\documentclass[11pt,a4paper,openany]{book}
% ============================================================
%  Preamble -- Topics in Mathematics with Applications in Finance
% ============================================================
\usepackage[T1]{fontenc}
\usepackage[utf8]{inputenc}
\usepackage{amsmath,amsthm,mathtools}
\usepackage{newpxtext,newpxmath} % provides the AMS symbols: do not load amssymb
\usepackage{bm}
\usepackage{microtype}
% Screen layout (read on a laptop, not printed): two parallel columns (paracol), the text column
% with the width it had on A4 (15 cm) and a 6.5 cm right column for notes, answers, feedback,
% examples, proofs, course notes, history and asides. Every right-column item starts level with the
% paragraph it belongs to; if the right column is still busy, the text waits (white space).
% The old A4 print layout was: a4paper,inner=2.4cm,outer=3.6cm,marginparwidth=3.0cm,marginparsep=0.3cm (twoside).
\usepackage[paperwidth=25cm,paperheight=29.7cm,top=2.6cm,bottom=2.6cm,left=2cm,textwidth=22cm,
  marginparsep=0.5cm,marginparwidth=6.5cm,headheight=23pt]{geometry}
\usepackage{paracol}
\setcolumnwidth{15cm/0.5cm,6.5cm}
% counters shared by the two columns (equations in a proof, examples, answers ... in the right column)
% (sectioning counters cannot be global with this kernel: they are copied to the right column by
% \synccounter before each switch, see \side@emit)
\globalcounter{equation}\globalcounter{footnote}\globalcounter{figure}
\usepackage{booktabs,array,tabularx,longtable}
\usepackage[dvipsnames]{xcolor}
\usepackage[shortlabels,inline]{enumitem}
\usepackage{tikz}
\usetikzlibrary{arrows.meta,positioning,decorations.pathreplacing,calc}
\usepackage{pgfplots}
\pgfplotsset{compat=1.18}
\usepackage[most]{tcolorbox}
\usepackage{listings}
\usepackage{fancyhdr}
\usepackage[colorlinks=true,linkcolor=NavyBlue,urlcolor=NavyBlue,citecolor=NavyBlue]{hyperref}
\usepackage[nameinlink]{cleveref}

\setlist{itemsep=2pt,topsep=4pt}
\setlength{\parskip}{3pt}
\setlength{\emergencystretch}{2em}
\hfuzz=5pt

% ---------- headers ----------
\pagestyle{fancy}
\fancyhf{}
\fancyhead[L]{\small\leftmark}
\fancyhead[R]{\small\makebox[\marginparwidth][r]{\rightmark\quad\thepage}}
\renewcommand{\headrulewidth}{0.3pt}
\fancypagestyle{plain}{\fancyhf{}\fancyfoot[C]{\small\thepage}\renewcommand{\headrulewidth}{0pt}}

% ---------- colours (bright on purpose: yellow notes, blue definitions, red theorems, green examples) ----------
\definecolor{defcol}{HTML}{1565C0}  % blue: definitions
\definecolor{thmcol}{HTML}{D32F2F}  % red: theorems, propositions, lemmas, corollaries
\definecolor{intcol}{HTML}{2E7D32}  % green: examples (also the green of the figures)
\definecolor{notecol}{HTML}{FBC02D} % yellow: margin notes and their answers
\definecolor{intucol}{HTML}{8E24AA} % purple: intuition
\definecolor{fincol}{HTML}{EF6C00}  % orange: finance
\definecolor{cscol}{HTML}{00897B}   % teal: case studies
\definecolor{keycol}{HTML}{C2185B}  % magenta: key formulas
\definecolor{exercol}{HTML}{3949AB} % indigo: exercises
\definecolor{srccol}{HTML}{546E7A}  % grey-blue: sources, course notes, remarks
\definecolor{histcol}{HTML}{B8860B} % bronze: history boxes (classical columns)
\definecolor{doubtcol}{HTML}{C62828}
\definecolor{ideacol}{HTML}{1565C0}
\definecolor{okcol}{HTML}{388E3C}
\colorlet{anscol}{notecol!85!black}

% ---------- table rows: a little air between lines in every table ----------
\AddToHook{env/tabular/begin}{\renewcommand{\arraystretch}{1.35}}
\AddToHook{env/tabularx/begin}{\renewcommand{\arraystretch}{1.35}}
\AddToHook{env/longtable/begin}{\renewcommand{\arraystretch}{1.35}}

% ---------- theorem-like environments (numbered by chapter), coloured headings ----------
\newtheoremstyle{coldef}{}{}{\normalfont}{}{\bfseries\color{defcol}}{.}{ }{}
\newtheoremstyle{colex}{}{}{\normalfont}{}{\bfseries\color{intcol}}{.}{ }{}
\newtheoremstyle{colexer}{}{}{\normalfont}{}{\bfseries\color{exercol}}{.}{ }{}
\newtheoremstyle{colthm}{}{}{\itshape}{}{\bfseries\color{thmcol}}{.}{ }{}
\newtheoremstyle{colrem}{}{}{\normalfont}{}{\itshape\color{srccol}}{.}{ }{}
\theoremstyle{coldef}
\newtheorem{definition}{Definition}[chapter]
\theoremstyle{colex}
\newtheorem{example}[definition]{Example}
\theoremstyle{colexer}
\newtheorem{exercise}{Exercise}[chapter]
\theoremstyle{colthm}
\newtheorem{theorem}[definition]{Theorem}
\newtheorem{proposition}[definition]{Proposition}
\newtheorem{lemma}[definition]{Lemma}
\newtheorem{corollary}[definition]{Corollary}
\theoremstyle{colrem}
\newtheorem{remark}[definition]{Remark}
\crefname{exercise}{Exercise}{Exercises}
\globalcounter{definition}\globalcounter{exercise}

% Right-column items are queued and written into the right column at the next point where the
% text column is between paragraphs at top level (end of the paragraph, end of a box, end of a
% list): there the columns are synchronised, the queued items are typeset in the right column,
% and the text continues from where it was.
\makeatletter
\newif\ifnote@inbox
\newif\ifside@inright
\def\note@queue{}
\def\side@pcname{paracol}
\newcommand{\side@emit}{\ifx\note@queue\@empty\else
  \let\note@tmp\note@queue\gdef\note@queue{}%
  \synccounter{chapter}\synccounter{section}\synccounter{subsection}\synccounter{subsubsection}%
  \switchcolumn*\side@inrighttrue\note@tmp\par\side@inrightfalse\switchcolumn\fi}
\newcommand{\note@flush}{\ifvmode\ifnote@inbox\else\ifside@inright\else
  \ifx\@currenvir\side@pcname\side@emit\fi\fi\fi\fi}
\newcommand{\side@put}[1]{\gappto\note@queue{#1}\note@flush}
\AddToHook{para/after}{\note@flush}
\tcbset{note queue/.code={\appto\kvtcb@afterbox{\note@flush}}}
% every chapter/appendix file is one paracol environment: it starts right after the \chapter
% line of the file (\begin{paracol}{2}) and is closed here (appendix A, only long tables, has none:
% longtable does not work inside paracol)
\AddToHook{include/end}{\par\note@flush\ifx\@currenvir\side@pcname\end{paracol}\fi}
\newcommand{\flushside}{\par\note@flush}
\makeatother
\tcbset{thmbox/.style={enhanced,breakable,colback=#1!7,colframe=#1!7,borderline west={3pt}{0pt}{#1},
  sharp corners,left=6pt,right=4pt,top=2pt,bottom=2pt,before skip=8pt,after skip=8pt,
  before upper={\csname note@inboxtrue\endcsname},note queue}}
\tcolorboxenvironment{definition}{thmbox=defcol}
\tcolorboxenvironment{theorem}{thmbox=thmcol}
\tcolorboxenvironment{proposition}{thmbox=thmcol}
\tcolorboxenvironment{lemma}{thmbox=thmcol}
\tcolorboxenvironment{corollary}{thmbox=thmcol}
% (examples are boxed in the right column, see \begin{example} below)
\tcolorboxenvironment{exercise}{thmbox=exercol}
\tcolorboxenvironment{remark}{thmbox=srccol}

% ---------- content boxes ----------
\tcbset{notebox/.style={enhanced,breakable,colback=#1!7,colframe=#1,
  fonttitle=\bfseries,coltitle=white,colbacktitle=#1,boxrule=0.6pt,arc=2pt,
  left=6pt,right=6pt,top=4pt,bottom=4pt,before skip=10pt,after skip=10pt,
  before upper={\csname note@inboxtrue\endcsname},note queue}}
% Box title: the box type in capitals, then the topic
\newcommand{\boxtitle}[2]{\MakeUppercase{#1}\ifx&#2&\else: #2\fi}
% Intuition: the "physicist's reading" of a result
\newtcolorbox{intuition}[1][]{notebox=intucol,title={\boxtitle{Intuition}{#1}}}
% Finance: where the maths meets the market
\newtcolorbox{finance}[1][]{notebox=fincol,title={\boxtitle{In finance}{#1}}}
% Case study: summary of an R case study / notebook
\newtcolorbox{casestudy}[1][]{notebox=cscol,title={\boxtitle{Case study}{#1}}}
% Summary / key formulas
\newtcolorbox{keyformulas}{notebox=keycol,fontupper=\small,title={\boxtitle{Key formulas}{}}}
% Sources: which lectures, videos and files a chapter is built on
\newtcolorbox{sources}{notebox=srccol,fontupper=\footnotesize,title={\boxtitle{Sources for this chapter}{}},before upper=\raggedright}
\newcolumntype{L}{>{\raggedright\arraybackslash}X}
% Note on course material (errata, missing lectures, differences 2013/2024): right column, see below
% Secondary material ("beyond the core"): can be skipped on a first reading
\newtcolorbox{secondary}[1][]{enhanced,breakable,colback=white,colframe=white,boxrule=0pt,
  borderline west={1.5pt}{0pt}{srccol!60},sharp corners,left=8pt,right=0pt,top=1pt,bottom=1pt,
  before skip=6pt,after skip=6pt,fontupper=\small,
  title={\footnotesize\sffamily$\diamond$\ \boxtitle{Beyond the core}{#1}},
  coltitle=srccol,colbacktitle=white,fonttitle=\bfseries,titlerule=0pt,toptitle=0pt,bottomtitle=0pt,
  before upper={\csname note@inboxtrue\endcsname},note queue}
% ============================================================
%  The right column: the reader's notes/answers/feedback in \tiny; examples, proofs, course notes,
%  history and asides in \scriptsize
% ============================================================
%  \note{...}     the reader's comment/doubt, NEW (not yet read by Claude)
%  \note*{...}    a request that has been carried out (green DONE tag)
%  \note+{...}    a comment that has been read and needs no answer (blue SEEN tag)
%  \begin{answer}{question}[topic] ... \end{answer}
%                 a doubt and its answer in one box: the question on top, then a dark yellow
%                 bar "ANSWER n: topic" and the answer; \answerpart{topic} starts a further
%                 answer in the same box (one per question); \label{ans:...} makes it citable
%  \feedback{...} a comment on one of the reader's own solutions (indigo)
%  \begin{history}[topic], \begin{aside}[topic], \begin{coursenote}[topic]
%  example and proof environments are also printed here (sideexample = example)
%  Inside a coloured box of the main text, margin items appear right after the box.
\makeatletter
\tcbset{sidebox/.style={enhanced,colback=#1!25,colframe=#1,boxrule=0.5pt,arc=1pt,
    left=3pt,right=3pt,top=2pt,bottom=2pt,before skip=0pt,after skip=0pt,
    fontupper=\tiny\sffamily\raggedright,fonttitle=\tiny\sffamily\bfseries,
    before upper={\note@inboxtrue}}}
\newcommand{\side@tag}[2]{\colorbox{#1}{\textcolor{white}{\bfseries\tiny #2}}\par}
\NewDocumentCommand{\note}{s t+ +m}{\side@put{\begin{tcolorbox}[sidebox=notecol]%
  \IfBooleanT{#1}{\side@tag{okcol}{DONE}}\IfBooleanT{#2}{\side@tag{defcol}{SEEN}}#3\end{tcolorbox}}}
\newcommand{\feedback}[1]{\side@put{\begin{tcolorbox}[sidebox=exercol]\side@tag{exercol}{FEEDBACK}#1\end{tcolorbox}}}
\newcounter{answer}[chapter]
\renewcommand{\theanswer}{\thechapter.\arabic{answer}}
\crefname{answer}{answer}{answers}
\globalcounter{answer}
\newcommand{\answerpart}[1]{\refstepcounter{answer}%
  \tcbsubtitle[colback=anscol,colframe=anscol,coltext=black,fontupper=\tiny\sffamily\bfseries,
    left=3pt,top=0.5pt,bottom=0.5pt,before skip=2pt,after skip=2pt]{ANSWER \theanswer\ifx&#1&\else: #1\fi}}
\NewDocumentEnvironment{answer}{m O{} +b}{\side@put{\begin{tcolorbox}[sidebox=notecol,colback=notecol!12]%
  #1\answerpart{#2}#3\end{tcolorbox}}}{}
% Content boxes of the right column (examples, proofs, course notes, history, asides): \scriptsize,
% one step above the \tiny of the reader's notes
\tcbset{sidecontent/.style={sidebox=#1,colback=#1!6,breakable,
    fontupper=\scriptsize\raggedright,fonttitle=\scriptsize\sffamily\bfseries,coltitle=white,colbacktitle=#1}}
\NewDocumentEnvironment{aside}{O{} +b}{\side@put{\begin{tcolorbox}[sidecontent=histcol,
  title={\boxtitle{Aside}{#1}}]#2\end{tcolorbox}}}{}
\NewDocumentEnvironment{coursenote}{O{} +b}{\side@put{\begin{tcolorbox}[sidecontent=srccol,
  title={\boxtitle{Course note}{#1}}]#2\end{tcolorbox}}}{}
% Examples and proofs keep their theorem-like look (numbering, labels, QED) but go to the right column
\let\origexample\example \let\endorigexample\endexample
\RenewDocumentEnvironment{example}{o +b}{\side@put{\begin{tcolorbox}[thmbox=intcol,
  fontupper=\scriptsize,before skip=0pt,after skip=0pt]%
  \IfValueTF{#1}{\begin{origexample}[#1]}{\begin{origexample}}#2\end{origexample}\end{tcolorbox}}}{}
\NewDocumentEnvironment{sideexample}{o +b}{\IfValueTF{#1}{\begin{example}[#1]}{\begin{example}}#2\end{example}}{}
\let\origproof\proof \let\endorigproof\endproof
\RenewDocumentEnvironment{proof}{o +b}{\side@put{\begin{tcolorbox}[thmbox=thmcol,colback=thmcol!3,
  fontupper=\scriptsize,before skip=0pt,after skip=0pt]%
  \IfValueTF{#1}{\begin{origproof}[#1]}{\begin{origproof}}#2\end{origproof}\end{tcolorbox}}}{}
\makeatother
% History: a fact from the history of finance/mathematics, framed like a classical
% temple front: a dentilled entablature (the title) on two Ionic columns (not breakable)
\newcommand{\histcolumn}[2]{% #1 = foot of the column, #2 = top of the capital (coordinate names)
  \fill[histcol] ([xshift=-6pt]#1) rectangle ([xshift=6pt,yshift=2.5pt]#1);
  \fill[histcol] ([xshift=-4.5pt,yshift=2.5pt]#1) rectangle ([xshift=4.5pt,yshift=4pt]#1);
  \fill[histcol!15] ([xshift=-3.5pt,yshift=4pt]#1) rectangle ([xshift=3.5pt,yshift=-5pt]#2);
  \foreach \dx in {-3.5,-1.2,1.2,3.5}
    \draw[histcol,line width=0.4pt] ([xshift=\dx pt,yshift=4pt]#1) -- ([xshift=\dx pt,yshift=-5pt]#2);
  \fill[histcol] ([xshift=-4.5pt,yshift=-5pt]#2) rectangle ([xshift=4.5pt,yshift=-2.5pt]#2);
  \foreach \dx in {-4.5,4.5} {
    \draw[histcol,line width=0.8pt,fill=white] ([xshift=\dx pt,yshift=-4.5pt]#2) circle (1.9pt);
    \fill[histcol] ([xshift=\dx pt,yshift=-4.5pt]#2) circle (0.6pt);}
  \fill[histcol] ([xshift=-6.5pt,yshift=-2.5pt]#2) rectangle ([xshift=6.5pt]#2);}
\makeatletter
\NewDocumentEnvironment{history}{O{} +b}{\side@put{\begin{tcolorbox}[histbox,title={\boxtitle{History}{#1}}]#2\end{tcolorbox}}}{}
\makeatother
\tcbset{histbox/.style={enhanced,colback=histcol!6,colframe=histcol,boxrule=0.8pt,sharp corners,
  fonttitle=\scriptsize\sffamily\bfseries,fontupper=\scriptsize\raggedright,coltitle=white,colbacktitle=histcol,
  left=18pt,right=18pt,top=6pt,bottom=3pt,before skip=0pt,after skip=0pt,
  before upper={\csname note@inboxtrue\endcsname},
  overlay={
    \draw[histcol,line width=2.5pt,dash pattern=on 2.5pt off 2pt]
      ([yshift=-1.25pt]title.south west) -- ([yshift=-1.25pt]title.south east);
    \coordinate (hfL) at ([xshift=9pt]frame.south west);
    \coordinate (htL) at ([xshift=9pt,yshift=-2.5pt]title.south west);
    \coordinate (hfR) at ([xshift=-9pt]frame.south east);
    \coordinate (htR) at ([xshift=-9pt,yshift=-2.5pt]title.south east);
    \histcolumn{hfL}{htL}\histcolumn{hfR}{htR}}}}
% The reader's own solution to an exercise, placed right below the exercise
%   \begin{mysolution}[optional topic] ... \end{mysolution}
\newtcolorbox{mysolution}[1][]{notebox=exercol,colback=exercol!5,fontupper=\small,
  title={\boxtitle{My solution}{#1}}}
% A request about the notes themselves (collected in the "Requests" chapter)
%   \request{done|open|noted}{where it was written}{the request}{what was done}
\makeatletter
\newcommand{\request}[4]{\ifstrequal{#1}{done}{\def\req@col{okcol}}{\ifstrequal{#1}{open}{\def\req@col{thmcol}}{\def\req@col{defcol}}}%
  \begin{tcolorbox}[sidebox=notecol,colback=notecol!12,fontupper=\small\sffamily\raggedright,
  before skip=5pt,after skip=5pt]%
  \side@tag{\req@col}{\MakeUppercase{#1}}%
  {\itshape #2.}\ #3%
  \tcbsubtitle[colback=anscol,colframe=anscol,coltext=black,fontupper=\scriptsize\sffamily\bfseries,
    left=3pt,top=0.5pt,bottom=0.5pt,before skip=2pt,after skip=2pt]{WHAT WAS DONE}#4\end{tcolorbox}}
\makeatother



% ---------- video timestamps ----------
%  \ts{24}{12}{23:10}  -> lecture 12 of the 2024 course, at 23:10
%  \ts{13}{18}{5:02}   -> lecture 18 of the 2013 course
\makeatletter % consecutive tags may break between them, but never before punctuation
\DeclareRobustCommand{\ts}[3]{\mbox{\footnotesize\sffamily\color{srccol}[$\blacktriangleright$\,20#1\,L#2\,@\,#3]}\@ifnextchar\ts{\allowbreak}{}}
\makeatother

% ---------- listings (short R snippets) ----------
\lstset{language=R,basicstyle=\ttfamily\footnotesize,keywordstyle=\color{defcol}\bfseries,
  commentstyle=\color{intcol}\itshape,stringstyle=\color{thmcol},frame=single,rulecolor=\color{black!20},
  breaklines=true,columns=fullflexible,showstringspaces=false}

% ---------- maths shortcuts ----------
\newcommand{\R}{\mathbb{R}}
\newcommand{\N}{\mathbb{N}}
\newcommand{\E}{\mathbb{E}}
\newcommand{\Prob}{\mathbb{P}}
\newcommand{\Q}{\mathbb{Q}}
\DeclareMathOperator{\Var}{Var}
\DeclareMathOperator{\Cov}{Cov}
\DeclareMathOperator{\Corr}{Corr}
\DeclareMathOperator{\tr}{tr}
\DeclareMathOperator{\rank}{rank}
\DeclareMathOperator{\diag}{diag}
\newcommand{\T}{^{\mathsf T}}
\newcommand{\dd}{\,\mathrm{d}}
\newcommand{\ind}{\mathbf{1}}
\newcommand{\eqd}{\stackrel{d}{=}}

% ---------- file references ----------
\DeclareRobustCommand{\file}[1]{\texttt{\small\detokenize{#1}}}


% Compile only some chapters while working on them, e.g.:
% \includeonly{chapters/ch01-markets-bonds}

\begin{document}

\frontmatter
\begin{titlepage}
  \centering
  \vspace*{3cm}
  {\Huge\bfseries Topics in Mathematics\\[4pt] with Applications in Finance\par}
  \vspace{1cm}
  {\Large Lecture notes\par}
  \vspace{0.5cm}
  {\large based on MIT 18.642 (Fall 2024) and MIT 18.S096 (Fall 2013)\par}
  \vspace{2cm}
  {\large Personal study notes of Alessio Martini\par}
  \vfill
  \begin{minipage}{0.85\textwidth}\small
    These notes are a personal, non-commercial study aid derived from the MIT OpenCourseWare
    courses \emph{18.642 Topics in Mathematics with Applications in Finance} (Fall 2024;
    instructors Peter Kempthorne, Vasily Strela, Jake Xia) and \emph{18.S096 Topics in
    Mathematics with Applications in Finance} (Fall 2013; instructors Peter Kempthorne,
    Choongbum Lee, Vasily Strela, Jake Xia), and on the guest lectures given in those courses.
    The original material is licensed under Creative Commons BY-NC-SA 4.0
    (\url{https://ocw.mit.edu/terms}); these notes are shared under the same licence.
    Lecture content is paraphrased and re-derived, not transcribed; any mistakes are mine.
  \end{minipage}
  \vspace{1cm}
\end{titlepage}

\chapter*{How to use these notes}
\addcontentsline{toc}{chapter}{How to use these notes}

Each chapter covers one \emph{topic} of the course. Chapters are ordered as the topics first
appear in the 2024 course (18.642); lectures from 2013 (18.S096) on the same topic are merged
into the same chapter, while 2013 lectures with no 2024 counterpart form
Part~\ref{part:2013}. Every chapter opens with a \textbf{Sources} box listing the lectures
(with their videos), slides, notes, case studies and problem sets it is built on.

\paragraph{Boxes.}
\begin{itemize}
  \item \textcolor{defcol}{\textbf{Definitions}} (blue), \textcolor{thmcol}{\textbf{theorems}} (red),
        \textcolor{intcol}{\textbf{examples}} (green), \textcolor{exercol}{\textbf{exercises}} and
        \textcolor{srccol}{\textbf{remarks}} are marked with a coloured bar on the left.
  \item \textcolor{keycol}{\textbf{Key formulas}} boxes collect the formulas to remember.
  \item \textcolor{intucol}{\textbf{Intuition}} boxes give the ``physicist's reading'' of a result.
  \item \textcolor{fincol}{\textbf{In finance}} boxes explain where the mathematics is used in practice.
  \item \textcolor{cscol}{\textbf{Case study}} boxes summarise the R case studies distributed with the course.
  \item \textcolor{srccol}{\textbf{Course notes}} flag errata in the slides, differences between the
        2013 and 2024 versions, and material that was not recorded.
  \item \textcolor{anscol}{\textbf{Answer}} boxes reply to the doubts in your margin notes (below).
\end{itemize}

\paragraph{Video timestamps.} A tag such as \ts{24}{12}{23:10} means: lecture~12 of the 2024
course, at minute 23:10 of the video, is where the point is explained in class. The 2024
lecture~1 is split into three videos, written \ts{24}{1-II}{12:36}.

\paragraph{Your annotations.} While reading, write directly in the chapter files. The macro \verb|\note{...}| prints a comment in the margin next to the text where it is written: a doubt, an idea, a reminder. For example \verb|\note{Why is duration divided by $1+y$?}| produces the box in the margin here. A doubt is answered in an \textbf{Answer} box right below the paragraph carrying the note; a request or reminder that has been carried out stays in the margin, marked DONE (\verb|\note*{...}|).\note{Why is the duration of a zero-coupon bond $T/(1+y)$ and not simply $T$? \emph{(Example: delete it from \file{main.tex}.)}}

\tableofcontents

\mainmatter
\part{The 2024 course (18.642)}\label{part:2024}
% !TEX root = ../main.tex
\chapter{Financial Markets, Instruments and Bond Mathematics}\label{ch:markets}

\begin{sources}
\small
\begin{tabularx}{\linewidth}{@{}l l L@{}}
\toprule
\textbf{Lecture} & \textbf{Speaker} & \textbf{Content}\\
\midrule
2024 L1-I   & V.~Strela, P.~Kempthorne & Organisation of the class, guest lectures, R/RStudio and financial data.\\
2024 L1-II  & J.~Xia    & Markets, products, participants, trader types, investment strategies, investment game.\\
2024 L1-III & V.~Strela & Compounding, discounting, zero-coupon and coupon bonds, yield, yield curve, duration.\\
2013 L1     & J.~Xia, V.~Strela & Same introduction with more detail on hedging, market making and prop trading; risk-aversion quiz; two intern projects (noisy delta, Kalman filter).\\
\bottomrule
\end{tabularx}

\smallskip
\textbf{Files.} Slides \file{mit18_642_f24_lec01_1.pdf} (intro, markets, bond math);
\file{mit18_642_f24_lec01_2.pdf} (R/RStudio) with the notebook \file{FM_Intro1.Rmd} in
\file{mit18_642_f24_rstudio.zip}; 2013 slides \file{MIT18_S096F13_lecnote1.pdf};
the 2024 \emph{Investment Game} page. No problem set of either year covers this chapter
directly (the 2024 problem sets start with linear algebra, \cref{ch:linalg}).
\end{sources}

\section*{Motivation}
The first lecture has two jobs: to give a vocabulary of markets, instruments and players that
every later application lecture takes for granted, and to introduce the single most important
idea in quantitative finance, \emph{pricing by no-arbitrage}, on the simplest possible
instruments: bonds. Everything in this chapter is elementary mathematically; what matters is
the way of reasoning. The discounting argument of \cref{sec:discounting} is the ancestor of
the replication argument behind Black--Scholes (\cref{ch:bs}) and the forward-rate formula
of \cref{sec:forward} is the seed of the HJM framework (\cref{ch:hjm}).
\note{What are the discounting argument? and the replication argument? and the forward-rate formula? and the HJM framework?}
\begin{answer}[Four arguments announced here]
The sentence is a roadmap; each item is developed later.
\begin{itemize}
  \item \textbf{Discounting argument} (\cref{sec:discounting}): the price today of \$1 paid in one
        year is fixed by comparison with riskless borrowing and lending, $x=1/(1+r)$; any other
        price gives a riskless profit.
  \item \textbf{Replication argument}: if a portfolio of traded assets pays exactly what a contract
        pays in every state of the world, the contract must cost what the portfolio costs.
        Discounting is the simplest case (the claim to \$1 is replicated by lending $1/(1+r)$);
        Black--Scholes (\cref{ch:bs}) replicates an option with a continuously rebalanced position
        of $\Delta$ shares plus cash.
  \item \textbf{Forward-rate formula} (\cref{prop:forward}): the interest fixed today for a loan
        from $t$ to $T$ is read off today's discount factors, $F_{tT}=\Lambda_t/\Lambda_T-1$.
  \item \textbf{HJM framework} (Heath--Jarrow--Morton, \cref{ch:hjm}): a model for the random
        evolution of the whole forward curve $f(t,T)$; no-arbitrage fixes its drift in terms of its
        volatility.
\end{itemize}
\end{answer}

% ------------------------------------------------------------------
\section{The course at a glance}
The course alternates \emph{mathematics lectures} (linear algebra, probability, stochastic
processes, regression, time series, volatility, stochastic calculus; P.~Kempthorne in 2024,
P.~Kempthorne and C.~Lee in 2013) with \emph{application lectures} by practitioners
(banks, hedge funds, exchanges, academia) \ts{24}{1-I}{4:05}. The mathematics is chosen
because it is used on trading floors, not for its own sake \ts{24}{1-I}{8:17}.
Illustrations and case studies are written in R; organisational details (problem sets, group
project, final paper, investment game, grading) are collected in \cref{app:logistics}.
\note{Nowadays everyone uses Python, so I will use it to illustrate the maths.}

\begin{history}[The course]
The class grew out of recruiting trips of J.~Xia and V.~Strela (then at Morgan Stanley) to MIT:
students with strong quantitative backgrounds asked more questions about finance than an
interview slot could answer, so a seminar was set up, later extended to a full term with
mathematics lectures by P.~Kempthorne \ts{24}{1-II}{4:14}. The 2013 edition (18.S096) was the
first to be recorded; the 2024 edition (18.642) keeps the same structure with new guest speakers.

\smallskip
{\footnotesize\textbf{Source.} J.~Xia's own account in the 2024 lecture (timestamp above); both
editions are on MIT OpenCourseWare (\url{https://ocw.mit.edu}).}
\end{history}

% ------------------------------------------------------------------
\section{Markets and products}\label{sec:products}

\subsection{Where trading happens}
\begin{itemize}
  \item \textbf{Exchanges}: centralised venues where standardised products (\emph{listed}
        securities, futures) trade under common rules \ts{24}{1-II}{8:23}.
  \item \textbf{Over-the-counter (OTC)}: bilateral contracts between two counterparties,
        not necessarily standardised and not subject to exchange rules; most swaps and many
        options trade OTC \ts{13}{1}{12:30}. OTC trading creates \emph{counterparty risk}
        (\cref{ch:counterparty}).
  \item \textbf{Electronic communication networks (ECN)}: electronic platforms acting as
        distributed exchanges \ts{24}{1-II}{9:25}.
\end{itemize}
\note{Difference between stock, equity, security? Define future, forward, option and swap. What are the exchange rules? Intuition behind counterparty risk?}
\begin{answer}[Securities, derivatives, exchanges, counterparty risk]
\begin{itemize}
  \item \textbf{Security}: any tradable financial claim (share, bond, note, \dots).
        \textbf{Equity}: ownership of a company, as opposed to debt. A \textbf{stock} (share) is the
        security that represents equity; in market jargon ``equities'' simply means stocks.
  \item \textbf{Forward}: private (OTC) agreement to buy or sell an asset at a date $T$ at a price
        $K$ fixed today; nothing is paid today, everything is settled at $T$.
        \textbf{Future}: the same, but standardised (size, dates) and exchange-traded, with gains and
        losses settled daily through margin accounts. \textbf{Option}: the right, not the
        obligation, to buy (call, payoff $(S_T-K)^+$) or sell (put, payoff $(K-S_T)^+$); the buyer
        pays a premium today. \textbf{Swap}: exchange of two streams of cash flows on a notional,
        e.g.\ fixed vs floating interest; it is a portfolio of forwards.
  \item \textbf{Exchange rules}: standardised contracts; only members trade directly (others go
        through brokers); orders are matched in an order book by price, then time; tick sizes,
        trading hours, circuit breakers and reporting duties. Trades are cleared by a central
        counterparty that collects margin from both sides, which is what removes bilateral
        counterparty risk.
  \item \textbf{Counterparty risk}: in an OTC contract your gain is only a promise of the other
        side. If they default when the contract is worth something to you, you lose that value
        (the cost of replacing the contract). It is like having lent money to them: the exposure is
        what they would owe you if they disappeared now, and it matters only when the contract has
        positive value to you. Collateral and clearing reduce it (\cref{ch:counterparty}).
\end{itemize}
\end{answer}

\subsection{What is traded}
\Cref{tab:products} lists the main product classes with the chapters in which their
mathematics appears.

\begin{table}[ht]
\centering\small
\caption{Main product classes.}\label{tab:products}
\begin{tabularx}{\linewidth}{@{}l X l@{}}
\toprule
\textbf{Class} & \textbf{What it is} & \textbf{Where}\\
\midrule
Money, FX & Currencies issued by different states; also gold and, arguably, crypto. & \ref{ch:timeseries}, \ref{ch:quanto}\\
Equity & Ownership of a company (stocks); baskets (indices such as the S\&P~500, ETFs). & \ref{ch:regression}, \ref{ch:portfolio}\\
Fixed income & Loans, securitised as bonds (government, corporate, municipal); interest-rate curves. & this ch., \ref{ch:rates}\\
Credit & Instruments whose value depends on default of an issuer. & \ref{ch:counterparty}, \ref{ch:quanto}\\
Real estate & Mortgages, mortgage-backed (MBS) and asset-backed securities (ABS). & --\\
Commodities & Metals, energy, agricultural products; mostly traded as futures. & \ref{ch:commodities}\\
Derivatives & Contracts whose payoff depends on another asset: forwards, futures, swaps, options, structured products. & \ref{ch:bs}, \ref{ch:rates}\\
\bottomrule
\end{tabularx}
\end{table}
\note*{Leave some space between lines in the table.}

A few points from the lectures deserve emphasis.
\begin{itemize}
  \item \textbf{Primary vs secondary market.} A company becomes listed through an
        \emph{initial public offering} (IPO), the primary market; afterwards its shares are
        traded among investors, the secondary market \ts{24}{1-II}{11:33}.
  \item \textbf{A bond is a securitised loan.} The issuer pays periodic interest
        (\emph{coupons}) and repays the \emph{principal} or \emph{notional} at maturity.
        Debt is ``more generic'' than equity: finance is fundamentally about someone who has
        money and someone who needs it \ts{13}{1}{14:36}. Since the issuer may fail to pay,
        a bond carries \emph{credit risk} \ts{24}{1-II}{12:36}; the 2008 crisis was a credit crisis
        that started in mortgages and spread through the derivatives built on them
        \ts{24}{1-II}{13:36}.
  \item \textbf{Asset-backed securities}: debt issued against a pool of assets (e.g.\ mortgages),
        whose risk depends on the income stream of the pool \ts{13}{1}{16:44}.
  \item \textbf{Derivatives}: an \emph{option} is the right but not the obligation to buy
        (call) or sell (put) at a fixed price; a \emph{forward} is an agreement today to trade
        at a future date; a \emph{future} is a standardised, exchange-traded forward; a
        \emph{swap} exchanges two streams of cash flows (e.g.\ fixed vs floating interest)
        \ts{24}{1-II}{15:41}.
\end{itemize}
\note{Is the primary market an actual market? what is a market?}
\begin{answer}[Primary market; what a market is]
The primary market is not a place but a process: the first sale of newly issued securities by
the issuer to investors (an IPO for shares, auctions or bank syndicates for bonds). The money goes to
the issuer, and investment banks organise the sale (underwriting, book-building). In the secondary
market investors trade existing securities among themselves and the issuer receives nothing.
In general a \emph{market} is any mechanism that brings buyers and sellers together and lets a price
form: an exchange with an order book, a network of dealers quoting prices (OTC), or an auction.
\end{answer}

\begin{finance}[Impossible events happen]
In April 2020 the front-month WTI crude-oil futures contract traded at a significantly
\emph{negative} price: holders of long contracts had to pay to get rid of them because nobody
could store the oil. Some brokers' systems had not been programmed for negative prices
\ts{24}{1-I}{14:30}. Negative interest rates, once considered impossible, were the norm in the
eurozone and Japan for years. Lesson: a model's hard constraints (``prices are positive'') are
assumptions, not facts.
\end{finance}

\begin{figure}[ht]
\centering
\begin{tikzpicture}
\begin{axis}[width=0.9\linewidth,height=5.2cm,xlabel={year},ylabel={WTI (\$ per barrel)},
  xmin=2011,xmax=2024.75,ymin=-45,ymax=130,xtick={2012,2014,...,2024},
  xticklabel style={/pgf/number format/1000 sep={}},grid=major,grid style={black!10}]
  \addplot[thin,black] table {data/fm_wti.dat};
  \addplot[thick,thmcol,domain=2011:2024.75,samples=2] {0};
  \node[font=\scriptsize,anchor=east,thmcol] at (axis cs:2020.1,-30) {20 Apr 2020: $-36.98$};
\end{axis}
\end{tikzpicture}
\caption{WTI crude oil (FRED series \texttt{DCOILWTICO}, weekly closes, daily in spring 2020),
2011--2024, with the zero line in red: re-drawn from the data of the course notebook
\file{FM_Intro1.Rmd} (\cref{sec:R}), shown in class \ts{24}{1-I}{14:54}.}\label{fig:wti}
\end{figure}
\note*{Here I am sure there were picture from the presentation. We should include them to make the presentation clearer and more interesting.}

% ------------------------------------------------------------------
\section{Market participants}\label{sec:participants}

\note{repeal? liabilities?}
\begin{itemize}
  \item \textbf{Commercial banks} take deposits and lend. \textbf{Investment banks} raise
        capital and trade. The two were separated in the US by the Glass--Steagall Act (1933);
        its repeal in 1999 is often blamed as a cause of the 2008 crisis \ts{13}{1}{17:44}.
        An investment bank typically has three client divisions \ts{24}{1-II}{17:45}:
        \emph{Equity} (stocks and equity derivatives), \emph{Fixed Income} (bonds, rates, FX,
        commodities and their derivatives) and the \emph{Investment Banking Division} (IBD:
        corporate finance, IPOs, bond issuance, M\&A advisory); plus asset and wealth
        management.
  \item \textbf{Dealers vs brokers.} A dealer (market maker) quotes a price and takes the other
        side of the trade, i.e.\ takes \emph{principal risk}; a broker only matches buyers and
        sellers and earns a commission \ts{13}{1}{23:00}.
  \item \textbf{Asset owners and managers}: mutual funds (usually long-only), insurance companies
        and pension funds (assets invested to meet liabilities), sovereign wealth funds,
        endowments, family offices \ts{13}{1}{24:04}.
  \item \textbf{Hedge funds} trade their own and their investors' capital to generate returns;
        \textbf{private equity} buys stakes in companies to improve and later sell them.
  \item \textbf{Retail investors}, \textbf{central banks} (policy-driven: set short rates,
        provide liquidity, intervene in FX) and \textbf{corporates} (mostly hedgers).
\end{itemize}
\begin{answer}[Repeal; liabilities]
To \emph{repeal} a law is to abolish it with a later law: the Gramm--Leach--Bliley Act (1999) removed
the Glass--Steagall separation, so commercial and investment banks could merge again (e.g.\
Citigroup). \emph{Liabilities} are what an entity owes: for a bank, its deposits; for a pension fund,
the pensions it has promised; for an insurer, future claims. ``Invested to meet liabilities'' means
choosing assets whose cash flows match those future payments (asset--liability management), e.g.\ long
bonds against long-dated pensions, with matching durations (\cref{sec:duration}).
\end{answer}
\note*{Central bank decides interest rate. Right now in 2024 the interest curve is inverted! The central bank of US is Federal Reserve Bank.}\note{what is a typical shape of interest curve?}
\begin{answer}[Central banks and the yield curve]
The central bank sets only the \emph{short-term} policy rate (in the US the federal funds target
range, decided by the FOMC; in the eurozone the ECB deposit rate). Longer yields are set by the
market: expectations of future policy rates plus a term premium. The US central bank is the
\emph{Federal Reserve System}: the Board of Governors in Washington plus twelve regional Federal
Reserve Banks. On the inversion: the 2-year vs 10-year Treasury spread was negative from July 2022
until early September 2024, the longest inversion on record; the September 2024 curve of
\cref{fig:yieldcurve} is still inverted at the short end.

\emph{Typical shape}: upward sloping and concave, steep over the first few years and flattening
beyond about ten years (1992 and 2021 in \cref{fig:yieldcurve}). The other shapes are flat,
inverted (short above long, when the market expects the central bank to cut) and humped (a
maximum at medium maturities).
\end{answer}

\paragraph{The driving force.} Every participant tries to maximise gains while limiting losses
under uncertainty. The fundamental trade is between \emph{lenders} (have money, need return) and
\emph{borrowers} (need money, bear risk) \ts{24}{1-II}{20:49}. The same pairing appears as
\emph{limited partners} (LP, the investors in a fund) vs \emph{general partners} (GP, the fund
manager). Whether markets are a zero-sum game depends on the product: many trades net to zero,
but markets exist because they transfer capital from those who have it to those who can use it
productively \ts{13}{1}{20:53}.
\note{Nice point: what is the point of a (financial) market?}
\begin{answer}[What financial markets are for]
\begin{itemize}
  \item \emph{Capital allocation}: savings flow to those who can use them productively;
  \item \emph{risk transfer}: hedgers pass risks to those willing to bear them for a premium;
  \item \emph{price discovery}: prices aggregate dispersed information into a public signal;
  \item \emph{liquidity}: assets can be turned into cash quickly, which makes holding them cheaper;
  \item \emph{moving consumption in time}: saving now to spend later, or the reverse.
\end{itemize}
Even when payoffs net to zero (derivatives), a trade can leave both sides better off, because
they start from different exposures and preferences.
\end{answer}

% ------------------------------------------------------------------
\section{Three types of traders}\label{sec:tradertypes}

\subsection{Hedgers}
Hedgers already have an exposure and want to \emph{reduce} it \ts{13}{1}{26:13}: a homeowner
with a floating-rate mortgage who locks in a fixed rate; an exporter with future euro revenues
who sells euros forward; a company in an M\&A deal hedging currency and rate exposure.
\note{Exposure = risk? what is M\&A?}
\begin{answer}[Exposure; M\&A]
\emph{Exposure} is the amount of value that moves when a risk factor moves: an exporter expecting
EUR~10 million of revenues has a EUR~10 million EUR/USD exposure. \emph{Risk} combines the exposure
with how much the factor moves: roughly, the size of the typical loss is exposure $\times$ volatility
of the factor. Exposure is thus a sensitivity (for derivatives, the delta). \emph{M\&A} = mergers and
acquisitions: one company buying another, or two merging. In a cross-border deal the buyer has to
pay in a foreign currency at a future closing date, hence the FX and rate hedges.
\end{answer}

\begin{example}[Samurai bonds and the carry trade {\ts{13}{1}{30:29}}]
Australian rates have historically been well above Japanese rates; the 2013 slide quotes a
10-year swap-rate difference of $4.5\% - 1\% = 3.5\%$. An Australian company can therefore
issue a bond in Japanese yen (JPY) (a \emph{Samurai bond}), paying $\approx 1\%$, and use the proceeds in
Australia where money earns $\approx 4.5\%$. Why doesn't everybody do this forever? Because the
Australian dollar (AUD) may weaken against the yen, wiping out the rate differential; and if everybody
holds the same position, the exit moves the market against them. To hedge, the company must
\emph{receive} JPY and \emph{pay} AUD in the future, which in practice means FX forwards,
interest-rate swaps and cross-currency (basis) swaps.
\end{example}
\note*{JPY = japanese yen, AUD = australian dollar.}

\begin{intuition}[Covered interest parity]
The yen/AUD trade is an example of a general principle you will see formalised in
\cref{ch:bs}: once the exchange-rate risk is hedged with a forward, the forward price itself
incorporates the interest-rate differential (\emph{covered interest parity}). A risk-free
``free lunch'' from the rate differential would be an arbitrage.
\end{intuition}

\subsection{Market makers}
A market maker quotes a \emph{bid} (price at which it buys) and an \emph{offer} or \emph{ask}
(price at which it sells), with bid $<$ offer, and aims to earn the spread rather than to
take a view \ts{24}{1-II}{21:56}. For transparent products (a liquid stock) the spread is tiny;
for opaque products (say a two-month call on Apple struck at 500) the market maker provides
liquidity and must then manage the residual inventory of risk \ts{13}{1}{33:41}. This is done by
balancing the \emph{Greeks} of the book \ts{13}{1}{34:46}, the partial derivatives of the
value $V$ of the option (or of the whole book) with respect to the price $S$ of the underlying,
time $t$ and the volatility $\sigma$ of the underlying:
\note*{def: V is the value of the option, S is the underlying price, t is time, sigma is volatility.}
\begin{center}
\begin{tabular}{@{}lll@{}}
\toprule
Greek & Definition & Meaning\\
\midrule
Delta $\Delta$ & $\partial V/\partial S$ & sensitivity to the underlying price\\
Gamma $\Gamma$ & $\partial^2 V/\partial S^2$ & curvature, change of delta\\
Theta $\Theta$ & $\partial V/\partial t$ & ``carry'' or time decay when nothing moves\\
Vega $\nu$ (or kappa) & $\partial V/\partial \sigma$ & sensitivity to volatility\\
\bottomrule
\end{tabular}
\end{center}
Beyond the Greeks, a dealer book is managed through tail-risk measures such as
\emph{Value at Risk} (\cref{ch:var}), capital usage, balance sheet and leverage. Before 2008 many
banks were levered about 40 times: with \$1 of capital and \$40 of exposure, a 2.5\% move wipes
out the capital \ts{13}{1}{35:47}.
\note{def of liquidity, residual inventory of risk. Dealer book = order book? tail-risk measures?}
\begin{answer}[Liquidity, inventory, dealer book, tail risk]
\begin{itemize}
  \item \textbf{Liquidity}: the ability to trade a large size quickly without moving the price much;
        measured by the bid--ask spread, the depth of the market and traded volume.
  \item \textbf{Residual inventory of risk}: after buying an option from a client the market maker
        holds a position it did not choose, and cannot sell the same option on at once (it is
        illiquid). It hedges what it can with liquid instruments (shares for delta, other options for
        vega); what is left is the residual risk it carries in inventory.
  \item \textbf{Dealer book $\neq$ order book}. An order book is the exchange's list of pending buy
        and sell orders at each price. A dealer's \emph{book} is its portfolio of positions, managed
        as a whole.
  \item \textbf{Tail-risk measures} describe the extreme losses of the profit-and-loss distribution:
        \emph{Value at Risk} (the loss exceeded with probability 1\%, say, over one day),
        \emph{Expected Shortfall} (the average loss beyond the VaR), stress tests (\cref{ch:var}).
\end{itemize}
\end{answer}

\begin{coursenote}[Vega or kappa?]
J.~Xia's anecdote \ts{13}{1}{6:15}: on his first day on the Morgan Stanley options desk he asked
for the ``vega report''; the training manual explained that the house term was \emph{kappa},
since vega is not a Greek letter. Both names are in use; this reflects how young quantitative
finance is as a discipline (the Black--Scholes model dates to 1973) \ts{13}{1}{9:19}.
\end{coursenote}

\subsection{Proprietary traders and portfolio managers}
Risk takers deliberately take positions to earn a return \ts{13}{1}{36:50}. Styles include:
\begin{itemize}
  \item \textbf{Directional}: long or short on a view (including ``gut'' trading).
  \note{what is gut trading?}
  \item \textbf{Arbitrage}: exploit violations of \emph{deterministic} price relationships,
        e.g.\ the same gold priced differently in two markets, or a spot price inconsistent with
        the forward price \ts{13}{1}{37:50}.
  \item \textbf{Value and relative value}: estimate the ``fair value'' of an instrument relative
        to others and trade on the difference over a longer horizon.
  \item \textbf{Systematic}: computer models such as trend following, momentum,
        statistical arbitrage across many stocks.
  \item \textbf{Fundamental / global macro}: trade balances, earnings, policy changes
        \ts{13}{1}{40:00}; \textbf{special situations}: distressed companies.
\end{itemize}
\begin{answer}[Gut trading]
Trading on instinct (``gut feeling'') and experience rather than on a model or a fixed rule: the
extreme form of \emph{discretionary} trading, as opposed to \emph{systematic} trading.
\end{answer}

\begin{definition}[Alpha and beta]\label{def:alphabeta}
Let $R_b$ be the return of a benchmark (e.g.\ an S\&P~500 index fund) and $R_a$ the return of a
portfolio. The linear regression
\[
  R_a = \alpha + \beta\, R_b + \varepsilon, \qquad \E[\varepsilon]=0,\ \Cov(\varepsilon, R_b)=0,
\]
decomposes the portfolio return into a part $\beta R_b$ explained by co-movement with the
benchmark (\emph{beta}, market exposure, cheap to obtain by buying the index) and an excess
return $\alpha$ (\emph{alpha}, the manager's skill) \ts{13}{1}{28:21}. The least-squares estimate is
$\hat\beta = \widehat{\Cov}(R_a,R_b)/\widehat{\Var}(R_b)$; regression is developed in
\cref{ch:regression} and the CAPM interpretation in \cref{ch:portfolio}.
\end{definition}
\note{def of what is a return. what is CAPM interpretation?}
\begin{answer}[Returns; CAPM]
Over one period, with price $P_t$ and dividend $D_t$, the \emph{simple return} is
$R_t=(P_t+D_t-P_{t-1})/P_{t-1}$ and the \emph{log return} is $r_t=\ln(1+R_t)$; log returns add up
over time and $r_t\approx R_t$ for small moves. The CAPM (\cref{ch:portfolio}) states that in
equilibrium expected \emph{excess} returns are proportional to beta with respect to the market
portfolio $M$: $\E[R_a]-r_f=\beta_a\bigl(\E[R_M]-r_f\bigr)$. Written for excess returns with the
market as benchmark, the regression above then has $\alpha=0$: a non-zero alpha is return that
market risk does not explain, i.e.\ skill or an unmodelled risk factor.
\end{answer}

\begin{coursenote}[Slide glitch]
On the 2013 slide the formula reads ``$R(a) = a + b \cdot R(b)$'': the Greek letters $\alpha,\beta$
were lost in the conversion, as J.~Xia remarks in class. In the same lecture a quote is given
as ``\$0.99 bid, \$0.95 offer''; the bid must be below the offer, and indeed later it becomes
``\$0.90 bid, \$0.95 offer''.
\end{coursenote}

% ------------------------------------------------------------------
\section{Where mathematics enters}\label{sec:mathuse}
Mathematics is used in three ways \ts{24}{1-II}{23:01}:
\begin{itemize}
  \item \textbf{Pricing models}: relative value and \emph{arbitrage-free} prices of complex
        instruments (options need PDEs; \cref{ch:bs}); a pricing model also delivers the risk
        parameters (Greeks) needed to hedge.
  \item \textbf{Risk management}: position sizing, diversification, when to take profit or cut
        losses, leverage, liquidity, counterparty risk; it counteracts greed and fear
        \ts{24}{1-II}{24:10}.
  \item \textbf{Trading strategies}: the ``holy grail'' or perpetual money machine. Such
        strategies, when they exist, are kept secret, and markets adapt: no model can be built
        and left alone forever \ts{13}{1}{42:03}. Firms protect them with confidentiality
        obligations on trade secrets that do not expire when an employee leaves, and with
        non-compete clauses (often a paid ``garden leave'' of some months) that keep departing
        traders and quants away from competitors.
\end{itemize}
\note*{Life-long no-disclosure contracts for trading strategies.}

\paragraph{Efficient markets vs behavioural finance.} Can one learn from historical patterns and
extrapolate? The \emph{efficient market hypothesis} says no: a \$100 bill lying on the ground would
have been picked up already. \emph{Behavioural finance} replies that people are not always
rational, which creates opportunities \ts{24}{1-II}{26:15}. Deterministic (physics-like)
relationships are nowadays almost always arbitraged away; what remains are \emph{statistical}
relationships, which hold only with some probability \ts{24}{1-II}{29:25}. Market cycles are long
and memories short: people over-weight recent experience \ts{13}{1}{51:23}.

\begin{example}[Risk-aversion quiz {\ts{13}{1}{44:04}}]\label{ex:quiz}
\emph{Loss framing.}
\begin{enumerate}[(A)]
  \item lose \$500 with probability $0.8$, win \$500 with probability $0.2$;
  \item lose \$280 for sure.
\end{enumerate}
\emph{Gain framing.}
\begin{enumerate}[(A)]
  \item win \$500 with probability $0.8$, lose \$500 with probability $0.2$;
  \item win \$280 for sure.
      \note{I dont understand what's the question. whats the best strategy? }
\end{enumerate}

\smallskip
In the loss framing $\E[A] = 0.8(-500)+0.2(500) = -300 < -280 = \E[B]$, so B has the higher
expectation; yet many people choose A, to avoid ``locking in'' a loss. In the gain framing
$\E[A] = +300 > 280 = \E[B]$, but most people choose B. A student's answer for B in class was
``higher Sharpe'': indeed $\operatorname{sd}(A) = 1000\sqrt{0.8\cdot 0.2} = 400$ while B is riskless.
\end{example}
\begin{answer}[What the quiz asks]
Each framing asks: \emph{which of the two would you choose?} It is a question about your
preferences, and there is no single right answer. Xia says so explicitly: he is not telling the class which
choice is right \ts{13}{1}{47:49}. A pure expected-value maximiser picks B in the loss framing
($-280>-300$) and A in the gain framing ($300>280$). Most people do the opposite in both. The
point is that the \emph{framing} (losing vs winning the same kind of bet) flips the choice, while a
rational choice should depend only on the payoffs, one's utility and one's wealth (see the
intuition box below).
\end{answer}
\note*{Brief def of Sharpe ratio, and mention the other metrics to assess an investment strategy.}

\paragraph{Sharpe ratio and other performance metrics.} For a strategy with return $R$ (per period),
risk-free rate $r_f$ and volatility $\sigma=\operatorname{sd}(R)$, the \emph{Sharpe ratio} is
\[
  \mathrm{SR}=\frac{\E[R]-r_f}{\sigma},
\]
the excess return per unit of risk. It is usually annualised, $\mathrm{SR}_{\text{year}}\approx\sqrt{252}\,\mathrm{SR}_{\text{day}}$ for
daily returns. It does not change under leverage, which is why it is the standard basis for
comparing strategies (\cref{ch10:prop-sharpe}). A riskless payoff has $\sigma=0$, hence the student's
``higher Sharpe'' for B. Other common metrics:
\begin{itemize}
  \item \textbf{Sortino ratio}: like Sharpe, but divides by the downside deviation, so upside
        volatility is not penalised;
  \item \textbf{maximum drawdown}: largest peak-to-trough loss of the cumulative P\&L;
        the \textbf{Calmar ratio} is annual return divided by maximum drawdown;
  \item \textbf{information ratio}: active return over a benchmark divided by the tracking error;
  \item \textbf{alpha and beta} with respect to a benchmark (\cref{def:alphabeta});
  \item \textbf{Value at Risk} and \textbf{Expected Shortfall}: tail losses (\cref{ch:var});
  \item \textbf{hit rate} (fraction of winning trades) and \textbf{profit factor} (gross gains over
        gross losses), close to the ratio $(G-L)/(G+L)$ of the Investment Game (\cref{ex:game});
  \item \textbf{turnover} and \textbf{capacity}: how much trading the strategy needs and how much
        capital it can absorb before costs eat the returns.
\end{itemize}

\begin{intuition}[Prospect theory: reflection effect and utility]
The pattern ``risk-seeking over losses, risk-averse over gains'' is the \emph{reflection effect}
of Kahneman and Tversky's prospect theory (not named in the lecture). In trading it appears as the
tendency to sell winners too early and hold losers too long, which is why discipline and
stop-loss rules are part of risk management \ts{13}{1}{46:10}. With a concave utility $u$
(risk aversion), B is preferred in the gain framing whenever $u(280) > 0.8\,u(500)+0.2\,u(-500)$;
the answer depends on wealth: a student with \$800 in the bank and an investor with \$100{,}000
should not answer the same way \ts{13}{1}{50:20}.
\end{intuition}
\note*{the titles of the boxes should be given more importance. here id call it something like: INTUITION: KT prospect theory, reflection effect and utility (this is just an idea)}

% ------------------------------------------------------------------
\section{Choosing an investment strategy}\label{sec:strategies}
J.~Xia's list of choices every investor faces \ts{24}{1-II}{27:19}:
\begin{center}\small
\begin{tabularx}{\linewidth}{@{}L@{\quad vs\quad}L@{}}
\toprule
Internal (direct investing) & external (fund of funds)\\
Public assets (equity, bonds, FX, commodities) & private assets (PE, VC, real estate, natural resources)\\
Passive (beta, benchmark) & active (alpha)\\
Systematic (quantitative rules) & discretionary (fundamental judgement)\\
Deterministic relationships & statistical relationships\\
Trend following (buy strength) & mean reversion (buy the dip)\\
Short term (fast, liquid) & long term (slow, illiquid)\\
Value: $V(0) \gg P$ & growth: $V(T) \gg V(0)$\\
History (repeating patterns) & future (extrapolation)\\
\bottomrule
\end{tabularx}
\end{center} \note*{i am pretty sure he said something more about each choice}
What Xia said about each choice:
\begin{itemize}
  \item \textbf{Internal vs external}: do it yourself, or pay other people to do it for you
        (invest in their funds, or in a fund of funds) \ts{24}{1-II}{27:44}.
  \item \textbf{Public vs private}: venture capital, private equity, buyouts, real-estate and
        natural-resource funds are not transparent, ``but they seem to make a lot of returns,
        too''. The two are very different markets and you need to pick one \ts{24}{1-II}{28:04}.
  \item \textbf{Passive vs active}: stock picking, or just buying the index. Buying the
        S\&P~500 over the last ten years beat most fund managers, ``so why shouldn't you just do the
        passive?'' \ts{24}{1-II}{28:23}.
  \item \textbf{Systematic vs discretionary}: decisions from quantitative rules or models, or case
        by case with your own judgement \ts{24}{1-II}{28:52}.
  \item \textbf{Deterministic vs statistical}: a deterministic relationship is like physics (A plus
        B always gives C); a statistical one gives C only with some probability. Deterministic
        arbitrage is nowadays very hard to find (the \$100 bills have been picked up), ``but sometimes,
        they still show up'' \ts{24}{1-II}{29:18}.
  \item \textbf{Trend following vs mean reversion}: when the market goes up, buy more, or sell
        (play the contrarian, buy the dip) \ts{24}{1-II}{29:50}.
  \item \textbf{Short vs long term}: short-term trading is fast, in liquid markets, and you can
        change your mind easily. Long-term decisions are slow: a start-up in a venture-capital
        portfolio may take ten years to prove itself \ts{24}{1-II}{30:10}.
  \item \textbf{Value vs growth}: Warren Buffett-style value investing wants the price well below
        today's value. Growth investing projects, say, NVIDIA's income over the next ten years and buys
        if the future value exceeds the current one. In recent years more and more people focus on
        growth, ``you need to figure out why'' \ts{24}{1-II}{30:37}.
  \item \textbf{History vs future}: believe that history repeats and follow the patterns, or refuse
        to simply extrapolate \ts{24}{1-II}{31:09}.
\end{itemize}
Here $P$ is the current price and $V(t)$ the (estimated) value of the company at time $t$: a value
investor buys what is cheap relative to today's value, a growth investor buys expected growth. \note{This is just a different time-scale approach: connection to renormalization.}
\begin{answer}[Value vs growth as time scales]
Partly right. Both styles bet that price converges to value; they differ in \emph{which} value and
over what horizon. Value bets that the gap $V(0)-P$ closes (a mean-reversion time, typically
years). Growth bets on the drift of $V(t)$ itself up to a distant $T$. So the choice is indeed about
which time scale carries your signal. The analogy with the renormalization group is loose, though:
there is no coarse-graining map and no fixed point here, only two different observables. The
solid physics connections are elsewhere: \emph{separation of time scales} (fast noise vs slow
fundamental drift, as in the Ornstein--Uhlenbeck model of \cref{ch:sde}) and the empirical
\emph{scaling laws} of returns across time scales (Mandelbrot's fat tails and multifractals;
R.~N.~Mantegna and H.~E.~Stanley, \emph{An Introduction to Econophysics}, 2000), where
renormalization-like ideas do appear.
\end{answer}
The key message \ts{24}{1-II}{31:28}: understand your own objectives and constraints (what the
money is for, how much you can afford to lose, your horizon), pick the area where you have an edge
over the rest of the market, and ``if math is your strength, make math your edge''. \note{what is an edge?}
\begin{answer}[Edge]
An \emph{edge} is a persistent advantage that makes your expected return, after costs, positive,
and that you can explain: ``you need to understand why you should make money. Other people
should give you the money'' \ts{24}{1-II}{32:10}. Typical sources:
\begin{itemize}
  \item better \emph{information} or better processing of public information;
  \item better \emph{models} (the quant's edge);
  \item \emph{speed} or lower \emph{costs};
  \item \emph{risk capacity} or \emph{patience}: holding what others must sell, or for longer than they can;
  \item being paid for providing a service (liquidity, insurance).
\end{itemize}
If you cannot say who is on the other side and why they lose, you probably have no edge.
\end{answer}

% ------------------------------------------------------------------
\section{Time value of money}\label{sec:tvm}

\subsection{Compounding}
Invest \$1 at a constant annual interest rate $r$ \ts{24}{1-III}{1:00}. \note{what are tipical values of r?} With annual compounding
(interest earns interest) it grows to $1+r$ after one year and to $(1+r)^n$ after $n$ years. If
interest is paid and reinvested $m$ times a year, each period earns $r/m$ and
\begin{equation}\label{eq:mcompound}
  1 \;\longmapsto\; \Bigl(1+\frac{r}{m}\Bigr)^{mn} \quad\text{after $n$ years.}
\end{equation}
\begin{answer}[Typical values of $r$]
A few percent per year, but the range is wide:
\begin{itemize}
  \item US policy rate (federal funds): $0$--$0.25\%$ in 2008--2015 and 2020--2022, then
        $5.25$--$5.50\%$ from July 2023 to September 2024;
  \item 10-year US Treasury yield: about $0.5\%$ at its 2020 low, about $5\%$ in October 2023,
        above $15\%$ in 1981 (when the policy rate was near $20\%$ to fight inflation);
  \item negative rates: the ECB deposit rate was $-0.5\%$ in 2019--2022, and the Bank of Japan's
        policy rate was $-0.1\%$ in 2016--2024;
  \item borrowers pay more than governments: mortgages, corporate bonds and credit cards carry a
        credit spread on top (from $1\%$ to more than $20\%$ a year for credit cards).
\end{itemize}
\end{answer}
Letting $m\to\infty$ (continuous compounding) \ts{24}{1-III}{3:02},
\begin{equation}\label{eq:contcompound}
  \lim_{m\to\infty}\Bigl(1+\frac{r}{m}\Bigr)^{mt} = e^{rt},
\end{equation}
since $\ln(1+r/m)^{mt} = mt\,\ln(1+r/m) = mt\,(r/m + O(m^{-2})) \to rt$.

\begin{history}[Compound interest and the number $e$]
V.~Strela credits Jacob Bernoulli with ``maybe the first mathematical-finance paper in history''
\ts{24}{1-III}{4:02}. The facts:
\begin{itemize}
  \item in 1683 Jacob Bernoulli (1655--1705) studied continuously compounded interest and asked for
        the limit of $(1+1/n)^n$ as $n\to\infty$, showing that it lies between 2 and 3;
  \item he published it in May 1690 in \emph{Acta Eruditorum}, in the paper ``Quaestiones nonnullae
        de usuris, cum solutione problematis de sorte alearum'' (``Some questions on interest,
        with the solution of a problem on games of chance''). Finance and probability were
        born together;
  \item the letter $e$ is Euler's: it first appears in a letter to Goldbach (1731) and in print in
        his \emph{Mechanica} (1736).
\end{itemize}
``First'' is generous: Fibonacci's \emph{Liber Abaci} (1202) already compares cash flows at
different dates with present-value calculations (W.~N.~Goetzmann).

\smallskip
{\footnotesize\textbf{References.} J.~Bernoulli, \emph{Acta Eruditorum}, May 1690, pp.~219--223;
J.~J.~O'Connor and E.~F.~Robertson, ``The number $e$'', MacTutor History of Mathematics
(\url{https://mathshistory.st-andrews.ac.uk/HistTopics/e/});
W.~N.~Goetzmann, ``Fibonacci and the Financial Revolution'', NBER Working Paper 10352 (2004).}
\end{history}\note*{I love history facts! can we implement a special box for them? like something that looks like old, antique. Furthermore, factcheck them and add references.}

\begin{remark}[Equivalent rates]
The same growth can be quoted with different conventions: a rate $r_m$ compounded $m$ times a year
and a continuously compounded rate $r_c$ are equivalent iff $(1+r_m/m)^m = e^{r_c}$, i.e.
$r_c = m\ln(1+r_m/m)$. A rate is meaningless without its compounding convention.
\end{remark}\note{is there any situation where is not allowed mathematically to use continuous compound? or where we get totally different results and we should be careful which one to use?}
\begin{answer}[When the compounding convention matters]
Mathematically nothing forbids continuous compounding. Every positive growth factor $G$ over a year
can be written as $e^{r_c}$ with $r_c=\ln G$. Once rates are converted with the formula above, all
conventions give the \emph{same} prices. Trouble comes from forgetting to convert:
\begin{itemize}
  \item \textbf{the size of the gap grows with $r$ and $t$}: $5\%$ annual is $r_c=4.88\%$, a small
        difference; $50\%$ annual is only $r_c=40.5\%$. Over 30 years at $5\%$, $1.05^{30}=4.32$ but
        $e^{1.5}=4.48$, so quoting one rate in the other convention misprices long bonds by several
        percent;
  \item \textbf{very negative rates}: $1+r_m/m$ must be positive, so a simple rate $r_m\le-m$ is
        meaningless (you would lose more than everything), while $e^{r_ct}$ is always positive;
  \item \textbf{contracts fix the convention}: coupons, money-market deposits and swaps specify
        compounding and day-count rules (\cref{ch:rates}), and a price must follow the contract;
  \item \textbf{continuous time}: in stochastic models (Black--Scholes, short-rate models)
        continuous compounding is the natural choice, because $e^{rt}$ is what $\dd B=rB\dd t$ produces.
\end{itemize}
\end{answer}

\subsection{Discounting by no-arbitrage}\label{sec:discounting}
How much should one pay today for the right to receive \$1 in one year \ts{24}{1-III}{5:03}? Call
the price $x$ and consider the strategy: \emph{borrow $x$ today and buy the claim}.\note{what does it mean buy the claim?} Today the
net cash flow is zero.\note{what is net cash flow?} In one year we receive \$1 and repay $x(1+r)$, so the final profit is
\[
  \Pi = 1 - x(1+r) \quad\text{with certainty.}
\]
If $\Pi>0$ we would do this on a huge scale (riskless profit from nothing); if $\Pi<0$ the
counterparty would do the opposite. Absence of arbitrage forces $\Pi = 0$, i.e.
$x = 1/(1+r)$ \ts{24}{1-III}{6:08}.
\begin{answer}[Claim; net cash flow]
A \emph{claim} is a contract that entitles its holder to a payment: here, the right to receive
\$1 in one year (a zero-coupon bond, \cref{sec:bonds}). \emph{Buying the claim} means paying its price $x$
today to become the one who will be paid. The \emph{net cash flow} at a date is money in minus money
out at that date. Today we receive $+x$ from the loan and pay $-x$ for the claim: net $0$, so the
strategy needs no capital of our own. In one year: $+1$ from the claim, $-x(1+r)$ to repay the loan.
\end{answer}

\begin{definition}[Discount factor]
The present value of \$1 received at time $t$ is the \emph{discount factor} $\Lambda_t$. With a
constant rate,
\[
  \Lambda_n = \frac{1}{(1+r)^n} \ \text{(annual compounding)},\qquad
  \Lambda_t = e^{-rt} \ \text{(continuous compounding)},
\]
and the present value of a cash flow $N$ at time $t$ is $N\Lambda_t$ \ts{24}{1-III}{8:09}.
\end{definition}\note{am i right if i say: given X money now, its value in the future is computet via the interest rate, while its value in the past is computed via the discount factor [we consider time as having an intrinsic value - different if we compound or not]. I AM RIGHT! look at: $X_{now} = X_{future} \cdot e^{-rt}$ and notice you can easily invert it to get $X_{future} = X_{now} \cdot e^{rt}$. [it is even simpler if you allow time to be negative, then you can use the same formula for both past and future]}

\begin{intuition}[No probabilities needed]
The argument never uses probabilities or preferences: only the existence of a riskless strategy
whose profit must vanish. The same logic, applied to a portfolio that is riskless only
\emph{instantaneously} (option + $\Delta$ shares), yields the Black--Scholes equation.
\end{intuition}

% ------------------------------------------------------------------
\section{Bonds}\label{sec:bonds}

\subsection{Zero-coupon and coupon bonds}
A \emph{zero-coupon} (discount) bond pays the notional $N$ at maturity $T$ and nothing else: it is
exactly ``the right to receive $N$ at $T$'' \ts{24}{1-III}{9:09}.  \note{id say that the contract is an obligation, for the other party to pay you in the future: a future/forward contract?} A \emph{coupon bond} pays
periodic coupons $C$ and the notional at maturity; US Treasuries pay coupons semi-annually
\ts{24}{1-III}{10:10}. \note{what about other countries and eu?} The cash flows from the buyer's point of view (the buyer lends $P$) are:

\begin{center}
\begin{tikzpicture}[x=1.05cm,y=0.5cm,>=Stealth]
  \draw[->] (0,0) -- (10.6,0) node[right]{$t$};
  \foreach \k in {1,...,9} \draw[->,thick,defcol] (\k,0) -- (\k,1.4) node[above]{\small $C$};
  \draw[->,thick,defcol] (10,0) -- (10,4) node[above]{\small $C+N$};
  \draw[->,thick,thmcol] (0,0) -- (0,-2.4) node[below]{\small $-P$};
  \foreach \k/\l in {0/0,1/1,2/2,9/n{-}1,10/n} \node[below,font=\scriptsize] at (\k,-0.1) {$\l$};
\end{tikzpicture}
\end{center}

\begin{proposition}[Bond prices at a flat rate]\label{prop:bondprice}
With annual coupons and a constant annual rate $r$,
\begin{align}
  P_{\text{zero}} &= \frac{N}{(1+r)^n},\label{eq:zcb}\\
  P_{\text{coupon}} &= \sum_{k=1}^{n}\frac{C}{(1+r)^k} + \frac{N}{(1+r)^n}
   = C\,\frac{(1+r)^n-1}{r(1+r)^n} + \frac{N}{(1+r)^n}.\label{eq:couponbond}
\end{align}
\end{proposition} \note{credo che il discount rate sulla parte del coupon debba essere divisa per n. no non e' vero, n e' il numero di anni, e noi calcoliamo il valore al tempo zero se riceviamo soldi al tempo k, dato un certo tasso di interesse.}
\begin{proof}
Each cash flow is discounted with its discount factor. With $v = 1/(1+r)$ the coupon part is the
geometric sum $C\sum_{k=1}^n v^k = C\,v\,\frac{1-v^n}{1-v} = C\,\frac{1-(1+r)^{-n}}{r}$, which equals
the stated expression.
\end{proof}

\begin{corollary}[Par bonds]\label{cor:par}
A bond trades \emph{at par} ($P=N$) iff its coupon rate equals the rate: $C = rN$.
\end{corollary} \note{does par stand for parity?}
\begin{proof}
Set $C = cN$ in \eqref{eq:couponbond}:
$P/N = \frac{c}{r}\bigl(1-(1+r)^{-n}\bigr) + (1+r)^{-n} = 1 + \bigl(\tfrac{c}{r}-1\bigr)\bigl(1-(1+r)^{-n}\bigr)$,
which equals $1$ iff $c=r$ (for $n\ge1$). The same identity shows $P>N$ (premium) iff $c>r$
and $P<N$ (discount) iff $c<r$.
\end{proof}

\subsection{Yield}
Prices of bonds with different coupons and maturities are not directly comparable. The market
quotes instead the \emph{yield} \ts{24}{1-III}{11:10}.

\begin{definition}[Yield to maturity]
The yield $y$ of a bond with market price $P$ is the constant discount rate that reproduces the
price: $P = \sum_k \mathrm{CF}_k (1+y)^{-t_k}$. For an $n$-year zero-coupon bond with $N=1$,
\[
  y = P^{-1/n} - 1 = \frac{1}{\sqrt[n]{P}} - 1 .
\]
For coupon bonds the equation is a polynomial in $(1+y)^{-1}$ and is inverted numerically
\ts{24}{1-III}{12:14}.
\end{definition}\note{C is the coupon, what is $F_k$? e' giusta la formula di P? perche' non usiamo continuous framework, sembra molto piu semplice-pulito? coupon non funzionerebbe! aggiungi valori tipici di yueld to maturity (e sopraa anche di coupon, prezzo di un bond, discount rate)}

Since every cash flow is positive, $P(y)$ is strictly decreasing and strictly convex in $y$, so the
yield is unique. Price decreases with yield and increases with coupon
(\cref{fig:priceyield}); by \cref{cor:par} the yield of a par bond equals its coupon.\note{l'idea e' che un bond verra' pagato una certa quantita X il giorno Y, e magari restituisce coupon C. Il problema e' trovare quale sia il prezzo adesso. [per me yield=interest rate=discount rate]}

\begin{figure}[ht]
\centering
\begin{tikzpicture}
\begin{axis}[width=0.8\linewidth,height=6cm,xlabel={yield $y$},ylabel={price of a 10-year bond ($N=1$)},
  domain=0.0005:0.30,samples=120,xmin=0,xmax=0.30,ymin=0,ymax=2.1,
  xtick={0,0.05,0.10,0.15,0.20,0.25,0.30},xticklabels={0\%,5\%,10\%,15\%,20\%,25\%,30\%},
  legend pos=north east,legend style={font=\small},grid=major,grid style={black!10}]
  \addplot[thick,fincol] {0.01*(1-(1+x)^(-10))/x + (1+x)^(-10)}; \addlegendentry{$C=1\%$}
  \addplot[thick,intcol] {0.05*(1-(1+x)^(-10))/x + (1+x)^(-10)}; \addlegendentry{$C=5\%$}
  \addplot[thick,defcol] {0.10*(1-(1+x)^(-10))/x + (1+x)^(-10)}; \addlegendentry{$C=10\%$}
  \addplot[only marks,mark=*,mark size=1.5pt] coordinates {(0.01,1) (0.05,1) (0.10,1)};
\end{axis}
\end{tikzpicture}
\caption{Price--yield curves of 10-year annual-coupon bonds, from \eqref{eq:couponbond}
(re-drawn from slide~22 of \file{lec01_1}). At $y=0$ the price is $1+10C$; the dots mark the
par points $y=C$.}\label{fig:priceyield}
\end{figure}

\subsection{The yield curve}
In reality bonds of different maturities have different yields: the \emph{term structure of
interest rates}, or \emph{yield curve} \ts{24}{1-III}{14:19}. Normally it slopes upward (lenders
want more yield for longer commitments); when short yields exceed long yields the curve is
\emph{inverted}. Inversions preceded the 2008 and 2020 recessions, which is why an inverted curve
is widely (though not infallibly) read as a recession signal \ts{24}{1-III}{16:24}. \note{what is the definition of recession?} The short end
is anchored by the central bank's policy rate \ts{24}{1-II}{14:41}; the long end reflects
expectations, risk premia, supply and demand.

\begin{figure}[ht]
\centering
\begin{tikzpicture}
\begin{axis}[width=0.8\linewidth,height=6.2cm,xlabel={maturity (years)},ylabel={yield (\%)},
  xmin=0,xmax=30,ymin=0,ymax=8,legend pos=south east,legend style={font=\small},
  grid=major,grid style={black!10},mark size=1.2pt]
  \addplot[thick,black!50,mark=*] coordinates {(0.25,2.95) (1,3.2) (5,5.3) (7,5.9) (10,6.35) (30,7.3)};
  \addlegendentry{11 Sep 1992}
  \addplot[thick,cscol,mark=*] coordinates {(0.25,3.95) (0.5,4.2) (1,4.15) (2,3.9) (3,3.95) (5,4.1) (7,4.25) (10,4.4) (20,4.7) (30,4.65)};
  \addlegendentry{12 Sep 2007}
  \addplot[thick,defcol,mark=*] coordinates {(0.25,0.05) (0.5,0.05) (1,0.07) (2,0.21) (3,0.43) (5,0.8) (7,1.1) (10,1.33) (20,1.85) (30,1.92)};
  \addlegendentry{10 Sep 2021}
  \addplot[very thick,fincol,mark=*] coordinates {(0.1,5.3) (0.25,5.2) (0.5,4.8) (1,4.35) (2,3.85) (3,3.75) (5,3.65) (7,3.7) (10,3.8) (20,4.2) (30,4.1)};
  \addlegendentry{3 Sep 2024}
\end{axis}
\end{tikzpicture}
\caption{US Treasury yield curves (values read approximately from slide~23 of \file{lec01_1};
source treasury.gov). 1992 and 2021: normal upward slope, at very different levels; 2007 and
2024: inverted short end.}\label{fig:yieldcurve}
\end{figure}\note{non capisco, se la maturity e' zero allora yield dovrebbe essere zero, no? non e' veramente a zero! perche' un prestito per un anno dovrebbe avere tasso di interesse piu alto che di due? cosi incentivi a comprare bond a scadenza breve, perche'?}\note{vorrei aggiungere yield curves of italian bonds}\note{fai il check con dati ufficiali di queste yield curves e aggiungi quella piu aggiornata di oggi}\note{i found this nice picture to see correlation with fed interest rate (?) https://www.gwkinvest.com/wp-content/uploads/Figure1-GP-June26.svg}

\paragraph{Bootstrapping.} ``Discount factors for all times can be bootstrapped from the yield
curve'' (slide 23). Concretely, suppose we observe par yields $c_1, c_2, \dots$ of annual-coupon
bonds with maturities $1,2,\dots$ years. A par bond with maturity $n$ has price $1$, so
\[
  1 = c_n\sum_{k=1}^{n}\Lambda_k + \Lambda_n
  \qquad\Longrightarrow\qquad
  \Lambda_n = \frac{1 - c_n\sum_{k=1}^{n-1}\Lambda_k}{1+c_n},
\]
which determines $\Lambda_1,\Lambda_2,\ldots$ recursively. Real curve construction (many instrument
types, interpolation, day-count conventions) is the subject of \cref{ch:rates}. \note{i need some more intuition on what is the discount factor and what is the yield, yield curve}

\subsection{Duration and convexity}\label{sec:duration}
A bondholder is exposed to yield changes, and the dependence is non-linear \ts{24}{1-III}{19:29}.
The sensitivities are derivatives, conventionally scaled by the price. \note{so they are dimensionless?}

\begin{definition}[Duration and convexity]
\[
  D = -\frac{1}{P}\frac{\dd P}{\dd y}, \qquad
  \mathcal C = \frac{1}{P}\frac{\dd^2 P}{\dd y^2},
\]
so that, by Taylor expansion,
\begin{equation}\label{eq:durconv}
  \frac{\Delta P}{P} \approx -D\,\Delta y + \tfrac12\,\mathcal C\,(\Delta y)^2 .
\end{equation}
\end{definition}\note{the convention for the signs is beacuse we expect D and C to be positive. look at fig 1.2}

\begin{coursenote}[Erratum in the slides]
On slide 24 of \file{lec01_1} convexity is written with the same letter and sign as duration,
``$D = -\frac1P\frac{d^2P}{dy^2}$''. The standard definition is $\mathcal C = +\frac1P\frac{d^2P}{dy^2}$,
which is positive for bonds with positive cash flows.
\end{coursenote}

\begin{proposition}\label{prop:duration}
\leavevmode
\begin{enumerate}[(i)]
  \item For cash flows $\mathrm{CF}_k$ at times $t_k$ discounted with annual compounding,
        \[ D = \frac{1}{1+y}\,\underbrace{\sum_k t_k\,\frac{\mathrm{CF}_k(1+y)^{-t_k}}{P}}_{D_{\mathrm{Mac}}}, \]
        where the \emph{Macaulay duration} $D_{\mathrm{Mac}}$ is the present-value-weighted average
        time of the cash flows; $D$ is also called \emph{modified duration}.
  \item Zero-coupon bond: $D_Z = \dfrac{T}{1+y}$, \quad $\mathcal C_Z = \dfrac{T(T+1)}{(1+y)^2}$.
  \item ``Continuous'' par bond (coupon paid continuously at rate $c=y$, continuous discounting,
        $N=1$): $D_p = \dfrac{1-e^{-yT}}{y}$.
\end{enumerate}
\end{proposition}
\begin{proof}
(i) Differentiate $P=\sum_k \mathrm{CF}_k(1+y)^{-t_k}$:
$\dd P/\dd y = -\sum_k t_k\,\mathrm{CF}_k(1+y)^{-t_k-1}$.

(ii) For $P=(1+y)^{-T}$: $P' = -T(1+y)^{-T-1}$ and $P'' = T(T+1)(1+y)^{-T-2}$; divide by $P$.

(iii) With coupon rate $c$ fixed, $P(y) = \int_0^T c\,e^{-ys}\dd s + e^{-yT} = c\,\frac{1-e^{-yT}}{y} + e^{-yT}$. Then
\[
P'(y) = c\,\frac{yTe^{-yT} - (1-e^{-yT})}{y^2} - Te^{-yT},
\]
and at $c=y$ (where $P=1$) the terms $\pm Te^{-yT}$ cancel, leaving $P'(y) = -(1-e^{-yT})/y$.
\end{proof}

\begin{intuition}[Duration as a centre of mass]
Macaulay duration is the ``centre of mass'' of the cash flows in time, with present values as
masses. A zero-coupon bond has all its mass at $T$; coupons pull the centre of mass earlier, so a
coupon bond has $D_{\mathrm{Mac}} < T$. For the continuous par bond $D_p = (1-e^{-yT})/y < T$ and
$D_p \to 1/y$ as $T\to\infty$: even a perpetual bond has finite duration because distant cash flows
weigh almost nothing.
\end{intuition}

\begin{example}[How good is the quadratic approximation?]
A 10-year zero at $y = 4\%$: $P = 1.04^{-10} = 0.6756$, $D = 10/1.04 = 9.615$,
$\mathcal C = 110/1.04^2 = 101.7$. If the yield rises by $\Delta y = 1\%$:
first order $-9.62\%$; adding convexity $+\tfrac12\cdot 101.7\cdot 0.0001 = +0.51\%$ gives $-9.11\%$;
the exact change is $(1.04/1.05)^{10}-1 = -9.13\%$.
\end{example}\note{we are still considering the yield to be fixed along the duration of the bond maturity, right?}

\begin{finance}[Why convexity is valuable]
Because $\mathcal C>0$, the price rises more when yields fall than it drops when yields rise by the
same amount. Two portfolios with the same price and duration but different convexity are
therefore not equivalent: the more convex one gains from large rate moves in either direction.
Trading desks hedge duration (first order, like delta) and monitor convexity (second order, like
gamma).
\end{finance}\note{first appearence of the Greeks, right? use the symbols $\Delta$, $\Gamma$, etc.}
\note{in general, i'd like to have graphs and conceptual maps when possible}

\subsection{Forward rates}\label{sec:forward}
\begin{definition}[Forward rate agreement]
A \emph{forward rate agreement} (FRA) fixes today the interest $F_{tT}$ to be paid at time $T$ on a
loan of $N$ that starts at a future date $t<T$ (slide 25): the lender pays $N$ at $t$ and receives
$N(1+F_{tT})$ at $T$.
\end{definition}

\begin{proposition}\label{prop:forward}
In the absence of arbitrage, $F_{tT} = \dfrac{\Lambda_t}{\Lambda_T} - 1$.
\end{proposition}\note{reminder: price now = Lambda payoff at maturity. Lambda serve a trasportare il valore lungo il tempo - nel passato}
\begin{proof}
Entering an FRA costs nothing today, so the present value of its cash flows must be zero:
$-N\Lambda_t + N(1+F_{tT})\Lambda_T = 0$. Equivalently (the form on the slide),
$1 = \frac{1}{\Lambda_t}(1+F_{tT})\Lambda_T$. If $F_{tT}$ were larger, one could lock in a riskless
profit by borrowing $N\Lambda_T(1+F_{tT})$ until $T$ \dots\ (\cref{ex:fra}).
\end{proof}

\begin{remark}
$F_{tT}$ is the interest over the whole period $[t,T]$; the annualised simple rate is
$F_{tT}/(T-t)$. In continuous time one defines the \emph{instantaneous forward rate}
$f(0,T) = -\partial_T \ln \Lambda_T$, so that $\Lambda_T = \exp\bigl(-\int_0^T f(0,s)\dd s\bigr)$: the
continuously compounded zero yield $-\frac1T\ln\Lambda_T$ is the \emph{average} of the forward rates
up to $T$. Modelling the dynamics of the whole curve $f(t,T)$ is the HJM framework (\cref{ch:hjm}).
\end{remark}\note{it feels that f plays the role of r, the interest rate. and Lambda is the correspective of the compount rate, right?}

% ------------------------------------------------------------------
\section{Working with data: R and RStudio}\label{sec:R}

\begin{casestudy}[\file{FM_Intro1.Rmd} --- collecting and plotting financial time series]
\textbf{Goal.} Show how easily financial data can be downloaded and displayed
\ts{24}{1-I}{10:22}. \textbf{Tools.} R packages \texttt{tidyquant} (downloads), \texttt{tidyverse}
(data handling), \texttt{fpp2} (forecasting); later lectures use \texttt{quantmod}, \texttt{xts}/\texttt{zoo}
(time series), \texttt{car} (regression), \texttt{rugarch} (GARCH), \texttt{fxregime} (regimes).
\textbf{Data (2011--2024).} From Yahoo Finance: S\&P~500 and VIX indices, stocks NVDA, GE, AAPL, GOOG,
AMZN, XOM, GME, AMC, the ARKK and XLF ETFs, BTC-USD. From FRED (St.\,Louis Fed, about 80{,}000
series): Treasury constant-maturity yields 3m, 1y, 5y, 10y (\texttt{DGS3MO}\dots\texttt{DGS10}),
Moody's Aaa/Baa corporate yields (\texttt{DAAA}, \texttt{DBAA}), WTI crude oil (\texttt{DCOILWTICO}),
Coinbase Bitcoin. \textbf{Steps.} Download, keep \texttt{(symbol, date, price)}, drop missing values,
merge, plot each series, save to \file{data_fm_intro1.Rdata} for later scripts.

\smallskip
\textbf{What to look at} \ts{24}{1-I}{12:25}: the VIX (``fear gauge'', the market's implied
volatility of the S\&P~500) spikes when the index drops, a relationship formalised by option
pricing (\cref{ch:volatility,ch:bs}); NVIDIA and Bitcoin raise the question of detecting bubbles;
WTI goes negative in April 2020; the Baa--Treasury spread is the credit spread modelled in
\cref{ch:regression}.
\end{casestudy}

\begin{figure}[ht]
\centering
\pgfplotsset{fmts/.style={width=0.9\linewidth,height=4.2cm,xmin=2011,xmax=2024.75,
  xtick={2012,2014,...,2024},xticklabel style={/pgf/number format/1000 sep={}},
  grid=major,grid style={black!10},every axis plot/.append style={thin}}}
\begin{tikzpicture}
\begin{axis}[fmts,ylabel={S\&P 500},xticklabels={},name=top]
  \addplot[defcol] table {data/fm_sp500.dat};
\end{axis}
\begin{axis}[fmts,ylabel={VIX},xlabel={year},at={(top.below south west)},anchor=north west,yshift=-4pt]
  \addplot[thmcol] table {data/fm_vix.dat};
\end{axis}
\end{tikzpicture}
\caption{S\&P~500 index and VIX, weekly closes 2011--2024, re-drawn from the data of
\file{FM_Intro1.Rmd}. Every VIX spike (2011, 2015, 2018, 2020, 2022) coincides with a drop of the index; the record daily close,
82.7, was on 16 March 2020.}\label{fig:spvix}
\end{figure}

\begin{figure}[ht]
\centering
\begin{tikzpicture}
\begin{semilogyaxis}[width=0.9\linewidth,height=5.5cm,xmin=2011,xmax=2024.75,
  xtick={2012,2014,...,2024},xticklabel style={/pgf/number format/1000 sep={}},xlabel={year},
  ylabel={price (\$, log scale)},grid=major,grid style={black!10},legend pos=south east,
  legend style={font=\small}]
  \addplot[thin,intcol] table {data/fm_nvda.dat}; \addlegendentry{NVIDIA (split-adjusted)}
  \addplot[thin,fincol] table {data/fm_btc.dat}; \addlegendentry{Bitcoin (BTC-USD)}
\end{semilogyaxis}
\end{tikzpicture}
\caption{NVIDIA and Bitcoin, weekly closes, from the data of \file{FM_Intro1.Rmd}. On a log scale
a straight line is constant exponential growth: NVIDIA rose about 500-fold in 2012--2024, the
kind of run that raises the bubble question \ts{24}{1-I}{14:36}.}\label{fig:nvdabtc}
\end{figure}

The notebook runs on RStudio Cloud (\url{https://posit.cloud}, free account, nothing to install) or
on desktop RStudio (faster, works offline); \file{lec01_2} gives step-by-step instructions, summarised
in \cref{app:rsetup}.

% ------------------------------------------------------------------
\section{Two examples of quant work (2013)}\label{sec:projects}
V.~Strela showed two projects done by a student intern at Morgan Stanley \ts{13}{1}{54:30}.

\paragraph{Noisy delta.} Prices of complex derivatives are often computed by Monte Carlo, so the
pricing function $\hat f(x)$ is only known up to statistical noise of standard deviation
$\operatorname{sd}(x)$. Delta is computed by a centred difference
$\hat\Delta_\varepsilon = [\hat f(x+\varepsilon)-\hat f(x-\varepsilon)]/(2\varepsilon)$. A small shift
$\varepsilon$ amplifies the noise, a large one introduces discretisation bias. The project chose
$\varepsilon$ to minimise the mean squared error $\E\bigl|f'(x)-\hat\Delta_\varepsilon\bigr|^2$ and found
a rule of the form $\varepsilon_{\text{opt}} \propto \bigl(\operatorname{sd}(x)/(d+|\Gamma(x)|)\bigr)^{1/3}$,
with $\Gamma$ the option's gamma and $c,d$ tuning constants (slide 19 of \file{lecnote1}); the resulting
delta is visibly smoother than with a constant shift.

\begin{intuition}[Where the cube root comes from]
Taylor-expanding, the centred difference has bias $\frac{\varepsilon^2}{6}f'''(x)$; with independent
noise at the two points its variance is $2\operatorname{sd}^2/(2\varepsilon)^2 = \operatorname{sd}^2/(2\varepsilon^2)$.
Hence
\begin{gather*}
  \mathrm{MSE}(\varepsilon) \approx \frac{\varepsilon^4 f'''(x)^2}{36} + \frac{\operatorname{sd}^2}{2\varepsilon^2},\\
  \frac{\dd\,\mathrm{MSE}}{\dd\varepsilon}=0 \iff \varepsilon^6 = \frac{9\,\operatorname{sd}^2}{f'''(x)^2}
  \iff \varepsilon_{\text{opt}} = \Bigl(\frac{3\operatorname{sd}}{|f'''(x)|}\Bigr)^{1/3}.
\end{gather*}
The same bias--variance trade-off appears in numerical differentiation of noisy experimental data.
The course's formula uses gamma as a proxy for the local curvature scale and regularises it with $d$;
the exact derivation is not in the material.
\end{intuition}

\paragraph{Kalman filter for FX quotes.} Five brokers quoted forward points for USD/RUB at random,
non-uniform times; the aim was to combine these noisy observations into a better estimate of the true
price with a Kalman filter \ts{13}{1}{57:41}. State-space models and the Kalman filter are part of
time-series analysis (\cref{ch:timeseries}).

% ------------------------------------------------------------------
\begin{keyformulas}
\begin{align*}
&\text{Compounding:} && (1+r)^n,\quad (1+r/m)^{mn},\quad e^{rt}\\
&\text{Discount factor:} && \Lambda_n=(1+r)^{-n},\quad \Lambda_t=e^{-rt}\\
&\text{Coupon bond:} && P = C\,\tfrac{1-(1+r)^{-n}}{r} + N(1+r)^{-n};\quad P=N \iff C=rN\\
&\text{Zero yield:} && y = P^{-1/n}-1\\
&\text{Bootstrapping:} && \Lambda_n = \bigl(1-c_n\textstyle\sum_{k<n}\Lambda_k\bigr)/(1+c_n)\\
&\text{Duration, convexity:} && D=-P'/P,\ \ \mathcal C=P''/P,\ \ \Delta P/P\approx -D\Delta y+\tfrac12\mathcal C\Delta y^2\\
&\text{Zero-coupon:} && D_Z=T/(1+y),\ \ \mathcal C_Z=T(T+1)/(1+y)^2\\
&\text{Continuous par bond:} && D_p=(1-e^{-yT})/y\\
&\text{Forward rate:} && F_{tT}=\Lambda_t/\Lambda_T-1\\
&\text{Alpha, beta:} && R_a=\alpha+\beta R_b+\varepsilon
\end{align*}
\end{keyformulas}\note{forse aiuterebbe aggiungere le unita' di misura, tipo $T^{-1}$ per r }

% ------------------------------------------------------------------
\section*{Exercises}
\addcontentsline{toc}{section}{Exercises}
No problem set covers this chapter; \cref{ex:game,ex:glossary} are the course's own activities,
the others are added for practice (marked \emph{not from the course}).

\begin{exercise}[2024 Investment Game (ungraded)]\label{ex:game}
With a hypothetical \$10{,}000, buy a single publicly traded stock, ETF, currency or bond (to be short,
buy an inverse ETF) at the close of Friday 13 September 2024 and sell at the close of Friday
15 November 2024, with one optional switch of the whole balance at the close of 11 October
\ts{24}{1-II}{32:28}. Track in a spreadsheet the number of shares, the daily P\&L in \$ and \%, the
total P\&L, the sum $G$ of the daily gains and the sum $L$ of the daily losses (in absolute value), and
the ratio $(G-L)/(G+L)$. Repeat the game on any two-month window of your choice.
\begin{enumerate}[(a)]
  \item Show that $G-L$ is the total P\&L and that $(G-L)/(G+L)\in[-1,1]$.
  \item Interpret the ratio: which values correspond to a ``clean'' trend, which to a noisy
        path with the same final P\&L?
\end{enumerate}
\end{exercise}
\note{a) easy to se, imposta le due disuguaglianze, sono ripettate trivialmente. b) il rateo e' profitto totale didiso per il valore assoluto di tutte le variazioni, quindi grandi variazioni implicano piccolo rateo e piccole variazioni grande rateo. inoltre, grande profitto implica grande rateo, piccolo profitto piccolo rateo. quindi questa e' una metrica di quanto profittevole e con bassa volatilita linvestimento e' stato}

\begin{exercise}[2013, optional homework]\label{ex:glossary}
Build your own financial glossary: list every term of this chapter you could not define precisely and
look it up \ts{13}{1}{53:27}.
\end{exercise}


\begin{exercise}[not from the course]
A 5-year bond pays annual coupons of 6\% on $N=100$.
\begin{enumerate}[(a)]
  \item Price it at $y=5\%$ and $y=7\%$.
  \item Compute its
      Macaulay and modified duration at $y=5\%$ and check them against the price change for $\Delta y=\pm 1$\,bp.
\end{enumerate}
\end{exercise}\note{$a. P_5 = 0.06*-5,526 + 100*1.276 = 127.3 ; P_5 = 0.06*-5,751 + 100*1.402 = 139.8 ; b. D_Mac = 0.06*100*(1*1.05^(-1) + 2*1.05^(-2) + 3*1.05^(-3) + 4*1.05^(-4) + 5*1.05^(-5)) / 127.3 = 0.592 ; D_mod = 0.592 / 1.05 = 0.564 ; TODO $}

\begin{exercise}[not from the course]
Par yields are $c_1 = 4\%$, $c_2 = 4.5\%$, $c_3 = 5\%$ (annual coupons). Bootstrap $\Lambda_1,\Lambda_2,\Lambda_3$,
the zero yields and the one-year forward rates $F_{1,2}$, $F_{2,3}$. Why is the zero curve above the par curve
when the curve slopes upward?
\end{exercise}

\begin{exercise}[not from the course]\label{ex:fra}
Complete the proof of \cref{prop:forward}: if the market quotes an FRA rate $F > \Lambda_t/\Lambda_T - 1$, write
down an explicit zero-cost strategy with zero-coupon bonds and the FRA that has a positive payoff at $T$.
\end{exercise}

\begin{exercise}[not from the course]
Compute the convexity of the continuous par bond of \cref{prop:duration}(iii) and its limit as $T\to\infty$.
\end{exercise}

% !TEX root = ../main.tex
\chapter{Linear Algebra and One-Period Financial Models}\label{ch:linalg}

\begin{sources}
\small
\begin{tabularx}{\linewidth}{@{}l l L@{}}
\toprule
\textbf{Lecture} & \textbf{Speaker} & \textbf{Content}\\
\midrule
2024 L2 & P.~Kempthorne & Vectors as portfolios, P\&L, arbitrage; matrices; stochastic matrices and Markov chains; single-period market model, completeness, pricing measure; linear systems; eigenvalues and diagonalisation.\\
2024 L4 (first 28 min) & P.~Kempthorne & Diagonalisation and state equations, symmetric matrices, SVD, Perron--Frobenius; equal-weighted portfolio case study and rebalancing.\\
2013 L2 & C.~Lee & Eigenvalues, symmetric matrices (proofs), SVD with a worked example, Perron--Frobenius and its link to the Ross recovery theorem.\\
\bottomrule
\end{tabularx}

\smallskip
\textbf{Files.} Slides \file{mit18_642_f24_lec02_1.pdf}; lecture note \emph{One-Period Financial Models}
\file{mit18_642_f24_lec02_2.pdf}; \file{mit18_642_f24_lec02_3.pdf} with
\file{mit18_642_f24_equal_weighted_portfolios.zip} (\file{Script_SP500.Rmd},
\file{Equal_Weighted_Portfolios.Rmd}, \file{SP500_data_subset.RData}); 2013 notes
\file{MIT18_S096F13_lecnote2.pdf}.
\textbf{Problem sets.} 2024 PS1 (all problems); 2013 PS1 (all problems).
\textbf{Not recorded.} 2024 L3 and 2013 L4, see \cref{sec:la-missing}.
\end{sources}

\section*{Motivation}
Linear algebra is the language in which portfolios, scenarios and covariance structures are
written. Three ideas of this chapter recur for the rest of the course:
\begin{itemize}
  \item a portfolio is a vector and its value is a dot product, so P\&L is linear in holdings;
  \item in a market with finitely many scenarios, \emph{no arbitrage} is equivalent to the
        existence of positive \emph{state prices}, and derivative prices are discounted expectations
        under a ``risk-neutral'' measure (\cref{sec:oneperiod}): this is the discrete skeleton of
        Black--Scholes;
  \item spectral decompositions (eigen-decomposition, SVD, Perron--Frobenius) describe long-run
        behaviour of Markov chains and the dominant directions of co-movement in data (PCA,
        \cref{ch:pca}).
\end{itemize}
The basic material (vector spaces, rank, determinants) is only recalled; a physicist can skim
\cref{sec:la-matrices} and focus on \cref{sec:la-portfolios,sec:markov,sec:oneperiod,sec:pf}.

% ------------------------------------------------------------------
\section{Vectors as portfolios}\label{sec:la-portfolios}

Let $\bm p(t)\in\R^{501}_{+}$ collect the closing prices on day $t$ of the 500 stocks of the S\&P~500
plus \emph{cash} as asset $0$ with $p_0(t)\equiv 1$ \ts{24}{2}{3:21}. A portfolio is a vector of holdings
$\bm q(t)$ (shares, and dollars of cash), and its value is the dot product
\[
  V_t = \bm q\T\bm p(t) = \sum_{j=0}^{500} q_j\,p_j(t).
\]

\paragraph{Timing and rebalancing.} The holdings $\bm q(t+1)$ carried through day $t+1$ are chosen at
the end of day $t$, using only information available then \ts{24}{2}{8:43}. Rebalancing by trading
$\delta_j(t)$ shares, $q_j(t+1) = q_j(t)+\delta_j(t)$, is \emph{self-financing} (no money added or
withdrawn) if $\sum_j \delta_j(t)\,p_j(t) = 0$ \ts{24}{2}{4:26}. Then $V_t = \bm q(t+1)\T\bm p(t)$ and
the daily P\&L is
\begin{equation}\label{eq:pnl}
  \mathrm{PnL}(t+1) = V_{t+1} - V_t = \bm q(t+1)\T\bigl[\bm p(t+1)-\bm p(t)\bigr],
\end{equation}
where the cash term drops out because $p_0$ does not change.

\begin{intuition}[Left-point sums]
Summing \eqref{eq:pnl} over days, the total gain is $\sum_t \bm q(t+1)\T\Delta\bm p(t)$: holdings decided
\emph{before} the price move multiply the price increment. This ``evaluate the integrand at the left end
point'' rule is exactly the discrete version of the Itô integral $\int \theta_t\,\dd S_t$ of
\cref{ch:stochcalc}; using $\bm q$ at the right end point would mean trading on tomorrow's prices.
\end{intuition}

\paragraph{Long/short and zero-cost portfolios.} Holdings may be negative. To \emph{sell short}
means to borrow shares through the broker (who must ``locate'' them), sell them now and buy
them back later \ts{24}{2}{11:44}; the sale proceeds are credited to the (margin) account, which allows long positions
worth more than 100\% of one's capital. For two portfolios $\bm q,\bm w$ the difference $\bm d=\bm q-\bm w$
is a long/short portfolio with $\mathrm{PnL}(\bm d)=\mathrm{PnL}(\bm q)-\mathrm{PnL}(\bm w)$.
A \emph{zero-cost} portfolio has $\bm d\T\bm p(t-1)=0$; an \emph{arbitrage} would be a zero-cost
portfolio whose P\&L is positive \ts{24}{2}{13:45}. The catch: at the time of the decision
$\bm p(t)$ is random, so the P\&L is random. A pure arbitrage requires a model of what $\bm p(t)$ can
be (\cref{sec:oneperiod}); strategies whose random P\&L merely tends to be positive are called
\emph{statistical arbitrage} \ts{24}{2}{16:49}.

\paragraph{Geometry.} $\|\bm v\| = \sqrt{\bm v\T\bm v}$, $\bm v\T\bm w = \|\bm v\|\|\bm w\|\cos\theta$;
$\bm v\T\bm w/\|\bm w\|$ is the length of the projection of $\bm v$ on the direction of $\bm w$
\ts{24}{2}{17:50}. Orthogonal projections are the geometry of least squares (\cref{ch:regression}).

% ------------------------------------------------------------------
\section{Matrices: a quick recap}\label{sec:la-matrices}
For $A=[\bm a_1\cdots\bm a_n]\in\R^{m\times n}$ the two readings of $A\bm v$ are
\begin{align*}
  A\bm v &= \textstyle\sum_{j=1}^n v_j\,\bm a_j && \text{(combination of columns)},\\
  (A\bm v)_i &= \operatorname{row}_i(A)\,\bm v && \text{(dot products with rows)};
\end{align*}
the first is the more useful one \ts{24}{2}{23:03}. Similarly $AB=[A\bm b_1\cdots A\bm b_p]$.
Remember $(AB)\T=B\T A\T$ and $AB\neq BA$ in general. Notation used later: $I_n=[\bm e_1\cdots\bm e_n]$,
$\bm 1$ the vector of ones, $J=\bm 1\bm 1\T$ the matrix of ones, $\diag(d_1,\dots,d_n)$.

\paragraph{Linear systems} \ts{24}{2}{1:14:25}. For $A\bm x=\bm b$ with $A\in\R^{m\times n}$: if
$m=n=\rank A$ then $\det A\ne0$ and $\bm x=A^{-1}\bm b$; otherwise the system is under-determined
($m<n$), over-determined ($m>n$, solved in the least-squares sense in \cref{ch:regression}) or
rank-deficient. Solvability: $\bm b\in\operatorname{col}(A)$; uniqueness: $\ker A=\{0\}$. These two
facts are exactly what decides \emph{replicability} of derivatives in \cref{sec:oneperiod}.

% ------------------------------------------------------------------
\section{Stochastic matrices and Markov chains}\label{sec:markov}

\begin{definition}[Stochastic matrix, Markov chain]
A square matrix $A$ with $a_{ij}\ge0$ whose \emph{columns} sum to $1$ is (column-)stochastic
\ts{24}{2}{29:29}. A Markov chain on states $\{1,\dots,m\}$ with stationary transition probabilities
$a_{ij}=\Prob(S_{t+1}=i\mid S_t=j)$ propagates the distribution
$\bm\pi(t)=(\Prob(S_t=1),\dots,\Prob(S_t=m))\T$ as
\[
  \bm\pi(t+1) = A\,\bm\pi(t),\qquad \bm\pi(t) = A^t\bm\pi(0).
\]
A \emph{stationary distribution} is a probability vector with $\bm\pi^*=A\bm\pi^*$, i.e.\ an
eigenvector of eigenvalue $1$ \ts{24}{2}{35:53}.
\end{definition}

\begin{coursenote}[Conventions]
The course uses column-stochastic matrices acting on column vectors. Many probability texts use the
transpose (row-stochastic $P$, row vectors, $\bm\pi\T(t+1)=\bm\pi\T(t)P$). Nothing changes except
transposes.
\end{coursenote}

\begin{proposition}\label{prop:stochastic}
If $A$ is column-stochastic then $1$ is an eigenvalue and every eigenvalue satisfies $|\lambda|\le1$.
\end{proposition}
\begin{proof}
$\bm 1\T A=\bm 1\T$, so $1$ is an eigenvalue of $A\T$, hence of $A$ (same characteristic polynomial).
If $A\bm v=\lambda\bm v$ then $|\lambda|\sum_i|v_i| = \sum_i|\sum_j a_{ij}v_j|\le\sum_j|v_j|\sum_i a_{ij}
= \sum_j|v_j|$.
\end{proof}

A stationary distribution need not be the limit of $\bm\pi(t)$: for $A=\bigl(\begin{smallmatrix}0&1\\1&0\end{smallmatrix}\bigr)$
the chain alternates deterministically and $\bm\pi(t)$ oscillates (a periodic chain, eigenvalue $-1$)
\ts{24}{2}{36:53}. If some power $A^k$ has all entries positive, the chain ``mixes'': by the
Perron--Frobenius theorem (\cref{sec:pf}) all other eigenvalues have $|\lambda|<1$ and
$\bm\pi(t)\to\bm\pi^*$ from any start \ts{24}{2}{37:54}.

\begin{example}[Two-state chain]\label{ex:twostate}
Let $p=\Prob(1\to2)$, $q=\Prob(2\to1)$, $0<p,q<1$:
$A=\bigl(\begin{smallmatrix}1-p & q\\ p & 1-q\end{smallmatrix}\bigr)$. From $\tr A=2-p-q$ and
$\det A=1-p-q$ the eigenvalues are $1$ and $\lambda_2=1-p-q$, with eigenvectors
$\bm\pi^*=\frac{1}{p+q}(q,p)\T$ and $(1,-1)\T$. Since $\bm\pi(0)-\bm\pi^*$ has zero sum it is a multiple of
$(1,-1)\T$, hence
\[
  \bm\pi(t) = \bm\pi^* + (1-p-q)^t\,\bigl(\bm\pi(0)-\bm\pi^*\bigr).
\]
Convergence is geometric with relaxation time $\tau=-1/\ln|1-p-q|$: the spectral gap $1-|\lambda_2|$
controls how fast the chain forgets its initial state.
\end{example}

\begin{finance}[Order-flow as a Markov chain]
In PS1 (2024) the side (buy/sell) of successive market orders is modelled as a two-state chain:
buys tend to follow buys and sells follow sells. The stationary distribution gives the long-run
fraction of buy orders; $\lambda_2$ measures the persistence of order flow \ts{24}{2}{38:56}.
\end{finance}

% ------------------------------------------------------------------
\section{The one-period market model}\label{sec:oneperiod}
This section follows the lecture note \emph{One-Period Financial Models} (\file{lec02_2}) and
\ts{24}{2}{41:03}--\ts{24}{2}{1:13:22}. It is the most important part of the chapter.

\subsection{Two assets, two states}\label{sec:binomial}
One period $[0,T]$ and two assets \ts{24}{2}{43:17}:
\begin{itemize}
  \item a \emph{bond} (risk-free asset) with $B_T = B_0(1+r_fT)$ known in advance;
  \item a \emph{stock} with known $S_0$ and random $S_T\in\{S_T^d, S_T^u\}$, $S_T^d<S_T^u$, depending on
        the state $\omega\in\{d,u\}$ of the economy at $T$.
\end{itemize}
No probabilities of $u$ and $d$ are specified: they will turn out to be irrelevant for pricing.
A portfolio $\bm\theta=(\theta_B,\theta_S)\T$ costs $V_0=\theta_BB_0+\theta_SS_0$ and pays
$V_T(\omega)=\theta_BB_T+\theta_SS_T(\omega)$.

A \emph{contingent claim} $C$ pays $C_T^u$ or $C_T^d$ at $T$. A portfolio \emph{replicates} $C$ if
\begin{equation}\label{eq:replic2}
  \begin{pmatrix} B_T & S_T^d\\ B_T & S_T^u\end{pmatrix}
  \begin{pmatrix}\theta_B\\ \theta_S\end{pmatrix}
  = \begin{pmatrix}C_T^d\\ C_T^u\end{pmatrix}.
\end{equation}
The determinant is $B_T(S_T^u-S_T^d)\neq0$, so every claim is replicable, uniquely:
\begin{equation}\label{eq:delta}
  \theta_S = \frac{C_T^u-C_T^d}{S_T^u-S_T^d},\qquad
  \theta_B = \frac{C_T^dS_T^u-C_T^uS_T^d}{B_T\,(S_T^u-S_T^d)} .
\end{equation}
$\theta_S$ is a finite-difference \emph{delta}, the hedge ratio of \cref{sec:tradertypes}. If the claim
trades at a price different from the cost of its replicating portfolio, buying the cheap one and
selling the expensive one is an arbitrage; so $C_0 = \theta_BB_0+\theta_SS_0$. Substituting
\eqref{eq:delta} and rearranging:
\begin{equation}\label{eq:rnbinomial}
  C_0 = \frac{1}{1+r_fT}\bigl[q_u\,C_T^u + q_d\,C_T^d\bigr],\qquad
  q_u = \frac{(1+r_fT)S_0-S_T^d}{S_T^u-S_T^d},\quad q_d = 1-q_u .
\end{equation}

\begin{proposition}[No arbitrage in the binomial model]\label{prop:binomial-na}
The two-asset model is arbitrage-free iff $S_T^d < S_0(1+r_fT) < S_T^u$, i.e.\ iff $0<q_u<1$.
\end{proposition}
The ``only if'' part is PS1 Problem~5 (\cref{ex:ps1-5}); both directions follow from the general
theorem below.

\begin{example}[Pricing a call]\label{ex:call}
Take the numbers of the lecture note: $B_0=100$, $B_T=103$ ($r_fT=3\%$), $S_0=100$, $S_T^u=130$,
$S_T^d=90$, and a call with strike $K=100$: $C_T^u=30$, $C_T^d=0$. Then
$\theta_S = 30/40 = 0.75$ shares, $\theta_B = -0.75\cdot 90/103 = -0.6553$ bonds, and
$C_0 = -65.53+75 = 9.466$. Equivalently $q_u = (103-90)/40 = 0.325$ and
$C_0 = 0.325\cdot 30/1.03 = 9.466$. Note that $q_u$ has nothing to do with how likely the up move
really is.
\end{example}

\begin{figure}[ht]
\centering
\begin{tikzpicture}[scale=0.036,>=Stealth]
  \draw[->,black!50] (-80,0) -- (150,0) node[right,black]{\small payoff if $d$};
  \draw[->,black!50] (0,-80) -- (0,150) node[above,black]{\small payoff if $u$};
  \draw[->,very thick,defcol] (0,0) -- (103,103) node[above right]{\small bond $(103,103)$};
  \draw[->,very thick,intcol] (0,0) -- (90,130) node[above]{\small stock $(90,130)$};
  \draw[->,thick,intcol,dashed] (0,0) -- (67.5,97.5) node[left,font=\scriptsize]{$0.75\cdot$stock\ };
  \draw[->,thick,defcol,dashed] (67.5,97.5) -- (0,30);
  \node[font=\scriptsize,defcol,anchor=east] at (30,68) {$-0.655\cdot$bond};
  \draw[->,very thick,thmcol] (0,0) -- (0,30) node[left]{\small call $(0,30)$};
  \draw[black!30] (-75,-75) -- (140,140);
  \node[font=\scriptsize,black!50,rotate=45] at (-45,-38) {riskless payoffs};
\end{tikzpicture}
\caption{The two-state model as a plane of payoff vectors (x: state $d$, y: state $u$), as in
Figure~1 of \file{lec02_2}. Bond and stock span $\R^2$, so every claim, e.g.\ the call of
\cref{ex:call}, is a combination of them \ts{24}{2}{55:18}.}\label{fig:payoffplane}
\end{figure}

\begin{intuition}[Why ``risk-neutral'']
From \eqref{eq:rnbinomial} applied to the stock itself,
$q_uS_T^u+q_dS_T^d=(1+r_fT)S_0$: under $Q=(q_u,q_d)$ \emph{every} asset earns the risk-free rate on
average, as it would in a world of risk-neutral investors. The real-world probabilities and
investors' risk preferences are already encoded in the \emph{current price} $S_0$; given $S_0$,
replication makes them irrelevant for the claim. This is why option prices do not depend on the
stock's expected return, a fact that will reappear as the absence of the drift $\mu$ from the
Black--Scholes formula.
\end{intuition}

\subsection{\texorpdfstring{$n$}{n} assets, \texorpdfstring{$m$}{m} states}
Generalise to $n$ assets and $m$ possible states $\omega_1,\dots,\omega_m$ at time $T$
\ts{24}{2}{1:06:57}:
\begin{itemize}
  \item $A_0=(P_0^1,\dots,P_0^n)\in\R^{1\times n}$: known prices at $t=0$;
  \item $A\in\R^{m\times n}$ with $A_{ij}=P_T^j(\omega_i)$: price of asset $j$ in state $i$. Column $j$ is the
        payoff vector of asset $j$, row $i$ the prices in scenario $i$;
  \item a portfolio $\bm\theta\in\R^n$ costs $V_0=A_0\bm\theta$ and pays $\bm V_T=A\bm\theta\in\R^m$.
\end{itemize}
Implicit assumptions: prices do not depend on quantities traded (unlimited liquidity), holdings
can be fractional, short selling is unrestricted, no transaction costs \ts{24}{2}{1:05:55}.

\begin{coursenote}[Notation]
The lecture note writes the portfolio as $\vec\pi$; here it is $\bm\theta$, to keep $\pi$ for
probability vectors.
\end{coursenote}

\begin{definition}[Replication, completeness, arbitrage, state prices]\label{def:onepmodel}
In the model $(A_0,A)$:
\begin{itemize}
  \item A claim with payoff $\bm C_T\in\R^m$ is \emph{replicable} if $A\bm\theta=\bm C_T$ for some $\bm\theta$.
        The market is \emph{complete} if every claim is replicable.
  \item An \emph{arbitrage} is a portfolio with $V_0\le 0$, $\bm V_T\ge0$ (componentwise) and
        $(V_0,\bm V_T)\neq(0,\bm 0)$: it never loses and either costs nothing and may gain, or pays you
        today and never costs anything later \ts{24}{2}{1:09:06}.
  \item A vector $\bm\psi\in\R^m$ with all $\psi_i>0$ is a vector of \emph{state prices} if
        $A_0=\bm\psi\T A$, i.e.\ $P_0^j=\sum_i\psi_iP_T^j(\omega_i)$ for every asset. Writing
        $\beta=\sum_i\psi_i$ and $q_i=\psi_i/\beta$, this reads
        \[
          P_0^j = \beta\,\E_Q\bigl[P_T^j\bigr],\qquad j=1,\dots,n,
        \]
        and $Q=(q_1,\dots,q_m)$ is a \emph{pricing measure} with discount factor $\beta$ \ts{24}{2}{1:11:16}.
\end{itemize}
\end{definition}

\begin{coursenote}[Definition of arbitrage]
The lecture note requires $V_0\le0$ and $\bm V_T\ge0$ with $V_T(\omega_i)>0$ for at least one state,
which omits the case ``negative cost, zero payoff''. When a riskless asset exists the two definitions
are equivalent (invest the negative cost in the bond); the version above makes the theorem below
true without that assumption.
\end{coursenote}

$\psi_i$ is the price of the \emph{Arrow--Debreu security} $\bm e_i$ paying \$1 in state $i$ and nothing
otherwise (when it is replicable). If a riskless bond exists, $B_0=\bm\psi\T\bm 1\,B_T$ gives
$\beta=1/(1+r_fT)$, the ordinary discount factor.

\begin{lemma}[Stiemke]\label{lem:stiemke}
Let $L\subset\R^k$ be a linear subspace. Then $L\cap\R^k_{+}=\{\bm 0\}$ iff there exists $\bm y\in\R^k$ with
all $y_i>0$ and $\bm y\perp L$.
\end{lemma}
\begin{proof}
($\Leftarrow$) If $\bm x\in L$, $\bm x\ge0$, $\bm x\ne0$, then $\bm y\T\bm x>0$, contradicting $\bm y\perp L$.
($\Rightarrow$) The simplex $\Delta=\{\bm x\ge0:\sum_ix_i=1\}$ is compact, convex and disjoint from the closed
convex set $L$. By the separating hyperplane theorem there are $\bm y$ and $c$ with $\bm y\T\bm l<c<\bm y\T\bm x$
for all $\bm l\in L$, $\bm x\in\Delta$. Since $L$ is a subspace, $\bm y\T\bm l$ bounded above on $L$ forces
$\bm y\T\bm l=0$; then $c>0$ and $y_i=\bm y\T\bm e_i>0$ for every $i$.
\end{proof}

\begin{theorem}[Fundamental theorem of asset pricing, one period]\label{thm:ftap}
In the one-period model $(A_0,A)$:
\begin{enumerate}[(i)]
  \item (Part 1) The model is arbitrage-free iff a state-price vector $\bm\psi>0$ exists.
  \item (Part 2) An arbitrage-free model is complete iff the state-price vector is unique.
  \item If $\bm C_T$ is replicable, its only arbitrage-free price is
        $C_0=\bm\psi\T\bm C_T=\beta\,\E_Q[\bm C_T]$, for any state-price vector $\bm\psi$.
\end{enumerate}
\end{theorem}
\begin{proof}
(i) Let $L=\{(-A_0\bm\theta,\,A\bm\theta):\bm\theta\in\R^n\}\subset\R^{1+m}$. By definition an arbitrage is a
non-zero element of $L\cap\R^{1+m}_+$. By \cref{lem:stiemke}, no arbitrage iff there is
$\bm y=(y_0,y_1,\dots,y_m)>0$ with $-y_0A_0\bm\theta+\sum_{i\ge1}y_i(A\bm\theta)_i=0$ for all $\bm\theta$, i.e.\
$A_0=\bm\psi\T A$ with $\psi_i=y_i/y_0>0$.

(ii) State prices solve $A\T\bm\psi=A_0\T$. The solution is unique iff $\ker A\T=\{0\}$ iff
$\rank A=m$ iff $\operatorname{col}(A)=\R^m$, i.e.\ iff the market is complete.

(iii) If $A\bm\theta=\bm C_T$ then $C_0$ must equal the cost of $\bm\theta$ (otherwise buy the cheaper, sell the
dearer), and $A_0\bm\theta=\bm\psi\T A\bm\theta=\bm\psi\T\bm C_T$.
\end{proof}

\begin{remark}
In the binomial model $A=\bigl(\begin{smallmatrix}B_T & S_T^d\\ B_T & S_T^u\end{smallmatrix}\bigr)$ is invertible:
$\bm\psi\T=A_0A^{-1}$ is unique and positive exactly under the condition of \cref{prop:binomial-na},
and \eqref{eq:rnbinomial} is (iii). With more states than independent assets ($\rank A<m$) the market
is incomplete: state prices are not unique and a non-replicable claim has an \emph{interval} of
arbitrage-free prices $\bigl(\inf_{\bm\psi}\bm\psi\T\bm C_T,\ \sup_{\bm\psi}\bm\psi\T\bm C_T\bigr)$.
\end{remark}

\paragraph{Finding an arbitrage.} In a complete market, solve $A_0=\bm\psi\T A$. If some
$\psi_i\le0$, the portfolio replicating the Arrow--Debreu security $\bm e_i$ ($A\bm\theta=\bm e_i$) costs
$A_0\bm\theta=\bm\psi\T\bm e_i=\psi_i\le0$ and pays \$1 in state $i$: an arbitrage (lecture note,
section 2.4).

\begin{intuition}[Pricing is a positive linear functional]
Replication makes the price a \emph{linear} function of the payoff vector, and no arbitrage makes it
\emph{positive} (non-negative payoffs have non-negative prices). A positive linear functional on $\R^m$
is represented by a positive vector, $\bm\psi$, which normalised becomes a probability measure. This is
the finite-dimensional Riesz representation theorem; in continuous time it becomes ``price =
discounted expectation under the risk-neutral measure'' (\cref{ch:bs}).
\end{intuition}

\begin{casestudy}[Visualising portfolios (\file{lec02_2}, section 3)]
For the binomial example of \cref{ex:call}, the R code (\texttt{ggplot2}) plots the payoff vectors
$(V_T^d,V_T^u)$ of many portfolios \ts{24}{2}{1:01:39}:
\begin{itemize}
  \item long-only portfolios with weights summing
      to 1 lie on the segment between the bond and the stock;
  \item weights in $[0,1]^2$ fill a parallelogram,
      coloured by the cost $V_0$;
  \item allowing shorting (weights in $[-2,2]^2$) fills a large parallelogram
      around the origin. In the last plot the zero-cost portfolios (white) lie on a line through the origin
      that avoids the positive quadrant: no zero-cost portfolio pays off in both states, i.e.\ no arbitrage
      \ts{24}{2}{1:04:48}.
\end{itemize}
\end{casestudy}

\begin{finance}[From the binomial tree to Black--Scholes]
The binomial model iterated over many periods is the Cox--Ross--Rubinstein tree; letting the step go
to zero gives Black--Scholes (\cref{ch:bs}). Part~2 of the theorem explains why stochastic-volatility
or jump models, which have more sources of randomness than traded assets, are incomplete: prices there
depend on a choice of pricing measure calibrated to market option prices.
\end{finance}

% ------------------------------------------------------------------
\section{Eigenvalues and diagonalisation}\label{sec:eigen}
$\lambda\in\mathbb C$ is an eigenvalue of $A\in\R^{n\times n}$ with eigenvector $\bm v\neq0$ if
$A\bm v=\lambda\bm v$, i.e.\ iff $\det(A-\lambda I)=0$, a polynomial equation of degree $n$: there is
always at least one (complex) eigenvalue \ts{13}{2}{6:17}. Eigenvectors are the directions in which $A$
acts as a pure rescaling.

If $A$ has $n$ linearly independent eigenvectors, $S=[\bm v_1\cdots\bm v_n]$ is invertible and
\[
  AS=S\Lambda,\qquad A=S\Lambda S^{-1},\qquad A^k=S\Lambda^kS^{-1},\qquad \Lambda=\diag(\lambda_1,\dots,\lambda_n).
\]
\paragraph{State equations} \ts{24}{4}{4:14}. For the linear dynamics $\bm u_t=A\bm u_{t-1}$ (the state
equation of a Kalman filter, \cref{ch:timeseries}), expand the initial state in the eigenbasis,
$\bm u_0=\sum_ic_i\bm v_i$ with $\bm c=S^{-1}\bm u_0$; then
\[
  \bm u_t = A^t\bm u_0 = \sum_{i=1}^n c_i\,\lambda_i^t\,\bm v_i .
\]
Modes with $|\lambda_i|<1$ die out, those with $|\lambda_i|>1$ explode, and if the largest eigenvalue is
$1$ and simple, $\bm u_t$ converges to a multiple of its eigenvector \ts{24}{4}{6:21}. (For a physicist:
normal-mode decomposition of a linear discrete-time system.)

\begin{coursenote}[Erratum in the slides]
Slide ``Eigenvalues and Eigenvectors'' of \file{lec02_1} states
$\max_{\|\bm v\|=1}\|A\bm v\|=\max|\lambda(A)|$ for every square $A$ (and it is repeated in class
\ts{24}{2}{1:17:34}). This holds for \emph{symmetric} (more generally normal) matrices, as in the 2013 notes
(Theorem~4 of \file{lecnote2}). In general $\max_{\|\bm v\|=1}\|A\bm v\|=\sigma_1(A)$, the largest singular value
(\cref{sec:svd}), and $\sigma_1\ge\max|\lambda|$: for $A=\bigl(\begin{smallmatrix}0&1\\0&0\end{smallmatrix}\bigr)$
all eigenvalues are $0$ but $\|A\bm e_2\|=1$.
\end{coursenote}

% ------------------------------------------------------------------
\section{Symmetric matrices}\label{sec:symmetric}
Most matrices in this course (covariance and correlation matrices, $A\T A$, Hessians) are symmetric,
and for them everything works \ts{13}{2}{17:01}.

\begin{theorem}[Spectral theorem]\label{thm:spectral}
Let $A=A\T\in\R^{n\times n}$. Then
\begin{enumerate}[(i)]
  \item all eigenvalues are real;
  \item eigenvectors of distinct eigenvalues
      are orthogonal;
  \item $A=U\Lambda U\T$ with $U$ orthogonal ($U\T U=I$) and $\Lambda$ real diagonal.
\end{enumerate}
\end{theorem}
\begin{proof}
(i) If $A\bm v=\lambda\bm v$, then $\bar{\bm v}\T A\bm v=\lambda\|\bm v\|^2$; conjugating and using $A=\bar A=A\T$,
the same quantity equals $\bar\lambda\|\bm v\|^2$, so $\lambda=\bar\lambda$ \ts{13}{2}{19:12}.

(ii) $\lambda_1\bm v_2\T\bm v_1=\bm v_2\T A\bm v_1=(A\bm v_2)\T\bm v_1=\lambda_2\bm v_2\T\bm v_1$, so
$\bm v_2\T\bm v_1=0$ if $\lambda_1\ne\lambda_2$.

(iii) Induction: take a unit eigenvector $\bm v_1$; its orthogonal complement $W$ is invariant, because
$\bm w\perp\bm v_1\Rightarrow (A\bm w)\T\bm v_1=\bm w\T A\bm v_1=\lambda_1\bm w\T\bm v_1=0$; restrict $A$ to $W$ (still
symmetric) and repeat (2013 notes, Proposition~3).
\end{proof}

Consequences used later:
\begin{itemize}
  \item \emph{Rayleigh quotient}: the top eigenvector attains
        \[ \max_{\|\bm v\|=1}\bm v\T A\bm v=\lambda_{\max},\qquad\text{and}\qquad \max_{\|\bm v\|=1}\|A\bm v\|=\max_i|\lambda_i|. \]
        This variational characterisation is the basis of PCA.
  \item $A$ is \emph{positive semi-definite} ($\bm v\T A\bm v\ge0$ for all $\bm v$) iff all $\lambda_i\ge0$ (2013 PS1 B-3).
        Every covariance matrix $\Sigma$ is PSD, because $\bm w\T\Sigma\bm w=\Var(\bm w\T\bm X)\ge0$ is the variance of
        the portfolio with weights $\bm w$ (2024 PS2 Problem~3, \cref{ch:pca}).
\end{itemize}

% ------------------------------------------------------------------
\section{Singular value decomposition}\label{sec:svd}

\begin{theorem}[SVD]\label{thm:svd}
Every $A\in\R^{m\times n}$ of rank $r$ can be written as $A=U\Sigma V\T$ with $U\in\R^{m\times m}$ and
$V\in\R^{n\times n}$ orthogonal and $\Sigma\in\R^{m\times n}$ diagonal with entries
$\sigma_1\ge\dots\ge\sigma_r>0=\sigma_{r+1}=\cdots$. Equivalently
\begin{equation}\label{eq:svdsum}
  A=\sum_{k=1}^r\sigma_k\,\bm u_k\bm v_k\T,
\end{equation}
a sum of $r$ rank-one matrices. The $\sigma_k$ (singular values) are the square roots of the non-zero
eigenvalues of $A\T A$ (equivalently of $AA\T$).
\end{theorem}
\begin{proof}[Construction]
$A\T A$ is symmetric PSD with rank $r$: by \cref{thm:spectral} it has orthonormal eigenvectors
$\bm v_1,\dots,\bm v_n$ with eigenvalues $\sigma_1^2\ge\dots\ge\sigma_r^2>0=\cdots$. Put
$\bm u_k=A\bm v_k/\sigma_k$ for $k\le r$. These are orthonormal:
$\bm u_j\T\bm u_k=\bm v_j\T A\T A\bm v_k/(\sigma_j\sigma_k)=\sigma_k^2\bm v_j\T\bm v_k/(\sigma_j\sigma_k)=\delta_{jk}$.
Complete $\{\bm u_k\}$ and $\{\bm v_k\}$ to orthonormal bases; then $U\T AV=\Sigma$ because $A\bm v_k=\sigma_k\bm u_k$
($k\le r$) and $A\bm v_k=0$ ($k>r$) \ts{13}{2}{33:05}.
\end{proof}

\begin{intuition}[Rotate, stretch, rotate]
Read $A\bm x=U\Sigma V\T\bm x$ from the right: $V\T$ rotates the input frame onto the axes
$\bm v_1,\dots,\bm v_n$, $\Sigma$ stretches axis $k$ by $\sigma_k$ (and kills the axes with $\sigma_k=0$), $U$
rotates onto the output frame $\bm u_1,\dots,\bm u_m$. An eigen-decomposition uses \emph{one} frame in which
$A$ scales; the SVD uses \emph{two} frames, $A\bm v_k=\sigma_k\bm u_k$, but exists for every matrix, even
rectangular \ts{13}{2}{27:52}.
\end{intuition}

\begin{example}[SVD by hand {\ts{13}{2}{47:05}}]
$A=\bigl(\begin{smallmatrix}3&2&2\\2&3&-2\end{smallmatrix}\bigr)$. Then $AA\T=\bigl(\begin{smallmatrix}17&8\\8&17\end{smallmatrix}\bigr)$
has eigenvalues $25$ and $9$, so $\sigma_1=5$, $\sigma_2=3$ ($A\T A$ is $3\times3$ with eigenvalues
$25,9,0$). Eigenvectors of $A\T A$: $\bm v_1=\frac1{\sqrt2}(1,1,0)\T$, $\bm v_2=\frac1{\sqrt{18}}(1,-1,4)\T$;
then $\bm u_1=A\bm v_1/5=\frac1{\sqrt2}(1,1)\T$, $\bm u_2=A\bm v_2/3=\frac1{\sqrt2}(1,-1)\T$, and
\[
  A=\begin{pmatrix}\tfrac1{\sqrt2}&\tfrac1{\sqrt2}\\[2pt] \tfrac1{\sqrt2}&-\tfrac1{\sqrt2}\end{pmatrix}
    \begin{pmatrix}5&0\\0&3\end{pmatrix}
    \begin{pmatrix}\tfrac1{\sqrt2}&\tfrac1{\sqrt2}&0\\[2pt] \tfrac1{\sqrt{18}}&-\tfrac1{\sqrt{18}}&\tfrac4{\sqrt{18}}\end{pmatrix}.
\]
The eigenvector of eigenvalue $0$ is not needed: this is the \emph{reduced} SVD, keeping only the $r$
non-zero singular values \ts{13}{2}{54:33}.
\end{example}

\paragraph{Why it matters for data.} Let $A$ be a data matrix with one row per company and one column per
date (e.g.\ $5\times365$). The reduced SVD stores $5\times5$, $5\times5$ and $5\times365$ matrices instead of a
$365\times365$ one, and if the data effectively live in a low-dimensional subspace, a few singular values
dominate \ts{13}{2}{56:35}. $AA\T$ (indexed by companies) measures how the companies' series move together,
so the leading singular vectors describe groups of co-moving stocks \ts{13}{2}{58:47}; with centred data
this is exactly principal component analysis (\cref{ch:pca}). The best rank-$k$ approximation of $A$ in the
Frobenius or operator norm is the truncation of \eqref{eq:svdsum} to its first $k$ terms (Eckart--Young
theorem, not proved in the course).

% ------------------------------------------------------------------
\section{The Perron--Frobenius theorem}\label{sec:pf}

\begin{theorem}[Perron--Frobenius]\label{thm:pf}
Let $A\in\R^{n\times n}$ have all entries $a_{ij}>0$. Then:
\begin{enumerate}[(i)]
  \item there is a real eigenvalue $\lambda_0>0$ with $|\lambda|<\lambda_0$ for every other eigenvalue $\lambda$;
  \item $\lambda_0$ has an eigenvector with all entries positive;
  \item $\lambda_0$ is simple: its eigenvectors form a one-dimensional space.
\end{enumerate}
\end{theorem}
\begin{proof}
Write $\bm x\ge0$ (resp.\ $>0$) for componentwise inequalities. Following the slides, let
$T=\{t\ge0:\ A\bm x\ge t\bm x\text{ for some }\bm x\ge0,\ \bm x\ne0\}$ and $\lambda_0=\sup T$.
$T$ contains $\min_ia_{ii}>0$ and is bounded (if $A\bm x\ge t\bm x$ with $\sum x_i=1$, look at the largest
component $x_k\ge1/n$: $t\le(A\bm x)_k/x_k\le n\max a_{ij}$); by compactness of the simplex the sup is attained
at some $\bm x^*$.

\emph{Step 1: $\lambda_0$ is an eigenvalue with a positive eigenvector.} Suppose $\bm y=A\bm x^*-\lambda_0\bm x^*\ge0$ were
non-zero. Since all $a_{ij}>0$, $A\bm y>0$, i.e.\ $A\bm z>\lambda_0\bm z$ for $\bm z=A\bm x^*>0$, and then
$A\bm z\ge(\lambda_0+\varepsilon)\bm z$ for small $\varepsilon>0$, contradicting the maximality of $\lambda_0$. So
$A\bm x^*=\lambda_0\bm x^*$, and $\bm x^*=A\bm x^*/\lambda_0>0$. The same argument shows: \emph{any} $\bm x\ge0$, $\bm x\neq0$ with
$A\bm x\ge\lambda_0\bm x$ is an eigenvector of $\lambda_0$.

\emph{Step 2: $|\lambda|\le\lambda_0$.} If $A\bm z=\lambda\bm z$ ($\bm z$ possibly complex), then
$|\lambda|\,|\bm z|=|A\bm z|\le A|\bm z|$ componentwise, where $|\bm z|=(|z_1|,\dots,|z_n|)$; hence $|\lambda|\in T$.

\emph{Step 3: strict inequality.} If $|\lambda|=\lambda_0$, Step 2 gives $A|\bm z|\ge\lambda_0|\bm z|$, so by Step 1
$A|\bm z|=\lambda_0|\bm z|$ and $|\bm z|>0$. Then $|\sum_ja_{ij}z_j|=\sum_ja_{ij}|z_j|$ for each $i$: equality in the
triangle inequality with positive weights forces all $z_j$ to have the same phase, $\bm z=e^{i\phi}|\bm z|$.
Substituting, $\lambda|\bm z|=A|\bm z|=\lambda_0|\bm z|$, so $\lambda=\lambda_0$.

\emph{Step 4: simplicity.} By Step 3 every eigenvector of $\lambda_0$ is a multiple of a positive vector. If
$\bm z>0$ were an eigenvector not proportional to $\bm x^*$, let $c=\min_ix^*_i/z_i$; then $\bm w=\bm x^*-c\bm z\ge0$ has a zero
component and $\bm w\ne0$, yet $\bm w=A\bm w/\lambda_0>0$: contradiction.
\end{proof}

\begin{coursenote}[On the proofs in the course material]
The slides prove Steps 1--2 as above and argue simplicity from ``$\bm x-\bm y$ has entries of both
signs''; this needs the fact, proved in Step~3, that every eigenvector of $\lambda_0$ has entries of one
sign. The 2013 notes only sketch the symmetric case. The theorem also holds (with
$|\lambda|\le\lambda_0$) for non-negative \emph{irreducible} matrices, the form needed for Markov chains
\ts{13}{2}{1:10:23}. The eigenvalue is in fact algebraically simple as well (not shown here).
\end{coursenote}

\paragraph{Applications in this course.}
\begin{itemize}
  \item \textbf{Markov chains}: a column-stochastic $A$ with some $A^k>0$ has $\lambda_0=1$ (\cref{prop:stochastic}),
        a unique positive stationary distribution, and $\bm\pi(t)\to\bm\pi^*$ geometrically (2024 PS1 Problem~4).
  \item \textbf{Covariance matrices} with all positive entries: the first principal component has all weights
        positive, a ``market factor'' (2024 PS2 Problem~3, \cref{ch:pca}).
  \item \textbf{Ross recovery theorem}: recovering real-world probabilities from option prices relies on
        Perron--Frobenius applied to a matrix of state prices \ts{13}{2}{1:01:54} (\cref{ch:ross}).
\end{itemize}

% ------------------------------------------------------------------
\section{Case study: equal-weighted portfolios}\label{sec:ewport}

\begin{casestudy}[\file{Script_SP500.Rmd} and \file{Equal_Weighted_Portfolios.Rmd}]
\textbf{Data.} \file{Script_SP500.Rmd} downloads (\texttt{tidyquant}) daily prices of the stocks
currently in the S\&P~500 from 2019-01-01 to 2024-08-31 and keeps those with no missing data:
$1426$ days $\times$ $488$ stocks, saved in \file{SP500_data_subset.RData}.

\textbf{Portfolio as a matrix--vector product.} Invest \$1000 equally at the first date:
$q_j=(1000/n)/p_j(0)$ shares, then \emph{never rebalance} (buy and hold). With $P$ the
$1426\times n$ matrix of prices, the value path is simply $\bm V=P\bm q$ \ts{24}{4}{23:31}. The same is done for
four ``hot'' stocks (AMZN, AAPL, NFLX, GOOG), \$250 each.

\textbf{Results} (from the knitted output). The \$1000 in the four stocks grows to \$3541.76 by
2024-08-30; the four \$250 stakes end at AAPL \$1516.55, GOOG \$790.26, NFLX \$655.08, AMZN \$579.87.
The concentrated portfolio beats the equal-weighted S\&P~500 portfolio but has much deeper drawdowns
\ts{24}{4}{24:33}.
\end{casestudy}

\paragraph{Weight drift and rebalancing.} Without rebalancing, the weight of asset $j$ at time $t$ is
\[
  w_j(t) = \frac{q_jp_j(t)}{\sum_kq_kp_k(t)} = \frac{w_j(0)\,G_j(t)}{\sum_kw_k(0)\,G_k(t)},\qquad G_j(t)=\frac{p_j(t)}{p_j(0)},
\]
so winners take over: Apple's weight goes from 25\% to $1516.55/3541.76\approx43\%$, and the portfolio
becomes less diversified and riskier \ts{24}{4}{25:37}. Rebalancing every day back to equal weights is
not obviously better: it systematically sells winners and buys losers, which hurts when there are
trends \ts{24}{4}{26:40}. One principled alternative is to rebalance so as to keep the \emph{risk} of the
portfolio at a target level (\cref{ch:portfolio}).

\begin{finance}[Survivorship bias (not discussed in class)]
The script uses the stocks that are in the S\&P~500 \emph{today} and have complete histories since 2019.
Companies that were in the index in 2019 but later dropped out (often because they performed badly) are
excluded, so the back-test of the equal-weighted portfolio is biased upward. Any historical study built on
current index membership has this problem.
\end{finance}

% ------------------------------------------------------------------
\section{Lectures without material}\label{sec:la-missing}
\begin{coursenote}[2024 L3 and 2013 L4]
\textbf{2024 L3}, \emph{Quantitative Equity Investing} (Jeff Shen, BlackRock): not recorded and no slides
published (the guest lecturer did not grant IP permission). \textbf{2013 L4}, \emph{Matrix Primer} (Morgan
Stanley Matrix team): no notes and no video; OCW only points to the Morgan Stanley Matrix microsite
(the bank's electronic trading and research platform for clients).
\end{coursenote}

% ------------------------------------------------------------------
\begin{keyformulas}
\begin{align*}
&\text{Portfolio value, P\&L:} && V_t=\bm q\T\bm p(t),\quad \mathrm{PnL}(t+1)=\bm q(t+1)\T\Delta\bm p(t)\\
&\text{Markov chain:} && \bm\pi(t)=A^t\bm\pi(0),\quad \bm\pi^*=A\bm\pi^*\\
&\text{Two-state chain:} && \lambda_2=1-p-q,\quad \bm\pi^*=(q,p)/(p+q)\\
&\text{Binomial replication:} && \theta_S=\tfrac{C^u-C^d}{S^u-S^d},\\
  &&& C_0=\tfrac{q_uC^u+q_dC^d}{1+r_fT},\\
  &&& q_u=\tfrac{(1+r_fT)S_0-S^d}{S^u-S^d}\\
&\text{State prices:} && A_0=\bm\psi\T A,\ \bm\psi>0 \iff \text{no arbitrage};\\
  &&& C_0=\bm\psi\T\bm C_T=\beta\,\E_Q[C_T]\\
&\text{Completeness:} && \rank A=m\iff \bm\psi\text{ unique}\\
&\text{Diagonalisation:} && A=S\Lambda S^{-1},\quad A^t\bm u_0=\textstyle\sum_ic_i\lambda_i^t\bm v_i\\
&\text{Spectral theorem:} && A=A\T\Rightarrow A=U\Lambda U\T,\ U\T U=I\\
&\text{SVD:} && A=U\Sigma V\T=\textstyle\sum_{k\le r}\sigma_k\bm u_k\bm v_k\T,\\
  &&& \sigma_k^2=\lambda_k(A\T A)
\end{align*}
\end{keyformulas}

% ------------------------------------------------------------------
\section*{Exercises}
\addcontentsline{toc}{section}{Exercises}

\subsection*{From 2024 Problem Set 1}

\begin{exercise}[PS1 (2024), Problem 1: Fibonacci matrix]
Let $A=\bigl(\begin{smallmatrix}1&1\\1&0\end{smallmatrix}\bigr)$, $F_0=0$, $F_1=1$, $F_{k+2}=F_{k+1}+F_k$ and
$\bm u_k=(F_{k+1},F_k)\T$.
\begin{enumerate}[(a)]
  \item Show $\bm u_{k+1}=A\bm u_k$ and $\bm u_k=A^k\bm u_0$.
  \item Find the eigenvalues $\lambda_1,\lambda_2$ of $A$ and show that $\bm s=(\lambda,1)\T$ is an eigenvector for $\lambda$.
  \item With $S=[\bm s_1\ \bm s_2]$ compute $S^{-1}$, verify $A=S\Lambda S^{-1}$, and show $\bm c=S^{-1}\bm u_0=a(1,-1)\T$; find $a$.
  \item Deduce $F_k=(\lambda_1^k-\lambda_2^k)/\sqrt5$ and the limit of $F_{k+1}/F_k$ (the golden ratio, also used in
      ``technical analysis'' of prices, e.g.\ Fibonacci retracements).
\end{enumerate}
\end{exercise}

\begin{exercise}[PS1 (2024), Problem 2]
Let $A=S\Lambda S^{-1}$ be diagonalisable. Under which conditions on the eigenvalues does $\lim_kA^k$
\begin{enumerate}[(a)]
  \item not exist,
  \item exist and equal the zero matrix,
  \item exist and be non-zero?
\end{enumerate}
\end{exercise}

\begin{exercise}[PS1 (2024), Problem 3: order-side Markov chain]
The side of successive market orders is a two-state chain (1 = buy, 2 = sell) with
$A=\bigl(\begin{smallmatrix}0.8&0.3\\0.2&0.7\end{smallmatrix}\bigr)$, $a_{ij}=\Prob(X_t=i\mid X_{t-1}=j)$.
\begin{enumerate}[(a)]
  \item Prove $\bm u_t=A\bm u_{t-1}$ and $\bm u_t=A^t\bm u_0$ for the distribution $\bm u_t$ of $X_t$.
  \item Starting from $X_0=1$, compute $\bm u_1,\dots,\bm u_5$ in R (\texttt{A \%*\% ut} in a loop) and increase $t$
      until convergence; repeat from $X_0=2$. Is the limit the same?
  \item What is the stationary distribution of the order side? Relate it to the eigenvector of eigenvalue 1
      returned by \texttt{eigen(A)} (which is normalised to unit Euclidean length, not to sum 1).
\end{enumerate}
\end{exercise}

\begin{exercise}[PS1 (2024), Problem 4]
For the matrix of the previous exercise:
\begin{enumerate}[(a)]
  \item find its eigenvalues explicitly;
  \item explain which case of
      Problem~2 applies;
  \item explain the link between the stationary distribution and the Perron--Frobenius
      theorem.
\end{enumerate}
\end{exercise}

\begin{exercise}[PS1 (2024), Problem 5; lecture note \file{lec02_2}, Exercise 1]\label{ex:ps1-5}
In the two-asset one-period model prove that no arbitrage requires
$S_T^d<S_0(1+r_fT)<S_T^u$. \emph{Hint:} if an inequality fails, build an explicit portfolio whose cost is
lower than its payoff in every state.
\end{exercise}

\subsection*{From the lecture note \emph{One-Period Financial Models}}

\begin{exercise}[\file{lec02_2}, Exercise 2]
With $B_0=100$, $B_T=103$, $S_0=100$, $S_T^u=130$, $S_T^d=90$, consider the put with strike $K=100$,
payoff $\max(0,K-S_T)$.
\begin{enumerate}[(a)]
  \item Find its replicating portfolio.
  \item Find its cost at $t=0$.
  \item Is this price
      arbitrage-free? Check put--call parity against \cref{ex:call}.
\end{enumerate}
\end{exercise}

\begin{exercise}[\file{lec02_2}, Exercise 3]
Starting from $C_0=\theta_BB_0+\theta_SS_0$, derive \eqref{eq:rnbinomial} and interpret
$(1+r_fT)^{-1}S_T^{u}$, $(1+r_fT)^{-1}S_T^{d}$ and $S_0(1+r_fT)$ (the forward value of the bond position
obtained by shorting one share). Show that $q_u+q_d=1$ and that $q_u,q_d>0$ under no arbitrage.
\end{exercise}

\subsection*{From 2013 Problem Set 1}

\begin{exercise}[PS1 (2013), A-1]
True or false:
\begin{enumerate}[(a)]
  \item row rank equals column rank;
  \item $\rank(AB)\le\rank(A)$;
  \item there are matrices with
      $BA=I$, $AC=I$ and $B\neq C$;
  \item full rank iff invertible;
  \item symmetric matrices are invertible;
  \item distinct eigenvalues imply invertible;
  \item diagonalisable iff invertible.
\end{enumerate}
\end{exercise}

\begin{exercise}[PS1 (2013), A-2 to A-4]
(A-2) Row and column rank of
$\left(\begin{smallmatrix}1&0&1&0&1\\1&1&1&0&0\\0&0&1&1&1\end{smallmatrix}\right)$ and
$\left(\begin{smallmatrix}1&1&0&0&0&0\\1&1&0&3&0&0\\0&0&2&1&-1&0\\0&0&4&0&-2&0\end{smallmatrix}\right)$.
(A-3) Determinant and inverse of $\left(\begin{smallmatrix}1&2\\4&-1\end{smallmatrix}\right)$ and
$\left(\begin{smallmatrix}-1&-2&3\\1&2&0\\4&6&3\end{smallmatrix}\right)$.
(A-4) Characteristic polynomial, eigenvalues and eigenvectors of
$\left(\begin{smallmatrix}-3&3&2\\1&-1&-2\\-1&-3&0\end{smallmatrix}\right)$.
\end{exercise}

\begin{exercise}[PS1 (2013), A-5 and A-6]
(A-5) Gram--Schmidt on $\bm v_1=(1,0,1,0,1)$, $\bm v_2=(1,1,1,0,0)$, $\bm v_3=(0,0,1,1,1)$; complete to an orthonormal basis of
$\R^5$ and find the matrix $U$ mapping it to the standard basis.
(A-6) SVD of $\left(\begin{smallmatrix}3&1&1\\-1&3&1\end{smallmatrix}\right)$, $\left(\begin{smallmatrix}1&1\\2&2\end{smallmatrix}\right)$ and
$\left(\begin{smallmatrix}1&0&1\\0&1&0\\1&0&1\end{smallmatrix}\right)$.
\end{exercise}

\begin{exercise}[PS1 (2013), B-1 and B-2]
(B-1) For an SVD $A=U\Sigma V\T$, true or false with explanation:
\begin{enumerate}[(a)]
  \item the number of non-zero diagonal entries of
      $\Sigma$ equals $\rank A$;
  \item $\Sigma$ is unique up to permutations;
  \item for square $A$ one can choose $U=V$;
  \item for symmetric $A$ one can choose $U=V$.
\end{enumerate}
(B-2) Prove that $AA\T$ and $A\T A$ have the same non-zero eigenvalues.
\end{exercise}

\begin{exercise}[PS1 (2013), B-3]
Prove that a symmetric PSD matrix can be written $D=U\T AU$ with $U$ orthogonal and $D$ diagonal with
non-negative entries.
\end{exercise}

\begin{exercise}[PS1 (2013), B-4: least squares via SVD]
Minimise $L=\|A\bm x-\bm w\|$ for $A\in\R^{m\times n}$.
\begin{enumerate}[(a)]
  \item Solve it when $A$ is diagonal.
  \item Using
      $A\bm x-\bm w=U(\Sigma V\T\bm x-U\T\bm w)$, reduce the general case to (a).
  \item Show that fitting a line $y=ax+b$ to
      points $(x_i,y_i)$ by least squares is a special case;
  \item same for a parabola $y=ax^2+bx+c$.
      (Continued in \cref{ch:regression}.)
\end{enumerate}
\end{exercise}

% !TEX root = ../main.tex
\chapter{Probability Theory}\label{ch:probability}
\begin{paracol}{2}

\begin{sources}
\begin{tabularx}{\linewidth}{@{}l l L@{}}
\toprule
\textbf{Lecture} & \textbf{Speaker} & \textbf{Content}\\
\midrule
2024 L4 (from 27:40) & P.~Kempthorne & Discrete, continuous and mixed variables; CDF and probability integral transform; moments, skewness, kurtosis; normal and lognormal laws, change of variables; expected call payoff; moment-generating functions; independence, covariance, random vectors and portfolio variance; start of PCA.\\
2024 L5 (to 1:12:25) & P.~Kempthorne & PCA of a random vector; $\chi^2$, $t$ and $F$ laws; weak law of large numbers, central limit theorem; \emph{Probability Theory for Asset Pricing}: Stein's lemma, expected-utility pricing, CARA utility, CAPM. (The rest of L5 is in \cref{ch:sp1}.)\\
2013 L3 & C.~Lee & pmf/pdf, mutual vs pairwise independence; why prices are lognormal; change of variables; exponential families; MGFs and the moment problem; weak LLN (casino and poker); CLT with proof.\\
\bottomrule
\end{tabularx}

\smallskip
\textbf{Files.} Slides \emph{Probability Theory} \file{mit18_642_f24_lec04_1.pdf}; slides
\emph{Probability Theory for Asset Pricing} \file{mit18_642_f24_lec04_2.pdf}; 2013 notes
\file{MIT18_S096F13_lecnote3.pdf}. Suggested readings on the 2024 Week~3 page: J.~Lintner (1965),
\emph{The valuation of risk assets and the selection of risky investments in stock portfolios and
capital budgets}, and \emph{Security prices, risk, and maximal gains from diversification}.
\textbf{Problem sets.} 2024 PS2 Problem~2; 2024 PS3 Problems~1--2; 2013 PS2 A-1, A-2, B-1 to B-4.
The other problems of these sets (random walks, Markov chains, martingales, PCA of yields) are in
\cref{ch:sp1,ch:pca}.
\textbf{Not recorded.} Nothing: all three lectures are on video.
\end{sources}

\section*{Motivation}
Every model in this course is a probability model. This chapter collects the probability that the
rest of the notes take for granted, organised around four questions that a quant asks about a
position:
\begin{itemize}
  \item \emph{What is the distribution of a price or a return?} Normal for log-returns, lognormal
        for prices, and in reality heavier tails than either (\cref{ch03:sec-normal,ch03:sec-heavy}).
        The distribution fixes the expected payoff of an option (\cref{ch03:sec-call}).
  \item \emph{How do assets move together?} Covariance matrices, portfolio variance $\bm a\T\Sigma\bm a$,
        diversification, the multivariate normal law and its principal components
        (\cref{ch03:sec-multi,ch03:sec-mvn,ch03:sec-pca}).
  \item \emph{What survives aggregation?} Laws of large numbers and the central limit theorem explain
        why a small edge repeated many times makes money and why log-returns over long horizons look
        Gaussian (\cref{ch03:sec-limits}).
  \item \emph{What should an asset cost?} With nothing more than expectations, covariances and one
        Gaussian identity (Stein's lemma), an expected-utility argument yields the capital asset
        pricing model (\cref{ch03:sec-capm}).
\end{itemize}
A physicist can skim \cref{ch03:sec-rv,ch03:sec-moments} and the definitions of
\cref{ch03:sec-chisq}; the less standard material is the option-payoff formula, the pitfalls of
moment-generating functions, conditional expectation, Stein's lemma and the CAPM derivation.

% ------------------------------------------------------------------
\section{Random variables and distribution functions}\label{ch03:sec-rv}

Financial random variables come in three kinds \ts{24}{4}{28:47}:
\begin{itemize}
  \item \emph{discrete}: default (1) or no default (0) of a counterparty, the FOMC decision on the Fed
        funds rate, the side (buy/sell) and the share size of the next market order in AAPL;\begin{answer}{what is FOMC?}[FOMC]
The \emph{Federal Open Market Committee}: the body of the US Federal Reserve that sets the target for the
Fed funds rate (the overnight rate between banks). It meets eight times a year; each decision (cut, hold,
hike, usually by 0.25\%) is a discrete random variable for markets. ``Open market'' because the Fed
implements the decision by buying and selling bonds in the open market.
\end{answer}
  \item \emph{continuous}: the value of an asset, the waiting time until the next market order
        \ts{24}{4}{29:49};
  \item \emph{mixed}: the value of a stock in six months, which is continuous except for an atom at
        $0$ (bankruptcy) or at the buyout price (the company is taken private) \ts{24}{4}{30:55}.
\end{itemize}

\begin{definition}[Distribution of a random variable]
Let $\mathcal X$ be the set of possible values (the \emph{sample space}) of $X$.
\begin{itemize}
  \item If $X$ is discrete, its \emph{probability mass function} (pmf) is $f_X(x)=\Prob(X=x)\ge0$ with
        $\sum_{x\in\mathcal X}f_X(x)=1$.
  \item If $X$ is continuous, its \emph{probability density function} (pdf) $f_X\ge0$ satisfies
        $\Prob(X\in[x,x+\dd x])=f_X(x)\dd x$ and $\int_{\mathcal X}f_X(x)\dd x=1$.
  \item The \emph{cumulative distribution function} (CDF) is $F_X(x)=\Prob(X\le x)$
        \ts{24}{4}{31:56}. It is non-decreasing, right-continuous, with $F_X(-\infty)=0$ and
        $F_X(+\infty)=1$, and it determines the distribution.
\end{itemize}
Then $\Prob(X\in A)=\sum_{x\in A}f_X(x)$ or $\int_Af_X(x)\dd x$, and $F_X(x)=\int_{-\infty}^xf_X(u)\dd u$ in
the continuous case, so $f_X=F_X'$.
\end{definition}

\begin{remark}
Both courses identify the sample space with the set of values of $X$ (2013 notes, section~1). The
measure-theoretic set-up (a probability space $(\Omega,\mathcal F,\Prob)$ and $X$ a measurable map
$\Omega\to\R$) is not needed here; it becomes useful for martingales and filtrations in \cref{ch:sp1}.
\end{remark}

\paragraph{Independence.} $X$ and $Y$ are \emph{independent} if
$\Prob(\{X\in A\}\cap\{Y\in B\})=\Prob(X\in A)\,\Prob(Y\in B)$ for all sets $A,B$; equivalently, the joint
pmf/pdf factorises, $f_{X,Y}(x,y)=f_X(x)f_Y(y)$ \ts{13}{3}{6:16}. For $n$ variables one must distinguish
\ts{13}{3}{7:17}: $X_1,\dots,X_n$ are \emph{mutually independent} if
$\Prob(\bigcap_i\{X_i\in A_i\})=\prod_i\Prob(X_i\in A_i)$ for all $A_1,\dots,A_n$, and \emph{pairwise
independent} if every pair is independent. Mutual implies pairwise, not conversely
(\cref{ch03:ex-pairwise}). \note+{interesting new subtelty}``Independent'' always means mutually independent.

\begin{proposition}[Probability integral transform]\label{ch03:prop-pit}
Let $X$ have a continuous CDF $F$. Then $U=F(X)$ is uniform on $(0,1)$ \ts{24}{4}{32:57}. Conversely, if
$U\sim\mathrm{Uniform}(0,1)$ and $F^{-1}(u)=\inf\{x:F(x)\ge u\}$ is the quantile function of an
arbitrary CDF $F$, then $F^{-1}(U)$ has CDF $F$.
\end{proposition}
\begin{proof}
If $F$ is strictly increasing, $\Prob(F(X)\le u)=\Prob(X\le F^{-1}(u))=F(F^{-1}(u))=u$ for $u\in(0,1)$. In
general, with $q=F^{-1}(u)$ continuity gives $F(q)=u$, and $\{F(X)\ge u\}=\{X\ge q\}$ (if $X\ge q$ then
$F(X)\ge F(q)=u$; if $X<q$ then $F(X)<u$ by definition of the infimum), so
$\Prob(F(X)\ge u)=1-F(q)=1-u$. For the converse, $F^{-1}(u)\le x\iff u\le F(x)$ by right-continuity, so
$\Prob(F^{-1}(U)\le x)=\Prob(U\le F(x))=F(x)$.
\end{proof}

\begin{finance}[Uses of the probability integral transform]
\begin{itemize}
  \item \emph{Model checking} \ts{24}{4}{33:57}: if a risk model predicts that tomorrow's return has CDF
        $F_t$, the numbers $u_t=F_t(x_t)$ computed from realised returns should look like an i.i.d.\
        uniform sample. A histogram of the $u_t$ piling up near $0$ and $1$ reveals tails that are too
        thin: this is how value-at-risk and density forecasts are back-tested (\cref{ch:var}).\note+{cool way to spot the heavy tails}
  \item \emph{Simulation}: $F^{-1}(U)$ turns uniform random numbers into draws from any distribution
        (inverse-transform sampling).
  \item \emph{Copulas} (not discussed in class): applying $F_i$ to each coordinate of a random vector
        separates the marginal laws from the dependence structure.\begin{answer}{non credo di aver capito quest'ultima}[Copulas]
Take two returns $(X,Y)$ and set $U=F_X(X)$, $V=F_Y(Y)$. By \cref{ch03:prop-pit} each of $U,V$ is uniform:
all the information about the \emph{marginal} laws (fat tails, skew) has been removed. What is left, the joint
law of $(U,V)$ on the unit square, is the pure \emph{dependence}: the \emph{copula} (Latin, ``link''). Sklar's
theorem: $F_{X,Y}(x,y)=C(F_X(x),F_Y(y))$. So one can model each asset's tails and their co-movement
separately, e.g.\ fat-tailed marginals glued with a Gaussian copula (the model blamed for CDO losses in 2008).
\end{answer}
\end{itemize}
\end{finance}

% ------------------------------------------------------------------
\section{Expectation, moments and the shape of a distribution}\label{ch03:sec-moments}

\begin{definition}[Moments, skewness, kurtosis]\label{ch03:def-moments}
For a random variable $X$ (sums replace integrals in the discrete case):
\begin{itemize}
  \item mean $\mu=\E[X]=\int x f_X(x)\dd x$; $k$-th moment $m_k=\E[X^k]$;
  \item variance $\Var(X)=\E[(X-\mu)^2]=m_2-m_1^2$ and standard deviation $\sigma=\sqrt{\Var(X)}$, which has the
        units of $X$ \ts{24}{4}{35:03};
  \item with the standardised variable $Z=(X-\mu)/\sigma$ \ts{24}{4}{36:07}, the \emph{skewness}
        $\gamma=\E[Z^3]$ and the \emph{kurtosis} $\kappa=\E[Z^4]$, both dimensionless \ts{24}{4}{37:12}.
\end{itemize}
$\gamma>0$ indicates a long right tail, $\gamma<0$ a long left tail. The normal law has $\gamma=0$ and
$\kappa=3$; $X$ is \emph{leptokurtic} (fat-tailed relative to the Gaussian) if $\kappa>3$
\ts{24}{4}{38:18}. The \emph{excess kurtosis} is $\kappa-3$.
\end{definition}

\begin{proposition}[Affine transformations]\label{ch03:prop-affine}
If $Y=a+bX$ with $b\neq0$, then $\E[Y]=a+b\mu$, $\Var(Y)=b^2\sigma^2$, $\gamma_Y=\operatorname{sign}(b)\,\gamma_X$
and $\kappa_Y=\kappa_X$.
\end{proposition}
\begin{proof}
$(Y-\E[Y])/\sigma_Y=b(X-\mu)/(|b|\sigma)=\operatorname{sign}(b)\,Z_X$; take third and fourth powers.
\end{proof}

\begin{coursenote}[Erratum (in class) and a caveat on the slides]
At \ts{24}{4}{1:02:42} it is said that linear transformations do not change the skewness but might change
the kurtosis. It is the other way round (\cref{ch03:prop-affine}): kurtosis is invariant, skewness
changes sign when $b<0$ (a short position has the opposite skew of the long one). Also, slide~6 of
\file{lec04_1} lists ``$\gamma=0$: $X$ is symmetric about $\mu$''. Symmetry implies $\gamma=0$, but not
conversely: $X\in\{-3,-1,2\}$ with probabilities $0.1,\,0.5,\,0.4$ has $\E[X]=\E[X^3]=0$ and is not symmetric.
Software differs on whether ``kurtosis'' means $\kappa$ or $\kappa-3$: check before comparing numbers.
\end{coursenote}

\begin{finance}[What the moments of returns mean]
For a return, the mean is the drift and the standard deviation the \emph{volatility} (quoted
annualised). Equity-index returns are typically negatively skewed (crashes are sudden, rallies are
gradual) and daily returns of almost all assets are strongly leptokurtic. Both facts matter for risk:
a model matched only to mean and variance underestimates the probability of large losses
(\cref{ch03:sec-heavy}), and they are visible in option markets as the volatility smile (\cref{ch:volatility}).
\end{finance}\begin{answer}{can you include some graphs of: a typical probability distribution of an equity index, of a daily return of an asset. what is the valatility smile?}[Graphs]
Done in \cref{ch03:fig-returns}, with the course's own data (daily closes of the S\&P~500 stocks,
2019--2024). The equal-weighted portfolio of all 487 stocks stands in for the index.
\answerpart{Volatility smile}
Take the Black--Scholes formula, which has one unknown, $\sigma$. From the market price of each option
solve for the $\sigma$ that reproduces it: the \emph{implied volatility}. If returns were normal (BS
true), it would be the same number for every strike. It is not: plotted against the strike it is a
curve, U-shaped (a ``smile'') or, for equity indices, falling from left to right (a ``skew'', or
``smirk''), sketched in panel (c). Options far from the money are dearer than BS says, because the
market knows the tails are fat (and, for indices, that crashes go down). Full story in \cref{ch:volatility}.
\end{answer}

\begin{figure}[!htb]
\centering
\pgfplotsset{retplot/.style={width=5.1cm,height=4.8cm,tick label style={font=\scriptsize},label style={font=\small},title style={font=\small},legend style={font=\scriptsize},xlabel={daily log-return},
  scaled x ticks=false,xticklabel style={/pgf/number format/fixed},legend style={font=\tiny,fill=white,fill opacity=0.8,text opacity=1}}}
\begin{tikzpicture}
\begin{axis}[retplot,title={(a) linear scale},xmin=-0.08,xmax=0.08,ymin=0,ymax=75,xtick={-0.05,0,0.05},legend pos=north west]
  \addplot[ybar interval,fill=defcol!35,draw=defcol!60] table[x=x,y=ew] {data/fm_dailyret.dat};
  \addplot[thmcol,thick,domain=-0.08:0.08,samples=120] {exp(-(x-0.00068)^2/(2*0.01386^2))/(0.01386*sqrt(2*pi))};
  \legend{S\&P~500 (EW),normal fit}
\end{axis}
\end{tikzpicture}\hfill
\begin{tikzpicture}
\begin{axis}[retplot,title={(b) log scale: the tails},xmin=-0.14,xmax=0.12,ymode=log,ymin=0.05,ymax=3000,xtick={-0.1,0,0.1},
  legend pos=north west]
  \addplot[only marks,mark size=0.9pt,defcol] table[x=x,y=ew] {data/fm_dailyret.dat};
  \addplot[only marks,mark size=0.9pt,intcol] table[x=x,y=aapl] {data/fm_dailyret.dat};
  \addplot[thmcol,thick,domain=-0.07:0.07,samples=120] {exp(-(x-0.00068)^2/(2*0.01386^2))/(0.01386*sqrt(2*pi))};
  \legend{S\&P~500 (EW),AAPL,normal (EW)}
\end{axis}
\end{tikzpicture}\hfill
\begin{tikzpicture}
\begin{axis}[width=5.1cm,height=4.8cm,tick label style={font=\scriptsize},label style={font=\small},title style={font=\small},legend style={font=\scriptsize},title={(c) implied volatility (schematic)},
  xlabel={strike $K/S_0$},ylabel={implied vol},xmin=0.7,xmax=1.3,ymin=0.1,ymax=0.4,ytick=\empty,
  legend pos=north east]
  \addplot[thmcol,thick,domain=0.7:1.3,samples=60] {0.18+1.2*(x-1)^2};
  \addplot[defcol,thick,domain=0.7:1.3,samples=60] {0.2-0.35*(x-1)+0.8*(x-1)^2};
  \addplot[black!50,dashed,domain=0.7:1.3] {0.15};
  \legend{smile,index skew,BS (flat)}
\end{axis}
\end{tikzpicture}
\caption{Daily log-returns, 2019-01-02 to 2024-08-30, from \file{SP500_data_subset.RData} (the data of the
case study in \cref{sec:ewport}). (a) The equal-weighted S\&P~500 portfolio against the normal with the same mean
and standard deviation (1.39\% a day): a sharper peak and fatter tails; skewness $-0.91$, kurtosis $17$. (b) On a log scale
the normal is a parabola and the data are not: both the index and Apple (kurtosis $8.5$) have 11 days beyond
$4\sigma$, where the normal would expect $0.09$. The extreme left points are March 2020. (c) A sketch, not data,
of the volatility smile and of the skew typical of index options (\cref{ch:volatility}).}\label{ch03:fig-returns}
\end{figure}

% ------------------------------------------------------------------
\section{Normal and lognormal distributions}\label{ch03:sec-normal}

\subsection{The normal distribution}
\begin{definition}[Normal law]
$X\sim N(\mu,\sigma^2)$, $\mu\in\R$, $\sigma>0$, if
\[
  f(x)=\frac{1}{\sigma\sqrt{2\pi}}\exp\Bigl(-\frac{(x-\mu)^2}{2\sigma^2}\Bigr),\qquad x\in\R .
\]
$\phi$ and $\Phi$ denote the density and CDF of the standard normal $N(0,1)$; $\Phi(-x)=1-\Phi(x)$.\begin{answer}{perche' non restiamo con una notazione unica? $\phi$ = $f_normal$ e $\Phi$ = $F_normal$}[Why $\phi,\Phi$]
Because $\phi$ and $\Phi$ are universal: every paper and the Black--Scholes formula use them, so you must
recognise them. They are also a special case, not a different notation: for $Z\sim N(0,1)$, $f_Z=\phi$
and $F_Z=\Phi$; for $X\sim N(\mu,\sigma^2)$, $F_X(x)=\Phi((x-\mu)/\sigma)$.
\end{answer}
\end{definition}
The density is symmetric about $\mu$, maximal at $\mu$, with inflection points at $\mu\pm\sigma$
\ts{24}{4}{39:29}; the probabilities of lying within $1,2,3$ standard deviations are
$0.6827$, $0.9545$, $0.9973$ \ts{24}{4}{40:31}. That $f$ integrates to $1$ is the Gaussian integral
(2013 PS2 B-2, \cref{ch03:ex-b2}); see \cref{ch03:fig-normal}.\note*{insert a graph of the distribution. in general, whenever it can be useful plot graphs and functions.}

\begin{figure}[!htb]
\centering
\begin{tikzpicture}
\begin{axis}[width=0.8\linewidth,height=5cm,axis lines=left,xmin=-3.6,xmax=3.6,ymin=0,ymax=0.45,
  xtick={-3,-2,-1,0,1,2,3},xticklabels={$\mu-3\sigma$,$\mu-2\sigma$,$\mu-\sigma$,$\mu$,$\mu+\sigma$,$\mu+2\sigma$,$\mu+3\sigma$},
  ytick=\empty,tick label style={font=\scriptsize},label style={font=\small},title style={font=\small},legend style={font=\scriptsize}]
  \addplot[fill=defcol!12,draw=none,domain=-3:3,samples=80] {exp(-x^2/2)/sqrt(2*pi)} \closedcycle;
  \addplot[fill=defcol!25,draw=none,domain=-2:2,samples=80] {exp(-x^2/2)/sqrt(2*pi)} \closedcycle;
  \addplot[fill=defcol!45,draw=none,domain=-1:1,samples=80] {exp(-x^2/2)/sqrt(2*pi)} \closedcycle;
  \addplot[very thick,defcol,domain=-3.6:3.6,samples=120] {exp(-x^2/2)/sqrt(2*pi)};
  \node[font=\scriptsize] at (axis cs:0,0.12) {68.3\%};
  \node[font=\scriptsize] at (axis cs:1.5,0.03) {95.4\%};
  \node[font=\scriptsize] at (axis cs:2.55,0.025) {99.7\%};
  \fill[thmcol] (axis cs:-1,0.242) circle (1.5pt) (axis cs:1,0.242) circle (1.5pt);
  \node[font=\scriptsize,thmcol,anchor=west] at (axis cs:1.05,0.26) {inflection points};
\end{axis}
\end{tikzpicture}
\caption{The normal density $N(\mu,\sigma^2)$ with the probability of lying within $1$, $2$, $3$ standard deviations
of the mean. The inflection points are at $\mu\pm\sigma$.}\label{ch03:fig-normal}
\end{figure}

\subsection{Why prices are not normal}
Suppose daily price changes $P_n-P_{n-1}$ were i.i.d.\ normal \ts{13}{3}{10:28}. Two things go wrong
\ts{13}{3}{11:30}:
\begin{itemize}
  \item the size of a move would not depend on the price level, whereas a \$50 stock typically moves
        about five times as many dollars as a \$10 stock: it is the \emph{percentage} change that is
        roughly stationary \ts{13}{3}{12:32};
  \item $P_n$ would be normal with variance growing like $n$, so it eventually becomes negative with
        substantial probability \ts{13}{3}{19:05}. Locally (short horizons, small moves) the additive model
        is harmless.
\end{itemize}
The remedy is to model \emph{log-returns} $r_n=\ln(P_n/P_{n-1})$ as i.i.d.\ normal. Then
$\ln P_n=\ln P_0+\sum_{i\le n}r_i$ is normal (sums of independent normals are normal,
\cref{ch03:sec-mgf}), and $P_n$ is \emph{lognormal}; positivity is automatic \ts{24}{4}{41:34}. In continuous
time the same construction with Brownian increments gives geometric Brownian motion
\ts{24}{4}{42:35} (\cref{ch:stochcalc,ch:sde}).

\subsection{Change of variables}
\begin{theorem}[Change of variables]\label{ch03:thm-cov}
Let $X$ have density $f_X$, let $g$ be differentiable and strictly monotone with inverse $h=g^{-1}$, and
$Y=g(X)$. If $g$ is increasing, $F_Y(y)=F_X(h(y))$ \ts{24}{4}{43:35}; in all cases
\[
  f_Y(y)=f_X\bigl(h(y)\bigr)\,\bigl|h'(y)\bigr| .
\]
\end{theorem}
\begin{proof}
If $g$ is increasing, $\{Y\le y\}=\{X\le h(y)\}$, so $F_Y=F_X\circ h$; differentiating with the chain rule and
$F_X'=f_X$ gives $f_Y=f_X(h)\,h'$ \ts{24}{4}{44:43}. If $g$ is decreasing, $F_Y(y)=\Prob(X\ge h(y))=1-F_X(h(y))$
and $f_Y=-f_X(h)h'=f_X(h)|h'|$.
\end{proof}
The 2013 notes (Theorem~1.3) call $h'$ the \emph{Jacobian} \ts{13}{3}{16:47}; in $\R^n$ it becomes
$f_{\bm Y}(\bm y)=f_{\bm X}(\bm h(\bm y))\,|\det D\bm h(\bm y)|$.

\begin{intuition}[Conservation of probability]
Probability is a conserved ``mass'': $f_Y(y)|\dd y|=f_X(x)|\dd x|$. Where $g$ stretches the axis the density
thins out, where it compresses the axis the density piles up. This is the same bookkeeping as for a
particle density under a change of coordinates.
\end{intuition}

\subsection{The lognormal distribution}
\begin{definition}[Lognormal law]
$Y$ is lognormal$(\mu,\sigma^2)$ if $\ln Y\sim N(\mu,\sigma^2)$, i.e.\ $Y=e^X$ with $X\sim N(\mu,\sigma^2)$
\ts{24}{4}{41:34}. By \cref{ch03:thm-cov} with $h(y)=\ln y$, $h'(y)=1/y$ \ts{24}{4}{45:45}:
\[
  f_Y(y)=\frac{1}{y\,\sigma\sqrt{2\pi}}\exp\Bigl(-\frac{(\ln y-\mu)^2}{2\sigma^2}\Bigr),\qquad y>0
\]
\ts{24}{4}{46:56}\ts{13}{3}{21:12}.
\end{definition}

\noindent Properties (the first three from monotonicity of $\ln$ and a derivative):
\begin{itemize}
  \item quantiles map to quantiles: $q_p(Y)=e^{\mu+\sigma z_p}$ with $z_p=\Phi^{-1}(p)$; the median is $e^\mu$;
  \item the mode is $e^{\mu-\sigma^2}$;
  \item $\E[Y^k]=e^{k\mu+k^2\sigma^2/2}$ for every real $k$ (\cref{ch03:sec-mgf}); in particular
        $\E [Y]=e^{\mu+\sigma^2/2}$, so mode $<$ median $<$ mean: the law is skewed to the right;
  \item products and powers: if $Y_1,Y_2$ are independent lognormal, $Y_1Y_2$ and $Y_1^c$ are lognormal.
\end{itemize}
The parameters $\mu,\sigma$ are the mean and standard deviation of $\ln Y$, \emph{not} of $Y$
\ts{13}{3}{23:15}; they are kept because the law is defined through the logarithm \ts{13}{3}{24:18}.

\begin{example}[Lognormal$(0.2,0.4^2)$, slides 15--16 of \file{lec04_1}]\label{ch03:ex-ln}
Let $X_0=1$ and $X_1$ lognormal with $\mu=0.2$, $\sigma=0.4$: a one-year log-return with mean $0.2$ and
standard deviation $0.4$ \ts{24}{4}{47:56}. Then (\cref{ch03:fig-lognormal}) the 5\%, 25\%, 50\%, 75\%, 95\%
quantiles are $0.633$, $0.933$, $1.221$, $1.600$, $2.358$ \ts{24}{4}{48:58}; the mode is $1.041$, the mean
$1.323$ and the standard deviation $0.551$. The quantiles are symmetric \emph{multiplicatively} about the
median: $1.221/0.633=2.358/1.221=1.93$.
\end{example}

\begin{figure}[ht]
\centering
\begin{tikzpicture}
\begin{axis}[width=0.85\linewidth,height=5.6cm,axis lines=left,xmin=0,xmax=3,ymin=0,ymax=1.12,
  xlabel={$y$},ylabel={$f_Y(y)$},tick label style={font=\small},label style={font=\small}]
  \addplot[very thick,defcol,domain=0.03:3,samples=150]
    {exp(-(ln(x)-0.2)^2/(2*0.16))/(x*0.4*sqrt(2*pi))};
  \draw[thmcol,thick] (axis cs:0.6326,0) -- (axis cs:0.6326,0.92) node[above,font=\scriptsize]{5\%};
  \draw[intcol,thick] (axis cs:0.9326,0) -- (axis cs:0.9326,0.92) node[above,font=\scriptsize]{25\%};
  \draw[black,thick] (axis cs:1.2214,0) -- (axis cs:1.2214,1.0) node[above,font=\scriptsize]{median};
  \draw[intcol,thick] (axis cs:1.5997,0) -- (axis cs:1.5997,0.92) node[above,font=\scriptsize]{75\%};
  \draw[thmcol,thick] (axis cs:2.3583,0) -- (axis cs:2.3583,0.92) node[above,font=\scriptsize]{95\%};
  \draw[fincol,thick,dashed] (axis cs:1.3231,0) -- (axis cs:1.3231,0.62);
  \node[font=\scriptsize,fincol,anchor=west] at (axis cs:1.33,0.06) {mean};
\end{axis}
\end{tikzpicture}
\caption{Density of the lognormal$(0.2,0.4^2)$ law with its 5, 25, 50, 75, 95\% quantiles (as on slide~16
of \file{lec04_1}) and its mean. The mode ($1.041$, the peak) lies left of the median, the mean right of
it.}\label{ch03:fig-lognormal}
\end{figure}

\begin{coursenote}[``An annualised return of about 20\%'']
In class the example is read as an asset with an annualised return of about 20\% and volatility 0.4
\ts{24}{4}{47:56}. Be careful which return: $0.2$ is the expected \emph{log}-return; the median gross
return is $e^{0.2}=1.221$ (+22.1\%) and the expected gross return is $e^{0.28}=1.323$ (+32.3\%). The gap
$\sigma^2/2=0.08$ is the convexity correction that reappears as the $-\sigma^2/2$ in the solution of the
geometric Brownian motion (\cref{ch:stochcalc}).
\end{coursenote}

\begin{finance}[Log-returns and the square-root-of-time rule]
Log-returns add over time: $\ln(S_T/S_0)=\sum_{\text{days}}r_i$. If daily log-returns are i.i.d.\ with mean
$\mu_d$ and variance $\sigma_d^2$, the $T$-day log-return has mean $T\mu_d$ and standard deviation
$\sigma_d\sqrt T$. This is why volatilities are quoted annualised and scaled by $\sqrt t$ (as in 2024 PS2
Problem~2(c), where $\sigma=\sqrt t\,\ln1.3$), and why the Black--Scholes model takes $S_T$ lognormal
(\cref{ch:bs}). For small moves the simple return $e^{r}-1\approx r$.
\end{finance}\begin{answer}{vorrei che la notazione fosse consistente per tutto il libro. S e' il prezzo della stock, giusto? cos'e' P allora? B e' il prezzo del bond?}[$S$, $P$, $B$]
Yes: $S$ = stock price, $B$ = bond (risk-free asset) price. $P$ is a \emph{generic} asset price, any asset;
above it comes from the 2013 lecture, which uses $P_n$ for the price on day $n$ of an unspecified asset
(as $P^j$ in \cref{def:onepmodel}). Here $P_n$ and $S_n$ are the same thing.
\end{answer}

\begin{exercise}[PS2 (2013), B-2]\label{ch03:ex-b2}
\leavevmode
\begin{enumerate}[(a)]
  \item Verify that $\phi(x)=\frac1{\sigma\sqrt{2\pi}}e^{-(x-\mu)^2/(2\sigma^2)}$ is a probability density (compute
      $\bigl(\int\phi\bigr)^2$ in polar coordinates).
  \item Compute the mean and variance of the lognormal law (it suffices to
      treat $\mu=0$; use (a) to simplify). (This is also Exercise~1.5(i) of \file{lecnote3}; Exercise~1.2 there asks
      to verify that the mean of $N(\mu,\sigma^2)$ is $\mu$.)
\end{enumerate}
\end{exercise}

\begin{exercise}[PS2 (2013), B-4: does time diversify risk?]
(From M.~Kritzman, \emph{Puzzles of Finance}, ch.~3; see also Kritzman and Rich, \emph{Beware of dogma: the truth
about time diversification}.) A stock has independent annual returns, each lognormal with $\mu=\ln1.1$ and
$\sigma=\ln1.2$. According to the author, the fraction of wealth lost with 0.1\% probability over $T$ years is
40.5\%, 58.43\%, 61.48\% and 73.92\% for $T=1,5,10,20$, and he concludes that if risk is measured by the
magnitude of potential loss, it \emph{increases} with the horizon.
\begin{enumerate}[(a)]
  \item Compute the fraction $h(T)$ of wealth
      lost with probability 0.1\% over $T$ years (in terms of $\Phi$).
  \item Find the $T$ that maximises $h(T)$. Is $h$
      increasing in $T$?
  \item Criticise the argument using (a) and (b).
\end{enumerate}
\end{exercise}

\begin{exercise}[\file{lec04_1}, slide 21]
\leavevmode
\begin{enumerate}[(a)]
  \item Find the skewness of a lognormal$(\mu,\sigma^2)$ variable (use the lognormal moments).
  \item Show directly
      from the density that the kurtosis of $N(\mu,\sigma^2)$ is 3.
\end{enumerate}
\end{exercise}

% ------------------------------------------------------------------
\section{Expected option payoffs}\label{ch03:sec-call}

A European call gives the right to buy the asset at time $T$ at the strike $K$; if $X$ is the price at $T$,
the payoff is the ``hockey stick'' $C=(X-K)^+=\max(0,X-K)$ \ts{24}{4}{50:01} (\cref{ch03:fig-hockey}).\note*{plot it}

\begin{figure}[!htb]
\centering
\begin{tikzpicture}
\begin{axis}[width=6.6cm,height=4.6cm,axis lines=middle,xmin=0,xmax=2.05,ymin=-0.1,ymax=1.1,
  xtick={1},xticklabels={$K$},ytick=\empty,xlabel={$X$},tick label style={font=\scriptsize},label style={font=\small},title style={font=\small},legend style={font=\scriptsize},title={call: $(X-K)^+$},
  xlabel style={at={(1,0)},anchor=north}]
  \addplot[very thick,intcol] coordinates {(0,0) (1,0) (2,1)};
\end{axis}
\end{tikzpicture}\hfill
\begin{tikzpicture}
\begin{axis}[width=6.6cm,height=4.6cm,axis lines=middle,xmin=0,xmax=2.05,ymin=-0.1,ymax=1.1,
  xtick={1},xticklabels={$K$},ytick={1},yticklabels={$K$},xlabel={$X$},tick label style={font=\scriptsize},label style={font=\small},title style={font=\small},legend style={font=\scriptsize},title={put: $(K-X)^+$},
  xlabel style={at={(1,0)},anchor=north}]
  \addplot[very thick,thmcol] coordinates {(0,1) (1,0) (2,0)};
\end{axis}
\end{tikzpicture}
\caption{Payoffs at expiry as functions of the price $X$ of the underlying: the call (the ``hockey stick'': flat
blade, then slope 1) and the put, which pays at most $K$.}\label{ch03:fig-hockey}
\end{figure}

\begin{theorem}[Expected call payoff]\label{ch03:thm-call}
Let $X$ have density $f$, CDF $F$ and $\E[|X|]<\infty$. For every constant $K$,
\[
  \E\bigl[(X-K)^+\bigr]=\int_K^\infty\!\!\int_x^\infty f(t)\dd t\dd x=\int_K^\infty\bigl[1-F(x)\bigr]\dd x .
\]
\end{theorem}
\begin{proof}
Following the slides \ts{24}{4}{51:05}, for $M>K$ integrate by parts with $u=x-K$, $v=-[1-F(x)]$:
\begin{align*}
  \int_K^M(x-K)f(x)\dd x&=-(M-K)\bigl[1-F(M)\bigr]\\
  &\quad+\int_K^M\bigl[1-F(x)\bigr]\dd x .
\end{align*}
The boundary term vanishes as $M\to\infty$ because $0\le(M-K)[1-F(M)]\le\int_M^\infty(x-K)f(x)\dd x$, the tail
of a convergent integral. (Alternatively: $(x-K)^+=\int_K^\infty\ind\{x>s\}\dd s$; take expectations and swap
the integrals.)
\end{proof}
This is the familiar fact that for a non-negative variable $\E[W]=\int_0^\infty\Prob(W>s)\dd s$
\ts{24}{4}{52:11}, applied to $W=(X-K)^+$. The slides assume finite variance; a finite mean is enough.

\begin{corollary}[Strike derivatives; not discussed in class]\label{ch03:cor-bl}
$C(K)=\E[(X-K)^+]$ satisfies $C'(K)=-\Prob(X>K)$ and $C''(K)=f(K)$.
\end{corollary}
So call prices across strikes determine the distribution of $X$ (under the pricing measure): this is
the Breeden--Litzenberger formula behind \cref{ch:duality,ch:ross}.\begin{answer}{questa parte non l'ho capita}[Breeden--Litzenberger]
Read $C(K)$ as the (undiscounted) price of a call, for every strike $K$. Differentiate twice in $K$: the
corollary gives $C''(K)=f(K)$. So if the market quotes calls at many strikes, the curvature of the price
curve \emph{is} the density of $S_T$ (under $Q$, \cref{ans:pricingQ}). Intuition: a ``butterfly''
(buy calls at $K-h$ and $K+h$, sell two at $K$) pays a small triangle around $K$, i.e.\ roughly
$h$ dollars if $S_T\approx K$; its price $\approx h^2C''(K)$ is the price of ``$S_T$ near $K$''.
\end{answer}

\begin{corollary}[Normal and lognormal payoffs, slide 18 of \file{lec04_1}]\label{ch03:cor-normal}
\leavevmode
\begin{enumerate}[(i)]
  \item If $X\sim N(\mu,\sigma^2)$ and $d=(\mu-K)/\sigma$, then
      \[
      \E\bigl[(X-K)^+\bigr]=(\mu-K)\,\Phi(d)+\sigma\,\phi(d)=(\mu-K)\,\Phi(d)+\frac{\sigma}{\sqrt{2\pi}}\,e^{-(K-\mu)^2/(2\sigma^2)} .
      \]
  \item If $X$ is lognormal$(\mu,\sigma^2)$ and $d=(\mu-\ln K)/\sigma$, then
      \[
      \E\bigl[(X-K)^+\bigr]=e^{\mu+\sigma^2/2}\,\Phi(d+\sigma)-K\,\Phi(d)
      \]
      \ts{24}{4}{53:14}.
\end{enumerate}
\end{corollary}\begin{answer}{non c'e' un errore nell'ultima uguaglianza di (i)?}[No error]
It is correct: $\phi(d)=e^{-d^2/2}/\sqrt{2\pi}$ with $d^2=(\mu-K)^2/\sigma^2$, and $(\mu-K)^2=(K-\mu)^2$.
\end{answer}
\begin{proof}
(i) Write $X=\mu+\sigma Z$: $(X-K)^+=(\mu-K+\sigma Z)\ind\{Z>-d\}$. Then $\Prob(Z>-d)=\Phi(d)$ and
$\E[Z\ind\{Z>a\}]=\int_a^\infty z\phi(z)\dd z=\phi(a)$, with $\phi(-d)=\phi(d)$.

(ii) is 2024 PS2 Problem~2(d)
(\cref{ch03:ex-ps2}).
\end{proof}

\begin{finance}[From expected payoffs to option prices]
A price is a \emph{discounted} expectation under a \emph{pricing} measure (\cref{thm:ftap}):\begin{answer}{cos'e' una pricing measure Q?}[Pricing measure]\label{ans:pricingQ}
The weights $q_i=\psi_i/\sum\psi$ of \cref{def:onepmodel}: positive numbers summing to 1, so they
\emph{look} like probabilities, chosen so that every price is a discounted average of payoffs,
$P_0=\beta\E_Q[P_T]$. They are not the real chances: they also include risk aversion (bad states get more
weight). Also called the \emph{risk-neutral} measure, see \cref{ch:bs}.
\end{answer}
$C_0=e^{-rT}\E_Q[(S_T-K)^+]$. In the Black--Scholes model, under $Q$, $\ln S_T\sim N(\ln S_0+(r-\sigma^2/2)T,\
\sigma^2T)$; inserting these parameters in \cref{ch03:cor-normal}(ii) gives $e^{\mu+\sigma^2T/2}=S_0e^{rT}$, $d=d_2$,
$d+\sigma\sqrt T=d_1$ and hence $C_0=S_0\Phi(d_1)-Ke^{-rT}\Phi(d_2)$ (\cref{ch:bs}). Part~(1) is the
\emph{Bachelier} model (1900), still used for interest rates and spreads, which can be negative. At the
money ($K=\mu$) it gives $\sigma/\sqrt{2\pi}\approx0.4\,\sigma$: the traders' rule of thumb ``an at-the-money
option is worth about 0.4 standard deviations of the underlying at expiry'' (not discussed in class).\begin{answer}{can you give a brief example?}[Rule of thumb]
A stock at \$100 with volatility 20\%/year; a 3-month at-the-money option. The standard deviation of $S_T$
is about $100\cdot0.2\cdot\sqrt{0.25}=\$10$, so the option is worth about $0.4\cdot10=\$4$. Black--Scholes
($r=0$) gives \$3.99.
\end{answer}
\end{finance}

\begin{exercise}[not from the course: puts and parity]
With the notation of \cref{ch03:thm-call}, show that $\E[(K-X)^+]=\int_{-\infty}^KF(x)\dd x$ and that
$\E[(X-K)^+]-\E[(K-X)^+]=\E[X]-K$. Interpret the second identity as put--call parity: a long call plus a short put is a
forward contract (\cref{sec:forward}).
\end{exercise}

% ------------------------------------------------------------------
\section{Moment-generating functions}\label{ch03:sec-mgf}

\begin{definition}[MGF, characteristic function]
The \emph{moment-generating function} of $X$ is $M_X(t)=\E[e^{tX}]$, $t\in\R$ \ts{24}{4}{53:14}\ts{13}{3}{36:02}.
It always satisfies $M_X(0)=1$, and the set where it is finite is an interval containing $0$ (by
convexity of $t\mapsto e^{tx}$). We say that $X$ \emph{has an MGF} if $M_X$ is finite on some
$(-\delta,\delta)$. The \emph{characteristic function} $\varphi_X(t)=\E[e^{itX}]=\E[\cos tX]+i\,\E[\sin tX]$
is finite for every $X$ and $t$, since $|e^{itX}|=1$ \ts{24}{4}{56:28}.
\end{definition}\begin{answer}{they are related by analitical continuation. posso dire che una e' la laplace transform e l'altra la fourier transform?}[Yes]
Yes: $M_X(t)$ is the (two-sided) Laplace transform of the density at $-t$, $\varphi_X(t)$ its Fourier
transform (up to the sign convention), and $\varphi_X(t)=M_X(it)$ when $M_X$ is finite near 0. As in
physics, the Fourier version always exists, the Laplace one only with enough decay.
\end{answer}

\begin{proposition}[Moments from the MGF]
If $M_X$ is finite on $(-\delta,\delta)$, all moments are finite, $\E[X^k]=M_X^{(k)}(0)$ and
$M_X(t)=\sum_{k\ge0}\frac{t^k}{k!}\,m_k$ for $|t|<\delta$ \ts{13}{3}{38:14}\ts{24}{4}{54:18}.
\end{proposition}
\begin{proof}
For $|s|<\delta$, $e^{s|x|}\le e^{sx}+e^{-sx}$, so $\E[e^{s|X|}]<\infty$; this dominates $\sum_k|tX|^k/k!$ for
$|t|\le s$, so expectation and series (and derivatives) can be exchanged.
\end{proof}

\paragraph{The MGF need not exist} \ts{24}{4}{55:21}. For the Cauchy density $f(x)=1/(\pi(1+x^2))$ (\cref{ch03:fig-laws}),\note*{plot it} $\E[|X|]=\infty$
and $M_X(t)=\infty$ for every $t\ne0$; its characteristic function is $e^{-|t|}$. The \emph{lognormal} law has
all moments but no MGF \ts{13}{3}{37:09}: for $t>0$ and every $k$, $e^{tY}\ge(tY)^k/k!$, hence
\[
  M_Y(t)\ \ge\ \frac{t^k\,\E[Y^k]}{k!}=\frac{t^k e^{k\mu+k^2\sigma^2/2}}{k!}\ \ge\ \exp\Bigl(\frac{k^2\sigma^2}{2}+k(\mu+\ln t)-k\ln k\Bigr)\xrightarrow[k\to\infty]{}\infty ,
\]
using $k!\le k^k$. For $t\le0$, instead, $0<e^{tY}\le1$ and $M_Y(t)$ is finite.

\begin{coursenote}[Erratum in the 2013 notes]
\file{lecnote3} (and the lecture, \ts{13}{3}{46:47}) state that for the lognormal law $\E[e^{tX}]$ does not
exist for all $t\neq0$. It is infinite for all $t>0$ but finite for all $t<0$; what fails is finiteness on a
neighbourhood of $0$, which is what the theorems below need. In the same paragraph,
$\frac{d^kM_X}{dx^k}(0)$ should read $\frac{d^kM_X}{dt^k}(0)$.
\end{coursenote}

\begin{theorem}[Uniqueness and continuity]\label{ch03:thm-mgf}
\leavevmode
\begin{enumerate}[(i)]
  \item If $M_X(t)=M_Y(t)<\infty$ for all $t$ in a neighbourhood of $0$, then $X$ and $Y$ have the same distribution
        \ts{24}{4}{56:28}\ts{13}{3}{39:18}.
  \item If $M_{X_n}(t)\to M_X(t)<\infty$ for all $t$ in a neighbourhood of $0$, then $X_n$ converges to $X$ \emph{in
        distribution}: $F_{X_n}(x)\to F_X(x)$ at every point where $F_X$ is continuous \ts{24}{4}{57:31}\ts{13}{3}{43:39}.
\end{enumerate}
The same statements hold for characteristic functions, with ``for all $t\in\R$'' and no existence
assumption. (Not proved in class; see Durrett, \emph{Probability: Theory and Examples}.)
\end{theorem}

\begin{coursenote}[Statement in the slides and notes]
Slide~19 of \file{lec04_1} and Theorem~2.2 of \file{lecnote3} assume equality or convergence ``for all $t$''.
The neighbourhood version above is the standard one and is the one needed in practice: the MGF of a
$\chi^2$ variable, for instance, is finite only for $t<1/2$ (\cref{ch03:sec-chisq}).
\end{coursenote}\begin{answer}{per ovviare al fatto che alcune distribuzioni non hanno i momenti ben definiti, perche' non aggiungiamo un weight esponenziale all'integrale cosi' abbiamo quantita' finite per caratterizzare le nostre distribuzioni heavy-tailed?}[It is done]
Good idea, and it is used. A damping $e^{-a|x|}$ makes every ``moment'' finite, but the numbers then
depend on $a$ and lose their meaning (mean, variance). The bounded weight $e^{itx}$ is the standard
answer: the characteristic function always exists and determines the law. Practitioners also use
quantile-based measures (median, interquartile range) and the tail index $\alpha$ of $\Prob(|X|>x)\sim x^{-\alpha}$.
In option pricing, damped transforms $e^{-ax}$ are the Carr--Madan (1999) method of pricing options by FFT.
\end{answer}

\begin{remark}[Moments do not determine a distribution]
If two variables have the same moments of all orders, they need not have the same law: the hole in the
argument ``same moments $\Rightarrow$ same MGF $\Rightarrow$ same law'' is that the MGF may not exist
\ts{13}{3}{41:26}\ts{13}{3}{42:28}. The standard example (Durrett, pp.~106--107, cited by \file{lecnote3}) is the
lognormal. Let $f$ be the lognormal$(0,1)$ density and, for $|a|\le1$, $f_a(y)=f(y)\,[1+a\sin(2\pi\ln y)]\ge0$.
With $y=e^u$, $\phi(u)e^{ku}=e^{k^2/2}\phi(u-k)$ and $v=u-k$,
\[
  \int_0^\infty y^kf(y)\sin(2\pi\ln y)\dd y=e^{k^2/2}\int_{\R}\phi(v)\sin(2\pi v+2\pi k)\dd v=0
  \qquad(k=0,1,2,\dots),
\]
since $\sin$ is odd. So every $f_a$ is a density with exactly the lognormal moments. (Not discussed in class.)
\end{remark}\begin{answer}{pero' non c'e' il problema che la pdf non e' definita positiva? inoltre, i momenti non si possono ricavare dalla characteristic function che esiste sempre?}[Positivity]
No: $|a\sin(\cdot)|\le1$, so $1+a\sin(2\pi\ln y)\ge0$ and $f_a\ge0$; it integrates to 1 (the $k=0$ case).
\answerpart{Moments from $\varphi$?}
Yes, $\E[X^k]=i^{-k}\varphi^{(k)}(0)$ when the moments exist. But the converse fails: the moments give only
the Taylor coefficients of $\varphi$ at 0, and for the lognormal that Taylor series has radius 0, so it
does not reconstruct $\varphi$. Different laws, different $\varphi$, same derivatives at 0.
\end{answer}

\paragraph{The normal MGF.} For $Z\sim N(0,1)$, complete the square, $tx-\tfrac12x^2=-\tfrac12(x-t)^2+\tfrac12t^2$
\ts{24}{4}{58:34}\ts{24}{4}{1:00:40}:
\[
  M_Z(t)=\int_{\R}\frac{e^{-\frac12(x-t)^2}}{\sqrt{2\pi}}\dd x\;e^{t^2/2}=e^{t^2/2}.
\]
Since $M_{a+bX}(t)=e^{at}M_X(bt)$ \ts{24}{4}{1:01:42}, for $X=\mu+\sigma Z\sim N(\mu,\sigma^2)$:
\begin{equation}\label{ch03:eq-normalmgf}
  M_X(t)=e^{\mu t+\sigma^2t^2/2}.
\end{equation}
Consequences:
\begin{itemize}
  \item \emph{Lognormal moments} \ts{24}{4}{1:02:42}: $\E[e^{kX}]=M_X(k)$, so $\E[Y^k]=e^{k\mu+k^2\sigma^2/2}$ for
        $Y=e^X$ (2024 PS2 Problem~2(a,b)).
  \item \emph{Sums}: if $X_1,\dots,X_n$ are independent, $M_{\sum X_i}(t)=\prod_iM_{X_i}(t)$, because the expectation of a
        product of independent variables is the product of expectations. Hence a sum of independent normals
        is $N(\sum\mu_i,\sum\sigma_i^2)$, by \cref{ch03:thm-mgf}(i).
  \item \emph{Normal moments}: $e^{t^2/2}=\sum_k t^{2k}/(2^kk!)$ gives $\E[Z^{2k}]=(2k)!/(2^kk!)$, e.g.\ $\E[Z^4]=3$:
        the normal kurtosis is $3$.
\end{itemize}

\begin{intuition}[Partition function and cumulants (not discussed in class)]
$M_X(t)$ is a Laplace transform, like a partition function $Z(\beta)$, and $K_X(t)=\ln M_X(t)$ plays the role of
a free energy. Its Taylor coefficients are the \emph{cumulants}: $K_X(t)=\mu t+\sigma^2\frac{t^2}{2}+\kappa_3\frac{t^3}{6}+
\kappa_4\frac{t^4}{24}+\cdots$ with $\kappa_3=\gamma\sigma^3$ and $\kappa_4=(\kappa-3)\sigma^4$. For independent variables
cumulants \emph{add}. The Gaussian is the ``free theory'': $K$ is exactly quadratic. For the standardised sum
$S_n/\sqrt n$ of i.i.d.\ variables the $k$-th cumulant is $n\kappa_k/n^{k/2}\to0$ for $k\ge3$: all non-Gaussian
features wash out, which is the central limit theorem in one line (\cref{ch03:thm-clt}).
\end{intuition}\begin{answer}{beautiful box. i want to know more.}[More on cumulants]
Three directions:
\begin{itemize}
  \item \emph{Large deviations} (Cram\'er): $\Prob(S_n/n\approx x)\approx e^{-nI(x)}$ with $I$ the Legendre
        transform of $K$, exactly free energy $\leftrightarrow$ entropy.
  \item \emph{Edgeworth expansion}: corrections to the CLT in powers of $\kappa_3/\sqrt n$, $\kappa_4/n$.
  \item \emph{Exponential tilting} $f_\theta(x)=e^{\theta x-K(\theta)}f(x)$ (the canonical ensemble): the Esscher
        transform, a way to build a pricing measure.
\end{itemize}
Reference: Feller, vol.~2, ch.~XVI; Dembo--Zeitouni, \emph{Large Deviations}.
\end{answer}

\paragraph{Exponential families (2013).} A family of pdfs or pmfs indexed by a parameter vector $\theta$ is an
\emph{exponential family} if \ts{13}{3}{26:24}
\[
  f(x\mid\theta)=h(x)\,c(\theta)\exp\Bigl(\sum_{i=1}^kw_i(\theta)\,t_i(x)\Bigr),
\]
with $h,t_i$ depending only on $x$ and $c,w_i$ only on $\theta$. The lognormal is one \ts{13}{3}{29:37}:
expanding the square,
\[
  f(x)=\frac1x\cdot\frac{e^{-\mu^2/(2\sigma^2)}}{\sigma\sqrt{2\pi}}\cdot\exp\Bigl(-\frac{(\ln x)^2}{2\sigma^2}+\frac{\mu\ln x}{\sigma^2}\Bigr),
\]
so $h(x)=1/x$, $c(\theta)=e^{-\mu^2/(2\sigma^2)}/(\sigma\sqrt{2\pi})$, $w_1=-1/(2\sigma^2)$, $t_1=(\ln x)^2$, $w_2=\mu/\sigma^2$,
$t_2=\ln x$. The point, stressed by P.~Kempthorne from the audience \ts{13}{3}{28:35}: the joint density of an
i.i.d.\ sample $\prod_jf(x_j\mid\theta)=\prod_jh(x_j)\,c(\theta)^n\exp\bigl(\sum_iw_i(\theta)\sum_jt_i(x_j)\bigr)$ has the
same form and depends on the data only through the \emph{sufficient statistics} $\sum_jt_i(x_j)$. Normal,
lognormal, exponential, Poisson and gamma laws are exponential families (2013 PS2 B-3); Cauchy and Student $t$
are not.

\begin{coursenote}[Erratum in the 2013 notes]
In the lognormal example of \file{lecnote3} (p.~3) the factor is printed $h(x)=\frac12$; it should be
$h(x)=\frac1x$, as in the displayed formula just above it.
\end{coursenote}

\begin{table}[ht]
\centering\small
\begin{tabularx}{\linewidth}{@{}l l L@{}}
\toprule
\textbf{Law} & \textbf{pmf / pdf} & \textbf{Where it appears}\\
\midrule
Bernoulli$(p)$ & $p^x(1-p)^{1-x}$, $x\in\{0,1\}$ & default indicators, side of an order, event contracts (\cref{ch:kalshi})\\
Poisson$(\lambda)$ & $\lambda^ke^{-\lambda}/k!$, $k\in\N$ & number of events (orders, jumps, defaults) in a period; 2013 PS2 A-2, B-1\\
Exponential$(\lambda)$ & $\lambda e^{-\lambda x}$, $x\ge0$ & waiting times between events (memoryless); 2013 PS2 A-1\\
Normal$(\mu,\sigma^2)$ & $\frac{1}{\sigma\sqrt{2\pi}}e^{-(x-\mu)^2/(2\sigma^2)}$ & log-returns, Brownian increments, regression errors\\
Lognormal$(\mu,\sigma^2)$ & $\frac{1}{x\sigma\sqrt{2\pi}}e^{-(\ln x-\mu)^2/(2\sigma^2)}$, $x>0$ & prices (Black--Scholes)\\
Gamma$(\alpha,\lambda)$ & $\frac{\lambda^\alpha}{\Gamma(\alpha)}x^{\alpha-1}e^{-\lambda x}$, $x>0$ & contains $\chi^2_n=\mathrm{Gamma}(n/2,1/2)$\\
$\chi^2_n$, $t_r$, $F_{m,n}$ & \cref{ch03:sec-chisq} & inference in regression (\cref{ch:regression}); $t$: heavy-tailed returns\\
Cauchy & $1/(\pi(1+x^2))$ & a law without mean, variance or MGF\\
\bottomrule
\end{tabularx}
\caption{Distributions used in the course.}\label{ch03:tab-laws}
\end{table}\note*{plot all the mentioned distributions}

\begin{figure}[!htb]
\centering
\pgfplotsset{lawplot/.style={width=5.3cm,height=4.2cm,tick label style={font=\scriptsize},label style={font=\small},title style={font=\small},legend style={font=\scriptsize},legend pos=north east}}
\begin{tikzpicture}
\begin{axis}[lawplot,title={Bernoulli$(0.3)$, Poisson$(3)$},ybar=0pt,bar width=4pt,xmin=-0.7,xmax=9.7,ymin=0]
  \addplot[fill=defcol!60,draw=none] coordinates {(0,0.7) (1,0.3)};
  \addplot[fill=thmcol!60,draw=none] coordinates {(0,0.0498) (1,0.1494) (2,0.2240) (3,0.2240) (4,0.1680) (5,0.1008) (6,0.0504) (7,0.0216) (8,0.0081) (9,0.0027)};
  \legend{Bernoulli,Poisson}
\end{axis}
\end{tikzpicture}\hfill
\begin{tikzpicture}
\begin{axis}[lawplot,title={exponential, gamma, $\chi^2$},xmin=0,xmax=8,ymin=0,ymax=1.05]
  \addplot[thick,defcol,domain=0:8,samples=100] {exp(-x)};
  \addplot[thick,intcol,domain=0.01:8,samples=100] {x*exp(-x)};
  \addplot[thick,thmcol,domain=0.05:8,samples=100] {x^0.5*exp(-x/2)/(2^1.5*0.886227)};
  \legend{Exp$(1)$,{Gamma$(2,1)$},$\chi^2_3$}
\end{axis}
\end{tikzpicture}\hfill
\begin{tikzpicture}
\begin{axis}[lawplot,title={normal, $t_3$, Cauchy},xmin=-5,xmax=5,ymin=0,ymax=0.42]
  \addplot[thick,defcol,domain=-5:5,samples=120] {exp(-x^2/2)/sqrt(2*pi)};
  \addplot[thick,intcol,domain=-5:5,samples=120] {0.367553*(1+x^2/3)^(-2)};
  \addplot[thick,thmcol,domain=-5:5,samples=120] {1/(pi*(1+x^2))};
  \legend{{$N(0,1)$},$t_3$,Cauchy}
\end{axis}
\end{tikzpicture}
\begin{tikzpicture}
\begin{axis}[lawplot,title={lognormal$(0,\sigma^2)$},xmin=0,xmax=4,ymin=0,ymax=1.7]
  \addplot[thick,defcol,domain=0.02:4,samples=120] {exp(-ln(x)^2/(2*0.0625))/(x*0.25*sqrt(2*pi))};
  \addplot[thick,intcol,domain=0.02:4,samples=120] {exp(-ln(x)^2/(2*0.25))/(x*0.5*sqrt(2*pi))};
  \addplot[thick,thmcol,domain=0.02:4,samples=120] {exp(-ln(x)^2/2)/(x*sqrt(2*pi))};
  \legend{$\sigma=0.25$,$\sigma=0.5$,$\sigma=1$}
\end{axis}
\end{tikzpicture}\hfill
\begin{tikzpicture}
\begin{axis}[lawplot,title={tails, log scale},xmin=0,xmax=8,ymode=log,ymin=1e-6,ymax=1,legend pos=south west]
  \addplot[thick,defcol,domain=0:5.3,samples=80] {exp(-x^2/2)/sqrt(2*pi)};
  \addplot[thick,intcol,domain=0:8,samples=80] {0.367553*(1+x^2/3)^(-2)};
  \addplot[thick,thmcol,domain=0:8,samples=80] {1/(pi*(1+x^2))};
  \legend{{$N(0,1)$},$t_3$,Cauchy}
\end{axis}
\end{tikzpicture}\hfill
\begin{tikzpicture}
\begin{axis}[lawplot,title={$F_{m,n}$},xmin=0,xmax=4,ymin=0,ymax=1]
  \addplot[thick,defcol,domain=0.01:4,samples=120] {10.368*x^1.5*(1+0.5*x)^(-7.5)};
  \legend{$F_{5,10}$}
\end{axis}
\end{tikzpicture}
\caption{The laws of \cref{ch03:tab-laws}. The $t_3$ and Cauchy densities look like the normal near $0$ but
their tails decay like powers ($x^{-4}$, $x^{-2}$), not like $e^{-x^2/2}$: on the log scale the gap at $x=5$ is a
factor $10^3$--$10^4$.}\label{ch03:fig-laws}
\end{figure}

\begin{exercise}[PS2 (2024), Problem 2: lognormal distribution]\label{ch03:ex-ps2}
Let $X$ be lognormal$(\mu,\sigma^2)$, i.e.\ $Y=\ln X\sim N(\mu,\sigma^2)$.
\begin{enumerate}[(a)]
  \item Prove that $\mu_*=\E[X]=e^{\mu+\sigma^2/2}$.
  \item Prove that $\sigma_*^2=\Var X=e^{2\mu+2\sigma^2}-e^{2\mu+\sigma^2}=\mu_*^2\,(e^{\sigma^2}-1)$.
      For (a) and (b) use the MGF of the normal law; no integrals are needed.
  \item A stock has price $S_0$ at $t=0$ and $S_0X$ at time $t$ (years), where $X$ is lognormal with $\mu=t\ln1.15$
      and $\sigma=\sqrt t\,\ln1.3$. Find the mean and standard deviation of the total return $X$ for $t=1/12,\,3/12,\,1$.
  \item Show that for a call with strike $K$ on a lognormal$(\mu,\sigma^2)$ price,
      $\E[(X-K)^+]=e^{\mu+\sigma^2/2}\Phi\bigl(\frac{\mu-\ln K}{\sigma}+\sigma\bigr)-K\Phi\bigl(\frac{\mu-\ln K}{\sigma}\bigr)$.
  \item For the stock of (c) with $S_0=\$100$, compute the expected payoff of the call for $T=1/12,\,3/12,\,1$ and
      strikes $K=\$100$ (at the money) and $K=\$120$ (out of the money). Comment on the effect of a longer time to
      expiration.
\end{enumerate}
\end{exercise}

\begin{exercise}[PS2 (2013), A-1: exponential distribution]
$X$ has density $f_X(x)=\lambda e^{-\lambda x}$, $x\ge0$.
\begin{enumerate}[(a)]
  \item Compute its MGF.
  \item Deduce mean and variance from the MGF.
  \item Compute $\Prob(X>t)$ for $t>0$.
  \item Prove the memoryless property $\Prob(X>s+t\mid X>s)=\Prob(X>t)$ for all $s,t>0$.
  \item If $X_1,\dots,X_n$ are i.i.d.\ exponential$(\lambda)$, prove that $\min\{X_1,\dots,X_n\}$ is exponential.
  \item Entering a bank you find all three tellers busy and nobody waiting. Service times of the three customers and
      yours are i.i.d.\ exponential. What is the probability that you are the last of the four to leave?
\end{enumerate}
\end{exercise}

\begin{exercise}[PS2 (2013), A-2: Poisson distribution]
$X$ has pmf $f_X(k)=\lambda^ke^{-\lambda}/k!$, $k=0,1,2,\dots$
\begin{enumerate}[(a)]
  \item Compute its MGF.
  \item Deduce mean and variance.
  \item Prove, with MGFs, that a sum of independent Poisson
      variables is Poisson.
\end{enumerate}
\end{exercise}

\begin{exercise}[PS2 (2013), B-1: from exponential waiting times to Poisson counts]
Let $X_1,X_2,\dots$ be i.i.d.\ exponential$(\lambda)$ and $\tau=\max\{t\ge0:X_1+\dots+X_t\le1\}$, the number of events with
exponential waiting times in an interval of length 1.
\begin{enumerate}[(a)]
  \item Compute $\Prob(\tau=0)$.
  \item Compute $\Prob(\tau=n)$ using
      that $S_n=X_1+\dots+X_n$ has CDF $F_{S_n}(x)=1-\sum_{k=0}^{n-1}\frac1{k!}e^{-\lambda x}(\lambda x)^k$ ($x>0$), and conclude that
      $\tau$ is Poisson.
\end{enumerate}
\end{exercise}

\begin{exercise}[PS2 (2013), B-3: exponential families]
Prove that the lognormal, Poisson and exponential distributions are exponential families, i.e.\ that their
densities or pmfs can be written as $h(x)c(\theta)\exp\bigl(\sum_iw_i(\theta)t_i(x)\bigr)$ with $c\ge0$, $h\ge0$.
\end{exercise}

% ------------------------------------------------------------------
\section{Several random variables}\label{ch03:sec-multi}

Independence of $X$ and $Y$ means that knowing one gives no information about the other; for densities
it is the factorisation $f_{X,Y}=f_Xf_Y$ \ts{24}{4}{1:03:47}.

\begin{definition}[Covariance, correlation]
$\Cov(X,Y)=\E[(X-\E[X])(Y-\E[Y])]=\E[XY]-\E[X]\,\E[Y]$ (so $\Cov(X,X)=\Var X$) and
\[
  \Corr(X,Y)=\frac{\Cov(X,Y)}{\sqrt{\Var X\,\Var Y}}=\E\Bigl[\frac{X-\mu_X}{\sigma_X}\cdot\frac{Y-\mu_Y}{\sigma_Y}\Bigr]
\]
\ts{24}{4}{1:04:52}\ts{24}{4}{1:05:53}. $X,Y$ are \emph{uncorrelated} if $\Cov(X,Y)=0$.
\end{definition}
Covariance is bilinear; by Cauchy--Schwarz $|\Corr(X,Y)|\le1$, with equality iff $Y=a+bX$ almost surely.
Independent variables are uncorrelated, but not conversely: correlation detects only \emph{linear}
dependence \ts{24}{4}{1:06:58}. Example: $X\sim N(0,1)$ and $Y=X^2$ have $\Cov(X,Y)=\E[X^3]=0$, although $Y$ is a
function of $X$.

\begin{proposition}[Mean and variance of a linear combination]\label{ch03:prop-portvar}
Let $\bm X=(X_1,\dots,X_n)\T$ have mean vector $\bm\mu=\E[\bm X]$ and covariance matrix
$\Sigma=\E[(\bm X-\bm\mu)(\bm X-\bm\mu)\T]$, $\sigma_{ij}=\Cov(X_i,X_j)$ \ts{24}{4}{1:07:59}. For a constant $\bm a\in\R^n$ and
$Y=\bm a\T\bm X$ \ts{24}{4}{1:09:07},
\[
  \E[Y]=\bm a\T\bm\mu,\qquad
  \Var Y=\bm a\T\Sigma\,\bm a=\sum_ia_i^2\Var X_i+2\sum_{i<j}a_ia_j\Cov(X_i,X_j) .
\]
More generally $\E[B\bm X+\bm b]=B\bm\mu+\bm b$ and $\Cov(B\bm X+\bm b)=B\Sigma B\T$.
\end{proposition}
\begin{proof}
$\Var Y=\E[(\bm a\T(\bm X-\bm\mu))^2]=\E[\bm a\T(\bm X-\bm\mu)(\bm X-\bm\mu)\T\bm a]=\bm a\T\Sigma\bm a$: constants come out of the
expectation, which acts entry by entry \ts{24}{4}{1:10:10}. The matrix version is identical.
\end{proof}
In particular $\Sigma$ is positive semi-definite, since $\bm a\T\Sigma\bm a=\Var(\bm a\T\bm X)\ge0$ (\cref{sec:symmetric}).
If $\bm X$ are asset returns and $\bm a$ portfolio weights, $Y$ is the portfolio return: its variance is a
quadratic form in the weights.

\paragraph{Diversification} \ts{24}{4}{1:11:13}\ts{24}{4}{1:12:24}. If the $X_i$ are uncorrelated with common variance
$\sigma^2$ and $a_i=1/n$, then $\sum a_i^2=1/n$ and $\Var Y=\sigma^2/n$ \ts{24}{4}{1:13:27}. With a common correlation
$\rho$ instead (not discussed in class),
\[
  \Var\Bigl(\frac1n\sum_iX_i\Bigr)=\frac{\sigma^2}{n}+\frac{n-1}{n}\,\rho\sigma^2\ \xrightarrow[n\to\infty]{}\ \rho\sigma^2 .
\]

\begin{finance}[Systematic and idiosyncratic risk]
Diversification removes the part of the risk that is specific to each asset, but not the part common to
all of them: with $\rho=0.3$ and $\sigma=30\%$, even an infinitely diversified equal-weighted portfolio keeps a
volatility of $\sqrt{0.3}\cdot30\%\approx16\%$. Only this common, \emph{systematic} risk should command a risk
premium in equilibrium: this is the content of the CAPM (\cref{ch03:sec-capm}) and the starting point of
factor models (\cref{ch:pca}).
\end{finance}

\subsection{Conditional distributions and conditional expectation}\label{ch03:sec-condexp}
This material is only used implicitly in the lectures; it is needed for martingales (\cref{ch:sp1}) and in
2024 PS3 Problem~2.

\begin{definition}[Conditional density and expectation]
If $(X,Y)$ has joint density $f_{X,Y}$, the conditional density of $Y$ given $X=x$ (for $f_X(x)>0$) is
$f_{Y|X}(y\mid x)=f_{X,Y}(x,y)/f_X(x)$, and $m(x)=\E[Y\mid X=x]=\int y\,f_{Y|X}(y\mid x)\dd y$. The
\emph{conditional expectation} $\E[Y\mid X]$ is the random variable $m(X)$. Similarly
$\Var(Y\mid X)=\E[(Y-m(X))^2\mid X]$. (Discrete case: replace densities by pmfs.)
\end{definition}

\begin{proposition}[Properties]\label{ch03:prop-condexp}
For integrable $Y$ and functions $g$:
\begin{enumerate}[(i)]
  \item tower property (law of iterated expectations): $\E\bigl[\E[Y\mid X]\bigr]=\E[Y]$;
  \item taking out what is known: $\E[g(X)Y\mid X]=g(X)\,\E[Y\mid X]$; if $X,Y$ are independent, $\E[Y\mid X]=\E[Y]$;
  \item best prediction: $\E[Y\mid X]$ minimises $\E[(Y-h(X))^2]$ over all functions $h$;
  \item law of total variance: $\Var Y=\E\bigl[\Var(Y\mid X)\bigr]+\Var\bigl(\E[Y\mid X]\bigr)$.
\end{enumerate}
\end{proposition}
\begin{proof}
(i) $\int m(x)f_X(x)\dd x=\iint y\,f_{X,Y}(x,y)\dd y\dd x=\E[Y]$.

(ii) Given $X=x$, $g(X)=g(x)$ is a constant; under
independence $f_{Y|X}=f_Y$.

(iii) $\E[(Y-h)^2]=\E[(Y-m)^2]+\E[(m-h)^2]+2\E[(Y-m)(m-h)(X)]$, and by (i)--(ii) the
cross term is $2\E\bigl[(m-h)(X)\,\E[Y-m(X)\mid X]\bigr]=0$.

(iv) Apply the same decomposition to
$Y-\E[Y]=(Y-m)+(m-\E[Y])$: the cross term vanishes, $\E[(Y-m)^2]=\E[\Var(Y\mid X)]$ by (i), and $\E[m(X)]=\E[Y]$.
\end{proof}

\begin{intuition}[Conditional expectation is a projection]
In the Hilbert space of square-integrable variables with $\langle U,V\rangle=\E[UV]$, $\E[Y\mid X]$ is the orthogonal
projection of $Y$ on the subspace of functions of $X$: (iii) is Pythagoras. Linear regression projects on the
smaller subspace of affine functions $a+bX$ (\cref{ch:regression}); for jointly normal variables the two
projections coincide (\cref{ch03:sec-mvn}). For a general information set $\mathcal G$ (e.g.\ the price history up
to time $t$) $\E[Y\mid\mathcal G]$ is defined by the same orthogonality, $\E[(Y-\E[Y\mid\mathcal G])\,W]=0$ for every
bounded $\mathcal G$-measurable $W$: this is the object in the definition of a martingale.
\end{intuition}

\begin{exercise}[not from the course]\label{ch03:ex-pairwise}
Let $X_1,X_2$ be independent with $\Prob(X_i=1)=\Prob(X_i=-1)=\frac12$, and $X_3=X_1X_2$. Show that $X_1,X_2,X_3$ are pairwise
independent but not mutually independent.
\end{exercise}

% ------------------------------------------------------------------
\section{The multivariate normal distribution}\label{ch03:sec-mvn}

\begin{definition}[Multivariate normal]
$\bm X\in\R^n$ is \emph{multivariate normal}, $\bm X\sim N_n(\bm\mu,\Sigma)$, if every linear combination $\bm a\T\bm X$ is
normal (a constant counting as $N(c,0)$); then $\bm\mu=\E[\bm X]$ and $\Sigma=\Cov\bm X$. Components of such a vector are
called \emph{jointly normal}.
\end{definition}

\begin{proposition}\label{ch03:prop-mvn}
\leavevmode
\begin{enumerate}[(i)]
  \item $\bm X\sim N_n(\bm\mu,\Sigma)$ iff $\E[e^{\bm t\T\bm X}]=\exp(\bm t\T\bm\mu+\frac12\bm t\T\Sigma\bm t)$ for all $\bm t\in\R^n$.
  \item $\bm X\eqd\bm\mu+A\bm Z$ with $\bm Z$ a vector of i.i.d.\ $N(0,1)$ and any $A$ with $AA\T=\Sigma$ (e.g.\ Cholesky, or
        $A=U\Lambda^{1/2}U\T$ from \cref{thm:spectral}).
  \item If $\Sigma$ is invertible, $\bm X$ has density
        $f(\bm x)=(2\pi)^{-n/2}(\det\Sigma)^{-1/2}\exp\bigl(-\frac12(\bm x-\bm\mu)\T\Sigma^{-1}(\bm x-\bm\mu)\bigr)$.
\end{enumerate}
\end{proposition}
\begin{proof}
(i) $\bm t\T\bm X\sim N(\bm t\T\bm\mu,\bm t\T\Sigma\bm t)$; evaluate \eqref{ch03:eq-normalmgf} at $1$. Conversely this MGF restricted to
$\bm t=s\bm a$ is the MGF of $N(\bm a\T\bm\mu,\bm a\T\Sigma\bm a)$.

(ii) $\E[e^{\bm t\T(\bm\mu+A\bm Z)}]=e^{\bm t\T\bm\mu}\prod_i\E[e^{(A\T\bm t)_iZ_i}]
=e^{\bm t\T\bm\mu+\frac12|A\T\bm t|^2}$, and $|A\T\bm t|^2=\bm t\T\Sigma\bm t$; conclude by the $n$-dimensional version of
\cref{ch03:thm-mgf}.

(iii) Change variables $\bm x=\bm\mu+A\bm z$ in the density $(2\pi)^{-n/2}e^{-|\bm z|^2/2}$ of $\bm Z$:
$|\det A|=(\det\Sigma)^{1/2}$ and $|\bm z|^2=(\bm x-\bm\mu)\T\Sigma^{-1}(\bm x-\bm\mu)$.
\end{proof}

\begin{theorem}[Properties of the multivariate normal]\label{ch03:thm-mvn}
Let $\bm X\sim N_n(\bm\mu,\Sigma)$.
\begin{enumerate}[(i)]
  \item $B\bm X+\bm b\sim N(B\bm\mu+\bm b,\,B\Sigma B\T)$; in particular all marginals are normal.
  \item Jointly normal variables are independent iff they are uncorrelated.
  \item Partition $\bm X=(\bm X_1,\bm X_2)$ with $\Sigma_{22}$ invertible. Then
        \[
          \bm X_1\mid\bm X_2=\bm x_2\ \sim\ N\bigl(\bm\mu_1+\Sigma_{12}\Sigma_{22}^{-1}(\bm x_2-\bm\mu_2),\ \Sigma_{11}-\Sigma_{12}\Sigma_{22}^{-1}\Sigma_{21}\bigr).
        \]
\end{enumerate}
\end{theorem}
\begin{proof}
(i) Every linear combination of $B\bm X+\bm b$ is a linear combination of $\bm X$ plus a constant; the moments follow from
\cref{ch03:prop-portvar}.

(ii) If the cross-covariances vanish, the MGF of \cref{ch03:prop-mvn}(i) factorises into the MGFs
of the blocks.

(iii) Let $\bm\varepsilon=\bm X_1-\bm\mu_1-\Sigma_{12}\Sigma_{22}^{-1}(\bm X_2-\bm\mu_2)$. By (i) $(\bm\varepsilon,\bm X_2)$ is
jointly normal, and $\Cov(\bm\varepsilon,\bm X_2)=\Sigma_{12}-\Sigma_{12}\Sigma_{22}^{-1}\Sigma_{22}=0$, so $\bm\varepsilon$ is independent of
$\bm X_2$ by (ii), with mean $\bm0$ and covariance $\Sigma_{11}-2\Sigma_{12}\Sigma_{22}^{-1}\Sigma_{21}+\Sigma_{12}\Sigma_{22}^{-1}\Sigma_{22}\Sigma_{22}^{-1}\Sigma_{21}
=\Sigma_{11}-\Sigma_{12}\Sigma_{22}^{-1}\Sigma_{21}$. Given $\bm X_2=\bm x_2$, $\bm X_1$ is a constant plus $\bm\varepsilon$.
\end{proof}

\noindent\textbf{Bivariate case.} If $(X,Y)$ is jointly normal with means $\mu_X,\mu_Y$, standard deviations
$\sigma,\tau$ and correlation $\rho$, then
\begin{equation}\label{ch03:eq-bivreg}
  Y=a+bX+\varepsilon,\qquad b=\frac{\rho\tau}{\sigma},\quad a=\mu_Y-b\mu_X,\quad \varepsilon\sim N\bigl(0,(1-\rho^2)\tau^2\bigr)\ \text{independent of }X .
\end{equation}
So $\E[Y\mid X]=\mu_Y+\rho\frac{\tau}{\sigma}(X-\mu_X)$ is exactly the regression line, and the conditional variance does
not depend on $X$. This decomposition is the hint of 2024 PS3 Problem~1 (\cref{ch03:ex-stein}).

\paragraph{Geometry.} The level sets of the density are the ellipsoids $(\bm x-\bm\mu)\T\Sigma^{-1}(\bm x-\bm\mu)=c$, centred at
$\bm\mu$, with principal axes along the eigenvectors of $\Sigma$ and semi-axes $\sqrt{c\lambda_i}$; in two dimensions they
are ellipses tilted upward for $\rho>0$ and downward for $\rho<0$ \ts{24}{5}{3:12}\ts{24}{5}{4:20}. Since
$(\bm X-\bm\mu)\T\Sigma^{-1}(\bm X-\bm\mu)=|\bm Z|^2$, this quadratic form has the $\chi^2_n$ law of \cref{ch03:sec-chisq}.

\begin{remark}[Normal marginals are not enough; not discussed in class]
Let $Z\sim N(0,1)$ and $S=\pm1$ with probability $\frac12$ each, independent of $Z$, and $W=SZ$. Then $W\sim N(0,1)$ and
$\Cov(Z,W)=\E[S]\,\E[Z^2]=0$, yet $Z$ and $W$ are dependent ($|W|=|Z|$) and $Z+W=0$ with probability $\frac12$, so
$(Z,W)$ is not jointly normal. Joint normality is a strong assumption about \emph{dependence}, and in
particular it makes simultaneous extreme moves of several assets very unlikely.
\end{remark}

% ------------------------------------------------------------------
\section{Principal components of a random vector}\label{ch03:sec-pca}

\begin{coursenote}[Scope]
This section is the probabilistic part of PCA (slides 26--32 of \file{lec04_1}, \ts{24}{4}{1:13:27} to
\ts{24}{5}{21:04}). Applications to equity indices and to the yield curve \ts{24}{5}{15:51}, announced in L5 for
S.~Andreev's lecture, are in \cref{ch:pca}.
\end{coursenote}

Let $\bm x\in\R^m$ have $\E[\bm x]=\bm\alpha$ and $\Cov\bm x=\Sigma$ \ts{24}{4}{1:14:27}. $\Sigma$ is symmetric and positive
semi-definite \ts{24}{4}{1:15:32}, so by \cref{thm:spectral}
\begin{align*}
  \Sigma&=\Gamma\Lambda\Gamma\T=\sum_{i=1}^m\lambda_i\bm\gamma_i\bm\gamma_i\T,\\
  \Lambda&=\diag(\lambda_1,\dots,\lambda_m),\quad
  \lambda_1\ge\dots\ge\lambda_m\ge0,\quad
  \Gamma=[\bm\gamma_1\cdots\bm\gamma_m],\ \Gamma\T\Gamma=I_m .
\end{align*}

\begin{definition}[Principal component variables]
$\bm p=\Gamma\T(\bm x-\bm\alpha)$, i.e.\ $p_i=\bm\gamma_i\T(\bm x-\bm\alpha)$ \ts{24}{4}{1:17:39}.
\end{definition}

\begin{proposition}\label{ch03:prop-pca}
\leavevmode
\begin{enumerate}[(i)]
  \item $\E[\bm p]=\bm0$ and $\Cov\bm p=\Gamma\T\Sigma\Gamma=\Lambda$: the PC variables are uncorrelated, with variances
        $\lambda_1\ge\dots\ge\lambda_m$ \ts{24}{4}{1:18:39}; and $\bm x=\bm\alpha+\Gamma\bm p$.
  \item Total variance: $\sum_i\Var(x_i)=\tr\Sigma=\tr(\Lambda\Gamma\T\Gamma)=\sum_k\lambda_k=\sum_k\Var(p_k)$ \ts{24}{5}{21:04}.
  \item Variational characterisation \ts{24}{5}{17:53}: $\bm\gamma_1$ maximises $\Var(\bm w\T\bm x)=\bm w\T\Sigma\bm w$ over unit vectors
        $\bm w$; $\bm\gamma_k$ maximises it over unit vectors orthogonal to $\bm\gamma_1,\dots,\bm\gamma_{k-1}$ \ts{24}{5}{18:55}; the
        maxima are $\lambda_1,\dots,\lambda_k$. The last component has the \emph{smallest} variance, and $\lambda_m=0$ iff
        $\bm\gamma_m\T(\bm x-\bm\alpha)=0$ almost surely, an exact linear relation among the $x_i$.
\end{enumerate}
\end{proposition}
\begin{proof}
(i) and (ii) are \cref{ch03:prop-portvar} and $\tr(AB)=\tr(BA)$. (iii) Write $\bm w=\Gamma\bm c$, $|\bm c|=|\bm w|=1$:
$\bm w\T\Sigma\bm w=\sum_i\lambda_ic_i^2\le\lambda_1$, with equality at $\bm c=\bm e_1$. Orthogonality to $\bm\gamma_1,\dots,\bm\gamma_{k-1}$
means $c_1=\dots=c_{k-1}=0$, and then the bound is $\lambda_k$. A variable with zero variance is a.s.\ equal to its mean.
\end{proof}

\begin{intuition}[Shift and rotate]
$\bm x\mapsto\bm p$ moves the origin to the mean and rotates the axes onto the principal axes of the
density ellipsoids; no distance is changed \ts{24}{5}{8:39}. $\Sigma$ is the ``inertia tensor'' of the cloud of
outcomes and PCA diagonalises it. A zero eigenvalue is a flat direction: before the euro the Deutsche Mark and
the French franc moved so closely that the smallest PC of currency returns was almost degenerate
\ts{24}{5}{19:59}.
\end{intuition}

\paragraph{Factor form} \ts{24}{5}{9:43}\ts{24}{5}{10:48}. Split $\Gamma=[\Gamma_1\,\Gamma_2]$ and $\bm p=(\bm p_1,\bm p_2)$ with the $K<m$
largest eigenvalues in the first block. Then
\[
  \bm x=\bm\alpha+\Gamma_1\bm p_1+\Gamma_2\bm p_2=\bm\alpha+B\bm f+\bm\varepsilon,\qquad B=\Gamma_1,\ \bm f=\bm p_1,\ \bm\varepsilon=\Gamma_2\bm p_2 ,
\]
like a factor model with $K$ factors, except that $\Cov\bm\varepsilon=\Gamma_2\Lambda_2\Gamma_2\T$ is not diagonal
\ts{24}{5}{11:48}. If $\sum_{k\le K}\lambda_k/\sum_k\lambda_k$ is close to $1$, the first $K$ components carry almost all the
variability.

\paragraph{Empirical PCA} \ts{24}{5}{12:48}\ts{24}{5}{13:48}. For data $X=[\bm x_1\cdots\bm x_T]\in\R^{m\times T}$ (one column per
date): row means $\bar{\bm x}=\frac1TX\bm1_T$; de-meaned matrix $X^*=X-\bar{\bm x}\bm1_T\T$; sample covariance
$\hat\Sigma=\frac1TX^*X^{*\mathsf T}=\hat\Gamma\hat\Lambda\hat\Gamma\T$; sample PCs $P=\hat\Gamma\T X^*$ ($m\times T$). Equivalently, from the
(reduced) SVD $X^*=VDU\T$ one gets $\hat\Lambda=D^2/T$, $\hat\Gamma=V$, $P=DU\T$ \ts{24}{5}{14:49} (\cref{ch03:ex-pcasvd}).

\begin{coursenote}[Erratum in the slides]
Slide~30 of \file{lec04_1} describes $U$ as an ``$(m\times T)$ row-orthonormal matrix, $UV'=I_m$''. For $X^*=VDU'$ to be
conformable, $U$ must be $T\times m$ with orthonormal columns, $U\T U=I_m$ (equivalently $U\T$ is $m\times T$ with
orthonormal rows); $UV'$ is not even defined unless $T=m$. The formulas $\hat\Lambda=D^2/T$, $\hat\Gamma=V$, $P=DU'$ are
correct (for $T\ge m$).
\end{coursenote}

\begin{exercise}[\file{lec04_1}, slide 30: PCA through the SVD]\label{ch03:ex-pcasvd}
Let $X^*$ be the de-meaned $m\times T$ data matrix and $X^*=VDU\T$ its reduced SVD ($V$ orthogonal $m\times m$, $D$ diagonal
with $d_1\ge\dots\ge d_m\ge0$, $U$ of size $T\times m$ with $U\T U=I_m$). Show that $\hat\Lambda=\frac1TD^2$, $\hat\Gamma=V$ and
$P=\hat\Gamma\T X^*=DU\T$.
\end{exercise}

% ------------------------------------------------------------------
\section{Sampling distributions: \texorpdfstring{$\chi^2$}{chi-square}, \texorpdfstring{$t$}{t} and \texorpdfstring{$F$}{F}}\label{ch03:sec-chisq}

These laws arise from sums of squares of Gaussians and reappear in every test of \cref{ch:regression}
\ts{24}{5}{21:04}.

\begin{definition}[$\chi^2$ distribution]
If $Z\sim N(0,1)$, $U=Z^2$ has the \emph{chi-squared law with 1 degree of freedom}, $\chi^2_1$. If $Z_1,\dots,Z_n$ are
i.i.d.\ $N(0,1)$, $V=Z_1^2+\dots+Z_n^2\sim\chi^2_n$ \ts{24}{5}{22:06}.
\end{definition}
\begin{proposition}
$\chi^2_1=\mathrm{Gamma}(\frac12,\frac12)$ and $\chi^2_n=\mathrm{Gamma}(\frac n2,\frac12)$, with
\[
  f_U(u)=\frac{u^{-1/2}e^{-u/2}}{\sqrt{2\pi}},\qquad
  f_V(v)=\frac{v^{n/2-1}e^{-v/2}}{2^{n/2}\Gamma(n/2)},\qquad M_V(t)=(1-2t)^{-n/2}\ \ (t<\tfrac12),
\]
$\E[V]=n$, $\Var V=2n$.
\end{proposition}
\begin{proof}
$F_U(u)=\Prob(-\sqrt u\le Z\le\sqrt u)=2\Phi(\sqrt u)-1$, so $f_U(u)=\phi(\sqrt u)/\sqrt u$ (the map $z\mapsto z^2$ is two-to-one,
so \cref{ch03:thm-cov} is applied on each branch). For the Gamma$(\alpha,\lambda)$ density,
$\int_0^\infty\frac{\lambda^\alpha}{\Gamma(\alpha)}x^{\alpha-1}e^{-(\lambda-t)x}\dd x=(1-t/\lambda)^{-\alpha}$ for $t<\lambda$, which is
$(1-2t)^{-1/2}$ for $\alpha=\lambda=\frac12$. The MGF of a sum of $n$ independent copies is the $n$-th power, $(1-2t)^{-n/2}$,
the MGF of Gamma$(\frac n2,\frac12)$; conclude by \cref{ch03:thm-mgf}. Mean $\alpha/\lambda=n$, variance $\alpha/\lambda^2=2n$.
\end{proof}

\begin{coursenote}[Erratum in the slides]
Slide~34 of \file{lec04_1} calls $\lambda$ in Gamma$(\alpha,\lambda)$ ``the scale parameter''. With the density
$\lambda^\alpha x^{\alpha-1}e^{-\lambda x}/\Gamma(\alpha)$ used on slide~33, $\lambda$ is the \emph{rate}; the scale is $1/\lambda$ ($=2$ for
$\chi^2$). R uses the same convention: \texttt{dgamma(x, shape = a, rate = l)}.
\end{coursenote}

\begin{definition}[Student $t$ and Fisher $F$]
Let $Z\sim N(0,1)$, $U\sim\chi^2_r$, $V\sim\chi^2_n$, $W\sim\chi^2_m$, independent. Then
$T=Z/\sqrt{U/r}\sim t_r$ (\emph{$t$ distribution with $r$ degrees of freedom}) \ts{24}{5}{23:11} and
$F=\frac{W/m}{V/n}\sim F_{m,n}$ \ts{24}{5}{24:13}.
\end{definition}
Properties:
\begin{itemize}
  \item $t_r$ has density $\frac{\Gamma((r+1)/2)}{\sqrt{r\pi}\,\Gamma(r/2)}(1+x^2/r)^{-(r+1)/2}$: a power-law tail $|x|^{-(r+1)}$, so moments of
        order $\ge r$ are infinite; $\Var T=r/(r-2)$ for $r>2$ and $\kappa=3+6/(r-4)$ for $r>4$. $t_1$ is the Cauchy law,
        and $t_r\to N(0,1)$ as $r\to\infty$ because $U/r\to1$.
  \item $t_r$ has heavier tails than the normal: $Z$ is divided by a random factor that is only \emph{around} $1$
        \ts{24}{5}{24:13}.
  \item $\E[F]=\E[W/m]\,\E[n/V]=n/(n-2)$ for $n>2$ (using $\E[1/V]=1/(n-2)$), and $T^2\sim F_{1,r}$ if $T\sim t_r$.
  \item $F$ statistics compare two independent variance estimates (analysis of variance, regression $F$ tests)
        \ts{24}{5}{25:16}. When the null hypothesis is false, the numerator is a \emph{non-central} $\chi^2$ and the
        statistic has a non-central $F$ law, which governs the power of the test \ts{24}{5}{26:16}.
\end{itemize}
Where they come from (preview of \cref{ch:regression}): for an i.i.d.\ $N(\mu,\sigma^2)$ sample, $(n-1)s^2/\sigma^2\sim\chi^2_{n-1}$
independently of $\bar X$, hence $\sqrt n(\bar X-\mu)/s\sim t_{n-1}$ (Student, 1908).

\subsection{Heavy tails}\label{ch03:sec-heavy}

\begin{figure}[ht]
\centering
\begin{tikzpicture}
\begin{axis}[width=0.8\linewidth,height=5.6cm,ymode=log,xmin=0,xmax=6,ymin=1e-9,ymax=1,
  axis lines=left,xlabel={$x$ (in standard deviations)},ylabel={density (log scale)},
  legend style={font=\scriptsize,draw=none,fill=none},legend pos=south west,
  tick label style={font=\small},label style={font=\small}]
  \addplot[very thick,defcol,domain=0:6,samples=121] {exp(-x^2/2)/sqrt(2*pi)};
  \addlegendentry{$N(0,1)$}
  \addplot[very thick,thmcol,dashed,domain=0:6,samples=121] {0.49007*(1+x^2/3)^(-3)};
  \addlegendentry{$t_5$ rescaled to variance 1}
\end{axis}
\end{tikzpicture}
\caption{Right tails of the standard normal and of a Student $t_5$ scaled to unit variance (density
$0.490\,(1+x^2/3)^{-3}$, kurtosis 9). On a log scale the Gaussian tail is a parabola, the $t$ tail a
logarithm.}\label{ch03:fig-tails}
\end{figure}

With equal variances the two laws of \cref{ch03:fig-tails} look alike near the centre, but
$\Prob(|X|>4)$ is $6.3\times10^{-5}$ for the normal and $3.6\times10^{-3}$ for the rescaled $t_5$, about 56 times more.
For daily returns with 252 trading days a year, a 4-standard-deviation day happens on average once every
63~years under the normal model and about once a year under the $t_5$ model.

\begin{proposition}[Normal variance mixtures; not discussed in class]\label{ch03:prop-mixture}
Let $X\mid V\sim N(\mu,\sigma^2V)$ with $V>0$ random and $\E[V^2]<\infty$. Then $\E[X]=\mu$, $\Var X=\sigma^2\E[V]$ and
\[
  \kappa_X=3\,\frac{\E[V^2]}{(\E[V])^2}=3\Bigl(1+\frac{\Var V}{(\E[V])^2}\Bigr)\ \ge\ 3,
\]
with equality iff $V$ is constant.
\end{proposition}
\begin{proof}
Condition on $V$: $\E[(X-\mu)^2\mid V]=\sigma^2V$ and $\E[(X-\mu)^4\mid V]=3\sigma^4V^2$; take expectations (tower property).
\end{proof}
The $t_r$ law is such a mixture: $T=Z\sqrt{r/U}$, so $T\mid U\sim N(0,r/U)$. \emph{Stochastic volatility} models, in
which the variance of returns is itself random, are of this type, and always produce kurtosis above 3. 2024
PS3 Problem~2 extends Stein's lemma (\cref{ch03:sec-capm}) to them (\cref{ch03:ex-sv}).

\begin{intuition}[Superstatistics]
A Gaussian whose variance fluctuates is like a Brownian particle in a medium whose temperature fluctuates
slowly: locally Gaussian, globally fat-tailed. Physicists call this superposition of statistics
\emph{superstatistics}; in finance the fluctuating temperature is the volatility, which clusters in time
(\cref{ch:volatility,ch:timeseries}).
\end{intuition}

\begin{finance}[Heavy tails in practice]
Daily returns of stocks and indices have sample kurtosis far above 3, and many studies find tails decaying
roughly like a power law, with exponent near 3 for individual stocks (the ``inverse cubic law''). Consequences:
normal-based value-at-risk understates extreme losses (\cref{ch:var}); sample moments of high order are
unstable or meaningless; averages converge more slowly than Gaussian intuition suggests, and not at all for a
Cauchy-like law (\cref{ch03:ex-cauchy}).
\end{finance}

\begin{exercise}[not from the course: a heavy-tailed mixture]
Let $X\mid V\sim N(0,V)$ with $V$ exponential with rate 1. Show that $M_X(t)=1/(1-t^2/2)$ for $|t|<\sqrt2$, identify the law of
$X$ (Laplace, density $\frac1{\sqrt2}e^{-\sqrt2|x|}$), and check \cref{ch03:prop-mixture}: $\Var X=1$, $\kappa=6$.
\end{exercise}

% ------------------------------------------------------------------
\section{Limit theorems}\label{ch03:sec-limits}

\begin{definition}[Modes of convergence]
$X_n\to X$ \emph{in probability} if $\Prob(|X_n-X|>\varepsilon)\to0$ for every $\varepsilon>0$; \emph{in distribution} if
$F_{X_n}(x)\to F_X(x)$ at every continuity point of $F_X$; \emph{almost surely} if $\Prob(X_n\to X)=1$. Almost sure
convergence implies convergence in probability, which implies convergence in distribution.
\end{definition}

\begin{lemma}[Markov and Chebyshev inequalities]\label{ch03:lem-cheb}
For $a>0$: $\Prob(|X|\ge a)\le\E[|X|]/a$. If $\Var X=\sigma^2<\infty$ and $\varepsilon>0$:
$\Prob(|X-\mu|\ge\varepsilon)\le\sigma^2/\varepsilon^2$.
\end{lemma}
\begin{proof}
$a\,\ind\{|X|\ge a\}\le|X|$ pointwise; take expectations. For Chebyshev apply Markov to $(X-\mu)^2$ and $a=\varepsilon^2$: the
parabola $(x-\mu)^2$ lies above the step function $\varepsilon^2\ind\{|x-\mu|\ge\varepsilon\}$ (\cref{ch03:fig-cheb}), and averaging
preserves the inequality \ts{24}{5}{29:29}\ts{24}{5}{30:30}.
\end{proof}

\begin{figure}[ht]
\centering
\begin{tikzpicture}
\begin{axis}[width=0.62\linewidth,height=4.8cm,axis lines=middle,xmin=-2.3,xmax=2.3,ymin=0,ymax=4.4,
  xtick={-1,1},xticklabels={$-\varepsilon$,$\varepsilon$},ytick={1},yticklabels={$\varepsilon^2$},
  xlabel={$d$},x label style={anchor=west},tick label style={font=\small}]
  \addplot[very thick,defcol,domain=-2.05:2.05,samples=83] {x^2};
  \draw[very thick,thmcol] (axis cs:-2.2,1) -- (axis cs:-1,1);
  \draw[very thick,thmcol] (axis cs:-1,0.02) -- (axis cs:1,0.02);
  \draw[very thick,thmcol] (axis cs:1,1) -- (axis cs:2.2,1);
  \draw[dashed,black!40] (axis cs:-1,0) -- (axis cs:-1,1);
  \draw[dashed,black!40] (axis cs:1,0) -- (axis cs:1,1);
  \node[defcol,font=\small,anchor=west] at (axis cs:1.55,3.6) {$d^2$};
  \node[thmcol,font=\small,anchor=south] at (axis cs:-1.6,1.05) {$\varepsilon^2\ind\{|d|\ge\varepsilon\}$};
\end{axis}
\end{tikzpicture}
\caption{The picture behind Chebyshev's inequality, as drawn in class \ts{24}{5}{29:29}: with $d=\bar X_n-\mu$,
$d^2\ge\varepsilon^2\ind\{|d|\ge\varepsilon\}$, so $\Var\bar X_n=\E[d^2]\ge\varepsilon^2\Prob(|d|\ge\varepsilon)$.}\label{ch03:fig-cheb}
\end{figure}

\begin{theorem}[Weak law of large numbers]\label{ch03:thm-wlln}
Let $X_1,X_2,\dots$ be i.i.d.\ with $\E[X_i]=\mu$ and $\Var X_i=\sigma^2<\infty$, and $\bar X_n=\frac1n(X_1+\dots+X_n)$. Then
$\bar X_n\to\mu$ in probability \ts{24}{5}{27:20}\ts{13}{3}{49:04}:
$\Prob(|\bar X_n-\mu|>\varepsilon)\le\dfrac{\sigma^2}{n\varepsilon^2}\to0$.
\end{theorem}
\begin{proof}
$\E[\bar X_n]=\mu$ by linearity. $\Var\bar X_n=\frac1{n^2}\sum_{i,j}\Cov(X_i,X_j)=\frac{\sigma^2}n$, since the cross terms vanish
\ts{13}{3}{58:43}. Apply \cref{ch03:lem-cheb} \ts{24}{5}{28:28}.
\end{proof}

\begin{remark}
\leavevmode
\begin{itemize}
  \item The proof only uses that the $X_i$ are pairwise \emph{uncorrelated} with equal means and variances.
  \item The conclusion holds with only $\E[|X_i|]<\infty$ (Khinchin's WLLN) \ts{24}{5}{31:38}, and the \emph{strong}
        law gives $\bar X_n\to\mu$ almost surely under the same assumption \ts{13}{3}{1:02:53}.
  \item Without a finite mean it fails: the mean of $n$ i.i.d.\ Cauchy variables is again Cauchy (\cref{ch03:ex-cauchy}).
\end{itemize}
\end{remark}

\begin{example}[How many observations?]\label{ch03:ex-samplesize}
To be 99\% sure that $|\bar X_n-\mu|<\varepsilon=0.1$
\ts{13}{3}{1:00:50}, Chebyshev requires
$\sigma^2/(n\varepsilon^2)\le0.01$; with $\sigma=1$ (our choice) this
is $n\ge10{,}000$. The CLT below suggests
$n\approx(2.576\,\sigma/\varepsilon)^2\approx664$.\sloppy Chebyshev is valid
for every distribution with finite variance, hence crude; sharper
exponential bounds (Chernoff, Hoeffding) use MGF-type arguments, as
mentioned in class \ts{13}{3}{1:01:53}.
\end{example}

\begin{finance}[The casino, the poker room and the trader (2013)]
Casino games against the house are designed so that the house has an edge of about 1--5\% (2013 notes)
while the variance of each game is large compared with the edge \ts{13}{3}{50:22}. If a player wins an
even-money bet with probability $0.48$ (the indicative figure quoted in class for blackjack
\ts{13}{3}{51:22}), the expected P\&L per \$1 bet is $0.48-0.52=-\$0.04$ with a standard deviation of about
\$1: over an evening the edge is invisible, but the casino plays an enormous $n$ and the LLN makes its
profit almost certain \ts{13}{3}{52:22}. In poker the players play against each other and the house simply
charges a fee per hand (the \emph{rake}) \ts{13}{3}{54:32}; a player whose edge over the others exceeds the rake
can let the LLN work for him. The same logic drives a hedge fund or a high-frequency trader: a small edge
per trade, many (nearly) independent trades \ts{13}{3}{55:32}. The practical difficulty is the variance:
drawdowns make it hard to tell a real edge from bad luck \ts{13}{3}{56:34}, and correlated trades reduce the
effective $n$.
\end{finance}

\begin{theorem}[Central limit theorem]\label{ch03:thm-clt}
Let $X_1,X_2,\dots$ be i.i.d.\ with mean $\mu$ and variance $\sigma^2\in(0,\infty)$. Then
\[
  Z_n=\frac{1}{\sqrt n}\sum_{i=1}^n(X_i-\mu)=\sqrt n\,(\bar X_n-\mu)\ \longrightarrow\ N(0,\sigma^2)\quad\text{in distribution}
\]
\ts{24}{5}{31:38}\ts{13}{3}{1:07:29}.
\end{theorem}
\begin{proof}[Proof when the MGF exists]
Both courses prove it this way \ts{24}{5}{33:48}\ts{13}{3}{1:09:46}. Take $\mu=0$ and let $M$ be the MGF of $X_i$, finite near $0$.
By independence the expectation of the product is the product of expectations \ts{24}{5}{34:51}:
\[
  M_{Z_n}(t)=\E\Bigl[\prod_{i=1}^ne^{tX_i/\sqrt n}\Bigr]=M\Bigl(\frac{t}{\sqrt n}\Bigr)^n .
\]
$M$ is smooth near $0$ with $M(0)=1$, $M'(0)=\E[X]=0$, $M''(0)=\E[X^2]=\sigma^2$, so by Taylor
$M(s)=1+\frac12\sigma^2s^2+o(s^2)$ as $s\to0$ \ts{24}{5}{35:54}. For fixed $t$,
\[
  \ln M_{Z_n}(t)=n\ln\Bigl(1+\frac{\sigma^2t^2}{2n}+o\bigl(\tfrac1n\bigr)\Bigr)\ \longrightarrow\ \frac{\sigma^2t^2}{2},
\]
as in $(1+r/n)^n\to e^r$, the limit behind continuous compounding \ts{24}{5}{37:01}. So $M_{Z_n}(t)\to e^{\sigma^2t^2/2}$, the MGF
of $N(0,\sigma^2)$, and \cref{ch03:thm-mgf}(ii) concludes.
\end{proof}

\begin{remark}
\leavevmode
\begin{itemize}
  \item The statement ``$\bar X_n\approx N(\mu,\sigma^2/n)$'' is written with quotes in class, because a limit law cannot
        depend on $n$; it is the practical approximation \ts{24}{5}{32:47}.
  \item The MGF is only a convenience: with characteristic functions the same proof works whenever $\sigma^2<\infty$.
  \item Identical distribution is not needed (Lindeberg--Lyapunov CLT) \ts{24}{5}{38:04}, and weak dependence is
        allowed (mixing sequences, martingale differences).
\end{itemize}
\end{remark}

To a student's question whether the CLT is of any use when returns
are neither independent nor identically distributed \ts{24}{5}{38:04},
the answer in class: what matters is that a quantity be a sum of many
small, roughly independent contributions. Projections of a
high-dimensional vector on a random direction, and the residuals of a
regression, are such sums and are often close to Gaussian
\ts{24}{5}{39:09}\ts{24}{5}{40:17}.

\begin{example}[Estimating a mean, 2013 notes Example 4.2]
To estimate an unknown mean from i.i.d.\ observations one uses
$\bar X_n$ \ts{13}{3}{1:16:15}: by the LLN it converges to $\mu$, and by
the CLT its error is approximately $N(0,\sigma^2/n)$ \ts{13}{3}{1:17:21},
which gives the usual interval $\bar X_n\pm1.96\,s/\sqrt n$.
\end{example}

\begin{coursenote}[Is the sample mean the best estimator?]
In the 2013 lecture the sample mean is called the maximum likelihood estimator \ts{13}{3}{1:18:23}. It is the MLE
for normal (and Poisson, exponential) data but not in general, and the notes rightly add that for some
distributions it is not the best choice: for Cauchy data it does not converge at all, while the sample
median does; for heavy-tailed data robust estimators are preferable (\cref{ch:regression}).
\end{coursenote}

\begin{finance}[Aggregation of returns]
A $T$-day log-return is a sum of $T$ daily log-returns, so by the CLT it is approximately normal for long
horizons even when daily returns are fat-tailed (as long as their variance is finite): return distributions
look more Gaussian at monthly than at daily frequency. Volatility clustering slows this convergence down. The
same theorem makes the binomial tree converge to the lognormal model (\cref{ch:bs}) and the random walk to
Brownian motion (\cref{ch:sp1}).
\end{finance}

\begin{exercise}[not from the course: no LLN for Cauchy]\label{ch03:ex-cauchy}
Assume that the standard Cauchy law has characteristic function $e^{-|t|}$. Show that if $X_1,\dots,X_n$ are i.i.d.\
standard Cauchy, $\bar X_n$ is again standard Cauchy. Why does this not contradict \cref{ch03:thm-wlln}? What does it
suggest about averaging returns with very heavy tails?
\end{exercise}

% ------------------------------------------------------------------
\section{Probability theory for asset pricing: from expected utility to the CAPM}\label{ch03:sec-capm}

The slides \file{lec04_2} derive the capital asset pricing model with nothing more than this chapter:
normal payoffs, covariances and a Gaussian identity \ts{24}{5}{40:17}. The development follows F.~de Jong and
B.~Rindi, \emph{The Microstructure of Financial Markets} (2009) \ts{24}{5}{1:11:22}.

\subsection{One risky asset}
One period, from $t_0$ (trading) to $t_1$ (end of period) \ts{24}{5}{41:21}. Before $t_0$ an agent is endowed
with $Q_a$ shares of a risky asset and $Q_f$ units of a risk-free asset. At $t_0$ the risk-free asset costs \$1 per
unit and the risky asset $p_a$ per share \ts{24}{5}{42:26}, so the initial wealth is $w_0=Q_ap_a+Q_f$. At $t_1$ the
risk-free asset pays $1+r_f$ per unit and the risky one pays $\tilde F\sim N(\mu,\sigma^2)$ per share \ts{24}{5}{43:28}.
(A normal payoff can be negative; the model is useful if this has negligible probability \ts{24}{5}{44:28}.)
The agent buys $X$ shares at $t_0$ ($X<0$: sells), financed with the risk-free asset, so the final wealth is
\ts{24}{5}{45:31}
\begin{equation}\label{ch03:eq-wealth}
  \tilde w=(Q_a+X)\tilde F+(Q_f-Xp_a)(1+r_f),\qquad \frac{\partial\tilde w}{\partial X}=\tilde F-p_a(1+r_f).
\end{equation}

\paragraph{Expected utility.} The agent chooses $X$ to maximise $\E[U(\tilde w)]$ for a utility function with
$U'>0$ \ts{24}{5}{46:37} (more wealth is better) and $U''<0$ \ts{24}{5}{47:37} (decreasing marginal utility). Why not
simply $U(w)=w$ \ts{24}{5}{48:39}? Because then there is no penalty for risk \ts{24}{5}{49:47}: with $U$ concave,
Jensen's inequality gives $\E[U(\tilde w)]\le U(\E[\tilde w])$, so the agent prefers a sure amount to a gamble with the
same mean. The ratio $A(w)=-U''(w)/U'(w)$ is the (Arrow--Pratt) coefficient of absolute risk aversion.

\begin{coursenote}[Shannon, Kelly and the log-utility debate]
The lecturer recalls a controversy between Claude Shannon and (he thinks) Milton Friedman on maximising the
expected logarithm of wealth \ts{24}{5}{49:47}. The well-documented debate is between Paul Samuelson, who argued
against maximising expected log-wealth, and the proponents of the \emph{Kelly criterion} (J.~Kelly,
C.~Shannon, E.~Thorp); see W.~Poundstone, \emph{Fortune's Formula} (2005). Log utility has $A(w)=1/w$.
\end{coursenote}

\subsection{Stein's lemma}
\begin{lemma}[Stein]\label{ch03:lem-stein}
Let $Y$ and $Y^*$ be jointly normal with $Y\sim N(\mu,\sigma^2)$, and let $g$ be differentiable with
$\E[|g(Y)(Y-\mu)|]<\infty$ and $\E[|g'(Y)|]<\infty$. Then \ts{24}{5}{50:58}\ts{24}{5}{52:01}
\[
  \E[g(Y)(Y-\mu)]=\sigma^2\,\E[g'(Y)]\qquad\text{and}\qquad \Cov\bigl(g(Y),Y^*\bigr)=\E[g'(Y)]\,\Cov(Y,Y^*).
\]
\end{lemma}
\noindent The proof is 2024 PS3 Problem~1 (\cref{ch03:ex-stein}). The first identity is a Gaussian integration
by parts, the second reduces to the first through the decomposition \eqref{ch03:eq-bivreg}.

\begin{intuition}[Gaussian integration by parts]
For a centred Gaussian, $\E[Y\,g(Y)]=\sigma^2\E[g'(Y)]$ is the identity behind Wick's theorem and the
Schwinger--Dyson equations: with $g(y)=y^{k-1}$ it gives $\E[Y^k]=(k-1)\sigma^2\E[Y^{k-2}]$, hence $\E[Y^4]=3\sigma^4$. The second
form says that the covariance of a nonlinear function $g(Y)$ with any jointly Gaussian $Y^*$ is that of its
linearisation, $g'\cdot\Cov(Y,Y^*)$, with the slope $g'$ averaged over the law of $Y$. In the pricing argument
$g=U'$: it converts ``covariance of marginal utility with the payoff'' into ``covariance of wealth with the
payoff''.
\end{intuition}

\begin{exercise}[PS3 (2024), Problem 1: Stein's lemma]\label{ch03:ex-stein}
Prove \cref{ch03:lem-stein} (with $X$ in place of $Y$ and $Y$ in place of $Y^*$): if $X,Y$ are jointly normal with
$X\sim N(\mu,\sigma^2)$ and $g$ is differentiable with $\E[g(X)(X-\mu)]$ and $\E[g'(X)]$ finite, then
$\E[g(X)(X-\mu)]=\sigma^2\E[g'(X)]$ and $\Cov(g(X),Y)=\E[g'(X)]\Cov(X,Y)$.
\emph{Hint:} integrate by parts; you may use that for jointly normal $(X,Y)$ with standard deviations $\sigma,\tau$ and
correlation $\rho$, $Y=a+bX+\varepsilon$ with $a=\mu_Y-b\mu_X$, $b=\rho\tau/\sigma$ and $\varepsilon\sim N(0,(1-\rho^2)\tau^2)$ independent of $X$.
\end{exercise}

\begin{exercise}[PS3 (2024), Problem 2: Stein's lemma for stochastic volatility]\label{ch03:ex-sv}
(From A.~Gron, B.~J{\o}rgensen, N.~Polson, \emph{Optimal portfolio choice and stochastic volatility}, Applied
Stochastic Models in Business and Industry, 2012; cited as a 2011 working paper in the problem set.)
$X$ has \emph{stochastic volatility} if $X\mid V\sim N(\mu,\sigma^2V)$ with $V>0$ random with density $p(V)$; the
law of $X$ is $p(X)=\int_0^\infty p(X\mid V)p(V)\dd V$ and has heavy tails. Assume $0<\E[V]<\infty$ and define the
\emph{size-biased} density $q(V)=V\,p(V)/\E[V]$ and the induced $q(X)=\int_0^\infty p(X\mid V)q(V)\dd V$. The
authors' theorem states
\[
  \Cov\bigl(g(X),X\bigr)=\E^Q[g'(X)]\,\Var X,\qquad(**)
\]
with $\E^Q$ the expectation under the measure induced by $q(V)$. Their proof reads
\begin{align*}
  \Cov(g(X),X)&\overset{(1)}{=}\E\bigl[g(X)(X-\E[X])\bigr]\\
  &\overset{(2)}{=}\E_V\bigl[\E_{X|V}[g(X)(X-\E[X])]\bigr]\\
  &\overset{(3)}{=}\E_V\bigl[\E_{X|V}[g(X)(X-\E[X\mid V])]\bigr].
\end{align*}
\begin{enumerate}[(a)]
  \item Show that $q$ is a probability density on $(0,\infty)$.
  \item Prove (1).
  \item Does (2) follow from the law of
      iterated expectations?
  \item Prove (3) (compare $\E[X\mid V]$ and $\E[X]$).
  \item Applying Stein's lemma conditionally
      on $V$ they obtain $\Cov(g(X),X)=\E_V\{\E_{X|V}[g'(X)]\,\sigma^2V\}=\E\bigl[g'(X)\tfrac{V}{\E[V]}\bigr]\sigma^2\E[V]$ $(*)$. Does $(*)$
      imply $(**)$? Explain.
\end{enumerate}
\end{exercise}

\subsection{First-order condition and the equilibrium price}
Assume one may differentiate under the expectation, which holds by
dominated convergence if, e.g., $U'$ is bounded, or for the exponential
utility below since normal variables have exponential moments
\ts{24}{5}{53:03}\ts{24}{5}{54:16}. By \eqref{ch03:eq-wealth} the first-order
condition is
\begin{align*}
  0&=\frac{\partial}{\partial X}\E[U(\tilde w)]
  =\E\bigl[U'(\tilde w)\bigl(\tilde F-p_a(1+r_f)\bigr)\bigr] \\
   &=\Cov\bigl(U'(\tilde w),\tilde F\bigr)
   +\E[U'(\tilde w)]\bigl(\mu-p_a(1+r_f)\bigr),
\end{align*}
using $\E[AB]=\Cov(A,B)+\E[A]\,\E[B]$ and that a constant shift does not
change a covariance \ts{24}{5}{55:19}. Now $\tilde w$ is normal (affine
in $\tilde F$), so \cref{ch03:lem-stein} with $g=U'$, $Y=\tilde w$,
$Y^*=\tilde F$ gives
\begin{align*}
  \Cov(U'(\tilde w),\tilde F)&=\E[U''(\tilde w)]\Cov(\tilde w,\tilde F)\\
  &=\E[U''(\tilde w)](Q_a+X)\sigma^2
\end{align*}
\ts{24}{5}{56:22}\ts{24}{5}{57:28}. Solving for $p_a$:
\begin{equation}\label{ch03:eq-price1}
  p_a=\frac{1}{1+r_f}\Bigl[\E[\tilde F]
  +\frac{\E[U''(\tilde w)]}{\E[U'(\tilde w)]}\,(Q_a+X)\Var(\tilde F)\Bigr].
\end{equation}
Since $U''<0<U'$ the second term is negative: the price is the
discounted expected payoff \emph{minus a risk adjustment}
\ts{24}{5}{58:29}.

\paragraph{CARA utility} \ts{24}{5}{59:30}. For $U(w)=c-e^{-Aw}$ with
$A>0$: $U'(w)=Ae^{-Aw}>0$ and $U''(w)=-A^2e^{-Aw}=-A\,U'(w)$, so the
ratio in \eqref{ch03:eq-price1} is the constant $-A$ and
\begin{equation}\label{ch03:eq-cara}
  p_a=\frac{1}{1+r_f}\Bigl[\E[\tilde F]-A\,(Q_a+X)\Var(\tilde F)\Bigr].
\end{equation}
The risk adjustment is proportional to the risk aversion $A$, to the
position held $Q_a+X$ and to the variance of the payoff
\ts{24}{5}{1:00:36}.

\begin{coursenote}[Erratum in the slides]
Slide~6 of \file{lec04_2} writes $U'(\tilde w)=AU(\tilde w)$ and $U''(\tilde w)=-A^2U(\tilde w)$. For $U=c-e^{-Aw}$ one has
$U'=Ae^{-Aw}=A(c-U)$ and $U''=-A^2e^{-Aw}$; even with $c=0$ the signs are wrong ($U'=-AU$, $U''=A^2U$). What the
derivation uses, $U''/U'=-A$, and hence \eqref{ch03:eq-cara}, are correct.
\end{coursenote}

\begin{remark}[CARA and normal = mean--variance; not discussed in class]\label{ch03:rem-meanvar}
Since $-A\tilde w$ is normal, \eqref{ch03:eq-normalmgf} at $t=1$ gives $\E[U(\tilde w)]=c-\exp\bigl(-A\,\E[\tilde w]+\frac{A^2}{2}\Var\tilde w\bigr)$.
Maximising expected utility is therefore equivalent to maximising the \emph{certainty equivalent}
$\E[\tilde w]-\frac A2\Var\tilde w$, a mean--variance criterion (\cref{ch:portfolio}). Its maximiser is the demand
$Q_a+X=\bigl(\E[\tilde F]-(1+r_f)p_a\bigr)/(A\sigma^2)$, and \eqref{ch03:eq-cara} is this demand curve solved for the price:
the price at which the agent is content to hold $Q_a+X$ shares.
\end{remark}

\paragraph{Expected return.} With $\tilde r=(\tilde F-p_a)/p_a$, \eqref{ch03:eq-price1} gives
\[
  \E[\tilde r]=r_f-\frac{\E[U''(\tilde w)]}{\E[U'(\tilde w)]}\,\frac{(Q_a+X)\Var(\tilde F)}{p_a}
  \ \overset{\text{CARA}}{=}\ r_f+A\,\frac{(Q_a+X)\Var(\tilde F)}{p_a}:
\]
the risk-free rate plus a risk premium that grows with risk aversion, position and variance.

\subsection{Many risky assets}
Now the agent holds $\bm Q\in\R^n$ shares of $n$ risky assets priced $\bm p$ at $t_0$, whose payoffs
$\tilde{\bm F}\sim N_n(\bm\mu_F,\Sigma_F)$ are jointly normal \ts{24}{5}{1:01:39}\ts{24}{5}{1:02:45}, and buys $\bm X$:
$\tilde w=(\bm Q+\bm X)\T\tilde{\bm F}+(Q_f-\bm X\T\bm p)(1+r_f)$, again normal. The $n$ first-order conditions are handled
exactly as before, with Stein's lemma applied to the jointly normal pair $(\tilde w,\tilde F_i)$:
\[
  p_i=\frac{1}{1+r_f}\Bigl[\E[\tilde F_i]+\frac{\E[U''(\tilde w)]}{\E[U'(\tilde w)]}\Cov(\tilde w,\tilde F_i)\Bigr],\qquad
  \Cov(\tilde w,\tilde F_i)=\bigl[\Sigma_F(\bm Q+\bm X)\bigr]_i .
\]
For CARA, $\bm p=\bigl[\bm\mu_F-A\,\Sigma_F(\bm Q+\bm X)\bigr]/(1+r_f)$. With $\tilde r_i=(\tilde F_i-p_i)/p_i$, so that
$\Cov(\tilde w,\tilde r_i)=\Cov(\tilde w,\tilde F_i)/p_i$,
\[
  \E[\tilde r_i]=r_f-\frac{\E[U''(\tilde w)]}{\E[U'(\tilde w)]}\Cov(\tilde w,\tilde r_i)\ \overset{\text{CARA}}{=}\ r_f+A\Cov(\tilde w,\tilde r_i)
\]
\ts{24}{5}{1:03:48}. The premium of asset $i$ is set by its covariance with the agent's \emph{total} wealth, not by its
own variance: an asset that hedges wealth ($\Cov<0$) is expected to earn less than $r_f$.

\subsection{The capital asset pricing model}
\paragraph{Market portfolio.} In equilibrium prices must be such that the assets are willingly held. With
identical agents (or a representative agent) holding no initial risky endowment, $\bm Q=\bm0$, the agent's
equilibrium purchase $\bm X$ is the whole supply of risky assets: the \emph{market portfolio}
\ts{24}{5}{1:03:48}\ts{24}{5}{1:04:53}. Its payoff, price and return are
\[
  \tilde F_m=\sum_iX_i\tilde F_i,\qquad p_m=\sum_iX_ip_i,\qquad
  \tilde r_m=\frac{\tilde F_m-p_m}{p_m}=\sum_i\omega_i\tilde r_i,\quad \omega_i=\frac{X_ip_i}{p_m},
\]
with $\omega_i$ the market-value weights \ts{24}{5}{1:05:56}. The key identity is
$\tilde w=\tilde F_m+(Q_f-p_m)(1+r_f)=p_m\tilde r_m+\text{const}$, hence $\Cov(\tilde w,\cdot)=p_m\Cov(\tilde r_m,\cdot)$.

\begin{theorem}[CAPM, security market line]\label{ch03:thm-capm}
In this equilibrium, with $\beta_i=\Cov(\tilde r_i,\tilde r_m)/\Var(\tilde r_m)$,
\[
  \E[\tilde r_i]=r_f+\beta_i\bigl(\E[\tilde r_m]-r_f\bigr),\qquad
  p_i=\frac{\E[\tilde F_i]}{1+r_f+\beta_i\bigl(\E[\tilde r_m]-r_f\bigr)} .
\]
\end{theorem}
\begin{proof}
Let $\theta=-\E[U''(\tilde w)]/\E[U'(\tilde w)]>0$. Averaging the expected returns with weights $\omega_i$
\ts{24}{5}{1:07:04}: $\E[\tilde r_m]=r_f+\theta\Cov(\tilde w,\tilde r_m)=r_f+\theta\,p_m\Var(\tilde r_m)$, so
$\theta=(\E[\tilde r_m]-r_f)/(p_m\Var\tilde r_m)$. Then for each asset
$\E[\tilde r_i]-r_f=\theta\Cov(\tilde w,\tilde r_i)=\theta\,p_m\Cov(\tilde r_m,\tilde r_i)=\beta_i(\E[\tilde r_m]-r_f)$ \ts{24}{5}{1:08:10}. Finally
$\E[\tilde r_i]=\E[\tilde F_i]/p_i-1$ gives the price \ts{24}{5}{1:10:18}.
\end{proof}

\begin{figure}[ht]
\centering
\begin{tikzpicture}
\begin{axis}[width=0.62\linewidth,height=5cm,axis lines=left,xmin=0,xmax=2.1,ymin=0,ymax=15,
  xlabel={$\beta_i$},ylabel={$\E[\tilde r_i]$ (\%)},xtick={0,0.5,1,1.5,2},ytick={2,8,14},
  tick label style={font=\small},label style={font=\small}]
  \addplot[very thick,defcol,domain=0:2,samples=2] {2+6*x};
  \addplot[only marks,mark=*,mark size=2pt,thmcol] coordinates {(0,2) (1,8)};
  \node[font=\scriptsize,anchor=west] at (axis cs:0.05,1.3) {$r_f$};
  \node[font=\scriptsize,anchor=north west] at (axis cs:1.02,8) {market portfolio};
  \node[font=\scriptsize,anchor=south east,defcol] at (axis cs:1.9,13.6) {SML};
  \addplot[only marks,mark=o,mark size=2pt,intcol] coordinates {(0.6,8)};
  \node[font=\scriptsize,anchor=south,intcol] at (axis cs:0.6,8.3) {$\alpha>0$};
\end{axis}
\end{tikzpicture}
\caption{Security market line with illustrative numbers $r_f=2\%$, $\E[\tilde r_m]=8\%$. In the CAPM every asset lies
on the line; an asset above it (open dot) would have a positive ``alpha''.}\label{ch03:fig-sml}
\end{figure}

\begin{coursenote}[Notation and assumptions in \file{lec04_2}]
The slides denote the market weights by $\mu_i$, which clashes with the means; here they are $\omega_i$. The step
``the agent's portfolio is the market portfolio'' is the market-clearing condition with identical agents,
stated on slide~12 as an assumption. The CAPM was formulated by W.~Sharpe (1964) and J.~Lintner (1965) from
mean--variance portfolio choice \ts{24}{5}{1:11:22}; by \cref{ch03:rem-meanvar} the CARA-normal route taken here is
equivalent. A Lintner paper was distributed on Canvas; the OCW Week~3 page lists both Lintner (1965) papers.
\end{coursenote}

\begin{finance}[Reading and using the CAPM]
\begin{itemize}
  \item $\beta_i$ is the slope of the regression of $\tilde r_i$ on $\tilde r_m$ \ts{24}{5}{1:09:14} (\cref{def:alphabeta},
        \cref{ch:regression}). High-beta stocks must offer higher expected returns; $\beta<1$ is ``defensive''.
  \item Only systematic risk is priced: the idiosyncratic part of $\Var\tilde r_i$ can be diversified away
        (\cref{ch03:sec-multi}) and earns nothing.
  \item Practice (not discussed in class): the cost of equity in corporate valuation is often set to
        $r_f+\beta\times$(equity risk premium); fund performance is measured by the \emph{alpha} over the SML.
        Empirically the SML is flatter than predicted (low-beta stocks earn more than the CAPM says), which
        motivated multi-factor models (\cref{ch:pca,ch:portfolio}).
\end{itemize}
\end{finance}

% ------------------------------------------------------------------
\begin{keyformulas}
\begin{align*}
&\text{PIT:} && F_X(X)\sim U(0,1)\ (F_X\text{ continuous}),\quad F^{-1}(U)\sim F\\
&\text{Change of variables:} && f_Y(y)=f_X(h(y))\,|h'(y)|,\quad h=g^{-1}\\
&\text{Lognormal:} && f_Y(y)=\tfrac{1}{y\sigma\sqrt{2\pi}}e^{-(\ln y-\mu)^2/(2\sigma^2)},\quad \E[Y^k]=e^{k\mu+k^2\sigma^2/2}\\
&\text{Call payoff:} && \E[(X-K)^+]=\textstyle\int_K^\infty[1-F(x)]\dd x\\
&\quad\text{normal:} && (\mu-K)\Phi(d)+\sigma\phi(d),\quad d=(\mu-K)/\sigma\\
&\quad\text{lognormal:} && e^{\mu+\sigma^2/2}\Phi(d+\sigma)-K\Phi(d),\quad d=(\mu-\ln K)/\sigma\\
&\text{MGFs:} && M_{a+bX}(t)=e^{at}M_X(bt),\quad N(\mu,\sigma^2)\!:e^{\mu t+\sigma^2t^2/2},\quad \chi^2_n\!:(1-2t)^{-n/2}\\
&\text{Portfolio:} && \E[\bm a]\T\bm X=\bm a\T\bm\mu,\quad \Var\bm a\T\bm X=\bm a\T\Sigma\bm a\\
&\text{Conditional normal:} && \E[Y\mid X]=\mu_Y+\rho\tfrac{\tau}{\sigma}(X-\mu_X),\quad \Var(Y\mid X)=(1-\rho^2)\tau^2\\
&\text{PCA:} && \bm p=\Gamma\T(\bm x-\bm\alpha),\ \Cov\bm p=\Lambda,\ \tr\Sigma=\textstyle\sum\lambda_k\\
&\text{Chebyshev, CLT:} && \Prob(|\bar X_n-\mu|>\varepsilon)\le\tfrac{\sigma^2}{n\varepsilon^2},\quad \sqrt n(\bar X_n-\mu)\to N(0,\sigma^2)\\
&\text{Stein:} && \Cov(g(Y),Y^*)=\E[g'(Y)]\Cov(Y,Y^*)\\
&\text{CARA price:} && p_a=\bigl[\E[\tilde F]-A(Q_a+X)\Var\tilde F\bigr]/(1+r_f)\\
&\text{CAPM:} && \E[\tilde r_i]=r_f+\beta_i(\E[\tilde r_m]-r_f),\quad \beta_i=\Cov(\tilde r_i,\tilde r_m)/\Var\tilde r_m
\end{align*}
\end{keyformulas}

% !TEX root = ../main.tex
\chapter{Stochastic Processes I: Random Walks, Markov Chains, Martingales}\label{ch:sp1}
\begin{paracol}{2}

\begin{sources}
\begin{tabularx}{\linewidth}{@{}l l L@{}}
\toprule
\textbf{Lecture} & \textbf{Speaker} & \textbf{Content}\\
\midrule
2024 L5 (from 1:12:25) & P.~Kempthorne & Definition of a martingale; the simple random walk and simulations of its paths; the martingale $S_n^2-n\sigma^2$.\\
2024 L6 (first 67 min) & P.~Kempthorne & Product and exponential martingales; filtrations, non-anticipating strategies and martingale transforms; stopping times; gambler's ruin (fair and biased) by optional stopping; Markov processes and chains, $n$-step transitions; credit-rating migration; an up/down Markov chain for AAPL.\\
2013 L5 & C.~Lee & Processes as distributions over paths; simple random walk, its properties, gambler's ruin by recursion; Markov chains, transition matrix, stationary distribution via Perron--Frobenius; martingales, stopping times, optional stopping theorem.\\
\bottomrule
\end{tabularx}

\smallskip
\textbf{Files.} Slides \file{mit18_642_f24_lec05.pdf} (\emph{Stochastic Processes I}, 44 slides in 88 pages of
incremental builds); 2013 lecture notes \file{MIT18_S096F13_lecnote5.pdf}. The R script
\file{MC_Example1.r} used for the AAPL example (slide 44) was distributed on Canvas only and is not among
the OCW files.
\textbf{Problem sets.} 2024 PS2 Problem~1, PS3 Problem~3; 2013 PS2 Problems A-3, A-4, B-5.
\textbf{Elsewhere.} The first 72 minutes of 2024 L5 (PCA, $\chi^2$/$t$/$F$ distributions, LLN, CLT,
a one-period CAPM) are in \cref{ch:probability}; the second half of 2024 L6 (regression) is in
\cref{ch:regression}.
\end{sources}

\section*{Motivation}
A price history is one realisation of a \emph{stochastic process}. Standing at time $t$ we know the
past path and want to say something about the future: whether it depends on the past (can prices be
predicted?), how the process behaves in the long run, and how likely extreme ``boundary'' events are
\ts{13}{5}{8:25}. This chapter introduces three families of discrete-time processes that answer these
questions in increasing generality and that return, in continuous time, for the rest of the course:
\begin{itemize}
  \item the \emph{random walk}, the simplest model of a price and the discrete skeleton of Brownian
        motion (\cref{ch:sp2});
  \item \emph{Markov chains}, where the future depends on the past only through the present state
        (credit ratings, market regimes); they extend the stochastic matrices of \cref{sec:markov};
  \item \emph{martingales}, the mathematical model of a fair game. Discounted prices are martingales
        under the risk-neutral measure, gains from trading are martingale transforms (discrete
        stochastic integrals, \cref{ch:stochcalc}), and the optional stopping theorem turns
        hard probability questions into one-line computations.
\end{itemize}
For a physicist: the random walk is diffusion, a Markov chain is a master equation relaxing to
equilibrium, and a martingale is a quantity that is conserved \emph{on average}; optional stopping is the
use of such a conservation law to solve a boundary-value problem.

% ------------------------------------------------------------------
\section{Stochastic processes}\label{ch04:sec-processes}

\begin{definition}[Stochastic process]
A \emph{stochastic process} is a family $\{X_t,\ t\in\mathcal T\}$ of random variables on a common
probability space $(\Omega,\mathcal F,\Prob)$, indexed by \emph{time}: $\mathcal T=\{0,1,2,\dots\}$
(discrete time) or $\mathcal T=[0,\infty)$ (continuous time) \ts{13}{5}{1:03}. Equivalently, each outcome
$\omega\in\Omega$ selects a whole \emph{path} $t\mapsto X_t(\omega)$, so a process is a probability
distribution on a space of paths \ts{13}{5}{3:09}.
\end{definition}
The first reading is the intuitive one (a random value at each time); the second is the one needed to do
mathematics: at time $t$ we observe the initial segment of a path drawn once and for all, and we reason
about the unobserved rest of it.

\begin{example}[Same marginals, different processes {\ts{13}{5}{4:10}}]\label{ch04:ex-threeproc}
For $t\ge0$ consider
\begin{enumerate}[(a)]
  \item $X_t=t$ with probability~1;
  \item $X_t=\varepsilon t$ for all $t$, with a
      \emph{single} sign $\varepsilon=\pm1$ drawn with probability $1/2$ each;
  \item $X_t=\varepsilon_tt$ with an
      independent sign $\varepsilon_t$ for every $t$.
\end{enumerate}
In (b) and (c) every $X_t$ has the same distribution
($\pm t$ with probability $1/2$), but in (b) one observation reveals the entire future, while in (c) the
past says nothing about the future beyond $X_t\in\{t,-t\}$ \ts{13}{5}{7:20}. A process is \emph{not}
determined by its one-time marginals: what matters is the joint law of the path.
\end{example}

% ------------------------------------------------------------------
\section{The random walk}\label{ch04:sec-rw}

\begin{definition}[Random walk]\label{ch04:def-rw}
Let $X_1,X_2,\dots$ be i.i.d.\ with mean $\mu$ and variance $\sigma^2$, and let $S_0=0$,
$S_n=X_1+\dots+X_n$. $\{S_n\}$ is a \emph{random walk} with steps $X_i$. When
$\Prob(X_i=+1)=\Prob(X_i=-1)=\tfrac12$ it is the (one-dimensional) \emph{simple random walk} (SRW)
\ts{13}{5}{12:35}, \ts{24}{5}{1:14:36}.
\end{definition}

The path can be drawn as the graph of $n\mapsto S_n$, or as the trajectory of a walker taking unit steps
left or right on a line \ts{13}{5}{13:40}. The 2013 course and the 2024 problem sets also use the
\emph{biased} walk ($\Prob(X=+1)=p\ne\tfrac12$) and walks with a zero step (2024 PS2 Problem~1,
\cref{ch04:ex-ps2-1}).

\begin{proposition}[Basic properties]\label{ch04:prop-rw}
For a random walk:
\begin{enumerate}[(i)]
  \item $\E[S_n]=n\mu$ and $\Var(S_n)=n\sigma^2$; for the SRW $\E[S_n]=0$, $\Var(S_n)=n$ \ts{24}{5}{1:15:36};
  \item (independent increments) for $0\le k_0\le k_1\le\dots\le k_r$ the increments
        $S_{k_1}-S_{k_0},\dots,S_{k_r}-S_{k_{r-1}}$ are mutually independent \ts{13}{5}{20:56};
  \item (stationary increments) $S_{k+h}-S_k\eqd S_h$ for all $k,h\ge0$ \ts{13}{5}{22:03};
  \item (CLT) $(S_n-n\mu)/(\sigma\sqrt n)\to N(0,1)$ in distribution.
\end{enumerate}
\end{proposition}
\begin{proof}
(i) Linearity of expectation, and additivity of variance for independent summands.

(ii) The increment over $(k_{i-1},k_i]$ is a function of the block $X_{k_{i-1}+1},\dots,X_{k_i}$;
functions of disjoint blocks of independent variables are independent.

(iii) $S_{k+h}-S_k=X_{k+1}+\dots+X_{k+h}$ is a sum of $h$ i.i.d.\ copies of $X$, like $S_h$.

(iv) Central limit theorem, \cref{ch:probability}.
\end{proof}
Properties (ii) and (iii) are exactly those that define Brownian motion in continuous time
\ts{13}{5}{22:03}: the random walk is ``the fundamental stochastic process'', and much of
\cref{ch:sp2} is its continuous limit.

\paragraph{The $\sqrt n$ scale.} By (iv) the SRW at time $n$ is typically at distance of order
$\sqrt n$ from the origin, not $n$: although $|S_n|$ could be as large as $n$, the probability that
$|S_n|\le 2\sqrt n$ tends to $0.954$ \ts{13}{5}{17:51}. Simulated paths (slides 4--8) become more
jagged as $n$ grows and stay mostly inside the envelopes $\pm\sqrt n$, $\pm2\sqrt n$, $\pm3\sqrt n$
\ts{24}{5}{1:16:43}; see \cref{ch04:fig-rw}. Nevertheless the envelopes are crossed again and again: with
probability one the SRW returns to $0$ infinitely often, and for every $c>0$ it exceeds $c\sqrt n$
infinitely often \ts{13}{5}{18:52} (a consequence of the law of the iterated logarithm,
$\limsup_n S_n/\sqrt{2n\ln\ln n}=1$ a.s., not proved in the course).

\begin{figure}[ht]
\centering
\begin{tikzpicture}[x=0.058cm,y=0.075cm,>=Stealth]
  \draw[black!40] (0,-40) rectangle (200,40);
  \draw[black!25] (0,0) -- (200,0);
  \foreach \yy in {-40,-20,0,20,40} \draw[black!40] (0,\yy) -- (-2,\yy) node[left,font=\scriptsize,black]{$\yy$};
  \foreach \xx in {0,50,100,150,200} \draw[black!40] (\xx,-40) -- (\xx,-42) node[below,font=\scriptsize,black]{$\xx$};
  \node[below,font=\small] at (100,-47) {step $n$};
  \node[rotate=90,font=\small] at (-17,0) {$S_n$};
  \draw[thmcol,thick,domain=0:200,samples=100] plot (\x,{sqrt(\x)});
  \draw[thmcol,thick,domain=0:200,samples=100] plot (\x,{-sqrt(\x)});
  \draw[thmcol,thick,dashed,domain=0:200,samples=100] plot (\x,{2*sqrt(\x)});
  \draw[thmcol,thick,dashed,domain=0:200,samples=100] plot (\x,{-2*sqrt(\x)});
  \pgfmathsetseed{3}
  \foreach \n [evaluate=\n as \s using {\p+2*round(rnd)-1}, remember=\s as \p (initially 0)] in {1,...,200}
    {\draw[defcol] (\n-1,\p) -- (\n,\s);}
  \pgfmathsetseed{17}
  \foreach \n [evaluate=\n as \s using {\p+2*round(rnd)-1}, remember=\s as \p (initially 0)] in {1,...,200}
    {\draw[intcol] (\n-1,\p) -- (\n,\s);}
  \pgfmathsetseed{29}
  \foreach \n [evaluate=\n as \s using {\p+2*round(rnd)-1}, remember=\s as \p (initially 0)] in {1,...,200}
    {\draw[fincol] (\n-1,\p) -- (\n,\s);}
  \pgfmathsetseed{41}
  \foreach \n [evaluate=\n as \s using {\p+2*round(rnd)-1}, remember=\s as \p (initially 0)] in {1,...,200}
    {\draw[cscol] (\n-1,\p) -- (\n,\s);}
  \node[thmcol,font=\scriptsize,anchor=west] at (201,{sqrt(200)}) {$\sqrt n$};
  \node[thmcol,font=\scriptsize,anchor=west] at (201,{2*sqrt(200)}) {$2\sqrt n$};
  \node[thmcol,font=\scriptsize,anchor=west] at (201,{-sqrt(200)}) {$-\sqrt n$};
  \node[thmcol,font=\scriptsize,anchor=west] at (201,{-2*sqrt(200)}) {$-2\sqrt n$};
\end{tikzpicture}
\caption{Four simulated paths of the simple random walk (200 steps) with the one- and two-standard-deviation
envelopes $\pm\sqrt n$, $\pm2\sqrt n$; slides 4--8 of \file{lec05} show the same picture for up to 10\,000
steps \ts{24}{5}{1:16:43}.}\label{ch04:fig-rw}
\end{figure}

\begin{coursenote}[Slips in the 2013 lecture]
At \ts{13}{5}{16:47} the value $S_t$ is said to be approximately normal ``with mean 0 and variance
$\sqrt t$''; the variance is $t$ and the \emph{standard deviation} is $\sqrt t$, which is the scale meant
in the rest of the discussion. (The $1/t$ proposed by a student is the variance of the average $S_t/t$.)
\end{coursenote}

\begin{intuition}[Diffusion]
$S_n$ is the position of a particle receiving independent unit kicks: $\langle S_n^2\rangle=n$ is
Einstein's $\langle x^2\rangle=2Dt$ with $D=\tfrac12$ (unit step, unit time). Rescaling space by
$\sqrt n$ and time by $n$, $t\mapsto S_{\lfloor nt\rfloor}/\sqrt n$ converges to Brownian motion
(Donsker's theorem, \cref{ch:sp2}); the jagged look of the long simulated paths is the nowhere-differentiability
of Brownian paths in the making.
\end{intuition}

\begin{finance}[Random walks as price models]
L.~Bachelier (1900), cited in the 2013 notes, modelled prices on the Paris exchange as a random walk;
the 2013 notes call the SRW ``a (rather inaccurate) model of stock price''. The course refines it in
two directions that appear in the problem sets: a drift, calibrated to an expected daily price change
(2024 PS2 Problem~1, where an intraday price moves in 78 five-minute steps), and a random walk for
the \emph{logarithm} of the price, $\log P_t=\log P_{t-1}+X_t$ (2024 PS3 Problem~3), which keeps prices
positive and makes returns, not price changes, i.i.d. Time-varying volatility comes later
(\cref{ch:volatility}).
\end{finance}

\begin{exercise}[PS2 (2024), Problem 1: random walk with bias]\label{ch04:ex-ps2-1}
Let $X_n$ be i.i.d.\ steps with $\Prob(X_n=+1)=p$, $\Prob(X_n=0)=1-(p+q)$, $\Prob(X_n=-1)=q$, where $p,q>0$ and $p+q<1$
(symmetric if $p=q$, biased upwards if $p>q$), and $S_n=\sum_{i=1}^nX_i$, $S_0=0$.
\begin{enumerate}[(a)]
  \item Derive the moment generating function $M_{X_i}(t)$ of a step and the mean and variance of $X_i$.
  \item Derive the moment generating function $M_{S_n}(t)$ and the mean and variance of $S_n$.
  \item With $\mu$, $\sigma^2$ the mean and variance of a step, prove that $Z_n=(S_n-n\mu)/\sqrt{n\sigma^2}$ converges in
      distribution to $N(0,1)$.
  \item Model the intraday price of a stock as such a walk with 78 five-minute steps (the six and a half hours of NYSE
      trading), with $p=0.52$ and $q=0.47$. Which step size gives an expected daily price change of $+\$0.50$? With that
      step size, what is the standard deviation of the daily price change?
\end{enumerate}
\end{exercise}

\begin{exercise}[\file{lecnote5}, Proposition 2.1]
Prove the independent- and stationary-increment properties of the random walk (\cref{ch04:prop-rw}) directly from
the definition, and check which of them survive if the steps are independent but not identically distributed.
\end{exercise}

% ------------------------------------------------------------------
\section{Markov chains}\label{ch04:sec-mc}

\subsection{The Markov property}
\begin{definition}[Markov process, Markov chain]\label{ch04:def-markov}
A process $\{X_t\}$ is a \emph{Markov process} if for all times $u<s<t$, given the present value $X_s$,
the future value $X_t$ is independent of the past values $X_u$ \ts{24}{6}{36:30}:
\[
  \bigl[X_t\mid X_r,\ r\le s\bigr]\ \eqd\ \bigl[X_t\mid X_s\bigr].
\]
A discrete-time Markov process with a finite or countable \emph{state space} $S$ is a \emph{Markov chain}:
\begin{equation}\label{ch04:eq-markov}
  \Prob(X_{n+1}=j\mid X_0=i_0,\dots,X_{n-1}=i_{n-1},X_n=i)=\Prob(X_{n+1}=j\mid X_n=i)
\end{equation}
for all states and all $n$ \ts{24}{6}{39:38}.
\end{definition}
``The effect of the past on the future is summarised by the current state'' \ts{13}{5}{33:32}: two paths
arriving at the same state at time $s$ have the same future distribution, however they got there
\ts{24}{6}{38:30}.

\begin{example}[Random walks are Markov]
If $S_{n+1}=S_n+X_{n+1}$ with $X_{n+1}$ independent of $(S_0,\dots,S_n)$, then
$\Prob(S_{n+1}\in\cdot\mid S_0,\dots,S_n)$ is the law of $S_n+X$, a function of $S_n$ only. For the SRW,
$\Prob(S_{n+1}=s\mid S_0,\dots,S_n)=\tfrac12$ if $s=S_n\pm1$ and $0$ otherwise \ts{13}{5}{37:45}. The
state space $\mathbb Z$ is infinite, so there is no finite transition matrix \ts{13}{5}{43:12}.
\end{example}

\begin{exercise}[PS2 (2013), A-3: identify Markov processes]
Which of the following are Markov processes?
\begin{enumerate}[(a)]
  \item The simple random walk $S_t$.
  \item $X_t=|S_t|$.
  \item $X_0=0$,
      $X_{t+1}=X_t+(Y_t-\tfrac1\lambda)$ with $Y_t$ i.i.d.\ exponential with parameter $\lambda$.
  \item $X_0=0$,
      $X_{t+1}=X_t+Z_tZ_{t-1}\cdots Z_0$ with $Z_t$ i.i.d.\ log-normal.
  \item $X_0=0$, $X_{t+1}=X_t+W_t+W_{t-1}+\dots+W_0$ with
      $W_t$ i.i.d.\ normal.
\end{enumerate}
\end{exercise}

\subsection{Transition matrices}
The one-step transition probabilities $\Prob(X_{n+1}=j\mid X_n=i)$ together with the initial
distribution $p_i=\Prob(X_0=i)$ specify the law of the whole process \ts{24}{6}{41:45}. The chain has
\emph{stationary} (time-homogeneous) transition probabilities if these do not depend on $n$
\ts{24}{6}{42:52}; then they form the \emph{transition matrix}
\[
  P=(P_{ij})_{i,j\in S},\qquad P_{ij}=\Prob(X_{n+1}=j\mid X_n=i),\qquad P_{ij}\ge0,\quad \sum_{j}P_{ij}=1,
\]
a \emph{row}-stochastic matrix (row $i$ = where you are, column $j$ = where you go) \ts{24}{6}{43:58}.
All chains below are time-homogeneous.

\begin{coursenote}[Row versus column conventions]
The 2024 slides use the row-stochastic $P$ above, acting on
\emph{row} vectors of probabilities from the right. The 2013 notes (and
\cref{sec:markov}) use the column-stochastic $A=P\T$ acting on column
vectors: $\bm\pi(t+1)=A\bm\pi(t)$ is the transpose of
$\bm p^{(t+1)}=\bm p^{(t)}P$. This chapter uses $P$. In the 2013
lecture the machine example below was first written in the row form
and then ``flipped'' \ts{13}{5}{51:41}.
\end{coursenote}

\begin{proposition}[Law of the path]\label{ch04:prop-path}
$\Prob(X_0=i_0,X_1=i_1,\dots,X_n=i_n)=p_{i_0}\,P_{i_0i_1}P_{i_1i_2}\cdots P_{i_{n-1}i_n}$.
\end{proposition}
\begin{proof}
Write $A=\{X_0=i_0,\dots,X_{n-1}=i_{n-1}\}$ and $B=\{X_n=i_n\}$. Then $\Prob(A\cap B)=\Prob(A)\Prob(B\mid A)$
(joint = marginal $\times$ conditional), $\Prob(B\mid A)=\Prob(X_n=i_n\mid X_{n-1}=i_{n-1})$ by the Markov
property, and this equals $P_{i_{n-1}i_n}$ by stationarity. Induction on $n$ \ts{24}{6}{45:04}, \ts{24}{6}{48:09}.
\end{proof}

\begin{theorem}[$n$-step transitions, Chapman--Kolmogorov]\label{ch04:thm-nstep}
Let $P^{(n)}_{ij}=\Prob(X_{m+n}=j\mid X_m=i)$ (independent of $m$). Then
\[
  P^{(n)}_{ij}=\sum_{k\in S}P_{ik}\,P^{(n-1)}_{kj},\quad P^{(0)}=I,\qquad\text{hence}\qquad P^{(n)}=P^n,
  \quad P^{(m+n)}=P^{(m)}P^{(n)} .
\]
The distribution of $X_n$ is the row vector $\bm p^{(n)}=\bm p\,P^n$, i.e.\
$\Prob(X_n=k)=\sum_jp_jP^{(n)}_{jk}$ \ts{24}{6}{52:29}.
\end{theorem}
\begin{proof}[Proof (first-step analysis)]
Condition on the first step \ts{24}{6}{49:16}:
\begin{align*}
P^{(n)}_{ij}&=\sum_k\Prob(X_1=k\mid X_0=i)\,\Prob(X_n=j\mid X_1=k,X_0=i)
 =\sum_kP_{ik}\,\Prob(X_n=j\mid X_1=k)\\
 &=\sum_kP_{ik}\,\Prob(X_{n-1}=j\mid X_0=k)=\sum_kP_{ik}P^{(n-1)}_{kj},
\end{align*}
using the Markov property and then stationarity. In matrix form
$P^{(n)}=PP^{(n-1)}$ \ts{24}{6}{51:21}; induction gives $P^n$.
\end{proof}
The 2013 notes decompose on the \emph{last} step instead, $P^{(n)}_{ij}=\sum_kP^{(n-1)}_{ik}P_{kj}$, and
the lecture reads the two-step matrix as a matrix product (``sum over all intermediate states'')
\ts{13}{5}{42:08}. First-step analysis is the more useful habit: it is how hitting probabilities and
expected hitting times are computed (\cref{ch04:sec-ruin}).

\begin{example}[A machine that breaks {\ts{13}{5}{44:15}}]\label{ch04:ex-machine}
A machine is working (W) or broken (B) each day. If working it breaks the next day with probability
$0.01$; if broken it is repaired by the next day with probability $0.8$:
\[
  P=\begin{pmatrix}0.99&0.01\\0.8&0.2\end{pmatrix},\qquad
  P^2=\begin{pmatrix}0.9881&0.0119\\0.952&0.048\end{pmatrix}.
\]
What is the state distribution after ten years, $\bm e_W P^{3650}$ \ts{13}{5}{46:26}? The lecture
``cheats'': after so long the distributions on day 3650 and 3651 should coincide, so the answer $\bm\pi$ should
satisfy $\bm\pi=\bm\pi P$ \ts{13}{5}{47:27}. This is the two-state chain of \cref{ex:twostate} with
$p=0.01$, $q=0.8$: $\bm\pi=(0.8,\,0.01)/0.81=(0.98765,\,0.01235)$ and the second eigenvalue is
$1-p-q=0.19$, so $P^n$ converges to the matrix with both rows equal to $\bm\pi$ like $0.19^n$ and
$P^{3650}$ equals it to machine precision. In the long run the machine is broken on $1.23\%$ of the days.
\end{example}

\subsection{Stationary distributions and convergence}
\begin{definition}[Stationary distribution]
A probability (row) vector $\bm\pi$ on $S$ is \emph{stationary} if $\bm\pi P=\bm\pi$, i.e.\
$\pi_j=\sum_k\pi_kP_{kj}$ for all $j$ \ts{13}{5}{53:49}. If $X_0\sim\bm\pi$ then $X_n\sim\bm\pi$ for every
$n$ \ts{13}{5}{54:52}. Equivalently, $\bm\pi\T$ is an eigenvector of eigenvalue $1$ of the column-stochastic
$A=P\T$.
\end{definition}

\begin{theorem}[Convergence to equilibrium]\label{ch04:thm-conv}
Let $S$ be finite with $m$ states and suppose some power $P^k$ has all entries positive (in particular, if
all $P_{ij}>0$). Then there is a unique stationary distribution $\bm\pi$, all $\pi_j>0$, and for every
initial distribution
\[
  \bigl\|\bm p\,P^n-\bm\pi\bigr\|_1\le 2\,(1-m\delta)^{\lfloor n/k\rfloor},\qquad
  \delta=\min_{i,j}(P^k)_{ij}>0 ;
\]
in particular $P^{(n)}_{ij}\to\pi_j$ for all $i,j$: the chain forgets its initial state geometrically fast.
\end{theorem}

\paragraph{The course's argument} (2013 notes, Theorem~3.3, for $k=1$ \ts{13}{5}{49:36}).
$A=P\T$ is column-stochastic with positive entries. By \cref{prop:stochastic}
its eigenvalues satisfy $|\lambda|\le1$ and $1$ is one of them; by
Perron--Frobenius (\cref{thm:pf}) the dominant eigenvalue is simple, has a
positive eigenvector, and all other eigenvalues have $|\lambda|<1$. The positive
eigenvector, normalised, is $\bm\pi$, and the other modes of $A^n\bm p\T$
decay (as in the state equations of \cref{sec:eigen}), so
$A^n\bm p\T\to\bm\pi\T$ \ts{13}{5}{52:41}.

\begin{proof}[A self-contained proof (not from the course)]
For two probability vectors $\bm\mu,\bm\nu$ let $\bm d=\bm\mu-\bm\nu$, so
$\sum_id_i=0$. Then, writing $Q=P^k$,
\begin{align*}
  (\bm dQ)_j&=\sum_id_iQ_{ij}=\sum_id_i\,(Q_{ij}-\delta)\\
  &\quad\Rightarrow\quad
  \|\bm dQ\|_1\le\sum_i|d_i|\sum_j(Q_{ij}-\delta)=(1-m\delta)\,\|\bm d\|_1,
\end{align*}
because $Q_{ij}-\delta\ge0$ and each row of $Q$ sums to $1$; similarly
$\|\bm dP\|_1\le\|\bm d\|_1$. Hence $\bm\mu\mapsto\bm\mu Q$ is a contraction
of the (complete) simplex in the $\ell^1$ distance, and by Banach's
fixed-point theorem has a unique fixed point $\bm\pi$. Since
$(\bm\pi P)Q=(\bm\pi Q)P=\bm\pi P$, $\bm\pi P$ is also a fixed point of
$Q$, so $\bm\pi P=\bm\pi$; any stationary distribution of $P$ is fixed
by $Q$, hence equals $\bm\pi$. Finally
\begin{align*}
  \|\bm pP^n-\bm\pi\|_1&=\|(\bm p-\bm\pi)P^n\|_1\\
  &\le(1-m\delta)^{\lfloor n/k\rfloor}\|\bm p-\bm\pi\|_1,
\end{align*}
with $\|\bm p-\bm\pi\|_1\le2$, and $\pi_j=(\bm\pi Q)_j\ge\delta>0$.
\end{proof}

\begin{intuition}[Relaxation to equilibrium]
$\bm p^{(n+1)}=\bm p^{(n)}P$ is a discrete-time master equation. The stationary distribution is its
equilibrium and $1-|\lambda_2|$ (the spectral gap) is the inverse relaxation time, as in \cref{ex:twostate}.
The contraction proof has a probabilistic meaning (coupling): run two copies of the chain from different
starts; at every block of $k$ steps they have probability at least $m\delta$ of landing in the same state,
after which they can be glued together. The initial condition is forgotten as soon as the copies have
met.
\end{intuition}

\begin{coursenote}[The puzzle left open in the 2013 lecture]
For the machine example the lecture sums the two components of $A\bm v=\lambda\bm v$ ($A$ column-stochastic)
and gets $\sum_iv_i=\lambda\sum_iv_i$, hence $\lambda=1$; then notices that the same argument seems to
force \emph{every} eigenvalue to be $1$, which is absurd, and leaves the point open \ts{13}{5}{56:57}. The
resolution: the identity only says that $\lambda=1$ \emph{or} $\sum_iv_i=0$. Eigenvectors of the other
eigenvalues have components summing to zero (like $(1,-1)\T$ in \cref{ex:twostate}), so they cannot be
normalised to probability vectors. For the Perron--Frobenius eigenvector, which is positive, the sum is
non-zero and the argument correctly gives $\lambda_0=1$.
\end{coursenote}

\paragraph{When the hypotheses fail.}
\begin{itemize}
  \item \emph{Periodicity.} The 2013 notes (Example~3.2) consider the walk $X_{n+1}=X_n\pm1\pmod N$ on the
        cycle $\mathbb Z_N$, with stationary distribution $\pi_i=1/N$ (the matrix is doubly stochastic, 2013 PS2
        A-4). If $N$ is even the chain alternates between even and odd states, no power of $P$ is positive and
        $P^{(n)}_{ij}$ does not converge: a stationary distribution need not be a limit (compare the chain
        $\bigl(\begin{smallmatrix}0&1\\1&0\end{smallmatrix}\bigr)$ in \cref{sec:markov}).
  \item \emph{Reducibility.} If the state space splits into pieces that do not communicate, the stationary
        distribution need not be unique (2013 PS2 A-4(a)). States $i,j$ \emph{communicate} if each can be
        reached from the other in some number of steps; a chain whose transition graph is strongly connected
        is \emph{irreducible} \ts{24}{6}{1:03:03}. For a finite irreducible chain the stationary distribution is
        unique and positive; if moreover it is aperiodic, some $P^k$ is positive and \cref{ch04:thm-conv}
        applies (standard facts, not proved in the course).
  \item \emph{Infinite state space.} ``A corresponding theorem is not true'' (2013 notes): the SRW on $\mathbb Z$
        has no stationary distribution. Indeed $\pi_j=\tfrac12(\pi_{j-1}+\pi_{j+1})$ makes $j\mapsto\pi_j$ linear;
        being non-negative on all of $\mathbb Z$ it is constant, and a constant cannot sum to $1$.
\end{itemize}

\paragraph{Estimating a chain from data (not discussed in class).} Given an observed path with $n_{ij}$
transitions $i\to j$, \cref{ch04:prop-path} gives the likelihood $\prod_{i,j}P_{ij}^{\,n_{ij}}$ (times
$p_{i_0}$). Maximising $\sum_{ij}n_{ij}\log P_{ij}$ subject to $\sum_jP_{ij}=1$ with Lagrange multipliers gives
the empirical frequencies $\hat P_{ij}=n_{ij}/\sum_kn_{ik}$; this is what the R package \texttt{markovchain}
uses in the case study below.

\begin{exercise}[PS2 (2013), A-4]
\leavevmode
\begin{enumerate}[(a)]
  \item Construct a finite Markov chain whose stationary distribution is not unique.
  \item A chain is \emph{doubly stochastic} if, in addition to $\sum_jp_{ij}=1$, also $\sum_ip_{ij}=1$ for every $j$ (the
      problem phrases this as ``each row sums to one'' for the transposed matrix $A$ of the 2013 notes). Prove that the uniform
      distribution $\pi_j=1/m$ is stationary for a doubly stochastic chain on $m$ states.
\end{enumerate}
\end{exercise}

\begin{exercise}[PS2 (2013), B-5: walk with reflecting barriers]
A person walks on the line, stepping right with probability $b$ and left with probability $1-b$ at each period. He
starts in one of the positions $1,\dots,m$; a step to $0$ (or to $m+1$) is instantly reflected back to $1$ (or to $m$).
Equivalently, in position $1$ (position $m$) he stays put with probability $1-b$ (probability $b$).
\begin{enumerate}[(a)]
  \item Model this as a Markov chain and write its transition matrix.
  \item Find the stationary distribution.
\end{enumerate}
\end{exercise}

\subsection{Markov chains in finance}

\paragraph{Credit ratings.} Rating agencies (S\&P, Moody's, Fitch) grade corporate and sovereign borrowers
from AAA (highest) through AA, A, BBB (investment grade) to BB, B, CCC, CC, C (speculative grade) and D
(default) \ts{24}{6}{52:29}. The rating determines the interest rate at which a company can borrow by
issuing bonds \ts{24}{6}{53:29}. Modelling the rating of an issuer as a Markov chain with a one-year
\emph{migration matrix}, default being an absorbing state ($P_{DD}=1$), the probability of default within
$n$ years for an issuer rated $i$ today is $(P^n)_{iD}$. Slide~43 shows such a matrix (attributed to the
CreditMetrics Technical Document): a AAA issuer keeps its rating over one year with probability
$43.78\%$, while a CCC issuer is upgraded to B with probability $41.53\%$, stays CCC with $32.46\%$ and
defaults with $19.15\%$ \ts{24}{6}{54:32}.

\begin{coursenote}[On the migration table of slide 43]
In the widely quoted one-year matrices (S\&P, Moody's, CreditMetrics) roughly 90\% of AAA issuers keep
their rating; the AAA row on the slide, which sends more than half of the AAA issuers to AA, is unusual.
The OCW version marks the table's source as unknown; use its numbers as an illustration only.
\end{coursenote}

\begin{example}[Default probabilities from a migration matrix (not from the course)]\label{ch04:ex-credit}
With three states, investment grade (IG), high yield (HY) and default (D), take
\[
  P=\begin{pmatrix}0.97&0.025&0.005\\0.05&0.90&0.05\\0&0&1\end{pmatrix}.
\]
The cumulative default probabilities $(P^n)_{\cdot D}$ are $0.5\%,\ 3.46\%,\ 8.46\%$ for an IG issuer and
$5\%,\ 20.75\%,\ 33.87\%$ for an HY issuer at $n=1,5,10$ years; they are not $n$ times the one-year
probability, because issuers migrate between grades before defaulting. Since D is absorbing and reachable from
everywhere, the unique stationary distribution is $\bm e_D$: every issuer eventually defaults, and the interesting
quantities are transient ones.
\end{example}

\begin{finance}[Uses and caveats of rating chains (not discussed in class)]
Migration matrices drive credit-portfolio models (CreditMetrics), regulatory capital, and the pricing of
credit-risky bonds and credit derivatives (\cref{ch:counterparty,ch:quanto}). The Markov assumption is only an
approximation: recently downgraded issuers are more likely to be downgraded again (``rating momentum''),
and migration rates vary with the business cycle, so the matrix is not time-homogeneous.
\end{finance}

\begin{casestudy}[Up and down days of AAPL (slide 44, \file{MC_Example1.r})]
\textbf{Goal.} Check whether the direction of daily price moves has short memory \ts{24}{6}{55:39}.
\textbf{Data.} AAPL daily closing prices $P_t$ downloaded with \texttt{quantmod} (2013/14 to September 2024);
the chain is fitted on the last four and a half years \ts{24}{6}{58:51}.
\textbf{Method.} A day is Up if $P_t>P_{t-1}$ and Down if $P_t<P_{t-1}$. The two-day state on day $t$ is one
of UU, UD, DU, DD (moves on days $t-1$ and $t$), which gives a 4-state chain whose transitions are partly
forced: from UU one can only go to UU or UD, and so on (\cref{ch04:fig-mc}). The package
\texttt{markovchain} (Spedicato, 2017) estimates the transition matrix by empirical frequencies and draws the
transition graph \ts{24}{6}{59:56}. The same is repeated with three-day states (8 states) \ts{24}{6}{1:02:01}.
\textbf{Results} (as reported in class; no knitted output is available). Two up days tend to be followed by a
third ($\hat P(\mathrm{UU}\to\mathrm{UU})>0.5$), while after two down days a reversal is more likely
($\hat P(\mathrm{DD}\to\mathrm{DD})<0.5$) \ts{24}{6}{1:00:56}. In the 8-state chain, three down days are
followed by an up day with probability close to $0.6$; the largest probability of a down day, $0.510$,
follows a mixed state with two up days preceded by a down day \ts{24}{6}{1:05:15}. All states communicate
(the graph is strongly connected) \ts{24}{6}{1:04:05}.
\textbf{Lesson.} Discretising returns into a few states and fitting a Markov chain is a quick exploratory
tool, easy to rerun on other tickers or periods \ts{24}{6}{1:06:18}. \emph{Our remark:} with about 1100
trading days each of the 8 states is visited roughly 140 times, so an estimated probability near $0.5$ has a
standard error of about $\sqrt{0.25/140}\approx0.04$: $0.510$ is indistinguishable from $0.5$, and even
$0.6$ is only about two standard errors away, before accounting for having looked at 8 states. Note also that
the slide's definitions leave days with $P_t=P_{t-1}$ unclassified; one inequality should be
non-strict, as the lecturer points out \ts{24}{6}{55:39}.
\end{casestudy}

\begin{figure}[ht]
\centering
\begin{tikzpicture}[>=Stealth,every node/.style={font=\small},
    st/.style={circle,draw=defcol,thick,fill=defcol!6,minimum size=9mm,inner sep=0pt}]
  % machine chain
  \node[st] (W) at (0,0) {W};
  \node[st] (B) at (2.8,0) {B};
  \draw[->,thick] (W) to[bend left=25] node[above,font=\scriptsize]{$0.01$} (B);
  \draw[->,thick] (B) to[bend left=25] node[below,font=\scriptsize]{$0.8$} (W);
  \draw[->,thick] (W) to[loop above,looseness=6] node[above,font=\scriptsize]{$0.99$} (W);
  \draw[->,thick] (B) to[loop above,looseness=6] node[above,font=\scriptsize]{$0.2$} (B);
  \node[font=\footnotesize] at (1.4,-1.6) {(a) machine chain, \cref{ch04:ex-machine}};
  % two-day chain
  \begin{scope}[xshift=6.3cm,yshift=0.9cm]
  \node[st] (UU) at (0,0) {UU};
  \node[st] (UD) at (3,0) {UD};
  \node[st] (DU) at (0,-2.2) {DU};
  \node[st] (DD) at (3,-2.2) {DD};
  \draw[->,thick,intcol] (UU) to[loop left,looseness=6] (UU);
  \draw[->,thick,thmcol] (UU) -- (UD);
  \draw[->,thick,intcol] (UD) to[bend left=12] (DU);
  \draw[->,thick,thmcol] (UD) -- (DD);
  \draw[->,thick,intcol] (DU) -- (UU);
  \draw[->,thick,thmcol] (DU) to[bend left=12] (UD);
  \draw[->,thick,intcol] (DD) -- (DU);
  \draw[->,thick,thmcol] (DD) to[loop right,looseness=6] (DD);
  \node[font=\footnotesize] at (1.5,-3.4) {(b) two-day up/down states};
  \end{scope}
\end{tikzpicture}
\caption{(a) The two-state machine chain of the 2013 lecture. (b) Allowed transitions of the two-day
states of the AAPL case study: the second letter of today's state is the first letter of tomorrow's; green
arrows are followed by an up day, red ones by a down day. The estimated probabilities were shown in class
\ts{24}{6}{1:00:56}.}\label{ch04:fig-mc}
\end{figure}

\begin{exercise}[Time to default (not from the course)]
For the rating chain of \cref{ch04:ex-credit}, let $Q$ be the $2\times2$ block of transitions among the non-default
states. Show that the expected number of years before default, starting from IG or HY, is the vector
$(I-Q)^{-1}\bm 1$ (condition on the first step), and compute it.
\end{exercise}

% ------------------------------------------------------------------
\section{Martingales}\label{ch04:sec-mart}

\subsection{Information and conditional expectation}
Let $X_1,X_2,\dots$ be the observed process. The \emph{information set} $\mathcal F_n$ collects what is known at
time $n$, i.e.\ the values of $X_1,\dots,X_n$ (formally, the $\sigma$-algebra they generate) \ts{24}{6}{8:32}.
A random variable $Y$ is \emph{known at time $n$}, written $Y\in\mathcal F_n$, iff $Y=f(X_1,\dots,X_n)$ for
some function $f$; the increasing family $\{\mathcal F_n\}$ is a \emph{filtration}, and one writes
$\E[Z\mid\mathcal F_n]=\E[Z\mid X_1,\dots,X_n]$ \ts{24}{6}{13:00}. We use the following properties of
conditional expectation (\cref{ch:probability}), valid for integrable $Z$:
\begin{itemize}
  \item \emph{taking out what is known}: $\E[YZ\mid\mathcal F_n]=Y\,\E[Z\mid\mathcal F_n]$ if $Y\in\mathcal F_n$;
  \item \emph{tower property}: $\E\bigl[\E[Z\mid\mathcal F_n]\mid\mathcal F_m\bigr]=\E[Z\mid\mathcal F_m]$ for
        $m\le n$; in particular $\E\bigl[\E[Z\mid\mathcal F_n]\bigr]=\E[Z]$;
  \item \emph{independence}: $\E[Z\mid\mathcal F_n]=\E[Z]$ if $Z$ is independent of $X_1,\dots,X_n$.
\end{itemize}

\subsection{Definition and first properties}
\begin{definition}[Martingale]\label{ch04:def-mart}
A sequence $M_n=f_n(X_1,\dots,X_n)$, $n\ge1$, with $M_0$ a constant, is a \emph{martingale} with respect
to $\{X_n\}$ (or to $\{\mathcal F_n\}$) if for all $n\ge1$
\begin{equation}\label{ch04:eq-mart}
  \E[|M_n|]<\infty\qquad\text{and}\qquad \E[M_n\mid\mathcal F_{n-1}]=M_{n-1}
\end{equation}
\ts{24}{5}{1:13:28}. It is a \emph{submartingale} if $\E[M_n\mid\mathcal F_{n-1}]\ge M_{n-1}$ and a
\emph{supermartingale} if $\le$ (terms not used in class).
\end{definition}
Standing at time $n-1$, the best prediction of the next value is the current value: the conditional
expectations of the future are ``flat'' \ts{24}{6}{2:05}. If $M_n$ is the balance of a gambler, the game is
\emph{fair}: whatever happened so far, the expected gain of the next round is zero \ts{13}{5}{1:01:24}. The
balance of a roulette player is a supermartingale, not a martingale: the game is designed so that
$\E[M_{n+1}\mid\mathcal F_n]<M_n$ \ts{13}{5}{1:02:24}.

\begin{coursenote}[2013 versus 2024 definitions]
The 2013 notes define a martingale with respect to its own past, $X_t=\E[X_{t+1}\mid X_0,\dots,X_t]$, and
omit the integrability condition. The 2024 definition allows the conditioning information to be generated
by a different, ``driving'' sequence $X_n$ (as in the examples below). A martingale with respect to
$\{X_n\}$ is also a martingale with respect to its own past, by the tower property, since
$M_0,\dots,M_{n-1}$ are functions of $X_1,\dots,X_{n-1}$.
\end{coursenote}

\begin{proposition}\label{ch04:prop-mart}
If $\{M_n\}$ is a martingale then $\E[M_n\mid\mathcal F_m]=M_m$ for all $m\le n$, and $\E[M_n]=M_0$ for all
$n$. The increments $D_k=M_k-M_{k-1}$ satisfy $\E[D_k\mid\mathcal F_{k-1}]=0$; if they are square-integrable
they are uncorrelated, and $\Var(M_n)=\sum_{k=1}^n\E[D_k^2]$.
\end{proposition}
\begin{proof}
By the tower property, $\E[M_n\mid\mathcal F_m]=\E\bigl[\E[M_n\mid\mathcal F_{n-1}]\mid\mathcal F_m\bigr]
=\E[M_{n-1}\mid\mathcal F_m]$, and induction on $n-m$ (2013 notes, Proposition~4.2). Taking $m=0$ gives
$\E[M_n]=M_0$. For $j<k$, $D_j\in\mathcal F_{k-1}$, so
$\E[D_jD_k]=\E\bigl[D_j\,\E[D_k\mid\mathcal F_{k-1}]\bigr]=0$; hence the variance of the sum is the sum of
the variances.
\end{proof}

\subsection{Examples}\label{ch04:sec-mart-ex}
The 2024 slides give four families; each one is used later.

\begin{example}[Sums: random walks {\ts{24}{5}{1:13:28}}]\label{ch04:ex-m1}
If $X_n$ are independent with $\E[X_n]=0$, then $S_n=X_1+\dots+X_n$ is a martingale:
$\E[S_n\mid\mathcal F_{n-1}]=S_{n-1}+\E[X_n]=S_{n-1}$. In particular the SRW is a martingale
\ts{13}{5}{1:01:24}. With non-zero means, the \emph{compensated} walk $S_n-\sum_{k\le n}\E[X_k]$ is a
martingale.
\end{example}

\begin{example}[Squares: $S_n^2-n\sigma^2$ {\ts{24}{5}{1:17:44}}]\label{ch04:ex-m2}
If moreover $\Var(X_n)=\sigma^2$, then $M_n=S_n^2-n\sigma^2$ is a martingale:
\[
  \E[S_n^2\mid\mathcal F_{n-1}]=S_{n-1}^2+2S_{n-1}\E[X_n]+\E[X_n^2]=S_{n-1}^2+\sigma^2 .
\]
Its paths look very different from those of $S_n$ (slides 10--12): bounded below by $-n\sigma^2$, unbounded
above, with a long right tail \ts{24}{6}{3:09}. It measures how much the walk has spread compared with its
expected spread; its continuous-time analogue $B_t^2-t$ encodes the quadratic variation of Brownian
motion, the origin of the extra term in It\^o's formula (\cref{ch:stochcalc}).
\end{example}

\begin{example}[Products {\ts{24}{6}{4:15}}]\label{ch04:ex-m3}
If $X_n\ge0$ are independent with $\E[X_n]=1$ and $M_0=1$, then $M_n=X_1X_2\cdots X_n$ is a martingale:
$\E[M_n\mid\mathcal F_{n-1}]=M_{n-1}\E[X_n]=M_{n-1}$. The 2024 slides take $X_n\in\{0.5,\,1.5\}$ with
probability $1/2$ each: the wealth of someone who stakes everything, every period, on a fair bet that gains
or loses 50\% \ts{24}{6}{5:20}. The 2013 lecture takes $X_n=2$ with probability $1/3$ and $X_n=\tfrac12$
with probability $2/3$, again with mean $1$ \ts{13}{5}{1:03:24}.
\end{example}

The simulations in class show a few paths growing to 5--10 times the initial wealth, while most lose
almost everything \ts{24}{6}{5:20}. The calculation behind the pictures (ours): $\E[M_n]=1$ for every $n$, but
$\log M_n=\sum_k\log X_k$ is a random walk with drift $\E[\log X]=\tfrac12\ln0.75=-0.144$ per step, so
$M_n\to0$ almost surely by the law of large numbers. After 50 steps the median wealth is
$0.75^{25}\approx7.5\times10^{-4}$, and $\Prob(M_{50}\ge1)=\Prob(\mathrm{Bin}(50,\tfrac12)\ge32)\approx0.032$;
the mean $1$ is carried by rare paths (the largest possible value is $1.5^{50}\approx6\times10^8$). In the 2013
example $\E[\log X]=\tfrac13\ln2-\tfrac23\ln2<0$ as well.

\begin{intuition}[Mean versus typical behaviour]
For multiplicative processes the mean and the typical value part ways: $\E[M_n]$ is dominated by
exponentially rare paths, much as a partition function can be dominated by rare configurations. The
typical growth rate is $\E[\log X]$, not $\log\E[X]$; by Jensen's inequality $\E[\log X]<\log\E[X]=0$ unless $X$ is
constant. In finance this is the gap between arithmetic and geometric mean returns (``volatility drag''),
and the reason for the debate mentioned in class \ts{24}{5}{49:47} on whether an investor should maximise
the expected \emph{logarithm} of wealth (the Kelly criterion, of which Paul Samuelson was the best-known
critic).
\end{intuition}

\begin{example}[Exponential martingales {\ts{24}{6}{6:28}}]\label{ch04:ex-m4}
Let $Y_n$ be i.i.d.\ with moment generating function $\phi(\lambda)=\E[e^{\lambda Y}]<\infty$. Then
$X_n=e^{\lambda Y_n}/\phi(\lambda)$ are independent with mean $1$, so by \cref{ch04:ex-m3}
\begin{equation}\label{ch04:eq-expmart}
  M_n=\frac{e^{\lambda S_n}}{\phi(\lambda)^n},\qquad S_n=Y_1+\dots+Y_n,
\end{equation}
is a martingale for every fixed $\lambda$ \ts{24}{6}{7:30}. If $\lambda_0\ne0$ solves $\phi(\lambda_0)=1$,
then $M_n=e^{\lambda_0S_n}$ is a martingale. Since $\phi$ is convex with $\phi(0)=1$ and $\phi'(0)=\E[Y]$, such
a root exists (with the sign opposite to $\E[Y]$) as soon as $\E[Y]\ne0$ and $\phi$ grows beyond $1$ on that side,
e.g.\ for bounded steps taking both signs.
\end{example}
For $\pm1$ steps, $\phi(\lambda)=pe^\lambda+qe^{-\lambda}$ and $\phi(\lambda)=1$ is a quadratic equation in
$e^\lambda$ with roots $1$ and $q/p$: hence $(q/p)^{S_n}$ is a martingale (\cref{ch04:fig-ruin}, left)
\ts{24}{6}{33:19}. This is the key to the biased gambler's ruin (\cref{ch04:sec-ruin}). The exponent
$\lambda_0$ is the ``adjustment coefficient'' of actuarial ruin theory, and \eqref{ch04:eq-expmart} is the
discrete ancestor of the exponential martingale $e^{\lambda B_t-\lambda^2t/2}$ of Brownian motion and of
the change of measure behind risk-neutral pricing (\cref{ch:bs}).

\paragraph{Markov versus martingale.} The two notions are independent \ts{13}{5}{1:05:32}. The SRW is both.
A biased walk is Markov but not a martingale. A martingale that is not Markov (our example): toss a fair
coin $\xi_1$ and let $M_n=\xi_1+c(\xi_1)\sum_{k=2}^n\xi_k$ with further fair signs $\xi_k$ and step size
$c=1$ if $\xi_1=+1$, $c=2$ if $\xi_1=-1$. It is a martingale (a martingale transform, see below), but
$M_3=-1$ can be reached with either step size, so the law of $M_4$ given $M_3=-1$ depends on the earlier path.
Processes of neither kind are the rule rather than the exception.

\begin{remark}[Prices or log-prices?]\label{ch04:rem-logprice}
If $\log P_t$ is a random walk with i.i.d.\ steps $X_t$ (2024 PS3 Problem~3), the price itself satisfies
$\E[P_t\mid\mathcal F_{t-1}]=P_{t-1}\E[e^{X}]$, so $P_t$ is a martingale iff $\E[e^X]=1$. For normal steps
$\E[e^X]=e^{\mu+\sigma^2/2}$, i.e.\ iff $\mu=-\sigma^2/2$: a martingale log-price ($\mu=0$) gives an
\emph{upward}-drifting price. This lognormal correction (2024 PS2 Problem~2, \cref{ch:probability}) is the same
$-\sigma^2/2$ that appears in the drift of $\log S_t$ under Black--Scholes.
\end{remark}

% ------------------------------------------------------------------
\section{Martingale transforms and stopping times}\label{ch04:sec-stopping}

\subsection{Trading strategies as martingale transforms}
\begin{definition}[Non-anticipating sequence, martingale transform]
A sequence $\{A_n,\ n\ge1\}$ is \emph{non-anticipating} (or predictable) with respect to $\{\mathcal F_n\}$ if
$A_n\in\mathcal F_{n-1}$ for every $n$: $A_n$ is decided with the information available at time $n-1$
\ts{24}{6}{9:44}. The \emph{martingale transform} of a martingale $\{M_n\}$ by $\{A_n\}$ is
\[
  \tilde M_n=M_0+A_1(M_1-M_0)+A_2(M_2-M_1)+\dots+A_n(M_n-M_{n-1})
  =M_0+\sum_{k=1}^nA_k\,\Delta M_k
\]
\ts{24}{6}{10:51}.
\end{definition}

\begin{theorem}[Martingale transform theorem]\label{ch04:thm-transform}
If $\{M_n\}$ is a martingale and $\{A_n\}$ is non-anticipating and bounded, $|A_n|\le C_n$ for constants
$C_n$, then $\{\tilde M_n\}$ is a martingale with respect to $\{\mathcal F_n\}$ \ts{24}{6}{13:00}.
\end{theorem}
\begin{proof}
$\tilde M_n$ is a function of $X_1,\dots,X_n$ and
$\E[|\tilde M_n|]\le|M_0|+\sum_{k\le n}C_k\bigl(\E[|M_k|]+\E[|M_{k-1}|]\bigr)<\infty$. Taking out what is known,
\[
  \E[\tilde M_n-\tilde M_{n-1}\mid\mathcal F_{n-1}]=A_n\,\E[M_n-M_{n-1}\mid\mathcal F_{n-1}]=A_n\cdot0=0 .
  \qedhere
\]
\end{proof}

\begin{finance}[You cannot beat a fair game]
Read $M_n$ as the (discounted) price of an asset and $A_k$ as the number of shares held from $k-1$ to $k$,
chosen using only the information available at $k-1$ \ts{24}{6}{10:51}. Then $\tilde M_n-M_0=\sum_kA_k\Delta M_k$ is
the cumulative trading gain: exactly the ``holdings times price increment, holdings decided first'' sum of
\eqref{eq:pnl} and the discrete version of the It\^o integral $\int\theta_t\dd S_t$ (\cref{ch:stochcalc}).
The theorem says that if prices are martingales, \emph{every} bounded trading strategy has zero expected gain.
This is the content of the fundamental theorem of asset pricing read backwards: absence of arbitrage is
equivalent to the existence of a measure $\Q$ under which discounted prices are martingales; in one period
this is the relation $P_0=\beta\,\E_Q[P_T]$ of \cref{thm:ftap}.
\end{finance}

\subsection{Stopping times}
\begin{definition}[Stopping time]\label{ch04:def-stop}
A random variable $\tau$ with values in $\{0,1,2,\dots\}\cup\{\infty\}$ is a \emph{stopping time} with respect to
$\{\mathcal F_n\}$ if $\{\tau\le n\}\in\mathcal F_n$ for every $n$ (equivalently $\{\tau=n\}\in\mathcal F_n$)
\ts{24}{6}{14:04}.
\end{definition}
Viewing $\omega$ as a path, the event $\{\omega:\tau(\omega)\le n\}$ is decided by the first $n$ steps of the path:
at every time we know whether we have already stopped \ts{24}{6}{16:16}. A stopping time is a rule to quit a
game that looks only at what has happened so far \ts{13}{5}{1:09:40}.

\begin{example}[Stopping times and non-stopping times {\ts{13}{5}{1:10:49}}]\label{ch04:ex-stop}
In the coin-toss game with balance $S_n$:
\begin{enumerate}[(a)]
  \item the first time the balance reaches \$100, $\tau=\min\{n:S_n=100\}$, is a stopping time; so is the first
        time it reaches either \$100 or $-\$50$, and any first hitting time of a set;
  \item the time of the first peak (the first local maximum of the balance) is \emph{not} a stopping time:
        to know that $n$ is a peak one must see $S_{n+1}<S_n$ \ts{13}{5}{1:11:50};
  \item one step after the first peak \emph{is} a stopping time, as a student remarks \ts{13}{5}{1:12:51}.
\end{enumerate}
\end{example}

\begin{definition}[Stopped and truncated processes]
For a stopping time $\tau$, the value at the stopping time is $X_\tau=\sum_{n\ge0}X_n\,\ind\{\tau=n\}$ (on
$\{\tau<\infty\}$) \ts{24}{6}{17:24}. Since $\tau$ can be infinite, one works with the truncated times
$n\wedge\tau=\min(n,\tau)$ \ts{24}{6}{18:28} and the \emph{stopped process} $\{X_{n\wedge\tau},\ n\ge0\}$, which
follows $X$ until $\tau$ and then stays frozen. $X_{n\wedge\tau}=\sum_{k<n}X_k\ind\{\tau=k\}+X_n\ind\{\tau\ge n\}$ is a
function of $X_0,\dots,X_n$.
\end{definition}

\begin{theorem}[Stopped martingales are martingales]\label{ch04:thm-stopped}
If $\{M_n\}$ is a martingale and $\tau$ a stopping time, then $\{M_{n\wedge\tau}\}$ is a martingale; in
particular $\E[M_{n\wedge\tau}]=M_0$ for every $n$.
\end{theorem}
\begin{proof}
Let $A_k=\ind\{\tau\ge k\}=1-\ind\{\tau\le k-1\}$. It is bounded and $A_k\in\mathcal F_{k-1}$, because
$\{\tau\le k-1\}\in\mathcal F_{k-1}$. The transform telescopes up to the stopping time:
\[
  M_0+\sum_{k=1}^nA_k(M_k-M_{k-1})=M_0+\sum_{k=1}^{n\wedge\tau}(M_k-M_{k-1})=M_{n\wedge\tau},
\]
which is a martingale by \cref{ch04:thm-transform} \ts{24}{6}{20:29}. Then apply \cref{ch04:prop-mart}.
\end{proof}

\begin{coursenote}[Erratum: slides 17, 20 and 23 of \file{lec05}]
\begin{itemize}
  \item Slide 23 (proof above): ``else replace $M_n$ by $M_n'=M_n-M_{n-1}$'' should read $M_n'=M_n-M_0$ (to reduce to
        $M_0=0$ one subtracts the initial value, not the previous one). In the displayed identity the right-hand side
        should have $n$ in place of the summation index $k$:
        $\sum_{k=1}^nA_k(M_k-M_{k-1})=M_\tau\ind(\tau\le n-1)+M_n\ind(\tau\ge n)=M_{n\wedge\tau}$.
  \item Slide 20 calls $\{0,1,2,\dots\}$ the ``state space'' of the process; in the definition of a stopping time it
        plays the role of the \emph{time index set}, in which $\tau$ takes values (together with $\infty$).
  \item Slide 17: in the special case the sum should read $S_n=\sum_{i=1}^nY_i$, and $\lambda_0$ must be a root of
        $\phi(\lambda_0)=1$ other than the trivial root $\lambda_0=0$.
\end{itemize}
\end{coursenote}

\begin{exercise}[Stopping times (not from the course)]
Let $\sigma,\tau$ be stopping times. Show that $\sigma\wedge\tau$, $\sigma\vee\tau$ and $\tau+1$ are stopping times, but that
$\tau-1$ in general is not. Deduce that ``one step after the first peak'' is a stopping time and ``the first peak'' is
not (\cref{ch04:ex-stop}).
\end{exercise}

\subsection{The optional stopping theorem}
\cref{ch04:thm-stopped} gives $\E[M_{n\wedge\tau}]=M_0$ at every fixed $n$. The \emph{optional stopping theorem}
passes to the limit $n\to\infty$, where $M_{n\wedge\tau}\to M_\tau$ on $\{\tau<\infty\}$.

\begin{theorem}[Optional stopping]\label{ch04:thm-ost}
Let $\{M_n\}$ be a martingale and $\tau$ a stopping time. Then $\E[M_\tau]=\E[M_0]$ in each of the following
cases:
\begin{enumerate}[(i)]
  \item $\tau\le T$ for a constant $T$ (2013 notes, Theorem~5.3, ``Doob's optional stopping theorem, weak form'')
        \ts{13}{5}{1:14:06};
  \item $\tau<\infty$ almost surely and the stopped process is bounded, $|M_{n\wedge\tau}|\le K$ for all $n$;
  \item $\E[\tau]<\infty$ and the increments are bounded, $|M_n-M_{n-1}|\le c$.
\end{enumerate}
\end{theorem}
\begin{proof}
(i) For $n=T$, $M_{T\wedge\tau}=M_\tau$, so $\E[M_\tau]=\E[M_0]$ by \cref{ch04:thm-stopped}. (This is the 2013 proof:
$M_\tau=M_0+\sum_{i=0}^{T-1}(M_{i+1}-M_i)\ind\{\tau\ge i+1\}$, and each term has zero mean after conditioning on
$\mathcal F_i$, because $\{\tau\ge i+1\}$ is decided by time $i$.)

(ii) Since $\tau<\infty$, $M_{n\wedge\tau}=M_\tau$ for all $n\ge\tau$, so $M_{n\wedge\tau}\to M_\tau$ almost surely; by
bounded convergence $\E[M_\tau]=\lim_n\E[M_{n\wedge\tau}]=\E[M_0]$.

(iii) $|M_{n\wedge\tau}|\le|M_0|+c\,\tau$, an integrable bound, so dominated convergence applies as in (ii).
\end{proof}

``A mortal being has no winning strategy'' in a fair game (2013 notes): whatever rule you use to decide when to
quit, as long as it does not peek into the future and you cannot play for an unbounded time with unbounded credit,
your expected final balance equals your initial one \ts{13}{5}{1:06:34}, \ts{13}{5}{1:15:08}. The hypotheses are
not decorations: 2013 notes, Exercise~5.4 (\cref{ch04:ex-note54}), asks why stopping the SRW when it first
reaches \$100 does not contradict the theorem although $\E[S_\tau]=100$.

\begin{finance}[The ``martingale'' betting system (not discussed in class)]
Bet \$1 on a fair coin; after each loss double the stake; stop at the first win. With $A_k=2^{k-1}$ on
$\{\tau\ge k\}$ and $\tau$ the first win, the net gain at $\tau$ is $2^{\tau-1}-(1+2+\dots+2^{\tau-2})=+1$ with
certainty, although $\E[\tau]=2<\infty$. Condition (iii) fails because the stakes, i.e.\ the increments of the
transformed process, are unbounded, and so does (ii), since the loss before the first win is unbounded. At any fixed
horizon $n$ the game is still fair: with probability $1-2^{-n}$ you have won \$1, with probability $2^{-n}$ you have lost
$2^n-1$, and the mean is $0$. This doubling scheme is the historical meaning of ``martingale'' in gambling (the
lecturer recalls the other meaning, a horse's bridle \ts{24}{5}{1:18:49}). Its trading version, doubling a losing
position to win it back, needs unlimited capital and is a classic route to ruin.
\end{finance}

\begin{exercise}[PS3 (2024), Problem 3: random walk for the log price]\label{ch04:ex-ps3-3}
Let $P_t$ be an asset price, $S_t=\log P_t$ and $X_t=S_t-S_{t-1}=\log(P_t/P_{t-1})$, with the $X_t$ i.i.d., $\E[X_t]=\mu$
and $\Var X_t=\sigma^2>0$.
\begin{enumerate}[(a)]
  \item For which $(\mu,\sigma^2)$ is $\{S_t\}$ a Markov process? Explain.
  \item For which $(\mu,\sigma^2)$ is $\{S_t\}$ a martingale? Explain.
  \item Let $S_0=0$ and $\tau=\min\{t:\ S_t>0.20\ \text{or}\ S_t<-0.20\}$, the first time the log price rises above
      $0.20$ or falls below $-0.20$. If the steps are i.i.d.\ $N(\mu,\sigma^2)$, which conditions on $\mu$ and $\sigma^2$ make
      $\Prob(S_\tau>0.20)>\tfrac12$? (The problem statement says ``$S_t$ are i.i.d.\ Normal''; the steps $X_t$ are meant.)
\end{enumerate}
\end{exercise}

\begin{exercise}[A martingale that is not Markov (not from the course)]
For the process $M_n=\xi_1+c(\xi_1)\sum_{k=2}^n\xi_k$ of \cref{ch04:sec-mart-ex}, write $M_n$ as a martingale
transform, and compute $\Prob(M_4=-3\mid M_3=-1)$ and $\Prob(M_4=-3\mid M_3=-1,\ M_1=-1)$ to show that $\{M_n\}$ is not
Markov.
\end{exercise}

\begin{exercise}[\file{lecnote5}, Exercise 5.4]\label{ch04:ex-note54}
In the coin-toss game the gambler stops the first time her balance reaches \$100. By construction $\E[S_\tau]=100$. Does
this contradict the optional stopping theorem? Which hypothesis of \cref{ch04:thm-ost} fails?
\end{exercise}

% ------------------------------------------------------------------
\section{Gambler's ruin}\label{ch04:sec-ruin}
Two players bet \$1 per round on a fair coin until one of them is broke. Let $S_n$ be my net gain after $n$
rounds, with my bankroll $B$ and my opponent's bankroll $A$ ($A,B$ positive integers) \ts{24}{6}{21:35},
\ts{13}{5}{23:03}. The game ends at
\[
  \tau=\min\{n:\ S_n=+A\ \text{or}\ S_n=-B\},
\]
a stopping time (a first hitting time), and I take all my opponent's money iff $S_\tau=+A$ \ts{24}{6}{22:46}.

\begin{lemma}\label{ch04:lem-finite}
For a walk with steps $\pm1$ and $\Prob(X=+1)=p\in(0,1)$, $\Prob(\tau>n(A+B))\le(1-p^{A+B})^n$. Hence
$\tau<\infty$ almost surely and $\E[\tau]<\infty$.
\end{lemma}
\begin{proof}
Split time into blocks of $A+B$ steps. If all steps of a block are $+1$ the walk, wherever it is in $(-B,A)$,
leaves the interval during that block. Blocks are independent and each has only $+1$ steps with probability
$p^{A+B}$, so surviving $n$ blocks has probability at most $(1-p^{A+B})^n$. Then
$\E[\tau]=\sum_{m\ge0}\Prob(\tau>m)\le(A+B)\sum_{n\ge0}(1-p^{A+B})^n<\infty$.
\end{proof}

\begin{proposition}[Fair gambler's ruin]\label{ch04:prop-ruin}
For the SRW,
\[
  \Prob(S_\tau=+A)=\frac{B}{A+B},\qquad \Prob(S_\tau=-B)=\frac{A}{A+B},\qquad \E[\tau]=AB .
\]
\end{proposition}
\begin{proof}
$S_n$ is a martingale and $|S_{n\wedge\tau}|\le\max(A,B)$, so by \cref{ch04:lem-finite} and optional stopping (b),
$0=\E[S_\tau]=A\,\Prob(S_\tau=A)-B\,\bigl(1-\Prob(S_\tau=A)\bigr)$ \ts{24}{6}{23:52}, \ts{24}{6}{24:54}, which gives the
first formula. For the second, $M_n=S_n^2-n$ is a martingale (\cref{ch04:ex-m2} with $\sigma^2=1$), so by
\cref{ch04:thm-stopped}
\[
  \E\bigl[S_{n\wedge\tau}^2\bigr]=\E[n\wedge\tau]\qquad\text{for every }n.
\]
As $n\to\infty$ the left side tends to $\E[S_\tau^2]$ by bounded convergence ($S^2_{n\wedge\tau}\le\max(A^2,B^2)$) and the
right side to $\E[\tau]$ by monotone convergence ($n\wedge\tau\uparrow\tau$). Hence \ts{24}{6}{29:13}
\[
  \E[\tau]=\E[S_\tau^2]=A^2\,\frac{B}{A+B}+B^2\,\frac{A}{A+B}=AB .\qedhere
\]
\end{proof}

With equal bankrolls each player wins with probability $\tfrac12$, by symmetry \ts{13}{5}{25:08}. Playing until I
win \$100 or lose \$50 ($A=100$, $B=50$), I win with probability $1/3$ \ts{13}{5}{26:08}, and the game lasts
$5000$ rounds on average. In a fair game the richer player is more likely to win everything
\ts{24}{6}{27:05}, and the expected duration is the product of the bankrolls \ts{24}{6}{29:13}.

\begin{coursenote}[Erratum: slides 24--26 of \file{lec05}]
\begin{itemize}
  \item ``Problem: Solve for $\Prob(\tau=+A)$'' (and the title of slide 25) should read $\Prob(S_\tau=+A)$: $\tau$ is a
        time, $S_\tau$ the level hit. A student pointed this out in class and the lecturer confirmed \ts{24}{6}{26:05}.
  \item ``$\lim_{n\to\infty}\E[S_{n\wedge\tau}]=S_\tau$ w.p.~1'' mixes a number and a random variable: the correct
        statement is $S_{n\wedge\tau}\to S_\tau$ almost surely (because $\Prob(\tau=\infty)=0$), and then
        $\E[S_\tau]=\lim_n\E[S_{n\wedge\tau}]=0$ by bounded convergence.
  \item Slide 26: ``$M_{n\wedge\tau}\le\max(A^2,B^2)+\tau$, bounded r.v.'' is not a valid justification, since $\tau$ is not
        bounded. The limit must be taken separately in the two terms of $M_{n\wedge\tau}=S^2_{n\wedge\tau}-n\wedge\tau$,
        with bounded and monotone convergence, as in the proof above.
\end{itemize}
\end{coursenote}

\paragraph{The 2013 route: first-step analysis.} Let $f(k)$ be the probability of hitting $+A$ before $-B$ when
starting from $k\in\{-B,\dots,A\}$ \ts{13}{5}{29:19}. Conditioning on the first toss and using the stationarity of
the increments,
\[
  f(k)=\tfrac12f(k+1)+\tfrac12f(k-1)\quad(-B<k<A),\qquad f(A)=1,\quad f(-B)=0
\]
\ts{13}{5}{30:20}. The increments $f(k+1)-f(k)$ are therefore constant: $f(-B+r)=\alpha r$ with
$\alpha=1/(A+B)$, and $f(0)=B/(A+B)$ (2013 notes, Section~2).

\begin{intuition}[Harmonic functions and electrostatics]
The recursion says that $f$ is \emph{discrete-harmonic}: its value is the average of its neighbours, i.e.\ its discrete
Laplacian vanishes. Equivalently $f(S_{n\wedge\tau})$ is a martingale, which is why the two methods agree. The hitting
probability is the solution of a Dirichlet problem, like the potential between two plates held at potentials $0$ and
$1$, hence linear. In continuous time the same idea links expectations of Brownian functionals to PDEs (backward
Kolmogorov and Feynman--Kac equations, used to price options in \cref{ch:bs}).
\end{intuition}

\begin{remark}[A student's symmetry argument {\ts{13}{5}{27:11}}]
For $A=100$, $B=50$ a student guessed $1/3$ as follows. Let $g(a,b)$ be the probability of gaining $a$ before losing
$b$. To gain $100$ before losing $50$, the walk must first reach $+50$ before $-50$ (probability $\tfrac12$); from $+50$
it must then gain $50$ before losing $100$, with probability $g(50,100)=1-g(100,50)$ by the reflection
$S\mapsto-S$. So $g=\tfrac12(1-g)$ and $g=1/3$. The lecturer was unsure it could be made rigorous; it can: restarting the
walk at a stopping time gives again a SRW independent of the past (the \emph{strong Markov property}, not proved in
this course).
\end{remark}

\paragraph{The biased walk.} Now $\Prob(X=+1)=p$, $\Prob(X=-1)=q=1-p$, $p\ne\tfrac12$ \ts{24}{6}{30:16}. First-step
analysis still works but is laborious \ts{24}{6}{31:16}; the exponential martingale solves the problem in two lines
\ts{24}{6}{32:19}.

\begin{proposition}[Biased gambler's ruin]\label{ch04:prop-biased}
With $r=q/p\ne1$,
\[
  \Prob(S_\tau=+A)=\frac{r^B-1}{r^{A+B}-1}.
\]
\end{proposition}
\begin{proof}
By \cref{ch04:ex-m4}, $M_n=r^{S_n}$ is a martingale with $M_0=1$. Since $-B\le S_{n\wedge\tau}\le A$, $M_{n\wedge\tau}$ lies
between $r^{-B}$ and $r^{A}$, and optional stopping (b) with \cref{ch04:lem-finite} gives \ts{24}{6}{34:22}
\[
  1=\E[r^{S_\tau}]=r^{A}\,\Prob(S_\tau=A)+r^{-B}\bigl(1-\Prob(S_\tau=A)\bigr)
  \quad\Longrightarrow\quad \Prob(S_\tau=A)=\frac{1-r^{-B}}{r^A-r^{-B}}=\frac{r^B-1}{r^{A+B}-1}.\qedhere
\]
\end{proof}
As $p\to\tfrac12$, $r=1+\varepsilon$ with $\varepsilon\to0$, and the ratio tends to $B\varepsilon/((A+B)\varepsilon)=B/(A+B)$,
recovering the fair case. The expected duration (not derived in class) follows from the martingale
$S_n-n(p-q)$ (\cref{ch04:ex-m1}) and optional stopping (c) (bounded increments, $\E[\tau]<\infty$):
$\E[S_\tau]=(p-q)\,\E[\tau]$ (Wald's identity), so
\[
  \E[\tau]=\frac{A\,\Prob(S_\tau=A)-B\,\Prob(S_\tau=-B)}{p-q}.
\]

\begin{figure}[ht]
\centering
\begin{tikzpicture}
\begin{axis}[name=phiax,width=0.47\linewidth,height=5.2cm,domain=-2:1,samples=100,ymin=0,ymax=3,xmin=-2,xmax=1,
  xlabel={$\lambda$},ylabel={$\phi(\lambda)$},title={\small MGF of a step, $p=0.8$},
  tick label style={font=\scriptsize},label style={font=\small}]
  \addplot[defcol,thick] {0.8*exp(x)+0.2*exp(-x)};
  \addplot[black!40,dashed] {1};
  \draw[thmcol,thick] (axis cs:-1.3863,0) -- (axis cs:-1.3863,3);
  \node[thmcol,font=\scriptsize,anchor=south west,align=left] at (axis cs:-1.35,1.6) {$\lambda_0=\ln(q/p)$\\$=-1.386$};
\end{axis}
\begin{axis}[at={(phiax.east)},anchor=west,xshift=1.4cm,width=0.47\linewidth,height=5.2cm,
  domain=0:40,samples=41,ymin=0,ymax=1,xmin=0,xmax=40,
  xlabel={my bankroll $B$ (with $A+B=40$)},ylabel={$\Prob(S_\tau=+A)$},title={\small Probability of winning all},
  tick label style={font=\scriptsize},label style={font=\small},
  legend style={font=\scriptsize,at={(0.03,0.97)},anchor=north west,draw=none,fill=none}]
  \addplot[intcol,thick] {((0.47/0.53)^x-1)/((0.47/0.53)^40-1)};
  \addplot[intcol!60,thick,dashed] {((0.49/0.51)^x-1)/((0.49/0.51)^40-1)};
  \addplot[black,thick] {x/40};
  \addplot[thmcol!60,thick,dashed] {((0.51/0.49)^x-1)/((0.51/0.49)^40-1)};
  \addplot[thmcol,thick] {((0.53/0.47)^x-1)/((0.53/0.47)^40-1)};
  \legend{$p=0.53$,$p=0.51$,$p=0.50$,$p=0.49$,$p=0.47$}
\end{axis}
\end{tikzpicture}
\caption{Left: the moment generating function $\phi(\lambda)=pe^\lambda+qe^{-\lambda}$ of a $\pm1$ step equals $1$ at
$\lambda=0$ and at $\lambda_0=\ln(q/p)$, as on slide 28 \ts{24}{6}{33:19}. Right: probability of winning the
opponent's whole bankroll, \cref{ch04:prop-biased}, as a function of my initial capital when the total capital is
$40$. A bias of a few percent per bet turns the straight line of the fair game into an S-shaped curve.}\label{ch04:fig-ruin}
\end{figure}

\begin{casestudy}[Probability surface in R (slides 30--31)]
The slides define \texttt{fcn.probA(A,B)} implementing \cref{ch04:prop-biased} with $p=0.53$, evaluate it on the grid
$A,B\in\{1,\dots,30\}$ with \texttt{outer} (after \texttt{Vectorize}), and draw the surface (\texttt{persp3D}) and its
contour lines \ts{24}{6}{35:30}. Some values: $A=B=10$ gives $0.769$ and $A=B=20$ gives $0.917$. Against an opponent of
unlimited wealth ($A\to\infty$) the probability of never being ruined is $1-r^B$: $0.45$ for $B=5$, $0.70$ for $B=10$ and
$0.97$ for $B=30$. This is why the contour lines on slide~31 become horizontal for large $A$.
\end{casestudy}

\begin{finance}[Edges and ruin (numbers ours)]
The limits $A\to\infty$ make the 2013 remark precise: ``if one has some advantage over an opponent in some game, then no
matter how small that advantage is, he/she will win in the long run''. With $p>q$ the probability of never being ruined
against an infinitely rich opponent is $1-(q/p)^B>0$; with $p<q$ ruin is certain. The casino is the infinitely rich
opponent with the edge: betting \$1 on red at European roulette ($p=18/37$) with \$100, aiming to reach \$200, succeeds
with probability $0.0045$ after about $3700$ bets on average; aiming for \$20 starting from \$10 succeeds with probability
$0.37$. For a trader, $p-q$ is the edge per trade and $B$ the capital in units of the bet size: a small edge protects
against ruin only if the bets are small relative to the capital.
\end{finance}

\begin{coursenote}[2013 versus 2024 treatment]
2013 derives the fair ruin probability by first-step analysis and then, at the end of the lecture, re-derives
$\Prob=1/3$ for $(100,-50)$ as an application of optional stopping, glossing over the fact that the hitting time is
unbounded \ts{13}{5}{1:16:09}. 2024 builds the machinery first (transforms, stopping times, stopped martingales), then
computes the ruin probability, the expected duration and the biased case, all by optional stopping. 2013 has instead
the link between stationary distributions and Perron--Frobenius, which 2024 does not discuss.
\end{coursenote}

\begin{exercise}[\file{lecnote5}, Exercise 5.5]
For positive integers $a,b$ the player stops the first time the balance equals $a$ or $-b$. Find the distribution of
$S_\tau$. (This is \cref{ch04:prop-ruin}; redo it without looking. The notes write ``$X_\tau=b$'' for the second
outcome; $-b$ is meant.)
\end{exercise}

\begin{exercise}[Biased ruin by first-step analysis (not from the course)]
For the biased walk let $f(k)=\Prob(\text{hit }A\text{ before }-B\mid S_0=k)$. Show that
$f(k)=pf(k+1)+qf(k-1)$, solve this linear recursion (characteristic roots $1$ and $q/p$) with $f(A)=1$, $f(-B)=0$, and
recover \cref{ch04:prop-biased}. Then use Wald's identity to show that the expected duration tends to $AB$ as
$p\to\frac12$.
\end{exercise}

% ------------------------------------------------------------------
\section{Looking ahead}\label{ch04:sec-ahead}
\begin{itemize}
  \item \emph{Random walk $\to$ Brownian motion} (\cref{ch:sp2}): independent stationary increments and the $\sqrt t$
        scale survive the limit; $S_n^2-n$ becomes $B_t^2-t$ and $e^{\lambda S_n}/\phi(\lambda)^n$ becomes
        $e^{\lambda B_t-\lambda^2t/2}$.
  \item \emph{Martingale transforms $\to$ It\^o integrals} (\cref{ch:stochcalc,ch:sde}): gains from trading, and the
        reason why hedged portfolios under the risk-neutral measure are martingales (\cref{ch:bs}).
  \item \emph{Markov chains $\to$ Markov diffusions}: Chapman--Kolmogorov becomes the Fokker--Planck equation, and
        first-step analysis becomes the backward equation behind option-pricing PDEs.
  \item \emph{Optional stopping} returns in first-passage problems: barrier options, default times in structural credit
        models, American-style early exercise.
  \item The \emph{Poisson process}, the continuous-time counting process of events with exponential waiting times, is
        not covered in these lectures; 2013 PS2 builds the Poisson distribution from exponential waiting times
        (\cref{ch:probability}).
\end{itemize}

% ------------------------------------------------------------------
\begin{keyformulas}
\begin{align*}
&\text{Random walk:} && \E[S_n]=n\mu,\quad \Var S_n=n\sigma^2,\\
& && (S_n-n\mu)/(\sigma\sqrt n)\to N(0,1)\\
&\text{Markov chain:} && \Prob(X_0=i_0,\dots,X_n=i_n)
  =p_{i_0}P_{i_0i_1}\cdots P_{i_{n-1}i_n},\\
& && P^{(n)}=P^n,\quad \bm p^{(n)}=\bm pP^n\\
&\text{Stationary dist.:} && \bm\pi P=\bm\pi;\quad P^k>0\ \Rightarrow\
  P^{(n)}_{ij}\to\pi_j\ \text{geometrically}\\
&\text{Martingale:} && \E[M_n\mid\mathcal F_{n-1}]=M_{n-1},\\
& && \E[M_n\mid\mathcal F_m]=M_m,\quad \E[M_n]=M_0\\
&\text{Examples:} && S_n,\quad S_n^2-n\sigma^2,\\
& && \textstyle\prod_kX_k\ (\E[X]=1),\quad e^{\lambda S_n}/\phi(\lambda)^n\\
&\text{Transform:} && \tilde M_n=M_0+\textstyle\sum_{k\le n}A_k\,\Delta M_k,\\
& && A_k\in\mathcal F_{k-1}\ \Rightarrow\ \text{martingale}\\
&\text{Opt. stopping:} && \E[M_\tau]=\E[M_0]\ \text{if $\tau$ bounded,}\\
& && \text{or $M_{n\wedge\tau}$ bounded, or $\E[\tau]<\infty$ \& bounded steps}\\
&\text{Fair ruin:} && \Prob(S_\tau=A)=\tfrac{B}{A+B},\quad \E[\tau]=AB\\
&\text{Biased ruin:} && \Prob(S_\tau=A)=\tfrac{(q/p)^B-1}{(q/p)^{A+B}-1},\\
& && \E[\tau]=\tfrac{\E[S_\tau]}{p-q}
\end{align*}
\end{keyformulas}

% !TEX root = ../main.tex
\chapter{Regression Analysis}\label{ch:regression}
\begin{paracol}{2}

\begin{sources}
\begin{tabularx}{\linewidth}{@{}l l L@{}}
\toprule
\textbf{Lecture} & \textbf{Speaker} & \textbf{Content}\\
\midrule
2024 L6 (from 1:07:22) & P.~Kempthorne & Multiple linear regression: set-up, polynomial, Fourier and autoregressive fits of the high-yield spread, the steps of model fitting, assumptions on the errors.\\
2024 L8 & P.~Kempthorne & OLS, hat matrix and geometry; normal linear model via moment-generating functions; $t$ and $F$ tests; prostate-cancer example (standardised covariates, $R^2$, diagnostics); Gauss--Markov; GLS.\\
2024 L11 & P.~Kempthorne & Recap, $t^2=F$, GLS, maximum likelihood, M-estimators (robust, quantile), ridge, principal-components regression, lasso; ETF case study; CAPM regressions (GE, all S\&P~500 stocks) and tests for a change of regime.\\
2013 L6 & P.~Kempthorne & The same theory: OLS, Gauss--Markov with proof, GLS, distribution theory via MGFs and the QR decomposition, maximum likelihood.\\
\bottomrule
\end{tabularx}

\smallskip
\textbf{Files.} Slides \file{mit18_642_f24_lec06_1.pdf} (used in L6, L8 and L11);
\file{mit18_642_f24_lec06_2.pdf} and \file{mit18_642_f24_lec06_3.pdf} (polynomial and Fourier fits
of the high-yield spread); \file{mit18_642_f24_lec06_4.pdf}, the knitted \file{ETF_Casestudy.rmd}
of \file{mit18_642_f24_etf_casestudy.zip} (which also contains \file{ETF_Casestudy_XLB.rmd/.pdf},
\file{spyders_symbols_labels.csv}, \file{casestudy_1_0_etfs.RData}); note \emph{Hypothesis Testing in
Regression Analysis} \file{mit18_642_f24_lec08.pdf}; 2013 notes \file{MIT18_S096F13_lecnote6.pdf};
2013 Case Study~1 \file{CaseStudy1.pdf} with \file{fm_casestudy_1_0.r} and
\file{fm_casestudy_1_rcode.r}; 2013 Case Study~2 \file{CaseStudy2.pdf} with
\file{fm_casestudy_2_rcode.r}, \file{fm_casestudy_fx_1.r} and \file{fred_fxrates.txt}. Installation
scripts and R instructions are described in \cref{app:rsetup}.
\textbf{Problem sets.} 2013 PS3 (all problems). Least squares via the SVD (2013 PS1, B-4) is in
\cref{ch:linalg}.
\textbf{Not published.} The note on CAPM regressions for all S\&P~500 stocks shown in 2024 L11; only
its GE part is in \file{lec08} (see \cref{ch05:capm}).
\end{sources}

\section*{Motivation}
Linear regression is the most used statistical tool in finance: every beta, hedge ratio, factor
exposure, CAPM test or ``is this strategy's alpha real?'' question is a regression. The mathematics
has two layers, and this chapter develops both:
\begin{itemize}
  \item \emph{Geometry} (no probability): least squares is the orthogonal projection of the data
        vector $\bm y\in\R^n$ on the column space of the design matrix $X$ (\cref{ch05:sec-ols}).
        This is linear algebra from \cref{ch:linalg}.
  \item \emph{Statistics}: under assumptions on the errors, the projection has optimality
        properties (Gauss--Markov, \cref{ch05:sec-gm}) and known sampling distributions
        (\cref{ch05:sec-normal}), which give $t$ and $F$ tests (\cref{ch05:sec-tests}), including
        tests for a change of regime (\cref{ch05:sec-chow}).
\end{itemize}
The rest is about what to do when the assumptions fail: diagnostics (\cref{ch05:sec-diag}), robust
and quantile estimators (\cref{ch05:sec-mle}), and shrinkage (\cref{ch05:sec-shrink}). The course's
message, repeated in every lecture, is methodological: propose a model, fit it, \emph{check the
assumptions}, revise \ts{24}{6}{1:18:03}. The case studies (sector ETFs, CAPM for GE and the S\&P~500,
the Chinese yuan) show what this looks like on real data.

% ------------------------------------------------------------------
\section{The linear regression model}\label{ch05:sec-model}

\paragraph{Data and goals.} We observe $n$ cases $i=1,\dots,n$; for each, a \emph{response}
(dependent variable) $y_i$ and $p$ \emph{explanatory} (independent) variables
$\bm x_i=(x_{i,1},\dots,x_{i,p})\T$ \ts{24}{6}{1:07:22}. In finance $y_i$ may be a stock return and
$\bm x_i$ market or macroeconomic factors \ts{13}{6}{2:04}. The goals are prediction, causal inference
(which factors really drive $y$?), approximation of an unknown function, and the discovery or validation
of functional relationships \ts{24}{6}{1:08:24}.

\begin{definition}[Linear model]\label{ch05:def-linmodel}
The \emph{general linear model} specifies the conditional distribution of $y_i$ given $\bm x_i$ as
\[
  y_i = \hat y_i + \varepsilon_i,\qquad \hat y_i=\beta_1x_{i,1}+\beta_2x_{i,2}+\dots+\beta_px_{i,p},
\]
with regression parameters $\bm\beta=(\beta_1,\dots,\beta_p)\T$ constant over the cases and a residual
(error) $\varepsilon_i$ varying over the cases. In matrix form
\[
  \bm y = X\bm\beta+\bm\varepsilon,\qquad
  \bm y=\begin{pmatrix}y_1\\ \vdots\\ y_n\end{pmatrix},\quad
  X=\begin{pmatrix}x_{1,1}&\cdots&x_{1,p}\\ \vdots&\ddots&\vdots\\ x_{n,1}&\cdots&x_{n,p}\end{pmatrix}\in\R^{n\times p}:
\]
cases are rows, variables are columns \ts{13}{6}{29:24}. An intercept is included by taking a first
column of ones.
\end{definition}

\begin{coursenote}[Erratum in the slides]
Slides 3 and 7 of \file{lec06_1} (and the 2013 notes) write the last term as $\beta_{i,p}x_{i,p}$: it must
be $\beta_px_{i,p}$, since the parameters do not depend on the case; Kempthorne corrects this in class
\ts{24}{6}{1:09:30}. In the same slides the bottom-right entry of $X$ is printed $x_{p,n}$ instead of
$x_{n,p}$.
\end{coursenote}

\paragraph{``Linear'' means linear in $\bm\beta$.} The columns of $X$ can be any functions of the
underlying variables \ts{24}{6}{1:10:32} \ts{13}{6}{6:19}:
\begin{itemize}
  \item polynomials, $x_{i,j}=x_i^{\,j}$ (a truncated Taylor expansion);
  \item Fourier terms, $x_{i,j}=\sin(jx_i)$ or $\cos(jx_i)$, for cyclical behaviour;
  \item time-series regressions, where $i$ indexes time and the regressors include lagged values
        $y_{i-1},y_{i-2},\dots$ of the response (autoregressions, \cref{ch:timeseries}) \ts{13}{6}{7:24}.
\end{itemize}
What matters is that $\hat y_i$ is linear in the parameters; this is what makes the estimation a problem
of linear algebra.

\begin{lab}[high-yield spread, functions of time against persistence]{ch05:lab-hyspread}
\textbf{Question.} Can the \emph{high-yield spread} be described by a function of calendar time, and what is
actually predictable in it? (Case study of \file{lec06_2} and \file{lec06_3}.)

\textbf{Data.} \file{data/fm_hyspread.dat}: daily spread 2010--2017, Moody's Baa corporate yield minus the 10-year US
Treasury yield, in percentage points \ts{24}{6}{1:11:33} (FRED series BAA10Y; 2001 days). Lower-rated borrowers pay
more than the Treasury, so it is always positive.

\textbf{In Python} (numpy, least squares with \texttt{np.linalg.lstsq}):
\begin{enumerate}[1.]
  \item Plot the spread. Regress it on powers of time, from the constant to an 8th-order polynomial, and record
        $R^2$ (rescale time to $[-1,1]$ first: why?).
  \item Regress it on Fourier terms $\sin(ju),\cos(ju)$, $j=1,\dots,J$, with $u$ running from 0 to $2\pi$ over the
        sample, for $J=1,\dots,4$.
  \item Regress today's spread on yesterday's, $y_t=a+by_{t-1}+\varepsilon_t$. Compare the three $R^2$, look at the
        fitted curves at the ends of the sample, and say which model you would use to forecast tomorrow.
\end{enumerate}
Reference script: \file{labs/ch05_hyspread.py} (\texttt{--plot} draws the fits).
\end{lab}
\begin{solution}
\textbf{Course results.} Higher orders follow the two big waves (peaks around 2012 and early 2016) better and
better, but the polynomials swing wildly at the ends of the sample, and nobody believes the spread
\emph{is} an 8th-order polynomial in calendar time \ts{24}{6}{1:12:37}. In class a third model was shown
(not in the files): the \emph{autoregression} of today's spread on yesterday's, whose fit is almost
indistinguishable from the data, with $R^2=0.995$ \ts{24}{6}{1:13:44}. (In \file{lec06_3} the legend colours do not
match the curves: the 1-cycle fit is drawn in blue, not green.)

\textbf{With the FRED data} (the course's exact window is not given, so the numbers differ a little):
\begin{center}\small
\begin{tabular}{@{}lcccccc@{}}
\toprule
model & poly 1 & poly 4 & poly 8 & Fourier 2 & Fourier 4 & AR(1)\\
\midrule
$R^2$ & 0.18 & 0.51 & 0.81 & 0.59 & 0.82 & 0.997\\
\bottomrule
\end{tabular}
\end{center}
Rescaling time avoids powers of numbers like 2017, whose columns are almost collinear (an ill-conditioned
$X\T X$). The AR(1) slope is $b=0.9997$: the spread is close to a random walk.

\textbf{Lesson.} A good in-sample fit is not a model. Deterministic functions of time only describe the
past; the autoregression exploits \emph{persistence} (tomorrow's spread is close to today's), which is
what is actually predictable. But with $b\approx1$ its $R^2$ says little: ``tomorrow = today'' already has
$R^2\approx0.997$; the honest comparison is on the daily \emph{changes} (\cref{ch:timeseries}).
\end{solution}

\paragraph{Steps for fitting a model} \ts{24}{6}{1:14:47} \ts{13}{6}{8:24}.
\begin{enumerate}[1.]
  \item Propose a model: the response and its \emph{scale} (e.g.\ log prices, so that errors are
        percentage moves \ts{13}{6}{9:25}), the explanatory variables (including functions of them), and
        assumptions on the distribution of $\bm\varepsilon$.
  \item Define a criterion for judging estimators (least squares, maximum likelihood, robust, Bayes,
        methods for missing data) \ts{13}{6}{23:07}.
  \item Characterise the best estimator under that criterion and compute it.
  \item Check the assumptions of step 1: residual analysis (the errors $\varepsilon_i$ are unobservable, the
        residuals $\hat\varepsilon_i=y_i-\hat y_i$ are their proxies), influence diagnostics, outlier detection
        \ts{13}{6}{26:16}.
  \item If necessary modify the model and/or the assumptions and go back to step 1.
\end{enumerate}
Step 5 is where the work is on hard problems: the model should be tailored to the process, not the
other way round \ts{13}{6}{12:31}.

\begin{definition}[Assumptions on the errors]\label{ch05:def-assumptions}
In increasing generality \ts{24}{6}{1:15:54} \ts{13}{6}{14:38}:
\begin{itemize}
  \item \emph{Normal linear model}: $\varepsilon_1,\dots,\varepsilon_n$ i.i.d.\ $N(0,\sigma^2)$.
  \item \emph{Gauss--Markov}: $\E[\bm\varepsilon]=\bm 0$ and $\Cov(\bm\varepsilon)=\sigma^2I_n$, i.e.\ zero mean,
        constant variance (homoscedastic), uncorrelated. Uncorrelated is weaker than independent
        \ts{13}{6}{15:40}.
  \item \emph{Generalised Gauss--Markov}: $\E[\bm\varepsilon]=\bm 0$, $\Cov(\bm\varepsilon)=\sigma^2\Sigma$ with a general
        positive definite $\Sigma$ (correlated and/or heteroscedastic errors).
  \item \emph{Non-Gaussian} errors: Laplace, Pareto, or \emph{contaminated normal}: a fraction $1-\delta$ of the
        $\varepsilon_i$ i.i.d.\ $N(0,\sigma^2)$ and a fraction $\delta$ from some contaminating distribution.
\end{itemize}
\end{definition}
Correlated residuals arise naturally when a straight line is fitted to a curved relationship
(residuals are systematically above, then below the line) or when unobserved factors persist in time
\ts{13}{6}{20:58}. Generalised Gauss--Markov is the typical situation of time-series regressions, and
heavy-tailed errors that of asset returns \ts{24}{6}{1:16:57}.

\begin{intuition}[A physicist's reading]
This is curve fitting with error bars. Least squares with equal error bars is $\chi^2$ minimisation;
Generalised Gauss--Markov is the case of a full covariance matrix of the measurements; the five steps are
the usual loop ``model, fit, look at the residuals (pulls), improve the model''. The one difference from a
lab: in finance there is no repeated experiment, and the ``error bars'' themselves must be estimated from
the same data.
\end{intuition}

% ------------------------------------------------------------------
\section{Ordinary least squares and its geometry}\label{ch05:sec-ols}

\subsection{Normal equations}
The \emph{least-squares criterion} is \ts{24}{8}{1:06}
\[
  Q(\bm\beta)=\sum_{i=1}^n\bigl(y_i-\bm x_i\T\bm\beta\bigr)^2=(\bm y-X\bm\beta)\T(\bm y-X\bm\beta)=\|\bm y-X\bm\beta\|^2,
\]
and the \emph{ordinary least-squares} (OLS) estimate $\hat{\bm\beta}$ minimises it. Differentiating,
\[
  \frac{\partial Q}{\partial\beta_j}=\sum_{i=1}^n2(-x_{i,j})\bigl(y_i-\bm x_i\T\bm\beta\bigr)=-2X_{[j]}\T(\bm y-X\bm\beta),
  \qquad
  \frac{\partial Q}{\partial\bm\beta}=-2X\T(\bm y-X\bm\beta),
\]
where $X_{[j]}$ is the $j$-th column of $X$ \ts{13}{6}{33:40}. The Hessian is
$\partial^2Q/\partial\bm\beta\,\partial\bm\beta\T=2X\T X$, positive semi-definite since
$\bm v\T X\T X\bm v=\|X\bm v\|^2\ge0$: $Q$ is convex, so every stationary point is a global minimum
\ts{24}{8}{5:16} \ts{13}{6}{32:34}. Setting the gradient to zero gives the \emph{normal equations}
\begin{equation}\label{ch05:eq-normal}
  X\T(\bm y-X\hat{\bm\beta})=\bm 0\iff X\T X\hat{\bm\beta}=X\T\bm y
  \quad\Longrightarrow\quad \hat{\bm\beta}=(X\T X)^{-1}X\T\bm y .
\end{equation}
$X\T X$ is invertible iff $X\bm v=\bm 0\Rightarrow\bm v=\bm 0$, i.e.\ iff $X$ has \emph{full column rank} $p$
\ts{24}{8}{7:33}: no explanatory variable may be a linear combination of the others \ts{13}{6}{37:50}.
Without full rank $\hat{\bm\beta}$ is not unique, but the fitted values are (see below), which is enough
for pure prediction \ts{13}{6}{37:50}.

\begin{remark}
Kempthorne jokes that there is nothing ``normal'' about the normal equations \ts{24}{8}{7:33}. A useful
way to remember the name: \eqref{ch05:eq-normal} says that the residual vector is \emph{normal}
(perpendicular) to every column of $X$.
\end{remark}

\subsection{The hat matrix}
The fitted values are $\hat{\bm y}=X\hat{\bm\beta}=H\bm y$ with the $n\times n$ \emph{hat matrix}
\ts{24}{8}{8:36} \ts{13}{6}{39:58}
\[
  H=X(X\T X)^{-1}X\T ,
\]
and the residuals are $\hat{\bm\varepsilon}=\bm y-\hat{\bm y}=(I_n-H)\bm y$.

\begin{proposition}[Properties of the hat matrix]\label{ch05:prop-hat}
Let $X\in\R^{n\times p}$ have full column rank.
\begin{enumerate}[(i)]
  \item $H\T=H$ and $H^2=H$; the same holds for $I_n-H$.
  \item $H$ is the orthogonal projection onto $\operatorname{col}(X)$, and $I_n-H$ the orthogonal projection
        onto $\operatorname{col}(X)^\perp$.
  \item $\tr H=\rank H=p$ and $\tr(I_n-H)=n-p$.
  \item $X\T\hat{\bm\varepsilon}=\bm 0$, hence $\hat{\bm\varepsilon}\perp\hat{\bm y}$ and
        $\|\bm y\|^2=\|\hat{\bm y}\|^2+\|\hat{\bm\varepsilon}\|^2$.
\end{enumerate}
\end{proposition}
\begin{proof}
(i) $H\T=X[(X\T X)^{-1}]\T X\T=H$ because $X\T X$ is symmetric;
\[H^2=X(X\T X)^{-1}(X\T X)(X\T X)^{-1}X\T=H\]
\ts{24}{8}{10:41}; $(I-H)^2=I-2H+H^2=I-H$ \ts{24}{8}{11:45}.

(ii) If $\bm v=X\bm a\in\operatorname{col}(X)$ then $H\bm v=X(X\T X)^{-1}X\T X\bm a=\bm v$; if $\bm w\perp\operatorname{col}(X)$ then
$X\T\bm w=\bm 0$ and $H\bm w=\bm 0$. Since $\R^n=\operatorname{col}(X)\oplus\operatorname{col}(X)^\perp$, $H$ is the orthogonal
projection, and $I-H$ the complementary one.

(iii) $\tr H=\tr\bigl((X\T X)^{-1}X\T X\bigr)=\tr I_p=p$, using $\tr(AB)=\tr(BA)$; for a projection, trace = rank.

(iv) is \eqref{ch05:eq-normal}; then $\hat{\bm y}\T\hat{\bm\varepsilon}=\hat{\bm\beta}\T X\T\hat{\bm\varepsilon}=0$ and Pythagoras
\ts{24}{8}{14:54}.
\end{proof}

\begin{figure}[ht]
\centering
\begin{tikzpicture}[scale=1.05,>=Stealth]
  \fill[defcol!8] (0,0) -- (6.4,0) -- (8.2,2.3) -- (1.8,2.3) -- cycle;
  \draw[defcol!60] (0,0) -- (6.4,0) -- (8.2,2.3) -- (1.8,2.3) -- cycle;
  \node[defcol,anchor=west] at (6.55,0.35) {$\operatorname{col}(X)$};
  \coordinate (O) at (1.2,0.6);
  \coordinate (Yh) at (5.0,1.4);
  \coordinate (Y) at (5.0,4.3);
  \draw[->,thick,black!60] (O) -- (3.4,0.6) node[below,black]{\small $X_{[1]}$};
  \draw[->,thick,black!60] (O) -- (2.3,1.85) node[left,black]{\small $X_{[2]}$};
  \draw[->,very thick,intcol] (O) -- (Y) node[above]{$\bm y$};
  \draw[->,very thick,defcol] (O) -- (Yh) node[below right]{$\hat{\bm y}=H\bm y$};
  \draw[->,very thick,thmcol] (Yh) -- (Y) node[midway,right]{$\hat{\bm\varepsilon}=(I-H)\bm y$};
  \draw[black!70] (5.0,1.7) -- (4.707,1.638) -- (4.707,1.338);
  \fill (O) circle (1.2pt) node[below left]{\small $\bm 0$};
\end{tikzpicture}
\caption{Least squares in $\R^n$: $\hat{\bm y}$ is the orthogonal projection of $\bm y$ on the column space of
$X$ (spanned here by two columns), and the residual is perpendicular to it
\ts{24}{8}{13:50} \ts{13}{6}{42:08}.}\label{ch05:fig-geometry}
\end{figure}

\begin{intuition}[Two pictures of the same fit]
In the usual scatter plot we look at $n$ points in $\R^{p}$ and a fitted hyperplane. In
\cref{ch05:fig-geometry} we look at \emph{one} point $\bm y$ in $\R^n$ (one axis per case) and a
$p$-dimensional subspace: least squares drops the perpendicular \ts{13}{6}{42:08}. The second picture is
the one in which all the distribution theory below becomes transparent: the residual lives in an
$(n-p)$-dimensional space, which is why ``$n-p$ degrees of freedom'' appear everywhere.
\end{intuition}

\begin{exercise}[PS3 (2013), Problem 1(a)--(c): the hat matrix]\label{ch05:ex-ps3-1ac}
\leavevmode
\begin{enumerate}[(a)]
  \item Prove that $H$ is a projection matrix: $H\T=H$ and $H^2=H$.
  \item The diagonal element $H_{ii}$ is the \emph{leverage} of case $i$. Show that $\partial\hat y_i/\partial y_i=H_{ii}$.
  \item Show that the average of the $H_{ii}$ is $p/n$. \emph{Hint:} $\tr(AB)=\tr(BA)$.
\end{enumerate}
\end{exercise}

\subsection{\texorpdfstring{$R^2$}{R-squared}}
Assume from now on that $X$ contains a column of ones (an intercept). Then $\bm 1\T\hat{\bm\varepsilon}=0$ by
\eqref{ch05:eq-normal}: residuals sum to zero and $\bar{\hat y}=\bar y$. Subtracting $\bar y\bm 1$ (which lies in
$\operatorname{col}(X)$) from both sides of Pythagoras gives the \emph{analysis of variance}
\[
  \underbrace{\sum_i(y_i-\bar y)^2}_{\text{TSS}}=\underbrace{\sum_i(\hat y_i-\bar y)^2}_{\text{ESS}}
  +\underbrace{\sum_i\hat\varepsilon_i^2}_{\text{RSS}} .
\]

\begin{definition}[Coefficient of determination]\label{ch05:def-r2}
$R^2=\text{ESS}/\text{TSS}=1-\text{RSS}/\text{TSS}$ (``multiple $R$-squared''), and the \emph{adjusted}
$R^2_{\text{adj}}=1-\dfrac{\text{RSS}/(n-p)}{\text{TSS}/(n-1)}$, which penalises the number of parameters.
\end{definition}

\begin{proposition}\label{ch05:prop-r2corr}
With an intercept, $R^2=\Corr(\bm y,\hat{\bm y})^2$, the squared sample correlation between observed and fitted
values \ts{24}{8}{59:31}.
\end{proposition}
\begin{proof}
Let $\tilde{\bm y}=\bm y-\bar y\bm 1$ and $\tilde{\hat{\bm y}}=\hat{\bm y}-\bar y\bm 1$. Since $\hat{\bm\varepsilon}\perp\hat{\bm y}$ and
$\hat{\bm\varepsilon}\perp\bm 1$, $\tilde{\bm y}\T\tilde{\hat{\bm y}}=(\tilde{\hat{\bm y}}+\hat{\bm\varepsilon})\T\tilde{\hat{\bm y}}=\|\tilde{\hat{\bm y}}\|^2$, so
$\Corr^2=\|\tilde{\hat{\bm y}}\|^4/(\|\tilde{\bm y}\|^2\|\tilde{\hat{\bm y}}\|^2)=\text{ESS}/\text{TSS}$.
\end{proof}
For simple regression ($p=2$, one slope) this is the familiar squared correlation of $x$ and $y$, and
$\hat\beta_1=\sum_i(x_i-\bar x)(y_i-\bar y)/\sum_i(x_i-\bar x)^2=\widehat{\Cov}(x,y)/\widehat{\Var}(x)$, the beta
of \cref{def:alphabeta}. It is worth calibrating the eye: the daily CAPM regression of GE on the market in
\cref{ch05:capm} has $R^2=0.30$, a clearly visible but very noisy linear cloud \ts{24}{11}{59:45}.

\begin{exercise}[not from the course: $R^2$ with orthogonal regressors]\label{ch05:ex-r2orth}
Let $\bm z_1,\dots,\bm z_m$ be centred and mutually orthogonal, and regress $\bm y$ on $\bm 1,\bm z_1,\dots,\bm z_m$.
\begin{enumerate}[(a)]
  \item Show that the multiple
      $R^2$ equals the sum of the $R^2$'s of the $m$ simple regressions of $\bm y$ on $\bm 1,\bm z_j$, and that each slope is the same in the
      simple and in the multiple regression.
  \item Check this on the ETF numbers ($0.545+0.006+0.079+0.012$ vs $0.642$) and
      explain why the order of the principal components by variance says nothing about their $R^2$.
\end{enumerate}
\end{exercise}

\subsection{The QR decomposition}\label{ch05:sec-qr}
Numerically one never inverts $X\T X$. The 2013 lecture \ts{13}{6}{1:08:26} uses the QR decomposition
\[
  X=QR,\qquad Q\in\R^{n\times p},\ Q\T Q=I_p,\qquad R\in\R^{p\times p}\ \text{upper triangular},
\]
obtained by Gram--Schmidt on the columns of $X$ \ts{13}{6}{1:09:26}. Column by column,
$X_{[1]}=Q_{[1]}r_{11}$, $X_{[2]}=Q_{[1]}r_{12}+Q_{[2]}r_{22}$, \dots, so
\[\begin{gathered}
  r_{11}=\|X_{[1]}\|,\quad Q_{[1]}=\frac{X_{[1]}}{r_{11}};\\
  r_{12}=Q_{[1]}\T X_{[2]},\quad r_{22}=\|X_{[2]}-r_{12}Q_{[1]}\|,\quad Q_{[2]}=\frac{X_{[2]}-r_{12}Q_{[1]}}{r_{22}};\ \dots
  \end{gathered}\]
(each column is orthogonalised against the previous ones) \ts{13}{6}{1:10:32}. Plugging $X=QR$ into the OLS
formulas \ts{13}{6}{1:11:35}:
\[
  \hat{\bm\beta}=R^{-1}Q\T\bm y,\qquad \Cov(\hat{\bm\beta})=\sigma^2(X\T X)^{-1}=\sigma^2R^{-1}(R^{-1})\T,\qquad H=QQ\T .
\]
$\hat{\bm\beta}$ is computed by back-substitution in $R\hat{\bm\beta}=Q\T\bm y$. The advantage (not stressed in
class): the condition number of $X\T X$ is the square of that of $X$, so forming $X\T X$ loses twice as many
digits when the regressors are nearly collinear. R's \texttt{lm()} works through a QR decomposition (the
\texttt{qr} component of a fitted \texttt{lm} object). The SVD route of 2013 PS1 B-4 (\cref{ch:linalg}) is
the most stable of all and is the one used for ridge and PC regression in \cref{ch05:sec-shrink}.

\subsection{Reparametrisation and standardised covariates}\label{ch05:sec-standard}
If $X'=XG$ with $G\in\R^{p\times p}$ invertible, then $\operatorname{col}(X')=\operatorname{col}(X)$, so $H$, $\hat{\bm y}$ and
$\hat{\bm\varepsilon}$ are unchanged and the new parameter is $\bm\beta'=G^{-1}\bm\beta$ (2013 PS3, \cref{ch05:ex-ps3-1de}).
The important special case is \emph{standardisation} \ts{24}{8}{54:05}: replace each covariate by its
$z$-score $(X_{[j]}-\bar x_j\bm 1)/s_j$. Then
\begin{itemize}
  \item fitted values, residuals, $R^2$, $F$ statistic, and the $t$ statistics and $p$-values of the slopes are
        identical \ts{24}{8}{57:24};
  \item the slope coefficient becomes $s_j\beta_j$: the effect on $y$ of a \emph{one-standard-deviation} move
        of $x_j$, which makes coefficients of variables in different units comparable \ts{24}{8}{58:27};
  \item the intercept changes meaning: it becomes $\bar y$, the prediction at the average covariate.
\end{itemize}
In the prostate example of \cref{ch05:sec-diag} the coefficient of \texttt{lcavol} goes from $0.5435$ to
$0.6080$ (its training-sample standard deviation is $1.119$) and that of \texttt{svi} from $0.877$ to
$0.349$, while the $t$ values stay $5.206$ and $2.949$.

% ------------------------------------------------------------------
\section{The Gauss--Markov theorem}\label{ch05:sec-gm}

\begin{definition}[Gauss--Markov assumptions]
The data $(\bm y,X)$ follow a linear model satisfying the Gauss--Markov assumptions if $\bm y$ is an
observation of a random vector $\bm Y$ with \ts{24}{8}{1:11:00} \ts{13}{6}{45:22}
\[
  \E[\bm Y\mid X,\bm\beta]=X\bm\beta,\qquad \Cov(\bm Y\mid X,\bm\beta)=\sigma^2I_n\quad(\sigma^2>0),
\]
equivalently $\bm Y=X\bm\beta+\bm\varepsilon$ with $\E[\bm\varepsilon]=\bm 0$, $\Cov(\bm\varepsilon)=\sigma^2I_n$. Everything is conditional
on $X$, which is treated as fixed.
\end{definition}

For a random vector $\bm Y$ with mean $\bm\mu$ and covariance $\Sigma$ and a constant matrix $A$,
\[
  \E[A\bm Y]=A\bm\mu,\qquad \Cov(A\bm Y)=A\,\Sigma\,A\T
\]
\ts{13}{6}{1:12:37}. With $A=(X\T X)^{-1}X\T$ this gives, under Gauss--Markov,
\begin{equation}\label{ch05:eq-covbeta}
  \E[\hat{\bm\beta}]=(X\T X)^{-1}X\T X\bm\beta=\bm\beta,\qquad
  \Cov(\hat{\bm\beta})=\sigma^2(X\T X)^{-1}X\T X(X\T X)^{-1}=\sigma^2(X\T X)^{-1}.
\end{equation}

\begin{theorem}[Gauss--Markov]\label{ch05:thm-gm}
For known constants $c_1,\dots,c_{p+1}$ let $\theta=c_1\beta_1+\dots+c_p\beta_p+c_{p+1}$ and
$\hat\theta=c_1\hat\beta_1+\dots+c_p\hat\beta_p+c_{p+1}$ with the OLS estimates. Then $\hat\theta$ is
\begin{enumerate}[(i)]
  \item unbiased;
  \item linear in $\bm y$, $\hat\theta=\sum_ib_iy_i$ with constants $b_i$ depending on $X$ only.
\end{enumerate}
Under the Gauss--Markov assumptions it has the smallest variance among \emph{all} linear unbiased estimators
of $\theta$: it is the \emph{best linear unbiased estimator} (BLUE) \ts{24}{8}{1:12:05} \ts{13}{6}{46:25}.
\end{theorem}
\begin{proof}
(2013 notes) \ts{13}{6}{51:42}. Without loss of generality $c_{p+1}=0$; write $\bm c=(c_1,\dots,c_p)\T$. Then
$\hat\theta=\bm c\T\hat{\bm\beta}=\bm c\T(X\T X)^{-1}X\T\bm y\equiv\bm d\T\bm y$ with $\bm d=X(X\T X)^{-1}\bm c\in\operatorname{col}(X)$, and
$\E[\hat\theta]=\bm c\T\bm\beta$ by \eqref{ch05:eq-covbeta}. Let $\tilde\theta=\bm b\T\bm y$ be any linear unbiased estimator and
$\bm f=\bm b-\bm d$, so $\tilde\theta=\hat\theta+\bm f\T\bm y$. Unbiasedness of both gives
$0=\E[\bm f\T\bm y]=\bm f\T X\bm\beta$ for \emph{all} $\bm\beta\in\R^p$, hence $X\T\bm f=\bm 0$: $\bm f\perp\operatorname{col}(X)$, in particular
$\bm f\perp\bm d$. Therefore
\[
  \Var(\tilde\theta)=\Var(\bm d\T\bm y)+\Var(\bm f\T\bm y)+2\Cov(\bm d\T\bm y,\bm f\T\bm y)
  =\Var(\hat\theta)+\sigma^2\|\bm f\|^2+2\sigma^2\bm d\T\bm f\ \ge\ \Var(\hat\theta),
\]
with equality iff $\bm f=\bm 0$ \ts{13}{6}{53:43}.
\end{proof}

\begin{coursenote}[Erratum in the 2013 proof]
Slide 17 of \file{lecnote6} writes $0=\E[\bm f\T\bm y]=\bm d\T\E[\bm y]=\bm f\T(X\bm\beta)$; the middle term should be
$\bm f\T\E[\bm y]$. The key point, easy to miss, is that unbiasedness must hold \emph{for every} $\bm\beta$, which is
what forces $\bm f\perp\operatorname{col}(X)$.
\end{coursenote}

\paragraph{Remarks.}
\begin{itemize}
  \item The theorem is very general \ts{24}{8}{1:13:07}: with $\bm c=\bm x_0$ it says that $\bm x_0\T\hat{\bm\beta}$ is the best
        linear unbiased estimate of the mean response at a new point $\bm x_0$; with $\bm c=\bm x_1-\bm x_2$, of the
        difference of mean responses of two cases. Since it holds for every $\bm c$, it is equivalent to the
        matrix statement $\Cov(\tilde{\bm\beta})-\Cov(\hat{\bm\beta})\succeq0$ for every linear unbiased $\tilde{\bm\beta}$.
  \item No normality is needed. Econometrics texts call this the BLUE property of OLS \ts{13}{6}{54:45}.
  \item ``Best'' only among \emph{linear unbiased} estimators. A biased estimator (ridge,
        \cref{ch05:sec-shrink}) or a non-linear one (robust M-estimators, \cref{ch05:sec-mle}) can have a smaller
        mean squared error, especially with heavy-tailed errors.
\end{itemize}

\paragraph{Design: where to put the $x$'s} \ts{24}{8}{27:37}. $\Var(\hat\beta_j)=\sigma^2C_{jj}$ with
$C_{jj}=[(X\T X)^{-1}]_{jj}$. If we control $X$ (an experiment), we should make $C_{jj}$ small. If $X\T X$ is
diagonal, $C_{jj}=1/\|X_{[j]}\|^2$: the longer the column, the more precise the estimate \ts{24}{8}{29:50}. For a
straight line $y=\beta_0+\beta_1x$,
\[
  \Var(\hat\beta_1)=\frac{\sigma^2}{\sum_i(x_i-\bar x)^2}:
\]
points far from $\bar x$ pin the slope down \ts{24}{8}{30:55}. In general the precision is governed by the
eigenvalues of $X\T X$ (optimal experimental design) \ts{24}{8}{31:59}.

\begin{intuition}[Lever arms]
A physicist measuring a spring constant spreads the loads over the widest possible range: an error
$\delta y$ at a point at distance $d$ from the centroid tilts the line by $\delta y/d$. The same ``lever arm'' will
reappear as the \emph{leverage} $H_{ii}$ of a case (\cref{ch05:sec-diag}): the points that pin the line down are
also the ones that can pull it away.
\end{intuition}

\begin{exercise}[not from the course: optimal design of a line fit]
For $y=\beta_0+\beta_1x+\varepsilon$ under Gauss--Markov with $n$ even and $x_i\in[-1,1]$ at your choice:
\begin{enumerate}[(a)]
  \item derive
      $\hat\beta_1=S_{xy}/S_{xx}$ and $\Var(\hat\beta_1)=\sigma^2/S_{xx}$ with $S_{xx}=\sum_i(x_i-\bar x)^2$;
  \item find the design minimising
      $\Var(\hat\beta_1)$ and compute the leverages $H_{ii}$ for it;
  \item explain why this design cannot detect curvature, and how the
      curvature test of \texttt{residualPlots()} would behave.
\end{enumerate}
\end{exercise}

% ------------------------------------------------------------------
\section{Generalised least squares}\label{ch05:sec-gls}
Suppose now $\bm Y=X\bm\beta+\bm\varepsilon$ with $\E[\bm\varepsilon]=\bm 0$ and $\Cov(\bm\varepsilon)=\sigma^2\Sigma$, where $\sigma^2$ is an unknown scale
and $\Sigma$ is a \emph{known} $n\times n$ positive definite matrix giving the relative variances and the
correlations of the observations \ts{24}{8}{1:13:07} \ts{13}{6}{54:45}. By the spectral theorem
(\cref{thm:spectral}) $\Sigma=U\Lambda U\T$ has the symmetric positive definite inverse square root
$\Sigma^{-1/2}=U\Lambda^{-1/2}U\T$. Transform the data \ts{24}{8}{1:15:23} \ts{13}{6}{55:47}:
\[\begin{gathered}
  \bm Y^*=\Sigma^{-1/2}\bm Y,\quad X^*=\Sigma^{-1/2}X,\quad \bm\varepsilon^*=\Sigma^{-1/2}\bm\varepsilon:\\
  \bm Y^*=X^*\bm\beta+\bm\varepsilon^*,\quad \Cov(\bm\varepsilon^*)=\sigma^2\Sigma^{-1/2}\Sigma\Sigma^{-1/2}=\sigma^2I_n .
\end{gathered}\]
The transformed model satisfies Gauss--Markov with the \emph{same} $\bm\beta$, so by \cref{ch05:thm-gm} the BLUE is
OLS on the starred data:
\begin{equation}\label{ch05:eq-gls}
  \hat{\bm\beta}_{\text{GLS}}=\bigl(X^{*\mathsf T}X^*\bigr)^{-1}X^{*\mathsf T}\bm Y^*=\bigl(X\T\Sigma^{-1}X\bigr)^{-1}X\T\Sigma^{-1}\bm Y,
  \qquad \Cov(\hat{\bm\beta}_{\text{GLS}})=\sigma^2\bigl(X\T\Sigma^{-1}X\bigr)^{-1}.
\end{equation}
Equivalently $\hat{\bm\beta}_{\text{GLS}}$ minimises the Mahalanobis-type criterion $(\bm y-X\bm\beta)\T\Sigma^{-1}(\bm y-X\bm\beta)$. If
$\Sigma=\diag(\sigma_1^2,\dots,\sigma_n^2)$ this is \emph{weighted least squares} with weights $1/\sigma_i^2$: noisy cases are
down-weighted \ts{24}{8}{1:16:24} \ts{13}{6}{57:55}. The main source of non-diagonal $\Sigma$ in finance is time
dependence: errors at nearby dates are more correlated than errors far apart \ts{24}{8}{1:14:15}; e.g.\ AR(1)
errors have $\Sigma_{ij}\propto\rho^{|i-j|}$ (\cref{ch:timeseries}).

\begin{intuition}[Whitening]
$\chi^2=\bm r\T C^{-1}\bm r$ with the covariance matrix $C$ of correlated measurements is exactly the GLS criterion.
Multiplying by $\Sigma^{-1/2}$ ``whitens'' the noise: in the new coordinates the noise is isotropic and the plain
Euclidean projection of \cref{ch05:fig-geometry} is optimal again.
\end{intuition}

\begin{remark}[Not discussed in class]
If the errors are correlated but OLS is used anyway, $\hat{\bm\beta}$ is still unbiased, but it is no longer best and its
true covariance is the ``sandwich'' $\sigma^2(X\T X)^{-1}X\T\Sigma X(X\T X)^{-1}$, not $\sigma^2(X\T X)^{-1}$: the standard
errors printed by \texttt{lm()} are wrong. In practice $\Sigma$ is unknown; it is estimated (feasible GLS) or the
standard errors are corrected (heteroscedasticity- and autocorrelation-consistent errors).
\end{remark}

% ------------------------------------------------------------------
\section{The normal linear model: distribution theory}\label{ch05:sec-normal}

\subsection{Multivariate normal vectors via MGFs}
For an $n$-variate random vector $\bm Y$ the moment-generating function is $M_{\bm Y}(\bm t)=\E[e^{\bm t\T\bm Y}]$,
$\bm t\in\R^n$ \ts{24}{8}{21:10}. In the normal linear model the $Y_i$ are independent $N(\mu_i,\sigma^2)$ with
$\bm\mu=X\bm\beta$ (independent but not identically distributed \ts{13}{6}{57:55}), so the expectation of the product
factorises \ts{24}{8}{22:13} \ts{13}{6}{1:02:09}:
\[
  M_{\bm Y}(\bm t)=\prod_{i=1}^n\E[e^{t_iY_i}]=\prod_{i=1}^ne^{t_i\mu_i+\frac12t_i^2\sigma^2}
  =\exp\Bigl(\bm t\T\bm\mu+\tfrac12\bm t\T\Sigma\bm t\Bigr),\qquad \Sigma=\sigma^2I_n .
\]

\begin{definition}[Multivariate normal]\label{ch05:def-mvn}
$\bm Y\sim N_n(\bm\mu,\Sigma)$, with $\bm\mu\in\R^n$ and $\Sigma$ symmetric positive semi-definite, if
$M_{\bm Y}(\bm t)=\exp(\bm t\T\bm\mu+\frac12\bm t\T\Sigma\bm t)$ for all $\bm t$. Then $\E[\bm Y]=\bm\mu$ and $\Cov\bm Y=\Sigma$; if
$\Sigma\succ0$ the density is $(2\pi)^{-n/2}(\det\Sigma)^{-1/2}\exp\bigl(-\frac12(\bm y-\bm\mu)\T\Sigma^{-1}(\bm y-\bm\mu)\bigr)$. Since an MGF
finite around $\bm 0$ determines the distribution, recognising this form is enough to identify a normal vector.
\end{definition}

\begin{coursenote}[Vocabulary]
In the 2024 L8 video Kempthorne often says ``multinomial'' \ts{24}{8}{20:07}: he means \emph{multinormal},
i.e.\ multivariate normal. Nothing in this chapter involves the multinomial distribution.
\end{coursenote}

\begin{proposition}[Linear maps of normal vectors]\label{ch05:prop-linmap}
If $\bm Y\sim N_n(\bm\mu,\Sigma)$ and $A\in\R^{m\times n}$ then $\bm Z=A\bm Y\sim N_m(A\bm\mu,A\Sigma A\T)$ \ts{24}{8}{31:59}
\ts{13}{6}{1:14:49}. In particular every sub-vector, and every component, is normal; and two jointly normal
sub-vectors with zero cross-covariance are independent.
\end{proposition}
\begin{proof}
$M_{\bm Z}(\bm\tau)=\E[e^{\bm\tau\T A\bm Y}]=M_{\bm Y}(A\T\bm\tau)=\exp\bigl(\bm\tau\T A\bm\mu+\frac12\bm\tau\T A\Sigma A\T\bm\tau\bigr)$. For a component take
$A=\bm e_j\T$ \ts{13}{6}{1:07:25}. If $\bm Z=(\bm Z_1,\bm Z_2)$ is normal with $\Cov(\bm Z_1,\bm Z_2)=0$, the quadratic form in the
exponent splits and the joint MGF is the product $M_{\bm Z_1}(\bm\tau_1)M_{\bm Z_2}(\bm\tau_2)$, which characterises
independence \ts{24}{8}{37:20}.
\end{proof}
Applying this with $A=(X\T X)^{-1}X\T$ \ts{24}{8}{24:21}: $\bm t=A\T\bm\tau=X(X\T X)^{-1}\bm\tau$ gives $\bm t\T\bm\mu=\bm\tau\T\bm\beta$ and
$\bm t\T(\sigma^2I)\bm t=\bm\tau\T\sigma^2(X\T X)^{-1}\bm\tau$, so $\hat{\bm\beta}\sim N_p(\bm\beta,\sigma^2(X\T X)^{-1})$ and
$\hat\beta_j\sim N(\beta_j,\sigma^2C_{jj})$ \ts{24}{8}{26:31} \ts{13}{6}{1:04:19}.

\begin{coursenote}[Errata in the MGF slides]
\begin{itemize}
  \item Slide 17 of \file{lec06_1}: the MGF of $\hat\beta_1$ is obtained with $\bm\tau=t(1,0,\dots,0)\T$, not $t(1,0,\dots,1)\T$.
  \item Slides 14--16 write the exponent as $\bm t\T\bm u+\dots$: $\bm u$ is the mean vector $\bm\mu$.
  \item Slide 18 (``Theorem*'') assumes $\rank A=m\le n$. The result holds for every $A$ if degenerate normals
        (singular $\Sigma$, no density) are allowed, as in \cref{ch05:def-mvn}; this is needed right away, because the
        stacked matrix $\bigl(\begin{smallmatrix}(X\T X)^{-1}X\T\\ I-H\end{smallmatrix}\bigr)$ used for independence has $n+p$
        rows but rank $n$, and $\Cov(\hat{\bm\varepsilon})$ is singular.
\end{itemize}
\end{coursenote}

\subsection{The fundamental theorem}
\begin{definition}[$\chi^2$, $t$ and $F$ distributions]
If $Z_1,\dots,Z_\nu$ are i.i.d.\ $N(0,1)$, $\sum_kZ_k^2\sim\chi^2_\nu$ \ts{13}{6}{1:16:57}. If $Z\sim N(0,1)$ and
$V\sim\chi^2_\nu$ are independent, $Z/\sqrt{V/\nu}\sim t_\nu$ (Student) \ts{24}{8}{38:21}. If $V_1\sim\chi^2_{\nu_1}$ and
$V_2\sim\chi^2_{\nu_2}$ are independent, $(V_1/\nu_1)/(V_2/\nu_2)\sim F_{\nu_1,\nu_2}$ \ts{24}{8}{45:38}. Note $t_\nu^2\sim F_{1,\nu}$.
\end{definition}

\begin{theorem}[Normal linear regression]\label{ch05:thm-normal}
Let $\bm y=X\bm\beta+\bm\varepsilon$ with $X\in\R^{n\times p}$ of rank $p$ and $\bm\varepsilon\sim N_n(\bm 0,\sigma^2I_n)$. Then
\ts{24}{8}{33:00} \ts{13}{6}{1:15:54}:
\begin{enumerate}[(i)]
  \item $\hat{\bm\beta}=(X\T X)^{-1}X\T\bm y\sim N_p\bigl(\bm\beta,\sigma^2(X\T X)^{-1}\bigr)$;
  \item $\hat{\bm\varepsilon}=(I_n-H)\bm y\sim N_n\bigl(\bm 0,\sigma^2(I_n-H)\bigr)$, a degenerate normal: $X\T\hat{\bm\varepsilon}=\bm 0$ identically
        \ts{24}{8}{34:06};
  \item $\hat{\bm\beta}$ and $\hat{\bm\varepsilon}$ are independent;
  \item $\text{RSS}=\hat{\bm\varepsilon}\T\hat{\bm\varepsilon}\sim\sigma^2\chi^2_{n-p}$; hence $\hat\sigma^2=\text{RSS}/(n-p)$ is unbiased;
  \item for each $j$,
        \begin{equation}\label{ch05:eq-tstat}
          \hat t_j=\frac{\hat\beta_j-\beta_j}{\hat\sigma\sqrt{C_{jj}}}\sim t_{n-p},\qquad C_{jj}=\bigl[(X\T X)^{-1}\bigr]_{jj}.
        \end{equation}
\end{enumerate}
\end{theorem}
\begin{proof}
(i) was shown above.

(ii) $\Cov\bigl((I-H)\bm y\bigr)=\sigma^2(I-H)(I-H)\T=\sigma^2(I-H)$ and $(I-H)X\bm\beta=\bm 0$.

(iii) $\hat{\bm\beta}$ and $\hat{\bm\varepsilon}$ are jointly normal (both linear in $\bm y$) with
$\Cov(\hat{\bm\beta},\hat{\bm\varepsilon})=(X\T X)^{-1}X\T(\sigma^2I)(I-H)=\sigma^2(X\T X)^{-1}(X\T-X\T H)=0$ because $X\T H=X\T$; apply
\cref{ch05:prop-linmap} (this is the ``joint MGF = product of MGFs'' argument of the 2024 slides).

(iv) Following the 2013 notes \ts{13}{6}{1:16:57}: let $X=QR$ and complete the columns of $Q$ to an orthonormal
basis of $\R^n$ by $W\in\R^{n\times(n-p)}$ (continue Gram--Schmidt on $[X\,|\,I_n]$). Then
$A=\bigl(\begin{smallmatrix}Q\T\\ W\T\end{smallmatrix}\bigr)$ is orthogonal and
\[
  \bm z=A\bm y=\begin{pmatrix}\bm z_Q\\ \bm z_W\end{pmatrix}\sim N_n\!\left(\begin{pmatrix}Q\T QR\bm\beta\\ W\T QR\bm\beta\end{pmatrix},\sigma^2AA\T\right)
  =N_n\!\left(\begin{pmatrix}R\bm\beta\\ \bm 0_{n-p}\end{pmatrix},\sigma^2I_n\right).
\]
So $\bm z_Q\sim N_p(R\bm\beta,\sigma^2I_p)$ and $\bm z_W\sim N_{n-p}(\bm 0,\sigma^2I_{n-p})$ are independent. Now
$\hat{\bm\beta}=R^{-1}Q\T\bm y=R^{-1}\bm z_Q$, and since $QQ\T+WW\T=A\T A=I_n$,
$\hat{\bm\varepsilon}=(I-QQ\T)\bm y=WW\T\bm y=W\bm z_W$. Hence $\text{RSS}=\|W\bm z_W\|^2=\|\bm z_W\|^2$ is $\sigma^2$ times a sum of $n-p$
independent squared standard normals. (This also re-proves (iii): $\hat{\bm\beta}$ and $\hat{\bm\varepsilon}$ are functions of the
independent blocks $\bm z_Q$ and $\bm z_W$.)

(v) $(\hat\beta_j-\beta_j)/(\sigma\sqrt{C_{jj}})\sim N(0,1)$ by (i), independent of $\text{RSS}/\sigma^2\sim\chi^2_{n-p}$ by (iii)--(iv);
dividing, $\sigma$ cancels and the ratio is \eqref{ch05:eq-tstat}.
\end{proof}

The unbiasedness of $\hat\sigma^2$ does not need normality \ts{24}{8}{36:15}: under Gauss--Markov,
$\E[\hat{\bm\varepsilon}\T\hat{\bm\varepsilon}]=\tr\Cov(\hat{\bm\varepsilon})=\sigma^2\tr(I_n-H)=\sigma^2(n-p)$ by \cref{ch05:prop-hat}(iii). The
$n-p$ appears because the residuals satisfy $p$ exact linear constraints $X\T\hat{\bm\varepsilon}=\bm 0$: there are not $n$
independent error terms \ts{24}{8}{35:14}.

\begin{coursenote}[Erratum: the missing square root]
Slide 20 of \file{lec06_1} and slide 34 of \file{lecnote6} write the denominator of $\hat t_j$ as
$\hat\sigma\,C_{j,j}$. It must be $\hat\sigma\sqrt{C_{jj}}$, the estimated standard deviation of $\hat\beta_j$; Kempthorne
corrects it at the start of L11 \ts{24}{11}{3:16}. Also, slide 38 of \file{lecnote6} writes
$\hat{\bm\beta}=(X\T X)^{-1}X\bm y$ for $(X\T X)^{-1}X\T\bm y$.
\end{coursenote}

\begin{intuition}[Why Student's $t$ has fat tails]
In $\hat t_j$ the unknown $\sigma$ is replaced by the random $\hat\sigma$. When $\hat\sigma$ happens to be small, $|\hat t_j|$ is
large; the denominator fluctuates around $1$ (in units of $\sigma$) and this mixing over scales fattens the tails
\ts{24}{8}{40:28}. W.~S.~Gosset, a statistician at the Guinness brewery, noticed that standardised means of
quality samples of size four were far more variable than $N(0,1)$ predicted; the brewery did not let him
publish under his name, hence ``Student'' \ts{24}{8}{41:29}. For $n-p\gtrsim30$ the
difference from the normal is small; with $n-p=63$ (prostate example) the two-sided 5\% critical value is
$2.00$ instead of $1.96$.
\end{intuition}

\begin{intuition}[Degrees of freedom, geometrically]
$\text{RSS}/\sigma^2$ is the squared length of a standard Gaussian vector projected on an $(n-p)$-dimensional subspace:
fitting $p$ parameters ``uses up'' $p$ dimensions. The familiar $1/(n-1)$ in the sample variance is the case
$p=1$ (projection on $\bm 1$), and $\chi^2/\mathrm{ndf}$ in a fit is $\hat\sigma^2/\sigma^2$.
\end{intuition}

A $(1-\alpha)$ confidence interval for $\beta_j$ follows immediately:
$\hat\beta_j\pm t_{n-p,1-\alpha/2}\,\hat\sigma\sqrt{C_{jj}}$, where $\hat\sigma\sqrt{C_{jj}}$ is the \emph{standard error}
$\mathrm{se}(\hat\beta_j)$ printed by R.

% ------------------------------------------------------------------
\section{Maximum likelihood and M-estimation}\label{ch05:sec-mle}

\subsection{Maximum likelihood in the normal model}
The \emph{likelihood} is the joint density of the data seen as a function of the unknown parameters,
$L(\bm\beta,\sigma^2)=p(\bm y\mid X,\bm\beta,\sigma^2)$; maximum likelihood (ML) estimates make the observed data ``most
likely'' \ts{24}{11}{10:53} \ts{13}{6}{1:20:05}. For the normal linear model
\[
  L(\bm\beta,\sigma^2)=\prod_{i=1}^n\frac{1}{\sqrt{2\pi\sigma^2}}e^{-\frac{(y_i-\bm x_i\T\bm\beta)^2}{2\sigma^2}},\qquad
  \log L=-\frac n2\log(2\pi\sigma^2)-\frac{1}{2\sigma^2}Q(\bm\beta)
\]
\ts{24}{11}{13:02}. Maximise in two steps: $\bm\beta$ enters only through $-Q(\bm\beta)$, so $\hat{\bm\beta}_{\text{ML}}=\hat{\bm\beta}_{\text{OLS}}$
for every $\sigma^2$; then $\partial\log L/\partial\sigma^2=-\frac{n}{2\sigma^2}+\frac{Q(\hat{\bm\beta})}{2\sigma^4}=0$ gives
\[
  \hat\sigma^2_{\text{ML}}=\frac{\text{RSS}}{n},\qquad \E[\hat\sigma^2]_{\text{ML}}=\frac{n-p}{n}\sigma^2\quad\text{(biased)}
\]
\ts{24}{11}{14:08}. So least squares is also ``doing something decent'' from the likelihood point of view when the
errors are Gaussian \ts{13}{6}{1:21:11}. ML estimators are asymptotically optimal (smallest variance as
$n\to\infty$) for any correctly specified model \ts{24}{11}{11:58}: if the error density were known and non-Gaussian,
ML would use it.

\subsection{Generalised M-estimators}
\begin{definition}[Generalised M-estimator]\label{ch05:def-mest}
For data $(\bm y,X)$ and the model $y_i=\bm x_i\T\bm\beta+\varepsilon_i$, choose $\hat{\bm\beta}$ to minimise
$Q(\bm\beta)=\sum_{i=1}^nh(y_i,\bm x_i,\bm\beta,\sigma^2)$ \ts{24}{11}{15:14}. The function $h$ defines the estimator:
\begin{enumerate}[(i)]
  \item least squares: $h=(y_i-\bm x_i\T\bm\beta)^2$;
  \item least absolute deviations (``MAD''): $h=|y_i-\bm x_i\T\bm\beta|$;
  \item maximum likelihood: $h=-\log p(y_i\mid\bm\beta,\bm x_i,\sigma^2)$ (so Laplace errors give LAD);
  \item robust M-estimator: $h=\chi(y_i-\bm x_i\T\bm\beta)$ with $\chi$ even and increasing on $(0,\infty)$;
  \item quantile estimator, for fixed $\tau\in(0,1)$: $h=\rho_\tau(y_i-\bm x_i\T\bm\beta)$ with the \emph{check loss}
        \[
          \rho_\tau(u)=\begin{cases}\tau\,|u|,&u\ge0,\\(1-\tau)\,|u|,&u<0.\end{cases}
        \]
\end{enumerate}
\end{definition}

\begin{figure}[ht]
\centering
\begin{tikzpicture}
\begin{axis}[width=0.62\linewidth,height=5.2cm,xmin=-2.5,xmax=2.5,ymin=0,ymax=2.6,
  axis lines=middle,xlabel={$u$},ylabel={$h(u)$},samples=101,
  legend style={font=\scriptsize,at={(1.02,1)},anchor=north west,draw=none},
  every axis x label/.style={at={(current axis.right of origin)},anchor=west},
  every axis y label/.style={at={(current axis.above origin)},anchor=south}]
  \addplot[thick,black!60,dashed,domain=-2.3:2.3]{x^2/2};
  \addlegendentry{$u^2/2$ (least squares)}
  \addplot[thick,defcol,domain=-2.4:2.4]{abs(x)};
  \addlegendentry{$|u|$ (LAD)}
  \addplot[thick,intcol,domain=-2.4:2.4]{(abs(x)<=1)*x^2/2+(abs(x)>1)*(abs(x)-0.5)};
  \addlegendentry{Huber, $c=1$}
  \addplot[very thick,thmcol,domain=-2.4:2.4]{(x>=0)*0.9*x+(x<0)*(-0.1)*x};
  \addlegendentry{$\rho_{0.9}$ (quantile)}
\end{axis}
\end{tikzpicture}
\caption{Loss functions of M-estimators. Quadratic loss lets large residuals dominate; LAD and Huber grow
linearly; the check loss $\rho_{0.9}$ penalises positive residuals nine times more than negative ones, so the fit
is pushed up to the 90th percentile.}\label{ch05:fig-losses}
\end{figure}

Robust estimators were designed for the contaminated-normal situation: normal errors most of the time, with a
small fraction of gross errors \ts{24}{11}{17:30}. When the error distribution is unknown they protect the fit
from a few extreme cases, which in finance means crash days. A standard example (not detailed in class) is
Huber's loss, quadratic for $|u|\le c$ and linear beyond.

\begin{proposition}[The check loss estimates quantiles]\label{ch05:prop-quantile}
Let $y_1,\dots,y_n\in\R$ and $f(c)=\sum_i\rho_\tau(y_i-c)$. Then $f$ is convex and piecewise linear, and $c$ minimises $f$ iff
\[
  \#\{i:y_i<c\}\le\tau n\le\#\{i:y_i\le c\},
\]
i.e.\ iff $c$ is a sample $\tau$-quantile. For $\tau=\frac12$ this is the median, and $\rho_{1/2}=\frac12|u|$, so
$\tau=\frac12$ gives LAD \ts{24}{11}{18:37}.
\end{proposition}
\begin{proof}
Each term is convex in $c$. The right derivative of $\rho_\tau(y_i-c)$ is $-\tau$ if $y_i>c$ and $1-\tau$ if $y_i\le c$,
so $f'_+(c)=-\tau\,\#\{y_i>c\}+(1-\tau)\,\#\{y_i\le c\}=\#\{y_i\le c\}-\tau n$; similarly $f'_-(c)=\#\{y_i<c\}-\tau n$. A convex
function is minimised exactly where $f'_-(c)\le0\le f'_+(c)$.
\end{proof}
With $c$ replaced by $\bm x_i\T\bm\beta$, the minimiser estimates the conditional $\tau$-quantile of $y$ given $\bm x$
(quantile regression): for $\tau=0.9$, the upper decile of outcomes \ts{24}{11}{19:41}. The objective is convex
(``a bowl with straight edges'') but not differentiable, so there are no normal equations to solve; one moves
downhill using one-sided derivatives, or solves a linear programme \ts{24}{11}{20:41}.

\begin{finance}[Heavy tails in practice]
In the GE CAPM regression of \cref{ch05:capm} the histogram of residuals is compared with two normal fits
\ts{24}{11}{1:00:53}: the ML fit (sample mean and standard deviation) and a robust one, whose scale is
estimated from the interquartile range ($\hat\sigma_{\text{rob}}=\mathrm{IQR}/1.349$, since the IQR of $N(0,\sigma^2)$ is
$1.349\sigma$) \ts{24}{11}{1:03:05}. The ML curve is too wide in the centre and too thin in the tails. A sharper
check uses \emph{fitted percentiles} $u_i=\Phi\bigl((\hat\varepsilon_i-\hat\mu)/\hat\sigma\bigr)$, which should be uniform on $(0,1)$ if the
model is right (probability integral transform) \ts{24}{11}{1:04:09}: with the ML scale they pile up near
$0.5$, with the robust scale they are uniform except in the extreme 1\% on each side
\ts{24}{11}{1:05:11}. Quantile regression is the natural tool for Value-at-Risk (\cref{ch:var}).
\end{finance}

\begin{exercise}[not from the course: quantile regression as maximum likelihood]
Show that minimising $\sum_i\rho_\tau(y_i-\bm x_i\T\bm\beta)$ is maximum likelihood for errors with the \emph{asymmetric Laplace} density
$f(u)=\frac{\tau(1-\tau)}{s}\exp\bigl(-\rho_\tau(u)/s\bigr)$, and check that $f$ is a density whose $\tau$-quantile is $0$. Which estimator of
\cref{ch05:def-mest} is ML for Laplace errors?
\end{exercise}

% ------------------------------------------------------------------
\section{Hypothesis testing}\label{ch05:sec-tests}

\subsection{Individual coefficients}
To test $H_0:\beta_j=b$ against $H_1:\beta_j\ne b$ use $\hat t=(\hat\beta_j-b)/\mathrm{se}(\hat\beta_j)$, which is $t_{n-p}$
under $H_0$ by \eqref{ch05:eq-tstat}; reject when $|\hat t|>c$ for a critical value $c$ \ts{24}{11}{5:26}. The
\emph{$p$-value} is the probability, computed under $H_0$, of a statistic at least as extreme as the one observed,
$P(|T_{n-p}|\ge|\hat t|)$ \ts{24}{8}{49:56}. The coefficient table of R's \texttt{summary(lm(\dots))} has columns
Estimate ($\hat\beta_j$), Std.\ Error ($\hat\sigma\sqrt{C_{jj}}$), $t$ value (test of $b=0$: the ``signal-to-noise
ratio'') and \texttt{Pr(>|t|)} \ts{24}{8}{48:48}. Conventionally $p<0.05$ is ``significant'' and $p<0.01$
``highly significant''. The case $b=0$ asks whether variable $j$ can be dropped from the model
\ts{24}{11}{5:26}.

\begin{example}[Is GE's beta equal to one? (\file{lec08}, section 1)]\label{ch05:ex-gebeta}
Daily excess returns of GE on the market excess return, $n=1676$ days (2017--2023): $\hat\alpha=-0.00054$ (se
$0.00051$, $p=0.285$), $\hat\beta=1.0830$ (se $0.0400$), $R^2=0.304$, $\hat\sigma=0.0207$. Testing $\beta=0$ is pointless
(a stock is obviously correlated with the market, $t=27.1$); the interesting null is $\beta=1$, ``same systematic
risk as the market'' \ts{24}{11}{59:45}:
\[
  \hat t=\frac{1.0830-1}{0.0400}=2.075,\qquad p=P(|T_{1674}|\ge2.075)=0.038,
\]
so GE's beta is (just) significantly above one at the 5\% level \ts{24}{11}{1:07:25}.
\end{example}

\begin{coursenote}[Errata in \file{lec08}, section 1]
On page 3 both hypotheses are labelled $H_0$; the alternative is $H_1:\beta_j\neq0$. On page 4,
``$(\hat\beta_2-\beta_2)\sim N(\beta_2,sd(\hat\beta_2))$'' should read $\hat\beta_2-\beta_2\sim N\bigl(0,sd(\hat\beta_2)^2\bigr)$, i.e.\
$\hat\beta_2\sim N(\beta_2,sd^2)$.
\end{coursenote}

\subsection{Nested models and the \texorpdfstring{$F$}{F} test}
\begin{lemma}[Quadratic forms of normal vectors]\label{ch05:lem-proj}
Let $\bm Y\sim N_n(\bm\mu,\sigma^2I_n)$ and let $P_1,P_2$ be orthogonal projections of ranks $r_1,r_2$ with $P_1P_2=0$ and
$P_1\bm\mu=P_2\bm\mu=\bm 0$. Then $\bm Y\T P_k\bm Y/\sigma^2\sim\chi^2_{r_k}$ ($k=1,2$), and the two quadratic forms are independent.
\end{lemma}
\begin{proof}
Write $P_k=U_kU_k\T$ with $U_k\in\R^{n\times r_k}$ having orthonormal columns (spectral theorem: the eigenvalues of a
projection are $0$ and $1$). Then $U_k\T\bm Y\sim N_{r_k}(U_k\T\bm\mu,\sigma^2I)$ with $U_k\T\bm\mu=U_k\T P_k\bm\mu=\bm 0$, and
$\bm Y\T P_k\bm Y=\|U_k\T\bm Y\|^2$. Moreover $U_1\T U_2=U_1\T(P_1P_2)U_2=0$, so $\Cov(U_1\T\bm Y,U_2\T\bm Y)=\sigma^2U_1\T U_2=0$ and
the two Gaussian vectors are independent (\cref{ch05:prop-linmap}).
\end{proof}

\begin{theorem}[$F$ test of a linear hypothesis]\label{ch05:thm-F}
In the normal linear model let $A\in\R^{q\times p}$ have rank $q$ and consider $H_0:A\bm\beta=\bm 0$. Let RSS be the
residual sum of squares of the unconstrained fit and $\text{RSS}_0$ that of the least-squares fit under the
constraint. Under $H_0$
\[
  F=\frac{(\text{RSS}_0-\text{RSS})/q}{\text{RSS}/(n-p)}\sim F_{q,\,n-p}.
\]
The most common case is $A=[0\ I_q]$: $H_0:\beta_{k+1}=\dots=\beta_p=0$ with $k=p-q$, i.e.\ the \emph{sub-model} using
only the first $k$ columns of $X$ is adequate \ts{24}{8}{43:35} \ts{24}{11}{6:33}.
\end{theorem}
\begin{proof}
The constrained fitted values range over $L_0=\{X\bm\beta:A\bm\beta=\bm 0\}$, a subspace of $\operatorname{col}(X)$ of dimension $p-q$
($X$ is injective). The constrained least-squares fit is the projection $H_0\bm y$ on $L_0$, so
$\text{RSS}_0=\bm y\T(I-H_0)\bm y$ and $\text{RSS}_0-\text{RSS}=\bm y\T(H-H_0)\bm y$. Since $L_0\subset\operatorname{col}(X)$, $HH_0=H_0H=H_0$; hence
$H-H_0$ is the orthogonal projection on $\operatorname{col}(X)\ominus L_0$ (rank $q$) and $(I-H)(H-H_0)=0$. Under $H_0$,
$\bm\mu=X\bm\beta\in L_0$, so $(H-H_0)\bm\mu=(I-H)\bm\mu=\bm 0$. By \cref{ch05:lem-proj}, $(\text{RSS}_0-\text{RSS})/\sigma^2\sim\chi^2_q$ and
$\text{RSS}/\sigma^2\sim\chi^2_{n-p}$ independently, and the ratio of the two over their degrees of freedom is $F_{q,n-p}$.
\end{proof}
Both numerator and denominator estimate $\sigma^2$ when $H_0$ is true \ts{24}{11}{6:33}; under the alternative the
numerator is inflated by $\|(H-H_0)\bm\mu\|^2/q$, so large $F$ is evidence against $H_0$ \ts{24}{8}{44:36}. An affine
hypothesis $A\bm\beta=\bm b$ reduces to this one by subtracting $X\bm\beta_b$ (any $\bm\beta_b$ with $A\bm\beta_b=\bm b$) from $\bm y$.

\begin{coursenote}[Erratum: denominator degrees of freedom]
Slide 20 of \file{lec06_1} gives $df_2=n-k$ for the sub-model $F$ test. It is $df_2=n-p$, the residual degrees of
freedom of the \emph{full} model, as in \file{lec08} (section 2) and in the formula on the same slide.
\end{coursenote}

\begin{proposition}[Wald form; $t^2=F$]\label{ch05:prop-wald}
For $H_0:A\bm\beta=\bm b$,
\[
  \text{RSS}_0-\text{RSS}=(A\hat{\bm\beta}-\bm b)\T\bigl[A(X\T X)^{-1}A\T\bigr]^{-1}(A\hat{\bm\beta}-\bm b).
\]
In particular, for a single coefficient ($A=\bm e_j\T$, $q=1$) $\text{RSS}_0-\text{RSS}=(\hat\beta_j-b)^2/C_{jj}$ and the $F$ statistic
equals $\hat t^{\,2}$ \ts{24}{11}{7:39}.
\end{proposition}
\begin{proof}
Minimising $\|\bm y-X\bm\beta\|^2$ subject to $A\bm\beta=\bm b$ with a Lagrange multiplier gives
$\hat{\bm\beta}_0=\hat{\bm\beta}-(X\T X)^{-1}A\T M^{-1}(A\hat{\bm\beta}-\bm b)$, $M=A(X\T X)^{-1}A\T$. As $X(\hat{\bm\beta}-\hat{\bm\beta}_0)\in\operatorname{col}(X)\perp\hat{\bm\varepsilon}$,
Pythagoras gives $\text{RSS}_0=\text{RSS}+(\hat{\bm\beta}-\hat{\bm\beta}_0)\T X\T X(\hat{\bm\beta}-\hat{\bm\beta}_0)$, and substituting,
$(\hat{\bm\beta}-\hat{\bm\beta}_0)\T X\T X(\hat{\bm\beta}-\hat{\bm\beta}_0)=(A\hat{\bm\beta}-\bm b)\T M^{-1}(A\hat{\bm\beta}-\bm b)$. For $q=1$, dividing by $\hat\sigma^2$ gives
$(\hat\beta_j-b)^2/(\hat\sigma^2C_{jj})=\hat t^{\,2}$.
\end{proof}
R's \texttt{car::linearHypothesis()} implements tests of linear hypotheses in this $F$ form \ts{24}{11}{1:08:29}.
For GE (\cref{ch05:ex-gebeta}) it gives $F=4.304=2.075^2$ ($p=0.038$) for $\beta=1$ and $F=1.144=1.070^2$ ($p=0.285$) for
$\alpha=0$: the same tests as the $t$ statistics.

\paragraph{Overall $F$ and ANOVA.} With an intercept, the test that all $p-1$ slopes vanish has
$\text{RSS}_0=\text{TSS}$, so
\[
  F=\frac{R^2/(p-1)}{(1-R^2)/(n-p)} .
\]
This is the last line of \texttt{summary.lm}: $F=13.4$ on $(8,63)$ degrees of freedom for the prostate data
($R^2=0.630$) and $F=271.9$ on $(4,607)$ for the ETF regression of \cref{ch05:sec-cases} ($R^2=0.642$). Adding
blocks of variables one at a time and recording the drop in RSS gives the sequential \emph{analysis-of-variance
table}, derived through the QR decomposition in 2013 PS3 (\cref{ch05:ex-ps3-3cd}).

\begin{finance}[Statistical versus practical significance; multiple testing]
With thousands of daily observations, tiny effects become ``significant'': in the ETF case study a principal
component explaining $0.6\%$ of the variance of the response has $p=0.052$ \ts{24}{11}{45:48}. Conversely, when the
same test is run on hundreds of stocks (is $\alpha=0$ for each S\&P~500 stock?), about 5\% of them will be
significant at the 5\% level even if every null is true \ts{24}{11}{1:16:47}. A handful of significant alphas is
therefore no evidence against the CAPM, and ``the best of 500 back-tests'' is not a strategy.
\end{finance}

\begin{exercise}[not from the course: $t^2=F$ on real numbers]
\leavevmode
\begin{enumerate}[(a)]
  \item From \cref{ch05:prop-wald}, show directly that for $H_0:\beta_j=b$ the $F$ statistic equals $\hat t^{\,2}$.
  \item Using only the
      GE output of \cref{ch05:ex-gebeta} (estimate, standard error, $n-p=1674$, RSS $=0.71819$), reproduce the RSS $0.72003$ of the
      model constrained by $\beta=1$ printed by \texttt{linearHypothesis}.
\end{enumerate}
\end{exercise}

% ------------------------------------------------------------------
\section{Testing for a change of regime}\label{ch05:sec-chow}
Regressions on time series are fitted over long samples, but the relationship may change: a firm changes business,
a currency changes its peg. \file{lec08} (section 3) shows how to test whether a regression is the same in two
periods \ts{24}{11}{1:09:33}.

\paragraph{Set-up.} Order the $n$ cases in time; period A is $i=1,\dots,n_A$, period B the remaining
$n_B=n-n_A$. Partition $\bm y=(\bm y_A,\bm y_B)$, $X=\bigl(\begin{smallmatrix}X_A\\X_B\end{smallmatrix}\bigr)$ and allow different parameters:
\[
  \bm y_A=X_A\bm\beta+\bm\varepsilon_A,\qquad \bm y_B=X_B\bm\gamma+\bm\varepsilon_B,\qquad \bm\varepsilon_A,\bm\varepsilon_B\ \text{independent},\ N(\bm 0,\sigma^2I),
\]
with $H_0:\bm\gamma=\bm\beta$ (no change) against $H_1:\bm\gamma\neq\bm\beta$ \ts{24}{11}{1:10:37}.

\paragraph{Approach 1: separate regressions.} In block form the unconstrained model is
\[
  \begin{pmatrix}\bm y_A\\ \bm y_B\end{pmatrix}=\begin{pmatrix}X_A&0\\0&X_B\end{pmatrix}\begin{pmatrix}\bm\beta\\ \bm\gamma\end{pmatrix}+\bm\varepsilon .
\]
The normal equations decouple, so $\hat{\bm\beta}=(X_A\T X_A)^{-1}X_A\T\bm y_A$, $\hat{\bm\gamma}=(X_B\T X_B)^{-1}X_B\T\bm y_B$ (separate OLS fits)
and $\text{RSS}=\text{RSS}_A+\text{RSS}_B$, with $n-2p$ residual degrees of freedom. The constrained model is the pooled regression
$\bm y=X\bm\beta+\bm\varepsilon$ with residual sum $\text{RSS}_0$. By \cref{ch05:thm-F} ($q=p$ constraints $\bm\gamma-\bm\beta=\bm 0$, $2p$ parameters)
\begin{equation}\label{ch05:eq-chow}
  F=\frac{(\text{RSS}_0-\text{RSS}_A-\text{RSS}_B)/p}{(\text{RSS}_A+\text{RSS}_B)/(n-2p)}\sim F_{p,\,n-2p}\quad\text{under }H_0 .
\end{equation}
In econometrics this is the \emph{Chow test} (the name is not used in the course).

\paragraph{Approach 2: parametrise the change} \ts{24}{11}{1:11:39}. Let $\bm\delta=\bm\gamma-\bm\beta$. Then
\[
  \begin{pmatrix}\bm y_A\\ \bm y_B\end{pmatrix}=\begin{pmatrix}X_A&0\\X_B&X_B\end{pmatrix}\begin{pmatrix}\bm\beta\\ \bm\delta\end{pmatrix}+\bm\varepsilon,\qquad H_0:\bm\delta=\bm 0 .
\]
The new design is the old one times $G=\bigl(\begin{smallmatrix}I&0\\-I&I\end{smallmatrix}\bigr)$ acting on the parameters, so it has the
same column space: same fit, same RSS, same $F$ (\cref{ch05:sec-standard}). The advantage is that the individual
$t$ statistics of $\hat{\bm\delta}$ tell \emph{which} coefficients changed. In R this is the model with interactions
between the regressors and a period indicator, \texttt{y $\sim$ x * ind.periodB}.

\begin{coursenote}[Erratum in \file{lec08}, section 3.1]
The $F$ statistic of Approach~1 is printed with denominator $\text{RSS}/(n-p)$. The unconstrained model has $2p$ parameters,
so it must be $\text{RSS}/(n-2p)$ as in \eqref{ch05:eq-chow}; the text right below the formula (and section 3.2) correctly
states $df_2=n-2p$.
\end{coursenote}

\begin{casestudy}[Did GE's CAPM change? (\file{lec08}, section 4)]
\textbf{Choosing the split.} Fit the CAPM regression of \cref{ch05:ex-gebeta} to 2017--2023 and plot the cumulative
sum of the residuals standardised to equal variance, $\hat\varepsilon_i/\sqrt{1-H_{ii}}$ (\cref{ch05:sec-diag}). Under the model
this should wander like a driftless random walk; instead it falls steadily to a minimum at 2018-12-27 and then
rises \ts{24}{11}{1:12:44}. Period A: up to 2018-12-27 (500 days); period B: after.

\textbf{Fit with interactions} (\texttt{Ra.RF $\sim$ Mkt.RF * ind.periodB}, $n=1676$):

\smallskip
\centering\small
\begin{tabular}{@{}lrrr@{}}
\toprule
 & Estimate & Std.\ error & $t$\\
\midrule
intercept ($\alpha_A$) & $-0.00288$ & $0.00092$ & $-3.13$\\
\texttt{Mkt.RF} ($\beta_A$) & $0.774$ & $0.110$ & $7.02$\\
\texttt{ind.periodB} ($\delta_\alpha$) & $0.00334$ & $0.00110$ & $3.04$\\
\texttt{Mkt.RF:ind.periodB} ($\delta_\beta$) & $0.354$ & $0.118$ & $2.99$\\
\bottomrule
\end{tabular}

\smallskip
\raggedright
So $\hat\alpha_B=0.00046$ and $\hat\beta_B=1.128$: the alpha turns from negative to (slightly) positive and the beta rises
from $0.77$ to $1.13$ \ts{24}{11}{1:13:44}. The $F$ test of $\delta_\alpha=\delta_\beta=0$ (\texttt{anova(lmfit0, lmfit0a)} or
\texttt{linearHypothesis}): RSS drops from $0.71819$ to $0.71020$,
$F=\frac{0.0079821/2}{0.71020/1672}=9.40$ on $(2,1672)$ degrees of freedom, $p=8.8\times10^{-5}$.

\textbf{Lesson.} CAPM parameters are not constant over long periods; practitioners estimate them on
\emph{rolling windows}, which raises the question of the window length \ts{24}{11}{1:13:44}.
\end{casestudy}

\begin{coursenote}[Choosing the break from the data (not discussed in class)]
Here the break date was chosen \emph{after} looking at the cumulative residuals, i.e.\ at the most favourable point
for rejection. The $F_{2,1672}$ distribution assumes the date is fixed in advance, so the reported $p$-value is too
optimistic. With an unknown break date one uses the supremum of the $F$ statistics over candidate dates (Andrews'
sup-$F$ test) or CUSUM tests, which formalise exactly the cumulative-residual plot used above. The
conclusion for GE would very likely survive, given $p\approx10^{-4}$.
\end{coursenote}

\paragraph{Open issues} (\file{lec08}, section 5). The test assumes i.i.d.\ normal errors with the same variance in
both periods. If the variances differ, fit the two periods separately by maximum likelihood and use a
generalised likelihood-ratio test instead of $F$. If the errors are dependent (check with the Durbin--Watson,
Ljung--Box, turning-point, difference-sign or McLeod--Li tests), time-series models are needed
(\cref{ch:timeseries}). If they are non-normal, large-sample theory still makes $\hat{\bm\beta}$ approximately normal and
$F$ tests are replaced by $\chi^2$ tests.

\begin{casestudy}[Exchange-rate regimes of the Chinese yuan (2013 Case Study 2)]
\textbf{Question.} The yuan (CNY) was pegged to the US dollar until July 2005; China then announced a managed float
against an undisclosed \emph{basket} of currencies, with daily moves up to 0.3\%. What is the implicit basket, does it
change, and how fast does the yuan move against it?

\textbf{Data and method.} \file{fm_casestudy_fx_1.r} downloads from FRED the daily dollar exchange rates listed in
\file{fred_fxrates.txt} (23 currencies plus trade-weighted dollar indices, 1999--2013).
\file{fm_casestudy_2_rcode.r} re-expresses eight of them in Swiss francs (a stable numeraire outside the basket), so
that the dollar itself becomes a regressor, and regresses daily log changes of CNY/SFR on those of USD, JPY, EUR,
GBP, KRW, MYR and THB against the franc (the Frankel--Wei regression).

\textbf{Results.}
\begin{itemize}
  \item 2001 -- June 2005: coefficient on USD $1.000$ (se $0.0005$, $t=1839$), all others insignificant,
        $R^2=0.9999$: a perfect dollar peg.
  \item July--December 2005: USD $0.19$ (not significant), Korean won $0.18$ ($t=4.9$), Malaysian ringgit $0.75$
        ($t=5.1$), $R^2=0.949$: the peg is gone.
  \item 2006--2013, year by year: USD weight $0.87$--$0.97$ ($t$ from 25 to 97), with small, changing, sometimes
        significant weights on MYR, KRW, EUR, GBP or THB; $R^2$ between $0.93$ and $0.998$.
  \item The intercept is a daily drift against the implied basket; annualised as $(1+\hat\alpha)^{252}-1$ it is negative in
        almost every year ($-2.9\%$ in 2006, $-6.1\%$ in 2007, $-7.3\%$ in 2008, $+0.2\%$ in 2009, then $-0.5\%$ to $-4.8\%$):
        fewer yuan per franc than the basket implies, i.e.\ a steady managed \emph{appreciation} of the yuan.
  \item Diagnostics for 2012: no heteroscedasticity, but residual tails heavier than normal.
\end{itemize}
\textbf{Lesson.} A regression coefficient can identify an exchange-rate regime, and fitting it period by period reveals
regime changes; the structural break of July 2005 is visible at a glance. \file{fm_casestudy_fx_1.r} also loads the R
package \texttt{fxregime}, which implements this regression together with formal tests for changes of regime.
\end{casestudy}

\begin{exercise}[not from the course: the Chow test]
\leavevmode
\begin{enumerate}[(a)]
  \item Show that the residual sum of squares of the block-diagonal design $\bigl(\begin{smallmatrix}X_A&0\\0&X_B\end{smallmatrix}\bigr)$ is
      $\text{RSS}_A+\text{RSS}_B$.
  \item Show that $\bigl(\begin{smallmatrix}X_A&0\\X_B&X_B\end{smallmatrix}\bigr)$ has the same column space, and relate the two
      parameter vectors.
  \item For GE, compute $\hat\alpha_B$, $\hat\beta_B$ and the $F$ statistic from the numbers in the case study of
      \cref{ch05:sec-chow}. Why can you not obtain the standard error of $\hat\beta_B$ from that table alone?
\end{enumerate}
\end{exercise}

% ------------------------------------------------------------------
\section{Regression diagnostics}\label{ch05:sec-diag}
Step 4 of the modelling loop. By \cref{ch05:thm-normal}(ii) the residuals are correlated and have unequal variances,
$\Var(\hat\varepsilon_i)=\sigma^2(1-H_{ii})$, even when the errors are i.i.d.\ \ts{24}{8}{1:01:33}. Diagnostics correct for this.

\paragraph{Leverage.} $H_{ii}$ is the \emph{leverage} of case $i$. Since $\hat y_i=\sum_jH_{ij}y_j$, $\partial\hat y_i/\partial y_i=H_{ii}$: the
fraction of its own fitted value that a case dictates (2013 PS3). From $H=H^2=H\T$,
$H_{ii}=\sum_jH_{ij}^2\ge H_{ii}^2$, so $0\le H_{ii}\le1$; with an intercept $H_{ii}\ge1/n$, and $\sum_iH_{ii}=p$, so the average is
$p/n$ \ts{24}{8}{1:05:44}. If $H_{ii}=1$ the fitted line passes through the point: that case alone determines its
fit \ts{24}{8}{1:04:44}. With an intercept,
\[
  H_{ii}=\frac1n+(\bm x_i-\bar{\bm x})\T(\mathcal X\T\mathcal X)^{-1}(\bm x_i-\bar{\bm x}),
\]
where $\mathcal X$ is the matrix of centred covariates (2013 PS3, \cref{ch05:ex-ps3-1de}): leverage grows with the
Mahalanobis distance of $\bm x_i$ from the centroid of the covariates (dividing by $n-1$ turns $\mathcal X\T\mathcal X$ into the sample
covariance matrix). High leverage is a property of the $x$'s only. R's \texttt{influence.measures()} flags cases with
$H_{ii}>3p/n$.

\paragraph{Residuals.} The \emph{standardised} residual (R: \texttt{rstandard}) is
$r_i=\hat\varepsilon_i/(\hat\sigma\sqrt{1-H_{ii}})$; the \emph{studentised} (deleted) residual (R: \texttt{rstudent}, or
\texttt{studres} in \texttt{MASS}) uses the estimate $\hat\sigma_{(i)}$ computed without case $i$,
$t_i=\hat\varepsilon_i/(\hat\sigma_{(i)}\sqrt{1-H_{ii}})$, and has exactly a $t_{n-p-1}$ distribution under the normal model, which makes
outliers easy to judge \ts{24}{8}{1:02:36}.

\begin{proposition}[Case deletion]\label{ch05:prop-deletion}
Let $\hat{\bm\beta}_{(i)}$, $\hat{\bm y}_{(i)}=X\hat{\bm\beta}_{(i)}$ and $\hat\sigma^2_{(i)}$ be computed with case $i$ deleted. Then
\[
  \hat{\bm\beta}_{(i)}=\hat{\bm\beta}-\frac{(X\T X)^{-1}\bm x_i\hat\varepsilon_i}{1-H_{ii}},\qquad
  \hat\sigma^2_{(i)}=\hat\sigma^2+\frac{1}{n-p-1}\Bigl(\hat\sigma^2-\frac{\hat\varepsilon_i^2}{1-H_{ii}}\Bigr),
\]
and \emph{Cook's distance} is
\[
  CD_i=\frac{\|\hat{\bm y}-\hat{\bm y}_{(i)}\|^2}{p\,\hat\sigma^2}=\frac{\hat\varepsilon_i^2}{p\,\hat\sigma^2}\,\frac{H_{ii}}{(1-H_{ii})^2}
  =\frac{r_i^2}{p}\,\frac{H_{ii}}{1-H_{ii}} .
\]
\end{proposition}
The proofs, via the Sherman--Morrison--Woodbury formula (deleting a row is a rank-one update of $X\T X$), are 2013 PS3
Problem~2 (\cref{ch05:ex-ps3-2a,ch05:ex-ps3-2bd}). Nothing needs to be refitted: all leave-one-out quantities come
from one fit. The last form of $CD_i$ says \emph{influence = outlyingness $\times$ leverage}: a large residual at a
low-leverage point, or a high-leverage point on the line, is harmless; both together move the fit. Since
$\|\hat{\bm y}-\hat{\bm y}_{(i)}\|^2=(\hat{\bm\beta}-\hat{\bm\beta}_{(i)})\T X\T X(\hat{\bm\beta}-\hat{\bm\beta}_{(i)})$, $CD_i$ measures how far $\hat{\bm\beta}$ moves in the metric of
its own confidence ellipsoid: values comparable with the median of $F_{p,n-p}$ (about $1$) are large
\ts{24}{8}{1:06:44}.

\paragraph{Graphical tools} \ts{24}{8}{1:08:51}. In R (\texttt{stats} and
\texttt{car}):
\begin{center}
\small
\begin{tabularx}{\linewidth}{@{}l L@{}}
\toprule
\texttt{plot(fit)} & residuals vs fitted (flat band), normal Q--Q,
scale--location ($\sqrt{|r_i|}$ vs fitted), residuals vs leverage
with Cook contours\\
\texttt{influence.measures()} & leave-one-out stats: DFBETAS,
DFFITS, COVRATIO, Cook's distance, hat values\\
\texttt{qqPlot()} & Q--Q plot of studentised residuals against
$t$ quantiles, with confidence band \ts{24}{8}{1:03:44}\\
\texttt{influencePlot()} & studentised residuals vs hat values,
bubble size $\propto$ Cook's distance\\
\texttt{residualPlots()} & residuals vs each regressor; curvature
test = $t$ test of added squared term; Tukey test = same with
$\hat y^2$ \ts{24}{8}{1:09:53}\\
\texttt{leveragePlots()} & added-variable plots: effect of each
regressor after adjusting for others\\
\texttt{durbinWatsonTest()} & autocorrelation of residuals
(time-ordered data)\\
\bottomrule
\end{tabularx}
\end{center}

\begin{casestudy}[Prostate cancer data (\file{lec06_1}, slides 21--33)]
\textbf{Data.} Not financial, deliberately: the \texttt{prostate} data of the \texttt{faraway} package (used in Hastie,
Tibshirani and Friedman's \emph{Elements of Statistical Learning}), 97 patients, response \texttt{lpsa} (log PSA), eight
covariates (log cancer volume \texttt{lcavol}, log prostate weight, age, \dots, seminal vesicle invasion \texttt{svi}).
A random 75\% training sample ($n=72$, \texttt{set.seed(1)}) is used. The experience transfers: only the variable
names change from one regression problem to another \ts{24}{8}{46:42}.

\textbf{Fit} ($p=9$ with intercept): $R^2=0.630$, $R^2_{\text{adj}}=0.583$, $\hat\sigma=0.724$ on 63 degrees of freedom,
$F=13.4$. Significant: \texttt{lcavol} ($t=5.21$) and \texttt{svi} ($t=2.95$); \texttt{lweight} borderline ($p=0.07$). The pairs
plot already shows the strong \texttt{lpsa}--\texttt{lcavol} relation \ts{24}{8}{48:48}. Standardising the covariates
leaves every $t$ value unchanged (\cref{ch05:sec-standard}) \ts{24}{8}{57:24}.

\textbf{Diagnostics.} Studentised residuals: unimodal, symmetric, within the Q--Q band except case 39 ($-3.06$). Hat
values: cases 32 ($0.41$) and 74 ($0.33$) exceed $3p/n=0.375$ or come close, against an average $p/n=0.125$. Cook's
distances of the flagged cases are only $0.07$--$0.13$. Curvature tests: none significant (smallest $p=0.10$ for
\texttt{svi}; Tukey test $p=0.69$) \ts{24}{8}{1:09:53}.

\textbf{Lesson.} A satisfactory model: no curvature, no dominant case; the scale--location plot hints at mild
non-constancy of the variance.
\end{casestudy}

\begin{exercise}[PS3 (2013), Problem 1(d)--(e): invariance and leverage]\label{ch05:ex-ps3-1de}
(d) Prove that $H$ is unchanged if $X$ is replaced by $X'=XG$ for any invertible $p\times p$ matrix $G$.
(e) Let $X$ be $n\times(p+1)$ with a first column of ones and $p$ covariates, and let $G$ be the identity except for the first
row, which is $(1,-\bar x_1,\dots,-\bar x_p)$, with $\bar x_j=\frac1n\sum_ix_{i,j}$. By (d) the models with $X$ and $X'=XG$ have the
same fitted values and residuals.
\begin{itemize}
  \item If $\bm\beta=(\beta_0,\dots,\beta_p)\T$ is the parameter for $X$, show that $\bm\beta'=G^{-1}\bm\beta$ is the parameter for $X'$; find
        $G^{-1}$ and the elements of $\bm\beta'$ explicitly.
  \item Show that $X'^{\mathsf T}X'=\bigl(\begin{smallmatrix}n&\bm 0_p\T\\ \bm 0_p&\mathcal X\T\mathcal X\end{smallmatrix}\bigr)$, where $\mathcal X$ has entries $x_{i,j}-\bar x_j$.
  \item Prove $H_{ij}=\frac1n+(\bm x_i-\bar{\bm x})\T[\mathcal X\T\mathcal X]^{-1}(\bm x_j-\bar{\bm x})$, where $\bm x_i=(x_{i,1},\dots,x_{i,p})\T$. Hence the
        leverage $H_{ii}$ increases with the (squared) Mahalanobis distance between $\bm x_i$ and $\bar{\bm x}$.
\end{itemize}
\end{exercise}

% ------------------------------------------------------------------
\section{Shrinkage and principal-components regression}\label{ch05:sec-shrink}
When regressors are many or strongly correlated (the four index ETFs below have pairwise correlations
$0.77$--$0.95$), $X\T X$ is ill-conditioned and $\Cov(\hat{\bm\beta})=\sigma^2(X\T X)^{-1}$ is huge in some directions. Three
remedies, all expressed through the SVD, were presented in L11.

\paragraph{Standardisation first} \ts{24}{11}{23:47}. Centre $\bm y$ and the columns of $X$ and scale the columns to unit
variance; the intercept is then estimated by $\bar y$. Otherwise a penalty on $\|\bm\beta\|$ would depend on the units:
measuring a covariate in cents instead of dollars would make its coefficient 100 times smaller and almost free
\ts{24}{11}{24:49}. Below $X$ is the standardised $n\times p$ matrix (no intercept column) with thin SVD
$X=UDV\T$, $U\T U=I_p$, $V$ orthogonal, $D=\diag(d_1\ge\dots\ge d_p)$ (\cref{thm:svd}), and $c_j=\bm u_j\T\bm y$.

\begin{coursenote}[Erratum and notation in the ridge slides]
Slide 41 of \file{lec06_1} rescales the columns by ``$S^{-1/2}X$'' with $S=\frac1n\diag(X\T X)$: since $S$ is $p\times p$ this must be
$XS^{-1/2}$ (multiplication on the right scales columns). On the same slide $H^*=I_n-\frac1n\bm 1\bm 1\T$ is the
\emph{centring} matrix, not the hat matrix.
\end{coursenote}

\subsection{Ridge regression}
\begin{definition}[Ridge]
For a complexity parameter $\lambda\ge0$ \ts{24}{11}{22:46},
\[
  \hat{\bm\beta}^{\text{ridge}}=\arg\min_{\bm\beta}\Bigl(\|\bm y-X\bm\beta\|^2+\lambda\|\bm\beta\|^2\Bigr)=(X\T X+\lambda I_p)^{-1}X\T\bm y .
\]
\end{definition}
The gradient condition is $-2X\T(\bm y-X\bm\beta)+2\lambda\bm\beta=\bm 0$; $X\T X+\lambda I$ is invertible for every $\lambda>0$ even when $X$ is
rank-deficient. The estimate is linear in $\bm y$ and equals OLS for $\lambda=0$. With the SVD
\ts{24}{11}{28:04}:
\[
  \hat{\bm\beta}^{\text{ridge}}=V(D^2+\lambda I)^{-1}DU\T\bm y,\qquad
  \hat{\bm y}^{\text{ridge}}=UD(D^2+\lambda I)^{-1}DU\T\bm y=\sum_{j=1}^p\frac{d_j^2}{d_j^2+\lambda}\,c_j\,\bm u_j ,
\]
whereas $\hat{\bm y}^{\text{LS}}=UU\T\bm y=\sum_jc_j\bm u_j$. Ridge shrinks the component of the fit along $\bm u_j$ by the factor
$d_j^2/(d_j^2+\lambda)$: little along the directions in which the regressors vary a lot (large $d_j$, the leading principal
components), much along the directions of little variation \ts{24}{11}{29:04}. It is biased,
$\E\hat{\bm\beta}^{\text{ridge}}=(X\T X+\lambda I)^{-1}X\T X\bm\beta$, but has smaller variance; $\lambda$ is chosen by cross-validation
(\texttt{cv.glmnet} in the case study).

\begin{proposition}[Ridge as a Bayes estimate]\label{ch05:prop-ridgebayes}
In the normal linear model with prior $\bm\beta\sim N_p(\bm 0,\tau^2I_p)$, the posterior mode (which, the posterior being Gaussian,
is also its mean) is $\hat{\bm\beta}^{\text{ridge}}$ with $\lambda=\sigma^2/\tau^2$ \ts{24}{11}{25:52}.
\end{proposition}
\begin{proof}
The posterior density is proportional to likelihood $\times$ prior,
\[\exp\bigl(-\|\bm y-X\bm\beta\|^2/(2\sigma^2)\bigr)\exp\bigl(-\|\bm\beta\|^2/(2\tau^2)\bigr);\]
maximising it means minimising $\|\bm y-X\bm\beta\|^2+(\sigma^2/\tau^2)\|\bm\beta\|^2$.
\end{proof}
The isotropic prior says that a priori no direction of $\R^p$ is preferred \ts{24}{11}{26:58}; a strong prior (small
$\tau$, large $\lambda$) pulls all coefficients towards zero.

\begin{exercise}[not from the course: ridge as augmented least squares]
\leavevmode
\begin{enumerate}[(a)]
  \item Show that $\hat{\bm\beta}^{\text{ridge}}$ is the OLS estimate for the augmented data $\tilde{\bm y}=(\bm y,\bm 0_p)$,
      $\tilde X=\bigl(\begin{smallmatrix}X\\ \sqrt\lambda I_p\end{smallmatrix}\bigr)$.
  \item Using the SVD, show that $\|\hat{\bm\beta}^{\text{ridge}}\|$ decreases in $\lambda$ and compute
      $\Cov(\hat{\bm\beta}^{\text{ridge}})$ under Gauss--Markov.
  \item Show that ridge has smaller variance than OLS in every direction:
      $\Cov(\hat{\bm\beta})-\Cov(\hat{\bm\beta}^{\text{ridge}})\succeq0$.
\end{enumerate}
\end{exercise}

\subsection{Principal-components regression}
The principal components of the (centred) regressors are $\bm z_j=X\bm v_j$, where $\bm v_j$ are the eigenvectors of the
sample covariance $S=\frac1nX\T X=\frac1nVD^2V\T$ \ts{24}{11}{30:05}. They are orthogonal, $\bm z_j=d_j\bm u_j$, with sample
variance $d_j^2/n$, and $\bm z_1$ is the unit-norm combination of the columns of $X$ with the largest variance
(\cref{ch:pca}). \emph{PC regression} fixes $m\le p$ and regresses $\bm y$ on $Z_m=[\bm z_1\cdots\bm z_m]=U_mD_m$
\ts{24}{11}{31:05}:
\[
  \hat{\bm\gamma}=(Z_m\T Z_m)^{-1}Z_m\T\bm y,\quad \hat\gamma_j=\frac{\bm z_j\T\bm y}{d_j^2},\qquad
  \hat{\bm\beta}^{\text{pc}}=V_m\hat{\bm\gamma},\qquad \hat{\bm y}^{\text{pc}}=U_mU_m\T\bm y=\sum_{j=1}^mc_j\bm u_j .
\]
Orthogonality makes each $\hat\gamma_j$ a separate simple regression, and
\[\Cov(\hat{\bm\beta}^{\text{pc}})=\sigma^2V_mD_m^{-2}V_m\T\preceq\sigma^2VD^{-2}V\T=\Cov(\hat{\bm\beta}):\]
variance is reduced by dropping the noisiest
directions \ts{24}{11}{32:13}.

\paragraph{Comparison} \ts{24}{11}{33:19}. The three fits use the same building blocks $c_j\bm u_j$:
\[
  \hat{\bm y}^{\text{LS}}=\sum_{j=1}^pc_j\bm u_j,\qquad
  \hat{\bm y}^{\text{ridge}}=\sum_{j=1}^p\frac{d_j^2}{d_j^2+\lambda}c_j\bm u_j,\qquad
  \hat{\bm y}^{\text{pc}}=\sum_{j=1}^mc_j\bm u_j .
\]
Least squares keeps all, ridge shrinks smoothly, PC regression keeps the first $m$ and discards the rest.

\begin{intuition}[Spectral filters]
This is the classic regularisation of an ill-posed inverse problem: ridge is Tikhonov regularisation (a smooth
filter $d^2/(d^2+\lambda)$ on the singular values), PC regression is the truncated SVD (a sharp cut-off). Both trade a
little bias for a large reduction of the noise amplified by small singular values.
\end{intuition}

A caveat, visible in the ETF case study \ts{24}{11}{46:50}: the principal components are ranked by how much the
\emph{regressors} vary, not by how much they explain $\bm y$. A low-variance component can be an important predictor,
so keeping ``the first $m$'' may throw away the useful direction.

\subsection{The lasso}
\begin{definition}[Lasso]
\[
  \hat{\bm\beta}^{\text{lasso}}=\arg\min_{\bm\beta}\Bigl(\|\bm y-X\bm\beta\|^2+\lambda\sum_{j=1}^p|\beta_j|\Bigr),
\]
equivalently: minimise $\|\bm y-X\bm\beta\|^2$ subject to $\|\bm\beta\|_1\le t$ \ts{24}{11}{33:19}. The constraint binds only for
$t<t_0=\|\hat{\bm\beta}^{\text{LS}}\|_1$; the \emph{regularisation path} is $\hat{\bm\beta}^{\text{lasso}}$ as a function of $\lambda\ge0$ (or of
$t\in[0,t_0]$).
\end{definition}
The $L_1$ ball is a diamond (a cross-polytope in $\R^p$) with corners on the axes. The elliptical contours of
$\|\bm y-X\bm\beta\|^2$, centred at $\hat{\bm\beta}^{\text{LS}}$, typically first touch it at a corner, where some coefficients are exactly
zero \ts{24}{11}{35:28} (\cref{ch05:fig-ridgelasso}). The lasso therefore does \emph{variable selection}: as $\lambda$ grows,
variables leave the model one after the other \ts{24}{11}{36:31}. There is no closed form; the path is computed by
coordinate descent (\texttt{glmnet}).

\begin{figure}[ht]
\centering
\begin{tikzpicture}[scale=0.95]
  \begin{scope}
    \draw[->,black!60] (-1.6,0) -- (2.8,0) node[right,black]{\small $\beta_1$};
    \draw[->,black!60] (0,-1.5) -- (0,4.3) node[above,black]{\small $\beta_2$};
    \fill[intcol!25] (0,0) circle (1);
    \draw[intcol!80!black,thick] (0,0) circle (1);
    \foreach \s in {2.093,1.35,0.6}
      \draw[thmcol,rotate around={-50:(0.6,2.2)}] (0.6,2.2) ellipse ({\s} and {0.5*\s});
    \fill[thmcol] (0.6,2.2) circle (1.5pt) node[right]{\small $\hat{\bm\beta}^{\text{LS}}$};
    \fill[black] (0.562,0.827) circle (1.8pt);
    \node[anchor=north west,font=\small] at (0.62,0.78) {$\hat{\bm\beta}^{\text{ridge}}$};
    \node[font=\small] at (0.6,-1.25) {ridge: $\beta_1^2+\beta_2^2\le t^2$};
  \end{scope}
  \begin{scope}[xshift=6.3cm]
    \draw[->,black!60] (-1.6,0) -- (2.8,0) node[right,black]{\small $\beta_1$};
    \draw[->,black!60] (0,-1.5) -- (0,4.3) node[above,black]{\small $\beta_2$};
    \fill[defcol!25] (1,0) -- (0,1) -- (-1,0) -- (0,-1) -- cycle;
    \draw[defcol,thick] (1,0) -- (0,1) -- (-1,0) -- (0,-1) -- cycle;
    \foreach \s in {2.519,1.6,0.7}
      \draw[thmcol,rotate around={-50:(0.6,2.2)}] (0.6,2.2) ellipse ({\s} and {0.5*\s});
    \fill[thmcol] (0.6,2.2) circle (1.5pt) node[right]{\small $\hat{\bm\beta}^{\text{LS}}$};
    \fill[black] (0,1) circle (1.8pt);
    \node[anchor=east,font=\small] at (-0.08,1.18) {$\hat{\bm\beta}^{\text{lasso}}$};
    \node[font=\small] at (0.6,-1.25) {lasso: $|\beta_1|+|\beta_2|\le t$};
  \end{scope}
\end{tikzpicture}
\caption{Constrained forms of ridge and lasso (same least-squares contours, $t=1$). The solution is where the smallest
contour touches the constraint region: a generic point of the circle for ridge, the corner $\beta_1=0$ of the diamond for
the lasso \ts{24}{11}{34:22}. (Contours computed for a quadratic form with axis ratio $2{:}1$.)}\label{ch05:fig-ridgelasso}
\end{figure}

\begin{remark}[Orthonormal design; not discussed in class]
If $X\T X=I_p$, then $\|\bm y-X\bm\beta\|^2=\|\bm y-X\hat{\bm\beta}\|^2+\|\bm\beta-\hat{\bm\beta}\|^2$ with $\hat{\bm\beta}=X\T\bm y$, and both problems decouple:
\[
  \hat\beta^{\text{ridge}}_j=\frac{\hat\beta_j}{1+\lambda},\qquad
  \hat\beta^{\text{lasso}}_j=\operatorname{sign}(\hat\beta_j)\bigl(|\hat\beta_j|-\lambda/2\bigr)_+ .
\]
Ridge rescales every coefficient; the lasso subtracts a constant and sets the small ones to exactly zero
(``soft thresholding'').
\end{remark}

Regularised models reappear in the 2013 lecture on regularised pricing and risk models (\cref{ch:regularized}), in
machine learning (\cref{ch:ml}), and PCA in \cref{ch:pca}.

% ------------------------------------------------------------------
\section{Case studies}\label{ch05:sec-cases}

\subsection{Sector ETFs regressed on index ETFs}
\begin{casestudy}[ETF case study (\file{lec06_4} = \file{ETF_Casestudy.pdf}, sector XLP)]
\textbf{Data.} Adjusted daily prices from Yahoo (2000-01-03 to 2024-09-19) of the 13 SPDR ETFs listed in
\file{spyders_symbols_labels.csv}: four index ETFs (SPY: S\&P~500, MDY: mid-caps, QQQ: Nasdaq-100, DIA: Dow Jones) and
nine sector ETFs (XLB materials, XLV health care, XLP consumer staples, XLY consumer discretionary, XLE energy, XLF
financials, XLI industrials, XLK technology, XLU utilities). Converted to weekly log returns (1289 weeks; saved in
\file{casestudy_1_0_etfs.RData}); the analysis uses 2013--2024 ($n=612$ weeks). Response: XLP (annualised mean 10.0\%,
volatility 13.7\%, against 13.7\% and 16.1\% for SPY).

\textbf{OLS} \ts{24}{11}{42:42}: $\hat y=0.0002+1.007\,\text{SPY}-0.368\,\text{MDY}-0.314\,\text{QQQ}+0.400\,\text{DIA}$, all
slopes with $|t|>4.3$; $R^2=0.642$, $\hat\sigma=0.0114$ (weekly), $F=271.9$. The mixed signs are a symptom of the collinear
regressors (index correlations $0.77$--$0.95$): XLP loads positively on the market but less on its growth and mid-cap
parts. Diagnostics: influential weeks 2016-11-11 (studentised residual $-4.8$), 2020-03-06 ($+4.4$), 2022-05-20
($-5.8$); the COVID week 2020-03-20 has hat value $0.135$ against $p/n=0.008$. The curvature tests are all significant
at 5\% ($p\approx0.02$--$0.04$): some non-linearity in the market response.

\textbf{$z$-scores.} Regressing on standardised index returns reproduces every $t$ value and $R^2$; only the scale of the
coefficients and the intercept change.

\textbf{PCA of the regressors} \ts{24}{11}{43:42}: the four standardised index returns have PC variances in proportions
$90.9\%$, $6.5\%$, $2.3\%$, $0.3\%$ (cumulative $97.4\%$ with two, $99.7\%$ with three). PC1 has equal loadings
($0.48$--$0.52$, the ``market''), PC2 contrasts QQQ ($+0.79$) with MDY and DIA.

\textbf{PC regressions} \ts{24}{11}{44:48}: simple regressions of XLP on each PC have $R^2=0.545$, $0.006$ ($p=0.052$),
$0.079$ and $0.012$. Their sum is $0.642$, the $R^2$ of the full regression, because the PCs are orthogonal
(\cref{ch05:ex-r2orth}). The third component, with only 2.3\% of the regressor variance, explains 13 times more of
XLP than the second: dropping it (PC123 would keep it, but a ``first two'' rule would not) changes the answer. Using PCs 1--3
gives coefficients very different from OLS; using the PCs with $|t|>\sqrt2$ (all four here) reproduces OLS exactly
\ts{24}{11}{51:03}.

\textbf{Ridge and lasso} (\texttt{glmnet}, $\lambda$ by 10-fold cross-validation over $10^{-5}$--$10^{3}$)
\ts{24}{11}{47:55}: coefficient traces shrink to zero, not monotonically for ridge; along the lasso path the third
variable (QQQ) leaves first, then MDY, then DIA \ts{24}{11}{50:01}. With the CV-optimal $\lambda$ both fits are
almost identical to OLS ($R^2=0.6417$ ridge, $0.6417$ lasso vs $0.6418$): with $n=612$ and 4 regressors there is little
to regularise \ts{24}{11}{51:03}.

\textbf{XLB version} (\file{ETF_Casestudy_XLB.pdf}, materials): $R^2=0.805$, loadings SPY $0.85$, MDY $0.50$, QQQ
$-0.38$, DIA not significant ($t=0.43$); PC1 alone gives $R^2=0.73$. The week 2020-03-20 is now very influential (Cook's
distance $0.99$), and all curvature tests are highly significant ($p<0.001$).

\textbf{Remark on the code.} Cumulative returns are plotted as $\prod_t(1+r_t)$ although $r_t$ are log returns; the correct
growth factor is $\exp(\sum_tr_t)$. The difference (about $-r_t^2/2$ per week) only affects the plots.
\end{casestudy}

\begin{finance}[Why regress a sector on the indices?]
The fitted coefficients are a \emph{hedge}: holding XLP and shorting $1.007$ SPY, $-0.368$ MDY, \dots\ leaves (in
sample) only the residual, i.e.\ the sector-specific part of the return, with most of the market risk removed
\ts{24}{11}{40:39}. The same regression, run with a hedge fund's returns as the response and liquid instruments as
regressors, is the basis of \emph{hedge-fund replication}: if most of a fund's return is explained by tradable factors,
those can be bought cheaply \ts{24}{11}{41:39}. (Regressing an index on its sectors would instead give $R^2\approx1$, as
noted by a student \ts{24}{11}{39:39}.)
\end{finance}

\subsection{CAPM regressions}\label{ch05:capm}
The CAPM (Sharpe 1964, Lintner 1965; theory in \cref{ch:portfolio}) states that in an efficient market in which all
investors hold mean--variance efficient portfolios and can borrow and lend at the risk-free rate $r_f$,
\[
  \E[R_j]=r_f+\beta_j\bigl(\E[R_M]-r_f\bigr),\qquad \beta_j=\frac{\Cov(R_j,R_M)}{\Var(R_M)},
\]
where $R_M$ is the return of the market portfolio \ts{24}{11}{54:12}. With excess returns $R^*_{j,t}=R_{j,t}-r_{f,t}$ and
$R^*_{M,t}$ observed over time, fit
\[
  R^*_{j,t}=\alpha_j+\beta_jR^*_{M,t}+\varepsilon_{j,t},\qquad \varepsilon_{j,t}\ \text{white noise},
\]
the alpha--beta regression of \cref{def:alphabeta} \ts{24}{11}{55:18}. Taking expectations, the CAPM holds iff
$\alpha_j=0$ \ts{24}{11}{56:27}. Two natural tests: $\alpha_j=0$ (is the asset efficiently priced?) and $\beta_j=1$ (does it carry
more or less systematic risk than the market?) \ts{24}{11}{1:06:14}.

\begin{casestudy}[CAPM and macro factors for GE and XOM, 2000--2013 (2013 Case Study 1)]
\textbf{Data.} \file{fm_casestudy_1_0.r} downloads daily prices of BAC, GE, JDSU, XOM and the S\&P~500 (Yahoo) and from FRED
the constant-maturity Treasury yields (3 months, 1, 5, 10 years), Moody's Aaa and Baa yields and the WTI oil price, for
2000-01-03 to 2013-05-31 (3373 days). \file{fm_casestudy_1_rcode.r} computes daily log returns and the daily risk-free
return $\log\bigl(1+r_{3M}\,\Delta t/360\bigr)$ with $\Delta t$ the number of calendar days since the previous close (money-market
convention).

\textbf{GE:} $\hat\alpha=-0.00013$ ($t=-0.56$), $\hat\beta=1.184$ ($t=66.6$), $R^2=0.568$: consistent with CAPM, with more market risk
than average. Influence: 129 of 3372 days exceed the leverage threshold; the plot of hat values over time shows they are
concentrated in the 2008 crisis, when market moves were extreme.

\textbf{Adding oil} (daily log change of WTI): for GE the oil coefficient is $-0.034$ ($t=-3.6$): GE falls when oil rises. For
XOM, simple CAPM $\hat\beta=0.830$, $R^2=0.451$ (lower than GE: an oil company tracks the market less); with oil,
$\hat\beta=0.782$ and oil coefficient $+0.132$ ($t=16.1$), $R^2$ up to $0.491$. The high-leverage days are those with extreme
market or oil moves, far from the centre of the $(R^*_M,\text{oil})$ scatter (the Mahalanobis interpretation of 2013
PS3).

\textbf{Lesson.} The market factor explains half of a large stock's daily variance; a second factor can be highly
significant while adding little $R^2$.
\end{casestudy}

\begin{casestudy}[CAPM for GE and for the S\&P~500 stocks (2024 L11)]
\textbf{GE 2017--2023} (\file{lec08}; excess return \texttt{Ra.RF} regressed on the market excess return
\texttt{Mkt.RF} of the Fama--French data): see \cref{ch05:ex-gebeta} and the regime test of \cref{ch05:sec-chow}. The residuals
are clearly non-normal \ts{24}{11}{1:00:53}; see the box on heavy tails in \cref{ch05:sec-mle}.

\textbf{All S\&P~500 stocks} (note shown in class, not published) \ts{24}{11}{1:14:45}: the same regression for the
380--400 constituents with a long enough history, grouped by sector. $R^2$ ranges over $(0,1)$, typically above $0.2$.
Alphas vary by sector (positive for construction and computer technology, negative for conglomerates)
\ts{24}{11}{1:15:47}, but the $p$-values of $\alpha_j=0$ are above $0.05$ for all but a handful of stocks, roughly the 5\%
expected by chance: the CAPM is not rejected stock by stock \ts{24}{11}{1:16:47}. Higher betas tend to go with higher
alphas \ts{24}{11}{1:17:53}. Betas are strongly sector-dependent: consumer staples lowest (the ten lowest betas, $0.37$
down to $0.18$, are mostly staples), technology and consumer discretionary highest \ts{24}{11}{1:18:59}. The largest
significant alpha, ENPH (energy), is $0.003$ per day, i.e.\ $(1.003)^{252}-1\approx113\%$ per year
\ts{24}{11}{1:20:01}.

\textbf{Next step.} Multi-factor models with the Fama--French factors, already downloaded with the data
\ts{24}{11}{1:22:11} (\cref{ch:pca}).
\end{casestudy}

\begin{finance}[Using alpha and beta]
Beta is the hedge ratio against the index and the measure of systematic risk that the CAPM says is rewarded;
alpha is the part of the return not explained by market exposure (\cref{def:alphabeta}). In practice betas are estimated
on rolling windows (\cref{ch05:sec-chow}) and often shrunk towards $1$, e.g.\ the widely quoted ``adjusted beta''
$\frac23\hat\beta+\frac13$ (not discussed in class): the same bias--variance trade as ridge regression, with the prior
centred at the market beta.
\end{finance}

% ------------------------------------------------------------------
\section{Differences between the 2013 and 2024 lectures}
\begin{coursenote}[2013 L6 vs 2024 L6, L8, L11]
The core (OLS, hat matrix, Gauss--Markov, GLS, normal distribution theory, ML, M-estimators) is the same material in
both years, with nearly identical slides. \textbf{2013} spends one lecture on it and adds the QR decomposition and
the full proofs of the Gauss--Markov and distribution theorems \ts{13}{6}{1:08:26}; its case studies are the CAPM with
oil (Case Study~1) and the yuan regimes (Case Study~2). \textbf{2024} spreads the topic over parts of three lectures and
adds real-data diagnostics (prostate data), hypothesis testing and regime-change tests (\file{lec08}), ridge, PC and lasso
regression, the ETF case study and the S\&P~500 CAPM study. The 2013 lecture ran out of time before M-estimation
\ts{13}{6}{1:21:11}; it is covered in the notes and in 2024 L11.
\end{coursenote}

% ------------------------------------------------------------------
\begin{keyformulas}
\begin{align*}
&\text{OLS:} && \hat{\bm\beta}=(X\T X)^{-1}X\T\bm y=R^{-1}Q\T\bm y,\\
  &&& H=X(X\T X)^{-1}X\T=QQ\T,\quad \hat{\bm\varepsilon}=(I-H)\bm y\\
&\text{Geometry:} && H^2=H=H\T,\ \tr H=p,\\
  &&& X\T\hat{\bm\varepsilon}=\bm 0,\ \text{TSS}=\text{ESS}+\text{RSS},\ R^2=\Corr(\bm y,\hat{\bm y})^2\\
&\text{Gauss--Markov:} && \E\hat{\bm\beta}=\bm\beta,\ \Cov\hat{\bm\beta}=\sigma^2(X\T X)^{-1};\\
  &&& \bm c\T\hat{\bm\beta}\ \text{is BLUE}\\
&\text{GLS:} && \hat{\bm\beta}_{\text{GLS}}=(X\T\Sigma^{-1}X)^{-1}X\T\Sigma^{-1}\bm y,\\
  &&& \Cov=\sigma^2(X\T\Sigma^{-1}X)^{-1}\\
&\text{Normal model:} && \hat{\bm\beta}\sim N_p(\bm\beta,\sigma^2(X\T X)^{-1})\\
  &&& \perp\!\!\!\perp\text{RSS}\sim\sigma^2\chi^2_{n-p},\ \hat\sigma^2=\tfrac{\text{RSS}}{n-p}\\
&t\text{ test:} && \hat t_j=(\hat\beta_j-b)/(\hat\sigma\sqrt{C_{jj}})\sim t_{n-p},\quad C_{jj}=[(X\T X)^{-1}]_{jj}\\
&F\text{ test:} && F=\tfrac{(\text{RSS}_0-\text{RSS})/q}{\text{RSS}/(n-p)}\sim F_{q,n-p}\\
&&& \text{RSS}_0-\text{RSS}=(A\hat{\bm\beta}-\bm b)\T[A(X\T X)^{-1}A\T]^{-1}(A\hat{\bm\beta}-\bm b)\\
&\text{Chow test:} && F=\tfrac{(\text{RSS}_0-\text{RSS}_A-\text{RSS}_B)/p}{(\text{RSS}_A+\text{RSS}_B)/(n-2p)}\\
  &&& \sim F_{p,n-2p}\\
&\text{ML:} && \hat{\bm\beta}_{\text{ML}}=\hat{\bm\beta},\quad \hat\sigma^2_{\text{ML}}=\text{RSS}/n\\
&\text{Diagnostics:} && \Var\hat\varepsilon_i=\sigma^2(1-H_{ii}),\ \ CD_i=\tfrac{r_i^2}{p}\tfrac{H_{ii}}{1-H_{ii}},\\
  &&& \hat{\bm\beta}_{(i)}=\hat{\bm\beta}-\tfrac{(X\T X)^{-1}\bm x_i\hat\varepsilon_i}{1-H_{ii}}\\
&\text{Shrinkage:} && \hat{\bm y}^{\text{ridge}}=\textstyle\sum_j\frac{d_j^2}{d_j^2+\lambda}c_j\bm u_j,\\
  &&& \hat{\bm y}^{\text{pc}}=\sum_{j\le m}c_j\bm u_j,\\
  &&& \lambda=\sigma^2/\tau^2\ \text{(Bayes)}\\
&\text{CAPM:} && R^*_{j,t}=\alpha_j+\beta_jR^*_{M,t}+\varepsilon_{j,t},\quad H_0:\alpha_j=0
\end{align*}
\end{keyformulas}

% ------------------------------------------------------------------
\section*{Exercises}
\addcontentsline{toc}{section}{Exercises}
Standalone exercises follow the section they practise; the connected series (2013 PS3, Problems~2 and~3, which build on one another) are kept together here.

\subsection*{From 2013 Problem Set 3}
Throughout, $\bm y=X\bm\beta+\bm\varepsilon$ with $X\in\R^{n\times p}$ of full rank $p\le n$, $\hat{\bm\beta}=(X\T X)^{-1}X\T\bm y$, $H=X(X\T X)^{-1}X\T$,
$\hat{\bm\varepsilon}=(I_n-H)\bm y$.

\begin{exercise}[PS3 (2013), Problem 2(a): Sherman--Morrison--Woodbury]\label{ch05:ex-ps3-2a}
Let $A$ be a $p\times p$ symmetric matrix of rank $p$ and $a,b$ two $q\times p$ matrices of rank $q<p$. Prove that, whenever the
inverses exist,
\[
  (A+a\T b)^{-1}=A^{-1}-A^{-1}a\T\bigl(I_q+bA^{-1}a\T\bigr)^{-1}bA^{-1}.
\]
\end{exercise}

\begin{exercise}[PS3 (2013), Problem 2(b)--(d): case-deletion influence]\label{ch05:ex-ps3-2bd}
(b) Using Sherman--Morrison--Woodbury, show that the least-squares estimate with case $i$ deleted is
$\hat{\bm\beta}_{(i)}=\hat{\bm\beta}-(X\T X)^{-1}\bm x_i\hat\varepsilon_i/(1-H_{ii})$, where $\bm x_i\T$ is the $i$-th row of $X$ and $\hat\varepsilon_i=y_i-\bm x_i\T\hat{\bm\beta}$.
(c) Cook's distance is $CD_i=\|\hat{\bm y}-\hat{\bm y}_{(i)}\|^2/(p\hat\sigma^2)$ with $\hat{\bm y}_{(i)}=X\hat{\bm\beta}_{(i)}$. Show that
$CD_i=\dfrac{\hat\varepsilon_i^2}{p\hat\sigma^2}\cdot\dfrac{H_{ii}}{(1-H_{ii})^2}$.
(d) Let $\hat\sigma^2_{(i)}$ be the unbiased estimate of $\sigma^2$ with case $i$ deleted. Show that
$\hat\sigma^2_{(i)}=\hat\sigma^2+\frac{1}{n-p-1}\bigl(\hat\sigma^2-\frac{\hat\varepsilon_i^2}{1-H_{ii}}\bigr)$.
\end{exercise}

\begin{exercise}[PS3 (2013), Problem 3(a)--(b): constrained least squares via QR]\label{ch05:ex-ps3-3ab}
Consider $H_0:\beta_{k+1}=\dots=\beta_p=0$. Let $X_I=[X_{[1]}\cdots X_{[k]}]$, $\hat{\bm\beta}_I=(X_I\T X_I)^{-1}X_I\T\bm y$ and
$\hat{\bm\beta}_0=(\hat{\bm\beta}_I,\bm 0_{p-k})$.
\begin{enumerate}[(a)]
  \item Show that $\hat{\bm\beta}_0$ minimises $(\bm y-X\bm\beta)\T(\bm y-X\bm\beta)$ subject to $\beta_j=0$ for $j>k$.
  \item With $X=QR$, show that $X_I=Q_IR_I$ where $Q_I$ is the first $k$ columns of $Q$ and $R_I$ the upper-left $k\times k$ block of $R$;
      verify $\hat{\bm\beta}_I=R_I^{-1}Q_I\T\bm y$ and $\hat{\bm y}_I=H_I\bm y$ with $H_I=Q_IQ_I\T$.
\end{enumerate}
\end{exercise}

\begin{exercise}[PS3 (2013), Problem 3(c)--(d): sequential ANOVA]\label{ch05:ex-ps3-3cd}
Let $A=\bigl(\begin{smallmatrix}Q\T\\ W\T\end{smallmatrix}\bigr)$ be orthogonal as in the proof of \cref{ch05:thm-normal} and $\bm z=A\bm y=(z_1,\dots,z_n)\T$.
(c) Prove that $\bm y\T\bm y=\sum_{i=1}^nz_i^2$, $\hat{\bm y}\T\hat{\bm y}=\sum_{i=1}^pz_i^2$, $\hat{\bm\varepsilon}\T\hat{\bm\varepsilon}=\sum_{i=p+1}^nz_i^2$, and for the constrained model
$\hat{\bm y}_I\T\hat{\bm y}_I=\sum_{i=1}^kz_i^2$, $\hat{\bm\varepsilon}_I\T\hat{\bm\varepsilon}_I=\sum_{i=k+1}^nz_i^2$.
(d) Using $\bm z_Q\sim N_p(R\bm\beta,\sigma^2I_p)$ independent of $\bm z_W\sim N_{n-p}(\bm 0,\sigma^2I_{n-p})$, deduce that
$SS_{\text{ERROR}}=\hat{\bm\varepsilon}\T\hat{\bm\varepsilon}\sim\sigma^2\chi^2_{n-p}$; that under $H_0$
$SS_{\text{REG}(k+1,\dots,p\mid1,\dots,k)}=\hat{\bm y}\T\hat{\bm y}-\hat{\bm y}_I\T\hat{\bm y}_I=\sum_{i=k+1}^pz_i^2\sim\sigma^2\chi^2_{p-k}$, independent of
$SS_{\text{ERROR}}$; and that $\hat F=\dfrac{SS_{\text{REG}}/(p-k)}{SS_{\text{ERROR}}/(n-p)}\sim F_{p-k,\,n-p}$. Arrange the results in the ANOVA
table with rows ``regression on $1,\dots,k$'' ($\hat{\bm y}_I\T\hat{\bm y}_I$, $k$ d.f.), ``regression on $k+1,\dots,p$ adjusting for
$1,\dots,k$'' ($p-k$ d.f., mean square $MS_0$, $\hat F=MS_0/MS_{\text{Error}}$), ``error'' ($n-p$ d.f.) and ``total'' ($\bm y\T\bm y$,
$n$ d.f.).
\end{exercise}

\printsolutions

% !TEX root = ../main.tex
\chapter{Linear Rates: Products and Models}\label{ch:rates}
\begin{paracol}{2}

\begin{sources}
\begin{tabularx}{\linewidth}{@{}l l L@{}}
\toprule
\textbf{Lecture} & \textbf{Speaker} & \textbf{Content}\\
\midrule
2024 L7 & A.~Gunstensen (Mizuho) & Linear rates products and the size of their markets; LIBOR and its replacement by SOFR; discounting = funding; deposits, SOFR futures and interest-rate swaps; swap valuation and par swap rates; yield-curve construction (knot points, CDF and spline interpolation, calibration, FOMC curves, turns); delta ladders and gamma; P\&L attribution; clearing and the CCP basis; electronic trading and price making; Q\&A on the rates outlook, research roles and municipal bonds.\\
\bottomrule
\end{tabularx}

\smallskip
\textbf{Files.} Slides \emph{Linear Rates Models} \file{mit18_642_f24_lec07.pdf} (45 slides). Slide~39 (fast
model pricing) and slides 40--44 (bonds) were skipped in class for lack of time \ts{24}{7}{1:14:29}; they are
covered in \cref{ch06:sec-etrading,ch06:sec-bonds}. The data file \file{fedbondyielddata_20230911.csv} and the
notebook \file{feddataanalysis.ipynb} are listed on OCW together with Lecture~9 (PCA of Treasury yield changes)
and are treated in \cref{ch:pca}.
\textbf{Problem sets.} None; all exercises at the end are \emph{not from the course}.
\textbf{2013.} No counterpart: the 2013 course has no lecture on linear rates products. The dynamics of the
whole curve are the subject of the 2013 HJM lecture (\cref{ch:hjm}).
\end{sources}

\section*{Motivation}
Bonds, deposits, futures and swaps are the largest and most liquid markets in the world, and their payoffs are
fixed or linear in interest rates: pricing them requires no model of rate \emph{dynamics}, only today's
\emph{discount curve}. \Cref{ch:markets} took that curve as given (a flat rate, or yields read off a screen,
\cref{fig:yieldcurve}). This chapter asks the practitioner's questions: which rate should one discount with,
which instruments reveal the curve, how does one fill the gaps between a few dozen quotes, how does one hedge
and explain the daily P\&L of a book of swaps? The lecturer announced a deliberately ``non-mathy'' talk: a
house is designed by thinking about rooms and how they connect, not by spending most of the time on hammers
\ts{24}{7}{1:01}. The mathematics is elementary (linear algebra, interpolation, Taylor expansions); what matters
is the set of conventions and design choices, and the consequences of each choice for hedging. The curve built
here is the input of every rates-option model; its PCA is in \cref{ch:pca} and its stochastic evolution in
\cref{ch:hjm}.

% ------------------------------------------------------------------
\section{Linear products and their markets}\label{ch06:sec-markets}

\begin{definition}[Linear product]
A rates product is \emph{linear} if its payoff is either fixed or a linear function of observed interest rates
\ts{24}{7}{2:06}: (vanilla) bonds, certificates of deposit, interest-rate futures, forward rate agreements,
interest-rate swaps. There is no optionality (no caps, swaptions, callable or range-accrual features).
\end{definition}
The PV of a linear product is a linear function of discount factors and forward rates, so it is fixed by
today's curve alone: no volatility enters, except through small corrections (futures margining, averaging,
payment delays) discussed below. Why should a quant care \ts{24}{7}{3:07}?
\begin{itemize}
  \item \textbf{Size and liquidity} (slide 4). Bonds outstanding about \$128T (ICMA, 2020), interest-rate swaps
        \$573T notional (BIS, 2023), SOFR futures open interest above \$10T. Daily volumes: US Treasuries
        about \$900B, swaps \$800B, SOFR futures about 4 million contracts. In a market this size, a tiny margin
        is a large business \ts{24}{7}{3:07}.
  \item \textbf{Hedging and inputs.} Linear products hedge the complex ones, and their curves are inputs of the
        complex models: one builds up from the simple stuff \ts{24}{7}{4:08}.
  \item \textbf{Post-2008 complexity.} Before the crisis vanilla modelling was considered trivial. Funding,
        clearing and margining (\cref{ch06:sec-funding,ch06:sec-clearing}) made even ``\$1 paid in ten
        years'' hard to value, with no unique answer \ts{24}{7}{4:08}. Regulation (including model-risk management)
        ended much of the business in opaque, highly levered securitisations that had been sold to buyers who
        could not assess them \ts{24}{7}{5:08}.
\end{itemize}

% ------------------------------------------------------------------
\section{Interest rates and reference rates}\label{ch06:sec-rates}

\subsection{Simple interest and accrual}
An interest rate is the rate at which one earns (or pays) interest on a principal \ts{24}{7}{7:10}. Money-market
instruments use \emph{simple} interest:
\begin{equation}\label{ch06:eq-simple}
  I = N\,r\,\Delta t,\qquad
  \mathrm{PV}(X \text{ paid after } \Delta t) = \frac{X}{1+r_d\,\Delta t},
\end{equation}
where $\Delta t$ is the \emph{accrual fraction} of the period and $r_d$ a discounting rate. Example from class:
\$1000 in one year discounted at 5\% is worth $1000/1.05=\$952.38$ today \ts{24}{7}{8:11}. The interest rate
$r$ of an investment is easy to observe (the rate a bank advertises on a CD); the discounting rate $r_d$ is not,
and it depends on who is doing the discounting (\cref{ch06:sec-funding}).

\begin{coursenote}[Day counts (not discussed in class)]
$\Delta t$ is computed from dates with a \emph{day-count convention}: ACT/360 (actual days$/360$; SOFR, money
markets, the floating legs of USD swaps), ACT/365F, 30/360 (bond-style), ACT/ACT (Treasuries). Dates themselves
are moved off holidays by a \emph{business-day convention} (e.g.\ ``modified following''). The lecturer's advice
to job candidates: after a year on a swaps desk you will know \emph{date arithmetic}, and at least 90\% of all
problems with swaps are date problems \ts{24}{7}{34:51}. The number of days between two dates has no unique answer;
it depends on the convention.
\end{coursenote}

\subsection{LIBOR}
The \emph{London Interbank Offered Rate} was created in the late 1960s as a reference for syndicated bank
lending: corporate loans are mostly floating-rate, resetting every one or three months, and a common index was
needed \ts{24}{7}{9:13}. Every day a panel of about 20 banks reported the rate at which they could borrow
unsecured for various terms; the rate was a trimmed average of the submissions \ts{24}{7}{10:16}. The first
swaps (a cross-currency swap involving IBM in 1981, then single-currency interest-rate swaps) were built on it, and
by LIBOR's cessation some \$350--400T of swaps referenced it \ts{24}{7}{11:16}.

Its weaknesses surfaced after 2008 (slide 7):
\begin{itemize}
  \item it was a \emph{quoted}, not a transacted rate: submitters estimated where they \emph{would} borrow;
  \item it was \emph{unsecured}, and after Lehman almost nobody lent unsecured term money: some tenors saw a
        handful of real transactions in six months while a rate had to be published every day
        \ts{24}{7}{13:26};
  \item hence it could be manipulated. A book can have millions of dollars of P\&L per basis point
        (bp $=0.01\%$) of a fixing, and a group of traders coordinated with their banks' submitters to nudge the
        fixings; this led to prison sentences and multi-billion-dollar fines \ts{24}{7}{12:21}.
\end{itemize}
Regulators (led from London) required a replacement ``anchored in actual transactions and hard to
manipulate'' \ts{24}{7}{13:26}. USD LIBOR ended in June 2023; sterling earlier, Canada and Mexico around 2024
\ts{24}{7}{18:40}. The migration of mortgages, bonds, loans and hundreds of trillions of swaps was enormous
and tedious, comparable to the euro conversion and Y2K \ts{24}{7}{15:31}.

\begin{finance}[Why banks liked an unsecured index]
A bank's floating-rate loans pay LIBOR plus a fixed spread. When the banking system is under stress, bank
credit spreads widen and LIBOR rises, so the bank's \emph{income} rises exactly when its own funding becomes
expensive: a natural hedge (a ``right-way'' feature) \ts{24}{7}{17:40}, \ts{24}{7}{19:44}. A secured rate
contains no bank credit risk and lacks this property. The flip side: LIBOR-based hedges of instruments that
depend on government rates suffer when the LIBOR--Treasury spread blows out. In 2011--12 this cost the
lecturer's employer of the time a write-off of several billion dollars on structured notes (\$8.4B quoted in
class), most of which came back later \ts{24}{7}{14:31}.
\end{finance}

\subsection{SOFR}
A \emph{repo} (repurchase agreement) is a collateralised loan: I deliver \$100M of Treasuries and receive, say,
\$98M cash, agreeing to buy the bonds back the next day; the interest paid is the \emph{repo rate} (reverse
repo for the other side) \ts{24}{7}{16:35}. This is a real market with real volume (about \$800B a day
according to slide~8).

\begin{definition}[SOFR]
The \emph{Secured Overnight Financing Rate} is the main USD reference rate: a volume-weighted median of the
rates on overnight repo transactions collateralised by US Treasuries, published every business day
\ts{24}{7}{15:31}.
\end{definition}

\begin{coursenote}[Erratum (spoken)]
In class SOFR is expanded as ``Securitized Overnight Funding Rate'' and described as an average. The name is
\emph{Secured Overnight Financing Rate}, computed by the Federal Reserve Bank of New York (published since April
2018, selected in 2017 by the Fed-convened Alternative Reference Rates Committee) as a volume-weighted
\emph{median} of transaction rates. Similarly ``OAS rate'' in the transcript at 8:11 should read OIS
(overnight index swap) rate, as on slide~6.
\end{coursenote}

SOFR has three inconvenient features \ts{24}{7}{16:35}: it is a \emph{daily} rate, while nobody wants to pay
interest daily; a term rate built from it is known only at the \emph{end} of the period (set in arrears), while
borrowers like to know their payment in advance; and it is \emph{secured}, while banks prefer to lend
unsecured (previous box) \ts{24}{7}{17:40}.

\paragraph{From daily fixings to term rates} \ts{24}{7}{20:49}. Over an accrual period of $d$ days, let $r_i$
be the SOFR fixing of business day $i$, applying for $d_i$ calendar days ($d_i=3$ for a Friday fixing, which
covers the weekend), $d=\sum_{i=1}^nd_i$. The two standard term rates (slide 9) are
\begin{equation}\label{ch06:eq-sofr}
  R^{\mathrm{comp}} = \frac{360}{d}\Biggl[\,\prod_{i=1}^{n}\Bigl(1+\frac{r_i d_i}{360}\Bigr)-1\Biggr],
  \qquad
  R^{\mathrm{avg}} = \frac1d\sum_{i=1}^{n} r_i d_i .
\end{equation}
$R^{\mathrm{comp}}$ is exactly what one earns by rolling an overnight deposit (interest on interest), and
$R^{\mathrm{comp}}\ge R^{\mathrm{avg}}$ when rates are non-negative (\cref{ch06:ex-sofr}). Most USD products pay
compounded SOFR: it is simpler to calibrate, and it needs no correction for rate volatility, which an average
does \ts{24}{7}{21:49}. The reason is \cref{ch06:prop-float} below: a compounded leg is replicated exactly by
rolling a deposit.

\begin{example}[Compounding vs averaging (not from class)]
A 91-day period with a fixing every day ($d_i=1$): 5.30\% for 40 days, then 5.05\% after a 25\,bp FOMC cut.
Then $R^{\mathrm{avg}}=5.1599\%$ and $R^{\mathrm{comp}}=5.1933\%$: a 3.3\,bp difference. It is ``a few basis
points running'', irrelevant for a retail borrower but twenty times the bid--offer of a swap market maker.
\end{example}

\paragraph{Other rates} \ts{24}{7}{21:49}--\ts{24}{7}{26:16} (slide 10).
\begin{itemize}
  \item \textbf{Term SOFR}: 1-, 3-, 6- and 12-month rates published by the CME by fitting a model of daily rates to
        SOFR futures prices and averaging the fitted daily rates over the term \ts{24}{7}{24:02}. It is set in
        advance, like LIBOR, and popular in corporate loans, whose borrowers do not want to handle daily
        fixings. To prevent it from becoming a new LIBOR, the Fed restricts dealers to trading it only to hedge
        their loan books. The market is therefore one-way, and term-SOFR swaps trade 4--5\,bp away from vanilla
        SOFR swaps, 20--30 times the bid--offer \ts{24}{7}{25:09}.
  \item \textbf{Fed Funds}: the older unsecured overnight interbank rate, still widely used; plus Prime,
        DISCO (a Fannie Mae discount-note rate, the lecturer's favourite name), MMD (a municipal benchmark) and many
        others \ts{24}{7}{26:16}.
\end{itemize}

% ------------------------------------------------------------------
\section{Discounting equals funding}\label{ch06:sec-funding}

Which rate should discount a future cash flow \ts{24}{7}{26:16}? Assume (slides 11--12) we own a fixed payment
$F$ at time $T$, worth $P(t)$ today, and that we bought it with borrowed money: in this world nobody pays out
of pocket, positions are funded. Consider one day, from $t_0$ to $t_1$, with no market move
\ts{24}{7}{27:21}:
\begin{align*}
  \text{interest paid on the funding:}\quad & I(t_1) = P(t_0)\,\frac{r_b}{360},\\
  \text{accretion of the claim:}\quad & P(t_1) = P(t_0)\Bigl(1+\frac{r_d}{360}\Bigr),\quad\text{i.e. }\
  \Delta P = P(t_0)\,\frac{r_d}{360},
\end{align*}
where $r_b$ is our borrowing rate and $r_d$ the rate used to discount $F$.

\begin{proposition}[Discounting rate $=$ funding rate]\label{ch06:prop-funding}
Let an agent borrow and lend at the rate $r_b$ over $[t,T]$ ($n$ days, daily compounding, rates constant as on
the slides). The only price of $F$ paid at $T$ that gives the agent no riskless trade is
$P(t)=F\,(1+r_b/360)^{-n}$: future cash flows must be discounted at the funding rate, $r_d=r_b$.
\end{proposition}
\begin{proof}
The funded position has a daily P\&L $P(r_d-r_b)/360$, known in advance. If $r_d>r_b$, buying the claim with
borrowed money earns a riskless carry; if $r_d<r_b$, selling it and lending the proceeds does. Summing over the
holding period, the terminal P\&L of ``borrow $P$, buy the claim'' is $F-P(1+r_b/360)^n$, which vanishes for
every holding period only if $P$ is discounted at $r_b$ \ts{24}{7}{28:24}.
\end{proof}

\begin{remark}
If the agent borrows at $r_b$ but can only lend at $r_l<r_b$, the argument only gives the band
$F(1+r_b/360)^{-n}\le P\le F(1+r_l/360)^{-n}$. Different agents have different funding rates (the lecturer's
example: borrowers at 4\%, 6\% and 8\% depending on their credit) and therefore different values for the same
contract: there is no single price \ts{24}{7}{9:13}. The lecturer's one-sentence summary of the lecture:
\emph{funding and discounting are the same thing} \ts{24}{7}{28:24}.
\end{remark}

\begin{intuition}[A stationary state]
Think of the funded position as a closed system in which ``nothing happens'': the claim drifts up at $r_d$,
the loan at $r_b$. Zero P\&L in the absence of market moves (a stationary state) requires equal drifts. The
discount curve is not a law of nature but the rate at which the holder's money actually grows or costs;
compare \cref{sec:discounting}, where the same argument with a single universal rate gave $1/(1+r)$.
\end{intuition}

\begin{finance}[Funding desks and OIS discounting]
Banks do not let each trader manage the funding of a thousand counterparty relationships: a central funding
desk (treasury) lends to all desks at a generic internal rate and manages the funding bases and small
collateral exposures itself \ts{24}{7}{9:13}. For collateralised trades the relevant funding rate is the rate
paid on the collateral, an overnight rate: hence the market standard of discounting with an OIS (now SOFR)
curve \ts{24}{7}{8:11}. The costs of funding uncollateralised trades are priced as funding valuation
adjustments, part of the XVA family of \cref{ch:counterparty}.
\end{finance}

% ------------------------------------------------------------------
\section{The instruments that reveal the curve}\label{ch06:sec-instruments}
To build a curve we start from what trades actively and is easy to observe (slide 13): whatever the model, it
must reprice these instruments, because they can be traded \ts{24}{7}{28:24}. The workhorses are cash deposits
(CDs), interest-rate futures and swaps \ts{24}{7}{29:26}. We write each price as a function of the discount
function $Z(t)$ (value today of \$1 paid at $t$; $\Lambda_t$ in \cref{ch:markets}), then invert.

\subsection{Deposits and certificates of deposit}
Short-term money can be invested at a fixed simple rate $r(t)$ for maturities up to about a year (CDs,
interbank deposits; also repo and Treasury bills, not discussed in class). Investing 1 today returns
$1+r(t)\Delta$ at $t$, so (slide 14)
\begin{equation}\label{ch06:eq-deposit}
  1=\bigl[1+r(t)\Delta\bigr]Z(t)\qquad\Longrightarrow\qquad Z(t)=\frac{1}{1+r(t)\Delta}
  \quad\text{\ts{24}{7}{29:26}.}
\end{equation}
Deposits are well observed but awkward for a swaps desk: they are \emph{funded} (a \$100M CD requires \$100M of
cash and balance sheet), whereas derivatives need little cash up front and allow more leverage
\ts{24}{7}{30:32}.

\subsection{Forward rates and FRAs}
A floating payment that sets in the future can be replicated with fixed flows \ts{24}{7}{39:04}. Let the rate
apply to $[t_1,t_2]$ with accrual fraction $\Delta$, and suppose zero-coupon bonds maturing at $t_1$ and $t_2$
can be traded.

\begin{proposition}[Implied forward rate]\label{ch06:prop-fwd}
The fair (zero-PV) rate fixed today for lending $N$ from $t_1$ to $t_2$ is
\begin{equation}\label{ch06:eq-fwd}
  f(t_1,t_2)=\frac1\Delta\Bigl(\frac{Z(t_1)}{Z(t_2)}-1\Bigr).
\end{equation}
\end{proposition}
\begin{proof}
Lending $N$ at $t_1$ and receiving $N(1+f\Delta)$ at $t_2$ has present value
$-NZ(t_1)+Nf\Delta\,Z(t_2)+NZ(t_2)$. Setting it to zero gives $1+f\Delta=Z(t_1)/Z(t_2)$. If the market rate were
higher, lend forward and short the replicating zeros (buy the $t_1$ zero, sell the $t_2$ zero) for a riskless
profit, and vice versa.
\end{proof}
This is \cref{prop:forward} of \cref{ch:markets} with $F_{t_1t_2}=f\Delta$.

\begin{coursenote}[Erratum: slide 20]
The slide states the problem for a period $[t_1,t_2]$ but writes the equations with $t_0,t_1$:
``$\mathrm{PV}=-NZ(t_0)+Nf\Delta Z(t_1)+NZ(t_1)=0$'' and $f=\frac1\Delta\bigl(\frac{Z(t_0)}{Z(t_1)}-1\bigr)$; it also
says ``borrow $N$ at $t_0$''. The consistent statement is \eqref{ch06:eq-fwd}: the start of the accrual period
plays the role of $t_0$ on the slide, the end that of $t_1$.
\end{coursenote}

\paragraph{FRAs (not discussed in class).} A forward rate agreement on notional $N$ and strike $K$ pays at
the end of the period $N(R-K)\Delta$, where $R$ is the rate fixed for $[t_1,t_2]$ (in the LIBOR world
usually settled at $t_1$ as $N(R-K)\Delta/(1+R\Delta)$, the same value). By \cref{ch06:prop-fwd} its value today
is $N(f-K)\Delta\,Z(t_2)$ and the fair strike is $K=f$. A swap is a strip of FRAs sharing one strike.

\subsection{SOFR futures}\label{ch06:sec-futures}
Short-rate futures are among the most liquid contracts in every currency \ts{24}{7}{30:32}. The US contracts
trade at the CME (slide 15) \ts{24}{7}{31:35}:
\begin{itemize}
  \item \textbf{3-month SOFR futures} settle on the \emph{compounded} SOFR over a period running from one
        \emph{IMM date} to the next (third Wednesday of March, June, September, December; the dates come from
        the old International Monetary Market in Chicago);
  \item \textbf{1-month SOFR futures} settle on the \emph{average} SOFR over a calendar month.
\end{itemize}
Prices are quoted as $F=100-R$, with $R$ the settlement rate in percent: a 5\% rate gives 95, a 6\% rate 94
\ts{24}{7}{22:58}. The convention makes the contract behave like a bond: when rates fall (the market ``rallies''),
the price rises \ts{24}{7}{31:35}. Futures prices estimate forward rates up to a small \emph{convexity adjustment}
(slide 15):
\begin{equation}\label{ch06:eq-futures}
  f(t_1,t_2) = 100 - F - \mathrm{Cvx},\qquad \mathrm{Cvx}>0 .
\end{equation}

\paragraph{Origin of the adjustment} \ts{24}{7}{32:39}. Futures are margined daily: a contract entered at no
cost (apart from an initial margin) pays the holder the daily price change. A long position gains when the price
rises, i.e.\ when rates fall, and the gain can then be reinvested only at the lower rate; when rates rise, the
losses must be funded at the higher rate. This adverse correlation between the margin flows and the rate at
which they are carried makes a long futures position slightly worse than a long FRA, so the futures price is
lower and the futures rate $100-F$ higher than the forward rate. The adjustment can be estimated from a model, or
simply fitted together with swaps and futures prices.

\begin{remark}[Size of the adjustment (not discussed in class)]
Daily settlement makes the futures rate the expectation of $R$ under the risk-neutral (bank-account) measure,
$R^{\mathrm{fut}}=\E^{Q}[R]$, while the forward rate is its expectation under the $t_2$-forward measure, whose
density with respect to $Q$ is $D/\E^Q[D]$ with $D=\exp(-\int_0^{t_2}r_s\dd s)$ (change of numeraire,
\cref{ch:bs}). Hence
\[
  R^{\mathrm{fut}}-f=\E^Q[R]-\frac{\E^Q[DR]}{\E^Q[D]}=-\frac{\Cov^Q(D,R)}{\E^Q[D]}>0,
\]
because high rates come with low discount factors. In the Ho--Lee model with normal short-rate volatility
$\sigma$, $\mathrm{Cvx}\approx\tfrac12\sigma^2t_1t_2$ (continuous compounding; e.g.\ Hull's
\emph{Options, Futures, and Other Derivatives}). With $\sigma=1\%$ per year: 0.6\,bp for a contract one year
out, 13\,bp at five years, 51\,bp at ten years. It is negligible at the front and essential at the back, which
is why long-dated curves use swaps rather than futures.
\end{remark}

\subsection{Interest-rate swaps}
\begin{definition}[Vanilla interest-rate swap]
An agreement between two counterparties to exchange periodic cash flows on a notional $N$ until maturity
(slide 16) \ts{24}{7}{33:47}: the \emph{fixed leg} pays a rate $K$ set on the trade date, the \emph{floating leg}
pays a floating index observed periodically (today compounded SOFR). The \emph{payer} pays fixed and receives
floating; the \emph{receiver} does the opposite. Notionals are not exchanged.
\end{definition}
Swaps are the most liquid rates derivative. They convert fixed-rate into floating-rate exposure and vice
versa, are off balance sheet, cost little up front, and can be customised to any schedule a counterparty will
accept \ts{24}{7}{33:47}.

\begin{figure}[ht]
\centering
\begin{tikzpicture}[x=2.3cm,y=0.55cm,>=Stealth]
  \draw[->] (0,0) -- (5.55,0) node[right]{\small $t$ (years)};
  \foreach \k in {1,...,5} {
    \draw[->,thick,intcol] (\k,0) -- (\k,2.4);
    \node[below,font=\scriptsize] at (\k,-0.1) {$t_{\k}$};
  }
  \node[font=\scriptsize,intcol,anchor=south] at (5,2.4) {$NK\Delta_i$};
  \foreach \k/\h in {0.25/1.3,0.5/1.5,0.75/1.6,1/1.8,1.25/1.9,1.5/2.0,1.75/2.1,2/2.15,2.25/2.2,2.5/2.25,2.75/2.3,3/2.3,
                     3.25/2.35,3.5/2.4,3.75/2.45,4/2.5,4.25/2.55,4.5/2.6,4.75/2.65,5/2.7}
    \draw[->,thick,thmcol] (\k,0) -- (\k,-\h);
  \node[below left,font=\scriptsize] at (0,-0.1) {$t_0$};
  \node[font=\scriptsize,thmcol,anchor=north] at (5,-2.8) {$NR_j\Delta_j$};
  \node[font=\small,intcol,anchor=west] at (0.05,2.9) {fixed leg: received};
  \node[font=\small,thmcol,anchor=west] at (0.05,-3.3) {floating leg: paid, rate unknown today};
\end{tikzpicture}
\caption{Cash flows of a 5-year receiver swap with annual fixed payments $NK\Delta_i$ and a floating leg drawn
with quarterly payments $NR_j\Delta_j$ (USD SOFR swaps actually pay the compounded rate annually on both legs).
The floating amounts are unknown today; with an upward-sloping forward curve they are expected to grow.}
\label{ch06:fig-swapflows}
\end{figure}

\paragraph{Rolling out the cash flows.} The logic is simple but error-prone (slide 17) \ts{24}{7}{35:54}:
today $\to$ trade date; $+$ spot rule (typically two business days) $\to$ effective date; $+$ tenor $\to$
unadjusted maturity; step back by the payment frequency $\to$ unadjusted roll dates; $+$ business-day adjustment
$\to$ accrual dates; $+$ reset rule $\to$ rate-fixing dates; $+$ payment rule $\to$ payment dates. Every step
depends on holiday calendars, which change (Juneteenth, a state funeral, new national holidays), so models must
treat them as data. The standard calendar source lists 294 distinct dollar calendars (every exchange and business
centre has its own) \ts{24}{7}{37:00}. With the wrong calendar you misprice, and by the winner's curse you get
the trade only when the mispricing is against you.

% ------------------------------------------------------------------
\section{Swap valuation}\label{ch06:sec-swapval}
Let the fixed leg pay at $t_1<\dots<t_n$ with accrual fractions $\Delta_i$ and the floating leg at
$s_1<\dots<s_m$ with fractions $\Delta_j$, both starting at $t_0=s_0$.

\paragraph{Fixed leg} \ts{24}{7}{37:00}. Discount each known payment (slide 18):
\begin{equation}\label{ch06:eq-fixed}
  \mathrm{PV}_{\mathrm{fix}}=\sum_{i=1}^{n}K\Delta_iN\,Z(t_i)=NK\,A,\qquad
  A=\sum_{i=1}^{n}\Delta_iZ(t_i),
\end{equation}
where $A$ is the \emph{annuity} (or level, or PV01 of the fixed leg per unit rate).

\paragraph{Floating leg} \ts{24}{7}{38:04}. Replace unknown rates by forwards (slide 19),
$\mathrm{PV}_{\mathrm{float}}=\sum_{j}f_j\Delta_jN\,Z(s_j)$. This needs two curves: an \emph{index} (projection)
curve for the forwards $f_j$ and a \emph{discount} curve for the $Z(s_j)$. For vanilla SOFR swaps both are the
SOFR curve, so there is only one curve. For compounded SOFR the replication is exact:

\begin{proposition}[Floating leg]\label{ch06:prop-float}
Assume the discount curve is the SOFR curve (\cref{ch06:prop-funding} with SOFR as the funding rate) and that
each floating period pays $NR_j\Delta_j$ at $s_j$, where $R_j$ is SOFR compounded over $[s_{j-1},s_j]$ as in
\eqref{ch06:eq-sofr}. Then the value of period $j$ is $N\bigl[Z(s_{j-1})-Z(s_j)\bigr]=Nf_j\Delta_jZ(s_j)$ with
$f_j$ the forward rate \eqref{ch06:eq-fwd}, and
\begin{equation}\label{ch06:eq-float}
  \mathrm{PV}_{\mathrm{float}}=N\bigl[Z(s_0)-Z(s_m)\bigr].
\end{equation}
\end{proposition}
\begin{proof}
Invest $N$ at $s_{j-1}$ in an overnight deposit and roll principal and interest every day. At $s_j$ it has become
$N\prod_i(1+r_id_i/360)=N(1+R_j\Delta_j)$, whatever the fixings. Hence the payment equals ``value of the rolled
deposit at $s_j$'' minus $N$. The first part costs $N$ at $s_{j-1}$, worth $NZ(s_{j-1})$ today; the second is a
fixed amount worth $NZ(s_j)$. The equality with $Nf_j\Delta_jZ(s_j)$ is \eqref{ch06:eq-fwd}. Summing over $j$,
the sum telescopes.
\end{proof}

\begin{intuition}[A floating-rate note is worth par]
Adding the notional at the end, $\mathrm{PV}_{\mathrm{float}}+NZ(s_m)=NZ(s_0)$: a note paying the funding rate
is worth par at each reset, because its coupons are exactly what the money would have earned anyway. So a
receiver swap is ``long a fixed-coupon bond, short a floating-rate note'': a bond position financed at the
floating rate. This is also why an arithmetic \emph{average} of SOFR needs a correction: it is not what a rolled
deposit earns, so the static replication fails and a small, volatility-dependent adjustment appears
\ts{24}{7}{21:49}.
\end{intuition}

\begin{corollary}[Par swap rate]\label{ch06:cor-par}
A swap traded at zero PV (the market standard) has the \emph{par swap rate} \ts{24}{7}{40:06}
\begin{equation}\label{ch06:eq-par}
  S=\frac{\sum_{j=1}^mf_j\Delta_jZ(s_j)}{\sum_{i=1}^n\Delta_iZ(t_i)}
   =\frac{Z(t_0)-Z(t_n)}{A}.
\end{equation}
The first form (slide 21) is a weighted average of forward rates; the second holds under the single-curve
assumptions of \cref{ch06:prop-float}. For a spot-starting swap $Z(t_0)\approx1$.
\end{corollary}
Par swap rates are quoted on every broker screen, and they are what the curve must reproduce. With annual
payments, $\Delta_i=1$ and $Z(t_0)=1$, \eqref{ch06:eq-par} is exactly the par-bond equation used for
bootstrapping in \cref{sec:bonds}: the swap rate is the coupon of a par bond priced off the SOFR curve.

\paragraph{Mark-to-market and DV01.} After inception, let $S_t$ and $A_t$ be the par rate and annuity of the
remaining schedule. Then
\begin{equation}\label{ch06:eq-mtm}
  V_{\mathrm{payer}}=N\bigl[(Z(t_0)-Z(t_n))-KA\bigr]=N\,A_t\,(S_t-K),\qquad V_{\mathrm{receiver}}=-V_{\mathrm{payer}},
\end{equation}
and to first order (holding $A_t$ fixed) a 1\,bp rise in the par rate changes the payer's value by
$\mathrm{DV01}\approx N A_t\times10^{-4}$.

\begin{example}[Valuing swaps (not from class)]
With the curve of \cref{ch06:ex-boot} below: a 5-year receiver swap on \$100M struck at $K=4.00\%$, when the
5-year par rate is $3.40\%$ and $A_5=4.5178$, is worth $100\text{M}\times0.006\times4.5178=\$2.71$M. A new
10-year par swap on \$100M has $A_{10}=8.3195$ and a DV01 of about \$83{,}200 per bp.
\end{example}

\begin{finance}[Why swaps were invented, and why they trade at par]
A comment from the audience \ts{24}{7}{41:14}: in the late 1980s swaps boomed because highly rated companies
could borrow long-term fixed cheaply, while weaker companies could only borrow short-term floating. Each borrowed
where it had a comparative advantage and they swapped the payments, both ending up with cheaper funding of the
type they wanted \ts{24}{7}{42:14} (\cref{ch06:ex-compadv}). Today swaps are traded at (or within a bid--offer of
0.1--0.2\,bp around) zero PV, i.e.\ at the par rate. Off-market swaps appear inside structured transactions,
where the upfront value is accounted for elsewhere or is where the P\&L is taken \ts{24}{7}{41:14}.
\end{finance}

\begin{remark}[Two curves in the LIBOR era]
With LIBOR the index (unsecured, bank credit) and the discount rate (collateral, OIS) differed, so forwards came
from a LIBOR curve and discount factors from an OIS curve, and the telescoping \eqref{ch06:eq-float} failed
(``multi-curve'' pricing). The SOFR transition has largely brought vanilla USD swaps back to a single curve.
\end{remark}

% ------------------------------------------------------------------
\section{Building the yield curve}\label{ch06:sec-curve}
The yield curve is the core valuation model of a rates desk: a function $t\mapsto Z(t)$ that feeds all other
pricing \ts{24}{7}{42:14}. It is built in three steps (slide 22) \ts{24}{7}{43:15}:
\begin{enumerate}[1.]
  \item parameterise it by \emph{knot points} $\{(t_k,Z(t_k))\}_{k=1}^{N}$;
  \item choose an \emph{interpolation} to get $Z(t)$ for $t\notin\{t_k\}$ (say a 4.5-year discount factor from
        annual knots);
  \item \emph{calibrate} the $Z(t_k)$ to market quotes.
\end{enumerate}

\subsection{Knot points}
The simplest and most common choice: the maturities of the calibration instruments (1, 2, 3, 5, 7, 10-year swaps
give knots at 1, 2, 3, 5, 7, 10 years) \ts{24}{7}{43:15}. Other choices are possible (e.g.\ FOMC meeting dates,
\cref{ch06:sec-fomc}), but the knots must be chosen together with the interpolation and the instruments: bad choices
give unstable hedges \ts{24}{7}{44:16}. Asked why the knots are not chosen to minimise interpolation error (e.g.\
Chebyshev nodes), the lecturer answers that one wants a one-to-one link between knots and calibration instruments:
with knots at 5, 10, 15, \dots\ years and instruments concentrated at the short end, the curve would be
over-determined in one region and under-determined in another. Moreover the hedges come out in terms of the input
instruments, which must be liquid and tradeable \ts{24}{7}{45:19}.

\subsection{Interpolation: local versus global}
\begin{definition}[CDF interpolation]\label{ch06:def-cdf}
\emph{Constant daily forward} (CDF) interpolation assumes that the daily forward rate
$f_d$, defined by $Z(d+1)=Z(d)/(1+f_d/360)$, is constant between consecutive knots \ts{24}{7}{44:16}.
\end{definition}

\begin{proposition}[CDF $=$ log-linear discount factors]\label{ch06:prop-cdf}
Under CDF interpolation, for $t_k\le t\le t_{k+1}$,
\begin{equation}\label{ch06:eq-cdf}
  \ln Z(t)=\frac{t_{k+1}-t}{t_{k+1}-t_k}\ln Z(t_k)+\frac{t-t_k}{t_{k+1}-t_k}\ln Z(t_{k+1}),
\end{equation}
equivalently the product $t\,r_z(t)$ of time and continuously compounded zero rate $r_z(t)=-\ln Z(t)/t$ is
linear between knots (slide 23). The interpolation is \emph{local}: $Z(t)$ depends only on the two neighbouring
knots.
\end{proposition}
\begin{proof}
$\ln Z(d)=\ln Z(t_k)-\sum_{d'=t_k}^{d-1}\ln(1+f_{d'}/360)$ is linear in the number of days when $f_{d'}$ is
constant; matching the two endpoints gives \eqref{ch06:eq-cdf}, and $t\,r_z(t)=-\ln Z(t)$. In continuous time
this is a piecewise-constant instantaneous forward $f(t)=-\dd\ln Z/\dd t$.
\end{proof}
CDF is very fast, local, and produces the characteristic staircase forward curve (\cref{ch06:fig-fwd}): the
simplest, least smooth method \ts{24}{7}{46:21}.

\begin{coursenote}[Erratum: figure on slide 23]
The chart next to the CDF bullets is titled ``Yield'', but a step function is what CDF produces for the
\emph{forward} rate. The zero yield $r_z(t)=\frac1t\int_0^tf(s)\dd s$ is continuous (of the form $a+b/t$ on each
interval). Read the vertical axis as the daily forward rate.
\end{coursenote}

\begin{definition}[Cubic spline]
Fit on each interval $[t_k,t_{k+1}]$ a cubic
\[
  x(t)=\sum_{j=0}^{3}\alpha_{kj}\,(t-t_k)^j ,
\]
where $x$ is the interpolated quantity (typically the zero rate), and impose continuity of $x,x',x''$ at the
interior knots (slide 24) \ts{24}{7}{46:21}. Two end conditions (e.g.\ $x''=0$ at both ends, a \emph{natural}
spline) close the system: $4(N-1)$ coefficients for $2(N-1)$ interpolation conditions, $2(N-2)$ smoothness
conditions and 2 end conditions.
\end{definition}

\begin{coursenote}[Erratum: slide 24]
The slide writes $x_i(t)=\sum_{i=0}^{3}\alpha_{ij}(t-t_i)^i$, using $i$ both for the interval and for the
summation. The correct form is the one above: interval index $k$, power index $j$.
\end{coursenote}

The spline is reasonably fast and very smooth, but \emph{global}: the curve near an 8-year point depends on all
knots, because moving a distant knot changes the shape everywhere through the smoothness constraints
\ts{24}{7}{47:24}. The lecturer's advice to future rates quants is to keep a laminated card reading ``fast, local
hedges \emph{or} smooth, global hedges'': one can be anywhere in between, but not have both. He has argued about
this with traders for 30 years. Traders want an 8-year trade to be hedged with 7- and 10-year swaps only, not with
5- or 15-year ones. Among smooth methods the cubic spline is the worst in this respect, precisely because it is the
smoothest. Tension splines and B-splines sit in between, and much interpolation technology from computer graphics
could be used \ts{24}{7}{46:21} (slide 29).

\begin{intuition}[The spline as an elastic beam]
A natural cubic spline minimises the bending energy $\int x''(t)^2\dd t$ among all interpolants: it is the
draftsman's elastic ruler pinned at the knots. Moving one pin bends the whole ruler. For equally spaced knots the
disturbance decays by the factor $2-\sqrt3\approx0.27$ per knot \emph{with alternating sign}, the root of
$\lambda^2+4\lambda+1=0$ in the tridiagonal system for the second derivatives. This is exactly the ringing
visible in the spline hedges of \cref{ch06:fig-ladder}. CDF corresponds to a ``string'' of independent segments
in $\ln Z$: no stiffness, no propagation, no smoothness.
\end{intuition}

\subsection{Calibration}
For each calibration instrument write the pricing error as a function of the curve (slide 25)
\ts{24}{7}{48:27}:
\[
  E_i(Z)=P_i(Z)-P_i^{\mathrm{obs}},\qquad i=1,\dots,N .
\]
\begin{itemize}
  \item With $N$ knots and $N$ instruments solve $E_i=0$ exactly with any root finder. When the interpolation is
        local and instruments are sorted by maturity, each new instrument introduces a single new unknown and the
        system can be \emph{bootstrapped} one equation at a time (as in \cref{sec:bonds}).
  \item Otherwise minimise $E=\sum_iE_i(Z)^2$, possibly with extra terms, e.g.\ smoothness penalties on an
        under-specified curve, or accept any curve that prices every instrument within its bid--offer (which is
        indistinguishable from an exact fit) \ts{24}{7}{49:28}.
\end{itemize}

\begin{example}[Bootstrapping a small swap curve (not from class)]\label{ch06:ex-boot}
Annual-pay par swap rates ($\Delta_i=1$, single curve, day counts ignored): 1Y 4.00\%, 2Y 3.60\%, 3Y 3.45\%,
5Y 3.40\%, 7Y 3.45\%, 10Y 3.55\%. For $T\le3$ every payment date is a knot and \eqref{ch06:eq-par} gives the
recursion $Z_T=(1-S_T\sum_{k<T}Z_k)/(1+S_T)$. For 5Y, CDF gives $Z_4=\sqrt{Z_3Z_5}$ and $Z_5$ solves the scalar
equation $1-Z_5=S_5\bigl(Z_1+Z_2+Z_3+\sqrt{Z_3Z_5}+Z_5\bigr)$; similarly for 7Y and 10Y. Result:
\begin{center}\small
\begin{tabular}{@{}lcccccc@{}}
\toprule
knot $t_k$ (years) & 1 & 2 & 3 & 5 & 7 & 10\\
\midrule
par rate (\%) & 4.00 & 3.60 & 3.45 & 3.40 & 3.45 & 3.55\\
$Z(t_k)$ & 0.96154 & 0.93184 & 0.90351 & 0.84640 & 0.78874 & 0.70466\\
zero rate $r_z$ (\%, cont.) & 3.922 & 3.530 & 3.382 & 3.335 & 3.390 & 3.500\\
forward on $(t_{k-1},t_k]$ (\%) & 3.922 & 3.138 & 3.088 & 3.265 & 3.528 & 3.757\\
\bottomrule
\end{tabular}
\end{center}
A natural cubic spline on zero rates, calibrated to the same six quotes, reprices them exactly too, but it
implies a 4-year par rate of 3.405\% against 3.419\% for CDF. For a maturity that is not quoted, the choice of
interpolation moves the price by more than a basis point, about ten times the bid--offer.
\end{example}

\begin{figure}[ht]
\centering
\begin{tikzpicture}
\begin{axis}[width=0.8\linewidth,height=6.2cm,xlabel={maturity (years)},ylabel={rate (\%)},
  xmin=0,xmax=10.3,ymin=2.9,ymax=4.05,xtick={0,1,2,3,5,7,10},
  legend pos=south east,legend style={font=\small,cells={anchor=west}},grid=major,grid style={black!10}]
  \addplot[very thick,thmcol,const plot] coordinates {(0,3.922) (1,3.1375) (2,3.0875) (3,3.2649) (5,3.5277) (7,3.7574) (10,3.7574)};
  \addlegendentry{forward, CDF}
  \addplot[thick,defcol,smooth] coordinates {(1.00,3.471) (1.25,3.260) (1.50,3.093) (1.75,2.993) (2.00,2.981) (2.25,3.031)
    (2.50,3.093) (2.75,3.146) (3.00,3.174) (3.25,3.186) (3.50,3.204) (3.75,3.228) (4.00,3.257) (4.25,3.289) (4.50,3.323)
    (4.75,3.358) (5.00,3.394) (5.25,3.429) (5.50,3.465) (5.75,3.500) (6.00,3.534) (6.25,3.566) (6.50,3.595) (6.75,3.621)
    (7.00,3.643) (7.25,3.662) (7.50,3.682) (7.75,3.701) (8.00,3.721) (8.25,3.740) (8.50,3.760) (8.75,3.779) (9.00,3.798)
    (9.25,3.817) (9.50,3.836) (9.75,3.855) (10.00,3.874)};
  \addlegendentry{forward, cubic spline on zeros}
  \addplot[thick,black!55,dashed,mark=none] coordinates {(1,3.9221) (2,3.5298) (3,3.3824) (4,3.353) (5,3.3354) (6,3.3674)
    (7,3.3903) (8,3.4362) (9,3.4719) (10,3.5004)};
  \addlegendentry{zero rate, CDF}
  \addplot[only marks,mark=*,mark size=1.6pt,fincol] coordinates {(1,4.00) (2,3.60) (3,3.45) (5,3.40) (7,3.45) (10,3.55)};
  \addlegendentry{par swap quotes}
\end{axis}
\end{tikzpicture}
\caption{Curves implied by the quotes of \cref{ch06:ex-boot}. CDF gives a staircase of forward rates with jumps at
the knots; the spline (drawn from the first knot; below 1 year both methods coincide) gives a smooth forward curve
that dips below the CDF forwards around 2 years. Both reprice the six quotes exactly. Continuously compounded
rates; the 1-year zero rate is $\ln1.04=3.922\%$.}
\label{ch06:fig-fwd}
\end{figure}

\subsection{Central-bank meeting curves, turns and data-driven curves}\label{ch06:sec-fomc}
\paragraph{FOMC curves.} The history of daily SOFR (slide 26) is not smooth at all: it looks like a staircase
\ts{24}{7}{49:28}. It is flat between FOMC meetings and jumps when the Fed changes its policy rate: near zero
after March 2020, a sequence of steps up in 2022--23, a plateau around 5.3\% until the 50\,bp cut of September 2024
(``the chart needs a 50\,bp drop at the right edge'') \ts{24}{7}{50:32}. The spike at the left edge is a repo
squeeze related to tax-payment flows (September 2019). This suggests CDF interpolation with knots at the
meeting dates: daily rates are constant between meetings, which is now almost universal in the industry
\ts{24}{7}{50:32}. Two complications (slides 27--28):
\begin{itemize}
  \item there are 8 FOMC meetings a year but only 4 quarterly futures, so the problem is under-determined. One
        adds 1-month futures or smoothness constraints, or uses \emph{FOMC swaps} (meeting-to-meeting OIS), which
        now trade and directly price the rate expected after a given meeting \ts{24}{7}{51:34};
  \item the meeting calendar is known only about two years ahead, and instruments are sparse beyond the front end.
        Hence a \emph{hybrid} curve: CDF on meeting dates at the front, spline at the back, which works well in
        practice \ts{24}{7}{50:32}.
\end{itemize}

\paragraph{Turns} \ts{24}{7}{52:37}. Banks report balance sheets and return on assets at quarter- and year-ends and
want small balance sheets on those dates, so borrowing across year-end costs a premium: a spike in the forward
curve every year. Curves include these \emph{turns} explicitly. Some hedge funds trade pairs of swaps ending just
before and just after year-end, isolating the turn and exploiting whoever models it badly; in size this is real
money \ts{24}{7}{53:40}.

\paragraph{Data-driven curves} \ts{24}{7}{53:40}. A USD curve has about 50 inputs of very different liquidity;
perhaps half a dozen are super-liquid. One can estimate historical factors (curve shapes) by PCA
(\cref{ch:pca}), drive the first few factors from the most liquid instruments and build the curve from them. This
works well for bond curves (slide 29).

% ------------------------------------------------------------------
\section{Risk: PV01, delta ladders and gamma}\label{ch06:sec-risk}
On a desk, ``risk'' first means: how much do I make or lose if rates move \ts{24}{7}{54:42}?

\begin{definition}[PV01 and delta ladder]
The \emph{PV01} (or DV01) of a position is its sensitivity to a 1\,bp move of the curve inputs, computed by a
two-sided finite difference (slide 30)
\[
  \frac{\partial\mathrm{PV}}{\partial r}\approx\frac{\mathrm{PV}(r+\epsilon)-\mathrm{PV}(r-\epsilon)}{2\epsilon},
  \qquad \epsilon=1\,\text{bp}.
\]
Bumping each curve input $x_k$ separately gives the vector of partial PV01s, the \emph{delta ladder} (partial
deltas, bucketed risk, delta strip) $\delta_k=\partial\mathrm{PV}/\partial x_k$.
\end{definition}
The ladder is what the desk hedges and what risk limits are set on. It is computed by brute force (bump the input,
rebuild the curve, revalue) or analytically, but the concept is the same \ts{24}{7}{54:42}. If the inputs are the
calibration quotes, the ladder translates directly into hedges:

\begin{proposition}[Hedging with the calibration instruments]\label{ch06:prop-hedge}
Let the curve be calibrated exactly to par swaps with quotes $q=(q_1,\dots,q_N)$ and annuities $A_k$. A payer
swap on instrument $k$, struck at today's quote, has $\partial\mathrm{PV}_k/\partial q_j=A_k\,\delta_{kj}$ per unit
notional. Hence a portfolio with ladder $\delta_k=\partial V/\partial q_k$ is first-order neutral to every quote
after adding payer swaps with notionals
\[
  x_k=-\frac{\delta_k}{A_k},\qquad k=1,\dots,N .
\]
\end{proposition}
\begin{proof}
Because the calibrated curve reprices instrument $k$ at par rate $q_k$ for every $q$, \eqref{ch06:eq-mtm} gives
$\mathrm{PV}_k(q)=(q_k-K)A_k(q)$. Differentiating, $\partial_{q_j}\mathrm{PV}_k=\delta_{kj}A_k+(q_k-K)\partial_{q_j}A_k$,
and the second term vanishes at $K=q_k$. The total first-order sensitivity to $q_j$ is $\delta_j+x_jA_j=0$.
\end{proof}

\begin{example}[An 8-year swap hedged with 7- and 10-year swaps (not from class)]
Take the curve of \cref{ch06:ex-boot} and a \$100M 8-year receiver swap at its par rate (3.492\% with CDF). Its
parallel PV01 is $-\$68{,}835$ per bp, equal to $-NA_8\times10^{-4}$ with $A_8=6.883$. The ladder with respect to
the six quotes depends on the interpolation (\cref{ch06:fig-ladder}):
\begin{itemize}
  \item \textbf{CDF}: $-\$39{,}982$ on the 7Y quote, $-\$28{,}840$ on the 10Y, a few dollars elsewhere. The hedge is
        to pay fixed on \$65.3M 7Y and \$34.7M 10Y: exactly what traders want;
  \item \textbf{cubic spline}: $-533$, $+3{,}244$, $-6{,}969$, $+19{,}314$, $-68{,}766$, $-15{,}138$ dollars per bp
        on 1Y, \dots, 10Y. The hedge requires positions of alternating sign in every maturity, including
        \$112M of 7Y and $-\$43$M of 5Y: a ``smeared hedge'' (slide 32).
\end{itemize}
In both cases the ladder sums to the parallel PV01.
\end{example}

\begin{figure}[ht]
\centering
\begin{tikzpicture}
\begin{axis}[width=0.85\linewidth,height=6cm,ybar,bar width=9pt,
  symbolic x coords={1Y,2Y,3Y,5Y,7Y,10Y},xtick=data,
  ylabel={hedge: payer notional (\$M)},xlabel={calibration swap},
  ymin=-55,ymax=125,legend pos=north west,legend style={font=\small},
  ymajorgrids,grid style={black!10},
  nodes near coords,every node near coord/.append style={font=\tiny},
  point meta=explicit symbolic]
  \addplot[fill=thmcol!70,draw=thmcol] coordinates {(1Y,0) [0] (2Y,0) [0] (3Y,0) [0] (5Y,0) [0] (7Y,65.3) [65.3] (10Y,34.7) [34.7]};
  \addlegendentry{CDF (local)}
  \addplot[fill=defcol!60,draw=defcol] coordinates {(1Y,5.5) [5.5] (2Y,-17.1) [$-$17.1] (3Y,24.9) [24.9] (5Y,-42.8) [$-$42.8] (7Y,112.3) [112.3] (10Y,18.2) [18.2]};
  \addlegendentry{cubic spline (global)}
\end{axis}
\end{tikzpicture}
\caption{Hedge of a \$100M 8-year receiver swap with the six calibration swaps, $x_k=-\delta_k/A_k$ from
\cref{ch06:prop-hedge}, for the two interpolations of \cref{ch06:ex-boot}. Both hedges total about \$100M, but
the spline one rings with alternating signs through the whole curve (``fast, local'' versus ``smooth, global''
\ts{24}{7}{47:24}).}
\label{ch06:fig-ladder}
\end{figure}

\paragraph{Gamma} \ts{24}{7}{55:47}. Second-order risk (gamma, or convexity) matters more or less depending on the
book. Slide 31 defines it as the change in delta for a 1\,bp move:
\[
  \Gamma=\frac{\mathrm{PV}(r+\epsilon)-2\,\mathrm{PV}(r)+\mathrm{PV}(r-\epsilon)}{\epsilon}
  =\epsilon\,\frac{\partial^2\mathrm{PV}}{\partial r^2}+O(\epsilon^3).
\]
Note the single $\epsilon$ in the denominator: this is $\Delta(r+\tfrac\epsilon2)-\Delta(r-\tfrac\epsilon2)$ with
$\Delta=\partial\mathrm{PV}/\partial r$, i.e.\ the second derivative times the bump, not the second derivative
itself. The full matrix of cross gammas $\partial^2\mathrm{PV}/\partial x_j\partial x_k$ needs bumps in pairs, of
order $N^2$ revaluations: $2{,}500$ for a 50-input curve, too expensive. Desks use a parallel gamma, a factor (PCA)
gamma, or scenario analysis instead \ts{24}{7}{55:47}.

\paragraph{Practical issues} \ts{24}{7}{56:49} (slide 32).
\begin{itemize}
  \item \emph{Bump size}: 1\,bp is a reporting convention. Monte Carlo pricers need larger bumps, or the noise
        swamps the difference (the ``noisy delta'' problem of \cref{ch:markets}); PDE pricers show grid effects
        for large bumps; too large a bump creates unnatural curve shapes.
  \item \emph{Sum of partials vs parallel shift}: by the chain rule they agree to first order, as in the example;
        the discrepancy comes from cross-gamma terms at finite $\epsilon$. Sequential bumps are one remedy.
  \item \emph{Interpolation}: the ladder of a position is a property of position \emph{and} curve model. Local
        interpolation concentrates it near the maturity, smooth interpolation smears it, which is a source of
        endless arguments \ts{24}{7}{56:49}.
\end{itemize}

% ------------------------------------------------------------------
\section{P\&L attribution}\label{ch06:sec-pnl}
Every morning a desk starts with positions and a market, and every evening it has made or lost money: why
\ts{24}{7}{56:49}? Skill, market direction, client flow earning the bid--offer? The attribution is also one of the
best checks of the model, and in some jurisdictions a regulatory requirement. Total P\&L is known exactly (slide 33),
$\mathrm{PnL}(t_0,t_1)=\mathrm{PV}(t_1)-\mathrm{PV}(t_0)$. It is split into three parts (slide 34)
\ts{24}{7}{57:58}:
\begin{itemize}
  \item \textbf{Portfolio change}: new trades, unwinds, novations, resets, exercises. Conceptually simple,
        mechanically the hardest (booking systems, operational flows).
  \item \textbf{Time}, or \emph{theta}: the change in value if the market ``does not move''.
  \item \textbf{Market moves}: to first order $\sum_i\frac{\partial\mathrm{PV}}{\partial x_i}\Delta x_i$ (the
        linear term on slide 34), plus $\frac12\sum_{i,j}\Gamma_{ij}\Delta x_i\Delta x_j$ and higher terms.
\end{itemize}
What remains is the \emph{unexplained} P\&L. A large or systematic unexplained part signals missing risk factors or
a bad model. The order in which the three steps are applied matters (the decomposition is path-dependent), so
conventions must be fixed and documented.

\paragraph{What does ``the market did not move'' mean?} \ts{24}{7}{58:59}. The curve is a function of maturity,
and when the calendar moves one must decide what is held fixed. There is more than one defensible answer
(not detailed in class):
\begin{itemize}
  \item \emph{Forwards realised}: yesterday's forward curve for fixed calendar dates is today's curve,
        $Z_{t_1}(T)=Z_{t_0}(T)/Z_{t_0}(t_1)$. Every funded position then earns exactly the funding rate: theta
        is pure carry.
  \item \emph{Constant curve by tenor}: the zero (or par) rate for each time-to-maturity is unchanged. On an
        upward-sloping curve a bond ``rolls down'' to lower yields and earns carry \emph{plus} roll-down
        (\cref{ch06:ex-theta}).
\end{itemize}
For a physicist: the first convention is Lagrangian (follow each cash flow), the second Eulerian (sit at a fixed
point of the maturity axis). Whatever the choice, the desk and its controllers must use the same one.

\begin{coursenote}[The spoken P\&L example]
In class the market term is illustrated with a long position of \$1B of bonds, a PV01 of ``\$10 million'' and
gains of \$50M for a 5\,bp move \ts{24}{7}{58:59}. The numbers were improvised: \$1B of 10-year bonds has a PV01
of roughly \$0.8--0.9M, so 5\,bp is worth about \$4M, and a long bond position gains when rates \emph{fall}.
\end{coursenote}

% ------------------------------------------------------------------
\section{Clearing, margin and the CCP basis}\label{ch06:sec-clearing}
A swap starts near zero PV, but as markets move its value drifts away from zero and becomes a credit exposure: if
a counterparty that owes me \$10M on a swap defaults, I lose it (slide 35) \ts{24}{7}{1:00:00}. Lehman, a huge
swap dealer, made this concrete in 2008: its counterparties had to replace hundreds of billions of hedges at once,
bid--offers widened about twentyfold for a month or two, and the one large dealer still trading normally captured
the flow \ts{24}{7}{1:01:03}. After the crisis regulators made \emph{central clearing} mandatory for vanilla
swaps: once two parties agree a trade, both assign their side to a central counterparty (CCP: LCH in London,
CME, JSCC in Japan) and post margin to it, so they no longer face each other's credit \ts{24}{7}{1:01:03}.

\begin{finance}[Fewer, larger risks]
Clearing replaces many small bilateral risks by concentrated exposure to a few enormous, deeply interconnected
CCPs. A failure of one of them would be a systemic catastrophe, which is why members post margin and default-fund
contributions and CCPs are heavily capitalised \ts{24}{7}{1:02:09}.
\end{finance}

\paragraph{The CCP basis} \ts{24}{7}{1:02:09}. A dollar owed by LCH is not the same as a dollar owed by CME:
positions at different CCPs cannot be offset operationally or for collateral. From about 2017 the fixed rate for
paying fixed at CME exceeded that at LCH. On a broker screen for 26 June 2017 (slide 36; source Tradition) the
CME--LCH spread grew with maturity, from 0.10\,bp at 1 year through 1.05 (5Y), 2.45 (10Y) and 3.40 (30Y) to
3.55\,bp at 50 years, in a market with a bid--offer of about 0.2\,bp \ts{24}{7}{1:03:10}. It persisted for years
(recently it has changed sign). Why not receive at CME, pay at LCH and retire \ts{24}{7}{1:03:10}?
\begin{itemize}
  \item \textbf{Double margin}: the two offsetting trades sit at different CCPs, so initial margin (1--2\% of
        notional, quoted in class) is posted twice. Its funding cost is the main driver
        \ts{24}{7}{1:04:16}.
  \item \textbf{Netting}: CME also clears futures, so a dealer long futures and short swaps there gets partial
        margin offsets that the arbitrage would destroy \ts{24}{7}{1:04:16}.
  \item \textbf{Preferences}: large US asset managers, one-way payers of fixed, preferred CME, a US-jurisdiction
        CCP not controlled by dealers, over LCH \ts{24}{7}{1:05:21}. (LCH Group is today majority-owned by the
        London Stock Exchange Group, with banks as minority shareholders; the perception quoted in class concerns
        its governance.)
\end{itemize}
The general lesson: margin and capital requirements let spreads persist between instruments that would otherwise
be arbitraged. To price correctly one must include how the trade is margined, how much capital it consumes, and
every other cost that used to be ignored. Estimating these costs and pushing them into the upfront price has been
one of the largest quant efforts of the last 15 years, made harder by the fact that many are portfolio effects
that must be allocated to individual trades \ts{24}{7}{1:06:21} (see \cref{ch:counterparty}).

% ------------------------------------------------------------------
\section{Electronic trading and price making}\label{ch06:sec-etrading}
Electronic trading is one of the hottest areas for quants \ts{24}{7}{1:07:23}. As a rule of thumb, equities are
about ten years ahead of FX, and FX ten years ahead of fixed income. Slide 37 (Flow Traders data, 2023) shows the
share of electronic trading: above 90\% for US cash equities and equity options, European ETFs and FX spot;
about 80\% for European government bonds and 65\% for US Treasuries; 35--45\% for investment-grade corporates;
22--27\% for high yield; about 12\% for munis and 10\% for emerging-market bonds. Small trades go electronic first:
\$1M of Treasuries does not move a market trading \$1T a day, whereas \$500M does and needs care
\ts{24}{7}{1:08:24}.

\paragraph{Participants} \ts{24}{7}{1:10:26}. \emph{Principal trading firms} (high-frequency shops) have no
customers, trade their own capital and are extremely fast; the lecturer met a physics PhD hired to optimise the
microwave links between Chicago and New York, and realised that banks will never beat them at latency
\ts{24}{7}{1:11:26}. \emph{Dealers} (banks) are fast (microseconds, not nanoseconds) and have customers, who pay
the bid--offer. \emph{Customers} want efficient execution, a view of the full market and confidence that they are
not being taken advantage of \ts{24}{7}{1:12:26}.

\paragraph{Price making} (slide 38) \ts{24}{7}{1:12:26}. Start from a mid price, from liquidity providers for
actively traded instruments or from a model for off-the-run ones, then \emph{skew} it according to:
\begin{itemize}
  \item the current position: be aggressive on risk-reducing trades, conservative on risk-increasing ones;
  \item the client tier: a valuable client relationship justifies a better price \ts{24}{7}{1:13:29};
  \item \emph{client toxicity}: the classic toxic client sends the same request to ten dealers at once. Whoever
        wins finds nine others hedging the same way, and the hedge has moved away (winner's curse)
        \ts{24}{7}{1:13:29};
  \item market volatility (widen when the market may move before the hedge is done) and short-term prediction
        models (market microstructure).
\end{itemize}
All of this requires tightly integrated systems. A credit check that answers after the client has gone is a lost
trade \ts{24}{7}{1:14:29}.

\paragraph{Fast model pricing (slide 39, not presented in class).} Model prices must be recomputed in microseconds.
One linearises the output around the last full valuation,
\[
  r(t)\approx r(t_0)+\sum_{i=1}^{n}\frac{\partial r}{\partial x_i}\,\Delta x_i ,
\]
splits the drivers into fast ones (fed through the linearisation in real time) and slow ones, and periodically
recomputes the full model. The slide omits the index on $\Delta x$; it must be $\Delta x_i$.

% ------------------------------------------------------------------
\section{Bonds revisited}\label{ch06:sec-bonds}
Slides 40--44 were not presented in class; they extend \cref{sec:bonds,sec:duration}.

\paragraph{Why bonds are harder than swaps.} Government bonds are the most liquid fixed-income instruments in many
markets. G20 government bonds are usually treated as default-free, while emerging-market bonds need a default model.
Despite their apparent simplicity bonds are harder than swaps: bonds are not fungible (each issue has its own
coupon, maturity, liquidity, repo specialness), whereas all swaps are generated from one curve (slide 40).

\paragraph{Price--yield with $f$ payments a year} (slide 41). Bonds trade and settle on price, but yield is the
better comparison across coupons and maturities. With $n$ remaining periods, coupon $C_i$ per period per unit
notional and $f$ coupons per year,
\[
  P_D(y)=\sum_{i=1}^{n}\frac{C_i}{(1+y/f)^{i}}+\frac{1}{(1+y/f)^{n}},
\]
which reduces to \eqref{eq:couponbond} for $f=1$. The slide's chart shows the familiar picture: premium when the
yield is below the coupon, discount when above.

\begin{proposition}[Modified and Macaulay duration]
Let $D_{\mathrm{mod}}=-\frac1P\frac{\partial P}{\partial y}$ and let $D_{\mathrm{Mac}}=\sum_i\frac{i}{f}\,
\frac{\mathrm{CF}_i(1+y/f)^{-i}}{P}$ be the PV-weighted average payment time in years. Then
$D_{\mathrm{Mac}}=D_{\mathrm{mod}}\,(1+y/f)$ (slide 42).
\end{proposition}
\begin{proof}
$\partial_y(1+y/f)^{-i}=-\frac{i}{f}(1+y/f)^{-i-1}$, so
$-P'(y)=(1+y/f)^{-1}\sum_i\frac if\mathrm{CF}_i(1+y/f)^{-i}=P\,D_{\mathrm{Mac}}/(1+y/f)$.
\end{proof}
(The slide writes the factor as $1+y/n$ with $n$ the compounding frequency, called $f$ on the previous slide.)
Duration is ``not generic but good to gain insight'': it assumes a flat move of a single yield, while swap risk is
a ladder.

\paragraph{Bonds against the swap curve: the z-spread} (slide 43). Swap and bond books hedge each other, so their
risks must be aggregated. Pricing a bond off the swap (SOFR) curve,
$P'_D=\sum_iC_iZ(t_i)+Z(t_n)$, will not reproduce its market price; the \emph{z-spread} $s$ is the parallel shift
of the continuously compounded discount rates that does:
\[
  P_D=\sum_{i=1}^{n}C_iZ(t_i)e^{-st_i}+Z(t_n)e^{-st_n}.
\]
It measures the bond's credit, liquidity and supply premium relative to swaps, and it lets the bond's rate risk be
expressed as a delta ladder on the swap curve plus a single spread risk.

\begin{example}[z-spread and DV01 (not from class)]
A 5-year bond with 4\% annual coupons trades at 101.50. Off the curve of \cref{ch06:ex-boot} it would be worth
$P'_D=102.71$; the z-spread is $s=25.6$\,bp. Its yield is 3.666\%, $D_{\mathrm{Mac}}=4.633$,
$D_{\mathrm{mod}}=4.469$, and its DV01 is $D_{\mathrm{mod}}\,P\times10^{-4}=0.0454$ per 100 of face value.
\end{example}

\paragraph{Bond flavours} (slide 44): government, asset-backed, corporate, municipal, high yield and more, each with
its own valuation and risk nuances. \emph{Municipal bonds} were discussed in class \ts{24}{7}{1:08:24}: a huge number
of small and heterogeneous issuers (states, towns, school boards, toll-road and bridge projects) and interest exempt
from federal income tax. At a tax rate $\tau$ an investor is indifferent between a muni yielding $y_m$ and a taxable
bond yielding $y_T$ when $y_m=(1-\tau)y_T$, so munis yield about 70\% of comparable taxable bonds
\ts{24}{7}{1:09:26}. Only taxpayers benefit, so the market is retail: the lecturer's former electronic muni platform
(acquired by TD) did 5{,}000--10{,}000 trades a day with an average size around \$28{,}000, absurdly small by
government-bond standards but profitable in volume. Emerging-market bonds are idiosyncratic (payment terms, credit
support) and still mostly traded by phone \ts{24}{7}{1:10:26}.

\begin{finance}[From the Q\&A]
\emph{Rates outlook} \ts{24}{7}{1:15:32}. The Fed balances inflation, pushed up by COVID-era stimulus and by the
energy shock after Russia's invasion of Ukraine, against employment. High rates discourage borrowing (business
start-ups, and above all housing, the largest purchase in most lives), and a cutting cycle had just begun.
\emph{Two kinds of fixed-income research} \ts{24}{7}{1:16:39}: client-facing economics research makes forecasts
and trade recommendations; internal quant research on pricing and hedging makes no predictions at all. It only
needs to know how to hedge (``if I knew where rates are going, I would be on my Caribbean island'').
\emph{Muni analytics} \ts{24}{7}{1:17:42}: revenue bonds depend on project cash flows (toll roads lost their revenue
in COVID); all munis depend on future tax policy, since the value of the exemption moves with the tax rate
(\cref{ch06:ex-muni}); and the practical problem is \emph{sourcing}: finding who holds the bonds a client wants,
which is why munis are still voice-driven \ts{24}{7}{1:18:43}.
\end{finance}

% ------------------------------------------------------------------
\begin{keyformulas}
\begin{align*}
&\text{Simple interest, deposit:} && I=Nr\Delta t,\\
& && Z(t)=1/\bigl(1+r(t)\Delta\bigr)\\
&\text{SOFR term rates:} && R^{\mathrm{comp}}=\tfrac{360}{d}\bigl[\textstyle\prod_i(1+\tfrac{r_id_i}{360})-1\bigr],\\
& && R^{\mathrm{avg}}=\tfrac1d\sum_ir_id_i\\
&\text{Discounting:} && r_d=r_b\ \text{(discount at the funding rate)}\\
&\text{Forward rate:} && f(t_1,t_2)=\tfrac1\Delta\bigl(Z(t_1)/Z(t_2)-1\bigr)\\
&\text{Futures:} && f=100-F-\mathrm{Cvx},\\
& && \mathrm{Cvx}\approx\tfrac12\sigma^2t_1t_2\ \text{(Ho--Lee)}\\
&\text{Swap legs:} && \mathrm{PV}_{\mathrm{fix}}=NK\textstyle\sum_i\Delta_iZ(t_i)=NKA,\\
& && \mathrm{PV}_{\mathrm{float}}=N\bigl[Z(t_0)-Z(t_n)\bigr]\\
&\text{Par rate, MTM:} && S=\bigl(Z(t_0)-Z(t_n)\bigr)/A,\\
& && V_{\mathrm{payer}}=NA(S-K),\quad \mathrm{DV01}\approx NA\cdot10^{-4}\\
&\text{CDF interpolation:} && \ln Z\ \text{linear between knots}\iff\\
& && f(t)\ \text{piecewise constant}\\
&\text{Calibration:} && E_i=P_i(Z)-P_i^{\mathrm{obs}}=0\quad\text{or}\quad\min\textstyle\sum_iE_i^2\\
&\text{PV01, gamma:} && \tfrac{\mathrm{PV}(r+\epsilon)-\mathrm{PV}(r-\epsilon)}{2\epsilon},\\
& && \Gamma=\tfrac{\mathrm{PV}(r+\epsilon)-2\mathrm{PV}(r)+\mathrm{PV}(r-\epsilon)}{\epsilon}\\
&\text{Hedge with inputs:} && x_k=-\delta_k/A_k\\
&\text{P\&L explain:} && \mathrm{PnL}=\text{portfolio}+\Theta\,\Delta t\\
& && +\textstyle\sum_i\delta_i\Delta x_i+\tfrac12\sum_{ij}\Gamma_{ij}\Delta x_i\Delta x_j+\text{unexplained}\\
&\text{Bonds:} && D_{\mathrm{Mac}}=D_{\mathrm{mod}}(1+y/f),\\
& && P_D=\textstyle\sum_iC_iZ(t_i)e^{-st_i}+Z(t_n)e^{-st_n}\\
&\text{Munis:} && y_m=(1-\tau)\,y_T
\end{align*}
\end{keyformulas}

% ------------------------------------------------------------------
\section*{Exercises}
\addcontentsline{toc}{section}{Exercises}
No problem set covers this lecture; all exercises are \emph{not from the course}.

\begin{exercise}[not from the course]\label{ch06:ex-sofr}
\leavevmode
\begin{enumerate}[(a)]
  \item Show that $R^{\mathrm{comp}}\ge R^{\mathrm{avg}}$ in \eqref{ch06:eq-sofr} whenever all $r_i\ge0$.
  \item For a constant rate $r$ fixed every day for $n$ days ($d_i=1$), show that
      $R^{\mathrm{comp}}-r=\frac{(n-1)r^2}{720}+O(r^3)$. Evaluate for $r=5.2\%$, $n=91$ and compare with the example of
      \cref{ch06:sec-rates}.
\end{enumerate}
\end{exercise}

\begin{exercise}[not from the course]
(a) Show that a note paying compounded SOFR in arrears on $N$,
plus $N$ at maturity, is worth $N$ at every reset date when
discounting is on the SOFR curve.

(b) Deduce \eqref{ch06:eq-float} again.

(c) USD SOFR swaps pay two business days after the end of each
period. Write the value of one floating period with this payment
lag and explain why the telescoping is no longer exact.
\end{exercise}

\bigskip

\begin{exercise}[not from the course]
Using the quotes of \cref{ch06:ex-boot}:

(a) bootstrap $Z_1,Z_2,Z_3$ by hand;

(b) show that CDF gives $Z_4=\sqrt{Z_3Z_5}$, write the scalar
equation for $Z_5$ and solve it numerically;

(c) compute the one-year forward rates $f(k-1,k)$,
$k=1,\dots,5$, and the 4-year par rate;

(d) prove \cref{ch06:prop-cdf} directly from the daily compounding
definition, and show that the forward curve jumps at a knot unless
the two adjacent segments have equal slopes of $\ln Z$.
\end{exercise}

\bigskip

\begin{exercise}[not from the course]
\leavevmode
\begin{enumerate}[(a)]
  \item Explain, without formulas, why the convexity adjustment makes
      the futures rate \emph{higher} than the forward rate, and what
      would change if margin balances earned a fixed rate instead of the
      overnight rate.
  \item With $\sigma=0.8\%$ per year, evaluate
      $\frac12\sigma^2t_1t_2$ for 3-month contracts expiring in 2, 5 and
      10 years. From which maturity does it exceed a bid--offer of
      0.25\,bp?
\end{enumerate}
\end{exercise}
\par\medskip

\begin{exercise}[not from the course]
A company has a 5-year \$50M floating-rate loan at
SOFR$+150$\,bp and wants to fix its cost.
\begin{enumerate}[(a)]
  \item Should it pay or receive fixed on a swap?
  \item With the 5-year par rate at 3.40\%, what all-in fixed
      rate does it lock in (ignore day-count and payment-date
      differences)?
  \item One year later the 4-year par rate is 3.90\% and
      $A_4=3.60$. What is the swap's value to the company,
      and where does the offsetting loss appear?
\end{enumerate}
\end{exercise}

\begin{exercise}[not from the course]\label{ch06:ex-compadv}
Firm A can borrow fixed at 4.00\% or floating at SOFR$+0.30\%$;
firm B at 5.20\% fixed or SOFR$+1.00\%$. A wants floating, B fixed.
\begin{enumerate}[(a)]
  \item Show that A has an absolute advantage in both markets but a
      comparative advantage in fixed.
  \item Design a swap through a bank earning 0.10\% that splits the
      remaining gain equally between A and B.
  \item Which risk does each firm take on in exchange (think of
      \cref{ch06:sec-clearing})?
\end{enumerate}
\end{exercise}
\par\medskip

\begin{exercise}[not from the course]
\leavevmode
\begin{enumerate}[(a)]
  \item Show that the gamma of slide 31 equals
      $\epsilon\,\partial^2\mathrm{PV}/\partial r^2+O(\epsilon^3)$.
  \item For a 10-year zero-coupon bond with continuously compounded
      rate 4\% and notional \$100M, compute the PV01 and this gamma for
      $\epsilon=1$\,bp.
  \item How many revaluations does a full central-difference
      cross-gamma matrix need for $N$ curve inputs?
  \item Why does a Monte Carlo pricer need a larger $\epsilon$
      (\cref{ch:markets})?
\end{enumerate}
\end{exercise}

\begin{exercise}[not from the course]\label{ch06:ex-theta}
Annual-compounding zero rates are 3\% (1 year) and 4\% (2 years).
\begin{enumerate}[(a)]
  \item Price a 2-year zero-coupon bond.
  \item Value it one year later if the forwards are realised, and if
      the curve by tenor is unchanged.
  \item Split the one-year return in (b) into carry and roll-down,
      and relate each to the two theta conventions of
      \cref{ch06:sec-pnl}.
\end{enumerate}
\end{exercise}

\begin{exercise}[not from the course]
Using the bond of the example in \cref{ch06:sec-bonds} (DV01
$0.0454$ per 100) and the 5-year swap annuity $A_5=4.5178$:
\begin{enumerate}[(a)]
  \item which swap and which notional make \$10M face of the bond
      DV01-neutral?
  \item Show that for small $s$,
      $s\approx(P'_D-P_D)/\sum_it_i\mathrm{CF}_iZ(t_i)$, and check it
      against $s=25.6$\,bp.
  \item Which risk remains after the hedge?
\end{enumerate}
\end{exercise}

\begin{exercise}[not from the course]\label{ch06:ex-muni}
\leavevmode
\begin{enumerate}[(a)]
  \item At which tax rate is a 2.8\% muni equivalent to a 4\%
      taxable bond?
  \item A 10-year muni is priced to yield 70\% of a 4\%
      taxable yield. If investors come to expect the tax rate to
      fall from 30\% to 20\%, estimate the price change with a
      modified duration of 8, all else equal. Connect this to the
      Q\&A remark on tax policy.
\end{enumerate}
\end{exercise}

% !TEX root = ../main.tex
\chapter{Principal Component Analysis and Factor Models}\label{ch:pca}
\begin{paracol}{2}

\begin{sources}
\begin{tabularx}{\linewidth}{@{}l l L@{}}
\toprule
\textbf{Lecture} & \textbf{Speaker} & \textbf{Content}\\
\midrule
2024 L9 & S.~Andreev (Two Sigma) & PCA as a rotation of coordinates; toy examples, robustness, PCA versus OLS; demeaning and variance normalisation; eigenvalue diagnostics, in- and out-of-sample variance explained; PCA on financial time series; US Treasury curve (level, slope, curvature, regime indicators); PC-neutral curves and flies, leverage.\\
2013 L15 & P.~Kempthorne & Linear factor models as cross-sectional and time-series regressions; Sharpe's single-index model, macroeconomic and fundamental (BARRA, Fama--French) models; factor-mimicking portfolios; statistical factor models: factor analysis (maximum likelihood, EM, likelihood-ratio test) and PCA; case study on Treasury yields.\\
\bottomrule
\end{tabularx}

\smallskip
\textbf{Files.} Slides \file{mit18_642_f24_lec09.pdf} (the European-indices example, slides 26--34, was
not covered in class); Jupyter notebook \file{feddataanalysis.ipynb} with the data file
\file{fedbondyielddata_20230911.csv} (distributed in 2024 week 5); readings
\file{mit18_642_f24_pset_analysis.pdf} (P.~Kempthorne, \emph{Principal Components Analysis in Finance:
Modeling the Dynamics of Interest Rates}, 2016) and \file{mit18_642_f24_pset_rcode.pdf} (its R code and a
5-year PCA of Treasury yields); 2013 slides \file{MIT18_S096F13_lecnote15.pdf} and case study
\file{MIT18_S096F13_CaseStudy7.pdf}.
\textbf{Problem sets.} 2024 PS2 Problem~3 and PS3 Problem~4; 2013 PS7 (all problems).
\end{sources}

\section*{Motivation}
Financial data come in large and highly redundant blocks: the yields of thirty Treasury maturities,
the returns of five hundred stocks, a dozen European equity indices. When the 10-year yield rises the
5- and 30-year yields almost always rise too, so thirty numbers per day carry perhaps three numbers'
worth of information. This chapter develops two complementary ways of exploiting the redundancy:
\begin{itemize}
  \item \textbf{Principal component analysis (PCA)}: a rotation of coordinates such that the first few
        axes carry most of the variance. It is model-free (unsupervised learning, no dependent
        variable), pure linear algebra, and fast \ts{24}{9}{3:05}.
  \item \textbf{Factor models}: the statistical model $\bm x_t=\bm\alpha+B\bm f_t+\bm\epsilon_t$ with a few
        common factors and uncorrelated asset-specific noise. It reduces the $m(m+1)/2$ entries of a
        covariance matrix to $O(mK)$ parameters, and is the backbone of risk models and of the
        CAPM/APT view of returns \ts{13}{15}{0:00}.
\end{itemize}
The 2024 lecture is the practitioner's view (a fixed-income quant at Two Sigma: workflow, pitfalls, and
how PCA turns into tradeable portfolios); the 2013 lecture is the statistician's view (models, estimation,
tests). Both analyse daily changes of US Treasury yields and find the same three factors: \emph{level},
\emph{slope} and \emph{curvature}. The mathematics needed is the spectral theorem, the SVD and
Perron--Frobenius of \cref{ch:linalg} (\cref{thm:spectral,thm:svd,thm:pf}) and least squares
(\cref{ch:regression}).

\begin{coursenote}[Notation across the sources]
The 2024 slides use a data matrix $X$ of size $N\times p$ (rows = observations) and loadings $w_k$, $W$;
the 2013 slides use $m$ assets, $T$ periods, a data matrix of size $m\times T$ (columns = periods) and
eigenvectors $\gamma_k$, $\Gamma$; the 2016 reading uses a tenors$\times$days matrix $A$. Here:
$p$ variables and $N$ observations with $X\in\R^{N\times p}$ in the PCA sections, $m$ assets and $T$ periods in
the factor-model sections, and $\gamma_k$, $\Gamma$ for eigenvectors (loadings) throughout.
\end{coursenote}

% ------------------------------------------------------------------
\section{Principal components of a random vector}\label{ch07:sec-pop}

Let $\bm x=(x_1,\dots,x_p)\T$ be a random vector with $\E[\bm x]=\bm\mu$ and $\Cov[\bm x]=\Sigma$. For any
weight vector $\bm w$, $\Var(\bm w\T\bm x)=\bm w\T\Sigma\bm w\ge0$: $\Sigma$ is symmetric positive semi-definite
(\cref{sec:symmetric}; 2024 PS2 Problem~3). If $\bm x$ are asset returns and $\bm w$ portfolio weights,
$\bm w\T\Sigma\bm w$ is the portfolio variance. By the spectral theorem (\cref{thm:spectral})
\[
  \Sigma=\Gamma\Lambda\Gamma\T=\sum_{k=1}^p\lambda_k\,\bm\gamma_k\bm\gamma_k\T,\qquad
  \lambda_1\ge\lambda_2\ge\dots\ge\lambda_p\ge0,\qquad \Gamma\T\Gamma=I_p .
\]

\subsection{Variance maximisation}

\begin{definition}[Principal components]\label{ch07:def-pca}
The \emph{first principal direction} is a unit vector $\bm w_1$ maximising $\Var(\bm w\T\bm x)=\bm w\T\Sigma\bm w$
subject to $\|\bm w\|=1$. The $k$-th principal direction $\bm w_k$ maximises the same variance subject to
$\|\bm w\|=1$ and $\bm w\perp\bm w_1,\dots,\bm w_{k-1}$ \ts{13}{15}{1:07:08}. The \emph{principal component
variables} (scores, PC factors) are $p_k=\bm w_k\T(\bm x-\bm\mu)$; the entries of $\bm w_k$ are the
\emph{loadings} of PC~$k$, and $\lambda_k/\sum_j\lambda_j$ is the \emph{fraction of variance explained}.
\end{definition}

\begin{theorem}[PCA is the spectral decomposition of $\Sigma$]\label{ch07:thm-pca}
With the notation above:
\begin{enumerate}[(i)]
  \item the $k$-th principal direction can be taken to be the eigenvector $\bm\gamma_k$, and
        $\Var(p_k)=\lambda_k$ (if the eigenvalues are distinct, $\bm w_k=\pm\bm\gamma_k$);
  \item $\bm p=\Gamma\T(\bm x-\bm\mu)$ has $\E[\bm p]=\bm 0$ and $\Cov[\bm p]=\Lambda$: the PCs are uncorrelated;
  \item $\bm x=\bm\mu+\Gamma\bm p=\bm\mu+\sum_{k=1}^p p_k\,\bm\gamma_k$;
  \item total variance is preserved: $\sum_i\Var(x_i)=\tr\Sigma=\tr\Lambda=\sum_k\lambda_k=\sum_k\Var(p_k)$.
\end{enumerate}
\end{theorem}
\begin{proof}
(i) Expand $\bm w=\sum_jc_j\bm\gamma_j$; then $\|\bm w\|^2=\sum_jc_j^2=1$ and $\bm w\T\Sigma\bm w=\sum_j\lambda_jc_j^2\le\lambda_1$,
with equality for $\bm c=\bm e_1$ (and only for $\bm c=\pm\bm e_1$ if $\lambda_1>\lambda_2$). Inductively, if
$\bm w_j=\bm\gamma_j$ for $j<k$, the constraint $\bm w\perp\bm\gamma_1,\dots,\bm\gamma_{k-1}$ forces $c_1=\dots=c_{k-1}=0$, so
$\bm w\T\Sigma\bm w=\sum_{j\ge k}\lambda_jc_j^2\le\lambda_k$, attained at $\bm\gamma_k$.

(ii) $\E[\bm p]=\Gamma\T(\E[\bm x]-\bm\mu)=\bm 0$ and $\Cov[\bm p]=\Gamma\T\Sigma\Gamma=\Gamma\T\Gamma\Lambda\Gamma\T\Gamma=\Lambda$
\ts{13}{15}{1:01:42}.

(iii) Multiply $\bm p=\Gamma\T(\bm x-\bm\mu)$ by $\Gamma=(\Gamma\T)^{-1}$.

(iv) $\tr(\Gamma\Lambda\Gamma\T)=\tr(\Lambda\Gamma\T\Gamma)=\tr\Lambda$ \ts{13}{15}{1:09:11}.
\end{proof}

Equivalently, with a Lagrange multiplier, $\nabla_{\bm w}[\bm w\T\Sigma\bm w-\lambda(\bm w\T\bm w-1)]=0$ gives
$\Sigma\bm w=\lambda\bm w$: the stationary points of the Rayleigh quotient $\bm w\T\Sigma\bm w/\bm w\T\bm w$ are the
eigenvectors and the stationary values the eigenvalues (\cref{sec:symmetric}).

\begin{intuition}[A rigid rotation of the data cloud]
$\bm x\mapsto\Gamma\T(\bm x-\bm\mu)$ moves the origin to the centre of mass and rotates the axes; distances are
unchanged \ts{13}{15}{1:02:46}. Andreev's toy example: a one-dimensional data set embedded in the plane
along the $x$-axis, the diagonal or the $y$-axis looks two-dimensional, but PCA immediately returns one
direction with all the variance and a second with none \ts{24}{9}{5:11}; add a little noise in the
perpendicular direction and $\lambda_2$ becomes small but non-zero \ts{24}{9}{7:22}. PCA returns two things:
\emph{directions} (eigenvectors) and their \emph{importance} (eigenvalues). For a physicist, $\Sigma$ is the
second-moment tensor of the cloud of points (unit masses); the inertia tensor $I=\tr(\Sigma)\,\mathbb 1-\Sigma$ has
the same eigenvectors, so the principal directions are the \emph{principal axes of inertia}, PC1 being the
axis with the smallest moment of inertia (the direction along which the cloud is most elongated).
\end{intuition}

\begin{remark}[Uniqueness]\label{ch07:rem-sign}
If $\bm\gamma_k$ is an eigenvector, so is $-\bm\gamma_k$: the sign of each PC is arbitrary (the 2024 slides:
``sign does not matter''), and different software may return different signs for the same data. If
$\lambda_k=\lambda_{k+1}$, any orthonormal basis of the eigenspace is equally valid: the directions are not
determined at all. This is the mathematical root of the robustness problems of \cref{ch07:sec-robust}.
\end{remark}

\begin{example}[Two variables with equal variances]\label{ch07:ex-2d}
For $\Sigma=\sigma^2\bigl(\begin{smallmatrix}1&\rho\\ \rho&1\end{smallmatrix}\bigr)$ with $\rho>0$: $\lambda_{1,2}=\sigma^2(1\pm\rho)$ with
$\bm\gamma_1=\frac1{\sqrt2}(1,1)\T$ (the sum) and $\bm\gamma_2=\frac1{\sqrt2}(1,-1)\T$ (the difference), whatever $\rho$ is;
PC1 explains the fraction $(1+\rho)/2$. As $\rho\to0$ the eigenvalues merge: an isotropic ``blob'' has no
preferred direction and moving a single data point can swing the estimated PC1 by any angle
\ts{24}{9}{12:44}. (The general bivariate case is 2013 PS7 Problem~1.)
\end{example}

\subsection{Geometry: best subspace, and PCA versus regression}\label{ch07:sec-geom}

\begin{proposition}[Best $K$-dimensional approximation]\label{ch07:prop-best}
Let $V\in\R^{p\times K}$ have orthonormal columns, so that $VV\T$ is the orthogonal projection onto their span.
Then
\[
  \E\bigl\|(\bm x-\bm\mu)-VV\T(\bm x-\bm\mu)\bigr\|^2=\tr\Sigma-\tr(V\T\Sigma V)\ \ge\ \sum_{j>K}\lambda_j ,
\]
with equality for $V=\Gamma_K=[\bm\gamma_1\cdots\bm\gamma_K]$. The first $K$ PCs span the $K$-dimensional subspace that
minimises the mean squared \emph{perpendicular} distance of the data from it.
\end{proposition}
\begin{proof}
With $\bm y=\bm x-\bm\mu$, $\|(I-VV\T)\bm y\|^2=\bm y\T(I-VV\T)\bm y$ since $I-VV\T$ is a projection; taking expectations gives
$\tr\Sigma-\tr(V\T\Sigma V)$. Now $\tr(V\T\Sigma V)=\sum_j\lambda_jc_j$ with $c_j=\|V\T\bm\gamma_j\|^2$. Each
$c_j\in[0,1]$ (a projection does not increase length) and $\sum_jc_j=\tr(V\T\Gamma\Gamma\T V)=\tr(V\T V)=K$. Hence
\[
  \sum_j\lambda_jc_j-\sum_{j\le K}\lambda_j=\sum_{j\le K}\lambda_j(c_j-1)+\sum_{j>K}\lambda_jc_j
  \le\lambda_K\Bigl(\sum_{j\le K}(c_j-1)+\sum_{j>K}c_j\Bigr)=0,
\]
and $V=\Gamma_K$ gives $c_j=1$ for $j\le K$, $0$ otherwise.
\end{proof}
For a sample this is the Eckart--Young statement of \cref{sec:svd}: the truncated SVD of the centred data
matrix is its best rank-$K$ approximation.

\paragraph{PCA is not regression} \ts{24}{9}{8:28}. OLS of $y$ on $x$ minimises \emph{vertical} squared
distances, OLS of $x$ on $y$ \emph{horizontal} ones, PCA \emph{perpendicular} ones (\cref{ch07:fig-ols}).
For standardised variables with correlation $\rho$ the three lines through the centre have slopes $\rho$,
$1/\rho$ and $1$ \ts{24}{9}{9:31}. OLS presupposes a dependent variable (a causal or predictive
direction); PCA treats all coordinates symmetrically and says nothing about causality. Use regression to
predict $y$ from $x$, PCA to describe the joint correlation structure \ts{24}{9}{10:34}. Two further
differences: PCA is equivariant under rotations of the data (rotate the cloud, the PCs rotate with it)
\ts{24}{9}{11:38}, but \emph{not} under rescaling of individual variables, whereas OLS predictions are
unaffected by changes of units. Hence the covariance-versus-correlation question of
\cref{ch07:sec-corr}.

\begin{figure}[ht]
\centering
\begin{tikzpicture}[scale=1.05,>=Stealth]
  \draw[->,black!50] (-2.9,0) -- (2.9,0) node[right,black]{\small $x$};
  \draw[->,black!50] (0,-2.9) -- (0,2.9) node[above,black]{\small $y$};
  \draw[black!30,dashed,rotate=45] (0,0) ellipse ({2*sqrt(1.6)} and {2*sqrt(0.4)});
  \foreach \p in {(0.24,0.54),(-0.06,-0.74),(-0.26,-0.95),(0.31,1.55),(-0.30,-0.63),(0.78,0.90),(0.36,-0.54),
    (0.21,0.89),(-1.24,-1.01),(-1.85,-2.14),(-1.79,-1.12),(-1.16,-0.28),(0.41,0.19),(-2.53,-1.83),(0.19,0.34),
    (-1.44,-1.15),(-0.84,-1.11),(1.41,0.18),(0.21,1.07),(-0.40,-0.21),(0.36,0.39),(-1.11,-0.44),(1.74,-0.32),
    (1.19,0.92),(1.08,-0.37),(0.32,0.85),(0.03,0.78),(0.17,0.84),(1.83,0.54),(0.46,-0.04),(0.38,-0.76),
    (-0.40,-0.29),(1.23,1.90),(-1.22,-1.31),(0.95,-1.18),(-0.27,-0.12),(1.63,1.70),(-0.12,-0.29),(-0.03,1.52)}
    \fill[black!45] \p circle (1.2pt);
  \draw[thick,thmcol] (-2.8,-1.68) -- (2.8,1.68) node[right,font=\scriptsize]{OLS $y\sim x$ (slope $\rho$)};
  \draw[thick,defcol] (-1.68,-2.8) -- (1.68,2.8) node[above,font=\scriptsize]{OLS $x\sim y$ (slope $1/\rho$)};
  \draw[very thick,intcol] (-2.55,-2.55) -- (2.55,2.55) node[right,font=\scriptsize]{PC1 (slope 1)};
  % residuals of one point
  \fill[black] (-0.47,1.72) circle (1.8pt) node[above left,font=\scriptsize]{$P$};
  \draw[thick,thmcol] (-0.47,1.72) -- (-0.47,-0.282);
  \draw[thick,defcol] (-0.47,1.72) -- (1.032,1.72);
  \draw[thick,intcol] (-0.47,1.72) -- (0.625,0.625);
\end{tikzpicture}
\caption{Forty standardised points with sample correlation $\rho=0.6$ and the $2\sigma$ ellipse. For the point
$P$ the three segments are the distances minimised by OLS of $y$ on $x$ (vertical), OLS of $x$ on $y$
(horizontal) and PCA (perpendicular); cf.\ slide 8 of \file{lec09} \ts{24}{9}{9:31}.}\label{ch07:fig-ols}
\end{figure}

\subsection{Positive correlations and the market factor}\label{ch07:sec-pf}
If all covariances are positive, $\Sigma_{ij}>0$ (typical for yields of different maturities or for stocks in
one market), the Perron--Frobenius theorem (\cref{thm:pf}) applies to $\Sigma$: $\lambda_1$ is simple and
$\bm\gamma_1$ can be chosen with \emph{all entries positive} (2024 PS2 Problem~3). PC1 is then a long-only
portfolio: the ``market'' or ``level'' factor, and all other PCs, being orthogonal to a positive vector,
have loadings of both signs: they are long/short spreads.

\begin{example}[Equicorrelated assets (not discussed in class)]\label{ch07:ex-equi}
For $p$ assets with common variance $\sigma^2$ and common correlation $\rho>0$,
$\Sigma=\sigma^2[(1-\rho)I+\rho\,\bm 1\bm 1\T]$. Then $\lambda_1=\sigma^2(1+(p-1)\rho)$ with $\bm\gamma_1=\bm1/\sqrt p$ (the
equal-weighted portfolio), and $\lambda_2=\dots=\lambda_p=\sigma^2(1-\rho)$. PC1 explains
$(1+(p-1)\rho)/p\to\rho$ as $p\to\infty$: even a modest average correlation produces a dominant first
component in a large universe, while the remaining $p-1$ eigenvalues are degenerate, so the individual
higher PCs are not identified (\cref{ch07:rem-sign}).
\end{example}

\begin{exercise}[PS2 (2024), Problem 3: principal components, special case]\label{ch07:ex-ps2}
A portfolio holds $n$ assets with fixed weights $\bm w=(w_1,\dots,w_n)\T$; over a given horizon the assets have
rates of return $\bm X=(X_1,\dots,X_n)\T$ with mean $\bm\mu$ and covariance $\Sigma=(\sigma_{ij})$, and the portfolio return
is $Y=\bm w\T\bm X$.
\begin{enumerate}[(a)]
  \item Show that $\E[Y]=\bm w\T\bm\mu$ and $\Var[Y]=\bm w\T\Sigma\bm w$.
  \item Prove that $\Sigma$ is positive semi-definite, for any random variables $X_1,\dots,X_n$.
  \item Suppose all entries of $\Sigma$ are positive. Prove that the largest eigenvalue $\lambda_{\max}$ has multiplicity one and
      an eigenvector $\bm v$ that can be chosen with all components positive.
  \item Write the first principal component variable $PC1$ in terms of $\bm X$, $\bm\mu$ and $\bm v$. What are its mean and
      variance?
\end{enumerate}
\end{exercise}

\begin{exercise}[not from the course]\label{ch07:ex-equi-proof}
Prove the claims of \cref{ch07:ex-equi}: for $\Sigma=\sigma^2[(1-\rho)I+\rho\bm1\bm1\T]$ find all eigenvalues and eigenvectors, and
the fraction of variance explained by PC1. For which $\rho$ is $\Sigma$ positive semi-definite? Relate the result to
\cref{thm:pf} when $\rho>0$.
\end{exercise}

% ------------------------------------------------------------------
\section{Sample PCA, SVD and practical choices}\label{ch07:sec-sample}

\subsection{Data matrix and singular value decomposition}
Stack $N$ observations (days) of $p$ variables (tenors, assets) in $X\in\R^{N\times p}$, row $i$ being $\bm x_i\T$.
With $\bar{\bm x}=X\T\bm1/N$, the centred matrix and sample covariance are
\[
  X_c=X-\bm 1\bar{\bm x}\T=\bigl(I_N-\tfrac1N\bm1\bm1\T\bigr)X,\qquad S=\frac{1}{N-1}X_c\T X_c
\]
(the 2013 slides divide by $N$; it does not matter for the eigenvectors) \ts{13}{15}{41:21}.

\begin{proposition}[PCA via the SVD]\label{ch07:prop-svd}
Let $X_c=UD\hat\Gamma\T$ be a thin SVD: $U\in\R^{N\times p}$ with $U\T U=I_p$, $D=\diag(d_1\ge\dots\ge d_p\ge0)$,
$\hat\Gamma\in\R^{p\times p}$ orthogonal. Then:
\begin{enumerate}[(i)]
  \item $S=\hat\Gamma\hat\Lambda\hat\Gamma\T$ with $\hat\lambda_k=d_k^2/(N-1)$: the right singular vectors are the loadings;
  \item the $N\times p$ matrix of PC scores $P=X_c\hat\Gamma=UD$ has zero column means and orthogonal columns;
        column $k$ (the time series of PC~$k$) has sample variance $\hat\lambda_k$;
  \item $X_c=P\hat\Gamma\T=\sum_kd_k\bm u_k\hat{\bm\gamma}_k\T$, and truncation after $K$ terms is the best rank-$K$
        approximation of $X_c$;
  \item $\rank X_c\le\min(N-1,p)$: with $N\le p$ observations at most $N-1$ eigenvalues are non-zero.
\end{enumerate}
\end{proposition}
\begin{proof}
(i) $X_c\T X_c=\hat\Gamma DU\T UD\hat\Gamma\T=\hat\Gamma D^2\hat\Gamma\T$.

(ii) $P=UD\hat\Gamma\T\hat\Gamma=UD$, so
$P\T P=D^2$; $\bm1\T P=\bm1\T X_c\hat\Gamma=\bm0$ because $\bm1\T X_c=\bm 0$.

(iii) is \cref{thm:svd} and the Eckart--Young
theorem.

(iv) $\bm1\T X_c=\bm0$ means the $N$ rows of $X_c$ are linearly dependent.
\end{proof}

Andreev computes the covariance matrix and diagonalises it (``we need the covariance matrix for other
things anyway''); the SVD gives the same result without forming $S$ \ts{24}{9}{18:02}, \ts{13}{15}{1:06:05}.
Numerically the SVD is preferable because forming $X_c\T X_c$ squares the condition number; in cost both
routes are $O(Np^2)$ (plus $O(p^3)$ for the eigen-decomposition).

\begin{coursenote}[Erratum: slides 10--11 of \file{lec09}]
\begin{itemize}
  \item ``$\Sigma=X^TX$ is the covariance matrix'' needs the factor $1/N$ (or $1/(N-1)$); without it the SVD
      relation reads $X\T X=W M\T MW\T$ with $M=\diag(\sqrt{(N-1)\lambda_k})$, not $\diag(\sqrt{\lambda_k})$.
  \item ``Computation of the covariance matrix is $O(N^2p^2)$'' is wrong: $X\T X$ costs $O(Np^2)$, the same order as
      the thin SVD; the advantage of the SVD is accuracy, not asymptotic cost.
  \item ``$W$ is the eigenvectors of
      $X$'': $W$ holds the right singular vectors of $X$, i.e.\ the eigenvectors of $X\T X$.
  \item Slide 11 calls
      $\Sigma=W\Lambda W\T$ the ``SVD decomposition'': it is the spectral decomposition (for a PSD matrix it coincides with the SVD).
\end{itemize}
\end{coursenote}

\begin{exercise}[\file{lecnote15}, slide 32: PCA via the SVD]\label{ch07:ex-lecsvd}
Let $X^*$ be the $m\times T$ row-centred data matrix, $\hat\Sigma_x=\frac1TX^*X^{*\prime}$, and $X^*=VDU'$ its SVD with $V$ ($m\times m$)
orthogonal, $U$ ($T\times m$) with orthonormal columns and $D=\diag(d_1\ge\dots\ge d_m\ge0)$. Show that
$\hat\Lambda=\frac1TD^2$, $\hat\Gamma=V$, and that the matrix of sample principal components is $P=\hat\Gamma'X^*=DU'$.
\end{exercise}

\subsection{A worked example: six days of Treasury yields}\label{ch07:sec-sixdays}
The 2016 reading (\file{pset_analysis}) runs PCA on the smallest possible data set: FRED constant-maturity
yields of 9 Treasury securities (3 months to 20 years) on the first six business days of 2001. Daily
changes in basis points, minus their row means, form the $9\times5$ matrix $A$ of \cref{ch07:tab-sixdays}
(tenors $\times$ days, the transpose of $X_c$, so the roles of $U$ and $\hat\Gamma$ are swapped: loadings are the
\emph{left} singular vectors of $A$).

\begin{table}[ht]
\centering\small
\caption{Six-day example: centred yield changes $A$ (bp) and the PCA of \file{pset_analysis}. Last rows:
PC variances $\hat\lambda_k=d_k^2/(n-1)$, $n=5$ days, and fractions of variance.}\label{ch07:tab-sixdays}
\begin{tabular}{@{}l rrrrr c rrr@{}}
\toprule
 & \multicolumn{5}{c}{$A$: centred daily changes (bp)} & & \multicolumn{3}{c}{loadings}\\
\cmidrule(lr){2-6}\cmidrule(l){8-10}
Tenor & Jan 03 & Jan 04 & Jan 05 & Jan 08 & Jan 09 & $s_j$ & PC1 & PC2 & PC3\\
\midrule
3M  & $-5.4$ & $-19.4$ & $-12.4$ & $19.6$ & $17.6$ & 17.70 & 0.38 & 0.53 & $-0.48$\\
6M  & $-4.6$ & $-14.6$ & $-12.6$ & $14.4$ & $17.4$ & 15.03 & 0.34 & 0.44 & $-0.05$\\
1Y  & $1.0$  & $-14.0$ & $-14.0$ & $9.0$  & $18.0$ & 14.12 & 0.36 & 0.26 & 0.23\\
2Y  & $9.6$  & $-10.4$ & $-16.4$ & $2.6$  & $14.6$ & 13.13 & 0.35 & $-0.03$ & 0.46\\
3Y  & $13.4$ & $-10.6$ & $-17.6$ & $1.4$  & $13.4$ & 13.99 & 0.37 & $-0.13$ & 0.43\\
5Y  & $18.6$ & $-11.4$ & $-15.4$ & $-0.4$ & $8.6$  & 14.03 & 0.35 & $-0.29$ & 0.12\\
7Y  & $20.8$ & $-11.2$ & $-14.2$ & $0.8$  & $3.8$  & 13.92 & 0.32 & $-0.37$ & $-0.23$\\
10Y & $20.8$ & $-12.2$ & $-11.2$ & $-0.2$ & $2.8$  & 13.37 & 0.30 & $-0.38$ & $-0.35$\\
20Y & $14.6$ & $-7.4$  & $-7.4$  & $0.6$  & $-0.4$ & 8.99  & 0.18 & $-0.28$ & $-0.36$\\
\midrule
\multicolumn{6}{@{}l}{$\hat\lambda_k$ (bp$^2$): 1323.9, 397.2, 34.2, 1.4, 0; total 1756.7} & & 75.4\% & 22.6\% & 1.9\%\\
\bottomrule
\end{tabular}
\end{table}

The singular values are $d=(72.77,\,39.86,\,11.70,\,2.39,\,0)$. Observations:
\begin{itemize}
  \item The total variance $1756.7$ equals the sum of the nine tenor variances $s_j^2$ (\cref{ch07:thm-pca}(iv)).
  \item PC5 has zero variance: five centred columns span at most four dimensions
        (\cref{ch07:prop-svd}(iv)). With real data one uses hundreds of days, yet the reading notes the
        ``curious fact'' that the loadings look almost the same as in the five-day toy.
  \item PC1 has positive, roughly equal loadings (a level shift); its standard deviation, 36.4 bp, exceeds
        that of every single tenor (at most 17.7 bp), because a unit-norm vector with nine positive entries
        of about $1/3$ adds up nine highly correlated moves.
  \item PC2 is short-minus-long (a change of slope), PC3 is negative at the two ends and positive in the
        middle (curvature; the overall sign is arbitrary). The daily scores (in bp) are $d_k\bm v_k$:
        \begin{align*}
          \text{PC1}&=(27.2,-37.6,-41.0,17.2,34.2)\ \text{(up, down, down, up, up)},\\
          \text{PC2}&=(-31.7,-4.5,3.1,18.5,14.6).
        \end{align*}
\end{itemize}

\begin{coursenote}[Erratum: dates in Table 1 of \file{pset_analysis}]
Table~1 of the reading labels the yield columns Jan 03, \dots, Jan 10 and the changes Jan 04, 05, 06, 07, 10;
but 6--7 January 2001 were a Saturday and a Sunday. The data are those of pages 5--6 of the reading: yields on
Jan 02, 03, 04, 05, 08, 09 and changes dated Jan 03, 04, 05, 08, 09, as used in \cref{ch07:tab-sixdays}.
\end{coursenote}

\subsection{Centring and scaling: covariance or correlation?}\label{ch07:sec-corr}
The typical workflow \ts{24}{9}{21:09}:
\begin{enumerate}[1.]
  \item transform the data to be as close as possible to multivariate
      normal (mean and covariance then describe the whole distribution: PCA only sees second moments);
  \item subtract the mean of each variable and possibly scale it to unit variance;
  \item compute the $p\times p$
      covariance (or correlation) matrix, symmetric and PSD by construction;
  \item diagonalise, order by decreasing
      eigenvalue, keep the first $K$ components.
\end{enumerate}

\paragraph{Centring is not optional.} PCA rotates about the origin \ts{24}{9}{25:28}. Feed it raw prices
(large positive numbers) and PC1 simply points from the origin to the centre of the cloud, capturing
nothing about co-movement \ts{24}{9}{26:31}. Slide 13 of \file{lec09} shows the same cloud before and after
centring. For daily yield changes the mean is tiny (a fraction of a basis point per day), so centring changes
little, but it remains good practice \ts{24}{9}{44:13}. In time series the \emph{rolling} mean can itself be
unstable, one more reason to work with changes rather than levels.

\paragraph{Scaling is a modelling choice} \ts{24}{9}{22:10}. Standardising $z_i=(x_i-\mu_i)/\sigma_i$ turns PCA of
$\Sigma$ into PCA of the correlation matrix $R=D_\sigma^{-1}\Sigma D_\sigma^{-1}$, $D_\sigma=\diag(\sigma_1,\dots,\sigma_p)$. Then
\begin{itemize}
  \item the total variance is $\tr R=p$ and every variable enters on the same footing;
  \item the eigenvectors of $R$ are \emph{not} the rescaled eigenvectors of $\Sigma$ (\cref{ch07:sec-geom}:
        PCA is not scale-equivariant);
  \item without scaling, a variable with much larger variance dominates PC1.
\end{itemize}
Andreev's toy example: the same data set looks one-dimensional (strongly elongated) in raw units and like
a structureless blob after normalising both axes \ts{24}{9}{23:12}. Rule of thumb \ts{24}{9}{24:24}: if the
coordinates have different units (temperature in Seattle versus pressure in Atlanta), standardise; if they
have the same units and their relative size is meaningful (grades on two quizzes, yield changes in bp),
covariance PCA keeps useful information about relative volatility. For Treasury yield changes Andreev uses
covariance PCA \ts{24}{9}{45:18}; how to normalise is ``very much the art'' and sometimes the research
question itself. For 2001--2005 Treasury changes the two choices give PC1/PC2/PC3 shares of
84.9/9.2/2.9\% (covariance) and 77.1/15.7/3.5\% (correlation); 2024 PS3 Problem~4 asks why.

For equity indices the slides use log returns $r_t=\ln(p_t/p_{t-1})$, which are closer to normal than price
changes; in that example covariance and correlation PCA give almost the same loadings because the index
variances are similar.

\subsection{How many components, and can we trust them?}\label{ch07:sec-robust}
PCA always returns an answer, even for pure noise \ts{24}{9}{12:44}. Andreev's advice: look at the
\emph{eigenvalues}, not only at the eigenvectors \ts{24}{9}{13:44}. A component is worth using if its eigenvalue
is well separated from the next one, ideally by an order of magnitude; with rolling PCA, plot the eigenvalue
(or variance-explained) time series on a log scale and check that the lines never cross \ts{24}{9}{14:50}. If
PC1 and PC2 swap places under small changes of the data, there is no stable structure and the PCA output
``obfuscates rather than clarifies'' \ts{24}{9}{15:54}. The standard display is the \emph{scree plot} of
$\hat\lambda_k$ against $k$ (\cref{ch07:fig-scree}).

\begin{intuition}[Degenerate perturbation theory (not discussed in class)]
Why do nearly equal eigenvalues make eigenvectors unstable? Estimation noise perturbs the covariance,
$\hat\Sigma=\Sigma+\delta\Sigma$ with $\delta\Sigma=O(N^{-1/2})$. First-order perturbation theory, exactly as for a
Hamiltonian in quantum mechanics, gives
\[
  \delta\lambda_k=\bm\gamma_k\T\,\delta\Sigma\,\bm\gamma_k,\qquad
  \delta\bm\gamma_k=\sum_{j\ne k}\frac{\bm\gamma_j\T\,\delta\Sigma\,\bm\gamma_k}{\lambda_k-\lambda_j}\,\bm\gamma_j .
\]
Eigenvalues are robust, but eigenvectors rotate by an amount inversely proportional to the \emph{gap}
$\lambda_k-\lambda_j$; for a degenerate level the eigenvectors are fixed by the perturbation, not by the data.
Andreev's ``move one point and watch PC1 swing'' demonstration \ts{24}{9}{11:38} is this formula at work.
\end{intuition}

\begin{figure}[ht]
\centering
\begin{tikzpicture}
\begin{axis}[width=0.72\linewidth,height=5.6cm,ymode=log,xlabel={component $k$},
  ylabel={fraction of variance},xmin=0.5,xmax=9.5,ymin=0.001,ymax=1.5,xtick={1,...,9},
  legend pos=north east,legend style={font=\small},grid=major,grid style={black!10},mark size=1.6pt]
  \addplot[thick,defcol,mark=*] coordinates {(1,0.8494) (2,0.0919) (3,0.0290) (4,0.0119) (5,0.00588)
    (6,0.00407) (7,0.00297) (8,0.00263) (9,0.00215)};
  \addlegendentry{covariance PCA}
  \addplot[thick,fincol,mark=square*] coordinates {(1,0.7708) (2,0.1571) (3,0.0355) (4,0.0175) (5,0.00806)
    (6,0.00481) (7,0.00245) (8,0.00207) (9,0.00167)};
  \addlegendentry{correlation PCA}
\end{axis}
\end{tikzpicture}
\caption{Scree plot (log scale) of the PCA of daily changes of 9 Treasury yields, 2001--2005 (2013 Case
Study~7; 2024 PS3 Problem~4). The first three eigenvalues are separated by roughly an order of magnitude
each, the pattern Andreev wants to see \ts{24}{9}{53:42}; standardising moves variance from PC1 to PC2.}\label{ch07:fig-scree}
\end{figure}

\paragraph{Effective dimension.} Slide 22 of \file{lec09} (after M.~Del Giudice, 2020) summarises a spectrum
in one number. With $\pi_k=\lambda_k/\sum_j\lambda_j$,
\[
  \mathrm{ED}=\frac{\bigl(\sum_k\lambda_k\bigr)^2}{\sum_k\lambda_k^2}=\frac{1}{\sum_k\pi_k^2}\in[1,p],
\]
the ``effective number of independent directions'' (the inverse participation ratio, familiar from
localisation physics): $\mathrm{ED}=1$ for a rank-one matrix such as the $4\times4$ matrix of ones, $2$ for two
decoupled blocks of perfectly correlated variables, $p$ for the identity. For the 2001--2005 Treasury covariance
PCA, ED $\approx1.37$.

\paragraph{Out-of-sample variance explained} \ts{24}{9}{20:07}. In finance PCA is calibrated on the past and
applied to the future. Let $W$ be the loadings estimated on a sample $X$ and let a new sample $Y$ have
covariance $\Sigma^Y=W^Y\Lambda^Y(W^Y)\T$. The PCs defined by $W$ have, in the new sample, covariance matrix
\begin{equation}\label{ch07:eq-oos}\begin{gathered}
  W\T\Sigma^YW=D\T\Lambda^YD,\qquad D=(W^Y)\T W\ \text{(orthogonal)},\\
  \text{share of PC }k=\frac{(D\T\Lambda^YD)_{kk}}{\tr\Lambda^Y}=\frac{\sum_jD_{jk}^2\lambda^Y_j}{\sum_j\lambda^Y_j}.\end{gathered}
\end{equation}
Since $\sum_jD_{jk}^2=1$, the out-of-sample variance of calibrated PC~$k$ is a weighted average of the new
eigenvalues: if the directions have rotated ($D\neq I$), variance leaks from the top factors to lower ones, and
by \cref{ch07:prop-best} the first $K$ calibrated PCs can never explain more than the first $K$ PCs of the new
sample itself. The off-diagonal entries of $D\T\Lambda^YD$ show that calibrated PCs are no longer uncorrelated
out of sample \ts{24}{9}{53:42}: in the Treasury example the correlations between next-day PC1, PC2, PC3 moves
are a few percent (slide 19: $7.1\%$ between PC1 and PC2), small enough for the factors to be called stable
\ts{24}{9}{54:47}.

\begin{coursenote}[A verbal slip]
At \ts{24}{9}{19:07} the speaker says that the ``product of eigenvalues'' is the variance of the data set; the
slide (and \cref{ch07:thm-pca}(iv)) has the correct statement: total variance $=\tr\Sigma=$ \emph{sum} of the
eigenvalues. (The product is $\det\Sigma$, the ``generalised variance''.)
\end{coursenote}

\begin{exercise}[not from the course]\label{ch07:ex-oos}
In \eqref{ch07:eq-oos} show that $\sum_k(D\T\Lambda^YD)_{kk}=\tr\Lambda^Y$ and that, for every $K$,
$\sum_{k\le K}(W\T\Sigma^YW)_{kk}\le\sum_{k\le K}\lambda^Y_k$. When does equality hold for $K=1$?
\end{exercise}

\subsection{PCA on financial time series}\label{ch07:sec-ts}
What is special about finance is that the data are \emph{time series} \ts{24}{9}{27:32}: each day is one
observation and the dimension is the number of variables observed that day. Consequences \ts{24}{9}{28:32}:
\begin{itemize}
  \item \textbf{Model changes, not levels} \ts{24}{9}{29:32}: levels (prices, yields) are non-stationary and not
        centred, while P\&L depends on changes after the bet is placed.
  \item \textbf{Hyperparameters} \ts{24}{9}{31:40}: the length of the look-back window, the weighting of old
        versus recent observations (rolling window, exponential weights), the set of variables included.
        More history or more variables is not always better: the marginal data may be noisier or less
        representative of where the bets will be placed.
  \item \textbf{Sampling frequency} \ts{24}{9}{33:46}: daily data are the default (long, clean histories);
        intraday data add observations but also microstructure noise. Volatility scales like
        $\sqrt{\text{time}}$, so the sampling interval sets its units: annualise daily numbers with $\sqrt{252}$
        before comparing \ts{24}{9}{34:51} (\cref{ch:volatility}).
  \item \textbf{Rolling PCA}: re-estimating every day turns loadings and eigenvalues into time series.
        Slowly varying loadings are reassuring; long stable periods interrupted by jumps signal regime
        changes \ts{24}{9}{30:33}.
\end{itemize}

% ------------------------------------------------------------------
\section{The Treasury yield curve: level, slope and curvature}\label{ch07:sec-yield}

\paragraph{The market.} The US Treasury funds the debt by regular auctions of bills and of 2, 3, 5, 7, 10,
(since 2020) 20 and 30-year notes and bonds \ts{24}{9}{38:00}; the most liquid points also trade as CME
futures (TU $\approx1.5$y, FV $\approx4$y, TY $\approx7$y, UXY $\approx10$y, US $\approx15$y, WN $\approx25$y equivalent maturities). The
market is the deepest in the world and allows large leverage \ts{24}{9}{35:55}; the Federal Reserve publishes
free, high-quality fitted curves \ts{24}{9}{36:58}. All these securities are loans to the same borrower, and
their prices depend on the expected path of the short rate set by the Fed: yield changes of different
maturities are therefore very highly correlated \ts{24}{9}{39:04}. In the 2024 data (par yields, daily changes)
the correlation is $98\%$ between 2y and 3y, $81\%$ between 2y and 10y, $68\%$ between 2y and 30y and $92\%$
between 10y and 30y (slide 17) \ts{24}{9}{40:07}: ``whenever you see a correlation matrix with high
concentrations of high correlations, it is a prime candidate for PCA''. The 2013 case study finds the
same pattern for 2001--2005: correlations above $0.9$ near the diagonal, decaying to $0.24$ between 3 months
and 20 years \ts{13}{15}{1:15:26}. Typical daily volatilities are 4--7 bp \ts{13}{15}{1:14:24}.

\paragraph{The yield curve} (\cref{sec:bonds}, \cref{fig:yieldcurve}) is usually upward sloping but can be
inverted, as in 2023--2024 \ts{24}{9}{41:11}; the Fed sets its short end (the 50 bp cut of September 2024
moved the leftmost point), the rest is market expectation \ts{24}{9}{43:12}. A movie of the curve would show
``strings flopping around'': mostly moving up and down, a little tilting, a little bending \ts{24}{9}{42:12}.

\paragraph{The three factors.} Running PCA on daily yield changes gives, in every period and data set
(\cref{ch07:fig-loadings}, \cref{ch07:tab-share}):
\begin{itemize}
  \item \textbf{PC1 = level}: all loadings of the same sign; yields move up or down together. It explains
        roughly 85--95\% of the variance \ts{24}{9}{47:24}, \ts{13}{15}{1:19:47}.
  \item \textbf{PC2 = slope}: loadings negative at the short end and positive at the long end (or vice versa):
        steepening versus flattening; 5--10\% \ts{24}{9}{48:24}.
  \item \textbf{PC3 = curvature}: same sign at both ends, opposite sign in the ``belly'': a second difference;
        1--3\% \ts{13}{15}{1:20:51}.
\end{itemize}
The top three factors explain 97--99\% of the daily variance.

\begin{figure}[ht]
\centering
\begin{tikzpicture}
\begin{axis}[width=0.52\linewidth,height=5.4cm,xmode=log,xmin=0.2,xmax=25,ymin=-0.6,ymax=0.6,
  xtick={0.25,0.5,1,2,5,10,20},xticklabels={3m,6m,1y,2y,5y,10y,20y},
  xlabel={tenor},ylabel={loading},title={\small (a) CMT yields, 2001--2005},
  grid=major,grid style={black!10},mark size=1.3pt,tick label style={font=\scriptsize},
  legend style={font=\scriptsize,at={(0.03,0.03)},anchor=south west}]
  \addplot[thick,defcol,mark=*] coordinates {(0.25,0.1088) (0.5,0.1581) (1,0.2531) (2,0.3905) (3,0.4197)
    (5,0.4206) (7,0.4009) (10,0.3697) (20,0.3101)};
  \addlegendentry{PC1 level}
  \addplot[thick,thmcol,mark=square*] coordinates {(0.25,-0.5225) (0.5,-0.4817) (1,-0.4107) (2,-0.2071)
    (3,-0.0873) (5,0.1082) (7,0.2289) (10,0.2837) (20,0.3621)};
  \addlegendentry{PC2 slope}
  \addplot[thick,intcol,mark=triangle*] coordinates {(0.25,0.5434) (0.5,0.2508) (1,-0.0514) (2,-0.4750)
    (3,-0.4022) (5,-0.0311) (7,0.1475) (10,0.2710) (20,0.3944)};
  \addlegendentry{PC3 curvature}
  \addplot[black!40,forget plot] coordinates {(0.2,0) (25,0)};
\end{axis}
\end{tikzpicture}\hfill
\begin{tikzpicture}
\begin{axis}[width=0.52\linewidth,height=5.4cm,xmin=0,xmax=31,ymin=-0.45,ymax=0.45,
  xlabel={maturity (years)},title={\small (b) par yields, Sep 2021--Aug 2023},
  grid=major,grid style={black!10},mark size=0.9pt,tick label style={font=\scriptsize}]
  \addplot[thick,defcol,mark=*] coordinates {(1,0.127) (2,0.185) (3,0.202) (4,0.208) (5,0.211) (6,0.212)
    (7,0.211) (8,0.209) (9,0.206) (10,0.203) (11,0.199) (12,0.195) (13,0.192) (14,0.188) (15,0.185) (16,0.181)
    (17,0.178) (18,0.176) (19,0.173) (20,0.171) (21,0.169) (22,0.167) (23,0.166) (24,0.164) (25,0.163)
    (26,0.162) (27,0.161) (28,0.160) (29,0.160) (30,0.159)};
  \addplot[thick,thmcol,mark=square*] coordinates {(1,-0.357) (2,-0.420) (3,-0.373) (4,-0.307) (5,-0.241)
    (6,-0.181) (7,-0.130) (8,-0.086) (9,-0.049) (10,-0.017) (11,0.010) (12,0.033) (13,0.052) (14,0.069)
    (15,0.084) (16,0.096) (17,0.107) (18,0.117) (19,0.126) (20,0.133) (21,0.140) (22,0.146) (23,0.151)
    (24,0.156) (25,0.161) (26,0.165) (27,0.168) (28,0.172) (29,0.175) (30,0.177)};
  \addplot[thick,intcol,mark=triangle*] coordinates {(1,0.273) (2,0.286) (3,0.260) (4,0.115) (5,-0.050)
    (6,-0.185) (7,-0.274) (8,-0.317) (9,-0.321) (10,-0.298) (11,-0.257) (12,-0.204) (13,-0.146) (14,-0.088)
    (15,-0.033) (16,0.017) (17,0.061) (18,0.097) (19,0.125) (20,0.146) (21,0.160) (22,0.166) (23,0.166)
    (24,0.160) (25,0.149) (26,0.133) (27,0.112) (28,0.087) (29,0.059) (30,0.028)};
  \addplot[black!40] coordinates {(0,0) (31,0)};
\end{axis}
\end{tikzpicture}
\caption{Loadings of the first three PCs of daily yield changes (covariance PCA). (a) 9 FRED constant-maturity
yields, 2001--2005, \texttt{princomp} output of 2013 Case Study~7. (b) Fed par yields at maturities 1--30 years,
504-day window ending 1 September 2023, computed from \file{fedbondyielddata_20230911.csv} as in
\file{feddataanalysis.ipynb} (signs fixed so that the 20--30y loadings sum to a positive number); PC1--PC3
explain 89.2\%, 8.3\% and 1.0\%. Compare slide 18 of \file{lec09}.}\label{ch07:fig-loadings}
\end{figure}

\begin{table}[ht]
\centering\small
\caption{Fraction of the variance of daily yield changes explained by the first three PCs.}\label{ch07:tab-share}
\setlength{\tabcolsep}{3pt}\begin{tabular}{@{}l l rrrr@{}}
\toprule
Data & Source & PC1 & PC2 & PC3 & top 3\\
\midrule
9 CMT yields, 6 days of 2001 & 2016 reading & 75.4 & 22.6 & 1.9 & 99.9\\
9 CMT yields, 2001--2005, covariance & Case Study 7 & 84.9 & 9.2 & 2.9 & 97.0\\
9 CMT yields, 2001--2005, correlation & PS7 (2013), PS3 (2024) & 77.1 & 15.7 & 3.5 & 96.3\\
9 CMT yields, 2009--2013, covariance & Case Study 7 & 88.6 & 6.8 & 1.6 & 96.9\\
30 par yields, 2-year window to Sep 2023 & computed from the 2024 data & 89.2 & 8.3 & 1.0 & 98.5\\
\bottomrule
\end{tabular}
\end{table}

\begin{intuition}[Modes of a vibrating string]
Level, slope and curvature are the analogues of the first normal modes of a string: 0, 1 and 2 sign changes
(nodes). This is partly a mathematical consequence of any smooth, strongly and positively correlated
curve (not discussed in class): for the correlation kernel $\rho_{ij}=e^{-\kappa|\tau_i-\tau_j|}$ on the same nine
tenors with $\kappa=0.1$, the first four eigenvectors have exactly 0, 1, 2, 3 sign changes and explain
66\%, 15\%, 9\%, 4\% (my computation; such kernels are ``oscillation matrices'' in the sense of
Gantmacher and Krein). The economics lies in the \emph{details}: the precise shapes, the explained
variances and how they change over time. The same holds for any Karhunen--Lo\`eve expansion of a random
field with a smooth covariance.
\end{intuition}

\paragraph{Reading the loadings as regime indicators} \ts{24}{9}{49:29}. The PC1 loading is small at the very
short end and levels off beyond a few years: short rates are anchored by the Fed, long rates by
expectations. The maturity at which PC1 ``levels off'' (in the notebook: the first tenor with loading above
0.15) was 1--3 years until the financial crisis, rose to 5--7 years when the Fed promised to keep rates
low for long (quantitative easing, forward guidance), fell back, jumped again with the COVID response and
decreased during the 2022 tightening \ts{24}{9}{50:33}, \ts{24}{9}{1:14:43}. Similarly the 2-year PC1 loading
collapsed after 2012 (the Fed was not moving) and rebounded when policy became uncertain; the
\emph{location of the belly} of PC3 (usually near 10y in the US) shifts with regimes \ts{24}{9}{51:34}.
Andreev's macro reading of curvature \ts{24}{9}{1:17:51}: the front end is driven by the central bank, the
intermediate part by inflation expectations, the long end by growth expectations; investors also differ by
segment (speculators and short-term funding at 2y, corporate issuance at 5--10y, pension funds and insurers
matching long liabilities at the long end) \ts{24}{9}{1:18:57}. In an emerging market such as Mexico the belly
sits nearer 5 years, reflecting shorter investment horizons. Once the level and slope are hedged out, these
small but structured features emerge from ``the big noise of rates going up and down'' \ts{24}{9}{1:19:59}.

\paragraph{The last PC is interesting too} \ts{13}{15}{1:21:57}. In the 2013 case study PC9, with the smallest
variance (0.8 bp/day), is essentially the 7-year yield against a weighted average of the 5- and 10-year: the
most stable linear combination of yields, i.e.\ the best hedge of the 7-year point with its neighbours.
Cumulating the scores shows the history of the factors \ts{13}{15}{1:23:00}: over 2001--2005 cumulative PC1
fell from 0 to about $-8$ and came back (rates fell, then rose), cumulative PC2 rose to about 6 and fell back
(the curve steepened, then flattened), curvature moved much less \ts{13}{15}{1:24:00}.

\begin{finance}[Nelson--Siegel--Svensson curves (not discussed in class)]
The Fed file \file{fedbondyielddata_20230911.csv} is the Gurkaynak--Sack--Wright data set: daily zero-coupon
yields (\texttt{SVENY01}--\texttt{30}, continuously compounded), par yields (\texttt{SVENPY}), instantaneous and
one-year forward rates (\texttt{SVENF}, \texttt{SVEN1F}) from 1961, and the parameters \texttt{BETA0}--\texttt{BETA3},
\texttt{TAU1}, \texttt{TAU2} of the fitted Svensson curve
\[
  y(n)=\beta_0+\beta_1\frac{1-e^{-n/\tau_1}}{n/\tau_1}
        +\beta_2\Bigl[\frac{1-e^{-n/\tau_1}}{n/\tau_1}-e^{-n/\tau_1}\Bigr]
        +\beta_3\Bigl[\frac{1-e^{-n/\tau_2}}{n/\tau_2}-e^{-n/\tau_2}\Bigr]
\]
(I checked that this reproduces \texttt{SVENY} to four decimals). The three Nelson--Siegel basis functions
are a parametric version of level ($\beta_0$), slope ($\beta_1$, which dies out with maturity) and curvature
($\beta_2$, a hump); $\beta_3$ adds a second hump. Because the published curves are smooth fits, their PCA
loadings are smoother than those of raw bond yields.
\end{finance}

\begin{casestudy}[Reading \file{pset_analysis} and its R code \file{pset_rcode}: five years of yields]
\textbf{Goal.} Reproduce the six-day example and extend it to 2001--2005. \textbf{Data.} FRED
constant-maturity yields \texttt{DGS3MO}, \dots, \texttt{DGS20} (script \file{fm_freddata.r}), 1247 daily changes.
\textbf{Method.} \texttt{svd()} of the centred days$\times$tenors matrix, $\hat\lambda_k=d_k^2/(n-1)$, bar plots of
loadings, daily and cumulative PC1/PC2 series, \texttt{pairs()} and interactive 3D plots of (PC1, PC2, PC3);
then the built-in \texttt{princomp()} and \texttt{screeplot()}, which give the same result up to the signs of
the loadings (\cref{ch07:rem-sign}). \textbf{Results.} Total variance 0.0306 (\%$^2$); PC1/PC2/PC3 explain
84.8/9.3/2.9\%; PC1 loadings 0.11 (3m) rising to 0.42 (3y--5y) and 0.31 (20y); cumulative PC1 falls
over the first half of the period and recovers, cumulative PC2 (the spread) rises then falls.
\textbf{Lesson.} Five days already reveal the structure that five years confirm.
\emph{Two small inconsistencies:} in the five-year section the text says $n$ is ``the number of columns
(days)'' but the days are now rows ($\hat\lambda_k=d_k^2/(\texttt{NROW(X)}-1)$ in the code); and the text calls PC2
``shorter minus longer'' although in this SVD its loadings are negative at the short end, i.e.\ long minus
short. Signs are arbitrary.
\end{casestudy}

\begin{casestudy}[\file{feddataanalysis.ipynb}: Andreev's notebook on the Fed data]
\textbf{Data.} Par yields \texttt{SVENPY01}--\texttt{30} from the first day with a 30-year yield (25 Nov 1985) to
1 Sep 2023, 9430 days, daily differences. \textbf{Method} (Python/pandas): 504-day (2-year) rolling covariance of
yield changes, \texttt{numpy.linalg.eigh} each day, signs normalised so that the 20--30y loadings sum to a
positive number; variance explained per factor (area and log plots); PC1 loading stability for 2, 5, 10, 30y;
``maturity for PC1 leveling'' (first tenor with loading $>0.15$), PC2 pivot and PC3 belly; next-day yield changes
projected on today's loadings (the out-of-sample check \eqref{ch07:eq-oos}); out-of-sample variance explained with
loadings two years old; a PC1/PC2-neutral 5s7s10s fly and a mean-reversion strategy on it.
\textbf{Results.} Correlation matrix of slide 17 (reproduced exactly); out-of-sample correlations between
next-day PC1, PC2, PC3 moves: $4.0\%$ (PC1/PC2), $-0.2\%$ (PC1/PC3), $-9.0\%$ (PC2/PC3) (the slide shows
$7.1\%$, $-1.6\%$, $-7.3\%$ from another run); strategy Sharpe ratio $0.10$ (slide: $0.09$).
\textbf{Lesson.} The last cell repeats the strategy with the signal \emph{not} delayed by one day and gets
spectacular, unrealistic returns: a textbook demonstration of look-ahead bias. Andreev recommends the
notebook as a template for a final project \ts{24}{9}{1:11:29}.
\end{casestudy}

% ------------------------------------------------------------------
\section{PC portfolios: curves, flies and hedging}\label{ch07:sec-pcport}

\subsection{Loadings as portfolio weights}
Why care about PC2 and PC3 if PC1 explains 90\% of the variance? Because they are where relative-value
trading happens \ts{24}{9}{55:49}. The six liquid tenors are so correlated that trading any of them outright
is essentially a bet on PC1 \ts{24}{9}{56:51}; PCA manufactures from them a handful of \emph{uncorrelated}
portfolios, i.e.\ more independent bets \ts{24}{9}{57:54}.

Positions in bonds are measured in \emph{dv01}: \$1 dv01 of a bond is the notional whose value changes by
\$1 for a 1 bp change of its yield (\cref{sec:duration}). Let $h_i$ be the dv01 held in tenor $i$ (positive =
long the bond, which gains when the yield falls). To first order the daily P\&L in dollars is
\[
  \mathrm{PnL}=-\bm h\T\Delta\bm y=-\sum_k(\bm h\T\bm\gamma_k)\,p_k,\qquad \Delta\bm y\ \text{in bp},
\]
using $\Delta\bm y=\sum_kp_k\bm\gamma_k$ (\cref{ch07:thm-pca}(iii), centring neglected). The \emph{exposure} of the portfolio to
PC~$k$ is $e_k=\bm h\T\bm\gamma_k$. Buying ``the PC1 portfolio'' means holding every bond in proportion to its PC1
loading \ts{24}{9}{58:59}; a portfolio with $e_1=0$ is \emph{PC1-neutral}, one with $e_1=e_2=0$ is PC1/PC2-neutral.

\begin{coursenote}[Erratum: the dv01 box on slide 20 of \file{lec09}]
The box defines ``buy \$1 dv01 of bond'' as the notional such that the position \emph{earns} \$1 when the yield
goes \emph{up} 1 bp. A long bond position loses when yields rise; the examples on the same slide (a steepener
sells the 10-year bond, a belly fly buys the 10-year) use the standard convention adopted here. The notebook
avoids the issue by computing returns directly in yield space ($+$weights$\cdot\Delta\bm y$), i.e.\ as positions in
yields rather than in bonds; any sign convention works if used consistently.
\end{coursenote}

\paragraph{PC-neutral curves} (``PC1 residuals''). Given two tenors $a,b$, the portfolio $\bm h=\bm e_a-\alpha\bm e_b$
(\$1 dv01 of $a$ against $\alpha$ dv01 of $b$) is PC1-neutral iff
\begin{equation}\label{ch07:eq-curve}
  \gamma_{a1}-\alpha\gamma_{b1}=0\iff\alpha=\frac{\gamma_{a1}}{\gamma_{b1}},
\end{equation}
where $\gamma_{ik}$ is the loading of tenor $i$ on PC~$k$. Example: the 2s10s \emph{steepener}, short \$1 dv01 of
10-year, long $\alpha$ dv01 of 2-year: it profits when the curve steepens and is insensitive to parallel (PC1) moves.
Such curve trades exist in the market as standard products \ts{24}{9}{1:00:00}.

\paragraph{PC-neutral flies} (``PC2 residuals''). Given three tenors, $\bm h=\alpha_a\bm e_a+\alpha_b\bm e_b-\bm e_c$ is neutral
to PC1 and PC2 iff
\begin{equation}\label{ch07:eq-fly}
  \begin{pmatrix}\gamma_{a1}&\gamma_{b1}\\ \gamma_{a2}&\gamma_{b2}\end{pmatrix}
  \begin{pmatrix}\alpha_a\\ \alpha_b\end{pmatrix}=\begin{pmatrix}\gamma_{c1}\\ \gamma_{c2}\end{pmatrix}.
\end{equation}
Example: the 2s10s30s ``belly'' fly buys \$1 dv01 of 10-year and sells $\alpha_2$ of 2-year and $\alpha_{30}$ of
30-year: a pure bet on curvature.

\begin{coursenote}[Erratum: slide 20 of \file{lec09}]
The fly system on the slide has $W_{b,0}$ in the lower-right entry; it must be $W_{b,1}$ (the PC2 loading of
tenor $b$), as in \eqref{ch07:eq-fly}. (The slides index PCs from 0.)
\end{coursenote}

\begin{example}[Hedge ratios and volatilities from the 2001--2005 PCA]\label{ch07:ex-hedges}
Rebuild $\Sigma=\Gamma\Lambda\Gamma\T$ from the loadings and standard deviations of 2013 Case Study~7 (this reproduces the
case-study volatilities and correlations to three digits) and measure risk as the standard deviation of daily P\&L
per \$1 dv01 of the reference bond:
\begin{center}\small
\begin{tabular}{@{}l l c c@{}}
\toprule
Trade (PC-neutrality) & Weights (dv01) & vol (\$/day) & outright vol\\
\midrule
2s10s curve (PC1) & $-1$ of 10y, $+0.947$ of 2y & 3.5 & 6.3 (10y)\\
2s5s10s fly (PC1, PC2) & $+1$ of 5y, $-0.423$ of 2y, $-0.690$ of 10y & 1.5 & 6.9 (5y)\\
5s7s10s fly (PC1, PC2) & $+1$ of 7y, $-0.367$ of 5y, $-0.667$ of 10y & 1.0 & 6.7 (7y)\\
\bottomrule
\end{tabular}
\end{center}
The 5s7s10s fly is the notebook's trade and is close to PC9 of the case study. To run it at the risk of an
outright 7-year position one needs about $6.4$ times the dv01 on the belly, plus the offsetting wings.
\end{example}

\subsection{Leverage, and why higher PCs are the interesting ones}
A PC2 portfolio moves perhaps a third as much as the PC1 portfolio per unit notional; scaling it up three
times gives the same daily standard deviation and a return stream uncorrelated with the market direction
\ts{24}{9}{1:01:03}. Far from being irrelevant, the higher PCs are ``where a lot of the money is made''
\ts{24}{9}{1:03:07}; beyond PC3, however, the leverage needed becomes too expensive. The strategy requires a
liquid market and cheap financing (repo); regulators (the CFTC and others) watch the resulting leveraged
long/short Treasury positions as a possible systemic risk \ts{24}{9}{1:02:05}. In equities PC1 is much less
dominant, so less leverage is needed; market-neutral long/short equity funds are precisely ``hedging out PC1''
\ts{24}{9}{1:04:09}.

\begin{center}\small
\begin{tabularx}{\linewidth}{@{}X X@{}}
\toprule
\textbf{Pros of PC-residual portfolios} & \textbf{Cons} (slide 21)\\
\midrule
Bets unrelated to the main market move (PC1). & Transaction cost per unit of volatility is much higher.\\
More effectively uncorrelated bets (diversity). & Weights are model-dependent and change over time.\\
Low volatility per notional: a mispricing may have a better signal-to-noise ratio. &
Much larger notionals, i.e.\ leverage and balance sheet; financing (overnight repo) can dry up in stress.\\
\bottomrule
\end{tabularx}
\end{center}

\paragraph{Model risk.} The hedge ratios are only as good as the loadings. When COVID hit and the Fed made an
emergency cut, the curve dynamics changed overnight \ts{24}{9}{1:07:21}; a leveraged PC2 position computed with
stale loadings is no longer PC1-neutral and its P\&L is dominated by unintended level risk, while there is no
history yet to recalibrate on \ts{24}{9}{1:08:21}. Possible answers: implied (option) markets, or calibrating
on past periods of a similar regime (e.g.\ all Fed cutting cycles), with its own pitfalls \ts{24}{9}{1:09:25}.
Event days also matter for risk: daily yield volatility is typically 5--6 bp, but the market may price twice
that on the day of a major inflation release \ts{24}{9}{1:06:18}.

\paragraph{Diversity of PC portfolios.} Slide 23 computes the effective dimension of the correlation matrix of
the returns of PC portfolios: PC1 alone 1, PC1+PC2 1.99, PC1+PC2+PC3 2.97 (almost exactly uncorrelated also
out of sample), and 5.28 when the residual portfolios left after removing three PCs are added. Higher-order
residuals also tend to \emph{mean-revert}. The notebook's strategy \ts{24}{9}{1:10:25}: build the PC1/PC2-neutral
5s7s10s fly every day from the current loadings; signal $s$ = z-score of its cumulative return against an
exponentially weighted mean (half-life 100 days) and standard deviation (half-life 60 days), delayed by one day;
position $=-s$. The Sharpe ratio is only about 0.1: the point is the method, not the strategy.

\begin{coursenote}[Differences between slides and speech]
At \ts{24}{9}{1:10:25} the speaker describes the strategy as ``momentum/mean reversion for PC1''; slide 24 and the
notebook apply mean reversion to the PC1/PC2-neutral fly.
\end{coursenote}

\subsection{Hedging a portfolio against PC1}
For an equity portfolio with weights $\bm V$ on $p$ indices (returns, so $\mathrm{PnL}=\bm V\T\bm r$), the exposure to
PC~$n$ is $\bm V\T\bm\gamma_n$. Projecting out the first component,
\[
  \bm V'=\bm V-(\bm V\T\bm\gamma_1)\,\bm\gamma_1,\qquad \bm V'^{\mathsf T}\bm\gamma_1=0,\qquad
  \Var(\bm V'^{\mathsf T}\bm r)=\bm V'^{\mathsf T}\Sigma\bm V'=\bm V\T\Sigma\bm V-\lambda_1(\bm V\T\bm\gamma_1)^2,
\]
the last identity because $(I-\bm\gamma_1\bm\gamma_1\T)\Sigma(I-\bm\gamma_1\bm\gamma_1\T)=\Sigma-\lambda_1\bm\gamma_1\bm\gamma_1\T$. If the expected return
(``alpha'') of the strategy is roughly orthogonal to PC1, hedging keeps the alpha and removes most of the
risk: higher Sharpe ratio.

\begin{coursenote}[Erratum: slide 32 of \file{lec09}]
``The stdev of the resulting portfolio $=V'^T\Sigma V'$'': that is the variance; the standard deviation is
$\sqrt{V'^T\Sigma V'}$.
\end{coursenote}

\begin{casestudy}[European stock indices (slides 26--34; not covered in class \ts{24}{9}{1:12:30})]
\textbf{Data.} Daily closes of AEX, DAX, FTSE MIB, IBEX, OMX, Euro Stoxx 50 and CAC, 4 Jan 2010 to 16 Sep 2016
(same time zone and currency), rebased to 100 EUR and converted to demeaned log returns. Correlations 72--98\%.
\textbf{PCA.} PC1 explains 89\% with nearly equal loadings: a ``Euro stocks'' factor. PC2 (about 5\%, half of the
rest) separates North (Netherlands, Germany, Sweden: positive) from South (Italy, Spain: negative), France near
zero. Covariance and correlation PCA agree. Rolling 6-month PCA: top three factors explain 90--98\% throughout,
PC1 loadings stable, PC2 fairly stable. \textbf{Hedging example} (thesis: German stocks outperform, 2010, target
100m EUR annualised volatility): buying the DAX (vol 20\%) needs 500m and returned 75m (Sharpe 0.75); the
PC1-neutral portfolio $\bm e_{\mathrm{DAX}}-0.377\,\bm\gamma_1$ (weights $+0.86$ DAX, about $-0.14$ each other index) has vol
6.7\%, needs 1{,}500m and returned about 230m (Sharpe 2.3). \textbf{Other uses:} North-versus-South trades;
Frankfurt opens before Madrid, so a PCA model can price Spain by proxy just before its open.
\emph{Slide typos:} the rebasing date is 4 Jan 2010 (not ``1/4/2016''), and the legend entry ``Swiss'' on slides 27
and 31 is the Euro Stoxx 50.
\end{casestudy}

\begin{finance}[Limitations (slide 35)]
PCA always gives an answer, so test robustness to the window and to the set of variables; it works best on data
close to multivariate normal; and it implies no causality. To predict a dependent variable, use regression
(\cref{ch:regression}). Andreev's closing advice \ts{24}{9}{1:22:03}: data that are hard to get (illiquid markets,
emerging-market curves) are where fewer competitors look, hence where quantitative modelling may still pay.
\end{finance}

% ------------------------------------------------------------------
\section{Linear factor models}\label{ch07:sec-lfm}

\subsection{Definition and covariance structure}
\begin{definition}[Linear factor model]\label{ch07:def-lfm}
For $m$ assets observed over $T$ periods \ts{13}{15}{1:06}, \ts{13}{15}{3:13}:
\begin{equation}\label{ch07:eq-lfm}
  x_{i,t}=\alpha_i+\beta_{1,i}f_{1,t}+\dots+\beta_{K,i}f_{K,t}+\epsilon_{i,t},\qquad
  \bm x_t=\bm\alpha+B\bm f_t+\bm\epsilon_t ,
\end{equation}
with intercepts $\bm\alpha\in\R^m$, \emph{factor loadings} $B\in\R^{m\times K}$ (row $i$ is $\bm\beta_i\T$), \emph{common factors}
$\bm f_t\in\R^K$, $K\ll m$, and \emph{specific factors} $\bm\epsilon_t$. Assumptions \ts{13}{15}{5:19}:
$\{\bm f_t\}$ is covariance stationary with $\E[\bm f_t]=\bm\mu_f$, $\Cov\bm f_t=\Omega_f$; $\{\bm\epsilon_t\}$ is white noise with
$\E[\bm\epsilon_t]=\bm0$, $\Cov\bm\epsilon_t=\Psi=\diag(\sigma_1^2,\dots,\sigma_m^2)$, $\Cov(\bm\epsilon_t,\bm\epsilon_{t'})=0$ for $t\neq t'$; and
$\Cov(\bm f_t,\bm\epsilon_{t'})=0$ for all $t,t'$.
\end{definition}

The key assumption is that $\Psi$ is \emph{diagonal}: all co-movement between assets goes through the factors.

\begin{proposition}[Moments of the factor model]\label{ch07:prop-lfm}
$\E[\bm x_t\mid\bm f_t]=\bm\alpha+B\bm f_t$, $\Cov[\bm x_t\mid\bm f_t]=\Psi$, and unconditionally
\begin{equation}\label{ch07:eq-lfmcov}
  \E[\bm x_t]=\bm\alpha+B\bm\mu_f,\qquad \Sigma_x=\Cov[\bm x_t]=B\,\Omega_f\,B\T+\Psi .
\end{equation}
\end{proposition}
\begin{proof}
By the law of total covariance, $\Cov[\bm x_t]=\E\bigl[\Cov(\bm x_t\mid\bm f_t)\bigr]+\Cov\bigl(\E[\bm x_t\mid\bm f_t]\bigr)=\Psi+B\Omega_fB\T$;
directly, $\Cov(B\bm f_t+\bm\epsilon_t)=B\Omega_fB\T+\Psi+2\Cov(B\bm f_t,\bm\epsilon_t)$ and the cross term vanishes \ts{13}{15}{7:25}.
\end{proof}

\paragraph{Parameter count} \ts{13}{15}{8:27}. An unrestricted covariance has $m(m+1)/2$ parameters; the factor
structure has $mK$ (loadings) $+K(K+1)/2$ ($\Omega_f$) $+m$ ($\Psi$) \ts{13}{15}{10:39}. For $m=500$ stocks and $K=10$
factors: $125{,}250$ versus $5{,}555$. Moreover each new asset adds only $K+1$ parameters. This is why factor models
are used for portfolio optimisation, where the (inverse) covariance matrix drives the allocation
(\cref{ch:portfolio}), and for risk management, where scenarios are naturally expressed as factor moves
\ts{13}{15}{16:54}.

\paragraph{Three regression views} (\cref{ch:regression}).
\emph{Cross-sectional}: fix $t$; \eqref{ch07:eq-lfm} is a regression of the $m$ returns on the columns of $B$ with
``coefficients'' $\bm f_t$ \ts{13}{15}{4:18}. \emph{Time series}: fix asset $i$; $\bm x_i=\bm1_T\alpha_i+F\bm\beta_i+\bm\epsilon_i$
with $F\in\R^{T\times K}$ the factor history and $\Cov\bm\epsilon_i=\sigma_i^2I_T$: an ordinary regression satisfying the
Gauss--Markov assumptions \ts{13}{15}{11:41}. \emph{Multivariate}: stacking all assets,
$X=\bm1_T\bm\alpha\T+FB\T+E$ with $X=[\bm x_1|\cdots|\bm x_m]\in\R^{T\times m}$ \ts{13}{15}{12:45} (the slides write $B$ for the
$K\times m$ matrix of the time-series form, the transpose of the cross-sectional $B$).
Whether factors or loadings are observed defines the three families below \ts{13}{15}{0:00}.

\subsection{Identification}\label{ch07:sec-ident}
When the factors are latent, the model is not identified \ts{13}{15}{42:40}: for any invertible $H\in\R^{K\times K}$,
$\bm f_t^*=H\bm f_t$ and $B^*=BH^{-1}$ give the same $\bm x_t$, with $\bm\mu_f^*=H\bm\mu_f$ and $\Omega_f^*=H\Omega_fH\T$
\ts{13}{15}{43:49}. One uses this freedom to normalise \ts{13}{15}{44:50}:
\begin{itemize}
  \item \emph{orthonormal factors}: if $\Omega_f=\Gamma_f\Lambda_f\Gamma_f\T$, take $H=\Lambda_f^{-1/2}\Gamma_f\T$, so that
        \[\Omega_f^*=\Lambda_f^{-1/2}\Gamma_f\T\Gamma_f\Lambda_f\Gamma_f\T\Gamma_f\Lambda_f^{-1/2}=I_K;\]
  \item \emph{zero-mean factors}: absorb $B\bm\mu_f$ into the intercept, $\bm\alpha^*=\bm\alpha+B\bm\mu_f$.
\end{itemize}
Then $\Sigma_x=BB\T+\Psi$ \ts{13}{15}{46:00}. An indeterminacy remains: $B\mapsto BQ$ with $Q$ orthogonal leaves $BB\T$
unchanged. This rotational freedom is exploited to make factors interpretable (\cref{ch07:sec-fa}).

\begin{coursenote}[Erratum: slide 22 of \file{lecnote15}]
The slide takes $H=\Gamma\Lambda^{-1/2}$; then $H\Omega_fH\T=\Gamma\Lambda^{-1/2}\Gamma\Lambda\Gamma\T\Lambda^{-1/2}\Gamma\T$, which is not $I_K$ in general.
The correct choice is $H=\Lambda^{-1/2}\Gamma\T$ (or $Q\Lambda^{-1/2}\Gamma\T$ for any orthogonal $Q$).
\end{coursenote}

% ------------------------------------------------------------------
\section{Macroeconomic and fundamental factor models}\label{ch07:sec-macrofund}

\subsection{Observed factors: Sharpe's single-index model}
The simplest factor model \ts{13}{15}{13:46}: excess returns regressed on the excess return $R_{M,t}$ of the
market index,
\[
  x_{i,t}=\alpha_i+\beta_iR_{M,t}+\epsilon_{i,t},\qquad
  \Sigma_x=\sigma_M^2\,\bm\beta\bm\beta\T+\Psi,\quad \sigma_M^2=\Var(R_{M,t}) .
\]
The covariance between two different assets is $\beta_i\beta_j\sigma_M^2$ \ts{13}{15}{14:48}: a rank-one matrix plus a
diagonal, trivial to build even for thousands of securities \ts{13}{15}{15:53}. $\beta_i$ measures the systematic
risk of asset $i$ (the CAPM of \cref{ch:portfolio} adds the equilibrium restriction $\alpha_i=0$).

\paragraph{Estimation} \ts{13}{15}{17:57}. The time-series regression of each asset satisfies Gauss--Markov, so OLS
gives BLUE estimates $(\hat\alpha_i,\hat\beta_i)$ (and the MLE under normality); unbiased residual variances are
$\hat\sigma_i^2=\hat{\bm\epsilon}_i\T\hat{\bm\epsilon}_i/(T-2)$; $\hat\sigma_M^2$ is the sample variance of $R_M$; and
$\hat\Sigma_x=\hat\sigma_M^2\hat{\bm\beta}\hat{\bm\beta}\T+\hat\Psi$.

\begin{intuition}[What PCA sees in a single-index world (not discussed in class)]
If all specific variances are equal, $\Psi=\psi I$, then $\Sigma_x=\sigma_M^2\bm\beta\bm\beta\T+\psi I$ has top eigenvector
$\bm\beta/\|\bm\beta\|$ with eigenvalue $\sigma_M^2\|\bm\beta\|^2+\psi$, and all other eigenvalues equal to $\psi$. PCA recovers the
market factor exactly and finds nothing else (a degenerate remainder, \cref{ch07:rem-sign}). With unequal
$\sigma_i^2$ the recovery is only approximate: this is the difference between PCA and factor analysis
(\cref{ch07:sec-pcafa}).
\end{intuition}

\subsection{Macroeconomic multifactor models}
Any observable series is a candidate factor, judged by how much it contributes to explaining returns and
covariances \ts{13}{15}{19:00}: market risk, inflation (CPI, PPI, commodity prices), industrial production, money
growth, interest rates, housing starts, unemployment. Chen, Roll and Ross (1986) used not the levels but the
\emph{surprises}: unanticipated changes, computed as innovations of a vector autoregression fitted to the macro
variables (\cref{ch:timeseries}) \ts{13}{15}{20:03}, \ts{13}{15}{21:07}. Since $F$ is observed, each asset's
loadings are again time-series OLS estimates, with $\hat\sigma_i^2=\hat{\bm\epsilon}_i\T\hat{\bm\epsilon}_i/(T-K-1)$,
$\hat\Omega_f$ the sample covariance of the factors and $\hat\Sigma_x=\hat B\hat\Omega_f\hat B\T+\hat\Psi$ \ts{13}{15}{22:10}.
If residual variances differ or residuals are correlated, generalised least squares is the efficient
estimator \ts{13}{15}{23:15}.

\begin{coursenote}[Erratum: slide 11 of \file{lecnote15}]
The estimator of the multifactor covariance is written $\hat\sigma_M^2\hat B\hat\Omega_f\hat B'+\hat\Psi$; the factor
$\hat\sigma_M^2$ is a leftover of the single-index slide. Correct: $\hat\Sigma_x=\hat B\hat\Omega_f\hat B\T+\hat\Psi$ (the factor
variances are already inside $\hat\Omega_f$).
\end{coursenote}

\subsection{Fundamental factor models}
Here the \emph{loadings} are observed asset attributes (industry, market capitalisation, dividend yield,
book-to-price, earnings-to-price, \dots) and the factor realisations are estimated. This is the BARRA approach
of Barr Rosenberg \ts{13}{15}{24:17}, still the standard of commercial risk models.

\paragraph{Industry factor model} \ts{13}{15}{29:35}. Split the $m$ stocks into $K$ industries and let
$\beta_{i,k}=1$ if stock $i$ is in industry $k$, $0$ otherwise: $B$ is a known, time-invariant matrix of
indicators \ts{13}{15}{30:37}. Oil stocks move together, utilities (regulated, low volatility) move together
and differently. For each $t$, estimate $\bm f_t$ by the cross-sectional regression of $\bm x_t$ on $B$
\ts{13}{15}{31:40}:
\begin{equation}\label{ch07:eq-fhat}
  \hat{\bm f}_t=(B\T B)^{-1}B\T\bm x_t,\qquad B\T B=\diag(m_1,\dots,m_K),
\end{equation}
where $m_k$ is the number of stocks in industry $k$: $\hat f_{k,t}$ is simply the \emph{average return of industry $k$}
\ts{13}{15}{32:43}. Then $\hat{\bm\epsilon}_t=\bm x_t-B\hat{\bm f}_t$, $\hat\Omega_f$ and $\hat\Psi$ are sample (co)variances of the
$\hat{\bm f}_t$ and of the residuals, and $\hat\Sigma=B\hat\Omega_fB\T+\hat\Psi$ \ts{13}{15}{33:48}.

\paragraph{Heteroscedasticity.} Stocks have different specific variances, so OLS in \eqref{ch07:eq-fhat} is
inefficient; the GLS estimator weights each stock by the inverse of its noise,
\[
  \hat{\bm f}_t=(B\T\Psi^{-1}B)^{-1}B\T\Psi^{-1}\bm x_t ,
\]
with $\Psi$ estimated in a first pass. Kempthorne stresses the analogy \ts{13}{15}{34:53}: mean-variance
optimisation weights assets by $\Sigma^{-1}\bm\mu$ (reward penalised by variance), GLS weights observations by the
inverse noise variance: in both cases noisy inputs must count less.

\begin{coursenote}[Erratum: slide 17 of \file{lecnote15}]
``$\hat\Sigma=B'\hat\Omega_fB+\hat\Psi$'' has the transposes swapped ($B'\hat\Omega_fB$ is not even defined, $B$ being $m\times K$);
correct: $\hat\Sigma=B\hat\Omega_fB\T+\hat\Psi$.
\end{coursenote}

\paragraph{Factor-mimicking portfolios} \ts{13}{15}{35:56}. The estimate $\hat{\bm f}_t=H\bm x_t$ is linear in the returns
\ts{13}{15}{37:00}, so each row of $H$ is a portfolio of the $m$ assets. Normalising the row weights to sum to one
(one unit of capital) gives a tradeable portfolio whose return tracks the factor \ts{13}{15}{39:04}.

\begin{proposition}[Mimicking portfolios have unit exposure]\label{ch07:prop-mimic}
Let $H=(B\T\Psi^{-1}B)^{-1}B\T\Psi^{-1}$ and let $\bm h_k\T$ be its $k$-th row. Then $HB=I_K$, so under \eqref{ch07:eq-lfm}
with $\bm\alpha=\bm0$ the portfolio $\bm h_k$ has return $f_{k,t}+\bm h_k\T\bm\epsilon_t$: unit exposure to factor $k$ and zero to
the others. Among all portfolios with these exposures ($B\T\bm h=\bm e_k$) it has the smallest specific variance
$\bm h\T\Psi\bm h$ (not discussed in class).
\end{proposition}
\begin{proof}
$HB=(B\T\Psi^{-1}B)^{-1}B\T\Psi^{-1}B=I_K$. For the minimum, the Lagrangian $\frac12\bm h\T\Psi\bm h-\bm\ell\T(B\T\bm h-\bm e_k)$ gives
$\bm h=\Psi^{-1}B\bm\ell$ and $B\T\Psi^{-1}B\bm\ell=\bm e_k$, i.e.\ $\bm h=\Psi^{-1}B(B\T\Psi^{-1}B)^{-1}\bm e_k=\bm h_k$ (Gauss--Markov in
portfolio language).
\end{proof}
In the industry model with OLS, $\bm h_k$ is the equal-weighted portfolio of industry $k$.

\paragraph{Fama--French factors} \ts{13}{15}{25:23}. Fama and French turned attributes into factor
\emph{returns} by sorting. At each $t$, rank the stocks by an attribute (market cap for \emph{size},
book-to-market for \emph{value versus growth}) \ts{13}{15}{26:25}, split into quintiles, and define the factor
as the return of the equal-weighted top quintile minus the bottom quintile (a zero-cost hedge portfolio)
\ts{13}{15}{27:32}. Then estimate each stock's loadings on these factors by time-series regressions. In the
published factors the sign is chosen so that the factor has a positive average return, e.g.\ small minus big
(SMB) and high minus low book-to-market (HML) \ts{13}{15}{28:35}. Value stocks are cheap relative to book
equity; growth stocks are expensive because the price anticipates future growth.

% ------------------------------------------------------------------
\section{Statistical factor models}\label{ch07:sec-stat}
Now neither factors nor loadings are observed: only the returns $\bm x_1,\dots,\bm x_T$ \ts{13}{15}{40:18}. Both
methods below model the sample covariance matrix.

\subsection{Factor analysis by maximum likelihood}\label{ch07:sec-fa}
Take the normalised model with Gaussian variables \ts{13}{15}{47:02}:
$\bm f_t\sim N_K(\bm0,I_K)$, $\bm\epsilon_t\sim N_m(\bm0,\Psi)$, independent over time, so that
$\bm x_t\sim N_m(\bm\alpha,\Sigma_x)$ i.i.d.\ with $\Sigma_x=BB\T+\Psi$ \ts{13}{15}{48:05}. The log-likelihood is
\ts{13}{15}{49:05}
\begin{equation}\label{ch07:eq-loglik}
  \ell(\bm\alpha,\Sigma_x)=-\frac{Tm}{2}\ln(2\pi)-\frac T2\ln\det\Sigma_x-\frac12\sum_{t=1}^T(\bm x_t-\bm\alpha)\T\Sigma_x^{-1}(\bm x_t-\bm\alpha),
\end{equation}
to be maximised over $\bm\alpha,B,\Psi$ subject to $\Sigma_x=BB\T+\Psi$. The maximiser of $\bm\alpha$ is $\bar{\bm x}$, and then
$\ell=-\frac T2\bigl[m\ln2\pi+\ln\det\Sigma_x+\tr(\Sigma_x^{-1}S)\bigr]$ with $S$ the (divisor-$T$) sample covariance.

\begin{coursenote}[Erratum: slide 24 of \file{lecnote15}]
The log-likelihood is written with $-\frac{TK}{2}\log(2\pi)-\frac K2\log|\Sigma|$; taking the logarithm of the likelihood on
slide 23 gives $-\frac{Tm}{2}\ln(2\pi)-\frac T2\ln|\Sigma_x|$, as in \eqref{ch07:eq-loglik}.
\end{coursenote}

\paragraph{The EM algorithm} \ts{13}{15}{51:06}. The likelihood is complicated because the factors are hidden;
if they were observed, estimating $B$ and $\Psi$ would be ordinary regression. The EM algorithm (Dempster, Laird
and Rubin 1977) alternates the two: estimate the hidden variables given the parameters (E-step), re-estimate the
parameters as if the hidden variables were known (M-step); each iteration increases the likelihood
\ts{13}{15}{52:10}. For centred data (not discussed in class; I verified numerically that the likelihood increases
monotonically):
\begin{align*}
  \text{E: }\ & \bm\beta=B\T\Sigma_x^{-1},\quad \E[\bm f_t\mid\bm x_t]=\bm\beta\bm x_t,\quad
     \E[\bm f_t\bm f_t\T\mid\bm x_t]=I-\bm\beta B+\bm\beta\bm x_t\bm x_t\T\bm\beta\T;\\
  \text{M: }\ & B_{\rm new}=S\bm\beta\T\bigl(I-\bm\beta B+\bm\beta S\bm\beta\T\bigr)^{-1},\quad
     \Psi_{\rm new}=\diag\bigl(S-B_{\rm new}\bm\beta S\bigr).
\end{align*}
Time dependence can be added through the linear state-space framework of the Kalman filter
(\cref{ch:timeseries}).

\paragraph{Factor scores} \ts{13}{15}{53:16}. With $\hat{\bm\alpha},\hat B,\hat\Psi$ in hand, estimate each $\bm f_t$ by the
cross-sectional GLS regression of $\bm x_t-\hat{\bm\alpha}$ on $\hat B$:
$\hat{\bm f}_t=(\hat B\T\hat\Psi^{-1}\hat B)^{-1}\hat B\T\hat\Psi^{-1}(\bm x_t-\hat{\bm\alpha})$. These estimated factor series can feed other
models, e.g.\ statistical factors of commodity returns used to explain stocks \ts{13}{15}{54:17}, and can be
rescaled into factor-mimicking portfolios.

\paragraph{Rotations and interpretation} \ts{13}{15}{56:23}. Factor analysis was born in psychometrics:
Spearman's analysis of intelligence tests, where one looks for the ``dimensions of intelligence''. The
orthogonal indeterminacy $B\mapsto BQ$ is used to choose rotated loadings that are easier to interpret (e.g.\ the
varimax criterion, which pushes each loading towards $0$ or $\pm$ large) \ts{13}{15}{57:30}.

\paragraph{How many factors? The likelihood-ratio test} \ts{13}{15}{58:36}. Compare the unrestricted MLE
($\tilde{\bm\alpha}=\bar{\bm x}$, $\tilde\Sigma=S$) with the $K$-factor MLE:
\[\begin{gathered}
  \mathrm{LR}(K)=2\bigl[\ell(\tilde{\bm\alpha},\tilde\Sigma)-\ell(\hat{\bm\alpha},\hat B\hat B\T+\hat\Psi)\bigr]\ \approx\ \chi^2_\nu,\\
  \nu=\frac{m(m+1)}2-\Bigl[mK+m-\frac{K(K-1)}2\Bigr]=\frac{(m-K)^2-(m+K)}2 ,
\end{gathered}\]
where the $K$-factor model has $mK+m$ parameters minus $K(K-1)/2$ for the rotational freedom. $H_0$ ($K$ factors
suffice) is rejected for large $\mathrm{LR}$. For $m=9$: $\nu=12,6,1$ for $K=3,4,5$, and $\nu<0$ for $K=6$, so at most five
factors can be fitted (R's \texttt{factanal} also applies Bartlett's small-sample correction). Unlike PCA, the ML
solution is scale-equivariant: analysing covariances or correlations gives equivalent answers (not discussed in
class).

\subsection{PCA as an approximate factor model; the principal factor method}\label{ch07:sec-pcafa}
Partition the PCA of $\Sigma_x$ into the first $K$ and the remaining $m-K$ components \ts{13}{15}{1:03:46}:
\[
  \bm x=\bm\alpha+\Gamma_1\bm p_1+\Gamma_2\bm p_2=\bm\alpha+B\bm f+\bm\epsilon,\qquad B=\Gamma_1,\ \bm f=\bm p_1,\ \bm\epsilon=\Gamma_2\bm p_2 .
\]
This looks like a factor model (and $B=\Gamma_1\Lambda_1^{1/2}$, $\bm f=\Lambda_1^{-1/2}\bm p_1$ gives orthonormal factors), but
$\Cov\bm\epsilon=\Gamma_2\Lambda_2\Gamma_2\T$ is in general \emph{not diagonal} \ts{13}{15}{1:04:59}: the ``specific'' parts are
correlated across assets. PCA is an approximate factor model, ``principal components'', not ``principal
factors''. The \emph{principal factor method} fixes this by iterating PCA on a reduced covariance:
\begin{enumerate}[1.]
  \item PCA of the sample covariance, $\hat\Sigma_x=\hat\Gamma\hat\Lambda\hat\Gamma\T$;
  \item initial estimates $\tilde B_0=\hat\Gamma_{(K)}\hat\Lambda_{(K)}^{1/2}$ (first $K$ columns and the leading $K\times K$ block),
        $\tilde\Psi_0=\diag(\hat\Sigma_x)-\diag(\tilde B_0\tilde B_0\T)$;
  \item eigen-decompose the reduced matrix $\hat\Sigma_x-\tilde\Psi_0$ (which under the model equals $BB\T$, of rank $K$) and
        repeat step 2 with its eigenvectors and eigenvalues, obtaining $\tilde B_1,\tilde\Psi_1$;
  \item iterate until $\tilde\Psi_s$ stops changing; use the last estimates.
\end{enumerate}

\begin{center}\small
\begin{tabularx}{\linewidth}{@{}l X X@{}}
\toprule
 & \textbf{PCA} & \textbf{Factor analysis (ML)}\\
\midrule
Model & none: exact rotation of the data & $\Sigma_x=BB\T+\Psi$, $\Psi$ diagonal, Gaussian likelihood\\
Output & ordered orthogonal directions, variances & loadings up to rotation, specific variances\\
Residuals & correlated across variables & uncorrelated by assumption\\
Scale & covariance and correlation PCA differ & scale-equivariant\\
Number of factors & scree plot, variance explained & likelihood-ratio test\\
Computation & one SVD & iterative (EM)\\
\bottomrule
\end{tabularx}
\end{center}

\begin{coursenote}[Erratum: slide 32 of \file{lecnote15}]
For the SVD $X^*=VDU'$ of the $m\times T$ centred matrix, $U$ must be $T\times m$ with orthonormal columns, $U'U=I_m$
(the slide says ``$U$: $(m\times T)$ row-orthonormal, $UV'=I_m$'', which is dimensionally inconsistent with
$X^*=VDU'$). With this correction the exercise on the slide (\cref{ch07:ex-lecsvd}) works out.
\end{coursenote}

\begin{casestudy}[2013 Case Study 7: factor modelling of Treasury yields]
\textbf{Data.} FRED constant-maturity yields, 9 tenors (3m to 20y), 3 Jan 2000 to 31 May 2013; daily changes (in
percentage points) over two periods \ts{13}{15}{1:11:14}: 2001--2005 and 2009--May 2013 (the section title says
``2009--2012'', the code uses 2009--2013).
\textbf{2001--2005} \ts{13}{15}{1:12:18}: all mean changes negative (rates fell on balance) \ts{13}{15}{1:13:23}; daily
volatilities 3.8 bp (3m) to 7.0 bp (3y); correlations decreasing away from the diagonal. PCA (\texttt{princomp}):
84.9\%, 9.2\%, 2.9\% \ts{13}{15}{1:18:43}; PC1 a weighted average peaking at 3--5y, PC2 long minus short, PC3
curvature, PC9 the 7y-versus-5y/10y hedge; the PC scores are exactly uncorrelated. Factor analysis
(\texttt{factanal}): 4 factors rejected ($\chi^2=26.8$, 6 df, $p=0.0002$), 5 factors rejected at 5\% ($\chi^2=4.33$, 1 df,
$p=0.037$), 6 factors cannot be fitted ($\nu<0$): no reduced-dimension model passes the test.
\textbf{2009--2013}: short rates at zero, so the 3m--1y yields hardly move (3m volatility 1.25 bp versus 6.6 bp at
7--20y); PCA 88.6\%, 6.8\%, 1.6\%, with PC1 loadings near zero at the short end; PC2 now opposes long and
\emph{medium} tenors. Factor analysis: 3 factors rejected ($\chi^2=231$, 12 df), 4 factors accepted ($\chi^2=5.2$, 6 df,
$p=0.52$); the 3m yield has a large specific variance (uniqueness 0.65).
\textbf{Lesson} \ts{13}{15}{1:25:03}: the structure depends on the period (monetary regime), so the calibration
window matters, and the fitted factors are only a starting point for modelling their dynamics (as in 2013 PS7
Problem~4). (Several uniquenesses equal 0.005, the lower bound used by \texttt{factanal}: near-``Heywood'' cases,
where a variable is almost entirely common factor; not discussed in class.)
\end{casestudy}

\begin{coursenote}[2013 versus 2024]
Both years analyse daily changes of US Treasury yields and find level, slope and curvature. 2013 (Kempthorne,
statistics) uses FRED constant-maturity yields at 9 tenors, fixed 5-year windows, and adds the factor-model and
factor-analysis machinery; 2024 (Andreev, practice) uses the Fed's fitted par yields at 30 maturities with a
rolling 2-year window, and focuses on stability, out-of-sample behaviour and PC-neutral trading. The 2024 course
has no lecture on macroeconomic or fundamental factor models; their portfolio-theory side reappears in
\cref{ch:portfolio}.
\end{coursenote}

% ------------------------------------------------------------------
\begin{keyformulas}
\begin{align*}
&\text{PCA:} && \Sigma=\Gamma\Lambda\Gamma\T,\quad
  \bm p=\Gamma\T(\bm x-\bm\mu),\quad \Cov\bm p=\Lambda,\\
& && \bm x=\bm\mu+\textstyle\sum_kp_k\bm\gamma_k\\
&\text{Variational:} && \bm\gamma_1=\arg\max_{\|\bm w\|=1}\bm w\T\Sigma\bm w,\\
& && \text{share}_k=\lambda_k/\tr\Sigma\\
&\text{Best subspace:} && \min_V\E[\|(I-VV\T)(\bm x-\bm\mu)\|^2]\\
& && =\textstyle\sum_{j>K}\lambda_j\ \ (V=\Gamma_K)\\
&\text{Sample PCA:} && X_c=UD\hat\Gamma\T,\quad
  \hat\lambda_k=d_k^2/(N-1),\quad P=X_c\hat\Gamma=UD\\
&\text{Correlation PCA:} && R=D_\sigma^{-1}\Sigma D_\sigma^{-1},\quad \tr R=p\\
&\text{Out of sample:} && W\T\Sigma^YW=D\T\Lambda^YD,\\
& && D=(W^Y)\T W;\quad \mathrm{ED}=(\textstyle\sum\lambda)^2/\textstyle\sum\lambda^2\\
&\text{PC exposures:} && e_k=\bm h\T\bm\gamma_k;\\
& && \text{curve } \alpha=\gamma_{a1}/\gamma_{b1};\\
& && \text{fly: } 2\times2\text{ system \eqref{ch07:eq-fly}}\\
&\text{PC1 hedge:} && \bm V'=\bm V-(\bm V\T\bm\gamma_1)\bm\gamma_1,\\
& && \Var=\bm V\T\Sigma\bm V-\lambda_1(\bm V\T\bm\gamma_1)^2\\
&\text{Factor model:} && \bm x_t=\bm\alpha+B\bm f_t+\bm\epsilon_t,\\
& && \Sigma_x=B\Omega_fB\T+\Psi\ \ (\Psi\text{ diagonal})\\
&\text{Single index:} && \Sigma_x=\sigma_M^2\bm\beta\bm\beta\T+\Psi\\
&\text{Factor estimates:} && \hat{\bm f}_t=(B\T\Psi^{-1}B)^{-1}B\T\Psi^{-1}\bm x_t,\\
& && HB=I_K\\
&\text{LR test df:} && \nu=\tfrac12\bigl[(m-K)^2-(m+K)\bigr]
\end{align*}
\end{keyformulas}

\begin{exercise}[not from the course]\label{ch07:ex-fa}
\leavevmode
\begin{enumerate}[(a)]
  \item Count the free parameters of the $K$-factor model $\Sigma_x=BB\T+\Psi$ and derive the degrees of freedom $\nu$ of the
      likelihood-ratio test.
  \item Check the values of $\nu$ reported in 2013 Case Study~7 and explain why a 6-factor model
      cannot be fitted to 9 yields.
  \item For the single-index model with $\Psi=\psi I$, show that PCA recovers $\bm\beta/\|\bm\beta\|$ as
      PC1; what happens if the $\sigma_i^2$ differ?
\end{enumerate}
\end{exercise}

% ------------------------------------------------------------------
\section*{Exercises}
\addcontentsline{toc}{section}{Exercises}
Standalone exercises follow the section they practise. The exercises that build on one another stay together here: PS3 Problem~4 with the exercises that reuse its data (PS7 Problem~4 and the PC-neutral hedges), and PS7 Problems~1--2.

\subsection*{From 2024 Problem Set 3}

\begin{exercise}[PS3 (2024), Problem 4; also PS7 (2013), Problem 3: PCA of yield changes]\label{ch07:ex-ps3}
Daily changes of the 9 FRED constant-maturity Treasury yields over 2001--2005 are analysed twice, with the sample
covariance matrix and with the sample correlation matrix. The standard deviations of the PC variables are
\begin{center}\small
\begin{tabular}{@{}l rrrrrrrrr@{}}
\toprule
 & 1 & 2 & 3 & 4 & 5 & 6 & 7 & 8 & 9\\
\midrule
covariance & 0.1618 & 0.0532 & 0.0299 & 0.0192 & 0.0135 & 0.0112 & 0.0096 & 0.0090 & 0.0081\\
correlation & 2.6338 & 1.1892 & 0.5652 & 0.3971 & 0.2694 & 0.2080 & 0.1485 & 0.1366 & 0.1225\\
\bottomrule
\end{tabular}
\end{center}
(proportions of variance 84.9, 9.2, 2.9\,\% and 77.1, 15.7, 3.5\,\% for the first three) and the first three loading
vectors are
\begin{center}\small
\begin{tabular}{@{}l rrr c rrr@{}}
\toprule
 & \multicolumn{3}{c}{covariance PCA} & & \multicolumn{3}{c}{correlation PCA}\\
\cmidrule(lr){2-4}\cmidrule(l){6-8}
Tenor & PC1 & PC2 & PC3 & & PC1 & PC2 & PC3\\
\midrule
3M  & 0.1088 & 0.5225 & 0.5434 & & $-0.2083$ & 0.6293 & 0.5831\\
6M  & 0.1581 & 0.4817 & 0.2508 & & $-0.2830$ & 0.5154 & $-0.0654$\\
1Y  & 0.2531 & 0.4107 & $-0.0514$ & & $-0.3369$ & 0.2698 & $-0.4066$\\
2Y  & 0.3905 & 0.2071 & $-0.4750$ & & $-0.3625$ & 0.0064 & $-0.4005$\\
3Y  & 0.4197 & 0.0873 & $-0.4022$ & & $-0.3664$ & $-0.0665$ & $-0.2719$\\
5Y  & 0.4206 & $-0.1082$ & $-0.0311$ & & $-0.3670$ & $-0.1619$ & 0.0036\\
7Y  & 0.4009 & $-0.2289$ & 0.1475 & & $-0.3608$ & $-0.2298$ & 0.1435\\
10Y & 0.3697 & $-0.2837$ & 0.2710 & & $-0.3525$ & $-0.2656$ & 0.2554\\
20Y & 0.3101 & $-0.3621$ & 0.3944 & & $-0.3290$ & $-0.3338$ & 0.4125\\
\bottomrule
\end{tabular}
\end{center}
(a) PC1 loadings are all positive in the covariance case and all negative in the correlation case. Is the
difference in sign meaningful? (Does the eigen-decomposition change if an eigenvector is multiplied by $-1$?)
(b) The PC1 loadings range more widely in the covariance case; the loading on the least variable change (3M) is
larger in magnitude in the correlation case, and those on the most variable changes are smaller. Explain why,
recalling that PC1 is the unit-norm weighted combination with the largest variance and that in the correlation
case each yield change is scaled to unit variance.
(c) Interpret the first three PCs of the correlation case and compare with the covariance case.
(d) If the goal is to model the dynamics of the whole term structure across all tenors, argue why correlation PCA
might be preferred.
(e) (2024 only) Compute the total variance of the correlation-case PCs. Is it an integer? Explain.
\end{exercise}

\begin{coursenote}[On the data of the previous exercise]
The problem sets say the changes are ``in basis point units''; the numbers are in percentage points (a PC1
standard deviation of 0.1618 is 16 bp per day, the tenor volatilities of Case Study 7 are 0.04--0.07). The 2013
version prints correlation-case standard deviations about $0.04\%$ smaller (their squares sum to $8.993$ instead of
$9$, a factor $(N-1)/N$ with $N\approx1247$ days, presumably a different divisor convention); the loadings coincide up
to signs.
\end{coursenote}

\subsection*{From 2013 Problem Set 7}

\begin{exercise}[PS7 (2013), Problem 1: PCA of a bivariate vector]\label{ch07:ex-ps7-1}
Let $\bm X=(X_1,X_2)\T$ have mean $\bm\alpha=(\alpha_1,\alpha_2)\T$ and covariance
$\Sigma=\bigl(\begin{smallmatrix}\sigma_1^2&\rho\sigma_1\sigma_2\\ \rho\sigma_1\sigma_2&\sigma_2^2\end{smallmatrix}\bigr)$.
\begin{enumerate}[(a)]
  \item Compute the eigenvalues $\lambda_1\ge\lambda_2\ge0$ of $\Sigma$.
  \item Compute orthonormal eigenvectors $\bm\gamma_1,\bm\gamma_2$.
  \item Verify $\Sigma=\sum_i\lambda_i\bm\gamma_i\bm\gamma_i\T$.
  \item For $p_i=\bm\gamma_i\T(\bm X-\bm\alpha)$ prove $\E[p_i]=0$, $\Var(p_i)=\lambda_i$ and $\Cov(p_1,p_2)=0$.
  \item Are the eigenvectors in (b) unique? If not, how are different solutions related?
\end{enumerate}
\end{exercise}

\begin{exercise}[PS7 (2013), Problem 2: PCA as a distance-preserving affine map]\label{ch07:ex-ps7-2}
A distance-preserving affine map of the plane can be written $T(\bm X)=\Phi\T(\bm X-\bm\mu)$ with $\Phi=[\bm\phi_1:\bm\phi_2]$
orthogonal and $\bm\mu\in\R^2$: it moves the origin to $\bm\mu$ and rotates the axes.
\begin{enumerate}[(a)]
  \item Find the unit vector $\bm\phi_1$ maximising $\Var(\bm\phi_1\T\bm X)$ and show that the maximum is $\lambda_1$ of the previous
      exercise.
  \item Find the unit vector $\bm\phi_2$ minimising $\Var(\bm\phi_2\T\bm X)$ and show that the minimum is $\lambda_2$.
  \item Give an equivalent definition of PCA in terms of affine transformations of the variables and their
      variances and expectations.
\end{enumerate}
\end{exercise}

\begin{exercise}[PS7 (2013), Problem 4: dynamics of the PC scores]\label{ch07:ex-ps7-4}
A VAR(3) with constant (\cref{ch:timeseries}) was fitted to the scores of the first three PCs of the
correlation-matrix analysis of \cref{ch07:ex-ps3} ($T=1245$ days). Estimates and $t$-values:
\begin{center}\small
\begin{tabular}{@{}l rr c rr c rr@{}}
\toprule
 & \multicolumn{2}{c}{eq.\ Comp.1} & & \multicolumn{2}{c}{eq.\ Comp.2} & & \multicolumn{2}{c}{eq.\ Comp.3}\\
\cmidrule(lr){2-3}\cmidrule(lr){5-6}\cmidrule(l){8-9}
Regressor & est. & $t$ & & est. & $t$ & & est. & $t$\\
\midrule
Comp.1, lag 1 & 0.0527 & 1.86 & & $-0.0203$ & $-1.65$ & & $-0.0029$ & $-0.48$\\
Comp.2, lag 1 & $-0.1340$ & $-2.08$ & & 0.1382 & 4.93 & & $-0.0143$ & $-1.03$\\
Comp.3, lag 1 & $-0.1322$ & $-1.00$ & & $-0.0833$ & $-1.45$ & & 0.0163 & 0.58\\
Comp.1, lag 2 & $-0.0530$ & $-1.87$ & & $-0.0072$ & $-0.59$ & & $-0.0135$ & $-2.23$\\
Comp.2, lag 2 & $-0.0542$ & $-0.84$ & & 0.0370 & 1.32 & & $-0.0037$ & $-0.27$\\
Comp.3, lag 2 & 0.1825 & 1.39 & & $-0.0195$ & $-0.34$ & & $-0.0523$ & $-1.86$\\
Comp.1, lag 3 & $-0.0169$ & $-0.60$ & & $-0.0085$ & $-0.69$ & & $-0.0019$ & $-0.32$\\
Comp.2, lag 3 & 0.0130 & 0.21 & & $-0.0403$ & $-1.46$ & & $-0.0030$ & $-0.22$\\
Comp.3, lag 3 & 0.0830 & 0.63 & & 0.0148 & 0.26 & & $-0.0712$ & $-2.53$\\
constant & $-0.0114$ & $-0.15$ & & 0.0121 & 0.38 & & 0.0002 & 0.01\\
\midrule
$R^2$; $F$-test $p$-value & 0.013 & 0.054 & & 0.029 & $4\cdot10^{-5}$ & & 0.013 & 0.059\\
\bottomrule
\end{tabular}
\end{center}
Residual standard errors are 2.60, 1.13, 0.56; residual correlations are below $0.025$ in absolute value; the
largest root of the characteristic polynomial has modulus 0.438.
(a)--(c) Interpret the estimated equations for Comp.1, Comp.2 and Comp.3.
(d) For each equation, is there evidence of mean reversion or of momentum in the PC variable? Explain.
\end{exercise}

\subsection*{Additional exercises}

\begin{exercise}[not from the course]\label{ch07:ex-fly}
Using the covariance-PCA loadings of \cref{ch07:ex-ps3}:
\begin{enumerate}[(a)]
  \item find the PC1-neutral 2s10s steepener and check
      that its exposure to PC2 is non-zero;
  \item find the PC1/PC2-neutral 2s5s10s fly (buy \$1 dv01 of 5y) and compute its
      exposure to PC3;
  \item prove the variance formula for the PC1-hedged portfolio $\bm V'=\bm V-(\bm V\T\bm\gamma_1)\bm\gamma_1$ and show
      that it never exceeds $\Var(\bm V\T\bm r)$.
\end{enumerate}
\end{exercise}

% !TEX root = ../main.tex
\chapter{Counterparty Risk, Margin and CVA}\label{ch:counterparty}
\begin{paracol}{2}

\begin{sources}
\begin{tabularx}{\linewidth}{@{}l l L@{}}
\toprule
\textbf{Lecture} & \textbf{Speaker} & \textbf{Content}\\
\midrule
2024 L10 & J.~Shepherd & \emph{Counterparty Risk Optimization} (Quantile Technologies, part of LSEG).
Types of risk and central clearing; coherent risk measures, VaR and expected shortfall (historical
and model-building estimates, time scaling, back-testing, Rockafellar--Uryasev); variation and initial
margin, ISDA SIMM; multilateral optimisation of initial margin in a network of banks, numerical issues,
``nice'' and ``fair'' solutions.\\
2013 L26 & Y.~Tang & \emph{Introduction to Counterparty Credit Risk} (Morgan Stanley, given remotely).
CVA, DVA, funding risk, wrong-way risk (Berkshire Hathaway's index puts), the CVA conundrum;
enterprise-level derivatives modelling; martingale measures, martingale testing, resampling and
interpolation. Last minutes: closing remarks of the instructors (\cref{ch08:sec-closing}).\\
\bottomrule
\end{tabularx}

\smallskip
\textbf{Files.} Slides \file{mit18_642_f24_lec10.pdf} (51 slides); 2013 slides
\file{MIT18_S096F13_lecnote26.pdf} (27 slides).
\textbf{Problem sets.} No problem set of either year covers this chapter; all exercises are marked
``not from the course''.
\textbf{Not covered in class.} 2024: slide 17 (FRTB) was skipped, slides 26--28 (SIMM practicalities)
and 48 (observed savings) only shown; 2013: slides 21--25 (extinguisher trade, ``bigger picture'')
were not reached for lack of time (summarised in \cref{ch08:rem-extinguisher}).
\end{sources}

\section*{Motivation}
An over-the-counter (OTC) derivative is a bilateral promise. Over the life of a swap sometimes you owe
your counterparty money and sometimes it owes you; if it defaults while it owes you, you lose part of the
receivable. Before 2008 this \emph{counterparty credit risk} was a side issue. During the crisis the
mark-to-market losses it caused were larger than the losses from actual defaults \ts{13}{26}{6:22}, and the
regulatory response (Basel~III capital for CVA, mandatory central clearing, margin rules for uncleared
derivatives) changed the whole derivatives business. The two guest lectures show the two sides of this change:
\begin{itemize}
  \item \textbf{2013 (Y.~Tang, Morgan Stanley)}: \emph{pricing} counterparty risk. The credit valuation
        adjustment (CVA) and its relatives (DVA, funding), wrong-way risk, and why all of them need models of
        the whole portfolio, not of one trade at a time.
  \item \textbf{2024 (J.~Shepherd, Quantile/LSEG)}: \emph{mitigating} it with collateral. Variation margin,
        initial margin computed as a value-at-risk, and how a network of banks can reduce the total initial
        margin they post by agreeing on new trades, found by convex optimisation.
\end{itemize}
The mathematics: option-like expectations under the pricing measure, coherent risk measures (VaR and
expected shortfall), quadratic forms of Gaussian vectors, convex and mixed-integer optimisation, and change of
numeraire. VaR is treated in more depth, from the 2013 lecture on the topic, in \cref{ch:var}.

% ==================================================================
\section{Counterparty risk in OTC derivatives}\label{ch08:sec-ccr}

\paragraph{Types of risk.} Trading derivatives creates several risks \ts{24}{10}{3:11}:
\emph{market} risk (values move with the market), \emph{credit} risk (a debtor does not pay an obligation),
\emph{operational and legal} risk (people, systems and events; the hardest to model, because it is mostly
people making mistakes), \emph{liquidity} risk (a transaction cannot be executed at the market price, typically
because the position is so large that selling it moves the price, or outgoing cash flows cannot be funded),
and \emph{counterparty} risk: the potential loss on the market value of transactions when the counterparty
defaults.

\paragraph{Counterparty versus credit risk.} Suppose you buy protection (a credit default swap, CDS) on a
JPMorgan bond from Goldman Sachs. You have \emph{credit} risk to JPMorgan: the value of the CDS depends on
JPMorgan defaulting. You have \emph{counterparty} risk to Goldman Sachs: whether it will actually pay what the
CDS is worth \ts{24}{10}{4:13}. For a corporate bond the economics are similar, but the risk to the issuer
is called \emph{issuer} risk \ts{13}{26}{3:10}. Counterparty risk has two components \ts{13}{26}{2:04}:
\begin{itemize}
  \item \emph{default risk}: at default we recover only a fraction $R$ of what we are owed;
  \item \emph{mark-to-market (MTM) risk}: even if nothing is due for ten years, a widening of the
        counterparty's credit spread lowers the value of the trade today, because a buyer of the trade would
        pay less for it \ts{13}{26}{3:10}.
\end{itemize}
What makes counterparty risk harder than bond credit risk is that the amount at risk is \emph{random and
two-sided}: the value of a swap can be positive or negative, and it depends on all trades with the same
counterparty.

\paragraph{Mitigation.} The standard tools (slide~4 of \file{lec10}) are \ts{24}{10}{5:15}: \emph{netting}
of offsetting payments; \emph{hedging} with new trades; \emph{collateralisation} (collecting margin, a kind of
insurance payment against the counterparty's default); contractual clauses such as periodic resets of the
mark-to-market value or break clauses; and \emph{central clearing counterparties} (CCPs).

\paragraph{Central clearing.} A CCP steps into the middle of bilateral trades by \emph{novation}: the trade
between A and B is replaced by two trades, A with the CCP and the CCP with B \ts{24}{10}{6:15}. By itself
this only changes the counterparty; the gain comes from \emph{multilateral netting}, since all of a firm's
positions now face the same counterparty and offset (\cref{ch08:fig-ccp}) \ts{24}{10}{7:20}. As a rough
rule, total counterparty risk scales like $\sum_{\text{counterparties}}\|\text{position}\|$, and netting reduces
the number of terms. Standard (vanilla) interest-rate swaps must be cleared by regulation, while more exotic
products such as swaptions are not eligible, so every bank ends up with a mixture of cleared and bilateral
trades \ts{24}{10}{8:25}.

\begin{figure}[ht]
\centering
\begin{tikzpicture}[>=Stealth,every node/.style={font=\small},
  pt/.style={circle,draw,thick,minimum size=7mm,inner sep=0pt}]
\begin{scope}
  \node[pt] (C) at (0,2.4) {C}; \node[pt] (D) at (2.4,2.4) {D};
  \node[pt] (A) at (0,0) {A};   \node[pt] (B) at (2.4,0) {B};
  \draw[->,thick,defcol] (C) -- node[left]{125} (A);
  \draw[->,thick,defcol] (A) -- node[below]{100} (B);
  \draw[->,thick,defcol] (B) -- node[right]{100} (D);
  \draw[->,thick,defcol] (D) -- node[above]{50} (C);
  \node at (1.2,-1.0) {bilateral: gross $375$};
\end{scope}
\begin{scope}[xshift=5.2cm]
  \node[pt] (C) at (0,2.4) {C}; \node[pt] (D) at (2.4,2.4) {D};
  \node[pt] (A) at (0,0) {A};   \node[pt] (B) at (2.4,0) {B};
  \node[circle,draw,thick,fill=okcol!15,minimum size=9mm,inner sep=0pt] (P) at (1.2,1.2) {\scriptsize CCP};
  \draw[thick,thmcol] (A) -- node[pos=0.45,below right,inner sep=1pt]{$+25$} (P);
  \draw[thick,black!30] (B) -- node[pos=0.45,below left,inner sep=1pt]{$0$} (P);
  \draw[thick,thmcol] (C) -- node[pos=0.45,above right,inner sep=1pt]{$-75$} (P);
  \draw[thick,thmcol] (D) -- node[pos=0.45,above left,inner sep=1pt]{$+50$} (P);
  \node at (1.2,-1.0) {after netting at the CCP: $150$};
\end{scope}
\end{tikzpicture}
\caption{Multilateral netting at a CCP (slide~5 of \file{lec10}). Left: four bilateral positions; an arrow
into a node counts $+$, out of a node $-$. Right: after novation each party faces only the CCP and its
positions net: A has $+125-100=+25$, B has $+100-100=0$ (no counterparty risk left), C has $+50-125=-75$,
D has $+100-50=+50$. The CCP itself is flat.}\label{ch08:fig-ccp}
\end{figure}

\begin{finance}[What a CCP costs (not discussed in class)]
Clearing concentrates risk in the CCP, which protects itself by collecting initial margin and contributions to a
default fund from its members. Splitting a portfolio between cleared and bilateral trades also destroys netting
\emph{between} the two sets. This is why ``move more trades into the CCP'' is a genuine design choice in the
optimisation of \cref{ch08:sec-opt}, not an automatic improvement.
\end{finance}

% ==================================================================
\section{Exposure, netting and collateral}\label{ch08:sec-exposure}

Fix a counterparty and a \emph{netting set} (all trades covered by one master agreement, whose values are
offset at default). Let $V(t)$ be the value \emph{to us} of the netting set computed with a
counterparty-default-free model, and $C(t)$ the value of the collateral we hold ($C<0$ if we have posted).

\begin{definition}[Exposure profiles]\label{ch08:def-exposure}
The \emph{exposure} at time $t$ is $E(t)=\bigl(V(t)-C(t)\bigr)^+$, the amount we would lose (before
recovery) if the counterparty defaulted at $t$. Then:
\begin{itemize}
  \item \emph{expected exposure} $\mathrm{EE}(t)=\E[E(t)]$, and \emph{expected positive exposure}
        $\mathrm{EPE}=\frac1T\int_0^T\mathrm{EE}(t)\dd t$, called \emph{mean positive exposure} (MPE) in the 2013
        slides;
  \item \emph{negative exposure} $\bigl(V(t)-C(t)\bigr)^-=\max(C(t)-V(t),0)$, what we owe; its mean is the MNE of
        the 2013 slides;
  \item \emph{potential future exposure} $\mathrm{PFE}_q(t)$, the $q$-quantile of $E(t)$ (typically $q=95\%$
        or $99\%$), used by banks to set credit limits (not discussed in class).
\end{itemize}
For pricing (CVA) the expectation is under the pricing measure $\Q$; for risk management under $\Prob$.
\end{definition}

The positive part is the key non-linearity: if $V>0$ at default we recover only $R\,V$, but if $V<0$ we
still owe the full amount to the defaulted counterparty. There is no ``default gain'', and this asymmetry
makes the loss an \emph{option-like} payoff \ts{13}{26}{19:10}. \Cref{ch08:fig-exposure} shows the typical
shapes: the uncertainty of $V(t)$ grows like $\sqrt t$ by diffusion, while for a swap the remaining cash flows,
and with them the sensitivity to rates, shrink linearly to maturity.

\begin{figure}[ht]
\centering
\begin{tikzpicture}
\begin{axis}[width=0.78\linewidth,height=5.2cm,xmin=0,xmax=10,ymin=0,ymax=2.8,
  xlabel={time $t$ (years)},ylabel={exposure / $\sigma$},
  tick label style={font=\scriptsize},label style={font=\small},
  legend style={font=\scriptsize,at={(0.02,0.98)},anchor=north west,draw=none,fill=none},
  legend cell align=left]
\addplot[defcol,thick,domain=0:10,samples=101]{0.3989*sqrt(x)*(1-x/10)};
\addlegendentry{swap, EE}
\addplot[defcol,thick,dashed,domain=0:10,samples=101]{1.645*sqrt(x)*(1-x/10)};
\addlegendentry{swap, PFE$_{95\%}$}
\addplot[thmcol,thick,domain=0:10,samples=101]{0.3989*0.5*sqrt(x)};
\addlegendentry{forward, EE}
\addplot[thmcol,thick,dashed,domain=0:10,samples=101]{1.645*0.5*sqrt(x)};
\addlegendentry{forward, PFE$_{95\%}$}
\end{axis}
\end{tikzpicture}
\caption{Stylised exposure profiles (not from the course). If $V(t)\sim N(0,s(t)^2)$ and there is no
collateral, then $\mathrm{EE}(t)=s(t)/\sqrt{2\pi}\approx0.40\,s(t)$ and $\mathrm{PFE}_{95\%}(t)=1.645\,s(t)$.
Swap: $s(t)=\sqrt t\,(1-t/T)$ with $T=10$ (diffusion times remaining duration, peak at $T/3$).
FX forward: $s(t)=0.5\sqrt t$ (a single exchange at maturity).}\label{ch08:fig-exposure}
\end{figure}

\begin{proposition}[Netting reduces exposure]\label{ch08:prop-netting}
For trades with values $V_1,\dots,V_n$ in one netting set,
$\bigl(\sum_iV_i\bigr)^+\le\sum_iV_i^+$ pathwise; hence $\mathrm{EE}_{\text{net}}(t)\le\sum_i\mathrm{EE}_i(t)$.
\end{proposition}
\begin{proof}
$x\mapsto x^+$ is subadditive: $(x+y)^+\le x^++y^+$ because the right side is $\ge x+y$ and $\ge0$.
Induct on $n$ and take expectations.
\end{proof}
If one trade goes from $0$ to $+\$100$M and an offsetting trade with the same counterparty goes to $-\$100$M, the
netted default loss is zero \ts{13}{26}{18:05}. The price of counterparty risk is therefore \emph{not} a sum of
per-trade quantities: to price it for one new trade you must know every other trade facing the same
counterparty. Ordinary derivative pricing never needs this \ts{13}{26}{7:29}.

\paragraph{Collateral.} Under a \emph{Credit Support Annex} (CSA) to the master agreement the parties exchange
collateral. \emph{Variation margin} (VM) is exchanged daily and equals the current value, $\mathrm{VM}(t)=V(t)$;
the CSA also fixes thresholds below which no collateral is called, minimum transfer amounts and eligible assets
(these details were not discussed in class, apart from the thresholds mentioned on the SIMM slides). With perfect
daily VM the exposure is no longer the \emph{level} of $V$ but its \emph{increment} between the last margin call
the defaulter honoured and the moment we have closed out or replaced the trades. That increment is covered by
\emph{initial margin} (\cref{ch08:sec-margin}).

\begin{intuition}[From $\sqrt{T}$ to $\sqrt{\text{MPOR}}$]
For a diffusive value process the typical size of $V(t)$ after 5 years is $\sigma\sqrt{5}$, while the typical move
over a 10-business-day close-out period is $\sigma\sqrt{10/250}$, about $9\%$ of it. Collateralisation replaces
a large, slowly growing exposure by a small, constantly renewed one. The price is liquidity: margin must be
paid in cash or high-quality securities, often at short notice.
\end{intuition}

\begin{exercise}[Netting and correlation (not from the course)]
Two trades with the same counterparty have values $V_1,V_2\sim N(0,\sigma^2)$ at time $t$, with correlation $\rho$, and
there is no collateral.
\begin{enumerate}[(a)]
  \item Compute $\mathrm{EE}$ with and without netting and show that the ratio is
      $\sqrt{(1+\rho)/2}$.
  \item What happens for $\rho=-1$ and $\rho=1$?
  \item Explain why the CVA of a new trade depends on the
      existing portfolio (the ``incremental CVA'').
\end{enumerate}
\end{exercise}

% ==================================================================
\section{CVA, DVA and funding: pricing counterparty risk}\label{ch08:sec-cva}

\subsection{A puzzle}
Tang opened with a question \ts{13}{26}{8:32}. You have an uncollateralised swap with value $0$ on day one
(ignoring counterparty risk). Later its value is $+\$100$M, and then the counterparty unexpectedly defaults with
$50\%$ recovery, so you receive $\$50$M in cash. Over the life of the trade, have you made $\$50$M, lost $\$50$M,
or are you flat?

A dealer has \emph{lost} $\$50$M \ts{13}{26}{12:45}. Its market risk is hedged, for example by an offsetting swap
with another dealer or an exchange (or the trade is simply intermediated back to back). When the client trade went
from $0$ to $+100$, the hedge went from $0$ to $-100$, and the hedge counterparty is solvent and must be paid in
full. The dealer receives $50$ and pays $100$. In general the loss at default is $(1-R)$ times the exposure. A
\emph{CVA desk} that had sold the trading desk protection on the counterparty (and itself bought CDS protection)
would have made the trading desk whole \ts{13}{26}{13:53}.

\subsection{Credit valuation adjustment}
\begin{definition}[CVA, 2013 slides]\label{ch08:def-cva}
The \emph{credit valuation adjustment} is the price of counterparty credit risk, an adjustment to the value given
by a counterparty-default-free model or a broker quote \ts{13}{26}{4:13}. With $\tau$ the counterparty's default
time, $T$ the final maturity of the netting set, $R_\tau$ the recovery rate, $V_{\tau-},C_{\tau-}$ the portfolio
value and collateral just before default, and $\beta_t=\exp\bigl(\int_0^tr_u\dd u\bigr)$ the money-market account,
\begin{equation}\label{ch08:eq-cva}
  \mathrm{CVA}_{\mathrm{MPE}} = -\E_0^{\Q}\!\left[\ind_{\{\tau\le T\}}\,
      \frac{(V_{\tau-}-C_{\tau-})^+\,(1-R_\tau)}{\beta_\tau}\right] \ts{13}{26}{15:57}.
\end{equation}
The adjusted value of the portfolio is $V_0+\mathrm{CVA}_{\mathrm{MPE}}$.
\end{definition}
The minus sign matters in practice: CVA on receivables is a cost, the amount charged to the counterparty.
Theoretical papers often define CVA as a positive number and subtract it; that is fine as long as you are
consistent \ts{13}{26}{14:55}. In \eqref{ch08:eq-cva}, $\ind_{\{\tau\le T\}}$ says there is no counterparty risk
if the counterparty defaults after the portfolio has matured; $(V-C)^+$ is the exposure; $(1-R)$ turns it into a
loss given default; and $\beta_\tau$ discounts \ts{13}{26}{17:01}.

\begin{proposition}[CVA as an integral of expected exposure (not derived in class)]\label{ch08:prop-cvaee}
Assume $\tau$ is independent of the market variables (no wrong-way risk, \cref{ch08:sec-wwr}), has a
deterministic hazard rate $h(t)$ with survival function $S(t)=\exp\bigl(-\int_0^th\bigr)$, and $R$ is constant.
Then
\[
  \mathrm{CVA} = -(1-R)\int_0^T \mathrm{EE}^*(t)\,h(t)S(t)\dd t,\qquad
  \mathrm{EE}^*(t)=\E^{\Q}\!\left[\frac{(V_t-C_t)^+}{\beta_t}\right].
\]
If moreover $h$ and $r$ are constant, $\mathrm{EE}(t)\approx\mathrm{EPE}$ is roughly flat, and the CDS spread obeys
the ``credit triangle'' $s\approx(1-R)h$, then
\begin{equation}\label{ch08:eq-cvabote}
  \mathrm{CVA}\approx-\,\mathrm{EPE}\times s\times D,\qquad D=\int_0^Te^{-(r+h)t}\dd t=\frac{1-e^{-(r+h)T}}{r+h},
\end{equation}
where $D$ is the risky duration (annuity) of the counterparty. This is the ``back-of-the-envelope'' formula of
the slides: $-\,$MPE $\times$ CDS par spread $\times$ duration \ts{13}{26}{13:53}.
\end{proposition}
\begin{proof}
Let $g(t)=\E^{\Q}\bigl[(V_t-C_t)^+/\beta_t\bigr]$. Since $\tau$ is independent of the market,
$\E[\ind_{\{\tau\le T\}}(1-R)(V_{\tau-}-C_{\tau-})^+/\beta_\tau]=(1-R)\E[\ind_{\{\tau\le T\}}g(\tau)]$ (paths are
continuous at a fixed time, so $V_{t-}=V_t$ a.s.), and $\tau$ has density $h(t)S(t)$. For the approximation, put
$\mathrm{EE}^*(t)=e^{-rt}\mathrm{EPE}$ and $hS(t)=he^{-ht}$, then use $(1-R)h=s$.
\end{proof}
The wider the spread (default more likely) and the longer the duration (more time to default), the larger the
charge \ts{13}{26}{14:55}.

\begin{example}[Size of CVA]
A 5-year uncollateralised swap book with $\mathrm{EPE}=\$5$M, counterparty spread $s=100$\,bp, $R=40\%$ (so
$h=1.67\%$) and $r=3\%$: $D=(1-e^{-0.0467\cdot5})/0.0467=4.46$ and $\mathrm{CVA}\approx-5\cdot0.01\cdot4.46
=-\$0.22$M. If the spread doubles, the CVA roughly doubles: a mark-to-market loss with no default at all.
\end{example}

\begin{intuition}[CVA is an option on a basket]
$(V_\tau-C_\tau)^+$ is the payoff of a call on the netting set with strike equal to the collateral, exercised at
the random time $\tau$. CVA is a default-probability-weighted strip of such options. Since a counterparty trades
many products (rates, FX, credit, equity, commodities, mortgages), it is an option on a \emph{basket of cross-asset
derivatives} \ts{13}{26}{20:11}. Pricing it needs a joint simulation of all risk factors of all trades: this is
what Tang calls enterprise-level modelling (\cref{ch08:sec-martingale}), and the reason his group at Morgan Stanley
was called Cross Asset Modeling.
\end{intuition}

\begin{finance}[CVA in the bank]
CVA changes the true economic price of a trade, so it is charged to clients and hedged dynamically by a CVA
desk, which buys CDS protection and hedges the market factors that drive $V$. It also enters regulatory capital:
Basel~III added a CVA capital charge, so trades with large CVA consume more capital and lower the return on
equity \ts{13}{26}{5:17}. In 2008, CVA mark-to-market losses were larger than the losses from actual defaults
\ts{13}{26}{6:22}.
\end{finance}

\begin{exercise}[CVA from exposure (not from the course)]
\leavevmode
\begin{enumerate}[(a)]
  \item Derive \cref{ch08:prop-cvaee} and the approximation \eqref{ch08:eq-cvabote}.
  \item For the exposure profile
      $\mathrm{EE}(t)=\sigma\sqrt t(1-t/T)/\sqrt{2\pi}$ of \cref{ch08:fig-exposure}, with $r=0$ and constant hazard rate $h$,
      compute CVA to first order in $h$ and compare with $\mathrm{EPE}\cdot s\cdot T$.
  \item Explain qualitatively how wrong-way
      risk changes the result.
\end{enumerate}
\end{exercise}

\subsection{DVA and first-to-default}
The same logic applied to \emph{our own} default gives a \emph{liability} CVA \ts{13}{26}{21:23}. With $\bar\tau$
our default time and $\bar R$ our recovery, the slides write
\[
  \mathrm{CVA}_{\mathrm{MNE}} = -\E_0^{\Q}\!\left[\ind_{\{\bar\tau\le T\}}\,
      \frac{(V_{\bar\tau-}-C_{\bar\tau-})^-\,(1-\bar R_{\bar\tau})}{\beta_{\bar\tau}}\right]
  = \E_0^{\Q}\!\left[\ind_{\{\bar\tau\le T\}}\,
      \frac{(C_{\bar\tau-}-V_{\bar\tau-})^+\,(1-\bar R_{\bar\tau})}{\beta_{\bar\tau}}\right]\ \ge 0 .
\]
This is a \emph{benefit}: if we default while owing money, we pay only a fraction of it. It is usually called
\emph{debit valuation adjustment} (DVA) \ts{13}{26}{22:27}.

\begin{coursenote}[Sign convention on slide~6 of \file{lecnote26}]
For the displayed formula to be a benefit, $(x)^-$ must be read as $\min(x,0)\le0$, as Tang explains
(``the positive sign becomes a negative sign''). With the other common convention $x^-=\max(-x,0)$ the
formula would carry the opposite sign. The right-hand side above avoids the ambiguity.
\end{coursenote}

Should the two adjustments know about each other? If we default first, the counterparty's later default no longer
matters to us, and vice versa. A \emph{bilateral} CVA puts first-to-default indicators $\ind_{\{\tau\le\min(\bar\tau,T)\}}$
and $\ind_{\{\bar\tau\le\min(\tau,T)\}}$ into the two terms (not written in class). Tang's view is that for pricing
this refinement does not need to be modelled; he gave references on the slides rather than discussing it
\ts{13}{26}{23:30}.

\begin{finance}[Why DVA is controversial (not discussed in class)]
DVA means that a bank books a gain when its own credit worsens, and it can hardly be hedged, since a bank cannot
sell protection on itself. Regulators therefore do not let DVA gains count towards capital, even though
accounting standards include DVA in fair value.
\end{finance}

\subsection{Funding and liquidity risk}
Second example \ts{13}{26}{23:30}: same trade, value $+\$100$M, uncollateralised, with market risk and counterparty
risk properly hedged. Is anything left? Students proposed interest-rate risk (hedged by assumption) and
``key-man'' risk \ts{13}{26}{25:40}. Tang's answer is \emph{funding liquidity risk} \ts{13}{26}{25:40}:
\begin{itemize}
  \item hedges are typically futures or collateralised trades. When our uncollateralised receivable is worth
        $+100$ the hedge is worth $-100$, and we must post $100$ of cash \emph{now}, while the receivable
        arrives only later. That cash must be borrowed at our unsecured funding spread \ts{13}{26}{26:47};
  \item the amount depends on market moves (\emph{contingent} liquidity risk), and failing to raise it is how
        firms die: ``you may end up like Lehman'' \ts{13}{26}{27:51};
  \item symmetrically, an uncollateralised \emph{payable} generates a funding \emph{benefit} (we receive
        collateral on its hedge), which can partially offset the funding need of the receivables
        \ts{13}{26}{27:51};
  \item further risks: unexpected tail events and (equity) capital \ts{13}{26}{28:54}.
\end{itemize}
In the spirit of \eqref{ch08:eq-cvabote}, a first estimate of the funding cost is $-s_F\times\mathrm{EPE}\times D$
and of the benefit $+s_F\times|\mathrm{ENE}|\times D$, with $s_F$ our own funding spread (a common simplification,
not from the lecture).

\begin{finance}[The XVA family (terminology not used in class)]
The adjustments to the default-free price are nowadays collectively called XVA: CVA (counterparty default),
DVA (own default), FVA (funding of uncollateralised positions, split into cost FCA and benefit FBA), KVA (cost
of regulatory capital held against the trade, Tang's ``capital''), and MVA (cost of funding the \emph{initial
margin} of \cref{ch08:sec-margin} over the life of the trade). The 2013 lecture covered CVA, DVA, funding and
capital. The 2024 lecture is about the world that the margin rules created, in which CVA is smaller because
trades are collateralised and the cost has moved into margin; its reading list cites J.~Gregory, \emph{The xVA
Challenge}.
\end{finance}

\subsection{Wrong-way risk}\label{ch08:sec-wwr}
\begin{definition}[Wrong-way risk]
\emph{Wrong-way risk} (WWR) is present when exposure to a counterparty is positively correlated with its
probability of default: the counterparty owes us more exactly when it is more likely to default
\ts{13}{26}{33:02}. (\emph{Right-way} risk is the opposite.) Then the independence assumption of
\cref{ch08:prop-cvaee} fails, and the relevant exposure is the \emph{conditional} one,
$\mathrm{CVA}=-(1-R)\int_0^T\E^{\Q}\bigl[(V_t-C_t)^+/\beta_t\,\big|\,\tau=t\bigr]\dd\Q(\tau\le t)$, which is
larger than $\mathrm{EE}^*(t)$ under WWR.
\end{definition}

\begin{casestudy}[Berkshire Hathaway's index puts (2013 slides 8--9)]
\textbf{The trade} \ts{13}{26}{29:54}. Selling a put below the current price is an alternative to buying the
stock: similar downside, no upside beyond the premium, and a steady income if repeated (``name your price and
get paid for it''). Around 2005--2006, as Tang recalls, Berkshire Hathaway sold long-dated (10--15 year) puts on
four indices (S\&P~500, FTSE~100, Euro Stoxx~50, Nikkei~225) and collected about $\$4$bn of premium \emph{without
posting collateral}. It was one of the largest derivative cash outflows of its time for the dealers, who were the
buyers \ts{13}{26}{30:55}. Had Berkshire posted collateral, the 2008 sell-off would have required many billions of
it (see page~16 of Berkshire's 2012 shareholder letter, cited on the slides).

\textbf{The dealer's risks} \ts{13}{26}{31:56}. Handling long-dated equity options was not the difficulty; many dealers
could not handle the \emph{counterparty} side: CVA on a large uncollateralised receivable; \emph{wrong-way} CVA,
because in a sell-off the puts gain value while an insurer with large equity holdings becomes more likely to
default \ts{13}{26}{33:02}; and \emph{wrong-way funding}, because the dealer paid $\$4$bn up front and had to fund
its collateralised hedges exactly in a crisis \ts{13}{26}{34:06}. The dealer therefore charged CVA, wrong-way CVA,
funding and wrong-way funding costs. It then had to risk-manage them: a static CDS hedge (fine for a bond) leaves
\emph{gap risk} because the exposure changes over time \ts{13}{26}{35:11}. Part of the risk was sold to investors
in Berkshire's bonds through a \emph{credit-linked note} paying a higher coupon, and through tranched portfolio
protection \ts{13}{26}{36:17}.

\textbf{Berkshire's side} \ts{13}{26}{37:18}. Selling puts is selling insurance on the equity market; the premium
is day-one cash, a cheap source of funding (Buffett compared it to borrowing at about 1\%); and there is no cash
outflow until maturity. So Berkshire suffered mark-to-market losses in 2008 but no liquidity drain.

\textbf{Feedback} \ts{13}{26}{38:25}. In the crisis, dealers' CVA hedging meant buying more and more CDS
protection on Berkshire, which widened Berkshire's CDS spread beyond what fundamentals suggested. Bond investors
then demanded higher coupons, raising Berkshire's funding cost, so Berkshire approached dealers about posting
collateral after all: ``no free lunch'' \ts{13}{26}{39:30}.

\textbf{Who loses?} Tang's answer: the end buyers of the puts, the hedgers who pay the premium and gain nothing if
the market does not fall far enough \ts{13}{26}{40:33}.
\end{casestudy}

\subsection{The CVA conundrum}
To hedge CVA on counterparty A you buy credit protection on A from counterparty B. Now you have counterparty risk
to B, so you buy protection on B from C, and so on: an infinite series \ts{13}{26}{41:41}. This also undermines
the theory: arbitrage pricing is replication with hedge instruments, and here replication needs infinitely many
of them \ts{13}{26}{43:52}. The practical way to stop the series at the first term is to buy the protection on A
\emph{fully collateralised}, typically from a dealer under a CSA: collateral is settled as the value moves, much
like daily settlement of futures, so the residual risk to B is negligible \ts{13}{26}{43:52}.

\begin{coursenote}[Slides 21--25 of \file{lecnote26}: the extinguisher (not presented)]\label{ch08:rem-extinguisher}
Tang ran out of time and referred students to the slides \ts{13}{26}{1:16:34}. The question: normally a dealer
\emph{charges} CVA; how can it structure a trade in which it \emph{pays} the counterparty for the latter's default
risk? Answer: a trade with positive value to the dealer that increases with the counterparty's default probability,
namely an \emph{extinguisher} (zero-recovery) swap, which knocks out at zero value if either party defaults. It is
worth something when the dealer expects to have more future \emph{payables} than receivables. Example: an
emerging-market sovereign issues a USD bond and swaps it into local currency with a dealer (receives USD, pays
local currency). Because of the interest-rate differential (the FX forward), the correlation between FX and
credit, and the jump of the FX rate at default, the dealer expects to owe money on the swap; the extinguisher
cancels this debt at the sovereign's default. In effect the dealer has bought protection on the sovereign from the
sovereign itself, and can sell it on to a third party. Issues: it reduces recovery for the other creditors; the
default event must reference the underlying bond; there are legal and franchise risks; and why would the sovereign
default, knowing it would lose its receivable on the swap? Slide~25 places all this in the ``bigger picture'' of a
bank's objectives: revenue, risk management (market, counterparty, liquidity, capital) and return on capital,
with the non-linear portfolio effects requiring enterprise-level models.
\end{coursenote}

% ==================================================================
\section{Risk measures: value at risk and expected shortfall}\label{ch08:sec-riskmeasures}
Initial margin is a risk measure of a portfolio's value change, and to optimise it we must know which risk
measures are well behaved. This section follows the first half of the 2024 lecture.

\subsection{Coherent risk measures}
A risk measure should capture both the \emph{size} of potential losses and their \emph{probability}
\ts{24}{10}{8:25}. Below $L$ denotes the \emph{loss} of a portfolio over a horizon (a random variable; gains are
negative losses) and $\rho(L)\in\R$ its risk.

\begin{definition}[Coherent risk measure, Artzner et al.\ 1999]\label{ch08:def-coherent}
$\rho$ is \emph{coherent} if for all losses $L,L_1,L_2$:
\begin{enumerate}[(i)]
  \item normalised: $\rho(0)=0$;
  \item positively homogeneous: $\rho(\lambda L)=\lambda\rho(L)$ for $\lambda\ge0$;
  \item translation invariant: $\rho(L+a)=\rho(L)+a$ for a deterministic amount $a\in\R$;
  \item monotone: $L_1\le L_2$ in all scenarios $\Rightarrow\rho(L_1)\le\rho(L_2)$;
  \item subadditive: $\rho(L_1+L_2)\le\rho(L_1)+\rho(L_2)$ \ts{24}{10}{9:29}.
\end{enumerate}
\end{definition}
Subadditivity is the important one: merging two portfolios should not create risk, and it may reduce it through
diversification \ts{24}{10}{10:31}.

\begin{coursenote}[Erratum: slide~6 of \file{lec10}]
The slide states homogeneity for all $\alpha$; it holds only for $\alpha\ge0$ (doubling a position doubles its
risk, but reversing it does not negate the risk). It also writes translation invariance as
$\rho(X+A)=\rho(X)+\alpha$ for ``a deterministic portfolio with guaranteed \emph{return} $\alpha$''. With $X$ a loss,
as in the rest of the slides, adding a guaranteed \emph{gain} $\alpha$ \emph{lowers} the risk:
$\rho(X-\alpha)=\rho(X)-\alpha$; equivalently, adding a deterministic \emph{loss} $a$ raises it by $a$, as in
(iii) above. Finally, the reference list dates Artzner--Delbaen--Eber--Heath to 2022; the paper appeared in
\emph{Mathematical Finance} in 1999, the year quoted in class.
\end{coursenote}

\subsection{VaR and ES}
\begin{definition}[VaR and ES]\label{ch08:def-vares}
For a confidence level $\beta\in(0,1)$ and a horizon (implicit in $L$):
\begin{align*}
  \mathrm{VaR}_\beta(L) &= \inf\{\ell\in\R:\ \Prob(L\le\ell)\ge\beta\}, \\
  \mathrm{ES}_\beta(L)  &= \frac{1}{1-\beta}\int_\beta^1\mathrm{VaR}_u(L)\dd u .
\end{align*}
If $L$ has a continuous density $f_L$ then $\int_{-\infty}^{\mathrm{VaR}_\beta}f_L=\beta$ and
\[
  \mathrm{ES}_\beta(L)=\E\bigl[L\,\big|\,L\ge\mathrm{VaR}_\beta\bigr]
  =\frac{1}{1-\beta}\int_{\mathrm{VaR}_\beta}^\infty x\,f_L(x)\dd x .
\]
\end{definition}
VaR answers ``what is the loss that I am $\beta$ sure not to exceed over the next 10 days?'' (e.g.\ $\$10$M at
$\beta=99\%$); ES answers ``if I am in the worst $1-\beta$ of cases, how much do I lose on average?''
\ts{24}{10}{10:31}. Clearly $\mathrm{ES}_\beta\ge\mathrm{VaR}_\beta$, since ES averages quantiles $\mathrm{VaR}_u\ge
\mathrm{VaR}_\beta$ \ts{24}{10}{12:35}. VaR was popularised by JPMorgan in the 1990s and ES appeared around 2000;
both are used heavily by regulators and in counterparty risk.

\begin{figure}[ht]
\centering
\begin{tikzpicture}
\begin{axis}[width=0.85\linewidth,height=5.2cm,xmin=-4,xmax=6.5,ymin=0,ymax=0.43,
  xlabel={loss $L$},ylabel={density},ytick=\empty,
  tick label style={font=\scriptsize},label style={font=\small},
  declare function={fb(\x)=0.97*exp(-\x*\x/2)/sqrt(2*pi)+0.03*exp(-(\x-5)*(\x-5)/0.18)/(0.3*sqrt(2*pi));
                    fr(\x)=exp(-\x*\x/(2*1.53983))/(1.24092*sqrt(2*pi));}]
\addplot[fill=ideacol,fill opacity=0.25,draw=none,domain=2.041:6.5,samples=120]{fb(x)}\closedcycle;
\addplot[fill=thmcol,fill opacity=0.25,draw=none,domain=2.041:6.5,samples=80]{fr(x)}\closedcycle;
\addplot[ideacol,thick,domain=-4:6.5,samples=200]{fb(x)};
\addplot[thmcol,thick,domain=-4:6.5,samples=200]{fr(x)};
\draw[dashed] (axis cs:2.041,0) -- (axis cs:2.041,0.30) node[above,font=\scriptsize]{$\mathrm{VaR}_{95\%}=2.04$};
\draw[thmcol,thick,->] (axis cs:2.56,0.13) node[above,font=\scriptsize]{ES $2.56$} -- (axis cs:2.56,0.005);
\draw[ideacol,thick,->] (axis cs:3.96,0.13) node[above,font=\scriptsize]{ES $3.96$} -- (axis cs:3.96,0.005);
\end{axis}
\end{tikzpicture}
\caption{Same VaR, different tails (in the spirit of slide~8 of \file{lec10}; numbers not from the course). Red:
$N(0,1.24^2)$. Blue: $97\%\ N(0,1)+3\%\ N(5,0.3^2)$, a book that usually makes small gains but occasionally
loses a lot. Both have $\mathrm{VaR}_{95\%}=2.04$, but the blue $\mathrm{ES}_{95\%}$ is $3.96$ against $2.56$:
VaR does not see what lies beyond the quantile.}\label{ch08:fig-fattail}
\end{figure}

\paragraph{Fat tails.} VaR is blind to the shape of the tail beyond the quantile (\cref{ch08:fig-fattail})
\ts{24}{10}{12:35}. Such a distribution is not artificial: a desk that sells out-of-the-money options collects a
small premium most of the time and occasionally suffers a large loss when the options end in the money
\ts{24}{10}{14:37}.

\paragraph{VaR is not subadditive.} Slide~9 \ts{24}{10}{14:37}: two independent portfolios A and B each lose $100$
with probability $4\%$ and $0$ otherwise. For each, $\mathrm{VaR}_{0.95}=0$ (the $4\%$ tail lies beyond the
quantile) and $\mathrm{ES}_{0.95}=(4\cdot100+1\cdot0)/5=80$. The sum A+B loses $200$ w.p.\ $0.16\%$, $100$ w.p.\
$7.68\%$, $0$ w.p.\ $92.16\%$, so $\mathrm{VaR}_{0.95}=100>0+0$ while $\mathrm{ES}_{0.95}=(0.16\cdot200+4.84\cdot100)/5
=103.2\le160$ \ts{24}{10}{15:39}. ES is subadditive in general (it is coherent; the slides point to a paper with
seven different proofs), but proving it takes some work. Conversely, for normal-like (elliptical) distributions
VaR \emph{is} subadditive, so in practice the failure is less dramatic than such contrived examples suggest
\ts{24}{10}{16:45}. For optimisation, however, it is fatal, because of the following fact.

\begin{proposition}[Subadditivity and convexity]\label{ch08:prop-convex}
Let $\rho$ be positively homogeneous. Then $\rho$ is subadditive iff it is convex.
\end{proposition}
\begin{proof}
If $\rho$ is subadditive, then for $\lambda\in[0,1]$:
$\rho(\lambda L_1+(1-\lambda)L_2)\le\rho(\lambda L_1)+\rho((1-\lambda)L_2)=\lambda\rho(L_1)+(1-\lambda)\rho(L_2)$.
If $\rho$ is convex: $\rho(L_1+L_2)=2\rho\bigl(\tfrac12L_1+\tfrac12L_2\bigr)\le\rho(L_1)+\rho(L_2)$.
\end{proof}
Without homogeneity the equivalence fails: $\sqrt{|x|}$ is subadditive but not convex. The slides (slide~33) state
``convexity is equivalent to sub-additivity''; the proposition makes precise when this is true.
Since VaR is not subadditive, it is not convex as a function of the portfolio, and non-convex functions are much
harder to minimise than convex ones \ts{24}{10}{16:45}.

\paragraph{Spectral risk measures.} A risk measure can be described by the weight it puts on each quantile,
$M_\phi(L)=\int_0^1\phi(u)\,\mathrm{VaR}_u(L)\dd u$ with $\phi\ge0$, $\int\phi=1$ \ts{24}{10}{17:45}. VaR puts all
its weight on $u=\beta$ (a delta function); ES puts a flat weight $1/(1-\beta)$ on $u>\beta$; an
\emph{exponential} spectral measure uses $\phi(u)\propto e^{-(1-u)/\gamma}$, weighting extreme losses most.
$M_\phi$ is coherent iff $\phi$ is non-decreasing (Acerbi, 2002): a sensible risk-averse measure never gives
worse losses less weight. This fails for VaR, whose weight drops to zero beyond $\beta$.

\begin{intuition}[Quantile versus conditional mean]
VaR is a quantile, like a median: it does not move if you drag the extreme tail further out. ES is the first moment
of the tail, so it does. Convexity is the property that makes a minimisation problem a bowl with one bottom (the
same role convexity of the free energy plays for thermodynamic stability); a risk measure without it can have
many local minima in portfolio space.
\end{intuition}

\begin{exercise}[VaR is not subadditive (not from the course)]
\leavevmode
\begin{enumerate}[(a)]
  \item Verify all the numbers of the A, B, A+B example (VaR and ES at $95\%$).
  \item Replace the $4\%$ loss probability by
      $p$ and find the range of $p$ for which $\mathrm{VaR}_{0.95}(A+B)>\mathrm{VaR}_{0.95}(A)+\mathrm{VaR}_{0.95}(B)$.
  \item Show that for jointly normal losses VaR is subadditive at any $\beta\ge\frac12$.
\end{enumerate}
\end{exercise}

\subsection{Normal model and time scaling}
\begin{proposition}[Normal VaR and ES]\label{ch08:prop-normal}
If $L\sim N(\mu,\sigma^2)$ and $z_\beta=N^{-1}(\beta)$, with $\varphi$ the standard normal density, then
\[
  \mathrm{VaR}_\beta=\mu+\sigma z_\beta,\qquad \mathrm{ES}_\beta=\mu+\sigma\,\frac{\varphi(z_\beta)}{1-\beta}
  =\mu+\sigma\,\frac{e^{-z_\beta^2/2}}{\sqrt{2\pi}\,(1-\beta)}\ts{24}{10}{18:45}.
\]
\end{proposition}
\begin{proof}
$L=\mu+\sigma Z$ and quantiles are equivariant under increasing affine maps. For ES,
$\E[Z\mid Z\ge z]=\frac{1}{1-\Phi(z)}\int_z^\infty x\varphi(x)\dd x=\frac{\varphi(z)}{1-\Phi(z)}$, since
$x\varphi(x)=-\varphi'(x)$.
\end{proof}
Numerically, $z_{0.99}=2.326$ and $\varphi(z_{0.99})/0.01=2.665$, while $\varphi(z_{0.975})/0.025=2.338$. Hence
under normality the $99\%$ VaR ($\mu+2.326\sigma$) and the $97.5\%$ ES ($\mu+2.338\sigma$) are almost the same
number: this is how the regulators' switch from 99\% VaR to 97.5\% ES was calibrated (slide~17).

\paragraph{Square root of time.} If daily changes $\Delta P_i$ are i.i.d.\ $N(0,\sigma^2)$, the $T$-day change has
standard deviation $\sigma\sqrt T$, so $T$-day VaR $=\sqrt T\times$ one-day VaR, and the same for ES
\ts{24}{10}{19:47}. Real returns are autocorrelated: a bad day tends to be followed by another.

\begin{proposition}[Autocorrelated daily changes]\label{ch08:prop-autocorr}
If $(\Delta P_i)$ is stationary with variance $\sigma^2$ and $\Corr(\Delta P_i,\Delta P_{i+k})=\rho^k$ (an AR(1)
structure), then
\[
  \Var\Bigl(\sum_{i=1}^T\Delta P_i\Bigr)=\sigma^2\Bigl[T+2\sum_{k=1}^{T-1}(T-k)\rho^k\Bigr],
\]
and for normal changes with zero mean the ratio of $T$-day to one-day VaR (or ES) is the square root of the bracket.
For $T=2$ it is $\sqrt{2(1+\rho)}$, as on slide~11.
\end{proposition}
\begin{proof}
Expand the variance of the sum: there are $T$ diagonal terms and $2(T-k)$ pairs at lag $k$.
\end{proof}
\begin{center}\small
\begin{tabular}{@{}lcccccc@{}}
\toprule
 & $T=1$ & $2$ & $5$ & $10$ & $50$ & $250$\\
\midrule
$\rho=0$    & 1.00 & 1.41 & 2.24 & 3.16 & 7.07 & 15.81\\
$\rho=0.05$ & 1.00 & 1.45 & 2.33 & 3.31 & 7.43 & 16.62\\
$\rho=0.1$  & 1.00 & 1.48 & 2.42 & 3.46 & 7.80 & 17.47\\
$\rho=0.2$  & 1.00 & 1.55 & 2.62 & 3.79 & 8.62 & 19.35\\
\bottomrule
\end{tabular}
\end{center}
The table of slide~11 is reproduced exactly by the formula above. For the EURO STOXX the observed lag-one
autocorrelation is about $0.1$ (between $0.05$ and $0.15$ for other indices) \ts{24}{10}{20:54}, so the usual
10-day rule $\sqrt{10}=3.16$ underestimates the true factor $3.46$ by about $10\%$; yet almost everybody uses
$\sqrt T$ \ts{24}{10}{21:59}. (The observed higher-lag autocorrelations on the slide are close to $0$ rather than
$\rho^k$; with only a lag-one correlation the factor is $\sqrt{10+2\cdot9\cdot0.1}=3.44$, practically the same.)

\begin{exercise}[Normal model and FRTB (not from the course)]
\leavevmode
\begin{enumerate}[(a)]
  \item Prove \cref{ch08:prop-normal}.
  \item Find the confidence level $\beta'$ at which the normal ES equals the normal
      $99\%$ VaR, and compare it with the $97.5\%$ used by FRTB.
  \item With $\bm\delta=(3000,4000,2000,1000)$ and the covariance
      matrix of slide~16 (rows $(1.02,0.97,0.80,0.94)$, $(0.97,1.08,0.82,0.96)$, $(0.80,0.82,0.89,0.76)$,
      $(0.94,0.96,0.76,0.96)$, all $\times10^{-4}$), reproduce $\sigma_P=96.6$, $\mathrm{VaR}=225$, $\mathrm{ES}=257$, and
      compute the 10-day values with and without the autocorrelation correction for $\rho=0.1$.
\end{enumerate}
\end{exercise}

\subsection{Estimating VaR and ES}\label{ch08:sec-estimation}

\begin{casestudy}[Historical simulation on four European indices (slides 12--15)]
\textbf{Data} \ts{24}{10}{21:59}. Daily closes of IBEX~35 (Spain), CAC~40 (France), DAX (Germany) and EURO STOXX~50,
2006--2023, containing two stress periods (the 2008 crisis and COVID in 2020). All four are in euros, so there is no FX
conversion; holidays differ across countries: days when all markets were closed were dropped, otherwise the
missing prices were forward-filled, which creates zero returns \ts{24}{10}{23:01}. (Shepherd's aside: in the
industry a surprising share of time goes into holiday calendars and FX.)

\textbf{Method} \ts{24}{10}{24:02}. With $v_i^j$ the level of index $j$ on day $i$ and positions
$\bm\delta=(3000,4000,2000,1000)$ in CAC, DAX, IBEX, EURO STOXX, the P\&L of day $i$ is
$\Delta P_i=\sum_j\delta_j\,(v_i^j-v_{i-1}^j)/v_{i-1}^j$. For a window of $n=500$ days (two years), sort the
losses: the $99\%$ VaR is the 5th worst ($1\%$ of $500$) and the $99\%$ ES the mean of the 4 worst; the $97.5\%$ VaR
is the 12th worst and the ES the mean of the 11 worst \ts{24}{10}{27:07}. (Some include the VaR scenario itself in
the average; the market convention is to exclude it \ts{24}{10}{28:10}.)

\textbf{Results} (slides 14--15; checked against the listed 12 largest losses):
\begin{center}\small
\begin{tabular}{@{}lcccc@{}}
\toprule
window & $\mathrm{VaR}_{99}$ & $\mathrm{ES}_{99}$ & $\mathrm{VaR}_{97.5}$ & $\mathrm{ES}_{97.5}$\\
\midrule
2022--2023 (recent)   & 320 & 377 & 227 & 303\\
2020--2021 (COVID)    & 450 & 815 & 306 & 534\\
2015--2016 (calm)     & 297 & 477 & 246 & 348\\
2007--2008 (crisis)   & 448 & 689 & 370 & 505\\
\bottomrule
\end{tabular}
\end{center}
\textbf{Lesson.} The choice of the historical window matters far more than any other modelling choice
\ts{24}{10}{26:07}. This is why regulators added \emph{stressed VaR}: besides the most recent window, VaR is also
computed on a fixed period known to be stressful (2008 or COVID) \ts{24}{10}{27:07}.
\end{casestudy}

\paragraph{Model building (variance--covariance).} Assume index returns $\Delta x_i\sim N(0,\sigma_i^2)$ with
correlations $\rho_{ij}$, covariance $C_{ij}=\rho_{ij}\sigma_i\sigma_j$, and a linear portfolio
$\Delta P=\sum_i\delta_i\Delta x_i$ with sensitivities $\delta_i=\partial P/\partial x_i$ \ts{24}{10}{28:10}. Then
$\Delta P\sim N(0,\sigma_P^2)$ with
\[
  \sigma_P^2=\sum_{i,j}\rho_{ij}\delta_i\delta_j\sigma_i\sigma_j=\bm\delta\T C\bm\delta,\qquad
  \mathrm{VaR}(T,\beta)=\sqrt T\,z_\beta\,\sigma_P,\qquad
  \mathrm{ES}(T,\beta)=\sqrt T\,\frac{\varphi(z_\beta)}{1-\beta}\,\sigma_P .
\]
With the 2022--2023 daily covariance matrix of slide~16 (entries about $10^{-4}$, i.e.\ daily volatilities of
about $1\%$) one gets $\sigma_P=96.6$, $\mathrm{VaR}(1,99\%)=225$ and $\mathrm{ES}(1,99\%)=257$ \ts{24}{10}{29:13},
well below the historical $320$ and $377$ for the same period: the empirical distribution has fatter tails than the
normal \ts{24}{10}{30:15}.

\begin{coursenote}[Erratum: ES formula on slide~16 of \file{lec10}]
The slide writes $\mathrm{ES}(T,\beta)=\sqrt T\,e^{-N^{-1}(\beta)}\sigma_P/(\sqrt{2\pi}(1-\beta))$. The exponent
must be $-[N^{-1}(\beta)]^2/2$, as on slide~11 and in \cref{ch08:prop-normal}. The quoted $\mathrm{ES}=257$ was
computed with the correct formula ($2.665\times96.6$); the misprinted one would give $376$.
\end{coursenote}

\subsection{Back-testing}
A VaR model can be tested by counting \emph{exceptions}, days on which the realised loss exceeded the VaR
\ts{24}{10}{30:15}. If the model is right, exceptions are independent Bernoulli trials with probability
$p=1-\beta$, so the number $N$ of exceptions in $n$ days is $\mathrm{Bin}(n,p)$ and
\[
  \Prob(N\ge m)=\sum_{k=m}^n\binom nk p^k(1-p)^{n-k}.
\]
The rule of thumb is to reject the model if this probability is below $5\%$ \ts{24}{10}{31:18}. In the example the
normal $\mathrm{VaR}(1,99\%)=225$ was exceeded on $12$ of the $500$ days, while $5$ were expected.

\begin{coursenote}[Erratum: the back-test on slide~16 and in class]
The slide's formula has $(1-p)^k$; the binomial probability needs $(1-p)^{n-k}$. More importantly, for
$n=500$, $m=12$, $p=0.01$ the correct value is $\Prob(N\ge12)=0.52\%$, not $5.2\%$. By the stated $5\%$ rule the
normal model is therefore \emph{rejected} (even at the $1\%$ level), whereas in class it was accepted as ``not
overly aggressive'' \ts{24}{10}{1:20:34}. At the $5\%$ level the model is accepted with up to $9$ exceptions
($\Prob(N\ge9)=6.7\%$) and rejected from $10$ on ($\Prob(N\ge10)=3.1\%$). The rejection also agrees with the
comparison with the historical VaR ($225$ against $320$). Note that this back-test is in-sample (covariance and
exceptions from the same window), which if anything favours the model.
\end{coursenote}

Back-testing VaR is easy; back-testing ES is not, and this is the main reason VaR is still used despite its flaws
\ts{24}{10}{32:23}. ES is a property of the whole tail of a distribution, while each day provides a single draw: one
realised loss cannot be compared with a tail mean without assumptions on the shape of the tail
\ts{24}{10}{1:17:27}. Regulation moved from VaR (Basel~I, 10-day $99\%$) to stressed VaR (Basel~2.5) to ES at
$97.5\%$ with variable liquidity horizons (FRTB, the Fundamental Review of the Trading Book). Banks still
back-test the VaR: they build one distributional model, compute capital from its ES, and accept the ES if the VaR of
the same model back-tests well \ts{24}{10}{1:16:23}. The ``model'' is the distributional assumption; VaR and ES are
two numbers read off it \ts{24}{10}{1:18:28}.

\begin{exercise}[Back-testing (not from the course)]
\leavevmode
\begin{enumerate}[(a)]
  \item Compute $\Prob(N\ge12)$ for $N\sim\mathrm{Bin}(500,0.01)$ exactly and with the Poisson approximation.
  \item Find the
      largest number of exceptions accepted at the $5\%$ level.
  \item Why is the in-sample back-test of the example biased in
      favour of the model?
  \item Explain why a single observed loss cannot ``test'' an ES forecast, and propose a test that
      uses the ES forecasts and realised losses of many days.
\end{enumerate}
\end{exercise}

\subsection{Expected shortfall as an optimisation problem}
Let $x\in X\subset\R^n$ be a portfolio, $y\in\R^m$ the random market change with density $p(y)$, and $f(x,y)$ the
loss. Write $\psi(x,\alpha)=\Prob(f(x,y)\le\alpha)$, so $\mathrm{VaR}_\beta(x)=\min\{\alpha:\psi(x,\alpha)\ge\beta\}$.

\begin{theorem}[Rockafellar--Uryasev]\label{ch08:thm-ru}
Let
\[
  F_\beta(x,\alpha)=\alpha+\frac{1}{1-\beta}\int_{\R^m}\bigl[f(x,y)-\alpha\bigr]^+p(y)\dd y .
\]
If $\psi(x,\cdot)$ is continuous, then $F_\beta(x,\cdot)$ is convex and continuously differentiable,
$\mathrm{VaR}_\beta(x)$ is a minimiser, and $\min_\alpha F_\beta(x,\alpha)=\mathrm{ES}_\beta(x)$. Consequently
\[
  \min_{x\in X}\mathrm{ES}_\beta(x)=\min_{(x,\alpha)\in X\times\R}F_\beta(x,\alpha),
\]
and if $f(\cdot,y)$ is convex and $X$ convex, this is a convex problem \ts{24}{10}{33:27}.
\end{theorem}
\begin{proof}
$\alpha\mapsto[f-\alpha]^+$ is convex, so $F_\beta$ is convex in $\alpha$, and
\[
  \frac{\partial F_\beta}{\partial\alpha}=1-\frac{\Prob(f>\alpha)}{1-\beta}=1+\frac{\psi(x,\alpha)-1}{1-\beta}
  =\frac{\psi(x,\alpha)-\beta}{1-\beta},
\]
which vanishes at $\alpha_\beta=\mathrm{VaR}_\beta(x)$. There
\[
  \int[f-\alpha_\beta]^+p\dd y=\int_{f\ge\alpha_\beta}f\,p\dd y-\alpha_\beta\bigl(1-\psi(x,\alpha_\beta)\bigr)
  =(1-\beta)\mathrm{ES}_\beta-\alpha_\beta(1-\beta),
\]
so $F_\beta(x,\alpha_\beta)=\alpha_\beta+\mathrm{ES}_\beta-\alpha_\beta=\mathrm{ES}_\beta$ (slides 18--20). If $f$ is jointly
convex in $(x,\alpha)$, so is $[f-\alpha]^+$ and hence $F_\beta$; minimising first over $\alpha$ and then over $x$ is
the same as minimising jointly.
\end{proof}
With $S$ equally likely scenarios $y_s$ (historical or simulated) and linear losses $f(x,y_s)=-x\T r_s$, the problem
becomes a \emph{linear program}:
\[
  \min_{x,\alpha,u}\ \alpha+\frac{1}{(1-\beta)S}\sum_{s=1}^Su_s\quad\text{s.t.}\quad u_s\ge -x\T r_s-\alpha,\ \ u_s\ge0,
  \ \ x\in X ,
\]
``basically a linear problem'' \ts{24}{10}{34:27}. No such representation exists for VaR. So ES is harder to
back-test but much easier to optimise, and for margin optimisation the latter matters more.

\begin{exercise}[ES as a linear program (not from the course)]
\leavevmode
\begin{enumerate}[(a)]
  \item Complete the proof of \cref{ch08:thm-ru} for a discrete distribution with $S$ equally likely scenarios, taking
      $\beta S$ to be an integer, and check that $\min_\alpha F$ is the average of the $(1-\beta)S$ largest losses.
  \item Write the LP that minimises the $97.5\%$ ES of a long-only fully invested portfolio of the four indices over the
      500 historical scenarios of 2022--2023.
  \item Why can the same trick not be used for VaR?
\end{enumerate}
\end{exercise}

% ==================================================================
\section{Margin: variation margin, initial margin and ISDA SIMM}\label{ch08:sec-margin}

\subsection{Why margin}
Consider a bank A that is long a swap with C and short an offsetting swap with B (slide~21). A receives cash flows
$\{T_i\}$ and pays $\{S_i\}$ on each swap, typically every three or six months; before the next payment date the
counterparty may default and not pay \ts{24}{10}{35:33}. To mitigate this the parties exchange \emph{variation
margin} daily, equal to the current value $V$ of the trades \ts{24}{10}{36:34}.

If C defaults it stops paying VM. A must replace or unwind the trade with C, which takes some days; meanwhile its
hedge with B keeps moving and A keeps paying VM to B. If the value of the defaulted trade to A rises in the
meantime, A has to pay more for the replacement than the collateral it holds, and that is a loss
\ts{24}{10}{37:35}. \emph{Initial margin} covers this loss.

\begin{definition}[Margin]
\emph{Margin} is collateral posted by one counterparty to the one that collects it, to cover potential losses on
derivatives (slide~22).
\begin{itemize}
  \item \emph{Variation margin} (VM) covers the current present value.
  \item \emph{Initial margin} (IM) covers the potential change of value between the counterparty's default and the
        time $t=\delta t$ at which the surviving party can exit the position. The period $\delta t$ is the
        \emph{margin period of risk} (MPOR).
\end{itemize}
IM is sized as a VaR (or ES) of the value change over the MPOR \ts{24}{10}{38:39}.
\end{definition}

Slide~23 shows it as a picture: a value path up to the default, a fan of possible paths over the MPOR, and a
distribution at close-out whose upper tail is covered by the IM collected and whose lower tail is the IM posted.
IM is held in segregated accounts and cannot be used for trading, so it is a real cost (the MVA of
\cref{ch08:sec-cva}) \ts{24}{10}{1:12:07}.

There are two kinds of IM \ts{24}{10}{38:39}. \emph{Cleared} IM is set by each CCP's own model and has existed for
decades. \emph{Bilateral} IM on uncleared trades was introduced by the post-crisis rules (Basel~III/BCBS--IOSCO) and
phased in from 2016; the industry standard model is ISDA's \emph{Standard Initial Margin Model} (SIMM).

\begin{exercise}[Tang's puzzle with collateral (not from the course)]
Redo the puzzle of \cref{ch08:sec-cva} (value $0\to+100$, default with $R=50\%$) when the trade is under a CSA with
daily VM and zero threshold, the counterparty stops paying VM at default, and the value moves from $100$ to $104$
during a 10-day MPOR. What is the loss with and without an initial margin of $5$? What part of the original $\$50$M
loss does VM remove, and what part does IM remove?
\end{exercise}

\subsection{The SIMM as a delta--gamma VaR}
SIMM is an exposure-based proxy for a 10-business-day ($\delta t=14/365$ in calendar time), $99\%$ VaR
\ts{24}{10}{39:41}. Let $P(t)$ be the portfolio value to the surviving party, the counterparty defaulting at $t=0$,
with $\mathrm{VM}(0)=P(0)$. Ignoring path dependence, cash flows in $(0,\delta t)$ and discounting, and letting only
the market state $\mathcal M$ move by $\Delta\mathcal M$ over the MPOR, a second-order expansion gives
\[
  P(\delta t)\approx P(0)+\bm\delta\T\Delta\mathcal M+\tfrac12\,\Delta\mathcal M\T\gamma\,\Delta\mathcal M ,
\]
with $\bm\delta$ the vector of deltas and $\gamma$ the (symmetric) gamma matrix. The loss is the \emph{increase}
$P(\delta t)-P(0)$, since the collateral is frozen at $P(0)$, and IM should satisfy
$\Prob\bigl(P(\delta t)-P(0)<\mathrm{IM}\bigr)=0.99$.

\begin{proposition}[Moments of the delta--gamma P\&L]\label{ch08:prop-simm}
If $\Delta\mathcal M\sim N(0,C)$, then $\Delta P=\bm\delta\T\Delta\mathcal M+\frac12\Delta\mathcal M\T\gamma\Delta\mathcal M$ has
\[
  \E[\Delta P]=\tfrac12\tr(\gamma C),\qquad \Var(\Delta P)=\bm\delta\T C\bm\delta+\tfrac12\tr\bigl((\gamma C)^2\bigr).
\]
Approximating $\Delta P$ by a normal with these moments,
$\mathrm{IM}\approx\frac12\tr(\gamma C)+2.326\sqrt{\bm\delta\T C\bm\delta+\frac12\tr((\gamma C)^2)}$.
\end{proposition}
\begin{proof}
Write $X=\Delta\mathcal M$. $\E[X\T\gamma X]=\sum_{ij}\gamma_{ij}C_{ij}=\tr(\gamma C)$. By Isserlis' (Wick's) theorem,
$\E[X_iX_jX_kX_l]=C_{ij}C_{kl}+C_{ik}C_{jl}+C_{il}C_{jk}$, so
\[\E[(X\T\gamma X)^2]=(\tr\gamma C)^2+2\tr(\gamma C\gamma C),\qquad \Var(X\T\gamma X)=2\tr((\gamma C)^2);\]
the factor
$\frac14$ gives $\frac12\tr((\gamma C)^2)$. The cross term $\Cov(\bm\delta\T X,X\T\gamma X)$ involves third moments of a
centred Gaussian and vanishes.
\end{proof}
Without gamma this is exactly the model-building VaR of \cref{ch08:sec-estimation} with $T=10$, $\beta=0.99$:
$\mathrm{IM}=2.326\sqrt{10}\sqrt{\bm\delta\T C_1\bm\delta}$ for a daily covariance $C_1$ \ts{24}{10}{39:41}. The
gamma (convexity) terms are the ``two additional terms'' of the slide: SIMM is ``basically VaR''.

\begin{coursenote}[Caveat on slide~25 of \file{lec10}]
The slide gives $\mathrm{IM}\approx2.326\sqrt{10}\sqrt{\bm\delta\T C\bm\delta}+\frac12\sqrt{10}\tr(\gamma C)
+\lambda\sqrt{\frac12\tr((\gamma C)^2)}$. Two remarks.
\begin{itemize}
  \item Splitting the square root of a sum into a delta term and a
      curvature term, with a separately calibrated multiplier $\lambda$, is a modelling choice (SIMM treats delta and
      curvature margins separately), not a consequence of the derivation; moreover $\Delta P$ is not Gaussian, so
      $\lambda$ absorbs the skew of the quadratic term.
  \item The time scaling is inconsistent: if $C$ is the \emph{daily}
      covariance, the 10-day covariance is $10C$, so the mean term is $\frac12\tr(\gamma\,10C)=5\tr(\gamma C)$, not
      $\frac12\sqrt{10}\tr(\gamma C)$, and the last term scales by $10$, not by $1$. Deltas scale like $\sqrt T$, gammas
      like $T$.
\end{itemize}
\end{coursenote}

\begin{exercise}[Delta--gamma margin (not from the course)]
\leavevmode
\begin{enumerate}[(a)]
  \item Prove $\Var(X\T\gamma X)=2\tr((\gamma C)^2)$ for $X\sim N(0,C)$ and symmetric $\gamma$.
  \item One risk factor with
      daily variance $c$, delta $\delta$ and gamma $\gamma$: write mean and variance of the 10-day $\Delta P$ and check how each
      term scales with the horizon.
  \item For a short option position ($\gamma<0$ in the value to the surviving party) discuss
      whether the normal approximation over- or underestimates the $99\%$ quantile.
\end{enumerate}
\end{exercise}

\subsection{SIMM in practice}
The formulation above has too many risk factors (so $C$, $\bm\delta$, $\gamma$ are huge) and $\gamma$ is hard to
compute. SIMM therefore simplifies further (slides 26--27): curvature is proxied from \emph{vega},
$\gamma\propto\mathcal V/(T\sigma_N)$; risk factors are grouped into \emph{buckets} (e.g.\ currencies) and \emph{risk
classes} (IR, FX, equity, commodity, qualifying and non-qualifying credit); sensitivities are delta, vega and
curvature; the results are aggregated up a tree with prescribed correlations, and the \emph{product classes}
(rates/FX, credit, equity, commodity) are simply added with no netting between them. Additional terms handle
thresholds and concentration, and when several regulators apply the posted IM is the maximum over them.

\paragraph{Calibration} \ts{24}{10}{40:43}. The correlations and risk weights, which cover all asset classes, are
recalibrated on historical data. In the lecture (October 2024) SIMM~2.7, due in early December 2024, was expected to
be the first calibration without the 2008 or COVID periods in its window, cutting SIMM IM by roughly $5$--$20\%$
and releasing on the order of $\$40$bn of margin. When the 2020 shocks arrived, the annual calibration made SIMM
underestimate risk for more than a year; ISDA, which owns the model and runs the calibration with the banks, then
moved to semi-annual recalibration \ts{24}{10}{1:13:11}.

\paragraph{Large positions} (slide~28). Big positions take longer to unwind, which is liquidity risk again. One
approach scales the MPOR with the position size, $\Delta T=\Delta T_0\max(1,N/N_0)$; since IM grows like
$N\sqrt{\Delta T}$, beyond $N_0$ one gets $\mathrm{IM}\propto N^{3/2}$. Another adds a concentration add-on
$W(\bm\delta)\,\bm\delta\T C\bm\delta$ with $W$ increasing (e.g.\ piecewise linear). Both keep the IM convex in the
position.

% ==================================================================
\section{Multilateral optimisation of initial margin}\label{ch08:sec-opt}
This is the main topic of the 2024 lecture and the business of Quantile.

\subsection{The problem}
Consider parties $\mathcal P=\{A,B,C,D,E,\dots\}$ (banks) with bilateral portfolios against each other and cleared
portfolios at one or more CCPs, all posting IM (slide~29) \ts{24}{10}{41:48}. The goal is to find a set of new
\emph{hedge trades} between pairs of parties that reduces the total IM in the system (slide~30)
\ts{24}{10}{42:50}. There are no hedge variables with the CCP: one can only book trades between two banks, but if
the chosen instrument is clearable (a vanilla swap) the trade is cleared and changes the cleared IM instead of the
bilateral one \ts{24}{10}{43:55}.

\begin{definition}[Margin minimisation problem, slide 31]\label{ch08:def-mmp}
Let $\mathcal H$ be the set of hedge instruments, $\mathcal M$ the market risk factors, $D^0_{p,q,m}$ the initial
sensitivity of party $p$'s portfolio facing $q$ to factor $m$, and $D_{h,p,q,m}$ the sensitivity of one unit of hedge
$h$. The unknowns are the amounts $x_{h,p,q}$ of hedge $h$ that $p$ trades with $q$. Solve
\begin{align*}
  \min_{x}\quad & \sum_{p\in\mathcal P}\Bigl(\sum_{q\ne p}\mathrm{SIMM}_{p,q}(D_{p,q})
                   +\sum_{c\in\mathcal{CCP}}\mathrm{IM}_{p,c}(D_{p,c})\Bigr)\\
  \text{s.t.}\quad & D_{p,q,m}=D^0_{p,q,m}+\textstyle\sum_{h\in\mathcal H}x_{h,p,q}D_{h,p,q,m} && \forall m\in\mathcal M,\\
  & x_{h,p,q}=-x_{h,q,p} && \forall h,\ p\ne q \quad\text{(symmetry)},\\
  & \textstyle\sum_{q}x_{h,p,q}=0 && \forall h,\ p \quad\text{(flatness)},\\
  & T^-_{p,q,m}\le\textstyle\sum_hx_{h,p,q}D_{h,p,q,m}\le T^+_{p,q,m} && \forall m,\ p\ne q\quad\text{(risk limits)}.
\end{align*}
\end{definition}
The objective is the total margin in the system: bilateral SIMM summed over all pairs, plus cleared IM summed over
CCPs \ts{24}{10}{43:55}. The new risk is the old risk plus the sum of the hedges times their risk
\ts{24}{10}{44:59}. The constraints:
\begin{itemize}
  \item \textbf{Symmetry}: if $p$ sells something to $q$, $q$ must buy exactly that from $p$. This was the
        constraint the banks worried about most: telling Goldman to buy 100 from JPMorgan while telling JPMorgan
        to sell 200 to Goldman would only be discovered when the trade is booked \ts{24}{10}{46:02}.
  \item \textbf{Flatness} (cash-flow neutrality): whatever $p$ sells to one party it buys from others, so each
        party's total position in every instrument, and hence its market risk, is unchanged. Only
        \emph{counterparty} risk is reshuffled, and nobody makes or loses money from the system itself. It is not
        strictly necessary, but relaxing it makes everything much harder \ts{24}{10}{47:02}.
  \item \textbf{Risk limits}: each party can bound how much the risk to each factor against each counterparty may
        change \ts{24}{10}{48:06}.
\end{itemize}

\begin{coursenote}[Notation on slide~31]
The slide writes $\mathrm{IM}_{p,ccp}(D_{p,q})$; the argument should be the cleared sensitivity $D_{p,ccp}$. The
risk update is written for bilateral risk only; for clearable hedge types $h\in\mathcal H_c$ the update goes to the
CCP, $D_{p,c,m}=D^0_{p,c,m}+\sum_{h\in\mathcal H_c}\sum_qx_{h,p,q}D_{h,m}$, which is how cleared IM enters the
optimisation \ts{24}{10}{42:50}.
\end{coursenote}

\subsection{Toy model: the median rule}
\begin{figure}[ht]
\centering
\begin{tikzpicture}[>=Stealth,every node/.style={font=\small},
  pt/.style={circle,draw,thick,minimum size=6mm,inner sep=0pt},
  lab/.style={fill=white,inner sep=1pt,font=\footnotesize}]
\foreach \dx/\ab/\bc/\ca/\ttl in {0/{-2}/{-9}/{-10}/{initial (margin 21)},
                                  4.2/{+9}/{+9}/{+9}/{hedges},
                                  8.4/{7}/{0}/{-1}/{result (margin 8)}}{
  \begin{scope}[xshift=\dx cm]
    \node[pt] (A) at (1.2,2.0) {A}; \node[pt] (B) at (0,0) {B}; \node[pt] (C) at (2.4,0) {C};
    \draw[->,thick,defcol] (A) -- node[lab]{$\ab$} (B);
    \draw[->,thick,defcol] (B) -- node[lab]{$\bc$} (C);
    \draw[->,thick,defcol] (C) -- node[lab]{$\ca$} (A);
    \node at (1.2,-0.7) {\ttl};
  \end{scope}}
\node at (3.3,0.9) {\Large $+$};
\node at (7.5,0.9) {\Large $=$};
\end{tikzpicture}
\caption{The toy model of slide~32. An arrow $p\to q$ labelled $x$ means that $p$'s position against $q$ is $x$
(and $q$'s against $p$ is $-x$). Moving $t=9$ units around the cycle $A\to B\to C\to A$ satisfies symmetry and
flatness by construction and reduces the total $|x|$ from $2+9+10=21$ to $7+0+1=8$.}\label{ch08:fig-triangle}
\end{figure}

\begin{proposition}[Median rule for a triangle]\label{ch08:prop-median}
Three parties, one risk factor, one hedge instrument, bilateral margin proportional to $|$net position$|$ (the
one-factor SIMM $2.326\sqrt{10}\,\sigma|D|$). Let $a_1,a_2,a_3$ be the initial positions oriented around the cycle
$A\to B\to C\to A$. Then every hedge satisfying symmetry and flatness adds the same amount $t$ to all three oriented
edges, and the total margin $\sum_k|a_k+t|$ is minimised at $t=-\operatorname{median}(a_1,a_2,a_3)$.
\end{proposition}
\begin{proof}
Flatness at A gives $x_{AB}+x_{AC}=0$, i.e.\ $x_{AB}=x_{CA}$ by symmetry; at B, $x_{BA}+x_{BC}=0$, i.e.\
$x_{BC}=x_{AB}$. So $x_{AB}=x_{BC}=x_{CA}=t$. The function $g(t)=\sum_k|a_k+t|$ is convex and piecewise linear with
slope $\#\{k:a_k+t>0\}-\#\{k:a_k+t<0\}$, which changes sign at $t=-\operatorname{median}(a)$.
\end{proof}
In the example of slide~32, $a=(-2,-9,-10)$, the median is $-9$, and moving $9$ units around the triangle turns the
positions into $(7,0,-1)$, cutting the margin from $21$ to $8$ \ts{24}{10}{48:06}. Any network can be decomposed into
triangles, and one can think of the optimisation as moving the median around each triangle. Since edges are shared
between triangles and there are additional constraints, this is only a heuristic and not the real solution method
\ts{24}{10}{49:07}.

\begin{coursenote}[Erratum: result panel of slide~32]
In the ``result'' triangle the labels on the A--B edge are $-7$ next to A and $7$ next to B. Adding the ``initial''
($-2$ at A, $2$ at B) and ``hedges'' ($9$ at A, $-9$ at B) panels gives $+7$ at A and $-7$ at B. The margin, which
depends only on $|x|$, is unaffected.
\end{coursenote}

\begin{intuition}[Hedges are circulations on a graph]
Think of $x_{h,\cdot,\cdot}$ as a \emph{flow} on the complete graph of parties: symmetry says it is an antisymmetric
edge function, and flatness says its divergence vanishes at every node (Kirchhoff's current law). Divergence-free
flows are circulations, i.e.\ the cycle space of the graph, of dimension $E-V+1$; for $n$ fully connected parties
this is $\binom n2-n+1=\binom{n-1}{2}$, spanned by triangles. For a triangle it is one number, the $t$ above. The
optimisation moves risk around loops without creating or destroying it at any node.
\end{intuition}

\begin{exercise}[Toy margin optimisation (not from the course)]
\leavevmode
\begin{enumerate}[(a)]
  \item Four parties on a cycle $A\to B\to C\to D\to A$ with oriented positions $a=(3,-5,8,-1)$: find the optimal
      circulation $t$ and the margin saving.
  \item Now add the chord A--C with position $2$: parametrise all hedges satisfying
      symmetry and flatness (the cycle space has dimension $2$) and minimise the total $\sum|x|$.
  \item For the fairness table
      of \cref{ch08:sec-opt}, check the least-squares criterion and propose a different fairness criterion (e.g.\ max--min
      of relative savings) that gives a different answer on some example.
\end{enumerate}
\end{exercise}

\subsection{Solving it in practice}
In reality one writes down the problem and hands it to a numerical solver \ts{24}{10}{50:07}. Some lessons:
\begin{itemize}
  \item \textbf{Convexity is everything.} Minimising ES is piecewise linear, SIMM is \emph{nearly} convex and
        concentration add-ons are \emph{potentially} convex (slide~33). Using general non-linear solvers did not
        work well; it is much easier to approximate the objective by a genuinely convex function and use a mature
        convex solver. Do not write your own \ts{24}{10}{51:07}. How hard a problem is depends on its structure
        and sparsity much more than on its size.
  \item \textbf{Modelling languages} separate the statement of a problem from the numerics. Quantile uses Gurobi;
        the slide's example is the integer program $\max 5x_1+8x_2$ s.t.\ $x_1+x_2\le5$, $3x_1+7x_2\le25$,
        $x_i\in\N$, solved in 10\,ms with optimum $31$ at $(3,2)$ \ts{24}{10}{52:08}.
  \item \textbf{Size.} The FX-delta problem (slide~36): 25 parties, 30 currencies, cleared and uncleared
        non-deliverable forwards (NDFs), one short maturity, about $18{,}000$ variables ($25\cdot24$ ordered pairs
        $\times30$ currencies) and $40{,}000$ constraints (pairwise risk limits, pairwise IM-increase limits, notional
        efficiency, symmetry, flatness, lot sizes), with SIMM plus cleared IM as a convex objective
        \ts{24}{10}{53:14}. NDFs are chosen because each has exposure to exactly one risk factor, one currency pair.
  \item \textbf{Linear algebra.} Interior-point methods take Newton steps: they find the zero of the gradient, and
        each step requires solving with the Hessian, naively $O(n^3)\approx10^{13}$ flops for $n\sim40{,}000$
        \ts{24}{10}{54:18}. In practice one does far better: the matrix is positive semidefinite (Cholesky), very
        sparse (well under $0.1\%$ non-zero) and structured. A few global rows (symmetry) and party rows (flatness)
        sit on top of a block-diagonal part, the $36{,}000$ risk constraints of individual $(p,q)$ pairs, which is
        purely diagonal when each hedge has one risk factor. This \emph{Dantzig--Wolfe} (block-angular) structure
        is exploited by the solvers \ts{24}{10}{55:18}.
  \item \textbf{Presolve} removes redundant rows ($x_1\ge1$ is implied by $x_1\ge2$) and here halved the problem;
        the barrier method then converged in 58 iterations and 11\,s, thanks to the quadratic convergence of Newton's
        method \ts{24}{10}{59:29}.
\end{itemize}

\paragraph{Nearly parallel risks.} The real difficulties are numerical \ts{24}{10}{56:22}. Two trades with almost
identical risk, e.g.\ exposures $3E$ and $E$ to two risk factors, give nearly singular systems (slide~39):
\[
  \begin{pmatrix}1&1\\ \tfrac13&\tfrac13\end{pmatrix}\binom{x_1}{y_1}=\binom{1}{\tfrac13}\ \text{has the solutions } y_1=1-x_1,
  \quad\text{but}\quad
  \begin{pmatrix}1&1\\ 0.3333&0.333\end{pmatrix}\binom{x_1}{y_1}=\binom{1}{b}
\]
has the unique solution $(0,1)$ for $b=0.333$, $(1,0)$ for $b=0.3333$, and $(-110,111)$ for $b=0.3$
\ts{24}{10}{57:25}. The determinant is $-3\times10^{-4}$ and the condition number about $7\times10^3$, so rounding in
the fourth digit moves the solution by $O(1)$ or even by hundreds. With a sign constraint $x_i\ge0$ the last system
is declared \emph{infeasible} purely because of numerical noise \ts{24}{10}{58:28}. Most of the practical work is
detecting and removing such situations, not solving convex problems in general.

\paragraph{Results.} On a 5-party and a 20-party system with one risk factor and no constraints, almost all margin
disappears, and every party ends up either long or short the factor against all its counterparties, never both:
all risk has been netted (slides 41--42) \ts{24}{10}{1:00:32}. With real constraints (slide~43) the picture is less
clean. The empty diagonal is ``margin against oneself'', and a block in the middle turned out to be the affiliates
of one large bank, which do not exchange margin with each other \ts{24}{10}{1:01:35}. A second application
(slide~44) is interest-rate margin: 15 parties, USD/EUR/GBP, swaptions and cleared swaps, about $5000$ variables
and $10{,}000$ constraints, where each trade loads on several risk factors (a swap expiring in 12 years is exposed to
the 10Y and 20Y swap rates).

\begin{exercise}[Rounding and integrality (not from the course)]
For the Gurobi example, $\max 5x_1+8x_2$ subject to $x_1+x_2\le5$, $3x_1+7x_2\le25$, $x_1,x_2\ge0$:
\begin{enumerate}[(a)]
  \item solve the LP
      relaxation (optimum $32.5$ at $(2.5,2.5)$);
  \item show that rounding the LP solution gives either an infeasible point or a
      poor one;
  \item find the integer optimum $31$ by enumeration. Relate this to the rounding of notionals in
      \cref{ch08:sec-opt}.
\end{enumerate}
\end{exercise}

\begin{exercise}[Nearly parallel risks (not from the course)]
For $A_\varepsilon=\bigl(\begin{smallmatrix}1&1\\ 1/3+\varepsilon&1/3\end{smallmatrix}\bigr)$ solve $A_\varepsilon(x,y)\T=(1,b)\T$
and show that $x=(b-1/3)/\varepsilon$. Compute the condition number of $A_\varepsilon$ for small $\varepsilon$ and explain the
three solutions of slide~39. How would you regularise the problem (e.g.\ a small penalty $\eta\|x\|^2$)?
\end{exercise}

\subsection{Nice and fair solutions}
A solution must be \emph{feasible} (all constraints hold) and \emph{optimal} (a good IM saving); in practice it
must also be \emph{nice} and \emph{fair} \ts{24}{10}{1:02:40}.

\paragraph{Nice.} Traders want round notionals (e.g.\ multiples of a million), few trades (\emph{notional
efficiency}), trade packages, and possibly more trades moved into CCPs for reasons beyond IM. The worst headache
came from NDFs: a USD notional and a EUR notional define the FX rate as their ratio, to 15 digits, but banks' booking
systems keep only 6--8 decimals, so the notionals had to be chosen as round dollar amounts that give exact euro
amounts at a 4- or 6-decimal rate \ts{24}{10}{1:03:41}. Rounding cannot be done afterwards. If a party sells 8 to one
counterparty and buys 4 from each of two others, rounding to the nearest 10 gives $10$, $0$, $0$ and breaks
flatness. Rounding must be part of the optimisation, which makes it a \emph{mixed-integer} problem, much harder:
the FX rounding step of slide~46 had about $10^4$ integer variables and took about a minute of branch-and-bound
\ts{24}{10}{1:04:43}.

\paragraph{Fair.} What does fairness mean for an optimiser \ts{24}{10}{1:05:47}? In the four-party example of
slide~47 three solutions give the same total saving:
\begin{center}\small
\begin{tabular}{@{}lcccc@{}}
\toprule
party & scenario 1 & scenario 2 & scenario 3 & single-party best\\
\midrule
A & 4 & 0 & 2 & 4\\
B & 4 & 4 & 4 & 4\\
C & 4 & 4 & 4 & 4\\
D & 0 & 4 & 2 & 4\\
\bottomrule
\end{tabular}
\end{center}
Scenarios 1 and 2 leave one party with nothing; scenario 3 splits the saving, but one could still ask why B and C get
their full 4 \ts{24}{10}{1:06:47}. One proposal is a two-step algorithm. First optimise for each party alone, keeping
all constraints, to get the best achievable $\mathrm{IM}_p^{\text{best}}$; then solve
$\min\sum_p(\mathrm{IM}_p-\mathrm{IM}_p^{\text{best}})^2$ \ts{24}{10}{1:07:53}. Here the squared shortfalls are
$16$, $16$ and $4+4=8$, so the algorithm picks scenario 3 \ts{24}{10}{1:09:00}. Fairness has no established definition
yet; Shepherd considers it the most interesting open problem in this area.

\paragraph{Network effect.} Across real runs, both the saving and the saving as a fraction of the initial IM grow with
the number of participating parties (slide~49): bigger networks benefit everyone \ts{24}{10}{1:10:04}.

\begin{finance}[Optimisation versus compression (not discussed in class)]
\emph{Portfolio compression} terminates redundant offsetting trades (e.g.\ circles of swaps A--B--C--A) and replaces
them with fewer trades of smaller gross notional and the same net risk. The IM optimisation of this lecture instead
\emph{adds} new trades that redistribute risk between pairs of counterparties. Both keep every party's market risk
unchanged (the flatness constraint) and exploit the same cycle structure of the network.
\end{finance}

% ==================================================================
\section{Enterprise-level modelling and martingale tools}\label{ch08:sec-martingale}
The second half of the 2013 lecture turns to the models behind CVA.

\subsection{Trade-level versus enterprise-level models}
A \emph{trade-level} model prices each trade independently (PV and Greeks), and the portfolio PV and Greeks are the
sums: this is what ``derivatives model'' usually means \ts{13}{26}{44:53}. \emph{Enterprise-level} models handle the
quantities that are \emph{not} linear in the trades: CVA, funding, capital (risk-weighted assets), liquidity
\ts{13}{26}{45:57}. They must simulate all trades and all market factors of a portfolio consistently, with a proper
joint distribution, while reusing the trade-level models to value each trade at future dates.

There is also feedback to the trade-level model \ts{13}{26}{47:02}. A cancellable swap facing a counterparty close to
default has an exercise boundary that ignores counterparty credit; cancelling early reduces or removes the CVA. In
practice the CVA desk pays the trading desk the CVA saving, which the trade model can take as an input to its
exercise decision \ts{13}{26}{48:05}.

Putting models together is numerically delicate: some trades are priced by PDE, some by simulation, and a joint
simulation introduces errors \ts{13}{26}{50:14}. Tang's remedy is built on martingales: the arbitrage-free price of
every traded quantity is a martingale under a suitable measure, and those conditions can be tested and enforced.

\subsection{Martingale measures and change of numeraire}
By Harrison and Pliska, if $Y_t$ is the price of a traded asset with no intermediate cash flows and $N_t>0$ is a
\emph{numeraire} (another such price), then $Y_t/N_t$ is a martingale under the measure $\Prob_N$ associated with
$N$ \ts{13}{26}{52:18}. The standard pairs (slides 13--15; see \cref{ch:bs,ch:hjm,ch:quanto}) are:
\begin{center}\small
\begin{tabular}{@{}lll@{}}
\toprule
numeraire & measure & martingale\\
\midrule
$\beta(t)=e^{\int_0^tr_u\dd u}$ & risk-neutral $\Q$ & $Y_t/\beta(t)$\\
$P(t,T)$ (zero-coupon bond) & $T$-forward $\Prob_T$ & forward price $F_Y(t,T)=Y_t/P(t,T)$\\
$P(t,T_{i+1})$ & $\Prob_{T_{i+1}}$ & forward LIBOR $L_i(t)=\frac{1}{\delta_i}\bigl(\frac{P(t,T_i)}{P(t,T_{i+1})}-1\bigr)$\\
$A_{i,N}(t)=\sum_{j=i+1}^N\delta_{j-1}P(t,T_j)$ & annuity $\Prob_{A_{i,N}}$ & swap rate $S_{i,N}(t)=\frac{P(t,T_i)-P(t,T_N)}{A_{i,N}(t)}$\\
$P^D(t,T)$ (domestic) & $\Prob_{T,D}$ & FX forward $F_{FX}(t,T)=\frac{S_{FX}(t)P^F(t,T)}{P^D(t,T)}$\\
\bottomrule
\end{tabular}
\end{center}
The pattern is always the same: write the quantity as a ratio of two traded assets \ts{13}{26}{55:28}. LIBOR is (a
multiple of) a portfolio of two bonds divided by the bond $P(t,T_{i+1})$ \ts{13}{26}{53:23}; the swap rate is a
portfolio of two bonds divided by the annuity \ts{13}{26}{54:27}; the FX forward is interest-rate parity, the foreign
bond converted into domestic currency divided by the domestic bond.

\begin{proposition}[Change of numeraire, slide 16]\label{ch08:prop-numeraire}
Let $N_1,N_2>0$ be numeraires and $Y_T$ an $\mathcal F_T$-measurable payoff. Then
\[
  \E_t^{N_1}[Y_T]=\frac{N_2(t)}{N_1(t)}\;\E_t^{N_2}\!\left[Y_T\,\frac{N_1(T)}{N_2(T)}\right],\qquad
  \text{i.e.}\quad \left.\frac{\dd\Prob_{N_1}}{\dd\Prob_{N_2}}\right|_{\mathcal F_T}=\frac{N_1(T)/N_1(0)}{N_2(T)/N_2(0)} .
\]
\end{proposition}
\begin{proof}
The price $X_t$ of the claim paying $X_T=Y_TN_1(T)$ does not depend on the measure used to compute it
\ts{13}{26}{57:32}:
$X_t=N_1(t)\E_t^{N_1}[X_T/N_1(T)]=N_2(t)\E_t^{N_2}[X_T/N_2(T)]$. Substitute $X_T$ and divide by $N_1(t)$.
\end{proof}
With $N_1=P(\cdot,T_i)$, $N_2=P(\cdot,T_{i+1})$ and $N_1(T)/N_2(T)=1+\delta_iL_i(T)$ for $T\le T_i$:
$\E_t^{T_i}[Y_T]=\E_t^{T_{i+1}}\bigl[Y_T(1+\delta_iL_i(T))\bigr]/(1+\delta_iL_i(t))$. This is the change of measure
behind the LIBOR market model \ts{13}{26}{56:29}.

\paragraph{Credit: the survival measure} \ts{13}{26}{58:38}. By analogy with the swap rate one expects the forward
CDS par spread $\tilde S_j(t,T_i,T_{i+1})$ of name $j$ to be a martingale under a ``risky annuity'' measure. The
catch is that the risky annuity $\tilde A_j$ (the value per unit spread of the premium leg, which pays only while
the name survives, i.e.\ a portfolio of zero-recovery defaultable bonds) drops to zero at default, and a numeraire
must stay positive \ts{13}{26}{59:40}. Sch\"onbucher's solution is to ignore the default states
\ts{13}{26}{1:00:45}: since $\ind_{\{t<\tau_j\}}\tilde S_j\tilde A_j$ is the value of the premium leg, a traded asset,
$\ind_{\{t<\tau_j\}}\tilde S_j\tilde A_j/\beta(t)$ is a $\Q$-martingale, and the \emph{survival measure}
\[
  \left.\frac{\dd\Prob_{\tilde A_j}}{\dd\Q}\right|_{t}=\frac{\ind_{\{t<\tau_j\}}\tilde A_j(t,T_i,T_{i+1})}{\beta(t)\,\tilde A_j(0,T_i,T_{i+1})}
\]
makes $\tilde S_j(t,T_i,T_{i+1})$ a martingale \ts{13}{26}{1:03:55}. $\Prob_{\tilde A_j}$ is absolutely continuous
with respect to $\Q$ but \emph{not equivalent} (it gives zero mass to default before $t$). One difficulty is traded
for another, but the second is easier, provided what happens at default is modelled separately
\ts{13}{26}{1:01:49}.

\begin{exercise}[Change of numeraire (not from the course)]
\leavevmode
\begin{enumerate}[(a)]
  \item Prove \cref{ch08:prop-numeraire} and deduce that $L_i(t)$ is a $\Prob_{T_{i+1}}$-martingale.
  \item Show that the swap
      rate $S_{i,N}(t)$ is a martingale under the annuity measure and that it is a weighted average of forward LIBORs with
      weights $\delta_jP(t,T_{j+1})/A_{i,N}(t)$.
  \item Verify \cref{ch08:prop-mginterp} and show by an example that ordinary
      linear interpolation in maturity $T$ is not a martingale in general.
\end{enumerate}
\end{exercise}

\subsection{Martingale testing, resampling and interpolation}
\begin{itemize}
  \item \textbf{Martingale testing}: check in the numerical implementation that known martingale relations such as
        $F(s)=\E_s[F(t)]$ hold; usually they fail slightly \ts{13}{26}{1:04:58}.
  \item \textbf{Martingale resampling}: transform the simulated variables (linearly or log-linearly) so that they
        satisfy the relations \emph{exactly}, keeping a variance metric fixed. The quadratic resampling of slide~19
        is $X=\frac{\mathrm{Std}(X)}{\mathrm{Std}(X_0)}\bigl(X_0-\E[X_0]\bigr)+\E[X]$, with $X_0$ the raw simulated values,
        $\E[X]$ the martingale target and $\mathrm{Std}(X)$ the target dispersion \ts{13}{26}{1:05:59}. This is moment
        matching, familiar from variance reduction in Monte Carlo.
  \item \textbf{Martingale interpolation}: models simulate only a few tenors (say 1Y and 5Y LIBOR). Interpolating a
        missing tenor (3Y) linearly in maturity does \emph{not} preserve the martingale property; interpolating
        linearly using the time-$s$ values as the independent variable does \ts{13}{26}{1:07:01}.
\end{itemize}

\begin{proposition}[Martingale interpolation]\label{ch08:prop-mginterp}
Let $M(\cdot,T_1)$ and $M(\cdot,T_2)$ be martingales, $M(s,T)=\E_s[M(t,T)]$ for $s\le t$, and let the target
$M(s,T_3)$ be known at a reference time $s$ (e.g.\ $s=0$, from today's curve). Define for $t\ge s$
\[
  M(t,T_3)=\frac{M(s,T_3)-M(s,T_2)}{M(s,T_1)-M(s,T_2)}\,M(t,T_1)+\frac{M(s,T_1)-M(s,T_3)}{M(s,T_1)-M(s,T_2)}\,M(t,T_2).
\]
Then $\E_s[M(t,T_3)]=M(s,T_3)$ \ts{13}{26}{1:10:06}.
\end{proposition}
\begin{proof}
The weights are $\mathcal F_s$-measurable and sum to $1$. Taking $\E_s$ replaces $M(t,T_k)$ by $a_k=M(s,T_k)$, giving
$\bigl[(a_3-a_2)a_1+(a_1-a_3)a_2\bigr]/(a_1-a_2)=a_3$ \ts{13}{26}{1:11:10}.
\end{proof}
Equivalently, $M(t,T_3)$ is read off the straight line through the points $(a_1,M(t,T_1))$ and $(a_2,M(t,T_2))$ at the
abscissa $a_3$: linear interpolation with the ``right'' independent variable.

\paragraph{The LIBOR market model} \ts{13}{26}{1:12:13}. Martingale modelling proceeds in three steps: find the
measure under which a quantity is a martingale, represent the martingale as a stochastic integral (log-normal, CEV or
stochastic volatility), then change to a common measure \ts{13}{26}{1:13:17}. For forward LIBORs,
$\dd L_i/L_i=\lambda_i(t)\dd W_t^{T_{i+1}}$ under $\Prob_{T_{i+1}}$. Under the terminal measure $\Prob_{T_N}$ (slide~20),
\[
  \frac{\dd L_i(t)}{L_i(t)}=-\sum_{j=i+1}^{N-1}\frac{\lambda_i\lambda_j\delta_jL_j(t)}{1+\delta_jL_j(t)}\,\rho_{ij}\dd t
  +\lambda_i\dd W_t^{T_N,L_i},\qquad \rho_{ij}=\Corr\bigl(\dd\ln L_i,\dd\ln L_j\bigr),
\]
the full-dimensional model with one Brownian motion per LIBOR and the correlations as inputs \ts{13}{26}{1:15:33}.
The drift follows from \cref{ch08:prop-numeraire} and Girsanov's theorem (\cref{ch:stochcalc}); interest-rate models
of this kind are the subject of \cref{ch:hjm}.

% ==================================================================
\section{Closing remarks of the 2013 course}\label{ch08:sec-closing}

\begin{coursenote}[Tang's conclusion and the end of 18.S096 (2013)]
Tang, who presented from New Jersey because snow had grounded his flight \ts{13}{26}{0:00}, concluded that banks need
enterprise-level models for the non-linear portfolio effects (CVA, funding, liquidity, capital), built on top of
trade-level models with martingale testing, resampling and interpolation. These risks have received much more
attention since the crisis \ts{13}{26}{1:16:34}.

The instructors then closed the course \ts{13}{26}{1:17:37}. A list of past final-paper topics was on the course site:
mostly Black--Scholes, its extensions and manipulations of the Black--Scholes equation, one statistical analysis of
commodity data (recommended as a direction), and numerical and Monte Carlo projects. Students were asked to reflect
on what they had learned and on what they would like to explore next. P.~Kempthorne described the course as
deliberately challenging, meant to give both the mathematical foundations and a real exposure to their applications
through practitioners; he said the final papers he had read so far were good and that the faculty wanted to remain
a resource after the course \ts{13}{26}{1:19:45}. The class would be repeated the following fall in a slightly
different form, students were invited to suggest topics (one value of the course being access to speakers from the
frontier of the industry), and to fill in the course review \ts{13}{26}{1:20:45}.
\end{coursenote}

\begin{coursenote}[2013 versus 2024]
The two lectures reflect a decade of regulation. In 2013 many OTC trades were uncollateralised, and the focus was
pricing and hedging CVA, DVA and funding at portfolio level. By 2024 standard swaps are centrally cleared and
bilateral trades between banks carry daily VM and regulatory IM, so counterparty risk is managed mainly through
collateral: the questions are how much margin (a VaR/ES) and how to reduce it (optimisation), and the cost of
counterparty risk has partly turned into the cost of funding margin. Both speakers stress the same point: the
relevant quantities are non-linear functions of the \emph{whole} portfolio, or even the whole network.
\end{coursenote}

\begin{coursenote}[Reading suggested on the slides]
2024: S.~Boyd and L.~Vandenberghe, \emph{Convex Optimization} (the slide misspells the first name as ``Steven'');
J.~Gregory, \emph{The xVA Challenge}; J.~Hull, \emph{Risk Management and Financial Institutions}; Artzner et al.\
(1999) on coherent risk measures; Rockafellar and Uryasev on optimising CVaR; the BCBS--IOSCO margin framework and
the ISDA SIMM~2.6 methodology; Embrechts and Wang, ``Seven proofs for the subadditivity of expected shortfall'';
``Subadditivity re-examined: the case for Value-at-Risk''. 2013: Harrison--Kreps (1979), Harrison--Pliska (1981),
Jamshidian (1997) on LIBOR and swap market models, Sch\"onbucher (2003) on survival measures, Tang and Li (2007),
Tang and Williams (2010) on funding benefit and cost, Zhu and Pykhtin (2007), ``A guide to modeling counterparty
credit risk''.
\end{coursenote}

% ==================================================================
\begin{keyformulas}
\begin{align*}
&\text{Exposure:} && E(t)=(V(t)-C(t))^+,\\
& && \mathrm{EE}(t)=\E[E(t)],\quad
  \textstyle(\sum_iV_i)^+\le\sum_iV_i^+\\
&\text{CVA:} && \mathrm{CVA}=-\E^{\Q}\bigl[\ind_{\{\tau\le T\}}(1-R)
  (V_{\tau-}-C_{\tau-})^+/\beta_\tau\bigr]
  \approx-\mathrm{EPE}\cdot s\cdot D\\
&\text{DVA:} && \mathrm{DVA}=\E^{\Q}\bigl[\ind_{\{\bar\tau\le T\}}(1-\bar R)
  (C_{\bar\tau-}-V_{\bar\tau-})^+/\beta_{\bar\tau}\bigr]\\
&\text{VaR, ES:} && \mathrm{VaR}_\beta=\inf\{\ell:\Prob(L\le\ell)\ge\beta\},\\
& && \mathrm{ES}_\beta=\tfrac1{1-\beta}\textstyle\int_\beta^1\mathrm{VaR}_u\dd u\\
&\text{Normal:} && \mathrm{VaR}_\beta=\mu+z_\beta\sigma,\\
& && \mathrm{ES}_\beta=\mu+\sigma\varphi(z_\beta)/(1-\beta),\\
& && \sigma_P^2=\bm\delta\T C\bm\delta\\
&\text{Time scaling:} && \text{factor }\sqrt{T+2\textstyle\sum_{k=1}^{T-1}(T-k)\rho^k}\\
& && (\sqrt T\text{ if }\rho=0)\\
&\text{Back-test:} && \Prob(N\ge m)=\textstyle\sum_{k\ge m}\binom nkp^k(1-p)^{n-k}\\
&\text{Rockafellar--Uryasev:} && \mathrm{ES}_\beta=\min_\alpha\bigl\{\alpha+\tfrac{1}{1-\beta}\E[(L-\alpha)^+]\bigr\}\\
&\text{Delta--gamma IM:} && \E[\Delta] P=\tfrac12\tr(\gamma C),\\
& && \Var\Delta P=\bm\delta\T C\bm\delta+\tfrac12\tr((\gamma C)^2)\\
&\text{Margin optim.:} && \min\textstyle\sum_p\bigl(\sum_q\mathrm{SIMM}_{p,q}
  +\sum_c\mathrm{IM}_{p,c}\bigr),\\
& && x_{h,p,q}=-x_{h,q,p},\ \sum_qx_{h,p,q}=0\\
&\text{Change of numeraire:} && \E_t^{N_1}[Y_T]=\tfrac{N_2(t)}{N_1(t)}
  \E_t^{N_2}\bigl[Y_TN_1(T)/N_2(T)\bigr]
\end{align*}
\end{keyformulas}

% !TEX root = ../main.tex
\chapter{Time Series Analysis}\label{ch:timeseries}

\begin{sources}
\small
\begin{tabularx}{\linewidth}{@{}l l L@{}}
\toprule
\textbf{Lecture} & \textbf{Speaker} & \textbf{Content}\\
\midrule
2024 L12 & P.~Kempthorne & Stationary processes; from prices to log returns (S\&P~500, Amazon, crude oil, 10-year yield); sample autocorrelations and the Box--Pierce test; Wold representation theorem; lag operator, impulse response; ARMA, AR($p$) and its stationarity condition, AR(1), Yule--Walker equations, MA($q$); differencing. The lecture ends at slide~22 of 29.\\
2013 L8 (from 26:19) & P.~Kempthorne & The same univariate theory (the first 26 minutes finish regression, \cref{ch:regression}); adds linear-trend reversion versus stochastic trend, ARIMA, maximum likelihood (LIML/FIML) and the information criteria AIC, BIC, HQ.\\
2013 L11 (from 28:35) & P.~Kempthorne & Multivariate time series: cross-covariance matrices, multivariate Wold, VAR($p$) and its VAR(1) form, equation-by-equation least squares and why it is optimal, maximum likelihood, model selection; VAR of US macroeconomic series (Case Study~5). The first 28 minutes are GARCH, \cref{ch:volatility}.\\
2013 L12 & P.~Kempthorne & The fitted macro VAR; cointegration, vector error-correction models, energy futures and crack spreads; linear state-space models (time-varying CAPM, time-varying regression, AR, MA, ARMA); the Kalman filter.\\
\bottomrule
\end{tabularx}

\smallskip
\textbf{Files.} Slides \file{mit18_642_f24_lec11-12.pdf} (29 slides, 73 pages with overlays); 2013 slides
\file{MIT18_S096F13_lecnote8.pdf}, \file{lecnote11.pdf}, \file{lecnote12.pdf}; 2013 Case Study~3
\file{CaseStudy3.pdf} (time series of the US 10-year Treasury yield, in R) and Case Study~5
\file{CaseStudy5.pdf} (VAR models of macroeconomic series, in R).
\textbf{Problem sets.} 2024 PS4 Problems~5--6 (Problems~1--4 are on Brownian motion, \cref{ch:sp2});
2013 PS4 (all six problems); 2013 PS6 Problems~1--2 (Problems~3--5 are portfolio theory,
\cref{ch:portfolio}). The exercises are at the end of the chapter.
\textbf{Not published or not covered.} Slide~5 of the 2024 lecture points to
\file{TimeSeries4plots.pdf} and \file{TimeSeries4acfplots.pdf}, which are not in the download, so the
exploratory plots of \cref{ch09:sec-logret} are described from the spoken commentary. The
crack-spread cointegration note of 2013 L12 exists only in the video. Slides 23--29 of 2024
(trend reversion, ARIMA, maximum likelihood, model selection, Dickey--Fuller test) were not reached in
class; 2013 L8 covered them up to model selection, and the Dickey--Fuller slide was not discussed in
either year.
\end{sources}

\section*{Motivation}
Prices, yields, spreads and macroeconomic indicators are observed \emph{in time}, and the observation
today carries information about tomorrow. Chapter~\ref{ch:sp1} modelled the simplest case, a random walk
with independent steps. Real data have more structure: yield changes are weakly autocorrelated, a credit
spread is pulled back to a long-run level, an interest rate reacts to unemployment and inflation with a
delay, two energy prices drift apart but never far. Time-series analysis supplies the vocabulary:
\begin{itemize}
  \item \emph{stationarity} and the autocorrelation function, to say what ``no structure left'' means;
  \item the \emph{Wold theorem} and the \emph{ARMA} family, to say what structure is possible in
        a stationary series, and \emph{differencing}/unit-root tests to deal with the series that are not
        (\cref{ch09:sec-stat,ch09:sec-wold,ch09:sec-arma,ch09:sec-nonstat});
  \item \emph{VAR} models and \emph{cointegration}, for several series at once
        (\cref{ch09:sec-var,ch09:sec-coint});
  \item \emph{state-space} models and the \emph{Kalman filter}, for hidden quantities such as a time-varying
        beta or a ``true'' price observed with noise (\cref{ch09:sec-ss}).
\end{itemize}
The mathematics is linear algebra and least squares (\cref{ch:linalg,ch:regression}): every model in
this chapter is fitted by regressing the series on its own past, and the stationarity conditions are
statements about the eigenvalues of a matrix. The case studies (10-year Treasury yield and a
three-variable macro VAR) show what the workflow looks like on real data.

% ==================================================================
\section{Stationarity and autocorrelation}\label{ch09:sec-stat}

A \emph{time series} is a stochastic process $\{X_t,\ t\in T\}$ of random variables indexed by the times
of observation, discrete ($T=\{1,2,\dots\}$, the case used almost throughout) or continuous
($T=[0,\infty)$) \ts{24}{12}{0:00}. Its probabilistic behaviour is fixed by the finite-dimensional
distributions $p(x_{t_1},\dots,x_{t_m})$ for every finite set of times \ts{24}{12}{1:09}\ts{13}{8}{27:23}.

\begin{definition}[Strict stationarity]\label{ch09:def-strict}
$\{X_t\}$ is \emph{strictly stationary} if $p(x_{t_1+\tau},\dots,x_{t_m+\tau})=p(x_{t_1},\dots,x_{t_m})$ for all
$\tau$, $m$ and $(t_1,\dots,t_m)$: the joint law is invariant under time translation \ts{24}{12}{2:10}.
\end{definition}

\begin{definition}[Covariance stationarity and ACF]\label{ch09:def-cov}
$\{X_t\}$ is \emph{covariance (weakly, second-order) stationary} if, for all $t$,
\[
  \E X_t=\mu,\qquad \Var X_t=\sigma_X^2<\infty,\qquad \Cov(X_t,X_{t+\tau})=\gamma(\tau)
\]
depend at most on the lag $\tau$ \ts{24}{12}{5:24}\ts{13}{8}{29:32}. The \emph{autocorrelation function}
(ACF) is $\rho(\tau)=\gamma(\tau)/\gamma(0)$ \ts{24}{12}{7:27}.
\end{definition}

Two facts are used constantly. First, $\gamma(-\tau)=\gamma(\tau)$ and $|\rho(\tau)|\le1$. Second, for
every $n$ the Toeplitz matrix $\bigl[\gamma(i-j)\bigr]_{i,j=1}^n$ is the covariance matrix of $(X_1,\dots,X_n)$, hence
positive semi-definite; not every even function $\gamma$ is an autocovariance function.
Strict stationarity together with a finite variance implies covariance stationarity; the converse fails
in general but holds for Gaussian processes, whose law is determined by the first two moments.

\begin{intuition}[Time-translation invariance]
Strict stationarity is the statistician's \emph{equilibrium}: the process looks the same whichever window
one watches, like the noise of a system that has relaxed to a steady state. It is what makes estimation
possible, because a single path can then be averaged over time in place of an ensemble average (the
ergodic hypothesis). Covariance stationarity asks only that the two-point correlation function depend on
time differences, $C(t,t')=\gamma(t-t')$; the Fourier transform of $\gamma$ is the power spectrum, which the
course does not use but which explains why cycles show up as oscillations in $\rho(\tau)$
(\cref{ch09:sec-ar2}). Parameters of a stationary model are constants that can be estimated
\ts{24}{12}{3:12}\ts{13}{8}{28:28}.
\end{intuition}

\subsection{Making financial series stationary}\label{ch09:sec-logret}
Prices are not stationary: they trend, and the size of their fluctuations grows with their level. The
standard remedy is to model the \emph{log return}. Let $Y_t$ be a positive price and
\begin{equation}\label{ch09:eq-logret}
  X_t=\log\frac{Y_t}{Y_{t-1}}=\log(1+R_t),\qquad R_t=\frac{Y_t-Y_{t-1}}{Y_{t-1}},\qquad
  Y_t=Y_0\exp\bigl(X_1+\dots+X_t\bigr),
\end{equation}
so $Y_t$ is an exponential of a random walk with steps $X_t$, and the modelling assumption is that
$\{X_t\}$ is covariance stationary \ts{24}{12}{9:33}\ts{24}{12}{10:33}. For small returns
$X_t\approx R_t-\tfrac12R_t^2$. Taking logs turns exponential growth into linear growth and stabilises
the variance, and differencing removes the linear trend; $X_t$ is then a candidate stationary series.

The lecture illustrated this on four series (the plots themselves are not in the download)
\ts{24}{12}{8:29}:
\begin{itemize}
  \item \emph{S\&P~500}, roughly 2011--2024. The index trends up with sharp COVID drops, so it is not
        stationary in any window; the monthly log returns fluctuate around a level close to zero with
        constant spread. Weekly and daily returns keep a stable mean, but the volatility is clearly not
        constant (extreme in 2020) \ts{24}{12}{11:42}\ts{24}{12}{13:48}. A fitted normal density
        underestimates both the mass near the mean and the tails: the returns are
        \emph{leptokurtic} (``slender''), as in \cref{ch03:sec-heavy}. Volatility clustering is the subject of
        \cref{ch:volatility}.
  \item \emph{Amazon}. The price shows strong exponential growth; on the log scale the variability is more
        homogeneous, and the monthly and weekly log returns look stationary with a better normal fit than the
        index. Daily returns show a few spikes that no Gaussian model would produce
        \ts{24}{12}{15:56}\ts{24}{12}{16:56}.
  \item \emph{Crude oil futures}. In April 2020 the front-month contract went \emph{negative}; a holder of
        100 contracts woke up owing the broker 100 times a negative price, and some brokers could not
        even report positions because their systems did not allow negative prices. The logarithm is
        undefined there, so the lecture's plots exclude those days \ts{24}{12}{18:00}\ts{24}{12}{19:03}.
        Log returns require positive prices (spreads and rates can be modelled directly).
  \item \emph{10-year US Treasury yield}, from above $3\%$ to below $1\%$ after COVID and rising again.
        Strictly, ``log return'' is the wrong name for a percentage change of a yield, but the same
        transformation gives unimodal, roughly symmetric increments, and the normal model fits better than
        for equity returns \ts{24}{12}{20:05}\ts{24}{12}{21:12}\ts{24}{12}{22:16}.
\end{itemize}
Modelling $Y_t$ as an exponential random walk means that $S_t=X_1+\dots+X_t$ is a sum of steps. In the
simplest model the steps are i.i.d.\ (\cref{ch04:sec-rw}); the models of this chapter let them depend
on their own past \ts{24}{12}{23:18}.

\subsection{Sample autocorrelation and tests of whiteness}\label{ch09:sec-sacf}
Given observations $x_1,\dots,x_T$ the \emph{sample autocorrelation at lag $k$} is
\begin{equation}\label{ch09:eq-sacf}
  \hat r_k=\frac{\sum_{t=1}^{T-k}(x_t-\bar x)(x_{t+k}-\bar x)}{\sum_{t=1}^{T}(x_t-\bar x)^2}
\end{equation}
\ts{24}{12}{24:25}, an estimate of $\rho(k)$. If the series is white noise, then $\hat r_k$ is
approximately $N(0,1/T)$ (Bartlett); the lecture writes the variance as $1/(T-k)$, which is the same for
$T\gg k$ \ts{24}{12}{25:35}. Hence the test of $H_0:\rho(k)=0$ rejects at the $5\%$ level when
\begin{equation}\label{ch09:eq-band}
  |\hat r_k|>1.96\sqrt{\tfrac1{T-k}},
\end{equation}
which is the pair of dashed lines that R's \texttt{acf} draws around zero \ts{24}{12}{26:42}. For
a daily series with $T=500$ the bands are $\pm0.088$; a single autocorrelation of $0.09$ is
already significant. The lecture saw a large negative monthly autocorrelation at lag~1 for the
S\&P~500 and, on weekly data, significant negative correlations around lag~5, which
a student read as \emph{mean reversion} (price changes that are too high tend to be reversed)
\ts{24}{12}{27:42}. Amazon showed no significant autocorrelation at any frequency, crude oil a negative
lag-2 and positive lag-3 weekly pattern, and Treasury yields strong low-lag dependence
\ts{24}{12}{29:50}\ts{24}{12}{30:58}.

\paragraph{Testing many lags at once.} If $\rho(1)=\dots=\rho(K)=0$ the $\sqrt T\hat r_j$ are
asymptotically independent $N(0,1)$, so
\begin{equation}\label{ch09:eq-bp}
  \mathrm{BP}=T\sum_{j=1}^{K}\hat r_j^{\,2}\ \approx\ \chi^2_K ,
\end{equation}
the \emph{Box--Pierce} statistic; reject the joint null when $\mathrm{BP}$ exceeds the $\chi^2_K$
quantile ($18.31$ for $K=10$ at $5\%$) \ts{24}{12}{36:10}\ts{24}{12}{37:18}. Two refinements, not
discussed in class: the Ljung--Box variant $T(T+2)\sum_j\hat r_j^2/(T-j)$ has a better small-sample
$\chi^2$ approximation; and when the test is applied to the residuals of a fitted ARMA$(p,q)$ model the
degrees of freedom drop to $K-p-q$.

\begin{finance}[Residual whiteness as a model check]
A time-series model writes $Y_t=\hat Y_t+e_t$. It is \emph{adequate} when the residuals have mean zero,
constant variance and no autocorrelation, i.e.\ they are white noise: nothing predictable is left
\ts{24}{12}{31:59}\ts{24}{12}{33:03}. If in addition they are Gaussian, the model delivers not only point
forecasts but valid \emph{prediction intervals} \ts{24}{12}{34:03}. The ACF of the residuals, and the
Box--Pierce test on it, are therefore the standard diagnostic: a model whose residuals still show
significant autocorrelation ``is not done yet'' \ts{24}{12}{38:24}. The flip side, developed in
\cref{ch:volatility}, is that uncorrelated does not mean independent: returns can be white noise while
their squares are strongly autocorrelated.
\end{finance}

% ==================================================================
\section{The Wold representation theorem}\label{ch09:sec-wold}

\begin{theorem}[Wold representation]\label{ch09:thm-wold}
Every zero-mean covariance-stationary series $\{X_t\}$ can be written
\begin{equation}\label{ch09:eq-wold}
  X_t=V_t+S_t,\qquad S_t=\sum_{i=0}^{\infty}\psi_i\,\eta_{t-i},
\end{equation}
where $\psi_0=1$, $\sum_i\psi_i^2<\infty$; $\{\eta_t\}$ is \emph{white noise}, i.e.\ $\E\eta_t=0$,
$\E\eta_t^2=\sigma^2$, $\E\eta_t\eta_s=0$ for $t\ne s$; $\{V_t\}$ is \emph{linearly deterministic} (a linear
combination of its own past with constant coefficients predicts it exactly); and
$\E(\eta_tV_s)=0$ for all $t,s$ \ts{24}{12}{38:24}\ts{13}{8}{31:40}.
\end{theorem}

The theorem says that once the deterministic part is removed, \emph{every} stationary series is a moving
average of white noise, possibly of infinite order. The white noise $\eta_t=X_t-P(X_t\mid X_{t-1},X_{t-2},\dots)$
is the \emph{innovation}: the part of $X_t$ that cannot be predicted linearly from the past, the ``news'' arriving
at $t$ \ts{13}{11}{43:32}\ts{13}{11}{44:38}. An example of a linearly deterministic process is a
sinusoid, $V_t=A_1\cos\omega t+A_2\sin\omega t$ with random $A_1,A_2$: it obeys the exact recursion
$V_t=2\cos\omega\,V_{t-1}-V_{t-2}$ (from $\cos\omega t=2\cos\omega\cos\omega(t-1)-\cos\omega(t-2)$), so two
observations fix it for ever; sums of sinusoids of different frequencies (a Fourier series) are also linearly
predictable \ts{24}{12}{40:30}\ts{24}{12}{41:32}\ts{24}{12}{42:32}.

\begin{proof}[Sketch of proof]
Work in the Hilbert space $L^2$ of square-integrable random variables with inner product $\langle U,W\rangle=\E UW$.
Let $H_t$ be the closed linear span of $\{X_s:s\le t\}$ and $P_{H}$ the orthogonal projection on $H$. Define
$\eta_t=X_t-P_{H_{t-1}}X_t$. Stationarity makes the time shift a unitary map, so $\sigma^2=\E\eta_t^2$ does not depend on
$t$ (we assume $\sigma^2>0$; otherwise $X$ is itself deterministic). Since $\eta_s\in H_s\subset H_{t-1}$ for $s<t$
and $\eta_t\perp H_{t-1}$, the $\eta_t$ are uncorrelated. Put $\psi_j=\E[X_t\eta_{t-j}]/\sigma^2$, independent of $t$; then
$\psi_0=\E[X_t\eta_t]/\sigma^2=\E[\eta_t^2]/\sigma^2=1$, because $X_t-\eta_t\in H_{t-1}$. The projection of $X_t$ on
$\mathrm{span}\{\eta_t,\dots,\eta_{t-J+1}\}$ is $\sum_{j<J}\psi_j\eta_{t-j}$, and Bessel's inequality gives
$\sigma^2\sum_j\psi_j^2\le\Var X_t<\infty$. Hence $S_t=\sum_{j\ge0}\psi_j\eta_{t-j}$ converges in $L^2$, and
$V_t=X_t-S_t$ is orthogonal to every $\eta_s$ with $s\le t$. One checks that $V_t\in\bigcap_sH_s$, the
\emph{remote past} of the process, so $V_t$ is determined exactly, by a linear combination, from
arbitrarily remote values: it is linearly deterministic.
\end{proof}

\begin{intuition}[What the theorem does and does not say]
The proof shows why the result holds: $\eta_t$ is the component of $X_t$ orthogonal to the past, and
expanding $X_t$ in the orthogonal family $\eta_t,\eta_{t-1},\dots$ is a Gram--Schmidt procedure. The
innovations are only \emph{uncorrelated}, not independent: the theorem describes the best \emph{linear}
prediction, and a process such as a GARCH model (\cref{ch:volatility}) is covered by it although its
innovations are dependent through their variances. Covariance stationarity of the sum
$S_t$ needs only $\sum\psi_i^2<\infty$; that is why the condition appears \ts{24}{12}{50:07}.
\end{intuition}

\subsection{Wold's theorem as a modelling recipe}\label{ch09:sec-wold-recipe}
The theorem suggests how to specify a stationary model for data $x_0,\dots,x_T$
\ts{24}{12}{44:41}\ts{13}{8}{34:48}. Fix a number $p$ of lags to represent the linearly deterministic part, and take an
estimation window of size $n$ (after de-trending so that the series has a constant mean level). Set
\[
  \bm y=\begin{pmatrix}y_1\\ \vdots\\ y_n\end{pmatrix},\qquad
  Z=\begin{pmatrix}1&y_0&y_{-1}&\cdots&y_{-(p-1)}\\ 1&y_1&y_0&\cdots&y_{-(p-2)}\\
  \vdots&\vdots&&&\vdots\\ 1&y_{n-1}&y_{n-2}&\cdots&y_{n-p}\end{pmatrix},
\]
so the first $p$ rows reach back before the start of the response. The linear projection on $p$ lags is the least-squares fit
\[
  \hat{\bm y}^{(p)}=Z(Z\T Z)^{-1}Z\T\bm y=\hat P(Y_t\mid Y_{t-1},\dots,Y_{t-p}),\qquad
  \hat{\bm\varepsilon}^{(p)}=\bm y-\hat{\bm y}^{(p)},
\]
the hat matrix of \cref{ch05:prop-hat}. The residual series is then modelled as a moving average
$\varepsilon_t^{(p)}=\sum_i\psi_i\eta_{t-i}$, giving estimates of the $\psi_i$ and of the innovations $\eta_t$
\ts{24}{12}{45:41}\ts{24}{12}{46:45}. Two diagnostics decide whether to stop:
\begin{itemize}
  \item the residual should be orthogonal to lags of $Y$ beyond $p$; if it is not, increase $p$;
  \item the estimated innovations should look like white noise (\cref{ch09:sec-sacf}); if not, change the
        specification of the moving average or add deterministic regressors \ts{24}{12}{47:50}.
\end{itemize}
As $p\to\infty$ the projection converges to $P(Y_t\mid Y_{t-1},Y_{t-2},\dots)$, but a sample of size $n$
supports only $p/n\to0$, so useful models have a modest $p$ or a moving average with a
parsimonious parametrisation \ts{24}{12}{48:55}\ts{13}{8}{38:55}. This is exactly why ARMA models
exist (\cref{ch09:sec-arma}).

\subsection{Lag operator, impulse response, invertibility}\label{ch09:sec-lag}
The \emph{lag} (backshift) operator $L$ shifts a series back one period: $LX_t=X_{t-1}$, $L^0=\mathrm{id}$,
$L^nX_t=X_{t-n}$; its inverse $L^{-n}X_t=X_{t+n}$ shifts forward \ts{24}{12}{52:09}\ts{13}{8}{40:01}. Polynomials in $L$ multiply
and factorise like ordinary polynomials in a complex variable $z$, and the Wold representation reads
\begin{equation}\label{ch09:eq-woldL}
  X_t=\psi(L)\eta_t+V_t,\qquad \psi(L)=\sum_{i\ge0}\psi_iL^i .
\end{equation}
The \emph{impulse response function} is $\mathrm{IR}(j)=\partial X_t/\partial\eta_{t-j}=\psi_j$, the effect of a
shock $j$ periods ago on the process today \ts{24}{12}{54:20}; graphed against $j$ it displays delayed
responses. If the Federal Reserve chairman changes interest rates and the variable of interest is GNP,
$\psi_j$ is the effect $j$ quarters later \ts{13}{8}{43:11}\ts{13}{8}{44:11}. The \emph{long-run cumulative response} is
$\sum_j\psi_j=\psi(1)$, the level to which the process eventually moves after one unit innovation
\ts{24}{12}{56:28}\ts{13}{8}{45:11}.

If $\psi(L)$ is \emph{invertible}, i.e.\ there is $\psi^{-1}(L)=\sum_i\psi^*_iL^i$ with
$\psi^{-1}(L)\psi(L)=L^0$, then for $V_t=0$
\begin{equation}\label{ch09:eq-arinf}
  \psi^{-1}(L)X_t=\eta_t\quad\Longleftrightarrow\quad X_t=\sum_{i\ge1}\pi_iX_{t-i}+\eta_t,\qquad \pi_i=-\psi^*_i,
\end{equation}
an infinite-order autoregression \ts{24}{12}{57:31}\ts{24}{12}{58:35}\ts{13}{8}{46:17}. The coefficients of the
least-squares projection on $p$ lags in \cref{ch09:sec-wold-recipe} converge to these $\pi_i$
\ts{13}{8}{47:20}. The basic tool for inverting is the geometric series: for a complex number $\lambda$ with
$|\lambda|>1$,
\begin{equation}\label{ch09:eq-geom}
  \bigl(1-\lambda^{-1}L\bigr)^{-1}=\sum_{i\ge0}\lambda^{-i}L^i,
\end{equation}
which converges (on bounded stationary series) precisely because the coefficients $\lambda^{-i}$ decay
\ts{24}{12}{1:06:02}.

% ==================================================================
\section{ARMA models}\label{ch09:sec-arma}

\begin{definition}[ARMA$(p,q)$]\label{ch09:def-arma}
$\{X_t\}$ follows an \emph{autoregressive moving-average} model of orders $(p,q)$ if
\begin{equation}\label{ch09:eq-arma}
  X_t-\mu=\phi_1(X_{t-1}-\mu)+\dots+\phi_p(X_{t-p}-\mu)+\eta_t+\theta_1\eta_{t-1}+\dots+\theta_q\eta_{t-q},
\end{equation}
where $\{\eta_t\}$ is white noise $\mathrm{WN}(0,\sigma^2)$. With the polynomials
$\phi(L)=1-\phi_1L-\dots-\phi_pL^p$ and $\theta(L)=1+\theta_1L+\dots+\theta_qL^q$ this reads
\[
  \phi(L)(X_t-\mu)=\theta(L)\eta_t,\qquad\text{and}\qquad X_t=\mu+\psi(L)\eta_t,\quad \psi(L)=\frac{\theta(L)}{\phi(L)}
\]
is the Wold decomposition, when $\phi(L)$ is invertible \ts{24}{12}{59:35}\ts{13}{8}{46:17}\ts{13}{8}{48:22}.
\end{definition}

The sign convention of the moving-average part is ``$+\theta_j$'' (as in R); some texts use ``$-\theta_j$''.
ARMA models are \emph{parsimonious}: $p+q$ parameters generate an infinite Wold expansion
\ts{24}{12}{1:00:39}. We first look at the two pure cases.

\subsection{Autoregressive models AR($p$)}\label{ch09:sec-ar}
Set $q=0$: $\phi(L)(X_t-\mu)=\eta_t$, that is,
\begin{equation}\label{ch09:eq-arp}
  X_t=c+\sum_{j=1}^p\phi_jX_{t-j}+\eta_t,\qquad c=\mu\,\phi(1)=\mu\Bigl(1-\sum_j\phi_j\Bigr),
\end{equation}
a \emph{linear regression of the series on its own lags}, with white-noise errors
\ts{24}{12}{1:01:47}\ts{24}{12}{1:02:49}\ts{13}{8}{50:34}. (Taking expectations of \eqref{ch09:eq-arp} gives $\mu=c+\mu\sum\phi_j$,
whence $c=\mu\phi(1)$.)

\begin{theorem}[Stationarity of AR($p$)]\label{ch09:thm-ar-stat}
The equation $\phi(L)(X_t-\mu)=\eta_t$ has a covariance-stationary solution of the form
$X_t-\mu=\sum_{i\ge0}\psi_i\eta_{t-i}$ if and only if every root $\lambda_j$ of the
\emph{characteristic equation} $\phi(z)=1-\phi_1z-\dots-\phi_pz^p=0$ lies outside the closed unit
disc, $|\lambda_j|>1$, $j=1,\dots,p$ \ts{24}{12}{1:03:51}\ts{13}{8}{52:35}.
\end{theorem}
\begin{proof}
Factorise $\phi(z)=\prod_{j=1}^p(1-z/\lambda_j)$ (the constant term is~1). If all $|\lambda_j|>1$,
each factor is inverted by the geometric series \eqref{ch09:eq-geom}, whose coefficients $\lambda_j^{-i}$ decay
geometrically, and the product $\phi^{-1}(L)=\prod_j(1-L/\lambda_j)^{-1}$ of absolutely summable series is
absolutely summable. Thus $\psi(L)=\phi^{-1}(L)$ has $\sum|\psi_i|<\infty$ and $X_t=\mu+\psi(L)\eta_t$ is
covariance stationary \ts{24}{12}{1:05:00}\ts{24}{12}{1:07:09}. Conversely, a root with $|\lambda|<1$ makes
the coefficients $\lambda^{-i}$ grow, so the expansion diverges (the process is explosive), and a root on
the circle $|\lambda|=1$ gives a unit root (the random walk is the case $p=1$, $\lambda=1$), whose variance
grows with $t$; in neither case is there a stationary solution driven by the past.
\end{proof}

\begin{remark}[Roots and eigenvalues]\label{ch09:rem-companion}
The condition is often stated on the reciprocals $1/\lambda_j$, the roots of $z^p-\phi_1z^{p-1}-\dots-\phi_p=0$,
which are the eigenvalues of the \emph{companion matrix}
\[
  F=\begin{pmatrix}\phi_1&\phi_2&\cdots&\phi_{p-1}&\phi_p\\ 1&0&\cdots&0&0\\ 0&1&&0&0\\ \vdots&&\ddots&&\vdots\\ 0&0&\cdots&1&0\end{pmatrix};
\]
stationarity is $|1/\lambda_j|<1$ for all $j$. Software conventions differ: R's \texttt{vars} package reports
the moduli of the eigenvalues (stationary if all $<1$), whereas \texttt{polyroot} returns roots of $\phi(z)$
(stationary if all $>1$) \ts{13}{12}{1:08}. The proof of the correspondence is given for the vector case in
\cref{ch09:prop-companion}.
\end{remark}

\begin{figure}[ht]
\centering
\begin{tikzpicture}[scale=1.15,>=Stealth]
  \draw[->] (-2.4,0)--(2.4,0) node[right] {$\mathrm{Re}\,z$};
  \draw[->] (0,-2.1)--(0,2.1) node[above] {$\mathrm{Im}\,z$};
  \draw[thick,cscol,fill=cscol!8] (0,0) circle (1);
  \node[cscol] at (0.05,0.5) {\footnotesize\begin{tabular}{c}roots here:\\ explosive\end{tabular}};
  \foreach \x/\y in {1.7/0.6,1.7/-0.6,-1.5/0,0/1.6}{\node[intcol] at (\x,\y) {\large$\times$};}
  \node[intcol,align=left,font=\footnotesize] at (1.55,1.55) {stationary\\[-2pt](all roots\\[-2pt]outside)};
  \node[thmcol] at (0.7,-0.45) {\large$\times$};
  \node[thmcol,font=\footnotesize] at (-1.05,-1.55) {one root inside};
  \draw[thmcol,->] (-1.0,-1.35)--(0.62,-0.5);
  \node[font=\footnotesize] at (1.2,-0.2) {$1$};
\end{tikzpicture}
\caption{Roots $\lambda_j$ of $\phi(z)=0$ in the complex plane. Green crosses: a stationary AR model, all roots
outside the unit circle. Red cross: a single root inside makes the process explosive.}\label{ch09:fig-roots}
\end{figure}

\subsubsection{AR(1)}\label{ch09:sec-ar1}
For $X_t-\mu=\phi(X_{t-1}-\mu)+\eta_t$ the characteristic equation is $1-\phi z=0$ with root $\lambda=1/\phi$:
the process is covariance stationary iff $|\phi|<1$ \ts{24}{12}{1:08:12}. Iterating,
$X_t-\mu=\sum_{j\ge0}\phi^j\eta_{t-j}$ (the Wold decomposition, $\psi_j=\phi^j$), and therefore
\begin{equation}\label{ch09:eq-ar1}
  \E X_t=\mu,\qquad \gamma(0)=\sigma_X^2=\frac{\sigma^2}{1-\phi^2},\qquad
  \gamma(j)=\phi^{\,j}\,\sigma_X^2,\qquad \rho(j)=\phi^{\,j}
\end{equation}
\ts{24}{12}{1:08:12}\ts{13}{8}{56:46}\ts{13}{8}{58:54}. (Multiply $X_t-\mu=\phi(X_{t-1}-\mu)+\eta_t$ by $X_{t-j}-\mu$
and take expectations, or sum the geometric series.) The variance exceeds the innovation variance $\sigma^2$ for every
$\phi\neq0$, and depends on $\phi$ only through $\phi^2$: positive and negative autocorrelation inflate it equally
\ts{13}{8}{57:50}. The regimes:
\begin{itemize}
  \item $0<\phi<1$: \emph{exponential mean reversion} to $\mu$, all correlations positive;
  \item $-1<\phi<0$: \emph{oscillating} exponential mean reversion, correlations alternate in sign
        \ts{24}{12}{1:09:15};
  \item $\phi=1$: the Wold decomposition does not exist, $X_t$ is a random walk, not stationary;
  \item $|\phi|>1$: explosive.
\end{itemize}

\begin{figure}[ht]
\centering
\begin{tikzpicture}
\begin{axis}[width=0.47\linewidth,height=4.6cm,xlabel={lag $j$},ylabel={$\rho(j)$},ymin=-1.05,ymax=1.05,
  xmin=-0.5,xmax=12.5,ytick={-1,-0.5,0,0.5,1},title={AR(1), $\phi=0.8$},title style={font=\footnotesize},grid=major,grid style={black!10}]
  \addplot+[ycomb,thick,mark=*,mark size=1.5pt,color=defcol] coordinates {(0,1.0) (1,0.8) (2,0.64) (3,0.512) (4,0.4096) (5,0.3277) (6,0.2621) (7,0.2097) (8,0.1678) (9,0.1342) (10,0.1074) (11,0.0859) (12,0.0687)};
\end{axis}
\end{tikzpicture}\hfill
\begin{tikzpicture}
\begin{axis}[width=0.47\linewidth,height=4.6cm,xlabel={lag $j$},ymin=-1.05,ymax=1.05,
  xmin=-0.5,xmax=12.5,ytick={-1,-0.5,0,0.5,1},title={AR(1), $\phi=-0.8$},title style={font=\footnotesize},grid=major,grid style={black!10}]
  \addplot+[ycomb,thick,mark=*,mark size=1.5pt,color=thmcol] coordinates {(0,1.0) (1,-0.8) (2,0.64) (3,-0.512) (4,0.4096) (5,-0.3277) (6,0.2621) (7,-0.2097) (8,0.1678) (9,-0.1342) (10,0.1074) (11,-0.0859) (12,0.0687)};
\end{axis}
\end{tikzpicture}
\caption{Theoretical autocorrelation $\rho(j)=\phi^j$ of an AR(1): geometric decay for $\phi>0$, geometric
decay with alternating sign for $\phi<0$.}\label{ch09:fig-acf-ar1}
\end{figure}

\begin{figure}[ht]
\centering
\begin{tikzpicture}
\begin{axis}[width=0.95\linewidth,height=5.2cm,xlabel={$t$},xmin=0,xmax=149,legend pos=south west,legend style={font=\footnotesize},
  grid=major,grid style={black!10}]
  \addplot[thick,color=defcol] coordinates {(0,0.0) (1,-2.556) (2,-1.754) (3,-2.059) (4,-2.203) (5,-2.088) (6,-3.795) (7,-3.457) (8,-3.804) (9,0.09) (10,0.302) (11,-0.096) (12,-0.363) (13,-0.976) (14,-1.885) (15,-1.993) (16,-1.212) (17,-1.269) (18,-0.121) (19,-0.303) (20,-0.233) (21,1.348) (22,1.691) (23,0.932) (24,0.609) (25,1.058) (26,2.835) (27,2.14) (28,1.575) (29,2.341) (30,1.104) (31,0.646) (32,1.432) (33,1.798) (34,1.619) (35,2.047) (36,-1.089) (37,0.096) (38,-0.878) (39,-2.415) (40,-1.776) (41,-0.809) (42,-1.133) (43,-2.039) (44,-1.707) (45,-1.504) (46,0.127) (47,0.856) (48,0.921) (49,1.895) (50,1.405) (51,0.268) (52,0.812) (53,1.273) (54,0.867) (55,-0.046) (56,0.19) (57,-2.332) (58,-1.292) (59,-0.607) (60,-2.155) (61,-1.77) (62,-2.469) (63,-1.341) (64,-3.174) (65,-3.613) (66,-2.361) (67,-0.851) (68,-2.881) (69,-2.947) (70,-2.177) (71,-2.46) (72,-0.5) (73,-1.616) (74,-1.019) (75,-1.915) (76,-0.222) (77,-0.21) (78,-0.551) (79,-2.186) (80,-0.177) (81,0.603) (82,1.266) (83,2.214) (84,2.231) (85,1.257) (86,0.268) (87,-0.572) (88,0.884) (89,-0.709) (90,-1.199) (91,-1.341) (92,-0.915) (93,-0.202) (94,-1.421) (95,-2.938) (96,-2.502) (97,-0.913) (98,-0.019) (99,0.2) (100,-0.147) (101,0.168) (102,-0.101) (103,0.732) (104,-0.173) (105,-0.012) (106,-0.121) (107,0.44) (108,0.599) (109,3.059) (110,4.099) (111,4.981) (112,2.194) (113,1.525) (114,0.687) (115,1.117) (116,-1.33) (117,0.044) (118,1.105) (119,-0.363) (120,-1.287) (121,-1.895) (122,-1.568) (123,-0.692) (124,1.46) (125,1.044) (126,1.655) (127,1.562) (128,3.087) (129,3.366) (130,4.23) (131,2.518) (132,1.948) (133,0.842) (134,2.221) (135,2.545) (136,1.858) (137,1.127) (138,1.443) (139,0.525) (140,-0.485) (141,0.069) (142,2.522) (143,1.898) (144,1.057) (145,-0.273) (146,-1.567) (147,-0.807) (148,0.165) (149,0.15)};
  \addplot[thick,color=thmcol] coordinates {(0,0.0) (1,-2.556) (2,-2.138) (3,-2.705) (4,-3.158) (5,-3.374) (6,-5.394) (7,-5.626) (8,-6.491) (9,-3.168) (10,-2.942) (11,-3.295) (12,-3.576) (13,-4.244) (14,-5.299) (15,-5.69) (16,-5.208) (17,-5.446) (18,-4.489) (19,-4.688) (20,-4.664) (21,-3.118) (22,-2.573) (23,-3.079) (24,-3.261) (25,-2.721) (26,-0.786) (27,-1.055) (28,-1.299) (29,-0.297) (30,-1.183) (31,-1.475) (32,-0.592) (33,-0.012) (34,0.08) (35,0.75) (36,-2.078) (37,-1.057) (38,-2.017) (39,-3.685) (40,-3.409) (41,-2.708) (42,-3.153) (43,-4.23) (44,-4.203) (45,-4.256) (46,-2.851) (47,-2.103) (48,-1.909) (49,-0.798) (50,-1.003) (51,-1.929) (52,-1.345) (53,-0.763) (54,-0.977) (55,-1.76) (56,-1.531) (57,-4.025) (58,-3.335) (59,-2.843) (60,-4.482) (61,-4.421) (62,-5.385) (63,-4.628) (64,-6.662) (65,-7.577) (66,-6.867) (67,-5.711) (68,-7.869) (69,-8.367) (70,-8.039) (71,-8.648) (72,-7.057) (73,-8.248) (74,-7.894) (75,-8.942) (76,-7.536) (77,-7.558) (78,-7.93) (79,-9.648) (80,-7.967) (81,-7.214) (82,-6.46) (83,-5.322) (84,-4.973) (85,-5.612) (86,-6.413) (87,-7.213) (88,-5.843) (89,-7.303) (90,-7.899) (91,-8.221) (92,-7.996) (93,-7.421) (94,-8.67) (95,-10.4) (96,-10.404) (97,-9.191) (98,-8.434) (99,-8.218) (100,-8.535) (101,-8.242) (102,-8.485) (103,-7.668) (104,-8.462) (105,-8.328) (106,-8.439) (107,-7.896) (108,-7.671) (109,-5.121) (110,-3.622) (111,-2.126) (112,-4.165) (113,-4.505) (114,-5.114) (115,-4.581) (116,-6.86) (117,-5.686) (118,-4.619) (119,-5.921) (120,-6.899) (121,-7.701) (122,-7.657) (123,-7.016) (124,-4.968) (125,-5.166) (126,-4.398) (127,-4.243) (128,-2.483) (129,-1.741) (130,-0.372) (131,-1.45) (132,-1.642) (133,-2.456) (134,-0.951) (135,-0.293) (136,-0.599) (137,-1.051) (138,-0.566) (139,-1.268) (140,-2.199) (141,-1.717) (142,0.746) (143,0.5) (144,-0.056) (145,-1.227) (146,-2.562) (147,-2.037) (148,-1.186) (149,-1.177)};
  \legend{{AR(1), $\phi=0.85$},{random walk ($\phi=1$)}}
\end{axis}
\end{tikzpicture}
\caption{Two paths driven by the same simulated white noise (seed~3). The AR(1) with $\phi=0.85$ keeps
returning to $\mu=0$; the random walk wanders (its variance grows like $t\sigma^2$).}\label{ch09:fig-ar1-rw}
\end{figure}

\begin{finance}[AR(1) in practice: mean reversion and Vasicek]\label{ch09:fin-vasicek}
The AR(1) is the discrete-time version of the Ornstein--Uhlenbeck (Vasicek) process
$\mathrm dX=-\kappa(X-\mu)\,\mathrm dt+\sigma_0\,\mathrm dW$ (\cref{ch:sde}): sampled every $h$, the exact
transition is an AR(1) with $\phi=e^{-\kappa h}$ and noise variance $\sigma_0^2(1-e^{-2\kappa h})/(2\kappa)$, consistent with
$\gamma(0)=\sigma_0^2/(2\kappa)$ in \eqref{ch09:eq-ar1}. The half-life of a shock is
$h\ln2/(-\ln\phi)$. Candidates for such a model are variables that should not wander too far: interest
rates, interest-rate spreads, real exchange rates and valuation ratios (dividend-to-price, earnings-to-price)
\ts{24}{12}{1:10:15}\ts{24}{12}{1:11:16}\ts{13}{8}{58:54}. The caveat of the 2013 lecture: this is usually true over short
periods; over long ones the level $\mu$ to which the series reverts drifts, so the model must be re-estimated
\ts{13}{8}{59:59}.
\end{finance}

\subsubsection{Yule--Walker equations}\label{ch09:sec-yw}
Multiply the AR($p$) equation $X_t-\mu=\sum_i\phi_i(X_{t-i}-\mu)+\eta_t$ by $X_{t-j}-\mu$ and take expectations. Since
$\E[\eta_t(X_{t-j}-\mu)]=0$ for $j\ge1$ and $=\sigma^2$ for $j=0$ (the noise is uncorrelated with the past),
\begin{equation}\label{ch09:eq-yw}
  \gamma(j)=\sum_{i=1}^p\phi_i\,\gamma(j-i)+\sigma^2\,\delta_{j0},\qquad j=0,1,2,\dots
\end{equation}
\ts{24}{12}{1:11:16}\ts{24}{12}{1:12:16}\ts{13}{8}{1:02:01} (the slide misprints the second lag as $X_{t-1}$; the
lecturer corrected it aloud). Dividing by $\gamma(0)$, $\rho(j)=\sum_i\phi_i\rho(j-i)$ for $j\ge1$. The equations
$j=1,\dots,p$ form the \emph{Yule--Walker system}
\begin{equation}\label{ch09:eq-ywsys}
  \begin{pmatrix}\gamma(1)\\ \gamma(2)\\ \vdots\\ \gamma(p)\end{pmatrix}=
  \begin{pmatrix}\gamma(0)&\gamma(1)&\cdots&\gamma(p-1)\\ \gamma(1)&\gamma(0)&\cdots&\gamma(p-2)\\
  \vdots&&\ddots&\vdots\\ \gamma(p-1)&\gamma(p-2)&\cdots&\gamma(0)\end{pmatrix}
  \begin{pmatrix}\phi_1\\ \phi_2\\ \vdots\\ \phi_p\end{pmatrix},
\end{equation}
using $\gamma(-j)=\gamma(j)$. The matrix is Toeplitz (constant along diagonals), a question the 2013 lecturer put to the
class \ts{13}{8}{1:03:02}. The zeroth equation then gives the noise variance
\begin{equation}\label{ch09:eq-ywsigma}
  \sigma^2=\gamma(0)-\sum_{j=1}^p\phi_j\gamma(j).
\end{equation}
\textbf{Estimation.} Replace the $\gamma(j)$ by sample autocovariances $\hat\gamma(j)$ and solve
\eqref{ch09:eq-ywsys}: these are the \emph{Yule--Walker estimates} \ts{24}{12}{1:13:20}. They are an application of the
\emph{method of moments}: equate sample and theoretical moments, and solve for the parameters. This is the tool of choice when the likelihood is hard to write
down or to optimise; econometrics is rich in such applications \ts{24}{12}{1:14:24}\ts{13}{8}{1:05:09}. (Two remarks beyond
the course: the system is solved efficiently by the Levinson--Durbin recursion, and with the divisor $T$ in $\hat\gamma$ the
fitted AR model is always stationary.)

\begin{example}[Yule--Walker with two lags]\label{ch09:ex-yw}
Suppose the sample autocorrelations of a series are $\hat\rho(1)=0.5$ and $\hat\rho(2)=0.2$. The $p=2$ system
$\bigl(\begin{smallmatrix}1&0.5\\0.5&1\end{smallmatrix}\bigr)\bigl(\begin{smallmatrix}\phi_1\\\phi_2\end{smallmatrix}\bigr)
=\bigl(\begin{smallmatrix}0.5\\0.2\end{smallmatrix}\bigr)$ has determinant $0.75$ and solution
$\phi_1=(0.5-0.5\cdot0.2)/0.75=8/15\approx0.533$ and $\phi_2=(0.2-0.25)/0.75=-1/15\approx-0.067$. The fitted process is
stationary: $\phi_1+\phi_2=0.467<1$, $\phi_2-\phi_1=-0.6<1$, $|\phi_2|<1$ (\cref{ch09:sec-ar2}). On a simulated AR(2)
with $\phi=(1,-0.5)$ and $n=2000$ the Yule--Walker estimates were $(1.017,-0.515)$ with $\hat\sigma^2=0.960$
(true value~1).
\end{example}

\subsubsection{Partial autocorrelation}\label{ch09:sec-pacf}
The \emph{partial autocorrelation} at lag $k$, $\alpha(k)$, is the correlation between $X_t$ and $X_{t-k}$ \emph{after
removing the linear effect of the intermediate values} $X_{t-1},\dots,X_{t-k+1}$
\ts{24}{12}{28:49}\ts{13}{12}{4:11}. Equivalently, it is the last coefficient $\phi_{kk}$ of the best linear predictor of
$X_t$ from $k$ lags, the coefficient of $X_{t-k}$ in the regression of $X_t$ on $X_{t-1},\dots,X_{t-k}$
\ts{13}{12}{5:13}: the ``incremental correlation'' of adding one more lag. Together the ACF and PACF identify the
order of a model:
\begin{center}\small
\begin{tabular}{@{}lll@{}}
\toprule
 & ACF & PACF\\
\midrule
AR($p$) & decays geometrically or as a damped sinusoid & \textbf{cuts off} after lag $p$\\
MA($q$) & \textbf{cuts off} after lag $q$ & decays\\
ARMA($p,q$) & decays & decays\\
\bottomrule
\end{tabular}
\end{center}
The AR row follows because the regression of an AR($p$) on more than $p$ lags has a zero coefficient on the extra
ones; the MA row is proved in \cref{ch09:sec-ma}.

\begin{casestudy}[10-year Treasury yield: ACF and PACF (2013 Case Study~3)]\label{ch09:cs-3a}
The case study downloads daily constant-maturity yields \texttt{DGS10} from FRED (3\,356 non-missing days, January~2000 to
May~2013; missing values are bond-market holidays), builds weekly (700) and monthly (161) closing series with
\texttt{to.weekly}, \texttt{to.monthly}, and plots ACF and PACF at each frequency. The high first-order
autocorrelation of the levels suggests a unit root at every frequency. The Augmented Dickey--Fuller tests of
\cref{ch09:sec-df} confirm it (p-values 0.057, 0.105 and 0.272 for daily, weekly and monthly levels, none rejecting a unit root
at $5\%$), and reject it for the first differences (p-values $0.01$ at all frequencies, the smallest value R reports).
For the monthly differences $y_0$ (160 values) the table of sample ACF and PACF is
\begin{center}\small
\begin{tabular}{@{}l rrrr@{}}
\toprule
lag & 1 & 2 & 3 & 4\\
\midrule
ACF & $0.035$ & $-0.211$ & $0.079$ & $-0.010$\\
PACF & $0.035$ & $-0.213$ & $0.100$ & $-0.068$\\
\midrule
last coeff.\ of AR($k$) fit by \texttt{lm} & $0.036$ & $-0.217$ & $0.104$ & $-0.072$\\
\bottomrule
\end{tabular}
\end{center}
The last row confirms the definition of the PACF: it is (nearly) the coefficient of the last lag in a regression on $k$ lags.
The small differences arise because \texttt{lm} conditions on the first $k$ cases, whereas \texttt{acf} uses all available
pairs at each lag. Only the lag-2 coefficient is significant ($t=-2.76$ in the AR(2) fit).
\end{casestudy}

\subsubsection{AR(2): roots, correlations and cycles}\label{ch09:sec-ar2}
The AR(2) model $X_t=\phi_0+\phi_1X_{t-1}+\phi_2X_{t-2}+\eta_t$ (2013 PS4 Problem~4, exercise \ref{ch09:ex-ar2})
is the simplest process that can oscillate. Write $\phi(z)=1-\phi_1z-\phi_2z^2=(1-G_1z)(1-G_2z)$, where
$G_i=1/\lambda_i$ are the roots of $G^2-\phi_1G-\phi_2=0$. Then
\[
  G_{1,2}=\frac{\phi_1\pm\sqrt{\phi_1^2+4\phi_2}}2,
\]
and the roots are real and distinct if $\phi_1^2+4\phi_2>0$, real and coincident if $\phi_1^2+4\phi_2=0$, and a
complex-conjugate pair if $\phi_1^2+4\phi_2<0$. Both $|G_i|<1$ if and only if the parameters lie in the
\emph{stationarity triangle}
\[
  \phi_1+\phi_2<1,\qquad \phi_2-\phi_1<1,\qquad |\phi_2|<1 .
\]
The autocorrelations obey $\rho(k)=\phi_1\rho(k-1)+\phi_2\rho(k-2)$ with $\rho(0)=1$ and
$\rho(1)=\phi_1/(1-\phi_2)$ (the first Yule--Walker equation). This is a linear difference equation whose solutions
are combinations of powers of the roots, $\rho(k)=C_1G_1^k+C_2G_2^k$ for $G_1\ne G_2$, and the two initial conditions give
\begin{equation}\label{ch09:eq-ar2rho}
  \rho(k)=\frac{G_1(1-G_2^2)\,G_1^{\,k}-G_2(1-G_1^2)\,G_2^{\,k}}{(G_1-G_2)(1+G_1G_2)} .
\end{equation}
For real roots the correlations decay geometrically in magnitude. For complex roots write
$G_{1,2}=d\,e^{\pm\mathrm i\omega}$ with $d=\sqrt{-\phi_2}$ and $\omega=\arccos\bigl(\phi_1/(2d)\bigr)$; then
\begin{equation}\label{ch09:eq-ar2osc}
  \rho(k)=d^{\,k}\,\frac{\sin(\omega k+F)}{\sin F},\qquad \tan F=\frac{1-\phi_2}{1+\phi_2}\tan\omega,
\end{equation}
a damped cosine with damping factor $d$, frequency $f_0=\omega/2\pi$ (period $1/f_0$) and phase $F$
(taken in the same quadrant as $\omega$). The 2013 problem set states the same result with $|\phi_1|$ and a factor $[\mathrm{sgn}\,\phi_1]^k$. A series that
oscillates around its mean with a typical period of $2\pi/\omega$ periods is the signature of complex roots.

\begin{example}[An AR(2) with a period of eight steps]\label{ch09:ex-ar2osc}
Take $\phi_1=1$, $\phi_2=-0.5$. Then $\phi_1^2+4\phi_2=-1<0$, the roots of $\phi(z)$ are $1\pm\mathrm i$ (modulus
$\sqrt2>1$: stationary), $d=1/\sqrt2$, $\omega=\arccos(1/\sqrt2)=\pi/4$: period exactly $8$, and $\tan F=3$. The
autocorrelations are $\rho=(1,\ 0.667,\ 0.167,\ -0.167,\ -0.25,\ -0.167,\ -0.042,\ 0.042,\dots)$, and the variance is
$\gamma(0)=\sigma^2(1-\phi_2)/[(1+\phi_2)((1-\phi_2)^2-\phi_1^2)]=2.4\,\sigma^2$, equal to $\sigma^2\sum_j\psi_j^2$ for the
Wold weights $\psi=(1,1,0.5,0,-0.25,-0.25,\dots)$ (checked numerically). Note $\psi(1)=1/(1-\phi_1-\phi_2)=2$: a unit shock
eventually raises the level by $2$ (\cref{fig:ar2}).
\end{example}

\begin{figure}[ht]
\centering
\begin{tikzpicture}
\begin{axis}[width=0.95\linewidth,height=5cm,xlabel={lag $k$},ylabel={$\rho(k)$},xmin=-0.5,xmax=20.5,ymin=-0.5,ymax=1.05,grid=major,grid style={black!10}]
  \addplot[dashed,black!50] coordinates {(0.0,1.0) (0.25,0.917) (0.5,0.8409) (0.75,0.7711) (1.0,0.7071) (1.25,0.6484) (1.5,0.5946) (1.75,0.5453) (2.0,0.5) (2.25,0.4585) (2.5,0.4204) (2.75,0.3856) (3.0,0.3536) (3.25,0.3242) (3.5,0.2973) (3.75,0.2726) (4.0,0.25) (4.25,0.2293) (4.5,0.2102) (4.75,0.1928) (5.0,0.1768) (5.25,0.1621) (5.5,0.1487) (5.75,0.1363) (6.0,0.125) (6.25,0.1146) (6.5,0.1051) (6.75,0.0964) (7.0,0.0884) (7.25,0.0811) (7.5,0.0743) (7.75,0.0682) (8.0,0.0625) (8.25,0.0573) (8.5,0.0526) (8.75,0.0482) (9.0,0.0442) (9.25,0.0405) (9.5,0.0372) (9.75,0.0341) (10.0,0.0313) (10.25,0.0287) (10.5,0.0263) (10.75,0.0241) (11.0,0.0221) (11.25,0.0203) (11.5,0.0186) (11.75,0.017) (12.0,0.0156) (12.25,0.0143) (12.5,0.0131) (12.75,0.012) (13.0,0.011) (13.25,0.0101) (13.5,0.0093) (13.75,0.0085) (14.0,0.0078) (14.25,0.0072) (14.5,0.0066) (14.75,0.006) (15.0,0.0055) (15.25,0.0051) (15.5,0.0046) (15.75,0.0043) (16.0,0.0039) (16.25,0.0036) (16.5,0.0033) (16.75,0.003) (17.0,0.0028) (17.25,0.0025) (17.5,0.0023) (17.75,0.0021) (18.0,0.002) (18.25,0.0018) (18.5,0.0016) (18.75,0.0015) (19.0,0.0014) (19.25,0.0013) (19.5,0.0012) (19.75,0.0011) (20.0,0.001)};
  \addplot[dashed,black!50] coordinates {(0.0,-1.0) (0.25,-0.917) (0.5,-0.8409) (0.75,-0.7711) (1.0,-0.7071) (1.25,-0.6484) (1.5,-0.5946) (1.75,-0.5453) (2.0,-0.5) (2.25,-0.4585) (2.5,-0.4204) (2.75,-0.3856) (3.0,-0.3536) (3.25,-0.3242) (3.5,-0.2973) (3.75,-0.2726) (4.0,-0.25) (4.25,-0.2293) (4.5,-0.2102) (4.75,-0.1928) (5.0,-0.1768) (5.25,-0.1621) (5.5,-0.1487) (5.75,-0.1363) (6.0,-0.125) (6.25,-0.1146) (6.5,-0.1051) (6.75,-0.0964) (7.0,-0.0884) (7.25,-0.0811) (7.5,-0.0743) (7.75,-0.0682) (8.0,-0.0625) (8.25,-0.0573) (8.5,-0.0526) (8.75,-0.0482) (9.0,-0.0442) (9.25,-0.0405) (9.5,-0.0372) (9.75,-0.0341) (10.0,-0.0313) (10.25,-0.0287) (10.5,-0.0263) (10.75,-0.0241) (11.0,-0.0221) (11.25,-0.0203) (11.5,-0.0186) (11.75,-0.017) (12.0,-0.0156) (12.25,-0.0143) (12.5,-0.0131) (12.75,-0.012) (13.0,-0.011) (13.25,-0.0101) (13.5,-0.0093) (13.75,-0.0085) (14.0,-0.0078) (14.25,-0.0072) (14.5,-0.0066) (14.75,-0.006) (15.0,-0.0055) (15.25,-0.0051) (15.5,-0.0046) (15.75,-0.0043) (16.0,-0.0039) (16.25,-0.0036) (16.5,-0.0033) (16.75,-0.003) (17.0,-0.0028) (17.25,-0.0025) (17.5,-0.0023) (17.75,-0.0021) (18.0,-0.002) (18.25,-0.0018) (18.5,-0.0016) (18.75,-0.0015) (19.0,-0.0014) (19.25,-0.0013) (19.5,-0.0012) (19.75,-0.0011) (20.0,-0.001)};
  \addplot+[ycomb,thick,mark=*,mark size=1.5pt,color=defcol] coordinates {(0,1.0) (1,0.6667) (2,0.1667) (3,-0.1667) (4,-0.25) (5,-0.1667) (6,-0.0417) (7,0.0417) (8,0.0625) (9,0.0417) (10,0.0104) (11,-0.0104) (12,-0.0156) (13,-0.0104) (14,-0.0026) (15,0.0026) (16,0.0039) (17,0.0026) (18,0.0007) (19,-0.0007) (20,-0.001)};
\end{axis}
\end{tikzpicture}
\caption{Autocorrelation of the AR(2) with $\phi=(1,-0.5)$: a damped cosine of period $8$ under the
envelope $\pm(1/\sqrt2)^k$ (dashed).}\label{fig:ar2}
\end{figure}

\begin{casestudy}[A cycle in the monthly yield changes (2013 Case Study~3)]\label{ch09:cs-3b}
Fitting AR(2) by least squares to the monthly differences of the 10-year yield gives
$\hat\phi_1=0.038$, $\hat\phi_2=-0.217$ (standard errors $0.079$, $0.079$), $R^2=0.048$. The roots of the characteristic
polynomial are the complex pair $0.088\pm2.145\,\mathrm i$, of squared modulus $4.61>1$: the fitted model is
stationary. With complex roots the series is cyclical, with period
$2\pi/\arccos\bigl(|\hat\phi_1|/(2\sqrt{-\hat\phi_2})\bigr)=4.107$, that is, ``just over $4$ months''; I recomputed this value
(and the roots) numerically from the printed coefficients. The function \texttt{ar()} selects the order with the AIC
(\cref{ch09:sec-modelsel}): \texttt{order = 7}. In the AR(7) fit the significant lags are $2$ ($t=-3.10$) and $7$
($t=-2.02$), $R^2=0.123$; the smallest \emph{squared} modulus of the characteristic roots is $1.4027$ (the case study calls it
the ``smallest root modulus'': the modulus is $\sqrt{1.4027}=1.184$), still outside the unit circle. Regression diagnostics
of the AR(7) fit (\texttt{influence.measures}) flag two high-leverage months, October and November 2008, ``the heart of the
financial crisis'' (\cref{ch05:sec-diag}).
\end{casestudy}

\subsection{Moving-average models MA($q$)}\label{ch09:sec-ma}
Set $p=0$: $X_t-\mu=\theta(L)\eta_t=\eta_t+\theta_1\eta_{t-1}+\dots+\theta_q\eta_{t-q}$
\ts{24}{12}{1:14:24}\ts{13}{8}{1:06:12}. Because $\eta$ is white noise, only overlapping terms contribute:
\begin{equation}\label{ch09:eq-ma}
  \begin{gathered}
  \E X_t=\mu,\qquad \gamma(0)=\sigma^2\bigl(1+\theta_1^2+\dots+\theta_q^2\bigr),\\
  \gamma(j)=\begin{cases}\sigma^2\bigl(\theta_j+\theta_{j+1}\theta_1+\dots+\theta_q\theta_{q-j}\bigr),&1\le j\le q,\\ 0,& j>q,\end{cases}
  \end{gathered}
\end{equation}
with $\theta_0=1$ \ts{24}{12}{1:16:26}. The ACF of an MA($q$) \emph{cuts off} at lag $q$ (a complete finite-memory
process: $X_t$ and $X_{t+j}$ share no noise term when $j>q$), which is the identification rule of the table above.

\paragraph{Invertibility.} The process is \emph{invertible} if all roots of $\theta(z)=0$ lie outside the closed unit
disc; then $\theta^{-1}(L)$ exists by \eqref{ch09:eq-geom} and $\eta_t=\theta^{-1}(L)(X_t-\mu)$ is an infinite
autoregression of the observations \ts{24}{12}{1:15:24}. Invertibility matters for two reasons. Identification: for
MA(1), $\rho(1)=\theta/(1+\theta^2)$, which has maximum $1/2$ at $\theta=1$ and is unchanged under
$\theta\mapsto1/\theta$; for instance $\theta=0.25$ and $\theta=4$ both give $\rho(1)=4/17=0.235$, and if $\sigma^2$ is replaced by $\theta^2\sigma^2$ the variance
$\gamma(0)=\sigma^2(1+\theta^2)$ is also unchanged: two different parameter sets with the same second-order structure, hence
indistinguishable from the ACF. Interpretation: only the invertible
representation has $\eta_t$ equal to the one-step forecast error of $X_t$ given its past, which is the Wold innovation; one
therefore always fits the invertible member of the pair.

\subsection{ARMA($p,q$) and the ARMA(1,1) model}\label{ch09:sec-arma11}
A general ARMA($p,q$) is covariance stationary when $\phi(z)$ has all roots outside the unit disc and is invertible when
$\theta(z)$ does. (If $\phi$ and $\theta$ share a root the factors cancel and the model is redundant, a point to remember
when over-fitting.) The Wold weights are the power-series coefficients of $\theta(z)/\phi(z)$, and the
autocovariances satisfy the Yule--Walker recursion $\gamma(j)=\sum_i\phi_i\gamma(j-i)$ for $j>q$.

For ARMA(1,1), $X_t-\phi X_{t-1}=\phi_0+\eta_t+\theta\eta_{t-1}$ (2013 PS4 Problem~6 and 2024 PS4 Problem~6,
exercise \ref{ch09:ex-arma11}), matching the powers in $\psi(L)=(1+\theta L)/(1-\phi L)$ gives
$\psi_0=1$ and $\psi_j=\phi^{j-1}(\phi+\theta)$ for $j\ge1$. Then $\E X_t=\phi_0/(1-\phi)$ and, summing $\sum_j\psi_j^2$,
\begin{equation}\label{ch09:eq-arma11}\begin{gathered}
  \gamma(0)=\sigma^2\,\frac{1+2\phi\theta+\theta^2}{1-\phi^2},\\
  \rho(1)=\frac{(1+\phi\theta)(\phi+\theta)}{1+2\phi\theta+\theta^2}=\phi+\frac{\theta\sigma^2}{\gamma(0)},\\
  \rho(k)=\phi\,\rho(k-1),\ k\ge2 .
\end{gathered}\end{equation}
Compare with the AR(1), where $\rho(k)=\phi^k$ from the very first lag: the ARMA(1,1) also decays geometrically at
rate $\phi$, but from lag~2 on, and the first correlation $\rho(1)$ is a free parameter that need not equal $\phi$. When
$\phi<0$ the ACF of an AR(1) alternates in sign from lag~1; for an ARMA(1,1) it may be positive at lag~1 and alternate
afterwards. Its distinctive pattern is \emph{small first correlation but long memory}: with $\phi=0.9$,
$\theta=-0.85$ one finds $\rho(1)=0.061$ and $\rho(10)=0.024$, while an AR(1) with $\phi=0.9$ has $\rho(1)=0.9$. An economic
index following such a process would show almost no day-to-day autocorrelation together with a slow, persistent
decay. This structure arises for the \emph{squared} residuals of a GARCH(1,1) model, which follow an ARMA(1,1) with
$\phi=\alpha_1+\beta_1$ and $\theta=-\beta_1$ (\cref{ch:volatility}, \ts{13}{11}{5:15}).
The comparison with MA(1) is the reverse: an MA(1) has a single non-zero correlation and then nothing
(2024 PS4 Problem~6(e)).

% ==================================================================
\section{Non-stationarity: differencing, ARIMA and unit-root tests}\label{ch09:sec-nonstat}

Many economic series behave like random walks. Box and Jenkins advocate removing trending behaviour by
\emph{differencing} \ts{24}{12}{1:17:28}\ts{13}{8}{1:07:14}. With $\Delta=1-L$,
\[
  \Delta X_t=X_t-X_{t-1},\qquad
  \Delta^2X_t=(1-L)^2X_t=X_t-2X_{t-1}+X_{t-2},\qquad
  \Delta^k=(1-L)^k=\sum_{j=0}^k\binom kj(-L)^j .
\]
A linear trend $a+bt$ becomes the constant $b$ under $\Delta$; a quadratic trend is removed by $\Delta^2$
\ts{24}{12}{1:18:30}\ts{13}{8}{1:08:20}. In practice one uses $k\le2$. Modelling $\Delta X_t$ or $\Delta^2X_t$ as stationary
means modelling the dynamics of the slope, or of the curvature, of the series, with a constant level and variability
\ts{24}{12}{1:19:32}.

\subsection{Two kinds of trend}\label{ch09:sec-trends}
\paragraph{Trend reversion.} Let $X_t=a+bt+\xi_t$, with $\xi_t=\phi\xi_{t-1}+\eta_t$, $|\phi|<1$
\ts{13}{8}{1:09:21}. Then $\E X_t=a+bt$ (not constant, hence not stationary), $\Var X_t=\sigma^2/(1-\phi^2)$, and
\[
  \Delta X_t=b+\xi_t-\xi_{t-1}=b+(1-L)\xi_t=b+(1-L)(1-\phi L)^{-1}\eta_t,
\]
which is stationary. A shock to $\xi$ dies out and $X$ returns to the deterministic line.

\paragraph{Stochastic trend.} The pure \emph{integrated process} of order one, $I(1)$, is $X_t=X_{t-1}+\eta_t$, that is,
$\Delta X_t=\eta_t$ \ts{13}{8}{1:10:22}. Given $X_0$, $X_t=X_0+TS_t$ with $TS_t=\sum_{j=1}^t\eta_j$ (the
\emph{stochastic trend}), a random walk with white-noise steps, and
\begin{equation}\label{ch09:eq-rwmoments}
  \begin{gathered}
  \Var X_t=t\sigma^2,\qquad \Cov(X_t,X_{t-j})=(t-j)\sigma^2,\\
  \Corr(X_t,X_{t-j})=\frac{t-j}{\sqrt{t(t-j)}}=\sqrt{1-\frac jt}\ \longrightarrow 1\quad(t\to\infty),
  \end{gathered}
\end{equation}
so the process is not stationary: variance and covariances depend on $t$, and consecutive values become
almost perfectly correlated. The trend cannot be forecast (the best forecast at any horizon is the last value, with error
variance growing like $h\sigma^2$).

The two models differ where it matters for finance: in the trend-reversion model the long-run cumulative response
$\psi(1)$ of $\Delta X_t$ to a shock is $\psi(1)=(1-1)/(1-\phi)=0$, the level effect is temporary; for the stochastic trend $\Delta X_t=\eta_t$ has
$\psi(1)=1$, the shock is permanent. Prices are usually modelled as the second kind (\cref{ch04:rem-logprice}),
returns as stationary. Differencing a trend-reversion series (``over-differencing'') creates an MA factor $(1-L)$ with a unit root, which is
not invertible; differencing a stochastic trend is the correct treatment.

\begin{definition}[ARIMA$(p,d,q)$]\label{ch09:def-arima}
$\{X_t\}$ is \emph{integrated ARMA}, ARIMA$(p,d,q)$, if $\Delta^dX_t$ is stationary (and lower-order differencing is not) and
follows an ARMA$(p,q)$ model (2024 slide~25, not discussed in class; \ts{13}{8}{1:11:26}).
\end{definition}
The practical questions (\ts{13}{8}{1:12:26}) are: determine the order $d$ of differencing needed to remove deterministic or
stochastic trends (\cref{ch09:sec-df}); estimate the parameters (\cref{ch09:sec-est}); choose among competing
$(p,d,q)$ (\cref{ch09:sec-modelsel}).

\subsection{Testing for a unit root: the Dickey--Fuller test}\label{ch09:sec-df}
Suppose $X_t=\phi X_{t-1}+\eta_t$ with $\eta_t$ white noise. The \emph{Dickey--Fuller} (DF) test has
\[
  H_0:\ \phi=1\ \text{(unit root, non-stationary)}\qquad\text{against}\qquad H_1:\ |\phi|<1\ \text{(stationary)} .
\]
Fit the AR(1) by least squares and form $t_{\phi=1}=(\hat\phi-1)/\mathrm{se}(\hat\phi)$. Under $H_1$,
$\sqrt T(\hat\phi-\phi)\xrightarrow{d}N(0,1-\phi^2)$ (the standard AR(1) result). Under $H_0$ the estimator is
\emph{superconsistent}, converging at the rate $1/T$, and $T(\hat\phi-1)$, like the $t$ statistic, has a non-standard
distribution, the Dickey--Fuller distribution (2024 slide~28, not discussed in either year). It can be derived in a few lines. For
$X_t=X_{t-1}+\eta_t$ and $X_0=0$,
\[
  \hat\phi-1=\frac{\sum_tX_{t-1}\eta_t}{\sum_tX_{t-1}^2},\qquad
  \sum_{t=1}^TX_{t-1}\eta_t=\tfrac12\Bigl(X_T^2-\sum_{t=1}^T\eta_t^2\Bigr),
\]
because $X_t^2=X_{t-1}^2+2X_{t-1}\eta_t+\eta_t^2$ telescopes. Dividing the numerator by $T\sigma^2$ and the
denominator by $T^2\sigma^2$, and using $X_{\lfloor rT\rfloor}/(\sigma\sqrt T)\to W(r)$ (Donsker's theorem; the
random walk converges to Brownian motion, \cref{ch:sp2}),
\begin{equation}\label{ch09:eq-df}
  T(\hat\phi-1)\ \xrightarrow{d}\ \frac{\tfrac12\bigl(W(1)^2-1\bigr)}{\int_0^1W(r)^2\,\mathrm dr}
  =\frac{\int_0^1W\,\mathrm dW}{\int_0^1W^2\,\mathrm dr},
\end{equation}
the last equality being Itô's formula $\int_0^1W\,\mathrm dW=\tfrac12(W(1)^2-1)$ (\cref{ch:stochcalc}). The distribution is skewed to the
left, so the usual normal or $t$ critical values reject the unit root too often. Simulating $20\,000$ random walks of length $500$,
the $5\%$ critical value of $t_{\phi=1}$ is $-1.95$ (no intercept) and $-2.89$ (with an intercept; the asymptotic value is
about $-2.86$), compared with $-1.645$ for the standard normal; the $1\%$ values are $-2.55$ and $-3.46$.

The test is often \emph{augmented} (ADF, Said--Dickey): lagged differences $\Delta X_{t-1},\dots,\Delta X_{t-k}$ are added to the
regression to absorb serial correlation, and $k$ is chosen by an information criterion. The null hypothesis is
\emph{non-stationarity}, so failing to reject does not prove a unit root: the power is low against $\phi$ near $1$. Other tests:
Phillips--Perron (PP) and Elliott--Rothenberg--Stock (ERS) for a unit-root null, and KPSS, whose null is stationarity
\ts{13}{8}{1:12:26}. The 10-year Treasury yield is the standard illustration (\cref{ch09:cs-3a}): the ADF test does not reject a
unit root for the daily, weekly or monthly \emph{levels} and rejects it for the \emph{differences}. The same test was used in the crack-spread example
(\cref{ch09:sec-crack}).

\subsection{Estimation of ARMA models}\label{ch09:sec-est}
For Gaussian innovations $\eta_t\sim N(0,\sigma^2)$ i.i.d., the parameters are estimated by maximum likelihood
(\cref{ch05:sec-mle}) (2024 slide~26; \ts{13}{8}{1:12:26}).
\begin{itemize}
  \item \emph{AR($p$).} Conditioning on the first $p$ observations, the log-likelihood is
        $\ell=-\frac n2\log(2\pi\sigma^2)-\frac1{2\sigma^2}\sum_{t=p+1}^{T}\bigl(x_t-c-\sum_j\phi_jx_{t-j}\bigr)^2$ with $n=T-p$:
        conditional maximum likelihood \emph{is} least squares of $x_t$ on its $p$ lags (\cref{ch09:sec-wold-recipe}). This is
        the \emph{limited-information} (LIML) method, and the resulting $\hat\sigma^2=\mathrm{RSS}/n$ is
        biased downwards by the same degrees-of-freedom argument as in regression \ts{13}{8}{6:10}\ts{13}{8}{7:11}. The
        \emph{full-information} (FIML) method adds the log-density of the first $p$ values under their stationary distribution,
        which matters when $T$ is small or the roots are close to the unit circle.
  \item \emph{MA and ARMA.} The likelihood is not a regression, because the past innovations are not observed. LIML sets the
        first $q$ innovations to zero and computes the others recursively from the data; the exact likelihood is obtained by
        writing the model in state-space form and applying the Kalman filter (\cref{ch09:sec-ss}), the
        \emph{prediction-error decomposition}. The maximisation is numerical, and stationarity and invertibility constrain the
        parameters; the standard trick is to reparametrise so that the optimiser works on an unconstrained
        scale (logarithms, logits), as in the GARCH fits of \cref{ch:volatility}.
  \item \emph{Method of moments}: Yule--Walker (\cref{ch09:sec-yw}) for AR.
\end{itemize}

\subsection{Model selection: AIC, BIC, HQ}\label{ch09:sec-modelsel}
Fit all ARMA$(p,q)$ with $0\le p\le p_{\max}$, $0\le q\le q_{\max}$, let $\tilde\sigma^2(p,q)$ be the Gaussian MLE of
$\sigma^2$, and choose $(p,q)$ to minimise (2024 slide~27; \ts{13}{8}{1:13:33})
\begin{align}
  \mathrm{AIC}(p,q)&=\log\tilde\sigma^2(p,q)+2\,\frac{p+q}n, &&\text{(Akaike)}\notag\\
  \mathrm{BIC}(p,q)&=\log\tilde\sigma^2(p,q)+\log n\,\frac{p+q}n, &&\text{(Bayes / Schwarz)}\label{ch09:eq-ic}\\
  \mathrm{HQ}(p,q)&=\log\tilde\sigma^2(p,q)+2\log(\log n)\,\frac{p+q}n. &&\text{(Hannan--Quinn)}\notag
\end{align}
Because the maximised Gaussian log-likelihood is $\hat\ell=-\frac n2\bigl(\log(2\pi\tilde\sigma^2)+1\bigr)$, each criterion is
$-2\hat\ell/n$ plus a penalty proportional to the number of parameters (up to an additive constant); the lecture ends on the question
of what penalty is appropriate \ts{13}{8}{1:14:33}. Adding one parameter lowers $\log\tilde\sigma^2$ and raises the penalty. The larger
model is preferred iff $2\Delta\hat\ell>c_n$, the likelihood-ratio statistic (asymptotically $\chi^2_1$ under the smaller model)
exceeds $c_n$, with $c_n=2$ (AIC), $\log n$ (BIC) or $2\log\log n$ (HQ):
\begin{center}\small
\begin{tabular}{@{}l rrr@{}}
\toprule
 & AIC & BIC & HQ\\
\midrule
threshold $c_n$ for $n=100$ & $2$ & $4.61$ & $3.05$\\
\quad equivalent $p$-value & $0.157$ & $0.032$ & $0.081$\\
threshold $c_n$ for $n=500$ & $2$ & $6.21$ & $3.65$\\
\quad equivalent $p$-value & $0.157$ & $0.013$ & $0.056$\\
\bottomrule
\end{tabular}
\end{center}
(For $\chi^2_1$ the AIC threshold $2$ corresponds to $|t|>\sqrt2=1.41$.) So AIC is a liberal criterion, which tends to overfit but
is the right choice if the true model is infinite-dimensional and one wants good forecasts; BIC assumes the true model has a finite number
of parameters that it will find with probability tending to one; HQ makes the same finite-model assumption but with a penalty
growing as slowly as consistency allows \ts{13}{8}{1:15:39}. The same ideas were used for regression in
\cref{ch:regression} and are used in Case Study~3 (\texttt{ar()} minimises the AIC over $p\le22$ and picks $p=7$, with $p=2$ at a
relative AIC of $1.61$) and for VAR models below.

% ==================================================================
\section{Multivariate time series and VAR models}\label{ch09:sec-var}

We now observe $m$ series together, $\bm X_t=(X_{1,t},\dots,X_{m,t})\T\in\R^m$: several exchange rates, a basket of stocks, or
macroeconomic variables \ts{13}{11}{30:44}. $\{\bm X_t\}$ is covariance stationary if every component is
\ts{13}{11}{31:47}.

\subsection{Second-order structure}\label{ch09:sec-multi}
With $\bm\mu=\E\bm X_t$ define the \emph{covariance matrix} $\Gamma_0=\E[(\bm X_t-\bm\mu)(\bm X_t-\bm\mu)\T]$ (symmetric, with the
variances of the components on the diagonal), the \emph{lag-$k$ cross-covariance matrix}
$\Gamma_k=\E[(\bm X_t-\bm\mu)(\bm X_{t-k}-\bm\mu)\T]$, whose $(j,j')$ entry is $\Cov(X_{j,t},X_{j',t-k})$, and the
\emph{correlation matrices} $R_k=D^{-1/2}\Gamma_kD^{-1/2}$ with $D=\mathrm{diag}(\Gamma_0)$
\ts{13}{11}{31:47}\ts{13}{11}{33:03}\ts{13}{11}{35:02}. The matrix $\Gamma_0$ is symmetric, but $\Gamma_k$ for $k\ne0$ is
\emph{not}, and $\Gamma_{-k}=\Gamma_k\T$: the lags of series $j'$ may be correlated with series $j$ without the converse
\ts{13}{11}{36:04}. If $[\Gamma_k]_{jj'}\ne0$ for some $k>0$ one says that $X_{j'}$ \emph{leads} $X_j$ (its past carries information
about $X_j$); if each leads the other there is \emph{feedback} \ts{13}{11}{37:08}\ts{13}{11}{38:13}. This is the idea of
\emph{Granger causality} (Granger and Engle received the Nobel prize for it): lags of one variable help predict another. It is
predictability, not causation in a philosophical sense.

The \emph{multivariate Wold theorem} holds verbatim with matrices:
$\bm X_t=\bm V_t+\sum_{k\ge0}\Psi_k\bm\eta_{t-k}$ with $\Psi_0=I_m$, $\sum_k\Psi_k\Psi_k\T$ convergent, $\bm V_t$ linearly
deterministic, and $\bm\eta_t$ multivariate white noise: $\E\bm\eta_t=\bm 0$, $\Var\bm\eta_t=\Sigma$ (positive semi-definite),
$\Cov(\bm\eta_t,\bm\eta_{t-k})=0$ for $k\ne0$, uncorrelated with $\bm V$ \ts{13}{11}{39:18}\ts{13}{11}{41:24}. The
innovation is $\bm\eta_t=\bm X_t-\E[\bm X_t\mid\mathcal F_{t-1}]$ in the Gaussian case, where $\mathcal F_{t-1}$ is the information available
before time $t$: what is new at $t$ \ts{13}{11}{42:25}\ts{13}{11}{43:32}.

\subsection{VAR($p$) models and their stationarity}\label{ch09:sec-var-def}
\begin{definition}[VAR($p$)]\label{ch09:def-var}
$\{\bm X_t\}$ follows a \emph{vector autoregression} of order $p$ if
\begin{equation}\label{ch09:eq-var}
  \bm X_t=\bm C+\Phi_1\bm X_{t-1}+\Phi_2\bm X_{t-2}+\dots+\Phi_p\bm X_{t-p}+\bm\eta_t,
\end{equation}
with $\bm C\in\R^m$, $\Phi_1,\dots,\Phi_p$ real $m\times m$ matrices and $\{\bm\eta_t\}$ multivariate white noise
$N(\bm 0,\Sigma)$ \ts{13}{11}{45:39}.
\end{definition}
Component $j$ is an AR($p$) of its own lags \emph{plus} lag regressors on all the other series:
$X_{j,t}=c_j+\sum_{k=1}^p[\Phi_k]_{jj}X_{j,t-k}+\sum_{j'\ne j}\sum_{k=1}^p[\Phi_k]_{jj'}X_{j',t-k}+\eta_{j,t}$: exchange rates
that co-move, macro variables that respond to each other \ts{13}{11}{46:43}\ts{13}{11}{47:43}.

\paragraph{VAR(1) form.} Any VAR($p$) is a VAR(1) in higher dimension \ts{13}{11}{48:51}. Stack
\[
  \bm Z_t=(\bm X_t\T,\bm X_{t-1}\T,\dots,\bm X_{t-p+1}\T)\T\in\R^{mp}.
\]
Then
\begin{equation}\label{ch09:eq-companion}
  \bm Z_t=\bm D+A\bm Z_{t-1}+\bm F_t,\qquad
  A=\begin{pmatrix}\Phi_1&\Phi_2&\cdots&\Phi_{p-1}&\Phi_p\\ I_m&0&\cdots&0&0\\ 0&I_m&&0&0\\ \vdots&&\ddots&&\vdots\\ 0&0&\cdots&I_m&0\end{pmatrix},\
  \bm D=\begin{pmatrix}\bm C\\ \bm 0\\ \vdots\\ \bm 0\end{pmatrix},\
  \bm F_t=\begin{pmatrix}\bm\eta_t\\ \bm 0\\ \vdots\\ \bm 0\end{pmatrix},
\end{equation}
with $A$ the ($mp\times mp$) \emph{companion matrix} \ts{13}{11}{50:56}\ts{13}{11}{52:00}.

\begin{proposition}[Stationarity of a VAR]\label{ch09:prop-companion}
The following are equivalent:
\begin{enumerate}[(i)]
  \item all eigenvalues of $A$ have modulus less than~1;
  \item all roots of
      $\det\bigl(I_m-\Phi_1z-\dots-\Phi_pz^p\bigr)=0$ lie outside the closed unit disc $|z|\le1$.
\end{enumerate}
Under (i)--(ii) the VAR has the unique
stationary solution $\bm X_t-\bm\mu=\sum_{k\ge0}\Psi_k\bm\eta_{t-k}$, where $\Psi_k$ is the upper-left $m\times m$ block of $A^k$, and
\begin{equation}\label{ch09:eq-varmean}
\bm\mu=\E\bm X_t=(I_m-\Phi_1-\dots-\Phi_p)^{-1}\bm C .
\end{equation}
\end{proposition}
\begin{proof}
Let $\lambda\ne0$ be an eigenvalue of $A$ with eigenvector $\bm v=(\bm v_1,\dots,\bm v_p)$. The block rows of $A\bm v=\lambda\bm v$
say $\sum_j\Phi_j\bm v_j=\lambda\bm v_1$ and $\bm v_{j-1}=\lambda\bm v_j$ for $j=2,\dots,p$, so $\bm v_j=\lambda^{p-j}\bm v_p$. Substituting,
$\bigl(\lambda^pI-\sum_j\Phi_j\lambda^{p-j}\bigr)\bm v_p=0$ with $\bm v_p\ne0$, i.e.\ $\det\bigl(I-\sum_j\Phi_jz^j\bigr)=0$ at $z=1/\lambda$;
the converse is the same computation backwards. Thus the eigenvalues of $A$ are the reciprocals of the roots, and (i)$\Leftrightarrow$(ii).
If (i) holds, $\bm Z_t-\bm\mu_Z=\sum_{k\ge0}A^k\bm F_{t-k}$ converges because the spectral radius of $A$ is $<1$ (so $\|A^k\|$ decays
geometrically), and reading off the first block gives the expansion. Taking expectations in \eqref{ch09:eq-var} gives $\bm\mu=\bm C+\sum\Phi_j\bm\mu$.
\end{proof}
The characteristic polynomial $\det(I_m-\Phi_1z-\dots-\Phi_pz^p)$ has degree $mp$, hence $mp$ roots
\ts{13}{11}{53:06}\ts{13}{11}{54:08}; the mean-adjusted series satisfies $\bm X_t-\bm\mu=\sum_j\Phi_j(\bm X_{t-j}-\bm\mu)+\bm\eta_t$
\ts{13}{11}{55:13}\ts{13}{11}{56:18}. The scalar case $m=1$ is \cref{ch09:thm-ar-stat}, and its proof-by-eigenvalue is
\cref{ch09:rem-companion}.

\begin{example}[Moments of a VAR(1)]\label{ch09:ex-var1}
For $\bm X_t=\bm C+A\bm X_{t-1}+\bm\eta_t$ with stationary $A$, let $\tilde{\bm X}_t=\bm X_t-\bm\mu$. Then
$\Gamma_0=A\Gamma_0A\T+\Sigma$, and with the vectorisation of \cref{ch09:sec-varest} this is solved by
$\mathrm{vec}\,\Gamma_0=(I-A\otimes A)^{-1}\mathrm{vec}\,\Sigma$; further, $\Gamma_h=A^h\Gamma_0$ for $h\ge0$. Take
$A=\bigl(\begin{smallmatrix}0.5&0.2\\-0.3&0.6\end{smallmatrix}\bigr)$, $\bm C=(1,0.5)\T$,
$\Sigma=\bigl(\begin{smallmatrix}1&0.3\\0.3&0.5\end{smallmatrix}\bigr)$. The eigenvalues of $A$ are $0.55\pm0.24\,\mathrm i$
(modulus $0.6<1$), $\bm\mu=(1.923,-0.192)\T$,
$\Gamma_0=\bigl(\begin{smallmatrix}1.443&0.244\\0.244&0.847\end{smallmatrix}\bigr)$ and
$\Gamma_1=A\Gamma_0=\bigl(\begin{smallmatrix}0.770&0.291\\-0.287&0.435\end{smallmatrix}\bigr)$, which is not symmetric (checked against the
iteration $\Gamma\leftarrow A\Gamma A\T+\Sigma$). This is the pattern that 2013 PS6 Problem~1 asks for with different numbers.
\end{example}

\subsection{VAR as a multivariate regression: estimation}\label{ch09:sec-varest}
Given $n$ observations $\bm x_1,\dots,\bm x_n$ and $p$ pre-sample values, define for each component $j$ the regression
$\bm y^{(j)}=Z\bm\beta^{(j)}+\bm\varepsilon^{(j)}$, where $\bm y^{(j)}=(x_{j,1},\dots,x_{j,n})\T$ and the \emph{same}
$n\times(mp+1)$ design matrix $Z$ has rows $(1,\bm x_{t-1}\T,\dots,\bm x_{t-p}\T)$ \ts{13}{11}{57:18}\ts{13}{11}{58:32}. The $\bm\beta^{(j)}$ contain the
$j$th rows of $\bm C,\Phi_1,\dots,\Phi_p$. Stacking columns,
\begin{equation}\label{ch09:eq-mvreg}
  Y=ZB+E,\qquad Y=[\bm y^{(1)}\cdots\bm y^{(m)}],\ B=[\bm\beta^{(1)}\cdots\bm\beta^{(m)}],
\end{equation}
a multivariate regression, called \emph{seemingly unrelated regressions} (SUR) in econometrics because the $m$ equations are
linked only through the correlation $\Sigma$ of their errors \ts{13}{11}{59:41}\ts{13}{11}{1:01:51}\ts{13}{11}{1:03:53}.

\paragraph{Component-wise least squares.} Fit each equation separately,
$\hat{\bm\beta}^{(j)}=(Z\T Z)^{-1}Z\T\bm y^{(j)}$, that is $\hat B=(Z\T Z)^{-1}Z\T Y$, take the residual matrix
$\hat E=Y-Z\hat B=(I_n-Z(Z\T Z)^{-1}Z\T)Y$, and estimate the innovation covariance by
$\hat\Sigma=\hat E\T\hat E/(n-mp-1)$ (unbiased entry by entry; the lecture notes divide by $n-pm$) \ts{13}{11}{1:04:58}.
Treating the equations separately looks wasteful, since the errors are correlated. It is not:

\begin{theorem}[Optimality of component-wise OLS]\label{ch09:thm-varols}
For a VAR($p$) with unrestricted coefficient matrices,
\begin{enumerate}[(i)]
  \item component-wise OLS equals the generalised least-squares (GLS) estimator for the
      whole system with any innovation covariance $\Sigma$, hence is BLUE (Gauss--Markov, \cref{ch05:thm-gm}, and \cref{ch05:sec-gls});
  \item with Gaussian
      innovations it is the maximum-likelihood estimator of $B$, and the MLE of $\Sigma$ is $\tilde\Sigma=\hat E\T\hat E/n$
      \ts{13}{11}{1:06:58}\ts{13}{11}{1:08:00}.
\end{enumerate}
\end{theorem}
\begin{proof}
Let $\otimes$ be the Kronecker product ($A\otimes B$ is the block matrix with blocks $a_{ij}B$; $(A\otimes B)\T=A\T\otimes B\T$ and
$(A\otimes B)(C\otimes D)=AC\otimes BD$) and $\mathrm{vec}$ the operator that stacks the columns of a matrix into one vector
\ts{13}{11}{1:09:04}\ts{13}{11}{1:10:04}. Put $\bm y=\mathrm{vec}\,Y$, $\mathcal X=I_m\otimes Z$, $\bm\beta^*=\mathrm{vec}\,B$,
$\bm\varepsilon=\mathrm{vec}\,E$. Then $\bm y=\mathcal X\bm\beta^*+\bm\varepsilon$ with $\Cov\bm\varepsilon=\Sigma\otimes I_n$, and
$(\Sigma\otimes I_n)^{-1}=\Sigma^{-1}\otimes I_n$. The GLS estimator is
$[\mathcal X\T(\Sigma^{-1}\otimes I_n)\mathcal X]^{-1}\mathcal X\T(\Sigma^{-1}\otimes I_n)\bm y$ with
\[
  \mathcal X\T(\Sigma^{-1}\otimes I_n)\mathcal X=(I\otimes Z\T)(\Sigma^{-1}\otimes I_n)(I\otimes Z)=\Sigma^{-1}\otimes Z\T Z,\qquad
  \mathcal X\T(\Sigma^{-1}\otimes I_n)=\Sigma^{-1}\otimes Z\T ,
\]
so it equals $[\Sigma\otimes(Z\T Z)^{-1}](\Sigma^{-1}\otimes Z\T)\bm y=[I_m\otimes(Z\T Z)^{-1}Z\T]\bm y=\mathrm{vec}\bigl((Z\T Z)^{-1}Z\T Y\bigr)$:
$\Sigma$ cancels and one recovers component-wise OLS \ts{13}{11}{1:11:06}\ts{13}{11}{1:12:08}. This proves (i).
For (ii), the Gaussian log-likelihood is
\[
  \ell(B,\Sigma)=-\tfrac{nm}2\log2\pi-\tfrac n2\log|\Sigma|-\tfrac12\mathrm{tr}\bigl[\Sigma^{-1}(Y-ZB)\T(Y-ZB)\bigr]
\]
\ts{13}{11}{1:13:12}\ts{13}{11}{1:14:12}. Writing $Y-ZB=\hat E+Z(\hat B-B)$ and using $Z\T\hat E=0$, the trace equals
$\mathrm{tr}[\Sigma^{-1}\hat E\T\hat E]+\mathrm{tr}[\Sigma^{-1}(\hat B-B)\T Z\T Z(\hat B-B)]$, and the second term is $\ge0$; so for \emph{every}
$\Sigma$ the maximum over $B$ is at $\hat B$ (a \emph{concentrated} likelihood, because $\hat B$ does not depend on $\Sigma$
\ts{13}{11}{1:15:16}\ts{13}{11}{1:16:19}). The remaining function is $-\frac n2[\log|\Sigma|+\mathrm{tr}(\Sigma^{-1}S)]$ with
$S=\hat E\T\hat E/n$. Let $u_1,\dots,u_m>0$ be the eigenvalues of $\Sigma^{-1/2}S\Sigma^{-1/2}$; then
$\log|\Sigma|+\mathrm{tr}(\Sigma^{-1}S)=\log|S|+\sum_i(u_i-\log u_i)\ge\log|S|+m$, since $u-\log u\ge1$, with equality iff all $u_i=1$,
i.e.\ $\Sigma=S$. (This is the Anderson--Olkin result quoted in the lecture \ts{13}{11}{1:16:19}.)
\end{proof}

\paragraph{Order selection and asymptotics.} As in the univariate case, fit $\mathrm{VAR}(p)$ for $0\le p\le p_{\max}$ and choose $p$ by
\begin{equation}\label{ch09:eq-icvar}
  \begin{gathered}
  \mathrm{AIC}(p)=\log|\tilde\Sigma(p)|+2\,\frac{pm^2}n,\qquad
  \mathrm{BIC}(p)=\log|\tilde\Sigma(p)|+\log n\,\frac{pm^2}n,\\
  \mathrm{HQ}(p)=\log|\tilde\Sigma(p)|+2\log(\log n)\,\frac{pm^2}n
  \end{gathered}
\end{equation}
\ts{13}{11}{1:17:24}. (The slides print $-\log|\tilde\Sigma(p)|$, a sign misprint: $\log|\tilde\Sigma|$ is the log
generalised variance, which decreases as $p$ increases, while the penalty increases; compare \cref{ch09:eq-ic}.) For a stationary VAR whose
innovations have bounded fourth moments and are independent over time, $Q=\mathrm{plim}\,Z\T Z/n$ exists and is non-singular and
$\sqrt n(\hat{\bm\beta}^*-\bm\beta^*)\to N(0,\Sigma\otimes Q^{-1})$; least-squares and Gaussian maximum-likelihood estimates are asymptotically
identical (2013 notes, last slide).

\subsection{Impulse responses and a macroeconomic VAR}\label{ch09:sec-irf}
The matrices $\Psi_k$ of the Wold expansion are the \emph{impulse response functions}: $[\Psi_k]_{jj'}$ is the effect on variable $j$, $k$ periods
later, of a unit innovation in variable $j'$ \ts{13}{11}{1:21:42}\ts{13}{11}{1:22:42}\ts{13}{11}{1:23:43}. Because $\Sigma$ is not
diagonal the innovations are correlated, and R's \texttt{irf} plots the \emph{orthogonalised} responses (shocks made uncorrelated through a Cholesky
factor of $\Sigma$, which depends on the ordering of the variables; a point the course does not discuss).

\begin{casestudy}[VAR of US macroeconomic series (2013 Case Study~5 and 2013 L11--L12)]\label{ch09:cs-5}
In the 1980s Litterman and Sims used vector autoregressions of nine macroeconomic aggregates (Treasury-bill rate, M1, GNP deflator, real GNP,
business fixed investment, unemployment, trade-weighted dollar, S\&P~500, commodity price index) to forecast the economy and to assess
the effect of policy variables such as the Federal Funds rate. The case study downloads monthly series from FRED with
\texttt{quantmod} (UNRATE, FEDFUNDS, TB3MS, CPIAUCSL, M1SL, GDPDEF, GDP, GPDI, TWEXBMTH, and the S\&P~500 from Yahoo) and fits VAR models with the
\texttt{vars} package to three of them: the unemployment rate, the Federal Funds rate and the CPI, over 1960--2000 (492 months)
\ts{13}{11}{1:18:29}\ts{13}{11}{1:19:33}\ts{13}{11}{1:20:38}.
\begin{itemize}
  \item \emph{Levels.} \texttt{VARselect} (up to 12 lags) returns $p=12,5,2,12$ for AIC, HQ, SC (=BIC), FPE. The VAR(2) chosen by SC has
        characteristic-root moduli $1.002,\,0.986,\,0.952,\,0.468,\,0.331,\,0.084$: the largest exceeds 1, so the fitted model is
        \emph{non-stationary} \ts{13}{12}{0:00}\ts{13}{12}{1:08}. (Here ``roots'' are eigenvalues of the companion matrix, stationary when $<1$;
        \cref{ch09:rem-companion}.)
  \item \emph{Differences.} Removing the random-walk component by first differences, the ACF and PACF show strong self-correlation and some cross-correlation between
        unemployment and Fed funds \ts{13}{12}{2:11}\ts{13}{12}{3:11}\ts{13}{12}{4:11}. \texttt{VARselect} gives $p=12,3,3,12$; the VAR(3) has
        all root moduli $<0.84$ (stationary). The dependent variable is the monthly change; the lecture's reading of the Fed funds equation
        \ts{13}{12}{6:15}\ts{13}{12}{7:21}:
        \[
          \Delta\mathrm{FF}_t=-0.71\,\Delta\mathrm{UN}_{t-1}+0.37\,\Delta\mathrm{FF}_{t-1}-0.18\,\Delta\mathrm{FF}_{t-2}+\dots\quad(t=-4.9,\ 8.0,\ -3.7).
        \]
        If unemployment rises, the Fed lowers rates the next period; a rate change tends to be followed by another in the same direction
        (the Fed does not want to shock the economy), and since $0.37\,x_{t-1}-0.18\,x_{t-2}=0.19\,x_{t-1}+0.18\,(x_{t-1}-x_{t-2})$ the
        negative second lag is an \emph{acceleration} term \ts{13}{12}{8:21}\ts{13}{12}{9:24}\ts{13}{12}{10:30}. The CPI equation is dominated by its own lags
        ($0.42$, $t=9.5$, and $0.29$ at lag~3, $t=6.7$), the unemployment equation by lags 2 and 3 of itself; $R^2$ are $0.12$ (unemployment),
        $0.23$ (Fed funds), $0.49$ (CPI); the residual correlation between unemployment and Fed funds is $-0.20$.
  \item \emph{Impulse responses} (bootstrap $95\%$ bands, 100 runs) \ts{13}{11}{1:22:42}. Levels VAR: after a rise in unemployment the Fed funds rate
        and the CPI decline (less employment, less pressure on prices); after a rise in the Fed funds rate unemployment and the CPI increase; after a
        rise in the CPI the Fed funds rate increases (consistent with policy against inflation). Differenced VAR: unemployment up $\Rightarrow$ Fed funds
        down over the next quarters; Fed funds up $\Rightarrow$ a modest increase in inflation, which persists for several quarters (the large third-lag
        partial autocorrelation of the CPI); CPI up $\Rightarrow$ modestly higher unemployment and higher Fed funds.
\end{itemize}
Three variables are far too few for a structural model of the economy: the lecture recommends ten to twelve
\ts{13}{12}{5:13}.
\end{casestudy}

% ==================================================================
\section{Cointegration}\label{ch09:sec-coint}

Differencing removes non-stationarity but discards information about the \emph{levels}. Cointegration is the framework that keeps it
\ts{13}{12}{10:30}\ts{13}{12}{11:36}.

\subsection{Definition and examples}
A process $\{\bm X_t\}$ is \emph{integrated of order $d$}, $I(d)$, if $(1-L)^d\bm X_t$ is stationary. If it has a VAR representation
$\Phi(L)\bm X_t=\bm\eta_t$ then $\Phi(L)=(1-L)^d\Phi^*(L)$ with $\Phi^*$ describing the stationary VAR($p-d$) of $\Delta^d\bm X_t$
\ts{13}{12}{12:43}\ts{13}{12}{13:43}. The point is that all the components may be $I(1)$ without the vector being ``jointly'' $I(1)$: linear
combinations of the levels may be stationary \ts{13}{12}{14:47}.

\begin{definition}[Cointegration]\label{ch09:def-coint}
An $m$-vector $\{\bm X_t\}$ of $I(1)$ series is \emph{cointegrated} if there is $\bm\beta\in\R^m$, $\bm\beta\ne\bm0$, with $\bm\beta\T\bm X_t$ an $I(0)$ (stationary)
process \ts{13}{12}{15:49}. The vector $\bm\beta$ is defined up to scale; normalising $\beta_1=1$,
$\bm\beta\T\bm X_t=u_t$ says $x_{1,t}=-(\beta_2x_{2,t}+\dots+\beta_mx_{m,t})+u_t$: a \emph{long-run equilibrium relationship}
$x_1=-(\beta_2x_2+\dots)$ with a stationary \emph{disequilibrium error} (cointegration residual) $u_t$ \ts{13}{12}{16:49}.
\end{definition}

\begin{finance}[Where cointegration appears]\label{ch09:fin-coint}
Examples listed in the lecture \ts{13}{12}{16:49}\ts{13}{12}{17:50}: the term structure of interest rates (the level of rates is non-stationary, the
\emph{spreads} between maturities are stationary: expectations hypothesis); purchasing-power parity (exchange rate and price levels of two countries: otherwise there would be arbitrage);
money demand (money, income, prices, interest rates); covered interest-rate parity (spot and forward rates); the law of one price (identical or equivalent
assets that must be valued alike: spot and futures, the same asset on two venues). In trading, a cointegrated pair is the basis of
\emph{convergence} (pairs) trades: two prices that wander individually but whose spread mean-reverts.
\end{finance}

\subsection{The error-correction (VECM) representation}\label{ch09:sec-vecm}
Let $\bm X_t$ follow a VAR($p$) \eqref{ch09:eq-var} and be $I(1)$. Subtract $\bm X_{t-1}$ from both sides, then add and subtract
$(\Phi_1-I)\bm X_{t-2}$, then $(\Phi_2+\Phi_1-I)\bm X_{t-3}$, and so on \ts{13}{12}{27:08}\ts{13}{12}{28:12}\ts{13}{12}{29:13}. The result is
\begin{equation}\label{ch09:eq-vecm}
  \begin{gathered}
  \Delta\bm X_t=\bm C+\Pi\,\bm X_{t-1}+\Gamma_1\Delta\bm X_{t-1}+\dots+\Gamma_{p-1}\Delta\bm X_{t-p+1}+\bm\eta_t,\\
  \Pi=\Phi_1+\dots+\Phi_p-I_m,\qquad \Gamma_k=-\sum_{j=k+1}^p\Phi_j,
  \end{gathered}
\end{equation}
the \emph{vector error-correction model} (VECM) \ts{13}{12}{30:14}\ts{13}{12}{31:14}. (Proof: for $p=2$, $\bm X_t=\bm C+\Phi_1\bm X_{t-1}+\Phi_2\bm X_{t-2}+\bm\eta_t$ is written
$\bm X_{t-2}=\bm X_{t-1}-\Delta\bm X_{t-1}$ in the last term, giving $\Delta\bm X_t=\bm C+(\Phi_1+\Phi_2-I)\bm X_{t-1}-\Phi_2\Delta\bm X_{t-1}+\bm\eta_t$; the general case
follows by induction.) Everything in \eqref{ch09:eq-vecm} is $I(0)$ except possibly $\Pi\bm X_{t-1}$, which must therefore be stationary too: it contains the
cointegrating combinations. Since $\bm X$ has unit roots, $\Pi$ is singular. Three cases \ts{13}{12}{32:19}\ts{13}{12}{33:19}:
\begin{itemize}
  \item $\mathrm{rank}\,\Pi=0$, $\Pi=0$: no cointegration; model the differences by a stationary VAR($p-1$);
  \item $\mathrm{rank}\,\Pi=m$: $\bm X$ is itself stationary (no unit roots);
  \item $0<r=\mathrm{rank}\,\Pi<m$: $\Pi=\alpha\beta\T$ with $\alpha,\beta$ of size $m\times r$ and full rank $r$. The $r$ columns of $\beta$ are linearly independent
        cointegrating vectors; the matrix $\alpha$ holds the \emph{adjustment} (error-correction) coefficients: the process is pulled back towards the equilibrium
        $\beta\T\bm X=0$ at speed $\alpha$.
\end{itemize}
The factorisation is not unique: $\alpha\beta\T=(\alpha G)(\beta G^{-\mathsf T})\T$ for any invertible $r\times r$ matrix $G$. The stationary
directions form an $r$-dimensional space; a choice of coordinates in it is convention \ts{13}{12}{34:22}.

\begin{example}[A cointegrated pair]\label{ch09:ex-coint}
Let $x_t=x_{t-1}+\varepsilon_{1t}$ (random walk), $u_t=0.7u_{t-1}+\varepsilon_{2t}$ and $y_t=2x_t+u_t$. Both $x$ and $y$ are $I(1)$, but $y_t-2x_t=u_t$ is
stationary: $\bm\beta=(-2,1)\T$. In error-correction form
$\Delta x_t=\varepsilon_{1t}$ and $\Delta y_t=2\Delta x_t+u_t-u_{t-1}=-0.3\,(y_{t-1}-2x_{t-1})+2\varepsilon_{1t}+\varepsilon_{2t}$. Hence
$\Delta\bm X_t=\alpha\,\bm\beta\T\bm X_{t-1}+\text{noise}$ with $\alpha=(0,-0.3)\T$, and $\Pi=\alpha\bm\beta\T=\bigl(\begin{smallmatrix}0&0\\0.6&-0.3\end{smallmatrix}\bigr)$
has rank one. In the VAR(1) in levels the matrix $I+\Pi$ has eigenvalues $1$ (the common stochastic trend) and $0.7$ (the stationary spread).
Only $y$ reacts to a disequilibrium: $x$ is ``weakly exogenous'' (the row of $\alpha$ for $x$ is zero). \Cref{ch09:fig-coint} shows a simulated path.
\end{example}

\begin{figure}[ht]
\centering
\begin{tikzpicture}
\begin{axis}[width=0.95\linewidth,height=5.4cm,xlabel={$t$},xmin=0,xmax=119,grid=major,grid style={black!10},legend style={font=\footnotesize,at={(0.5,1.03)},anchor=south,legend columns=3}]
  \addplot[thick,color=defcol] coordinates {(0,0.0) (1,0.3) (2,0.03) (3,-0.86) (4,-1.32) (5,-2.31) (6,-2.25) (7,-0.91) (8,-1.4) (9,-2.02) (10,-1.53) (11,-1.18) (12,-1.07) (13,-2.0) (14,-2.03) (15,-1.34) (16,-2.68) (17,-3.14) (18,-5.04) (19,-6.33) (20,-8.17) (21,-8.41) (22,-9.67) (23,-9.4) (24,-9.24) (25,-9.43) (26,-11.95) (27,-12.49) (28,-12.54) (29,-12.42) (30,-13.95) (31,-14.43) (32,-15.41) (33,-16.22) (34,-15.16) (35,-15.96) (36,-16.0) (37,-15.11) (38,-15.7) (39,-15.81) (40,-15.7) (41,-15.63) (42,-16.86) (43,-16.78) (44,-15.42) (45,-16.97) (46,-16.11) (47,-15.99) (48,-16.63) (49,-14.63) (50,-13.87) (51,-15.07) (52,-15.0) (53,-14.42) (54,-14.61) (55,-13.92) (56,-13.99) (57,-13.32) (58,-11.89) (59,-12.56) (60,-12.36) (61,-12.82) (62,-12.69) (63,-13.88) (64,-14.46) (65,-14.66) (66,-13.76) (67,-12.61) (68,-13.94) (69,-14.73) (70,-14.08) (71,-16.08) (72,-16.54) (73,-16.64) (74,-15.38) (75,-14.69) (76,-15.02) (77,-15.39) (78,-15.64) (79,-14.11) (80,-14.54) (81,-14.84) (82,-14.49) (83,-14.61) (84,-14.81) (85,-15.92) (86,-15.94) (87,-16.38) (88,-15.21) (89,-14.56) (90,-14.58) (91,-13.92) (92,-14.26) (93,-13.2) (94,-13.21) (95,-12.63) (96,-13.92) (97,-13.57) (98,-15.26) (99,-17.29) (100,-17.6) (101,-18.5) (102,-18.33) (103,-16.09) (104,-16.92) (105,-17.54) (106,-17.34) (107,-16.85) (108,-17.02) (109,-17.23) (110,-16.53) (111,-16.01) (112,-17.04) (113,-17.12) (114,-17.08) (115,-18.14) (116,-17.88) (117,-18.74) (118,-17.76) (119,-17.57)};
  \addplot[thick,color=thmcol] coordinates {(0,0.0) (1,0.65) (2,-0.27) (3,-2.02) (4,-4.04) (5,-6.28) (6,-5.45) (7,-3.76) (8,-3.66) (9,-5.69) (10,-3.76) (11,-3.35) (12,-2.37) (13,-4.08) (14,-5.04) (15,-2.61) (16,-4.45) (17,-5.68) (18,-9.82) (19,-12.57) (20,-16.87) (21,-16.52) (22,-19.47) (23,-18.92) (24,-19.05) (25,-19.63) (26,-25.2) (27,-25.13) (28,-25.27) (29,-24.41) (30,-27.59) (31,-29.06) (32,-31.15) (33,-33.0) (34,-30.71) (35,-32.43) (36,-32.52) (37,-31.42) (38,-32.71) (39,-31.55) (40,-31.75) (41,-32.15) (42,-34.13) (43,-33.01) (44,-31.33) (45,-34.41) (46,-32.93) (47,-33.53) (48,-33.91) (49,-29.73) (50,-28.02) (51,-30.79) (52,-30.17) (53,-29.29) (54,-29.62) (55,-28.79) (56,-29.37) (57,-26.82) (58,-24.2) (59,-25.24) (60,-24.82) (61,-25.98) (62,-25.93) (63,-27.76) (64,-29.1) (65,-29.53) (66,-27.65) (67,-24.62) (68,-27.04) (69,-28.65) (70,-27.94) (71,-32.82) (72,-32.98) (73,-32.62) (74,-30.39) (75,-28.79) (76,-29.16) (77,-29.66) (78,-29.94) (79,-27.57) (80,-27.71) (81,-29.48) (82,-28.32) (83,-28.46) (84,-28.56) (85,-29.98) (86,-29.67) (87,-31.91) (88,-30.84) (89,-28.92) (90,-29.64) (91,-28.17) (92,-28.24) (93,-27.21) (94,-28.24) (95,-26.37) (96,-28.59) (97,-27.82) (98,-30.97) (99,-35.42) (100,-36.69) (101,-38.14) (102,-38.05) (103,-34.13) (104,-34.91) (105,-35.87) (106,-34.98) (107,-34.5) (108,-35.0) (109,-35.73) (110,-34.47) (111,-32.89) (112,-35.16) (113,-34.78) (114,-34.34) (115,-35.19) (116,-35.83) (117,-36.99) (118,-35.24) (119,-34.95)};
  \addplot[thick,color=intcol] coordinates {(0,0.0) (1,0.05) (2,-0.32) (3,-0.29) (4,-1.4) (5,-1.66) (6,-0.95) (7,-1.94) (8,-0.85) (9,-1.64) (10,-0.7) (11,-0.99) (12,-0.23) (13,-0.08) (14,-0.98) (15,0.06) (16,0.91) (17,0.6) (18,0.25) (19,0.08) (20,-0.53) (21,0.29) (22,-0.12) (23,-0.12) (24,-0.56) (25,-0.77) (26,-1.3) (27,-0.16) (28,-0.2) (29,0.44) (30,0.31) (31,-0.2) (32,-0.33) (33,-0.57) (34,-0.39) (35,-0.5) (36,-0.53) (37,-1.2) (38,-1.32) (39,0.07) (40,-0.36) (41,-0.88) (42,-0.41) (43,0.55) (44,-0.48) (45,-0.46) (46,-0.7) (47,-1.55) (48,-0.64) (49,-0.46) (50,-0.28) (51,-0.65) (52,-0.18) (53,-0.45) (54,-0.4) (55,-0.95) (56,-1.39) (57,-0.17) (58,-0.43) (59,-0.12) (60,-0.11) (61,-0.34) (62,-0.54) (63,-0.0) (64,-0.18) (65,-0.22) (66,-0.14) (67,0.61) (68,0.83) (69,0.81) (70,0.23) (71,-0.67) (72,0.1) (73,0.65) (74,0.37) (75,0.59) (76,0.88) (77,1.11) (78,1.33) (79,0.66) (80,1.37) (81,0.21) (82,0.67) (83,0.76) (84,1.06) (85,1.87) (86,2.2) (87,0.85) (88,-0.42) (89,0.2) (90,-0.47) (91,-0.34) (92,0.27) (93,-0.8) (94,-1.82) (95,-1.12) (96,-0.76) (97,-0.68) (98,-0.45) (99,-0.83) (100,-1.49) (101,-1.14) (102,-1.38) (103,-1.95) (104,-1.06) (105,-0.78) (106,-0.3) (107,-0.81) (108,-0.96) (109,-1.27) (110,-1.42) (111,-0.88) (112,-1.08) (113,-0.55) (114,-0.18) (115,1.09) (116,-0.07) (117,0.48) (118,0.28) (119,0.19)};
  \legend{$x_t$ (random walk),$y_t=2x_t+u_t$,$u_t=y_t-2x_t$ (stationary)}
\end{axis}
\end{tikzpicture}
\caption{A simulated cointegrated pair (seed~7). Each of $x$ and $y$ wanders like a random walk, but $y-2x$ stays around zero.}\label{ch09:fig-coint}
\end{figure}

\subsection{Estimation and tests for the cointegration rank}\label{ch09:sec-cointest}
Sims, Stock and Watson (1990) and Park and Phillips (1989) prove that least squares applied to the original VAR in \emph{levels}
gives consistent estimates with the same asymptotic distributions as maximum likelihood, and the reduced-rank constraints on $\Pi$ hold
asymptotically: no new tool is needed for estimation \ts{13}{12}{34:22}\ts{13}{12}{35:33}\ts{13}{12}{36:40}. For inference on the \emph{rank}
$r$, Johansen (1991) developed maximum likelihood by reduced-rank regression for Gaussian innovations, following Banerjee and Hendry's use of the
concentrated likelihood \ts{13}{12}{37:43}. It gives likelihood-ratio tests of the number of cointegrating vectors: Johansen's \emph{trace} statistic
$-T\sum_{i>r}\log(1-\hat\lambda_i)$ (a sum over the smallest eigenvalues, testing rank $\le r$) and his \emph{maximum-eigenvalue} statistic
$-T\log(1-\hat\lambda_{r+1})$ (rank $r$ against $r+1$), where the $\hat\lambda_i$ are squared sample canonical correlations between
$\Delta\bm X_t$ and $\bm X_{t-1}$ after adjusting for the lagged differences (the slides describe them as the eigenvalues of $\hat\Pi$; I give the standard
definition). A cruder route, not in the course, is the Engle--Granger two-step method: regress one series on the others and apply an ADF test
to the residual (with different critical values).

\subsection{Energy futures and crack spreads}\label{ch09:sec-crack}
The lecture's example of cointegration uses futures on the CME from the Energy Information Administration: WTI crude oil (CL, quoted in dollars per barrel), RBOB
gasoline and heating oil (both in cents per gallon), from 2006 to 2013, first-month contracts. A refinery ``cracks'' crude oil into gasoline, heating oil, jet fuel;
multiplying the product prices by $42$ gallons per barrel puts inputs and outputs in the same units, and the products stay above crude, so the
\emph{crack spreads} (product minus crude) represent the value added of refining \ts{13}{12}{18:54}\ts{13}{12}{19:55}\ts{13}{12}{20:56}\ts{13}{12}{22:01}. The class
listed what drives the spread: cost of refining, supply and demand, seasonality (heating oil dearer in winter), and the refiner's profit
\ts{13}{12}{23:02}\ts{13}{12}{24:05}; the spreads look stationary \ts{13}{12}{25:05}. The case note (not in the download) tests the levels with the Augmented Dickey--Fuller test:
p-values $0.164$ for the first crude contract and $0.327$ for heating oil, so a unit root cannot be rejected \ts{13}{12}{38:47}\ts{13}{12}{39:49}. In the Johansen
procedure the test of rank $0$ is almost significant at the $10\%$ level and the test of rank $\le1$ is not, so the rank is at most one, with cointegrating vector $(1,\,1.3,\,-1.7)$ on crude, RBOB and heating oil, so that
``crude plus gasoline minus heating oil'' is the stationary combination (ignoring seasonality) \ts{13}{12}{41:55}\ts{13}{12}{43:00}. The first-month contract was chosen because
refiners hedge on the futures market, so the analysis addresses hedging refinery risk; longer maturities would be used depending on what one hedges \ts{13}{12}{44:07}.
The video jumps here between explanations, so the numbers are reported as spoken; commodity models are the subject of \cref{ch:commodities}.

% ==================================================================
\section{Linear state-space models and the Kalman filter}\label{ch09:sec-ss}

Many models of economics and finance, including ARMA models, regressions with time-varying coefficients and models with hidden factors, can be written in
one common form that has an exact and cheap recursive solution, the Kalman filter \ts{13}{12}{45:11}\ts{13}{12}{46:12}.

\subsection{The model}\label{ch09:sec-ss-model}
Let $\bm y_t\in\R^k$ be the observation at time $t$, $\bm s_t\in\R^m$ the (hidden) \emph{state}, $\bm\varepsilon_t\in\R^k$ the observation error and
$\bm\eta_t\in\R^n$ the state innovation. The \emph{linear Gaussian state-space model} is
\begin{equation}\label{ch09:eq-ss}
  \bm s_{t+1}=T_t\bm s_t+R_t\bm\eta_t,\quad \bm\eta_t\sim N(\bm0,Q_t);\qquad
  \bm y_t=Z_t\bm s_t+\bm\varepsilon_t,\quad \bm\varepsilon_t\sim N(\bm0,H_t),
\end{equation}
the \emph{state} (transition) equation and the \emph{observation} (measurement) equation, with i.i.d.\ noises independent of one another;
$T_t$ ($m\times m$) is the transition matrix, $R_t$ ($m\times n$, often columns of the identity) selects which components are shocked, $Z_t$ ($k\times m$) the observation
matrix, $Q_t,H_t$ positive definite \ts{13}{12}{47:18}\ts{13}{12}{48:22}. Jointly,
$\bigl(\bm s_{t+1};\bm y_t\bigr)=\bigl(\begin{smallmatrix}T_t\\ Z_t\end{smallmatrix}\bigr)\bm s_t+\bm u_t$ with
$\bm u_t=(R_t\bm\eta_t;\bm\varepsilon_t)\sim N(\bm0,\Omega_t)$, $\Omega_t=\mathrm{blockdiag}(R_tQ_tR_t\T,\,H_t)$ \ts{13}{12}{49:25}\ts{13}{12}{50:28}. Often the
model is time-invariant. The state is not unique: replacing $\bm s_t$ by $M\bm s_t$ ($M$ invertible) and $T_t$ by $MT_tM^{-1}$, $Z_t$ by $Z_tM^{-1}$ gives an equivalent model
\ts{13}{12}{1:02:59}.

\subsection{Examples}\label{ch09:sec-ss-ex}
\paragraph{Time-varying CAPM.} Extend the CAPM regression (\cref{ch05:capm}) so that alpha and beta drift as Gaussian random walks:
\[
  r_t=\alpha_t+\beta_tr_{m,t}+\varepsilon_t,\quad \varepsilon_t\sim N(0,\sigma_\varepsilon^2);\qquad
  \alpha_{t+1}=\alpha_t+\xi_t,\quad\beta_{t+1}=\beta_t+\zeta_t,
\]
with $\xi_t\sim N(0,\sigma_\xi^2)$, $\zeta_t\sim N(0,\sigma_\zeta^2)$ and $\varepsilon,\xi,\zeta$ mutually independent, $r_t$ the excess return of the asset and $r_{m,t}$ that of the market
\ts{13}{12}{48:22}\ts{13}{12}{51:33}. In state-space form $\bm s_t=(\alpha_t,\beta_t)\T$, $T_t=R_t=I_2$, $Q=\mathrm{diag}(\sigma_\xi^2,\sigma_\zeta^2)$,
$Z_t=(1,\ r_{m,t})$ (a known $1\times2$ row), $H=\sigma_\varepsilon^2$ \ts{13}{12}{52:35}\ts{13}{12}{53:36}.

\paragraph{Regression with time-varying coefficients.} $y_t=\bm x_t\T\bm\beta_t+\varepsilon_t$ with $\beta_{j,t+1}=\beta_{j,t}+\eta_{j,t}$, $\eta_{j,t}\sim N(0,\sigma_j^2)$, is
$\bm s_t=\bm\beta_t$, $T_t=R_t=I_p$, $Z_t=\bm x_t\T$, $Q=\mathrm{diag}(\sigma_1^2,\dots,\sigma_p^2)$, $H=\sigma^2$ \ts{13}{12}{54:37}. If all $\sigma_j^2=0$ the coefficients are constant
and one obtains ordinary linear regression; applying the filter of the next subsection observation by observation is then \emph{recursive least squares}, updating the estimate as
data arrive \ts{13}{12}{55:37}.

\paragraph{AR($p$).} With $\bm s_t=(y_t,y_{t-1},\dots,y_{t-p+1})\T$,
$\bm s_{t+1}=F\bm s_t+\bm e_1\eta_t$, where $F$ is the companion matrix of \cref{ch09:rem-companion}, $\bm e_1=(1,0,\dots,0)\T$, and $y_t=\bm e_1\T\bm s_t$ (no measurement
error, $\varepsilon_t=0$) \ts{13}{12}{56:43}\ts{13}{12}{57:46}\ts{13}{12}{58:46}.

\paragraph{ARMA($p,q$) (Harvey's representation).} Older estimation methods for moving-average models were hard to implement; the state-space form made exact maximum likelihood routine
\ts{13}{12}{59:48}\ts{13}{12}{1:00:55}. Put $m=\max(p,q+1)$, $\phi_j=0$ for $j>p$, $\theta_j=0$ for $j>q$, $\theta_0=1$; then $\{y_t\}$ is ARMA($m,m-1$). Define the state by $s_{1,t}=y_t$ and
\[
  s_{k,t}=\sum_{i=k}^{m}\phi_i\,y_{t+k-1-i}+\sum_{j=k-1}^{m-1}\theta_j\,\eta_{t+k-1-j},\qquad k=2,\dots,m .
\]
Rewriting the model equation in terms of $s_{1,t}$ and $s_{2,t}$ and repeating gives
\begin{equation}\label{ch09:eq-harvey}
  \bm s_{t+1}=T\bm s_t+\bm R\,\eta_{t+1},\quad y_t=\bm e_1\T\bm s_t,\qquad
  T=\begin{pmatrix}\phi_1&1&0&\cdots&0\\ \phi_2&0&1&&\vdots\\ \vdots&&&\ddots&0\\ \phi_{m-1}&0&\cdots&0&1\\ \phi_m&0&\cdots&0&0\end{pmatrix},\quad
  \bm R=\begin{pmatrix}1\\ \theta_1\\ \vdots\\ \theta_{m-1}\end{pmatrix}
\end{equation}
\ts{13}{12}{1:01:55}\ts{13}{12}{1:06:01}. \emph{Check for ARMA(1,1)} ($m=2$): the first row of $\bm s_{t+1}=T\bm s_t+\bm R\eta_{t+1}$ is
$y_{t+1}=\phi y_t+s_{2,t}+\eta_{t+1}$ and the second is $s_{2,t+1}=\theta\eta_{t+1}$, so $s_{2,t}=\theta\eta_t$ and $y_{t+1}=\phi y_t+\eta_{t+1}+\theta\eta_t$, as required. I also
verified numerically that for ARMA(2,1) with $\phi=(0.5,-0.3)$, $\theta=0.4$ the state-space variances ($\gamma(0)=1.893$, $\gamma(1)=1.036$ in units of $\sigma^2$, from the discrete Lyapunov equation
$P=TPT\T+\bm R\bm R\T$) agree with the Wold-weight sums. The MA($q$) is the case $p=0$. (The lecture notes give a different MA($q$) form, with the state $(\eta_{t-1},\dots,\eta_{t-q})$ and
$y_t=\sum_j\theta_j s_j+\eta_t$; there the measurement error and the state innovation are the \emph{same} variable, which violates the independence assumption of \eqref{ch09:eq-ss}, so I use the ARMA form above.)

\subsection{The Kalman filter}\label{ch09:sec-kalman}
Let $\mathcal F_t=\{\bm y_1,\dots,\bm y_t\}$. The \emph{Kalman filter} is the recursive computation of the conditional laws of the state and of the next observation given the information so far
\ts{13}{12}{1:07:03}\ts{13}{12}{1:08:10}. Because everything is Gaussian these are determined by means and covariances:
\begin{align*}
  \bm s_{t|t}&=\E[\bm s_t\mid\mathcal F_t], & \bm s_{t|t-1}&=\E[\bm s_t\mid\mathcal F_{t-1}], & \bm y_{t|t-1}&=\E[\bm y_t\mid\mathcal F_{t-1}],\\
  P_{t|t}&=\Cov(\bm s_t\mid\mathcal F_t), & P_{t|t-1}&=\Cov(\bm s_t\mid\mathcal F_{t-1}), & F_t&=\Cov(\bm y_t\mid\mathcal F_{t-1}),
\end{align*}
where the lecture writes $\Omega_s(t\mid t)$, $\Omega_s(t\mid t-1)$, $\Omega_y(t\mid t-1)$ \ts{13}{12}{1:09:12}. The \emph{innovation} (one-step forecast error) is
$\bm v_t=\bm y_t-\bm y_{t|t-1}=\bm y_t-Z_t\bm s_{t|t-1}$ \ts{13}{12}{1:10:12}. The filter alternates two steps.
\begin{enumerate}[1.]
  \item \textbf{Prediction} \ts{13}{12}{1:11:18}. From $(\bm s_{t|t},P_{t|t})$ predict the next state and observation:
        \begin{equation}\label{ch09:eq-kfpred}
          \begin{gathered}
          \bm s_{t+1|t}=T_t\bm s_{t|t},\qquad P_{t+1|t}=T_tP_{t|t}T_t\T+R_tQ_tR_t\T,\\
          \bm y_{t+1|t}=Z_{t+1}\bm s_{t+1|t},\qquad F_{t+1}=Z_{t+1}P_{t+1|t}Z_{t+1}\T+H_{t+1}.
          \end{gathered}
        \end{equation}
  \item \textbf{Correction (filtering)} \ts{13}{12}{1:12:21}. When $\bm y_t$ is observed, update:
        \begin{equation}\label{ch09:eq-kfupd}
          \bm s_{t|t}=\bm s_{t|t-1}+G_t\bm v_t,\qquad P_{t|t}=P_{t|t-1}-G_tF_tG_t\T,\qquad G_t=P_{t|t-1}Z_t\T F_t^{-1},
        \end{equation}
        $G_t$ being the \emph{filter gain}. The uncertainty on the state decreases by the information in the innovation (\ts{13}{12}{1:13:23}). (The slides write the gain with the index
        $(t-1\mid t)$; it must be $(t\mid t-1)$.)
\end{enumerate}
Two further operations reuse the same quantities. \emph{Forecasting} $h$ steps ahead iterates the prediction step without corrections \ts{13}{12}{1:14:24}. \emph{Smoothing} improves the estimate of a past
state using the whole sample $\mathcal F_T$, by the backward recursion
$\bm s_{t|T}=\bm s_{t|t}+S_t(\bm s_{t+1|T}-\bm s_{t+1|t})$ with $S_t=P_{t|t}T_t\T P_{t+1|t}^{-1}$ (and the analogous recursion for the covariance) \ts{13}{12}{1:14:24}.

\begin{proof}[Derivation of the correction step]
Conditionally on $\mathcal F_{t-1}$ the pair $(\bm s_t,\bm y_t)$ is jointly Gaussian with means $(\bm s_{t|t-1},\,Z_t\bm s_{t|t-1})$ and covariance blocks
$\Var\bm s_t=P_{t|t-1}$, $\Cov(\bm s_t,\bm y_t)=P_{t|t-1}Z_t\T$ (because $\bm\varepsilon_t$ is independent of $\bm s_t$), $\Var\bm y_t=F_t=Z_tP_{t|t-1}Z_t\T+H_t$. Conditioning a Gaussian on part of
its coordinates (\cref{ch03:sec-mvn}) gives the mean $\bm s_{t|t-1}+P_{t|t-1}Z_t\T F_t^{-1}(\bm y_t-Z_t\bm s_{t|t-1})$ and the covariance
$P_{t|t-1}-P_{t|t-1}Z_t\T F_t^{-1}Z_tP_{t|t-1}=P_{t|t-1}-G_tF_tG_t\T$, which are \eqref{ch09:eq-kfupd}. The prediction step is the linear transformation of a Gaussian by the
state equation.
\end{proof}

\paragraph{Likelihood.} By the chain rule $p(\bm y_1,\dots,\bm y_T;\theta)=p(\bm y_1;\theta)\prod_{t\ge2}p(\bm y_t\mid\mathcal F_{t-1};\theta)$, and
$\bm y_t\mid\mathcal F_{t-1}\sim N(\bm y_{t|t-1},F_t)$. So the filter delivers the exact log-likelihood of all parameters $\theta$ in $T_t,Z_t,R_t,Q_t,H_t$ as a by-product
(\emph{prediction-error decomposition}) \ts{13}{12}{1:15:25}:
\begin{equation}\label{ch09:eq-kfll}
  \ell(\theta)=-\frac{kT}2\log2\pi-\frac12\sum_{t=1}^T\log|F_t|-\frac12\sum_{t=1}^T\bm v_t\T F_t^{-1}\bm v_t .
\end{equation}
It is maximised by Newton-type methods or by the EM algorithm of Dempster, Laird and Rubin (\cref{ch07:sec-fa}). As $T\to\infty$ the MLE is consistent and asymptotically normal
with covariance the inverse Fisher information (2013 notes, last slide). The initial state is modelled as $N(\bm s_{1|0},P_{1|0})$ with pre-specified moments: a large variance (diffuse) if one knows nothing, or the
stationary distribution of the state if the model is stationary \ts{13}{12}{1:16:25}. The observations are jointly Gaussian and the recursions provide a full specification of the model: means and
covariances characterise the whole law.

\begin{example}[The local-level model: Kalman filter, EWMA and ARIMA(0,1,1)]\label{ch09:ex-locallevel}
The simplest model with a hidden state is $y_t=\mu_t+\varepsilon_t$, $\mu_{t+1}=\mu_t+\eta_t$: a slowly moving true level observed with noise. Here $T=Z=R=1$,
$Q=\Var\eta$, $H=\Var\varepsilon$. The recursions \eqref{ch09:eq-kfpred}--\eqref{ch09:eq-kfupd} become $P_{t+1|t}=P_{t|t}+Q$, $G_t=P_{t|t-1}/(P_{t|t-1}+H)$ and
$m_t=(1-G_t)\,m_{t-1}+G_t\,y_t$, where $m_t=\mu_{t|t}$. At the steady state $P=P_{t|t-1}$ solves $P=(1-K)P+Q$ with $K=P/(P+H)$, that is $P^2-QP-QH=0$:
\[
  P=\frac{Q+\sqrt{Q^2+4QH}}2,\qquad K=\frac P{P+H},\qquad P_{t|t}=(1-K)P .
\]
For $Q=1$, $H=4$: $P=(1+\sqrt{17})/2=2.562$, $K=0.390$, $P_{t|t}=1.562$, while the raw observation has error variance $H=4$
(a Monte Carlo over $2\,000$ paths gives a filter mean-squared error of $1.558$). The filter is an \emph{exponentially weighted moving average} (EWMA) of the observations with
weight $K$ and decay factor $1-K=0.610$. The same object appears in ARIMA language: $\Delta y_t=\eta_{t-1}+\varepsilon_t-\varepsilon_{t-1}$ has $\gamma(0)=Q+2H$, $\gamma(1)=-H$, so
$\rho(1)=-H/(Q+2H)=-4/9$, an MA(1) whose invertible parameter is $\theta=-0.6096$ (from $\theta/(1+\theta^2)=\rho(1)$), equal to $-(1-K)$. Thus the local-level model \emph{is} ARIMA(0,1,1),
and the steady-state Kalman filter is the exponentially weighted forecast (\cref{ch09:sec-arma11,ch09:sec-nonstat}). The same recursion applied to squared returns is the RiskMetrics
volatility estimator, a GARCH(1,1) with $\alpha_0=0$ and $\alpha_1+\beta_1=1$ (\cref{ch:volatility}) \ts{13}{11}{27:33}. \Cref{ch09:fig-kf} shows a simulated path.
\end{example}

\begin{figure}[ht]
\centering
\begin{tikzpicture}
\begin{axis}[width=0.95\linewidth,height=5.4cm,xlabel={$t$},xmin=0,xmax=59,grid=major,grid style={black!10},legend style={font=\footnotesize,at={(0.5,1.03)},anchor=south,legend columns=3}]
  \addplot[only marks,mark=*,mark size=1.2pt,color=black!45] coordinates {(0,-0.47) (1,2.96) (2,1.74) (3,2.07) (4,2.5) (5,-0.47) (6,3.05) (7,1.59) (8,3.53) (9,-0.35) (10,4.44) (11,3.38) (12,2.49) (13,3.97) (14,4.85) (15,6.38) (16,0.47) (17,5.18) (18,4.19) (19,3.76) (20,1.07) (21,4.08) (22,-0.56) (23,2.42) (24,1.97) (25,-4.64) (26,-2.52) (27,-5.72) (28,-4.11) (29,-1.53) (30,0.38) (31,-7.45) (32,-5.79) (33,-2.85) (34,-0.38) (35,-3.0) (36,-4.79) (37,-1.66) (38,-2.49) (39,-3.68) (40,-4.39) (41,-1.9) (42,-0.96) (43,-3.05) (44,-3.97) (45,-6.93) (46,-7.07) (47,-3.56) (48,-5.82) (49,-9.06) (50,-2.13) (51,-2.91) (52,0.87) (53,-2.82) (54,-2.91) (55,-6.22) (56,-0.06) (57,3.03) (58,-5.51) (59,-4.82)};
  \addplot[thick,color=defcol] coordinates {(0,0.03) (1,1.39) (2,2.62) (3,2.11) (4,1.81) (5,1.28) (6,1.85) (7,1.8) (8,2.54) (9,0.7) (10,2.26) (11,2.17) (12,2.85) (13,2.71) (14,2.33) (15,2.79) (16,3.62) (17,3.42) (18,3.26) (19,3.95) (20,3.08) (21,1.56) (22,1.96) (23,1.29) (24,-0.63) (25,-1.45) (26,-1.91) (27,-3.11) (28,-4.6) (29,-4.56) (30,-3.67) (31,-3.9) (32,-4.64) (33,-4.26) (34,-3.54) (35,-3.84) (36,-3.29) (37,-2.25) (38,-2.46) (39,-3.27) (40,-2.92) (41,-2.68) (42,-1.58) (43,-2.86) (44,-3.52) (45,-4.36) (46,-6.1) (47,-5.97) (48,-5.44) (49,-6.18) (50,-4.8) (51,-3.97) (52,-3.35) (53,-2.94) (54,-1.99) (55,-3.32) (56,-2.71) (57,-2.1) (58,-3.87) (59,-3.53)};
  \addplot[thick,color=thmcol] coordinates {(0,-0.33) (1,1.28) (2,1.48) (3,1.72) (4,2.02) (5,1.05) (6,1.83) (7,1.73) (8,2.44) (9,1.35) (10,2.55) (11,2.88) (12,2.72) (13,3.21) (14,3.85) (15,4.84) (16,3.13) (17,3.93) (18,4.03) (19,3.93) (20,2.81) (21,3.31) (22,1.79) (23,2.04) (24,2.01) (25,-0.59) (26,-1.34) (27,-3.05) (28,-3.47) (29,-2.71) (30,-1.5) (31,-3.83) (32,-4.59) (33,-3.91) (34,-2.53) (35,-2.71) (36,-3.52) (37,-2.8) (38,-2.68) (39,-3.07) (40,-3.59) (41,-2.93) (42,-2.16) (43,-2.51) (44,-3.08) (45,-4.58) (46,-5.55) (47,-4.77) (48,-5.18) (49,-6.7) (50,-4.91) (51,-4.13) (52,-2.18) (53,-2.43) (54,-2.62) (55,-4.02) (56,-2.48) (57,-0.32) (58,-2.35) (59,-3.31)};
  \legend{observations $y_t$,true level $\mu_t$,filtered $\mu_{t|t}$}
\end{axis}
\end{tikzpicture}
\caption{Local-level model with $Q=1$, $H=4$ (seed~11). The filter (red) tracks the hidden level (blue) much more closely than the raw observations (dots).}\label{ch09:fig-kf}
\end{figure}

\begin{finance}[Where state-space models appear]
Time-varying betas and alphas (the CAPM example above) and, in general, any regression whose coefficients are believed to drift; dynamic factor models, where the factors of
\cref{ch:pca} follow their own time-series model; a ``true'' price observed through noisy quotes, as in the USD/RUB forward-point example of \cref{ch:markets}, where five
brokers quote at random and a Kalman filter estimates the underlying price \ts{13}{1}{57:41}; and exact maximum-likelihood estimation of ARMA models
(\cref{ch09:sec-est}). The filter is cheap: each step costs a few matrix products, and it delivers the likelihood as well as the forecast.
\end{finance}

% ==================================================================
\section{Notes on the sources}\label{ch09:sec-notes}

\begin{coursenote}[2013 versus 2024, and errata]
\begin{itemize}
  \item \textbf{Overlap.} 2024 L12 and 2013 L8 use the same slides for univariate theory (\file{lec11-12} vs.\ \file{lecnote8}); the 2013 lecture goes further (trend reversion, ARIMA,
        estimation, model selection, with the information-criteria discussion at the end). Multivariate series, cointegration, state space and the Kalman filter are 2013-only.
  \item \textbf{Slide misprints.}
  \begin{itemize}
    \item Yule--Walker slide: the second lag is written $X_{t-1}$ instead of $X_{t-2}$ (corrected aloud, \ts{24}{12}{1:12:16}), and the sum in the estimator of $\sigma^2$
        starts at $j-1$ instead of $j=1$.
    \item 2013 L11, model selection for VAR: the criteria are printed with $-\log|\tilde\Sigma(p)|$; it should be $+\log|\tilde\Sigma(p)|$ (\cref{ch09:eq-icvar}).
    \item 2013 L12, Kalman filter: the gain $G_t$ is written with covariances at $(t-1\mid t)$; it should be $(t\mid t-1)$.
    \item The notes divide $\tilde\Sigma$ for the VAR by $n-pm$; the unbiased
        divisor is $n-mp-1$ when an intercept is estimated.
  \end{itemize}
  \item \textbf{Case Study 3.} The printed ``smallest root modulus $1.4027$'' is the smallest \emph{squared} modulus (\cref{ch09:cs-3b}). The AR(2) cycle of ``just over four months'' is
        reproduced ($4.107$).
  \item \textbf{Conventions.} $\mathrm{Var}(\hat r_k)\approx1/(T-k)$ (2024) versus $1/T$ (Bartlett): equivalent for $T\gg k$. R's \texttt{vars} reports moduli of companion eigenvalues, whereas
        \texttt{polyroot} returns roots of $\phi(z)$ (\cref{ch09:rem-companion}). ``Log returns'' of yields is a loose name for relative changes of the yield.
  \item \textbf{Not in the download.} The plots \file{TimeSeries4plots.pdf}, \file{TimeSeries4acfplots.pdf} and the crack-spread cointegration note; the numbers of \cref{ch09:sec-crack} are as spoken.
\end{itemize}
\end{coursenote}

\begin{intuition}[How the pieces fit together]
The chapter is one idea seen from four sides. Stationarity plus the Wold theorem say that a stationary series is a linear filter of white noise. ARMA models parametrise that filter with few
parameters; the AR part is a regression on the past (so least squares of \cref{ch:regression} applies) and the roots of its polynomial decide stability. Non-stationary series are
brought to that setting by differencing, and unit-root tests decide when. In several dimensions the polynomial becomes a matrix (VAR), stability is a condition on the eigenvalues of
the companion matrix, and cointegration is the case where a rank-deficient matrix $\Pi$ ties non-stationary series together. State-space models recast all of it as a hidden Markov-type
Gaussian system, in which the Kalman filter is the exact recursive regression on the past. Forward links: volatility (GARCH is ARMA for squared returns, \cref{ch:volatility}), portfolio
theory (\cref{ch:portfolio}), continuous time (the AR(1) is the Ornstein--Uhlenbeck process, \cref{ch:sde}; the Dickey--Fuller limit uses Brownian motion, \cref{ch:sp2}).
\end{intuition}

% ==================================================================
\begin{keyformulas}
\begin{align*}
&\text{Stationarity:} && \E X_t=\mu,\quad \gamma(\tau)=\Cov(X_t,X_{t+\tau}),\quad\rho=\gamma/\gamma(0)\\
&\text{White noise test:} && \mathrm{BP}=T\textstyle\sum_{j\le K}\hat r_j^2\sim\chi^2_K,\quad |\hat r_k|>1.96/\sqrt T\\
&\text{Wold:} && X_t=V_t+\psi(L)\eta_t,\quad \psi_0=1,\ \textstyle\sum\psi_j^2<\infty\\
&\text{ARMA:} && \phi(L)(X_t-\mu)=\theta(L)\eta_t,\quad \psi(L)=\theta(L)/\phi(L)\\
&\text{Stationary if:} && \text{roots of }\phi(z)\text{ outside }|z|\le1\\
& && (\Leftrightarrow\ \text{eigenvalues of companion matrix }|\lambda|<1)\\
&\text{AR(1):} && \gamma(0)=\sigma^2/(1-\phi^2),\quad \rho(j)=\phi^j\\
&\text{Yule--Walker:} && \gamma(j)=\textstyle\sum_{i\le p}\phi_i\gamma(j-i)+\sigma^2\delta_{j0}\\
&\text{MA($q$):} && \gamma(0)=\sigma^2\textstyle\sum_{i\le q}\theta_i^2,\quad \gamma(j)=0\ (j>q)\\
&\text{ARMA(1,1):} && \gamma(0)=\sigma^2\tfrac{1+2\phi\theta+\theta^2}{1-\phi^2},\quad \rho(k)=\phi\rho(k-1)\ (k\ge2)\\
&\text{Differencing:} && \Delta=1-L;\ \text{ARIMA}(p,d,q):\ \Delta^dX\sim\text{ARMA}(p,q)\\
&\text{Dickey--Fuller:} && T(\hat\phi-1)\to\tfrac{\frac12(W(1)^2-1)}{\int_0^1W^2\dd r}\ \text{under }\phi=1\\
&\text{Information criteria:} && \log\tilde\sigma^2+c_n\tfrac{p+q}n,\quad c_n=2,\ \log n,\ 2\log\log n\\
&\text{VAR($p$):} && \bm X_t=\bm C+\textstyle\sum_j\Phi_j\bm X_{t-j}+\bm\eta_t,\quad \bm\mu=(I-\textstyle\sum\Phi_j)^{-1}\bm C\\
&\text{VAR(1) moments:} && \Gamma_0=A\Gamma_0A\T+\Sigma,\quad \Gamma_h=A^h\Gamma_0\\
&\text{Estimation:} && \hat B=(Z\T Z)^{-1}Z\T Y,\quad \tilde\Sigma=\hat E\T\hat E/n\\
&\text{VECM:} && \Delta\bm X_t=\bm C+\Pi\bm X_{t-1}+\textstyle\sum_k\Gamma_k\Delta\bm X_{t-k}+\bm\eta_t\\
& && \Pi=\textstyle\sum_j\Phi_j-I=\alpha\beta\T,\quad \mathrm{rank}\,\Pi=r\\
&\text{State space:} && \bm s_{t+1}=T\bm s_t+R\bm\eta_t,\quad \bm y_t=Z\bm s_t+\bm\varepsilon_t\\
&\text{Kalman:} && \bm s_{t|t}=\bm s_{t|t-1}+G_t\bm v_t,\quad G_t=P_{t|t-1}Z\T F_t^{-1}\\
& && P_{t|t}=P_{t|t-1}-G_tF_tG_t\T,\quad P_{t+1|t}=TP_{t|t}T\T+RQR\T
\end{align*}
\end{keyformulas}

% ==================================================================
\section*{Exercises}
\addcontentsline{toc}{section}{Exercises}

\subsection*{From 2013 Problem Set 4 (and 2024 Problem Set 4, Problems 5--6)}
The 2013 PS4 has six problems on univariate models; 2024 PS4 Problem~5 is 2013 Problem~1(a)--(e), and 2024 Problem~6 is 2013 Problem~6 plus a part (e).

\begin{exercise}[PS4 (2013), Problem 1; PS4 (2024), Problem 5: covariance-stationary AR(2)]\label{ch09:ex-ps4-1}
Let $X_t=\phi_0+\phi_1X_{t-1}+\phi_2X_{t-2}+\eta_t$, $\eta_t\sim\mathrm{WN}(0,\sigma^2)$, be covariance stationary.
\begin{enumerate}[(a)]
  \item Show that $\mu=\E X_t=\phi_0/(1-\phi_1-\phi_2)$.
  \item Show that $\gamma(k)=\phi_1\gamma(k-1)+\phi_2\gamma(k-2)$, $k\ge1$.
  \item Show that $\rho_k=\phi_1\rho_{k-1}+\phi_2\rho_{k-2}$.
  \item The Yule--Walker equations for $k=1,2$ are $\rho_1=\phi_1+\phi_2\rho_1$ and $\rho_2=\phi_1\rho_1+\phi_2$ (using $\rho_0=1$, $\rho_{-k}=\rho_k$): solve for $\rho_1,\rho_2$ in terms of $\phi_1,\phi_2$.
  \item Solve for $\phi_1,\phi_2$ in terms of $\rho_1,\rho_2$.
  \item (2013 only) Derive complete formulas for $\rho_k$, $k>2$, using Problem~4(b) below.
\end{enumerate}
\end{exercise}

\begin{exercise}[PS4 (2013), Problem 2: difference equations of an AR($p$)]\label{ch09:ex-ps4-2}
Let $\phi(L)X_t=\phi_0+\eta_t$ be an AR($p$) with $\phi(L)=1-\phi_1L-\dots-\phi_pL^p$, and consider the homogeneous equation $\phi(L)g(t)=0$.
\begin{enumerate}[(a)]
  \item Show $\mu=\phi_0/(1-\sum_i\phi_i)$, so $\mu=0$ iff $\phi_0=0$
      (``de-meaning'' removes the constant).
  \item For $\gamma(t)=\Cov(X_t,X_0)$ prove $\phi(L)\gamma(t)=0$ for all $t\ge1$.
  \item Same for $\rho(t)=\Corr(X_t,X_0)$.
  \item Given $x_0,\dots,x_t$, let $\hat x_t(h)=\E[X_{t+h}\mid x_t,\dots,x_0]$ for $h>0$
      and $\hat x_t(h)=x_{t+h}$ for $h\le0$; show $\hat x_t(h)=\phi_0+\sum_i\phi_i\hat x_t(h-i)$ for every $h>0$.
  \item For the de-meaned process show that the forecast function satisfies $\phi(L)\hat x_t(h)=0$ in $h$.
\end{enumerate}
\end{exercise}

\begin{exercise}[PS4 (2013), Problem 3: solutions of the difference equations]\label{ch09:ex-ps4-3}
Consider $\phi(L)g(t)=0$, $t=0,1,2,\dots$
\begin{enumerate}[(a)]
  \item Show that $g(t)=Ce^{-\lambda t}$ solves it if $z=e^{\lambda}$ is a root of $\phi(z)=1-\sum_i\phi_iz^i=0$.
  \item Show that $g(t)=CG^t$ solves it if $G^{-1}$ is a root.
  \item If the $p$ roots
      $z_i=G_i^{-1}$ are distinct, show that $g(t)=\sum_iC_iG_i^t$ is a solution.
  \item If all roots are outside the unit circle, show every such solution is bounded; since the autocovariance is a solution, this is the necessary and sufficient
      condition for covariance stationarity.
\end{enumerate}
\end{exercise}

\begin{exercise}[PS4 (2013), Problem 4: the AR(2) process]\label{ch09:ex-ar2}
\leavevmode
\begin{enumerate}[(a)]
  \item Solve for the two roots $z_1,z_2$ of the characteristic equation and give the conditions under which they are distinct and real, coincident and real, or complex conjugates.
  \item Prove $\rho_k=C_1G_1^k+C_2G_2^k$
      with $G_i=z_i^{-1}$; solve for $C_1,C_2$ when the roots are distinct and show that $\rho_k=\dfrac{G_1(1-G_2^2)G_1^k-G_2(1-G_1^2)G_2^k}{(G_1-G_2)(1+G_1G_2)}$.
  \item For distinct real roots show that the autocovariance decreases geometrically in magnitude.
  \item Under what conditions on $\phi_1$ is it always positive, and when does it alternate in sign?
  \item \emph{Optional.} For complex roots write $G_{1,2}=de^{\pm\mathrm i2\pi f_0}$ and show
      $\rho_k=[\mathrm{sgn}(\phi_1)]^k|d|^k\sin(2\pi f_0k+F)/\sin F$ with $|d|=\sqrt{-\phi_2}$, $2\pi f_0=\arccos\bigl(|\phi_1|/(2\sqrt{-\phi_2})\bigr)$ and $\tan F=\frac{1+d^2}{1-d^2}\tan(2\pi f_0)$.
\end{enumerate}
\end{exercise}

\begin{exercise}[PS4 (2013), Problem 5: the MA(1) process]\label{ch09:ex-ps4-5}
$X_t=\eta_t+\theta\eta_{t-1}$.
\begin{enumerate}[(a)]
  \item Derive the ACF.
  \item Compute the first-order autocorrelation for $\theta=0.5$ and $\theta=2.0$.
  \item Solve the formula for $\rho_1$ for $\theta$ in terms of $\rho_1$; is the solution unique?
  \item For each of the two models decide whether the process is invertible and, if so, write it as an infinite-order autoregression.
\end{enumerate}
\end{exercise}

\begin{exercise}[PS4 (2013), Problem 6; PS4 (2024), Problem 6: ARMA(1,1)]\label{ch09:ex-arma11}
For the covariance-stationary $(1-\phi L)X_t=\phi_0+(1+\theta L)\eta_t$:
\begin{enumerate}[(a)]
  \item Prove $\mu=\phi_0/(1-\phi)$.
  \item Prove $\Var X_t=\gamma(0)=\sigma^2\bigl[1+\frac{(\phi+\theta)^2}{1-\phi^2}\bigr]$
      (equivalently $\sigma^2(1+2\phi\theta+\theta^2)/(1-\phi^2)$).
  \item Prove that $\rho_1=\phi+\theta\sigma^2/\gamma(0)$ and $\rho_k=\phi\rho_{k-1}$ for $k>1$.
  \item Compare the ACF with that of the AR(1) with the same $\phi_0,\phi$: which pattern of the ACF is impossible for an AR(1)?
      What behaviour would an economic index following such a process show?
  \item (2024 only) Compare with the ACF of an MA(1) with the same parameters; what is the pattern of an MA(1) and how does it change for MA($q$), $q>1$?
      Try to obtain (b)--(c) directly from the covariance recursion rather than via the Wold weights used in \cref{ch09:sec-arma11}.
\end{enumerate}
\end{exercise}

\subsection*{From 2013 Problem Set 6 (Problems 1--2)}
\begin{exercise}[PS6 (2013), Problem 1: moments of a VAR(1)]\label{ch09:ex-ps6-1}
Let $\bm X_t=(X_{1,t},X_{2,t})\T$ follow $X_{1,t}=0.3+0.8X_{1,t-1}+\eta_{1,t}$, $X_{2,t}=0.2+0.6X_{1,t-1}+0.4X_{2,t-1}+\eta_{2,t}$, with $(\eta_{1,t},\eta_{2,t})\T$ i.i.d.\ $N(\bm0,\Sigma)$, $\Sigma=\bigl(\begin{smallmatrix}3&-1\\-1&3\end{smallmatrix}\bigr)$.
Compute
\begin{enumerate}[(a)]
  \item $\bm\mu=\E\bm X_t$,
  \item $\Gamma(0)=\Cov\bm X_t$,
  \item $\Gamma(1)=\Cov(\bm X_t,\bm X_{t-1})$, and
  \item derive a formula for $\Gamma(h)$, $h\ge1$. (See \cref{ch09:ex-var1} for the method.)
\end{enumerate}
\end{exercise}

\begin{exercise}[PS6 (2013), Problem 2: an AR(2) as a difference of two MA processes]\label{ch09:ex-ps6-2}
With $\{\eta_t\}$ i.i.d.\ $\mathrm{WN}(0,\sigma^2)$ let $w_t=5(1-0.5L)^{-1}\eta_t$, $v_t=4(1-0.4L)^{-1}\eta_t$ and $x_t=w_t-v_t$.
\begin{enumerate}[(a)]
  \item Solve for the coefficients $\psi_i$ in $x_t=\eta_t+\sum_{i\ge1}\psi_i\eta_{t-i}$.
  \item Prove $\{x_t\}$ is an AR(2) process.
  \item Solve for $\phi_1,\phi_2$ in $x_t=\phi_1x_{t-1}+\phi_2x_{t-2}+\eta_t$.
  \item Prove that any stationary AR(2) process can be written as the difference of two (possibly infinite-order) moving-average processes driven by the same $\{\eta_t\}$.
\end{enumerate}
\end{exercise}

\subsection*{Extra exercises}
\begin{exercise}[Box--Pierce by hand (not from the course)]\label{ch09:ex-x-bp}
A daily return series with $T=500$ has $\hat r_1=0.09$, $\hat r_2=-0.05$, $\hat r_3=0.04$, $\hat r_4=0.02$, $\hat r_5=-0.06$.
\begin{enumerate}[(a)]
  \item Which lags are individually significant at $5\%$?
  \item Compute BP for $K=5$ and compare with $\chi^2_5$ (critical value $11.07$).
  \item Explain how the two answers can disagree and which you would report.
  \item If the numbers came from the residuals of an ARMA(1,1) fit, what would the degrees of freedom be?
\end{enumerate}
\end{exercise}

\begin{exercise}[Stationarity of a VAR (not from the course)]\label{ch09:ex-x-companion}
\leavevmode
\begin{enumerate}[(a)]
  \item Write the companion matrix of the univariate AR(2) $X_t=0.5X_{t-1}+0.3X_{t-2}+\eta_t$, find its eigenvalues ($0.852,-0.352$) and the roots of $\phi(z)$; verify $\lambda=1/z$.
  \item For the VAR(1) with $A=\bigl(\begin{smallmatrix}0.9&0.4\\0&0.5\end{smallmatrix}\bigr)$ compute $\Gamma_0$ from $\mathrm{vec}\,\Gamma_0=(I-A\otimes A)^{-1}\mathrm{vec}\,\Sigma$ for $\Sigma=I$ and the impulse responses $\Psi_k$;
      which shock has the longest-lasting effect on which variable?
\end{enumerate}
\end{exercise}

\begin{exercise}[Dickey--Fuller by simulation (not from the course)]\label{ch09:ex-x-df}
\leavevmode
\begin{enumerate}[(a)]
  \item Simulate $10^4$ random walks of length $500$ and compute $t_{\phi=1}$ from the regression of $X_t$ on $X_{t-1}$ with and without intercept; check the $5\%$ quantiles $-1.95$ and $-2.89$.
  \item Repeat with an AR(1) with $\phi=0.95$ and estimate the power of the test at the $5\%$ level; explain why
      non-rejection on a persistent yield series is weak evidence for a unit root.
  \item Verify the identity $\sum_{t=1}^TX_{t-1}\eta_t=\frac12(X_T^2-\sum\eta_t^2)$ and show that its expectation is zero; then estimate $\E[T(\hat\phi-1)]$ by simulation and explain why the ratio in \eqref{ch09:eq-df} is nevertheless biased downwards (the denominator is positively correlated with the numerator's negative part).
\end{enumerate}
\end{exercise}

\begin{exercise}[Cointegration and the VECM (not from the course)]\label{ch09:ex-x-coint}
For the pair of \cref{ch09:ex-coint}
\begin{enumerate}[(a)]
  \item rewrite the system as a VAR(1) in levels and verify that $I+\Pi$ has eigenvalues $1$ and $0.7$.
  \item Show that $\Pi$ is not invertible and identify $\alpha$ and $\beta$; check that another factorisation $(\alpha G)(\beta G^{-\mathsf T})\T$ gives the same $\Pi$.
  \item Simulate the system, regress $y_t$ on $x_t$ and apply an ADF test to the residuals (Engle--Granger); do the same for two independent random walks and comment on ``spurious regression''.
  \item What changes in the VECM if both $\varepsilon_{1t}$ and $u_{t-1}$ enter the equation of $x_t$?
\end{enumerate}
\end{exercise}

\begin{exercise}[Kalman filter by hand (not from the course)]\label{ch09:ex-x-kf}
Take the local-level model with $Q=1$, $H=4$ and initial state $N(0,10)$.
\begin{enumerate}[(a)]
  \item Run the filter for the observations $y_1=1.5$, $y_2=0.8$, $y_3=2.2$: compute $G_t$, $\mu_{t|t}$, $P_{t|t}$.
  \item Compute the two-step forecast of $y_5$ and its variance.
  \item Show that $P_{t|t-1}$ converges monotonically to
      $(Q+\sqrt{Q^2+4QH})/2$.
  \item Show that $G_t\to0$ as $Q\to0$ and interpret it (the level is constant and the filter reduces to the sample mean).
\end{enumerate}
\end{exercise}

\begin{exercise}[Time-varying beta (not from the course)]\label{ch09:ex-x-beta}
In the state-space model of the time-varying CAPM, write out the Kalman recursions for $\bm s_t=(\alpha_t,\beta_t)\T$.
\begin{enumerate}[(a)]
  \item Show that if $\sigma_\xi=\sigma_\zeta=0$ the recursion reproduces recursive least squares for the simple regression.
  \item Explain how the ratio
      $\sigma_\zeta^2/\sigma_\varepsilon^2$ sets the effective window length of the estimate, and relate it to the local-level example.
  \item With daily data for one stock and the S\&P~500, estimate $\sigma_\zeta^2$ by maximising \eqref{ch09:eq-kfll}.
\end{enumerate}
\end{exercise}

% !TEX root = ../main.tex
\chapter{Portfolio Theory and Portfolio Management}\label{ch:portfolio}
\begin{paracol}{2}

\begin{sources}
\begin{tabularx}{\linewidth}{@{}l l L@{}}
\toprule
\textbf{Lecture} & \textbf{Speaker} & \textbf{Content}\\
\midrule
2013 L14 & P.~Kempthorne & Markowitz mean--variance analysis (two assets, then $m$ assets with Lagrange multipliers); the risk-free asset, the market portfolio, Tobin's separation theorem and the capital market line; von Neumann--Morgenstern utility; portfolio constraints (long-only, holding, turnover, benchmark, tracking error, factor, integer); estimation of means and covariances; alternative risk measures (MAD, semi-variance, VaR, CVaR, coherent measures, skewness and kurtosis). Live demonstration of Case Study~6 (simulation and sector ETFs).\\
2013 L16 & J.~Xia & Blackboard lecture: goals of portfolio management, the return--volatility map of asset classes, portfolio theory through special cases of two assets, Sharpe ratio, beta, the Kelly formula, the 60/40 portfolio, risk parity and leverage, the ``free lunch'' of rebalancing, crowd synchronisation (Millennium Bridge, metronomes) and the limits of the theory. Long question session.\\
2024 L13 & J.~Xia & The same programme in a slide version with new material from the speaker's own research: the endowment model, the loss of meaning of volatility as risk, the gain--loss ratio $(G-L)/(G+L)$ for sizing, a feedback model of crowd behaviour and the origin of power laws. Questions on taxes, expected versus worst loss, stop losses.\\
\bottomrule
\end{tabularx}

\smallskip
\textbf{Files.} 2024 slides \file{mit18_642_f24_lec13.pdf} (38 pages). 2013 slides \file{MIT18_S096F13_lecnote14.pdf}
(34 pages) and \file{lecnote16.pdf} (2 pages: only the outline, the lecture was given on the blackboard);
2013 Case Study~6 \file{CaseStudy6.pdf} (portfolio theory with US sector ETFs, 14 pages); the R output
\file{Smltn_TwoAst.pdf} (simulation of two-asset portfolios, 6 pages) and the four sector-ETF output files
\file{ETF_pridA_15.pdf}, \file{ETF_pridA_30.pdf} (2009--2013, maximum allocation $15\%$ and $30\%$),
\file{ETF_pridB_15.pdf}, \file{ETF_pridB_30.pdf} (2003--2006); the names ``periodA'' and ``periodB'' are those of
the case study text.
\textbf{Problem sets.} 2013 PS6 Problems~3--5 (two-asset and $m$-asset minimum-variance portfolios,
equicorrelation); Problems~1--2 of the same set are time series (\cref{ch:timeseries}). The exercises are at the
end of the chapter.
\textbf{Not covered or not available.} The R code that produced the case-study output is not in the download (only
its printed output), so the numbers of \cref{ch10:sec-case} are quoted, and checked by an independent computation
from the printed input tables. The video shown by the speaker in both years and the 2013 blackboard are not in the
notes; \file{lecnote16} contains only the outline of the lecture. The second half of 2013 L14 mentions ``sector
market-neutral models'' with no data (\cref{ch10:sec-case}).
\end{sources}

\section*{Motivation}
Every investor faces the same question: given a list of things one could hold, how much of each? The 2013 and 2024
lectures answer it twice. Kempthorne (2013 L14) gives the mathematical theory that earned Harry Markowitz a Nobel
prize: describe a portfolio by its expected return and its variance, and find the portfolios that offer the
largest return for a given variance. The theory is a piece of applied linear algebra (\cref{ch:linalg}): a
quadratic form to be minimised under linear constraints, solved with Lagrange multipliers, whose solution is
governed by the inverse covariance matrix $\Sigma^{-1}$. Xia (2013 L16, 2024 L13), who has managed money for
a bank and for the Harvard endowment, then asks what survives in practice: the inputs are estimated with large
errors, volatility is a poor measure of risk, and markets are populated by agents who react to each other.

The chapter follows this arc.
\begin{itemize}
  \item \cref{ch10:sec-problem}: the objectives, constraints and horizon that define a portfolio problem.
  \item \cref{ch10:sec-two,ch10:sec-mvo,ch10:sec-rf}: Markowitz theory for two assets, for $m$ assets, and with a
        riskless asset (efficient frontier, tangency portfolio, capital market line), with complete derivations;
        \cref{ch10:sec-capm} links it to beta and the CAPM (\cref{ch03:thm-capm}).
  \item \cref{ch10:sec-rebal,ch10:sec-rp}: diversification, rebalancing, leverage and risk parity.
  \item \cref{ch10:sec-utility,ch10:sec-kelly}: utility theory as the foundation of mean--variance analysis, and
        sizing by the Kelly criterion and by Xia's gain--loss ratio.
  \item \cref{ch10:sec-constr,ch10:sec-est,ch10:sec-risk}: constraints, estimation error, and other risk measures.
  \item \cref{ch10:sec-crowd}: the ``beyond Markowitz'' ideas of the 2024 lecture (crowds, power laws).
  \item \cref{ch10:sec-case}: the 2013 case study and the notes on sources and errata.
\end{itemize}
\emph{Notation.} The slides call the vector of expected returns $\bm\alpha$; here it is $\bm\mu$, to keep
$\alpha$ for the regression intercept of \cref{def:alphabeta}. Vectors are columns, $\bm 1=(1,\dots,1)\T$,
$\Sigma$ is the (positive definite) covariance matrix of the returns $\bm R$ of $m$ risky assets and
$\sigma_i^2=\Sigma_{ii}$.

% ==================================================================
\section{The portfolio construction problem}\label{ch10:sec-problem}

Both Xia lectures start with the same exercise: each student writes on a sheet the composition of a
hypothetical portfolio (\$10{,}000 in 2024; in 2013 anything from a \$1{,}000 saving to a fund's start-up capital), in percent, in five minutes and
without being told the goal \ts{24}{13}{3:14}\ts{13}{16}{3:10}. The answers become data for the rest of the hour.
In 2024 they were mostly index funds (S\&P~500, QQQ), Treasuries and cash, one option strategy, no cryptocurrency
(in 2021 the class had picked Dogecoin) \ts{24}{13}{9:26}\ts{24}{13}{11:35}; in 2013 they ranged from bonds and
S\&P index funds to timberland, rare coins and lottery tickets \ts{13}{16}{10:16}\ts{13}{16}{11:16}.

\subsection{Objective, loss tolerance, horizon}
Before choosing what to hold one must say what the portfolio is \emph{for}. The considerations listed by Xia are:
\begin{itemize}
  \item a \emph{return objective} (10\%, 50\%, 100\%?) and a \emph{time horizon}, which he deliberately left out
        of the exercise \ts{24}{13}{4:15}\ts{24}{13}{18:01};
  \item a \emph{loss tolerance}: how much of \$10{,}000 can be lost before one stops, and drawdown limits
        \ts{24}{13}{15:54};
  \item an \emph{edge}: why should this investment beat the alternatives \ts{24}{13}{5:17};
  \item \emph{constraints}: liquidity and spending needs, exclusions (fossil fuels, tobacco), inflation, absolute
        versus relative returns against a benchmark (the excess return being the \emph{alpha} of
        \cref{def:alphabeta}) \ts{24}{13}{14:48}\ts{24}{13}{16:56};
  \item \emph{asset--liability matching}: after 2008 many asset owners learned that when the portfolio draws down
        the liabilities dominate the balance sheet and spending must be cut \ts{24}{13}{18:01};
  \item \emph{peer pressure and career risk}: managers are compared quarter to quarter although the owner's
        horizon is long \ts{24}{13}{19:01}.
\end{itemize}
The 2013 version adds the individual's life cycle: income peaks around age 50, spending follows a different
curve, and the portfolio must bridge the gap; the same question is asked of a pension fund, an endowment and a
trading desk (how much capital to allocate to each trader, or to each of 30 strategies in a new hedge fund)
\ts{13}{16}{14:32}\ts{13}{16}{17:43}.
In 2024 the objective is written (slide~7) as $\max R/\sigma$ subject to constraints \ts{24}{13}{14:48}.

\subsection{Layers of decisions and the risk--return map}
The slides separate three layers \ts{24}{13}{7:22}: choice of \emph{markets and instruments};
\emph{data, signals, models, factors, predictions and strategies}; and finally \emph{sizing, capital allocation,
portfolio construction, optimisation and risk management}. Only the last is the topic of this chapter. It comes
down to comparing investments through their return and their risk, first taken to be the volatility
\ts{24}{13}{8:25}. On the map with volatility on the horizontal axis (2024 slide~6, 2013 blackboard),
cash sits on the vertical axis with zero volatility and a positive return (over 4\% in October 2024,
\ts{24}{13}{11:35}); bonds have a little more of both, stocks more, an index fund less volatility than a single stock,
commodities more, private equity and venture capital (illiquid, with a wide range of outcomes) are far up to the
right \ts{24}{13}{12:39}\ts{24}{13}{13:44}\ts{13}{16}{24:06}. Two extremes anchor the picture: a lottery ticket has a
return near $-100\%$ and a small volatility, and the fair coin flip has zero return and $100\%$ volatility
\ts{13}{16}{23:02}. A rational investor discards every point that has another point above and to the left of
it, which is the idea of the efficient frontier \ts{13}{16}{26:13}\ts{13}{16}{27:18}.

\begin{finance}[The endowment model]
A university endowment is a perpetual portfolio. It spends about $5\%$ of its value each year, and inflation is
about $3\%$, so the nominal return target is about $8\%$ (more precisely $1.05\times1.03-1\approx8.2\%$),
over horizons ``longer than ten years'' \ts{24}{13}{21:07}\ts{24}{13}{22:12}. The slide adds that liquidity
must cover $30$--$40\%$ of the annual operating budget, and Xia says that roughly $40\%$ of a university's budget is now
paid from endowment spending, as research grants shrink \ts{24}{13}{22:12}. An $8\%$ target is not easy to obtain
consistently when bonds yield little, so endowments hold a long list of strategies \ts{24}{13}{23:15}: cash,
government and corporate bonds and credit, macro and relative-value hedge funds, quantitative funds (trend
following/CTA, statistical arbitrage), long/short equity, multi-manager platforms, private equity, real estate and
natural resources, and newer assets (crypto, intellectual property, legal claims). The ``endowment model'' is
characterised by external rather than internal managers, public and private markets, generalists and specialists,
absolute and relative returns, mostly active management with some passive beta, and peer comparison
\ts{24}{13}{25:20}\ts{24}{13}{26:27}. Its central skill is manager selection (``can you predict fund performance,
luck or skill?''). Endowments are tax-exempt, so taxes are ignored; a family office would tilt to
private equity to defer realised gains \ts{24}{13}{1:13:53}\ts{24}{13}{1:14:54}.
\end{finance}

\subsection{The classic problem}
Assume the objectives and the loss tolerance are given and that one can predict the return, volatility and
correlation matrix of every investment. What fraction should go into each? The classic problem is to maximise the
portfolio return and minimise its risk \ts{24}{13}{27:30}. Three observations of the lecture frame the rest of
the chapter \ts{24}{13}{28:30}\ts{24}{13}{29:35}. First, the problem is one of \emph{sizing} at every level: for asset
classes it is asset allocation, and because ten to fifteen classes make optimisation hard, practitioners often
reduce them to three to five \emph{risk factors} (equity, bonds, inflation, currency, credit; see
\cref{ch:pca}). Second, risk management is the same problem seen from the other side, with the added concerns of
concentration, illiquidity and unwanted risks. Third, everything depends on predictions (``capital market
assumptions'') which are unreliable (\cref{ch10:sec-est}).

% ==================================================================
\section{Two assets: diversification and the feasible set}\label{ch10:sec-two}

Consider two risky assets with returns $R_1,R_2$, expected returns $\mu_1,\mu_2$, volatilities
$\sigma_1,\sigma_2$ and correlation $\rho$, and the fully invested portfolios with weights $(1-w,w)$, $0\le w\le1$
(no short sales) \ts{13}{14}{5:21}\ts{24}{13}{30:39}:
\begin{gather}\label{ch10:eq-two}
  R_w=(1-w)R_1+wR_2,\qquad \mu_w=(1-w)\mu_1+w\mu_2,\notag\\
  \sigma_w^2=(1-w)^2\sigma_1^2+w^2\sigma_2^2+2w(1-w)\rho\sigma_1\sigma_2 .
\end{gather}
The \emph{feasible set} is the curve $\{(\sigma_w,\mu_w):0\le w\le1\}$ in the volatility--return plane
\ts{13}{14}{7:28}. Since $\mu_w$ is affine in $w$ and $\sigma_w^2$ is quadratic, $\sigma_w^2$ is a quadratic
function of $\mu_w$: the curve is one branch of a hyperbola (a parabola in the $(\mu,\sigma^2)$ plane)
\ts{13}{14}{27:32}. Its shape is governed by $\rho$, illustrated by the special cases of both Xia lectures.

\begin{itemize}
  \item $\rho=1$: $\sigma_w=(1-w)\sigma_1+w\sigma_2$ is linear in $w$, hence so is $\mu$ as a function of $\sigma$:
        a straight segment between the two assets and no diversification gain
        \ts{24}{13}{31:39}\ts{13}{16}{34:34}.
  \item $\rho=-1$: $\sigma_w=|(1-w)\sigma_1-w\sigma_2|$, two segments meeting on the return axis at
        $w=\sigma_1/(\sigma_1+\sigma_2)$, where the portfolio has \emph{zero} variance
        \ts{24}{13}{32:39}\ts{13}{14}{17:06}. A portfolio without risk must earn the riskless rate, otherwise
        there is arbitrage; hedges of a security by a future exploit exactly this \ts{13}{14}{18:06}\ts{13}{14}{19:06}.
  \item $\rho=0$: the curve bends to the left of the straight line, and the variance cannot be reduced to zero.
  \item $\sigma_1=0$ (asset~1 is cash with return $r_0$): $\sigma_w=w\sigma_2$ and
        $\mu_w=r_0+(\mu_2-r_0)\sigma_w/\sigma_2$, a straight line with slope $(\mu_2-r_0)/\sigma_2$, the
        \emph{capital allocation line} of \cref{ch10:sec-rf} \ts{24}{13}{33:41}\ts{24}{13}{34:45}\ts{13}{16}{35:37}.
\end{itemize}

\begin{proposition}[Minimum-variance portfolio of two assets]\label{ch10:prop-two}
Let $\sigma_1^2+\sigma_2^2-2\rho\sigma_1\sigma_2>0$ (that is, not $\rho=1$ with $\sigma_1=\sigma_2$). The variance
\eqref{ch10:eq-two} is minimised over $w\in\R$ at
\[
  w^*=\frac{\sigma_1^2-\rho\sigma_1\sigma_2}{\sigma_1^2+\sigma_2^2-2\rho\sigma_1\sigma_2},\qquad
  \sigma^{*2}=\frac{\sigma_1^2\sigma_2^2(1-\rho^2)}{\sigma_1^2+\sigma_2^2-2\rho\sigma_1\sigma_2}.
\]
The weight $w^*$ lies in $[0,1]$ (the minimum is attained by a long-only portfolio) iff
$\rho\le\min(\sigma_1/\sigma_2,\sigma_2/\sigma_1)$; otherwise the minimum over $[0,1]$ is at the endpoint with smaller
volatility. If $\sigma_1=\sigma_2=\sigma$ then $w^*=\tfrac12$ and $\sigma^{*2}=\sigma^2(1+\rho)/2$.
\end{proposition}
\begin{proof}
Differentiating \eqref{ch10:eq-two}: $\tfrac12\,\dd\sigma_w^2/\dd w=-(1-w)\sigma_1^2+w\sigma_2^2+(1-2w)\rho\sigma_1\sigma_2$,
which vanishes at $w^*$; the second derivative $\sigma_1^2+\sigma_2^2-2\rho\sigma_1\sigma_2$ is positive, so this is the
global minimum. Substituting gives $\sigma^{*2}$ (use $1-w^*=(\sigma_2^2-\rho\sigma_1\sigma_2)/(\dots)$).
The numerator of $w^*$ is $\ge0$ iff $\rho\le\sigma_1/\sigma_2$, and that of $1-w^*$ is $\ge0$ iff
$\rho\le\sigma_2/\sigma_1$.
\end{proof}
For $\rho=0$ the formula gives $w^*=\dfrac{1/\sigma_2^2}{1/\sigma_1^2+1/\sigma_2^2}$: the assets are weighted
\emph{inversely to their variances}, as computed by Kempthorne on the board \ts{13}{14}{12:52}\ts{13}{14}{13:54}
(Exercise~\ref{ch10:ex-ps6-3} extends it).
Two consequences of \cref{ch10:prop-two} explain the pictures of the lectures.
\begin{itemize}
  \item Adding a little of the riskier
      asset~2 to asset~1 \emph{reduces} the variance as long as $\rho<\sigma_1/\sigma_2$
      (the derivative of $\sigma_w^2$ at $w=0$ is $2\sigma_1(\rho\sigma_2-\sigma_1)$, negative iff $\rho<\sigma_1/\sigma_2$). Even for $\rho=0$ diversification lowers the risk and, when asset~2 has the higher return,
      also \emph{raises} the return: every portfolio to the left of the minimum-variance point is dominated
      \ts{13}{14}{14:57}\ts{13}{14}{15:59}. For $\rho$ larger than the ratio of volatilities the minimum is the safer asset alone.
  \item The equal-volatility, zero-correlation case gives $\sigma^*=\sigma/\sqrt2$: ``a significant reduction of
      the risk by half-and-half'' \ts{13}{16}{31:31}\ts{13}{16}{33:33}; for three such assets it is $\sigma/\sqrt3$
      \ts{13}{16}{41:28}, and for $m$ assets $\sigma/\sqrt m$ (Exercise~\ref{ch10:ex-ps6-4}).
\end{itemize}

\begin{example}[The two assets of the 2013 simulation]\label{ch10:ex-two}
Take $\mu=(0.15,0.25)$, $\sigma=(0.20,0.30)$, the values printed in the plots of Case Study~6
(\cref{ch10:sec-case} explains why these are not the values read out in class). Then $\sigma_1/\sigma_2=2/3$ and
\begin{center}
\begin{tabular}{@{}lrrrrrr@{}}
\toprule
$\rho$ & $-1$ & $-0.4$ & $0$ & $0.4$ & $0.8$ & $1$\\
\midrule
$w^*$ (asset~2) & 0.400 & 0.360 & 0.308 & 0.195 & $-0.235$ & $-2$\\
$\sigma^*$ (long-only) & 0 & 0.130 & 0.166 & 0.192 & 0.200 & 0.200\\
$\mu^*$ (long-only) & 0.190 & 0.186 & 0.181 & 0.170 & 0.150 & 0.150\\
\bottomrule
\end{tabular}
\end{center}
For $\rho=0.8>2/3$ and $\rho=1$ the unconstrained $w^*$ is negative (a short position in asset~2), and the
long-only minimum is asset~1 alone. This is why in the last panel of \file{Smltn_TwoAst} (correlation $0.8$) the
minimum-variance portfolio (cyan curve) coincides with asset~1. All values were computed from
\cref{ch10:prop-two}.
\end{example}

\begin{figure}[ht]
\centering
\begin{tikzpicture}
\begin{axis}[width=0.86\linewidth,height=6.6cm,xlabel={volatility $\sigma_w$},ylabel={expected return $\mu_w$},
  xmin=-0.005,xmax=0.36,ymin=0.135,ymax=0.265,xtick={0,0.1,0.2,0.3},scaled ticks=false,/pgf/number format/fixed,grid=major,grid style={black!10},
  legend style={font=\footnotesize,at={(0.5,1.03)},anchor=south,legend columns=5},
  tick label style={font=\small},label style={font=\small}]
  \addplot[thick,thmcol] coordinates {(0.2,0.15) (0,0.19) (0.3,0.25)};
  \addplot[thick,cscol] coordinates {(0.2000,0.1500) (0.1845,0.1550) (0.1702,0.1600) (0.1575,0.1650) (0.1467,0.1700) (0.1383,0.1750) (0.1327,0.1800) (0.1304,0.1850) (0.1315,0.1900) (0.1358,0.1950) (0.1432,0.2000) (0.1531,0.2050) (0.1652,0.2100) (0.1789,0.2150) (0.1940,0.2200) (0.2101,0.2250) (0.2270,0.2300) (0.2446,0.2350) (0.2626,0.2400) (0.2811,0.2450) (0.3000,0.2500)};
  \addplot[thick,intcol] coordinates {(0.2000,0.1500) (0.1906,0.1550) (0.1825,0.1600) (0.1759,0.1650) (0.1709,0.1700) (0.1677,0.1750) (0.1664,0.1800) (0.1671,0.1850) (0.1697,0.1900) (0.1741,0.1950) (0.1803,0.2000) (0.1879,0.2050) (0.1970,0.2100) (0.2072,0.2150) (0.2184,0.2200) (0.2305,0.2250) (0.2433,0.2300) (0.2568,0.2350) (0.2707,0.2400) (0.2852,0.2450) (0.3000,0.2500)};
  \addplot[thick,fincol] coordinates {(0.2000,0.1500) (0.1965,0.1550) (0.1940,0.1600) (0.1925,0.1650) (0.1920,0.1700) (0.1927,0.1750) (0.1944,0.1800) (0.1971,0.1850) (0.2008,0.1900) (0.2054,0.1950) (0.2110,0.2000) (0.2173,0.2050) (0.2243,0.2100) (0.2320,0.2150) (0.2404,0.2200) (0.2492,0.2250) (0.2586,0.2300) (0.2684,0.2350) (0.2786,0.2400) (0.2891,0.2450) (0.3000,0.2500)};
  \addplot[thick,defcol] coordinates {(0.2,0.15) (0.3,0.25)};
  \legend{$\rho=-1$,$\rho=-0.4$,$\rho=0$,$\rho=0.4$,$\rho=1$}
  \addplot[only marks,mark=*,mark size=2pt,black,forget plot] coordinates {(0.2,0.15) (0.3,0.25)};
  \node[font=\footnotesize,anchor=east] at (axis cs:0.19,0.150) {asset 1};
  \node[font=\footnotesize,anchor=west] at (axis cs:0.305,0.25) {asset 2};
  \addplot[only marks,mark=square*,mark size=2pt,intcol,forget plot] coordinates {(0.1664,0.1808)};
  \draw[->,intcol] (axis cs:0.235,0.163) -- (axis cs:0.170,0.1790);
  \node[font=\footnotesize,anchor=west,intcol] at (axis cs:0.235,0.160) {min.\ variance, $\rho=0$};
\end{axis}
\end{tikzpicture}
\caption{Feasible set of the portfolios $(1-w)R_1+wR_2$, $0\le w\le1$, for $\mu=(0.15,0.25)$, $\sigma=(0.20,0.30)$
and several correlations (\cref{ch10:ex-two}; compare slides~15--16 of \file{lec13} and the top-right panels of
\file{Smltn_TwoAst}). For $\rho=-1$ the curve reaches the return axis. The lower part of each curve, below the
minimum-variance point, is dominated.}\label{ch10:fig-two}
\end{figure}

\begin{intuition}[Correlation as an angle]
Write $\sigma_w^2=|(1-w)\bm r_1+w\bm r_2|^2$ with vectors $\bm r_1,\bm r_2$ of lengths $\sigma_1,\sigma_2$ and
angle $\theta$, $\cos\theta=\rho$: the volatility of a portfolio is the \emph{length} of a weighted sum of two
vectors. Perfect correlation means parallel vectors (lengths add), $\rho=-1$ antiparallel (lengths cancel), and $\rho=0$
orthogonal (Pythagoras, $\sigma_w^2=(1-w)^2\sigma_1^2+w^2\sigma_2^2$). Diversification is the triangle inequality
being strict. The minimum-variance portfolio is the foot of the perpendicular from the origin to the segment
joining the tips of the two vectors.
\end{intuition}


% ==================================================================
\section{Markowitz mean--variance optimisation}\label{ch10:sec-mvo}

Now let there be $m$ risky assets with return vector $\bm R$, $\E[\bm R]=\bm\mu$, $\Cov\bm R=\Sigma$ (positive
definite) in a \emph{single period}. A portfolio is a vector $\bm w$ of fractions of the wealth invested in each
asset, $\bm w\T\bm 1=1$ (one unit of capital). Its return $R_{\bm w}=\bm w\T\bm R$ has
\begin{equation}\label{ch10:eq-mv}
  \mu_{\bm w}=\bm w\T\bm\mu,\qquad \sigma_{\bm w}^2=\bm w\T\Sigma\bm w
\end{equation}
(\cref{ch03:prop-portvar}) \ts{13}{14}{3:11}\ts{13}{14}{4:16}. Investors prefer higher $\mu_{\bm w}$ and lower
$\sigma^2_{\bm w}$, so portfolios are compared through the pair $(\mu_{\bm w},\sigma_{\bm w}^2)$
\ts{13}{14}{23:14}. Markowitz's idea (1952) is to pose this as a convex optimisation problem.

\begin{definition}[Markowitz problem I]\label{ch10:def-markowitz}
For a target return $\mu_0$, choose $\bm w\in\R^m$ (short positions allowed) to
\[
  \text{minimise } \tfrac12\bm w\T\Sigma\bm w\quad\text{subject to}\quad \bm w\T\bm\mu=\mu_0,\ \ \bm w\T\bm 1=1 .
\]
\end{definition}
The factor $\frac12$ only simplifies the derivatives \ts{13}{14}{24:14}. The objective is strictly convex ($\Sigma\succ0$)
and the constraints are linear, so a stationary point of the Lagrangian is the unique global minimum
\ts{13}{14}{25:20}.

\begin{theorem}[Efficient portfolios without a riskless asset]\label{ch10:thm-mvo}
Put
\[
  a=\bm\mu\T\Sigma^{-1}\bm\mu,\qquad b=\bm\mu\T\Sigma^{-1}\bm 1,\qquad c=\bm 1\T\Sigma^{-1}\bm 1,\qquad
  \Delta=ac-b^2 .
\]
If $\bm\mu$ is not a multiple of $\bm 1$ then $\Delta>0$ and problem~I has the unique solution
\begin{equation}\label{ch10:eq-w0}
  \bm w_0=\lambda_1\Sigma^{-1}\bm\mu+\lambda_2\Sigma^{-1}\bm 1,\qquad
  \begin{pmatrix}\lambda_1\\ \lambda_2\end{pmatrix}=\begin{pmatrix}a&b\\ b&c\end{pmatrix}^{-1}\begin{pmatrix}\mu_0\\ 1\end{pmatrix}
  =\frac1\Delta\begin{pmatrix}c\mu_0-b\\ a-b\mu_0\end{pmatrix},
\end{equation}
whose variance is
\begin{equation}\label{ch10:eq-frontier}
  \sigma_0^2=\bm w_0\T\Sigma\bm w_0=\frac{c\mu_0^2-2b\mu_0+a}{\Delta}
  =\frac1c+\frac{c}{\Delta}\Bigl(\mu_0-\frac bc\Bigr)^2 .
\end{equation}
\end{theorem}
\begin{proof}
The Lagrangian is $L=\frac12\bm w\T\Sigma\bm w+\lambda_1(\mu_0-\bm w\T\bm\mu)+\lambda_2(1-\bm w\T\bm 1)$
\ts{13}{14}{24:14}. Setting $\partial L/\partial\bm w=\Sigma\bm w-\lambda_1\bm\mu-\lambda_2\bm 1=\bm 0$ gives
$\bm w=\lambda_1\Sigma^{-1}\bm\mu+\lambda_2\Sigma^{-1}\bm 1$. Substituting in the two constraints gives the linear
system $\mu_0=a\lambda_1+b\lambda_2$, $1=b\lambda_1+c\lambda_2$, that is \eqref{ch10:eq-w0} \ts{13}{14}{26:26}.
The matrix $\begin{psmallmatrix}a&b\\b&c\end{psmallmatrix}$ is the Gram matrix of $\bm\mu,\bm 1$ for the inner
product $\langle\bm x,\bm y\rangle=\bm x\T\Sigma^{-1}\bm y$; its determinant $\Delta$ is positive by the strict
Cauchy--Schwarz inequality, because $\bm\mu$ and $\bm 1$ are linearly independent. Finally
$\sigma_0^2=\bm w_0\T\Sigma\bm w_0=\lambda_1\bm w_0\T\bm\mu+\lambda_2\bm w_0\T\bm 1=\lambda_1\mu_0+\lambda_2$, which
equals \eqref{ch10:eq-frontier}; completing the square in $\mu_0$ gives the second form.
\end{proof}

Thus $\sigma_0^2$ is a parabola in the target return \ts{13}{14}{27:32}: the smallest variance obtainable
with mean $\mu_0$, and the set of all pairs $(\sigma_0,\mu_0)$ is a hyperbola in the volatility--return plane.

\begin{corollary}[Global minimum-variance portfolio and the frontier]\label{ch10:cor-gmv}
The parabola \eqref{ch10:eq-frontier} is smallest at $\mu_0=b/c$, where
\[
  \bm w_g=\frac{\Sigma^{-1}\bm 1}{\bm 1\T\Sigma^{-1}\bm 1},\qquad \sigma_g^2=\frac1c,\qquad \mu_g=\frac bc .
\]
The \emph{efficient frontier} is the upper branch $\mu_0\ge b/c$ (portfolios below it are dominated by the
portfolio with the same variance and larger return); the whole curve is a hyperbola with asymptotes
$\mu=b/c\pm\sqrt{\Delta/c}\,\sigma$. Along the efficient branch $\mu$ is an increasing concave function of $\sigma$, and no portfolio lies to the left of the hyperbola.
\end{corollary}

\subsection{The three equivalent problems}
The slides pose the same optimisation in three ways (slide~7) \ts{13}{14}{28:36}\ts{13}{14}{29:42}:
\begin{itemize}
  \item \textbf{Problem I}: minimise the variance for a given mean $\mu_0$;
  \item \textbf{Problem II}: maximise $\bm w\T\bm\mu$ for a given variance $\sigma_0^2$, with $\bm w\T\bm 1=1$;
  \item \textbf{Problem III}: for a risk-aversion parameter $\lambda\ge0$ maximise
        $\bm w\T\bm\mu-\frac\lambda2\bm w\T\Sigma\bm w$ with $\bm w\T\bm 1=1$.
\end{itemize}
All three have the same first-order conditions, and the \emph{efficient frontier} is the trace of the solutions
when $\mu_0$ (I), $\sigma_0^2$ (II) or $\lambda$ (III) varies. Problem~III can be solved in one line, and it shows the
structure of the frontier. The condition $\bm\mu-\lambda\Sigma\bm w-\gamma\bm 1=\bm 0$ (with multiplier $\gamma$
for the budget) gives $\bm w=\lambda^{-1}\Sigma^{-1}(\bm\mu-\gamma\bm 1)$, and $\bm 1\T\bm w=1$ fixes $\gamma=(b-\lambda)/c$:
\begin{equation}\label{ch10:eq-twofund}
  \bm w(\lambda)=\bm w_g+\frac1\lambda\,\bm v,\qquad \bm v=\Sigma^{-1}\Bigl(\bm\mu-\frac bc\bm 1\Bigr),\qquad
  \bm 1\T\bm v=0 .
\end{equation}
Every efficient portfolio is the minimum-variance portfolio plus a multiple $1/\lambda$ of one fixed
\emph{zero-cost} (long/short) portfolio $\bm v$. Hence $\mu=b/c+(a-b^2/c)/\lambda$ and
$\sigma^2=1/c+(a-b^2/c)/\lambda^2$, since $\bm v\T\Sigma\bm v=a-b^2/c$ and $\bm w_g\T\Sigma\bm v=0$; eliminating $\lambda$
recovers \eqref{ch10:eq-frontier}. Consequently any two distinct efficient portfolios span the whole frontier (the
\emph{two-fund theorem}, not stated in class). The risk aversion $\lambda$ is the Arrow--Pratt index of
\cref{ch10:sec-utility}: a very risk-averse investor ($\lambda\to\infty$) holds $\bm w_g$.

\begin{example}[Nine sector ETFs, unconstrained (computed from the rounded case-study inputs)]\label{ch10:ex-etf-unc}
Take the annualised means, volatilities and correlations of the nine SPDR sector ETFs for 2009--2013 printed in
Case Study~6 (\cref{ch10:tab-etf}). Then $a=3.44$, $b=11.44$, $c=88.04$, $\Delta=171.9$. The global minimum-variance
portfolio has $\sigma_g=1/\sqrt c=10.7\%$ and $\mu_g=b/c=13.0\%$, better in volatility than every single ETF
(the lowest is consumer staples, $12\%$, but its mean is $15\%$), and the asymptotic slope is
$\sqrt{\Delta/c}=1.40$. But the weights, in the order XLB, XLV, XLP, XLY, XLE, XLF, XLI, XLK, XLU, are
\[
  \bm w_g\approx(0.09,\ 0.14,\ 0.94,\ -0.17,\ -0.15,\ -0.02,\ -0.27,\ 0.26,\ 0.18),
\]
with $94\%$ in consumer staples and a $27\%$ short position in industrials: an unconstrained optimum is not
what an investor would hold. This is the reason for the constraints of \cref{ch10:sec-constr}.
\end{example}

% ==================================================================
\section{A riskless asset: tangency portfolio and capital market line}\label{ch10:sec-rf}

Add a riskless asset $0$ with constant return $R_0\equiv r_0$ (zero variance, uncorrelated with everything)
\ts{13}{14}{31:55}. The investor puts $\bm w$ in the risky assets and $1-\bm w\T\bm 1$ in the riskless one, which
may be negative (borrowing at $r_0$). There is no budget constraint left on $\bm w$, and
\begin{equation}\label{ch10:eq-mvrf}
  \mu_{\bm w}=\bm w\T\bm\mu+(1-\bm w\T\bm 1)r_0=r_0+\bm w\T(\bm\mu-r_0\bm 1),\qquad \sigma_{\bm w}^2=\bm w\T\Sigma\bm w .
\end{equation}
\emph{Geometric picture.} Two assets are enough: combining the riskless asset with a risky portfolio~$P$ gives the
line through $(0,r_0)$ and $(\sigma_P,\mu_P)$; the higher the slope, the better. The best line is the tangent to
the efficient frontier of the risky assets \ts{24}{13}{35:50}\ts{13}{14}{33:01}, which lies above the frontier
because we can now hold a fraction of a risky portfolio in cash \ts{13}{14}{35:05}. The mathematics confirms it.

\begin{theorem}[Optimal portfolios with a riskless asset; Tobin separation]\label{ch10:thm-tobin}
Let $\bm\eta=\bm\mu-r_0\bm 1$ (excess returns) and $d^2=\bm\eta\T\Sigma^{-1}\bm\eta=a-2r_0b+r_0^2c$. The problem
$\min\frac12\bm w\T\Sigma\bm w$ subject to $\mu_{\bm w}=\mu_0$ has the solution
\begin{equation}\label{ch10:eq-rfsol}
  \bm w_0=\lambda_1\Sigma^{-1}\bm\eta,\qquad \lambda_1=\frac{\mu_0-r_0}{d^2},\qquad
  \sigma_0=\frac{|\mu_0-r_0|}{d}.
\end{equation}
Assume $\bm 1\T\Sigma^{-1}\bm\eta=b-r_0c>0$ (i.e.\ $r_0<\mu_g$) and define the \emph{market (tangency) portfolio}
\begin{equation}\label{ch10:eq-market}
  \bm w_M=\frac{\Sigma^{-1}\bm\eta}{\bm 1\T\Sigma^{-1}\bm\eta},\qquad
  \mu_M=\frac{a-r_0b}{b-r_0c},\qquad \sigma_M=\frac{d}{b-r_0c}.
\end{equation}
Then every optimal portfolio with $\mu_0>r_0$ invests in the risky assets in the proportions of $\bm w_M$: it holds
the fraction
\begin{equation}\label{ch10:eq-wm}
  m_{\bm w}=\bm w_0\T\bm 1=\frac{\mu_0-r_0}{\mu_M-r_0}
\end{equation}
of its wealth in $\bm w_M$ and $1-m_{\bm w}$ in the riskless asset \ts{13}{14}{41:30}\ts{13}{14}{42:31}.
Its return and risk are $\mu_P=r_0+m_{\bm w}(\mu_M-r_0)$ and $\sigma_P=m_{\bm w}\sigma_M$.
\end{theorem}
\begin{proof}
The Lagrangian $L=\frac12\bm w\T\Sigma\bm w+\lambda_1[(\mu_0-r_0)-\bm w\T\bm\eta]$ gives
$\Sigma\bm w=\lambda_1\bm\eta$, i.e.\ $\bm w=\lambda_1\Sigma^{-1}\bm\eta$, and the constraint
$\bm w\T\bm\eta=\mu_0-r_0$ gives $\lambda_1d^2=\mu_0-r_0$ \ts{13}{14}{35:05}\ts{13}{14}{36:08}. The variance is
$\lambda_1^2\bm\eta\T\Sigma^{-1}\bm\eta=\lambda_1^2d^2=(\mu_0-r_0)^2/d^2$. Since $\bm w_0=\lambda_1(b-r_0c)\bm w_M$,
the risky weights are proportional to $\bm w_M$ with total weight $m_{\bm w}=\lambda_1(b-r_0c)$; \eqref{ch10:eq-market}
follows from $\mu_M=\bm w_M\T\bm\mu$, $\sigma_M^2=\bm\eta\T\Sigma^{-1}\bm\eta/(\bm 1\T\Sigma^{-1}\bm\eta)^2$, and
$\mu_M-r_0=d^2/(b-r_0c)$ gives \eqref{ch10:eq-wm}.
\end{proof}

\begin{corollary}[Capital market line and price of risk]\label{ch10:cor-cml}
The optimal portfolios lie on the half-line, the \emph{capital market line} (CML)
\begin{equation}\label{ch10:eq-cml}
  \mu_P=r_0+\frac{\mu_M-r_0}{\sigma_M}\,\sigma_P,\qquad \sigma_P\ge0,
\end{equation}
which is tangent to the efficient frontier at $M$. Its slope $(\mu_M-r_0)/\sigma_M=d$ is the \emph{Sharpe ratio} of
the market portfolio, the ``market price of risk'': the compensation per unit of volatility. Investors of all risk
appetites hold the \emph{same} risky portfolio and differ only in how much of it (Tobin, 1958)
\ts{13}{14}{43:42}\ts{13}{14}{44:47}. If borrowing at $r_0$ is possible, $m_{\bm w}>1$ (i.e.\ $\mu_0>\mu_M$) gives points of the CML
beyond $M$, a \emph{leveraged} market portfolio \ts{13}{14}{46:52}.
\end{corollary}

\begin{definition}[Sharpe ratio]
For a portfolio with return $R_P$, the Sharpe ratio is $\mathrm{SR}_P=(\mu_P-r_0)/\sigma_P$
\ts{24}{13}{36:54}\ts{13}{16}{44:21}.
\end{definition}

\begin{proposition}[The tangency portfolio maximises the Sharpe ratio]\label{ch10:prop-sharpe}
For every risky portfolio $\bm w\neq\bm 0$, $(\bm w\T\bm\eta)/\sqrt{\bm w\T\Sigma\bm w}\le d$, with equality iff $\bm w$
is a positive multiple of $\Sigma^{-1}\bm\eta$, i.e.\ of $\bm w_M$. The Sharpe ratio does not change when a portfolio
is mixed with the riskless asset or levered at rate $r_0$.
\end{proposition}
\begin{proof}
Cauchy--Schwarz for the inner product defined by $\Sigma$:
\[(\bm w\T\bm\eta)^2=\bigl(\bm w\T\Sigma\cdot\Sigma^{-1}\bm\eta\bigr)^2
\le(\bm w\T\Sigma\bm w)(\bm\eta\T\Sigma^{-1}\bm\eta).\] Scaling $\bm w\to t\bm w$ multiplies $\mu-r_0$ and $\sigma$ by $t$.
\end{proof}

The last statement is the answer to Xia's question in 2013 (``when you lever up, does the Sharpe ratio change?''):
with $w_1=-1$ in cash and $w_2=2$ in the risky asset the volatility is $2\sigma_2$ and the excess return
$2(R_2-R_1)$, so the ratio $(R_2-R_1)/\sigma_2$ is the same at every point of the line
\ts{13}{16}{56:12}\ts{13}{16}{57:16}\ts{13}{16}{1:01:28}; the 2024 slide~18 draws the same picture,
$\sigma=2\sigma_p$, $R=2R_p-R_f$ \ts{24}{13}{38:56}. Leverage moves the investor along the line, it does not improve it.

\begin{example}[Two assets and a riskless rate]\label{ch10:ex-tobin}
Take $\bm\mu=(0.15,0.25)$, $\Sigma=\mathrm{diag}(0.04,0.09)$ (uncorrelated assets of \cref{ch10:ex-two}) and
$r_0=0.05$. Then $\bm\eta=(0.10,0.20)$, $\Sigma^{-1}\bm\eta=(2.5,\ 2.22)$, $\bm w_M=(0.529,\ 0.471)$,
$\mu_M=19.7\%$, $\sigma_M=17.6\%$ and $d=\sqrt{0.10^2/0.04+0.20^2/0.09}=0.833$; the individual Sharpe ratios are only
$0.50$ and $0.67$. An investor who wants $\mu_0=12\%$ puts $(0.12-0.05)/(0.197-0.05)=47.6\%$ of the wealth in $M$ and
$52.4\%$ in cash; one who wants $\mu_0=20\%$ borrows $2\%$ of the wealth. \sloppy Twice levered ($m_{\bm w}=2$) gives
$\mu=34.4\%$ at $\sigma=35.3\%$, still on the line (\cref{ch10:fig-cml}).
\end{example}

\begin{figure}[ht]
\centering
\begin{tikzpicture}
\begin{axis}[width=0.86\linewidth,height=7cm,xlabel={volatility $\sigma$},ylabel={expected return $\mu$},
  xmin=0,xmax=0.45,ymin=0,ymax=0.40,xtick={0,0.1,0.2,0.3,0.4},scaled ticks=false,/pgf/number format/fixed,grid=major,grid style={black!10},
  tick label style={font=\small},label style={font=\small}]
  \addplot[thick,defcol,dashed] coordinates {(0.3672,0.0900) (0.3354,0.1000) (0.3046,0.1100) (0.2751,0.1200) (0.2474,0.1300) (0.2220,0.1400) (0.2000,0.1500) (0.1825,0.1600) (0.1709,0.1700) (0.1664,0.1808)};
  \addplot[very thick,defcol] coordinates {(0.1664,0.1808) (0.1697,0.1900) (0.1803,0.2000) (0.1970,0.2100) (0.2184,0.2200) (0.2433,0.2300) (0.2707,0.2400) (0.3000,0.2500) (0.3306,0.2600) (0.3622,0.2700) (0.3946,0.2800) (0.4276,0.2900) (0.4610,0.3000)};
  \addplot[thick,thmcol] coordinates {(0,0.05) (0.40,0.3833)} node[pos=0.22,above,sloped,font=\footnotesize]{CML, slope $0.833$};
  \addplot[only marks,mark=*,mark size=2.2pt,thmcol] coordinates {(0.17647,0.19706) (0,0.05) (0.3529,0.3441)};
  \addplot[only marks,mark=square*,mark size=2pt,black] coordinates {(0.1664,0.1808) (0.2,0.15) (0.3,0.25)};
  \node[font=\footnotesize,anchor=east,thmcol] at (axis cs:0.168,0.225) {$M$};
  \node[font=\footnotesize,anchor=south west,thmcol] at (axis cs:0.005,0.055) {$r_0$};
  \node[font=\footnotesize,anchor=north west,thmcol] at (axis cs:0.36,0.335) {leveraged};
  \node[font=\footnotesize,anchor=south west] at (axis cs:0.203,0.153) {asset 1};
  \node[font=\footnotesize,anchor=north west] at (axis cs:0.305,0.245) {asset 2};
  \node[font=\footnotesize,anchor=east] at (axis cs:0.150,0.225) {min.\ variance};\draw[->] (axis cs:0.120,0.215) -- (axis cs:0.1640,0.1830);
    \node[font=\footnotesize,defcol,anchor=east] at (axis cs:0.445,0.19) {frontier (solid: efficient)};
\end{axis}
\end{tikzpicture}
\caption{Frontier of the two assets of \cref{ch10:ex-two} for $\rho=0$ with short sales allowed (upper branch solid,
dominated branch dashed) and the capital market line of \cref{ch10:ex-tobin} with $r_0=5\%$, tangent at $M$.
Points on the CML beyond $M$ are levered.}\label{ch10:fig-cml}
\end{figure}

\begin{intuition}[Why one fund is enough]
The direction $\Sigma^{-1}\bm\eta$ solves the \emph{normal equations} $\Sigma\bm w=\bm\eta$, the same form as
$X\T X\hat{\bm\beta}=X\T\bm y$ of \cref{ch:regression}: it is the portfolio whose covariances with the assets are
proportional to their excess returns, which is exactly the CAPM relation \eqref{ch10:eq-sml} of
\cref{ch10:sec-capm}. Once the best ``direction'' in asset space is fixed, scaling it up or down is only a question of
appetite, like rescaling a normal mode; the riskless asset supplies the one extra degree of freedom (the
length) that the fully invested constraint of \cref{ch10:sec-mvo} did not leave.
\end{intuition}

\begin{coursenote}[Naming and the case $r_0>\mu_g$]
Kempthorne calls $\bm w_M$ ``the market portfolio'' because it is the \emph{fully invested} optimal
portfolio in the sense of $\bm w_M\T\bm 1=1$ \ts{13}{14}{39:15}; the same portfolio is called the ``tangency portfolio''
elsewhere. That it coincides with the cap-weighted market of all assets is a statement of \emph{equilibrium}
(Sharpe 1964, Lintner 1965), proved in \cref{ch03:thm-capm} and used in \cref{ch10:sec-capm}; here it is only the
name of a solution. The slides implicitly assume $b-r_0c>0$; if instead $r_0>\mu_g$ the tangency to the frontier
touches the lower, dominated branch and $\bm w_M$ has a negative Sharpe ratio; the optimal portfolio for $\mu_0>r_0$
is then $-\bm w_M$ scaled (short the ``market''), which is unrealistic and signals an unreasonable $r_0$. The
slides also assume that the borrowing and lending rates coincide; with a higher borrowing rate the CML has a
kink at $M$ and only the segment up to $M$ is tangent to the frontier (not discussed in class).
The key papers listed on slide~16 of \file{lecnote14} are Markowitz (1952), Tobin (1958), Sharpe (1964),
Lintner (1965) and Fama (1970) \ts{13}{14}{47:54}; Kempthorne remarks that virtually all of them have a Nobel prize,
and that Fama had been awarded it two weeks earlier \ts{13}{14}{48:55}.
\end{coursenote}

% ==================================================================
\section{Beta, the security market line and the CAPM}\label{ch10:sec-capm}

The mean--variance solution also contains the CAPM, by a route different from the expected-utility derivation of
\cref{ch03:sec-capm}. This derivation is not given in either lecture, but it explains the slides on
$\alpha,\beta$ (2024 slide~17, 2013 L16).
\begin{proposition}[Security market line from the tangency portfolio]\label{ch10:prop-sml}
Let $M$ be the tangency portfolio \eqref{ch10:eq-market}. For every asset $i$ (and every portfolio),
\begin{equation}\label{ch10:eq-sml}
  \mu_i-r_0=\beta_i(\mu_M-r_0),\qquad \beta_i=\frac{\Cov(R_i,R_M)}{\sigma_M^2}.
\end{equation}
\end{proposition}
\begin{proof}
By definition $\Sigma\bm w_M=\bm\eta/(\bm 1\T\Sigma^{-1}\bm\eta)$, i.e.\ $\Cov(R_i,R_M)=[\Sigma\bm w_M]_i=
\eta_i\,\kappa$ with $\kappa=1/(\bm 1\T\Sigma^{-1}\bm\eta)$. Multiplying by $w_{M,i}$ and summing gives
$\sigma_M^2=\kappa(\mu_M-r_0)$, hence $\eta_i=\Cov(R_i,R_M)/\kappa=\Cov(R_i,R_M)(\mu_M-r_0)/\sigma_M^2$.
\end{proof}
Under the assumptions of Sharpe and Lintner every investor holds $M$, so $M$ is the market portfolio of all risky
assets and \eqref{ch10:eq-sml} is the CAPM (\cref{ch03:thm-capm}). The regression of $R_i-r_0$ on the market
excess return has slope $\beta_i$ and, by \eqref{ch10:eq-sml}, zero intercept: in the language of \cref{def:alphabeta}
\[
  R_P-r_0=\alpha+\beta\,(R_{BM}-r_0)+\varepsilon,\qquad \beta=\frac{\Cov(R_P,R_{BM})}{\sigma_{BM}^2}=\rho\,\frac{\sigma_P}{\sigma_{BM}}
\]
(2024 slide~17 \ts{24}{13}{37:54}), where a nonzero $\alpha$ measures a departure from the security market line, and
$\sigma_P^2=\beta^2\sigma_{BM}^2+\sigma_\varepsilon^2$ splits the risk into a systematic and an idiosyncratic part
(regression, \cref{ch05:capm}). In the two-asset special case of the lectures (riskless asset plus one risky
asset $2$, $\sigma_P=w_2\sigma_2$) the correlation with asset~2 is $1$ and $\beta=\sigma_P/\sigma_2=w_2$
\ts{13}{16}{1:02:34}\ts{13}{16}{1:03:38}: the 60/40 investor asked how much equity beta the portfolio has needs the
same computation with both assets (\cref{ch10:sec-rp}).

\begin{finance}[Alpha, beta and the Sharpe ratio in practice]
Beta is cheap to buy (an index fund); alpha is what a manager is paid for, and the Sharpe ratio $\mathrm{SR}$
measures return per unit of total risk while the \emph{information ratio} $\alpha/\sigma_\varepsilon$ measures
skill per unit of active risk (not discussed in class). A portfolio manager benchmarked to an index is judged on
$\alpha$ and on the tracking error (\cref{ch10:sec-constr}); a fund with a high Sharpe ratio can be levered to
any desired volatility, so Sharpe ratios are the fair basis for comparing funds of different volatilities.
\end{finance}


% ==================================================================
\section{Diversification and rebalancing: the ``free lunch''}\label{ch10:sec-rebal}

\subsection{How much can diversification remove?}
For the equally weighted portfolio of $n$ assets, with $\bar\sigma^2$ the average variance and $\bar c$ the average
pairwise covariance,
\begin{equation}\label{ch10:eq-ew}
  \Var(R_{\rm ew})=\frac1{n^2}\Bigl(\sum_i\sigma_i^2+\sum_{i\ne j}\Sigma_{ij}\Bigr)=\frac{\bar\sigma^2}n+\Bigl(1-\frac1n\Bigr)\bar c .
\end{equation}
The first term (idiosyncratic risk) vanishes like $1/n$, the second does not: what survives is the average
covariance, the \emph{systematic} risk that no diversification removes (\cref{ch07:ex-equi} for the spectral
version). For uncorrelated assets of equal volatility, $\sigma/\sqrt n$; Exercises~\ref{ch10:ex-ps6-4} and
\ref{ch10:ex-ps6-5} ask for the minimum-variance portfolios in the uncorrelated, equicorrelated and general cases.

\subsection{Rebalancing}
Xia's second illustration of diversification in both years is a two-year, two-asset example
\ts{24}{13}{40:00}\ts{13}{16}{52:00}. Asset~1 doubles in year~1 and halves in year~2; asset~2 halves and then doubles.
Alone, each asset ends where it started (return $0$). Start with $50\%$ in each \ts{24}{13}{41:06}.
\begin{itemize}
  \item Year~1: $0.5\cdot(+100\%)+0.5\cdot(-50\%)=+25\%$, so the value is $1.25$, of which
        $1.00$ ($80\%$) is now asset~1 and $0.25$ ($20\%$) asset~2 \ts{24}{13}{42:13}.
  \item \emph{Without rebalancing}, year~2 gives $0.8\cdot(-50\%)+0.2\cdot(+100\%)=-20\%$, and the two-year return is
        $1.25\times0.8-1=0$: no gain from a diversified portfolio of perfectly negatively correlated assets.
  \item \emph{With rebalancing} to $50/50$ at the end of year~1 (sell $0.375$ of asset~1, buy asset~2: $0.625$ each),
        year~2 gives $0.5\cdot(-50\%)+0.5\cdot(+100\%)=+25\%$: two years at $+25\%$, total $1.25^2=1.5625$
        \ts{24}{13}{42:13}\ts{13}{16}{54:05}\ts{13}{16}{55:08}.
\end{itemize}
Rebalanced, the portfolio has \emph{zero} volatility and a steady $25\%$ return, ``a straight line''.

\begin{figure}[ht]
\centering
\begin{tikzpicture}
\begin{axis}[width=0.62\linewidth,height=5.2cm,xlabel={year},ylabel={value of \$1},xmin=-0.1,xmax=2.1,ymin=0.3,ymax=2.2,
  xtick={0,1,2},grid=major,grid style={black!10},legend style={font=\footnotesize,at={(0.5,1.03)},anchor=south,legend columns=2},
  tick label style={font=\small},label style={font=\small}]
  \addplot[thick,defcol,mark=*] coordinates {(0,1) (1,2) (2,1)};
  \addplot[thick,intcol,mark=*] coordinates {(0,1) (1,0.5) (2,1)};
  \addplot[thick,thmcol,mark=square*,dashed] coordinates {(0,1) (1,1.25) (2,1.0)};
  \addplot[very thick,cscol,mark=square*] coordinates {(0,1) (1,1.25) (2,1.5625)};
  \legend{asset 1,asset 2,50/50 no rebalancing,50/50 rebalanced}
\end{axis}
\end{tikzpicture}
\caption{The two-year example of Xia (2024 slide~20, 2013 L16): each asset ends at $1$, the buy-and-hold $50/50$
portfolio also, the rebalanced one at $1.5625$.}\label{ch10:fig-rebal}
\end{figure}

Two objections are discussed in class \ts{24}{13}{43:16}\ts{24}{13}{44:18}. Is this not a bet on mean reversion,
losing money in a trend? Yes: rebalancing sells the winner and buys the loser, which pays when returns alternate
and loses when they trend (the case of \cref{sec:ewport}, where daily rebalancing to equal weights was not obviously
better). The rule of thumb is therefore \emph{consistency of views}: the equal weights expressed the belief that the
two assets have the same distribution and the same risk; unless the view has changed, letting the weights drift
means holding a non-equal portfolio by accident. ``The only free lunch is diversification, but you have to
rebalance to eat it.''

\begin{proposition}[Excess growth of a constant mix; not discussed in class]\label{ch10:prop-growth}
Let asset $i$ follow a geometric Brownian motion with $\dd S_i/S_i=\mu_i\dd t+\sigma_i\dd W_i$ and
$\Cov(\dd W_i,\dd W_j)=\rho_{ij}\dd t$. Asset $i$ alone grows at the logarithmic rate $\mu_i-\frac12\sigma_i^2$.
A portfolio rebalanced \emph{continuously} to the weights $\bm w$ has
\begin{equation}\label{ch10:eq-growth}
  g(\bm w)=\bm w\T\bm\mu-\tfrac12\bm w\T\Sigma\bm w,\qquad
  g(\bm w)-\sum_iw_i\bigl(\mu_i-\tfrac12\sigma_i^2\bigr)=\tfrac12\Bigl(\sum_iw_i\sigma_i^2-\bm w\T\Sigma\bm w\Bigr)\ge0 .
\end{equation}
\end{proposition}
\begin{proof}
The wealth $V$ of the constant mix satisfies $\dd V/V=\bm w\T\bm\mu\dd t+\bm w\T\bm\sigma\dd\bm W$ (its return is
$\bm w\T$ times the vector of returns, \cref{ch:stochcalc}), so by It\^o's formula
$\dd\log V=(\bm w\T\bm\mu-\tfrac12\bm w\T\Sigma\bm w)\dd t+\dots$. The inequality is $\Var(\sum_iw_iR_i)\le\sum_iw_i\Var R_i$,
i.e.\ convexity of the variance.
\end{proof}
The gain is the ``volatility harvesting'' or ``diversification return''. Two uncorrelated assets with $\sigma=20\%$
and the same $\mu$, kept at $50/50$, grow at $\mu-1\%$ per year, while each alone grows at $\mu-2\%$: the
rebalanced portfolio gains $\sigma^2/4=1\%$ per year over the average of its components. The example above is the
extreme case: $\rho=-1$, growth rates of $0\%$ for each asset, $25\%$ for the mix.

% ==================================================================
\section{Leverage, the 60/40 portfolio and risk parity}\label{ch10:sec-rp}

Two-asset portfolios are important in practice, because managers are often benchmarked to a mix of equities and
bonds, the most famous being the \emph{60/40} portfolio: $60\%$ in equities and $40\%$ in bonds. Any fund is still
compared with it \ts{13}{16}{41:28}\ts{13}{16}{42:34}\ts{13}{16}{47:42}.

\subsection{Weights are not risks}
Equities are much more volatile than bonds, so a $60/40$ portfolio is really an equity portfolio in risk terms.
The idea of \emph{risk parity} (the name was coined in 2005 by Edward Qian, as Xia recalls) is to allocate
\emph{risk} rather than capital: equal risk from each component, instead of weights by market value
\ts{13}{16}{49:54}\ts{13}{16}{50:55}\ts{24}{13}{38:56}. Before 2000, Xia recalls, a passive $60/40$ was the norm; the
2000 equity peak and bond rally and the 2008 crash, when government bonds rallied while equities fell, made
investors question it \ts{13}{16}{48:48}\ts{13}{16}{49:54}. To define the risk contribution of an asset one
needs a decomposition of $\sigma_P=\sqrt{\bm w\T\Sigma\bm w}$; it comes from Euler's theorem, since $\sigma_P$ is
homogeneous of degree one in $\bm w$.

\begin{definition}[Risk contributions]\label{ch10:def-rc}
The \emph{risk contribution} of asset $i$ to a portfolio $\bm w$ is
\[
  \mathrm{RC}_i=w_i\frac{\partial\sigma_P}{\partial w_i}=\frac{w_i(\Sigma\bm w)_i}{\sigma_P}=\frac{w_i\Cov(R_i,R_P)}{\sigma_P},
  \qquad \sum_i\mathrm{RC}_i=\sigma_P .
\]
A portfolio has \emph{equal risk contributions} (risk parity, ERC) if $\mathrm{RC}_i=\mathrm{RC}_j$ for all $i,j$, that is
$w_i(\Sigma\bm w)_i=$ const.
\end{definition}
The identity $\sum_i\mathrm{RC}_i=\sigma_P$ is Euler's theorem, or simply $\sum_iw_i(\Sigma\bm w)_i=\bm w\T\Sigma\bm w=\sigma_P^2$.

\begin{proposition}[Risk parity for two assets and for equal correlations]\label{ch10:prop-rp}
\leavevmode
\begin{enumerate}[(i)]
  \item For two assets with volatilities $\sigma_1,\sigma_2$ and \emph{any} correlation $\rho>-1$ the risk-parity
      weights are $w_1:w_2=1/\sigma_1:1/\sigma_2$ (inverse volatility).
  \item The same holds for $m$ assets with equal
      pairwise correlations, whatever the volatilities.
  \item If, in addition, all assets have the \emph{same Sharpe
      ratio}, this risk-parity portfolio is the tangency portfolio $\bm w_M$ of \cref{ch10:thm-tobin}.
\end{enumerate}
\end{proposition}
\begin{proof}
(i) If $w_1\sigma_1=w_2\sigma_2=:\kappa$ then $w_1(\Sigma\bm w)_1=w_1(w_1\sigma_1^2+w_2\rho\sigma_1\sigma_2)=\kappa^2(1+\rho)$ and
likewise for asset~2.

(ii) With $\Sigma=D C D$, $D=\diag(\sigma_i)$, $C=(1-\rho)I+\rho\bm 1\bm 1\T$, and
$w_i=\kappa/\sigma_i$, i.e.\ $D\bm w=\kappa\bm 1$: $w_i(\Sigma\bm w)_i=w_i\sigma_i(C D\bm w)_i=\kappa^2(C\bm 1)_i=\kappa^2(1+(m-1)\rho)$,
independent of $i$.

(iii) $\mu_i-r_0=s\sigma_i$ gives $\bm\eta=sD\bm 1$ and $\Sigma^{-1}\bm\eta=sD^{-1}C^{-1}\bm 1=
s\,D^{-1}\bm 1/(1+(m-1)\rho)$, proportional to $1/\sigma_i$.
\end{proof}

\begin{example}[60/40 versus risk parity, hypothetical numbers]\label{ch10:ex-rp}
Let equities have volatility $15\%$ and bonds $5\%$, uncorrelated, and let both have a Sharpe ratio of $0.3$
(excess returns $4.5\%$ and $1.5\%$). These numbers are \emph{not} from the course; they are chosen so that the
volatility ratio is $3$, which reproduces the $25/75$ equity/bond risk-parity split quoted in class
\ts{13}{16}{59:21}. Then:
\begin{center}
\begin{tabular}{@{}lrrrrr@{}}
\toprule
portfolio (eq./bond) & $\sigma_P$ & excess ret. & Sharpe & eq.\ risk & bond risk\\
\midrule
$60/40$ & $9.2\%$ & $3.3\%$ & $0.358$ & $95.3\%$ & $4.7\%$\\
risk parity $25/75$ & $5.3\%$ & $2.25\%$ & $0.424$ & $50\%$ & $50\%$\\
risk parity levered $1.74\times$ & $9.2\%$ & $3.9\%$ & $0.424$ & $50\%$ & $50\%$\\
\bottomrule
\end{tabular}
\end{center}
The $60/40$ portfolio takes $95\%$ of its risk from equities. Risk parity has a higher Sharpe ratio (here it is
the tangency portfolio by \cref{ch10:prop-rp}(iii), $0.3\sqrt2=0.424$) but a lower return; to get it back, one
\emph{levers} it: borrowing $74\%$ of the capital at the riskless rate brings the volatility up to that of the $60/40$
portfolio and the excess return to $3.9\%$ against $3.3\%$
\ts{13}{16}{56:12}\ts{13}{16}{59:21}\ts{24}{13}{40:00}. Values computed with the formulas above.
\end{example}
That is the ``levered risk parity'' of 2024 slide~19: a levered bond allocation to achieve an equal risk
contribution from equities and bonds, so that the leverage is what buys back the return. Because Sharpe ratios
are invariant to leverage (\cref{ch10:prop-sharpe}), the gain comes from choosing the point of maximal Sharpe ratio.

\begin{finance}[Limits of risk parity]
Risk parity uses \emph{no return forecasts}, only volatilities and correlations; this is attractive because,
a remark from the audience (apparently by Kempthorne, whom the transcript does not name) put it, returns cannot be
estimated accurately even with more data, whereas volatilities can (\cref{ch10:sec-est}) \ts{13}{16}{1:25:00}. But it does depend on a forecast of volatility from history: asked why risk-parity
funds lost heavily in May--June 2013, Xia explained that they overweight low-volatility bonds, which sold off when
the Fed chairman announced the tapering of quantitative easing \ts{13}{16}{1:14:26}. He then asked whether this
episode proves the approach wrong, or the 2008 crisis proves it right, and concluded that the evidence is
inconclusive: one is ``driving while looking in the rear-view mirror''
\ts{13}{16}{1:15:30}. Leverage carries liquidity risk, not treated in the lecture \ts{13}{16}{59:21}.
\end{finance}

\begin{intuition}[Leverage and beta]
Levering a portfolio by a factor $\ell$ multiplies its beta to every factor, and its volatility, by $\ell$, and
the excess return by $\ell$: the Sharpe ratio is unchanged, like the direction of a vector when it is rescaled.
For $\sigma_P=w_2\sigma_2$ in the two-asset special case, $\beta=\sigma_P/\sigma_2$
\ts{13}{16}{1:03:38}: ``a higher volatility of the portfolio means a higher beta to the risky asset''.
\end{intuition}

% ==================================================================
\section{Utility theory: why mean and variance?}\label{ch10:sec-utility}

Markowitz reduces a portfolio to its mean and variance. When is that legitimate? Kempthorne's answer is the
\emph{von Neumann--Morgenstern theory} of decisions under uncertainty: a rational agent has a utility function
$u$ of wealth and chooses the portfolio $P$ maximising the \emph{expected utility} $\E[u(W)]$ of the terminal
wealth $W=W_0(1+R_P)$ \ts{13}{14}{48:55}\ts{13}{14}{50:00}\ts{13}{14}{51:01}. Sensible preferences give
$u'>0$ (more is better) and $u''<0$ (decreasing marginal utility, risk aversion) \ts{13}{14}{52:01}. The
Arrow--Pratt coefficients are
\begin{align*}
  \lambda_A(W)&=-\frac{u''(W)}{u'(W)}\quad\text{(absolute risk aversion)},\\
  \lambda_R(W)&=-\frac{W\,u''(W)}{u'(W)}\quad\text{(relative)}.
\end{align*}
(Slide~19 of \file{lecnote14}; the same $\lambda_A$ appears in \cref{ch03:sec-capm}, where the CARA case is treated.)

\begin{table}[ht]
\centering
\begin{tabular}{@{}llccl@{}}
\toprule
& $u(W)$ & $\lambda_A(W)$ & $\lambda_R(W)$ & remark\\
\midrule
linear & $a+bW$, $b>0$ & $0$ & $0$ & risk neutral\\
quadratic & $W-\frac12\lambda W^2$, $W<1/\lambda$ & $\dfrac\lambda{1-\lambda W}$ & $\dfrac{\lambda W}{1-\lambda W}$ & mean--variance is exact\\
exponential & $1-e^{-\lambda W}$ & $\lambda$ & $\lambda W$ & CARA\\
power & $W^{1-\lambda}$, $0<\lambda<1$ & $\lambda/W$ & $\lambda$ & CRRA\\
logarithmic & $\ln W$ & $1/W$ & $1$ & CRRA, Kelly\\
\bottomrule
\end{tabular}
\caption{The utility functions of slide~20 of \file{lecnote14} \ts{13}{14}{54:08}, with their risk-aversion coefficients (computed here).
Quadratic utility has \emph{increasing} absolute risk aversion and is not increasing beyond $W=1/\lambda$.}\label{ch10:tab-utility}
\end{table}

\subsection{From expected utility to mean--variance}
\begin{proposition}[Three routes to mean--variance analysis]\label{ch10:prop-mv-utility}
\leavevmode
\begin{enumerate}[(i)]
  \item For any smooth $u$, expanding around $w_*=\E[W]$ gives to second order
      $\E[u(W)]\approx u(w_*)+\tfrac12u''(w_*)\Var W$, i.e.\ up to a positive factor and a constant,
      $\E[W]-\frac12\lambda_A(w_*)\Var W$ \ts{13}{14}{53:06}\ts{13}{14}{54:08}.
  \item For quadratic utility the expected utility is exactly $\mu_W-\frac\lambda2(\mu_W^2+\sigma_W^2)$, a function of the mean
      and variance of $W$ alone \ts{13}{14}{55:09}.
  \item If $W\sim N(\mu_W,\sigma_W^2)$, for \emph{any} increasing concave $u$, $\E[u(W)]$ is increasing in $\mu_W$ and
      decreasing in $\sigma_W^2$, so that only the pair $(\mu_W,\sigma_W)$ matters and the efficient frontier is the set of
      candidates for the optimum; for exponential utility, $\E[u(W)]=1-\exp\{-\lambda(\mu_W-\frac\lambda2\sigma_W^2)\}$.
\end{enumerate}
\end{proposition}
\begin{proof}
(i) Taylor's formula at $w_*$, with $\E[W-w_*]=0$: the first-order term drops, and $u''=-\lambda_Au'$.

(ii) Compute
$\E[W-\frac\lambda2W^2]$.

(iii) Write $\E[u(W)]=\int u(x)\varphi(x;\mu_W,\sigma_W^2)\dd x$; the Gaussian density satisfies the
heat equation $\partial_{\sigma^2}\varphi=\frac12\partial_x^2\varphi$ and $\partial_\mu\varphi=-\partial_x\varphi$, so
integrating by parts, $\partial_{\sigma^2}\E[u]=\frac12\E[u''(W)]<0$ and $\partial_\mu\E[u]=\E[u'(W)]>0$. For the exponential,
use the moment generating function of the normal (\cref{ch03:eq-normalmgf}).
\end{proof}
\begin{intuition}[Variance as the heat-equation time]
For a Gaussian, the variance plays the role of the time in the diffusion equation: $\E[u(W)]$ is $u$ smoothed by
a Gaussian of width $\sigma_W$, and smoothing a concave function can only lower its value.
\end{intuition}
So mean--variance analysis is the right thing to do if the utility is quadratic, \emph{or} if returns are Gaussian
(Kempthorne, 2013: ``the imagination that returns are Gaussian''), at least to first orders otherwise
\ts{13}{14}{55:09}\ts{13}{14}{56:11}. With exponential utility and normal returns, wealth
$W=W_0(1+\bm w\T\bm R)$ has certainty equivalent $W_0\bigl(1+\bm w\T\bm\mu-\frac\lambda2\bm w\T\Sigma\bm w\bigr)$ with
$\lambda=\lambda_AW_0$, exactly \emph{problem III} of \cref{ch10:sec-mvo}; with a riskless asset the solution is
$\bm w=\lambda^{-1}\Sigma^{-1}\bm\eta$, i.e.\ the fraction $m_{\bm w}=(\mu_M-r_0)/(\lambda\sigma_M^2)$ of the wealth in the
market portfolio (from \cref{ch10:thm-tobin}): the risk aversion fixes only \emph{how much}, in accord with
Tobin's separation.

\subsection{Higher moments}
When returns are not Gaussian one goes on in the Taylor series (slides 32--33 of \file{lecnote14})
\ts{13}{14}{56:11}:
\[
  \E[u(W)]=u(w_*)+\tfrac12u''(w_*)\Var W+\tfrac1{3!}u'''(w_*)\,m_3+\tfrac1{4!}u^{(4)}(w_*)\,m_4+\cdots,\qquad m_k=\E[(W-w_*)^k],
\]
with $m_3$ and $m_4$ the third and fourth central moments (the slides call them ``Skew'' and ``Kurtosis''). For
the utilities of \cref{ch10:tab-utility} (exponential, power, log) $u'''>0$ and $u^{(4)}<0$: investors like positive
skewness and dislike large fourth moments. The corresponding optimisation
$\max\ \E[R_P]-\lambda_1\Var R_P+\lambda_2m_3-\lambda_3m_4$ subject to $\bm 1\T\bm w=1$ has no closed form; the slides mention
multi-objective methods and \emph{polynomial goal programming}, and warn that skewness and kurtosis are hard
to estimate (they are sensitive to outliers and need large samples).

% ==================================================================
\section{Sizing: the Kelly criterion and the gain--loss ratio}\label{ch10:sec-kelly}

\subsection{Kelly}
The \emph{Kelly formula} is mentioned by Xia and, at the end of the lecture, by Kempthorne \ts{13}{16}{45:21}\ts{13}{16}{1:26:07}. Consider repeated
independent bets in which a stake gains $g$ per unit staked with probability $p$ and loses $l$ per unit with
probability $q=1-p$. Betting a fraction $f$ of the wealth each time multiplies the wealth by $1+fg$ or $1-fl$, and by the
law of large numbers (\cref{ch03:sec-limits}) $\frac1n\log W_n\to G(f)=p\log(1+fg)+q\log(1-fl)$ almost surely. The
fraction that maximises the long-run growth rate is
\begin{equation}\label{ch10:eq-kelly}
  f^*=\frac pl-\frac qg=\frac{pg-ql}{gl},
\end{equation}
since $G'(f)=pg/(1+fg)-ql/(1-fl)=0$ gives $pg(1-fl)=ql(1+fg)$, i.e.\ $f\,gl(p+q)=pg-ql$. For an even-money bet
($g=l=1$), $f^*=p-q=2p-1$: with probability $\frac12$ one bets nothing, with $p=1$ everything, and with $p=0$ one bets
everything \emph{on the other side}, as Xia says \ts{13}{16}{45:21}\ts{13}{16}{46:23}. It is expected-log utility
(\cref{ch10:tab-utility}), the ``Kelly criterion'' of Kelly, Shannon and Thorp, criticised by Samuelson in the
debate recalled in the case of \cref{ch03:sec-capm} and in \emph{Fortune's Formula}; the co-lecturer, speaking from the audience,
stresses that, unlike the single-period theory of this chapter, it deals with \emph{multi-period} investment
\ts{13}{16}{1:26:07}\ts{13}{16}{1:27:08}. For one asset with Gaussian excess return $\mu$ and variance
$\sigma^2$ per period the growth rate of a fraction $f$ is $r_0+f\mu-\frac12f^2\sigma^2$
(\cref{ch10:prop-growth}), maximised at $f^*=\mu/\sigma^2$, the Sharpe ratio divided by the volatility
(not discussed in class); Kelly is problem III with $\lambda=1$.

\subsection{Xia's gain--loss ratio}
Xia argues that volatility is a poor risk measure because it is symmetric; a long out-of-the-money call
\emph{wants} high volatility, a short option position wants low volatility \ts{24}{13}{47:28}, and the Sharpe ratio
inherits the flaw; the Sortino ratio separates the downside but does not tell how much to invest
\ts{24}{13}{48:30}. He proposes to summarise the distribution of the profit and loss $R$ of an investment
(per period, or of the terminal outcome) by the \emph{expected gain} and \emph{expected loss}, both positive numbers,
\[
  G=\E[R^+],\qquad L=\E[R^-],\qquad \text{so}\qquad G-L=\E[R],\quad G+L=\E[|R|] ,
\]
and to rank investments by the ratio
\begin{equation}\label{ch10:eq-gl}
  \frac{G-L}{G+L}=1-\frac{2L}{G+L}=\frac{\E[R]}{\E[|R|]}\in[-1,1],
\end{equation}
which he interprets directly as a position size (2024 slide~24) \ts{24}{13}{48:30}\ts{24}{13}{49:33}. In the
$(L,G)$ plane every investment is a point; the $45^\circ$ line is $G=L$ (zero expected return) and the
ratio is the slope of the ray from the origin: one wants points far above the diagonal
(slide~24: ``$G\gg L$'') and maximises $G$ for given $L$, which replaces the volatility budget by a
\emph{loss} budget \ts{24}{13}{50:37}\ts{24}{13}{51:37}. For a discrete payoff, $G$ and $L$ are sums of probability
times payoff over winning and losing outcomes; in the continuous case they are integrals of the payoff against the
density over $R>0$ and $R<0$, or against a target return other than $0$ \ts{24}{13}{50:37}. Estimation needs only
``half distributions'', no more work than a mean and a standard deviation \ts{24}{13}{52:41}\ts{24}{13}{53:47};
the gain and loss of a combination of two investments depends on how they move together and is found by
Monte Carlo simulation of the joint distribution \ts{24}{13}{54:54}. The \emph{2024 Investment Game} of \cref{ex:game}
is the empirical version: for the chosen stock, sum the daily gains and losses over two months and compute
$(G-L)/(G+L)$; picking a very volatile stock raises $G-L$ but also the denominator, so the ratio ranks
investments of different volatilities on a common scale \ts{24}{13}{51:37}.

\begin{proposition}[Properties of the ratio \eqref{ch10:eq-gl}; not proved in class]\label{ch10:prop-gl}
\leavevmode
\begin{enumerate}[(i)]
  \item For an even-money bet ($R=\pm1$ with probabilities $p,q$) the ratio is $p-q$, the Kelly fraction \eqref{ch10:eq-kelly};
      for $g\neq l$ the two differ (for $g=2$, $l=1$, $p=\frac12$: Kelly $f^*=0.25$, ratio $=1/3$). So ``the same as Kelly
      for binary bets'' (slide~24) holds for even-money bets, or as an approximation.
  \item If $R\sim N(m,s^2)$ the ratio equals $S/h(S)$ with $S=m/s$ and $h(S)=2\varphi(S)+S(2\Phi(S)-1)$, an increasing function
      of the Sharpe ratio: for symmetric returns it gives the same ranking as the Sharpe ratio.
  \item For skewed payoffs the two rankings differ: $R\equiv2$ w.p.\ $\frac12$ and $0$ w.p.\ $\frac12$ has mean $1$ and standard deviation $1$,
      so does $N(1,1)$, and the ratios are $1$ (never a loss) and $0.857$.
\end{enumerate}
\end{proposition}
\begin{proof}
(i) $G=p$, $L=q$, $G+L=1$. For the Kelly value use \eqref{ch10:eq-kelly}.

(ii) For $R=m+sZ$,
$\E[|R|]=s\bigl[2\varphi(S)+S(2\Phi(S)-1)\bigr]$ (fold the normal), and $h'(S)=2\Phi(S)-1$, so
$(S/h)'=(h-Sh')/h^2=2\varphi(S)/h^2>0$.

(iii) $G=1$, $L=0$; for $N(1,1)$, $h(1)=1.1666$ and $1/h(1)=0.857$.
\end{proof}

% ==================================================================


\section{Portfolio constraints}\label{ch10:sec-constr}

Real portfolios are subject to constraints, which Kempthorne lists (slides 22--25 of \file{lecnote14})
\ts{13}{14}{57:16}\ts{13}{14}{58:20}. With $\bm w$ the target portfolio, $\bm w_0$ the current one, $\bm x=\bm w-\bm w_0$ the
trade and $\bm w_B$ a benchmark portfolio:
\begin{itemize}
  \item \textbf{Long-only}: $w_j\ge0$. \textbf{Holding constraints}: $L_j\le w_j\le U_j$, for instance a cap
        on each position; in a medium-frequency equity strategy, a position must not exceed a day's trading
        volume, or it could not be liquidated \ts{13}{14}{58:20}.
  \item \textbf{Turnover}: $|x_j|\le U_j$ for each asset and $\sum_j|x_j|\le U_*$ overall, because one cannot
        trade arbitrarily much in one period.
  \item \textbf{Benchmark exposure}: $\sum_i|w_i-w_{B,i}|\le U_B$, to remain close to an index (S\&P~500,
        NASDAQ-100, Russell) while trying to beat it \ts{13}{14}{59:30}.
  \item \textbf{Tracking error}: $\mathrm{TE}_P=R_P-R_B=(\bm w-\bm w_B)\T\bm R$, so
        $\Var(\mathrm{TE}_P)=(\bm w-\bm w_B)\T\Sigma(\bm w-\bm w_B)\le\bar\sigma^2_{\mathrm{TE}}$
        \ts{13}{14}{1:00:31}.
  \item \textbf{Factor constraints}: for a factor model $R_{i,t}=\alpha_i+\sum_k\beta_{i,k}f_{k,t}+\epsilon_{i,t}$
        (\cref{ch07:sec-lfm}), the exposure to factor $k$ is $\sum_i\beta_{i,k}w_i$ and one imposes
        $|\sum_i\beta_{i,k}w_i|\le U_k$; many strategies \emph{neutralise} all exposures, $\sum_i\beta_{i,k}w_i=0$
        \ts{13}{14}{1:01:33}\ts{13}{14}{1:02:33}.
  \item \textbf{Minimum transaction or holding size} and \textbf{integer constraints}: equities trade in round lots
        and one share of a stock worth several hundred dollars is a significant amount compared with one of a
        \$50 stock, which matters for small portfolios but not for large ones \ts{13}{14}{1:02:33}\ts{13}{14}{1:03:33}.
\end{itemize}
Constraints of the first four types are linear or quadratic in $\bm w$: slide~25 writes general linear constraints
$A_w\bm w\le\bm u_w$, $A_x\bm x\le\bm u_x$, $A_B(\bm w-\bm w_B)\le\bm u_B$ and quadratic ones $\bm w\T Q_w\bm w\le q_w$,
$\bm x\T Q_x\bm x\le q_x$, $(\bm w-\bm w_B)\T Q_B(\bm w-\bm w_B)\le q_B$. The optimisation is then the same as before with
additional Lagrange multipliers, that is a \emph{convex quadratic programme} \ts{13}{14}{1:03:33}. Integer constraints make
it non-convex.

\begin{proposition}[Structure of the constrained solution; not discussed in class]\label{ch10:prop-constr}
Consider the problem I with additional linear inequality constraints.
\begin{enumerate}[(i)]
  \item The optimal value $\sigma_0^2(\mu_0)$ is
      never smaller than the unconstrained one \eqref{ch10:eq-frontier}, and it increases as constraints are added or tightened.
  \item The Karush--Kuhn--Tucker conditions with a fixed set of active constraints are \emph{linear} in
      $(\bm w,\text{multipliers})$ and in $\mu_0$; therefore the optimal weights are piecewise affine in the target return
      $\mu_0$, and the pieces end when a new constraint becomes active or a multiplier changes sign.
\end{enumerate}
\end{proposition}
\begin{proof}
(i) A smaller feasible set has a larger minimum.

(ii) With the active constraints imposed as equalities,
the first-order conditions are the linear system of \cref{ch10:thm-mvo} enlarged by the active rows, whose
solution is affine in the right-hand side $(\mu_0,1,\text{bounds})$; strict convexity makes it unique.
\end{proof}
That is exactly what the plots of the case study show (\cref{ch10:sec-case}): weights that grow along straight
lines in the target return, kinks when a weight reaches its cap, and new assets entering. Below the first
kink (only the constraints $w_j\ge0$ active) the solution is a \emph{scaled} fixed portfolio: ``they enter in the same
proportion'', the de-levered tangency portfolio of \cref{ch10:thm-tobin} computed on the assets that enter, i.e.\ the
\emph{long-only} tangency portfolio, not the unconstrained one, which has short positions (Case Study~6, and Kempthorne
\ts{13}{14}{1:09:02}).

\begin{finance}[Benchmarked management]
A manager measured against an index does not solve problem~I but the same quadratic problem in the \emph{active
weights} $\bm x_B=\bm w-\bm w_B$: maximise $\bm x_B\T\bm\mu-\frac\lambda2\bm x_B\T\Sigma\bm x_B$, the tracking-error
variance being the risk. The efficient frontier is then drawn in the plane (tracking error, active return), and its
slope is the information ratio. Factor-neutrality constraints remove market, sector or style bets so that the
return is pure stock selection.
\end{finance}

% ==================================================================
\section{Estimation error and the fragility of the optimiser}\label{ch10:sec-est}

Everything above takes $\bm\mu$ and $\Sigma$ as known, ``the problem where we assume basically certainty''
\ts{13}{14}{2:07}. In practice they are estimated, and the estimates are noisy \ts{13}{14}{20:08}. Xia's first
objection to modern portfolio theory is exactly that: the capital market assumptions of return,
volatility and correlation are not reliable, and the mean--variance solution is very sensitive to them, is
``unstable'' and has ``unlimited'' solutions unless bounded by artificial constraints
\ts{24}{13}{45:22}\ts{24}{13}{46:27}.

\subsection{How noisy is a sample mean?}
Suppose returns are i.i.d.\ Gaussian with annual mean $\mu$ and volatility $\sigma$ and are observed for $T$ years, at any
frequency. The sample mean has standard error $\sigma/\sqrt T$ \emph{independent of the sampling frequency}, since only
the total time span matters, whereas the sample volatility, computed from $n$ increments, has relative standard error
$1/\sqrt{2n}$ (\cref{ch03:sec-limits}, \cref{ch:volatility}). For 500 weekly returns ($T=9.6$ years, $n=500$):
\begin{center}
\begin{tabular}{@{}lccc@{}}
\toprule
& true value & standard error & sample values in \file{Smltn_TwoAst}\\
\midrule
mean of asset 1 & $15\%$ & $\sigma_1/\sqrt T=6.4\%$ & $0.085$--$0.21$\\
mean of asset 2 & $25\%$ & $\sigma_2/\sqrt T=9.7\%$ & $0.177$--$0.343$\\
vol.\ of asset 1 & $20\%$ & $3.2\%$ of itself ($0.6\%$) & $0.191$--$0.200$\\
vol.\ of asset 2 & $30\%$ & $3.2\%$ of itself ($0.9\%$) & $0.287$--$0.291$\\
\bottomrule
\end{tabular}
\end{center}
This is what Kempthorne shows by simulation: sample means may be far from the population values while volatilities
and correlations are estimated much better \ts{13}{14}{21:08}\ts{13}{14}{22:11}, and what the audience remark of
2013 L16 says about returns versus volatilities \ts{13}{16}{1:25:00}. In finance ``more data'' means longer periods
(the mean does not improve by sampling faster), while regime changes limit how far back one can go: the choice of
the estimation period is itself an issue \ts{13}{14}{1:20:06}. With the five-year sector-ETF sample of
Case Study~6 the mean of a sector with $\sigma=20\%$ has a standard error of $9\%$, larger than most differences
between the sector means ($10\%$ to $24\%$).

\subsection{What the optimiser does with the errors}
Because $\bm w\propto\Sigma^{-1}\bm\eta$ and highly correlated assets make $\Sigma$ nearly singular, errors in
the means are amplified in the weights. On the rounded nine-ETF inputs of \cref{ch10:tab-etf} (computed here), the
tangency portfolio for $r_0=0$ is extremely leveraged, and \emph{very small changes of the means, of one
percentage point per asset}, change it a lot:
\begin{center}
\footnotesize
\begin{tabular}{@{}lrrrrrrrrrr@{}}
\toprule
weights & XLB & XLV & XLP & XLY & XLE & XLF & XLI & XLK & XLU & $\sum|\Delta w|$\\
\midrule
base means & $-0.32$ & $0.59$ & $0.58$ & $1.49$ & $-0.06$ & $-0.36$ & $-0.45$ & $-0.10$ & $-0.35$ & \\
XLY mean $+1$ pp & $-0.34$ & $0.63$ & $0.53$ & $1.73$ & $-0.05$ & $-0.39$ & $-0.54$ & $-0.18$ & $-0.38$ & $0.58$\\
$\pm1$ pp draw A & $-0.32$ & $0.75$ & $0.28$ & $1.19$ & $-0.09$ & $-0.35$ & $-0.29$ & $0.07$ & $-0.24$ & $1.24$\\
$\pm1$ pp draw B & $-0.57$ & $0.68$ & $0.42$ & $2.07$ & $-0.05$ & $-0.55$ & $-0.19$ & $-0.33$ & $-0.48$ & $1.91$\\
\bottomrule
\end{tabular}
\end{center}
(Draw A is $+,+,-,-,+,+,+,+,+$ pp and draw B $-,-,-,+,-,-,+,-,-$ pp for XLB to XLU; the weights sum to $1$.) For comparison,
the sum of absolute weights is $4.3$: a $\pm1$pp error in each mean, well inside the statistical
error, moves $29\%$ and $44\%$ of that gross exposure (draws A and B). In this respect the caps of \cref{ch10:sec-constr} act as
regularisation, ``artificial constraints'' that bound the solutions \ts{24}{13}{45:22}. The minimum-variance portfolio
$\bm w_g$ does not use $\bm\mu$ at all, which is why it is popular.

\subsection{Remedies}
The slides list the standard alternatives \ts{13}{14}{1:20:06} (slide~27 of \file{lecnote14}): the sample means and
covariance are the least-squares, unbiased and, under Gaussian assumptions, maximum likelihood estimates, but
one may instead use \emph{exponentially weighted moving averages} (\cref{ch:volatility}), \emph{dynamic factor
models}, and optimisation with a simple model for $\Sigma$, notably Sharpe's \emph{single-index model}
$R_i=\alpha_i+\beta_iR_m+\epsilon_i$ (with uncorrelated $\epsilon_i$), which gives $\Sigma=\sigma_m^2\bm\beta\bm\beta\T+
\diag(\sigma^2_{\epsilon,i})$, or the \emph{constant correlation model}, $\Sigma_{ij}=\bar\rho\sigma_i\sigma_j$. A factor model
estimates $2m+1$ parameters instead of $m(m+1)/2$ (for $m=500$ stocks, $1001$ instead of $125{,}250$), and ``results in
more precise inputs to the optimisation'' \ts{13}{14}{1:20:06}; see \cref{ch07:sec-lfm,ch:regularized} (shrinkage,
regularised risk models) and \cref{ch05:sec-shrink}. The returns are the weak point: forecasting them is the actual
job of an active manager, as an audience member observed in 2013 \ts{13}{16}{1:22:59}\ts{13}{16}{1:24:00}.

\begin{finance}[Why the frontier is estimated with errors: ``error maximisation'']
An optimiser puts most weight on the assets whose estimated means are too high and whose estimated variances are too
low, which is precisely where estimation errors are largest: the out-of-sample performance of the
mean--variance solution is typically worse than the in-sample frontier suggests (not discussed in class).
\end{finance}

% ==================================================================
\section{Other risk measures}\label{ch10:sec-risk}

Volatility is only one measure of risk, and Kempthorne closes 2013 L14 by showing how the machinery extends
\ts{13}{14}{1:20:06}\ts{13}{14}{1:21:13} (slides 29--33). Let $R_P=\bm w\T\bm R$.
\begin{itemize}
  \item \textbf{Mean absolute deviation}, $\mathrm{MAD}=\E[|\bm w\T(\bm R-\bm\mu)|]$, leads to a linear programme with linear or
        quadratic constraints. Where the variance is the natural scale for Gaussian returns, MAD is natural for
        heavier-tailed, e.g.\ exponential (Laplace), returns \ts{13}{14}{1:21:13}.
  \item \textbf{Semi-variance}, $\E[\min(R_P-\E[R_P],0)^2]$, the downside variance, introduced by Markowitz himself:
        it controls downside rather than upside risk \ts{13}{14}{1:21:13}. Its square root, relative to a target,
        is the denominator of the Sortino ratio \ts{24}{13}{48:30}.
  \item \textbf{Value at risk}, $\mathrm{VaR}_{1-\epsilon}(R_P)=\min\{r:\Prob(R_P\le-r)\le\epsilon\}$ with $\epsilon=0.05,0.01,0.001$,
        popularised by JP~Morgan's RiskMetrics \ts{13}{14}{1:22:18}. It is standard for cash instruments (stocks, bonds) but,
        as Kempthorne says, for derivatives with nonlinear payoffs ``things get complicated'' \ts{13}{14}{1:24:25}.
  \item \textbf{Conditional VaR} (expected shortfall), $\mathrm{CVaR}_{1-\epsilon}(R_P)=\E[-R_P\mid-R_P\ge\mathrm{VaR}_{1-\epsilon}(R_P)]$, the average
        loss beyond the VaR, which is where VaR is silent; expected to enter bank regulation
        \ts{13}{14}{1:23:22}. The optimisation of CVaR is a linear programme (Rockafellar and Uryasev, 2000, cited on the slide).
\end{itemize}
Definitions with the loss convention, the coherence axioms and the estimation of VaR and ES are in
\cref{ch08:sec-riskmeasures,ch:var}; here only the portfolio-theory consequences are needed.

\begin{proposition}[VaR is not subadditive; variance is not monotone; not stated in class]\label{ch10:prop-var}
Two independent positions each lose $100$ with probability $4\%$ and otherwise $0$. At the $95\%$ level,
$\mathrm{VaR}=0$ for each (the loss occurs with probability $4\%<5\%$), but for the portfolio of both
$\Prob(\text{loss}>0)=1-0.96^2=7.84\%>5\%$, so $\mathrm{VaR}_{95\%}=100>0+0$: diversification \emph{increases} VaR.
The expected shortfalls are $80$ for each and $103.2$ for the portfolio ($\le160$), as CVaR is subadditive. Likewise
the variance violates monotonicity: $R=0$ and $R'\in\{0,1\}$ with probability $\frac12$ each satisfy $R\le R'$ but
$\Var R=0<\Var R'=\frac14$, whereas monotonicity requires $s(R)\ge s(R')$.
\end{proposition}
\noindent (Checked by direct computation: $\mathrm{ES}_{95\%}(\text{portfolio})=(0.16\%\cdot200+4.84\%\cdot100)/5\%=103.2$.)

\begin{coursenote}[Erratum: translation invariance on slide~31 of \file{lecnote14}]
The slide, written for \emph{returns}, defines a coherent measure by: monotonicity ($R_P\le R_{P'}$ w.p.~1 $\Rightarrow s(R_P)\ge s(R_{P'})$),
subadditivity, positive homogeneity, and ``translational invariance $s(R_P+a)\le s(R_P)-a$''. The standard axiom
(Artzner--Delbaen--Eber--Heath 1999) is an \emph{equality}, $s(R_P+a)=s(R_P)-a$: adding a sure gain $a$ lowers the risk by
exactly $a$ (in the loss convention of \cref{ch08:def-coherent}, $\rho(L+a)=\rho(L)+a$). The slide's ``coherent'' verdicts
(variance not, VaR not, CVaR yes) are correct \ts{13}{14}{1:24:25}.
\end{coursenote}

\begin{finance}[Expected loss, worst loss, stop losses, drawdowns]
Asked for the best risk measure, Xia answers the \emph{expected loss} $L$, arguing that gains and losses are
usually not symmetric; one should also know the worst case but cannot size a portfolio on it (one would never invest) \ts{24}{13}{1:14:54}\ts{24}{13}{1:15:58}. Stop losses on trading platforms are set tighter than the
expected loss (cut the position by half after a small loss, another half after the next), which gives the
trader an asymmetric, option-like payoff and forces a distribution of the trades that fits it; and hedge-fund managers,
who take a share of profits but not of the losses, hold ``a free call option'' on the fund
\ts{24}{13}{1:16:58}\ts{24}{13}{1:18:00}\ts{24}{13}{1:19:01}. The loss tolerance of \cref{ch10:sec-problem} is often expressed as a
\emph{maximum drawdown} $\max_{s\le t}\bigl(\max_{u\le s}V_u-V_s\bigr)/\max_{u\le s}V_u$, the largest peak-to-trough fall of
the portfolio value $V$ (not defined in class; see the drawdowns of the concentrated portfolio in \cref{sec:ewport}).
\end{finance}

% ==================================================================
\section{Beyond Markowitz: crowds and power laws (2024)}\label{ch10:sec-crowd}

In 2024 Xia devotes the second half of the lecture to three ideas from his own research (slides marked
``sections 4--7 highlight my research work''): the loss of meaning of volatility as risk (\cref{ch10:sec-kelly}), the
unpredictability of markets because of \emph{crowd behaviour}, and the origin of \emph{power-law} distributions.
They are not tested, and are heuristic, but they answer his ``MPT issue number 2'': how do we predict the future from
the past \ts{24}{13}{54:54}? The scientific method (observe, quantify, extract a pattern, build a model, predict,
repeat and calibrate, control) applies to physics and engineering, but in finance the historical patterns may not
be repeatable because humans adapt and herd \ts{24}{13}{56:01}\ts{24}{13}{57:01}. In 2013 the same point was
made with videos and no formulas: a bridge whose walkers synchronise their steps, and a bank of metronomes on a
moving plate that synchronise \ts{13}{16}{1:05:51}\ts{13}{16}{1:08:02}; in 2024 with a video of birds
\ts{24}{13}{58:03} and the London Millennium Bridge \ts{24}{13}{59:05}. The lesson for portfolio management: an
investor who finds the ``optimal'' strategy finds that everyone else implements it too, and the collective system
is not optimal but fragile; individual managers also have an incentive to herd, since underperforming the peers is
punished more than being wrong with the crowd (career risk)
\ts{13}{16}{1:10:12}\ts{13}{16}{1:16:33}\ts{13}{16}{1:17:37}. Participants change the market they are trying to optimise \ts{13}{16}{1:12:19}.

\subsection{A feedback model of crowd behaviour}
Slide~28 (drawn on slide~34 with heterogeneous agents) describes the interaction loop of $N$ agents, an external force
$E$ and an observation $O$ (the price) \ts{24}{13}{1:00:05}:
\begin{equation}\label{ch10:eq-crowd}
  \dd S_i=C_i\dd E+B_i\dd O+\varepsilon_i,\qquad \dd O=\sum_{i=1}^NA_i\dd S_i ,
\end{equation}
where $S_i$ is the action of agent $i$, $A_i$ its power (influence on the observation), $C_i$ its response to the
external force, $B_i$ its \emph{reactive function} (how strongly it reacts to the observation) and $\varepsilon_i$
noise. In the rational state $B_i$ takes low values $\{B_L\}$; when $\dd O$ moves a lot, agents switch to a reactive
state with high values $\{B_H\}$, and $N_H$ counts them \ts{24}{13}{1:01:08}. The \emph{order parameter}
\[
  R(t)=\frac{\bigl|\sum_i\dd S_i(t)\bigr|}{\sum_i|\dd S_i(t)|}\in[0,1]
\]
is $1$ if all agents act in the same direction (synchronised) and near $0$ if they are not \ts{24}{13}{1:01:08}.
The slides plot $R$ against $N_H/N$ (rising steeply from $\approx0.1$ at $N_H=0$ to $1$ at $N_H=N$) and against the noise
amplitude (falling from $1$ to near $0$: noise desynchronises, as in coupled oscillators)
\ts{24}{13}{1:02:15}, and a simulated ``bubble'': an external force that rises, then falls and oscillates makes the
number of reactive agents rise and fall and the observation follow; if the force stays high forever, all agents end up
reactive and the observation explodes, a \emph{tipping point} when $A\times B_{\max}>1$ (slide~30)
\ts{24}{13}{1:03:21}\ts{24}{13}{1:04:28}.

\begin{proposition}[Mean-field reading of the tipping point; my derivation, not shown in class]\label{ch10:prop-tipping}
Suppose all agents react with the same $B_i=B$. Substituting the first equation of \eqref{ch10:eq-crowd} into the second,
\begin{align*}
  \dd O&=\sum_iA_iC_i\dd E+B\Bigl(\sum_iA_i\Bigr)\dd O+\sum_iA_i\varepsilon_i,\quad\text{so}\\
  \dd O&=\frac{N}{1-\bar A B},\quad N=\sum_iA_iC_i\dd E+\sum_iA_i\varepsilon_i,\quad \bar A=\sum_iA_i .
\end{align*}
The response of the price to any external push or noise is amplified by the factor $1/(1-\bar AB)$, which is finite
and positive iff $\bar AB<1$ and diverges as $\bar AB\to1$: with $\bar A B\ge1$ the loop gain exceeds $1$ and the
feedback is unstable. The switch from $B_L$ to $B_H$ (a higher $B$ when many agents are reactive) moves the system
towards the threshold, which is the tipping point of slide~30.
\end{proposition}
\noindent This is the analogue of the gain of a linear amplifier with positive feedback (Barkhausen criterion), or of the
critical coupling of a mean-field ferromagnet, $\chi=1/(1-J\bar A)$, the susceptibility diverging at the critical
temperature. The lecture presents the mechanism qualitatively; the numerical scenarios on the slides are simulations
whose parameters are not given.

\subsection{Power laws}
Why are wealth, city sizes, venture-capital returns or viral posts so unequal, unlike heights?
\ts{24}{13}{1:05:33}\ts{24}{13}{1:06:33}. The $80$--$20$ rule, the Matthew effect and ``winner takes all'' describe distributions
that are \emph{scale-free}: the top $20\%$ of the people own $80\%$ of the wealth, and within them again the top $20\%$
own $80\%$, and so on. In symbols $p(bx)=g(b)p(x)$ for all $b>0$ holds (with smooth $p$) iff $p$ is a power law, $p(x)=Cx^{-a}$
(with $g(b)=b^{-a}$; the ``only if'' is the standard functional-equation argument, $p(bx)=g(b)p(x)$ forces $g(bb')=g(b)g(b')$) \ts{24}{13}{1:05:33}.
The three usual descriptions are related as follows (slide~33) \ts{24}{13}{1:09:39}\ts{24}{13}{1:10:39}:
\begin{align*}
  \text{Pareto (tail):}&\quad \Prob[X>x]\sim x^{-k}, & \text{Zipf (rank $r$):}&\quad \E[x_{(r)}]\sim r^{-b},\\
  \text{density:}&\quad f(x)=Cx^{-a}, & a&=1+k=1+\frac1b .
\end{align*}
(Differentiating the tail gives $a=1+k$; the $r$-th largest of $N$ satisfies $r/N=\Prob[X>x_{(r)}]$, so $x_{(r)}\sim r^{-1/k}$ and $b=1/k$.)
A Gaussian has no such structure; heights are Gaussian because they are set once, whereas in social
phenomena ``the rich get richer'': agents have \emph{different powers} $A(i,t)$ and gain power by betting in the
right direction of the observation, so the power of the crowd concentrates in a few ``super agents'' (governments and central banks, the Fed being the example named in class; in the old days, some hedge funds) whose actions dominate the system
\ts{24}{13}{1:08:38}\ts{24}{13}{1:10:39}\ts{24}{13}{1:12:49}. He also mentions entropy and stability of such systems as a topic
``beyond the scope''.

\begin{example}[The $80$--$20$ rule and the Pareto exponent]\label{ch10:ex-pareto}
If wealth has the tail $\Prob[X>x]=(x_m/x)^k$ with $k>1$, the top fraction $q$ of the population holds the fraction
$q^{1-1/k}$ of the total wealth (integrate $x$ over the tail). The $80$--$20$ rule, $0.2^{1-1/k}=0.8$, gives
$k=\ln5/\ln4=1.161$, $a=2.16$, $b=0.861$; and the top $20\%$ of the top $20\%$ hold $80\%$ of the $80\%$, i.e.\ $64\%$ of all
wealth (top $4\%$), the self-similarity of the lecture. Note $k$ close to $1$: the mean is barely finite
(for $k\le1$ it is infinite).
\end{example}

\begin{finance}[What the crowd material means for practice]
Xia's conclusions of the 2024 lecture (slide~35) \ts{24}{13}{1:11:45}: portfolio management is about \emph{sizing}; clarify
objectives and know the loss tolerance; diversification helps but requires rebalancing; volatility is not risk and the
Sharpe ratio captures neither skewness nor sizing (use the gain--loss ratio); capital market assumptions are
not reliable; crowd behaviour can push markets to extremes; watch the powerful super agents; understand the key drivers of the
portfolio. The 2013 version ends similarly: ``solving the problem'' is not done, since \emph{when you participate
in the market you change the market}, and quantitative finance ``is not like solving physics problems''
\ts{13}{16}{1:11:12}\ts{13}{16}{1:12:19}. A questioner asked whether this is game theory; Xia agreed that it has
a lot to do with it, but that the markets have no clearly defined rules \ts{13}{16}{1:13:22}. Whether a
computer can beat a human in investment, he ``cannot confidently say it will not happen''; a colleague added that
managers also manage people \ts{13}{16}{1:18:37}.
\end{finance}


% ==================================================================
\section{Case study: simulation and US sector ETFs (2013 Case Study~6)}\label{ch10:sec-case}

\begin{casestudy}[\file{CaseStudy6.pdf}, \file{Smltn_TwoAst.pdf}, \file{ETF_prid*.pdf}]
\textbf{Goal.} Kempthorne's live demonstration of the theory in R \ts{13}{14}{8:32}\ts{13}{14}{1:04:54}:
\begin{itemize}
  \item the feasible set of two
      assets for several correlations;
  \item mean--variance optimal portfolios of the nine US sector ETFs, with a riskless
      asset, with the constraint that no asset exceeds $30\%$ or $15\%$ of the capital.
\end{itemize}

\textbf{Part 1: two-asset simulation.} Two Gaussian assets with $\mu=(0.15,0.25)$, $\sigma=(0.20,0.30)$ (see the erratum
below), 500 weekly returns simulated for $\rho=-0.8,-0.4,0,0.4,0.8$. Each page of \file{Smltn_TwoAst} has four
panels: cumulative returns of the two assets (blue: asset~2, green: asset~1); the feasible set of
\cref{ch10:sec-two}; the scatter plot of the weekly returns; and the same cumulative returns with the cumulative return
of the minimum-variance portfolio (cyan) \ts{13}{14}{8:32}\ts{13}{14}{10:44}. In the PDF as distributed only the two asset
points are drawn in the second panel and the scatter panel is blank, so the curves of \cref{ch10:fig-two} were recomputed.
The sample statistics printed above the panels are:
\begin{center}
\small
\begin{tabular}{@{}lrrrrr@{}}
\toprule
$\rho$ & $-0.8$ & $-0.4$ & $0$ & $0.4$ & $0.8$\\
\midrule
sample means & $0.115,\ 0.341$ & $0.144,\ 0.343$ & $0.183,\ 0.334$ & $0.210,\ 0.321$ & $0.085,\ 0.177$\\
sample vols & $0.191,\ 0.291$ & $0.194,\ 0.291$ & $0.200,\ 0.289$ & $0.200,\ 0.287$ & $0.196,\ 0.287$\\
sample corr.\ & $-0.784$ & $-0.359$ & $0.031$ & $0.395$ & $0.791$\\
\bottomrule
\end{tabular}
\end{center}
Sample volatilities and correlations are close to the population values; the sample means are off by about one
standard error (\cref{ch10:sec-est}); for $\rho=0.8$ both are well below the population values ($0.085$ and $0.177$ against $0.15$ and $0.25$), an unlucky
draw, and the cumulative-return curves of the last panel look very different from the others. Lessons observed live: with
$\rho\le0.4$ the minimum-variance portfolio has a return and a volatility better than asset~1 alone; with
$\rho=-1$ the variance would be zero \ts{13}{14}{11:47}\ts{13}{14}{17:06}; every portfolio to the left of the minimum-variance point is
suboptimal \ts{13}{14}{14:57}.

\textbf{Part 2: nine sector ETFs.} The SPDR sector funds XLB (materials), XLV (health care), XLP (consumer staples),
XLY (consumer discretionary), XLE (energy), XLF (financials), XLI (industrials), XLK (technology) and XLU (utilities), for two
periods: 2009--2013 (``periodA'', up to a week before the lecture; $251$ observations on the time axis, i.e.\ weekly
data over about $4.8$ years, an inference: the frequency is not stated) and 2003--2006 (``periodB'', $209$ observations).
\cref{ch10:tab-etf} gives the annualised return and volatility of each. The correlations are high, with the
average pairwise correlation $0.74$ in 2009--2013 (range $0.57$--$0.93$) and $0.55$ in 2003--2006 (range
$0.16$--$0.85$): there is much less diversification available in the aftermath of the crisis.

\textbf{Method.} For a grid of $101$ target returns (``target return index'' in the plots) the R script minimises the variance
of a portfolio of the nine ETFs and the riskless asset (return $\approx0$: the frontier of plot~4 starts at the origin),
subject to full investment, $w_j\ge0$ and $w_j\le30\%$ or $15\%$. It prints the allocations for three target volatilities and
draws six plots: \emph{1}~risk--return scatter with the riskless asset; \emph{2}~optimal allocations versus target
return (lines); \emph{3}~the same as a stacked bar chart (dark red at the top is cash); \emph{4}~the efficient frontier over
the ETFs; \emph{5}~the frontier with the target portfolio for a $10\%$ volatility (circle); \emph{6}~cumulative returns of the
ETFs with that optimal portfolio in bold blue \ts{13}{14}{1:04:54}--\ts{13}{14}{1:11:11}.
\end{casestudy}

\begin{table}[ht]
\centering
\begin{tabular}{@{}llrrrr@{}}
\toprule
& & \multicolumn{2}{c}{2009--2013} & \multicolumn{2}{c}{2003--2006}\\
\cmidrule(lr){3-4}\cmidrule(lr){5-6}
ETF & sector & return & vol. & return & vol.\\
\midrule
XLB & materials & 0.16 & 0.24 & 0.16 & 0.18\\
XLV & health care & 0.17 & 0.15 & 0.07 & 0.12\\
XLP & consumer staples & 0.15 & 0.12 & 0.08 & 0.09\\
XLY & consumer discretionary & 0.24 & 0.21 & 0.14 & 0.14\\
XLE & energy & 0.15 & 0.25 & 0.26 & 0.20\\
XLF & financials & 0.13 & 0.33 & 0.15 & 0.13\\
XLI & industrials & 0.18 & 0.23 & 0.15 & 0.14\\
XLK & technology & 0.18 & 0.19 & 0.12 & 0.18\\
XLU & utilities & 0.10 & 0.16 & 0.19 & 0.13\\
\bottomrule
\end{tabular}
\caption{Annualised return and volatility of the nine sector ETFs printed in \file{CaseStudy6}. (The correlation
matrices are printed there as well; the computations of \cref{ch10:ex-etf-unc} and of this section use them.)}\label{ch10:tab-etf}
\end{table}

\paragraph{Results (2009--2013).} With a $30\%$ cap and a target volatility of $9.9\%$ the printed optimum is XLV $15.7\%$,
XLP $30\%$, XLY $22.6\%$, cash $31.7\%$, for a return of $12.4\%$; for a target volatility of $15.3\%$ it is fully
invested: XLV $30\%$, XLP $22.7\%$, XLY $30\%$, XLK $17.3\%$, for a return of $18.7\%$; for $0.9\%$ volatility it
holds $93.5\%$ in cash. The comments of the case study, confirmed by the plots \ts{13}{14}{1:07:59}--\ts{13}{14}{1:11:11}:
\begin{itemize}
  \item As the target return rises from zero only XLY, XLP and XLV enter, in \emph{fixed proportions} (the de-levered
        long-only tangency portfolio, \cref{ch10:prop-constr}; the unconstrained one of \cref{ch10:sec-est} is very
        different); allocations grow linearly (plot~2).
  \item The cap binds first for consumer staples; XLY and XLV are increased; when all three are at $30\%$ technology
        (higher return) enters and the allocations to XLY and XLP fall to make room. No allocation is ever given to financials.
  \item In the efficient frontier of plot~4, every ETF except consumer discretionary is dominated by an optimal
        portfolio, and the frontier shows the CML-like straight segment up to the first kink.
  \item With a $15\%$ cap the constraints bind sooner: XLV, XLP, XLY, XLK, then XLU and XLI, and finally XLB; the frontier is
        \emph{below} the $30\%$ one (\cref{ch10:prop-constr}(i)), because to reach a higher return one has to
        hold assets with a lower return per unit of risk \ts{13}{14}{1:12:14}\ts{13}{14}{1:14:24}.
  \item In 2003--2006 the optimal portfolios are completely different: utilities, energy, financials, consumer
        staples and industrials, XLU being the first to hit the cap, and technology never enters ($15\%$ cap:
        XLP, XLY, XLE, XLF, XLI, XLU at the cap, XLB in addition at the top). The ranking of sectors in one period says
        little about the next.
\end{itemize}

\begin{figure}[ht]
\centering
\begin{tikzpicture}
\begin{axis}[width=0.9\linewidth,height=7.4cm,xlabel={volatility},ylabel={expected return},
  xmin=0.05,xmax=0.36,ymin=0.04,ymax=0.32,xtick={0.1,0.15,0.2,0.25,0.3,0.35},scaled ticks=false,/pgf/number format/fixed,grid=major,grid style={black!10},
  legend style={font=\footnotesize,at={(0.5,1.03)},anchor=south,legend columns=3},
  tick label style={font=\small},label style={font=\small}]
  \addplot[thick,thmcol,dashed] coordinates {(0.05,0.0927) (0.16,0.297)};
  \addplot[very thick,defcol] coordinates {(0.005,0.0065) (0.02,0.0258) (0.04,0.0516) (0.06,0.0774) (0.08,0.1030) (0.10,0.1281) (0.12,0.1530) (0.13,0.1654) (0.14,0.1767) (0.15,0.1867) (0.16,0.1911)};
  \addplot[very thick,intcol] coordinates {(0.005,0.0065) (0.02,0.0258) (0.04,0.0515) (0.06,0.0765) (0.08,0.0991) (0.10,0.1183) (0.12,0.1350) (0.13,0.1433) (0.14,0.1514) (0.15,0.1590) (0.16,0.1664)};
  \legend{unconstrained CML (Sharpe $1.85$),cap $30\%$,cap $15\%$}
  \addplot[only marks,mark=*,mark size=1.8pt,black,forget plot] coordinates {(0.24,0.16) (0.15,0.17) (0.12,0.15) (0.21,0.24) (0.25,0.15) (0.33,0.13) (0.23,0.18) (0.19,0.18) (0.16,0.10)};
  \node[font=\scriptsize,anchor=north east] at (axis cs:0.236,0.157) {XLB};
  \node[font=\scriptsize,anchor=south east,fill=white,fill opacity=0.8,inner sep=1pt] at (axis cs:0.146,0.173) {XLV};
  \node[font=\scriptsize,anchor=north west] at (axis cs:0.123,0.146) {XLP};
  \node[font=\scriptsize,anchor=south] at (axis cs:0.21,0.243) {XLY};
  \node[font=\scriptsize,anchor=north] at (axis cs:0.25,0.147) {XLE};
  \node[font=\scriptsize,anchor=north] at (axis cs:0.33,0.127) {XLF};
  \node[font=\scriptsize,anchor=south west] at (axis cs:0.226,0.185) {XLI};
  \node[font=\scriptsize,anchor=south] at (axis cs:0.19,0.183) {XLK};
  \node[font=\scriptsize,anchor=north] at (axis cs:0.16,0.097) {XLU};
\end{axis}
\end{tikzpicture}
\caption{Nine sector ETFs, 2009--2013 (inputs of \cref{ch10:tab-etf} and the printed correlations; the risk-free rate is
taken to be $0$). Dots: individual ETFs. Blue and green: efficient frontier with cash and long-only positions capped at $30\%$
and $15\%$ (plots~4 of \file{ETF_pridA_30} and \file{ETF_pridA_15}, recomputed here). The dashed line is the
CML of the \emph{unconstrained} problem, with a Sharpe ratio of $1.85$, reached only with large short positions
(\cref{ch10:sec-est}).}\label{ch10:fig-etf}
\end{figure}

\begin{coursenote}[Independent check of the case study]
The R code is not available, so the printed allocations were re-derived from the printed means, volatilities and
correlations (all rounded to two decimals) by maximising the return at a given volatility, with $r_0=0$,
$w_j\in[0,c]$ and full investment or less, with a general-purpose solver. At a target volatility of $15.3\%$ ($30\%$ cap) it gives XLV $30\%$, XLP $22.9\%$, XLY $30\%$,
XLK $17.1\%$ and a return of $18.8\%$ (printed: $30$, $22.7$, $30$, $17.3$, $18.7$); at $9.9\%$ it gives XLV $21.8\%$, XLP $30\%$, XLY $18.7\%$,
cash $29.5\%$ and $12.7\%$ (printed: $15.7$, $30$, $22.6$, $31.7$, $12.4$). The sets of assets, the order in which the caps bind
and the frontiers agree with the printed ones; the few percentage points of difference in XLV and XLY
come from the rounding of the inputs (a half-point change in a mean moves weights by several points), a small illustration of the sensitivity of \cref{ch10:sec-est}. For the $15\%$ cap at $9.9\%$ volatility the recomputation gives XLV, XLP, XLY, XLK at $15\%$,
XLU $3.6\%$, XLI $1.6\%$, cash $34.8\%$ (printed: XLU $6.9\%$, cash $33.1\%$, return $11.7\%$ in both).
\end{coursenote}

\begin{finance}[Sector market-neutral models and diversification]
In the second demonstration of 2013 L14 Kempthorne replaces the ETFs by hedge-fund style strategies: multi-factor pricing
models trading long/short, market-neutral, within each of the nine sectors, over five years
\ts{13}{14}{1:15:29}\ts{13}{14}{1:16:33}. Their total returns are modest ($20$--$60\%$ over five years) but they are
\emph{little correlated} with each other, so the optimiser gets a large benefit from combining them; the model for the utilities sector
receives the largest weight, then industrials and energy, and a target volatility of $10\%$ is reached \ts{13}{14}{1:17:35}\ts{13}{14}{1:18:36}.
The data behind the plots are not in the files, so only this qualitative description is given. The lesson is the one
of \cref{ch10:sec-two}: low correlation makes the frontier bend to the left, and diversification is worth most for
uncorrelated bets.
\end{finance}

% ==================================================================
\section{Notes on the sources, differences between the years and errata}\label{ch10:sec-notes}

\begin{coursenote}[2013 versus 2024]
The 2013 lectures are split: Kempthorne gives the mathematics (his lecture has a full derivation, a case study,
utility theory, constraints and risk measures), Xia the practice. The 2024 lecture has only Xia, and it is a
different lecture from his 2013 one. It keeps the special cases of two assets (slides 14--20; Xia says he skips the mathematics that he covered in earlier years \ts{24}{13}{40:00}) and the rebalancing example, and it adds: the endowment model, MPT
issue~1 (volatility as risk) with the gain--loss ratio, MPT issue~2 (predicting the future) with the crowd model, and
power laws. It drops the Kelly formula, the beta discussion and the Millennium Bridge \emph{and} metronome videos of 2013 (2024 shows
a video of birds and mentions the bridge). The 2024 slides point to the 2013 video of Xia's lecture as
``one of the top five most viewed courses'' of OpenCourseWare \ts{24}{13}{1:06}. What Kempthorne derives (Markowitz problem, tangency
portfolio, Tobin, utility, constraints, risk measures) is not covered in 2024 by anyone: the derivations of this chapter
follow \file{lecnote14}.
\end{coursenote}

\begin{coursenote}[Erratum: parameters of the two-asset simulation]
Kempthorne says at the start of the demonstration that asset~1 has mean $15\%$ and volatility $25\%$, asset~2 mean $20\%$ and
volatility $30\%$ \ts{13}{14}{5:21}, and page~2 of \file{CaseStudy6} prints the same numbers (with typos:
``$E(R_1)=0.20=\alpha_2$'' and ``$\sqrt{\Var R_1}=0.30=\sigma_2$'' should refer to $R_2$, and ``$\sigma_w^2=\Var(R_2)$'' should be
$\Var(R_w)$). But the plots of \file{Smltn_TwoAst} are titled ``mean returns $0.15,0.25$, asset volatilities $0.20,0.30$'',
Kempthorne himself reads ``mean returns $15\%$, $25\%$, volatility $20$ to $30$'' when he shows them \ts{13}{14}{10:44}, and the
sample statistics ($0.19$--$0.20$ and $0.29$ for the volatilities, cumulative return $\approx3.3$ for asset~2, i.e.\ about $0.34$ per year over $9.6$ years) fit
these values, not the printed ones. This chapter uses the values of the plots.
\end{coursenote}

\begin{coursenote}[Errata and remarks on \file{CaseStudy6}]
\begin{itemize}
  \item Section~2.3 refers to ``Plot~4 in the two files \file{ETF_S_1_periodA_30.pdf} and \file{ETF_S_1_periodA_30.pdf}'': the second
      should be the $15\%$ file (\file{ETF_pridA_15}); the case study cites the file names of the original R output
      (\texttt{ETFS\_1\_periodA\_30}), which differ from the OCW names \file{ETF_pridA_30}.
  \item Section~3.3 (2003--2006, $15\%$ cap) says both ``technology
      never enters the optimal allocation'' and ``the allocations to XLK increase until its limit is reached'': the second
      bullet is copied from Section~2.3; the printed table has XLK at $0.000$ throughout.
  \item The bars of plot~3 use very
      similar colours for consumer discretionary and energy (both green) and the line plots repeat colours (XLB and XLU are both
      black in the first plot of each file), so the ETFs are identified through the tables.
\end{itemize}
\end{coursenote}

\begin{coursenote}[Slides and lectures: small corrections]
\begin{itemize}
  \item \file{lecnote14}, slide~19 writes $\E[u(W)]\propto\E[W-\frac12\lambda(W-w_*)^2]\approx\E[W]-\frac12\lambda\Var W$: the
      $\propto$ hides an additive constant and the positive factor $u'(w_*)$; the content is that of \cref{ch10:prop-mv-utility}(i).
  \item Slide~31, translation invariance: see the erratum in \cref{ch10:sec-risk}.
  \item Slide~24 of \file{lec13} asserts the ratio $(G-L)/(G+L)$ is ``the same as
      the Kelly criterion for binary bets'': true only for even-money bets (\cref{ch10:prop-gl}(i)).
  \item In 2013 L16 Xia says that after year~1 the winning
      asset ``becomes 75'' when rebalancing: the correct amounts are $62.5$ in each asset (\cref{ch10:sec-rebal}); the 2024 slide has it right.
  \item The transcript of 2013 L14 says Fama ``was just awarded, what, two weeks ago'': the 2013 economics prize
      (Fama, Hansen, Shiller) was announced on 14 October 2013, ten days before the lecture of 24 October (the case-study date).
  \item In 2013 L16 two remarks made ``from the audience'' (about estimating returns versus volatility, and about the Kelly criterion, Shannon and Samuelson)
      are by a co-lecturer, in all likelihood Kempthorne; the transcript does not identify the speaker.
  \item \file{lecnote16} contains only the outline (five items: portfolio construction, portfolio theory in special cases, risk parity, limitations, what we learned): Xia used a single slide and the blackboard, so 2013 L16
      has no formulas apart from those reconstructed here from the transcript and from the 2024 slides.
\end{itemize}
\end{coursenote}

% ==================================================================
\begin{keyformulas}
\begin{align*}
&\text{Portfolio:} && \mu_{\bm w}=\bm w\T\bm\mu,\quad \sigma_{\bm w}^2=\bm w\T\Sigma\bm w,\quad \bm w\T\bm 1=1\\
&\text{Two assets:} && \sigma_w^2=(1-w)^2\sigma_1^2+w^2\sigma_2^2+2w(1-w)\rho\sigma_1\sigma_2\\
& && w^*=\tfrac{\sigma_1^2-\rho\sigma_1\sigma_2}{\sigma_1^2+\sigma_2^2-2\rho\sigma_1\sigma_2},\quad
  \sigma^{*2}=\tfrac{\sigma_1^2\sigma_2^2(1-\rho^2)}{\sigma_1^2+\sigma_2^2-2\rho\sigma_1\sigma_2}\\
&\text{Constants:} && a=\bm\mu\T\Sigma^{-1}\bm\mu,\ b=\bm\mu\T\Sigma^{-1}\bm 1,\ c=\bm 1\T\Sigma^{-1}\bm 1,\ \Delta=ac-b^2\\
&\text{Optimal weights:} && \bm w_0=\lambda_1\Sigma^{-1}\bm\mu+\lambda_2\Sigma^{-1}\bm 1,\quad
  (\lambda_1,\lambda_2)=\tfrac1\Delta(c\mu_0-b,\ a-b\mu_0)\\
&\text{Frontier:} && \sigma_0^2=\tfrac{c\mu_0^2-2b\mu_0+a}\Delta=\tfrac1c+\tfrac c\Delta(\mu_0-\tfrac bc)^2\\
&\text{Min.\ variance:} && \bm w_g=\tfrac{\Sigma^{-1}\bm 1}c,\quad\sigma_g^2=\tfrac1c,\quad \mu_g=\tfrac bc\\
&\text{Risk aversion:} && \bm w(\lambda)=\bm w_g+\tfrac1\lambda\Sigma^{-1}(\bm\mu-\tfrac bc\bm 1)\\
&\text{With riskless asset:} && \bm w_0=\tfrac{\mu_0-r_0}{d^2}\Sigma^{-1}\bm\eta,\quad \bm\eta=\bm\mu-r_0\bm 1,\quad d^2=\bm\eta\T\Sigma^{-1}\bm\eta\\
&\text{Tangency portfolio:} && \bm w_M=\tfrac{\Sigma^{-1}\bm\eta}{\bm 1\T\Sigma^{-1}\bm\eta},\quad \mu_M-r_0=\tfrac{d^2}{b-r_0c},\quad
  \sigma_M=\tfrac d{b-r_0c}\\
&\text{Tobin, CML:} && m_{\bm w}=\tfrac{\mu_0-r_0}{\mu_M-r_0},\quad \mu_P=r_0+\tfrac{\mu_M-r_0}{\sigma_M}\sigma_P\\
&\text{Sharpe ratio:} && \mathrm{SR}=\tfrac{\mu_P-r_0}{\sigma_P}\le d\ \text{(max.\ at }\bm w_M),\ \text{leverage-invariant}\\
&\text{SML:} && \mu_i-r_0=\beta_i(\mu_M-r_0),\quad \beta_i=\tfrac{\Cov(R_i,R_M)}{\sigma_M^2}=\rho_{iM}\tfrac{\sigma_i}{\sigma_M}\\
&\text{Equal weights:} && \Var R_{\rm ew}=\tfrac{\bar\sigma^2}n+(1-\tfrac1n)\bar c\\
&\text{Constant mix:} && g(\bm w)=\bm w\T\bm\mu-\tfrac12\bm w\T\Sigma\bm w\ \ (\text{log growth})\\
&\text{Risk contribution:} && \mathrm{RC}_i=\tfrac{w_i(\Sigma\bm w)_i}{\sigma_P},\ \textstyle\sum_i\mathrm{RC}_i=\sigma_P;\ \text{2 assets: }w_i\propto1/\sigma_i\\
&\text{Utility:} && \lambda_A=-\tfrac{u''}{u'},\ \lambda_R=-\tfrac{Wu''}{u'};\quad \E[u]\approx u(w_*)+\tfrac12u''\Var W;\\
& && \text{CARA+normal: CE}=\mu_W-\tfrac\lambda2\sigma_W^2\\
&\text{Kelly:} && f^*=\tfrac pl-\tfrac qg;\ \ \text{Gaussian: }f^*=\tfrac{\mu}{\sigma^2}\\
&\text{Gain--loss ratio:} && \tfrac{G-L}{G+L}=\tfrac{\E[R]}{\E[|R|]}\in[-1,1],\ G=\E[R^+],\ L=\E[R^-]\\
&\text{Constraints:} && \mathrm{TE}^2=(\bm w-\bm w_B)\T\Sigma(\bm w-\bm w_B),\ \ \textstyle\sum_i\beta_{ik}w_i=0\\
&\text{Estimation:} && \mathrm{se}(\hat\mu)=\tfrac\sigma{\sqrt T},\quad \tfrac{\mathrm{se}(\hat\sigma)}\sigma\approx\tfrac1{\sqrt{2n}}\\
&\text{Crowd loop:} && \dd O=\tfrac{\sum A_iC_i\dd E+\sum A_i\varepsilon_i}{1-\bar AB}\ \ (\text{mean field})\\
&\text{Power law:} && \Prob[X>x]\sim x^{-k},\ f\sim x^{-a},\ a=1+k=1+\tfrac1b
\end{align*}
\end{keyformulas}

% ==================================================================
\section*{Exercises}
\addcontentsline{toc}{section}{Exercises}

\subsection*{From 2013 Problem Set 6 (Problems 3--5; Problems 1--2 are in \cref{ch:timeseries})}
In all three problems the portfolio $\bm w$ is fully invested, $\sum_jw_j=1$, with no short sales, $w_j\ge0$.

\begin{exercise}[PS6 (2013), Problem 3: two risky assets]\label{ch10:ex-ps6-3}
Two assets with returns $\bm R=(R_1,R_2)\T$, means $(\alpha_1,\alpha_2)$, covariance
$\Sigma=\begin{psmallmatrix}\sigma_1^2&\rho\sigma_1\sigma_2\\ \rho\sigma_1\sigma_2&\sigma_2^2\end{psmallmatrix}$ and a portfolio
$R_{\bm w}=w_1R_1+w_2R_2$.
\begin{enumerate}[(a)]
  \item Prove that $\Var(R_{\bm w})\le\max(\sigma_1^2,\sigma_2^2)$ for every such portfolio.
  \item For $\sigma_1=\sigma_2$ and $\rho=0$ find the minimum-variance portfolio $\bm w^*$ and its variance.
  \item For $\sigma_1=\sigma_2$ and arbitrary $\rho$ find $\bm w^*$ and $\Var(R_{\bm w^*})$, and graph the minimal variance as a function of
      $\rho\in[-1,1]$.
\end{enumerate}
\end{exercise}

\begin{exercise}[PS6 (2013), Problem 4: $m>2$ assets, equal volatilities]\label{ch10:ex-ps6-4}
Now $m>2$ assets with covariance matrix $\Sigma=(\rho_{ij}\sigma_i\sigma_j)$.
\begin{enumerate}[(a)]
  \item Prove that $\Var(R_{\bm w})\le\max_j\sigma_j^2$ for every long-only fully invested portfolio.
  \item If $\sigma_1=\dots=\sigma_m$ and all $\rho_{ij}=0$ ($i\neq j$), find the minimum-variance portfolio, express its variance as a function
      of $m$ and find its limit as $m\to\infty$.
  \item If $\sigma_1=\dots=\sigma_m$ and $\rho_{ij}\equiv\rho$ for all $i\neq j$ (\emph{equicorrelation}): first show that
      $\rho\ge-1/(m-1)$ (hint: $\bm 1$ is an eigenvector of $\Sigma$; compute its eigenvalue and impose that $\Sigma$ be positive
      semi-definite). Then find the minimum-variance portfolio and its variance, graph the variance as a function of $\rho\in[-\frac1{m-1},1]$,
      find the limit as $m\to\infty$ and compare it with (b), commenting on the ability to diversify away risk by adding equicorrelated assets.
\end{enumerate}
\end{exercise}

\begin{exercise}[PS6 (2013), Problem 5: unequal variances and the general case]\label{ch10:ex-ps6-5}
\leavevmode
\begin{enumerate}[(a)]
  \item For $m=2$, $\sigma_1\neq\sigma_2$ and $\rho=0$, find the minimum-variance portfolio and its variance.
  \item For $m>2$ and $\Sigma=\diag(\sigma_1^2,\dots,\sigma_m^2)$ (arbitrary variances, zero correlations) find the minimum-variance
      portfolio and its variance, and express the variance in terms of $m$ and $\tilde\sigma^2=\bigl(\frac1m\sum_j\sigma_j^{-2}\bigr)^{-1}$.
  \item For $m>2$ and an arbitrary positive definite $\Sigma$, find the minimum-variance portfolio and its variance.
\end{enumerate}
\end{exercise}

\begin{exercise}[2024 Investment Game and the gain--loss ratio (ungraded)]\label{ch10:ex-game}
The 2024 lecture (slide~36) repeats the game of \cref{ex:game}: $\$10{,}000$ in one stock or ETF from the close of 13~September to the
close of 15~November, one optional switch on 11~October, and the final calculation of $G$, $L$ and $(G-L)/(G+L)$ from the daily
profits and losses. Carry it out on a two-month window of your choice (for example for the sector ETFs of \cref{ch10:tab-etf}), compute the
Sharpe ratio of the same daily P\&L, and compare the ranking of two investments by the two measures; see \cref{ch10:prop-gl}.
\end{exercise}

\subsection*{Extra exercises (not from the course)}

\begin{exercise}[not from the course: risk parity]
Two assets have volatilities $16\%$ and $6\%$ and correlation $0.3$.
\begin{enumerate}[(a)]
  \item Find the risk-parity weights and check that the risk
      contributions of \cref{ch10:def-rc} are equal.
  \item What fraction of the risk of the $60/40$ portfolio comes from the first asset?
  \item By what factor must the
      risk-parity portfolio be levered to have the volatility of the $60/40$ portfolio?
\end{enumerate}
\end{exercise}

\begin{exercise}[not from the course: the frontier]
\leavevmode
\begin{enumerate}[(a)]
  \item Verify that $\bm w_g$ of \cref{ch10:cor-gmv} minimises $\bm w\T\Sigma\bm w$ under $\bm w\T\bm 1=1$ and that the frontier
      \eqref{ch10:eq-frontier} has asymptotes $\mu=b/c\pm\sqrt{\Delta/c}\,\sigma$.
  \item Show that the difference of two efficient portfolios is proportional to $\bm v$ of
      \eqref{ch10:eq-twofund} and deduce the two-fund theorem.
  \item Show that, if $r_0<\mu_g$, the tangency portfolio maximises the
      Sharpe ratio over the fully invested portfolios, and find the tangency point of the line through $(0,r_0)$ with the hyperbola.
\end{enumerate}
\end{exercise}

\begin{exercise}[not from the course: CARA utility]
An investor with $u(W)=1-e^{-AW}$ and initial wealth $W_0$ can invest in $m$ assets with normal returns and in a riskless asset. Show that the
optimal amounts invested in the risky assets are $A^{-1}\Sigma^{-1}\bm\eta$ (in dollars) and identify the fraction of the wealth held in the
tangency portfolio.
\end{exercise}

\begin{exercise}[not from the course: Kelly]
A bet wins $g=2$ per unit staked with probability $p=0.6$ and loses $l=1$ otherwise.
\begin{enumerate}[(a)]
  \item Compute the Kelly fraction and the
      optimal growth rate $G(f^*)$.
  \item What growth rate does the half-Kelly fraction $f^*/2$ give, as a percentage of $G(f^*)$?
  \item Show that for small edges
      half-Kelly gives $3/4$ of the optimal growth rate.
\end{enumerate}
\end{exercise}

\begin{exercise}[not from the course: gain--loss ratio]
Investment~A pays $+3$ with probability $0.4$ and $-1$ with probability $0.6$. Investment~B pays $+1$ or $-1$
with probability $0.5$ each plus a fixed $0.3$. Compute for each the mean, the standard deviation, the Sharpe ratio, and
$(G-L)/(G+L)$; comment on the difference in rankings and on the Kelly fractions.
\end{exercise}

\begin{exercise}[not from the course: VaR and CVaR]
Modify \cref{ch10:prop-var} to three independent positions with default probability $3\%$ and loss $100$: compute the
$95\%$ VaR and expected shortfall of one position and of the portfolio of three, and check subadditivity of ES.
\end{exercise}

\begin{exercise}[not from the course: rebalancing]
Two independent assets have geometric Brownian motions with $\sigma_1=30\%$, $\sigma_2=10\%$ and the same $\mu$. Using
\cref{ch10:prop-growth}, find the constant-mix weight that maximises the log growth rate and the value of the ``excess growth'' at
that weight. Compare with equal weights.
\end{exercise}

\begin{exercise}[not from the course: crowd model and power laws]
\leavevmode
\begin{enumerate}[(a)]
  \item In \cref{ch10:prop-tipping}, take $N$ agents with $A_i=A/N$ and $B_i\in\{B_L,B_H\}$, a fraction $x$ in the reactive state; write the loop
      gain and find the critical $x$ for the tipping point.
  \item For a Pareto tail $\Prob[X>x]=(x_m/x)^k$, $k>1$, prove that the top fraction $q$
      holds $q^{1-1/k}$ of the total, and compute $k$ for a ``$90$--$10$'' rule.
\end{enumerate}
\end{exercise}

% !TEX root = ../main.tex
\chapter{Stochastic Processes II: Brownian Motion}\label{ch:sp2}
\begin{paracol}{2}

\begin{sources}
\begin{tabularx}{\linewidth}{@{}l l L@{}}
\toprule
\textbf{Lecture} & \textbf{Speaker} & \textbf{Content}\\
\midrule
2024 L14 & P.~Kempthorne & History; Brownian motion with diffusion coefficient $\sigma^2$: increments, Markov property, covariance; transition density and the diffusion equation; Brownian motion as limit of normalised random walks (invariance principle); reflection principle, maximum, first hitting time; reflected and absorbed Brownian motion; Brownian bridge and the empirical distribution function; Brownian motion with drift; gambler's ruin by infinitesimal one-step analysis; non-differentiability; quadratic variation; log-price quadratic variation; Laplace increments at exponential times; filtrations and adapted processes.\\
2013 L17 & C.~Lee & Continuous-time processes and why their law is hard to describe; the existence theorem for standard Brownian motion; random-walk limit; Brown, Einstein and Bachelier; zero crossings, $\sqrt t$ envelope, nowhere differentiability; distribution of the maximum by the reflection principle; quadratic variation and $(\dd B)^2=\dd t$; Brownian motion with drift; \emph{why $S_t\ne e^{B_t}$}, Taylor expansion and the first appearance of It\^o's lemma.\\
\bottomrule
\end{tabularx}

\smallskip
\textbf{Files.} 2024 slides \file{mit18_642_f24_lec14_1.pdf} (77 pages with overlays, 20 slides) and the figure decks
\file{lec14_2} (two-dimensional Brownian motion, 201 animation frames), \file{lec14_3} (normal densities $p(x\mid t)$,
$t=1,\dots,10$), \file{lec14_4} (limiting random walk, 201 frames), \file{lec14_5} (reflected Brownian motion, three trials),
\file{lec14_6} ($\Prob[M(t)>x]$), \file{lec14_7} (hitting-time density), \file{lec14_8} (gambler's ruin with drift, four
panels), \file{lec14_9} (quadratic variation, 201 frames); 2013 note \file{MIT18_S096F13_lecnote17.pdf}.
\textbf{Problem sets.} 2024 PS4 Problems~1--2 (conditional moments and scaling of Brownian motion; Problems~3--4 are the
maximum-likelihood estimation of a geometric Brownian motion, \cref{ch:volatility}); 2024 PS5 Problems~1--2 (L\'evy processes:
Brownian motion with drift and the Poisson process, with simulations; Problems~3--5 are volatility estimation,
\cref{ch:volatility}). 2013 PS8 Problem~A-2 is 2024 PS4 Problem~1 and Problem~B-1 is on two-dimensional Brownian motion; PS8
belongs to \cref{ch:stochcalc}, see the course note in \cref{ch11:sec-ahead}. The exercises are at the end of the chapter.
\textbf{Not covered.} The construction of Brownian motion (the existence theorem is only quoted in both years), the
proofs of Donsker's theorem and of nowhere differentiability, and the law of the iterated logarithm.
The R code that generated the figure decks is not part of the download; the figures of this chapter are recomputed from formulas.
\end{sources}

\section*{Motivation}
Chapter~\ref{ch:sp1} studied random walks, martingales and stopping times in discrete time. Prices, however,
change at every instant a trade takes place, and the mathematics of derivative pricing (\cref{ch:stochcalc,ch:sde,ch:bs}) is
carried by a \emph{continuous-time} model. Its basic building block is \emph{Brownian motion}, the continuous-time random walk.
It was found three times, in three unrelated places: by the botanist Robert Brown (1827), watching pollen grains jitter in
water; by Louis Bachelier (1900), as a model of prices on the Paris exchange; by Albert Einstein (1905), who explained
Brown's observation by the bombardment of the grain by water molecules and obtained the diffusion law
$\langle x^2\rangle=2Dt$ \ts{24}{14}{0:00}\ts{13}{17}{18:21}\ts{13}{17}{20:29}. Norbert Wiener, an MIT mathematician,
gave it its rigorous foundation in 1923, which is why it is also called the \emph{Wiener process}
\ts{24}{14}{1:04}\ts{13}{17}{13:43}.

The reason for its centrality is a \emph{universality} statement, made precise in \cref{ch11:sec-donsker}: any random walk with
finite-variance steps, watched from far away, looks like Brownian motion. For a physicist this is the central limit theorem in
disguise; for a trader it says that a price built from many small, roughly independent trades is described by two numbers, a
drift and a volatility. This chapter develops the calculus of \emph{sample paths} that the next chapters use:

\begin{itemize}
  \item the definition and the elementary properties (Gaussian, Markov, martingale), and the diffusion equation
        (\cref{ch11:sec-def});
  \item Brownian motion as the scaling limit of the random walk (\cref{ch11:sec-donsker});
  \item the strange geometry of the paths: continuous but nowhere differentiable, with \emph{finite quadratic variation}
        $[B]_T=T$, the fact that later becomes It\^o's lemma (\cref{ch11:sec-paths});
  \item the reflection principle, giving the law of the running maximum and of the first hitting time
        (\cref{ch11:sec-reflection}), with reflected and absorbed Brownian motion;
  \item Brownian motion with drift and the gambler's ruin problem, geometric Brownian motion, and why the
        price is \emph{not} $e^{B_t}$ (\cref{ch11:sec-drift,ch11:sec-gbm});
  \item further processes built from $B$: bridge, random-time sampling, two dimensions, L\'evy processes, filtrations
        (\cref{ch11:sec-variants}).
\end{itemize}

% ==================================================================
\section{Definition and elementary properties}\label{ch11:sec-def}

\subsection{Brownian motion}
In discrete time the law of a random walk is specified through its increments: ``each step is $\pm1$ with probability
$\tfrac12$''. In continuous time there is no next step to describe, so one lists \emph{properties} of increments over all
intervals and asserts that a process with these properties exists \ts{13}{17}{4:10}\ts{13}{17}{5:18}\ts{13}{17}{9:28}.

\begin{definition}[Standard Brownian motion]\label{ch11:def-bm}
A stochastic process $\{B(t),\,t\ge0\}$ is a \emph{standard Brownian motion} (Wiener process) if
\begin{enumerate}[(i)]
  \item $B(0)=0$ and $t\mapsto B(t)$ is continuous (with probability one);
  \item (stationary Gaussian increments) for $0\le s<t$, $B(t)-B(s)\sim N(0,\,t-s)$;
  \item (independent increments) for $t_1\le t_2\le t_3\le t_4$ the increments $B(t_2)-B(t_1)$ and $B(t_4)-B(t_3)$ are
        independent; more generally the increments over non-overlapping intervals are mutually independent
        \ts{24}{14}{3:08}\ts{24}{14}{4:13}\ts{24}{14}{5:17}\ts{13}{17}{7:21}.
\end{enumerate}
\end{definition}

The 2024 slides call $\sigma^2$ the \emph{diffusion coefficient} and write $B(t+\Delta t)-B(t)\sim N(0,\sigma^2\Delta t)$
already in the definition; from the slide on drift onwards $B$ is standard and $\sigma B$ is used. We follow the second convention:
\emph{Brownian motion with diffusion coefficient $\sigma^2$} is $\sigma B(t)$, and one may start it at any level $x$ by
adding $x$ (``Brownian motion started at $x$'') \ts{24}{14}{9:36}. Property (ii) is the one to remember: over a time span
$h$ the process moves by a centred normal variable whose \emph{variance is the time elapsed}, so its typical size is
$\sqrt h$ \ts{13}{17}{12:40}.

\begin{coursenote}[Existence, and why the definition is indirect]
Both lecturers stress that (i)--(iii) are a \emph{list of demands} on a probability law on the space of continuous
paths, and that the theorem ``such a law exists'' is hard and not proved in the course \ts{13}{17}{6:20}\ts{13}{17}{10:36}\ts{13}{17}{11:37}.
The sample space is the space of paths $\omega:[0,\infty)\to\R$ itself, with $\Prob$ a measure on it (Wiener measure); a
path is an outcome, and there is no ``density at $\omega$'' \ts{13}{17}{5:18}. One standard route (Wiener 1923; L\'evy 1948)
builds $B$ as a uniform limit of random piecewise-linear functions obtained by refining a dyadic grid; another is
Donsker's theorem below, where $B$ is \emph{defined} as the limit of rescaled random walks (this is how the 2024 lecture
introduces it, \ts{24}{14}{17:22}). What (i)--(iii) do determine is the joint law of $(B(t_1),\dots,B(t_m))$ for every finite
set of times, by writing $B(t_k)$ as a sum of independent normal increments; continuity of the paths is the extra input
that makes the law unique on path space.
\end{coursenote}

\begin{proposition}[Gaussian characterisation]\label{ch11:prop-gauss}
A process with continuous paths and $B(0)=0$ is a standard Brownian motion if and only if it is a centred Gaussian process
with covariance function $\Cov(B(s),B(t))=\min(s,t)$.
\end{proposition}
\begin{proof}
($\Rightarrow$) Take $s\le t$. Then $B(t)=B(s)+[B(t)-B(s)]$ with independent summands, so
$\E[B(s)B(t)]=\E[B(s)^2]+\E[B(s)]\,\E[B(t)-B(s)]=s+0$ \ts{24}{14}{7:33}\ts{24}{14}{8:34}. Any vector $(B(t_1),\dots,B(t_m))$ is a linear
image of independent normal increments, hence jointly Gaussian.
($\Leftarrow$) For $s<t$ the variable $B(t)-B(s)$ is centred normal with variance $t-2s+s=t-s$. For
$s_1<t_1\le s_2<t_2$, $\Cov(B(t_1)-B(s_1),B(t_2)-B(s_2))=t_1-t_1-s_1+s_1=0$ (write out the four
minima); jointly Gaussian and uncorrelated implies independent.
\end{proof}
Note the meaning of the covariance: $B(s)$ and $B(t)$ are correlated to the extent that the two time windows $[0,s]$ and $[0,t]$
overlap, $\Corr(B(s),B(t))=\sqrt{s/t}$ for $s<t$ \ts{24}{14}{8:34}.

\begin{proposition}[Elementary properties]\label{ch11:prop-elem}
Let $B$ be a standard Brownian motion and $\mathcal F_t=\sigma(B(u),u\le t)$.
\begin{enumerate}[(i)]
  \item \emph{Mean and variance.} $\E[B(t)]=0$, $\Var B(t)=t$: the standard deviation grows like $\sqrt t$
        \ts{24}{14}{6:28}\ts{24}{14}{7:33}.
  \item \emph{Markov property.} For $h>0$, $B(t+h)=B(t)+[B(t+h)-B(t)]$ where the bracket is $N(0,h)$ and independent of
        $\mathcal F_t$; hence the law of $B(t+h)$ given $\mathcal F_t$ is $N(B(t),h)$, which depends on the past only through $B(t)$
        \ts{24}{14}{5:17}\ts{24}{14}{6:28}\ts{24}{14}{11:54}.
  \item \emph{Martingales.} $B(t)$, $B(t)^2-t$ and, for every real $\lambda$, $\exp(\lambda B(t)-\tfrac12\lambda^2t)$ are
        martingales with respect to $\{\mathcal F_t\}$.
  \item \emph{Symmetries.} $-B$; the shifted process $B(t_0+t)-B(t_0)$; the rescaled process $c^{-1/2}B(ct)$ for $c>0$; and the
        time-inverted process $tB(1/t)$ ($t>0$, and $0$ at $t=0$) are again standard Brownian motions.
\end{enumerate}
\end{proposition}
\begin{proof}
(i) is property (ii) of \cref{ch11:def-bm} with $s=0$.

(ii) is property (iii) of \cref{ch11:def-bm} together with the definition of $\mathcal F_t$.

(iii) For $s<t$, using independence of
$B(t)-B(s)$ from $\mathcal F_s$: $\E[B(t)\mid\mathcal F_s]=B(s)$;
$\E[B(t)^2\mid\mathcal F_s]=B(s)^2+2B(s)\E[B(t)-B(s)]+\E[(B(t)-B(s))^2]=B(s)^2+(t-s)$; and with the moment generating
function of a normal variable, $\E[e^{\lambda(B(t)-B(s))}]=e^{\lambda^2(t-s)/2}$, so
$\E[e^{\lambda B(t)-\lambda^2t/2}\mid\mathcal F_s]=e^{\lambda B(s)}e^{\lambda^2(t-s)/2}e^{-\lambda^2t/2}=e^{\lambda B(s)-\lambda^2s/2}$.

(iv) Each candidate is continuous (at $t=0$ for the inverted one use $B(u)/u\to0$, which is the law of large numbers for
$B$), centred Gaussian, and one checks that its covariance is $\min(s,t)$; then apply \cref{ch11:prop-gauss}.
For the rescaling, $\Cov(c^{-1/2}B(cs),c^{-1/2}B(ct))=c^{-1}\min(cs,ct)=\min(s,t)$; this is the content of 2024 PS4
Problem~2 (\cref{ch11:ex-scaling}); for the inversion, $st\min(1/s,1/t)=\min(t,s)$.
\end{proof}

Property (c) is the continuous-time image of \cref{ch04:ex-m1,ch04:ex-m2,ch04:ex-m4}: the same three martingales that
appeared for the random walk (sum, square minus $n\sigma^2$, exponential) survive the limit. The variance $t$ that
appears in $B^2-t$ is the \emph{quadratic variation} of \cref{ch11:sec-paths}. The martingale $\exp(\lambda B-\lambda^2t/2)$ is
the ancestor of the change of measure behind risk-neutral pricing (\cref{ch:bs}).

\begin{intuition}[The physicist's reading]\label{ch11:int-physics}
Standard Brownian motion is the position of a free overdamped particle whose velocity has been eliminated: the Langevin
equation $\dot x=\sqrt{2D}\,\xi(t)$ with white noise $\xi$ has solution $x(t)=\sqrt{2D}\,B(t)$, so
$\langle x^2\rangle=2Dt$ (Einstein's law) and $\sigma^2=2D$. The white noise $\xi=\dot B$ does not exist as a function
(\cref{ch11:sec-paths}); only its integral does. The law of the whole path is the \emph{Wiener measure}, the measure of the
Euclidean path integral of a free particle,
$\propto\exp\bigl(-\tfrac12\int_0^t\dot x^2\dd s\bigr)\prod\dd x(s)$, formal but useful; the transition density below is
its propagator (the heat kernel), and $\min(s,t)$ is the Green function of $-\dd^2/\dd t^2$ with Dirichlet condition at $0$.
\end{intuition}

\begin{finance}[Bachelier's model]
Bachelier's thesis modelled the price itself as $S(t)=S(0)+\mu t+\sigma B(t)$, ``arithmetic'' Brownian motion; the
2013 note credits him with the first use of the process to value stocks and options \ts{13}{17}{13:43}. Its defect is that the price becomes
negative with positive probability, so equities are modelled through $\log S$ or through relative changes
(\cref{ch11:sec-gbm}). The arithmetic model survives where levels can be negative or small compared with their moves, in
particular for interest rates and spreads, where ``normal volatility'' is quoted in basis points per year.
\end{finance}

% ------------------------------------------------------------------
\subsection{Transition density and the diffusion equation}\label{ch11:sec-diffusion}
Let $B^x$ be Brownian motion started at $x$, i.e.\ $x+B$. For $\tau>0$ and diffusion coefficient $\sigma^2$, the density of
$B^x(s+\tau)$ given $B^x(s)=x$ is the Gaussian
\begin{equation}\label{ch11:eq-kernel}
  p(y,\tau\mid x)=\frac1{\sqrt{2\pi\sigma^2\tau}}\exp\Bigl(-\frac{(y-x)^2}{2\sigma^2\tau}\Bigr),
\end{equation}
that is, $\Prob[B(t)\le c\mid B(s)=x]=\Phi\bigl((c-x)/(\sigma\sqrt{t-s})\bigr)$ in terms of the standard normal cdf $\Phi$ and pdf
$\varphi$ \ts{24}{14}{9:36}\ts{24}{14}{10:42}\ts{24}{14}{11:54}\ts{24}{14}{13:03}.

\begin{proposition}[Diffusion equation]\label{ch11:prop-heat}
The kernel \eqref{ch11:eq-kernel} satisfies
\begin{equation}\label{ch11:eq-heat}
  \frac{\partial p}{\partial\tau}=\frac{\sigma^2}2\,\frac{\partial^2p}{\partial y^2},
  \qquad \int_{-\infty}^{\infty}p(y,\tau\mid x)\dd y=1\ \ (\tau>0),\qquad p(y,\tau\mid x)\xrightarrow[\tau\downarrow0]{}\delta(y-x),
\end{equation}
and \eqref{ch11:eq-heat} has no other solution with these properties. The same holds in the variable $x$ (backward equation).
\end{proposition}
\begin{proof}
Put $u=y-x$. Then $\partial_\tau p=p\,\bigl(-\tfrac1{2\tau}+\tfrac{u^2}{2\sigma^2\tau^2}\bigr)$ and
$\partial_y^2p=p\,\bigl(-\tfrac1{\sigma^2\tau}+\tfrac{u^2}{\sigma^4\tau^2}\bigr)$, so $\partial_\tau p=\tfrac{\sigma^2}2\partial_y^2p$.
Integrating \eqref{ch11:eq-heat} over $y$ gives $\partial_\tau\int p=\tfrac{\sigma^2}2[\partial_yp]_{-\infty}^{\infty}=0$: probability
is conserved \ts{24}{14}{16:17}. The Fourier transform in $y$ turns the equation into $\partial_\tau\hat p=-\tfrac12\sigma^2k^2\hat p$,
whose solution with $\hat p(k,0)=e^{-ikx}$ is $e^{-ikx-\sigma^2k^2\tau/2}$, the transform of \eqref{ch11:eq-kernel}; this gives uniqueness
among tempered solutions \ts{24}{14}{17:22}. The backward version follows from $p$ depending on $y-x$ only.
\end{proof}
The semigroup property (Chapman--Kolmogorov, \cref{ch04:thm-nstep}) is the convolution identity
$\int p(z,\tau_2\mid y)\,p(y,\tau_1\mid x)\dd y=p(z,\tau_1+\tau_2\mid x)$, the statement that a sum of independent normals is normal
with added variances. In other disciplines \eqref{ch11:eq-heat} is the \emph{heat equation}, in population dynamics the
\emph{Fokker--Planck} or \emph{Kolmogorov forward} equation \ts{24}{14}{17:22}; it returns in the valuation of derivatives
(\cref{ch:bs}), where the Black--Scholes PDE is a backward equation of this type.

\begin{figure}[ht]
\centering
\begin{tikzpicture}
\begin{axis}[width=0.9\linewidth,height=5cm,xlabel={$y$},ylabel={$p(y,t\mid0)$},xmin=-6,xmax=6,ymin=0,ymax=0.45,
  legend style={at={(0.5,1.03)},anchor=south,legend columns=3,font=\footnotesize},grid=major,grid style={black!10}]
  \addplot[thick,color=defcol,domain=-6:6,samples=100] {1/sqrt(2*pi*1)*exp(-x^2/(2*1))};
  \addplot[thick,color=intcol,domain=-6:6,samples=100] {1/sqrt(2*pi*2.5)*exp(-x^2/(2*2.5))};
  \addplot[thick,color=thmcol,domain=-6:6,samples=100] {1/sqrt(2*pi*6)*exp(-x^2/(2*6))};
  \addplot[only marks,mark=*,mark size=1.6pt,color=defcol] coordinates {(-1,0.24197) (1,0.24197)};
  \addplot[only marks,mark=*,mark size=1.6pt,color=intcol] coordinates {(-1.5811,0.15316) (1.5811,0.15316)};
  \addplot[only marks,mark=*,mark size=1.6pt,color=thmcol] coordinates {(-2.4495,0.09946) (2.4495,0.09946)};
  \legend{{$t=1$},{$t=2.5$},{$t=6$}}
\end{axis}
\end{tikzpicture}
\caption{Diffusion of the density of $B(t)$ (the panels of \file{lec14_3}, which shows $t=1,1.5,\dots,10$). The
density drops at the centre and rises in the tails; the dots mark the inflection points $y=\pm\sqrt t$, where the curvature
changes sign, so that $\partial_tp<0$ inside and $>0$ outside \ts{24}{14}{14:08}\ts{24}{14}{15:14}.}\label{ch11:fig-heat}
\end{figure}


% ==================================================================
\section{Brownian motion as the limit of random walks}\label{ch11:sec-donsker}

Let $X_1,X_2,\dots$ be i.i.d.\ with $\E[X_i]=0$, $\Var X_i=1$, and $S_n=X_1+\dots+X_n$
(\cref{ch04:sec-rw}). By the central limit theorem (\cref{ch03:thm-clt}), $\Prob[S_n/\sqrt n\le c]\to\Phi(c)$, i.e.\ $\Prob[S_n\le x]\approx\Phi(x/\sqrt n)$
\ts{24}{14}{17:22}\ts{24}{14}{18:27}. To speak about a \emph{path} rather than one time, rescale space by $\sqrt n$ and time by $n$.

\begin{definition}[Normalised random walk]\label{ch11:def-bn}
$B_n(t)=S_{\lfloor nt\rfloor}/\sqrt n$ for $t\ge0$, where $\lfloor nt\rfloor$ is the largest integer $\le nt$; the process jumps at
$t=k/n$ by $X_k/\sqrt n$ and is constant in between \ts{24}{14}{19:29}. The piecewise-linear version, which joins the points
$(k/n,S_k/\sqrt n)$ (the version of the 2013 note, ``connect the dots''), is continuous and differs from it by at most
$\max_{k\le nT}|X_k|/\sqrt n$ on $[0,T]$ \ts{13}{17}{14:56}\ts{13}{17}{16:01}.
\end{definition}

\begin{coursenote}[Erratum in the 2013 note]
\file{lecnote17} defines the piecewise-linear path by $Z(t/n)=Y_t$ for a walk with variance-one steps; the rescaling by $1/\sqrt n$ is missing (the
values would be of order $\sqrt n$ and no limit could exist). The correct definition is $Z(t/n)=Y_t/\sqrt n$, as in the 2024 slides.
\end{coursenote}

\begin{theorem}[Donsker's invariance principle]\label{ch11:thm-donsker}
As $n\to\infty$ the process $B_n$ converges in distribution (weakly, as a random element of the space of paths on $[0,T]$
with the uniform metric) to a standard Brownian motion $B$, whatever the distribution of the steps $X_i$
\ts{24}{14}{20:44}. In particular, for every continuous functional $F$ of the path (e.g.\ $\sup_{t\le1}$, $\int_0^1(\cdot)^2\dd t$,
the value at time $1$), $F(B_n)\xrightarrow{d}F(B)$.
\end{theorem}
\begin{proof}[What is proved here]
\emph{Finite-dimensional convergence.} For $0\le t_1<\dots<t_m$ the increments $B_n(t_j)-B_n(t_{j-1})$ are sums of
$\lfloor nt_j\rfloor-\lfloor nt_{j-1}\rfloor\sim n(t_j-t_{j-1})$ i.i.d.\ steps divided by $\sqrt n$, over disjoint blocks, hence independent; by the central limit
theorem each converges to $N(0,t_j-t_{j-1})$, and by independence jointly, which is exactly the joint law of the increments of $B$
\ts{13}{17}{17:19}. This gives $\Var B(t)=\lim_n\Var B_n(t)=t$ and $\E[B(t)]=0$, the mean and variance of the slides
\ts{24}{14}{19:29}. \emph{Tightness}, that no path of $B_n$ oscillates wildly between grid times so that the limit is a
continuous process, requires a maximal inequality (Kolmogorov or Doob) and is not proved in the course; it is what turns
convergence of the finite-dimensional laws into convergence of laws on path space \ts{13}{17}{6:20}.
\end{proof}

\begin{figure}[ht]
\centering
\begin{tikzpicture}
\begin{axis}[width=0.47\linewidth,height=4.8cm,xlabel={$t$},ylabel={$B_n(t)$},xmin=0,xmax=1,title={$n=20$, $\pm1$ steps},title style={font=\footnotesize},grid=major,grid style={black!10}]
  \addplot[thick,const plot,color=defcol] coordinates {(0.000,0.000) (0.050,-0.224) (0.100,-0.447) (0.150,-0.224) (0.200,-0.447) (0.250,-0.224) (0.300,0.000) (0.350,0.224) (0.400,0.000) (0.450,-0.224) (0.500,-0.447) (0.550,-0.671) (0.600,-0.447) (0.650,-0.224) (0.700,-0.447) (0.750,-0.224) (0.800,-0.447) (0.850,-0.224) (0.900,0.000) (0.950,0.224) (1.000,0.447)};
\end{axis}
\end{tikzpicture}\hfill
\begin{tikzpicture}
\begin{axis}[width=0.47\linewidth,height=4.8cm,xlabel={$t$},xmin=0,xmax=1,title={$n=150$},title style={font=\footnotesize},
  legend style={at={(0.5,1.25)},anchor=south,legend columns=2,font=\scriptsize},grid=major,grid style={black!10}]
  \addplot[thick,color=defcol] coordinates {(0.000,0.000) (0.007,0.082) (0.013,0.163) (0.020,0.082) (0.027,0.163) (0.033,0.082) (0.040,0.163) (0.047,0.245) (0.053,0.163) (0.060,0.245) (0.067,0.163) (0.073,0.082) (0.080,0.000) (0.087,0.082) (0.093,0.000) (0.100,-0.082) (0.107,-0.163) (0.113,-0.245) (0.120,-0.327) (0.127,-0.408) (0.133,-0.327) (0.140,-0.408) (0.147,-0.327) (0.153,-0.245) (0.160,-0.327) (0.167,-0.408) (0.173,-0.490) (0.180,-0.572) (0.187,-0.490) (0.193,-0.572) (0.200,-0.490) (0.207,-0.408) (0.213,-0.327) (0.220,-0.408) (0.227,-0.490) (0.233,-0.408) (0.240,-0.490) (0.247,-0.408) (0.253,-0.327) (0.260,-0.245) (0.267,-0.327) (0.273,-0.245) (0.280,-0.163) (0.287,-0.082) (0.293,-0.163) (0.300,-0.245) (0.307,-0.163) (0.313,-0.245) (0.320,-0.327) (0.327,-0.408) (0.333,-0.327) (0.340,-0.408) (0.347,-0.327) (0.353,-0.408) (0.360,-0.490) (0.367,-0.408) (0.373,-0.327) (0.380,-0.245) (0.387,-0.163) (0.393,-0.245) (0.400,-0.327) (0.407,-0.245) (0.413,-0.163) (0.420,-0.082) (0.427,-0.163) (0.433,-0.245) (0.440,-0.163) (0.447,-0.245) (0.453,-0.327) (0.460,-0.245) (0.467,-0.327) (0.473,-0.245) (0.480,-0.327) (0.487,-0.408) (0.493,-0.327) (0.500,-0.408) (0.507,-0.490) (0.513,-0.572) (0.520,-0.653) (0.527,-0.735) (0.533,-0.653) (0.540,-0.572) (0.547,-0.653) (0.553,-0.735) (0.560,-0.653) (0.567,-0.572) (0.573,-0.653) (0.580,-0.735) (0.587,-0.816) (0.593,-0.898) (0.600,-0.816) (0.607,-0.735) (0.613,-0.653) (0.620,-0.735) (0.627,-0.653) (0.633,-0.572) (0.640,-0.653) (0.647,-0.572) (0.653,-0.653) (0.660,-0.735) (0.667,-0.653) (0.673,-0.572) (0.680,-0.653) (0.687,-0.572) (0.693,-0.490) (0.700,-0.572) (0.707,-0.653) (0.713,-0.572) (0.720,-0.490) (0.727,-0.408) (0.733,-0.490) (0.740,-0.572) (0.747,-0.490) (0.753,-0.572) (0.760,-0.490) (0.767,-0.408) (0.773,-0.490) (0.780,-0.572) (0.787,-0.490) (0.793,-0.572) (0.800,-0.490) (0.807,-0.408) (0.813,-0.327) (0.820,-0.245) (0.827,-0.327) (0.833,-0.245) (0.840,-0.327) (0.847,-0.408) (0.853,-0.327) (0.860,-0.245) (0.867,-0.327) (0.873,-0.408) (0.880,-0.327) (0.887,-0.245) (0.893,-0.163) (0.900,-0.082) (0.907,-0.163) (0.913,-0.245) (0.920,-0.163) (0.927,-0.245) (0.933,-0.327) (0.940,-0.408) (0.947,-0.490) (0.953,-0.572) (0.960,-0.490) (0.967,-0.408) (0.973,-0.490) (0.980,-0.572) (0.987,-0.490) (0.993,-0.572) (1.000,-0.490)};
  \addplot[thick,color=thmcol] coordinates {(0.000,0.000) (0.007,-0.106) (0.013,-0.204) (0.020,-0.309) (0.027,-0.230) (0.033,-0.259) (0.040,-0.339) (0.047,-0.431) (0.053,-0.397) (0.060,-0.483) (0.067,-0.587) (0.073,-0.537) (0.080,-0.635) (0.087,-0.661) (0.093,-0.661) (0.100,-0.698) (0.107,-0.702) (0.113,-0.593) (0.120,-0.635) (0.127,-0.738) (0.133,-0.888) (0.140,-0.905) (0.147,-0.933) (0.153,-0.912) (0.160,-0.950) (0.167,-0.989) (0.173,-1.048) (0.180,-1.090) (0.187,-1.077) (0.193,-1.108) (0.200,-1.100) (0.207,-0.945) (0.213,-0.905) (0.220,-1.034) (0.227,-0.893) (0.233,-0.864) (0.240,-0.941) (0.247,-0.867) (0.253,-0.866) (0.260,-0.916) (0.267,-0.968) (0.273,-1.049) (0.280,-1.045) (0.287,-0.957) (0.293,-0.984) (0.300,-0.950) (0.307,-0.796) (0.313,-0.896) (0.320,-0.875) (0.327,-0.900) (0.333,-0.986) (0.340,-1.070) (0.347,-1.071) (0.353,-1.036) (0.360,-1.079) (0.367,-1.075) (0.373,-1.016) (0.380,-0.991) (0.387,-0.921) (0.393,-0.883) (0.400,-0.843) (0.407,-0.694) (0.413,-0.643) (0.420,-0.632) (0.427,-0.574) (0.433,-0.650) (0.440,-0.677) (0.447,-0.612) (0.453,-0.561) (0.460,-0.434) (0.467,-0.434) (0.473,-0.553) (0.480,-0.394) (0.487,-0.305) (0.493,-0.391) (0.500,-0.279) (0.507,-0.276) (0.513,-0.424) (0.520,-0.459) (0.527,-0.500) (0.533,-0.370) (0.540,-0.435) (0.547,-0.456) (0.553,-0.478) (0.560,-0.509) (0.567,-0.584) (0.573,-0.566) (0.580,-0.478) (0.587,-0.427) (0.593,-0.503) (0.600,-0.597) (0.607,-0.587) (0.613,-0.645) (0.620,-0.696) (0.627,-0.833) (0.633,-0.674) (0.640,-0.599) (0.647,-0.679) (0.653,-0.605) (0.660,-0.495) (0.667,-0.690) (0.673,-0.735) (0.680,-0.767) (0.687,-0.714) (0.693,-0.723) (0.700,-0.742) (0.707,-0.747) (0.713,-0.596) (0.720,-0.419) (0.727,-0.462) (0.733,-0.538) (0.740,-0.318) (0.747,-0.398) (0.753,-0.445) (0.760,-0.442) (0.767,-0.402) (0.773,-0.318) (0.780,-0.286) (0.787,-0.358) (0.793,-0.316) (0.800,-0.296) (0.807,-0.143) (0.813,-0.144) (0.820,-0.253) (0.827,-0.338) (0.833,-0.220) (0.840,-0.264) (0.847,-0.436) (0.853,-0.483) (0.860,-0.483) (0.867,-0.386) (0.873,-0.469) (0.880,-0.415) (0.887,-0.350) (0.893,-0.407) (0.900,-0.422) (0.907,-0.278) (0.913,-0.137) (0.920,-0.067) (0.927,-0.040) (0.933,0.053) (0.940,0.041) (0.947,0.033) (0.953,-0.037) (0.960,-0.036) (0.967,-0.043) (0.973,0.183) (0.980,0.168) (0.987,0.271) (0.993,0.379) (1.000,0.364)};
  \legend{{$\pm1$ steps},{Gaussian steps}}
\end{axis}
\end{tikzpicture}
\caption{Normalised random walks $B_n(t)=S_{\lfloor nt\rfloor}/\sqrt n$ (two independent simulations, seeds fixed). Already at
$n=150$ the Bernoulli and the Gaussian walk are visually indistinguishable, the point of the animation \file{lec14_4}
(there up to $n=2000$) \ts{24}{14}{21:46}\ts{24}{14}{22:47}.}\label{ch11:fig-donsker}
\end{figure}

\begin{intuition}[Universality: the CLT as a fixed point]
Only the first two moments of the steps survive in the limit; all the rest (skewness, kurtosis, lattice structure) is
washed out at rate $1/\sqrt n$. This is the renormalisation-group picture: $S_n/\sqrt n$ is the coarse-graining of
$S_n$ by a factor $n$, and the Gaussian is the fixed point of this flow, with $\sigma^2$ the only relevant coupling. The
consequence for modelling is what the lecture emphasised: one may use Brownian motion when the microscopic steps are
non-Gaussian, provided they have finite variance and are not too dependent \ts{24}{14}{22:47}\ts{24}{14}{23:52}. The 2013
lecture gives the two physical examples: the pollen grain receives a huge number of tiny left/right pushes from
water molecules (Einstein), and a stock price is pushed up or down by a huge number of tiny buy and sell orders
\ts{13}{17}{17:19}\ts{13}{17}{18:21}\ts{13}{17}{21:30}\ts{13}{17}{23:45}. The limit fails when the variance is infinite (stable L\'evy
processes with jumps) or when the steps are strongly dependent (long memory, fractional Brownian motion).
\end{intuition}

\begin{finance}[Sampling frequency and the model]
An analyst observes prices on a grid of spacing $h$: if the increments are like a random walk, the sampled process at fine
resolution looks like Brownian motion (\cref{ch11:thm-donsker} with $n\sim1/h$) \ts{13}{17}{17:19}. Two warnings. At very short horizons
(ticks) market microstructure (bid--ask bounce) makes the increments negatively correlated, so the limit is not
Brownian; and price jumps have heavy tails, so ``finite variance'' can fail (\cref{ch03:sec-heavy}). Universality is a
statement about the \emph{long-horizon} behaviour of the sum of many small steps.
\end{finance}

\paragraph{Where the theorem is used.} Because $F(B_n)\to F(B)$ for continuous $F$, asymptotic distributions of statistics
of random walks are the distributions of functionals of Brownian motion. For the simple random walk one even has an exact
discrete analogue of the maximum formula of \cref{ch11:sec-reflection}: $\Prob[\max_{k\le n}S_k\ge a]=2\Prob[S_n>a]+\Prob[S_n=a]$
(checked exactly for $n=100$, $a=10$: both sides equal $0.3197$), so $\Prob[\max_{k\le n}S_k\ge a\sqrt n]\to2(1-\Phi(a))$, the
Brownian result. Similarly $n^{-2}\sum_kS_{k-1}^2\to\int_0^1B(t)^2\dd t$ and $n^{-1}\sum_kS_{k-1}X_k\to\int_0^1B\dd B$,
which is how the Dickey--Fuller limit of \cref{ch09:eq-df} (\cref{ch09:sec-df}) is obtained; the Kolmogorov--Smirnov statistic
of \cref{ch11:sec-bridge} is another example.

% ==================================================================
\section{Sample paths: roughness and quadratic variation}\label{ch11:sec-paths}

\subsection{Size, zeros and the envelope}
Almost every Brownian path has the following features (2013 lecture): it \emph{crosses the $t$-axis infinitely often}, in
fact it has zeros accumulating at $t=0$ and at every later zero, and it neither drifts off nor stays near zero; at time $t_0$
it is $N(0,t_0)$, so it stays ``around'' the curves $\pm\sqrt t$ (typical size $\sqrt{t_0}$, never $100$ times that in
practice) \ts{13}{17}{25:48}\ts{13}{17}{26:48}. The precise version of ``it does not deviate too much'' is the law of the
iterated logarithm (not proved in the course): $\limsup_{t\to\infty}B(t)/\sqrt{2t\ln\ln t}=1$ almost surely, so the
$\sqrt t$ envelope is exceeded infinitely often, but only by slowly growing factors (compare \cref{ch04:sec-rw}).

\subsection{Non-differentiability}
The path is continuous by definition, but it is nowhere differentiable, which was ``really surprising'' in the words of the 2013 lecture,
and matters because the whole of classical calculus is built on derivatives \ts{13}{17}{27:48}\ts{13}{17}{28:53}.

\begin{proposition}[No derivative]\label{ch11:prop-nodiff}
Fix $t_0\ge0$. With probability one, $B$ is not differentiable at $t_0$. In fact, with probability one there is no
$t_0$ at which $B$ is differentiable (Paley--Wiener--Zygmund, 1933).
\end{proposition}
\begin{proof}[Proof of the first statement]
The heuristic of the slides \ts{24}{14}{1:02:33}\ts{24}{14}{1:03:39}: the difference quotient $D(\Delta t)=[B(t_0+\Delta t)-B(t_0)]/\Delta t$
is $N(0,1/\Delta t)$, so for every fixed $M<\infty$, $\Prob(|D|\le M)=\Prob(|N(0,1)|\le M\sqrt{\Delta t})\le M\sqrt{2\Delta t/\pi}\to0$:
the slope does not settle down; it is of order $(\Delta t)^{-1/2}$ (equivalently $\Delta B=O_P(\sqrt{\Delta t})\gg\Delta t$).
To make this a proof, note that a finite derivative at $t_0$ forces, for some integers $K,m$, $|B(t_0+s)-B(t_0)|\le Ks$ for all
$0<s<1/m$; the event $E_{K,m}$ so defined satisfies $\Prob(E_{K,m})\le\Prob(|B(t_0+\delta)-B(t_0)|\le K\delta)\le K\sqrt{2\delta/\pi}$
for every $\delta<1/m$, hence $\Prob(E_{K,m})=0$; take the countable union over $K,m$ \ts{13}{17}{43:15}\ts{13}{17}{44:19}\ts{13}{17}{45:23}.
The statement for all $t_0$ simultaneously needs a covering argument and is quoted.
\end{proof}

\begin{coursenote}[Erratum in the 2013 proof of non-differentiability]
The written proof (\file{lecnote17}, Proposition~2.4) bounds the event by $\Prob(M(\delta)\le A\delta)$ for the maximum
$M(\delta)=\max_{s\le\delta}B(s)$, and computes it as ``$2(1-\Phi(A\sqrt\delta))$, which tends to zero''. That value tends to $1$, not
$0$. By \cref{ch11:prop-max} the correct value is $\Prob(M(\delta)\le A\delta)=2\Phi(A\sqrt\delta)-1$, which does tend to $0$; the
lecture itself stumbles over the direction of the inequality at this point \ts{13}{17}{46:23}. Both versions are statements for a fixed $t_0$; the slide of 2024 (``not differentiable at $t$, for any $t\ge0$'') is likewise a statement for each fixed $t$.
The argument above avoids the maximum altogether.
\end{coursenote}

\subsection{Quadratic variation}
Let $\Pi=\{0=t_0<t_1<\dots<t_n=T\}$ be a partition of $[0,T]$ with mesh $|\Pi|=\max_j(t_j-t_{j-1})$ and define the
$\Pi$-quadratic variation
\ts{24}{14}{1:03:39}\ts{24}{14}{1:04:44}\ts{24}{14}{1:05:45}
\[
  QV_\Pi([0,T])=\sum_{j=1}^{n}\bigl[B(t_j)-B(t_{j-1})\bigr]^2 .
\]

\begin{theorem}[Quadratic variation of Brownian motion]\label{ch11:thm-qv}
$QV_\Pi([0,T])\to T$ as $|\Pi|\to0$ in $L^2$ (hence in probability), and almost surely along any sequence of partitions with
$\sum_k|\Pi_k|<\infty$, for instance the dyadic partitions \ts{24}{14}{1:06:57}\ts{13}{17}{47:24}.
\end{theorem}
\begin{proof}
Write $\Delta_j=t_j-t_{j-1}$ and $\Delta_jB=B(t_j)-B(t_{j-1})\sim N(0,\Delta_j)$, independent over $j$. Then
$\E[(\Delta_jB)^2]=\Delta_j$, $\Var[(\Delta_jB)^2]=3\Delta_j^2-\Delta_j^2=2\Delta_j^2$ (the fourth moment of $N(0,\Delta_j)$ is $3\Delta_j^2$), so
\[
  \E[Q]V_\Pi=\sum_j\Delta_j=T,\qquad
  \Var\,QV_\Pi=2\sum_j\Delta_j^2\le2|\Pi|\sum_j\Delta_j=2|\Pi|\,T\xrightarrow[|\Pi|\to0]{}0 .
\]
For equally spaced points $QV_\Pi\eqd(T/n)\,\chi^2_n$ \ts{24}{14}{1:08:03}, mean $T$ and variance $2T^2/n$. Chebyshev's inequality gives
convergence in probability; if $\sum_k|\Pi_k|<\infty$ then $\sum_k\Prob(|QV_{\Pi_k}-T|>\varepsilon)\le\sum_k2|\Pi_k|T/\varepsilon^2<\infty$ and
Borel--Cantelli gives almost sure convergence.
\end{proof}

\begin{coursenote}[Erratum and remarks on the proof in the notes and slides]
\begin{itemize}
  \item The 2013 proof invokes the law of large numbers for the sum of $n$ i.i.d.\ variables $(B(t_{i+1})-B(t_i))^2$ of mean $T/n$; the
      summands change with $n$ (a triangular array), so the classical law does not apply verbatim, and the mean and variance computation above (Chebyshev)
      is the clean way; the quantity converges in $L^2$ and, for nested partitions, almost surely. The lecture's own computation
      had a slip corrected by a student (the mean of $X_i^2$ is $T/n$, not $T^2/n^2$) \ts{13}{17}{56:07}.
  \item Slide~19 of \file{lec14\_1} (PDF pages 73--74) says ``$QV([0,T])=1$ for virtually every path''; the value is $T$ (it equals $1$ only for $T=1$), as the
      claim on the same slide states.
  \item On the slide of the Brownian bridge (slide~13, \cref{ch11:sec-bridge}) ``Moments of $B(t)$'' should read
      ``moments of $X(t)$''.
  \item Slide~2 dates Einstein's diffusion paper to 1904; it appeared in 1905 (as in the 2013 note and in
      our Motivation).
  \item The diffusion equation on slide~6 (PDF page 15) is written $\partial_tp=\tfrac12\partial_y^2p$, which is the case $\sigma=1$ announced there
      (``easy to verify for $s=0$ and $\sigma=1$''); for a general diffusion coefficient the factor is $\tfrac12\sigma^2$, \cref{ch11:prop-heat}.
\end{itemize}
\end{coursenote}

The theorem says that the \emph{sum of squared increments is not random and does not depend on the path}: $[B]_T=T$ for
almost every path \ts{24}{14}{1:06:57}\ts{24}{14}{1:09:17}. Compare a continuously differentiable function $f$: by the mean value theorem
$\sum_j(f(t_j)-f(t_{j-1}))^2=\sum_jf'(s_j)^2\Delta_j^2\le\max|f'|^2\,|\Pi|\,T\to0$, so its quadratic variation is zero
\ts{13}{17}{48:29}\ts{13}{17}{49:39}\ts{13}{17}{50:42}\ts{13}{17}{52:44}. Brownian paths ``bounce back and forth too much'': the squared moves
do not vanish as the grid is refined. Two corollaries.

\begin{corollary}[Infinite total variation]\label{ch11:cor-tv}
Almost surely $\sum_j|B(t_j)-B(t_{j-1})|\to\infty$ as $|\Pi|\to0$.
\end{corollary}
\begin{proof}
$\sum_j\Delta_jB^2\le\max_j|\Delta_jB|\cdot\sum_j|\Delta_jB|$. The left side tends to $T>0$ and, by uniform continuity of the
path, $\max_j|\Delta_jB|\to0$; hence the sum of absolute increments diverges.
\end{proof}
So one cannot define $\int f\dd B$ pathwise as a Riemann--Stieltjes integral, and the stochastic (It\^o) integral of \cref{ch:stochcalc} is needed.

The rule of thumb is written
\begin{equation}\label{ch11:eq-dB2}
  (\dd B)^2=\dd t,\qquad \dd B\,\dd t=0,\qquad (\dd t)^2=0,
\end{equation}
where $\dd B\sim\sqrt{\dd t}$ carries the whole effect at second order: the increments of a Brownian path are of order
$\sqrt{\Delta t}$, their squares of order $\Delta t$ (summing to a finite number), while a product $\Delta B\,\Delta t$
is $O(\Delta t^{3/2})$ and vanishes on summation \ts{13}{17}{53:49}. This is the reason why the Taylor expansion of a function of $B$ has a second-order term (\cref{ch11:sec-gbm}).

\begin{figure}[ht]
\centering
\begin{tikzpicture}
\begin{axis}[width=0.95\linewidth,height=5cm,xlabel={$t$},ylabel={$QV_n(t)$},xmin=0,xmax=1,ymin=0,ymax=1.4,
  legend style={at={(0.5,1.03)},anchor=south,legend columns=3,font=\footnotesize},grid=major,grid style={black!10}]
  \addplot[thick,color=defcol] coordinates {(0.000,0.000) (0.000,0.000) (0.020,0.083) (0.020,0.083) (0.040,0.214) (0.040,0.214) (0.060,0.217) (0.060,0.217) (0.080,0.224) (0.080,0.224) (0.100,0.228) (0.100,0.228) (0.120,0.229) (0.120,0.229) (0.140,0.311) (0.140,0.311) (0.160,0.312) (0.160,0.312) (0.180,0.327) (0.180,0.327) (0.200,0.547) (0.200,0.547) (0.220,0.548) (0.220,0.548) (0.240,0.551) (0.240,0.551) (0.260,0.552) (0.260,0.552) (0.280,0.561) (0.280,0.561) (0.300,0.584) (0.300,0.584) (0.320,0.587) (0.320,0.587) (0.340,0.591) (0.340,0.591) (0.360,0.593) (0.360,0.593) (0.380,0.611) (0.380,0.611) (0.400,0.612) (0.400,0.612) (0.420,0.612) (0.420,0.612) (0.440,0.659) (0.440,0.659) (0.460,0.665) (0.460,0.665) (0.480,0.671) (0.480,0.671) (0.500,0.671) (0.500,0.671) (0.520,0.677) (0.520,0.677) (0.540,0.752) (0.540,0.752) (0.560,0.753) (0.560,0.753) (0.580,0.755) (0.580,0.755) (0.600,0.775) (0.600,0.775) (0.620,0.790) (0.620,0.790) (0.640,0.792) (0.640,0.792) (0.660,0.808) (0.660,0.808) (0.680,0.814) (0.680,0.814) (0.700,0.815) (0.700,0.815) (0.720,0.824) (0.720,0.824) (0.740,0.984) (0.740,0.984) (0.760,1.004) (0.760,1.004) (0.780,1.023) (0.780,1.023) (0.800,1.078) (0.800,1.078) (0.820,1.080) (0.820,1.080) (0.840,1.090) (0.840,1.090) (0.860,1.094) (0.860,1.094) (0.880,1.117) (0.880,1.117) (0.900,1.117) (0.900,1.117) (0.920,1.117) (0.920,1.117) (0.940,1.157) (0.940,1.157) (0.960,1.168) (0.960,1.168) (0.980,1.168) (0.980,1.168) (1.000,1.193)};
  \addplot[thick,color=thmcol] coordinates {(0.000,0.000) (0.010,0.009) (0.020,0.025) (0.030,0.038) (0.040,0.045) (0.050,0.052) (0.060,0.058) (0.070,0.065) (0.080,0.083) (0.090,0.097) (0.100,0.103) (0.110,0.110) (0.120,0.116) (0.130,0.133) (0.140,0.146) (0.150,0.151) (0.160,0.154) (0.170,0.171) (0.180,0.180) (0.190,0.193) (0.200,0.199) (0.210,0.216) (0.220,0.224) (0.230,0.234) (0.240,0.245) (0.250,0.249) (0.260,0.260) (0.270,0.272) (0.280,0.280) (0.290,0.288) (0.300,0.307) (0.310,0.313) (0.320,0.323) (0.330,0.328) (0.340,0.335) (0.350,0.344) (0.360,0.352) (0.370,0.362) (0.380,0.385) (0.390,0.402) (0.400,0.410) (0.410,0.423) (0.420,0.433) (0.430,0.438) (0.440,0.452) (0.450,0.464) (0.460,0.474) (0.470,0.483) (0.480,0.494) (0.490,0.505) (0.500,0.509) (0.510,0.515) (0.520,0.526) (0.530,0.530) (0.540,0.537) (0.550,0.552) (0.560,0.558) (0.570,0.574) (0.580,0.583) (0.590,0.598) (0.600,0.607) (0.610,0.623) (0.620,0.634) (0.630,0.645) (0.640,0.657) (0.650,0.664) (0.660,0.674) (0.670,0.687) (0.680,0.695) (0.690,0.702) (0.700,0.713) (0.710,0.729) (0.720,0.738) (0.730,0.745) (0.740,0.753) (0.750,0.764) (0.760,0.776) (0.770,0.789) (0.780,0.800) (0.790,0.812) (0.800,0.822) (0.810,0.831) (0.820,0.841) (0.830,0.850) (0.840,0.853) (0.850,0.874) (0.860,0.881) (0.870,0.882) (0.880,0.896) (0.890,0.898) (0.900,0.911) (0.910,0.925) (0.920,0.933) (0.930,0.939) (0.940,0.953) (0.950,0.963) (0.960,0.968) (0.970,0.988) (0.980,0.997) (0.990,1.013) (1.000,1.031)};
  \addplot[dashed,black!70] coordinates {(0,0) (1,1)};
  \legend{{$n=50$ steps},{$n=1000$ steps},{$T=t$}}
\end{axis}
\end{tikzpicture}
\caption{Cumulative quadratic variation $QV_n(t)=\sum_{k\le nt}(\Delta B_k)^2$ of a Gaussian random walk
$\Delta B_k=X_k/\sqrt n$, at two resolutions. At $t=1$ the values are $1.19$ and $1.03$ (theory: standard deviation
$\sqrt{2/n}$, i.e.\ $0.20$ and $0.045$); the path converges to the straight line $t$, as in \file{lec14\_9}
(there $n$ up to $2000$) \ts{24}{14}{1:08:03}\ts{24}{14}{1:09:17}.}\label{ch11:fig-qv}
\end{figure}

\begin{finance}[Realised variance and stochastic volatility]
If the log price is $\log(P_t/P_0)=\mu t+\sigma B(t)$, the sum of squared log returns over a period $[0,T]$ (the \emph{realised
variance}) estimates $\sigma^2T$ \emph{without knowing the drift and without error as the sampling becomes finer}: with
$n$ returns its relative standard deviation is $\sqrt{2/n}$ (about $8.9\%$ for $n=252$ daily returns of one year, so
$8.9\%/2\approx4.5\%$ in volatility terms). This is the theoretical basis for measuring volatility by high-frequency data
(\cref{ch:volatility}). The lecture proposes plotting the cumulative squared returns of an asset against time: for a Brownian model the
graph is a straight line of slope $\sigma^2$; slope changes and curvature reveal periods of low and high volatility and
suggest a model of dynamic volatility \ts{24}{14}{1:10:21}\ts{24}{14}{1:11:22}\ts{24}{14}{1:12:22}. Drift, by contrast, cannot be
learned by refining the sampling: $\hat\mu$ has standard error $\sigma/\sqrt T$ whatever $n$ is (2024 PS4 Problem~3,
\cref{ch:volatility}).
\end{finance}


% ==================================================================
\section{The reflection principle, the maximum and hitting times}\label{ch11:sec-reflection}

Suppose Brownian motion models a stock price during a trading day (9:30 to 16:00). Traders care about the \emph{highest
and lowest} values reached, not only the closing value \ts{13}{17}{31:00}: barrier options pay if a level is touched, a stop-loss triggers when it is
reached, a firm defaults when its value hits a threshold. The two objects are the running maximum and the first hitting time.
For a standard Brownian motion $B$ and levels $a>0$ define
\begin{equation}\label{ch11:eq-defM}
  M(t)=\max_{0\le u\le t}B(u),\qquad \tau_a=\min\{u\ge0:\ B(u)=a\}.
\end{equation}
The maximum exists because a continuous function attains its maximum on the compact set $[0,t]$ \ts{24}{14}{28:05}\ts{13}{17}{31:00}.
Since the path is continuous and starts at $0<a$, hitting $a$ before $t$ and having a maximum $\ge a$ are the same event:
\begin{equation}\label{ch11:eq-tauM}
  \{\tau_a\le t\}=\{M(t)\ge a\}.
\end{equation}
$\tau_a$ is a \emph{stopping time} (\cref{ch04:def-stop}): whether $\tau_a\le t$ is decided by the path up to time $t$
\ts{24}{14}{24:53}\ts{13}{17}{33:04}.

\begin{lemma}[Reflection principle]\label{ch11:lem-reflection}
Let $\widetilde B(u)=B(u)$ for $u\le\tau_a$ and $\widetilde B(u)=a-\bigl(B(u)-a\bigr)=2a-B(u)$ for $u>\tau_a$: the path is
reflected in the level $a$ from the first hitting time on. Then $\widetilde B$ is again a standard Brownian motion
\ts{24}{14}{24:53}\ts{24}{14}{25:57}.
\end{lemma}
The proof uses the \emph{strong Markov property}: at the stopping time $\tau_a$ the process $B(\tau_a+u)-B(\tau_a)$ is a new
Brownian motion independent of the past $\mathcal F_{\tau_a}$; the Markov property of \cref{ch11:prop-elem}(ii) still holds when the fixed time $s$ is replaced by a
stopping time, because $\tau_a$ can be approximated by discrete stopping times (\cref{ch04:thm-stopped}) (2013 note: ``the
distribution of $B(t)-B(\tau_a)$ is not affected by the fact that we conditioned on $\tau_a<t$'') \ts{13}{17}{35:54}. Since $-B$ is
also a Brownian motion, gluing the old path up to $\tau_a$ to the reflected new increment $-(B(\tau_a+u)-a)$ gives
a process with the same law. The lemma is the statement that after touching the level the path is equally likely to go
up or down \ts{24}{14}{24:53}\ts{24}{14}{25:57}\ts{24}{14}{26:57}; \cref{ch11:fig-refl} shows one path and its reflection (the deck
\file{lec14\_5} shows it for $x=50$ and $10\,000$ steps).

\begin{figure}[ht]
\centering
\begin{tikzpicture}
\begin{axis}[width=0.95\linewidth,height=5.4cm,xlabel={$t$},xmin=0,xmax=1,ymin=-0.7,ymax=1.5,
  legend style={at={(0.5,1.03)},anchor=south,legend columns=3,font=\footnotesize},grid=major,grid style={black!10}]
  \addplot[thick,color=defcol] coordinates {(0.000,0.000) (0.005,-0.042) (0.010,0.003) (0.015,0.076) (0.020,0.149) (0.025,0.277) (0.030,0.250) (0.035,0.289) (0.040,0.263) (0.045,0.162) (0.050,0.112) (0.055,0.122) (0.060,0.057) (0.065,0.044) (0.070,0.123) (0.075,0.163) (0.080,0.204) (0.085,0.228) (0.090,0.216) (0.095,0.084) (0.100,0.154) (0.105,0.047) (0.110,0.062) (0.115,0.055) (0.120,0.065) (0.125,0.082) (0.130,0.059) (0.135,0.122) (0.140,0.032) (0.145,0.088) (0.150,-0.032) (0.155,0.052) (0.160,0.016) (0.165,0.042) (0.170,0.149) (0.175,-0.004) (0.180,-0.026) (0.185,0.014) (0.190,0.078) (0.195,0.174) (0.200,0.217) (0.205,0.253) (0.210,0.261) (0.215,0.279) (0.220,0.284) (0.225,0.355) (0.230,0.366) (0.235,0.413) (0.240,0.527) (0.245,0.610) (0.250,0.586) (0.255,0.579) (0.260,0.631) (0.265,0.585) (0.270,0.456) (0.275,0.384) (0.280,0.422) (0.285,0.407) (0.290,0.382) (0.295,0.435) (0.300,0.361) (0.305,0.521) (0.310,0.522) (0.315,0.499) (0.320,0.643) (0.325,0.566) (0.330,0.531) (0.335,0.525) (0.340,0.457) (0.345,0.433) (0.350,0.474) (0.355,0.486) (0.360,0.483) (0.365,0.553) (0.370,0.686) (0.375,0.797) (0.380,0.648) (0.385,0.745) (0.390,0.800) (0.395,0.829) (0.400,0.760) (0.405,0.728) (0.410,0.713) (0.415,0.614) (0.420,0.702) (0.425,0.765) (0.430,0.764) (0.435,0.697) (0.440,0.792) (0.445,0.837) (0.450,0.912) (0.455,0.883) (0.460,1.052) (0.465,1.042) (0.470,1.181) (0.475,1.039) (0.480,1.061) (0.485,1.066) (0.490,0.888) (0.495,0.842) (0.500,0.756) (0.505,0.656) (0.510,0.603) (0.515,0.520) (0.520,0.543) (0.525,0.576) (0.530,0.595) (0.535,0.537) (0.540,0.502) (0.545,0.621) (0.550,0.597) (0.555,0.595) (0.560,0.554) (0.565,0.591) (0.570,0.654) (0.575,0.710) (0.580,0.740) (0.585,0.790) (0.590,0.775) (0.595,0.819) (0.600,0.749) (0.605,0.756) (0.610,0.817) (0.615,0.746) (0.620,0.801) (0.625,0.823) (0.630,0.737) (0.635,0.616) (0.640,0.499) (0.645,0.425) (0.650,0.448) (0.655,0.465) (0.660,0.552) (0.665,0.713) (0.670,0.679) (0.675,0.535) (0.680,0.429) (0.685,0.286) (0.690,0.237) (0.695,0.199) (0.700,0.289) (0.705,0.182) (0.710,0.164) (0.715,0.146) (0.720,0.137) (0.725,0.096) (0.730,0.245) (0.735,0.232) (0.740,0.235) (0.745,0.284) (0.750,0.300) (0.755,0.425) (0.760,0.339) (0.765,0.325) (0.770,0.289) (0.775,0.283) (0.780,0.346) (0.785,0.303) (0.790,0.164) (0.795,0.168) (0.800,0.094) (0.805,0.137) (0.810,0.153) (0.815,0.301) (0.820,0.339) (0.825,0.280) (0.830,0.305) (0.835,0.381) (0.840,0.386) (0.845,0.384) (0.850,0.365) (0.855,0.336) (0.860,0.270) (0.865,0.356) (0.870,0.384) (0.875,0.282) (0.880,0.358) (0.885,0.360) (0.890,0.376) (0.895,0.387) (0.900,0.471) (0.905,0.314) (0.910,0.312) (0.915,0.352) (0.920,0.343) (0.925,0.226) (0.930,0.137) (0.935,0.152) (0.940,0.253) (0.945,0.116) (0.950,0.075) (0.955,0.098) (0.960,0.116) (0.965,0.156) (0.970,0.239) (0.975,0.295) (0.980,0.193) (0.985,0.174) (0.990,0.098) (0.995,0.100) (1.000,0.209)};
  \addplot[thick,dashed,color=thmcol] coordinates {(0.390,0.800) (0.390,0.800) (0.395,0.771) (0.400,0.840) (0.405,0.872) (0.410,0.887) (0.415,0.986) (0.420,0.898) (0.425,0.835) (0.430,0.836) (0.435,0.903) (0.440,0.808) (0.445,0.763) (0.450,0.688) (0.455,0.717) (0.460,0.548) (0.465,0.558) (0.470,0.419) (0.475,0.561) (0.480,0.539) (0.485,0.534) (0.490,0.712) (0.495,0.758) (0.500,0.844) (0.505,0.944) (0.510,0.997) (0.515,1.080) (0.520,1.057) (0.525,1.024) (0.530,1.005) (0.535,1.063) (0.540,1.098) (0.545,0.979) (0.550,1.003) (0.555,1.005) (0.560,1.046) (0.565,1.009) (0.570,0.946) (0.575,0.890) (0.580,0.860) (0.585,0.810) (0.590,0.825) (0.595,0.781) (0.600,0.851) (0.605,0.844) (0.610,0.783) (0.615,0.854) (0.620,0.799) (0.625,0.777) (0.630,0.863) (0.635,0.984) (0.640,1.101) (0.645,1.175) (0.650,1.152) (0.655,1.135) (0.660,1.048) (0.665,0.887) (0.670,0.921) (0.675,1.065) (0.680,1.171) (0.685,1.314) (0.690,1.363) (0.695,1.401) (0.700,1.311) (0.705,1.418) (0.710,1.436) (0.715,1.454) (0.720,1.463) (0.725,1.504) (0.730,1.355) (0.735,1.368) (0.740,1.365) (0.745,1.316) (0.750,1.300) (0.755,1.175) (0.760,1.261) (0.765,1.275) (0.770,1.311) (0.775,1.317) (0.780,1.254) (0.785,1.297) (0.790,1.436) (0.795,1.432) (0.800,1.506) (0.805,1.463) (0.810,1.447) (0.815,1.299) (0.820,1.261) (0.825,1.320) (0.830,1.295) (0.835,1.219) (0.840,1.214) (0.845,1.216) (0.850,1.235) (0.855,1.264) (0.860,1.330) (0.865,1.244) (0.870,1.216) (0.875,1.318) (0.880,1.242) (0.885,1.240) (0.890,1.224) (0.895,1.213) (0.900,1.129) (0.905,1.286) (0.910,1.288) (0.915,1.248) (0.920,1.257) (0.925,1.374) (0.930,1.463) (0.935,1.448) (0.940,1.347) (0.945,1.484) (0.950,1.525) (0.955,1.502) (0.960,1.484) (0.965,1.444) (0.970,1.361) (0.975,1.305) (0.980,1.407) (0.985,1.426) (0.990,1.502) (0.995,1.500) (1.000,1.391)};
  \addplot[thin,black!60] coordinates {(0,0.8) (1,0.8)};
  \addplot[thin,dotted,black!60] coordinates {(0.390,-0.7) (0.390,0.8)};
  \node[font=\footnotesize,anchor=north west] at (axis cs:0.395,-0.42) {$\tau_a$};
  \node[font=\footnotesize,anchor=south west] at (axis cs:0.01,0.82) {level $a=0.8$};
  \legend{{path $B$},{reflected path $\widetilde B$ after $\tau_a$}}
\end{axis}
\end{tikzpicture}
\caption{A simulated Brownian path ($200$ steps, $a=0.8$) and its reflection in the level $a$ after the first hitting time
$\tau_a$. The two paths have the same probability; the original ends below $a$, its reflection above.}\label{ch11:fig-refl}
\end{figure}

\begin{proposition}[Distribution of the maximum]\label{ch11:prop-max}
For $t>0$ and $a>0$,
\begin{equation}\label{ch11:eq-max}
  \Prob[M(t)\ge a]=2\,\Prob[B(t)>a]=2\Bigl[1-\Phi\Bigl(\frac a{\sqrt t}\Bigr)\Bigr].
\end{equation}
Equivalently $M(t)\eqd|B(t)|$, with density $\sqrt{2/(\pi t)}\,e^{-m^2/(2t)}$ on $m>0$, and $\E[M(t)]=\sqrt{2t/\pi}$.
\end{proposition}
\begin{proof}
Two views, both used in the lectures.
\emph{Path pairing} \ts{24}{14}{28:05}\ts{24}{14}{29:20}\ts{24}{14}{30:27}. Let $A_{\mathrm{END}}=\{B(t)>a\}$ and $A_{\mathrm{MAX}}=\{M(t)\ge a\}$.
By continuity $A_{\mathrm{END}}\subset A_{\mathrm{MAX}}$ (a path ending above $a$ has crossed it, so $\tau_a<t$). Every
$\omega_0\in A_{\mathrm{END}}$ is paired with its reflection $\omega_1$ (\cref{ch11:lem-reflection}), which
ends below $a$ but has still maximum $\ge a$, and has the same probability; conversely every path in $A_{\mathrm{MAX}}\setminus A_{\mathrm{END}}$
arises this way. Hence $\Prob[A_{\mathrm{MAX}}]=\Prob[A_{\mathrm{MAX}}\cap A_{\mathrm{END}}]+\Prob[A_{\mathrm{MAX}}\cap A_{\mathrm{END}}^c]
=\Prob[A_{\mathrm{END}}]+\Prob[A_{\mathrm{END}}]$ (the event $B(t)=a$ has probability $0$).
\emph{Conditioning on $\tau_a$} \ts{13}{17}{34:09}\ts{13}{17}{36:56}\ts{13}{17}{40:08}. Given $\tau_a<t$, the strong Markov property says $B(t)-B(\tau_a)=B(t)-a$
is symmetric, so $\Prob[B(t)>a\mid\tau_a<t]=\Prob[B(t)<a\mid\tau_a<t]=\tfrac12$. Therefore
$\Prob[M(t)\ge a]=\Prob[\tau_a<t]=2\Prob[B(t)>a,\ \tau_a<t]=2\Prob[B(t)>a]$, the last step because $B(t)>a$ already
implies $\tau_a<t$.
The rest follows from $\Prob[B(t)>a]=1-\Phi(a/\sqrt t)$, by differentiating in $a$, and by $\E[M]=\int_0^\infty\Prob[M>a]\dd a$.
\end{proof}

\begin{remark}[Joint law of maximum and endpoint; not discussed in class]\label{ch11:rem-joint}
The same reflection gives, for $m>0$ and $x\le m$, $\Prob[M(t)\ge m,\ B(t)\le x]=\Prob[B(t)\ge2m-x]=1-\Phi\bigl((2m-x)/\sqrt t\bigr)$
(reflect the part of the path after $\tau_m$: the event $\{B(t)\le x\}$ becomes $\{\widetilde B(t)\ge2m-x\}$).
For $x=m$ this is \eqref{ch11:eq-max}. By symmetry, the minimum satisfies $\Prob[\min_{u\le t}B(u)\le-a]=2[1-\Phi(a/\sqrt t)]$
as well.
\end{remark}

\begin{finance}[Barrier events: a path is twice as likely to touch a level as to end beyond it]
Let $\log(S_t/S_0)=\sigma B(t)$ (no drift, $\sigma=30\%$ per year). The probability that the price is below $S_0e^{-0.2}\approx0.82\,S_0$ \emph{at the end of
one year} is $\Phi(-0.2/0.3)=25.2\%$; the probability that the stock \emph{touches} that level at some time during the year is
$2\Phi(-0.2/0.3)=50.5\%$ (\cref{ch11:prop-max} applied to $-B$). For a knock-out or knock-in
barrier option, or a stop-loss, the touching probability is what matters, and ``double the terminal probability'' is
the first-order rule; a drift or discrete (say daily) monitoring changes the number, the latter by lowering it
(a path can cross between two observations without being seen) \ts{13}{17}{31:00}. Range-based volatility estimators
combine the maximum and the minimum of the log price over a day: with the range $R=\max-\min$ of a driftless Brownian path one has
$\E[R^2]=4\ln2\cdot\sigma^2t$ (Parkinson 1980, derived from the joint law of \cref{ch11:rem-joint}; it is the subject of 2024 PS5 Problem~3 and
of \cref{ch:volatility}); in the 2013 lecture Kempthorne, from the audience, remarks that he uses the range of the process as a more precise measure of volatility than close-to-close prices
\ts{13}{17}{39:05}\ts{13}{17}{40:08}.
\end{finance}

\subsection{First hitting time}
By \eqref{ch11:eq-tauM} and \cref{ch11:prop-max},
\begin{equation}\label{ch11:eq-hitcdf}
  \Prob[\tau_a\le t]=2\Bigl[1-\Phi\Bigl(\frac a{\sqrt t}\Bigr)\Bigr],\qquad
  f_{\tau_a}(t)=\frac{a}{\sqrt{2\pi}}\,t^{-3/2}\,e^{-a^2/(2t)},\quad0<t<\infty,
\end{equation}
\sloppy the density being the derivative in $t$ (using $\frac{\dd}{\dd t}\Phi(a t^{-1/2})=-\tfrac12 a t^{-3/2}\varphi(a t^{-1/2})$)
\ts{24}{14}{31:35}\ts{24}{14}{32:37}\ts{24}{14}{33:44}. This is the L\'evy distribution.

\begin{proposition}[Properties of $\tau_a$]\label{ch11:prop-tau}
\leavevmode
\begin{enumerate}[(i)]
  \item $\tau_a<\infty$ almost surely: a Brownian path reaches every level.
  \item $\E[\tau_a]=\infty$: the tail is $\Prob[\tau_a>t]=2\Phi(a/\sqrt t)-1\sim a\sqrt{2/(\pi t)}$ and $\int^\infty t^{-1/2}\dd t=\infty$.
  \item The density is unimodal with mode at $t=a^2/3$, and its median is $a^2/\bigl(\Phi^{-1}(3/4)\bigr)^2\approx2.198\,a^2$.
  \item Scaling: $\tau_a\eqd a^2\tau_1$.
\end{enumerate}
\end{proposition}
\begin{proof}
(i) $\Prob[\tau_a\le t]\to2[1-\Phi(0)]=1$ as $t\to\infty$.

(ii) as stated.

(iii) $\frac{\dd}{\dd t}\ln f=-\tfrac{3}{2t}+\tfrac{a^2}{2t^2}=0$ at $t=a^2/3$;
the median solves $2[1-\Phi(a/\sqrt t)]=\tfrac12$, i.e.\ $a/\sqrt t=\Phi^{-1}(3/4)=0.6745$.

(iv) from \eqref{ch11:eq-hitcdf}, which depends on $a/\sqrt t$ only.
\end{proof}

\begin{figure}[ht]
\centering
\begin{tikzpicture}
\begin{axis}[width=0.47\linewidth,height=5cm,xlabel={$t$},ylabel={$\Prob[M(t)>x]$},xmin=0,xmax=5,ymin=0,ymax=0.7,
  title={Maximum exceeds $x$ before $t$},title style={font=\footnotesize},
  legend style={at={(0.03,0.97)},anchor=north west,font=\scriptsize},grid=major,grid style={black!10}]
  \addplot[thick,color=defcol] coordinates {(0.040,0.000) (0.141,0.008) (0.242,0.042) (0.344,0.088) (0.445,0.134) (0.546,0.176) (0.647,0.214) (0.749,0.248) (0.850,0.278) (0.951,0.305) (1.052,0.330) (1.153,0.352) (1.255,0.372) (1.356,0.390) (1.457,0.407) (1.558,0.423) (1.660,0.438) (1.761,0.451) (1.862,0.464) (1.963,0.475) (2.064,0.486) (2.166,0.497) (2.267,0.507) (2.368,0.516) (2.469,0.525) (2.571,0.533) (2.672,0.541) (2.773,0.548) (2.874,0.555) (2.976,0.562) (3.077,0.569) (3.178,0.575) (3.279,0.581) (3.380,0.587) (3.482,0.592) (3.583,0.597) (3.684,0.602) (3.785,0.607) (3.887,0.612) (3.988,0.617) (4.089,0.621) (4.190,0.625) (4.291,0.629) (4.393,0.633) (4.494,0.637) (4.595,0.641) (4.696,0.644) (4.798,0.648) (4.899,0.651) (5.000,0.655)};
  \addplot[thick,color=intcol] coordinates {(0.040,0.000) (0.141,0.000) (0.242,0.000) (0.344,0.001) (0.445,0.003) (0.546,0.007) (0.647,0.013) (0.749,0.021) (0.850,0.030) (0.951,0.040) (1.052,0.051) (1.153,0.063) (1.255,0.074) (1.356,0.086) (1.457,0.098) (1.558,0.109) (1.660,0.121) (1.761,0.132) (1.862,0.143) (1.963,0.153) (2.064,0.164) (2.166,0.174) (2.267,0.184) (2.368,0.194) (2.469,0.203) (2.571,0.212) (2.672,0.221) (2.773,0.230) (2.874,0.238) (2.976,0.246) (3.077,0.254) (3.178,0.262) (3.279,0.269) (3.380,0.277) (3.482,0.284) (3.583,0.291) (3.684,0.297) (3.785,0.304) (3.887,0.310) (3.988,0.317) (4.089,0.323) (4.190,0.329) (4.291,0.334) (4.393,0.340) (4.494,0.345) (4.595,0.351) (4.696,0.356) (4.798,0.361) (4.899,0.366) (5.000,0.371)};
  \addplot[thick,color=thmcol] coordinates {(0.040,0.000) (0.141,0.000) (0.242,0.000) (0.344,0.000) (0.445,0.000) (0.546,0.000) (0.647,0.000) (0.749,0.001) (0.850,0.001) (0.951,0.002) (1.052,0.003) (1.153,0.005) (1.255,0.007) (1.356,0.010) (1.457,0.013) (1.558,0.016) (1.660,0.020) (1.761,0.024) (1.862,0.028) (1.963,0.032) (2.064,0.037) (2.166,0.041) (2.267,0.046) (2.368,0.051) (2.469,0.056) (2.571,0.061) (2.672,0.066) (2.773,0.072) (2.874,0.077) (2.976,0.082) (3.077,0.087) (3.178,0.092) (3.279,0.098) (3.380,0.103) (3.482,0.108) (3.583,0.113) (3.684,0.118) (3.785,0.123) (3.887,0.128) (3.988,0.133) (4.089,0.138) (4.190,0.143) (4.291,0.148) (4.393,0.152) (4.494,0.157) (4.595,0.162) (4.696,0.166) (4.798,0.171) (4.899,0.175) (5.000,0.180)};
  \legend{{$x=1$},{$x=2$},{$x=3$}}
\end{axis}
\end{tikzpicture}\hfill
\begin{tikzpicture}
\begin{axis}[width=0.47\linewidth,height=5cm,xlabel={$t$},ylabel={$f(t\mid x)$},xmin=0,xmax=5,ymin=0,ymax=0.5,
  title={Density of the hitting time $\tau_x$},title style={font=\footnotesize},
  legend style={at={(0.97,0.97)},anchor=north east,font=\scriptsize},grid=major,grid style={black!10}]
  \addplot[thick,color=defcol] coordinates {(0.050,0.002) (0.134,0.195) (0.218,0.395) (0.302,0.459) (0.386,0.456) (0.469,0.428) (0.553,0.393) (0.637,0.358) (0.721,0.326) (0.805,0.297) (0.889,0.271) (0.973,0.249) (1.057,0.229) (1.141,0.211) (1.225,0.196) (1.308,0.182) (1.392,0.170) (1.476,0.159) (1.560,0.149) (1.644,0.140) (1.728,0.132) (1.812,0.124) (1.896,0.117) (1.980,0.111) (2.064,0.106) (2.147,0.100) (2.231,0.096) (2.315,0.091) (2.399,0.087) (2.483,0.083) (2.567,0.080) (2.651,0.077) (2.735,0.073) (2.819,0.071) (2.903,0.068) (2.986,0.065) (3.070,0.063) (3.154,0.061) (3.238,0.059) (3.322,0.057) (3.406,0.055) (3.490,0.053) (3.574,0.051) (3.658,0.050) (3.742,0.048) (3.825,0.047) (3.909,0.045) (3.993,0.044) (4.077,0.043) (4.161,0.042) (4.245,0.041) (4.329,0.039) (4.413,0.038) (4.497,0.037) (4.581,0.036) (4.664,0.036) (4.748,0.035) (4.832,0.034) (4.916,0.033) (5.000,0.032)};
  \addplot[thick,color=intcol] coordinates {(0.050,0.000) (0.134,0.000) (0.218,0.001) (0.302,0.006) (0.386,0.019) (0.469,0.035) (0.553,0.052) (0.637,0.068) (0.721,0.081) (0.805,0.092) (0.889,0.100) (0.973,0.106) (1.057,0.111) (1.141,0.113) (1.225,0.115) (1.308,0.116) (1.392,0.115) (1.476,0.115) (1.560,0.114) (1.644,0.112) (1.728,0.110) (1.812,0.108) (1.896,0.106) (1.980,0.104) (2.064,0.102) (2.147,0.100) (2.231,0.098) (2.315,0.095) (2.399,0.093) (2.483,0.091) (2.567,0.089) (2.651,0.087) (2.735,0.085) (2.819,0.083) (2.903,0.081) (2.986,0.079) (3.070,0.077) (3.154,0.076) (3.238,0.074) (3.322,0.072) (3.406,0.071) (3.490,0.069) (3.574,0.067) (3.658,0.066) (3.742,0.065) (3.825,0.063) (3.909,0.062) (3.993,0.061) (4.077,0.059) (4.161,0.058) (4.245,0.057) (4.329,0.056) (4.413,0.055) (4.497,0.054) (4.581,0.053) (4.664,0.052) (4.748,0.051) (4.832,0.050) (4.916,0.049) (5.000,0.048)};
  \addplot[thick,color=thmcol] coordinates {(0.050,0.000) (0.134,0.000) (0.218,0.000) (0.302,0.000) (0.386,0.000) (0.469,0.000) (0.553,0.001) (0.637,0.002) (0.721,0.004) (0.805,0.006) (0.889,0.009) (0.973,0.012) (1.057,0.016) (1.141,0.019) (1.225,0.022) (1.308,0.026) (1.392,0.029) (1.476,0.032) (1.560,0.034) (1.644,0.037) (1.728,0.039) (1.812,0.041) (1.896,0.043) (1.980,0.044) (2.064,0.046) (2.147,0.047) (2.231,0.048) (2.315,0.049) (2.399,0.049) (2.483,0.050) (2.567,0.050) (2.651,0.051) (2.735,0.051) (2.819,0.051) (2.903,0.051) (2.986,0.051) (3.070,0.051) (3.154,0.051) (3.238,0.051) (3.322,0.051) (3.406,0.051) (3.490,0.051) (3.574,0.050) (3.658,0.050) (3.742,0.050) (3.825,0.049) (3.909,0.049) (3.993,0.049) (4.077,0.048) (4.161,0.048) (4.245,0.047) (4.329,0.047) (4.413,0.047) (4.497,0.046) (4.581,0.046) (4.664,0.045) (4.748,0.045) (4.832,0.044) (4.916,0.044) (5.000,0.044)};
  \legend{{$x=1$},{$x=2$},{$x=3$}}
\end{axis}
\end{tikzpicture}
\caption{Left: $\Prob[M(t)>x]=2[1-\Phi(x/\sqrt t)]$ (file \file{lec14\_6}); at $t=5$ the values are $0.655$, $0.371$, $0.180$.
Right: the hitting-time density \eqref{ch11:eq-hitcdf} (file \file{lec14\_7}); the peak of $x=1$ is at $t=1/3$ with height
$0.463$, and levels farther away are reached later, with a much more spread-out distribution
\ts{24}{14}{31:35}\ts{24}{14}{33:44}.}\label{ch11:fig-max}
\end{figure}

\begin{intuition}[A path reaches everything, but takes an infinite mean time]
Recurrence and infinite mean go together: Brownian motion in one dimension reaches every level with probability one,
but the waiting time has the heavy tail $t^{-1/2}$. It is the continuous version of the simple random walk, which returns
to $0$ a.s.\ while the expected return time is infinite (\cref{ch04:sec-ruin}, \cref{ch04:thm-ost}: optional stopping fails for $\tau_a$ because
$B(\tau_a)=a\ne0=\E[B(0)]$ and the stopped process is not uniformly bounded).
The scaling $\tau_a\eqd a^2\tau_1$ is diffusion again: time to travel $a$ grows like $a^2$.
\end{intuition}

\subsection{Reflected and absorbed Brownian motion}\label{ch11:sec-absorbed}
Two modifications of $B$ appear in the lecture \ts{24}{14}{33:44}\ts{24}{14}{34:50}.

\begin{definition}[Reflected and absorbed Brownian motion]\label{ch11:def-refabs}
\emph{Reflected Brownian motion} at zero is $R(t)=|B(t)|$ (the path is folded onto the positive half line).
\emph{Absorbed Brownian motion} started at $x>0$ is $A(t)=x+B(t)$ for $t\le\tau$ and $A(t)=0$ for $t>\tau$, where $\tau$ is
the first hitting time of $0$.
\end{definition}

\begin{proposition}[Moments of $R$ and law of $A$]\label{ch11:prop-refabs}
\leavevmode
\begin{enumerate}[(i)]
  \item $\E[R(t)]=\sqrt{2t/\pi}$ and $\Var R(t)=(1-2/\pi)\,t$.
  \item The law of $A(t)$ has an atom at $0$ of mass $2[1-\Phi(x/\sqrt t)]$ (the probability of having been absorbed), and on $y>0$ the
      density
      \begin{equation}\label{ch11:eq-absorbed}
      p_A(y,t\mid x)=p(y,t\mid x)-p(-y,t\mid x)=\frac1{\sqrt{2\pi t}}\Bigl[e^{-(y-x)^2/2t}-e^{-(y+x)^2/2t}\Bigr].
      \end{equation}
\end{enumerate}
\end{proposition}
\begin{proof}
(i) $R(t)\eqd M(t)$ by \cref{ch11:prop-max}, so $\E[R]=\sqrt{2t/\pi}$; $\E[R^2]=\E[B^2]=t$ gives the variance.

(ii) Absorption at $0$ from $x$ is the hitting of the level $-x$ by $B$ from $0$, so $\Prob[\tau\le t]=2[1-\Phi(x/\sqrt t)]$
by \eqref{ch11:eq-hitcdf}. For $y>0$, $\Prob[x+B(t)\in\dd y,\ \tau\le t]$ is, by reflection in $0$, the density of a path ending at $-y$, that is
$p(-y,t\mid x)$; subtracting this from the free density gives \eqref{ch11:eq-absorbed} (the \emph{method of images}, with an image
charge of opposite sign at $-x$).
\end{proof}
Check: $\int_0^\infty p_A\dd y=\Phi(x/\sqrt t)-[1-\Phi(x/\sqrt t)]=1-2[1-\Phi(x/\sqrt t)]$, so density plus atom sum to $1$; and
$\E[A(t)]=x$ for all $t$ (absorbed Brownian motion is the martingale $B(t\wedge\tau)$, the loss of mass to $0$ being
compensated by the surviving paths). \emph{Not to be confused:} the slides use the word ``reflected'' both for the folded process $|B|$ and, in the section
on the reflection principle, for the path reflected \emph{in a level after the hitting time}; only the second is used in \cref{ch11:lem-reflection}.

\begin{finance}[A price that can go bankrupt]
The lecture proposes absorbed Brownian motion as the price of an asset which is positive as long as the firm is viable, but stays at $0$
once bankrupt \ts{24}{14}{34:50}. Structural credit models (Merton, Black--Cox; cf.\ \cref{ch:counterparty}) do exactly this: the
probability of default before $t$ is a hitting probability, and the survival probability with $\sigma$-volatile assets started $x$ above the barrier
is $2\Phi(x/(\sigma\sqrt t))-1$, a straightforward application of \eqref{ch11:eq-hitcdf}.
\end{finance}


% ==================================================================
\section{Brownian motion with drift and the gambler's ruin problem}\label{ch11:sec-drift}

\begin{definition}[Brownian motion with drift]\label{ch11:def-drift}
$X(t)=\mu t+\sigma B(t)$ for a standard Brownian motion $B$, a \emph{drift} $\mu\in\R$ and a \emph{volatility} $\sigma>0$
\ts{24}{14}{43:19}\ts{13}{17}{47:24}. $X$ has continuous paths, independent increments and
$X(t)-X(s)\sim N(\mu(t-s),\sigma^2(t-s))$, so $\E[X(t)]=\mu t$ and $\Var X(t)=\sigma^2t$; it reduces to $B$ for $\mu=0,\sigma=1$.
\end{definition}

\begin{proposition}[Drift wins in the long run]\label{ch11:prop-drift}
For every $\varepsilon>0$, with probability tending to one $(\mu-\varepsilon)t\le X(t)\le(\mu+\varepsilon)t$ as $t\to\infty$; indeed
$X(t)/t\to\mu$ (almost surely). Yet at moderate times the noise $\sigma\sqrt t$ dominates the drift $\mu t$ as long as $t<(\sigma/\mu)^2$.
\end{proposition}
\begin{proof}
$X(t)/t-\mu=\sigma B(t)/t$ has variance $\sigma^2/t\to0$ (convergence in probability; the almost sure version is the law of large numbers for $B$). The
two terms $\mu t$ and $\sigma\sqrt t$ are equal at $t=(\sigma/\mu)^2$.
\end{proof}
The last remark of the proposition is the reason why drifts are so hard to estimate: for a stock with $\mu=10\%$ and $\sigma=30\%$ the time
scale is $(\sigma/\mu)^2=9$ years, and for it the estimate $\hat\mu=X(T)/T$ has standard error $\sigma/\sqrt T$ (\cref{ch:volatility}).

\paragraph{One-step analysis.} Look at the increment $\Delta X=X(t+\Delta t)-X(t)=\mu\Delta t+\sigma\Delta B$ with
$\Delta B\sim N(0,\Delta t)$ independent of the past. Then
\begin{equation}\label{ch11:eq-onestep}
  \E[\Delta] X=\mu\,\Delta t,\qquad \E[(\Delta X)^2]=\sigma^2\Delta t+(\mu\Delta t)^2=\sigma^2\Delta t+o(\Delta t),\qquad
  \E[|\Delta X|^c]=o(\Delta t)\ \ (c>2)
\end{equation}
\ts{24}{14}{44:20}\ts{24}{14}{45:27}\ts{24}{14}{46:33}\ts{24}{14}{47:36}: to order $\Delta t$, only the mean and the variance of the
increment matter. If $f$ is twice differentiable, Taylor's formula and \eqref{ch11:eq-onestep} give
\begin{equation}\label{ch11:eq-generator}
  \E\bigl[f(x+\Delta X)\bigr]=f(x)+\Bigl(\mu f'(x)+\tfrac12\sigma^2f''(x)\Bigr)\Delta t+o(\Delta t),
\end{equation}
so $\mathcal Lf=\mu f'+\tfrac12\sigma^2f''$ is the \emph{generator} of the process. This expansion is why the lecture cares about
the powers of the increment: the value of a derivative on an underlying $X$ is a function $f(X)$, and its dynamics are read
from the Taylor series \ts{24}{14}{47:36}\ts{24}{14}{48:45}. It is the seed of It\^o's lemma (\cref{ch11:sec-gbm}).

\subsection{Gambler's ruin for Brownian motion}\label{ch11:sec-ruin}
Let $X$ be a Brownian motion with drift $\mu$ and volatility $\sigma$, started at $x\in(a,b)$, and let
$\tau=\min\{u:X(u)=a\text{ or }b\}$. Find
\[
  u(x)=\Prob\bigl[X(\tau)=b\mid X(0)=x\bigr],
\]
the probability of reaching the upper level first. This generalises the random walk of \cref{ch04:sec-ruin} to continuous
time and continuous state \ts{24}{14}{48:45}\ts{24}{14}{49:47}\ts{24}{14}{50:51}.

\paragraph{Infinitesimal derivation (the lecture).} Wait a short time $\Delta t$ so small that hitting a level in
the meantime is negligible. By the Markov property, $u(x)=\E[u(x+\Delta X)]$; expanding with \eqref{ch11:eq-generator},
\begin{gather*}
  u(x)=u(x)+\bigl(\mu u'(x)+\tfrac12\sigma^2u''(x)\bigr)\Delta t+o(\Delta t)\\
  \Longrightarrow\ \mu\,u'(x)+\tfrac12\sigma^2u''(x)=0,\qquad u(a)=0,\ u(b)=1
\end{gather*}
\ts{24}{14}{51:55}\ts{24}{14}{52:59}\ts{24}{14}{54:04}\ts{24}{14}{55:16}. This is a first-order equation in $v=u'$: $v'/v=-2\mu/\sigma^2$, so $u'(x)=C\,e^{-2\mu x/\sigma^2}$ and
\begin{equation}\label{ch11:eq-ruin}
  u(x)=\frac{g(x)-g(a)}{g(b)-g(a)},\quad g(x)=e^{-2\mu x/\sigma^2}\ \ (\mu\ne0);\qquad
  u(x)=\frac{x-a}{b-a}\ \ (\mu=0)
\end{equation}
\ts{24}{14}{56:16}\ts{24}{14}{57:18}\ts{24}{14}{58:18}\ts{24}{14}{59:20}. For $\mu=0$, $g(x)=x$ is the other solution of the equation
(the general solution is $Ax+B$); with no drift, the probability of success is linear in the distance travelled, exactly as for the fair gambler
(\cref{ch04:prop-ruin}).

\paragraph{Rigorous derivation by a martingale.} The infinitesimal argument assumes that $u$ is twice differentiable, which
one does not want to take for granted. The proof by optional stopping needs nothing of the kind. By
\cref{ch11:prop-elem}(iii) with $\lambda=\theta\sigma$, the process $e^{\theta X(t)}=e^{\theta\mu t+\theta\sigma B(t)}$ is a martingale iff $\theta\mu=-\theta^2\sigma^2/2$, i.e.
$\theta=-2\mu/\sigma^2$. Hence $M(t)=g(X(t))=e^{-2\mu X(t)/\sigma^2}$ is a martingale, exactly the continuous analogue of
$r^{S_n}$ with $r=q/p$ of \cref{ch04:ex-m4}. The stopped martingale $M(t\wedge\tau)$ is bounded, and $\tau<\infty$ a.s., so
optional stopping (\cref{ch04:thm-ost}(ii)) gives $g(x)=\E[g(X(\tau))]=u\,g(b)+(1-u)\,g(a)$, which is \eqref{ch11:eq-ruin}. It is also the exact limit
of the discrete formula \cref{ch04:prop-biased}: for steps $\pm\sigma\sqrt h$ with $p=\tfrac12(1+\mu\sqrt h/\sigma)$, one has
$r=q/p\approx e^{-2\mu\sqrt h/\sigma}$, and $r^{S_n}\to e^{-2\mu X/\sigma^2}$. A numerical check (not in the course): for
$\mu=0.5$, $\sigma=1$, $a=0$, $b=3$, $x=1$ formula \eqref{ch11:eq-ruin} gives $u=0.66524$, and the exact ruin probability
of the biased walk with $h=10^{-4}$ is $0.66524$.

\begin{figure}[ht]
\centering
\begin{tikzpicture}
\begin{axis}[width=0.9\linewidth,height=5.6cm,xlabel={starting level $x$ ($a=0$, $b=1$, $\sigma=1$)},ylabel={$u(x)=\Prob[\text{hit }b\text{ before }a]$},
  xmin=0,xmax=1,ymin=0,ymax=1.02,legend style={at={(0.5,1.03)},anchor=south,legend columns=5,font=\footnotesize},grid=major,grid style={black!10}]
  \addplot[thick,color=thmcol,domain=0:1,samples=60] {(exp(-2*(-3)*x)-1)/(exp(-2*(-3))-1)};
  \addplot[thick,color=fincol,domain=0:1,samples=60] {(exp(-2*(-1)*x)-1)/(exp(-2*(-1))-1)};
  \addplot[thick,color=black!70,domain=0:1,samples=2] {x};
  \addplot[thick,color=intcol,domain=0:1,samples=60] {(exp(-2*1*x)-1)/(exp(-2*1)-1)};
  \addplot[thick,color=defcol,domain=0:1,samples=60] {(exp(-2*3*x)-1)/(exp(-2*3)-1)};
  \legend{{$\mu=-3$},{$\mu=-1$},{$\mu=0$},{$\mu=1$},{$\mu=3$}}
\end{axis}
\end{tikzpicture}
\caption{Probability of hitting $b=1$ before $a=0$ as a function of the starting point, for several drifts (the four panels of
\file{lec14\_8} show $b=3$ and $b=1$ with $\mu=\pm(10^{-4},0.5,\dots,3)$). With no drift the probability is the straight line
$x$; a positive drift bends the curve above it (the barrier $b$ is reached first from almost everywhere), a negative drift below it
\ts{24}{14}{1:00:20}\ts{24}{14}{1:01:30}\ts{24}{14}{1:02:33}.}\label{ch11:fig-ruin}
\end{figure}

\begin{proposition}[Consequences of \eqref{ch11:eq-ruin}; not discussed in class]\label{ch11:prop-ruincons}
\leavevmode
\begin{enumerate}[(i)]
  \item Symmetry: $u_\mu(x)+u_{-\mu}(a+b-x)=1$.
  \item If $\mu>0$, the probability that $X$ ever falls to a level $d$ below its start is $e^{-2\mu d/\sigma^2}$,
      so the running minimum of $X$ is exponentially distributed with rate $2\mu/\sigma^2$; if $\mu\le0$ every level below is
      reached with probability one.
  \item If $\mu=0$, $\E[\tau]=(x-a)(b-x)/\sigma^2$.
\end{enumerate}
\end{proposition}
\begin{proof}
(i) reflect the interval, $x\mapsto a+b-x$, which changes the sign of the drift and swaps the two barriers.

(ii) Let $b\to\infty$ in
$1-u=(g(b)-g(x))/(g(b)-g(a))$: for $\mu>0$, $g(b)\to0$ and the limit is $g(x)/g(a)=e^{-2\mu(x-a)/\sigma^2}$; for $\mu<0$, $g(b)\to\infty$ and the limit is $1$.

(iii) $X(t)^2-\sigma^2t$ is a martingale (\cref{ch11:prop-elem}(iii)); optional stopping gives $\sigma^2\E[\tau]=\E[X(\tau)^2]-x^2
=\frac{(b-x)a^2+(x-a)b^2}{b-a}-x^2=(x-a)(b-x)$.
\end{proof}

\begin{finance}[How likely is a losing streak? (numbers not from the course)]
Take the parameters used in 2024 PS5 for a simulation, $\mu=0.30$ and $\sigma=0.40$, for the \emph{log} value of a strategy (the drift
already includes any $-\sigma^2/2$ correction). The probability that its value ever falls by $0.2$ (about $18\%$ in level terms) below its running
starting point, however long one waits, is $e^{-2\cdot0.30\cdot0.2/0.16}=e^{-0.75}=47\%$, even though the strategy has a positive expected return of
$30\%$ per year. Positive drift protects in the long run but not from drawdowns: the minimum of the path is exponential with mean $\sigma^2/(2\mu)=0.27$.
\end{finance}

% ==================================================================
\section{Geometric Brownian motion and the first look at It\^o's lemma}\label{ch11:sec-gbm}

The price of a stock cannot be $S_0+\sigma B_t$: it would become negative. What is reasonable, the lecture argues, is that the
\emph{percentage} change is normal, $\Delta S/S=\mu\Delta t+\sigma\Delta B$, written $\dd S_t=S_t(\mu\dd t+\sigma\dd B_t)$
\ts{13}{17}{1:00:19}\ts{13}{17}{1:01:24}. The natural question, asked in the 2013 lecture, is whether the solution of $\dd S=S\dd B$ is $S_t=e^{B_t}$, as in classical calculus; the answer is no
\ts{13}{17}{1:01:24}\ts{13}{17}{1:02:28}. Here is why.

\paragraph{The Taylor argument.} Let $f$ be a smooth function (in the 2013 lecture a call payoff, $f(B_T)=(B_T-s_0)^+$, or the price of any derivative on the underlying)
and $B$ a Brownian motion \ts{13}{17}{1:03:29}\ts{13}{17}{1:04:29}\ts{13}{17}{1:05:33}. If $B$ were differentiable, $\dd f(B_t)=f'(B_t)\frac{\dd B}{\dd t}\dd t$; it is not
\ts{13}{17}{1:06:36}, so one keeps the increment $\dd B$ and writes Taylor's formula
$f(B_{t+x})-f(B_t)=f'(B_t)\,\Delta B+\tfrac12f''(B_t)\,\Delta B^2+\dots$ with $\Delta B=B_{t+x}-B_t$
\ts{13}{17}{1:07:39}\ts{13}{17}{1:08:44}\ts{13}{17}{1:09:46}\ts{13}{17}{1:10:48}. In calculus the term $\Delta B^2$ is negligible; here, by
\eqref{ch11:eq-dB2}, $\Delta B^2\approx\Delta t$ is of the same order as the time step, so it cannot be dropped:
\begin{equation}\label{ch11:eq-ito}
  \dd f(B_t)=f'(B_t)\dd B_t+\tfrac12f''(B_t)\dd t\qquad\text{(It\^o's lemma, heuristic form)}
\end{equation}
\ts{13}{17}{1:11:55}\ts{13}{17}{1:13:01}\ts{13}{17}{1:15:12}. The lecturer's summary: if one thing is to be remembered it is \eqref{ch11:eq-ito}, if two, that $(\dd B)^2=\dd t$
\ts{13}{17}{1:14:05}. Proofs and generalisations (functions of $t$ and $X$, stochastic integrals) are in \cref{ch:stochcalc}.

\paragraph{Application to the exponential.} With $f(x)=e^{x}$, \eqref{ch11:eq-ito} gives $\dd e^{B_t}=e^{B_t}\dd B_t+\tfrac12e^{B_t}\dd t$: the
process $e^{B}$ has a \emph{drift} $\tfrac12S\dd t$ that $\dd S=S\dd B$ does not have. Indeed $\E[e^{B_t}]=e^{t/2}$ (moment generating function of the normal), whereas
$\dd S=S\dd B$ means that $S$ is a martingale with constant mean $S_0$; the same computation is the exponential martingale of
\cref{ch11:prop-elem}(iii), and the correction $-\tfrac12\sigma^2t$ in the exponent is the same lognormal correction as in \cref{ch04:rem-logprice}.
The solution of $\dd S=S(\mu\dd t+\sigma\dd B)$ is therefore
\begin{equation}\label{ch11:eq-gbm}
  S_t=S_0\exp\Bigl[\bigl(\mu-\tfrac12\sigma^2\bigr)t+\sigma B_t\Bigr],
\end{equation}
\emph{geometric Brownian motion} (GBM): $\log S_t$ is a Brownian motion with drift $\mu-\tfrac12\sigma^2$ and volatility $\sigma$, so $S_t$ is lognormal
(\cref{ch03:sec-normal}) with
\[
  \E[S_t]=S_0e^{\mu t},\qquad \operatorname{median}(S_t)=S_0e^{(\mu-\sigma^2/2)t},\qquad \Var S_t=S_0^2e^{2\mu t}\bigl(e^{\sigma^2t}-1\bigr).
\]

\paragraph{The mechanism in discrete time (our check).} The discrete product of one-step relative changes,
$\prod_i(1+\sigma\Delta_iB)$ over a grid of mesh $h$, has mean exactly $\prod(1+0)=1$ and equals
$\exp\bigl(\sum_i\ln(1+\sigma\Delta_iB)\bigr)=\exp\bigl(\sigma B_t-\tfrac12\sigma^2\sum_i(\Delta_iB)^2+\dots\bigr)$; the sum of
squares converges to $t$ by \cref{ch11:thm-qv}, so the product converges to $\exp(\sigma B_t-\tfrac12\sigma^2t)$: the extra
factor $e^{-\sigma^2t/2}$ is the quadratic variation at work. In a simulation with $2\times10^5$ paths of $400$ $\pm\sqrt{h}$ steps ($t=1$, $\sigma=1$) the
mean of the product is $1.005$, whereas the mean of $e^{B_1}$ is $1.656$ (theory: $1$ and $e^{1/2}=1.649$).

\begin{finance}[Arithmetic versus geometric growth]
With $\mu=8\%$ and $\sigma=20\%$ the expected value of an investment grows as $e^{0.08\cdot10}=2.23$ over ten years, but the median as
$e^{(0.08-0.02)10}=1.82$, and the probability of ending below the initial value is $\Phi(-0.06\cdot\sqrt{10}/0.2)=17\%$. The gap between the mean and the typical
path grows with $\sigma^2t$, exactly as for the multiplicative martingales of \cref{ch04:sec-mart} (volatility drag). GBM is the model of Black, Scholes and Merton
(\cref{ch:bs}); Brownian motion with drift, on the other hand, is the model of the \emph{logarithm} of the price, which is why the
gambler's-ruin formulas of \cref{ch11:sec-ruin} apply to log-prices and log-levels.
\end{finance}

% ==================================================================
\section{Processes built from Brownian motion}\label{ch11:sec-variants}

\subsection{The Brownian bridge and the empirical distribution function}\label{ch11:sec-bridge}
\begin{definition}[Brownian bridge]\label{ch11:def-bridge}
$X(t)=B(t)-tB(1)$, $0\le t\le1$ \ts{24}{14}{35:57}\ts{24}{14}{37:05}.
\end{definition}
The path is pinned at both ends: $X(0)=0=X(1)$, since it is $B$ minus the straight line from $0$ to $B(1)$ that it would have followed on average. It is a centred Gaussian
process with
\begin{equation}\label{ch11:eq-bridge}
  \Var X(t)=t(1-t),\qquad \Cov(X(s),X(t))=\min(s,t)-st\quad(s,t\in[0,1]),
\end{equation}
from $\Cov(X(s),X(t))=s-st-st+st$ for $s\le t$. Since $\Cov(X(t),B(1))=t-t=0$ and the pair is jointly Gaussian, $X$ is independent of $B(1)$;
consequently $X$ has the law of Brownian motion \emph{conditioned} to return to $0$ at time $1$ (the endpoint is
irrelevant for the shape of the fluctuations). The variance is largest in the middle, $\tfrac14$ at $t=\tfrac12$.

\paragraph{Where it comes from: the empirical distribution function.} Let $U_1,\dots,U_n$ be i.i.d.\ uniform on $(0,1)$ and
$\hat F_n(t)=\#\{U_i\le t\}/n$. Each indicator $\ind\{U_i\le t\}$ is Bernoulli with mean $t$ and variance $t(1-t)$, so $\E[\hat F_n(t)]=t$ and
$\Var\hat F_n(t)=t(1-t)/n$; by the central limit theorem
$\sqrt n[\hat F_n(t)-t]\to N(0,t(1-t))$ \ts{24}{14}{38:05}\ts{24}{14}{39:08}\ts{24}{14}{40:09}. The whole process $\sqrt n[\hat F_n(t)-t]$ converges to
the Brownian bridge (Donsker, 1952): variance $t(1-t)$ and covariance $\min(s,t)-st$ are those of \eqref{ch11:eq-bridge}, checked directly
on the multinomial counts.

\begin{finance}[Testing a return distribution: probability integral transform and Kolmogorov--Smirnov]
If $Y_1,\dots,Y_n$ are i.i.d.\ with a \emph{continuous} cdf $F$, then $U_i=F(Y_i)$ are i.i.d.\ uniform (\cref{ch03:prop-pit}), so
the statistic $D_n=\sup_y|\hat F_n(y)-F(y)|=\sup_t|\hat F_n^U(t)-t|$ has a distribution that does not depend on $F$: $\sqrt nD_n\to\sup_{t}|X(t)|$, whose law is
$\Prob[\sup|X|\le x]=1-2\sum_{k\ge1}(-1)^{k-1}e^{-2k^2x^2}$ (Kolmogorov); the $5\%$ critical value is $1.358$.
To test whether a set of returns is compatible with an assumed model, transform the data by the model's cdf and test uniformity
\ts{24}{14}{41:10}\ts{24}{14}{42:14}\ts{24}{14}{43:19}. (If the parameters of $F$ are estimated from the same data, the critical values are smaller, as in the Lilliefors
correction.) This is a model-selection tool; residual diagnostics of \cref{ch:timeseries} play the same role for time series.
\end{finance}

\subsection{Sampling at random times: Laplace increments}\label{ch11:sec-laplace}
Observe $B$ at $t=0,1,2,\dots$: the increments are i.i.d.\ $N(0,1)$ (with drift and volatility, $N(\mu,\sigma^2)$)
\ts{24}{14}{1:13:28}. Suppose instead the observation times are random, with i.i.d.\ exponential gaps $T\sim\mathrm{Exp}(\lambda)$ (rate $\lambda$)
independent of $B$ \ts{24}{14}{1:14:35}. Then the increment $B(T)$ has density
\begin{equation}\label{ch11:eq-laplace}
  f(x)=\int_0^\infty\lambda e^{-\lambda t}\frac{e^{-x^2/2t}}{\sqrt{2\pi t}}\dd t=\sqrt{\tfrac\lambda2}\;e^{-\sqrt{2\lambda}\,|x|},
\end{equation}
the \emph{Laplace} (two-sided exponential) density, instead of a bell curve: a \emph{normal variance mixture} in
the sense of \cref{ch03:prop-mixture}, $B(T)=\sqrt T\,Z$, with variance $\E[T]=1/\lambda$ and excess kurtosis $3\E[T^2]/(\E[T])^2-3=3$ (checked by numerical integration of \eqref{ch11:eq-laplace} at two points). Its tails are heavier than the normal ones, and the
lecture reports that a Laplace model fits S\&P~500 returns much better \ts{24}{14}{1:15:44}. The economic reading: the useful clock for prices
is not calendar time but a random \emph{business time} (trades, information arrivals), and the process is a Brownian motion \emph{subordinated} to a
random clock. Time-changed and stochastic-volatility models take this idea further (\cref{ch:volatility}).

\subsection{Two-dimensional Brownian motion}\label{ch11:sec-2d}
The lecture opens with the animation of a particle moving in the plane, where each coordinate takes independent Gaussian steps,
no direction being preferred \ts{24}{14}{1:04}\ts{24}{14}{2:04}\ts{24}{14}{3:08}. Formally $\bm B(t)=(B_1(t),B_2(t))$ with independent standard Brownian motions has density
$(2\pi t)^{-1}e^{-|\bm y|^2/2t}$, invariant under rotations; the distance $|\bm B(t)|$ has the Rayleigh law
$\Prob[|\bm B(t)|<\varepsilon]=1-e^{-\varepsilon^2/(2t)}$ (for $t=\varepsilon=1$: $39.3\%$) and $\E[|\bm B(t)|]=\sqrt{\pi t/2}$.
The path in the plane is fractal-looking and its two projections are the one-dimensional processes of the rest of the chapter (file \file{lec14\_2}, $2000$ steps).
Correlated Brownian motions, obtained from independent ones by a matrix (\cref{ch:linalg}), are the basis of models with several risk factors.

\subsection{Filtrations and adapted processes}\label{ch11:sec-adapted}
Let $X$ be a stochastic process, for instance a price. The \emph{filtration} $\{\mathcal F_t\}$ generated by $X$ is the increasing family of $\sigma$-fields of
events determined by $\{X(u),\,0\le u\le t\}$: $\mathcal F_t\subset\mathcal F_{t'}$ for $t<t'$; it is the information available at time $t$
\ts{24}{14}{1:16:44}\ts{24}{14}{1:17:48}. A process $Y$ is \emph{adapted} to $X$ if $Y(t)$ depends only on $\{X(u),u\le t\}$, equivalently
$\mathcal G_t\subset\mathcal F_t$ for the filtration $\mathcal G$ of $Y$: the transformation is not forward-looking. For example if $Y(t)$
is the cash flow of a strategy that trades the security $X$, adaptedness says that one cannot trade on tomorrow's price
\ts{24}{14}{1:17:48}\ts{24}{14}{1:18:55}. Adaptedness is what makes the gains from trading a martingale when $X$ is one: the continuous-time
analogue of the martingale transform of \cref{ch04:thm-transform} (with constant-mean processes, martingales, ``applied again and again'' in the
lecture \ts{24}{14}{1:18:55}\ts{24}{14}{1:19:59}) is the stochastic integral $\int Y\dd X$ of \cref{ch:stochcalc}, and the price of an option is a conditional
expectation with respect to $\mathcal F_t$ (\cref{ch:bs}). All the martingale statements of \cref{ch11:prop-elem} hold with respect to the filtration of $B$.

\subsection{L\'evy processes: Brownian motion and the Poisson process}\label{ch11:sec-levy}
2024 PS5 opens with a general definition (Problems~1--2, \cref{ch11:ex-levy-bm,ch11:ex-levy-poisson}); the topic was not lectured on.

\begin{definition}[L\'evy process]\label{ch11:def-levy}
$\{L(t),\,t\ge0\}$ is a \emph{L\'evy process} if
\begin{enumerate}[(i)]
  \item $L(0)=0$;
  \item stationary increments: $L(t)-L(s)\eqd L(t-s)$ for $s\le t$;
  \item independent increments over disjoint intervals;
  \item continuity in probability: $\lim_{s\to t}\Prob[|L(t)-L(s)|>\varepsilon]=0$ for all $\varepsilon>0$ and $t$.
\end{enumerate}
\end{definition}
(The problem set indexes $t$ over all of $\R$; only $t\ge0$ is used.) It generalises Definition~\ref{ch11:def-bm}, which asks in
addition that the increments be normal and the paths continuous: continuity in probability is much weaker (it does not forbid jumps
at unpredictable times). Two examples are the subject of the problem set. \emph{Brownian motion with drift} (\cref{ch11:def-drift})
is a L\'evy process: its increments are $N(\mu h,\sigma^2h)$. The \emph{Poisson process} $N$ with intensity $\lambda$ has
independent increments with $N(t)-N(s)\sim\mathrm{Poisson}(\lambda(t-s))$; it is continuous in probability, since
$\Prob[N(t)-N(s)\ge1]=1-e^{-\lambda(t-s)}\to0$, although \emph{every} path is a pure-jump staircase.
The two are the extreme cases of the L\'evy--It\^o decomposition (not covered): every L\'evy process is the sum of a drift, a Brownian motion,
and a superposition of jumps (compensated when small).

\begin{table}[ht]
\centering\small
\begin{tabularx}{\linewidth}{@{}l L L L@{}}
\toprule
& \textbf{paths} & \textbf{increment over $h$} & \textbf{quadratic variation on $[0,T]$}\\
\midrule
$X=\mu t+\sigma B$ & continuous, nowhere differentiable & $N(\mu h,\sigma^2h)$ & $\sigma^2T$, deterministic\\
Poisson $N$, rate $\lambda$ & piecewise constant, unit jumps & $\mathrm{Poisson}(\lambda h)$ & $N(T)$ (each jump adds $1^2$), random\\
\bottomrule
\end{tabularx}
\caption{Brownian motion with drift and the Poisson process, the two building blocks of jump--diffusion models.}\label{ch11:tab-levy}
\end{table}

\begin{intuition}[Two universal limits]
Sums of many small independent increments give Brownian motion (\cref{ch11:thm-donsker}). Sums of many \emph{rare} unit
increments give the Poisson process: a walk that steps up by $1$ with probability $\lambda/n$ at each of $n$ sub-intervals has
$\mathrm{Bin}(n,\lambda/n)\to\mathrm{Poisson}(\lambda)$ increments (the law of small numbers). The compensated process
$N(t)-\lambda t$ is a martingale with variance $\lambda t$, and $(N(t)-\lambda t)/\sqrt\lambda$ looks Brownian when the rate is large, a bridge between the two limits. In finance the two blocks together give the
\emph{jump--diffusion} of Merton (1976): Brownian noise for the day-to-day movements, Poisson jumps for crashes, defaults and announcements.
Realised variance (\cref{ch11:sec-paths}) measures both: the quadratic variation of the diffusive part plus the sum of the squared jumps.
\end{intuition}

% ==================================================================
\section{Looking ahead and the two lectures compared}\label{ch11:sec-ahead}
\begin{itemize}
  \item $(\dd B)^2=\dd t$ and \eqref{ch11:eq-ito} are the beginning of \emph{stochastic calculus}: It\^o integrals $\int\theta\dd B$ for adapted $\theta$, It\^o's lemma for
        $f(t,X_t)$, geometric Brownian motion as the solution of an SDE (\cref{ch:stochcalc,ch:sde}), Girsanov's theorem and risk-neutral pricing (\cref{ch:bs}).
  \item The martingale $e^{\lambda B-\lambda^2t/2}$ and the gambler's-ruin martingale $e^{-2\mu X/\sigma^2}$ are the continuous-time versions
        of \cref{ch04:ex-m4}; the reflection principle and hitting-time density are the working tools for barrier options and
        structural credit models.
  \item The diffusion equation \eqref{ch11:eq-heat} is the seed of the Black--Scholes PDE, the Fokker--Planck equation and the Feynman--Kac formula of the next chapters.
  \item The Dickey--Fuller distribution of \cref{ch09:sec-df} is the law of $\int B\dd B/\int B^2\dd t$.
\end{itemize}

\begin{coursenote}[2013 and 2024 compared; problem-set overlap]
Both years cover the same core: definition, random-walk limit, non-differentiability, reflection principle and the maximum, quadratic variation, drift. They differ
in emphasis. \emph{2013} (C.~Lee) is rigorous about the definition (existence theorem, path space), proves non-differentiability by the maximum, and spends the second half of the lecture on
\emph{quadratic variation as the reason for It\^o's lemma}, ending with the question $S_t=e^{B_t}$?; he does not treat the bridge,
Laplace increments, reflected and absorbed Brownian motion, the hitting-time density, or gambler's ruin with drift. \emph{2024}
(P.~Kempthorne) is more applied and computational: transition density and diffusion equation, the normalised random walk simulated
for Bernoulli and Gaussian steps, path pairing for the reflection principle, hitting-time density,
the bridge and the empirical distribution function, gambler's ruin by infinitesimal analysis, quadratic variation of log-prices, subordination and adapted processes; it
stops before It\^o's lemma, which is in the next lecture (\cref{ch:stochcalc}). The quadratic-variation theorem is stated for arbitrary partitions in 2024 and proved for
equal spacing in 2013; \cref{ch11:thm-qv} above proves it in general. Problem sets: 2024 PS4 Problem~1 is 2013 PS8 Problem~A-2 (same three parts, with $\sigma=2$ in part~(b)
in 2024); 2013 PS8 Problem~B-1 (density and rotation invariance of a two-dimensional Brownian motion, and $\Prob[\bm B(t)\in D]$ for a disc $D$) belongs to \cref{ch:stochcalc}, together with the rest of
PS8; \cref{ch11:sec-2d} gives the results it needs. 2024 PS4 Problems~3--4 (MLE for the geometric Brownian motion) are treated in \cref{ch:volatility}.
\end{coursenote}

% ==================================================================
\begin{keyformulas}
\begin{align*}
&\text{Brownian motion:} && B(0)=0,\ B(t)-B(s)\sim N(0,t-s)\text{ independent}\\
& && \Cov(B(s),B(t))=\min(s,t)\\
&\text{Martingales:} && B(t),\quad B(t)^2-t,\quad e^{\lambda B(t)-\lambda^2t/2}\\
&\text{Density:} && p(y,\tau\mid x)=\tfrac{1}{\sqrt{2\pi\sigma^2\tau}}e^{-(y-x)^2/(2\sigma^2\tau)}\\
& && \partial_\tau p=\tfrac{\sigma^2}{2}\partial_y^2p\\
&\text{Donsker:} && B_n(t)=S_{\lfloor nt\rfloor}/\sqrt n\ \Rightarrow\ B(t)\quad(\E[X]=0,\ \Var X=1)\\
&\text{Scaling:} && \{B(ct)\}\eqd\{\sqrt c\,B(t)\}\\
&\text{Paths:} && \text{continuous, nowhere differentiable; }\ \sum(\Delta B)^2\to T,\ \ (\dd B)^2=\dd t\\
&\text{Maximum:} && \Prob[M(t)\ge a]=2\Prob[B(t)>a]=2[1-\Phi(a/\sqrt t)]\\
&\text{Hitting time:} && f_{\tau_a}(t)=\tfrac{a}{\sqrt{2\pi}}t^{-3/2}e^{-a^2/(2t)},\quad \E[\tau_a]=\infty\\
&\text{Reflected:} && \E[|B(t)|]=\sqrt{2t/\pi},\quad \Var|B(t)|=(1-\tfrac2\pi)t\\
&\text{Absorbed at 0:} && p_A(y)=p(y,t\mid x)-p(-y,t\mid x)\\
& && \Prob[A(t)=0]=2[1-\Phi(x/\sqrt t)]\\
&\text{Bridge:} && X(t)=B(t)-tB(1),\quad \Cov(X(s),X(t))=\min(s,t)-st\\
&\text{Drift:} && X=\mu t+\sigma B,\quad \mathcal Lf=\mu f'+\tfrac12\sigma^2f''\\
&\text{Ruin:} && \Prob[X(\tau)=b\mid X(0)=x]=\tfrac{e^{-2\mu x/\sigma^2}-e^{-2\mu a/\sigma^2}}{e^{-2\mu b/\sigma^2}-e^{-2\mu a/\sigma^2}}\\
& && \bigl(=\tfrac{x-a}{b-a}\text{ if }\mu=0\bigr)\\
&\text{It\^o (heuristic):} && \dd f(B)=f'(B)\dd B+\tfrac12f''(B)\dd t\\
&\text{GBM:} && S_t=S_0e^{(\mu-\sigma^2/2)t+\sigma B_t},\quad \E[S_t]=S_0e^{\mu t}\\
&\text{Laplace:} && B(T),\ T\sim\mathrm{Exp}(\lambda):\ f(x)=\sqrt{\lambda/2}\,e^{-\sqrt{2\lambda}|x|}
\end{align*}
\end{keyformulas}

% ==================================================================
\section*{Exercises}
\addcontentsline{toc}{section}{Exercises}

\subsection*{From the problem sets}
\begin{exercise}[PS4 (2024), Problem 1; PS8 (2013), Problem A-2: conditional moments]\label{ch11:ex-cond}
\leavevmode
\begin{enumerate}[(a)]
  \item Let $B(t)$ be a standard Brownian motion and $t>s\ge0$ fixed. Compute $\E[B(t)\mid B(s)]$ and $\Var[B(t)\mid B(s)]$.
  \item Let $X(t)$ be a Brownian motion with drift $\mu$ and volatility $\sigma=2$ (the 2013 version has no volatility). Compute $\E[X(t)\mid X(s)]$ and $\Var[X(t)\mid X(s)]$.
  \item For a standard Brownian motion and fixed reals $t>s>0$, compute $\E\bigl[\exp\bigl(\sigma(B(t)-B(s))\bigr)\bigr]$, for a real constant $\sigma$ (the multiplier of the increment inside the exponential).
\end{enumerate}
\end{exercise}

\begin{exercise}[PS4 (2024), Problem 2: Brownian scaling]\label{ch11:ex-scaling}
Let $B$ be a standard Brownian motion. Prove:
\begin{enumerate}[(a)]
  \item if $B(0)=0$, then for every $t>0$, $\{B(st),s\ge0\}\eqd\{t^{1/2}B(s),s\ge0\}$ as processes;
  \item more precisely, the two families have the same finite-dimensional distributions: for $s_1<\dots<s_n$,
      $(B(s_1t),\dots,B(s_nt))\eqd(t^{1/2}B(s_1),\dots,t^{1/2}B(s_n))$.
\end{enumerate}
\end{exercise}

\begin{exercise}[PS5 (2024), Problem 1: Brownian motion with drift as a L\'evy process]\label{ch11:ex-levy-bm}
Let $W$ be a Brownian motion with drift $\mu\in\R$ and volatility $\sigma>0$.
\begin{enumerate}[(a)]
  \item Prove that $W$ is a L\'evy process in the sense of \cref{ch11:def-levy}.
  \item Simulate the path of $W$ with $T=1$ year, $\mu=0.30$, $\sigma=0.40$ and $m=252$ increments, $t_i=iT/m$, and plot five simulated paths (the process values at the time increments).
\end{enumerate}
\end{exercise}

\begin{exercise}[PS5 (2024), Problem 2: the Poisson process]\label{ch11:ex-levy-poisson}
A stochastic process $\{N(t),t\ge0\}$ is a Poisson process of intensity (jump rate) $\lambda$ if for any $0=t_0<t_1<\dots<t_n$ the increments $N(t_k)-N(t_{k-1})$ are independent,
and $N(t)-N(s)$ is Poisson with mean $\lambda(t-s)$ for all $0\le s<t$.
\begin{enumerate}[(a)]
  \item Prove that $N$ is a L\'evy process.
  \item Simulate the path of a Poisson process with $T=1$ year, $\lambda=24$ events per year and $m=252$ increments, $t_i=iT/m$, and plot five simulated paths.
\end{enumerate}
\end{exercise}

\subsection*{Additional exercises (not from the course)}
\begin{exercise}[not from the course: time inversion]\label{ch11:ex-inversion}
Prove that $W(t)=tB(1/t)$ ($t>0$), $W(0)=0$, is a standard Brownian motion (check that it is centred Gaussian with covariance $\min(s,t)$; for continuity at $0$ use
$B(u)/u\to0$ as $u\to\infty$). Deduce that $B(t)/t\to0$ as $t\to\infty$ from the continuity of $W$ at $0$.
\end{exercise}

\begin{exercise}[not from the course: reflection and the minimum]\label{ch11:ex-refl}
Let $m(t)=\min_{u\le t}B(u)$. Show $\Prob[m(t)\le-a]=2\Prob[B(t)<-a]$ and that $M(t)-B(t)\eqd M(t)\eqd|B(t)|$ at a fixed time $t$ (hint: use \cref{ch11:rem-joint} or the
time-reversed path $B(t-u)-B(t)$). What is the probability that a driftless log price with $\sigma=25\%$ never falls below $-10\%$ during one year?
\end{exercise}

\begin{exercise}[not from the course: a hitting time with infinite mean]\label{ch11:ex-tau}
Show $\Prob[\tau_a<\infty]=1$ and $\E[\tau_a]=\infty$ from \eqref{ch11:eq-hitcdf}. Why does optional stopping ($\E[B(\tau_a)]=0$)
fail? Use $\tau_a\eqd a^2\tau_1$ to show that the median of $\tau_a$ scales like $a^2$; compute the median of $\tau_1$.
\end{exercise}

\begin{exercise}[not from the course: absorbed Brownian motion]\label{ch11:ex-abs}
For absorbed Brownian motion $A$ started at $x>0$ (\cref{ch11:def-refabs}) show that \eqref{ch11:eq-absorbed} and the mass $2[1-\Phi(x/\sqrt t)]$ at $0$ sum to $1$, and that $\E[A(t)]=x$.
Is $A(t)^2-t\wedge\tau$ a martingale?
\end{exercise}

\begin{exercise}[not from the course: exit time and ruin]\label{ch11:ex-exit}
For $X=\mu t+\sigma B$ with $\mu\ne0$ compute $\E[\tau]$ for the exit time of $(a,b)$ (hint: the martingale $X(t)-\mu t$ and \eqref{ch11:eq-ruin}). Check that the limit $\mu\to0$ is $(x-a)(b-x)/\sigma^2$.
\end{exercise}

\begin{exercise}[not from the course: geometric Brownian motion]\label{ch11:ex-gbm}
For $S_t$ in \eqref{ch11:eq-gbm} compute $\E[S_t]$, $\Var S_t$, and $\Prob[S_t<S_0]$. With $\mu=0.08$, $\sigma=0.2$ find the horizon at which the probability of a loss falls below $5\%$.
Show that $\E[e^{\lambda B(t)}]=e^{\lambda^2t/2}$ and use it to explain why $e^{B}$ is not a martingale.
\end{exercise}

\begin{exercise}[not from the course: the Laplace mixture]\label{ch11:ex-laplace}
Derive \eqref{ch11:eq-laplace} (hint: $\int_0^\infty t^{-1/2}e^{-\lambda t-x^2/2t}\dd t=\sqrt{\pi/\lambda}\,e^{-\sqrt{2\lambda}|x|}$), verify that the density integrates to one, and compute its variance and excess kurtosis.
\end{exercise}

% !TEX root = ../main.tex
\chapter{Data Science and AI for Biomedical Portfolios}\label{ch:biomed}
\begin{paracol}{2}

\begin{sources}
\begin{tabularx}{\linewidth}{@{}l l L@{}}
\toprule
\textbf{Lecture} & \textbf{Speaker} & \textbf{Content}\\
\midrule
2024 L18 & A.~W.~Lo (MIT Sloan; guest) & Biomedicine at an inflection point (gene therapy); the ``valley of death''; risk versus uncertainty; Sharpe ratio and the Madoff example; one anti-cancer drug as an investment; a portfolio of 150 independent projects; rare diseases, correlation and BridgeBio; a fusion ``mega fund''; questions (patents, expertise, COVID-19 vaccines).\\
\bottomrule
\end{tabularx}

\smallskip
\textbf{Files.} Transcript only: \file{Lecture 18 - Applying Data Science and Artificial Intelligence to Managing Biomedical Portfolios.txt}. The same text exists as
\file{18642-lecture-18-version-5_transcript.pdf} in \file{18.642-fall-2024/static_resources}. There are no slides, no
problem set and no other file for this lecture; everything said about the slides below is therefore unknown, and the
chapter is built on the spoken text.
\textbf{Problem sets.} None. The exercises at the end are all \emph{not from the course}.
\textbf{Not from the lecture.} Every derivation marked ``Reconstruction'' (correlated projects, binomial and
beta-binomial counts, securitisation, ML for trial success) is my own mathematics behind the numbers quoted by
Lo; the numerics were checked by a short Python script. The 2024 L1 slides list a past final-project paper,
\emph{A Megafund Approach to Pediatric Oncology Drug Development}, which uses exactly this framework.
\end{sources}

\section*{Motivation}
Prof.\ Lo is a financial economist (director of the MIT Laboratory for Financial Engineering, principal investigator
at CSAIL) \ts{24}{18}{0:00}\ts{24}{18}{2:07}. His message is not about biology: it is that a \emph{financing structure} can change the
risk--reward profile of an investment, and that the required mathematics is a few lines of probability. A single
drug-development project is a lottery ticket that nobody wants: costly, slow, and very likely to fail. A
portfolio of many such tickets whose outcomes are only weakly correlated has a much better ratio of expected
return to risk, and that is what makes it fundable. The chapter follows the lecture and then adds the
mathematics that the lecture only quotes:
\begin{itemize}
  \item the Sharpe ratio as the quantity to be raised, and the two ways to raise it (\cref{ch12:sec-sharpe});
  \item the single-project calculation with its numbers, and the $\sqrt N$ law for independent projects
        (\cref{ch12:sec-single,ch12:sec-portfolio});
  \item \emph{what breaks when projects are correlated}: the variance of a sum, the diversification limit,
        counts of successes, and parameter uncertainty as a hidden common factor
        (\cref{ch12:sec-corr,ch12:sec-counts}); this is the point the lecturer stresses with rare diseases and
        with the fusion question;
  \item a sketch of securitisation (debt and equity tranches) and of predicting trial success from data
        (\cref{ch12:sec-tranche,ch12:sec-ml});
  \item the second example of the lecture, fusion energy, and a numerical erratum (\cref{ch12:sec-fusion}).
\end{itemize}
Portfolio variance is developed in general in \cref{ch:portfolio}; the covariance formula for sums is in
\cref{ch:probability}.

% ==================================================================
\section{The setting: a valley of death}\label{ch12:sec-setting}

The lecture opens with two clinical stories, both single-gene diseases treated by a one-time gene therapy
(a correct copy of the gene is carried by a virus into the cells): Leber's congenital amaurosis, a blindness
caused by one DNA ``typo'' \ts{24}{18}{4:07}, and Canavan disease, in which the myelin insulation of nerve
cells degrades and children rarely live beyond about ten years \ts{24}{18}{6:13}\ts{24}{18}{9:19}. Lo uses them, together with
the ``convergence'' of life sciences, physical sciences and engineering and the various ``omics''
\ts{24}{18}{13:36}, to argue that science is advancing fast. The bottleneck he identifies is \emph{economics}
\ts{24}{18}{14:39}.

The specific obstacle is the \emph{valley of death}: the gap between academic experiments and the first
human clinical trials, where money is hard to raise \ts{24}{18}{15:42}. He explains it with the distinction, familiar in economics, between
\begin{itemize}
  \item \emph{risk}: randomness whose probability law can be parametrised (normal, gamma, $\chi^2$), and
  \item \emph{uncertainty}: ``unknown unknowns'', with no data from which to assess probabilities
        \ts{24}{18}{16:47}.
\end{itemize}
Drug development has three features that investors must digest: it costs hundreds of millions to billions
of dollars per approved drug, the probability of success is low, and the process takes 10 to 15 years
\ts{24}{18}{17:50}. Compared with writing an app, all three are extreme.

\begin{intuition}[Risk versus uncertainty, in physicist's terms]
Risk is a stochastic model with known parameters (a Hamiltonian with a heat bath at known temperature);
uncertainty is not knowing the Hamiltonian. \cref{ch12:sec-counts} shows that the two are not unrelated:
uncertainty about a parameter that all projects share acts exactly like a common random field coupling them,
and it limits diversification just as correlation does.
\end{intuition}

% ==================================================================
\section{Sharpe ratio: what investors want}\label{ch12:sec-sharpe}

Lo's classroom exercise shows four cumulative-return curves, unnamed, from 2008 onward, and asks which one
to choose \ts{24}{18}{18:54}\ts{24}{18}{19:59}. The four assets are revealed as Treasury bills
(safe, about zero real return), the S\&P~500, Pfizer, and the Fairfield Sentry fund, the feeder of Madoff's
Ponzi scheme, which was chosen by most because it looked like the best trade-off \ts{24}{18}{20:04}\ts{24}{18}{21:04}. The moral
is that investors want a high expected return per unit of risk, formalised by the Sharpe ratio.

\begin{definition}[Sharpe ratio]\label{ch12:def-sharpe}
For an investment with return $R$, risk-free return $r_f$ over the same horizon, the \emph{Sharpe ratio} is
\[
  \mathrm{SR}=\frac{\E[R]-r_f}{\sqrt{\Var R}}.
\]
\end{definition}
The lecture quotes $\mathrm{SR}\approx0.33$ for the S\&P~500 and $0.34$ for Pfizer, and an ``order of magnitude
higher'' figure, on paper, for the Ponzi scheme \ts{24}{18}{22:07}. Lo adds that the Sharpe ratio of
biomedical investing has been declining for decades, not because the numerator is unattractive but
because risk and uncertainty grow faster as the science becomes more sophisticated \ts{24}{18}{23:07}.

There are, jokes Lo, two ways to increase a ratio: raise the numerator or lower the denominator
(or both) \ts{24}{18}{24:11}. Better science and engineering act on both. What is under-used in biomedicine is
\emph{new financing and business models}, and this is where mathematicians can contribute: constructing
structures whose risk--reward trade-off differs from that of the underlying assets
\ts{24}{18}{25:14}.

\begin{finance}[Why the Madoff curve fooled people]
A smooth, upward curve has a tiny sample standard deviation, so its Sharpe ratio is enormous; a Sharpe ratio
that looks too good is itself a warning sign. In the lecture the point is behavioural (``we all want
high-yielding, low-risk assets'') and it sets up the target: not a lottery ticket, but a pooled structure with a
respectable Sharpe ratio \ts{24}{18}{22:07}.
\end{finance}

% ==================================================================
\section{One anti-cancer drug as an investment}\label{ch12:sec-single}

Lo poses the question as a fund. Raise \$200 million, wait 10 years, and with probability $95\%$ return nothing
\ts{24}{18}{25:14}. The back-of-the-envelope numbers for an oncology project are \ts{24}{18}{26:19}\ts{24}{18}{27:22}:
\begin{itemize}
  \item cost of the clinical trials: about \$200 million, up front;
  \item 10 years to complete them;
  \item probability of success: a bit under $5\%$ historically ($3.4\%$ in oncology), rounded to $p=5\%$;
  \item if successful, revenues of \$2 billion per year in years 11--20 (the patent lifetime is 20 years, so the
        exclusivity runs until then), equivalent to a single payment of \$12.3 billion in year~10.
\end{itemize}
The equivalence is a present-value computation (\cref{sec:discounting}). The lecture does not state the discount
rate, but the numbers are reproduced by $r=10\%$:
\begin{equation}\label{ch12:eq-annuity}
  2\times\sum_{k=1}^{10}\frac{1}{1.1^{\,k}}=2\times\frac{1-1.1^{-10}}{0.1}=2\times6.1446=12.29\approx12.3
  \quad\text{(billion, at year 10).}
\end{equation}
(This is a \emph{Reconstruction}: only the \$12.3 billion is quoted.)

\begin{example}[Reconstruction: mean and volatility of the single project]\label{ch12:ex-single}
Let $G$ be the gross ten-year return, $G=12.3/0.2=61.5$ with probability $p=0.05$ and $0$ otherwise. Then
\[
  \E G=pG_0=3.075,\qquad \Var G=G_0^2\,p(1-p),\qquad \sqrt{\Var G}=61.5\sqrt{0.0475}=13.40,
\]
with $G_0=61.5$. The annualised expected return is $3.075^{1/10}-1\approx11.9\%$, the ``about $12\%$'' of the lecture. The
annualised standard deviation, obtained by scaling the ten-year standard deviation by $\sqrt{10}$, is
$13.40/\sqrt{10}=4.24$, that is $424\%$, the ``$423.5\%$'' of the lecture \ts{24}{18}{28:23} (the small
difference comes from rounding of the inputs). The Sharpe ratio with $r_f\approx0$ is $0.12/4.24\approx0.03$,
``close to zero'': hardly anyone would invest.
\end{example}
Lo's conclusion is that few investors take this bet \ts{24}{18}{28:23}. The scaling $\sqrt{10}$ used for the
volatility is the usual rule for independent yearly increments; how the expected return is annualised is a convention
(the simple average $207.5\%/10\approx21\%$ would give another number), so the pair $(12\%,423.5\%)$ is best taken as
the lecturer's stated input to the next step.

% ==================================================================
\section{A portfolio of independent projects}\label{ch12:sec-portfolio}

Instead of one project, invest in $N=150$ at once: this needs $150\times\$200\text{M}=\$30$ billion, a sum that
Lo defers (``assume we have it'') \ts{24}{18}{29:26}. If the projects are independent and identically distributed
the expected return of the portfolio equals that of each project, while
the standard deviation is divided by $\sqrt N$: $423.5\%/\sqrt{150}=34.6\%$ \ts{24}{18}{30:30}. The Sharpe ratio is
then about $12\%/34.6\%\approx0.35$, ten times the single-project value.

\begin{proposition}[Reconstruction: equal-weight portfolio of $N$ i.i.d.\ projects]\label{ch12:prop-iid}
Let $R_1,\dots,R_N$ be i.i.d.\ with mean $\mu$ and variance $\sigma^2$ (returns per dollar invested) and let
$R_P=\frac1N\sum_i R_i$. Then $\E R_P=\mu$, $\Var R_P=\sigma^2/N$ and $\mathrm{SR}_N=\sqrt N\,\mathrm{SR}_1$.
\end{proposition}
\begin{proof}
Linearity of $\E$ gives the mean. For the variance, independence gives
$\Var\bigl(\sum_iR_i\bigr)=\sum_i\Var R_i=N\sigma^2$, and dividing the sum by $N$ divides the variance by $N^2$.
\end{proof}
With the numbers of \cref{ch12:ex-single}: $\mathrm{SR}_1=0.0283$, $\mathrm{SR}_{150}=0.0283\sqrt{150}=0.347$.
Lo phrases the step as ``the variance of the sum is the sum of the variances'' and ``the risk goes down by the square
root of the number of projects''; the second phrase refers to the standard deviation \ts{24}{18}{29:26}\ts{24}{18}{30:30}.

\begin{intuition}[Central-limit bookkeeping]
Ten thousand independent spins have relative fluctuations of order $10^{-2}$ of their mean: this is the same
$1/\sqrt N$ law that makes pressure well defined. The economics of ``a \$30 billion fund'' is the statement that
the fund is a large-$N$ system whose fluctuations are small, provided the couplings between the components are
small; the rest of the chapter is about that proviso.
\end{intuition}

Lo then asks where \$30 billion would come from, and answers that if the fund is attractive enough, small
investors can supply it in aggregate (part of a retirement account or of a summer-internship salary)
\ts{24}{18}{30:30}. This is the second half of the strategy: the same mathematics that makes the fund attractive
also justifies broad retail participation, and \cref{ch12:sec-tranche} sketches how such capital can be
structured.

% ==================================================================
\section{Correlation: the real constraint}\label{ch12:sec-corr}

The i.i.d.\ assumption is flagged by Lo himself (``I'll come back to that assumption later'')
\ts{24}{18}{29:26}. He returns to it with \emph{rare diseases} \ts{24}{18}{31:30}. There are roughly $12{,}000$
of them, each with a small patient population, but more than 30 million Americans are affected in
aggregate. Development is comparatively cheap because trials can be small (he mentions a
Leber's-amaurosis approval on 30--40 patients) and the target is a known gene \ts{24}{18}{32:31}\ts{24}{18}{33:36}. Above all
the diseases are so different that the successes and failures of different programmes are barely
correlated. To see why this matters, recall the variance of a sum.

\begin{proposition}[Variance of a portfolio of $N$ correlated projects]\label{ch12:prop-corr}
Let $R_1,\dots,R_N$ have common variance $\sigma^2$ and common pairwise correlation $\rho$. Then
\begin{equation}\label{ch12:eq-varcorr}
  \Var R_P=\frac{1}{N^2}\Bigl(\sum_i\Var R_i+\sum_{i\neq j}\Cov(R_i,R_j)\Bigr)
  =\sigma^2\Bigl(\frac1N+\frac{N-1}{N}\,\rho\Bigr),
\end{equation}
so that
\begin{equation}\label{ch12:eq-sharpecorr}
  \mathrm{SR}_N=\frac{\mathrm{SR}_1}{\sqrt{\,1/N+(1-1/N)\rho\,}}\ \xrightarrow[N\to\infty]{}\ \frac{\mathrm{SR}_1}{\sqrt\rho}.
\end{equation}
\end{proposition}
\begin{proof}
The double sum $\sum_{i\neq j}$ has $N(N-1)$ ordered pairs, each equal to $\rho\sigma^2$. Hence
$\Var R_P=\frac{1}{N^2}\bigl(N\sigma^2+N(N-1)\rho\sigma^2\bigr)$, which is~\eqref{ch12:eq-varcorr}.
The Sharpe ratio uses the same expected return, whence~\eqref{ch12:eq-sharpecorr}.
\end{proof}
This is the statement of the lecture: with $N$ projects there are $N(N-1)/2$ pairwise covariances, ``a lot of them'',
and in a typical investment portfolio they are almost all positive \ts{24}{18}{33:36}\ts{24}{18}{34:38}. Only when they vanish
does the risk fall as $1/\sqrt N$. The relevant limit of \eqref{ch12:eq-varcorr} is the
\emph{undiversifiable} floor $\sigma^2\rho$: no number of projects removes the common component.

\begin{table}[ht]
\centering\small
\caption{Reconstruction: standard deviation and Sharpe ratio of an equal-weight portfolio for the single-project
inputs of \cref{ch12:ex-single} (mean $12\%$, standard deviation $423.5\%$, $r_f=0$), computed from
\eqref{ch12:eq-varcorr}.}\label{ch12:tab-corr}
\begin{tabular}{@{}lrrrr@{}}
\toprule
& \multicolumn{2}{c}{$N=15$} & \multicolumn{2}{c}{$N=150$}\\
\cmidrule(lr){2-3}\cmidrule(l){4-5}
$\rho$ & std.\ dev. & SR & std.\ dev. & SR\\
\midrule
$0$    & $109\%$ & $0.11$ & $34.6\%$ & $0.35$\\
$0.01$ & $117\%$ & $0.10$ & $54.6\%$ & $0.22$\\
$0.02$ & $124\%$ & $0.10$ & $69.0\%$ & $0.17$\\
$0.05$ & $143\%$ & $0.08$ & $100\%$ & $0.12$\\
$0.10$ & $169\%$ & $0.07$ & $138\%$ & $0.09$\\
\bottomrule
\end{tabular}
\end{table}

\begin{figure}[ht]
\centering
\begin{tikzpicture}
\begin{axis}[width=0.82\linewidth,height=6.2cm,xmode=log,xlabel={number of projects $N$},
  ylabel={Sharpe ratio $\mathrm{SR}_N$},xmin=1,xmax=1000,ymin=0,ymax=1.0,
  legend pos=north west,legend cell align=left,legend style={font=\footnotesize},grid=both,
  grid style={black!10},samples=50,domain=1:1000,every axis plot/.append style={thick}]
\addplot[defcol]{0.0283*sqrt(x)};
\addlegendentry{$\rho=0$}
\addplot[thmcol]{0.0283/sqrt(1/x+(1-1/x)*0.02)};
\addlegendentry{$\rho=0.02$}
\addplot[fincol]{0.0283/sqrt(1/x+(1-1/x)*0.05)};
\addlegendentry{$\rho=0.05$}
\addplot[intcol]{0.0283/sqrt(1/x+(1-1/x)*0.2)};
\addlegendentry{$\rho=0.2$}
\end{axis}
\end{tikzpicture}
\caption{Reconstruction: Sharpe ratio of an equal-weight portfolio of $N$ projects with pairwise correlation $\rho$,
from \eqref{ch12:eq-sharpecorr} with $\mathrm{SR}_1=0.0283$. For $\rho=0$ it grows as $\sqrt N$; for $\rho>0$ it saturates
at $\mathrm{SR}_1/\sqrt\rho$.}\label{ch12:fig-sharpe}
\end{figure}

\begin{coursenote}[How many projects, how much money?]
Lo's claim is that because rare-disease programmes are nearly uncorrelated, one needs ``10 or 20''
projects rather than 150, hence \$400--500 million rather than \$30 billion \ts{24}{18}{34:38}. That is not obtained by scaling the
\$200 million per project of the oncology example: 15 such projects would cost \$3 billion. It also assumes that rare-disease
projects are cheaper and have a different profile than the oncology numbers used above, which the lecture does
not quantify. The reported input is a Monte Carlo simulation with industry numbers published in his paper,
giving about $22\%$ per year and a Sharpe ratio above~1 \ts{24}{18}{35:38}; the parameters are not in the
lecture, so \cref{ch12:tab-corr} should \emph{not} be read as a reproduction of that result: with the oncology
inputs and $N=15$ a Sharpe ratio of $0.11$ is all one gets, even with independence.
\end{coursenote}

\begin{finance}[The BridgeBio story]\label{ch12:fin-bridge}
A former student of Lo (a biotech venture capitalist) reran the simulations with his own parameters,
then with further modifications; they ran about 35 versions over six months, after which he quit his job to
found a company \ts{24}{18}{36:43}. BridgeBio Pharma came out of stealth in 2017 with a private-equity investor
(KKR) taking part in a biotech for the first time; it raised about \$40 million after a friends-and-family round,
then \$135 million, then \$300 million (about \$400--500 million in total), had about 15 projects, went public and raised a further
\$750 million, and traded at about \$4.6 billion of market capitalisation on the day of the lecture
\ts{24}{18}{37:45}\ts{24}{18}{38:48}. By his account about 20 projects were in the clinic, two drugs had been approved in 2021
and the Canavan gene therapy shown earlier belongs to the company \ts{24}{18}{39:52}. Lo's remark, worth
keeping: finance cannot turn a bad drug into a good one, but the wrong financing can destroy a good one; science
comes first \ts{24}{18}{41:57}. (Facts about the company are as stated in class; they were not checked.)
\end{finance}

% ==================================================================
\section{Counting successes}\label{ch12:sec-counts}

For a project financed at year 0 the natural random variable is a Bernoulli success indicator
$X_i\in\{0,1\}$ with $\Prob(X_i=1)=p$. This section is a \emph{Reconstruction} (not from the course) of what
the number of successes looks like when the indicators are independent or correlated.

\paragraph{Independent projects.}
$K=\sum_iX_i\sim\mathrm{Bin}(N,p)$. The probability of at least one success is $1-(1-p)^N$: with $p=5\%$ it is
$54\%$ for $N=15$ and $99.95\%$ for $N=150$ (the chance of nothing is $0.95^{150}=4.6\times10^{-4}$). If every success returns
$G_0=61.5$ times a single project's cost and every project has the same cost, the portfolio (total cost $N$ units) returns
$G_0K$ units, so the portfolio is at least at break-even when $K\ge N/G_0$. For $N=150$ this means $K\ge3$
($150/61.5=2.44$), which has probability $98.2\%$; for $N=15$ it means $K\ge1$, probability $54\%$.

\paragraph{Correlated outcomes.}
Two Bernoulli variables with the same $p$ have $\Cov(X_i,X_j)=\rho\,p(1-p)$, and then
\begin{equation}\label{ch12:eq-varK}
  \Var K=Np(1-p)\bigl(1+(N-1)\rho\bigr),
\end{equation}
by the same counting argument as in \cref{ch12:prop-corr}: $\rho>0$ inflates the fluctuations of $K$ by
$1+(N-1)\rho$ compared with the binomial. The distribution of $K$ also acquires a fat right tail and a fat mass at
$0$. Two standard ways to produce correlated Bernoulli indicators:

\begin{enumerate}[(a)]
  \item \emph{One-factor Gaussian model} (reconstruction; the model behind many credit-risk portfolios). Let
        $Z_i=\sqrt a\,Y+\sqrt{1-a}\,\varepsilon_i$ with $Y,\varepsilon_i$ i.i.d.\ $N(0,1)$ and
        declare project $i$ a success if $Z_i<c=\Phi^{-1}(p)$. Given the common factor $Y$ the projects are independent with
        success probability $q(Y)=\Phi\bigl((c-\sqrt a\,Y)/\sqrt{1-a}\bigr)$, so
        $\Prob(K\ge k)=\E_Y\bigl[\Prob(\mathrm{Bin}(N,q(Y))\ge k)\bigr]$, a one-dimensional integral. The Bernoulli
        correlation $\rho$ is smaller than the latent correlation $a$ (it is $0.026$ for $a=0.1$ and $0.058$ for $a=0.2$).
  \item \emph{Uncertainty about $p$} (beta-binomial). The lecturer's distinction between risk and uncertainty has an exact
        counterpart. Suppose the projects share an unknown success rate $\Pi\sim\mathrm{Beta}(\alpha,\beta)$ with
        $\E\Pi=p$ and, given $\Pi$, are independent Bernoulli($\Pi$). Then $\Cov(X_i,X_j)=\Var\Pi$ and
        \begin{equation}\label{ch12:eq-rhobeta}
          \rho=\frac{\Var\Pi}{p(1-p)}=\frac{1}{\alpha+\beta+1}.
        \end{equation}
        Indeed $\Var\Pi=p(1-p)/(\alpha+\beta+1)$ for a beta distribution. The parameter $\alpha+\beta$ is the
        ``effective sample size'' of the prior information about $p$: if the historical rate ($3.4\%$ against $5\%$)
        were known to only about 50 observations, $\alpha+\beta=49$ and $\rho=0.02$, a correlation that no amount of
        diversification can remove (\cref{ch12:tab-corr} with $\rho=0.02$).
\end{enumerate}

\begin{table}[ht]
\centering\small
\caption{Reconstruction: probability that the portfolio of $N$ projects with $p=5\%$ has at least $k$ successes,
independent case ($\rho=0$), one-factor Gaussian model with latent correlation $a$ and beta-binomial
($\alpha+\beta=49$, $\rho=0.02$). $N=150,k=3$ is break-even in the sense of the text.}\label{ch12:tab-counts}
\begin{tabular}{@{}lrrrr@{}}
\toprule
& $N=15,\ k\ge1$ & $N=150,\ k\ge1$ & $N=150,\ k\ge3$ & $N=150,\ k\ge4$\\
\midrule
independent          & $0.537$ & $0.9995$ & $0.982$ & $0.945$\\
Gaussian, $a=0.1$    & $0.482$ & $0.968$  & $0.821$ & $0.730$\\
Gaussian, $a=0.2$    & $0.432$ & $0.899$  & $0.689$ & $0.601$\\
Gaussian, $a=0.5$    & $0.300$ & $0.616$  & $0.419$ & $0.368$\\
beta-binomial        & $0.491$ & $0.970$  & $0.838$ & $0.753$\\
\bottomrule
\end{tabular}
\end{table}

\begin{intuition}[What correlation does to the tails]
The mean number of successes ($7.5$ out of 150) is unchanged; what changes is the probability of the unlucky
scenario in which \emph{everything fails together}. That scenario is invisible in the average and decisive for
someone who has to repay creditors. This is why the lecture insists on the near-independence of rare diseases
\ts{24}{18}{33:36} and why the portfolio manager's key parameter is a correlation, not a mean.
\end{intuition}

% ==================================================================
\section{From portfolio to structure: a securitisation sketch}\label{ch12:sec-tranche}

In the lecture the \$30 billion is obtained by appeal to the investing public \ts{24}{18}{30:30} and, for fusion, to
retail money \ts{24}{18}{57:45}. The natural instrument, which Lo calls a ``mega fund'' \ts{24}{18}{54:42} and
which the past final project of 2024 L1 develops for paediatric oncology, is a pool of projects funded by
tranches of different seniority. He does not describe tranches in this lecture; the following is a generic
\emph{Reconstruction}.

\begin{example}[Reconstruction: one senior debt tranche on the 150-project pool]\label{ch12:ex-tranche}
Pool of $N=150$ projects, all financed at year 0 by \$30 billion, with $K$ successes and terminal cash flow
$V=12.3\,K$ (billion dollars) at year 10. Suppose the whole \$30 billion is raised as zero-coupon debt at $5\%$,
so the face value due at year 10 is $30\times1.05^{10}=48.9$. It is repaid if and only if
$12.3K\ge48.9$, i.e.\ $K\ge4$. With independent projects $\Prob(K\ge4)=0.945$; in the one-factor model the same probability is
$0.73$ for $a=0.1$ and $0.60$ for $a=0.2$ (\cref{ch12:tab-counts}). A structure with a junior equity tranche that
absorbs the first losses, and only a fraction of the capital raised as senior debt, would push the
senior repayment probability far higher; but the equity holders take the correlation risk, and the
example shows that its price depends mostly on $\rho$, not on $N$.
\end{example}
The moral of \cref{ch12:ex-tranche} is that securitisation transforms one risk profile into several
(the same principle as in structured credit), and that a rating of the senior tranche is a statement about the joint
law of the project outcomes. The correlation input is exactly what the two rare-disease and fusion arguments in the
lecture are about.

% ==================================================================
\section{Where the data science and AI enter}\label{ch12:sec-ml}

Despite the title, the lecture says almost nothing about machine learning: artificial intelligence appears in
passing as a competitor for investors' money \ts{24}{18}{1:16:35}, and Lo's affiliation with CSAIL
opens the talk \ts{24}{18}{0:00}. His actual ``data science'' is a \emph{simulation}
of the portfolio with parameters supplied by industry experts \ts{24}{18}{35:38}. What follows is a
short \emph{Reconstruction} (not from the course) of how the data would enter a real model.

\begin{enumerate}[(a)]
  \item \emph{Success probability from features.} Model the outcome $Y_i\in\{0,1\}$ of project $i$ with covariates $x_i$
        (indication, mechanism, phase, sponsor, biomarker) as $\Prob(Y_i=1\mid x_i)=\sigma(w\T x_i)$, with
        $\sigma(u)=1/(1+e^{-u})$, and fit $w$ by minimising the log-loss
        $-\sum_i\bigl[Y_i\log\sigma(w\T x_i)+(1-Y_i)\log(1-\sigma(w\T x_i))\bigr]$ with a penalty on $\|w\|$
        (regularisation, \cref{ch:regularized}; more general classifiers, \cref{ch:ml}). The output is a
        project-specific $p_i$ in place of the pooled $p=5\%$; the portfolio mean is then $\sum_ip_iG_i$.
  \item \emph{Shrinkage and uncertainty.} With few historical trials, the estimate of $p$ is imprecise, and a Bayesian update
        of a $\mathrm{Beta}(\alpha,\beta)$ prior with $s$ successes in $n$ trials gives
        $\mathrm{Beta}(\alpha+s,\beta+n-s)$ with mean $(\alpha+s)/(\alpha+\beta+n)$. By \eqref{ch12:eq-rhobeta}
        the residual uncertainty about $p$ is the common factor that limits diversification; more data
        lowers $\rho$.
  \item \emph{Dependence between projects.} Correlations (shared mechanisms, shared regulators, shared
        macro conditions for funding) are estimated from the same features, e.g.\ by treating shared
        attributes as factors, as in \cref{ch:pca}; the sample correlation of a handful of binary outcomes
        is very noisy, so the structure of the factors matters more than the numbers.
\end{enumerate}

\begin{coursenote}[What the lecture actually contains]
No data set, code or slide of the lecture is available. The only quantitative statements are those reproduced in
\cref{ch12:sec-single,ch12:sec-portfolio}: $\$200$M, 10 years, $5\%$, \$12.3B, $12\%$, $423.5\%$, $150$, $34.6\%$, Sharpe
$\approx0.34$--$0.35$, the rare-disease reduction to 10--20 projects and \$400--500M, and a Sharpe ratio above 1 from a
Monte Carlo of a published paper. Everything else in this chapter is derived here.
\end{coursenote}

% ==================================================================
\section{Second example: a fusion mega fund}\label{ch12:sec-fusion}

The second half of the lecture applies the same idea to energy. The argument is that the ``energy transition''
is not a transition: consumption of oil, gas and coal is still growing as more people join the middle class,
so low-carbon, on-demand sources are needed in addition \ts{24}{18}{42:02}\ts{24}{18}{43:03}. The candidates listed are fusion, fission and
geothermal; fusion promises more energy per kilogram of fuel than fission and no long-lived waste, but requires the three
Lawson conditions of temperature, pressure (density) and confinement time \ts{24}{18}{47:17}\ts{24}{18}{48:20}. Lo points to
a 20-tesla magnet built at MIT with Commonwealth Fusion Systems in 2021 and a laboratory reaction with net energy
gain \ts{24}{18}{49:23}\ts{24}{18}{50:27}.

His question, asked with Dennis Whyte, is whether finance can pay for fusion. The idea is to \emph{decompose} ``a
commercially viable power plant'' into a portfolio of enabling technologies (sustained
plasma at about 100 million degrees, ultra-clean vacuum systems, magnetic confinement, and so on), each with a value of its
own, for instance a vacuum system also useful in drug development, and to treat the portfolio as in the biomedical
case, with perhaps 150 components \ts{24}{18}{51:34}\ts{24}{18}{52:37}\ts{24}{18}{53:39}. He adds a
board of scientific and business advisers whose participation would follow once serious capital is deployed
\ts{24}{18}{54:42}.

A student asks whether the components would not be highly correlated. The answer is twofold: complete reactors from
different companies are correlated, but not perfectly, because they use different approaches (magnetic versus
inertial); and, more importantly, technologies such as heating, vacuum and magnet shaping have little to do with each
other \ts{24}{18}{56:44}\ts{24}{18}{57:45}. In the language of \cref{ch12:prop-corr}: the argument is that
$\rho$ is small, and it should be treated as the key empirical input, since the floor $\mathrm{SR}_1/\sqrt\rho$
is not eliminated by adding components that share a common ``fusion is feasible'' factor (compare \eqref{ch12:eq-rhobeta}).

\begin{coursenote}[Erratum: the retail-funding arithmetic]
Lo asks for a one-time \$1,500 from each of $3\%$ of 130 million households, ``more than \$30 billion''
\ts{24}{18}{56:44}\ts{24}{18}{57:45}. The product is
\[
  130\times10^{6}\times0.03\times1500=\$5.85\times10^{9},
\]
about \$5.9 billion, not \$30 billion; reaching \$30 billion at \$1,500 each needs about $15\%$ of households.
(He also says once that fusion is ``when we split an atom'': that is fission; fusion joins nuclei. Both are slips of
speech in the transcript; the intended meaning is clear.)
\end{coursenote}

\section{The rest of the lecture}\label{ch12:sec-rest}

\paragraph{The Harvey Lodish story} \ts{24}{18}{58:45}. As a young faculty member in 1983 Lodish was
approached to help develop a treatment for Gaucher's disease, a genetic enzyme deficiency; the company became Genzyme,
which got the first enzyme-replacement therapy approved in 1991 and was sold to Sanofi for \$20 billion in 2014.
In 2002 his unborn grandson was diagnosed with the same disease, and the drug saved him. The point Lo draws is that finance need not be
zero-sum: ``doing well by doing good'' \ts{24}{18}{1:02:57}.

\paragraph{Questions.}
\begin{itemize}
  \item \emph{Patents} \ts{24}{18}{1:05:05}: a patent gives exclusive use of the compound for 20 years (which is why revenues
        run over years 11--20 in \cref{ch12:sec-single}); afterwards anyone can manufacture cheaply, like aspirin.
        Producing the pill is cheap; the expensive part is the trials that lead to approval.
  \item \emph{How to find and evaluate opportunities} \ts{24}{18}{1:07:08}\ts{24}{18}{1:12:24}: expertise. Neil Kumar (BridgeBio's
        co-founder, a chemical-engineering PhD) worked at a consultancy on pharma strategy, then at a biotech venture
        firm, which had passed on single-drug projects as too risky; the simulations gave him the conviction to fund
        them himself.
  \item \emph{How can students invest?} \ts{24}{18}{1:08:11}: BridgeBio is public, at very high risk; private companies are
        harder; MIT is rich in opportunities for investors; Lo also teaches health-care financing (15.482, with an online
        version in MIT's Open Learning Library) \ts{24}{18}{1:11:21}.
  \item \emph{Crises} \ts{24}{18}{1:15:32}: COVID-19 temporarily removed the valley of death: the time between the
        sequencing of the virus and the first human injection was 63 days, against 3--5 years normally; afterwards the valley returned
        because vaccines are not attractive to investors compared with, for example, AI \ts{24}{18}{1:17:39}\ts{24}{18}{1:18:40}.
\end{itemize}

% ==================================================================
\begin{keyformulas}
\begin{itemize}
  \item Sharpe ratio: $\mathrm{SR}=(\E R-r_f)/\sqrt{\Var R}$.
  \item Single project (lecture): cost \$200M, 10 years, $p=5\%$, payoff \$12.3B; $\E G=3.075$ per dollar,
        $\sqrt{\Var G}=61.5\sqrt{p(1-p)}=13.4$ over 10 years ($424\%$ per year after dividing by $\sqrt{10}$).
  \item $N$ i.i.d.\ projects: $\Var R_P=\sigma^2/N$, $\mathrm{SR}_N=\sqrt N\,\mathrm{SR}_1$.
  \item Common correlation $\rho$: $\Var R_P=\sigma^2\bigl(\tfrac1N+\tfrac{N-1}{N}\rho\bigr)$,
        $\mathrm{SR}_N\to\mathrm{SR}_1/\sqrt\rho$.
  \item Successes: $\Var K=Np(1-p)\bigl(1+(N-1)\rho\bigr)$; $\Prob(K\ge1)=1-(1-p)^N$ if independent.
  \item Unknown $p\sim\mathrm{Beta}(\alpha,\beta)$: $\rho=1/(\alpha+\beta+1)$.
\end{itemize}
\end{keyformulas}

\section*{Exercises}
\addcontentsline{toc}{section}{Exercises}

\begin{exercise}[not from the course]
Verify \eqref{ch12:eq-annuity}: find the discount rate for which \$2 billion in each of years 11--20 is worth
\$12.3 billion at year 10, and compute the value at year 0 (in units of the \$200M cost) of the project of
\cref{ch12:ex-single}, using its expected payoff.
\end{exercise}

\begin{exercise}[not from the course]
Prove \eqref{ch12:eq-varK}. Then compute the standard deviation of $K$ for $N=150$, $p=0.05$ and $\rho\in\{0,0.02,0.05\}$
and compare it with the mean $7.5$.
\end{exercise}

\begin{exercise}[not from the course]
With $R_P=\frac1N\sum_iR_i$ and general covariance matrix $\Sigma$, the variance is $\frac{1}{N^2}\mathbf 1\T\Sigma\mathbf 1$
(\cref{ch:portfolio}). Show that for the common-correlation matrix $\Sigma=\sigma^2[(1-\rho)I+\rho\,\mathbf 1\mathbf 1\T]$ this reproduces
\eqref{ch12:eq-varcorr}, and find its eigenvalues. Which eigenvector carries the ``undiversifiable'' part?
\end{exercise}

\begin{exercise}[not from the course]
Show that in the beta-binomial model of \cref{ch12:sec-counts}(b), $\rho=1/(\alpha+\beta+1)$, and that after observing
$s$ successes in $n$ historical trials of a common process the posterior correlation between two future projects is
$1/(\alpha+\beta+n+1)$. What sample size $n$ makes it fall below $0.005$ starting from $\alpha+\beta=49$?
\end{exercise}

\begin{exercise}[not from the course]
In \cref{ch12:ex-tranche} let the senior debt be only a fraction $f$ of the \$30 billion, with face value $30f\times1.05^{10}$.
Find, for the independent model, the smallest $f$ such that the senior tranche is repaid with probability at least $0.99$.
(The rest of the capital is equity and absorbs the first losses.)
\end{exercise}

\begin{exercise}[not from the course]
Erratum check: with the numbers of \cref{ch12:sec-fusion}, what fraction of the 130 million households must contribute \$1,500
to raise \$30 billion? What one-time contribution per household would be needed from $3\%$ of them?
\end{exercise}

\begin{exercise}[not from the course]
Simulate the one-factor Gaussian model of \cref{ch12:sec-counts}(a) for $N=150$, $p=0.05$, $a=0.1$ and estimate
$\Prob(K\ge3)$ by Monte Carlo with $10^5$ paths; compare with \cref{ch12:tab-counts}.
\end{exercise}

% !TEX root = ../main.tex
\chapter{Volatility Modeling}\label{ch:volatility}
\begin{paracol}{2}

\begin{sources}
\begin{tabularx}{\linewidth}{@{}l l L@{}}
\toprule
\textbf{Lecture} & \textbf{Speaker} & \textbf{Content}\\
\midrule
2024 L19 & P.~Kempthorne & Definition and annualisation of volatility; historical estimators (simple, exponential, RiskMetrics); geometric Brownian motion and maximum likelihood; Garman--Klass and the other open-high-low-close (OHLC) estimators; case study of the S\&P~500 and of ARKK (naive, ARIMA and seasonal ARIMA models of rolling volatility); Poisson jump diffusion; Laplace law; ARCH, GARCH, maximum likelihood, stylised facts, $t$ innovations, euro/dollar case study.\\
2013 L9 (from 19:52) & P.~Kempthorne & The same programme (slides \file{lecnote9}): historical volatility, time scales in GBM (trading versus calendar days), Garman--Klass, jump diffusion, ARCH, GARCH, fits to the euro/dollar rate. The first 20 minutes finish univariate time series (information criteria, Dickey--Fuller test, Case Study~3), see \cref{ch:timeseries}.\\
2013 L11 (first 28 min) & P.~Kempthorne & Recap and completion: GARCH$(1,1)$ as ARMA for squared returns, long-run variance, likelihood and parameter constraints, residual diagnostics, $t$ innovations and their degrees of freedom, value at risk, volatility clustering, mean reversion and persistence, RiskMetrics. The rest of the lecture is multivariate time series (\cref{ch:timeseries}).\\
\bottomrule
\end{tabularx}

\smallskip
\textbf{Files.} 2024 slides \file{mit18_642_f24_lec17_1.pdf} (``Volatility Modeling'', 31 slides, 36 pages with overlays) and \file{mit18_642_f24_lec17_2.pdf} (35 pages, R case study on the S\&P~500); zip \file{mit18_642_f24_rstudio_pset5.zip} (case studies of the S\&P~500 and of ARKK, R Markdown source and knitted PDF); 2013 slides \file{MIT18_S096F13_lecnote9.pdf} (31 pages), 2013 Case Study~4 \file{CaseStudy4.pdf} (volatility of exchange-rate returns, in R). The case study loads a workspace built by \file{fm_casestudy_fx_1.r} from the Federal Reserve data \file{fred_fxrates.txt} (setup files: appendix~\ref{app:rsetup}).
\textbf{Problem sets.} 2024 PS4 Problems~3--4 (Problems~1--2 and 5--6 are in \cref{ch:sp2} and \cref{ch:timeseries}); 2024 PS5 Problems~3--5 (Problems~1--2, L\'evy processes, are in \cref{ch:sp2}); 2013 PS5, all four problems (its Problems~1--2 are the same as 2024 PS4 Problems~3--4). Standalone exercises follow the section they practise; the connected ones are at the end of the chapter.
\textbf{Not published or not covered.} The papers of Garman--Klass, Parkinson and Yang--Zhang, and the RiskMetrics technical document, are distributed on Canvas or on the web and are not in the download; the auxiliary R function file \file{test_vol1b.r} used by Case Study~4 is also missing, so the two-time-scale fits of \cref{ch13:sec-gbm} are described from the printed output. The 2024 lecture reached the end of the slides; slide~31 (extensions of GARCH) is only a list of names. Stochastic-volatility and implied-volatility models are named but not developed in either year.
\end{sources}

\section*{Motivation}
Prices move, and \emph{how much} they move is the central risk quantity of finance: it sets the width of a confidence
band for tomorrow's P\&L (\cref{ch:var}), the price of an option (\cref{ch:bs}), the size of a margin
requirement (\cref{ch:counterparty}) and the weights of a portfolio (\cref{ch:portfolio}). Two facts make it
interesting.
\begin{itemize}
  \item Volatility is \emph{not observed}. The price is; the standard deviation of its increments has to be estimated,
        and the quality of the estimate depends dramatically on the information used (closing prices only, or the
        whole intraday range) and on the model assumed (constant volatility, or volatility that moves).
  \item Volatility is \emph{not constant}. Calm and turbulent weeks alternate, so that squared returns are strongly
        autocorrelated although returns themselves are not: white noise that is not independent
        (\cref{ch09:sec-sacf}). The models of Engle and Bollerslev, ARCH and GARCH, put this in one line and, with
        one extra parameter, produce fat tails, mean reversion of volatility and a forecasting formula.
\end{itemize}
For a physicist the situation is that of a noisy detector whose noise level drifts: one measures a
\emph{scale parameter}, one has to choose between averaging over a long window (small statistical error, large
bias if the scale drifts) and a short one, and the statistics of the scale itself becomes a subject of
study. The chapter follows the lectures: historical estimators (\cref{ch13:sec-hist}), the constant-volatility
benchmark (\cref{ch13:sec-gbm}), range-based estimators and the S\&P~500 case study (\cref{ch13:sec-ohlc}),
mixtures for fat tails (\cref{ch13:sec-mix}), ARCH and GARCH (\cref{ch13:sec-arch,ch13:sec-garch}), estimation and the
euro/dollar case study (\cref{ch13:sec-mle}), and a summary of stylised facts (\cref{ch13:sec-facts}).

% ==================================================================
\section{Volatility and its historical estimation}\label{ch13:sec-hist}

\subsection{Definition and annualisation}\label{ch13:sec-def}
The lectures define \emph{volatility} as the annualised standard deviation of the change in price (or value) of a
financial security \ts{24}{19}{1:00}\ts{13}{9}{21:57}. It is a risk measure and an input to option pricing.
Given prices $P_t$ the \emph{log return} over one period is $R_t=\log(P_t/P_{t-1})$; log returns add over
successive periods whereas percentage returns multiply \ts{24}{19}{4:13}, as in \cref{ch09:eq-logret}. If $\{R_t\}$ is
covariance stationary \ts{24}{19}{5:16}, the \emph{period volatility} is $\sigma=\sqrt{\Var R_t}$, and the
\emph{annualised volatility} is
\begin{equation}\label{ch13:eq-annual}
  \mathrm{vol}=\sqrt N\,\sigma,\qquad N=252\ (\text{daily}),\ 52\ (\text{weekly}),\ 12\ (\text{monthly}),
\end{equation}
where $252$ is the number of business days in a year \ts{24}{19}{2:00}\ts{13}{9}{25:10}. The justification is
the variance of a sum: the return over a year is the sum of $N$ period returns, and if these are
\emph{uncorrelated} with a common variance, $\Var\bigl(\sum_{t=1}^NR_t\bigr)=N\sigma^2$. A daily standard deviation of
$1\%$ is thus a volatility of $1\%\times\sqrt{252}=15.9\%$. Scaling all horizons to one year makes assets and
maturities comparable. Orders of magnitude quoted in 2013: bond yields move by less than $5$ percentage points a
year, currencies by about $10\%$, equity indices by $30\%$--$40\%$ in a nervous year \ts{13}{9}{23:03}.

\begin{remark}[What the square-root rule needs]\label{ch13:rem-sqrt}
Only \emph{uncorrelatedness} is used, not independence, so \eqref{ch13:eq-annual} holds for the
\emph{unconditional} variance of ARCH and GARCH returns, which are uncorrelated but dependent
(\cref{ch13:prop-mds}). It fails when returns are autocorrelated (\cref{ch08:prop-autocorr}) and, as shown in
\cref{ch13:sec-forecast}, it fails for the \emph{conditional} variance over a horizon: after a turbulent day
the volatility expected over the next month is smaller than $\sqrt{21}$ times today's volatility.
\end{remark}

\subsection{Sample volatility and its precision}\label{ch13:sec-sample}
Given $R_1,\dots,R_T$ the sample estimate is
\begin{equation}\label{ch13:eq-samplevol}
  \hat\sigma=\sqrt{\frac1{T-1}\sum_{t=1}^T(R_t-\bar R)^2},\qquad \bar R=\frac1T\sum_{t=1}^TR_t,
\end{equation}
which uses the unbiased estimator of the variance, hence the divisor $T-1$ \ts{24}{19}{5:16}. How precise is it?

\begin{proposition}[Precision of the sample volatility]\label{ch13:prop-precision}
Let $R_t$ be i.i.d.\ with kurtosis $\kappa$ (\cref{ch03:def-moments}). Then $\hat\sigma^2$ is asymptotically normal with
$\Var\hat\sigma^2\approx(\kappa-1)\sigma^4/T$. For Gaussian returns the result is exact in the form
$(T-1)\hat\sigma^2/\sigma^2\sim\chi^2_{T-1}$ (\cref{ch03:sec-chisq}), so $\Var\hat\sigma^2=2\sigma^4/(T-1)$, and by the delta method
\[
  \frac{\operatorname{sd}(\hat\sigma)}{\sigma}\approx\frac1{\sqrt{2(T-1)}} .
\]
\end{proposition}
\begin{proof}
Put $\mu=\E[R]$ and $Y_t=(R_t-\mu)^2$, i.i.d.\ with mean $\sigma^2$ and variance $\E[(R-\mu)^4]-\sigma^4=(\kappa-1)\sigma^4$.
Then $\hat\sigma^2=\frac{T}{T-1}\bigl[\overline Y-(\bar R-\mu)^2\bigr]$, and $(\bar R-\mu)^2=O(1/T)$ does not affect the leading term, so
the central limit theorem applied to $\overline Y$ gives $\Var\hat\sigma^2\approx(\kappa-1)\sigma^4/T$. For $\kappa=3$ this is $2\sigma^4/T$.
The $\chi^2$ statement is the classical result for the sample variance of a normal sample, and $\Var\chi^2_{T-1}=2(T-1)$. Finally
$\hat\sigma=g(\hat\sigma^2)$ with $g(x)=\sqrt x$, $g'(\sigma^2)=1/(2\sigma)$, so $\Var\hat\sigma\approx\Var(\hat\sigma^2)/(4\sigma^2)=\sigma^2/(2(T-1))$.
\end{proof}
With $T=21$ daily returns (one month) the relative standard error of $\hat\sigma$ is $15.8\%$: a monthly estimate
of $12\%$ annual volatility is uncertain by about $\pm1.9$ percentage points; with $T=252$ it is $4.5\%$. For a
heavy-tailed series with $\kappa=6$ the variance of $\hat\sigma^2$ is $2.5$ times larger than in the Gaussian case.
(This is the quantitative content of the ``trade-off'' below and of 2013 PS5 Problem~4.)

\begin{exercise}[PS5 (2013), Problem 4: are year-to-year differences in volatility significant?]\label{ch13:ex-ps5-4}
The annual sample variances of the daily log returns of the S\&P~500 close are (year: daily variance, number of days)
2006: $3.981\times10^{-5}$, 251; 2007: $1.019\times10^{-4}$, 251; 2008: $6.677\times10^{-4}$, 253; 2009: $2.950\times10^{-4}$, 252; 2010: $1.295\times10^{-4}$, 252; 2011: $2.164\times10^{-4}$, 252; 2012: $6.459\times10^{-5}$, 250.
Model the returns of one year as an i.i.d.\ sample from $N(\mu,\sigma^2)$ and let $\hat\sigma^2$ be the unbiased sample variance.
\begin{enumerate}[(a)]
  \item Prove that $\hat\sigma^2\sim\frac{\sigma^2}{n-1}\chi^2_{n-1}$.
  \item Invert the $\chi^2$ percentiles $q_{0.025},q_{0.975}$ to obtain a two-sided $95\%$ confidence interval for $\sigma^2$ (the problem gives the percentiles for $249$ to $252$ degrees of freedom): compute it for 2008, express it as an annualised volatility ($\sqrt{253}\,\sigma$),
      and say whether the sample volatility of any other year falls inside it.
  \item Under the null hypothesis of equal daily variances in two years, show that $S=\hat\sigma^2_{2008}/\hat\sigma^2_{2007}$ has an $F$ distribution (state its degrees of freedom); the problem gives $q_{0.025}=0.7804$, $q_{0.975}=1.2815$ for $F_{252,250}$. Compute $S$ for 2008 against 2007, find the
      $P$-value (the $\alpha$ at which the test is on the boundary), and repeat for 2008 against 2006.
\end{enumerate}
\end{exercise}

\subsection{Historical predictors: simple and exponential averages}\label{ch13:sec-ewma}
To \emph{predict} volatility one needs a single-period sample estimate $\hat\sigma_t^2$ for each period
$t$. If $t$ indexes months and the data are daily, it is the sample variance of the daily returns of month $t$; if $t$
indexes days and the data are daily, it is simply $R_t^2$ (the mean is negligible); with intraday data the daily value is the sum of
squared intraday returns (the \emph{realised variance}) \ts{24}{19}{7:28}\ts{13}{9}{24:06}. Four predictors
of $\sigma^2_{t+1}$ are then defined \ts{24}{19}{6:23}\ts{13}{9}{27:15}:
\begin{align}
  \text{historical average:}&\quad \tilde\sigma^2_{t+1}=\frac1t\sum_{j=1}^t\hat\sigma_j^2 &&\text{(all data)},\notag\\
  \text{simple moving average:}&\quad \tilde\sigma^2_{t+1}=\frac1m\sum_{j=0}^{m-1}\hat\sigma^2_{t-j} &&\text{(last $m$)},\notag\\
  \text{exponential moving average (EMA):}&\quad \tilde\sigma^2_{t+1}=(1-\beta)\hat\sigma^2_t+\beta\,\tilde\sigma^2_t,\quad 0\le\beta\le1
  &&\text{(all data)},\label{ch13:eq-ema}\\
  \text{EWMA over the last $m$:}&\quad \tilde\sigma^2_{t+1}=\frac{\sum_{j=0}^{m-1}\beta^j\hat\sigma^2_{t-j}}{\sum_{j=0}^{m-1}\beta^j}
  &&\text{(last $m$)}.\notag
\end{align}
The EMA is a recursion and can be applied over the whole history. Unrolling it gives the same weights as the last
line for $m=\infty$.

\begin{proposition}[Weights of the exponential average]\label{ch13:prop-ewma}
If the recursion \eqref{ch13:eq-ema} runs from the remote past, $\tilde\sigma^2_{t+1}=(1-\beta)\sum_{j\ge0}\beta^j\hat\sigma^2_{t-j}$.
The weights sum to $1$; the mean age of the data is $\beta/(1-\beta)$, the half-life of a weight is $\ln2/\ln(1/\beta)$, and if the
$\hat\sigma_t^2=R_t^2$ come from i.i.d.\ zero-mean Gaussian returns, $\tilde\sigma^2$ is unbiased with
\[
  \Var\tilde\sigma^2_{t+1}=2\sigma^4\,\frac{1-\beta}{1+\beta},\qquad\text{i.e. an effective sample size }\
  n_{\mathrm{eff}}=\frac{1+\beta}{1-\beta}.
\]
\end{proposition}
\begin{proof}
Iterate the recursion $n$ times: $\tilde\sigma^2_{t+1}=(1-\beta)\sum_{j<n}\beta^j\hat\sigma^2_{t-j}+\beta^n\tilde\sigma^2_{t+1-n}$, and let $n\to\infty$.
The mean age is $(1-\beta)\sum_jj\beta^j=\beta/(1-\beta)$. The weight falls to one half when $\beta^j=\frac12$. For independent
$\hat\sigma^2_t$ with variance $2\sigma^4$ the variance of the weighted sum is $2\sigma^4\sum_jw_j^2=2\sigma^4(1-\beta)^2/(1-\beta^2)$, and a
simple average of $n$ observations has variance $2\sigma^4/n$.
\end{proof}
\begin{center}\small
\begin{tabular}{@{}lccc@{}}
\toprule
 & $\beta=0.90$ & $\beta=0.94$ & $\beta=0.97$\\
\midrule
mean age (days) & 9.0 & 15.7 & 32.3\\
half-life (days) & 6.6 & 11.2 & 22.8\\
$n_{\mathrm{eff}}$ (days) & 19 & 32.3 & 65.7\\
\bottomrule
\end{tabular}
\end{center}
\begin{figure}[ht]
\centering
\begin{tikzpicture}
\begin{axis}[width=0.86\linewidth,height=5.4cm,xmin=-1,xmax=61,ymin=0,ymax=0.065,
  axis lines=left,xlabel={age $j$ (days)},ylabel={weight $w_j$},
  legend style={font=\scriptsize,draw=none,fill=none},legend pos=north east,
  tick label style={font=\small},label style={font=\small},scaled y ticks=false]
  \addplot[very thick,defcol,domain=0:60,samples=61,ycomb,mark=none] {0.06*0.94^x};
  \addlegendentry{EMA $\beta=0.94$}
  \addplot[very thick,thmcol,dashed,domain=0:60,samples=61] {0.03*0.97^x};
  \addlegendentry{EMA $\beta=0.97$}
  \addplot[very thick,intcol,const plot,domain=0:21,samples=22] {1/21};
  \addlegendentry{simple average, $m=21$}
\end{axis}
\end{tikzpicture}
\caption{Weights given to the squared return $j$ days ago. The simple average of $21$ days is flat and then cuts off,
so that a spike is kept for exactly $21$ days and then dropped; the exponential averages fade out smoothly.}\label{ch13:fig-ewma}
\end{figure}
\begin{finance}[RiskMetrics]\label{ch13:fin-riskmetrics}
The exponential average is the workhorse of the RiskMetrics methodology that J.P.~Morgan published in the 1990s
\ts{24}{19}{10:39}.\begin{aside}[RiskMetrics and the lecturer]
Kempthorne was in the Sloan School's sponsoring research centre at the time.
\end{aside} Its technical document
tabulates decay factors for foreign exchange, zero-coupon yields, equity indices and other assets, chosen by comparing the
mean squared error of the resulting volatility forecasts: they are generally well above $0.9$, more like $0.95$,
which says that volatility changes slowly \ts{24}{19}{11:41}; the 2013 lecture translates this into an averaging
period of roughly 45 to 90 days \ts{13}{9}{26:12}. (The published values are $\beta=0.94$ for daily and
$0.97$ for monthly data; this is not stated in class.) The point of a benchmark is not optimality but agreement:
everybody knows what is measured \ts{24}{19}{12:44}. The methodology is adopted without a long-run mean of
volatility, which is why it corresponds to $\alpha_0=0$, $\alpha_1+\beta_1=1$ in a GARCH$(1,1)$ (\cref{ch13:sec-igarch}).
\end{finance}

\subsection{Regression forecasts, trade-offs and evaluation}\label{ch13:sec-regr}
A more flexible predictor regresses the target on the last $p$ single-period estimates,
$\tilde\sigma^2_{t+1}=\gamma_{1,t}\hat\sigma_t^2+\gamma_{2,t}\hat\sigma^2_{t-1}+\dots+\gamma_{p,t}\hat\sigma^2_{t-p+1}+u_t$, fitted on all data or
on a rolling window of the last $m$ \ts{24}{19}{9:36}. It resembles but is not the autoregression of the series $\hat\sigma^2_t$, since the
coefficients are re-estimated as the window moves.

\begin{coursenote}[Erratum in the slides]
Slide~5 of \file{lec17_1} (and slide~7 of the 2013 deck) writes the coefficient of the second lag as $\gamma_{1,t}$, the same symbol as
that of the first lag; it should be $\gamma_{2,t}$, as above.
\end{coursenote}

Every choice (window length $m$, decay $\beta$, number of lags) balances two effects: more data reduces the statistical error of the estimate
(\cref{ch13:prop-precision}), but data close to time $t$ describe the current volatility better if it changes
\ts{24}{19}{9:36}\ts{13}{9}{28:15}. The choice is made \emph{out of sample}, separately for asset classes, sampling frequencies and forecast horizons,
with error measures such as the mean squared error, the mean absolute error and the mean absolute percentage error
\begin{equation}\label{ch13:eq-mape}
  \mathrm{MSE}=\frac1n\sum e_t^2,\qquad \mathrm{MAE}=\frac1n\sum|e_t|,\qquad
  \mathrm{MAPE}=\frac{100}n\sum\Bigl|\frac{e_t}{y_t}\Bigr|,\qquad e_t=y_t-\hat y_t,
\end{equation}
\ts{24}{19}{9:36}\ts{13}{9}{29:17}. (MAPE requires a target that never comes close to zero.)


% ==================================================================
\section{The Gaussian benchmark: geometric Brownian motion}\label{ch13:sec-gbm}

\subsection{The model and its returns}
In continuous time the standard model of a price is the \emph{geometric Brownian motion}
\begin{equation}\label{ch13:eq-gbm}
  \dd S(t)=\mu S(t)\dd t+\sigma S(t)\dd W(t),
\end{equation}
with $\sigma$ the volatility and $\mu$ the mean return per unit time, $W$ a standard Brownian motion: its increments over
an interval are $N(0,\text{length})$ and increments over disjoint intervals are independent
\ts{24}{19}{12:44}\ts{13}{9}{31:23}. Dividing by $S$ shows that the relative increment is a
drift plus a random walk with Gaussian steps. The theory is the subject of \cref{ch:sp2,ch:stochcalc}; what is needed here is its consequence for sampled data. If the price is observed at
$t_0<t_1<\dots<t_n$ with $\Delta_j=t_j-t_{j-1}$, the log returns $R_j=\log(S(t_j)/S(t_{j-1}))$ are \emph{independent} with
\begin{equation}\label{ch13:eq-gbmret}
  R_j\sim N(\mu_*\Delta_j,\ \sigma^2\Delta_j),\qquad \mu_*=\mu-\tfrac12\sigma^2 ,
\end{equation}
that is, $\log S(t)$ is a Brownian motion with drift $\mu_*$ and volatility $\sigma$ \ts{24}{19}{14:00}\ts{13}{9}{33:36}. The
drift of the \emph{log} price differs from the drift of the price by $-\sigma^2/2$: expanding $\log S$ to second order,
$\dd\log S=\dd S/S-\tfrac12(\dd S/S)^2$, and $(\dd W)^2=\dd t$ gives the extra term \ts{24}{19}{15:05}
(the same lognormal correction as in \cref{ch04:rem-logprice}; the derivation is \cref{ch:stochcalc}). In working
with stochastic differential equations one must not stop at first order.

\subsection{Maximum likelihood, and what more data can and cannot buy}
For equally spaced data with $\Delta_j\equiv1$ the log-likelihood of $R_1,\dots,R_n$ is that of i.i.d.\ normals, and the
maximum-likelihood estimators are \ts{24}{19}{15:05}\ts{13}{9}{34:41}
\begin{equation}\label{ch13:eq-gbmmle}
  \hat\mu_*=\bar R=\frac1n\sum_{t=1}^nR_t,\qquad \hat\sigma^2=\frac1n\sum_{t=1}^n(R_t-\bar R)^2,
\end{equation}
the sample mean and the sample variance with divisor $n$ (not $n-1$); the derivation is 2013 PS5 Problem~1(a)
(\cref{ch13:ex-ps5-1}). If the $\Delta_j$ vary, ``exercise'' (\cref{ch13:ex-ps5-2}) \ts{24}{19}{17:18}. A question with a striking
answer is what happens when the number $n$ of observations in a fixed window $[0,T]$ grows, that is when the
sampling frequency increases \ts{13}{9}{35:46}. The exercises show the following. The maximum-likelihood
estimator of the drift is always the cumulative return divided by the length of the period, so its variance is
$\sigma^2/T$ \emph{whatever the sampling}: no amount of intraday data helps to measure a drift. The estimator of the
variance, on the other hand, becomes arbitrarily precise, $\Var\hat\sigma^2=O(1/n)$, and even for irregular sampling the
variance of the estimator is the same for every spacing of the same $n+1$ points (2024 PS4 Problem~4(c)).
Physically: over a fixed observation time one can determine the noise amplitude of a diffusion as well as one wishes, but the
mean velocity only as well as $\sigma/\sqrt T$. This asymmetry underlies the estimators of \cref{ch13:sec-ohlc}: extra
information inside the day sharpens $\sigma$ but not $\mu$.

\subsection{Which clock? Trading time versus calendar time}\label{ch13:sec-clock}
Case Study~4 of 2013 fits \eqref{ch13:eq-gbmret} to daily EUR/USD and GBP/USD returns from 1999 to September 2013 (3\,708 returns) on two time
scales \ts{13}{9}{36:48}. In \emph{trading time} every return is one period, $\Delta_j=1$. In \emph{clock (calendar) time}
$\Delta_j$ is the number of calendar days between the two quotes, so that the Monday return, which spans a weekend, has a variance
three times larger.
\begin{secondary}[The fitted numbers]
The annualised estimates printed by the case study are

\begin{center}\small
\begin{tabular}{@{}lcccc@{}}
\toprule
 & \multicolumn{2}{c}{trading days} & \multicolumn{2}{c}{clock time}\\
 & $\hat\mu_*$ (per year) & $\hat\sigma$ & $\hat\mu_*$ (per day) & $\hat\sigma$\\
\midrule
EUR/USD & $0.00926$ & $0.1025$ & $2.5\times10^{-5}$ & $0.1228$\\
GBP/USD & $-0.00185$ & $0.0946$ & $-5.1\times10^{-6}$ & $0.1119$\\
\bottomrule
\end{tabular}
\end{center}
\end{secondary}
The volatility depends on the model of time: $10.25\%$ against $12.2\%$ for the euro \ts{13}{9}{40:02}.
\begin{secondary}[Diagnostics: Q--Q plot and fitted percentiles]
The
diagnostic used to decide is the \emph{normal Q--Q plot} (sorted returns against the expected order statistics
$x_j=\E[X_{[j]}]$ of a standard normal sample; a straight line means Gaussian) and the histogram of the \emph{fitted percentiles}
$\hat F(R_j)$, which is uniform on $(0,1)$ if the model is right \ts{13}{9}{42:08}. Both show
\emph{leptokurtic} returns (slender centre, fat tails): the extreme fitted percentile bins are over-populated, the middle
under-populated \ts{13}{9}{45:14}. In calendar time the $1\%$ percentile is exceeded about as often as it should, so the
tail fit is slightly better, but not the centre \ts{13}{9}{46:20}.
\end{secondary}

\begin{intuition}[Subordinated clocks]\label{ch13:int-clock}
If a Brownian motion looks wrong when observed against calendar time, it may still be a Brownian motion against another
clock $\tau(t)$: $\log S(t)=B(\tau(t))$. Such processes are called \emph{subordinated}. Natural clocks are
cumulative volume or number of trades, a proxy for the flow of information reaching the market
\ts{13}{9}{41:04}. In physics the same idea appears as ``operational time'' in anomalous diffusion. The Laplace law of
\cref{ch13:sec-laplace} is the simplest example, with an exponentially distributed clock.
\end{intuition}

% ==================================================================
\section{Range-based estimators from open, high, low and close}\label{ch13:sec-ohlc}

\subsection{The setting and the notion of efficiency}\label{ch13:sec-eff}
A daily \emph{bar} reports four prices: open $O$, high $H$, low $L$ and close $C$. Close-to-close volatility uses one number per day and
ignores the rest \ts{24}{19}{18:22}\ts{13}{9}{47:20}. Garman and Klass (1980) asked how much precision the other three
prices buy. Work with $X(t)=\log S(t)$, a Brownian motion \emph{without drift} ($\mu_*=0$) with variance $\sigma^2$ per day, and
let $f\in(0,1)$ be the fraction of the day during which the market is closed \ts{24}{19}{19:27}. Day~0 closes at
$t=0$ (price $C_0$), day~1 opens at $t=f$ ($O_1$) and closes at $t=1$ ($C_1$), and
\[
  H_1=\max_{f\le t\le1}X(t),\qquad L_1=\min_{f\le t\le1}X(t),
\]
all in logarithms. The three increments $O_1-C_0\sim N(0,f\sigma^2)$, $C_1-O_1\sim N(0,(1-f)\sigma^2)$ and their sum
$C_1-C_0\sim N(0,\sigma^2)$ describe the overnight gap, the trading session and the whole day; the first two are independent.

Every estimator below is an unbiased estimator of $\sigma^2$, and estimators are compared by their variance.

\begin{definition}[Efficiency]\label{ch13:def-eff}
The \emph{efficiency} of an unbiased estimator $\hat\sigma^2$ of $\sigma^2$ is $\mathrm{eff}(\hat\sigma^2)=\Var(\hat\sigma^2_0)/\Var(\hat\sigma^2)$, where
$\hat\sigma_0^2=(C_1-C_0)^2$ is the close-to-close estimator \ts{24}{19}{25:52}. An estimator of efficiency $e$ based on $n$
days is as precise as the close-to-close estimator based on $en$ days.
\end{definition}
Since $C_1-C_0=\sigma Z$ with $Z\sim N(0,1)$, $\hat\sigma_0^2=\sigma^2Z^2$ with $\sigma^2Z^2\sim\sigma^2\chi^2_1$ has mean $\sigma^2$ and
variance $\sigma^4\Var\chi^2_1=2\sigma^4$ \ts{24}{19}{21:32}\ts{13}{9}{50:30}. The same holds for
the two normalised gap and session estimators,
\begin{equation}\label{ch13:eq-gk1}
  \hat\sigma_1^2=\frac{(O_1-C_0)^2}{f},\qquad \hat\sigma_2^2=\frac{(C_1-O_1)^2}{1-f},\qquad
  \E[\hat\sigma_i^2]=\sigma^2,\quad\Var\hat\sigma_i^2=2\sigma^4\quad(i=0,1,2),
\end{equation}
because $(O_1-C_0)/\sqrt f$ and $(C_1-O_1)/\sqrt{1-f}$ are again $\sigma Z$ \ts{24}{19}{22:41}\ts{13}{9}{51:33}.

\begin{lemma}[Combining independent unbiased estimators]\label{ch13:lem-comb}
Let $V_1,V_2$ be independent unbiased estimators of the same parameter with variances $v_1,v_2$. Among the estimators
$aV_1+(1-a)V_2$ the variance is minimal for $a=v_2/(v_1+v_2)$ (inverse-variance weighting), it equals $v_1v_2/(v_1+v_2)$, and
efficiencies \emph{add}: $\mathrm{eff}=\mathrm{eff}_1+\mathrm{eff}_2$.
\end{lemma}
\begin{proof}
$\Var=a^2v_1+(1-a)^2v_2$ is a parabola with minimum at $a=v_2/(v_1+v_2)$, where it equals $v_1v_2/(v_1+v_2)$. Hence $1/\Var=1/v_1+1/v_2$,
and efficiency is $2\sigma^4$ times $1/\Var$.
\end{proof}
Since $\hat\sigma^2_1$ and $\hat\sigma^2_2$ are independent \ts{24}{19}{23:44}, the average
$\hat\sigma^2_*=\tfrac12\hat\sigma^2_1+\tfrac12\hat\sigma^2_2$ has $\Var=\sigma^4$ and efficiency $2$: the opening price alone
doubles the information of the close. In 2013 the lecturer stressed the reading of this number: the opening price buys as much as a
second day of data \ts{13}{9}{52:40}. Weighting the gap and the session equally is right because both are
$\sigma Z$ in law, whatever the value of $f$ (which drops out of \eqref{ch13:eq-gk1}).

\subsection{Parkinson's high--low estimator}\label{ch13:sec-parkinson}
Parkinson (1980) proposed to use the range $H-L$ of the Brownian path over the day \ts{24}{19}{25:52}\ts{13}{9}{55:46}:
\begin{equation}\label{ch13:eq-parkinson}
  \hat\sigma_3^2=\frac{(H_1-L_1)^2}{4\ln2}\qquad(f=0).
\end{equation}
It is unbiased because for a driftless Brownian motion of variance $\sigma^2$ on a unit interval, $\E[(H-L)^2]=4\ln2\,\sigma^2$.
The proof needs the joint law of the maximum and the minimum, obtained by the reflection principle (Feller, 1951) and not
reproduced here; a check by simulation is given below. The same law gives the fourth moment $\E[(H-L)^4]=9\zeta(3)\sigma^4$, whence
$\Var\hat\sigma_3^2=\bigl(9\zeta(3)-16\ln^22\bigr)\sigma^4/(16\ln^22)=0.4073\,\sigma^4$.
A single day's range thus carries more than four days of close-to-close information: ``the high and low have a lot
of information in them'' \ts{24}{19}{26:55}. Note that the estimator is positive whenever the range is positive, even on a day with a zero
close-to-close return.

\begin{coursenote}[Numerical discrepancy in the slides: the efficiency of Parkinson's estimator]\label{ch13:cn-park}
Slide~10 of \file{lec17_1} (and slide~13 of the 2013 deck) states $\mathrm{eff}(\hat\sigma_3^2)\approx5.2$. The exact value from
the moments above is $2/0.4073=4.91$, and a simulation of $2\times10^5$ Brownian paths (Brownian-bridge extremes, $400$ steps) gives $\Var\hat\sigma_3^2=0.4074$ and
efficiency $4.91$. The lecture slides also state $a\approx0.17$ as the optimal weight for the combination
$\hat\sigma_4^2=a\hat\sigma_1^2+(1-a)\hat\sigma_3^2$ of the gap estimator and the Parkinson estimator; by \cref{ch13:lem-comb} the optimal weight is $1/(1+\mathrm{eff}_3)=0.169$ with the exact value ($0.161$ with $5.2$), and the combination has
efficiency $1+4.91=5.91$, not $6.2$. So $a\approx0.17$ is consistent with the exact value and $6.2$ with the quoted $5.2$; the
statistical conclusions are unaffected, and the figures $7.4$ and $8.4$ below are reproduced exactly. This is my check, not a statement of the course.
\end{coursenote}

\begin{exercise}[PS5 (2024), Problem 3: the extreme-value estimator of Parkinson (1980)]\label{ch13:ex-ps5-p3}
Read M.~Parkinson, ``The extreme value method for estimating the variance of the rate of return'', \emph{J.\ Business} 53(1) (1980) 61--65. (b) Explain the definition of the estimator $D_x$ of Section~II. (c) Explain the definition of $D_l$ of Section~III. (d) Using the results of Section~IV,
explain why the improvement obtained with $D_l$ as an estimate of the variance $V$ (Section~V) could matter particularly in studies of the dependence of $V$ on time and on price, if there is any. (Part (a), ``read the article'', has no answer.) The estimator of \cref{ch13:sec-parkinson} is the one in question.
\end{exercise}

\subsection{Garman--Klass, and the ``best analytic'' estimator}\label{ch13:sec-gk}
Garman and Klass (1980) found the estimator of minimal variance among those that are \emph{scale invariant}, i.e.\ that
scale as $\lambda^2$ when the price is multiplied by $\lambda$, hence expressible in the normalised prices
$u=H-O$, $d=L-O$, $c=C-O$ (logarithms) \ts{24}{19}{27:57}\ts{13}{9}{56:47}:
\begin{equation}\label{ch13:eq-gkbest}
  \hat\sigma_{**}^2=0.511\,(u-d)^2-0.019\,\bigl\{c(u+d)-2ud\bigr\}-0.383\,c^2,\qquad \mathrm{eff}(\hat\sigma^2_{**})\approx7.4 .
\end{equation}
The derivation minimises a quadratic form over the coefficients of an analytic scale-invariant combination: ``some quadratic
forms and minimising them''. The coefficients are close to those of the simpler estimator used in practice
\begin{equation}\label{ch13:eq-gk}
  \hat\sigma^2_{\mathrm{GK}}=\tfrac12(u-d)^2-(2\ln2-1)\,c^2,\qquad 2\ln2-1=0.386,
\end{equation}
which is the one implemented in the case study (formula of \cref{ch13:sec-gkcs} below, with $\ln(H/L)=u-d$, $\ln(C/O)=c$).
It is unbiased because $\E\bigl[\tfrac12(H-L)^2\bigr]=2\ln2\,\sigma^2$ (from $\E[(H-L)^2]=4\ln2\,\sigma^2$) and $\E[c^2]=\sigma^2$: the two terms give
$(2\ln2)-(2\ln2-1)=1$ times $\sigma^2$. This explains where the mysterious constant $2\ln2-1$ comes from: it is the coefficient that cancels the
bias of the range term. Simulation: efficiency $7.40$ for both \eqref{ch13:eq-gkbest} and \eqref{ch13:eq-gk}.
Why does one bother with efficiency? With an efficient estimator one can use a shorter window and still be precise, which allows
tracking changes of volatility with less lag \ts{24}{19}{30:03}\ts{13}{9}{58:56}.

For $0<f<1$ the opening price may differ from the previous close, and the \emph{composite} estimator combines the overnight gap with the
best analytic estimator of the session, normalised by its length $1-f$:
\begin{equation}\label{ch13:eq-gkcomp}
  \hat\sigma^2_{G K}=a\,\frac{(O_1-C_0)^2}{f}+(1-a)\,\frac{\hat\sigma^2_{**}}{1-f},\qquad a\approx0.12,\qquad
  \mathrm{eff}\approx8.4 .
\end{equation}
Both terms are unbiased for $\sigma^2$ and independent (disjoint increments), with efficiencies $1$ and $7.4$, so
\cref{ch13:lem-comb} gives the optimal weight $a=1/(1+7.4)=0.119\approx0.12$ and the efficiency $1+7.4=8.4$ \ts{24}{19}{31:13}\ts{13}{9}{58:56}, exactly the values of
the slides.
\begin{secondary}[Gap plus range]
(In the same way the composite of the gap and the session range, whose second term is
$(H-L)^2/(4\ln2\,(1-f))$, has $a=0.17$ and efficiency $5.9$.)
\end{secondary}

\begin{exercise}[(not from the course): unbiasedness of a range estimator]\label{ch13:ex-extra2}
Using $\E[(H-L)^2]=4\ln2\,\sigma^2$ and $\E[c^2]=\sigma^2$, verify that the Garman--Klass estimator \eqref{ch13:eq-gk} is unbiased, and that with a drift $\mu_*\ne0$ the close-to-close estimator $c^2$ has expectation $\sigma^2+\mu_*^2$ (units of a day). Why does this not contradict the estimator being unbiased for a driftless process?
\end{exercise}

\subsection{Rogers--Satchell and Yang--Zhang}\label{ch13:sec-rsyz}
\begin{secondary}[Estimators robust to drift and opening jumps]
All estimators so far assume no drift, and the composite one assumes that the only jump is the opening gap. Two later refinements are
listed in the slides as extensions and used in the case study.
\begin{itemize}
  \item \emph{Rogers--Satchell} (1991) allow a non-zero drift (but no opening jump):
        \begin{equation}\label{ch13:eq-rs}
          \hat\sigma^2_{\mathrm{RS}}=u(u-c)+d(d-c)=\ln\frac HC\ln\frac HO+\ln\frac LC\ln\frac LO .
        \end{equation}
        It is unbiased for every drift. (Simulation, with a drift of $1.5\sigma$ per day, far larger than any real one: mean of
        \eqref{ch13:eq-rs} $=1.002\,\sigma^2$, whereas close-to-close gives $3.26\,\sigma^2$, Parkinson $1.85\,\sigma^2$ and \eqref{ch13:eq-gk} $1.30\,\sigma^2$.) At zero drift its
        efficiency is $5.97$, below that of Garman--Klass: robustness against drift is paid for with variance.
  \item \emph{Yang--Zhang} (2000) allow both a drift and opening gaps. Over $n$ days,
        \begin{equation}\label{ch13:eq-yz}
          \hat\sigma^2_{\mathrm{YZ}}=\hat\sigma^2_{o}+k\,\hat\sigma^2_{c}+(1-k)\,\hat\sigma^2_{\mathrm{RS}},\qquad
          k=\frac{0.34}{1.34+\frac{n+1}{n-1}},
        \end{equation}
        where $\hat\sigma^2_o$ is the sample variance of the overnight log returns $\ln(O_i/C_{i-1})$, $\hat\sigma_c^2$ the sample variance of the
        open-to-close returns $\ln(C_i/O_i)$ and $\hat\sigma^2_{\mathrm{RS}}$ the mean of \eqref{ch13:eq-rs} over the $n$ days
        \ts{24}{19}{31:13}. It is a weighted average of the overnight, open-to-close and Rogers--Satchell variances with the weights that minimise the estimation error,
        and, because the first two are \emph{sample} variances, it does not depend on the drift. For $n=21$, $k=0.139$. It is the
        version recommended ``these days'' \ts{24}{19}{31:13}.
\end{itemize}
\end{secondary}

\begin{coursenote}[Names and efficiencies in the case study]\label{ch13:cn-names}
The R function \texttt{volatility()} of the package \texttt{TTR} (through \texttt{tidyquant}) offers six estimators:
\texttt{close}, \texttt{parkinson}, \texttt{garman.klass}, \texttt{rogers.satchell}, \texttt{gk.yz} and \texttt{yang.zhang}. The one labelled
``Garman--Klass Yang and Zhang'' (\texttt{gk.yz}) is \eqref{ch13:eq-gk} plus the squared overnight return
$\ln^2(O_i/C_{i-1})$, i.e.\ the composite \eqref{ch13:eq-gkcomp} with unit weights and no normalisation by $f$ and $1-f$; it is
unbiased if the overnight variance is proportional to $f$. The case study writes that it ``has an efficiency of 8 times''
the close-to-close estimator. Under the model of \cref{ch13:sec-eff} its efficiency is $2/\bigl(2f^2+(1-f)^2\,2/7.4\bigr)$: it is $8.4$ at its best, for $f\approx0.11$, $7.4$ for $f=0$,
$3.5$ for $f=0.5$ and only $1.8$ for an index that is closed $17.5$ hours out of $24$ ($f=0.73$) if overnight variance were proportional to closed time (which is not so: the overnight variance of equity indices is a much smaller share, and this is why Yang--Zhang re-weight; not discussed in class). Parameter names: TTR's \texttt{alpha}=$1.34$ enters as
$k=(\alpha-1)/(\alpha+(n+1)/(n-1))$, identical to \eqref{ch13:eq-yz}.
\end{coursenote}

\subsection{Checking the efficiencies by simulation}
\begin{secondary}[Monte Carlo check]
The claims of the previous subsections can be verified with a Monte Carlo experiment: $2\times10^5$ Brownian paths on $[0,1]$ ($\sigma=1$, $f=0$), with the high and low of each of $400$ steps drawn exactly from the Brownian bridge so that discretisation does not bias the range.
\begin{center}\small
\begin{tabular}{@{}lcccc@{}}
\toprule
estimator & mean, $\mu_*=0$ & efficiency, $\mu_*=0$ & mean, $\mu_*=1.5\sigma$ & \\
\midrule
close-to-close $c^2$ & $1.00$ & $1.00$ & $3.26$ & \\
Parkinson \eqref{ch13:eq-parkinson} & $1.00$ & $4.91$ & $1.85$ & \\
Rogers--Satchell \eqref{ch13:eq-rs} & $1.00$ & $5.97$ & $1.00$ & \\
Garman--Klass \eqref{ch13:eq-gk} & $1.00$ & $7.40$ & $1.30$ & \\
best analytic \eqref{ch13:eq-gkbest} & $1.00$ & $7.40$ & $1.29$ & \\
\bottomrule
\end{tabular}
\end{center}
With real data the high and low are recorded from discrete trades and quotes, so range estimators are biased downwards; this is
not discussed in the course.
\end{secondary}

\subsection{Case study: volatility of the S\&P~500 and of ARKK (2024)}\label{ch13:sec-gkcs}
\begin{casestudy}[Historical volatility with R: S\&P~500 (lecture) and ARKK (problem set)]\label{ch13:cs-ohlc}
\textbf{Goal.} Compare the six estimators on real bars and ask which is more \emph{predictable}
\ts{24}{19}{32:16}.
\textbf{Lesson.} The efficient estimator (Yang--Zhang) is also the more predictable, the estimates of the same asset differ noticeably with the
estimator, and volatility levels differ across assets by an order of magnitude; the conclusions above are only descriptive.
\end{casestudy}
\begin{secondary}[The S\&P~500 and ARKK case study in detail]
\textbf{Data.} Daily OHLC bars from Yahoo Finance through \texttt{tidyquant::tq\_get}:
the S\&P~500 index from 2014-01-02 to 2024-11-04 (2\,729 days) and the ARKK exchange-traded fund (Cathie Wood) from 2014-10-31 to 2024-10-30
(2\,516 days) \ts{24}{19}{34:26}. \textbf{Method.} Rolling windows of $n=21$ days and $N=252$, with the formulas
\begin{gather*}
  \hat\sigma_{\mathrm{cc}}=\sqrt{\tfrac N{n-2}\textstyle\sum_{i=1}^{n-1}(r_i-\bar r)^2},\ \ r_i=\ln\tfrac{C_{i+1}}{C_i},\\
  \hat\sigma_{\mathrm{P}}=\sqrt{\tfrac N{4n\ln2}\textstyle\sum_{i=1}^n\ln^2\tfrac{H_i}{L_i}},\\
  \hat\sigma_{\mathrm{GK}}=\sqrt{\tfrac Nn\textstyle\sum_{i=1}^n\bigl[\tfrac12\ln^2\tfrac{H_i}{L_i}-(2\ln2-1)\ln^2\tfrac{C_i}{O_i}\bigr]},\\
  \hat\sigma_{\mathrm{RS}}=\sqrt{\tfrac Nn\textstyle\sum_{i=1}^n\bigl[\ln\tfrac{H_i}{C_i}\ln\tfrac{H_i}{O_i}+\ln\tfrac{L_i}{C_i}\ln\tfrac{L_i}{O_i}\bigr]},
\end{gather*}
plus \texttt{gk.yz} and Yang--Zhang \eqref{ch13:eq-yz}; \emph{close-to-close} uses the unbiased variance of the $n-1$ returns in a window of $n$ prices.
\textbf{Results (S\&P~500).} All estimators show the same pattern: annualised volatility mostly between $5\%$ and $25\%$, around $12\%$--$15\%$ in
calm periods, with spikes and a huge one in the COVID crash of spring 2020 (peak close-to-close $0.98$, Parkinson $0.53$, GK $0.52$, Yang--Zhang $0.70$)
\ts{24}{19}{34:26}. The range estimators are smoother and different in level: over 2023--24 the Yang--Zhang volatility
is generally below the close-to-close one, though sometimes far above it (scatter plot of the pairs) \ts{24}{19}{36:32}. A close-to-close
window keeps a spike until it drops out $21$ days later, so it stays elevated longer than the others; estimators that fall faster after a spike
are probably better for predicting realised volatility \ts{24}{19}{38:43}.

\medskip
\textbf{Modelling the rolling volatility, $\approx450$ days of 2023--24} (\texttt{forecast}/\texttt{fpp2} package of Hyndman and Athanasopoulos, \emph{Forecasting: Principles and Practice}; Chapter~3 for the naive method and residual diagnostics, Chapter~8 for automatic ARIMA and seasonal ARIMA) \ts{24}{19}{40:47}:
\begin{center}\small
\begin{tabular}{@{}lcccc@{}}
\toprule
S\&P~500, 2023--24 & model & Ljung--Box $Q^*$ (df, $p$) & MAPE\\
\midrule
close-to-close & naive & $31.2$ (10, $5\times10^{-4}$) & $3.67$\\
 & ARIMA$(3,1,0)$ & $16.2$ (7, $0.023$) & $3.74$\\
 & seasonal, $[21]$ & $89.1$ (34, $8\times10^{-7}$) & (meaningless)\\
Yang--Zhang & naive & $88.1$ (10, $10^{-14}$) & $2.31$\\
 & ARIMA$(1,1,1)$ & $9.5$ (8, $0.31$) & $2.23$\\
 & ARIMA$(2,1,3)(2,0,1)_{[21]}$ & $31.8$ (34, $0.58$) & $1.92$\\
\bottomrule
\end{tabular}
\end{center}
The \emph{naive} model $\hat y_t=y_{t-1}$ is a random walk (efficient-market flavour: the best forecast is where things are now); its residuals are the first
differences and the Ljung--Box test (\cref{ch09:eq-bp}) rejects whiteness \ts{24}{19}{41:49}. \texttt{auto.arima}
picks a model by the corrected AIC, $\mathrm{AICc}=\mathrm{AIC}+2k(k+1)/(n-k-1)$ with $k$ parameters (for ARIMA$(3,1,0)$: $-3303.73$ against $-3303.64$, $n\approx450$)
\ts{24}{19}{45:02}; the coefficients are judged by their ratio to the standard error: for
close-to-close ARIMA$(3,1,0)$ the AR(3) coefficient has $0.126/0.047\approx2.7$ (``almost 3''), the AR(2) $1.4$
\ts{24}{19}{47:12}. For Yang--Zhang the fit is an ARIMA$(1,1,1)$ with $\hat\phi=0.696$ and $\hat\theta=-0.446$, many standard errors from zero, and the
residuals are close to white noise, except for a large autocorrelation at \emph{lag $21$}, the length of the window
\ts{24}{19}{48:14} (explained in \cref{ch13:prop-lag21}). A \emph{seasonal} ARIMA with period $21$ removes it
\ts{24}{19}{50:17}: with a seasonal MA coefficient of $-0.78$ ($13$ standard errors) the Ljung--Box test no longer rejects, and the training MAPE falls to $1.92$
\ts{24}{19}{57:58}. The forecast plot shows a point forecast (solid line) and Gaussian prediction intervals that widen with the horizon, with the volatility
reverting to a mean level after a low period \ts{24}{19}{59:01}.

\textbf{ARKK (problem set).} The same code gives a very different picture: close-to-close volatility of ARKK ranges from $10\%$ to $120\%$, with peaks of
about $1.2$ in 2020 and $1.0$ in 2022, three to four times the index. Results: close-to-close naive $Q^*=35.3$ (df~10), MAPE $3.50$; ARIMA$(1,1,1)$ with
$\hat\phi=0.876$, $\hat\theta=-0.799$, MAPE $3.55$; Yang--Zhang naive $Q^*=13.5$ ($p=0.19$, already white), MAPE $2.01$; ARIMA$(1,1,1)$
$\hat\phi=0.780$, $\hat\theta=-0.697$, MAPE $1.99$; seasonal ARIMA$(0,1,1)(2,0,2)_{[21]}$ with drift, $Q^*=37.5$ ($p=0.44$), MAPE $1.78$. For the
close-to-close series \texttt{auto.arima} on the differences selects no seasonal term at all and the Ljung--Box test still rejects
($Q^*=124.6$, df~40).
The scatter of YZ against close-to-close for ARKK lies around the diagonal, with YZ below it at high volatility.
\end{secondary}

\begin{coursenote}[Erratum: stale numbers in a note of the slides]\label{ch13:cn-mape}
Page~31 of \file{lec17_2} says ``the MAPE of the seasonal ARIMA model of the YZ volatility is $1.746$ versus $3.58$ for the naive
model of the close volatility'', but the printed accuracy tables on pages~15 and 30 of the same file give $1.917$ and $3.666$. The note is
hard-coded text left over from another run; in the knitted S\&P~500 PDF of the problem-set zip, which is a re-run with 2\,726 days, the note is
generated from the output ($1.905$ against $3.626$). The values quoted above are those of the lecture PDF.
\end{coursenote}

\begin{secondary}[Why lag 21: the window itself]
\begin{proposition}[Why the residuals of a rolling volatility have an autocorrelation at the window length]\label{ch13:prop-lag21}
Let $x_t$ be i.i.d.\ with variance $s^2$ and let $v_t=\frac1m\sum_{j=0}^{m-1}x_{t-j}$ be the rolling mean of window $m$. Then the first differences (the
residuals of the naive model) satisfy $\Delta v_t=(x_t-x_{t-m})/m$, so
$\rho_{\Delta v}(1)=\dots=\rho_{\Delta v}(m-1)=0$ and $\rho_{\Delta v}(m)=-\tfrac12$.
\end{proposition}
\begin{proof}
$\Delta v_t=v_t-v_{t-1}=(x_t-x_{t-m})/m$: all terms but two cancel. The variance is $2s^2/m^2$; for lags $1\le h<m$ the two
differences $(x_t-x_{t-m})$ and $(x_{t+h}-x_{t+h-m})$ share no term; for $h=m$ they share $x_t$ with opposite signs, covariance $-s^2/m^2$.
\end{proof}
This is what the lecture describes as observations \emph{entering and leaving the window} \ts{24}{19}{49:15}; a simulation with $4\times10^5$
i.i.d.\ chi-squared terms gives $\rho(21)=-0.501$. It also shows that the good performance of the seasonal model is partly mechanical: the term $x_{t-m}$ leaving the window is known
from the past, so a model with a seasonal moving-average term can predict it, and the variance of the one-step error falls to that of $x_t/m$,
half of $\Var(\Delta v_t)$; the root-mean-square error should drop by the factor $1/\sqrt2=0.71$, and it drops from $0.0051$ to $0.0039$ for Yang--Zhang. Part of the
``predictability of volatility'' in the case study is the predictability of the estimator's own window, not of the market (our remark).
\end{secondary}

\begin{exercise}[PS5 (2024), Problem 4: volatility case studies]\label{ch13:ex-ps5-p4}
The R Markdown files
\begin{itemize}\item \file{Case_Study_Estimating__SP500_Historical_Volatility.Rmd}, \item \file{Case_Study_Estimating__ARKK_Historical_Volatility.Rmd}\end{itemize}
produce the S\&P~500 and ARKK case studies (\cref{ch13:cs-ohlc}), which use the naive, ARIMA and seasonal ARIMA methods of Hyndman and Athanasopoulos,
\emph{Forecasting: Principles and Practice} (2nd ed.), Chapter~3 (naive method, residual diagnostics) and Chapter~8 (automatic ARIMA).
\begin{enumerate}[(a)]
  \item Comment on the S\&P~500 results you find interesting or significant.
  \item The same for ARKK.
  \item Compare the volatility findings for ARKK and the S\&P~500.
\end{enumerate}
\end{exercise}

\begin{exercise}[PS5 (2024), Problem 5, optional (extra credit): your own stock]\label{ch13:ex-ps5-p5}
Choose a stock symbol and edit the ARKK file to run the same analysis.
\begin{enumerate}[(a)]
  \item Produce the PDF of the case study.
  \item Comment on results you find interesting.
  \item Compare your stock with the S\&P~500 and ARKK.
\end{enumerate}
\end{exercise}


% ==================================================================
\section{Fat tails from mixtures: jumps and the Laplace law}\label{ch13:sec-mix}
Geometric Brownian motion has returns that are exactly normal, in contrast with the histograms of \cref{ch13:sec-clock}.
Two extensions keep independent returns but give them heavier tails; both are \emph{normal variance mixtures} in the sense of
\cref{ch03:prop-mixture}, which will be used repeatedly: if $X\mid V\sim N(\mu,\sigma^2V)$, the kurtosis is $3\,\E[V^2]/(\E[V])^2\ge3$.

\subsection{Poisson jump diffusion}\label{ch13:sec-pjd}
The \emph{Poisson jump diffusion} adds to \eqref{ch13:eq-gbm} a compound Poisson process \ts{24}{19}{1:00:08}\ts{13}{9}{58:56}:
\begin{equation}\label{ch13:eq-pjd}
  \frac{\dd S(t)}{S(t)}=\mu\dd t+\sigma\dd W(t)+\gamma\sigma Z(t)\dd\pi(t),
\end{equation}
where $\pi$ is a Poisson process of rate $\lambda$ (the jumps arrive at random, unpredictable times), the marks $Z(t)$ are i.i.d.\ $N(0,1)$ and
$\gamma$ sets the scale of the jumps in units of the diffusion volatility. For the distribution of the return we treat the jump as acting additively on the log-price, the convention of the
lecture (the exact jump of $\log S$ would be $\log(1+\gamma\sigma Z)$). Over a period $\Delta$ the number $N$ of jumps is Poisson$(\lambda\Delta)$, and given $N=k$ the centred return is the sum of the
Gaussian diffusion increment and $k$ independent Gaussian jumps:
\[
  R-m\mid N=k\ \sim\ N\bigl(0,\ \sigma^2(\Delta+k\gamma^2)\bigr).
\]
So the return has a \emph{Poisson mixture of Gaussian} density $\sum_{k\ge0}e^{-\lambda\Delta}\frac{(\lambda\Delta)^k}{k!}\,\varphi_{\sigma^2(\Delta+k\gamma^2)}(x-m)$ \ts{24}{19}{1:02:28}\ts{13}{9}{1:00:00}.

\begin{proposition}[Variance and kurtosis of the jump diffusion]\label{ch13:prop-pjd}
The variance of $R$ is $\sigma^2\Delta(1+\lambda\gamma^2)$ and its kurtosis is
\[
  \kappa=3+\frac{3\gamma^4\lambda}{\Delta\,(1+\lambda\gamma^2)^2}\ >3 .
\]
\end{proposition}
\begin{proof}
Apply \cref{ch03:prop-mixture} with $V=\Delta+N\gamma^2$: $\E[V]=\Delta(1+\lambda\gamma^2)$ and, since $\Var N=\lambda\Delta$, $\Var V=\gamma^4\lambda\Delta$. Then
$\kappa=3\bigl(1+\Var V/(\E[V])^2\bigr)$.
\end{proof}
For $\lambda=0.3$ per day, $\gamma=3$ and $\Delta=1$ day, $\kappa=8.3$ (simulation: $8.32$). The excess kurtosis decays as $1/\Delta$: on longer horizons many jumps are added
and the central limit theorem restores normality, in accordance with the empirical observation that monthly returns are closer to Gaussian than daily ones.
\begin{secondary}[Estimation by the EM algorithm]
Estimation by maximum likelihood is hard because the likelihood is a product of infinite series, but the model has a natural latent variable, the number of jumps in each period; the
\emph{EM algorithm} of Dempster, Laird and Rubin then has closed-form steps and returns, as a by-product, the posterior number of jumps per period
\ts{24}{19}{1:01:19}\ts{13}{9}{1:04:18} (Pickard, Kempthorne and Zakaria, 1987).
\end{secondary}
A finding reported in class: when fitted to stock returns, the jump rate came out very high, a jump every
couple of days, so the model was not detecting rare crashes but mimicking a fat-tailed distribution by a mixture of a plain normal, a normal with one jump, with two jumps, and so on
\ts{24}{19}{1:02:28}.

\subsection{Brownian motion observed at random times: the Laplace law}\label{ch13:sec-laplace}
If a Brownian motion is observed not at equally spaced times but after random, \emph{exponentially distributed} increments of time, the increments of the process follow the
\emph{Laplace} (two-sided exponential) law \ts{24}{19}{1:03:32}
\begin{equation}\label{ch13:eq-laplace}
  f(x\mid\mu,b)=\frac1{2b}\,e^{-|x-\mu|/b},\qquad \E[X]=\mu,\quad \Var X=2b^2,\quad \kappa=6 .
\end{equation}
\begin{history}[Laplace's first law of errors]
Laplace himself introduced it in 1774 as the ``first law of errors''.

\textbf{References.} P.-S.~Laplace, ``Mémoire sur la probabilité des causes par les événements'', \emph{Mémoires de l'Académie royale des sciences présentés par divers savans} 6 (1774) 621--656; S.~M.~Stigler, \emph{The History of Statistics}, Harvard University Press, 1986, ch.~3.
\end{history}

\begin{proposition}\label{ch13:prop-laplace}
Let $\tau$ be exponential with mean $s$ and $X=\mu+\sigma B_\tau$ with $B$ a standard Brownian motion independent of $\tau$. Then $X\sim\mathrm{Laplace}(\mu,b)$ with $b^2=\sigma^2s/2$.
\end{proposition}
\begin{proof}
Given $\tau$, $X\sim N(\mu,\sigma^2\tau)$, so the characteristic function is $\E[e^{iu(X-\mu)}]=\E[e^{-\sigma^2u^2\tau/2}]=\bigl(1+\sigma^2su^2/2\bigr)^{-1}$
(Laplace transform of the exponential law). The Laplace density has characteristic function $e^{iu\mu}/(1+b^2u^2)$; the two agree for $b^2=\sigma^2s/2$. As a mixture
$V=\tau$ has $\E[V^2]/(\E[V])^2=2$, whence $\kappa=6$ (\cref{ch03:prop-mixture}) and $\Var X=\sigma^2s=2b^2$.
\end{proof}
It is a mixture with an exponentially distributed clock, the simplest ``subordinated'' process (\cref{ch13:int-clock}), and it has
$P(|X-\mu|>4\sqrt{\Var X})=e^{-4\sqrt2}=0.35\%$ against $6.3\times10^{-5}$ for the normal.

\begin{casestudy}[Laplace against normal on S\&P~500 returns]\label{ch13:cs-laplace}
The lecture fits both laws to $n=1000$ daily S\&P~500 returns, 2014-10-10 to 2018-10-02 \ts{24}{19}{1:05:34}. Method of moments for the Laplace law:
$\hat\mu=\bar R=0.000427$, $\hat b=\sqrt{\hat\Var/2}=0.005652$; the implied annualised volatility is $\sqrt{252}\sqrt2\,\hat b=12.69\%$ (a check: $\sqrt2\times0.005652\times\sqrt{252}=0.1269$). The
histogram is peaked with exponentially decaying tails and the Laplace curve follows it much better than the Gaussian (drawn in red for comparison). The \emph{Q--Q plot} (sorted observations against theoretical quantiles) is
close to a straight line for the Laplace law and clearly curved for the normal \ts{24}{19}{1:06:37}. So the Laplace law is a better model of independent daily returns;
what neither can explain is that returns are not independent.
\end{casestudy}

% ==================================================================
\section{ARCH models}\label{ch13:sec-arch}
The models so far have independent returns and hence a volatility that is the same in every period. The data say otherwise (\cref{ch13:sec-facts}): large
returns tend to follow large returns. Engle (1982) proposed to let the volatility of period $t$ depend on the returns of the
preceding periods \ts{24}{19}{1:07:45}\ts{13}{9}{1:05:24}.

\begin{definition}[ARCH$(p)$]\label{ch13:def-arch}
Let $y_t=\log(S_t/S_{t-1})$ and $\mathcal F_t$ the information available at time $t$. The \emph{autoregressive conditionally heteroskedastic} model of order $p$ is
\begin{equation}\label{ch13:eq-arch}
  y_t=\mu_t+\varepsilon_t,\qquad \varepsilon_t=\sigma_tz_t,\qquad
  \sigma_t^2=\alpha_0+\alpha_1\varepsilon_{t-1}^2+\dots+\alpha_p\varepsilon_{t-p}^2,
\end{equation}
where $\mu_t=\E[y_t\mid\mathcal F_{t-1}]$ is the conditional mean, $\{z_t\}$ are i.i.d.\ with $\E[z_t]=0$, $\Var z_t=1$, independent of $\mathcal F_{t-1}$, and
$\alpha_0>0$, $\alpha_i\ge0$. Then $\sigma_t^2=\Var(y_t\mid\mathcal F_{t-1})$ is the \emph{conditional variance}.
\end{definition}
The name says it all: \emph{conditional heteroskedasticity} means that the variance given the past changes over time; \emph{autoregressive} refers to its dependence on lagged squares. If $\varepsilon_{t-1}$ was
large, the variance of $\varepsilon_t$ is large and $\varepsilon_t$ is likely to be large too. The constraints have a reason: if some $\alpha_i$ were negative, a large past shock could make $\sigma_t^2$ negative
\ts{13}{9}{1:08:30}\ts{24}{19}{1:12:05}.
\begin{secondary}[Why $\alpha_0>0$]
(The slides write $\alpha_j\ge0$ for $j=0,\dots,p$, i.e.\ allow $\alpha_0=0$, in which case the variance can shrink to zero and the model degenerates; we take $\alpha_0>0$.)
\end{secondary}

\begin{proposition}[ARCH innovations are uncorrelated but not independent]\label{ch13:prop-mds}
Let $\varepsilon_t=\sigma_tz_t$ with $\sigma_t>0$ $\mathcal F_{t-1}$-measurable, $\E[\sigma_t^2]<\infty$, and $z_t$ as in \cref{ch13:def-arch}. Then
\begin{enumerate}[(i)]
  \item $\E[\varepsilon_t\mid\mathcal F_{t-1}]=0$ and $\E[\varepsilon_t^2\mid\mathcal F_{t-1}]=\sigma_t^2$;
  \item $\Cov(\varepsilon_t,\varepsilon_s)=0$ for $t\ne s$: the innovations form a martingale difference sequence and are uncorrelated;
  \item for ARCH$(1)$ with $\E[\varepsilon^4]<\infty$, $\Cov(\varepsilon_t^2,\varepsilon_{t-1}^2)=\alpha_1\Var\varepsilon_t^2>0$: the squares are correlated, so the $\varepsilon_t$ are not independent.
\end{enumerate}
\end{proposition}
\begin{proof}
(i) $z_t$ is independent of $\sigma_t$: $\E[\sigma_tz_t\mid\mathcal F_{t-1}]=\sigma_t\E[z_t]=0$, $\E[\sigma_t^2z_t^2\mid\mathcal F_{t-1}]=\sigma_t^2$.

(ii) For $s<t$, $\varepsilon_s$ is $\mathcal F_{t-1}$-measurable, so $\E[\varepsilon_t\varepsilon_s]=\E\bigl[\varepsilon_s\E[\varepsilon_t\mid\mathcal F_{t-1}]\bigr]=0$.

(iii) $\E[\varepsilon_t^2\varepsilon_{t-1}^2]=\E[\sigma_t^2\varepsilon_{t-1}^2]=\alpha_0\E[\varepsilon^2]+\alpha_1\E[\varepsilon^4]$, while $\E[\varepsilon_t^2]\E[\varepsilon^2_{t-1}]=(\E[\varepsilon^2])(\alpha_0+\alpha_1\E[\varepsilon^2])$; subtract.
\end{proof}
This is the precise sense of the remark of \cref{ch09:sec-sacf} that residual whiteness is not the end of the story: the ACF of the returns is flat (the martingale
difference), that of the squares is not \ts{24}{19}{1:09:55}. For the S\&P~500 the autocorrelations of the squared, de-meaned log returns
are significantly positive at (about) the first $10$ lags \ts{24}{19}{1:09:55}; for the euro/dollar rate the ACF and PACF of the squares are large and of the returns marginal
(\cref{ch13:cs-fx}) \ts{13}{9}{1:18:02}.

\subsection{ARCH is an autoregression for the squared returns}\label{ch13:sec-arch-ar}
Put $u_t=\varepsilon_t^2-\sigma_t^2$ and add it to both sides of \eqref{ch13:eq-arch}:
\begin{equation}\label{ch13:eq-archar}
  \varepsilon_t^2=\alpha_0+\alpha_1\varepsilon_{t-1}^2+\dots+\alpha_p\varepsilon_{t-p}^2+u_t,\qquad
  \E[u_t\mid\mathcal F_{t-1}]=0,\quad \Var[u_t\mid\mathcal F_{t-1}]=(\kappa_z-1)\sigma_t^4 ,
\end{equation}
where $\kappa_z=\E[z^4]$ (equal to $3$ for Gaussian $z$, so that the conditional variance of $u_t$ is $2\sigma_t^4$)
\ts{24}{19}{1:09:55}\ts{13}{9}{1:09:31}. The squared innovations follow an AR$(p)$ whose noise $u_t$ has mean zero and a variance that
changes with time. This yields a quick test for ARCH effects.

\begin{coursenote}[Erratum: the noise of the AR equation is not white noise with variance $2\sigma^4$]\label{ch13:cn-u}
Slides~21 and~25 of \file{lec17_1} (slides~20 and~24 of 2013) write $\Var[u_t\mid\mathcal F_t]=\Var(\varepsilon_t^2)=2\sigma_t^4$, and later
$u_t\sim\mathrm{WN}(0,2\sigma^4)$. Two corrections. First, the conditional variance is taken given $\mathcal F_{t-1}$ (given $\mathcal F_t$ the variable $u_t$ is known), and it equals
$(\kappa_z-1)\sigma_t^4$, i.e.\ $2\sigma_t^4$ only for Gaussian $z_t$. Second, $\sigma_t^2$ is random, so $u_t$ is not white noise in the strong sense: it is a martingale difference
(hence uncorrelated), and it has a constant \emph{unconditional} variance $(\kappa_z-1)\E[\sigma_t^4]$ only if $\E[\varepsilon^4]<\infty$ (\cref{ch13:thm-kurt}). The ARMA representation of GARCH below is therefore an
ARMA with uncorrelated, not independent, innovations; the second-moment formulas of \cref{ch09:sec-arma11} apply verbatim.
\end{coursenote}

\paragraph{Lagrange-multiplier test for ARCH effects.} To test $H_0:\alpha_1=\dots=\alpha_p=0$, fit the mean model, take the residuals $\hat\varepsilon_t=y_t-\hat\mu_t$, regress
$\hat\varepsilon_t^2$ on a constant and $\hat\varepsilon^2_{t-1},\dots,\hat\varepsilon^2_{t-p}$ (an AR$(p)$ fit to the squared residuals), and compute
\begin{equation}\label{ch13:eq-lm}
  \mathrm{LM}=nR^2\ \overset{H_0}{\approx}\ \chi^2_p ,
\end{equation}
{\emergencystretch=3em with $R^2$ the coefficient of determination of that regression \ts{24}{19}{1:09:55}\ts{13}{9}{1:10:34}. Under $H_0$ the squares $\hat\varepsilon_t^2$ have
no linear dependence on the past and $R^2$ is small; for $p=1$, $R^2$ is the squared sample correlation between $\hat\varepsilon_t^2$ and $\hat\varepsilon^2_{t-1}$, so
$nR^2\approx n\hat r_1^2$ with $\sqrt n\hat r_1\to N(0,1)$, as in \cref{ch09:eq-bp}.
\begin{secondary}[McLeod--Li and quasi-maximum likelihood]
For general $p$ the statistic is asymptotically equivalent to the
Box--Pierce statistic of the squared residuals (the McLeod--Li test; this remark is not from the class). The regression uses least squares, whose assumptions are not those of a Gaussian likelihood; its estimates are
\emph{quasi-maximum-likelihood} estimates, consistent but not efficient \ts{13}{9}{1:11:40}.
\end{secondary}\par}

\begin{finance}[Volatility clustering and prediction]
An ARCH model predicts nothing about the \emph{sign} of tomorrow's return but a lot about its size, and size is what matters for risk: the conditional variance $\sigma_t^2$ that it produces is a one-step-ahead
volatility forecast that reacts to the latest shock. This is the improvement over the historical averages of \cref{ch13:sec-ewma}, whose weights are chosen by hand
instead of being estimated.
\end{finance}

\begin{exercise}[PS5 (2013), Problem 3: properties of the ARCH$(1)$ model]\label{ch13:ex-ps5-3}
Let $y_t=\mu_t+\varepsilon_t$, $\varepsilon_t=z_t\sigma_t$, $\sigma_t^2=\alpha_0+\alpha_1\varepsilon_{t-1}^2$, with $z_t$ i.i.d., $\E[z_t]=0$, $\Var z_t=1$, $\E[z_t^3]=0$, $\E[z_t^4]=\kappa_z$ ($=3$ if Gaussian). Prove:
\begin{enumerate}[(a)]
  \item $\E[\varepsilon_t^2]=\alpha_0/(1-\alpha_1)$;
  \item $\E[\varepsilon_t^3]=0$;
  \item $\E[\varepsilon_t^4]=\kappa_z\alpha_0^2(1+\alpha_1)/[(1-\alpha_1)(1-\kappa_z\alpha_1^2)]$;
  \item which constraints on $\alpha_0,\alpha_1$ does (c) require for the fourth moment to be finite and stationary;
  \item the kurtosis of $\varepsilon_t$: for Gaussian $z_t$, is the unconditional distribution of $\varepsilon_t$ heavier-tailed than the Gaussian?
      (The problem as printed writes the constant $\kappa_z$ in (c) without naming it; it is the kurtosis of $z_t$ stated just before.)
\end{enumerate}
\end{exercise}


% ==================================================================
\section{GARCH models}\label{ch13:sec-garch}

\subsection{Definition and parsimony}\label{ch13:sec-garch-def}
Data often need a large ARCH order to reproduce the persistence of volatility (a hint of the infinite-order autoregression of Wold's theorem, \cref{ch09:sec-wold}). Bollerslev (1986)
added lagged \emph{variances} to the right-hand side \ts{24}{19}{1:12:05}\ts{13}{9}{1:14:50}\ts{13}{11}{1:07}.

\begin{definition}[GARCH$(p,q)$]\label{ch13:def-garch}
With $y_t=\mu_t+\varepsilon_t$ and $\varepsilon_t=\sigma_tz_t$ as in \cref{ch13:def-arch},
\begin{equation}\label{ch13:eq-garch}
  \sigma_t^2=\alpha_0+\sum_{i=1}^p\alpha_i\varepsilon_{t-i}^2+\sum_{j=1}^q\beta_j\sigma_{t-j}^2,\qquad\alpha_0>0,\ \alpha_i\ge0,\ \beta_j\ge0 .
\end{equation}
The most used case is the \emph{GARCH$(1,1)$},
\begin{equation}\label{ch13:eq-garch11}
  \sigma_t^2=\alpha_0+\alpha_1\varepsilon_{t-1}^2+\beta_1\sigma_{t-1}^2 ,
\end{equation}
parsimonious (three parameters, plus those of the mean) and adequate for many financial series; in practice the $(1,1)$ order is often the one selected among the GARCH orders
\ts{24}{19}{1:13:09}\ts{13}{9}{1:16:58}.
\end{definition}

\begin{proposition}[GARCH$(1,1)$ is an ARCH$(\infty)$ with geometric weights]\label{ch13:prop-archinf}
If $\beta_1<1$ and the process is stationary (\cref{ch13:thm-stat}),
\[
  \sigma_t^2=\frac{\alpha_0}{1-\beta_1}+\alpha_1\sum_{j\ge0}\beta_1^j\,\varepsilon_{t-1-j}^2 .
\]
\end{proposition}
\begin{proof}
Iterate \eqref{ch13:eq-garch11}: $\sigma_t^2=\alpha_0\sum_{j<n}\beta_1^j+\alpha_1\sum_{j<n}\beta_1^j\varepsilon^2_{t-1-j}+\beta_1^n\sigma^2_{t-n}$; the last term tends to $0$ in expectation
because $\E[\sigma^2_{t-n}]$ is finite and constant and $\beta_1^n\to0$.
\end{proof}
So the conditional variance is an exponentially weighted average of past squared innovations plus a constant: the same structure as the EMA of
\cref{ch13:prop-ewma}, with the decay factor $\beta_1$ and the weight $\alpha_1/(1-\beta_1)$ on the average. Two parameters $(\alpha_1,\beta_1)$ replace the ten or more coefficients of a high-order ARCH,
which is the parsimony argument of the lecture: the Wold theorem promises that some infinite-order model fits, and GARCH gives a parametrisation of the infinite coefficient sequence by two numbers
\ts{13}{9}{1:20:10}.

\begin{figure}[ht]
\centering
\begin{tikzpicture}
\begin{axis}[width=0.96\linewidth,height=5.6cm,xmin=0,xmax=201,ymin=-4.6,ymax=4.6,
  axis lines=left,xlabel={day $t$},ylabel={return (\%)},
  legend style={font=\scriptsize,draw=none,fill=white,fill opacity=0.85,text opacity=1,at={(0.5,1.02)},anchor=south,legend columns=2},
  tick label style={font=\small},label style={font=\small}]
  \addplot[thick,defcol,ycomb,mark=none] coordinates {(1,-1.40) (2,0.10) (3,-1.09) (4,-1.23) (5,1.50) (6,0.87) (7,-1.05) (8,-2.02) (9,1.34) (10,1.38) (11,0.03) (12,-0.58) (13,1.29) (14,3.22) (15,2.38) (16,-0.03) (17,0.64) (18,-1.69) (19,-4.11) (20,0.55) (21,-0.34) (22,4.14) (23,0.34) (24,3.12) (25,2.82) (26,-2.99) (27,-0.86) (28,-0.54) (29,-3.04) (30,1.27) (31,0.10) (32,-1.81) (33,0.44) (34,-0.09) (35,0.80) (36,-0.48) (37,-0.33) (38,0.11) (39,0.59) (40,-0.05) (41,1.66) (42,0.50) (43,-0.06) (44,-0.63) (45,-2.40) (46,-1.21) (47,0.24) (48,0.34) (49,1.01) (50,-0.44) (51,1.60) (52,-2.83) (53,0.77) (54,-2.07) (55,2.06) (56,1.26) (57,-1.65) (58,-1.27) (59,1.54) (60,0.32) (61,-0.14) (62,-1.06) (63,0.61) (64,-0.56) (65,-0.53) (66,-1.07) (67,0.20) (68,-0.21) (69,-0.74) (70,-0.85) (71,-0.29) (72,-0.95) (73,1.17) (74,0.54) (75,-0.54) (76,1.55) (77,-1.41) (78,0.70) (79,0.27) (80,-0.52) (81,1.15) (82,0.63) (83,0.29) (84,1.92) (85,0.65) (86,0.64) (87,-0.10) (88,-0.67) (89,0.70) (90,-0.94) (91,-0.93) (92,-0.03) (93,-1.89) (94,0.19) (95,0.11) (96,0.43) (97,-1.17) (98,-0.76) (99,-1.42) (100,-1.46) (101,1.03) (102,0.09) (103,-1.09) (104,-1.42) (105,-0.55) (106,-0.83) (107,-0.42) (108,0.15) (109,1.00) (110,0.42) (111,1.14) (112,1.27) (113,0.34) (114,-1.50) (115,-2.04) (116,0.02) (117,1.98) (118,1.59) (119,1.35) (120,-0.02) (121,-0.63) (122,0.06) (123,-1.38) (124,0.25) (125,0.62) (126,-1.01) (127,1.72) (128,-0.77) (129,0.90) (130,3.13) (131,-0.07) (132,0.23) (133,1.04) (134,-1.47) (135,1.42) (136,0.51) (137,-0.22) (138,0.03) (139,-0.40) (140,-0.33) (141,0.51) (142,-0.22) (143,0.29) (144,0.40) (145,-0.62) (146,0.75) (147,-0.25) (148,0.91) (149,-0.34) (150,-0.21) (151,0.03) (152,0.29) (153,0.18) (154,0.41) (155,0.19) (156,-0.07) (157,-0.37) (158,0.13) (159,-0.58) (160,0.42) (161,-0.63) (162,-0.36) (163,0.62) (164,-0.63) (165,-0.23) (166,-0.58) (167,0.17) (168,-0.27) (169,0.39) (170,0.18) (171,0.06) (172,-0.10) (173,-0.09) (174,0.14) (175,-0.37) (176,0.24) (177,0.18) (178,-0.08) (179,-0.08) (180,0.15) (181,0.00) (182,-0.28) (183,0.61) (184,0.44) (185,-0.16) (186,0.77) (187,0.20) (188,0.24) (189,0.63) (190,-0.19) (191,0.24) (192,0.34) (193,-0.03) (194,0.28) (195,-0.70) (196,-0.09) (197,-0.46) (198,0.48) (199,0.33) (200,-1.34)};
  \addlegendentry{$\varepsilon_t$}
  \addplot[thick,thmcol,no marks] coordinates {(1,1.000) (2,1.056) (3,0.989) (4,1.003) (5,1.033) (6,1.099) (7,1.071) (8,1.066) (9,1.220) (10,1.229) (11,1.242) (12,1.158) (13,1.100) (14,1.121) (15,1.531) (16,1.643) (17,1.525) (18,1.434) (19,1.456) (20,1.965) (21,1.830) (22,1.700) (23,2.131) (24,1.976) (25,2.125) (26,2.197) (27,2.281) (28,2.131) (29,1.981) (30,2.116) (31,2.007) (32,1.859) (33,1.833) (34,1.705) (35,1.582) (36,1.495) (37,1.399) (38,1.306) (39,1.217) (40,1.154) (41,1.078) (42,1.162) (43,1.099) (44,1.028) (45,0.988) (46,1.245) (47,1.234) (48,1.154) (49,1.084) (50,1.073) (51,1.016) (52,1.101) (53,1.422) (54,1.349) (55,1.446) (56,1.522) (57,1.479) (58,1.489) (59,1.451) (60,1.451) (61,1.354) (62,1.261) (63,1.232) (64,1.168) (65,1.108) (66,1.052) (67,1.053) (68,0.989) (69,0.930) (70,0.912) (71,0.908) (72,0.861) (73,0.877) (74,0.920) (75,0.886) (76,0.856) (77,0.970) (78,1.033) (79,0.998) (80,0.941) (81,0.903) (82,0.939) (83,0.910) (84,0.863) (85,1.051) (86,1.010) (87,0.973) (88,0.914) (89,0.892) (90,0.875) (91,0.887) (92,0.895) (93,0.843) (94,1.030) (95,0.968) (96,0.910) (97,0.870) (98,0.914) (99,0.900) (100,0.981) (101,1.051) (102,1.047) (103,0.981) (104,0.996) (105,1.055) (106,1.007) (107,0.986) (108,0.937) (109,0.883) (110,0.901) (111,0.861) (112,0.903) (113,0.957) (114,0.907) (115,1.000) (116,1.175) (117,1.097) (118,1.234) (119,1.275) (120,1.277) (121,1.190) (122,1.132) (123,1.058) (124,1.101) (125,1.033) (126,0.992) (127,0.994) (128,1.106) (129,1.068) (130,1.047) (131,1.464) (132,1.361) (133,1.269) (134,1.236) (135,1.260) (136,1.273) (137,1.200) (138,1.122) (139,1.049) (140,0.992) (141,0.938) (142,0.900) (143,0.851) (144,0.809) (145,0.778) (146,0.768) (147,0.774) (148,0.739) (149,0.770) (150,0.741) (151,0.708) (152,0.675) (153,0.654) (154,0.631) (155,0.623) (156,0.604) (157,0.584) (158,0.580) (159,0.563) (160,0.584) (161,0.584) (162,0.606) (163,0.599) (164,0.617) (165,0.634) (166,0.615) (167,0.627) (168,0.606) (169,0.592) (170,0.588) (171,0.573) (172,0.556) (173,0.542) (174,0.530) (175,0.521) (176,0.526) (177,0.521) (178,0.515) (179,0.506) (180,0.499) (181,0.494) (182,0.487) (183,0.491) (184,0.529) (185,0.539) (186,0.530) (187,0.583) (188,0.569) (189,0.558) (190,0.586) (191,0.571) (192,0.560) (193,0.557) (194,0.542) (195,0.538) (196,0.578) (197,0.561) (198,0.569) (199,0.577) (200,0.571)};
  \addlegendentry{$\pm\sigma_t$}
  \addplot[thick,thmcol,no marks,forget plot] coordinates {(1,-1.00) (2,-1.06) (3,-0.99) (4,-1.00) (5,-1.03) (6,-1.10) (7,-1.07) (8,-1.07) (9,-1.22) (10,-1.23) (11,-1.24) (12,-1.16) (13,-1.10) (14,-1.12) (15,-1.53) (16,-1.64) (17,-1.52) (18,-1.43) (19,-1.46) (20,-1.97) (21,-1.83) (22,-1.70) (23,-2.13) (24,-1.98) (25,-2.12) (26,-2.20) (27,-2.28) (28,-2.13) (29,-1.98) (30,-2.12) (31,-2.01) (32,-1.86) (33,-1.83) (34,-1.71) (35,-1.58) (36,-1.50) (37,-1.40) (38,-1.31) (39,-1.22) (40,-1.15) (41,-1.08) (42,-1.16) (43,-1.10) (44,-1.03) (45,-0.99) (46,-1.25) (47,-1.23) (48,-1.15) (49,-1.08) (50,-1.07) (51,-1.02) (52,-1.10) (53,-1.42) (54,-1.35) (55,-1.45) (56,-1.52) (57,-1.48) (58,-1.49) (59,-1.45) (60,-1.45) (61,-1.35) (62,-1.26) (63,-1.23) (64,-1.17) (65,-1.11) (66,-1.05) (67,-1.05) (68,-0.99) (69,-0.93) (70,-0.91) (71,-0.91) (72,-0.86) (73,-0.88) (74,-0.92) (75,-0.89) (76,-0.86) (77,-0.97) (78,-1.03) (79,-1.00) (80,-0.94) (81,-0.90) (82,-0.94) (83,-0.91) (84,-0.86) (85,-1.05) (86,-1.01) (87,-0.97) (88,-0.91) (89,-0.89) (90,-0.88) (91,-0.89) (92,-0.90) (93,-0.84) (94,-1.03) (95,-0.97) (96,-0.91) (97,-0.87) (98,-0.91) (99,-0.90) (100,-0.98) (101,-1.05) (102,-1.05) (103,-0.98) (104,-1.00) (105,-1.05) (106,-1.01) (107,-0.99) (108,-0.94) (109,-0.88) (110,-0.90) (111,-0.86) (112,-0.90) (113,-0.96) (114,-0.91) (115,-1.00) (116,-1.18) (117,-1.10) (118,-1.23) (119,-1.27) (120,-1.28) (121,-1.19) (122,-1.13) (123,-1.06) (124,-1.10) (125,-1.03) (126,-0.99) (127,-0.99) (128,-1.11) (129,-1.07) (130,-1.05) (131,-1.46) (132,-1.36) (133,-1.27) (134,-1.24) (135,-1.26) (136,-1.27) (137,-1.20) (138,-1.12) (139,-1.05) (140,-0.99) (141,-0.94) (142,-0.90) (143,-0.85) (144,-0.81) (145,-0.78) (146,-0.77) (147,-0.77) (148,-0.74) (149,-0.77) (150,-0.74) (151,-0.71) (152,-0.68) (153,-0.65) (154,-0.63) (155,-0.62) (156,-0.60) (157,-0.58) (158,-0.58) (159,-0.56) (160,-0.58) (161,-0.58) (162,-0.61) (163,-0.60) (164,-0.62) (165,-0.63) (166,-0.61) (167,-0.63) (168,-0.61) (169,-0.59) (170,-0.59) (171,-0.57) (172,-0.56) (173,-0.54) (174,-0.53) (175,-0.52) (176,-0.53) (177,-0.52) (178,-0.52) (179,-0.51) (180,-0.50) (181,-0.49) (182,-0.49) (183,-0.49) (184,-0.53) (185,-0.54) (186,-0.53) (187,-0.58) (188,-0.57) (189,-0.56) (190,-0.59) (191,-0.57) (192,-0.56) (193,-0.56) (194,-0.54) (195,-0.54) (196,-0.58) (197,-0.56) (198,-0.57) (199,-0.58) (200,-0.57)};
\end{axis}
\end{tikzpicture}
\caption{A simulated GARCH$(1,1)$ path, Gaussian $z_t$, $\alpha_1=0.12$, $\beta_1=0.85$, long-run volatility $1\%$ per day (200 days). Returns (bars) and the conditional
volatility $\pm\sigma_t$: calm and turbulent stretches alternate, and the amplitude of the returns follows $\sigma_t$, though the returns are uncorrelated (lag-1 autocorrelation $0.06$, of the squares $0.15$).}\label{ch13:fig-garchpath}
\end{figure}

\subsection{ARMA representation of the squared innovations}\label{ch13:sec-arma}
\begin{proposition}[GARCH is ARMA for $\varepsilon^2$]\label{ch13:prop-arma}
Let $u_t=\varepsilon_t^2-\sigma_t^2$, so that $\E[u_t\mid\mathcal F_{t-1}]=0$. A GARCH$(p,q)$ satisfies, with $\alpha_i=0$ for $i>p$ and $\beta_j=0$ for $j>q$,
\begin{equation}\label{ch13:eq-armasq}
  \varepsilon_t^2=\alpha_0+\sum_{i=1}^{\max(p,q)}(\alpha_i+\beta_i)\varepsilon_{t-i}^2+u_t-\sum_{j=1}^q\beta_ju_{t-j},
\end{equation}
an ARMA$(\max(p,q),q)$ for $\varepsilon_t^2$. For the GARCH$(1,1)$:
$\varepsilon_t^2=\alpha_0+(\alpha_1+\beta_1)\varepsilon_{t-1}^2+u_t-\beta_1u_{t-1}$, an ARMA$(1,1)$ with $\phi=\alpha_1+\beta_1$ and $\theta=-\beta_1$.
\end{proposition}
\begin{proof}
From $\sigma_{t-j}^2=\varepsilon_{t-j}^2-u_{t-j}$ substitute in \eqref{ch13:eq-garch}: $\varepsilon_t^2-u_t=\alpha_0+\sum_i\alpha_i\varepsilon^2_{t-i}+\sum_j\beta_j(\varepsilon_{t-j}^2-u_{t-j})$; collect terms.
\end{proof}
{\emergencystretch=3em This is the statement announced in \cref{ch09:sec-arma11} \ts{24}{19}{1:13:09}\ts{13}{11}{4:15}.
\begin{secondary}[A slip in the 2013 transcript]
(the transcript of 2013 writes $-\beta_1u_t$ for $-\beta_1u_{t-1}$).
\end{secondary}
With the AR polynomial $A(L)=1-\sum_i(\alpha_i+\beta_i)L^i$ and the MA polynomial $B(L)=1-\sum_j\beta_jL^j$
one has
\[
  A(L)\varepsilon^2_t=\alpha_0+B(L)u_t \qquad \ts{13}{11}{6:19}.
\]
}

\begin{secondary}[The root condition]
\begin{lemma}[Stationarity of the ARMA]\label{ch13:lem-roots}
For $c_i=\alpha_i+\beta_i\ge0$ the roots of $A(z)=1-\sum_ic_iz^i$ lie outside the closed unit disc iff $\sum_ic_i<1$.
\end{lemma}
\begin{proof}
If $\sum c_i<1$, then for $|z|\le1$, $|\sum c_iz^i|\le\sum c_i<1$, so $A(z)\ne0$. If $\sum c_i\ge1$, then $A(0)=1>0\ge A(1)$, and $A$ is continuous, so it has a root in $(0,1]$.
\end{proof}
\end{secondary}
Consequently the ARMA of \eqref{ch13:eq-armasq} is covariance stationary iff $\sum(\alpha_i+\beta_i)<1$; for the $(1,1)$ case, $\alpha_1+\beta_1<1$
\ts{24}{19}{1:13:09}\ts{13}{11}{6:19}. When $\E[\varepsilon^4]<\infty$ (\cref{ch13:thm-kurt}) the noise $u_t$ has constant variance and is uncorrelated, and the ARMA formulas of
\cref{ch09:sec-arma11} give the autocorrelation function of the squares:
\begin{equation}\label{ch13:eq-acfsq}
  \rho_{\varepsilon^2}(1)=\frac{\alpha_1(1-\alpha_1\beta_1-\beta_1^2)}{1-2\alpha_1\beta_1-\beta_1^2},\qquad
  \rho_{\varepsilon^2}(k)=(\alpha_1+\beta_1)^{k-1}\rho_{\varepsilon^2}(1),\ k\ge2 .
\end{equation}
\begin{secondary}[Checking the formula]
(Insert $\phi=\alpha_1+\beta_1$, $\theta=-\beta_1$ in \cref{ch09:eq-arma11}: $1+\phi\theta=1-\alpha_1\beta_1-\beta_1^2$, $\phi+\theta=\alpha_1$, $1+2\phi\theta+\theta^2=1-2\alpha_1\beta_1-\beta_1^2$.) The lag-one
correlation does not depend on the kurtosis of $z_t$. For $\alpha_1=0.10$, $\beta_1=0.85$: $\rho(1)=0.179$, $\rho(5)=0.146$ (simulation, $2\times10^6$ days: $0.181$ and
$0.146$); for $\alpha_1=0.05$, $\beta_1=0.85$ one gets $\rho(1)=0.061$ with slow decay, the ``small first correlation but long memory'' example of \cref{ch09:sec-arma11}.
\end{secondary}
The slow geometric decay at rate $\alpha_1+\beta_1$, near $1$, is what the significant autocorrelations of squared returns at many lags (the case studies) require, and what a low-order ARCH
cannot produce.

\subsection{Stationarity and long-run variance}\label{ch13:sec-longrun}
\begin{theorem}[Stationarity of the GARCH$(1,1)$]\label{ch13:thm-stat}
Let $\alpha_0>0$ and $\phi=\alpha_1+\beta_1$. There is a covariance stationary solution of \eqref{ch13:eq-garch11} iff $\phi<1$, and then
\begin{equation}\label{ch13:eq-longrun}
  \sigma_*^2\;:=\;\E[\varepsilon_t^2]=\E[\sigma_t^2]=\frac{\alpha_0}{1-\alpha_1-\beta_1}.
\end{equation}
More generally, a GARCH$(p,q)$ is covariance stationary iff $\sum_i(\alpha_i+\beta_i)<1$, with $\sigma_*^2=\alpha_0/\bigl[1-\sum_{i\le\max(p,q)}(\alpha_i+\beta_i)\bigr]$.
\end{theorem}
\begin{proof}
Put $a_t=\alpha_1z_t^2+\beta_1$, an i.i.d.\ sequence with $\E[a_t]=\phi$; then \eqref{ch13:eq-garch11} reads $\sigma_t^2=\alpha_0+a_{t-1}\sigma_{t-1}^2$, where $a_{t-1}$ is independent of $\sigma_{t-1}^2$.
\emph{Necessity.} If a stationary solution with $\E[\sigma_t^2]=m<\infty$ exists, taking expectations gives $m=\alpha_0+\phi m$, so $m=\alpha_0/(1-\phi)$, which is positive only if $\phi<1$.
\emph{Sufficiency.} Iterating $n$ times,
$\sigma_t^2=\alpha_0\bigl(1+a_{t-1}+a_{t-1}a_{t-2}+\dots+a_{t-1}\cdots a_{t-n+1}\bigr)+a_{t-1}\cdots a_{t-n}\,\sigma^2_{t-n}$. If $\phi<1$ the series with $n=\infty$
has non-negative terms and expectation $\alpha_0\sum_k\phi^k=\alpha_0/(1-\phi)<\infty$, hence converges almost surely and in $L^1$, and it defines a stationary solution (a function of
$z_{t-1},z_{t-2},\dots$); the remainder term has expectation $\phi^n\E[\sigma^2_{t-n}]\to0$, so this is the only solution with finite mean. For the general $(p,q)$ case take expectations in
\eqref{ch13:eq-garch} (with $\E[\varepsilon^2]=\E[\sigma^2]$) and use \cref{ch13:lem-roots} for the equivalence with the ARMA polynomial \ts{13}{11}{7:23}.
\end{proof}
Taking expectations of \eqref{ch13:eq-garch11} is how the lecture obtains the \emph{long-run variance} $\sigma_*^2=\alpha_0+(\alpha_1+\beta_1)\sigma_*^2$ \ts{24}{19}{1:14:13}. Conversely
$\alpha_0=(1-\alpha_1-\beta_1)\sigma_*^2$: one can parametrise the model by $(\sigma_*^2,\alpha_1,\beta_1)$, and $\sigma_*^2$ is then simply the unconditional variance of the returns (which the
sample variance estimates).

\begin{coursenote}[Erratum: typo in the general long-run variance; constraints]\label{ch13:cn-longrun}
Slide~26 of \file{lec17_1} (slide~25 of 2013) writes $\sigma_*^2=\alpha_0+\bigl(\sum_1^p\alpha_j+\sum_1^q\beta_1\bigr)\sigma_*^2$; the last coefficient should read $\beta_j$. Slide~27 has a typo
in the constraint $\sum\alpha_u$ for $\sum\alpha_i$. The strict stationarity of the process is a weaker condition than covariance stationarity (Nelson, 1990; not discussed in class):
the GARCH$(1,1)$ has a strictly stationary solution iff $\E\ln(\alpha_1z_t^2+\beta_1)<0$, and by Jensen $\E[\ln a_t]\le\ln\phi$, so $\phi<1$ suffices but $\phi\ge1$ is not excluded: the model may be
strictly stationary with infinite variance.
\end{coursenote}

\subsection{Persistence, mean reversion and half-life}\label{ch13:sec-persist}
Subtracting $\sigma_*^2$ from \eqref{ch13:eq-garch11} written with $\alpha_0=(1-\phi)\sigma_*^2$ gives the \emph{mean-reverting form}
\begin{equation}\label{ch13:eq-meanrev}
  \sigma_t^2-\sigma_*^2=\alpha_1\bigl(\varepsilon_{t-1}^2-\sigma_*^2\bigr)+\beta_1\bigl(\sigma_{t-1}^2-\sigma_*^2\bigr),\qquad
  \varepsilon_t^2-\sigma_*^2=\phi\,\bigl(\varepsilon_{t-1}^2-\sigma_*^2\bigr)+u_t-\beta_1u_{t-1},
\end{equation}
the second being the ARMA of \cref{ch13:prop-arma} \ts{24}{19}{1:14:13}\ts{13}{11}{25:25}. It says that the deviation of the squared return from its long-run
average is a fraction $\phi$ of the previous deviation plus noise: a discrete Ornstein--Uhlenbeck process (an AR(1), \cref{ch09:sec-ar1}; the continuous-time version is \cref{ch:sde}) with MA(1) noise. Because $0<\phi<1$ volatility \emph{reverts to the mean}: the
excess of variance over $\sigma_*^2$ shrinks by the factor $\phi$ per period \ts{13}{11}{26:27}. The sum $\alpha_1+\beta_1$ is therefore called the \emph{persistence} of the model \ts{13}{11}{27:33}; the
\emph{half-life} of a variance shock is $\ln\frac12/\ln\phi$ periods:
\begin{center}\small
\begin{tabular}{@{}lcccccc@{}}
\toprule
persistence $\phi$ & $0.90$ & $0.95$ & $0.97$ & $0.99$ & $0.997$\\
\midrule
half-life (days) & $6.6$ & $13.5$ & $22.8$ & $69$ & $231$\\
\bottomrule
\end{tabular}
\end{center}
In the ARMA language the persistence is the AR root; the parameters $\alpha_1$ (reaction to news) and $\beta_1$ (memory) have distinct roles: for the same $\phi$ a larger $\alpha_1$ gives a more
jagged volatility with larger kurtosis (\cref{ch13:thm-kurt}), a larger $\beta_1$ a smoother one.

\begin{finance}[Persistence of the euro/dollar volatility]\label{ch13:fin-persist}
The Gaussian GARCH$(1,1)$ fitted to EUR/USD returns in the 2013 case study (\cref{ch13:cs-fx}) has $\hat\alpha_1=0.0278$, $\hat\beta_1=0.9696$, $\hat\omega=1.19\times10^{-7}$: persistence $0.9973$,
half-life $259$ trading days (about a year), long-run variance $\hat\omega/(1-\hat\phi)=4.46\times10^{-5}$, i.e.\ a daily volatility $0.668\%$ and an annual $10.6\%$, remarkably close to the $10.25\%$ of the constant-volatility fit
of \cref{ch13:sec-clock}. The case study concludes that the model is ``very close to being non-stationary'' and that there is practically no tendency to revert to a specific volatility level within a horizon of
a few months \ts{13}{11}{27:33}: volatility is a slowly wandering quantity, a feature that risk managers exploit by using an exponential average with a slow decay.
\end{finance}

\subsection{Forecasting the variance}\label{ch13:sec-forecast}
\begin{theorem}[Variance forecasts of a GARCH$(1,1)$]\label{ch13:thm-forecast}
Let $\phi=\alpha_1+\beta_1<1$. For every $h\ge1$,
\begin{equation}\label{ch13:eq-forecast}
  \E\bigl[\varepsilon_{t+h}^2\mid\mathcal F_t\bigr]=\E\bigl[\sigma^2_{t+h}\mid\mathcal F_t\bigr]=\sigma_*^2+\phi^{\,h-1}\bigl(\sigma_{t+1}^2-\sigma_*^2\bigr),
\end{equation}
where $\sigma_{t+1}^2=\alpha_0+\alpha_1\varepsilon_t^2+\beta_1\sigma_t^2$ is known at time $t$. The expected total variance of the next $H$ periods is
\begin{equation}\label{ch13:eq-hday}
  \E\Bigl[\sum_{h=1}^H\varepsilon^2_{t+h}\,\Big|\,\mathcal F_t\Bigr]=H\sigma_*^2+\bigl(\sigma_{t+1}^2-\sigma_*^2\bigr)\frac{1-\phi^H}{1-\phi}.
\end{equation}
\end{theorem}
\begin{proof}
By the tower property $\E[\varepsilon^2_{t+h}\mid\mathcal F_t]=\E\bigl[\E[\varepsilon^2_{t+h}\mid\mathcal F_{t+h-1}]\mid\mathcal F_t\bigr]=\E[\sigma^2_{t+h}\mid\mathcal F_t]$. Write $m_h$ for the deviation
$\E[\sigma^2_{t+h}\mid\mathcal F_t]-\sigma_*^2$. Taking $\E[\cdot\mid\mathcal F_t]$ in \eqref{ch13:eq-meanrev} for $h\ge2$ gives $m_h=\alpha_1m_{h-1}+\beta_1m_{h-1}=\phi\,m_{h-1}$, and $m_1=\sigma^2_{t+1}-\sigma_*^2$; induct.
The second formula is the geometric sum $\sum_{h=1}^H\phi^{h-1}$.
\end{proof}
Two consequences. The forecast decays geometrically to the long-run variance at the rate $\phi$ per period: in the short run the variance moves incrementally, in the long run it returns to $\sigma_*^2$
\ts{13}{11}{26:27}. And, since the $\varepsilon_t$ are uncorrelated, \eqref{ch13:eq-hday} is the variance of the $H$-period return (the conditional variance of a sum of martingale differences is the sum of the conditional variances of the terms), so
the annualised volatility expected over the next $H$ days is $\sqrt{(252/H)\times}$ the right-hand side of \eqref{ch13:eq-hday}.

\begin{example}[After a turbulent day]\label{ch13:ex-forecast}
Let $\sigma_*=1\%$ per day (annualised $15.9\%$), $\alpha_1=0.06$, $\beta_1=0.91$, so $\phi=0.97$ and the half-life is $22.8$ days, and let today's news raise the forecast for tomorrow to $\sigma_{t+1}=2\%$
($31.7\%$ annualised). Then the annualised volatility expected over the next $H$ days is $31.7\%$ for $H=1$ and
\begin{center}\small
\begin{tabular}{@{}lccccc@{}}
\toprule
$H$ (days) & $5$ & $10$ & $21$ & $63$ & $252$\\
\midrule
GARCH$(1,1)$ & $31.0\%$ & $30.2\%$ & $28.6\%$ & $24.4\%$ & $18.8\%$\\
$\sqrt H$ rule from $\sigma_{t+1}$ & $31.7\%$ & $31.7\%$ & $31.7\%$ & $31.7\%$ & $31.7\%$\\
\bottomrule
\end{tabular}
\end{center}
For a $21$-day horizon the standard deviation of the return is $8.26\%$ against $\sqrt{21}\times2\%=9.17\%$ from the naive scaling, an overstatement of $11\%$; for a year the naive rule is wrong by a
factor $1.7$. This is the failure of the square-root-of-time rule announced in \cref{ch13:rem-sqrt}, and it is relevant for the multi-day value at risk of \cref{ch:var} (cf.\ \cref{ch08:prop-autocorr}, where the discrepancy has the opposite sign, being due to autocorrelation of returns).
\end{example}

\begin{figure}[ht]
\centering
\begin{tikzpicture}
\begin{axis}[name=L,width=0.49\linewidth,height=5.6cm,xmin=1,xmax=150,ymin=0,ymax=4.3,
  axis lines=left,xlabel={horizon $h$ (days)},ylabel={$\E_t[\sigma^2_{t+h}]/\sigma_*^2$},
  legend style={font=\scriptsize,draw=none,fill=white,fill opacity=0.85,text opacity=1,legend pos=north east},
  tick label style={font=\small},label style={font=\small},title style={font=\small},title={variance forecast}]
  \addplot[very thick,defcol,domain=1:150,samples=100] {1+3*0.97^(x-1)};
  \addlegendentry{$\phi=0.97$}
  \addplot[very thick,defcol,domain=1:150,samples=100,forget plot] {1-0.75*0.97^(x-1)};
  \addplot[very thick,thmcol,dashed,domain=1:150,samples=100] {1+3*0.99^(x-1)};
  \addlegendentry{$\phi=0.99$}
  \addplot[very thick,thmcol,dashed,domain=1:150,samples=100,forget plot] {1-0.75*0.99^(x-1)};
  \draw[gray,dotted] (axis cs:1,1)--(axis cs:150,1);
\end{axis}
\begin{axis}[at={(L.east)},anchor=west,xshift=1.1cm,width=0.49\linewidth,height=5.6cm,xmin=1,xmax=252,ymin=10,ymax=35,
  axis lines=left,xlabel={horizon $H$ (days)},ylabel={annualised volatility (\%)},
  legend style={font=\scriptsize,draw=none,fill=white,fill opacity=0.85,text opacity=1,legend pos=south west},
  tick label style={font=\small},label style={font=\small},title style={font=\small},title={average volatility over $H$ days}]
  \addplot[very thick,defcol,domain=1:252,samples=120] {sqrt(252*(1+3*(1-0.97^x)/(x*0.03)))};
  \addlegendentry{$\phi=0.97$}
  \addplot[very thick,thmcol,dashed,domain=1:252,samples=120] {sqrt(252*(1+3*(1-0.99^x)/(x*0.01)))};
  \addlegendentry{$\phi=0.99$}
  \draw[gray,dotted] (axis cs:1,31.75)--(axis cs:252,31.75);
  \draw[gray,dotted] (axis cs:1,15.87)--(axis cs:252,15.87);
  \node[font=\scriptsize,gray,anchor=south west] at (axis cs:100,32) {$\sqrt H$ rule};
  \node[font=\scriptsize,gray,anchor=north west] at (axis cs:100,15.5) {long-run level};
\end{axis}
\end{tikzpicture}
\caption{Left: expected conditional variance $\E_t[\sigma^2_{t+h}]$ in units of $\sigma_*^2$ after a shock ($\sigma_{t+1}=2\sigma_*$) and after a calm day ($\sigma_{t+1}=\sigma_*/2$), for two persistences, from \eqref{ch13:eq-forecast}. Right:
annualised volatility expected over the next $H$ days after the shock of \cref{ch13:ex-forecast}, from \eqref{ch13:eq-hday}; the dotted lines are the $\sqrt H$ extrapolation of today's volatility and the long-run level $15.9\%$.}\label{ch13:fig-forecast}
\end{figure}

\begin{exercise}[(not from the course): forecasting]\label{ch13:ex-extra1}
For a GARCH$(1,1)$ with $\alpha_1=0.08$, $\beta_1=0.90$, long-run daily volatility $1\%$, and $\sigma_{t+1}=1.6\%$, compute the expected annualised volatility over the next $21$ and $252$ days from \eqref{ch13:eq-hday}, and the half-life of a variance shock. Check by simulation.
\end{exercise}

\subsection{Fourth moment and kurtosis}\label{ch13:sec-kurt}
GARCH reproduces the fat tails of returns even when the innovations $z_t$ are Gaussian, and one can say by how much.
\begin{theorem}[Kurtosis of the GARCH$(1,1)$]\label{ch13:thm-kurt}
Let $\kappa_z=\E[z_t^4]<\infty$, $\phi=\alpha_1+\beta_1<1$, and assume $\Delta:=1-\phi^2-(\kappa_z-1)\alpha_1^2>0$. Then $\E[\varepsilon_t^4]<\infty$ and the kurtosis of $\varepsilon_t$ is
\begin{equation}\label{ch13:eq-kurt}
  \kappa_\varepsilon=\frac{\E[\varepsilon_t^4]}{(\E[\varepsilon_t^2])^2}=\kappa_z\,\frac{1-\phi^2}{1-\phi^2-(\kappa_z-1)\alpha_1^2}\ >\ \kappa_z .
\end{equation}
If $\Delta\le0$ the fourth moment is infinite. For Gaussian $z$ the condition is $(\alpha_1+\beta_1)^2+2\alpha_1^2<1$, i.e.\ $3\alpha_1^2+2\alpha_1\beta_1+\beta_1^2<1$.
\end{theorem}
\begin{proof}
With $a_t=\alpha_1z_t^2+\beta_1$ as before, $\E[a_t]=\phi$ and $\E[a_t^2]=\alpha_1^2\kappa_z+2\alpha_1\beta_1+\beta_1^2=\phi^2+(\kappa_z-1)\alpha_1^2$, so $\Delta=1-\E[a_t^2]$.
\emph{Finite fourth moment iff $\E[a^2]<1$:} the series of \cref{ch13:thm-stat} has terms $\alpha_0\prod_{j\le k}a_{t-j}$ of $L^2$ norm $\alpha_0(\E[a^2])^{k/2}$, summable iff $\E[a^2]<1$ (Minkowski's inequality gives sufficiency);
if $\E[\sigma^4]<\infty$ and stationary, the identity below forces $\E[a^2]<1$.
Squaring $\sigma_t^2=\alpha_0+a_{t-1}\sigma_{t-1}^2$ and using independence of $a_{t-1}$ and $\sigma_{t-1}^2$,
\[
  \E[\sigma^4]=\alpha_0^2+2\alpha_0\phi\,\sigma_*^2+\E[a^2]\,\E[\sigma^4]\ \Longrightarrow\
  \E[\sigma^4]=\frac{\alpha_0^2(1+\phi)}{(1-\phi)\,\Delta},
\]
because $\alpha_0^2+2\alpha_0\phi\cdot\alpha_0/(1-\phi)=\alpha_0^2(1+\phi)/(1-\phi)$. Since $\varepsilon_t=\sigma_tz_t$ with $z_t$ independent of $\sigma_t$, $\E[\varepsilon^4]=\kappa_z\E[\sigma^4]$ and
$\kappa_\varepsilon=\kappa_z\E[\sigma^4]/\sigma_*^4=\kappa_z\dfrac{(1+\phi)(1-\phi)^2}{(1-\phi)\Delta}=\kappa_z\dfrac{1-\phi^2}{\Delta}>\kappa_z$.
\end{proof}
The mechanism is that of \cref{ch03:prop-mixture}: given the past, $\varepsilon_t$ is normal with variance $\sigma_t^2$, and this variance fluctuates, so the marginal law is a mixture of
Gaussians with kurtosis $\kappa_z\E[\sigma^4]/(\E[\sigma^2])^2\ge\kappa_z$, strictly unless $\sigma_t$ is constant \ts{13}{11}{24:19}; the slides describe GARCH as a ``stochastic mixture of Gaussians with higher kurtosis''.
For $\beta_1=0$ one recovers the ARCH$(1)$ formulas of 2013 PS5 Problem~3 (\cref{ch13:ex-ps5-3}): $\E[\varepsilon^4]=\kappa_z\alpha_0^2(1+\alpha_1)/\bigl[(1-\alpha_1)(1-\kappa_z\alpha_1^2)\bigr]$, finite iff $\alpha_1<1/\sqrt3=0.577$ for Gaussian $z$.

\begin{secondary}[Kurtosis for typical parameters]
\begin{center}\small
\begin{tabular}{@{}lccc@{}}
\toprule
$(\alpha_1,\beta_1)$ & $\phi$ & $1-\phi^2-2\alpha_1^2$ & kurtosis (Gaussian $z$)\\
\midrule
$(0.05,\,0.90)$ & $0.95$ & $0.0925$ & $3.16$\\
$(0.10,\,0.85)$ & $0.95$ & $0.0775$ & $3.77$ (simulation $3.79$)\\
$(0.0278,\,0.9696)$, EUR/USD fit & $0.9973$ & $0.0038$ & $4.22$\\
$(0.10,\,0.88)$ & $0.98$ & $0.0196$ & $6.06$\\
$(0.09,\,0.90)$ & $0.99$ & $0.0037$ & $16.1$\\
$(0.12,\,0.87)$ & $0.99$ & $-0.0089$ & $\infty$\\
\bottomrule
\end{tabular}
\end{center}
\end{secondary}
Persistence close to $1$ makes the fourth moment barely exist, and then sample kurtosis is an unreliable estimator of the theoretical one: sample fourth moments of such a process converge very slowly.
\begin{secondary}[Power-law tails]
(A tail-index refinement, not discussed: $\P(|\varepsilon|>x)\sim x^{-\xi}$ with $\xi$ solving $\E[a_t^{\xi/2}]=1$, so GARCH tails are power laws, as in the empirical ``inverse cubic'' law of \cref{ch03:sec-heavy}.)
\end{secondary}

\subsection{The RiskMetrics limit: IGARCH}\label{ch13:sec-igarch}
The RiskMetrics exponential average with decay factor $\beta$ is the GARCH$(1,1)$ with $\alpha_0=0$, $\alpha_1=1-\beta$, $\beta_1=\beta$, that is $\phi=1$:
\begin{equation}\label{ch13:eq-igarch}
  \sigma_{t+1}^2=(1-\beta)\,\varepsilon_t^2+\beta\,\sigma_t^2,
\end{equation}
an \emph{integrated} GARCH (IGARCH). It has no long-run variance, and by \eqref{ch13:eq-forecast} with $\phi=1$ its forecast is flat, $\E_t[\sigma^2_{t+h}]=\sigma^2_{t+1}$ for all $h$, so it scales volatility with $\sqrt H$
\ts{13}{11}{27:33}. The lecturer describes the advantage: one tracks a possibly non-stationary volatility without presuming that a long-run average consistent with the past exists.
\begin{secondary}[Two caveats on IGARCH]
Two caveats, not
discussed: as a \emph{generating} process the IGARCH with $\alpha_0=0$ is degenerate, since $\ln\sigma_t^2=\ln\sigma_0^2+\sum_{j\le t}\ln a_j$ has i.i.d.\ increments of strictly negative mean
$\E[\ln a]<\ln\E[a]=0$ (Jensen), so $\sigma_t^2\to0$ almost surely; for $\alpha_1=0.06$ and Gaussian $z$ the drift is $\E[\ln a]=-0.0032$ per day, slow enough to be invisible in a sample of a few hundred days. And the
recursion \eqref{ch13:eq-igarch} is exactly the steady-state Kalman filter of a local-level model (\cref{ch09:sec-ss}), which is a better way to think of it: a \emph{filter} for the current level of volatility, not a
model of its evolution.
\end{secondary}


% ==================================================================
\section{Estimation of ARCH and GARCH models}\label{ch13:sec-mle}

\subsection{The likelihood}\label{ch13:sec-lik}
A GARCH model specifies the law of $y_t$ \emph{given the past}, and this is exactly what the likelihood needs.
Write $\theta=(c,\alpha_0,\alpha_1,\dots,\alpha_p,\beta_1,\dots,\beta_q)$ for the parameters, $y_t=c+\varepsilon_t$ (a constant mean; \cref{ch13:sec-fxcs} uses an AR mean),
and $\varepsilon_t=\sigma_tz_t$. The joint density factorises into conditional densities (the chain rule of probability) \ts{24}{19}{1:14:13}\ts{13}{11}{8:26}.
\begin{proposition}[Gaussian log-likelihood of a GARCH model]\label{ch13:prop-lik}
If $z_t$ is $N(0,1)$ then $y_t\mid\mathcal F_{t-1}\sim N(c,\sigma_t^2)$, where $\sigma_t^2=\sigma_t^2(\theta)$ is computed recursively from
\eqref{ch13:eq-garch} and the data, and
\begin{equation}\label{ch13:eq-loglik}
  \ell(\theta)=\sum_{t=1}^T\log p(y_t\mid\mathcal F_{t-1};\theta)=-\frac12\sum_{t=1}^T\Bigl[\log(2\pi)+\log\sigma_t^2+\frac{\varepsilon_t^2}{\sigma_t^2}\Bigr],\qquad\varepsilon_t=y_t-c .
\end{equation}
\end{proposition}
\begin{proof}
$p(y_1,\dots,y_T)=\prod_tp(y_t\mid y_{t-1},\dots,y_1)$ and, given $\mathcal F_{t-1}$, both $c$ and $\sigma_t^2$ are known numbers, so the conditional law is the stated normal. Take logarithms.
\end{proof}
Two points.
\begin{secondary}[Starting values]
First, the recursion needs starting values ($\varepsilon_0^2$ and $\sigma_0^2$): common choices are the sample variance for both, or the long-run variance $\alpha_0/(1-\phi)$; the effect fades out
geometrically at rate $\beta_1$ (\cref{ch13:prop-archinf}) but is felt for a long time when $\beta_1\approx1$ (not discussed in class).
\end{secondary}
Second, the likelihood is \emph{not} a quadratic function of $\theta$
and has no closed-form maximiser (unlike the regression of \cref{ch:regression}); the maximum is found numerically. The two terms have a clear meaning: $\log\sigma_t^2$ penalises a large predicted variance,
$\varepsilon_t^2/\sigma_t^2$ penalises a variance too small for the return that occurred, and the maximum balances them, exactly as in the Gaussian estimate of the scale of an i.i.d.\ sample
(where $\sigma_t^2\equiv\hat\sigma^2$ gives \eqref{ch13:eq-gbmmle}).

\paragraph{Constraints and reparametrisation.} The parameters are constrained: $\alpha_0>0$, $\alpha_i,\beta_j\ge0$ and $\sum_i\alpha_i+\sum_j\beta_j<1$ (\cref{ch13:thm-stat}) \ts{13}{11}{9:30}. The lecturer's
remark is practical: optimisation algorithms behave well on \emph{unconstrained} parameters, so one transforms them, the logarithm for a positive parameter, the logit
$\log\frac x{1-x}$ for a parameter in $(0,1)$ \ts{13}{11}{10:39}. For the GARCH$(1,1)$ one can write
\[
  \alpha_0=e^{a},\qquad \alpha_1=\phi\,s,\quad\beta_1=\phi\,(1-s),\quad \phi=\frac1{1+e^{-u}},\quad s=\frac1{1+e^{-v}},
\]
with $(a,u,v)\in\R^3$ free; then $\alpha_0>0$, $\alpha_1,\beta_1>0$ and $\alpha_1+\beta_1=\phi\in(0,1)$ hold automatically. (This particular map is our example.)
\begin{secondary}[Variance targeting]
A second trick, not discussed, is \emph{variance targeting}:
replace $\alpha_0$ by $(1-\phi)\hat s^2$, with $\hat s^2$ the sample variance of the returns, which removes one parameter and enforces the long-run level \eqref{ch13:eq-longrun}.
\end{secondary}

\begin{remark}[Estimation problems]\label{ch13:rem-ident}
If $\alpha_1=0$ the parameter $\beta_1$ disappears from the model (since $\sigma_t^2$ is constant), so the likelihood is flat in $\beta_1$ and the estimators of $\alpha_1,\beta_1$ are badly behaved near this boundary; and when $\alpha_1+\beta_1\approx1$ the
asymptotic theory is delicate as well. Standard errors in \cref{ch13:sec-fxcs} are those printed by the software. When $z_t$ is not Gaussian the Gaussian likelihood \eqref{ch13:eq-loglik} is still a sensible objective: its maximiser is a
\emph{quasi-maximum-likelihood} estimator, which is consistent under conditions weaker than normality (Bollerslev and Wooldridge, 1992; not discussed) but no longer efficient, and it is the natural counterpart of
the remark on the LM test in \cref{ch13:sec-arch-ar}.
\end{remark}

\begin{exercise}[(not from the course): estimation]\label{ch13:ex-extra4}
Simulate $2000$ returns of a Gaussian GARCH$(1,1)$ with $(\alpha_0,\alpha_1,\beta_1)=(2\times10^{-6},0.08,0.90)$, write the log-likelihood \eqref{ch13:eq-loglik} in the reparametrisation of \cref{ch13:sec-lik} and maximise it; compare the estimates with the truth over $200$ replications and comment on the bias and dispersion of $\hat\alpha_1+\hat\beta_1$.
\end{exercise}

\subsection{Innovations with the Student $t$ distribution}\label{ch13:sec-t}
The Gaussian assumption leaves two symptoms in the residuals of the case study: standardised residuals with values six standard deviations from the mean, and a bend of the Q--Q plot at the extremes
\ts{13}{11}{14:47}. A natural heavier-tailed alternative is the Student $t$ law. Because $\Var z_t$ must be $1$, one takes the \emph{standardised} $t$ with $\nu>2$ degrees of freedom:
\begin{equation}\label{ch13:eq-tstd}
  z=\sqrt{\frac{\nu-2}{\nu}}\,T_\nu,\qquad
  f_\nu(z)=\frac{\Gamma\bigl(\tfrac{\nu+1}2\bigr)}{\Gamma\bigl(\tfrac\nu2\bigr)\sqrt{\pi(\nu-2)}}\Bigl(1+\frac{z^2}{\nu-2}\Bigr)^{-\frac{\nu+1}2},
\end{equation}
where $T_\nu$ has the usual $t$ density $\propto(1+x^2/\nu)^{-(\nu+1)/2}$ \ts{13}{11}{19:08}. The scaling is needed because $\Var T_\nu=\nu/(\nu-2)$; this is why the case-study code multiplies by $\sqrt{(\nu-2)/\nu}$.

\begin{proposition}[Moments and tails of the $t$ law]\label{ch13:prop-t}
Let $T_\nu=Z/\sqrt{V/\nu}$ with $Z\sim N(0,1)$ independent of $V\sim\chi^2_\nu$. Then $\Var T_\nu=\nu/(\nu-2)$ for $\nu>2$; the kurtosis is finite for $\nu>4$ and equals
\[
  \kappa_t=3+\frac6{\nu-4}\qquad(\nu>4),
\]
and $\Prob(|T_\nu|>x)\sim c_\nu x^{-\nu}$ as $x\to\infty$: a power-law tail, with only the moments of order $<\nu$ finite.
\end{proposition}
\begin{proof}
Conditionally on $V$, $T_\nu\sim N(0,\nu/V)$, so $T_\nu$ is a normal variance mixture (\cref{ch03:prop-mixture}) with mixing variable $\nu/V$. Now $\E[V^{-1}]=1/(\nu-2)$ and $\E[V^{-2}]=1/((\nu-2)(\nu-4))$ for the $\chi^2_\nu$ law
(they follow from $\E[V^{k}]=2^k\Gamma(\nu/2+k)/\Gamma(\nu/2)$ for $k=-1,-2$, and are infinite for $\nu\le2$, resp.\ $\nu\le4$). Hence $\Var T_\nu=\nu\E[V^{-1}]=\nu/(\nu-2)$, $\E[T_\nu^4]=3\nu^2\E[V^{-2}]=3\nu^2/((\nu-2)(\nu-4))$ and
$\kappa_t=\E[T^4]/(\Var T)^2=3(\nu-2)/(\nu-4)=3+6/(\nu-4)$. The tail follows from the density \eqref{ch13:eq-tstd}, which decays like $|z|^{-(\nu+1)}$.
\end{proof}
So $\nu=10$ has $\kappa_t=4$, $\nu=5$ has $9$, $\nu=30$ has $3.23$; as $\nu\to\infty$ the law tends to the normal. The value $\nu=30$ (or $25$--$40$, the answers of the students in class) is when the density and the distribution function are
visually indistinguishable from the normal \ts{13}{11}{19:08}. What differs is the \emph{far tail}, best seen on a logarithmic scale \ts{13}{11}{20:13}, \cref{ch13:fig-tails}.

\begin{figure}[ht]
\centering
\begin{tikzpicture}
\begin{axis}[width=0.9\linewidth,height=6cm,xmin=0,xmax=6,ymin=-9.2,ymax=0,
  axis lines=left,xlabel={$x$ (standard deviations)},ylabel={$\log_{10}\Prob(X>x)$},
  legend style={font=\scriptsize,draw=none,fill=white,fill opacity=0.85,text opacity=1,at={(0.5,1.03)},anchor=south,legend columns=3},
  tick label style={font=\small},label style={font=\small}]
  \addplot[very thick,defcol,no marks] coordinates {(0.00,-0.301) (0.25,-0.397) (0.50,-0.511) (0.75,-0.645) (1.00,-0.800) (1.25,-0.976) (1.50,-1.175) (1.75,-1.397) (2.00,-1.643) (2.25,-1.913) (2.50,-2.207) (2.75,-2.526) (3.00,-2.870) (3.25,-3.239) (3.50,-3.633) (3.75,-4.053) (4.00,-4.499) (4.25,-4.971) (4.50,-5.469) (4.75,-5.993) (5.00,-6.543) (5.25,-7.119) (5.50,-7.721) (5.75,-8.350) (6.00,-9.006)};
  \addlegendentry{normal}
  \addplot[very thick,thmcol,dashed,no marks] coordinates {(0.00,-0.301) (0.25,-0.406) (0.50,-0.531) (0.75,-0.676) (1.00,-0.839) (1.25,-1.017) (1.50,-1.206) (1.75,-1.404) (2.00,-1.608) (2.25,-1.815) (2.50,-2.023) (2.75,-2.231) (3.00,-2.437) (3.25,-2.640) (3.50,-2.839) (3.75,-3.034) (4.00,-3.224) (4.25,-3.410) (4.50,-3.591) (4.75,-3.767) (5.00,-3.938) (5.25,-4.104) (5.50,-4.266) (5.75,-4.423) (6.00,-4.576)};
  \addlegendentry{$t_{10}$}
  \addplot[very thick,intcol,dotted,no marks] coordinates {(0.00,-0.301) (0.25,-0.455) (0.50,-0.608) (0.75,-0.762) (1.00,-0.915) (1.25,-1.069) (1.50,-1.222) (1.75,-1.376) (2.00,-1.529) (2.25,-1.683) (2.50,-1.836) (2.75,-1.990) (3.00,-2.144) (3.25,-2.297) (3.50,-2.451) (3.75,-2.604) (4.00,-2.758) (4.25,-2.911) (4.50,-3.065) (4.75,-3.218) (5.00,-3.372) (5.25,-3.526) (5.50,-3.679) (5.75,-3.833) (6.00,-3.986)};
  \addlegendentry{Laplace}
\end{axis}
\end{tikzpicture}
\caption{One-sided tail probability of three laws with unit variance, against the distance $x$ from the mean in standard deviations. Up to about $x=2$ the three curves are close, and the $t_{10}$ is even
\emph{below} the normal near $x=1.6$ (its centre is more peaked, ``leptokurtic''); beyond $x\approx2.5$ the normal falls off super-exponentially, the Laplace exponentially and the $t_{10}$ as a power law.
At $x=6$ the normal has $10^{-9}$, the $t_{10}$ $2.7\times10^{-5}$ and the Laplace $10^{-4}$.}\label{ch13:fig-tails}
\end{figure}

\begin{coursenote}[What ``a six-sigma move'' means]\label{ch13:cn-6sd}
The 2013 lecture says that with $10$ degrees of freedom a six-standard-deviation move ``occurs about 1 in 10\,000 times'' \ts{13}{11}{21:15}. The figure corresponds to $|T_{10}|>6$ for the \emph{unstandardised} $t_{10}$ ($\Prob=1.3\times10^{-4}$, i.e.\ one
time in $7\,600$). In units of the standard deviation of the fitted law, $|z|>6$ means $|T_{10}|>6/\sqrt{0.8}=6.7$, with probability $5.3\times10^{-5}$ (one in $19\,000$; one-sided $2.7\times10^{-5}$); for the normal it is $2.0\times10^{-9}$. The order of magnitude and the message are unchanged: with fat tails such moves are not
impossible, one per few decades of daily data.
\end{coursenote}

\paragraph{Likelihood with $t$ innovations.} With $\varepsilon_t=\sigma_tz_t$ and the density \eqref{ch13:eq-tstd}, the conditional density of $y_t$ is $f_\nu(\varepsilon_t/\sigma_t)/\sigma_t$, and
\begin{equation}\label{ch13:eq-loglikt}
  \ell(\theta,\nu)=T\log\frac{\Gamma\bigl(\frac{\nu+1}2\bigr)}{\Gamma\bigl(\frac\nu2\bigr)\sqrt{\pi(\nu-2)}}-\frac12\sum_{t=1}^T\log\sigma_t^2-\frac{\nu+1}2\sum_{t=1}^T\log\Bigl(1+\frac{\varepsilon_t^2}{(\nu-2)\sigma_t^2}\Bigr).
\end{equation}
Compared with \eqref{ch13:eq-loglik}, the quadratic penalty $\varepsilon_t^2/\sigma_t^2$ is replaced by a logarithm that grows slowly: an extreme return is punished much less, so the fit is less driven by a few outliers (the
volatilities of the two fits in the case study are nearly identical for this reason and because the parameters of $\sigma_t^2$ are largely determined by the bulk of the data). The degrees of freedom $\nu$ is estimated like any other parameter, most simply by \emph{profiling}: fit the model for a grid of $\nu$ and
compare the maximised log-likelihoods \ts{13}{11}{23:17}\ts{24}{19}{1:19:31}. The tail heaviness is inherited by the returns twice, from $z_t$ and from the fluctuation of $\sigma_t$: by \cref{ch13:thm-kurt} the GARCH$(1,1)$ with $t_\nu$ innovations has kurtosis
$\kappa_z(1-\phi^2)/[1-\phi^2-(\kappa_z-1)\alpha_1^2]$ with $\kappa_z=3+6/(\nu-4)$.

\begin{exercise}[(not from the course): kurtosis with $t$ innovations]\label{ch13:ex-extra3}
For which $(\alpha_1,\beta_1)$ is the kurtosis of a GARCH$(1,1)$ with $t_6$ innovations finite? Compare with the Gaussian case for $\alpha_1+\beta_1=0.95$. Prove that the maximum $\alpha_1$ for an ARCH$(1)$ with fourth moment decreases from $1/\sqrt3$ to $1/\sqrt{\kappa_z}$.
\end{exercise}

\subsection{Diagnostics and model choice}\label{ch13:sec-diag}
After fitting, the standardised residuals $\hat z_t=\hat\varepsilon_t/\hat\sigma_t$ should behave like i.i.d.\ draws of the assumed law with unit variance. Hence \ts{24}{19}{1:14:13}\ts{13}{11}{11:39}:
\begin{itemize}
  \item $\hat z_t$ should be uncorrelated (Ljung--Box/Box--Pierce test of \cref{ch09:eq-bp}) and, the point of the exercise, \emph{$\hat z_t^2$ should be uncorrelated}: the GARCH structure should have absorbed the dependence of the squares.
  \item Normality (or fit of the assumed $t$): normal Q--Q plot, the Jarque--Bera test $\mathrm{JB}=T\bigl[\hat S^2/6+(\hat K-3)^2/24\bigr]\approx\chi^2_2$ (skewness $\hat S$, kurtosis $\hat K$), the Shapiro--Wilk test, the test based on the uniformity of the fitted percentiles $\hat F(y_t)$
        (as in \cref{ch13:sec-clock}) and the Kolmogorov--Smirnov distance between empirical and fitted distribution functions.
  \item Model selection among orders and among innovation laws by the Akaike and Bayesian information criteria (\cref{ch09:eq-ic}); a model with more parameters, such as a free $\nu$, must gain enough log-likelihood to pay for them.
\end{itemize}
(Strictly the residual tests use estimated parameters, so their null distributions are only approximate; not discussed.)

% ------------------------------------------------------------------
\subsection{Case study: EUR/USD returns, ARCH and GARCH fits (2013)}\label{ch13:sec-fxcs}
\begin{casestudy}[EUR/USD returns with ARCH, GARCH and $t$ GARCH models]\label{ch13:cs-fx}
\textbf{Lesson.} Volatility clustering is strongly present, a parsimonious GARCH$(1,1)$ describes it better than a high-order ARCH, the Gaussian assumption is inadequate in the tails
but the volatility estimates are robust to it, and the $t$ law with about $10$ degrees of freedom fits the standardised residuals.
\end{casestudy}
\begin{secondary}[The EUR/USD case study in detail]
\textbf{Data.} The daily returns of the euro in dollars, from the Federal Reserve series used in \cref{ch13:sec-clock} (3\,708 log returns, 1999 to September 2013). The lecture of 2024 shows the same pages of the note for the
dollar--pound series; the numbers below are those of the note, which uses EUR/USD \ts{13}{9}{1:16:58}\ts{24}{19}{1:15:15}.
\textbf{Step 1: dependence in the squares.} The ACF and PACF of the returns show only marginal dependence, those of the \emph{squared} returns have highly significant, slowly decaying autocorrelations and a very significant PACF at low lags, which
suggests an autoregression for the squares, i.e.\ ARCH \ts{13}{9}{1:18:02}\ts{24}{19}{1:15:15}.
\textbf{Step 2: Gaussian ARCH$(p)$ and GARCH$(1,1)$ by maximum likelihood} (R function \texttt{tseries::garch}):
\begin{center}\small
\footnotesize\setlength{\tabcolsep}{3.5pt}
\begin{tabular}{@{}lcccccc@{}}
\toprule
model & $\hat\alpha_0$ & $\hat\alpha_1$ & other & JB & LB $p$ & floor\\
\midrule
ARCH$(1)$ & $4.05\times10^{-5}$ & $0.0285$ & & $626.6$ & $0.79$ & $10.1\%$\\
ARCH$(2)$ & $3.62\times10^{-5}$ & $0.0264$ & $\hat\alpha_2=0.1001$ & $459.2$ & $0.76$ & $9.6\%$\\
ARCH$(10)$ & $2.06\times10^{-5}$ & $0.0035$ & $\sum_{i}\hat\alpha_i=0.504$ & $368.4$ & $0.85$ & $7.2\%$\\
GARCH$(1,1)$ & $1.31\times10^{-7}$ & $0.0278$ & $\hat\beta_1=0.9692$ & $77.1$ & $0.0043$ & $3.3\%$\\
\bottomrule
\end{tabular}
\end{center}
(JB is the Jarque--Bera statistic of the residuals, LB $p$ the Ljung--Box $p$-value for $\hat z^2$; the last column is the annualised lower bound of the fitted volatility, $\sqrt{252\hat\alpha_0}$ for ARCH, $\sqrt{252\hat\alpha_0/(1-\hat\beta_1)}$ for GARCH; the estimated
$\alpha_i$ of ARCH$(10)$ are all significant except the first and the last, and the sum $0.504$ is the persistence.) The volatility paths of the four fits \ts{13}{11}{2:09}\ts{13}{9}{1:19:09}
show the following. (i) Every ARCH model has a \emph{hard lower bound} $\sqrt{\hat\alpha_0}$, lower for higher orders, which is fixed by the sample; (ii) the GARCH$(1,1)$ with three parameters is smoother and explores a much wider range of levels than
the ten-parameter ARCH$(10)$ \ts{13}{9}{1:17:21}\ts{24}{19}{1:17:21}; (iii) its persistence is $\hat\alpha_1+\hat\beta_1=0.9970$: it is ``very close to being non-stationary'', with the long-run variance
$\hat\alpha_0/(1-\hat\phi)=4.34\times10^{-5}$ (daily volatility $0.66\%$, $10.5\%$ annual). Residual checks: the Jarque--Bera test rejects normality for all fits, less strongly for GARCH$(1,1)$; the Ljung--Box tests of the squared residuals no longer reject for the ARCH fits
and marginally still reject ($p=0.004$) for the GARCH$(1,1)$.

\textbf{Step 3: AR mean, $t$ innovations.} A mean model is fitted first: an AR$(7)$ is chosen by AIC, with only the seventh lag significant ($t=2.0$, $R^2=0.004$), and residual standard error $0.00645$. In the AR$(7)$--GARCH$(1,1)$ fit of the
package \texttt{rugarch} \ts{13}{11}{13:44}:
\begin{center}\small
\begin{tabular}{@{}lcccc@{}}
\toprule
innovations & $\hat\omega$ & $\hat\alpha_1$ & $\hat\beta_1$ & persistence\\
\midrule
Gaussian & $1.192\times10^{-7}$ & $0.02777$ & $0.96956$ & $0.99733$\\
$t_{10}$ (fixed) & $1.285\times10^{-7}$ & $0.02920$ & $0.96810$ & $0.99730$\\
\bottomrule
\end{tabular}
\end{center}
The Q--Q plot of the AR$(7)$ residuals (constant variance) bends strongly at both ends, with values up to $6$ standard deviations \ts{13}{11}{14:47}; those of the Gaussian GARCH residuals bend much less but still bend beyond $\pm2.5$ \ts{13}{11}{15:54}, and those of the $t_{10}$ fit are along the line.
The two fitted volatility paths are practically identical (differences at most $0.0002$ in the daily $\sigma_t$, about $0.3$ percentage points in annualised volatility), and the persistence is the same to four digits \ts{13}{11}{16:59}.
\textbf{Step 4: choosing $\nu$.} Fitting for $\nu=4,\dots,30$ with the parameters of the mean and of $\sigma_t^2$ re-estimated each time gives the profile log-likelihood below (the note prints it with the sign of a negative log-likelihood, and plots ``Minus Log Likelihood'');
its maximum is at $\nu=10$ \ts{13}{11}{23:17} (the 2024 lecture reads $9$ from the same plot \ts{24}{19}{1:19:31}: the difference is $0.16$ log-likelihood units).
\begin{center}\small
\begin{tabular}{@{}lccccccccc@{}}
\toprule
$\nu$ & $4$ & $5$ & $6$ & $8$ & $10$ & $12$ & $15$ & $20$ & $30$\\
\midrule
$\ell_{\max}-\ell(\nu)$ & $33.1$ & $15.3$ & $7.1$ & $1.0$ & $0$ & $0.6$ & $2.4$ & $5.6$ & $10.3$\\
\bottomrule
\end{tabular}
\end{center}
\end{secondary}
\begin{secondary}[A confidence interval for $\nu$]
The profile above is one number per $\nu$ from the note. Our own reading, not made in class: since $2[\ell_{\max}-\ell(\nu)]$ is approximately $\chi^2_1$, the values of $\nu$ with $\ell_{\max}-\ell<1.92$ form an approximate $95\%$ confidence interval, $\nu\in[8,14]$; the Gaussian
limit is rejected (for $\nu=30$, $2\times10.3=20.7$, and the normal, $\nu=\infty$, is worse still).
\end{secondary}

\begin{coursenote}[The $9$ degrees of freedom and ``spherical symmetry'']\label{ch13:cn-nine}
In 2024 the lecturer read $9$ degrees of freedom off the plot and gave a heuristic explanation: a $t$ law arises as the projection of a spherically symmetric law in a space of dimension equal to the degrees of freedom, so
$\nu\approx9$ would mean that the volatility is stable over about nine periods \ts{24}{19}{1:20:34}. Take this as a suggestive statement only. What is true and used above is different: $t_\nu$ is a Gaussian mixture with an inverse-chi-squared variance (\cref{ch13:prop-t}), i.e.\ the return is Gaussian with a variance that is random and
slowly changing, which is the same idea as the GARCH itself. In the case study the estimate of $\nu$ is also very flat around its maximum (interval $[8,14]$ above), so $9$ and $10$ are equivalent.
\end{coursenote}

\subsection{Value at risk with $t$ innovations}\label{ch13:sec-var}
The conditional distribution of the next return is $c+\sigma_{t+1}z$ with $z$ standardised. The one-day value at risk at level $p$ (the loss exceeded with probability $p$; the notion of \cref{ch:var}) is
\begin{equation}\label{ch13:eq-var}
  \mathrm{VaR}_p=-\bigl(\hat\mu_{t+1}+q_p\,\hat\sigma_{t+1}\bigr),\qquad q_p=\sqrt{\tfrac{\nu-2}{\nu}}\;t_{\nu,p}\ \ (\text{Student}),\qquad q_p=\Phi^{-1}(p)\ \ (\text{Gaussian}),
\end{equation}
with $t_{\nu,p}$ the $p$-quantile of $T_\nu$ \ts{13}{11}{18:05}. The lecture's observation is that at the usual levels of $2.5\%$ or $5\%$ the two multipliers are barely different, and that the difference appears in the far tail \ts{13}{11}{19:08}:
\begin{center}\small
\begin{tabular}{@{}lcccc@{}}
\toprule
level $p$ & $5\%$ & $2.5\%$ & $1\%$ & $0.1\%$\\
\midrule
Gaussian, $-q_p$ & $1.645$ & $1.960$ & $2.326$ & $3.090$\\
standardised $t_{10}$, $-q_p$ & $1.621$ & $1.993$ & $2.472$ & $3.706$\\
ratio & $0.99$ & $1.02$ & $1.06$ & $1.20$\\
\bottomrule
\end{tabular}
\end{center}
At $5\%$ the $t$ quantile is actually \emph{smaller}, because the standardised $t$ has a more peaked centre (\cref{ch13:fig-tails}). In the case study the $2.5\%$ limits of the AR$(7)$--GARCH$(1,1)$--$t_{10}$ model were exceeded by $107$ of the $3\,708$ returns on the lower side ($2.89\%$) and by $92$ on the upper side
($2.48\%$), against $92.7$ expected on each side \ts{13}{11}{27:33}.
\begin{secondary}[Is 2.89\% compatible with 2.5\%?]
(With this frequency the lower exceedances are consistent with $2.5\%$: the one-sided binomial $p$-value is $0.076$; our computation.)
\end{secondary}

\begin{coursenote}[A detail of the case-study code]\label{ch13:cn-varcode}
In the R code the limits are first computed around the fitted conditional mean and then \emph{overwritten} by limits around the constant long-term mean $\hat c_0/(1-\sum\hat\phi_i)$ of the AR part (the second assignment replaces the first).
The reported exceedance frequencies are therefore those of the second version. Since the conditional mean varies very little ($R^2=0.4\%$) the difference is immaterial, but a VaR should be centred on $\hat\mu_{t+1}$ as in \eqref{ch13:eq-var}.
\end{coursenote}

\begin{finance}[From GARCH to a risk number]
A risk desk fits a GARCH$(1,1)$ (or applies RiskMetrics) to its P\&L, obtains $\hat\sigma_{t+1}$, multiplies by a quantile and reports the VaR; multi-day VaR uses the forecasts of \cref{ch13:sec-forecast} rather than the $\sqrt H$ rule. What can go wrong: the
volatility estimate reacts with a lag to a regime change; the quantile is model dependent far in the tail; and the independence of $z_t$ is assumed. The backtest is the exceedance count of the preceding paragraph.
\end{finance}

% ==================================================================
\section{Stylised facts and extensions}\label{ch13:sec-facts}
The lecture ends with a list of empirical regularities of returns and volatility, attributed to Engle, Bollerslev and Nelson (1994) \ts{24}{19}{1:21:36}\ts{13}{11}{24:19}:
\begin{center}\small
\begin{tabular}{@{}p{3.3cm}p{5.0cm}p{4.4cm}@{}}
\toprule
stylised fact & evidence in this chapter & ingredient of the models\\
\midrule
returns nearly uncorrelated, squares correlated & ACF/PACF of the euro returns and their squares; S\&P~500 squares significant at about $10$ lags & ARCH: $\varepsilon_t$ martingale difference, $\varepsilon_t^2$ an AR (\cref{ch13:prop-mds})\\
volatility clustering & large $\varepsilon_t^2$ follow large ones (GARCH path of \cref{ch13:fig-garchpath}) & $\alpha_1>0$, persistence $\phi$ near $1$\\
fat tails, leptokurtosis & Q--Q plots; kurtosis of the euro residuals; Laplace better than normal & GARCH is a stochastic mixture of Gaussians; $t$ innovations\\
mean reversion and persistence of volatility & $\hat\alpha_1+\hat\beta_1=0.997$, half-life about a year & $\sigma_t^2-\sigma_*^2$ decays at rate $\phi$\\
volatility depends on the estimator and on the time scale & trading against clock time; range estimators & \cref{ch13:sec-clock,ch13:sec-ohlc}\\
\bottomrule
\end{tabular}
\end{center}
The slides close with a list of extensions, only named \ts{13}{11}{28:35}: EGARCH (Nelson, 1991/1992), TGARCH (Glosten, Jagannathan and Runkle, 1993), PGARCH (Ding, Engle and Granger), GARCH-in-mean, and non-Gaussian innovations
(Student $t$, generalised error distribution, GED). The lecturer's advice is to master the specification under Gaussian and $t$ assumptions first; the extensions are reading for a final paper \ts{13}{11}{28:35}\ts{24}{19}{1:21:36}.
\begin{secondary}[The extensions in formulas]
To make the names concrete (the formulas are not in the lectures):
\begin{align}
  \text{GJR/TGARCH:}&\quad \sigma_t^2=\alpha_0+\bigl(\alpha_1+\gamma\,\ind_{\{\varepsilon_{t-1}<0\}}\bigr)\varepsilon_{t-1}^2+\beta_1\sigma_{t-1}^2,\notag\\
  \text{EGARCH:}&\quad \log\sigma_t^2=\omega+\beta\log\sigma_{t-1}^2+\alpha\,|z_{t-1}|+\gamma\,z_{t-1},\notag\\
  \text{GARCH-in-mean:}&\quad y_t=c+\lambda\,\sigma_t^2+\varepsilon_t .\notag
\end{align}
The parameter $\gamma$ captures the \emph{leverage effect}, the observation that a negative return raises future volatility more than a positive return of the same size (a fall in the price of a firm raises its leverage); EGARCH has no positivity constraints because it models $\log\sigma_t^2$. GARCH-in-mean lets the
expected return depend on the risk. All these keep the GARCH$(1,1)$ recursion as their core.
\end{secondary}

\begin{finance}[Realised, implied and stochastic volatility]
Volatility is also \emph{implied} by option prices: the number that, inserted into a formula such as that of \cref{ch:bs}, reproduces the market price. The implied volatility is a forward-looking quantity and typically exceeds the subsequently realised one (a risk premium); the models of this chapter
estimate the \emph{physical} volatility from past returns. Neither year develops implied volatility or \emph{stochastic-volatility} models (where $\sigma_t$ has its own noise, not a function of the past returns as in GARCH); we only mention them, in the words of the course notes, as topics for further study.
\end{finance}
\begin{intuition}[Physicist's summary]
A GARCH$(1,1)$ is a damped memory kernel for the energy of the noise: $\sigma_t^2$ is a leaky integrator of past $\varepsilon^2$ with a relaxation time $-1/\ln\phi$, driven by the noise itself (multiplicative noise).
Everything in the chapter follows from this: the stationary level is the balance of injection and decay, the fat tails are the fluctuations of the local noise temperature, and IGARCH (RiskMetrics) is the same integrator with no leak.
\end{intuition}

% ==================================================================
\section*{Summary of results}
\begin{keyformulas}
\begin{align*}
&\text{Annualised volatility:} && \sqrt N\sigma,\ N=252/52/12;\quad \operatorname{sd}(\hat\sigma)/\sigma\approx1/\sqrt{2(T-1)}\\
&\text{EWMA:} && \tilde\sigma^2_{t+1}=(1-\beta)\hat\sigma^2_t+\beta\tilde\sigma^2_t,\quad n_{\mathrm{eff}}=\tfrac{1+\beta}{1-\beta}\\
& && \text{half-life }\ln2/\ln(1/\beta),\quad\text{mean age }\beta/(1-\beta)\\
&\text{GBM returns:} && R_j\sim N(\mu_*\Delta_j,\sigma^2\Delta_j),\quad \mu_*=\mu-\tfrac12\sigma^2\\
&\text{GBM MLE ($\Delta_j=1$):} && \hat\mu_*=\bar R,\quad \hat\sigma^2=\tfrac1n\textstyle\sum(R_j-\bar R)^2\\
&\text{Efficiencies ($\mu_*=0$):} && \text{close-to-close }1,\ \text{Parkinson }4.9,\ \text{Rogers--Satchell }6.0\\
& && \text{Garman--Klass }7.4,\ \text{with opening jump }8.4\\
&\text{Parkinson:} && \hat\sigma^2=(H-L)^2/(4\ln2)\\
&\text{Garman--Klass:} && \hat\sigma^2=\tfrac12(u-d)^2-(2\ln2-1)c^2\\
&\text{Jump diffusion:} && \Var R=\sigma^2\Delta(1+\lambda\gamma^2),\quad \kappa=3+\tfrac{3\gamma^4\lambda}{\Delta(1+\lambda\gamma^2)^2}\\
&\text{Laplace:} && f=\tfrac1{2b}e^{-|x-\mu|/b},\quad \Var=2b^2,\quad \kappa=6\\
&\text{GARCH$(1,1)$:} && \sigma_t^2=\alpha_0+\alpha_1\varepsilon_{t-1}^2+\beta_1\sigma_{t-1}^2\\
&\text{ARMA form:} && \varepsilon_t^2=\alpha_0+(\alpha_1+\beta_1)\varepsilon_{t-1}^2+u_t-\beta_1u_{t-1}\\
&\text{Stationarity, long run:} && \phi=\alpha_1+\beta_1<1,\quad \sigma_*^2=\tfrac{\alpha_0}{1-\phi}\\
&\text{Forecast:} && \E_t[\sigma^2_{t+h}]=\sigma_*^2+\phi^{h-1}(\sigma^2_{t+1}-\sigma_*^2)\\
& && \text{half-life }\ln(1/2)/\ln\phi\\
&\text{ACF of squares:} && \rho(1)=\tfrac{\alpha_1(1-\alpha_1\beta_1-\beta_1^2)}{1-2\alpha_1\beta_1-\beta_1^2},\quad \rho(k)=\phi^{k-1}\rho(1)\\
&\text{Kurtosis:} && \kappa_\varepsilon=\kappa_z\tfrac{1-\phi^2}{1-\phi^2-(\kappa_z-1)\alpha_1^2}\\
& && \text{(finite iff the denominator is positive)}\\
&\text{Gaussian log-likelihood:} && \ell=-\tfrac12\textstyle\sum_t[\log2\pi+\log\sigma_t^2+\varepsilon_t^2/\sigma_t^2]\\
&\text{Standardised $t$:} && z=\sqrt{\tfrac{\nu-2}{\nu}}T_\nu,\quad\kappa_t=3+\tfrac6{\nu-4}\\
&\text{Value at risk:} && \mathrm{VaR}_p=-(\mu_{t+1}+q_p\sigma_{t+1})\\
&\text{RiskMetrics/IGARCH:} && \alpha_0=0,\ \alpha_1=1-\beta,\ \beta_1=\beta\quad(\phi=1)
\end{align*}
\end{keyformulas}

% ==================================================================
\section*{Exercises}
\addcontentsline{toc}{section}{Exercises}
Standalone exercises follow the section they practise. 2013 PS5 Problems~1--2 are a connected series (Problem~2 generalises the equal sampling intervals of Problem~1 and they share the set-up below), so they stay together here.

\subsection*{From 2013 Problem Set 5 (its Problems 1 and 2 are Problems 3 and 4 of 2024 Problem Set 4)}
In the first two problems $X(t)$ follows $\dd X/X=\mu\dd t+\sigma\dd W$ and the price is sampled on $[0,T]$, so that the log returns $Y_i=\log(X_i/X_{i-1})$ are independent $N(\mu_*h_i,\sigma^2h_i)$ with $\mu_*=\mu-\tfrac12\sigma^2$ (this is given).

\begin{exercise}[PS5 (2013), Problem 1; PS4 (2024), Problem 3: equal sampling intervals]\label{ch13:ex-ps5-1}
Sample at equal increments $h=T/n$, $t_i=ih$.
\begin{enumerate}[(a)]
  \item Show that the maximum-likelihood estimators are $\hat\mu_*=\frac1n\sum Y_i$ and $\hat\sigma^2=\frac1n\sum(Y_i-\hat\mu_*)^2$ (the 2024 version normalises by $h$:
      $\hat\mu_*=\frac1{hn}\sum Y_i$, $\hat\sigma^2=\frac1{hn}\sum(Y_i-h\hat\mu_*)^2$).
  \item Derive the distribution of $\hat\mu_*$, with its mean and variance.
  \item Derive the distribution of $\hat\sigma^2$, with its mean and variance.
  \item Let $n\to\infty$ on the fixed period: what are the limiting distributions of $\hat\mu_{*,n}$ and $\hat\sigma^2_n$?
  \item An estimator sequence $\hat\theta_n$ is weakly consistent if $\lim_n\Prob(|\hat\theta_n-\theta|>\epsilon)=0$ for every $\epsilon>0$ (the 2013 text writes the condition with a small typo; the 2024 hint is Markov's inequality). Which of the two sequences is weakly consistent?
\end{enumerate}
\end{exercise}

\begin{exercise}[PS5 (2013), Problem 2; PS4 (2024), Problem 4: variable sampling intervals]\label{ch13:ex-ps5-2}
Now the increments $h_i>0$ are arbitrary with $\sum_ih_i=T$.
\begin{enumerate}[(a)]
  \item Derive the MLE of $\mu_*$ and its distribution for a fixed set $\{h_i\}$.
  \item Derive the MLE of $\sigma^2$ and its distribution.
  \item With only $n+1$ price points allowed (including $X_0$ and $X_T$), prove that: the MLE of $\sigma^2$
      depends on the spacing, but its \emph{variance} is the same for every spacing; and the MLE of $\mu_*$ is the same, with the same variance, for every spacing. (In the 2024 statement the parts of Problem~4 are mislabelled 3(a)--3(c).)
\end{enumerate}
\end{exercise}

% !TEX root = ../main.tex
\chapter{Event Contracts and a Regulated Exchange}\label{ch:kalshi}
\begin{paracol}{2}

\begin{sources}
\begin{tabularx}{\linewidth}{@{}l l L@{}}
\toprule
\textbf{Lecture} & \textbf{Speaker} & \textbf{Content}\\
\midrule
2024 L20 & T.~Mansour (Kalshi), guest; introduced by P.~Kempthorne & Why an exchange for events; the three-year regulatory road and the lawsuit against the regulator; betting market versus exchange; arbitrage against other assets; cash collateral and margin; who trades (institutions, retail); revenue model; market makers and quoted spreads; career advice. Questions from students and from J.~Xia and V.~Strela.\\
\bottomrule
\end{tabularx}

\smallskip
\textbf{Files.} None: only the video transcript exists for this lecture (no slides, no notes, no
problem set). Everything attributed to the lecture below is taken from the spoken words, and the
tags \ts{24}{20}{0:00}~point into the video. The company facts (volumes, fees, dates, investors)
are the speaker's own statements in late 2024 and are reported as such, without independent
verification. \textbf{Not from the course.} The mathematics of event contracts (Sections~\ref{ch14:sec-binary} to
\ref{ch14:sec-info}) is a reconstruction added for these notes: the lecture contains no formulas.
Every such section says so, and all exercises are marked ``(not from the course)''.
\end{sources}

\section*{Motivation}
Most of this course prices contracts whose payoff is a smooth function of an underlying: a bond, a
forward, a call. The guest lecture is about the most primitive contract one can imagine, one that pays a
fixed amount if a stated event happens and nothing otherwise. The founder of Kalshi describes the idea
that started the company: if there are markets for equities, credit, commodities and currencies, why not a
market on \emph{any question about the future} \ts{24}{20}{4:13}? The lecture itself is a story (the idea, the
regulator, the product, the business), so the notes do two things. Sections~\ref{ch14:sec-story}
and~\ref{ch14:sec-market} report faithfully what was said. Sections~\ref{ch14:sec-binary} to
\ref{ch14:sec-info} supply the mathematics a physicist would want next to the story: an event contract
is a digital option, its price is a (risk-neutral) probability, related contracts must obey
no-arbitrage constraints, a market maker earns a spread and pays for adverse selection, a trader can size a
bet with the Kelly rule, and a market forecast can be scored and compared with the information behind it.
It is a short chapter; the finance content is light and the value is in seeing the one-period model of
\cref{sec:oneperiod} at its smallest scale.

% ==================================================================
\section{The story of the exchange (as told in class)}\label{ch14:sec-story}

\textbf{The speaker.} Tarek Mansour took this course as an MIT student, worked at Goldman Sachs,
Citadel and Bridgewater, and about three years before the lecture co-founded Kalshi with Luana Lopes Lara,
also from MIT \ts{24}{20}{0:00}\ts{24}{20}{1:02}. He credits the course for showing that every
mathematical concept has a financial application, and recalls a term paper on reinforcement learning applied
to the stock market and the use of Markov chains on a trading desk \ts{24}{20}{39:55}\ts{24}{20}{40:57}.

\textbf{The idea.} In 2016 many market participants wanted a hedge against Brexit and had only proxies:
stocks, options, other traditional instruments, indirect and expensive because traded over the counter
\ts{24}{20}{3:12}\ts{24}{20}{4:13}. The proposal: a financial market that prices any event of economic,
social or cultural relevance. He argued that such a market could be the largest of all, simply because it is
broader than any other \ts{24}{20}{4:13}\ts{24}{20}{5:15}.

\textbf{The regulatory road.} Such markets had historically been illegal or unregulated in the US. The
founders phoned essentially every securities and derivatives lawyer in the country; the answers ranged from
laughter to ``the CFTC will never allow it'' \ts{24}{20}{5:15}\ts{24}{20}{6:22}. They kept going
because none of the objections was, in his words, intellectually satisfying: manipulation is a solvable
problem, as it was for stock markets, and the ``gambling'' objection does not distinguish event bets from grain
futures, where some trade to hedge and some to speculate \ts{24}{20}{6:22}. It took two and a half to three years
to get the exchange and the clearinghouse regulated, after which Kalshi launched what he calls the first
legal prediction market in the US \ts{24}{20}{6:22}\ts{24}{20}{7:24}.

\textbf{The lawsuit.} The CFTC did not allow the company to list an election market. Kalshi sued its own
regulator, and, according to the speaker, won ``earlier this fall'' (2024), which made it possible for the
first time in a century to trade the US election legally; the significance is that the ruling widens the universe of
things counted as financial instruments \ts{24}{20}{7:24}\ts{24}{20}{8:26}. He reports more than one
billion (dollars) of volume in the last month, plans to add categories (finance, sports, ``every piece of
news'') and to make the contracts tradable from brokerage accounts such as Schwab, Robinhood and Fidelity
\ts{24}{20}{8:26}\ts{24}{20}{9:32}. The company had raised more than \$100 million from several investors
\ts{24}{20}{10:33}. After the election the volume stayed higher than before it; categories mentioned as growing are
crypto, economics, cabinet appointments, weather and entertainment \ts{24}{20}{21:03}\ts{24}{20}{22:04}.

\begin{coursenote}[Dates and legal status]
Statements such as ``first legal prediction market'', ``first time in 100 years'' and the outcome of the
lawsuit are the speaker's account at the date of the lecture. The notes neither confirm nor extend them;
for anything with legal weight consult primary sources.
\end{coursenote}

% ==================================================================
\section{How the market is organised (as told in class)}\label{ch14:sec-market}

\textbf{Betting market versus exchange.} A student asks how this differs from a betting market
\ts{24}{20}{11:33}. The answer has two parts. First, structure: it is an \emph{exchange} with an order book,
counterparties trade against each other, there is no bookie who sets prices, and an institution hedging
Brexit would never accept bookie prices, while an over-the-counter quote from a bank is negotiated bilaterally
and, in his view, poor \ts{24}{20}{11:33}\ts{24}{20}{12:36}. Second, nature of the risk: rolling a die at a casino
creates an artificial risk for entertainment, whereas Brexit is a \emph{natural} risk that some people carry
and may wish to transfer, by hedging, or increase, by speculating \ts{24}{20}{12:36}\ts{24}{20}{13:39}. That,
he says, is why it is a financial market. He also expects betting markets and exchanges eventually to converge.

\textbf{Arbitrage.} Asked about arbitrage between betting markets and Kalshi, he says Kalshi is the only
legal venue in the US, so the arbitrage that matters is against traditional assets: during the election many
hedge funds traded the election market against crypto and the S\&P~500 \ts{24}{20}{13:39}. Section~\ref{ch14:sec-noarb}
formalises what such cross-asset and cross-contract consistency means.

\textbf{Collateral and margin.} In a futures market a participant posts a small fraction of the notional, say
1\% to 10\%, with the clearinghouse, which then settles gains and calls for more margin after losses
\ts{24}{20}{18:57}. Kalshi is \emph{fully cash-collateralised}: the money is put up front, positions are marked to market
in real time, and a participant cannot lose more than was posted \ts{24}{20}{17:57}\ts{24}{20}{18:57}. Leverage
is expected to be introduced, at the earliest at the end of the following year, under strict regulatory
requirements aimed at preventing anything like 2008; for products with strong correlation some collateral is
already returned \ts{24}{20}{18:57}\ts{24}{20}{20:01}. The economics of the last statement is explored in
\cref{ch14:ex-portfolio-margin}.

\textbf{Who trades and how the business earns.} By volume, 30--40\% is institutional and the rest retail, with an
expected drift toward 50/50. Institutions trade much more volume but can be charged less; retail trades less
and can be charged more, and the speaker calls the institution-first model of the large exchanges a strategic
mistake \ts{24}{20}{22:04}\ts{24}{20}{23:11}. The company was not profitable at the time, by choice: the
strategy is growth first \ts{24}{20}{26:17}\ts{24}{20}{27:20}. Revenue sources listed:
\begin{itemize}
  \item trading fees, of the order of
      \$1--\$2 per \$100 traded;
  \item a separate market-making arm, competing with other firms, among them SIG;
  \item in the future, data, currently free and open \ts{24}{20}{27:20}\ts{24}{20}{28:22}.
\end{itemize}

\textbf{Market makers.} Kalshi imposes obligations on market makers; his example is the election market, where
they had to quote a spread of \$0.01 with at least \$1 million on each side. The tight spread meant that a
very large trade (he mentions a hypothetical \$100 million) could be executed within about one cent of
slippage \ts{24}{20}{28:22}\ts{24}{20}{29:23}. Incentives are given at the start of a market, after which
flow and the bid--ask spread sustain the market maker \ts{24}{20}{29:23}. In the CFTC world, he adds, there is
little competition between exchanges (CME has interest-rate swaps and the S\&P, ICE energy), and he expects
Kalshi to dominate event-related products \ts{24}{20}{29:23}\ts{24}{20}{30:27}.

\textbf{Advice.} The closing part of the lecture is career advice: perseverance, taking more risk, not
over-weighting authority figures, taking courses that connect theory with practice
\ts{24}{20}{30:27}\ts{24}{20}{34:45}\ts{24}{20}{38:51}. It contains nothing to derive and is not
summarised further.

% ==================================================================
\section{Event contracts as digital options}\label{ch14:sec-binary}

\begin{coursenote}[Reconstruction]
From here to Section~\ref{ch14:sec-info} the material is \emph{not from the course}: the lecture states
no formulas. It is written to connect what Mansour describes with the tools of \cref{ch:probability} and
\cref{sec:oneperiod}.
\end{coursenote}

\begin{definition}[Event contract]\label{ch14:def-event}
Let $A$ be an event decided at a settlement date $T$. The \emph{event contract} on $A$ (``Yes'' contract)
pays $\ind_A$ units of cash at $T$: one unit if $A$ occurs and zero otherwise. The ``No'' contract pays
$\ind_{A^c}=1-\ind_A$.
\end{definition}

This is an Arrow--Debreu security for the two-state space $\{A,A^c\}$: in the one-period model of
\cref{sec:oneperiod} it is the asset whose payoff vector is $(1,0)$. If the price is $\pi_A$ today and the cash
account earns rate $r$ over $[0,T]$, the fundamental theorem \cref{thm:ftap} says that absence of arbitrage is
equivalent to the existence of a pricing measure $\Q$ with
\begin{equation}\label{ch14:eq-price}
  \pi_A \;=\; e^{-rT}\,\E^{\Q}[\ind_A] \;=\; e^{-rT}\,\Q(A),
  \qquad 0\le \pi_A\le e^{-rT}.
\end{equation}
Since exchange contracts are cash-collateralised and run for weeks or months at interest rates of a few
per cent, the discount factor is close to $1$ and the quoted price in $[0,1]$ is read as a probability,
$\pi_A\approx \Q(A)$. (Whether the collateral earns interest is a product detail the lecture does not
mention; we ignore it and set $r=0$ from now on.)

\begin{proposition}[Event contract as a digital option]\label{ch14:prop-digital}
Let $S_T$ be a terminal price and $C(K)$ the price of a European call with strike $K$ and payoff
$(S_T-K)^+$. A contract paying $\ind_{\{S_T>K\}}$ (a ``digital'' or ``binary'' call) has price
\[
  D(K)=-\,\frac{\partial C}{\partial K}(K)\quad(r=0),
\]
the negative slope of the call price as a function of strike.
\end{proposition}
\begin{proof}
The call spread that is long one call at $K$ and short one at $K+h$ pays $(S_T-K)^+-(S_T-K-h)^+$, which equals $0$
for $S_T\le K$, $h$ for $S_T\ge K+h$, and lies between $0$ and $h$ in between. Dividing by $h$ and letting $h\to0$
the payoff converges pointwise to $\ind_{\{S_T>K\}}$ (at $S_T=K$ it is a matter of convention). Its price is
$(C(K)-C(K+h))/h\to -C'(K)$.
\end{proof}

So event contracts are the elementary bricks out of which vanilla options are built: a call is the integral of digitals,
$C(K)=\int_K^\infty D(k)\dd k$ (for $r=0$). \Cref{ch:bs} computes $D(K)$ in the Black--Scholes model. What is new
in an event exchange is that the event need not be a function of a traded price at all (weather, a cabinet appointment,
an election), so there is no underlying with which to replicate; the price is then set by trading alone.

\begin{intuition}[Physicist's reading]
A contract on $A$ measures $\ind_A$, a projector in the algebra of observables, and the price is its
``expectation value'' in the state $\Q$. The set of admissible prices for several events is the set of
expectation values of a commuting family of projectors, that is, the convex hull of the vertices $\ind_\omega$. Section~\ref{ch14:sec-noarb} is
just a description of that polytope.
\end{intuition}

\begin{finance}[Risk-neutral versus real-world probability]
By \eqref{ch14:eq-price} the price gives $\Q(A)$, not the true probability $\Prob(A)$. If the
pricing kernel is $m=\dd\Q/\dd\Prob$ then $\Q(A)=\E^{\Prob}[m\ind_A]$. Example:
$\Prob(A)=0.5$ and $m(A)=0.8$, $m(A^c)=1.2$ (the event is a ``good'' state in which one is willing to pay less
for a payoff, e.g.\ a payoff correlated with market gains); then $\E^{\Prob}[m]=1$ and the price is
$0.5\cdot0.8=0.40$, well below the true $0.50$. For an event that is unrelated to aggregate wealth, $m$ is constant,
$\Q=\Prob$, and the price is the physical probability; this is the case in which prediction markets can be
read directly as forecasts. The lecture does not discuss this point; it does say that hedge funds traded the election
contracts against crypto and the S\&P \ts{24}{20}{13:39}, which is exactly a bet that the election
outcome is correlated with those assets.
\end{finance}

% ==================================================================
\section{No-arbitrage constraints between contracts}\label{ch14:sec-noarb}

Let $\Omega$ be a finite set of mutually exclusive, exhaustive outcomes (e.g.\ which candidate wins) and, with $r=0$,
suppose there are contracts $\ind_{\omega}$ for each $\omega\in\Omega$ with prices $\pi_\omega$.

\begin{proposition}[Prices sum to one]\label{ch14:prop-sum}
If there is no arbitrage and contracts on all $\omega\in\Omega$ trade at the same price for buying and selling, then $\sum_\omega\pi_\omega=1$.
More generally, for events $A\subseteq B$: $0\le\pi_A\le\pi_B\le1$, and for disjoint $A,B$: $\pi_{A\cup B}=\pi_A+\pi_B$.
\end{proposition}
\begin{proof}
A portfolio holding one of every contract pays $\sum_\omega \ind_\omega=1$ in every state. If its cost $\sum\pi_\omega$
were less than $1$ it would be an arbitrage (buy it), if more than $1$ the reverse portfolio would be (sell it). The
other statements follow the same way: $\ind_B-\ind_A\ge0$, so its price cannot be negative, and
$\ind_{A\cup B}=\ind_A+\ind_B$ by linearity.
\end{proof}

With a bid--ask spread the equality becomes inequalities. Let $a_\omega$ be the best ask (price at which one buys) and $b_\omega$ the best bid.
Buying all contracts costs $\sum a_\omega$ and pays exactly $1$; selling all yields $\sum b_\omega$. Hence, before fees,
\begin{equation}\label{ch14:eq-arbband}
  \sum_{\omega} b_\omega \;\le\; 1 \;\le\; \sum_{\omega} a_\omega
\end{equation}
is necessary for absence of arbitrage. Similarly the ``No'' contract on $A$ is the complement of the ``Yes'':
buying No at $a^{\mathrm{No}}$ is equivalent to selling Yes at price $1-a^{\mathrm{No}}$, so on a well-organised book
$b^{\mathrm{Yes}}=1-a^{\mathrm{No}}$ and $a^{\mathrm{Yes}}=1-b^{\mathrm{No}}$: Yes and No are two views of the same order book.

\begin{example}[Three candidates]\label{ch14:ex-three}
Asks of $0.48,\ 0.36,\ 0.19$ on three exhaustive candidates sum to $1.03\ge1$ (allowed), while bids of
$0.47,\ 0.35,\ 0.18$ sum to $1.00$. Suppose instead the asks were $0.45,\ 0.36,\ 0.16$: they sum to $0.97<1$. Buying one of each costs
$0.97$ and pays $1$ for sure, a profit of $3$ cents per set, before fees. Fees of the order of \$1--\$2 per \$100 traded
\ts{24}{20}{27:20} (which is $1$ to $2$ cents per dollar of price paid, so $\approx 1$ to $2$ cents on a $0.97$ cost) would eat
most of this margin: this is why the prices need not be exactly consistent, and an arbitrage is real only if it survives
costs.
\end{example}

\begin{proposition}[Monotone thresholds and the implied density]\label{ch14:prop-thresh}
Let $X$ be a real quantity (an inflation print, say) and $P(k)$ the price of ``$X>k$''. Then $P$ is non-increasing in $k$,
and $P(k)-P(k')=\Q(k<X\le k')\ge0$ for $k<k'$. A market with a ladder of thresholds $k_1<\dots<k_n$ therefore
implies a discrete risk-neutral distribution for $X$.
\end{proposition}
\begin{proof}
$\{X>k'\}\subseteq\{X>k\}$; apply \cref{ch14:prop-sum}.
\end{proof}

\begin{remark}
Applying this to $X=S_T$ recovers \cref{ch14:prop-digital}: $P(K)=D(K)=-C'(K)$, so $-P'(k)$ is the risk-neutral
density of $S_T$, a fact used in option-pricing theory for the implied distribution.
\end{remark}

\begin{finance}[Portfolio margin and correlated products]\label{ch14:fin-margin}
With full cash collateral a seller of a contract that pays $1$ must lock $1-\pi$ on top of the premium $\pi$ received, so the locked
amount per sold contract equals the maximum possible payout, $1$. Consider sold contracts on two disjoint,
exhaustive events $A$ and $B$ (so $B=A^c$): exactly one of the two contracts pays, so the joint liability is $1$, not $2$.
Collateral of $1$ suffices, and $1$ can be released. This is the elementary form of the statement, heard in the lecture, that where two
products are strongly correlated some collateral is returned \ts{24}{20}{20:01}. (The lecture gives no formula; the general
rule is that the collateral of a portfolio is its worst-case loss $-\min_\omega \text{payoff}(\omega)$ minus premiums, a linear
program over the finite state space.)
\end{finance}

% ==================================================================
\section{Market making, spread and adverse selection}\label{ch14:sec-mm}

The lecture says that market makers commit to a tight spread (one cent, with \$1 million a side, in the election market) and are paid
by flow and by the spread \ts{24}{20}{28:22}\ts{24}{20}{29:23}. A minimal model that shows what this spread has to cover is
the following (a one-period version of the Glosten--Milgrom model, \emph{not from the course}).

\paragraph{Model.} An event $A$ has prior probability $\tfrac12$ (and $r=0$), so the outcome $V=\ind_A\in\{0,1\}$.
A market maker quotes a bid $b$ and an ask $a$ and faces one arriving trader. With probability $\alpha$ the trader is
\emph{informed}: he knows $V$ and buys if $V=1$, sells if $V=0$. With probability $1-\alpha$ the trader is
uninformed and buys or sells with probability $\tfrac12$ each. The market maker breaks even on each side (competition), that
is, sets each quote equal to the conditional expectation of $V$ given the trade.

\begin{proposition}[Spread equals the informed fraction]\label{ch14:prop-gm}
In this model $a=\tfrac{1+\alpha}{2}$, $b=\tfrac{1-\alpha}{2}$, so the spread is $a-b=\alpha$, and the market maker's expected profit is zero.
\end{proposition}
\begin{proof}
$\Prob(\text{buy}\mid V{=}1)=\alpha+\tfrac{1-\alpha}{2}$ and $\Prob(\text{buy}\mid V{=}0)=\tfrac{1-\alpha}{2}$. By Bayes with prior
$\tfrac12$, $a=\E[V\mid\text{buy}]=\dfrac{\alpha+(1-\alpha)/2}{\alpha+(1-\alpha)}=\dfrac{1+\alpha}{2}$. By symmetry
$b=\E[V\mid\text{sell}]=\tfrac{1-\alpha}2$.
\end{proof}

\begin{example}
For $\alpha=0.1$: $a=0.55$, $b=0.45$. Against an uninformed trader (probability $0.9$) the market maker gains, per unit, half the spread, $0.05$, on average
(she sells at $0.55$ a contract worth $0.5$ in expectation). Against an informed trader (probability $0.1$) she
loses $0.45$: she sells at $0.55$ a contract that pays $1$. The expected profit is $0.9(0.05)-0.1(0.45)=0$, confirming
\cref{ch14:prop-gm}. If the required spread is one cent, as in the election market, the model says the market maker
believed that $\alpha\approx0.01$: at most one trade in a hundred was against someone who knew the answer, and
large trades could be absorbed because the fraction of informed flow was so small.
\end{example}

The model captures only adverse selection. Two further costs, both standard, complete the picture:
\emph{inventory risk} (a market maker who has sold $n$ Yes contracts is short $n\ind_A$ and carries the variance
$n^2p(1-p)$ until she hedges, by buying back or by trading a correlated asset, as the funds arbitraging the election
against the S\&P did \ts{24}{20}{13:39}) and \emph{fees}. A
market maker's spread over a lifetime of the contract must cover both, which is why exchanges incentivise it
at the start of a market, when informed fraction is high and volume low \ts{24}{20}{29:23}.

\begin{intuition}[Where the spread comes from]
The spread is not a fee for making the trade, but insurance against the person on the other side knowing more than you.
Doubling the informed fraction doubles the spread. Near settlement, when everyone knows the answer, $\alpha\to1$ and
both quotes converge to $0$ and $1$: the market ``goes to the boundary'' and quoting a tight two-sided market stops being possible.
\end{intuition}

\begin{figure}[ht]
\centering
\begin{tikzpicture}
\begin{axis}[width=0.72\linewidth,height=5.2cm,xlabel={informed fraction $\alpha$},ylabel={quote},
  xmin=0,xmax=1,ymin=0,ymax=1,legend pos=north west,legend cell align=left,grid=both,grid style={black!10}]
\addplot[domain=0:1,thick,samples=2,color=red!70!black]{(1+x)/2};
\addplot[domain=0:1,thick,samples=2,color=blue!70!black]{(1-x)/2};
\legend{ask $a$,bid $b$}
\end{axis}
\end{tikzpicture}
\caption{Quotes of the break-even market maker of \cref{ch14:prop-gm}. The gap between the lines is the informed fraction $\alpha$.}
\label{ch14:fig-quotes}
\end{figure}

% ==================================================================
\section{Sizing a bet: the Kelly criterion}\label{ch14:sec-kelly}

A participant who believes $\Prob(A)=p$ while the contract trades at $c<1$ must decide how much to buy. Let wealth be $W$ and
invest a fraction $f\in[0,1)$ of it in Yes contracts at price $c$ (so the number of contracts is $fW/c$). Ignoring fees, wealth after
settlement is
\[
  W\bigl(1+fb\bigr)\ \text{ if }A,\qquad W(1-f)\ \text{ if }A^c,\qquad b=\frac{1-c}{c},
\]
where $b$ is the net odds per unit staked.

\begin{proposition}[Kelly fraction for a binary contract]\label{ch14:prop-kelly}
The expected log growth $g(f)=p\ln(1+fb)+(1-p)\ln(1-f)$ is strictly concave and maximised at
\[
  f^\star=\frac{p-c}{1-c}\qquad(\text{if }p>c;\ f^\star=0\text{ otherwise}).
\]
\end{proposition}
\begin{proof}
$g'(f)=\dfrac{pb}{1+fb}-\dfrac{1-p}{1-f}=0$ gives $pb(1-f)=(1-p)(1+fb)$, i.e.\ $f(pb+(1-p)b)=pb-(1-p)$, so
$f^\star=p-\dfrac{1-p}{b}=p-\dfrac{(1-p)c}{1-c}=\dfrac{p-c}{1-c}$. Concavity: $g''<0$ since both logarithms are concave.
\end{proof}

\begin{example}\label{ch14:ex-kelly}
Let $c=0.55$ and $p=0.65$: $b=0.818$, $f^\star=0.10/0.45=0.222$ and $g(f^\star)=0.0206$ per bet (about $2\%$ of wealth per bet, in
the log sense). Staking half of that, $f=0.111$, still gives $g=0.0153$ (three quarters of the optimal growth) with a much smaller
variance, whereas staking double, $f=0.444$, gives $g=-0.0041<0$: an \emph{over-bet} with a positive edge loses money in the long run.
(Numbers checked numerically.)
\end{example}

\begin{figure}[ht]
\centering
\begin{tikzpicture}
\begin{axis}[width=0.75\linewidth,height=5.4cm,xlabel={fraction staked $f$},ylabel={$g(f)$},
  xmin=0,xmax=0.5,ymin=-0.03,ymax=0.03,grid=both,grid style={black!10},axis lines=left]
\addplot[domain=0:0.5,samples=80,thick,color=blue!70!black]{0.65*ln(1+x*0.81818)+0.35*ln(1-x)};
\draw[dashed] (axis cs:0,0) -- (axis cs:0.5,0);
\draw[dotted] (axis cs:0.2222,-0.03) -- (axis cs:0.2222,0.0206);
\node[anchor=south west,font=\small] at (axis cs:0.23,0.0206) {$f^\star=0.222$};
\end{axis}
\end{tikzpicture}
\caption{Expected log growth for $c=0.55$, $p=0.65$ (\cref{ch14:ex-kelly}). The curve crosses zero at about twice the Kelly fraction.}
\label{ch14:fig-kelly}
\end{figure}

\begin{finance}[Edge and price]
Only the gap $p-c$ matters, scaled by the maximum loss $1-c$: an event contract is the cleanest setting for Kelly, because the outcome
is a genuine Bernoulli variable and the odds are read directly from the price. The estimate of $p$ is itself uncertain, and
$f^\star$ is sensitive to it ($\partial f^\star/\partial p=1/(1-c)$), which is why practitioners use a fraction of Kelly.
\end{finance}

% ==================================================================
\section{Scoring rules, calibration and aggregation}\label{ch14:sec-info}

\paragraph{Brier score.} A forecaster reports $q\in[0,1]$ for an event with outcome $\ind_A$, and is scored by
$(q-\ind_A)^2$ (lower is better). If the true probability is $p$, the expected score is
\[
  \E\,(q-\ind_A)^2 = p(1-q)^2+(1-p)q^2 = p(1-p)+(q-p)^2 .
\]
It is minimised at $q=p$: the Brier score is a \emph{proper} scoring rule, honesty is optimal. For $p=0.6$ the expected scores are
$0.24$ at $q=0.6$, $0.25$ at $q=0.5$ and $0.28$ at $q=0.8$ (verified numerically). The irreducible term $p(1-p)$ is the variance of the event;
only the excess $(q-p)^2$ measures forecasting error. A market price is a forecast, and it is \emph{calibrated} if among all events priced at $q$ a fraction $q$ occur;
comparing a collection of prices with outcomes in this way is how one tests the claim, implicit in the lecture, that an exchange for events makes the
world ``smarter about the future'' \ts{24}{20}{32:38}.

\paragraph{Aggregation of information.} Suppose each of $n$ traders holds an independent private signal $s_i$ about $A$, with likelihood
ratio $\lambda_i=\Prob(s_i\mid A)/\Prob(s_i\mid A^c)$. By Bayes' rule in odds form,
\[
  \frac{\Prob(A\mid s_1,\dots,s_n)}{\Prob(A^c\mid s_1,\dots,s_n)}=\frac{\Prob(A)}{\Prob(A^c)}\prod_i\lambda_i,
\]
so in \emph{log-odds} the evidence adds: $\operatorname{logit}p=\operatorname{logit}p_0+\sum_i\ln\lambda_i$, with $\operatorname{logit}p=\ln\frac p{1-p}$.
With prior $\tfrac12$ and $\lambda=(2,3,1.5)$ the posterior odds are $9$ and $p=0.9$, higher than any single trader would say
(each alone would have $2/3$, $3/4$, $3/5$). The purpose of a market is to make prices carry the pooled log-odds, and
trading is the mechanism: each trader whose belief differs from the price trades, and by \cref{ch14:prop-kelly}
the size of the trade grows with the edge, so the price moves toward the belief-weighted pool. That this works
in practice is an empirical claim, not a theorem: it needs enough liquidity, the market-maker obligations of the previous sections, and traders who cannot manipulate at little cost
(the ``manipulation'' objection heard by the founders \ts{24}{20}{6:22}).

\begin{keyformulas}
\begin{itemize}
  \item Event contract: pays $\ind_A$; price $\pi_A=e^{-rT}\Q(A)\approx\Q(A)$ (\cref{ch14:eq-price}); digital call $D(K)=-\partial C/\partial K$.
  \item No arbitrage: $\sum_\omega\pi_\omega=1$; with a spread, $\sum b_\omega\le1\le\sum a_\omega$ (\cref{ch14:eq-arbband}); thresholds: $P(k)$ non-increasing.
  \item Yes and No: $b^{\mathrm{Yes}}=1-a^{\mathrm{No}}$.
  \item Market maker with informed fraction $\alpha$ (prior $\tfrac12$): $a=\tfrac{1+\alpha}2$, $b=\tfrac{1-\alpha}2$, spread $\alpha$.
  \item Kelly for a contract bought at $c$ with belief $p$: $f^\star=(p-c)/(1-c)$.
  \item Brier: $\E[(q-\ind_A)^2]=p(1-p)+(q-p)^2$.
  \item Independent evidence: $\operatorname{logit}p=\operatorname{logit}p_0+\sum_i\ln\lambda_i$.
  \item From the lecture: cash collateral $100\%$, fee of order $\$1$--$\$2$ per $\$100$, election quotes $\$0.01$ wide with $\$1$M per side.
\end{itemize}
\end{keyformulas}

% ==================================================================
\section*{Exercises}
\addcontentsline{toc}{section}{Exercises}

\begin{exercise}[not from the course]\label{ch14:ex-digital}
In the Black--Scholes model with $r=0$, $C(K)=S\,\Phi(d_1)-K\,\Phi(d_2)$ (\cref{ch:bs}), where $\Phi$ is the standard normal
distribution function. Show that $D(K)=-C'(K)=\Phi(d_2)$, and interpret it as $\Q(S_T>K)$.
\end{exercise}

\begin{exercise}[not from the course]\label{ch14:ex-arb}
Four candidates have asks $0.40,\ 0.30,\ 0.15,\ 0.10$ and bids $0.38,\ 0.28,\ 0.13,\ 0.08$. Is there an arbitrage without fees? With a fee of $1$ cent per
contract bought, how much must the sum of asks fall below $1$ for buying one of each to be profitable?
\end{exercise}

\begin{exercise}[not from the course]\label{ch14:ex-gm}
Generalise \cref{ch14:prop-gm} to a prior $\Prob(V=1)=\pi_0\neq\tfrac12$ and derive the ask and bid. Show that the spread
is $\alpha$ times a factor depending on $\pi_0$, and identify the factor.
\end{exercise}

\begin{exercise}[not from the course]\label{ch14:ex-kellyfee}
A trader believes $p=0.70$ and can buy Yes at $c=0.60$, paying a fee of $0.01$ per contract at purchase. Compute the Kelly fraction with and without the fee (treat the fee as an
increase of the price, $c\to0.61$). By how much does the fee reduce the optimal log growth $g(f^\star)$?
\end{exercise}

\begin{exercise}[not from the course]\label{ch14:ex-portfolio-margin}
Two independent events $A,B$ with $\Prob=\tfrac12$ each. Under full cash collateral, a trader sells Yes on both at price $0.5$. Compute the
locked collateral and the maximum loss of the portfolio. Do the same when $B=A^c$ (compare \cref{ch14:fin-margin}). Explain why the
difference in the required collateral is due to the correlation.
\end{exercise}

\begin{exercise}[not from the course]\label{ch14:ex-brier}
Show that the logarithmic score $-\ln q$ (if $A$) and $-\ln(1-q)$ (if $A^c$) is also proper, and compute its expected value for true probability $p$, in terms of the binary entropy and the Kullback--Leibler divergence.
\end{exercise}

% !TEX root = ../main.tex
\chapter{Black--Scholes Formula and Risk-Neutral Valuation}\label{ch:bs}
\begin{paracol}{2}

\begin{sources}
\begin{tabularx}{\linewidth}{@{}l l L@{}}
\toprule
\textbf{Lecture} & \textbf{Speaker} & \textbf{Content}\\
\midrule
2024 L21 & V.~Strela & Update of the yield curve after the September 2024 Fed meeting; two-horse race (the bookie who sets odds from the pool); forwards, calls, puts; one-step replication of a forward and of a call; risk-neutral probability; lognormal dynamics, why $\dd W\sim\sqrt{\dd t}$, heuristic It\^o formula; replication argument and the Black--Scholes PDE; boundary conditions and the closed-form formula; the risk-neutral density; put--call parity on Bloomberg screens (AAPL, IBM) and the implied-volatility smile; the digital option and the call spread.\\
2013 L19 & V.~Strela & The same material (nearly the same slides), with the discrete one-step market in more detail (bond units $b$, the interest-rate quiz on the forward), suggestions for final-paper topics (heat equation, formula by expectation), an IBM option screen of March 2006, put--call parity, and the digital option as an interview question.\\
\bottomrule
\end{tabularx}

\smallskip
\textbf{Files.} Slides \file{mit18_642_f24_lec21.pdf} (26 slides) and 2013 slides \file{MIT18_S096F13_lecnote19.pdf}
(25 slides); no problem set, case study or code accompanies this lecture in either year.
\textbf{Not recorded.} 2024 L22, \emph{The Spectrum of Systematic Trading Strategies in Liquid Instruments}
(R.~Garon, Millennium Management), has no video and no slides: the guest did not grant permission to publish them, so there
is nothing to summarise here; the next recorded lecture is L23 (machine learning, \cref{ch:ml}).
\textbf{Exercises.} None of the course problem sets deals with this lecture; the exercises (each after the section it practises)
are all marked ``(not from the course)''.
\end{sources}

\section*{Motivation}
Chapter~\ref{ch:linalg} priced a call in a one-period market with two states and found that the real-world
probability of ``up'' did not matter (\cref{ex:call}). This chapter takes that idea to its limit. The
question is: \emph{what is the price today of a payoff $f(S_T)$ that is fixed by the price of a stock at a
future date $T$?} The answer, due to Black, Scholes and Merton (1973), is
\begin{itemize}
  \item \emph{model-dependent but preference-free}: once one specifies how the stock moves, the price is
        determined, and neither the expected return of the stock nor the risk appetite of investors enters;
  \item \emph{constructive}: the price is the cost of a self-financing portfolio in stock and cash that
        \emph{replicates} the payoff, and that portfolio is also the hedge;
  \item \emph{an expectation}: the price equals a discounted expectation of the payoff under a fictitious
        \emph{risk-neutral} measure $\Q$, under which every traded asset grows at the risk-free rate;
  \item \emph{a PDE}: the same price solves the Black--Scholes equation, a heat equation in disguise.
\end{itemize}
The lecture is deliberately placed \emph{before} the stochastic-calculus lectures
(\cref{ch:stochcalc,ch:sde}): It\^o's formula is introduced heuristically, in the way a physicist would expand
a function of a Brownian path, and the rigorous statements are deferred. Read this chapter as the
\emph{application} that motivates the formalism of the next ones, and as the continuous-time counterpart of
\cref{sec:oneperiod}. Sections~\ref{ch15:sec-bookie}--\ref{ch15:sec-onestep} are the discrete world,
\cref{ch15:sec-tree,ch15:sec-lognormal} the passage to continuous time,
\cref{ch15:sec-pde,ch15:sec-rn,ch15:sec-formula} the two routes to the formula, and
\cref{ch15:sec-parity,ch15:sec-digital,ch15:sec-greeks,ch15:sec-smile} what one does with it in practice.

% ==================================================================
\section{Where the lecture starts: the yield curve}\label{ch15:sec-curve}
As promised in the first lecture, the class opens with an update of the US Treasury curve of \cref{fig:yieldcurve}.
Slide~2 overlays the curve of 3 September 2024 (just before the course began) on that of 15 November 2024
\ts{24}{21}{0:00}. The Fed had cut in the meantime by 75 basis points in total (in the calendar: 50~bp in September and
25~bp in November, a detail not spoken in class), and the short end of the curve moved down by about that
amount while the long end moved by only about 50~bp \ts{24}{21}{1:07}: not a parallel shift but a steepening.

The lecturer then puts the current shape next to three historical snapshots \ts{24}{21}{2:14}: the
curves of 2007 (just before the crisis) and 2019 were both \emph{inverted} and a recession followed, while by
September 2021 the curve had a normal shape and the recession had been short. He calls this ``non-scientific
prediction'': a serious analysis would decompose the curves with principal components (\cref{ch:pca}), and the
term structure is, in his words, in the centre of finance because the discount factors of \cref{sec:discounting}
are read off it. For the present chapter only one consequence matters: discounting uses a rate $r$, which the
lecture takes \emph{constant}. Everything that follows is an exercise in a single-rate world; making $r$ stochastic
is the subject of the short-rate and HJM models of \cref{ch:hjm}.

% ==================================================================
\section{Setting the odds: a bookie's version of hedging}\label{ch15:sec-bookie}
The lecture introduces risk-neutral valuation with an example that has no mathematics in it
\ts{24}{21}{4:21}\ts{13}{19}{0:00}. A bookie knows that horse~1 wins a race with probability $20\%$ and horse~2 with
probability $80\%$. The public does not know this, and bets \$10\,000 on horse~1 and \$50\,000 on horse~2
\ts{24}{21}{5:24}\ts{13}{19}{1:01}; the pool is \$60\,000 (a commission, say $1\%$, is what the bookie keeps in
the end and is ignored below). The bookie must announce the odds ``$k$ to one'': a winning stake of $s$ returns
$s+ks$.

\paragraph{Odds from the ``true'' probabilities: $4$ to $1$.} Odds $4:1$ equal the true ratio $0.8:0.2$. If horse~1
wins, the bookie collects \$60\,000 and pays $10\,000+4\cdot10\,000=\$50\,000$: profit \$10\,000. If horse~2 wins, he pays
$50\,000+\tfrac14\cdot50\,000=\$62\,500$: loss \$2\,500 \ts{24}{21}{6:28}. The expected profit is
$0.2\cdot10\,000+0.8\cdot(-2\,500)=0$, but only on average: the bookie carries the risk of the one race that is run.

\paragraph{Odds from the pool: $5$ to $1$.} Now ignore the true probabilities and set the odds by the money
staked, $50\,000/10\,000=5$. If horse~1 wins the bookie pays $10\,000+5\cdot10\,000=\$60\,000$; if horse~2 wins he pays
$50\,000+\tfrac15\cdot50\,000=\$60\,000$. Either way he pays exactly what he collected \ts{24}{21}{8:37}\ts{13}{19}{3:06}: the book is
\emph{hedged} and any commission is riskless profit. The pool proportions
\[
  q_1=\frac{10\,000}{60\,000}=\frac16,\qquad q_2=\frac{50\,000}{60\,000}=\frac56
\]
play the role of a \emph{risk-neutral} probability: whichever the true probabilities $(p_1,p_2)$, the bookie's
expected profit with odds set from $q$ is zero \emph{and its variance is zero}. The price of a bet on horse~1 (its stake,
relative to the payout) is $q_1=1/6$, not $p_1=1/5$.

\begin{intuition}[What the example teaches]
The bookie prices by \emph{what is already priced by the market} (the pool) and not by what he believes. The
numbers $q_i$ are not forecasts: they are the prices of the two Arrow--Debreu securities of \cref{def:onepmodel}
(``\$1 if horse $i$ wins'') normalised to sum to one. Three lessons carry over to derivatives \ts{24}{21}{9:42}:
\begin{itemize}
  \item an arbitrage-free price is the cost of a portfolio that reproduces the payoff in every outcome;
  \item that cost can be written as an expectation, but with the ``pool'' probabilities $q$ and not the
      true ones;
  \item the pool probabilities need not coincide with reality, and so a model of real-world
      probabilities is not needed to price and hedge.
\end{itemize}
\end{intuition}

\begin{coursenote}[Transcript slips]
In the 2013 transcript the bookie's payout on the second horse is given as ``\$12.25'' for the quarter of \$50\,000;
it is \$12\,500 (so the total payout is \$62\,500 and the loss \$2\,500, as on the slide). Also in 2013,
to the quiz ``which interest rate does the forward strike imply?'' a student answers ``2.5'', which the lecturer first
repeats and then corrects on the fly: $80\to88.41$ is about $10\%$ over the two years of \cref{ch15:sec-payoffs}, i.e.\ $5\%$ per year
continuously compounded, $80\,e^{0.05\cdot2}=88.41$ \ts{13}{19}{16:42}.
\end{coursenote}

\begin{exercise}[not from the course: the bookie]\label{ch15:ex-bookie}
Pools of $A$ on horse~1 and $B$ on horse~2. Show that odds set to $B/A$ to $1$ on horse~1 (and $A/B$ to $1$ on horse~2) give the bookie the same profit whichever horse wins, and identify the ``risk-neutral probabilities'' $q_i$.
Compute the standard deviation of the bookie's profit in the example if the odds are set from the true probabilities $(0.2,0.8)$, and compare with the pool odds.
\end{exercise}

% ==================================================================
\section{Forwards, calls and puts}\label{ch15:sec-payoffs}
Throughout, $S_t$ is the price at time $t$ of a stock (no dividends), $T$ the \emph{expiry} of the contract, $K$ its
\emph{strike} and $\tau=T-t$ the time to expiry. The three simplest derivatives pay, at $T$ \ts{24}{21}{10:47}\ts{13}{19}{4:06},
\[
  \text{forward: } S_T-K,\qquad \text{call: } (S_T-K)^+,\qquad \text{put: } (K-S_T)^+ ,
\]
with $x^+=\max(x,0)$. A \emph{forward contract} is an obligation to buy at $T$ at the price $K$ fixed today; $K$ is chosen so that
the contract costs nothing to enter \ts{24}{21}{11:54}\ts{13}{19}{5:07}. A \emph{call option} is the right, not the obligation, to
buy at $K$, and a \emph{put option} the right to sell at $K$ \ts{24}{21}{12:55}. The call is insurance against
the stock going up if one is short, and the put is insurance against a fall: its owner is protected below $K$. Unlike a forward,
an option costs a premium today, because it can only win. A call (put) is \emph{in the money} if $S>K$ ($S<K$), \emph{at the money} if $S=K$, and
\emph{out of the money} otherwise; the lecture always means \emph{European} options, exercisable only at $T$
(American options, exercisable any time, have no closed form and are computed on trees or finite-difference grids).

\begin{figure}[ht]
\centering
\begin{tikzpicture}
\begin{axis}[name=A,width=0.48\linewidth,height=6.2cm,xmin=0,xmax=200,ymin=-90,ymax=125,axis lines=left,
  xlabel={$S$},tick label style={font=\footnotesize},label style={font=\small},title={\small call and forward},
  title style={yshift=-3pt},legend style={at={(0.97,0.03)},anchor=south east,font=\scriptsize,draw=none,fill=none}]
  \addplot[thick,black!45] coordinates {(0,0) (80,0) (200,120)};
  \addplot[very thick,defcol] coordinates {(2,0.0000) (8,0.0000) (14,0.0000) (20,0.0000) (26,0.0004) (32,0.0077) (38,0.0574) (44,0.2499) (50,0.7561) (56,1.7730) (62,3.4588) (68,5.8914) (74,9.0629) (80,12.9014) (86,17.3009) (92,22.1466) (98,27.3322) (104,32.7679) (110,38.3825) (116,44.1224) (122,49.9485) (128,55.8331) (134,61.7571) (140,67.7071) (146,73.6744) (152,79.6530) (158,85.6391) (164,91.6300) (170,97.6241) (176,103.6203) (182,109.6178) (188,115.6161) (194,121.6151) (200,127.6144)};
  \addplot[very thick,dashed,thmcol] coordinates {(0,-88.41) (200,111.59)};
  \addplot[very thick,dotted,intcol] coordinates {(2,-77.9967) (8,-71.9967) (14,-65.9967) (20,-59.9967) (26,-53.9967) (32,-47.9967) (38,-41.9967) (44,-35.9967) (50,-29.9967) (56,-23.9967) (62,-17.9967) (68,-11.9967) (74,-5.9967) (80,0.0033) (86,6.0033) (92,12.0033) (98,18.0033) (104,24.0033) (110,30.0033) (116,36.0033) (122,42.0033) (128,48.0033) (134,54.0033) (140,60.0033) (146,66.0033) (152,72.0033) (158,78.0033) (164,84.0033) (170,90.0033) (176,96.0033) (182,102.0033) (188,108.0033) (194,114.0033) (200,120.0033)};
  \legend{payoff $(S-80)^+$,{call price $C(t,S)$},forward payoff $S-K$,{forward price $F(t,S)$}}
\end{axis}
\begin{axis}[at={(A.east)},anchor=west,xshift=1.1cm,width=0.48\linewidth,height=5.6cm,xmin=0,xmax=200,ymin=-5,ymax=80,axis lines=left,
  xlabel={$S$},tick label style={font=\footnotesize},label style={font=\small},title={\small put},title style={yshift=-3pt},
  legend style={at={(0.97,0.97)},anchor=north east,font=\scriptsize,draw=none,fill=none}]
  \addplot[thick,black!45] coordinates {(0,80) (80,0) (200,0)};
  \addplot[very thick,defcol] coordinates {(2,70.3870) (8,64.3870) (14,58.3870) (20,52.3870) (26,46.3874) (32,40.3947) (38,34.4444) (44,28.6369) (50,23.1431) (56,18.1600) (62,13.8458) (68,10.2784) (74,7.4499) (80,5.2884) (86,3.6879) (92,2.5336) (98,1.7192) (104,1.1549) (110,0.7694) (116,0.5094) (122,0.3355) (128,0.2201) (134,0.1441) (140,0.0941) (146,0.0614) (152,0.0400) (158,0.0261) (164,0.0170) (170,0.0111) (176,0.0073) (182,0.0048) (188,0.0031) (194,0.0021) (200,0.0014)};
  \legend{payoff $(80-S)^+$,{put price $P(t,S)$}}
\end{axis}
\end{tikzpicture}
\caption{Payoffs at expiry (grey or dashed) and prices two years before expiry (coloured), as functions of the stock price,
for $K=80$, $r=5\%$, $\sigma=20\%$, $\tau=2$ years (call, put) and $K=88.41=80e^{0.1}$ (forward), the numbers of slides~3--6.
The volatility is not stated on the slides; $20\%$ reproduces the digital curve of slide~26 (\cref{ch15:sec-digital}).
The price curves are the Black--Scholes formula derived below; the forward price is $S-Ke^{-r\tau}$.}\label{ch15:fig-payoffs}
\end{figure}

\Cref{ch15:fig-payoffs} anticipates the goal of the lecture \ts{24}{21}{13:55}: to compute the curved lines, the prices today
of the options, and the straight line for the forward. It also states the three messages of the lecture, which are exactly
the content of the box below \ts{24}{21}{15:01}\ts{13}{19}{9:22}.

\begin{keyformulas}
\textbf{Three messages of the lecture} (slide~7).
\begin{itemize}
  \item Given the current price of the stock \emph{and a model for its dynamics}, the price of a derivative is
        not uncertain.
  \item The price does not depend on the risk preferences of market participants (nor on the expected return of the stock).
  \item The mathematics of stochastic processes computes the price \emph{and the risks} (hedge ratios) of the derivative.
\end{itemize}
\end{keyformulas}

% ==================================================================
\section{One step: replication and the risk-neutral probability}\label{ch15:sec-onestep}
The lecture builds the theory in a market with one period of length $\dd t$ and two states \ts{24}{21}{16:04}\ts{13}{19}{11:27}. It is
the market of \cref{sec:binomial} with continuous compounding: the stock moves from $S_0$ to $S_1$ (\emph{up}, with real-world
probability $p$) or $S_2<S_1$ (\emph{down}), a riskless account grows from $B$ to $Be^{r\dd t}$, and a derivative $f$ pays $f_1$ or $f_2$.

\subsection{The forward}\label{ch15:sec-fwd1}
A forward on the stock pays $S-K_0$ at $\dd t$ and costs nothing today. \emph{The naive approach} sets the price
equal to the expected payoff with the real-world $p$,
$p(S_1-K_0)+(1-p)(S_2-K_0)=0$, i.e.\ $K_0=pS_1+(1-p)S_2$ (equal to $\frac12(S_1+S_2)$ if $p=\frac12$)
\ts{24}{21}{17:05}\ts{13}{19}{13:31}. It is wrong, because nobody knows $p$, and it is exploitable, as follows
\ts{24}{21}{19:08}\ts{13}{19}{14:33}.

\begin{proposition}[Forward price]\label{ch15:prop-fwd}
In the absence of arbitrage, $K_0=S_0e^{r\dd t}$: the strike of a zero-cost forward is the price of the stock capitalised at the
risk-free rate, whatever the real-world dynamics.
\end{proposition}
\begin{proof}
At time~$0$, borrow $S_0$ from the bank, buy the stock and sell the forward with strike $K_0$: the net cost is zero. At $\dd t$ deliver
the stock in exchange for $K_0$ and repay $S_0e^{r\dd t}$: the payoff is $K_0-S_0e^{r\dd t}$ in both states, a certain
amount that cost nothing. If it were positive this would be an arbitrage; if negative, the opposite position would
be \ts{24}{21}{20:11}\ts{13}{19}{15:35}. Hence $K_0=S_0e^{r\dd t}$.
\end{proof}
So the stock itself (financed with a loan) \emph{replicates} the forward: it has exactly the same payoff, and no assumption on $p$, $S_1$, $S_2$
or on how the stock ``really'' behaves was used. In the lecture's example, $S_0=80$, $T=2$ years and $K=88.41$: $88.41/80=e^{0.1}$, so $r=5\%$
continuously compounded \ts{13}{19}{16:42}. With $r=0$ one gets $K_0=S_0$, and one can still ask for a probability that gives the
same price as the expectation: the number $q$ with $K_0=qS_1+(1-q)S_2$ satisfies
\begin{equation}\label{ch15:eq-q}
  q=\frac{S_0e^{r\dd t}-S_2}{S_1-S_2}\in(0,1)\quad\text{iff}\quad S_2<S_0e^{r\dd t}<S_1,\qquad
  qS_1+(1-q)S_2=S_0e^{r\dd t}.
\end{equation}
This $q$ (called $p_n$ or $p_{rn}$ on the slides, $q$ in 2013) is the \emph{risk-neutral probability}: under it the stock earns exactly the
risk-free rate on average \ts{24}{21}{23:25}. The condition $0<q<1$ is the no-arbitrage condition of \cref{prop:binomial-na}
with $1+r_fT$ replaced by $e^{r\dd t}$.

\subsection{A general claim}\label{ch15:sec-claim1}
For a call (or any claim) the stock alone is no longer enough, but a mix of stock and cash is. Look for a stock position $a$ (in units of stock)
and a cash amount $B$ (in the riskless account, growing to $Be^{r\dd t}$) with \ts{24}{21}{25:30}--\ts{24}{21}{28:33}\ts{13}{19}{19:51}
\begin{equation}\label{ch15:eq-repl}
  f_1=aS_1+Be^{r\dd t},\qquad f_2=aS_2+Be^{r\dd t}.
\end{equation}
Subtracting, $a=\dfrac{f_1-f_2}{S_1-S_2}$, and then $B=e^{-r\dd t}\dfrac{S_1f_2-S_2f_1}{S_1-S_2}$. By no arbitrage $f_0=aS_0+B$
\ts{24}{21}{29:42}\ts{13}{19}{20:55}, and substituting one finds
\begin{equation}\label{ch15:eq-f0}
  f_0=e^{-r\dd t}\bigl[q\,f_1+(1-q)\,f_2\bigr],\qquad q=\frac{S_0e^{r\dd t}-S_2}{S_1-S_2},
\end{equation}
the same $q$ as in \eqref{ch15:eq-q}. Indeed
$aS_0+B=e^{-r\dd t}\bigl[\frac{(f_1-f_2)S_0e^{r\dd t}+S_1f_2-S_2f_1}{S_1-S_2}\bigr]
=e^{-r\dd t}\bigl[f_1\frac{S_0e^{r\dd t}-S_2}{S_1-S_2}+f_2\frac{S_1-S_0e^{r\dd t}}{S_1-S_2}\bigr]$
\ts{24}{21}{31:51}\ts{13}{19}{21:59}. This is \eqref{eq:rnbinomial} of \cref{sec:binomial} again.

\begin{proposition}[One-step risk-neutral pricing]\label{ch15:prop-onestep}
In the one-step market with $S_2<S_0e^{r\dd t}<S_1$, every claim $f$ is replicable, and its unique arbitrage-free price is
$f_0=e^{-r\dd t}\E_\Q[f_1]$ with $\Q=(q,1-q)$ given by \eqref{ch15:eq-q}. Equivalently, both $e^{-rt}S_t$ and $e^{-rt}f_t$ are martingales under $\Q$
(a two-point statement of the first fundamental theorem, \cref{thm:ftap}).
\end{proposition}
The price of the call is thus computed by taking \emph{expectations of its payoff}, $C_0=e^{-r\dd t}q(S_1-K)$ if $S_2<K<S_1$
(slide~12) \ts{24}{21}{31:51}. The bookie's pool proportions play the role of $q$; a physicist would say that $\Q$ is the measure that makes the
\emph{discounted} price a martingale, i.e.\ an unbiased random walk.

\begin{example}[Numbers of \cref{ex:call} in continuous compounding]\label{ch15:ex-call1}
Take $S_0=100$, $S_1=130$, $S_2=90$, $K=100$ and $e^{r\dd t}=1.03$ (the value used in \cref{ex:call}). Then $a=30/40=0.75$,
$B=-0.75\cdot90/1.03=-65.53$, $q=(103-90)/40=0.325$, and $f_0=0.75\cdot100-65.53=0.325\cdot30/1.03=9.466$, in agreement with
\cref{ex:call}. The real-world probability $p$ never appears.
\end{example}

\begin{intuition}[Replicating portfolio versus riskless portfolio]\label{ch15:int-riskless}
Equation~\eqref{ch15:eq-repl} says that the cash-plus-stock portfolio $aS+B$ has the same value as $f$ in both states. Turn it around
\ts{24}{21}{32:52}: the portfolio ``one unit of the derivative, minus $a$ units of stock'' pays $f_i-aS_i=Be^{r\dd t}$ in both
states: it is \emph{riskless}, and must therefore grow at the risk-free rate. For a physicist this is a statement that the hedged
combination has zero ``noise'', i.e.\ $a=\Delta f/\Delta S$ is a discrete derivative that cancels the fluctuation. In continuous time the same argument is
used in the form ``riskless portfolio'' (\cref{ch15:sec-pde}). Which of the two forms is more convenient depends on the problem \ts{24}{21}{32:52}.
\end{intuition}

\begin{exercise}[not from the course: one-step market]\label{ch15:ex-onestep}
$S_0=50$, $S_1=60$, $S_2=42$, $r\dd t=\ln1.02$.
\begin{enumerate}[(a)]
  \item Compute $q$, check the no-arbitrage condition and find the fair forward strike.
  \item Price a call ($K=50$) and a put ($K=50$) by replication and by $\E_\Q$ and verify parity.
  \item Let the real-world probability of ``up'' be $p=0.9$: compute $e^{-r\dd t}\E_{\Prob}[\text{call}]$ and the profit of a trader who sells the call at that price and hedges with $a$ shares.
\end{enumerate}
\end{exercise}


% ==================================================================
\section{Many steps: the binomial tree and its limit}\label{ch15:sec-tree}
This section is not spelled out in the lecture, but it is the bridge between the one-step market of
\cref{ch15:sec-onestep} and the continuous model of the slides (the second bullet of 2013 slide~21, ``tree methods, equivalent to
explicit scheme''); it is also the way most exotic options are still priced.

Split $[0,T]$ into $n$ steps of length $\dd t=T/n$. In each step the stock is multiplied by $u$ or by $d$; the
Cox--Ross--Rubinstein (CRR) choice is
\begin{equation}\label{ch15:eq-crr}
  u=e^{\sigma\sqrt{\dd t}},\qquad d=\frac1u=e^{-\sigma\sqrt{\dd t}},\qquad
  q=\frac{e^{r\dd t}-d}{u-d}.
\end{equation}
Each node of the tree is a copy of the one-step market, so \cref{ch15:prop-onestep} applies at every node and prices are computed by
\emph{backward induction} from the payoff at the leaves,
\begin{equation}\label{ch15:eq-backward}
  f(S,t)=e^{-r\dd t}\bigl[q\,f(uS,t+\dd t)+(1-q)\,f(dS,t+\dd t)\bigr],
  \qquad f(S,T)=\text{payoff}(S).
\end{equation}
Unrolling it, $f(S_0,0)=e^{-rT}\E_\Q[\text{payoff}(S_T)]$ with $S_T=S_0u^{J}d^{n-J}$ and $J\sim\mathrm{Bin}(n,q)$. Since
the number $q$ was fixed by the requirement that the \emph{discounted} stock be a martingale, this is the multi-period FTAP of \cref{thm:ftap}
and the martingale $e^{-rt}S_t$ of \cref{ch:sp1}.

\begin{proposition}[Two limits of the tree]\label{ch15:prop-limit}
Let $\dd t\to0$ with \eqref{ch15:eq-crr}.
\begin{enumerate}[(i)]
  \item \emph{Distribution.} $q=\frac12+\frac{r-\sigma^2/2}{2\sigma}\sqrt{\dd t}+O(\dd t)$ and
        $\ln(S_T/S_0)\Rightarrow N\bigl((r-\tfrac12\sigma^2)T,\ \sigma^2T\bigr)$. The price of a call
        converges to $e^{-rT}\E_\Q[(S_T-K)^+]$ with $S_T$ lognormal, i.e.\ to the Black--Scholes value of \cref{ch15:sec-formula}.
  \item \emph{Equation.} If $f$ is smooth, \eqref{ch15:eq-backward} implies
        $f_t+\tfrac12\sigma^2S^2f_{SS}+rSf_S-rf=o(1)$, the Black--Scholes PDE of \cref{ch15:sec-pde}.
\end{enumerate}
\end{proposition}
\begin{proof}
(i) Expand $u=1+\sigma\sqrt{\dd t}+\tfrac12\sigma^2\dd t+O(\dd t^{3/2})$, $d=1-\sigma\sqrt{\dd t}+\tfrac12\sigma^2\dd t+O(\dd t^{3/2})$, so
$u-d=2\sigma\sqrt{\dd t}\,(1+O(\dd t))$ and $e^{r\dd t}-d=\sigma\sqrt{\dd t}+(r-\tfrac12\sigma^2)\dd t+O(\dd t^{3/2})$, which gives the
expansion of $q$. A step of $\ln S$ is $\sigma\sqrt{\dd t}\,\xi$ with $\xi=\pm1$, $\Prob_\Q(\xi=1)=q$; its mean is
$\sigma\sqrt{\dd t}(2q-1)=(r-\tfrac12\sigma^2)\dd t+O(\dd t^{3/2})$ and its variance $\sigma^2\dd t\,[1-(2q-1)^2]=\sigma^2\dd t\,(1+O(\dd t))$.
Summing $n=T/\dd t$ independent steps, the mean tends to $(r-\tfrac12\sigma^2)T$ and the variance to $\sigma^2T$; the steps are bounded by $\sigma\sqrt{\dd t}\to0$,
so the Lindeberg central limit theorem (\cref{ch:probability}) applies. The payoff is continuous with linear growth and
$\E_\Q[S_T]=S_0e^{rT}$ exactly for every $n$, which supplies the uniform integrability needed to pass to the limit in the expectation.

(ii) By the definition of $q$, $q(u-1)+(1-q)(d-1)=e^{r\dd t}-1$ exactly, and $q(u-1)^2+(1-q)(d-1)^2=\sigma^2\dd t\,(1+O(\sqrt{\dd t}))$. Taylor-expanding
$f(uS,t+\dd t)$ and $f(dS,t+\dd t)$ to second order in $(u-1)S$, $(d-1)S$ and first order in $\dd t$ and inserting in \eqref{ch15:eq-backward} gives
$e^{r\dd t}f=f+f_t\dd t+(e^{r\dd t}-1)Sf_S+\tfrac12\sigma^2S^2f_{SS}\dd t+o(\dd t)$; subtract $f$ and divide by $\dd t$.
\end{proof}

\Cref{ch15:tab-crr} and \cref{ch15:fig-crr} show the convergence for the option of \cref{ch15:fig-payoffs} at the money:
$S_0=K=80$, $T=2$, $r=5\%$, $\sigma=20\%$, for which the formula of \cref{ch15:sec-formula} gives $C=12.9014$. The error decays like $1/n$ but alternates in sign
with the parity of $n$, because at the money the strike sits on a node for even $n$ and in the middle of a lattice cell for odd $n$: the column $n\cdot\text{error}$ is
stable near $-2.26$ for even $n$ and near $+1.73$ for odd $n$ (e.g.\ $n=51,101,1001$), so the oscillation does not die out, only its amplitude shrinks. The single-step value $14.51$ is more than $12\%$ off.

\begin{table}[ht]
\centering\small
\begin{tabular}{r r r r r}
\toprule
$n$ & tree price & error & $n\cdot$error & $q$\\
\midrule
1 & 14.5106 & $+1.6092$ & $+1.61$ & 0.6132\\
2 & 11.8731 & $-1.0283$ & $-2.06$ & 0.5775\\
5 & 13.2506 & $+0.3492$ & $+1.75$ & 0.5481\\
10 & 12.6786 & $-0.2228$ & $-2.23$ & 0.5338\\
50 & 12.8564 & $-0.0451$ & $-2.25$ & 0.5150\\
100 & 12.8789 & $-0.0226$ & $-2.26$ & 0.5106\\
1000 & 12.8992 & $-0.0023$ & $-2.26$ & 0.5034\\
5000 & 12.9010 & $-0.0005$ & $-2.27$ & 0.5015\\
\bottomrule
\end{tabular}
\caption{CRR tree price of the European call $S_0=K=80$, $T=2$, $r=5\%$, $\sigma=20\%$ (Black--Scholes: 12.9014). Computed for these notes.}\label{ch15:tab-crr}
\end{table}

\begin{figure}[ht]
\centering
\begin{tikzpicture}
\begin{axis}[width=0.85\linewidth,height=5.4cm,axis lines=left,xmin=0,xmax=60,ymin=11.5,ymax=14.7,
  xlabel={number of steps $n$},ylabel={price},tick label style={font=\small},label style={font=\small},
  legend style={at={(0.97,0.97)},anchor=north east,font=\scriptsize,draw=none,fill=none}]
  \addplot[only marks,mark=*,mark size=1.3pt,defcol] coordinates {(1,14.5106) (2,11.8731) (3,13.4811) (4,12.3590) (5,13.2506) (6,12.5342) (7,13.1508) (8,12.6240) (9,13.0952) (10,12.6786) (11,13.0599) (12,12.7153) (13,13.0354) (14,12.7416) (15,13.0175) (16,12.7614) (17,13.0038) (18,12.7768) (19,12.9930) (20,12.7892) (21,12.9842) (22,12.7993) (23,12.9770) (24,12.8078) (25,12.9710) (26,12.8150) (27,12.9658) (28,12.8211) (29,12.9614) (30,12.8264) (31,12.9575) (32,12.8311) (33,12.9541) (34,12.8352) (35,12.9511) (36,12.8389) (37,12.9484) (38,12.8422) (39,12.9460) (40,12.8451) (41,12.9438) (42,12.8478) (43,12.9418) (44,12.8502) (45,12.9400) (46,12.8524) (47,12.9384) (48,12.8545) (49,12.9369) (50,12.8564) (51,12.9355) (52,12.8581) (53,12.9342) (54,12.8597) (55,12.9330) (56,12.8612) (57,12.9319) (58,12.8626) (59,12.9308) (60,12.8638)};
  \addplot[thick,thmcol,dashed] coordinates {(0,12.9014) (60,12.9014)};
  \legend{CRR tree,Black--Scholes}
\end{axis}
\end{tikzpicture}
\caption{Convergence of the binomial price to the Black--Scholes value (same parameters as \cref{ch15:tab-crr}); the alternation between odd and even $n$ is visible.}\label{ch15:fig-crr}
\end{figure}

\begin{intuition}[Why $\sigma\sqrt{\dd t}$ and where $-\frac12\sigma^2$ comes from]
The steps of the tree are not independent of the time step: if the space step were $\propto\dd t$ the process would freeze,
and if it were larger than $\sqrt{\dd t}$ the variance would explode. The unique scaling that leaves a non-trivial limit is the diffusive one,
$\text{step}\sim\sqrt{\dd t}$, as in Donsker's theorem (\cref{ch11:thm-donsker}); this is what the lecturer proves
in the next section by counting variances. The drift correction $-\tfrac12\sigma^2$ in the exponent is the same lognormal correction as in
\cref{ch04:rem-logprice}: the tree has $\E_\Q[S_T]=S_0e^{rT}$ (the martingale property is exact), but $\ln S_T$ has mean
$(r-\tfrac12\sigma^2)T$ because $\E[e^{X}]=e^{\E[X]+\Var X/2}$ for a normal variable.
\end{intuition}

\begin{exercise}[not from the course: the tree]\label{ch15:ex-tree}
\leavevmode
\begin{enumerate}[(a)]
  \item Show that the CRR tree prices a forward $(S_T-K)$ exactly, for every $n$.
  \item Write a program that prices the call of \cref{ch15:tab-crr} and reproduces the $1/n$ error and the odd--even oscillation; then choose the
      strike lattice-aligned ($K=S_0u^{k}$) and observe the effect.
  \item Show that with $u=e^{\sigma\sqrt{\dd t}+\nu\dd t}$, $d=e^{-\sigma\sqrt{\dd t}+\nu\dd t}$ (drift $\nu$ in the tree) one still converges to Black--Scholes provided $q$ is defined by the martingale condition.
\end{enumerate}
\end{exercise}

% ==================================================================
\section{Lognormal dynamics and a heuristic It\^o formula}\label{ch15:sec-lognormal}

\subsection{The model}
The lecture models the stock by the \emph{lognormal} (geometric Brownian motion, \cref{ch11:eq-gbm}) dynamics
\ts{24}{21}{33:56}\ts{13}{19}{22:04}
\begin{equation}\label{ch15:eq-gbm}
  \dd S_t=\mu S_t\dd t+\sigma S_t\dd W_t,
\end{equation}
where $\mu$ is the (real-world) drift, $\sigma$ the \emph{volatility}, and $\dd W$ the increment of a Brownian motion: normal, with
$\E[\dd] W=0$ and $\operatorname{sd}(\dd W)=\sqrt{\dd t}$, independent for disjoint time intervals \ts{13}{19}{23:04}. The proportional change $\dd S/S$
has a deterministic part and a noise part. The lecture says clearly that the choice is an \emph{example}: the method needs only that the stock
has some dynamics, and other choices (stochastic volatility, jumps) lead to other equations of the same type
\ts{24}{21}{33:56}.

\subsection{Why the standard deviation is $\sqrt{\dd t}$}
In the 2024 lecture Strela stops to answer a question that ``always bothers'' him \ts{24}{21}{36:07}. Take a random walk $W_n=\sum_{i=1}^n\xi_i$ with steps $\xi_i=\pm1$ with
probability $\frac12$, uncorrelated: $\E[W_n]=0$ (a martingale) and $\Var W_n=\sum_i\E[\xi_i^2]=n$
\ts{24}{21}{37:10}. Now put the $n$ steps in a finite time $T$, $\dd t=T/n$, and let each step have size $\pm\delta$:
$\Var W=n\delta^2=T\delta^2/\dd t$. For the limit $\dd t\to0$ to be neither $0$ nor $\infty$ one needs $\delta^2/\dd t$ to stay finite, i.e.\ $\delta\propto\sqrt{\dd t}$
\ts{24}{21}{40:26}. That is the content of Donsker's theorem (\cref{ch11:thm-donsker}); the same
argument fixed the size of the steps of the tree \eqref{ch15:eq-crr}.

\subsection{Taylor expansion with a $\sqrt{\dd t}$ noise}
Let $f(S,t)$ be smooth. In a deterministic world the second-order Taylor formula is
$\dd f=f_t\dd t+f_S\dd S+\tfrac12f_{tt}\dd t^2+f_{tS}\dd t\dd S+\tfrac12f_{SS}\dd S^2+\dots$ \ts{24}{21}{43:35}. With $\dd S$ from \eqref{ch15:eq-gbm} the terms have very
different sizes \ts{24}{21}{44:45}: $\dd t\dd S\sim\dd t^{3/2}$ and $\dd t^2$ are negligible, but
\[
  (\dd S)^2=\sigma^2S^2(\dd W)^2+O(\dd t^{3/2}),\qquad(\dd W)^2\sim\dd t,
\]
is of the \emph{same} order as $\dd t$ and must be kept. One writes $(\dd W)^2=\dd t$ (\cref{ch11:eq-dB2}) and obtains, to order $\dd t$,
\ts{24}{21}{45:45}
\begin{equation}\label{ch15:eq-ito}
  \dd f=\Bigl(f_t+\mu Sf_S+\tfrac12\sigma^2S^2f_{SS}\Bigr)\dd t+\sigma Sf_S\dd W,
\end{equation}
\emph{It\^o's formula} for a function of $S$ and $t$ (slide~15). The lecturer calls it a ``quick and dirty derivation'' \ts{24}{21}{47:46}; the
rigorous statement, with the It\^o integral and the proof, is in \cref{ch:stochcalc}.

\begin{coursenote}[A student's objection]
A student asks how $(\dd W)^2=\dd t$ can hold when $\dd W$ is random \ts{24}{21}{55:16}. The answer given (by the lecturer and by the instructor, P.~Kempthorne) is that
the notation is sloppy but reflects a theorem: the \emph{quadratic variation} of Brownian motion over $[0,t]$ is exactly $t$,
$\sum_j(\Delta_jW)^2\to t$ almost surely as the mesh goes to zero (\cref{ch11:thm-qv}) \ts{24}{21}{56:21}. Each squared increment is random, of mean $\dd t$
and standard deviation $\sqrt2\,\dd t$, but their \emph{sum} over $n\sim1/\dd t$ intervals has mean $t$ and standard deviation
$\sqrt{2t\,\dd t}\to0$: a law of large numbers. Inside an integral $(\dd W)^2$ can therefore be replaced by $\dd t$.
\end{coursenote}

\begin{intuition}[Physicist's reading]
Equation~\eqref{ch15:eq-ito} is the Taylor expansion of a function of a Brownian path, with the ``ordinary'' second-order term promoted to first order
because the path is fractal ($\Delta W\sim\sqrt{\Delta t}$). Its physical avatar is the Langevin equation with multiplicative noise interpreted in the It\^o
(left-point) sense: $\tfrac12\sigma^2S^2f_{SS}$ is the ``noise-induced drift'' that a Stratonovich convention would absorb in the definition of the drift
(\cref{ch:stochcalc}). The same extra term makes $\ln S$ drift at $\mu-\tfrac12\sigma^2$ rather than $\mu$: taking $f=\ln S$ in \eqref{ch15:eq-ito}
gives $\dd\ln S=(\mu-\tfrac12\sigma^2)\dd t+\sigma\dd W$, and hence \cref{ch11:eq-gbm},
$S_T=S_t\exp\bigl[(\mu-\tfrac12\sigma^2)\tau+\sigma(W_T-W_t)\bigr]$.
\end{intuition}

% ==================================================================
\section{The Black--Scholes equation}\label{ch15:sec-pde}
We now do in continuous time what \cref{ch15:sec-claim1} did in one step \ts{24}{21}{47:46}\ts{13}{19}{23:04}. \textbf{Assumptions} (implicit in the lecture): the
stock follows \eqref{ch15:eq-gbm} with \emph{constant} $\mu$, $\sigma$; the interest rate $r$ is constant and the same for borrowing and lending; the stock
pays no dividend; trading is continuous, without transaction costs, with unrestricted short selling and fractional positions; the claim is European.

\subsection{Derivation by replication}
Let a portfolio hold $a_t=a(S_t,t)$ shares of the stock and a cash amount $B_t$ in the money-market account, which grows as
$\dd B$ (due to interest) $=rB\dd t$ \ts{24}{21}{49:52}. Its value is $V=aS+B$. The portfolio is \emph{self-financing}: no money is added or withdrawn,
so the change of value is only due to price moves and interest,
\begin{equation}\label{ch15:eq-selffin}
  \dd V=a\,\dd S+rB\dd t=(a\mu S+rB)\dd t+a\sigma S\dd W .
\end{equation}
(On the slides: $\dd f=a\,\dd S+\dd B$; the lecture writes $\dd B=rB\dd t$ and $B=f-aS$, the same thing.) We want the portfolio to \emph{replicate} the claim:
$V_t=f(S_t,t)$ at all times, with $f(S,T)$ the payoff. Equate \eqref{ch15:eq-ito} and \eqref{ch15:eq-selffin} \ts{24}{21}{51:59}\ts{13}{19}{25:12}.

\emph{Noise terms} (coefficients of $\dd W$): $\sigma Sf_S=a\sigma S$, so
\begin{equation}\label{ch15:eq-hedge}
  a=\frac{\partial f}{\partial S}\quad(\text{the \emph{delta}})\ts{24}{21}{57:27}.
\end{equation}
\emph{Drift terms} (coefficients of $\dd t$), with $B=f-aS=f-Sf_S$ from $V=f$:
$f_t+\mu Sf_S+\tfrac12\sigma^2S^2f_{SS}=\mu Sf_S+r(f-Sf_S)$. The real-world drift $\mu$ cancels and one is left with
\begin{equation}\label{ch15:eq-bs}
  \boxed{\ \frac{\partial f}{\partial t}+\frac12\sigma^2S^2\frac{\partial^2f}{\partial S^2}+rS\frac{\partial f}{\partial S}-rf=0\ }
\end{equation}
\ts{24}{21}{58:33}\ts{13}{19}{26:20}, the \emph{Black--Scholes equation} (slide~17), a backward parabolic PDE with final condition $f(S,T)=\text{payoff}(S)$. It is the result of Black and Scholes (1973); Scholes and Merton received the 1997 Nobel prize in economics (Black had died in 1995), and all three were associated with MIT \ts{24}{21}{1:01:47}.
If the claim traded at any price other than $V_t$, buying the cheaper and selling the dearer of $f$ and $V$ would be an arbitrage, which is why $f=V$.

The lecture derives the same equation in the ``riskless portfolio'' form \ts{24}{21}{49:52}. Let $\Pi=f-aS$ with $a$ held fixed over the interval
$[t,t+\dd t]$. By \eqref{ch15:eq-ito}, $\dd\Pi=\dd f-a\dd S=(f_t+\tfrac12\sigma^2S^2f_{SS}+(f_S-a)\mu S)\dd t+(f_S-a)\sigma S\dd W$.
Choosing $a=f_S$ kills \emph{both} $\mu$ and the noise: $\dd\Pi=(f_t+\tfrac12\sigma^2S^2f_{SS})\dd t$ is deterministic, hence must equal $r\Pi\dd t=r(f-Sf_S)\dd t$ by no arbitrage,
which is \eqref{ch15:eq-bs} again. The delta-hedged position (long the option, short $f_S$ shares) earns the risk-free rate.

\begin{keyformulas}
\textbf{What the equation says} (slide~18; \ts{24}{21}{1:00:39}).
\begin{itemize}[nosep]
  \item Every tradable claim on $S$, whatever its payoff, satisfies \eqref{ch15:eq-bs}; the payoff enters only through the final (and boundary) conditions.
  \item The real-world drift $\mu$ does not appear: only $\sigma$ and $r$ do. (Compare \cref{sec:binomial}: the probability of ``up'' does not appear.)
  \item The replicating portfolio is also a \emph{hedging strategy}: to hedge a short position in the claim, hold $a=\partial f/\partial S$ shares and put the rest, $f-Sf_S$, in the bank.
        For a traded, liquid stock this is what a derivatives desk does all day \ts{13}{19}{28:31}.
  \item With another model for $S$ the equation changes (different coefficients, possibly more state variables), but the derivation is identical.
\end{itemize}
\end{keyformulas}

\subsection{Final and boundary conditions}
As for any PDE one needs conditions \ts{24}{21}{1:03:55}\ts{13}{19}{30:40}. The \emph{final} condition is the payoff, $C(S,T)=(S-K)^+$ for a call and $P(S,T)=(K-S)^+$ for a put.
The \emph{boundary} conditions follow from the economic sense of the options \ts{24}{21}{1:04:56}\ts{13}{19}{31:49}:
\begin{align}
  \text{call:}\quad & C(0,t)=0,\qquad C(S,t)\sim S-Ke^{-r(T-t)}\ \ (S\to\infty),\label{ch15:eq-bccall}\\
  \text{put:}\quad  & P(0,t)=Ke^{-r(T-t)},\qquad P(S,t)\to0\ \ (S\to\infty).\label{ch15:eq-bcput}
\end{align}
The stock price $S=0$ is absorbing for geometric Brownian motion (a stock worth zero stays worthless), so the call is worthless and the put is a
certain payment of $K$ at $T$, worth $Ke^{-r(T-t)}$ today. For $S\to\infty$ the call is exercised almost surely and is worth the forward, $S-Ke^{-r(T-t)}$.

\begin{coursenote}[Erratum: the call at infinity]
Slide~19 of 2024 (slide~18 of 2013) states $C(\infty,t)\cong S$. To leading order this is right, but the correct asymptotics is
\eqref{ch15:eq-bccall}: the call behaves as the \emph{forward} $S-Ke^{-r(T-t)}$, as the 2013 lecturer says in words (``discounted $S$ minus $K$'') \ts{13}{19}{31:49}. The put condition
$P(0,t)=Ke^{-r(T-t)}$ is correct on both sets of slides.
\end{coursenote}

\subsection{Reduction to the heat equation}\label{ch15:sec-heat}
Slide~18 states that ``by a change of variables'' \eqref{ch15:eq-bs} becomes the heat equation $u_\tau=u_{xx}$; the 2024 lecturer calls finding
the change of variables a good exercise \ts{24}{21}{1:02:51}, and in 2013 he proposed it as a final-paper topic \ts{13}{19}{30:40}. Here it is
(the slide's normalisation $u_\tau=u_{xx}$ differs from ours by a rescaling $\tau\to\sigma^2\tau/2$).

Put $\tau=T-t$ (so $\partial_t=-\partial_\tau$) and $y=\ln S$ (so $Sf_S=f_y$ and $S^2f_{SS}=f_{yy}-f_y$). Then \eqref{ch15:eq-bs} reads
\[
  f_\tau=\tfrac12\sigma^2f_{yy}+\bigl(r-\tfrac12\sigma^2\bigr)f_y-rf .
\]
Remove the discounting with $f=e^{-r\tau}w$, which gives $w_\tau=\tfrac12\sigma^2w_{yy}+\bigl(r-\tfrac12\sigma^2\bigr)w_y$, and remove the drift with a
\emph{Galilean shift}, $x=y+(r-\tfrac12\sigma^2)\tau$, $w(y,\tau)=u(x,\tau)$: since $w_\tau=u_\tau+(r-\tfrac12\sigma^2)u_x$, $w_y=u_x$, $w_{yy}=u_{xx}$, one gets
\begin{equation}\label{ch15:eq-heat}
\begin{gathered}
  f(S,t)=e^{-r\tau}\,u(x,\tau),\qquad x=\ln S+\bigl(r-\tfrac12\sigma^2\bigr)\tau,\\
  \frac{\partial u}{\partial\tau}=\frac{\sigma^2}2\frac{\partial^2u}{\partial x^2},\qquad u(x,0)=\text{payoff}(e^x).
\end{gathered}
\end{equation}
(One checks it on exponentials: $f=e^{-r\tau}e^{\alpha x+\sigma^2\alpha^2\tau/2}$ satisfies \eqref{ch15:eq-bs} for every $\alpha$; verified symbolically for these notes.)
By \cref{ch11:prop-heat} the solution with initial datum $u_0$ is $u(x,\tau)=\int p(y,\tau\mid x)\,u_0(y)\dd y$, the Gaussian kernel \eqref{ch11:eq-kernel}
with variance $\sigma^2\tau$: that is, $u(x,\tau)=\E\bigl[u_0(x+\sigma\sqrt\tau Z)\bigr]$ with $Z\sim N(0,1)$. Written in the original variables,
\begin{equation}\label{ch15:eq-pde-solution}
  f(S,t)=e^{-r\tau}\,\E\Bigl[\text{payoff}\bigl(Se^{(r-\sigma^2/2)\tau+\sigma\sqrt\tau Z}\bigr)\Bigr],
\end{equation}
which is exactly the risk-neutral expectation of \cref{ch15:sec-rn}: the two routes to the price (PDE and expectation) give the same answer, as they must (this is the
Feynman--Kac formula, \cref{ch:sde}).

\begin{intuition}[Diffusion of a payoff backwards in time]
The pricing PDE is the heat equation run \emph{backwards}: the payoff at $T$ is the ``initial temperature profile'' and the price today is that profile after
diffusing for a time $\tau$ with diffusivity $\sigma^2/2$, in the frame of the log-price moving with velocity $r-\tfrac12\sigma^2$ and with a decay $e^{-r\tau}$ (discounting).
The hockey stick smoothed by a Gaussian is precisely the curved price line of \cref{ch15:fig-payoffs}; the smoothing width
$\sigma\sqrt\tau$ is the size of the ``typical move'' of the log-price until expiry, so an option far from the strike is the payoff, and an option near
the strike is worth more than its intrinsic value because the price may still cross $K$ (the ``time value'' or convexity value).
\end{intuition}

\begin{exercise}[not from the course: the heat equation]\label{ch15:ex-heat}
\leavevmode
\begin{enumerate}[(a)]
  \item Verify the change of variables of \cref{ch15:sec-heat} by direct substitution.
  \item Solve $u_\tau=\frac12\sigma^2u_{xx}$ for $u(x,0)=\ind\{x>\ln K\}$ and show that $f=e^{-r\tau}\Phi(d_2)$ (the digital).
  \item The slide's normalisation is $u_\tau=u_{xx}$: find the rescaling of $\tau$.
\end{enumerate}
\end{exercise}


% ==================================================================
\section{Risk-neutral valuation}\label{ch15:sec-rn}
The second route to the price is the one the lecturer wants to stress: ``mostly about risk-neutral pricing, and Black--Scholes is an illustration of the method''
\ts{24}{21}{0:00}. In discrete time the price was $f_0=e^{-r\dd t}\E_\Q[f_1]$ with the special probability $q$. The continuous-time statement is
\ts{24}{21}{1:06:01}\ts{13}{19}{33:55}
\begin{equation}\label{ch15:eq-rnval}
  f_t=e^{-r(T-t)}\,\E_\Q\bigl[f_T\mid S_t\bigr],\qquad
  \text{where under }\Q:\quad \dd S=rS\dd t+\sigma S\dd W^{\Q},
\end{equation}
so that $S_t=e^{-r(T-t)}\E_\Q[S_T\mid S_t]$: under $\Q$ the \emph{discounted} stock price $e^{-rt}S_t$ is a martingale (slides 13, 21).

\paragraph{How to find $\Q$.} The lecturer's argument \ts{24}{21}{1:07:01}: take the expectation of \eqref{ch15:eq-gbm} (the
$\dd W$ term has mean zero), $\dd\,\E[S_t]=\mu\,\E[S_t]\dd t$, so $\E[S_T]=S_te^{\mu(T-t)}$ for the real-world dynamics. In the risk-neutral
world the stock must grow at the risk-free rate, $\E_\Q[S_T]=S_te^{r(T-t)}$ (this was the property $qS_1+(1-q)S_2=S_0e^{r\dd t}$ of the tree), which forces $\mu\to r$.
The volatility is \emph{unchanged}: the size of the fluctuations is a pathwise property (the quadratic variation, \cref{ch11:thm-qv}), the same under every equivalent measure, whereas the drift
can be changed at will. Changing only the drift of a Brownian motion is the content of Girsanov's theorem (\cref{ch:sde}): $W^{\Q}_t=W_t+\lambda t$ with
$\lambda=(\mu-r)/\sigma$, the \emph{market price of risk} (the Sharpe ratio of the stock; this identification is not made in the lecture). In this language
\begin{equation}\label{ch15:eq-girsanov}
  \dd S=\mu S\dd t+\sigma S\dd W=rS\dd t+\sigma S\bigl(\dd W+\lambda\dd t\bigr)=rS\dd t+\sigma S\dd W^{\Q}.
\end{equation}
The lognormal density of the terminal price under $\Q$ is (slide~21)
\begin{equation}\label{ch15:eq-density}
  p_\Q(S_T\mid S_t)=\frac1{S_T\,\sigma\sqrt{2\pi(T-t)}}\exp\left[-\frac{\bigl(\ln(S_T/S_t)-(r-\sigma^2/2)(T-t)\bigr)^2}{2\sigma^2(T-t)}\right],
\end{equation}
the density of the lognormal variable \cref{ch11:eq-gbm} with $\mu\to r$ \ts{24}{21}{1:09:15}. Pricing a claim then means ``integrate the payoff against this density and discount'', instead of solving
a PDE \ts{24}{21}{1:09:15}; this is easy for a lognormal density and typically not for other models. The equivalence of \eqref{ch15:eq-rnval} with the PDE is the Feynman--Kac
formula (\cref{ch:sde}), already visible in \eqref{ch15:eq-pde-solution}.

\begin{coursenote}[Erratum: the density on slide 21 (2024) and 20 (2013)]
The pre-factor of the printed density is wrong in both years. In 2024 it reads $S_t/(\sigma S_T\sqrt{2\pi T})$, in 2013 $1/(\sigma S\sqrt{2\pi T})$. The correct normalisation, which makes the density integrate to one
in $S_T$, is $1/\bigl(S_T\,\sigma\sqrt{2\pi(T-t)}\bigr)$, as in \eqref{ch15:eq-density}: the $S_T$ in the denominator is the Jacobian of $S_T=e^{X}$ and the time is the \emph{time to expiry}
$T-t$, as in the exponent, not $T$. The exponent on the slide is correct.
\end{coursenote}

\begin{intuition}[Why ``neutral'']
There is no claim that investors are risk-neutral. Real-world prices already contain risk premia: $\mu>r$ is the reward for bearing the risk of the stock. The claim is that the \emph{price of a derivative relative to the
stock} does not depend on those premia because the derivative can be replicated by the stock and the bank; so one may compute the price in any world where
the stock earns $r$, and the simplest is the one where nobody cares about risk. Changing $\mu\to r$ moves the drift of the log-price, not its width: a physicist would say that $\Q$ is the same
diffusion in a different, tilted, ``potential'' (the tilt is the drift), and the Radon--Nikodym factor $\exp(-\lambda W_T-\tfrac12\lambda^2T)$ is a
Boltzmann-like reweighting of the paths.
\end{intuition}

\begin{example}[The wrong expectation]\label{ch15:ex-wrong}
For $S_0=K=80$, $T=2$, $r=5\%$, $\sigma=20\%$ the correct price is $12.90$. Suppose the stock has real-world drift $\mu=12\%$. Discounting the \emph{real-world} expectation
of the payoff at the risk-free rate gives $e^{-rT}\E[(S_T-K)^+]=22.17$, a $72\%$ overestimate: it is not an arbitrage price, because one cannot replicate the option with it. A Monte Carlo estimate under $\Q$ with $2\cdot10^5$
paths gives $12.955\pm0.040$ (one standard error), consistent with $12.901$. (Computed for these notes.)
\end{example}

\begin{finance}[Where the approach is used]
For a desk the two languages are complementary. The PDE (finite-difference grids) is the tool for low-dimensional problems with early exercise (American options), barriers and local-vol models; the expectation
(Monte Carlo) for high-dimensional or path-dependent payoffs; trees are equivalent to an explicit finite-difference scheme (\cref{ch15:prop-limit}(ii)) \ts{24}{21}{1:04:56}\ts{13}{19}{34:59}.
Slide~22 of 2024 lists what the whole business uses: PDE, numerical methods, stochastic calculus, simulation, statistics, ``much, much more''.
\end{finance}

% ==================================================================
\section{The Black--Scholes formula}\label{ch15:sec-formula}

\begin{theorem}[Black--Scholes formula]\label{ch15:thm-bs}
In the model \eqref{ch15:eq-gbm}, with $\tau=T-t$, the European call and put are worth
\begin{align}
  C(S,t)&=S\,\Phi(d_1)-Ke^{-r\tau}\,\Phi(d_2)=e^{-r\tau}\bigl[e^{r\tau}S\,\Phi(d_1)-K\,\Phi(d_2)\bigr],\label{ch15:eq-call}\\
  P(S,t)&=Ke^{-r\tau}\,\Phi(-d_2)-S\,\Phi(-d_1)=e^{-r\tau}\bigl[K\,\Phi(-d_2)-e^{r\tau}S\,\Phi(-d_1)\bigr],\label{ch15:eq-put}
\end{align}
where $\Phi$ is the standard normal cdf and
\begin{equation}\label{ch15:eq-d12}
  d_1=\frac{\ln(S/K)+(r+\tfrac12\sigma^2)\tau}{\sigma\sqrt\tau},\qquad d_2=d_1-\sigma\sqrt\tau=\frac{\ln(S/K)+(r-\tfrac12\sigma^2)\tau}{\sigma\sqrt\tau}.
\end{equation}
\end{theorem}
The right-hand forms are the ones on slide~20 \ts{24}{21}{1:04:56}\ts{13}{19}{32:52}. In the lecture the formula is quoted as ``the solution of the PDE in terms of the error function, which has to be tabulated''; here are two derivations.

\begin{proof}[Proof 1: risk-neutral expectation]
By \eqref{ch15:eq-rnval}, $C=e^{-r\tau}\E[(S_T-K)^+]$ with $S_T=Se^{(r-\sigma^2/2)\tau+\sigma\sqrt\tau Z}$, $Z\sim N(0,1)$ (the marginal of \eqref{ch15:eq-density}). The event $\{S_T>K\}$ is
$\{Z>-d_2\}$ because $\ln(S_T/K)=\sigma\sqrt\tau\,(Z+d_2)$. So
\[
  \E\bigl[(S_T-K)^+\bigr]=\E\bigl[S_T\ind\{Z>-d_2\}\bigr]-K\,\Prob(Z>-d_2),\qquad\Prob(Z>-d_2)=\Phi(d_2).
\]
For the first term use the identity, for real $b,c$,
\[
  \E\bigl[e^{bZ}\ind\{Z>c\}\bigr]=\int_c^\infty e^{bz}\frac{e^{-z^2/2}}{\sqrt{2\pi}}\dd z=e^{b^2/2}\int_c^\infty\frac{e^{-(z-b)^2/2}}{\sqrt{2\pi}}\dd z=e^{b^2/2}\,\Phi(b-c)
\]
(completing the square). With $b=\sigma\sqrt\tau$, $c=-d_2$ one has $b-c=d_2+\sigma\sqrt\tau=d_1$ and
$\E[S_T\ind\{Z>-d_2\}]=Se^{(r-\sigma^2/2)\tau}e^{\sigma^2\tau/2}\Phi(d_1)=Se^{r\tau}\Phi(d_1)$. Multiplying by $e^{-r\tau}$ gives \eqref{ch15:eq-call}. The put follows in the same way with
$\{Z<-d_2\}$, or from \cref{ch15:thm-parity} below.
\end{proof}
This is also \cref{ch03:cor-normal}(ii) with $\mu_{\ln}=\ln S+(r-\tfrac12\sigma^2)\tau$ and $\sigma_{\ln}=\sigma\sqrt\tau$, exactly as anticipated in the finance box of that chapter; and it is the
two-line computation suggested in 2013 as a final-paper topic \ts{13}{19}{34:59}.

\begin{proof}[Proof 2: solving the PDE]
By \eqref{ch15:eq-pde-solution} the solution of \eqref{ch15:eq-bs} with final condition $(S-K)^+$ and the boundary conditions \eqref{ch15:eq-bccall} is the same Gaussian integral as in Proof~1; uniqueness of
the bounded-growth solution of the heat equation (\cref{ch11:prop-heat}) shows that it is \emph{the} solution. One can also verify directly that \eqref{ch15:eq-call} solves \eqref{ch15:eq-bs}, using the identity
$S\varphi(d_1)=Ke^{-r\tau}\varphi(d_2)$ (\cref{ch15:sec-greeks}), and that $C\to(S-K)^+$ as $\tau\to0$ (for $S>K$, $d_{1,2}\to+\infty$; for $S<K$, $d_{1,2}\to-\infty$).
\end{proof}

\begin{example}[The at-the-money call of the slides]\label{ch15:ex-atm}
For $S=K=80$, $\tau=2$, $r=5\%$, $\sigma=20\%$: $d_1=0.14/0.2828=0.4950$, $d_2=0.2121$, $\Phi(d_1)=0.6897$, $\Phi(d_2)=0.5840$,
$S\Phi(d_1)=55.175$, $Ke^{-r\tau}\Phi(d_2)=72.39\cdot0.5840=42.274$ and $C=12.901$. By parity the put is $5.288$.
The two terms are the price of \emph{receiving the stock if $S_T>K$} and of \emph{paying $K$ if $S_T>K$}, the number $\Phi(d_2)=0.584$ being the risk-neutral probability of finishing in the money.
\end{example}

\begin{coursenote}[On the parameters of the slides]
The slides label the plots with $S=80$, $K=80$, $T=2$ years (forward: $K=88.41$) and never state $r$ or $\sigma$. The forward strike identifies $r=5\%$ (continuously compounded, $80e^{0.1}=88.41$), and the digital of slide~26 tends to $0.905=e^{-0.1}$ for large $S$ and takes the value $\approx0.53$ at $S=80$;
with $\sigma=20\%$ formula \eqref{ch15:eq-digital} gives $0.528$. All figures in this chapter that reproduce the slides use $r=5\%$, $\sigma=20\%$; this identification is ours.
\end{coursenote}

\begin{intuition}[Reading the formula]
\begin{itemize}[nosep]
  \item $e^{-r\tau}\Phi(d_2)=$ price of the \emph{cash-or-nothing} (digital) option, $\Phi(d_2)=\Q(S_T>K)$; $S\Phi(d_1)$ is the price of the \emph{asset-or-nothing} option, which pays $S_T$ if $S_T>K$.
        $\Phi(d_1)$ is a probability too, but under the measure in which the stock, not the bank, is the numeraire: the probability that the option is exercised as seen by a holder of the stock.
  \item As $\sigma\sqrt\tau\to0$ the option becomes the forward $(S-Ke^{-r\tau})^+$: the time value is the smoothing of the hockey stick over a width $\sigma\sqrt\tau$ (\cref{ch15:fig-payoffs}).
  \item At the money forward ($Se^{r\tau}=K$), $d_1=-d_2=\tfrac12\sigma\sqrt\tau$ and $C=S\,[2\Phi(\tfrac12\sigma\sqrt\tau)-1]\approx0.4\,S\sigma\sqrt\tau$, the Bachelier rule of thumb quoted in \cref{ch:probability} (not discussed in class).
  \item The price depends on $S,K,\tau,r,\sigma$ but not on $\mu$. The only quantity that is not observed is $\sigma$, which is what traders quote instead of a price (\cref{ch15:sec-smile}).
\end{itemize}
\end{intuition}

\begin{exercise}[not from the course: the formula by expectation]\label{ch15:ex-formula}
Derive the put formula \eqref{ch15:eq-put} directly from $P=e^{-r\tau}\E_\Q[(K-S_T)^+]$ and check parity. Evaluate the call and the put for $S=100$, $K=90$, $\tau=0.5$, $r=3\%$, $\sigma=25\%$.
Show that $C$ is increasing and convex in $S$, increasing in $\sigma$ and $\tau$, and that $(S-Ke^{-r\tau})^+\le C\le S$ (the no-arbitrage bounds).
\end{exercise}

\begin{exercise}[not from the course: risk-neutral drift]\label{ch15:ex-girsanov}
Apply \eqref{ch15:eq-ito} to $e^{-rt}S_t$ to show that $\dd(e^{-rt}S_t)=(\mu-r)e^{-rt}S_t\dd t+\sigma e^{-rt}S_t\dd W$ and, with $W^{\Q}_t=W_t+\lambda t$, that it is a martingale for $\lambda=(\mu-r)/\sigma$. Compute
$\E_{\Prob}[C_T]$ and $\E_{\Q}[C_T]$ for the call of \cref{ch15:ex-atm} with $\mu=12\%$ and explain why $e^{-rT}\E_{\Prob}[C_T]$ is not the price.
\end{exercise}

% ==================================================================
\section{Put--call parity: a model-free replication}\label{ch15:sec-parity}
The lecture illustrates replication with a relation that requires \emph{no} assumption on the stock dynamics \ts{24}{21}{1:10:22}\ts{13}{19}{38:07}.

\begin{theorem}[Put--call parity]\label{ch15:thm-parity}
For European options on a non-dividend-paying asset with the same strike $K$ and expiry $T$, in the absence of arbitrage,
\begin{equation}\label{ch15:eq-parity}
  C(t)-P(t)=S_t-Ke^{-r(T-t)}.
\end{equation}
\end{theorem}
\begin{proof}
The payoffs at $T$ satisfy $(S_T-K)^+-(K-S_T)^+=S_T-K$ (the ``hockey stick'' and its reflection add up to a straight line \ts{24}{21}{1:11:25}). So buying a call and selling a put
replicates a forward with strike $K$ \emph{statically}, i.e.\ with no trading before $T$: it is a forward contract, and by
\cref{ch15:prop-fwd} (applied over $[t,T]$) it is worth $S_t-Ke^{-r(T-t)}$ (stock, minus the bond that pays $K$). The two portfolios have the same payoff in every state, so they have the same price. \ts{24}{21}{1:12:31}
\end{proof}
If \eqref{ch15:eq-parity} failed, say $C-P>S-Ke^{-r\tau}$, one would sell the call, buy the put and the stock, and borrow $Ke^{-r\tau}$, locking in the difference at no risk. The remark of the lecture is that the relation is
\emph{independent of the dynamics}: one can price a put from the observed calls without a model, which matters in the real market where the stock is \emph{not} lognormal
\ts{24}{21}{1:12:31}\ts{13}{19}{38:07}. It is also the reason why the implied volatility of a call and of a put with the same strike and expiry must coincide.

\subsection{Parity on Bloomberg screens}
The slides check the relation on real quotes (mid $=$ average of bid and ask). The tables below use numbers read from the screens for these notes; a reading error of a cent is possible.
\begin{table}[ht]
\centering\small
\begin{tabular}{r r r r r r}
\toprule
$K$ & call mid & put mid & $C-P$ & $S-Ke^{-r\tau}$ & difference\\
\midrule
75 & 6.900 & 0.300 & 6.600 & 6.555 & $+0.045$\\
80 & 2.950 & 1.300 & 1.650 & 1.583 & $+0.067$\\
85 & 0.750 & 4.150 & $-3.400$ & $-3.390$ & $-0.010$\\
90 & 0.125 & 8.800 & $-8.675$ & $-8.362$ & $-0.313$\\
\bottomrule
\end{tabular}
\caption{IBM options expiring 22 April 2006, screen of 8 March 2006 (2013 slide~23), $S=81.14$; the discount uses $r=4.5\%$ and $\tau=45/365$ (our assumptions: the slide does not
show a rate, so this is a rough choice, and dividends are ignored). Equity options are American, and the deep in-the-money put ($K=90$, wide spread) carries an early-exercise premium that parity does not include.}\label{ch15:tab-ibm}
\end{table}

\begin{table}[ht]
\centering\small
\begin{tabular}{r r r r r r}
\toprule
$K$ & call mid & put mid & $C-P$ & $S-PV(D)-Ke^{-r\tau}$ & difference\\
\midrule
90 & 23.63 & 0.26 & 23.37 & 23.14 & $+0.23$\\
100 & 14.00 & 0.77 & 13.23 & 13.15 & $+0.08$\\
110 & 6.10 & 2.83 & 3.27 & 3.17 & $+0.10$\\
120 & 1.62 & 8.35 & $-6.73$ & $-6.81$ & $+0.08$\\
\bottomrule
\end{tabular}
\caption{AAPL options expiring 20 January 2017 ($\tau=81$ days), screen of 31 October 2016 (2024 slide~24), $S=113.54$, $r=0.83\%$ and the dividend of \$0.57 shown on the screen (its present value is
subtracted from $S$: parity for a dividend-paying stock is $C-P=S-PV(D)-Ke^{-r\tau}$, ``not discussed'' in class).}\label{ch15:tab-aapl}
\end{table}
On slide~24 the lecturer plots the call, the put and their difference against strike: the difference is a straight line of slope $-e^{-r\tau}\approx-1$, while the implied volatility
of the same options is a U-shaped curve (see \cref{ch15:sec-smile}). A straight line is what parity requires, and ``if it wouldn't be, people would jump on it immediately'' \ts{24}{21}{1:13:36}.
In 2013 the same idea is shown on the IBM screen of 2006 (slide~23), at the end of the lecture: far from the strike a call is worth $S-K$ ($S-K=81.14-55=26.14$, against a quoted mid of $26.50$), and between
strikes the prices change by 5 when the strike changes by 5 \ts{13}{19}{36:03}.

\begin{finance}[Replication as a way of thinking]
The lecturer's closing message \ts{24}{21}{1:14:40}\ts{13}{19}{48:38}: ``what you try to do is look at the payout of a portfolio and replicate it with the instruments at hand, or at least approximate it,
because if you did, you know the price now, or you know an arbitrage''. Conversely one builds a complicated portfolio of listed options and looks for a mispricing with respect to a replicating one: nowadays this
is ``practically impossible'' for vanilla options, but was done ``25 years ago'' with more complex options in FX (a story he leaves to J.~Xia). The hedge of a book of exotics is usually dynamic and needs rebalancing
\ts{13}{19}{49:42}; static replications like parity and the call spread below are the exceptions, and the most robust.
\end{finance}

\begin{exercise}[not from the course: put--call parity]\label{ch15:ex-parity}
Prove parity when the stock pays a known cash dividend $D$ at $t_D\in(t,T)$: $C-P=S-De^{-r(t_D-t)}-Ke^{-r\tau}$. Using the numbers of \cref{ch15:tab-ibm} choose an interest rate $r$ that best fits the four strikes by least squares, and comment on the residuals.
Design an arbitrage (list the trades and the guaranteed profit) if the call in the $K=100$ row of \cref{ch15:tab-aapl} were quoted at 14.50 instead of 14.00.
\end{exercise}

% ==================================================================
\section{Digital options and call spreads}\label{ch15:sec-digital}
The lecture ends with a challenge \ts{24}{21}{1:15:43}, which in 2013 was posed as an interview question \ts{13}{19}{43:24}: \emph{price and hedge a digital (all-or-nothing) option}, which pays $1$
if $S_T>K$ and $0$ otherwise, given that the prices of calls of \emph{all strikes} are known. Slide~26 of 2024 compares the Black--Scholes digital with a ``call spread'' built from the same
AAPL data.

\paragraph{Replication.} A call spread that is long a call with strike $K-\varepsilon$ and short a call with strike $K+\varepsilon$ pays
$(S_T-K+\varepsilon)^+-(S_T-K-\varepsilon)^+$, a ramp from $0$ (for $S_T\le K-\varepsilon$) to $2\varepsilon$ (for $S_T\ge K+\varepsilon$). A student proposes it and the lecturer confirms:
$\pm\tfrac12$ around $K$ gives a ramp up to $1$, $\pm\tfrac14$ needs two spreads, and in general the payoff of $\frac1{2\varepsilon}$ spreads is a ramp of unit height and width $2\varepsilon$
\ts{13}{19}{44:31}. Letting $\varepsilon\to0$, the digital price is
\begin{equation}\label{ch15:eq-digital-limit}
  D(K)=\lim_{\varepsilon\to0}\frac{C(K-\varepsilon)-C(K+\varepsilon)}{2\varepsilon}=-\frac{\partial C}{\partial K}
\end{equation}
\emph{whatever the model} \ts{13}{19}{46:37}, since $C(K)=e^{-r\tau}\E_\Q[(S_T-K)^+]$ and $\partial_K(S_T-K)^+=-\ind\{S_T>K\}$; this is the Breeden--Litzenberger relation of \cref{ch03:cor-bl} ($D=e^{-r\tau}\Q(S_T>K)$, while $\partial^2_KC=e^{-r\tau}p_\Q(K)$). The lecturer notes that the answer is a
derivative of the call price with respect to $K$, and that it should not be surprising: the derivative of a ramp is a step \ts{13}{19}{46:37}; the same idea works for any payoff, e.g.\ a box. In practice the dealer
sells the digital to a client and hedges with two liquid calls on the exchange, and digitals are less liquid than calls \ts{13}{19}{47:37}.

In Black--Scholes the digital costs
\begin{equation}\label{ch15:eq-digital}
  D=e^{-r\tau}\,\Phi(d_2),
\end{equation}
which for $S=K=80$, $\tau=2$, $r=5\%$, $\sigma=20\%$ is $0.5284$; the call spread with half-width $\varepsilon=10,5,2,1,0.1$ costs $0.5291,\ 0.5286,\ 0.52846,\ 0.52843,\ 0.528423$ (computed for these notes).
For large $S$, $D\to e^{-r\tau}=0.905$, the value of a certain \$1 at $T$, as in the left graph of slide~26 \ts{24}{21}{1:16:44}.

\begin{figure}[ht]
\centering
\begin{tikzpicture}
\begin{axis}[width=0.9\linewidth,height=5.6cm,axis lines=left,xmin=20,xmax=140,ymin=0,ymax=1.08,
  xlabel={$S$},tick label style={font=\footnotesize},label style={font=\small},
  legend style={at={(0.03,0.97)},anchor=north west,font=\scriptsize,draw=none,fill=none}]
  \addplot[thick,black!45] coordinates {(20,0) (80,0) (80,1) (140,1)};
  \addplot[very thick,defcol] coordinates {(20,0.0000) (23,0.0000) (26,0.0001) (29,0.0003) (32,0.0011) (35,0.0030) (38,0.0070) (41,0.0142) (44,0.0259) (47,0.0431) (50,0.0666) (53,0.0967) (56,0.1331) (59,0.1753) (62,0.2220) (65,0.2722) (68,0.3244) (71,0.3772) (74,0.4295) (77,0.4802) (80,0.5284) (83,0.5736) (86,0.6153) (89,0.6534) (92,0.6877) (95,0.7183) (98,0.7453) (101,0.7691) (104,0.7897) (107,0.8076) (110,0.8230) (113,0.8362) (116,0.8474) (119,0.8568) (122,0.8649) (125,0.8716) (128,0.8773) (131,0.8820) (134,0.8859) (137,0.8892) (140,0.8920)};
  \addplot[thick,dashed,thmcol] coordinates {(20,0.0000) (23,0.0000) (26,0.0001) (29,0.0005) (32,0.0016) (35,0.0040) (38,0.0088) (41,0.0172) (44,0.0302) (47,0.0489) (50,0.0738) (53,0.1050) (56,0.1421) (59,0.1844) (62,0.2309) (65,0.2803) (68,0.3313) (71,0.3827) (74,0.4334) (77,0.4824) (80,0.5291) (83,0.5729) (86,0.6134) (89,0.6504) (92,0.6838) (95,0.7138) (98,0.7405) (101,0.7640) (104,0.7846) (107,0.8025) (110,0.8180) (113,0.8314) (116,0.8429) (119,0.8527) (122,0.8610) (125,0.8681) (128,0.8741) (131,0.8791) (134,0.8834) (137,0.8869) (140,0.8899)};
  \legend{payoff $\ind\{S>K\}$,{Black--Scholes $D(t,S)$},call spread $\varepsilon=10$}
\end{axis}
\end{tikzpicture}
\caption{The digital option ($K=80$, $\tau=2$, $r=5\%$, $\sigma=20\%$) and the call spread with strikes $70$ and $90$. The spread is already close to the digital price; for $\varepsilon\to0$ it converges to it.}\label{ch15:fig-digital}
\end{figure}

\paragraph{With market data.} The right graph of slide~26 is the call spread built from the AAPL quotes (slide~24) as a function of $K-S$. From the numbers of \cref{ch15:tab-aapl}, $(C(110)-C(120))/10=0.448$
and $(C(105)-C(125))/20=0.451$ estimate the digital with $K=115$ (probability of finishing in the money, discounted). A flat-volatility Black--Scholes digital with the implied volatility of the $K=115$ call ($19.8\%$) would give only $0.413$: the
difference is caused by the \emph{skew} of the smile ($\partial\sigma/\partial K<0$ there) and is a first-principles example of ``the world is not Black--Scholes'' (\cref{ch15:sec-smile}). The calculation is ours (not in class).

\begin{intuition}[Static versus dynamic replication]
Any smooth payoff $g(S_T)$ is a continuous portfolio of calls, $g(S_T)=g(0)+g'(0)S_T+\int_0^\infty g''(K)(S_T-K)^+\dd K$ (a Taylor formula with integral remainder in $K$; not discussed in class), so its price
is a fixed combination of call prices with no model at all. The Black--Scholes dynamic hedge is what one needs when the instruments are \emph{not} available at all strikes.
\end{intuition}

\begin{exercise}[not from the course: digital options]\label{ch15:ex-digital}
\leavevmode
\begin{enumerate}[(a)]
  \item Prove $D=-\partial C/\partial K$ for any model where $C(K)=e^{-r\tau}\E_\Q[(S_T-K)^+]$.
  \item Show that a call spread with strikes $K$ and $K+\varepsilon$, scaled by $1/\varepsilon$, \emph{super-replicates} the digital,
      so that the dealer's ask for the digital is $\le[C(K)-C(K+\varepsilon)]/\varepsilon$, and find the sub-replicating spread giving a bid.
  \item In a model with a smile, $\sigma=\sigma(K)$, show that $D=e^{-r\tau}\Phi(d_2)-\mathcal V\,\sigma'(K)$ and explain the sign
      of the correction for the AAPL data of \cref{ch15:tab-smile}.
\end{enumerate}
\end{exercise}


% ==================================================================
\section{Greeks and hedging}\label{ch15:sec-greeks}
The replication argument produced one hedge ratio, $a=\partial f/\partial S$ \eqref{ch15:eq-hedge}. Traders call the derivatives of the price with respect to its inputs the \emph{Greeks}; they measure
the risks of a position (``the mathematical apparatus allows one to compute the current price of a derivative and its risks'', slide~7). The lecture only uses the delta; the rest of this section is not
discussed in class but follows from \cref{ch15:thm-bs} and completes the picture.

\begin{proposition}[Greeks of the European call]\label{ch15:prop-greeks}
With $\varphi=\Phi'$ the standard normal density and $\tau=T-t$,
\begin{equation}\label{ch15:eq-greeks}
\begin{gathered}
  \Delta=\frac{\partial C}{\partial S}=\Phi(d_1),\qquad
  \Gamma=\frac{\partial^2C}{\partial S^2}=\frac{\varphi(d_1)}{S\sigma\sqrt\tau},\qquad
  \mathcal V=\frac{\partial C}{\partial\sigma}=S\varphi(d_1)\sqrt\tau,\\
  \Theta=\frac{\partial C}{\partial t}=-\frac{S\varphi(d_1)\sigma}{2\sqrt\tau}-rKe^{-r\tau}\Phi(d_2),\qquad
  \rho=\frac{\partial C}{\partial r}=K\tau e^{-r\tau}\Phi(d_2).
\end{gathered}
\end{equation}
For the put, $\Delta_P=\Delta-1$ and $\Gamma,\mathcal V$ are the same; $\Theta_P=-\frac{S\varphi(d_1)\sigma}{2\sqrt\tau}+rKe^{-r\tau}\Phi(-d_2)$ and $\rho_P=-K\tau e^{-r\tau}\Phi(-d_2)$.
The price satisfies \eqref{ch15:eq-bs} in the form $\Theta+\tfrac12\sigma^2S^2\Gamma+rS\Delta=rC$.
\end{proposition}
\begin{proof}
First the identity $S\varphi(d_1)=Ke^{-r\tau}\varphi(d_2)$: since $d_2=d_1-\sigma\sqrt\tau$,
$\varphi(d_2)=\varphi(d_1)\exp(\sigma\sqrt\tau\,d_1-\tfrac12\sigma^2\tau)$ and $\sigma\sqrt\tau\,d_1=\ln(S/K)+(r+\tfrac12\sigma^2)\tau$, so $\varphi(d_2)=\varphi(d_1)\,(S/K)\,e^{r\tau}$.
Now $\partial_Sd_1=\partial_Sd_2=1/(S\sigma\sqrt\tau)$ and
$\partial_SC=\Phi(d_1)+\bigl[S\varphi(d_1)-Ke^{-r\tau}\varphi(d_2)\bigr]\partial_Sd_1=\Phi(d_1)$, the bracket vanishing by the identity. Differentiating once more gives $\Gamma$.
For $\mathcal V$: $\partial_\sigma C=S\varphi(d_1)\partial_\sigma d_1-Ke^{-r\tau}\varphi(d_2)\partial_\sigma d_2=S\varphi(d_1)\,\partial_\sigma(d_1-d_2)=S\varphi(d_1)\sqrt\tau$, by the identity and $d_1-d_2=\sigma\sqrt\tau$. The theta and rho are computed in the same way (or theta from the PDE, once $\Delta,\Gamma$ are known);
the put values follow from parity $P=C-S+Ke^{-r\tau}$.
\end{proof}

Numbers for $S=K=100$, $\tau=1$, $r=5\%$, $\sigma=20\%$ ($d_1=0.35$, $d_2=0.15$; call $10.4506$, put $5.5735$; checked against finite differences and against the PDE, residual $10^{-16}$, for these notes):
\begin{center}\small
\begin{tabular}{l r r r r r}
\toprule
 & $\Delta$ & $\Gamma$ & $\mathcal V$ (per unit $\sigma$) & $\Theta$ (per year) & $\rho$ (per unit $r$)\\
\midrule
call & $0.6368$ & $0.01876$ & $37.52$ & $-6.414$ & $53.23$\\
put  & $-0.3632$ & $0.01876$ & $37.52$ & $-1.658$ & $-41.89$\\
\bottomrule
\end{tabular}
\end{center}
A one-point (1\%) rise in volatility adds about $0.375$ to the price, and one day of time decay of the call costs about $6.41/365=0.018$.

\begin{figure}[ht]
\centering
\begin{tikzpicture}
\begin{axis}[name=A,width=0.48\linewidth,height=5.2cm,xmin=60,xmax=140,ymin=0,ymax=1.05,axis lines=left,xlabel={$S$},
  tick label style={font=\footnotesize},label style={font=\small},title={\small delta $\Phi(d_1)$},title style={yshift=-3pt},
  legend style={at={(0.03,0.97)},anchor=north west,font=\scriptsize,draw=none,fill=none}]
  \addplot[very thick,defcol] coordinates {(60,0.0138) (62,0.0207) (64,0.0300) (66,0.0420) (68,0.0572) (70,0.0759) (72,0.0981) (74,0.1239) (76,0.1533) (78,0.1861) (80,0.2219) (82,0.2604) (84,0.3009) (86,0.3431) (88,0.3862) (90,0.4298) (92,0.4733) (94,0.5162) (96,0.5580) (98,0.5983) (100,0.6368) (102,0.6733) (104,0.7075) (106,0.7394) (108,0.7688) (110,0.7958) (112,0.8203) (114,0.8426) (116,0.8626) (118,0.8805) (120,0.8965) (122,0.9106) (124,0.9230) (126,0.9339) (128,0.9434) (130,0.9517) (132,0.9589) (134,0.9651) (136,0.9704) (138,0.9750) (140,0.9789)};
  \addplot[very thick,thmcol,dashed] coordinates {(60,0.0000) (62,0.0000) (64,0.0000) (66,0.0000) (68,0.0001) (70,0.0003) (72,0.0009) (74,0.0023) (76,0.0051) (78,0.0105) (80,0.0199) (82,0.0352) (84,0.0584) (86,0.0912) (88,0.1349) (90,0.1898) (92,0.2550) (94,0.3286) (96,0.4078) (98,0.4892) (100,0.5695) (102,0.6454) (104,0.7147) (106,0.7757) (108,0.8276) (110,0.8704) (112,0.9046) (114,0.9313) (116,0.9515) (118,0.9664) (120,0.9772) (122,0.9847) (124,0.9900) (126,0.9935) (128,0.9959) (130,0.9974) (132,0.9984) (134,0.9990) (136,0.9994) (138,0.9997) (140,0.9998)};
  \legend{$\tau=1$,$\tau=0.25$}
\end{axis}
\begin{axis}[at={(A.east)},anchor=west,xshift=1.1cm,width=0.48\linewidth,height=5.2cm,xmin=60,xmax=140,ymin=0,ymax=0.07,axis lines=left,xlabel={$S$},
  tick label style={font=\footnotesize},label style={font=\small},title={\small gamma},title style={yshift=-3pt},
  legend style={at={(0.03,0.97)},anchor=north west,font=\scriptsize,draw=none,fill=none}]
  \addplot[very thick,defcol] coordinates {(60,0.0029) (62,0.0040) (64,0.0053) (66,0.0068) (68,0.0084) (70,0.0102) (72,0.0120) (74,0.0138) (76,0.0156) (78,0.0172) (80,0.0186) (82,0.0198) (84,0.0207) (86,0.0214) (88,0.0217) (90,0.0218) (92,0.0216) (94,0.0212) (96,0.0206) (98,0.0197) (100,0.0188) (102,0.0177) (104,0.0165) (106,0.0153) (108,0.0141) (110,0.0129) (112,0.0117) (114,0.0106) (116,0.0095) (118,0.0085) (120,0.0075) (122,0.0066) (124,0.0058) (126,0.0051) (128,0.0044) (130,0.0039) (132,0.0033) (134,0.0029) (136,0.0025) (138,0.0021) (140,0.0018)};
  \addplot[very thick,thmcol,dashed] coordinates {(60,0.0000) (62,0.0000) (64,0.0000) (66,0.0000) (68,0.0001) (70,0.0002) (72,0.0004) (74,0.0010) (76,0.0019) (78,0.0036) (80,0.0060) (82,0.0095) (84,0.0139) (86,0.0191) (88,0.0247) (90,0.0301) (92,0.0349) (94,0.0385) (96,0.0404) (98,0.0407) (100,0.0393) (102,0.0365) (104,0.0327) (106,0.0282) (108,0.0236) (110,0.0192) (112,0.0151) (114,0.0116) (116,0.0087) (118,0.0063) (120,0.0045) (122,0.0031) (124,0.0022) (126,0.0014) (128,0.0009) (130,0.0006) (132,0.0004) (134,0.0002) (136,0.0001) (138,0.0001) (140,0.0001)};
  \legend{$\tau=1$,$\tau=0.25$}
\end{axis}
\end{tikzpicture}
\caption{Delta and gamma of a call with $K=100$, $r=5\%$, $\sigma=20\%$. Near expiry delta approaches the step function of the payoff and gamma concentrates around the strike, which makes short-dated at-the-money options the
hardest to hedge.}\label{ch15:fig-greeks}
\end{figure}

\paragraph{Delta hedging.} To be short one call and hedged one holds $\Delta=\Phi(d_1)$ shares (about $64$ shares per $100$ calls in the example above) and a negative cash position; the combination has
no first-order sensitivity to $S$. The residual risk is second order: over a small move $\delta S$ the hedged book changes by $\Theta\,\delta t+\tfrac12\Gamma\,\delta S^2$,
and the PDE $\Theta+\tfrac12\sigma^2S^2\Gamma=r(C-S\Delta)$ says that, on average, the gain from gamma ($\tfrac12\Gamma\,\delta S^2\approx\tfrac12\Gamma\sigma^2S^2\delta t$) pays for the time decay: a long-gamma
position bleeds theta and earns it back through movement, exactly at the rate that makes the hedged book earn $r$. This is the origin of the trader's statement that an option trade is a bet on realised versus implied volatility (\cref{ch:volatility}).

\begin{table}[ht]
\centering\small
\begin{tabular}{r r r r}
\toprule
rebalancings $n$ & mean error & std.\ dev. & std.\ dev.\ $\times\sqrt n$\\
\midrule
1 & $-0.208$ & 6.109 & 6.11\\
2 & $-0.088$ & 4.379 & 6.19\\
4 & $-0.092$ & 3.225 & 6.45\\
12 & $-0.040$ & 1.902 & 6.59\\
52 & $-0.014$ & 0.927 & 6.69\\
252 & $-0.005$ & 0.426 & 6.76\\
\bottomrule
\end{tabular}
\caption{Error of discrete delta hedging of a short call ($S_0=K=100$, $T=1$, $r=5\%$, $\sigma=20\%$, price $10.4506$ received) in a market whose real-world drift is $\mu=10\%\ne r$:
terminal value of the hedge portfolio minus the payoff, $2\cdot10^4$ simulated paths, hedge ratio $\Phi(d_1)$ recomputed at $n$ equally spaced dates. Computed for these notes.}\label{ch15:tab-hedge}
\end{table}

\paragraph{What discrete rebalancing costs.} The derivation \eqref{ch15:eq-selffin} assumed that the hedge is adjusted \emph{continuously}. \Cref{ch15:tab-hedge} shows what happens with $n$ rebalancings: the hedge error
has a small mean (below $0.21$ here, and shrinking with $n$) and standard deviation proportional to $n^{-1/2}$ (the last column is roughly constant), i.e.\ proportional to $\sqrt{\dd t}$: the risk is
not eliminated for a finite $n$, only made small, and every rebalancing has a transaction cost, so the desk chooses a finite frequency \ts{13}{19}{28:31}\ (``every time $\dd t$ we will have to rebalance'').

\begin{finance}[What breaks the assumptions]
Discrete hedging and transaction costs (above); jumps and gaps, for which no hedge in the stock alone is perfect and the market is incomplete (\cref{sec:oneperiod}: a non-replicable claim has an interval of arbitrage-free prices)
\ts{24}{21}{1:14:40}; stochastic volatility, so that $\sigma$ itself needs hedging with other options; funding and credit spreads (\cref{ch:counterparty}); dividends and borrow costs. Each is a chapter of its own in a derivatives library;
Black--Scholes remains the language in which prices are \emph{quoted}.
\end{finance}

\begin{exercise}[not from the course: Greeks]\label{ch15:ex-greeks}
Prove $S\varphi(d_1)=Ke^{-r\tau}\varphi(d_2)$, derive $\Delta,\Gamma,\mathcal V,\Theta$ and check that they satisfy the Black--Scholes PDE. For an at-the-money-forward option show that $C\approx0.4\,S\sigma\sqrt\tau$; prove in general that $\mathcal V=S^2\sigma\tau\,\Gamma$. What is the $\Delta$ of a call with $S$ very large? Of a call very close to expiry, at the money?
\end{exercise}

\begin{exercise}[not from the course: delta hedging by simulation]\label{ch15:ex-hedge}
Reproduce \cref{ch15:tab-hedge}: simulate GBM with drift $\mu$, sell a call at the Black--Scholes price, rebalance the delta at $n$ dates, and record the terminal error. Check that the distribution does not depend on $\mu$ (for large $n$),
that its width scales as $n^{-1/2}$, and add a proportional transaction cost $c|\Delta a|S$ to find the $n$ that minimises the cost-plus-risk for your favourite risk penalty. What changes if you hedge with a wrong $\sigma$ (say $15\%$ when the true one is $20\%$)?
\end{exercise}

% ==================================================================
\section{Implied volatility and the smile}\label{ch15:sec-smile}
The formula has one input that is not observed, $\sigma$. Since $\partial C/\partial\sigma=\mathcal V>0$ and $C\to(S-Ke^{-r\tau})^+$ as $\sigma\to0$, $C\to S$ as $\sigma\to\infty$, the map $\sigma\mapsto C_{BS}(\sigma)$ is a strictly increasing
bijection from $(0,\infty)$ onto $\bigl((S-Ke^{-r\tau})^+,S\bigr)$; hence every market price in that interval determines a unique \emph{implied volatility} $\sigma_{\mathrm{imp}}$ (found numerically, e.g.\ by Newton's method with the vega as derivative). Parity implies that call and put
of the same strike and expiry have the same $\sigma_{\mathrm{imp}}$.

In Black--Scholes $\sigma$ is one constant for the stock, so the implied volatility should be the \emph{same for all strikes and expiries}. The lecture points out that it is not
\ts{24}{21}{1:13:36}\ts{13}{19}{37:04}: plotted against strike it is a curve, a ``smile'' (in 2013 the word is ``skewed, or rather smiled''), shown in the lower right graph of slide~24. \Cref{ch15:tab-smile} and
\cref{ch15:fig-smile} report the numbers of that screen.

\begin{table}[ht]
\centering\small
\begin{tabular}{r r r r}
\toprule
$K$ & call mid & screen implied vol.\ (\%) & recomputed (\%)\\
\midrule
90 & 23.63 & 32.53 & 33.12\\
100 & 14.00 & 23.71 & 24.13\\
110 & 6.10 & 20.47 & 20.69\\
120 & 1.62 & 19.06 & 19.08\\
130 & 0.31 & 19.41 & 19.37\\
140 & 0.09 & 21.45 & 21.79\\
\bottomrule
\end{tabular}
\caption{AAPL calls of \cref{ch15:tab-aapl} ($S=113.54$, $\tau=81$ days). ``Recomputed'' inverts the European Black--Scholes formula with $S-PV(D)$, $r=0.83\%$ (our computation); the small differences from the Bloomberg values reflect the
American-style and dividend conventions of the screen, and are largest for deep in-the-money and far out-of-the-money quotes.}\label{ch15:tab-smile}
\end{table}

\begin{figure}[ht]
\centering
\begin{tikzpicture}
\begin{axis}[width=0.85\linewidth,height=5.6cm,axis lines=left,xmin=70,xmax=150,ymin=15,ymax=55,
  xlabel={strike $K$ ($S=113.54$)},ylabel={implied volatility (\%)},tick label style={font=\small},label style={font=\small}]
  \addplot[very thick,defcol,mark=*,mark size=1.4pt] coordinates {(75,51.08) (80,45.65) (82.5,43.27) (85,38.82) (90,32.53) (92.5,29.02) (95,26.92) (97.5,25.21) (100,23.71) (105,22.04) (110,20.47) (115,19.57) (120,19.06) (125,18.91) (130,19.41) (135,20.25) (140,21.45) (145,22.37)};
  \draw[thmcol,dashed] (axis cs:113.54,15) -- (axis cs:113.54,55) node[pos=0.93,right,font=\scriptsize]{spot};
\end{axis}
\end{tikzpicture}
\caption{Implied volatility of the AAPL calls, screen of 31 October 2016 (values read from 2024 slide~24; the quote at $K=87.5$ is missing). A constant-volatility model would give a horizontal line.}\label{ch15:fig-smile}
\end{figure}

The curve decreases steeply for low strikes and turns up slightly for high ones: options in the wings are relatively \emph{expensive}. Two consequences. First, the lognormal distribution is not the distribution the market uses: returns have
heavier tails and skew (\cref{ch03:sec-heavy}), volatility moves with the price, and the wings price in crash risk. Second, prices of \emph{non-vanilla} claims cannot be read off one volatility number: as in the digital
of \cref{ch15:sec-digital}, the skew enters. The lecturer's conclusion \ts{24}{21}{1:14:40}\ is that Black--Scholes is a ``nice example'' of the method, but more complicated dynamics are needed for stocks: \emph{stochastic volatility}
(the volatility itself follows an SDE; \cref{ch:volatility,ch:sde}) or \emph{jump-diffusion}. With jumps ``life becomes more complicated'': the market is incomplete and martingale
techniques for continuous paths no longer suffice. The risk-neutral \emph{method} is unchanged: find the measure(s) $\Q$ under which discounted traded assets are martingales, and price by discounted expectation; what changes is that (for stochastic volatility and jumps) $\Q$
is no longer unique (\cref{thm:ftap}, part~2), and calibration to the observed option prices (the smile) is what selects it.

\begin{exercise}[not from the course: implied volatility]\label{ch15:ex-iv}
Implement Newton's method for $\sigma_{\mathrm{imp}}$ with the vega as derivative, starting from $\sigma_0=\sqrt{2|\ln(S/K)+r\tau|/\tau}$. Recompute the last column of \cref{ch15:tab-smile}. Why does the iteration fail for deep out-of-the-money options and what is a safer alternative (bisection, Brent)?
\end{exercise}

% ==================================================================
\section{The two lectures compared, and errata}\label{ch15:sec-compare}
\begin{coursenote}[2013 versus 2024]
The core slides are the same: the 2024 deck is the 2013 deck (V.~Strela was the speaker in both) with five changes.
\begin{itemize}[nosep]
  \item \emph{Opening.} 2024 starts with the update of the yield curve (\cref{ch15:sec-curve}); 2013 goes straight to the bookie.
  \item \emph{More derivation.} The 2024 lecture spends about 15 minutes on why $\operatorname{sd}(\dd W)=\sqrt{\dd t}$ and on the Taylor expansion behind It\^o's formula and on a student's objection to $(\dd W)^2=\dd t$
        (\cref{ch15:sec-lognormal}); 2013 assumes the students know It\^o from the lecture that precedes it (``did you talk about It\^o already? OK, cool'').
  \item \emph{Discrete market.} 2013 keeps the number $b$ of bond units and $B_0$ separate (its slide~10 has $b=\frac{S_1f_2-S_2f_1}{(S_1-S_2)B_0e^{r\dd t}}$), 2024 uses the cash amount $B$ directly; the risk-neutral probability is called $q$ in 2013 and $p_{rn}$ in 2024.
  \item \emph{Examples.} 2024 replaces the March 2006 IBM screen (2013 slide~23) with two recent Bloomberg screens (AAPL 2016 with the implied-volatility smile, and IBM 2015), adds the put--call parity plot and the real-data
        call spread (slide~26); 2013 discusses parity and the smile only orally, at the end, and treats the digital as an interview question with the audience answering.
  \item \emph{Suggestions.} 2013 explicitly proposes as final-paper topics the reduction to the heat equation and the derivation of the formula as a lognormal expectation
        \ts{13}{19}{30:40}, both carried out above; 2024 announces the next two (guest) lectures, one of them by J.~Hull on machine learning (\cref{ch:ml}).
\end{itemize}
In 2013 the lecturer also remarks that the two-year rate was then about $32.5$ basis points, so that the $5\%$ of the example is illustrative and not realistic \ts{13}{19}{17:43}.
\end{coursenote}

\begin{coursenote}[Errata found in the material]
\begin{itemize}
  \item \emph{Density of $S_T$} (2024 slide~21, 2013 slide~20): the pre-factor should be $1/(S_T\sigma\sqrt{2\pi(T-t)})$; see \cref{ch15:sec-rn}.
  \item \emph{Boundary condition of the call} (slide~19, 2013 slide~18): $C(S,t)\to S-Ke^{-r(T-t)}$ as $S\to\infty$, not only $\cong S$; see \cref{ch15:sec-pde}.
  \item \emph{Spoken slips} (2013 transcript): the bookie's payout ``\$12.25'' is \$12\,500 (\cref{ch15:sec-bookie}). (2024 transcript): ``our strike has to be 0'' after the argument for zero interest rates means $K=S_0$ ($K=S_0e^{r\dd t}$ in general) \ts{24}{21}{22:20};
        ``$S_2$ less than 0'' in the call diagram is $S_2<K$ \ts{24}{21}{26:32}.
\end{itemize}
\end{coursenote}

% ==================================================================
\begin{keyformulas}
\textbf{One step and tree} (continuous compounding, up/down factors $u,d$):
\begin{gather*}
  q=\frac{e^{r\dd t}-d}{u-d},\quad qu+(1-q)d=e^{r\dd t},\quad f_0=e^{-r\dd t}\bigl[qf_u+(1-q)f_d\bigr],\quad a=\frac{f_u-f_d}{S_u-S_d},\\
  \text{CRR: }u=e^{\sigma\sqrt{\dd t}}=1/d,\quad q=\tfrac12+\tfrac{r-\sigma^2/2}{2\sigma}\sqrt{\dd t}+O(\dd t).
\end{gather*}
\textbf{Model and PDE:}
\begin{gather*}
  \dd S=\mu S\dd t+\sigma S\dd W,\qquad \dd f=\bigl(f_t+\mu Sf_S+\tfrac12\sigma^2S^2f_{SS}\bigr)\dd t+\sigma Sf_S\dd W,\\
  f_t+\tfrac12\sigma^2S^2f_{SS}+rSf_S-rf=0,\qquad f(S,T)=\text{payoff},\qquad a=f_S .
\end{gather*}
\textbf{Risk-neutral valuation:}
\[
  f_t=e^{-r(T-t)}\E_\Q[f_T],\qquad \dd S=rS\dd t+\sigma S\dd W^{\Q},\qquad W^{\Q}_t=W_t+\tfrac{\mu-r}{\sigma}t .
\]
\textbf{Black--Scholes formula} ($\tau=T-t$):
\begin{gather*}
  C=S\Phi(d_1)-Ke^{-r\tau}\Phi(d_2),\qquad P=Ke^{-r\tau}\Phi(-d_2)-S\Phi(-d_1),\\
  d_{1,2}=\frac{\ln(S/K)+(r\pm\frac12\sigma^2)\tau}{\sigma\sqrt\tau},\qquad C-P=S-Ke^{-r\tau}.
\end{gather*}
\textbf{Greeks:}
\begin{gather*}
  \Delta_C=\Phi(d_1),\quad\Gamma=\frac{\varphi(d_1)}{S\sigma\sqrt\tau},\quad\mathcal V=S\varphi(d_1)\sqrt\tau,\quad S\varphi(d_1)=Ke^{-r\tau}\varphi(d_2),\\
  \Theta+\tfrac12\sigma^2S^2\Gamma+rS\Delta=rC .
\end{gather*}
\textbf{Digital and call spread:} $D=e^{-r\tau}\Phi(d_2)=-\partial C/\partial K=\lim_{\varepsilon\to0}\dfrac{C(K-\varepsilon)-C(K+\varepsilon)}{2\varepsilon}$.
\end{keyformulas}

% !TEX root = ../main.tex
\chapter{Introduction to Machine Learning}\label{ch:ml}

\begin{sources}
\small
\begin{tabularx}{\linewidth}{@{}l l L@{}}
\toprule
\textbf{Lecture} & \textbf{Speaker} & \textbf{Content}\\
\midrule
2024 L23 & J.~Hull (Rotman School, Toronto; guest) & What machine learning is; unsupervised, supervised and reinforcement learning; training/validation/test sets and overfitting; neural networks, activation functions, gradient descent, backpropagation; two applications (learning the Black--Scholes--Merton price, explaining the movements of the implied-volatility surface); reinforcement learning, exploration versus exploitation, the game of Nim, $Q$-learning, deep $Q$-learning; hedging derivatives by reinforcement learning; questions from the class.\\
\bottomrule
\end{tabularx}

\smallskip
\textbf{Files.} Slides \file{mit18_642_f24_lec23.pdf} (40 slides, no overlays); transcript
\file{Lecture 23 - Introduction to Machine Learning.txt} (about 75 minutes).
\textbf{Problem sets.} None; all exercises are marked ``(not from the course)''.
\textbf{Not covered in 2013.} The 2013 course had no machine-learning lecture.
\textbf{Not in the lecture.} The bias--variance decomposition, the convergence analysis of gradient descent,
the derivation of backpropagation, $k$-means and the Bellman equation are \emph{not} derived in the lecture, which is a
survey by an applied researcher. They are derived here so that the slides can be read as mathematics;
each such passage is marked \emph{(not discussed in class)}. The numerical experiments marked
``our reconstruction'' were run for these notes and are \emph{not} Hull's results.
\end{sources}

\section*{Motivation}
Every chapter so far fitted a model whose functional form was decided in advance: a linear factor model
(\cref{ch:regression,ch:pca}), an ARMA process (\cref{ch:timeseries}), a Gaussian return distribution
(\cref{ch:probability}). Machine learning starts from the opposite end. One has a large data set and a flexible
family of functions with many parameters, and asks the computer to find the function that predicts well
\emph{on data it has not seen}. The mathematics is that of the earlier chapters, used with a different emphasis:
\begin{itemize}
  \item least squares becomes one loss function among many, minimised by \emph{gradient descent} instead of by a closed
        formula (\cref{ch05:sec-model});
  \item overfitting, which in \cref{ch05:sec-shrink} was tamed by a ridge penalty chosen by cross-validation, becomes the
        central practical issue, handled by a train/validation/test split and by stopping the optimiser early;
  \item the flexible family is the \emph{neural network}, a composition of linear maps and fixed nonlinearities whose
        gradient is computed by the chain rule (``backpropagation'');
  \item decisions taken in sequence, such as hedging an option over its life, lead to \emph{reinforcement learning},
        dynamic programming without a model.
\end{itemize}
John Hull, known for his derivatives textbook, gave this lecture in 2024 and uses his own examples: a network that
imitates the Black--Scholes--Merton formula, a network that predicts how implied volatility moves when the
stock moves, and reinforcement-learning hedging. \Cref{ch16:sec-what} gives the vocabulary; \cref{ch16:sec-sup}
develops the theory of generalisation; \cref{ch16:sec-unsup} treats clustering; \cref{ch16:sec-gd,ch16:sec-nn}
develop gradient descent and neural networks; \cref{ch16:sec-bsm,ch16:sec-vol} are the two finance applications; and
\cref{ch16:sec-rl} covers reinforcement learning. The regularised models of \cref{ch:regularized} are the classical
counterpart of what is said here about overfitting.

% ==================================================================
\section{What machine learning is}\label{ch16:sec-what}
\subsection{Vocabulary, and three kinds of learning}
Hull introduces machine learning (ML) as a branch of artificial intelligence: give a program a lot of data and let it
find relationships between variables and make predictions \ts{24}{23}{7:17}. Some techniques date to the 1950s; what changed
is computer speed and the falling cost of storing data. The terminology differs from statistics, and the lecture insists
that this is a real obstacle to newcomers \ts{24}{23}{11:39}--\ts{24}{23}{12:44}.
\begin{table}[ht]
\centering\small
\begin{tabular}{@{}l l l@{}}
\toprule
\textbf{Statistics} & \textbf{Machine learning} & \textbf{In these notes}\\
\midrule
independent variables, regressors & \emph{features} & $\bm x\in\R^d$\\
dependent variable, response & \emph{target} (or \emph{label} for a class) & $y$\\
coefficients $\bm\beta$ & \emph{weights} (and \emph{biases}) & $\bm\theta$\\
fitting, estimation & \emph{training} & minimise the loss\\
sample, one pass over the data & \emph{epoch} (one pass of the optimiser) & \\
nonlinear link / basis function & \emph{activation function} & $\varphi$\\
step size of a numerical optimiser & \emph{learning rate} & $\eta$\\
tuning parameter (e.g.\ ridge $\lambda$) & \emph{hyperparameter} & $\lambda,\ \alpha,\ \varepsilon$\\
\bottomrule
\end{tabular}
\caption{A small dictionary. Hull's book has a ten-page glossary; these are the terms used in the lecture.}\label{ch16:tab-dict}
\end{table}

The lecture classifies algorithms in three families \ts{24}{23}{11:39}.
\begin{description}
  \item[Unsupervised learning] looks for structure in features alone, e.g.\ divides customers into clusters
        \ts{24}{23}{14:53}. ``No supervision'' means that nobody tells the algorithm what to predict \ts{24}{23}{13:47}.
        See \cref{ch16:sec-unsup}.
  \item[Supervised learning] predicts a numerical target (\emph{regression}) or a class (\emph{classification}) from features.
        Linear regression (\cref{ch:regression}) is the classical example. See \cref{ch16:sec-sup,ch16:sec-nn}.
  \item[Reinforcement learning] learns a rule for a \emph{sequence} of decisions in an environment that changes
        unpredictably. See \cref{ch16:sec-rl}.
\end{description}

\subsection{Digitisation versus learning; statistics versus ML}
Hull's periodisation of industrial revolutions (steam, then electricity and mass production, then computers and
digitisation, roughly 1950--2000, and now the fourth, AI) is used to separate two things \ts{24}{23}{7:17}. If loan officers
apply known rules, a computer can \emph{automate} them, which is digitisation. If the rules are unknown, ML can \emph{infer} them
from past applications and decisions; and one can go a step further and improve the rules by using data on how the loans
turned out \ts{24}{23}{8:24}--\ts{24}{23}{9:24}.

The second contrast concerns workflow \ts{24}{23}{9:24}--\ts{24}{23}{10:29}. Traditional statistics: formulate a hypothesis, collect data, test.
Machine learning: collect data, try different models, find patterns or build a predictive tool.
\begin{intuition}[The physicist's reading]
A physicist knows both modes: one fits a theory-driven model (Hooke's law) and one also explores a data set for unexpected
structure. What makes the second mode legitimate is precisely what the lecture goes on to stress, namely that the
model is judged on data that played no role in selecting it (the validation and test sets of
\cref{ch16:sec-sup}). Searching a data set for patterns without such a safeguard is what physicists call the
look-elsewhere effect: with enough candidate models, one always fits. ML's answer is procedural (sample splitting),
where classical hypothesis testing corrects the significance level.
\end{intuition}

% ==================================================================
\section{Supervised learning: generalisation and the bias--variance trade-off}\label{ch16:sec-sup}
\subsection{Loss, training set, validation set, test set}
\begin{definition}[Supervised learning problem]\label{ch16:def-sup}
Given pairs $(\bm x_i,y_i)$, $i=1,\dots,n$ (features and target), choose a function $h_{\bm\theta}$ in a parametrised
family $\mathcal H$ so as to make the \emph{empirical risk}
\[
  \hat R(\bm\theta)=\frac1n\sum_{i=1}^n\ell\bigl(y_i,h_{\bm\theta}(\bm x_i)\bigr)
\]
small, where $\ell$ is a loss function. The goal is a small \emph{expected} loss on new pairs, not on the $n$ given ones.
\end{definition}
The choice $\ell(y,\hat y)=(y-\hat y)^2$ with a linear $h_{\bm\theta}(\bm x)=\bm\beta\T\bm x$ is ordinary least squares
(\cref{ch05:sec-model}); other losses (absolute, check loss) are the M-estimators of \cref{ch05:sec-mle}. Hull's lecture
uses squared error (``mse'') and, in one plot, absolute error (see the erratum in \cref{ch16:sec-bsm}).

Because there is no hypothesis, the data are divided into three sets \ts{24}{23}{15:58}--\ts{24}{23}{17:03}, typically about 60\%, 20\% and 20\%.
\begin{itemize}
  \item the \emph{training set} determines the parameters $\bm\theta$ of each candidate model;
  \item the \emph{validation set} compares candidate models (and picks hyperparameters such as the degree of a polynomial or
        the number of layers);
  \item the \emph{test set} is kept aside and used once, to report the out-of-sample performance of the model finally chosen
        \ts{24}{23}{18:06}.
\end{itemize}
The rule of thumb is to increase the complexity of the model until performance on the validation set starts
to deteriorate \ts{24}{23}{18:06}. The three-way split is the simple, data-rich analogue of the cross-validation used in
\cref{ch05:sec-shrink} to choose the ridge parameter, and the validation error plays the role of the criterion that
\cref{ch:regularized} will minimise in a more structured setting.

\paragraph{Hull's salary example} \ts{24}{23}{19:09}--\ts{24}{23}{21:21}. Salaries in a profession are regressed on age with polynomials
(slide~11, whose raw data are not published). A fifth-degree polynomial follows the training points closely but does poorly
on the validation points (\emph{overfitting}); a straight line misses the rise-and-fall pattern (\emph{underfitting}); a
quadratic, in which salary rises and then tails off with age, is the best model. \Cref{ch16:fig-bv} reproduces the
mechanism, not the data.

\subsection{The bias--variance decomposition (not discussed in class)}
Overfitting and underfitting are two faces of one identity. Fix a point $\bm x$ and write $y=f(\bm x)+\varepsilon$ with
$\E\varepsilon=0$, $\Var\varepsilon=\sigma^2$, $\varepsilon$ independent of the training sample; $\hat h$ is the
fitted function, a random function through the training sample. Expanding the square around $\E\hat h(\bm x)$ (the cross terms vanish by
independence and by $\E[\hat h-\E\hat h]=0$),
\begin{equation}\label{ch16:eq-bv}
  \E\bigl[(y-\hat h(\bm x))^2\bigr]=\underbrace{\sigma^2}_{\text{irreducible noise}}
  +\underbrace{\bigl(f(\bm x)-\E\hat h(\bm x)\bigr)^2}_{\text{bias}^2}
  +\underbrace{\Var\hat h(\bm x)}_{\text{variance}} .
\end{equation}
For a concrete family the last two terms can be computed exactly. Take the linear-in-parameters case of
\cref{ch05:sec-model}: a fixed design, $\bm y=\bm f+\bm\varepsilon\in\R^n$, $\Cov\bm\varepsilon=\sigma^2I$, and the fit
$\hat{\bm y}=H\bm y$ with $H$ the hat matrix (a rank-$p$ orthogonal projection, \cref{ch05:prop-hat}); $p$ is the number of
fitted parameters (for a polynomial of degree $d$, $p=d+1$).

\begin{proposition}[Training and test error of a projection fit]\label{ch16:prop-bv}
Let $b^2=\frac1n\|(I-H)\bm f\|^2$ be the squared bias averaged over the design points. Then
\[
  \underbrace{\frac1n\E\|\bm y-\hat{\bm y}\|^2}_{\text{training error}}=\sigma^2+b^2-\frac{p}{n}\sigma^2,\qquad
  \underbrace{\frac1n\E\|\bm y^*-\hat{\bm y}\|^2}_{\text{test error at the same }\bm x_i,\ \text{fresh noise}}=\sigma^2+b^2+\frac{p}{n}\sigma^2 ,
\]
and the variance term of \eqref{ch16:eq-bv}, averaged over the design points, is $\frac pn\sigma^2$.
\end{proposition}
\begin{proof}
Since $I-H$ is a projection of trace $n-p$: $\bm y-\hat{\bm y}=(I-H)\bm f+(I-H)\bm\varepsilon$, and the cross term has zero mean, so
$\E\|\bm y-\hat{\bm y}\|^2=\|(I-H)\bm f\|^2+\sigma^2\tr(I-H)=nb^2+\sigma^2(n-p)$. For the fit error,
$\hat{\bm y}-\bm f=-(I-H)\bm f+H\bm\varepsilon$ gives $\E\|\hat{\bm y}-\bm f\|^2=nb^2+\sigma^2p$; the variance is
$\E\|H\bm\varepsilon\|^2=\sigma^2\tr H=p\sigma^2$. Finally $\bm y^*-\hat{\bm y}=\bm\varepsilon^*+(\bm f-\hat{\bm y})$ with
$\bm\varepsilon^*$ independent of $\hat{\bm y}$, so the test error is $\sigma^2+\frac1n\E\|\hat{\bm y}-\bm f\|^2$.
\end{proof}
The gap test $-$ training $=2\sigma^2p/n$ is the \emph{optimism} of the training error (it is the correction in Mallows' $C_p$).
More parameters lower the bias and raise the variance linearly in $p$; the training error decreases with $p$
\emph{always}, so it cannot choose the model. This is why the lecture selects on a validation set.

\begin{example}[Polynomials of degree $d$ fitted to $\sin 2\pi x$]\label{ch16:ex-poly}
Take $n=20$ equispaced points $x_i\in[0,1]$, $f(x)=\sin(2\pi x)$, $\sigma=0.3$ (\emph{our reconstruction}; Hull's salary
data are not available). The bias $b^2$ is the residual of the projection of $\bm f$ on the polynomials, computed
numerically (Legendre basis, python); the variance is exactly $\sigma^2(d+1)/n$.
Test error and training error follow from \cref{ch16:prop-bv}: the test error is minimal at $d=3$ (value $0.114$, of which $\sigma^2=0.09$ is noise), and it then rises slowly with
$d$, while the training error keeps falling ($0.078$ at $d=3$, $0.050$ at $d=8$).
\end{example}
\begin{figure}[ht]
\centering
\begin{tikzpicture}
\begin{axis}[width=0.92\linewidth,height=6.2cm,xlabel={degree $d$ of the polynomial},ylabel={mean squared error},xmin=-0.3,xmax=8.3,ymin=0,ymax=0.62,
  xtick={0,...,8},grid=major,grid style={black!10},legend pos=north east,legend style={font=\footnotesize,fill=white,draw=black!30}]
  \addplot[thick,color=defcol,mark=*,mark size=1.6pt] coordinates {(0,0.5695) (1,0.3303) (2,0.3348) (3,0.1140) (4,0.1185) (5,0.1170) (6,0.1215) (7,0.1260) (8,0.1305)};
  \addplot[thick,color=thmcol,mark=square*,mark size=1.6pt] coordinates {(0,0.5605) (1,0.3123) (2,0.3078) (3,0.0780) (4,0.0735) (5,0.0630) (6,0.0585) (7,0.0540) (8,0.0495)};
  \addplot[thick,dashed,color=intcol,mark=triangle*,mark size=1.6pt] coordinates {(0,0.4750) (1,0.2313) (2,0.2313) (3,0.0060) (4,0.0060) (5,0.0000) (6,0.0000) (7,0.0000) (8,0.0000)};
  \addplot[thick,dashed,color=cscol,mark=diamond*,mark size=1.8pt] coordinates {(0,0.0045) (1,0.0090) (2,0.0135) (3,0.0180) (4,0.0225) (5,0.0270) (6,0.0315) (7,0.0360) (8,0.0405)};
  \legend{test error,training error,$b^2$ (bias$^2$),variance $\sigma^2(d+1)/n$}
\end{axis}
\end{tikzpicture}
\caption{Expected training and test errors for polynomial regression of degree $d$ ($n=20$, $\sigma^2=0.09$; our reconstruction of the
mechanism of slide~11). The test error is $\sigma^2+b^2+\sigma^2(d+1)/n$; the training error is $\sigma^2+b^2-\sigma^2(d+1)/n$.
The minimum of the test error (degree~3) is the model the validation set would select; the training error alone would choose the largest degree.}\label{ch16:fig-bv}
\end{figure}

\begin{finance}[Overfitting is the everyday enemy of backtests]
Financial data are short and noisy (a few thousand daily returns, a low signal-to-noise ratio). A model with many free
parameters will always fit the past; the validation and test sets must respect time order, otherwise
information from the future leaks into the fit (\cref{ch:timeseries}). The lecture's advice to keep a test set
untouched until the end is, for a trading strategy, the discipline of an out-of-sample period used once.
\end{finance}
\begin{remark}[Forward reference]
The classical route to a controlled bias--variance trade-off in a regression is not to shrink the number of parameters but the
size of the coefficients: ridge regression, whose effective number of parameters is $\sum_jd_j^2/(d_j^2+\lambda)<p$
(\cref{ch05:sec-shrink}). \Cref{ch:regularized} uses the same idea to calibrate pricing and risk models.
\end{remark}

% ==================================================================
\section{Unsupervised learning: clustering}\label{ch16:sec-unsup}
The lecture spends one slide on unsupervised learning \ts{24}{23}{14:53}--\ts{24}{23}{15:58}: give the software a list of
customers with, say, ten attributes (annual spend, number of products bought, location, \dots) and ask it to group them so that
customers in a cluster have similar feature values. The value is in finding clusters one had not thought of.
No algorithm is named. The standard one is $k$-means, derived here \emph{(not discussed in class)}.

\begin{definition}[$k$-means]\label{ch16:def-kmeans}
Given points $\bm x_1,\dots,\bm x_n\in\R^d$ and $K$, choose an assignment $c:\{1..n\}\to\{1..K\}$ and centres
$\bm\mu_1,\dots,\bm\mu_K$ to minimise
\[
  J(c,\bm\mu)=\sum_{i=1}^n\|\bm x_i-\bm\mu_{c(i)}\|^2 .
\]
\end{definition}
\begin{proposition}[Lloyd's algorithm]\label{ch16:prop-lloyd}
Alternate
\begin{enumerate}[1.]
  \item $c(i)\leftarrow\arg\min_k\|\bm x_i-\bm\mu_k\|^2$ and
  \item $\bm\mu_k\leftarrow$ mean of the points assigned to $k$.
\end{enumerate}
Each step does not increase $J$, so $J$ decreases monotonically and the algorithm stops after finitely many steps at a local
minimum (in general not the global one).
\end{proposition}
\begin{proof}
For fixed centres, $J$ is a sum over $i$ of terms each minimised by the nearest centre, so step 1 cannot increase $J$.
For a fixed assignment and one cluster $C_k$, $\sum_{i\in C_k}\|\bm x_i-\bm m\|^2$ is a convex quadratic in $\bm m$ with
gradient $2\sum_{i\in C_k}(\bm m-\bm x_i)$, which vanishes at the mean, so step 2 cannot increase $J$.
There are finitely many assignments and $J$ is non-increasing, so the sequence of assignments cannot cycle except through
equal values, and the algorithm stops.
\end{proof}
\begin{example}[Six points]\label{ch16:ex-kmeans}
Take $(0,0),(0,1),(1,0)$ and $(5,5),(5,6),(6,5)$ with initial centres $(0,0)$ and $(5,5)$. The assignment is already
stable, the means are $(\tfrac13,\tfrac13)$ and $(\tfrac{16}3,\tfrac{16}3)$, and $J=\tfrac43+\tfrac43=\tfrac83\approx2.667$ (checked in python).
\end{example}
\begin{intuition}[Link with principal components]
$k$-means minimises a within-cluster sum of squares, exactly as PCA minimises a residual sum of squares over subspaces
(\cref{ch07:sec-geom}): PCA compresses to a linear subspace, $k$-means to $K$ points. It is a known result (Ding and He, 2004; not
proved here) that a continuous relaxation of the $k$-means objective is solved by the leading principal components of the data,
so PCA is often used to reduce the features before clustering. Unlike PCA, $J$ has many local minima; one restarts from different
initial centres.
\end{intuition}

% ==================================================================
\section{Optimisation: gradient descent}\label{ch16:sec-gd}
\subsection{The algorithm and Hull's example}
Training a model means minimising $\hat R(\bm\theta)$, a function of many parameters, for which there is no closed formula
(least squares is the exception). Hull describes the procedure in words \ts{24}{23}{31:08}: choose starting values and a \emph{learning rate}, compute the gradient,
which shows the direction in which the parameters can be improved, take a step against it, and repeat. In symbols,
\begin{equation}\label{ch16:eq-gd}
  \bm\theta_{k+1}=\bm\theta_k-\eta\,\nabla\hat R(\bm\theta_k).
\end{equation}
The first slide on the subject (slide~16) minimises the one-variable function $y=x^2-8x+20$
starting at $x_0=1$ \ts{24}{23}{34:17}; the minimum is at $x=4$ (there $y=4$). The gradient is $2x-8=-6$ at $x_0$, and the iterates printed on the slide, $2.2$ and $2.92$,
correspond to a learning rate $\eta=0.2$: $x_1=1+0.2\cdot6=2.2$, $x_2=2.2+0.2\cdot3.6=2.92$.

\begin{coursenote}[A slip in the spoken example]
In the spoken introduction Hull says the function is ``$6x^2-8x+20$'' with minimiser $x=4$ \ts{24}{23}{34:17}. That cannot be right: he also says the gradient is
$2x-8$, which is the derivative of $x^2-8x+20$, the function on the slide. (For $6x^2-8x+20$ the minimiser would be $2/3$.) The
slide is correct. Also, on slide~16 the vertical lines of the picture lie at about $1, 2.4, 2.9, 3.3$; the labels $2.2$ and $2.92$ in the
text should be trusted, since $2.2$ and $2.92$ are exactly what \eqref{ch16:eq-gd} gives with $\eta=0.2$.
\end{coursenote}

\subsection{Convergence, learning rate, scaling (not discussed in class)}
Near a minimum every smooth function is a quadratic, so the quadratic case explains the behaviour of the method.
Let $f(\bm\theta)=\tfrac12(\bm\theta-\bm\theta^*)\T A(\bm\theta-\bm\theta^*)+f^*$ with $A$ symmetric positive definite, eigenvalues
$0<\lambda_{\min}\le\dots\le\lambda_{\max}$. Then $\nabla f=A(\bm\theta-\bm\theta^*)$ and \eqref{ch16:eq-gd} gives, for the error $\bm e_k=\bm\theta_k-\bm\theta^*$,
\[
  \bm e_{k+1}=(I-\eta A)\bm e_k,\qquad \bm e_k=(I-\eta A)^k\bm e_0 .
\]
In the eigenbasis of $A$ each component is multiplied by $(1-\eta\lambda_j)$ at every step. Hence:
\begin{proposition}[Gradient descent on a quadratic]\label{ch16:prop-gd}
The iteration converges for every $\bm e_0$ iff $|1-\eta\lambda_j|<1$ for all $j$, i.e.\ $0<\eta<2/\lambda_{\max}$. The error is
multiplied by $r(\eta)=\max_j|1-\eta\lambda_j|$ per step; $r$ is minimal, and equal to $\frac{\kappa-1}{\kappa+1}$
with $\kappa=\lambda_{\max}/\lambda_{\min}$ the condition number, at $\eta=2/(\lambda_{\max}+\lambda_{\min})$.
For $\eta>1/\lambda_{\max}$ the components along the top eigenvector change sign at each step (oscillation).
\end{proposition}
\begin{proof}
The first two statements are the geometric-series criterion applied to each eigen-component. The maximum of $|1-\eta\lambda|$ over $[\lambda_{\min},\lambda_{\max}]$ is attained at an end point;
$r$ is decreasing in $\eta$ while $1-\eta\lambda_{\min}$ dominates and increasing when $\eta\lambda_{\max}-1$ dominates, so the minimum is where
$1-\eta\lambda_{\min}=\eta\lambda_{\max}-1$, which gives the stated $\eta$ and $r=\frac{\lambda_{\max}-\lambda_{\min}}{\lambda_{\max}+\lambda_{\min}}$.
\end{proof}
For Hull's example $A=2$ (one dimension), so $\lambda=2$, $e_k=(1-2\eta)^ke_0=(0.6)^k\cdot(-3)$ for $\eta=0.2$: convergence is geometric, with $1-2\eta=0.6$. The
stability condition is $\eta<1$; $\eta=\tfrac12$ reaches the minimum in one step; $\eta>\tfrac12$ overshoots and alternates around $4$; $\eta>1$ diverges
(\cref{ch16:fig-gd}). This is Hull's warning about the learning rate \ts{24}{23}{36:26}--\ts{24}{23}{37:28}: too small, and one never reaches the bottom in reasonable time; too big, and the iterates
oscillate from one side of the valley to the other.

\begin{figure}[ht]
\centering
\begin{tikzpicture}
\begin{axis}[width=0.47\linewidth,height=5.4cm,xlabel={$x$},ylabel={$y=x^2-8x+20$},xmin=0,xmax=8,ymin=0,ymax=22,grid=major,grid style={black!10}]
  \addplot[thick,color=black!60,smooth] coordinates {(0,20) (0.5,16.25) (1,13) (1.5,10.25) (2,8) (2.5,6.25) (3,5) (3.5,4.25) (4,4) (4.5,4.25) (5,5) (5.5,6.25) (6,8) (6.5,10.25) (7,13) (7.5,16.25) (8,20)};
  \addplot[only marks,mark=*,mark size=2pt,color=thmcol] coordinates {(1.000,13.000) (2.200,7.240) (2.920,5.166) (3.352,4.420)};
  \node[font=\footnotesize,anchor=south west,color=thmcol] at (axis cs:1.1,13.2) {$x_0=1$};
  \node[font=\footnotesize,anchor=south west,color=thmcol] at (axis cs:2.3,7.4) {$2.2$};
  \node[font=\footnotesize,anchor=south west,color=thmcol] at (axis cs:3.0,5.4) {$2.92$};
\end{axis}
\end{tikzpicture}\hfill
\begin{tikzpicture}
\begin{axis}[width=0.47\linewidth,height=5.4cm,xlabel={step $k$},ylabel={$x_k$},xmin=0,xmax=8,ymin=-2.6,ymax=10,xtick={0,2,4,6,8},grid=major,grid style={black!10},
  legend columns=2,legend style={font=\scriptsize,draw=black!30,at={(0.5,-0.32)},anchor=north}]
  \addplot[thick,color=defcol,mark=*,mark size=1.2pt] coordinates {(0,1) (1,1.3) (2,1.57) (3,1.813) (4,2.032) (5,2.229) (6,2.406) (7,2.565) (8,2.709)};
  \addplot[thick,color=intcol,mark=*,mark size=1.2pt] coordinates {(0,1) (1,2.2) (2,2.92) (3,3.352) (4,3.611) (5,3.767) (6,3.86) (7,3.916) (8,3.95)};
  \addplot[thick,color=fincol,mark=*,mark size=1.2pt] coordinates {(0,1) (1,6.4) (2,2.08) (3,5.536) (4,2.771) (5,4.983) (6,3.214) (7,4.629) (8,3.497)};
  \addplot[thick,color=thmcol,mark=*,mark size=1.2pt] coordinates {(0,1) (1,7.3) (2,0.37) (3,7.993) (4,-0.392) (5,8.832) (6,-1.315) (7,9.846) (8,-2.431)};
  \draw[dashed,black!50] (axis cs:0,4) -- (axis cs:8,4);
  \legend{$\eta=0.05$,$\eta=0.2$,$\eta=0.9$,$\eta=1.05$}
\end{axis}
\end{tikzpicture}
\caption{Left: the first three steps of gradient descent from $x_0=1$ with $\eta=0.2$ on the function of slide~16. Right: the iterates for four learning rates
($x_{k+1}-4=(1-2\eta)(x_k-4)$; dashed line: the minimiser $x=4$). $\eta=0.05$ is too small, $\eta=0.9$ oscillates but converges, and $\eta=1.05$ diverges.}\label{ch16:fig-gd}
\end{figure}

The role of $\kappa$ explains Hull's third remark, that the method \emph{works best when the variables have been scaled}
\ts{24}{23}{37:28}. For a least-squares loss $\tfrac12\|\bm y-X\bm\beta\|^2$ the matrix $A$ is $X\T X$, and its condition number is the ratio of the squared singular values of $X$. A
regressor measured in units 100 times larger changes the corresponding diagonal entry of $X\T X$ by $10^4$ and can make $\kappa$ huge; then the stable step size is dictated by
the stiff direction while the flat direction crawls (the rate is $\frac{\kappa-1}{\kappa+1}\approx1-2/\kappa$). Standardising the columns, as done before ridge in \cref{ch05:sec-shrink}, is the remedy in both cases.

\begin{remark}[Not covered: stochastic gradients]
With $n$ in the millions one computes the gradient of \eqref{ch16:eq-gd} on a random subset (\emph{mini-batch}) of the data at each step, and an \emph{epoch} is one pass over all batches. Hull uses
``epoch'' for one training iteration \ts{24}{23}{41:44}; in his examples the whole training set is used at every step of an epoch, as far as the slides say. The lecture also mentions ``clever ways'' of setting
the learning rate \ts{24}{23}{37:28} without naming them (adaptive schemes such as Adam belong there).
\end{remark}

\subsection{Early stopping is a form of ridge regression (not discussed in class)}
Hull's stopping rule, ``continue until the validation results diverge from the training results'' \ts{24}{23}{36:26}, has a precise meaning in the linear model. Let $X=UDV\T$ be the thin SVD as in \cref{ch05:sec-shrink}
and run \eqref{ch16:eq-gd} on $\tfrac12\|\bm y-X\bm\beta\|^2$ from $\bm\beta_0=\bm 0$. Along the $j$-th right singular vector the error contracts by $(1-\eta d_j^2)$ per step, so after $k$ steps the fitted values are
\[
  \hat{\bm y}_k=\sum_{j=1}^p\bigl[1-(1-\eta d_j^2)^k\bigr]c_j\bm u_j,\qquad\text{ridge: }\ \hat{\bm y}^{\text{ridge}}=\sum_{j=1}^p\frac{d_j^2}{d_j^2+\lambda}c_j\bm u_j,\quad c_j=\bm u_j\T\bm y .
\]
Both filters go from $0$ (for $d_j\to0$) to $1$ (for large $d_j$): directions of large variance are fitted first, small-variance directions last and only if we keep going. The filter $1-e^{-\eta kd^2}$ (the small-$\eta$ limit of the first) and
$d^2/(d^2+\lambda)$ agree qualitatively with $\lambda\approx1/(\eta k)$. For $\eta=0.05$ and $d^2\in\{0.1,1,10\}$ the two filters are $(0.049,0.401,0.999)$ for gradient descent after $k=10$ steps and $(0.048,0.333,0.833)$ for ridge with
$\lambda=2$; after $k=200$ steps and $\lambda=0.1$ they are $(0.633,1.000,1.000)$ and $(0.500,0.909,0.990)$ (python). Stopping early therefore shrinks small-variance directions like a ridge penalty, and the number of epochs is the hyperparameter chosen on
the validation set instead of $\lambda$.

% ==================================================================
\section{Neural networks}\label{ch16:sec-nn}
\subsection{Architecture}
A supervised model that predicts the output from the input ``in one go'' (linear regression is the model to keep in mind) is replaced by a sequence of stages
\ts{24}{23}{22:21}--\ts{24}{23}{24:34}: the input layer holds the known variables, $v_{0,1},v_{0,2},\dots$; a first hidden layer $v_{1,1},v_{1,2},\dots$ is computed from it, a second hidden layer from the first, and so on up to the output layer (slide~13).
Each neuron is a weighted combination of all the neurons of the previous layer, to which an \emph{activation function} is applied \ts{24}{23}{25:41}.

\begin{definition}[Feedforward neural network]\label{ch16:def-nn}
Let $L\ge1$, layer widths $m_0=d,m_1,\dots,m_L$, weight matrices $W^{(l)}\in\R^{m_l\times m_{l-1}}$, bias vectors $\bm b^{(l)}\in\R^{m_l}$ and activation functions $\varphi_l:\R\to\R$ applied componentwise. With $\bm a^{(0)}=\bm x$,
\[
  \bm z^{(l)}=W^{(l)}\bm a^{(l-1)}+\bm b^{(l)},\qquad \bm a^{(l)}=\varphi_l\bigl(\bm z^{(l)}\bigr),\qquad l=1,\dots,L,
\]
and the network output is $h_{\bm\theta}(\bm x)=\bm a^{(L)}$, $\bm\theta=\{W^{(l)},\bm b^{(l)}\}_l$. The parameters are the \emph{weights}
\ts{24}{23}{26:48}--\ts{24}{23}{27:51}.
\end{definition}
Hull's slide writes a neuron as a ``weighted average'' followed by $f$; the constant term $\bm b^{(l)}$ (the bias) is left implicit. The number of parameters is the sum of $m_l(m_{l-1}+1)$; for the Black--Scholes network below (5 inputs, three hidden layers of 20 neurons,
one output) it is $20\cdot6+2\cdot20\cdot21+21=981$, to be estimated from $6{,}000$ observations. In the volatility-surface network (3 inputs) it is $941$.

\paragraph{Activation functions} (slide~14) \ts{24}{23}{28:54}--\ts{24}{23}{30:01}: identity $f(y)=y$; sigmoid $f(y)=1/(1+e^{-y})\in(0,1)$; hyperbolic tangent
$f(y)=(e^{2y}-1)/(e^{2y}+1)\in(-1,1)$; ReLU $f(y)=\max(y,0)$. Two identities are used below:
\[
  \sigma'(y)=\sigma(y)\bigl(1-\sigma(y)\bigr)\le\tfrac14,\qquad
  \tanh y=\frac{e^{2y}-1}{e^{2y}+1}=2\sigma(2y)-1,\quad(\tanh)'=1-\tanh^2 .
\]
(The slide's formula for tanh is correct; both identities were confirmed numerically.) If all activations are the identity the network collapses to $\bm x\mapsto W^{(L)}\cdots W^{(1)}\bm x+\text{const}$, a linear model, which is Hull's remark
\ts{24}{23}{28:54}: it is why the nonlinearity is essential, and why the identity is nevertheless used on the \emph{last} layer, where it lets the output range over $\R$. A sigmoid output would confine the prediction to $(0,1)$, which is impossible for an option price
\ts{24}{23}{39:36}.
\begin{figure}[ht]
\centering
\begin{tikzpicture}
\begin{axis}[width=0.8\linewidth,height=5.2cm,xlabel={$y$},xmin=-4,xmax=4,ymin=-1.15,ymax=2.2,grid=major,grid style={black!10},
  legend columns=3,legend pos=north west,legend style={font=\footnotesize,fill=white,draw=black!30}]
  \addplot[thick,color=defcol,smooth] coordinates {(-4.00,0.0180) (-3.50,0.0293) (-3.00,0.0474) (-2.50,0.0759) (-2.00,0.1192) (-1.50,0.1824) (-1.00,0.2689) (-0.50,0.3775) (0.00,0.5000) (0.50,0.6225) (1.00,0.7311) (1.50,0.8176) (2.00,0.8808) (2.50,0.9241) (3.00,0.9526) (3.50,0.9707) (4.00,0.9820)};
  \addplot[thick,color=intcol,smooth] coordinates {(-4.00,-0.9993) (-3.50,-0.9982) (-3.00,-0.9951) (-2.50,-0.9866) (-2.00,-0.9640) (-1.50,-0.9051) (-1.00,-0.7616) (-0.50,-0.4621) (0.00,0.0000) (0.50,0.4621) (1.00,0.7616) (1.50,0.9051) (2.00,0.9640) (2.50,0.9866) (3.00,0.9951) (3.50,0.9982) (4.00,0.9993)};
  \addplot[thick,color=thmcol] coordinates {(-4,0) (0,0) (2.2,2.2)};
  \legend{sigmoid,$\tanh$,ReLU}
\end{axis}
\end{tikzpicture}
\caption{The activation functions of slide~14 (the ReLU is cut at the upper edge of the frame).}\label{ch16:fig-act}
\end{figure}
\begin{remark}[Not discussed: what the hidden layers buy]
A network with one hidden layer of sufficiently many sigmoid (or tanh) units can approximate any continuous function on a compact set to arbitrary accuracy (universal approximation theorem, Cybenko 1989, Hornik 1991); the theorem says nothing about how many units are
needed or how to find the weights. The lecture's point is more modest: several stages make the model ``a lot more flexible'' than one written-down equation \ts{24}{23}{24:34}. A network with no hidden layer, an identity output and squared loss \emph{is} the linear regression of
\cref{ch:regression}; with a sigmoid output and the cross-entropy loss $\ell=-y\log p-(1-y)\log(1-p)$, $p=\sigma(\bm w\T\bm x)$, it is logistic regression, the natural model for the loan-acceptance example of \cref{ch16:sec-what}, with the remarkably simple gradient
$\partial\ell/\partial z=p-y$ (from $\sigma'=\sigma(1-\sigma)$).
\end{remark}

\subsection{Backpropagation}
Gradient descent needs $\partial\ell/\partial\theta$ for \emph{every} weight, and there may be many. The lecture says only that each output is a function of a function of a function, so the gradient follows from the chain rule, with shortcuts that avoid computing each derivative from scratch
\ts{24}{23}{32:16}--\ts{24}{23}{33:16}, and that this is called backpropagation \ts{24}{23}{37:28}. The derivation is short \emph{(not discussed in class)}.

For one training pair and loss $\ell=\tfrac12\|\bm a^{(L)}-\bm y\|^2$ define the \emph{error signals} $\bm\delta^{(l)}=\partial\ell/\partial\bm z^{(l)}$.
\begin{proposition}[Backpropagation]\label{ch16:prop-bp}
\[
  \bm\delta^{(L)}=(\bm a^{(L)}-\bm y)\odot\varphi_L'(\bm z^{(L)}),\qquad
  \bm\delta^{(l)}=\bigl(W^{(l+1)\mathsf T}\bm\delta^{(l+1)}\bigr)\odot\varphi_l'\bigl(\bm z^{(l)}\bigr),
\]
\[
  \frac{\partial\ell}{\partial W^{(l)}}=\bm\delta^{(l)}\bm a^{(l-1)\mathsf T},\qquad \frac{\partial\ell}{\partial\bm b^{(l)}}=\bm\delta^{(l)} ,
\]
where $\odot$ is the componentwise product. One forward pass (to store the $\bm z^{(l)},\bm a^{(l)}$) and one backward pass give the whole gradient at a cost proportional to the number of weights.
\end{proposition}
\begin{proof}
The last formula is the chain rule for $\bm a^{(L)}=\varphi_L(\bm z^{(L)})$. Since $\bm z^{(l+1)}$ depends on $\bm z^{(l)}$ only through $\bm a^{(l)}=\varphi_l(\bm z^{(l)})$,
$\delta^{(l)}_j=\sum_i\delta^{(l+1)}_i\,\partial z^{(l+1)}_i/\partial z^{(l)}_j=\sum_i\delta^{(l+1)}_iW^{(l+1)}_{ij}\varphi_l'(z^{(l)}_j)$. Finally $z^{(l)}_i=\sum_jW^{(l)}_{ij}a^{(l-1)}_j+b^{(l)}_i$ gives $\partial\ell/\partial W^{(l)}_{ij}=\delta^{(l)}_ia^{(l-1)}_j$ and $\partial\ell/\partial b_i^{(l)}=\delta^{(l)}_i$.
\end{proof}
The name is the direction of the recursion: the error signal is propagated from the output layer to the input layer, and each layer's derivative reuses the one above. (The formulas were checked against a finite-difference derivative for a $3\to4\to1$ sigmoid network: the derivative of a weight agrees to
$8$ digits.) For sigmoid layers $\varphi'=a(1-a)\le\tfrac14$, so the product of $\varphi'$ along many layers shrinks the signal, the \emph{vanishing-gradient} problem, one reason why ReLU ($\varphi'\in\{0,1\}$) is popular in deep networks. Hull reports that all the
activation functions he tried behaved similarly well, and that sigmoid worked best in his two applications \ts{24}{23}{1:12:20}.

% ==================================================================
\section{Application: learning the Black--Scholes--Merton price}\label{ch16:sec-bsm}
\begin{casestudy}[A network that imitates Black--Scholes--Merton (slides 18--21)]
\textbf{Goal.} Test whether a network can recover a known pricing function from noisy data \ts{24}{23}{38:30}.
\textbf{Data.} $10{,}000$ European call prices from the Black--Scholes--Merton formula, with parameters drawn at random: $S\in[40,60]$, $K\in[0.5S,1.5S]$, $r\in[0,5\%]$, $\sigma\in[10\%,40\%]$, $T\in[0.25,2]$ years. To each price a
normal error with mean $0$ and standard deviation $0.15$ is added. The split is $6{,}000$ training, $2{,}000$ validation, $2{,}000$ test observations \ts{24}{23}{40:41}.
\textbf{Method.} Three hidden layers of $20$ neurons, sigmoid activation everywhere except the last layer, where the identity is used because a sigmoid output would lie in $(0,1)$ and prices do not \ts{24}{23}{39:36}. Training error and validation error are recorded as the number of epochs grows.
\textbf{Result.} The validation error stops improving and slowly worsens after a while while the training error keeps decreasing (slide~19); on the 50-epoch moving average (slide~20) the minimum is found and training is stopped after $2{,}575$ epochs
\ts{24}{23}{40:41}--\ts{24}{23}{41:44}. With only $10{,}000$ observations the network imitates the formula very well, in spite of the noise \ts{24}{23}{42:46}.
\textbf{Lesson.} The noise variance cannot be fitted away: the best possible mean squared error is $0.15^2=0.0225$ and the flexible model reaches it only if it stops before it starts fitting the noise. Early stopping is the safeguard (\cref{ch16:sec-gd}).
\end{casestudy}

\begin{coursenote}[Erratum: the axis of slides 19--20 is not the mean squared error]
The titles of slides 19 and 20 say ``mse'', but the vertical axis is labelled \emph{Mean Absolute Error} and the curves level off near $0.12$. For a noise $N(0,0.15^2)$ the smallest attainable mean absolute error is $0.15\sqrt{2/\pi}=0.1197$, while the smallest
attainable mean squared error is $0.0225$. So the plotted quantity is the mean \emph{absolute} error, and it is the title that is wrong (the transcript also says only ``mse''). Slide~26, by contrast, does show a mean squared error (values around $6.5\times10^{-5}$).
\end{coursenote}

\paragraph{Our reconstruction.} To see the mechanics we repeated the experiment (numpy, hand-written backpropagation as in \cref{ch16:prop-bp}, inputs standardised, full-batch Adam with learning rate $10^{-2}$; the slides do not say which optimiser
was used). After $4{,}000$ epochs the test-set mean squared error against the noisy prices is $0.0259$ against a floor of $0.0225$, and the root-mean-square distance between the network and the \emph{noise-free} Black--Scholes--Merton price is $0.054$ (the prices range up to about $25$). At $4{,}000$ epochs the validation
error was still decreasing, so no overfitting phase was reached in this short run; Hull's $2{,}575$ epochs are specific to his optimiser and learning rate, not a universal number.

\begin{finance}[Why imitate a formula that one already has?]
Hull's answer \ts{24}{23}{42:46}--\ts{24}{23}{44:53}: some derivatives have no formula and are valued by Monte Carlo simulation, which can take minutes, too slow with a client on the phone or for scenario analysis, where a portfolio is revalued thousands of times.
The recipe is
\begin{enumerate}[1.]
  \item run the simulation many times to build a table relating price to inputs, without adding any noise;
  \item train a network to reproduce the table;
  \item price by a forward pass through the network, which is nearly instantaneous.
\end{enumerate}
The Black--Scholes exercise is the test bed, because there the right answer is known.
The price paid is the training set: the surrogate is reliable only inside the range of inputs it was trained on.
\end{finance}

% ==================================================================
\section{Application: how the implied-volatility surface moves}\label{ch16:sec-vol}
The second application \ts{24}{23}{45:53}--\ts{24}{23}{50:07} concerns hedging with the delta of \cref{ch:bs} and the volatility surface of \cref{ch:volatility}. The standard delta hedge takes $\partial c/\partial S$ with the implied volatility held fixed. If the implied volatility tends to move with the asset price
(when the stock rises the firm is less levered, so volatility falls \ts{24}{23}{46:57}), the hedge that minimises the variance of the hedged position is different.

\begin{proposition}[Minimum-variance delta (not discussed in class)]\label{ch16:prop-mvdelta}
Let an option have value $f(S,\sigma)$ with $\delta=\partial f/\partial S$ and $\nu=\partial f/\partial\sigma$ (vega), and to first order $\Delta f\approx\delta\,\Delta S+\nu\,\Delta\sigma$ over a short interval. Write $\Delta\sigma=b\,\Delta S+\epsilon$ with
$b=\Cov(\Delta\sigma,\Delta S)/\Var(\Delta S)$ the regression slope (\cref{ch:regression}) and $\epsilon$ uncorrelated with $\Delta S$. The number of shares $h$ minimising $\Var(\Delta f-h\Delta S)$ is
\[
  \delta_{\mathrm{MV}}=\delta+\nu\,b .
\]
\end{proposition}
\begin{proof}
$\Delta f-h\Delta S=(\delta+\nu b-h)\Delta S+\nu\epsilon$, the two terms are uncorrelated, so the variance is $(\delta+\nu b-h)^2\Var(\Delta S)+\nu^2\Var\epsilon$, minimal at $h=\delta+\nu b$.
\end{proof}
Hull and White (2017) proposed an analytic model for $\E[\Delta\sigma\mid\Delta S]$ (its form is not on the slides). The lecture's neural network replaces it by a learned function
\ts{24}{23}{48:01}: \emph{inputs} the daily asset-price return, the moneyness (measured by delta) and the time to maturity; \emph{output} the change in implied volatility. For a call, $\nu>0$ and typically $b<0$, so $\delta_{\mathrm{MV}}<\delta$.

\begin{casestudy}[S\&P~500 options (slides 24--27)]
\textbf{Data.} S\&P~500 call options, 2014--2019, $100$ options sampled at random per day, $125{,}700$ observations, split $60/20/20\%$ ($75{,}420$, $25{,}140$, $25{,}140$).
\textbf{Method.} Three hidden layers of $20$ neurons; several activation functions tried, sigmoid best \ts{24}{23}{48:01}. Training stopped after $5{,}826$ epochs, where the validation mse turned up (slide~26: mse of about $6.5\times10^{-5}$ on the training and $6.8\times10^{-5}$ on the validation set).
\textbf{Results.} On the test set the network gave only a modest $11\%$ improvement over the Hull--White analytic model \ts{24}{23}{49:02}. But when the VIX index of day $t$ was added as a feature, a further improvement of over $60\%$ was found: the behaviour of the surface differs
in high-volatility and low-volatility environments \ts{24}{23}{49:02}--\ts{24}{23}{50:07}. The slides do not define ``improvement''; presumably it is the relative reduction of the mean squared error of the predicted volatility change.
\textbf{Lesson.} The 60\% is a result one is unlikely to find by ``playing around'' \ts{24}{23}{50:07}: the network was used to \emph{discover} that the regime matters, and the regime variable was then added by hand.
\end{casestudy}

% ==================================================================
\section{Reinforcement learning}\label{ch16:sec-rl}
\subsection{The model and the exploration--exploitation trade-off}
Reinforcement learning (RL) finds a strategy for a \emph{series} of decisions in a changing environment \ts{24}{23}{50:07}--\ts{24}{23}{51:11}; software of this kind beats the best human chess and Go players. The general model (slide~30) has a state $S_t$, an action $A_t$ taken in it,
a reward $R_{t+1}$ and a new state $S_{t+1}$ \ts{24}{23}{52:21}: $S_0\xrightarrow{A_0}(R_1,S_1)\xrightarrow{A_1}(R_2,S_2)\cdots$.

\begin{definition}[Value of an action]\label{ch16:def-Q}
Fix a policy $\pi$ (a rule $S\mapsto A$). With discount factor $\gamma\in[0,1]$ the \emph{return} is $G_t=\sum_{k\ge0}\gamma^kR_{t+k+1}$, and
$Q^\pi(s,a)=\E[G_t\mid S_t=s,A_t=a]$ is the value of taking $a$ in $s$ and following $\pi$ afterwards. The optimal $Q^*=\max_\pi Q^\pi$ satisfies the Bellman equation
\[
  Q^*(s,a)=\E\Bigl[R_{t+1}+\gamma\max_{a'}Q^*(S_{t+1},a')\ \Bigm|\ S_t=s,A_t=a\Bigr].
\]
\end{definition}
The Bellman equation is dynamic programming (the backward induction of \cref{ch:bs} on a tree) but written for an unknown environment, so that the expectation is replaced by observed trials. The lecture does not use these formulas; in Nim below $\gamma=1$ and the only reward is $\pm1$ at the end.

At each trial the agent must choose between \emph{exploitation}, taking the best action found so far, and \emph{exploration}, taking a random one to learn something new \ts{24}{23}{53:22}. With probability $\varepsilon$ it explores and with probability $1-\varepsilon$ it
exploits (\emph{$\varepsilon$-greedy}); $\varepsilon$ starts at $1$ and is multiplied at each trial by a decay factor such as $0.995$ (a hyperparameter set by trial and error, slide~31), so that after enough trials the agent exploits only \ts{24}{23}{54:25}. With the decay $0.9995$ used
in the Nim tables, $\varepsilon_n=0.9995^n$ is $0.607,\ 0.082,\ 0.0067$ after $1{,}000$, $5{,}000$, $10{,}000$ trials.

\subsection{Nim and the updating rule}
The game (slide~32) \ts{24}{23}{54:25}--\ts{24}{23}{55:31}: a pile of $n$ matches, two players alternately take $1$, $2$ or $3$; whoever must take the last match loses. Hull takes $n=8$, reward $+1$ for a win and $-1$ for a loss, an opponent who plays at random, state $S$ = matches left when it is the agent's turn, action
$A$ = number taken, and $Q(S,A)$ initialised at $0$.

\begin{proposition}[Perfect play in Nim]\label{ch16:prop-nim}
The player to move with $s$ matches loses against perfect play iff $s\equiv1\pmod 4$.
\end{proposition}
\begin{proof}
Induction on $s$. With $s=1$ the mover must take the last match and loses. If $s\not\equiv1$, i.e.\ $s-1\equiv1,2,3\pmod 4$ ($s\ge2$), the mover takes $(s-1)\bmod4\in\{1,2,3\}$ and leaves $4m+1$ matches to the opponent, who loses by the induction hypothesis. If $s\equiv1$ and $s>1$,
whatever $a\in\{1,2,3\}$ the mover takes leaves $s-a\equiv1-a\not\equiv1\pmod4$, a winning position for the opponent.
\end{proof}
This is the ``$4n+1$'' strategy of the lecture \ts{24}{23}{56:34}--\ts{24}{23}{57:39}: five matches is lost, nine is lost, and so on. Starting from $8$ the right move is to take $3$. The point of the exercise is to see whether a machine that knows nothing finds it
\ts{24}{23}{57:39}.

\paragraph{Updating} (slides 33--34) \ts{24}{23}{58:40}--\ts{24}{23}{1:00:45}. In the first trial of the slides the agent takes $1$, the opponent $3$, the agent $1$, the opponent $3$, and the agent wins. The pairs $(8,1)$ and $(4,1)$ visited are updated with the rule
\begin{equation}\label{ch16:eq-qupdate}
  Q(S,A)^{\text{new}}=Q(S,A)^{\text{old}}+\alpha\bigl[G-Q(S,A)^{\text{old}}\bigr],
\end{equation}
where $\alpha$ (e.g.\ $0.05$) is a hyperparameter: with $G=1$ and $Q^{\text{old}}=0$ the new value is $0.05$, as on the slide. It is an exponentially weighted moving average of the gains, and a stochastic-approximation step for the mean of $G$. There are two choices for the gain $G$:
\begin{itemize}
  \item the final reward of the trial (\emph{Monte Carlo method}): $G=R_{\text{end}}=\pm1$;
  \item the value of the next state under the best decision, $G=R+\gamma\max_{a'}Q(S',a')$ (\emph{temporal-difference learning}, or \emph{$Q$-learning}), which is the Bellman equation with the expectation replaced by one observation.
\end{itemize}
Monte Carlo waits until the game is over and needs the whole trajectory. Temporal-difference learning \emph{bootstraps}: it updates from a neighbouring estimate immediately. Under $\varepsilon$-greedy play the Monte Carlo average estimates the value of the
current exploratory policy, whereas the temporal-difference target estimates the value of the greedy one, which is why the latter learns faster while $\varepsilon$ is large. (That $Q$-learning converges to $Q^*$ if $\alpha$ decays suitably and every pair is visited infinitely often is a theorem of Watkins and Dayan; not shown here.)

\begin{example}[Exact solution against a random opponent (our reconstruction)]\label{ch16:ex-nim}
Against an opponent who picks uniformly among the legal moves, the optimal values follow from backward induction, $Q^*(s,a)=\frac1{n}\sum_{b}V^*(s-a-b)$ with $V^*(s)=\max_aQ^*(s,a)$, $V^*(0)$ read as $+1$ when the opponent took the last match
and $Q^*(s,s)=-1$ when the agent did. The values (python) are:
\begin{center}\small
\begin{tabular}{@{}c|ccccccccc@{}}
$a\backslash s$ & 1 & 2 & 3 & 4 & 5 & 6 & 7 & 8\\ \hline
1 & $-1$ & $1$ & $0$ & $\tfrac13$ & $\tfrac13$ & $1$ & $\tfrac79$ & $\tfrac79$\\
2 &  & $-1$ & $1$ & $0$ & $\tfrac13$ & $\tfrac13$ & $1$ & $\tfrac79$\\
3 &  &  & $-1$ & $1$ & $0$ & $\tfrac13$ & $\tfrac13$ & $1$\\
\end{tabular}
\end{center}
At $s=8$ the best move is $3$ with value $1$ (the opponent, left with $5$ and moving at random, always ends in a lost position for himself, since the agent then plays perfectly), while $1$ and $2$ give $7/9\approx0.778$, positive because the opponent is not perfect. Compare slides 35--36:
after $10{,}000$ trials the slide's row $s=8$ reads $(0.674,0.793,1.000)$ for Monte Carlo and $(0.813,0.910,1.000)$ for temporal-difference, so the action $3$ is found and the other two values are noisy estimates of $0.778$. Column $7$ is $0$ everywhere in the slides, correctly: after the agent's first move and the opponent's reply at most $6$ matches remain, so state $7$ is never met.
Some entries, e.g.\ $Q(3,1)=0.618$ in the temporal-difference table against the exact $0$, are stale: once $\varepsilon$ is small the agent rarely tries a bad action again, and its value keeps the noise of the few early trials.
\end{example}
\begin{figure}[ht]
\centering
\begin{tikzpicture}
\begin{axis}[width=0.92\linewidth,height=5.4cm,xlabel={number of trials},ylabel={mean $Q(8,3)$},xmin=0,xmax=10000,scaled x ticks=false,xtick={0,2000,4000,6000,8000,10000},ymin=0,ymax=1.08,grid=major,grid style={black!10},
  legend pos=south east,legend style={font=\footnotesize,fill=white,draw=black!30}]
  \addplot[thick,color=thmcol,mark=*,mark size=1.2pt] coordinates {(500,0.259) (1000,0.385) (1500,0.522) (2000,0.641) (2500,0.721) (3000,0.781) (3500,0.836) (4000,0.863) (4500,0.892) (5000,0.916) (5500,0.935) (6000,0.947) (6500,0.965) (7000,0.967) (7500,0.980) (8000,0.981) (8500,0.985) (9000,0.991) (9500,0.991) (10000,0.995)};
  \addplot[thick,color=intcol,mark=square*,mark size=1.2pt] coordinates {(500,0.942) (1000,0.999) (1500,1.000) (2000,1.000) (2500,1.000) (3000,1.000) (3500,1.000) (4000,1.000) (4500,1.000) (5000,1.000) (5500,1.000) (6000,1.000) (6500,1.000) (7000,1.000) (7500,1.000) (8000,1.000) (8500,1.000) (9000,1.000) (9500,1.000) (10000,1.000)};
  \legend{Monte Carlo,temporal difference}
\end{axis}
\end{tikzpicture}
\caption{Learning Nim with $\varepsilon_0=1$, decay $0.9995$, $\alpha=0.05$ (parameters of slides 35--36). The curves are the value of the winning move $Q(8,3)$, averaged over 100 independent runs of our own simulation. The temporal-difference rule learns the right value in about a thousand
trials, the Monte Carlo rule in about ten times more (compare the tables of slides 35 and 36).}\label{ch16:fig-nim}
\end{figure}
In our simulation of $200$ independent runs, the fraction of runs in which the greedy move at $s=8$ is $3$ after $1{,}000$, $5{,}000$, $10{,}000$ trials is $62.5\%,\,95\%,\,100\%$ for the Monte Carlo rule and $100\%$ at all three horizons for temporal-difference learning. This matches the lecture's remark that with Monte Carlo the machine is not yet sure after $1{,}000$ trials, is
fairly sure after $5{,}000$ and certain after $10{,}000$, and that temporal difference is faster \ts{24}{23}{1:01:52}--\ts{24}{23}{1:02:54}. Hull adds that when the opponent also learns, moves $1$ and $2$ from $8$ turn negative \ts{24}{23}{1:03:55}: against a perfect opponent $Q(8,1)$ and $Q(8,2)$ are $-1$, because $7$ and $6$ leave a winning position.

\subsection{Deep $Q$-learning}
When there are many states or actions the table is filled only partly \ts{24}{23}{1:03:55}--\ts{24}{23}{1:05:00}, so $Q$ is estimated as a function by a network $Q_{\bm\theta}(S,A)$, trained to minimise the sum of squared differences between the network and the target \ts{24}{23}{1:05:00} (slide~37):
\[
  L(\bm\theta)=\sum_{\text{observed }(S,A,R,S')}\Bigl(R+\gamma\max_{a'}Q_{\bm\theta}(S',a')-Q_{\bm\theta}(S,A)\Bigr)^2 ,
\]
where, as in \eqref{ch16:eq-qupdate}, the target is held fixed when differentiating (a ``semi-gradient''; common implementations keep a lagged copy of $\bm\theta$ inside the target for stability; not discussed in class). This is nonlinear least squares, i.e.\ \cref{ch16:sec-gd,ch16:sec-nn}
with a moving target. The table of Nim is the special case with one parameter per pair $(S,A)$ and a constant step: \eqref{ch16:eq-qupdate} is one gradient step on the square $\tfrac12(G-Q(S,A))^2$ with learning rate $\alpha$.

\subsection{Hedging derivatives by reinforcement learning}
Hull's research question \ts{24}{23}{1:05:00}--\ts{24}{23}{1:06:05}: the traditional hedge uses the Greek letters, i.e.\ derivatives of the portfolio value with respect to a variable over a very short horizon. RL instead looks several periods ahead and chooses today's hedge to work well over the whole
horizon. \textbf{Findings} (slide~39) \ts{24}{23}{1:06:05}--\ts{24}{23}{1:10:14}:
\begin{itemize}
  \item for vanilla options it does about as well as Greek-letter hedging when there are no transaction costs, but it \emph{saves transaction costs}, in some cases significantly;
  \item for exotic options such as barrier options it produces better results even with no transaction costs;
  \item the user \emph{chooses the objective function}: e.g.\ VaR$_{95}$ or CVaR$_{95}$ of the hedging loss, which the Greeks cannot do;
  \item it is robust and works well in stressed periods; JP Morgan is the main competitor doing this research;
  \item disadvantage: it is much more computationally demanding, minutes rather than nothing for a Greek, but computers are getting faster.
\end{itemize}
\begin{intuition}[Why looking ahead saves costs (own illustration, not from the lecture)]
Hull's example \ts{24}{23}{1:07:09}: the delta is $0.5$, rises to $0.7$, and there is a 50\% chance that at the next hedging date it will be back at $0.5$. With a proportional cost $c$ per share traded, rebalancing at once costs $0.2c$ now and, with probability $\tfrac12$, another $0.2c$
later: expected cost $0.3c$. Waiting costs $0$ now and $0.2c$ only if the delta stays at $0.7$: expected cost $0.1c$. The price of waiting is a hedge that is $0.2$ shares short for one period, which adds a variance $0.2^2\Var(\Delta S)=0.04\Var(\Delta S)$ to the hedged position. Waiting is better
when $0.2c$ exceeds the value put on this variance, i.e.\ when transaction costs are large, which is Hull's conclusion \ts{24}{23}{1:08:13}. A Greek-letter hedger, who looks one period ahead, always rebalances; a multi-period agent can trade the two effects off.
\end{intuition}
\begin{coursenote}[VaR and CVaR: a misstatement in the transcript]
Explaining the objective functions, Hull says that with VaR$_{95}$ one wants ``the 95th percentile point of the loss distribution to be as high as possible'' \ts{24}{23}{1:09:13}. A risk manager wants it \emph{as low as possible}: the objective is to \emph{minimise} the
95th percentile of the loss (as for the CVaR$_{95}$ described in the next sentence, the average loss beyond that percentile, ``as low as possible''). The definitions and the properties of both measures are in \cref{ch:var}.
\end{coursenote}

% ==================================================================
\section{The questions at the end}
The remaining minutes are questions \ts{24}{23}{1:11:20}--\ts{24}{23}{1:15:24}. Asked whether the choice of activation matters, Hull answers that sigmoid has worked well, perhaps by chance, that the alternatives tried gave results close to each other, and that the size of the model (neurons per layer, number of layers) matters more.
Asked whether a two-layer network can be extended to three layers without starting again, he believes one must start from scratch. (This is the practice in his examples; in the wider literature weights of a smaller network are sometimes reused as an initialisation, not discussed.) He points to the Python code for the book's examples on his
Toronto website \ts{24}{23}{1:13:20}, and, asked how the objective function is chosen, says that it depends on the problem: to match an exotic option price one minimises the difference between the network and target price \ts{24}{23}{1:14:20}--\ts{24}{23}{1:15:24}.

\begin{finance}[Summary: where the tools sit in the course]
\begin{itemize}
  \item \emph{Regression and shrinkage} (\cref{ch:regression}): squared loss, linear model, ridge penalty; the network is the nonlinear generalisation and early stopping is the ridge penalty in disguise (\cref{ch16:sec-gd}).
  \item \emph{PCA} (\cref{ch:pca}): dimension reduction on the features, also the linear cousin of clustering.
  \item \emph{Volatility and Black--Scholes} (\cref{ch:volatility,ch:bs}): the network learns the pricing function and the response of the surface; the minimum-variance delta is a regression coefficient.
  \item \emph{Regularised models} (\cref{ch:regularized}): the finance-industry treatment of the same bias--variance balance, for pricing and risk model calibration.
  \item \emph{VaR} (\cref{ch:var}): the loss functionals used as objectives in RL hedging.
\end{itemize}
\end{finance}

\begin{keyformulas}
\begin{align*}
&\text{Empirical risk:} && \hat R(\bm\theta)=\tfrac1n\textstyle\sum_i\ell(y_i,h_{\bm\theta}(\bm x_i))\\
&\text{Bias--variance:} && \E(y-\hat h)^2=\sigma^2+(f-\E\hat h)^2+\Var\hat h\\
&\text{Projection fit:} && \text{test}=\sigma^2+b^2+\tfrac pn\sigma^2,\quad\text{train}=\sigma^2+b^2-\tfrac pn\sigma^2\\
&\text{$k$-means:} && J=\textstyle\sum_i\|\bm x_i-\bm\mu_{c(i)}\|^2\ \text{(Lloyd: monotone decrease)}\\
&\text{Gradient descent:} && \bm\theta_{k+1}=\bm\theta_k-\eta\nabla\hat R(\bm\theta_k)\\
&\text{Quadratic case:} && \bm e_{k+1}=(I-\eta A)\bm e_k;\ 0<\eta<\tfrac2{\lambda_{\max}};\ r_{\min}=\tfrac{\kappa-1}{\kappa+1}\\
&\text{Early stopping:} && \text{filter }1-(1-\eta d^2)^k\ \leftrightarrow\ \text{ridge }\tfrac{d^2}{d^2+\lambda},\ \lambda\approx\tfrac1{\eta k}\\
&\text{Network layer:} && \bm z^{(l)}=W^{(l)}\bm a^{(l-1)}+\bm b^{(l)},\quad\bm a^{(l)}=\varphi(\bm z^{(l)})\\
&\text{Activations:} && \sigma(y)=\tfrac1{1+e^{-y}},\ \sigma'=\sigma(1-\sigma),\ \tanh y=2\sigma(2y)-1\\
&\text{Backpropagation:} && \bm\delta^{(l)}=(W^{(l+1)\mathsf T}\bm\delta^{(l+1)})\odot\varphi'(\bm z^{(l)})\\
& && \partial\ell/\partial W^{(l)}=\bm\delta^{(l)}\bm a^{(l-1)\mathsf T}\\
&\text{Min-variance delta:} && \delta_{\mathrm{MV}}=\delta+\nu\,\Cov(\Delta\sigma,\Delta S)/\Var(\Delta S)\\
&\text{Bellman:} && Q^*(s,a)=\E[R+\gamma\max_{a'}Q^*(S',a')]\\
&\text{$Q$ update:} && Q\leftarrow Q+\alpha[G-Q]\\
& && G=R_{\text{end}}\ (\text{MC}),\quad G=R+\gamma\max_{a'}Q(S',a')\ (\text{TD})\\
&\text{Exploration:} && \varepsilon_n=\varepsilon_{n-1}\cdot\text{decay},\quad \varepsilon_n=0.9995^n
\end{align*}
\end{keyformulas}

% ==================================================================
\section*{Exercises}
\addcontentsline{toc}{section}{Exercises}

\begin{exercise}[not from the course: gradient descent on an ill-conditioned quadratic]
Let $f(\bm\theta)=\tfrac12(\theta_1^2+10\theta_2^2)$.
\begin{enumerate}[(a)]
  \item For which $\eta$ does gradient descent converge?
  \item Find the $\eta$ that minimises the contraction factor and this factor.
  \item Starting at $(1,1)$, how many steps are needed to bring $\|\bm\theta\|$ below $10^{-3}$ with that $\eta$,
      and with $\eta=0.05$?
  \item Rescale $\theta_2=\theta_2'/\sqrt{10}$ and repeat: what has changed?
\end{enumerate}
\end{exercise}
\begin{exercise}[not from the course: backpropagation]
A network has one input, one hidden layer with two $\tanh$ neurons, and a linear output, $h=w_1\tanh(u_1x+c_1)+w_2\tanh(u_2x+c_2)+d$, with loss $\ell=\tfrac12(h-y)^2$. Using \cref{ch16:prop-bp}, compute $\partial\ell/\partial u_1$, $\partial\ell/\partial c_1$, $\partial\ell/\partial w_1$ and check with a finite-difference in
python for random values of the parameters.
\end{exercise}
\begin{exercise}[not from the course: optimism of any linear smoother]
Let $\hat{\bm y}=S\bm y$ for a fixed matrix $S$ (ridge regression, \cref{ch05:sec-shrink}, is an example with $S=X(X\T X+\lambda I)^{-1}X\T$), $\bm y=\bm f+\bm\varepsilon$, $\Cov\bm\varepsilon=\sigma^2I$. Show that $\frac1n\E\|\bm y^*-\hat{\bm y}\|^2-\frac1n\E\|\bm y-\hat{\bm y}\|^2=\frac{2\sigma^2}{n}\tr S$, where
$\bm y^*=\bm f+\bm\varepsilon^*$ is an independent copy, so that $\tr S$ plays the role of the number of parameters in \cref{ch16:prop-bv}. Evaluate $\tr S$ for ridge in terms of the singular values.
\end{exercise}
\begin{exercise}[not from the course: $k$-means]
\leavevmode
\begin{enumerate}[(a)]
  \item Show that the mean minimises $\sum_{i\in C}\|\bm x_i-\bm m\|^2$ and that the minimum equals $\frac1{2|C|}\sum_{i,j\in C}\|\bm x_i-\bm x_j\|^2$.
  \item Take the four points $(0,0),(4,0),(0,1),(4,1)$ and $K=2$. Show that ``left pair/right pair'' and ``bottom pair/top pair'' are both stable
      under Lloyd's algorithm, and compute $J$ for each ($1$ and $16$): a local minimum need not be a global one.
\end{enumerate}
\end{exercise}
\begin{exercise}[not from the course: Nim]
\leavevmode
\begin{enumerate}[(a)]
  \item Verify \cref{ch16:prop-nim} for $s\le13$ by listing the winning moves.
  \item Verify by hand the exact values $Q^*(5,\cdot)=(\tfrac13,\tfrac13,0)$ and $Q^*(7,\cdot)=(\tfrac79,1,\tfrac13)$ of \cref{ch16:ex-nim}.
  \item For which starting number of matches $n\le13$ is there a first move that wins even if the opponent plays perfectly?
\end{enumerate}
\end{exercise}
\begin{exercise}[not from the course: one $Q$-learning step]
In Nim with $\alpha=0.05$, $\gamma=1$, let $Q(8,3)=0.5$ and $Q(4,\cdot)=(0.1,0.2,0.4)$. A trial goes $8\xrightarrow{3}5\to4$ (opponent takes $1$). Compute the updated $Q(8,3)$ by
\begin{enumerate}[(a)]
  \item the temporal-difference rule and
  \item the Monte Carlo rule, if the trial ends with a win.
\end{enumerate}
\end{exercise}
\begin{exercise}[not from the course: minimum-variance delta]
A call has $\delta=0.60$ and vega $\nu=25$ (price change per unit change of $\sigma$ in absolute terms, e.g.\ $0.01$ of volatility is $0.25$). Suppose over a day $\Delta\sigma=-0.008\,\Delta S+\epsilon$ ($\Delta S$ in dollars). Compute $\delta_{\mathrm{MV}}$.
How does the variance of the hedged position change when $h=\delta_{\mathrm{MV}}$ replaces $h=\delta$, if $\Var(\Delta S)=1$ and $\Var(\epsilon)=10^{-5}$?
\end{exercise}

% !TEX root = ../main.tex
\chapter{Stochastic Calculus (It\^o Integral and It\^o's Lemma)}\label{ch:stochcalc}
\begin{paracol}{2}

\begin{sources}
\begin{tabularx}{\linewidth}{@{}l l L@{}}
\toprule
\textbf{Lecture} & \textbf{Speaker} & \textbf{Content}\\
\midrule
2024 L24 & P.~Kempthorne & Brownian motion with drift and the one-step analysis; filtrations; the It\^o integral built from constant, step and adapted step integrands; the worked case $\int B\,\dd B$; the general integral, its martingale property, the isometry and the Riemann-sum (left-point) representation; Gaussian integrals; It\^o's formula for $f(B_t)$ and $f(t,B_t)$, integrals by anti-derivatives, the PDE condition for martingales, the exponential martingale, ruin with drift; standard processes (last slide of the lecture).\\
2024 L25 (to 18:00) & P.~Kempthorne & Recap of It\^o's formula; geometric Brownian motion as an SDE, the log price and the drift correction $\mu-\tfrac12\sigma^2$, mean and variance of the price, annualised return, simulation of 50 paths. (From 19:01 the lecture derives the Black--Scholes equation and solves the heat equation; that part belongs to \cref{ch:sde,ch:bs}, only the derivation by It\^o's lemma is summarised here.)\\
2013 L18 & C.~Lee & Taylor expansion and $(\dd B)^2=\dd t$; It\^o's lemma for $f(t,B_t)$ and for $f(t,X_t)$; examples $B^2$, $e^{\mu t+\sigma B}$, geometric Brownian motion without drift; integration as the inverse of differentiation, left-point sums, adapted processes; normality, isometry and martingale property; equivalent measures, Radon--Nikod\'ym derivative and Girsanov's theorem (with a comment by P.~Kempthorne at 1:10:39).\\
\bottomrule
\end{tabularx}

\smallskip
\textbf{Files.} Slides \file{mit18_642_f24_lec19_1.pdf} (26 slides, 84 pages with overlays; this is Lecture~24, not 19);
\file{mit18_642_f24_lec19_2.pdf} (7 pages: the R simulation of Brownian motion with drift, of $e^{W}$ and of the cumulative
average log return); 2013 note \file{MIT18_S096F13_lecnote18.pdf} (7 pages).
\textbf{Problem sets.} 2013 PS8, all problems (A-1 to A-5, B-1 to B-3); Problem~A-2 also opens 2024 PS4 and is treated
in \cref{ch:sp2}, where its solution is developed. 2024 PS5 (L\'evy processes, simulations) belongs to \cref{ch:sp2}; there is no 2024 problem set on
stochastic calculus. Each exercise follows the section it practises.
\textbf{Not covered in class.} 2024 does not treat change of measure, 2013 does not construct the integral (the isometry is
stated without proof); the construction of \cref{ch17:sec-general}, the proof sketch of It\^o's lemma with control of the
remainder, the multi-dimensional formula and the Stratonovich integral are additions of these notes (marked as such).
\end{sources}

\section*{Motivation}
A stock price is not differentiable in time, and a Brownian path is not differentiable either
(\cref{ch:sp2}). Yet finance needs to speak about \emph{integrals} such as the gain $\int\theta_t\,\dd S_t$ of a trading
strategy that holds $\theta_t$ units of an asset whose price is $S_t$, about the \emph{differential} of an option price
$G(S_t,t)$ when $S$ is random, and about the solution of $\dd S=\mu S\,\dd t+\sigma S\,\dd B$. Ordinary calculus fails because
the squared increments of $B$ do not vanish faster than $\dd t$: $(\dd B)^2=\dd t$
(\cref{ch11:thm-qv,ch11:eq-dB2}). Stochastic calculus is the calculus that keeps this one extra term and is otherwise
as close to ordinary calculus as possible. In physical language: it is the calculus of a particle
driven by white noise, and the extra term is the It\^o (or ``noise-induced'') drift that a physicist has to decide
about when writing a Langevin equation.
\begin{itemize}
  \item The \emph{It\^o integral} $\int_0^tf\,\dd B$ evaluates the integrand at the \emph{left} end of each small
        interval, so that the integrand never looks into the future (\cref{ch17:sec-simple,ch17:sec-general}); the price is
        that classical rules change, e.g.\ $\int_0^tB\,\dd B=\tfrac12(B_t^2-t)$ (\cref{ch17:sec-BdB}), which is the identity the Dickey--Fuller
        distribution of \cref{ch09:sec-df} is built on.
  \item \emph{It\^o's lemma} is the chain rule of this calculus, $\dd f(B_t)=f'(B_t)\dd B_t+\tfrac12f''(B_t)\dd t$
        (\cref{ch17:sec-lemma}). It converts every question about functions of Brownian motion into a
        drift and a martingale part (\cref{ch17:sec-examples}) and solves geometric Brownian motion
        (\cref{ch17:sec-gbm}).
  \item With several noise sources the same idea gives the product rule and the multi-dimensional formula
        (\cref{ch17:sec-multi}); and Girsanov's theorem (\cref{ch17:sec-girsanov}) says that a drift can be removed
        by changing the probability measure, which is the mathematical heart of risk-neutral pricing
        (\cref{ch:bs}).
\end{itemize}
We use the background of \cref{ch:sp1} (martingales, conditional expectation) and \cref{ch:sp2} (definition of Brownian
motion, quadratic variation, filtrations). The stochastic differential equations built on this calculus are the subject of
\cref{ch:sde}.

% ==================================================================
\section{From Brownian motion with drift to It\^o processes}\label{ch17:sec-drift}

Let $B$ be a standard Brownian motion and put
\begin{equation}\label{ch17:eq-bmdrift}
  X(t)=\mu t+\sigma B(t),\qquad t\ge0,
\end{equation}
Brownian motion with drift $\mu$ and variance parameter $\sigma^2$ (\cref{ch11:def-drift}) \ts{24}{24}{1:05}.
The increments are independent and $\Var[X(t)-X(s)]=\sigma^2|t-s|$. The lecture looks at a small step,
$X(t+\Delta t)-X(t)=\mu\Delta t+\sigma\Delta B$ with $\Delta B=B(t+\Delta t)-B(t)$ \ts{24}{24}{4:09}; the increment
$\Delta X$ is \emph{exactly} $N(\mu\Delta t,\sigma^2\Delta t)$ and
\begin{equation}\label{ch17:eq-onestep}
  \E[(\Delta X)^2]=\sigma^2\Delta t+(\mu\Delta t)^2=\sigma^2\Delta t+o(\Delta t),\qquad
  \E[|\Delta X|^c]=o(\Delta t)\quad(c>2)
\end{equation}
\ts{24}{24}{6:19}: only the \emph{second} moment of the increment is of the order of the time step, all higher
ones are negligible. (The lecture slide writes $\E[(\Delta X)^c]$ without absolute value; for odd $c$ that is a different
quantity, but it is also $o(\Delta t)$.) This is the whole content of ``$(\dd B)^2=\dd t$'' in probabilistic language.

In differential notation \eqref{ch17:eq-bmdrift} reads $\dd X=\mu\,\dd t+\sigma\,\dd B$, and its solution is
\begin{equation}\label{ch17:eq-sol}
  X_t=X_0+\int_0^t\mu\,\dd s+\int_0^t\sigma\,\dd B_s=X_0+\mu t+\sigma\bigl(B_t-B_0\bigr)
\end{equation}
\ts{24}{24}{7:24}. The differential form is only shorthand for the integral form. The lecture
wants to allow a drift and a volatility that depend on time and on the state,
$\dd X_t=\mu(X_t,t)\,\dd t+\sigma(X_t,t)\,\dd B_t$: an \emph{It\^o process}, the model for a price, an interest rate or the value of a derivative
that depends on time to maturity and on the level of the underlying \ts{24}{24}{9:32}. The second integral
in \eqref{ch17:eq-sol} is then no longer a trivial difference of the endpoints; a student asks exactly this and the lecture answers that
It\^o integrals are the missing definition \ts{24}{24}{10:36}.

\begin{intuition}[Why no ordinary integral: white noise]
If $B$ were differentiable, $\xi=\dd B/\dd t$ would be a noise and $\dd X=\mu\,\dd t+\sigma\xi\,\dd t$ a Langevin equation. The path of $B$
is nowhere differentiable, and by \cref{ch11:cor-tv} its total variation is infinite, so $\int f\,\dd B$ cannot be defined path by
path as a Riemann--Stieltjes integral. Physicists write $\langle\xi(t)\xi(s)\rangle=\delta(t-s)$: $\xi$ is a generalised function, and a
Langevin equation with \emph{multiplicative} noise $\sigma(X)\xi$ is ambiguous until one says at which point of
the small interval $\sigma$ is evaluated. Stochastic calculus resolves the ambiguity by a rule; the It\^o rule is
``at the beginning'' (\cref{ch17:sec-simple,ch17:sec-BdB}), the alternative ``at the middle'' gives the Stratonovich integral
that obeys the ordinary chain rule.
\end{intuition}

% ==================================================================
\section{Filtrations, adapted processes and integrals of step processes}\label{ch17:sec-simple}

\subsection{The probability set-up}\label{ch17:sec-setup}
As in the lecture, let $(\Omega,\{\mathcal F_t\},\Prob)$ be the probability model of the Brownian motion: $\Omega$ the set of
paths $\omega$, $\mathcal F_t$ the $\sigma$-field of events determined by the path up to time $t$, i.e.\ by $\{B_s,s\le t\}$
(the information available at $t$; the filtration was introduced in \cref{ch11:sec-adapted}), and $B_t=B_t(\omega)$ a
specific path \ts{24}{24}{12:42}. Below $B$ is a Brownian motion
\emph{with respect to} this filtration: $B_t-B_s$ is $N(0,t-s)$ and independent of $\mathcal F_s$ for $s<t$.

\begin{definition}[Adapted process]\label{ch17:def-adapted}
A process $\{\Delta_t\}$ is \emph{adapted} (to $B$, or to the filtration $\{\mathcal F_t\}$) if for every $t$ the random variable $\Delta_t$
depends only on $\{B_s:s\le t\}$, i.e.\ is $\mathcal F_t$-measurable \ts{13}{18}{40:10}.
\end{definition}
Adapted means \emph{non-anticipating}: a strategy that decides the position at $t$ from what has been observed until $t$.
The 2013 lecture gives examples for a process $X$ \ts{13}{18}{42:26}: $\Delta_t=X_t$ is adapted;
$\Delta_t=\min\{X_t,c\}$ is adapted; $\Delta_t=X_{t+1}$ is not (it looks one unit ahead); $\Delta_t=\max_{0\le s\le T}X_s$, for a fixed
$T>t$, is not (it uses the future until $T$). In the same spirit $\Delta_t=X_{t/2}$ is adapted and $\Delta_t=X_{2t}$ is not. If $\tau$ is a stopping time (\cref{ch04:def-stop}),
the \emph{stopped} process $X_{t\wedge\tau}$ is adapted; the note writes $X_\tau$, see the erratum in \cref{ch17:sec-compare}.

\begin{finance}[Trading strategies must be adapted]
If $S_t$ is a price and $\Delta_t$ the number of units held at $t$, the strategy makes sense only if $\Delta$ is adapted: one cannot
buy more just before a jump that has not yet happened. The 2013 note's example is the strategy $\Delta_t=\pm1$ (buy or sell one
share): it is admissible only if the sign is decided from past information \ts{13}{18}{39:05}. The discrete counterpart is the
daily P\&L of \cref{eq:pnl}, where the holdings $\bm q(t+1)$ are chosen at the end of day $t$: the
sum $\sum_tq(t+1)\Delta p(t)$ is a left-point sum.
\end{finance}

\begin{exercise}[PS8 (2013), Problem A-3: adapted processes]\label{ch17:ex-a3}
Let $X_t$ be a given stochastic process. Which of the following processes $Y_t$ are adapted to $X_t$?
\begin{enumerate}[(a)]
  \item $Y_t=X_{T-t}$ for $t\le T$, for a fixed $T$.
  \item $Y_t=\max_{0\le s\le2t}X_s$.
  \item $Y_t=|\{i\in[0,t]:X_i\ge0\}|$ (the amount of time in $[0,t]$ during which $X$ is non-negative; the printed set-notation is ambiguous).
\end{enumerate}
\end{exercise}

\subsection{Integrals of step processes}\label{ch17:sec-stepint}
The lecture builds the integral $I(f)=\int_0^Tf\,\dd B$ in stages of increasing generality
\ts{24}{24}{19:15}. For each path $\omega$ one wants $I(f)(\omega)$, a \emph{random variable} with a mean and a
variance.
\begin{enumerate}[(a)]
  \item \emph{Constant integrand} $f\equiv1$: $I(f)=\int_0^T\dd B_s=B_T-B_0=B_T$, with mean $0$ and variance $T$
        \ts{24}{24}{21:17}. For an indicator $f=\ind(a<s\le b)$ the integral is $B_b-B_a$, with variance $b-a$
        \ts{24}{24}{25:29}.
  \item \emph{Deterministic step function} $f(s)=\sum_{i=0}^{n-1}a_i\,\ind(t_i<s\le t_{i+1})$ with constants $a_i$ and
        $0=t_0<t_1<\dots<t_n=T$:
        \begin{equation}\label{ch17:eq-stepint}
          I(f)=\sum_{i=0}^{n-1}a_i\bigl(B_{t_{i+1}}-B_{t_i}\bigr),\qquad
          \E[I(f)]=0,\quad \Var I(f)=\sum_ia_i^2(t_{i+1}-t_i)
        \end{equation}
        \ts{24}{24}{26:33}. The $a_i$ can be read as the number of shares held during $(t_i,t_{i+1}]$ of an asset
        following Brownian motion, and $I(f)$ as the gain or loss. The slides write the last time as $t_{n+1}=T$; with $n$ intervals it is $t_n=T$.
  \item \emph{Adapted step process}: the levels $a_i=a_i(\omega)$ are random, \emph{$\mathcal F_{t_i}$-measurable} and with $\E[a_i^2]<\infty$
        \ts{24}{24}{28:40}. The integral is again \eqref{ch17:eq-stepint}, now with random $a_i$
        \ts{24}{24}{31:54}.
\end{enumerate}
Case (c) is the only new idea, and it is the one that fixes the meaning of the integral: \emph{the level $a_i$ is known at the left
end $t_i$ of the interval on which it is used}. We call such an $f=\sum_ia_i\ind(t_i<s\le t_{i+1})$ a \emph{simple process}
(``non-anticipatory'' in the lecture) and write $I_t(f)=\sum_ia_i(B_{t\wedge t_{i+1}}-B_{t\wedge t_i})$ for the integral up to time $t$.

\begin{proposition}[Simple integrands: mean, martingale, isometry]\label{ch17:prop-simple}
Let $f$ be a simple process. Then $I_t(f)$ is linear in $f$, and
\begin{enumerate}[(i)]
  \item $\{I_t(f)\}_{t\le T}$ is a martingale, in particular $\E[I_t(f)]=0$;
  \item (\emph{It\^o isometry}) $\E\bigl[I_T(f)^2\bigr]=\E\int_0^Tf(s)^2\,\dd s=\sum_i\E[a_i^2](t_{i+1}-t_i)$; more generally
        $\E[(I_t-I_s)^2\mid\mathcal F_s]=\E[\int_s^tf^2\,\dd u\mid\mathcal F_s]$.
\end{enumerate}
\end{proposition}
\begin{proof}
Write $\Delta_iB=B_{t_{i+1}}-B_{t_i}$, independent of $\mathcal F_{t_i}$ with mean $0$ and variance $\Delta t_i=t_{i+1}-t_i$.
(i) Take $s<t$; refining the grid we may assume $s$ and $t$ are grid points. Then $I_t-I_s=\sum_{s\le t_i<t}a_i\Delta_iB$ and, for each such $i$,
by the tower property, $\E[a_i\Delta_iB\mid\mathcal F_s]=\E\bigl[a_i\,\E[\Delta_iB\mid\mathcal F_{t_i}]\bigm|\mathcal F_s\bigr]=0$, because $a_i$ is
$\mathcal F_{t_i}$-measurable and $\Delta_iB$ is independent of $\mathcal F_{t_i}$ with mean $0$
(this is the answer given in the lecture to a student who asked how the mean can vanish when $a_i$ is random
\ts{24}{24}{35:04}). (The products are integrable since $\E[|a_i\Delta_iB|]\le(\E[a_i^2]\Delta t_i)^{1/2}$.)

(ii) Expand $I_T^2=\sum_ia_i^2(\Delta_iB)^2+2\sum_{i<j}a_i\Delta_iB\,a_j\Delta_jB$. For $i<j$ the factor $a_i\Delta_iB\,a_j$ is $\mathcal F_{t_j}$-measurable
and $\Delta_jB$ is independent of it with mean $0$, so the cross terms have mean $0$. For the diagonal terms $a_i$ and $\Delta_iB$ are independent, so
$\E[a_i^2(\Delta_iB)^2]=\E[a_i^2]\Delta t_i$. The conditional version is the same computation given $\mathcal F_s$.
\end{proof}

The isometry is the identity that carries everything that follows: \emph{the variance of a stochastic integral is the integrated
mean square of the integrand}. It is the quadratic variation $(\dd B)^2=\dd t$ at work: the cross terms cancel because increments are
independent, and each squared increment is replaced by its mean $\Delta t_i$.

\begin{example}[The step approximation of $\int B\,\dd B$]\label{ch17:ex-stepBdB}
Take the grid $t_i=iT/n$ and $a_i=B_{t_i}$ (adapted, with $\E[a_i^2]=t_i$). Then $I=\sum_{i=0}^{n-1}B_{t_i}(B_{t_{i+1}}-B_{t_i})$ has mean $0$ and, by the isometry,
\[
  \Var I=\sum_i\E[B_{t_i}^2]\,\Delta t_i=\sum_{i=0}^{n-1}\frac{iT}{n}\cdot\frac Tn=\frac{T^2(n-1)}{2n}\ \xrightarrow[n\to\infty]{}\ \frac{T^2}2=\int_0^Ts\,\dd s
\]
\ts{24}{24}{33:58}. The limit $n\to\infty$ is the lecture's fourth case and is the subject of \cref{ch17:sec-BdB}.
\end{example}

% ==================================================================
\section{The It\^o integral in general}\label{ch17:sec-general}

\subsection{Construction (not given in class)}\label{ch17:sec-constr}
The lecture states the definition and refers to Steele's book, chapter~6, for ``some very strong mathematical arguments'', adding that it is ``not
really a fun read'' \ts{24}{24}{43:35}. The idea is short enough to record, because it is what the isometry is for.
Let $\mathcal H^2$ be the set of processes $f(\omega,s)$, $0\le s\le T$, that are adapted ($f(\cdot,s)$ is $\mathcal F_s$-measurable
for every $s$) and satisfy
\begin{equation}\label{ch17:eq-h2}
  \|f\|^2:=\E\int_0^Tf(\omega,s)^2\,\dd s=\int_\Omega\Bigl[\int_0^Tf(\omega,s)^2\,\dd s\Bigr]\dd\Prob(\omega)<\infty
\end{equation}
\ts{24}{24}{45:40}. This is the squared norm of $L^2(\dd\Prob\times\dd t)$, the space of integrands. The steps are:
\begin{enumerate}[1.]
  \item \emph{Approximation.} Every $f\in\mathcal H^2$ is the limit, in this norm, of simple processes $f_n$ (for a continuous adapted bounded $f$ take
        $f_n=\sum f(t_i)\ind(t_i<s\le t_{i+1})$ on finer and finer grids and use bounded convergence; the general case is a density argument).
  \item \emph{Cauchy property.} By linearity and \cref{ch17:prop-simple}(ii), $\E[(I(f_n)-I(f_m))^2]=\|f_n-f_m\|^2\to0$: the integrals $I(f_n)$ form a Cauchy
        sequence in $L^2(\dd\Prob)$.
  \item \emph{Definition.} $I(f):=\lim_nI(f_n)$ in $L^2(\dd\Prob)$; the limit does not depend on the approximating sequence
        (interlace two sequences and use step~2).
\end{enumerate}
This is why the isometry $\|I(f)\|_{L^2(\dd\Prob)}=\|f\|_{L^2(\dd\Prob\times\dd t)}$ (the slide writes the two \emph{squared} norms, and
$\Var I(f)=\|I(f)\|^2$) is highlighted in the lecture: it is an isometry between the space of integrands and the space of integrals, it ``formalises the convergence of
sequences of It\^o integrals'' and lets one speak about closed subspaces and limits \ts{24}{24}{51:02}\ts{13}{18}{49:10}.
The integral up to $t<T$ is the integral of the truncated integrand $f\ind(0\le s\le t)$ \ts{24}{24}{46:46}.

\begin{theorem}[Properties of the It\^o integral]\label{ch17:thm-ito-int}
Let $f,g\in\mathcal H^2$ and $X_t=\int_0^tf(\omega,s)\,\dd B_s$.
\begin{enumerate}[(i)]
  \item \emph{Linearity:} $\int(af+bg)\,\dd B=a\int f\,\dd B+b\int g\,\dd B$.
  \item \emph{Martingale:} $\{X_t\}$ is a martingale with respect to $\{\mathcal F_t\}$ (it has a continuous version), $\E[X_t]=0$ and
        $\E[X_{t+k}\mid\mathcal F_t]=X_t$ \ts{24}{24}{47:51}\ts{13}{18}{52:22}.
  \item \emph{Isometry:} $\E[X_t^2]=\E\int_0^tf^2\,\dd s$ and
        $\Var[X_{t+k}\mid\mathcal F_t]=\E\bigl[\int_t^{t+k}f^2\,\dd s\bigm|\mathcal F_t\bigr]$ \ts{24}{24}{48:57}.
  \item \emph{Deterministic integrands are Gaussian:} if $f=f(s)$ is a deterministic (continuous) function, $\{X_t\}$ is a mean-zero \emph{Gaussian}
        process with independent increments and covariance
        \[
          \Cov(X_s,X_t)=\int_0^{s\wedge t}f(u)^2\,\dd u
        \]
        \ts{24}{24}{56:28}\ts{13}{18}{46:59}.
\end{enumerate}
\end{theorem}
\begin{proof}
(i)--(iii) hold for simple processes (\cref{ch17:prop-simple}), and both sides are continuous for $L^2$-limits: $L^2$-limits preserve expectations,
second moments and conditional expectations ($\E[\,\cdot\mid\mathcal F_s]$ is a contraction on $L^2$). (iv) For a step function
$X_t=\sum a_i(B_{t\wedge t_{i+1}}-B_{t\wedge t_i})$ is a linear combination of independent Gaussian increments, hence Gaussian, and
$L^2$-limits of Gaussian variables are Gaussian (characteristic functions converge). The covariance follows from the isometry,
$\Cov(X_s,X_t)=\E[X_s(X_s+(X_t-X_s))]=\E[X_s^2]$ for $s\le t$ because the increment $X_t-X_s$ is orthogonal to $X_s$.
\end{proof}

The sum $I(f)=\sum a_i\Delta_iB$ of a \emph{deterministic} integrand is a sum of independent normals, and the theorem says that this survives in the limit: the
2013 lecture presents it as ``the sum of normal variables is normal'' \ts{13}{18}{46:59}. For a \emph{random} adapted integrand
$X_t$ is a martingale with a known variance but in general \emph{not} Gaussian; the standard example is
$\int_0^tB\,\dd B=\tfrac12(B_t^2-t)$ (\cref{ch17:sec-BdB}), which is a shifted $\chi^2$ variable, skewed to the right.

\subsection{The Riemann-sum representation and the left point}\label{ch17:sec-riemann}
\begin{theorem}[Riemann representation]\label{ch17:thm-riemann}
For a continuous $f:\R\to\R$ and the partitions $t_i=iT/n$ of $[0,T]$,
\[
  \lim_{n\to\infty}\sum_{i=1}^{n}f\bigl(B_{t_{i-1}}\bigr)\bigl(B_{t_i}-B_{t_{i-1}}\bigr)=\int_0^Tf(B_s)\,\dd B_s
\]
in probability \ts{24}{24}{55:23}\ts{13}{18}{33:47}.
\end{theorem}
\begin{proof}[Sketch]
The sum is $I(f_n)$ for the simple process $f_n(s)=f(B_{t_{i-1}})$ on $(t_{i-1},t_i]$. If $f$ is bounded, $f_n(s)\to f(B_s)$ for every $s$ by continuity of the path,
so $\|f_n-f(B)\|^2\to0$ by bounded convergence and $I(f_n)\to I(f(B))$ in $L^2$ by the isometry. For unbounded $f$ one truncates $f$ where the path exceeds a level and lets the level grow:
the paths are bounded on $[0,T]$, so the truncation is inactive with probability close to one (hence convergence in probability only).
\end{proof}
In ordinary Riemann integration it does not matter which point of each small interval is used; \emph{here it does}. Evaluating at
the right end, or at the middle, gives different limits, because $\sum(\Delta B)^2$ does not vanish
\ts{13}{18}{35:48}; the explicit case is in \cref{ch17:sec-BdB}. The lecture also mentions the general statement that any intermediate point
$t_i^*\in[t_{i-1},t_i]$ gives the same limit for a deterministic continuous $f$ (\emph{Gaussian integrals}, slide 12, after Steele) \ts{24}{24}{57:30};
this is right, because for deterministic $f$ the difference between two choices is a sum of independent terms of variance $\sum(f(t_i^*)-f(t_{i-1}))^2\Delta t_i\to0$.

\begin{finance}[The left point is the no-lookahead rule]
The It\^o rule is the mathematical form of ``you cannot see the future'': the position over $(t_i,t_{i+1}]$ is decided at $t_i$. The
2013 lecture says that this intuition is \emph{built into the theory}, which is why it fits finance so well; the same sum with the position taken from the right end
would use the price move itself to decide how much to hold \ts{13}{18}{37:01}. The
martingale property of \cref{ch17:thm-ito-int}(ii) is the statement that the gains of an adapted strategy in a fair game (a martingale price)
are themselves a fair game, the continuous-time version of the martingale transform of \cref{ch04:thm-transform}.
\end{finance}

\begin{remark}[Integrability conditions; not discussed in class]\label{ch17:rem-local}
The last slide of L24 defines a \emph{standard process} $X_t=x_0+\int_0^ta(\omega,s)\,\dd s+\int_0^tb(\omega,s)\,\dd B_s$ with $a,b$ adapted and
$\Prob(\int_0^t|a|\,\dd s<\infty)=\Prob(\int_0^tb^2\,\dd s<\infty)=1$, ``otherwise the integrals may explode'' \ts{24}{24}{1:21:11}. The weaker condition
$\int_0^Tb^2\,\dd s<\infty$ almost surely (instead of finite expectation) still defines $\int b\,\dd B$ (by stopping when $\int b^2$ exceeds a level), but the result is then only a
\emph{local} martingale: mean zero and the isometry may fail. Whenever we assert a martingale property below we check $\E\int f^2<\infty$.
\end{remark}

% ==================================================================
\section{The basic computation: $\int B\,\dd B$}\label{ch17:sec-BdB}

Classical calculus would give $\int_0^TB\,\dd B=\tfrac12B_T^2$. The lecture's fourth case computes the limit of the step approximations of
\cref{ch17:ex-stepBdB} and finds something else \ts{24}{24}{37:10}.

\begin{theorem}[Ito integral of Brownian motion]\label{ch17:thm-BdB}
For every $T>0$,
\begin{equation}\label{ch17:eq-BdB}
  \int_0^TB_s\,\dd B_s=\tfrac12B_T^2-\tfrac12T\qquad\text{almost surely.}
\end{equation}
\end{theorem}
\begin{proof}
For a partition $0=t_0<\dots<t_n=T$ use the algebraic identity, ``trivial but incredibly useful'' in the lecture,
\[
  B_{t_i}\bigl(B_{t_{i+1}}-B_{t_i}\bigr)=\tfrac12\bigl(B_{t_{i+1}}^2-B_{t_i}^2\bigr)-\tfrac12\bigl(B_{t_{i+1}}-B_{t_i}\bigr)^2 ,
\]
which is $b(c-b)=\tfrac12(c^2-b^2)-\tfrac12(c-b)^2$. Summing over $i$ the first term telescopes:
\[
  \sum_{i=0}^{n-1}B_{t_i}\bigl(B_{t_{i+1}}-B_{t_i}\bigr)=\tfrac12B_T^2-\tfrac12\sum_{i=0}^{n-1}\bigl(B_{t_{i+1}}-B_{t_i}\bigr)^2
\]
\ts{24}{24}{40:21}. By \cref{ch11:thm-qv} the last sum, the quadratic variation of the partition, tends to $T$ as the mesh goes to zero
(in $L^2$, hence in probability; almost surely for dyadic partitions), and by \cref{ch17:thm-riemann} the left side tends to $\int_0^TB\,\dd B$ \ts{24}{24}{42:33}.
\end{proof}
So the correction $-\tfrac12T$ ``drags'' on the integral because of the quadratic variation \ts{24}{24}{42:33}; \eqref{ch17:eq-BdB}
also appears in 2013 as the example of a `violation' of the fundamental theorem of calculus \ts{13}{18}{18:49}.
Checks that do not need the proof: the mean is $\E\tfrac12(B_T^2-T)=0$, as required for a martingale (\cref{ch17:thm-ito-int}(ii));
the variance is $\tfrac14\Var(B_T^2)=\tfrac14\cdot2T^2=\tfrac12T^2$, which is the isometry value $\int_0^Ts\,\dd s$ of \cref{ch17:ex-stepBdB} \ts{24}{24}{35:04};
the random variable $\tfrac12(B_T^2-T)$ is a shifted, rescaled $\chi^2_1$, hence \emph{not} Gaussian (skewness $2\sqrt2\approx2.83$). The 2013 note verifies the martingale property directly by
computing $\E[\int_{t_1}^{t_2}B\,\dd B\mid\mathcal F_{t_1}]=0$ (\file{lecnote18}, Example~2.4).

\begin{corollary}[Used in the Dickey--Fuller distribution]\label{ch17:cor-df}
For a standard Brownian motion $W$, $\int_0^1W\,\dd W=\tfrac12\bigl(W(1)^2-1\bigr)$. This is the numerator of the limiting distribution of the unit-root estimator
$T(\hat\phi-1)$ in \cref{ch09:sec-df} (there obtained from the discrete telescoping identity $X_t^2=X_{t-1}^2+2X_{t-1}\eta_t+\eta_t^2$, the finite-$n$ version of the proof above).
\end{corollary}

\begin{figure}[ht]
\centering
\begin{tikzpicture}
\begin{axis}[width=0.95\linewidth,height=5.2cm,xlabel={$s$},xmin=0,xmax=1,grid=major,grid style={black!10},
  legend style={at={(0.5,1.03)},anchor=south,legend columns=2,font=\footnotesize}]
  \addplot[thick,color=defcol] coordinates {(0.0000,0.0000) (0.0078,-0.1232) (0.0156,-0.2517) (0.0234,-0.2291) (0.0312,-0.1628) (0.0391,-0.0687) (0.0469,-0.0301) (0.0547,-0.2288) (0.0625,-0.2774) (0.0703,-0.3722) (0.0781,-0.4260) (0.0859,-0.6291) (0.0938,-0.3400) (0.1016,-0.2979) (0.1094,-0.3106) (0.1172,-0.3417) (0.1250,-0.4502) (0.1328,-0.4472) (0.1406,-0.3885) (0.1484,-0.4632) (0.1562,-0.4399) (0.1641,-0.4302) (0.1719,-0.4800) (0.1797,-0.4677) (0.1875,-0.3809) (0.1953,-0.3197) (0.2031,-0.1277) (0.2109,-0.1700) (0.2188,-0.1815) (0.2266,-0.0575) (0.2344,0.2062) (0.2422,0.1416) (0.2500,0.2716) (0.2578,0.0598) (0.2656,0.1384) (0.2734,0.2804) (0.2812,0.1440) (0.2891,0.1783) (0.2969,0.2344) (0.3047,0.3237) (0.3125,0.3700) (0.3203,0.3875) (0.3281,0.5103) (0.3359,0.6271) (0.3438,0.6298) (0.3516,0.5137) (0.3594,0.6232) (0.3672,0.4292) (0.3750,0.3974) (0.3828,0.3802) (0.3906,0.5184) (0.3984,0.5302) (0.4062,0.5513) (0.4141,0.3311) (0.4219,0.3885) (0.4297,0.3675) (0.4375,0.4374) (0.4453,0.3439) (0.4531,0.2546) (0.4609,0.1643) (0.4688,0.0984) (0.4766,0.0843) (0.4844,0.0752) (0.4922,-0.0017) (0.5000,0.0494) (0.5078,0.0549) (0.5156,0.2017) (0.5234,0.2235) (0.5312,0.2691) (0.5391,0.2395) (0.5469,0.2686) (0.5547,0.2059) (0.5625,0.0829) (0.5703,0.0933) (0.5781,0.0661) (0.5859,-0.0341) (0.5938,-0.0859) (0.6016,-0.2074) (0.6094,-0.1317) (0.6172,-0.0400) (0.6250,0.1375) (0.6328,0.1200) (0.6406,0.0297) (0.6484,0.0718) (0.6562,0.0739) (0.6641,-0.0627) (0.6719,-0.0381) (0.6797,-0.1069) (0.6875,-0.0963) (0.6953,-0.0038) (0.7031,0.1150) (0.7109,0.1828) (0.7188,0.1325) (0.7266,0.1869) (0.7344,0.2413) (0.7422,0.2290) (0.7500,0.3717) (0.7578,0.4868) (0.7656,0.4584) (0.7734,0.3444) (0.7812,0.4190) (0.7891,0.4836) (0.7969,0.5691) (0.8047,0.4646) (0.8125,0.5343) (0.8203,0.6967) (0.8281,0.5767) (0.8359,0.4207) (0.8438,0.4948) (0.8516,0.5692) (0.8594,0.5157) (0.8672,0.6352) (0.8750,0.6546) (0.8828,0.6514) (0.8906,0.7721) (0.8984,0.9235) (0.9062,0.8794) (0.9141,0.8084) (0.9219,0.8341) (0.9297,0.8367) (0.9375,0.8805) (0.9453,0.9166) (0.9531,0.8504) (0.9609,0.8484) (0.9688,0.8561) (0.9766,0.6402) (0.9844,0.6973) (0.9922,0.8059) (1.0000,1.0320)};
  \addplot[thick,color=thmcol] coordinates {(0.0000,0.0000) (0.1250,0.0000) (0.1250,-0.4502) (0.2500,-0.4502) (0.2500,0.2716) (0.3750,0.2716) (0.3750,0.3974) (0.5000,0.3974) (0.5000,0.0494) (0.6250,0.0494) (0.6250,0.1375) (0.7500,0.1375) (0.7500,0.3717) (0.8750,0.3717) (0.8750,0.6546) (1.0000,0.6546)};
  \legend{path $B_s$,step integrand $f_8(s)$}
\end{axis}
\end{tikzpicture}
\caption{The step approximation of \cref{ch17:ex-stepBdB}: a simulated path of $B$ on $[0,1]$ (blue) and the adapted step integrand $f_8(s)=B_{t_i}$ on $(t_i,t_{i+1}]$ (red, $n=8$). Each step is
frozen at the value of the path at its \emph{left} end.}\label{ch17:fig-steps}
\end{figure}

\subsection{Left, middle and right points: It\^o, Stratonovich, backward}\label{ch17:sec-strat}
Evaluate the integrand at another point of each interval and the limit changes; for $\int B\,\dd B$ this can be done exactly.
With $\Delta_iB=B_{t_{i+1}}-B_{t_i}$:
\begin{align*}
  \text{left point (It\^o):}\quad&\sum_iB_{t_i}\Delta_iB=\tfrac12B_T^2-\tfrac12\sum_i(\Delta_iB)^2\ \to\ \tfrac12B_T^2-\tfrac12T,\\
  \text{mid-point (Stratonovich):}\quad&\sum_i\tfrac12\bigl(B_{t_i}+B_{t_{i+1}}\bigr)\Delta_iB=\tfrac12\sum_i\bigl(B_{t_{i+1}}^2-B_{t_i}^2\bigr)=\tfrac12B_T^2\quad\text{exactly},\\
  \text{right point (backward):}\quad&\sum_iB_{t_{i+1}}\Delta_iB=\tfrac12B_T^2+\tfrac12\sum_i(\Delta_iB)^2\ \to\ \tfrac12B_T^2+\tfrac12T.
\end{align*}
The three limits differ by $\pm\tfrac12T$, the quadratic variation, \cref{ch17:fig-sums}. The mid-point rule, the \emph{Stratonovich} integral $\int B\circ\dd B=\tfrac12B_T^2$,
obeys the ordinary rules of calculus (chain rule without the second-order term); it is the limit of physically smooth noise, which is why it is the
convention of physicists who start from a Langevin equation with a finite correlation time. The price is that it is not
a martingale and the integrand looks into the future: the integrand is evaluated at a point that is not yet known. The right-point sums are
the ``equivalent version'' of the 2013 remark, in which
one gets $\dd f=f'(B)\dd B-\tfrac12f''(B)\dd t$, i.e.\ formally $(\dd B)^2=-\dd t$ \ts{13}{18}{37:01}; this is
correct once the integral is understood as the backward one (from the identity above, $\int B\,\dd^{\rm back}B=\tfrac12B^2+\tfrac12t$, so $\dd(\tfrac12B^2)=B\,\dd^{\rm back}B-\tfrac12\dd t$).
The lecture calls it ``just cool stuff'' and moves on; It\^o's is the convention for finance.

\begin{figure}[ht]
\centering
\begin{tikzpicture}
\begin{axis}[width=0.95\linewidth,height=6cm,xmode=log,log basis x=2,xlabel={number of subintervals $n$},xmin=1.7,xmax=5000,ymin=-0.25,ymax=1.25,
  ytick={0.0325,0.5325,1.0325},yticklabels={{$\tfrac12(B_1^2-1)$},{$\tfrac12B_1^2$},{$\tfrac12(B_1^2+1)$}},yticklabel style={font=\footnotesize},
  xtick={2,8,32,128,512,2048},xticklabels={2,8,32,128,512,2048},
  legend style={at={(0.5,1.03)},anchor=south,legend columns=3,font=\footnotesize},grid=major,grid style={black!10}]
  \addplot[thick,color=defcol,mark=*,mark size=1.3pt] coordinates {(2,0.0485) (4,0.2010) (8,-0.0403) (16,-0.0138) (32,-0.0054) (64,-0.0497) (128,-0.0924) (256,-0.0335) (512,-0.0036) (1024,0.0145) (2048,0.0341) (4096,0.0262)};
  \addplot[thick,color=intcol,mark=square*,mark size=1.3pt] coordinates {(2,0.5325) (4,0.5325) (8,0.5325) (16,0.5325) (32,0.5325) (64,0.5325) (128,0.5325) (256,0.5325) (512,0.5325) (1024,0.5325) (2048,0.5325) (4096,0.5325)};
  \addplot[thick,color=thmcol,mark=triangle*,mark size=1.6pt] coordinates {(2,1.0164) (4,0.8640) (8,1.1053) (16,1.0788) (32,1.0704) (64,1.1147) (128,1.1573) (256,1.0985) (512,1.0685) (1024,1.0505) (2048,1.0308) (4096,1.0387)};
  \addplot[dashed,black!60,forget plot] coordinates {(1.7,0.0325) (5000,0.0325)};
  \addplot[dashed,black!60,forget plot] coordinates {(1.7,0.5325) (5000,0.5325)};
  \addplot[dashed,black!60,forget plot] coordinates {(1.7,1.0325) (5000,1.0325)};
  \legend{left point (It\^o),mid-point (Stratonovich),right point}
\end{axis}
\end{tikzpicture}
\caption{One simulated Brownian path on $[0,1]$ with $B_1=1.032$, discretised at $n$ equally spaced points ($n=2,\dots,4096$, sub-sampled from a single
fine path). The left-point sums converge (slowly, at rate $n^{-1/2}$) to the It\^o value $\tfrac12(B_1^2-1)=0.032$, the mid-point sums are \emph{exactly} the Stratonovich value $\tfrac12B_1^2=0.532$, and the
right-point sums tend to $\tfrac12(B_1^2+1)=1.032$; the three limits are the tick labels of the vertical axis.}\label{ch17:fig-sums}
\end{figure}

\begin{exercise}[PS8 (2013), Problem B-2: It\^o isometry]\label{ch17:ex-b2}
Use the It\^o isometry to calculate $\E[X_t^2]$, where $X_t=\int_0^tB_s\,\dd B_s$.
\end{exercise}

\begin{exercise}[not from the course: Stratonovich]\label{ch17:ex-x1}
Show that the mid-point sums $\sum_i\bigl[\tfrac12(B_{t_i}+B_{t_{i+1}})\bigr]^{2}\Delta_iB$ converge to $\tfrac13B_T^3$ as the mesh tends to zero, and compare with the left-point (It\^o) limit $\int_0^TB^2\,\dd B=\tfrac13B_T^3-\int_0^TB_s\,\dd s$ from \eqref{ch17:eq-antider}. (Hint: compare each term with $\tfrac13(B_{t_{i+1}}^3-B_{t_i}^3)$.)
\end{exercise}

% ==================================================================
\section{It\^o's lemma}\label{ch17:sec-lemma}

\subsection{One variable}\label{ch17:sec-lemma1}
\begin{theorem}[It\^o's formula for $f(B_t)$]\label{ch17:thm-ito1}
If $f:\R\to\R$ has a continuous second derivative, then for every $t\ge0$
\begin{equation}\label{ch17:eq-ito1}
  f(B_t)=f(0)+\int_0^tf'(B_s)\,\dd B_s+\frac12\int_0^tf''(B_s)\,\dd s,
\end{equation}
written in shorthand $\dd f(B_t)=f'(B_t)\,\dd B_t+\tfrac12f''(B_t)\,\dd t$.
\end{theorem}
\ts{24}{24}{57:30}\ts{24}{25}{0:00}\ts{13}{18}{0:00}\ \noindent
The first integral is a stochastic integral (mean zero, a martingale when $\E\int f'(B)^2<\infty$), the second an ordinary Riemann integral: a
random \emph{adapted drift} \ts{24}{24}{58:36}. It is the chain rule with one extra term, the \emph{It\^o correction} $\tfrac12f''\,\dd t$.
(Both years also record the classical formula, obtained if $\dd B/\dd t$ existed, $\dd f=f'(B)\tfrac{\dd B}{\dd t}\dd t$, which ``makes no sense'', and the
attempt $\dd f=f'\dd B$ that ``at least makes sense'' but ``does not quite work'' \file{lecnote18}, \S1.)

\begin{proof}[Proof (sketch with the error control; the lecture only Taylor-expands)]
Take a partition $t_i=it/n$ and telescope, $f(B_t)-f(0)=\sum_i[f(B_{t_{i+1}})-f(B_{t_i})]$ \ts{24}{24}{59:45}. Taylor's theorem with the Lagrange remainder gives, with $\Delta_iB=B_{t_{i+1}}-B_{t_i}$
and some $\xi_i$ between $B_{t_i}$ and $B_{t_{i+1}}$,
\[
  f(B_{t_{i+1}})-f(B_{t_i})=f'(B_{t_i})\Delta_iB+\tfrac12f''(B_{t_i})(\Delta_iB)^2+\tfrac12\bigl[f''(\xi_i)-f''(B_{t_i})\bigr](\Delta_iB)^2 .
\]
Sum over $i$ and treat the three sums.
(i) $\sum f'(B_{t_i})\Delta_iB\to\int_0^tf'(B)\,\dd B$ by \cref{ch17:thm-riemann}.

(ii) Write $(\Delta_iB)^2=\Delta t_i+[(\Delta_iB)^2-\Delta t_i]$. The sum $\tfrac12\sum f''(B_{t_i})\Delta t_i$ is a Riemann sum of the
continuous function $s\mapsto f''(B_s)$ and tends to $\tfrac12\int_0^tf''(B_s)\,\dd s$. The fluctuation $D_n=\tfrac12\sum f''(B_{t_i})[(\Delta_iB)^2-\Delta t_i]$ is a sum of martingale differences ($\Delta_iB$ is independent of $\mathcal F_{t_i}$ and
$\E[(\Delta_iB)^2-\Delta t_i]=0$), so for bounded $f''$
\begin{align*}
\E[D_n^2]&=\tfrac14\sum_i\E\bigl[f''(B_{t_i})^2\bigr]\Var\bigl[(\Delta_iB)^2\bigr]=\tfrac12\sum_i\E\bigl[f''(B_{t_i})^2\bigr]\Delta t_i^2\\
&\le\tfrac12\|f''\|_\infty^2\,t\,\max_i\Delta t_i\to0 ,
\end{align*}
using $\Var[(\Delta_iB)^2]=2\Delta t_i^2$. This is the point where ``$(\dd B)^2=\dd t$'' enters: \emph{the squared increment can be replaced by its mean because the difference has variance of order $\Delta t^2$, whose sum
is negligible}.

(iii) The remainder is bounded by $\tfrac12\varpi(\max_i|\Delta_iB|)\sum(\Delta_iB)^2$, where $\varpi$ is the modulus of continuity of $f''$ on the (compact) range of the path; the sum is bounded
in probability (it converges to $t$) and $\max_i|\Delta_iB|\to0$ by uniform continuity of the path, so the remainder tends to $0$ in probability. For unbounded $f''$ one stops the path when it leaves a large interval.
Collecting the three limits gives \eqref{ch17:eq-ito1}.
\end{proof}

\begin{intuition}[Where the extra term comes from]
Expand $f(B+\Delta B)-f(B)=f'\Delta B+\tfrac12f''\Delta B^2+\tfrac16f'''\Delta B^3+\dots$. Since $\Delta B\sim\sqrt{\Delta t}$, the terms are
of order $\Delta t^{1/2},\Delta t,\Delta t^{3/2},\dots$: after summing $1/\Delta t$ steps, the first one is a random martingale, the second contributes a finite drift (its fluctuations are
negligible, as shown above), and all others vanish. This is the 2013 argument: ``$\dd t^2$ and $\dd t\,\dd B$ are insignificant, the only term that matters is $\dd B^2=\dd t$'', with the
heuristic $\dd B\approx\sqrt{\dd t}$ \ts{13}{18}{4:34}. A physicist recognises it: taking expectations kills the martingale,
$\dd\E[f(B_t)]/\dd t=\tfrac12\E[f''(B_t)]$, i.e.\ $\tfrac12\partial_x^2$ is the \emph{generator} of Brownian motion and the density of $B_t$ solves the diffusion equation
of \cref{ch11:prop-heat}. It\^o's formula is the pathwise version of the heat equation.
\end{intuition}

\subsection{Functions of time and space; the multiplication table}\label{ch17:sec-lemma2}
For $f\in C^{1,2}(\R_+\times\R)$ (continuous derivatives $f_t,f_x,f_{xx}$) the Taylor expansion has the second-order terms
$\tfrac12f_{tt}\Delta t^2+f_{tx}\Delta t\Delta B+\tfrac12f_{xx}\Delta B^2$; only the last one survives
\ts{24}{24}{1:06:15}\ts{24}{25}{1:03}\ts{13}{18}{2:17}. We summarise this in the
\emph{multiplication table} of It\^o calculus
\begin{equation}\label{ch17:eq-table}
  (\dd B)^2=\dd t,\qquad \dd B\,\dd t=0,\qquad (\dd t)^2=0 ,
\end{equation}
(the same as \cref{ch11:eq-dB2}; the products $\dd B\,\dd t\sim\dd t^{3/2}$ and $\dd t^2$ are negligible against $\dd t$).

\begin{theorem}[It\^o's formula for $f(t,B_t)$]\label{ch17:thm-ito2}
For $f\in C^{1,2}(\R_+\times\R)$,
\begin{equation}\label{ch17:eq-ito2}
  f(t,B_t)-f(0,0)=\int_0^t\Bigl[f_t+\tfrac12f_{xx}\Bigr](s,B_s)\,\dd s+\int_0^tf_x(s,B_s)\,\dd B_s ,
\end{equation}
i.e.\ $\dd f=\bigl(f_t+\tfrac12f_{xx}\bigr)\dd t+f_x\,\dd B$.
\end{theorem}
\ts{24}{24}{1:08:20}\ts{13}{18}{7:37}\ \noindent
The proof is that of \cref{ch17:thm-ito1} with the Taylor expansion in two variables (the additional terms $f_t\Delta t$ sum to a Riemann integral). The 2013 lecture notes that the second
variable is just a placeholder: one writes $f_x$ for the derivative in the second argument and then \emph{substitutes} $x=B_t$ \ts{13}{18}{9:48}.

\begin{theorem}[It\^o's formula for standard processes]\label{ch17:thm-ito3}
Let $\dd X_t=a_t\,\dd t+b_t\,\dd B_t$, with $a,b$ adapted and $\int_0^t|a|\,\dd s<\infty$, $\int_0^tb^2\,\dd s<\infty$ almost surely (a \emph{standard process}, \cref{ch17:rem-local}). For $f\in C^{1,2}$,
\begin{equation}\label{ch17:eq-ito3}
  \dd f(t,X_t)=f_t\,\dd t+f_x\,\dd X_t+\tfrac12f_{xx}\,(\dd X_t)^2=\Bigl(f_t+a_tf_x+\tfrac12b_t^2f_{xx}\Bigr)\dd t+b_tf_x\,\dd B_t ,
\end{equation}
where $(\dd X_t)^2=b_t^2\,\dd t$ by \eqref{ch17:eq-table}.
\end{theorem}
\ts{24}{24}{1:21:11}\ts{24}{25}{0:00}\ts{13}{18}{11:53}
\begin{proof}[Idea]
As above, $\dd f=f_t\dd t+f_x\dd X+\tfrac12f_{xx}(\dd X)^2$, and substituting $\dd X=a\,\dd t+b\,\dd B$,
$(\dd X)^2=a^2\dd t^2+2ab\,\dd t\,\dd B+b^2\dd B^2=b^2\dd t$: only $b^2\dd B^2$ survives, the extra term is $\tfrac12b^2f_{xx}\dd t$. The 2013 lecture stresses that ``it all started from one fact, the quadratic variation'' and
that the formula looks complicated only because one forgets where the terms come from \ts{13}{18}{14:58}.
The note's proof (Theorem~1.1) has a stray ``$\dots\dd t\,\dd B_T$'' for the discarded terms, which are $ab\,f_{xx}\,\dd t\,\dd B$ and $\tfrac12a^2f_{xx}(\dd t)^2$.
The rigorous version needs the same estimates as in \cref{ch17:thm-ito1} with $f''(B_{t_i})$ replaced by $f_{xx}(t_i,X_{t_i})b_{t_i}^2$.
\end{proof}

The lecture writes the last result for a ``Markov'' process $\dd X=\mu(X_t,t)\dd t+\sigma(X_t,t)\dd B_t$ and $Y_t=G(X_t,t)$ as
\[
  \dd G=\Bigl[\mu\,G_x+G_t+\tfrac12\sigma^2G_{xx}\Bigr]\dd t+\sigma\,G_x\,\dd B_t
\]
\ts{24}{24}{1:22:14}\ts{24}{25}{5:16}; this is \eqref{ch17:eq-ito3} with $a=\mu(X,t)$, $b=\sigma(X,t)$.

\subsection{Integrals by anti-derivatives}\label{ch17:sec-antider}
Solving \eqref{ch17:eq-ito1} for the stochastic integral gives the lecture's method to compute $\int f(B_s)\,\dd B_s$: if $F'=f$ (choose $F(0)=0$), then
\begin{equation}\label{ch17:eq-antider}
  \int_0^tf(B_s)\,\dd B_s=F(B_t)-\tfrac12\int_0^tf'(B_s)\,\dd s
\end{equation}
\ts{24}{24}{1:01:58}\ts{24}{25}{1:03}. For $F(x)=x$ ($f=1$, $f'=0$) it returns $B_t$; for $F(x)=\tfrac12x^2$ it returns $\tfrac12B_t^2-\tfrac12t$, the result of \cref{ch17:sec-BdB}
\ts{24}{24}{1:05:13}. Another example (ours): for $F=f=e^x-1$,
\[
  \int_0^te^{B_s}\,\dd B_s=e^{B_t}-1-\tfrac12\int_0^te^{B_s}\,\dd s ,
\]
an exact identity between a martingale and a functional of the path, checked pathwise in a simulation below.

\begin{coursenote}[Erratum: the It\^o integral is not a Gaussian process]
Slide~13 of \file{lec19\_1} (and the spoken commentary \ts{24}{24}{58:36}) describes $\int_0^tf'(B_s)\,\dd B_s$ in \eqref{ch17:eq-ito1} as a ``zero-mean adaptive \emph{Gaussian} process''. Only the zero mean
(indeed the martingale property) is correct; for a random integrand the integral is in general not Gaussian, as the
example $\int_0^tB\,\dd B=\tfrac12(B_t^2-t)$ shows. Gaussianity holds only for \emph{deterministic} integrands (\cref{ch17:thm-ito-int}(iv)).
\end{coursenote}

\begin{coursenote}[Erratum: slide 20 and the upper limit]
Slide~20 of \file{lec19\_1} defines the standard process by $X_t=x_0+\int_0^ta\,\dd s+\int_0^Tb\,\dd B_s$, with the upper limit $T$ in the stochastic integral; it should be $t$, as in the first integral and as in the definition on slide~9.
\end{coursenote}

\begin{remark}[Stratonovich form; not discussed in class]\label{ch17:rem-strat}
If one wants a chain rule without the second-order term, use the mid-point integral. The two are related by
$\int\sigma(X)\circ\dd B=\int\sigma(X)\,\dd B+\tfrac12\int\sigma'(X)\sigma(X)\,\dd t$ (for $X$ a solution of the It\^o equation with volatility $\sigma$): for $\sigma(x)=x$, $X=B$, this is $\int B\circ\dd B=\int B\,\dd B+\tfrac12t$,
which is the pair of formulas of \cref{ch17:sec-strat}. Consequently a Langevin equation with multiplicative noise has a \emph{noise-induced drift} $\tfrac12\sigma\sigma'$ if read in It\^o's convention and none in Stratonovich's;
in finance the It\^o reading is used because prices must be martingales when there is no drift.
\end{remark}

\begin{exercise}[PS8 (2013), Problem A-4: It\^o's formula]\label{ch17:ex-a4}
Use It\^o's formula to compute the differentials of the following ($B_t$ is a Brownian motion):
\begin{enumerate}[(a)]
  \item $f(t,B_t)=B_t^3$;
  \item $\sin B_t$;
  \item $\cos(t^3+B_t^2)$;
  \item $e^{B_t^2}$;
  \item $f(t,B_t)=\int B_t^2\,\dd B_t$;
  \item $f(t,B_t)=\int B_t\,\dd t$;
  \item $f(t,X_t)=X_t^2$, where $\dd X_t=\mu\,\dd t+\sigma\,\dd B_t$ with $\mu,\sigma$ constants. (Parts (e) and (f) are printed as indefinite integrals; interpret them as the processes $\int_0^tB_s^2\,\dd B_s$ and $\int_0^tB_s\,\dd s$.)
\end{enumerate}
\end{exercise}

\begin{exercise}[PS8 (2013), Problem B-3: integration by parts]\label{ch17:ex-b3}
\leavevmode
\begin{enumerate}[(a)]
  \item Prove the integration by parts formula for the It\^o integral, $\int_0^th(s)\,\dd B_s=h(t)B_t-\int_0^th'(s)B_s\,\dd s$ (for a deterministic differentiable $h$).
  \item Using (a), prove that for each fixed $T$ the random variable $\int_0^TB_s\,\dd s$ is normally distributed.
\end{enumerate}
\end{exercise}

% ==================================================================
\section{Examples: martingales, moments, exponentials}\label{ch17:sec-examples}

\begin{example}[The square]\label{ch17:ex-square}
$f(x)=x^2$: $\dd(B_t^2)=2B_t\,\dd B_t+\dd t$ \ts{13}{18}{18:49}\ts{24}{25}{1:03}. Hence $B_t^2-t=2\int_0^tB\,\dd B$ is a martingale and $\E[B_t^2]=t$, in
agreement with \cref{ch11:prop-elem}(iii). For $X$ with $\dd X=\mu\,\dd t+\sigma\,\dd B$ the same formula gives $\dd(X^2)=(2\mu X+\sigma^2)\,\dd t+2\sigma X\,\dd B$; the lecture takes $\mu=0$, $\sigma=1$, $G=x^2$ and
concludes $\int_0^tX\,\dd X=\tfrac12(X_t^2-t)$ \ts{24}{25}{1:03}.
\end{example}

\begin{example}[Moments of Brownian motion]\label{ch17:ex-moments}
For $f(x)=x^n$, $\dd(B^n)=nB^{n-1}\dd B+\tfrac12n(n-1)B^{n-2}\dd t$. Taking expectations (the stochastic integral has mean zero because $\E\int B^{2n-2}\dd s<\infty$),
$m_n(t)=\E[B_t^n]$ satisfies $m_n(t)=\tfrac12n(n-1)\int_0^tm_{n-2}(s)\,\dd s$. Starting from $m_2(s)=s$ this gives $m_4(t)=\tfrac12\cdot4\cdot3\int_0^ts\,\dd s=3t^2$ and $m_6(t)=\tfrac12\cdot6\cdot5\int_0^t3s^2\,\dd s=15t^3$. In general $\E[B_t^{2k}]=(2k-1)!!\,t^k$,
odd moments vanish. A Monte Carlo check with $2\times10^4$ paths of $1000$ steps gave $\E[B_1^4]\approx2.99$ (theory $3$).
\end{example}

\begin{example}[Exponentials and geometric Brownian motion without drift]\label{ch17:ex-exp}
Let $f(t,x)=\exp(\mu t+\sigma x)$. Then $f_t=\mu f$, $f_x=\sigma f$, $f_{xx}=\sigma^2f$ and
\begin{equation}\label{ch17:eq-exp}
  \dd f(t,B_t)=\bigl(\mu+\tfrac12\sigma^2\bigr)f\,\dd t+\sigma f\,\dd B_t
\end{equation}
\ts{13}{18}{22:03}. (The student's question at that point, why $\tfrac12$ multiplies $\sigma^2$ and not $\sigma'^2$: in the general formula
\eqref{ch17:eq-ito3} the $\sigma$ is the volatility of the argument process $X$, here $X=B$ has volatility $1$; the $\sigma$ of \eqref{ch17:eq-exp} sits inside the function $f$.)
The model question of the 2013 lecture: $S_t=e^{\sigma B_t}$ is \emph{not} the solution of $\dd S=\sigma S\,\dd B$ (a process with no drift), because by \eqref{ch17:eq-exp} with $\mu=0$ it has drift $\tfrac12\sigma^2S$
\ts{13}{18}{24:18}. The drift is removed by choosing $\mu=-\tfrac12\sigma^2$: the process
\begin{equation}\label{ch17:eq-expmart}
  S_t=\exp\bigl(\sigma B_t-\tfrac12\sigma^2t\bigr)\quad\text{satisfies}\quad\dd S_t=\sigma S_t\,\dd B_t,\qquad S_0=1,
\end{equation}
\ts{13}{18}{26:24}\ts{24}{24}{1:12:36}, and is a martingale with $\E[S_t]=1$ (both years; the lecture of 2024 obtains it as the case $f(t,x)=e^{\sigma x-\sigma^2t/2}$ with $f_t=-\tfrac12\sigma^2f$, $f_x=\sigma f$, $f_{xx}=\sigma^2f$). Compare
\cref{ch04:ex-m4}: this is its continuous-time limit.
In discrete form the statement is the exact identity $S_t-1=\sigma\int_0^tS\,\dd B$; we checked it pathwise for $\sigma=1$ with left-point sums: the root-mean-square difference between
$S_1-1$ and $\sum S_{t_i}\Delta_iB$ over $2000$ paths is $0.295,\ 0.092,\ 0.028,\ 0.009$ for $n=10,10^2,10^3,10^4$ steps, decreasing like $n^{-1/2}$ (the strong order of the Euler scheme, see \cref{ch:sde}).
\end{example}

\begin{example}[Time and a drift: $t^2+X^2$]\label{ch17:ex-tx}
Let $\dd X=\mu\,\dd t+\sigma\,\dd B$ and $f(t,x)=t^2+x^2$. Then $f_t=2t$, $f_x=2x$, $f_{xx}=2$ and
\[
  \dd f(t,X_t)=2t\,\dd t+2X_t\,\dd X_t+(\dd X_t)^2=\bigl(2t+2\mu X_t+\sigma^2\bigr)\dd t+2\sigma X_t\,\dd B_t .
\]
\end{example}
\begin{coursenote}[Erratum in \file{lecnote18}, Example 1.3(iv)]
The note writes $\dd f=2t\,\dd t+2X_t\,\dd X_t+(\dd X_t)^2=2t\,\dd t+2(\mu\,\dd t+\sigma\,\dd B_t)+\sigma^2\dd t=(2t+2\mu+\sigma^2)\dd t+2\sigma\,\dd B_t$: the factor $X_t$ in
$2X_t\,\dd X_t$ has been dropped in the second step. The correct result is the one above.
\end{coursenote}

\subsection{Martingales from the PDE condition}\label{ch17:sec-pde}
Equation \eqref{ch17:eq-ito2} shows at once when $f(t,B_t)$ is a martingale: the $\dd t$-term must vanish.

\begin{theorem}[PDE condition for martingales]\label{ch17:thm-pde}
Let $f\in C^{1,2}$ satisfy the backward heat equation
\begin{equation}\label{ch17:eq-pdecond}
  f_t+\tfrac12f_{xx}=0 ,
\end{equation}
and suppose $\E\int_0^T\bigl(f_x(s,B_s)\bigr)^2\dd s<\infty$. Then $M_t=f(t,B_t)=f(0,0)+\int_0^tf_x\,\dd B$ is a (square-integrable) martingale with $\E[|M_T-M_0|^2]=\E\int_0^Tf_x^2\,\dd s$.
\end{theorem}
\ts{24}{24}{1:10:26}\ts{24}{25}{2:06}\ts{13}{18}{54:35}\ \noindent
This is the connection ``martingales $\leftrightarrow$ PDEs'' of the lecture: a solution of the PDE gives a martingale, and conversely the
martingale condition reads $f_t=-\tfrac12f_{xx}$ \ts{24}{24}{1:11:31}. Both years extract the general principle:
\emph{in an SDE written as ``something $\times\dd t$ plus something $\times\dd B$'', the first term is the drift (the trend) and the second the noise; a process with no $\dd t$-term is a martingale, a process with a drift is not}
\ts{13}{18}{54:35}.
The 2013 lecture points at the use in pricing: a hedged portfolio with no noise term must have zero drift (relative to the risk-free rate), otherwise it would be an
arbitrage \ts{13}{18}{56:45}, made precise in \cref{ch:bs}.

Examples (checked by direct differentiation): $f=x$; $f=x^2-t$ ($f_t=-1$, $f_{xx}=2$); $f=x^3-3tx$ ($f_t=-3x$, $f_{xx}=6x$); $f=x^4-6tx^2+3t^2$ ($f_t=-6x^2+6t$, $f_{xx}=12x^2-12t$);
in general $f=t^{n/2}H_n(x/\sqrt t)$ with the Hermite polynomials (own remark). And the exponential martingale, the lecture's Example~4:
$f(t,x)=e^{\sigma x-\sigma^2t/2}$ has $f_t=-\tfrac12\sigma^2f$ and $\tfrac12f_{xx}=\tfrac12\sigma^2f$, so \eqref{ch17:eq-pdecond} holds; and
$\E\int_0^t\sigma^2f(s,B_s)^2\,\dd s=\sigma^2\int_0^te^{\sigma^2s}\dd s<\infty$ ($\E[e^{2\sigma B_s-\sigma^2s}]=e^{\sigma^2s}$), so $f(t,B_t)$ is a martingale with $\E[f(t,B_t)]=f(0,0)=1$ for all $t$
\ts{24}{24}{1:12:36}. The variable $f(t,B_t)=e^{\sigma B_t}\,e^{-\sigma^2t/2}$ is a scaled lognormal, and
$\E[e^{\sigma B_t}]=e^{\sigma^2t/2}$ (moment generating function, \cref{ch03:sec-normal}; it was a problem-set question in the class) \ts{24}{24}{1:14:45}.

\begin{remark}[The integrability condition matters; not discussed in class]
A solution of \eqref{ch17:eq-pdecond} with a stochastic integral that is only a \emph{local} martingale, i.e.\ with $\E\int f_x^2=\infty$, need not be a martingale. The condition of the theorem
is the analogue for continuous time of the bounded-increments condition in optional stopping (\cref{ch04:sec-stopping}).
\end{remark}

\subsection{Brownian motion with drift and ruin}\label{ch17:sec-ruin}
Let $X_t=\mu t+\sigma B_t$ and $\tau=\inf\{t:X_t=A\text{ or }X_t=-B'\}$ with $A,B'>0$ (the lecture writes $B$ for the lower barrier; we write $B'$ to avoid a clash with the Brownian motion). We want $\Prob(X_\tau=A)$
\ts{24}{24}{1:15:46}. Look for $h$ such that $M_t=h(X_t)$ is a martingale, so that the stopped process $M_{t\wedge\tau}$ is one too, with boundary values $h(A)=1$, $h(-B')=0$. Then
$\Prob(X_\tau=A)=\E[M_\tau]=M_0=h(0)$ \ts{24}{24}{1:16:48}. By \eqref{ch17:eq-ito3}, $\dd h(X_t)=\bigl(\mu h'+\tfrac12\sigma^2h''\bigr)(X_t)\,\dd t+\sigma h'(X_t)\,\dd B_t$, so the condition
is $\mu h'+\tfrac12\sigma^2h''=0$; the lecture obtains it through $f(t,x)=h(\mu t+\sigma x)$ and \eqref{ch17:eq-pdecond}, where $f_t=\mu h'$ and $f_{xx}=\sigma^2h''$
\ts{24}{24}{1:19:04}. It is a first-order equation for $h'$, $h''=-(2\mu/\sigma^2)h'$, whose solution with the boundary conditions is
\[
  h(y)=\frac{e^{-2\mu y/\sigma^2}-e^{2\mu B'/\sigma^2}}{e^{-2\mu A/\sigma^2}-e^{2\mu B'/\sigma^2}},\qquad \Prob(X_\tau=A)=h(0)
\]
\ts{24}{24}{1:21:11}. This is the formula \eqref{ch11:eq-ruin} of \cref{ch11:sec-ruin} with $a=-B'$, $b=A$ and starting point $0$; there it was obtained by the infinitesimal argument and by
optional stopping applied to $e^{-2\mu X/\sigma^2}$, which is the exponential martingale \eqref{ch17:eq-expmart} with $\theta=-2\mu/\sigma^2$.
It\^o's lemma gives the same result in one line, and shows why $\mu h'+\tfrac12\sigma^2h''=0$ is the right equation: the operator $\mu\partial_x+\tfrac12\sigma^2\partial_x^2$ is the \emph{generator} of $X$.

\begin{exercise}[PS8 (2013), Problem A-1: identify martingales]\label{ch17:ex-a1}
Which of the following are martingales (with respect to their natural filtration)?
\begin{enumerate}[(a)]
  \item A simple random walk.
  \item $X_t=|S_t|$ where $S_t$ is a simple random walk.
  \item $X_0=0$ and $X_{t+1}=X_t+(Y_t-\tfrac1\lambda)$ for $t\ge0$, where the $Y_t$ are i.i.d.\ exponential
      random variables with parameter $\lambda$.
  \item $X_0=0$ and $X_{t+1}=X_t+Z_tZ_{t-1}\cdots Z_0$ for $t\ge0$, where the $Z_t$ are i.i.d.\ lognormal.
  \item $X_t=B^{(1)}_tB^{(2)}_t$, the product of two independent Brownian motions.
\end{enumerate}
\end{exercise}

\begin{exercise}[PS8 (2013), Problem A-2]\label{ch17:ex-a2}
This problem (conditional mean and variance of Brownian motion, of Brownian motion with drift, and $\E\exp(\sigma(B(t)-B(s)))$) is 2024 PS4 Problem~1 and is stated with the other Brownian-motion exercises in \cref{ch11:ex-cond}.
\end{exercise}

% ==================================================================
\section{Geometric Brownian motion}\label{ch17:sec-gbm}

The standard model of a price is that its \emph{relative} change is a Brownian motion with drift,
\begin{equation}\label{ch17:eq-gbm-sde}
  \dd P_t=\mu P_t\,\dd t+\sigma P_t\,\dd B_t,\qquad \mu\in\R,\ \sigma>0,
\end{equation}
\ts{24}{25}{5:16}\ts{13}{18}{25:22}. Since $P_t$ appears on both sides one is tempted to think that $\ln P_t$ has differential $\dd P/P$; It\^o's lemma shows it does not. Apply \eqref{ch17:eq-ito3} to
$G(P,t)=\ln P$, with $G_P=1/P$, $G_t=0$, $G_{PP}=-1/P^2$ and $a=\mu P$, $b=\sigma P$:
\begin{equation}\label{ch17:eq-logp}
  \dd\ln P_t=\Bigl[\frac1P\,\mu P+\frac12\Bigl(-\frac1{P^2}\Bigr)\sigma^2P^2\Bigr]\dd t+\frac1P\,\sigma P\,\dd B_t=\Bigl(\mu-\tfrac12\sigma^2\Bigr)\dd t+\sigma\,\dd B_t
\end{equation}
\ts{24}{25}{6:19}\ts{24}{24}{1:22:14}. The log price is a Brownian motion with drift $\mu-\tfrac12\sigma^2$ and variance rate $\sigma^2$: the drift on the log scale has a \emph{drag} $-\tfrac12\sigma^2$ due to the quadratic variation
\ts{24}{25}{6:19}. Integrating,
\begin{equation}\label{ch17:eq-gbm}
  P_t=P_0\exp\Bigl[\bigl(\mu-\tfrac12\sigma^2\bigr)t+\sigma B_t\Bigr].
\end{equation}
Conversely, \eqref{ch17:eq-gbm} solves \eqref{ch17:eq-gbm-sde}: apply \eqref{ch17:eq-exp} with drift parameter $\mu-\tfrac12\sigma^2$, which gives the drift $\bigl(\mu-\tfrac12\sigma^2+\tfrac12\sigma^2\bigr)P=\mu P$. (That there is no other solution is the uniqueness theorem of \cref{ch:sde}; \cref{ch11:eq-gbm} derived the same solution.)
Equation \eqref{ch17:eq-expmart} is the special case $\mu=0$ and $P_0=1$, i.e.\ the martingale.

\subsection{Distribution, mean and annualised return}\label{ch17:sec-gbmdist}
Over $(t,T]$, with $\tau=T-t$, the log return is a sum of independent normal increments,
\begin{equation}\label{ch17:eq-logret}
  \ln\frac{P_T}{P_t}\sim N\Bigl(\bigl(\mu-\tfrac12\sigma^2\bigr)\tau,\ \sigma^2\tau\Bigr),\qquad
  \E[P_T\mid P_t]=P_te^{\mu\tau},\qquad \Var[P_T\mid P_t]=P_t^2e^{2\mu\tau}\bigl(e^{\sigma^2\tau}-1\bigr)
\end{equation}
\ts{24}{25}{7:23}, by the lognormal moments (\cref{ch03:sec-normal}); the slide's ``$\Var[P_t\mid P_t]$'' should read $\Var[P_T\mid P_t]$.
The parameter $\mu$ is the \emph{expected rate of return} (mean growth of the price), while the typical (median) growth rate is $\mu-\tfrac12\sigma^2$. The lecture's numerical example:
$\mu=0.30$ per year, $\sigma=0.40$, $\tau=\tfrac12$ year. Then
\[
  \E[P_T\mid P_t]=P_t\,e^{0.15}=1.1618\,P_t,\qquad \operatorname{sd}[P_T\mid P_t]=1.1618\sqrt{e^{0.08}-1}\,P_t=0.3353\,P_t
\]
\ts{24}{25}{9:32}: an expected gain of $16.2\%$ rather than the $15\%$ one would obtain by naively halving $\mu$ (the convexity of the exponential, Jensen), and a standard deviation larger than a naive reading of $\sigma$ suggests.
The continuously compounded annualised rate of return $r=\tfrac1\tau\ln(P_T/P_t)$ satisfies
\begin{equation}\label{ch17:eq-rate}
  r\sim N\Bigl(\mu-\tfrac12\sigma^2,\ \frac{\sigma^2}{\tau}\Bigr)
\end{equation}
\ts{24}{25}{11:37}. Hence over long horizons the realised rate of return converges (almost surely, by the strong law $B_T/T\to0$) to the constant $\mu-\tfrac12\sigma^2$, and the variance of the average rate goes to zero as $1/\tau$:
the growth of the price is ``at a constant rate with vanishing variability'' \ts{24}{25}{12:40}. Note that this constant is $\mu-\tfrac12\sigma^2$, \emph{not} $\mu$: if $\mu<\tfrac12\sigma^2$ then $P_t\to0$
almost surely although $\E[P_t]=P_0e^{\mu t}$ grows (own remark; for $\mu=0.05$, $\sigma=0.4$ the typical rate is $-3\%$ per year). The mean is carried by rare, enormous outcomes.

\begin{figure}[ht]
\centering
\begin{tikzpicture}
\begin{axis}[width=0.95\linewidth,height=6.2cm,xlabel={$t$ (years)},ylabel={$X(t)=e^{W(t)}$},xmin=0,xmax=4,ymin=0,ymax=18,
  legend style={at={(0.5,1.03)},anchor=south,legend columns=3,font=\footnotesize},grid=major,grid style={black!10}]
  \addplot[thin,color=defcol!80,forget plot] coordinates {(0.0,1.0) (0.083,1.045) (0.167,1.005) (0.25,0.985) (0.333,0.913) (0.417,0.989) (0.5,0.986) (0.583,0.886) (0.667,0.902) (0.75,0.9) (0.833,0.927) (0.917,0.804) (1.0,0.719) (1.083,0.602) (1.167,0.624) (1.25,0.668) (1.333,0.606) (1.417,0.607) (1.5,0.656) (1.583,0.69) (1.667,0.686) (1.75,0.727) (1.833,0.963) (1.917,1.05) (2.0,1.114) (2.083,1.216) (2.167,1.274) (2.25,1.367) (2.333,1.288) (2.417,1.134) (2.5,1.377) (2.583,1.285) (2.667,1.029) (2.75,1.009) (2.833,0.816) (2.917,0.857) (3.0,0.814) (3.083,0.712) (3.167,0.634) (3.25,0.575) (3.333,0.579) (3.417,0.735) (3.5,0.785) (3.583,0.857) (3.667,0.805) (3.75,0.781) (3.833,0.831) (3.917,0.868) (4.0,0.816)};
  \addplot[thin,color=thmcol!80,forget plot] coordinates {(0.0,1.0) (0.083,0.933) (0.167,0.867) (0.25,0.881) (0.333,0.868) (0.417,0.936) (0.5,1.033) (0.583,1.107) (0.667,1.08) (0.75,1.1) (0.833,1.001) (0.917,1.138) (1.0,1.256) (1.083,1.468) (1.167,1.534) (1.25,1.849) (1.333,1.778) (1.417,2.043) (1.5,1.934) (1.583,1.808) (1.667,2.006) (1.75,1.621) (1.833,1.682) (1.917,1.458) (2.0,1.74) (2.083,1.741) (2.167,2.077) (2.25,2.217) (2.333,2.12) (2.417,1.824) (2.5,2.214) (2.583,2.084) (2.667,2.092) (2.75,2.063) (2.833,2.237) (2.917,2.313) (3.0,2.247) (3.083,2.872) (3.167,3.029) (3.25,3.67) (3.333,3.458) (3.417,4.137) (3.5,4.576) (3.583,4.924) (3.667,4.598) (3.75,5.28) (3.833,6.178) (3.917,6.582) (4.0,7.014)};
  \addplot[thin,color=intcol!80,forget plot] coordinates {(0.0,1.0) (0.083,1.244) (0.167,1.298) (0.25,1.246) (0.333,1.367) (0.417,1.544) (0.5,1.259) (0.583,1.372) (0.667,1.301) (0.75,1.32) (0.833,1.277) (0.917,1.197) (1.0,1.281) (1.083,1.425) (1.167,1.627) (1.25,1.676) (1.333,1.578) (1.417,1.544) (1.5,1.568) (1.583,1.285) (1.667,1.347) (1.75,1.287) (1.833,1.279) (1.917,1.325) (2.0,1.72) (2.083,2.188) (2.167,2.533) (2.25,2.983) (2.333,2.876) (2.417,2.955) (2.5,2.974) (2.583,3.162) (2.667,3.065) (2.75,3.156) (2.833,2.893) (2.917,2.939) (3.0,3.013) (3.083,2.739) (3.167,2.752) (3.25,2.837) (3.333,2.832) (3.417,2.775) (3.5,2.847) (3.583,3.618) (3.667,3.824) (3.75,4.023) (3.833,4.471) (3.917,5.493) (4.0,5.242)};
  \addplot[thin,color=fincol!80,forget plot] coordinates {(0.0,1.0) (0.083,0.99) (0.167,0.968) (0.25,1.02) (0.333,0.955) (0.417,1.049) (0.5,1.24) (0.583,1.442) (0.667,1.229) (0.75,1.323) (0.833,1.359) (0.917,1.405) (1.0,1.286) (1.083,1.478) (1.167,1.509) (1.25,1.632) (1.333,2.245) (1.417,2.489) (1.5,2.217) (1.583,2.078) (1.667,2.218) (1.75,2.12) (1.833,2.029) (1.917,2.189) (2.0,1.918) (2.083,1.984) (2.167,1.837) (2.25,1.431) (2.333,1.454) (2.417,1.448) (2.5,1.519) (2.583,1.6) (2.667,1.622) (2.75,1.642) (2.833,1.576) (2.917,1.563) (3.0,1.632) (3.083,1.293) (3.167,1.137) (3.25,1.56) (3.333,1.55) (3.417,1.541) (3.5,1.715) (3.583,1.573) (3.667,1.487) (3.75,1.538) (3.833,1.364) (3.917,1.588) (4.0,2.018)};
  \addplot[thin,color=cscol!80,forget plot] coordinates {(0.0,1.0) (0.083,0.882) (0.167,0.984) (0.25,1.067) (0.333,1.19) (0.417,1.082) (0.5,1.069) (0.583,1.111) (0.667,1.318) (0.75,1.406) (0.833,1.663) (0.917,1.628) (1.0,1.615) (1.083,1.349) (1.167,1.363) (1.25,1.614) (1.333,1.226) (1.417,1.507) (1.5,1.608) (1.583,1.428) (1.667,1.56) (1.75,1.61) (1.833,1.669) (1.917,1.814) (2.0,1.775) (2.083,2.166) (2.167,2.142) (2.25,2.178) (2.333,2.177) (2.417,2.86) (2.5,2.285) (2.583,1.97) (2.667,2.211) (2.75,1.911) (2.833,2.266) (2.917,2.366) (3.0,2.12) (3.083,2.158) (3.167,2.16) (3.25,1.953) (3.333,1.667) (3.417,1.938) (3.5,1.742) (3.583,1.628) (3.667,1.495) (3.75,1.2) (3.833,1.007) (3.917,1.051) (4.0,1.052)};
  \addplot[thick,black,dotted,domain=0:4,samples=41] {exp(0.30*x)};
  \addplot[thick,black,domain=0:4,samples=41] {exp(0.38*x)};
  \addplot[thick,black!70,dashed,domain=0:4,samples=41] {exp(0.30*x+0.8*sqrt(x))};
  \addplot[thick,black!70,dashed,domain=0:4,samples=41,forget plot] {exp(0.30*x-0.8*sqrt(x))};
  \legend{median,mean,$\pm2$ sd band}
\end{axis}
\end{tikzpicture}
\caption{The simulation of \file{lec19\_2} in the setting of the lecture (\ts{24}{25}{13:42}): $W(t)=0.30\,t+0.40\,B(t)$ on $[0,4]$ and $X=e^{W}$, five paths (different random numbers from the lecture's, $m=1008$ steps of which every $21$st is drawn).
Dotted: the median $e^{0.30t}$; solid: the mean $e^{(0.30+0.08)t}$ (the drift of $X$ in \eqref{ch17:eq-gbm-sde} is $0.30+\tfrac12\cdot0.40^2$); dashed: the band $\exp(0.30t\pm2\cdot0.40\sqrt t)$, the image of the $\pm2$ standard-deviation band of $W$.
The band is far above the typical path, as in the lecture's picture.}\label{ch17:fig-gbm}
\end{figure}

\begin{coursenote}[Two conventions for $\mu$ in the lecture]
Slide~25 of \file{lec19\_1} and \file{lec19\_2} simulate $W$, a Brownian motion with drift $0.30$, volatility $0.40$ and $T=4$ years, and then $X=e^{W}$. In this convention $0.30$ is the drift of the \emph{log} price, and $X$ solves
\eqref{ch17:eq-gbm-sde} with $\mu=0.30+\tfrac12(0.40)^2=0.38$, so that $\E[X_4]=e^{1.52}=4.57$ and $\operatorname{median}X_4=e^{1.2}=3.32$. In the numerical example of the previous slides $0.30$ is the $\mu$ of \eqref{ch17:eq-gbm-sde}: then $\E[P_T]/P_t=e^{0.30\tau}$ and the median is $e^{0.22\tau}$.
Both are right, but they are different processes; the lecture does not remark on it.
\end{coursenote}

\paragraph{The simulation.}
The lecture simulates 5 and then 50 paths, draws the mean and $\pm2$ standard-deviation bands of $W$ and their exponentials, and then the \emph{cumulative average log return} $R(t)=\ln[X(t)/X(0)]/t=W(t)/t$, whose bands are $\mu\pm2\sigma/\sqrt t$
\ts{24}{25}{13:42}. Two observations of the lecture:
\begin{itemize}
  \item among the 50 exponentials some paths reach $10$--$20$ times the initial value, ``extraordinary
      returns'' of the type of a few stocks over some years, consistent with a plain geometric Brownian motion \ts{24}{25}{15:50};
  \item $R(t)$ is very dispersed for short horizons and stabilises around the
      constant level as the horizon grows \ts{24}{25}{16:54}.
\end{itemize}
Our own run (numpy, not the R code of \file{lec19\_2}, so the paths differ) of $50$ paths: $\Prob(X_4>10)=1-\Phi\bigl((\ln10-1.2)/0.8\bigr)=8.4\%$, so $4.2$ of $50$ paths are expected above $10$; the run had one (maximum $17.1$),
and its sample mean of $X_4$ was $3.48$, well below the theoretical $4.57$, because the mean is dominated by rare large paths and $50$ paths are too few to contain them; the median $2.79$ versus $3.32$. The average log return $R(4)$ had mean $0.245$ and standard deviation
$0.183$ over the 50 paths (theory $0.30$ and $0.40/2=0.20$).

\begin{finance}[It\^o's lemma and the Black--Scholes equation]\label{ch17:fin-bs}
The second half of L25 (from 19:01) derives the pricing equation with the tools of this chapter; the argument is developed in \cref{ch:sde,ch:bs}, here is the core. Let $P$ follow \eqref{ch17:eq-gbm-sde} and let
$G_t=G(P_t,t)$ be the price of a derivative that depends only on the asset price and time. By \eqref{ch17:eq-ito3},
\[
  \dd G=\Bigl(G_P\,\mu P+G_t+\tfrac12\sigma^2P^2G_{PP}\Bigr)\dd t+G_P\,\sigma P\,\dd B ,
\]
\ts{24}{25}{19:01}. Short one derivative and hold $G_P$ units of the asset, $V=-G+G_PP$. The noise terms cancel in an infinitesimal time (the $\dd B$ of $G$ is $G_P\sigma P\,\dd B$, that of $G_PP$ is the same),
leaving $\dd V=-\bigl(G_t+\tfrac12\sigma^2P^2G_{PP}\bigr)\dd t$, a riskless change \ts{24}{25}{22:15}. By no arbitrage a riskless portfolio earns the risk-free rate $r$, $\dd V=rV\,\dd t$
\ts{24}{25}{24:16}, whence the \emph{Black--Scholes equation}
\begin{equation}\label{ch17:eq-bs}
  G_t+rPG_P+\tfrac12\sigma^2P^2G_{PP}-rG=0 ,
\end{equation}
valid for \emph{any} contingent claim; different claims differ by the terminal condition at maturity, e.g.\ $G(P,T)=(P-K)^+$ for a call \ts{24}{25}{25:20}. The delta $G_P$ changes with time and the hedge must be rebalanced: this
is the answer given in class to the student's question about the changing amount of asset \ts{24}{25}{29:28}. Note that the drift $\mu$ of the asset has disappeared from \eqref{ch17:eq-bs}: this is the key fact behind risk-neutral pricing (\cref{ch17:sec-girsanov}).
\end{finance}

% ==================================================================
\section{Several dimensions, covariation and the product rule}\label{ch17:sec-multi}

This section is not covered in the lectures (only the two-dimensional Brownian motion of 2013 PS8 Problem~B-1 touches it, see \cref{ch11:sec-2d}); it is the natural completion of the formalism and is marked as our addition.

\subsection{Covariation}
For two processes the multiplication table \eqref{ch17:eq-table} generalises: the \emph{covariation} of $X$ and $Y$ over $[0,t]$ is $[X,Y]_t=\lim\sum_i(\Delta_iX)(\Delta_iY)$ as the mesh goes to zero (in probability). Its differential is written
$\dd X\,\dd Y$. If $\dd X=a\,\dd t+b\,\dd B$ and $\dd Y=a'\,\dd t+b'\,\dd B$ (same $B$) then $\dd X\,\dd Y=bb'\,\dd t$.
For \emph{independent} Brownian motions $B^{(1)},B^{(2)}$,
\begin{equation}\label{ch17:eq-indep}
  \dd B^{(1)}\dd B^{(2)}=0 ,\qquad\text{because}\quad \E\Bigl[\Bigl(\sum_i\Delta_iB^{(1)}\Delta_iB^{(2)}\Bigr)^2\Bigr]=\sum_i\Delta t_i^2\to0
\end{equation}
(the cross terms vanish by independence and mean zero; the diagonal terms are $\E[\Delta B^{(1)2}]\E[\Delta B^{(2)2}]=\Delta t^2$). If $B^{(2)}=\rho B^{(1)}+\sqrt{1-\rho^2}\,\widetilde B$ with $\widetilde B$ independent of $B^{(1)}$ (correlated Brownian motions, \cref{ch:linalg}), then $\dd B^{(1)}\dd B^{(2)}=\rho\,\dd t$.
A numerical check ($2\times10^5$ paths, $200$ steps, $t=1$): the root-mean-square of $\sum\Delta B^{(1)}\Delta B^{(2)}$ was $0.0707$, the theoretical $1/\sqrt{200}=0.0707$.

\subsection{It\^o's formula in several variables}
\begin{theorem}[Multi-dimensional It\^o formula; not discussed in class]\label{ch17:thm-multi}
Let $B=(B^{(1)},\dots,B^{(m)})$ be independent Brownian motions and $X=(X_1,\dots,X_n)$ with $\dd X_i=a_i\,\dd t+\sum_{k=1}^mb_{ik}\,\dd B^{(k)}$ (standard processes). For $f\in C^{1,2}(\R_+\times\R^n)$,
\begin{equation}\label{ch17:eq-multi}
  \dd f(t,X_t)=\Bigl[f_t+\sum_ia_if_{x_i}+\tfrac12\sum_{i,j}(bb\T)_{ij}f_{x_ix_j}\Bigr]\dd t+\sum_{i,k}b_{ik}f_{x_i}\,\dd B^{(k)}_t ,
\end{equation}
i.e.\ $\dd f=f_t\,\dd t+\sum_if_{x_i}\dd X_i+\tfrac12\sum_{i,j}f_{x_ix_j}\,\dd X_i\,\dd X_j$ with $\dd X_i\,\dd X_j=(bb\T)_{ij}\,\dd t$.
\end{theorem}
The proof is the Taylor expansion to second order with \eqref{ch17:eq-table} and \eqref{ch17:eq-indep}. The two consequences we need:

\begin{corollary}[Product rule and integration by parts]\label{ch17:cor-product}
For standard processes $X,Y$, $\dd(XY)=X\,\dd Y+Y\,\dd X+\dd X\,\dd Y$. In particular for $h$ deterministic and continuously differentiable, $\dd(hB)=h\,\dd B+B\,h'\,\dd t$.
\end{corollary}
\begin{proof}
Take $f(x,y)=xy$: $f_{xx}=f_{yy}=0$, $f_{xy}=1$ in \eqref{ch17:eq-multi}. In the second statement $\dd h\,\dd B=h'\,\dd t\,\dd B=0$.
\end{proof}
The extra term $\dd X\,\dd Y$ is absent in classical calculus. For $X=Y=B$ it gives $\dd(B^2)=2B\,\dd B+\dd t$, \cref{ch17:ex-square}; for correlated Brownian motions $\dd(B^{(1)}B^{(2)})=B^{(1)}\dd B^{(2)}+B^{(2)}\dd B^{(1)}+\rho\,\dd t$, so $B^{(1)}B^{(2)}-\rho t$
is a martingale (it is a sum of two It\^o integrals with $\E\int B^2<\infty$); for $\rho=0$ the product of independent Brownian motions is a martingale.

\subsection{Two- and $n$-dimensional Brownian motion}
Let $\bm B=(B^{(1)},\dots,B^{(n)})$ be an $n$-dimensional Brownian motion (independent coordinates); for $n=2$ this is the process of 2013 problem B-1. Apply \eqref{ch17:eq-multi} to a function $f(\bm x)$ of $\bm x\in\R^n$:
\begin{equation}\label{ch17:eq-nD}
  \dd f(\bm B_t)=\nabla f(\bm B_t)\cdot\dd\bm B_t+\tfrac12\Delta f(\bm B_t)\,\dd t ,\qquad\Delta=\sum_i\partial_{x_i}^2 .
\end{equation}
\begin{itemize}
  \item \emph{Harmonic functions.} If $\Delta f=0$ then $f(\bm B_t)$ is a (local) martingale: the martingales of an $n$-dimensional Brownian motion are the harmonic functions. The generator is $\tfrac12\Delta$, as in the heat kernel of \cref{ch11:sec-2d}.
  \item \emph{The squared distance.} For $f=|\bm x|^2$, $\dd|\bm B|^2=2\bm B\cdot\dd\bm B+n\,\dd t$, so $|\bm B_t|^2-nt$ is a martingale and $\E[|\bm B_t|^2]=nt$ (for $n=2$: $2t$; in the physicist's language $\langle r^2\rangle=2dDt$ with diffusion coefficient $D=\tfrac12$ in $d=n$ dimensions).
  \item \emph{The radial part} (Bessel process). For $f=|\bm x|$: $\nabla f=\bm x/|\bm x|$, $\Delta f=(n-1)/|\bm x|$, so $R_t=|\bm B_t|$ satisfies $\dd R=\dd W+\dfrac{n-1}{2R}\,\dd t$ with $\dd W=\dfrac{\bm B}{|\bm B|}\cdot\dd\bm B$,
        and $(\dd W)^2=\sum_iB_i^2/|\bm B|^2\,\dd t=\dd t$; by L\'evy's characterisation (not proved here) $W$ is a Brownian motion. The drift $(n-1)/(2R)$ is the ``centrifugal'' push of a random walk in the plane away from the origin. As a check for $n=2$: $\E[R_t]=\sqrt{\pi t/2}$ from the Rayleigh law of \cref{ch11:sec-2d}, and $\tfrac{\dd}{\dd t}\E[R_t]=\E\tfrac1{2R_t}=\tfrac12\sqrt{\pi/(2t)}$, the derivative of $\sqrt{\pi t/2}$ (using $\E 1/R_t=\sqrt{\pi/(2t)}$).
\end{itemize}

\begin{exercise}[PS8 (2013), Problem B-1: two-dimensional Brownian motion]\label{ch17:ex-b1}
Let $Y(t)=(X_1(t),X_2(t))$ where $X_1,X_2$ are independent Brownian motions (a $2$-dimensional Brownian motion).
\begin{enumerate}[(a)]
  \item For fixed $t$, find the probability density function of $Y(t)$.
  \item Let $D_\rho=\{x\in\R^2:|x|<\rho\}$ and compute $\Prob(Y(t)\in D_\rho)$.
\end{enumerate}
\end{exercise}

\begin{exercise}[not from the course: exit from a disc]\label{ch17:ex-x4}
For a $2$-dimensional Brownian motion started at the origin, let $\tau=\inf\{t:|\bm B_t|=\rho\}$. Use $|\bm B_t|^2-2t$ (\cref{ch17:sec-multi}) and optional stopping (\cref{ch04:sec-stopping}) to find $\E[\tau]$.
\end{exercise}

\begin{exercise}[not from the course: discounted price]\label{ch17:ex-x2}
Let $\dd P=rP\,\dd t+\sigma P\,\dd B$. Show that $\dd(e^{-rt}P_t)=\sigma e^{-rt}P_t\,\dd B_t$ using the product rule, so that the discounted price is a martingale, and compute $\Var P_T$.
\end{exercise}

% ==================================================================
\section{Change of measure and Girsanov's theorem}\label{ch17:sec-girsanov}

This is the last topic of the 2013 lecture; 2024 did not treat it \ts{13}{18}{58:55}. The question is: \emph{can a process with drift be viewed as a process without drift?} Both are probability distributions over paths, and a typical path of the
first drifts along the line $\mu t$ while a typical path of the second stays around the axis \ts{13}{18}{1:00:57}. The answer of Girsanov's theorem is yes, at the price of changing the probability measure, and the consequence
is that ``almost any question about Brownian motion with drift may be rephrased as a parallel question about standard Brownian motion'' (the quotation from Steele in the 2013 note).

\subsection{Equivalent measures and the Radon--Nikod\'ym derivative}
Two probability measures $\Prob,\widetilde\Prob$ on the same space $\Omega$ are \emph{equivalent} if $\Prob(A)>0\iff\widetilde\Prob(A)>0$ for all events $A$: they agree about which events are impossible, although not on their probabilities
\ts{13}{18}{1:03:11}. In the finite example of the lecture, $\Omega=\{1,2,3\}$, $\Prob=(\tfrac13,\tfrac13,\tfrac13)$ and $\widetilde\Prob=(\tfrac23,\tfrac16,\tfrac16)$: for $A=\{1,2\}$ one has $\Prob(A)=\tfrac23$ and $\widetilde\Prob(A)=\tfrac56$, different, but both positive
(the transcript garbles the numbers of $\widetilde\Prob$; these are consistent with the values $\tfrac23,\tfrac56$ quoted). If instead $\widetilde\Prob=(\tfrac23,\tfrac13,0)$ the measures are not equivalent: $\Prob(\{3\})=\tfrac13$ but $\widetilde\Prob(\{3\})=0$
\ts{13}{18}{1:06:22}. Equivalent measures are related by a positive density:

\begin{theorem}[Radon--Nikod\'ym; quoted]\label{ch17:thm-rn}
$\widetilde\Prob$ is equivalent to $\Prob$ if and only if there is a random variable $Z>0$ with $\E[Z]=1$ such that $\widetilde\Prob(A)=\E[Z\,\ind_A]$ for all $A$; then $\widetilde\E[Y]=\E[ZY]$. One writes $Z=\dd\widetilde\Prob/\dd\Prob$, the \emph{Radon--Nikod\'ym derivative}.
\end{theorem}
\ts{13}{18}{1:04:17}\ \noindent
In the finite example $Z=(2,\tfrac12,\tfrac12)$. A mapping of paths in general destroys equivalence: the law of the \emph{square} $B(t)^2$ is not equivalent to that of $B(t)$, because $B^2$ never takes negative values and $B$ does \file{lecnote18}, \S3 (this is essentially PS8 Problem~A-5).

\begin{exercise}[PS8 (2013), Problem A-5: non-equivalent laws]\label{ch17:ex-a5}
Let $B(t)$ be a Brownian motion. Prove that the probability distributions of $B(t)$ and of $B(t)^2$ are not equivalent (two probability distributions $\Prob$ and $\Q$ are equivalent if for every set $X$, $\Prob(X)>0$ if and only if $\Q(X)>0$).
\end{exercise}

\subsection{Girsanov's theorem}
\begin{theorem}[Girsanov, simple version]\label{ch17:thm-girsanov1}
Let $X$ be a Brownian motion with drift $\mu$ (and unit variance) on $[0,T]$ under $\Prob$, and define $\widetilde\Prob$ by
\begin{equation}\label{ch17:eq-Z}
  \frac{\dd\widetilde\Prob}{\dd\Prob}=Z=\exp\Bigl(-\mu X_T+\tfrac12\mu^2T\Bigr).
\end{equation}
Then $\widetilde\Prob$ is a probability measure, equivalent to $\Prob$, and under $\widetilde\Prob$ the process $X$ is a standard Brownian motion \emph{without} drift.
\end{theorem}
\begin{proof}
Take $0=t_0<\dots<t_n=T$ and the increments $x_i=X_{t_i}-X_{t_{i-1}}$; under $\Prob$ they are independent $N(\mu\Delta_i,\Delta_i)$ with $\Delta_i=t_i-t_{i-1}$, with joint density $\prod_i(2\pi\Delta_i)^{-1/2}e^{-(x_i-\mu\Delta_i)^2/2\Delta_i}$. Since $Z=\prod_ie^{-\mu x_i+\mu^2\Delta_i/2}$,
the density of the increments under $\widetilde\Prob$, i.e.\ of $Z\,\dd\Prob$, is the product of
\[
  (2\pi\Delta_i)^{-1/2}\exp\Bigl[-\frac{(x_i-\mu\Delta_i)^2}{2\Delta_i}-\mu x_i+\frac{\mu^2\Delta_i}2\Bigr]=(2\pi\Delta_i)^{-1/2}\exp\Bigl[-\frac{x_i^2}{2\Delta_i}\Bigr],
\]
because $-\tfrac{(x-\mu\Delta)^2}{2\Delta}=-\tfrac{x^2}{2\Delta}+\mu x-\tfrac12\mu^2\Delta$. This is a probability density (so $\E[Z]=1$) and says that the increments are independent $N(0,\Delta_i)$ under $\widetilde\Prob$. The finite-dimensional distributions of $X$ are those of Brownian motion, and the paths are continuous.
\end{proof}
The path picture of a physicist: with $\Prob[x(\cdot)]\propto e^{-\frac12\int(\dot x-\mu)^2}$ and $\widetilde\Prob[x(\cdot)]\propto e^{-\frac12\int\dot x^2}$, the ratio is $e^{\mu(x_T-x_0)-\mu^2T/2}=1/Z$: reweighting paths by an exponential in the endpoint (an exponential tilting, or Cameron--Martin) shifts the drift.

\begin{coursenote}[Erratum: the sign in the note's Theorem 3.2]
\file{lecnote18} (Theorem~3.2) states $\dd\widetilde\Prob/\dd\Prob=e^{-\mu\omega(T)-\mu^2T/2}$ for $X$ a Brownian motion with drift $\mu$ under $\Prob$. The correct exponent is $-\mu X_T\,\mathbf{+}\,\tfrac12\mu^2T$, as in \eqref{ch17:eq-Z}. With the sign of the note,
$\E[Z]=\E[e^{-\mu X_T-\mu^2T/2}]=e^{-\mu^2T/2}e^{-\mu^2T/2}=e^{-\mu^2T}\neq1$ (as $\E[e^{-\mu X_T}]=e^{-\mu^2T+\mu^2T/2}$), so $Z$ would not be a density (Monte Carlo, $\mu=0.5$, $T=1$: $0.7788=e^{-0.25}$ against $1.0000$ for the correct sign; with the correct sign $\E[Z\,\ind_{X_T>1}]=0.1586=1-\Phi(1)$ and $\E[ZX_T^2]=1.00$, the values of a standard normal).
The general theorem (\cref{ch17:thm-girsanov2}) below, which the note states correctly, is consistent with it (take $\theta=\mu$ and start from $X-\mu t$).
\end{coursenote}

\begin{theorem}[Girsanov, general version]\label{ch17:thm-girsanov2}
Let $B$ be a Brownian motion under $\Prob$ on $[0,T]$, $\theta$ a bounded adapted process, and
\begin{equation}\label{ch17:eq-Zt}
  Z_T=\exp\Bigl(-\int_0^T\theta_u\,\dd B_u-\tfrac12\int_0^T\theta_u^2\,\dd u\Bigr),\qquad \widetilde\Prob=Z_T\cdot\Prob .
\end{equation}
Then $Y_t=B_t+\int_0^t\theta_u\,\dd u$ is a Brownian motion under $\widetilde\Prob$.
\end{theorem}
\ts{13}{18}{1:10:39}
\begin{proof}[Proof sketch (not given in class)]
By \eqref{ch17:eq-ito3}, $\dd Z_t=-\theta_tZ_t\,\dd B_t$, so $Z$ is a positive martingale with $Z_0=1$, $\E[Z_T]=1$ (the integrand is square-integrable since $\theta$ is bounded), and by Bayes $\widetilde\E[G\mid\mathcal F_s]=\E[Z_tG\mid\mathcal F_s]/Z_s$.
For real $\lambda$ put $N_t=Z_t\exp(\lambda Y_t-\tfrac12\lambda^2t)$. Then
$\ln N_t=\int_0^t(\lambda-\theta)\,\dd B-\tfrac12\int_0^t(\lambda-\theta)^2\dd u$ (the terms $\lambda\int\theta$ and $-\tfrac12\lambda^2t$ complete the square), so $N$ is again a martingale of the form \eqref{ch17:eq-Zt}. Then for $s<t$
\[
  \widetilde\E\bigl[e^{\lambda(Y_t-Y_s)}\bigm|\mathcal F_s\bigr]=\frac{\E[N_t\mid\mathcal F_s]\,e^{\lambda^2t/2}e^{-\lambda Y_s}}{Z_s}=e^{\lambda^2(t-s)/2} .
\]
This is the moment generating function of $N(0,t-s)$ for every $\lambda$, so $Y_t-Y_s$ is $N(0,t-s)$ under $\widetilde\Prob$ and independent of $\mathcal F_s$: $Y$ is a Brownian motion.
\end{proof}
For constant $\theta=\mu$: $Z_T=e^{-\mu B_T-\mu^2T/2}$ and $Y=B+\mu t$ is Brownian under $\widetilde\Prob$; equivalently $B=Y-\mu t$ is a Brownian motion with drift $-\mu$ under $\widetilde\Prob$. The general theorem lets the drift removed be random (adapted).
The 2013 lecture states the simple version and emphasises how to use it: an expectation of a functional of a process with drift under $\Prob$ becomes an expectation of a functional of an undrifted process under $\widetilde\Prob$, $\E[F]=\widetilde\E[F/Z]$ (the lecture writes ``$Z$ tilde'', as it swaps the roles of the two measures)
\ts{13}{18}{1:14:50}. A martingale is a fair game; the change of measure converts an unfair game into a fair one, and this is why it matters for pricing.

\begin{finance}[Risk-neutral measure; not discussed in class]
Let the real-world price follow \eqref{ch17:eq-gbm-sde} and let $r$ be the risk-free rate. Take the constant $\theta=(\mu-r)/\sigma$ (the \emph{market price of risk}) and $Y=B+\theta t$, a Brownian motion under $\widetilde\Prob=\Q$. Then
$\dd P=\mu P\,\dd t+\sigma P(\dd Y-\theta\,\dd t)=rP\,\dd t+\sigma P\,\dd Y$: under $\Q$ the price grows at the risk-free rate, and $e^{-rt}P_t$ is a martingale. The derivative price is the discounted $\Q$-expectation of the payoff (\cref{ch:bs});
this is the same reason why $\mu$ disappeared from \eqref{ch17:eq-bs}. Both measures are equivalent, so ``impossible under one'' means ``impossible under the other'', which is exactly what no-arbitrage requires (\cref{thm:ftap}).
\end{finance}

\begin{remark}[Why equivalence is a surprise]
The two lectures record the same astonishment \ts{13}{18}{1:09:36}: paths with and without drift look completely different, since as $t\to\infty$ one hugs the line $\mu t$ and the other the axis. Both statements are true: on a \emph{finite} interval $[0,T]$ the measures are equivalent, the density $Z$ being an
explicit, finite, positive number; on $[0,\infty)$ they are mutually singular, since $\{\lim X_t/t=\mu\}$ has probability $1$ under $\Prob$ and $0$ under $\widetilde\Prob$ (strong law of large numbers), and $Z_T=e^{-\mu X_T+\mu^2T/2}\to0$ almost surely under $\Prob$ as $T\to\infty$.
\end{remark}

\begin{remark}[The discrete picture]
P.~Kempthorne's comment during the 2013 lecture \ts{13}{18}{1:10:39}: a gambler's fortune against a casino with unfavourable odds is a random walk with a negative drift; taking out the expected loss turns it into a martingale, to which martingale theory (convergence theorems) applies. In the discrete case the change of measure is a likelihood ratio: a
walk with up-probability $p=\tfrac12(1+\varepsilon)$ against the fair one has, over $n$ steps with $S_n=\sum\xi_i$, $\prod_i\bigl[(2p)^{-(1+\xi_i)/2}(2q)^{-(1-\xi_i)/2}\bigr]=\exp\bigl(-\varepsilon S_n+\tfrac12n\varepsilon^2+O(n\varepsilon^3)\bigr)$;
with steps $\pm\sqrt h$ and $\varepsilon=\mu\sqrt h$ this tends to $e^{-\mu X_T+\mu^2T/2}$, the density \eqref{ch17:eq-Z} (own computation).
\end{remark}

\begin{exercise}[not from the course: Girsanov by simulation]\label{ch17:ex-x3}
For $\mu=0.5$, $T=1$ let $X_T\sim N(\mu,1)$ under $\Prob$ and $Z=e^{-\mu X_T+\mu^2/2}$. Compute (exactly, and by simulation) $\E[Z]$, $\E[Z\ind_{\{X_T>1\}}]$ and $\E[Z\,X_T^2]$, and explain the results using \cref{ch17:thm-girsanov1}.
\end{exercise}

% ==================================================================
\section{Errata and the two lectures compared}\label{ch17:sec-compare}

\begin{coursenote}[2013 and 2024 compared]
\emph{Order of presentation.} 2013 (C.~Lee) starts from the \emph{differentiation}: Taylor expansion, $(\dd B)^2=\dd t$ and the lemma for $f(t,B_t)$ and $f(t,X_t)$ come first; the integral is \emph{defined as the inverse} of differentiation
(``a very stupid definition'' whose existence is left open), and only afterwards is it linked to left-point Riemann sums. 2024 (P.~Kempthorne) goes the other way, from integrals of step processes through the isometry to the limit and then derives the formula from Taylor's theorem; the two constructions agree.
\emph{Content only in 2013:} adapted processes with examples, the ``normal'' theorem for deterministic integrands, the martingale-versus-drift dictionary and its use in arbitrage arguments, change of measure and Girsanov. \emph{Content only in 2024:} the explicit step-process cases
including $\int B\,\dd B$ as a limit, the isometry as an isometry of $L^2$ spaces, Gaussian integrals with covariance, the anti-derivative method, the PDE condition and the ruin problem, geometric Brownian motion with numerical example and simulation, the derivation of the Black--Scholes equation.
\emph{Common:} both give the isometry and the martingale property without proof, both stop at the level of ``reasonable functions'' (the 2013 note writes $g\in L^2$; the 2024 slides $\E\int f^2<\infty$).
\end{coursenote}

\begin{coursenote}[Further errata and misprints]
\begin{itemize}
  \item \file{lecnote18} Theorem~2.3 writes the integrability condition as ``$\int\!\!\int_0^tg^2(t,B_t)\,\dd t\,\dd B_t<\infty$''; it should be $\E\int_0^tg(s,B_s)^2\,\dd s<\infty$, the condition \eqref{ch17:eq-h2}.
  \item \file{lecnote18} Example~2.2(iv) claims that $X_\tau$ is an adapted process for a stopping time $\tau$; $X_\tau$ is a single random variable, not a process. The intended statement is that the stopped process $t\mapsto X_{t\wedge\tau}$ is adapted.
  \item \file{lecnote18} Example~1.3(i) calls $\mu t+\sigma B_t$ Brownian motion with ``variance $\sigma$''; the variance parameter is $\sigma^2$.
  \item Slide~7 of \file{lec19\_1} defines the step process with a last time $t_{n+1}=T$ (it is $t_n=T$); slide~10 writes $\Var[I(f)]=\|I(f)\|_{L^2}$, where the norm has to be squared (the slide says ``(squared) norm'' below).
  \item In the 2013 lecture the differential is stated aloud as $\dd f=f'\,\dd B+f''\,\dd t$ \ts{13}{18}{1:11}: the coefficient of the last term is $\tfrac12$, correctly written in the note.
  \item The 2024 transcript (spoken) says that the expectation of the increment powers $c>2$ is ``comparable to'' $\sigma^2\Delta t$ \ts{24}{24}{7:24}; the slide correctly says $o(\Delta t)$.
\end{itemize}
\end{coursenote}

\paragraph{Looking ahead.} \Cref{ch:sde} takes \eqref{ch17:eq-ito3} to solve stochastic differential equations: existence and uniqueness, the Ornstein--Uhlenbeck and Vasicek models, and the numerical (Euler--Maruyama) scheme whose strong order $n^{-1/2}$ we saw in \cref{ch17:ex-exp}; the Feynman--Kac formula
links the PDE condition of \cref{ch17:thm-pde} to expectations. \Cref{ch:bs} uses \eqref{ch17:eq-bs} and the change of measure of \cref{ch17:sec-girsanov} to price options.

% ==================================================================
\begin{keyformulas}
\begin{align*}
&\text{Table:} && (\dd B)^2=\dd t,\quad \dd B\,\dd t=0,\quad(\dd t)^2=0,\quad \dd B^{(1)}\dd B^{(2)}=\rho\,\dd t\\
&\text{It\^o integral:} && \E\!\int_0^tf\,\dd B=0,\qquad \E\Bigl(\int_0^tf\,\dd B\Bigr)^2=\E\!\int_0^tf^2\,\dd s\\
& && \int_0^tB\,\dd B=\tfrac12\bigl(B_t^2-t\bigr)\quad(\text{left point})\\
&\text{It\^o's lemma:} && \dd f(t,X_t)=\Bigl(f_t+af_x+\tfrac12b^2f_{xx}\Bigr)\dd t+bf_x\,\dd B,\qquad\dd X=a\,\dd t+b\,\dd B\\
&\text{Product rule:} && \dd(XY)=X\,\dd Y+Y\,\dd X+\dd X\,\dd Y\\
&\text{Martingale test:} && f_t+\tfrac12f_{xx}=0\ \Rightarrow\ f(t,B_t)\text{ martingale}\quad(\text{e.g.\ }e^{\sigma B-\sigma^2t/2},\ B^2-t)\\
&\text{GBM:} && \dd P=\mu P\,\dd t+\sigma P\,\dd B\ \Longrightarrow\ P_t=P_0e^{(\mu-\sigma^2/2)t+\sigma B_t}\\
& && \E[P_t]=P_0e^{\mu t},\qquad r=\tfrac1\tau\ln\tfrac{P_T}{P_t}\sim N\bigl(\mu-\tfrac12\sigma^2,\sigma^2/\tau\bigr)\\
&\text{Girsanov:} && Z=e^{-\mu X_T+\mu^2T/2}\ (X\text{ has drift }\mu\text{ under }\Prob)\\
& && \Rightarrow\ X\text{ is driftless under }Z\cdot\Prob
\end{align*}
\end{keyformulas}

% !TEX root = ../main.tex
\chapter{Stochastic Differential Equations}\label{ch:sde}
\begin{paracol}{2}

\begin{sources}
\begin{tabularx}{\linewidth}{@{}l l L@{}}
\toprule
\textbf{Lecture} & \textbf{Speaker} & \textbf{Content}\\
\midrule
2024 L25 (from 19:01) & P.~Kempthorne & The first 18 minutes (recap of It\^o's formula, geometric Brownian motion) are in \cref{ch:stochcalc}. Then: the Black--Scholes equation by hedging (slides 2--4, also \cref{ch:stochcalc,ch:bs}); the heat equation: linearity and superposition, the Dirac initial condition, general initial data, the indicator example, solution by similarity/invariance and by step functions (slides 5--12); SDEs: Lipschitz and growth conditions, existence and uniqueness, coefficient matching for geometric Brownian motion and its ``paradox'', the Ornstein--Uhlenbeck process and Vasicek's model, the product-process solution, systems of SDEs and the position--velocity example (slides 13--20); numerical methods listed on slide~21.\\
2013 L21 & C.~Lee & The goal of an SDE and the existence/uniqueness theorem; coefficient matching for geometric Brownian motion and (by a product guess) for the Ornstein--Uhlenbeck process; numerical methods: finite differences for ODEs, Monte Carlo along a fixed Brownian path, the tree method; the heat equation, superposition, the Dirac delta and the Gaussian kernel, the link with the Black--Scholes equation.\\
\bottomrule
\end{tabularx}

\smallskip
\textbf{Files.} Slides \file{mit18\_642\_f24\_lec24.pdf} (21 slides, this is Lecture~25, not 24; slides 2--4 Black--Scholes, 5--12 heat equation, 13--20 SDEs, 21 numerical methods and references);
2013 note \file{MIT18\_S096F13\_lecnote21.pdf} (7 pages).
\textbf{Problem sets.} 2013 PS9 (\file{pset9.pdf}: Problems A-1, A-2, A-3; there is no Part~B), and Exercise~3.2 of \file{lecnote21}. 2024 has no problem set on this topic. The exercises are at the end of the chapter.
\textbf{Additions of these notes} (marked as such where they occur): the proof of existence and uniqueness by Picard iteration, generators, the Feynman--Kac formula, the Fokker--Planck equation, the exact AR(1) sampling of the Ornstein--Uhlenbeck process, the Vasicek bond price, the Euler--Maruyama and Milstein convergence experiments, and the CIR process. The lectures state only that the Black--Scholes PDE and expectations are two faces of one object; the Feynman--Kac formula that proves it is here.
\end{sources}

\section*{Motivation}
\Cref{ch:stochcalc} taught us to differentiate functions of Brownian motion. This chapter runs the calculus backwards, exactly as an ordinary differential equation is the reverse of differentiation: \emph{given} the local rules ``expected change per unit time $=\mu$, variance of the change per unit time $=\sigma^2$'', find the process. A stochastic differential equation (SDE) is the law of motion of a random quantity: a stock price, a short interest rate, the velocity of a particle in a fluid. Four questions organise the chapter.
\begin{itemize}
  \item \emph{Does a solution exist, and is it unique?} Yes under a Lipschitz and a growth condition (\cref{ch18:sec-exist}).
  \item \emph{Can we solve it in closed form?} Rarely, but the two most important examples can: geometric Brownian motion, the model of a price, and the Ornstein--Uhlenbeck (Vasicek) process, the model of a mean-reverting quantity such as an interest rate (\cref{ch18:sec-cm,ch18:sec-ou}). Sampled at discrete times the Ornstein--Uhlenbeck process is exactly the AR(1) model of \cref{ch:timeseries}.
  \item \emph{How does an SDE connect to partial differential equations?} Through the generator of the process: the Feynman--Kac formula turns a PDE into an expectation and back, and the Fokker--Planck equation is the PDE for the density (\cref{ch18:sec-fk}). Both lectures end with the simplest of these PDEs, the heat equation, which is the Black--Scholes equation in disguise (\cref{ch18:sec-heat}).
  \item \emph{What if there is no closed form?} Simulate: the Euler--Maruyama scheme, Monte Carlo, trees (\cref{ch18:sec-num}).
\end{itemize}
For a physicist most of this is familiar under other names: an SDE with additive noise is a Langevin equation, the Ornstein--Uhlenbeck process is the velocity of a Brownian particle, and the Fokker--Planck equation and Feynman--Kac formula are the Schr\"odinger-type equation and the path integral of the same diffusion. We use It\^o's formula \eqref{ch17:eq-ito3} of \cref{ch:stochcalc} throughout.

% ==================================================================
\section{Stochastic differential equations: existence and uniqueness}\label{ch18:sec-exist}

\subsection{The definition}
Let $B$ be a standard Brownian motion and $\mu(t,x),\sigma(t,x)$ two given functions of time and space. The SDE
\begin{equation}\label{ch18:eq-sde}
  \dd X_t=\mu(t,X_t)\,\dd t+\sigma(t,X_t)\,\dd B_t,\qquad X_0=x_0
\end{equation}
\ts{24}{25}{1:04:34}\ts{13}{21}{0:00} is a \emph{shorthand for an integral equation}, exactly as in It\^o's formula. A \emph{solution} is a continuous adapted process $X$ such that for all $t\ge0$
\begin{equation}\label{ch18:eq-sdeint}
  X_t=x_0+\int_0^t\mu(s,X_s)\,\dd s+\int_0^t\sigma(s,X_s)\,\dd B_s ,
\end{equation}
the first integral being an ordinary (pathwise) integral and the second an It\^o integral (\cref{ch:stochcalc}). Note the difference with all of \cref{ch:stochcalc}: there the integrands were given adapted processes; here the integrand $\sigma(s,X_s)$ depends on the unknown $X$ itself, so the object to be found appears on both sides.
\begin{intuition}[Reading an SDE]
Over a short time $h$, \eqref{ch18:eq-sde} says $X_{t+h}-X_t\approx\mu(t,X_t)\,h+\sigma(t,X_t)\,(B_{t+h}-B_t)$: conditionally on the present the change is normal with mean $\mu h$ and variance $\sigma^2h$. So $\mu$ is the local drift (mean velocity) and $\sigma^2$ the local variance rate, the \emph{diffusion coefficient} of a physicist ($\sigma^2=2D$). In the language of \cref{ch:stochcalc} the equation is a Langevin equation $\dot X=\mu(t,X)+\sigma(t,X)\,\xi_t$ with white noise $\xi=\dot B$, in the It\^o convention: the noise multiplies $\sigma$ evaluated at the \emph{start} of each step.
\end{intuition}
The 2013 lecture stresses that this is the reverse of what had been done so far \ts{13}{21}{0:00}\ts{13}{21}{1:06}: before, a function was differentiated to obtain an equation; now the equation is given and the process has to be found by going backwards. Even for ordinary or partial differential equations one does not expect a closed form; what one \emph{can} always say is whether a solution exists.

\subsection{Existence and uniqueness}
\begin{theorem}[Existence and uniqueness]\label{ch18:thm-eu}
Fix $T>0$ and suppose the coefficients of \eqref{ch18:eq-sde} satisfy, for some constant $K$ and all $t\in[0,T]$, $x,y\in\R$, the \emph{space-variable Lipschitz condition}
\begin{equation}\label{ch18:eq-lip}
  |\mu(t,x)-\mu(t,y)|^2+|\sigma(t,x)-\sigma(t,y)|^2\le K|x-y|^2
\end{equation}
and the \emph{spatial growth condition}
\begin{equation}\label{ch18:eq-growth}
  |\mu(t,x)|^2+|\sigma(t,x)|^2\le K\bigl(1+|x|^2\bigr).
\end{equation}
Then there is a continuous adapted solution $X$ on $[0,T]$ with $\sup_{0\le t\le T}\E X_t^2<\infty$. Moreover if $X$ and $Y$ are two continuous solutions with this $L^2$ bound, then $\Prob\bigl(X_t=Y_t\ \text{for all }t\in[0,T]\bigr)=1$.
\end{theorem}
\ts{24}{25}{1:04:34}\ts{24}{25}{1:05:34}\ts{24}{25}{1:06:36}\ts{13}{21}{2:13}\ts{13}{21}{3:18}\ts{13}{21}{4:27}
Both lectures state the theorem without proof, refer to the literature (Steele; Oksendal), and comment on the hypotheses. The Lipschitz condition says that $\mu$ and $\sigma$ cannot change too fast when the space variable moves a little: at most linearly in the distance \ts{13}{21}{4:27}. The growth condition says that the coefficients cannot blow up too fast at infinity \ts{24}{25}{1:05:34}\ts{13}{21}{5:31}. ``Essentially unique'' means unique up to a null set: the two solutions have the same paths with probability one \ts{24}{25}{1:06:36}. The 2013 speaker adds the advice that matters in practice: do not worry about the technical conditions; \emph{a solution exists, but do not expect a closed form} \ts{13}{21}{5:31}.

\begin{proof}[Sketch by Picard iteration (not given in class)]
Put $X^{(0)}_t=x_0$ and $X^{(n+1)}_t=x_0+\int_0^t\mu(s,X^{(n)}_s)\dd s+\int_0^t\sigma(s,X^{(n)}_s)\dd B_s$. Using $(a+b)^2\le2a^2+2b^2$, Cauchy--Schwarz for the $\dd s$ integral, the It\^o isometry (\cref{ch17:thm-ito-int}) for the noise integral, and \eqref{ch18:eq-lip} (which contains the Lipschitz bound of $\mu$ and of $\sigma$ separately), the differences $\varphi_n(t)=\E|X^{(n+1)}_t-X^{(n)}_t|^2$ satisfy
\[
  \varphi_n(t)\le C\int_0^t\varphi_{n-1}(s)\,\dd s,\qquad C=2K(T+1),
\]
hence $\varphi_n(t)\le\varphi_0(T)(Ct)^n/n!$ by induction, and $\varphi_0(T)<\infty$ by \eqref{ch18:eq-growth}. The series $\sum_n\sqrt{\varphi_n}$ converges, so $X^{(n)}$ is Cauchy in $L^2$ uniformly in $t$; its limit solves \eqref{ch18:eq-sdeint}. (Continuity of paths needs Doob's maximal inequality; we skip it.) For uniqueness, if $X,Y$ are two $L^2$-bounded solutions then $D(t)=\E|X_t-Y_t|^2$ satisfies $D(t)\le C\int_0^tD(s)\,\dd s$; Gronwall's lemma gives $D\equiv0$, so $X_t=Y_t$ a.s.\ for each $t$, and by continuity for all $t$ simultaneously.
\end{proof}
The two conditions are \emph{sufficient, not necessary}, and each has its own job: the Lipschitz condition prevents non-uniqueness and the growth condition prevents explosion. Two ordinary differential equations (the case $\sigma=0$) show what goes wrong.
\begin{example}[Why both conditions matter; not discussed in class]\label{ch18:ex-fail}
\leavevmode
\begin{enumerate}[(a)]
  \item $\dot x=x^2$, $x(0)=1$ has the solution $x(t)=1/(1-t)$, which explodes at $t=1$: $\mu(x)=x^2$ is locally Lipschitz but violates \eqref{ch18:eq-growth}.
  \item $\dot x=3x^{2/3}$, $x(0)=0$ has the two solutions $x\equiv0$ and $x=t^3$: $x^{2/3}$ is not Lipschitz at $0$. The stochastic analogue is the SDE of PS9 Problem~A-2, $\dd X=\tfrac13X^{1/3}\dd t+X^{2/3}\dd B$ (real cube roots): with $X_0=a>0$ the process $X_t=(a^{1/3}+B_t/3)^3$ is a solution, and it is the only one \emph{until $X$ first hits $0$}, where the coefficients stop being Lipschitz. For $a=0$ both $X\equiv0$ and $X_t=(B_t/3)^3$ solve the equation (apply It\^o's formula to $f(B)=(B/3)^3$): uniqueness fails exactly where \eqref{ch18:eq-lip} fails.
\end{enumerate}
\end{example}
Under only \emph{local} Lipschitz continuity one still has uniqueness and existence up to an explosion time. The solution of \eqref{ch18:eq-sde} under the conditions of \cref{ch18:thm-eu} is a \emph{Markov process} (quoted): given $X_t$, the future depends on the past only through the present, because the future noise is independent of $\mathcal F_t$. This is what makes the generator of \cref{ch18:sec-fk} meaningful.

\begin{coursenote}[Misprints in the 2013 note, Theorem 1.1]
\file{lecnote21} introduces the equation ``for given functions $a$ and $b$'' although the coefficients are called $\mu,\sigma$; writes the solution as ``$X(T)=\int_0^T\dots$'', omitting the initial value $x_0$ (it is \eqref{ch18:eq-sdeint} with $t=T$); and leaves the placeholder ``such that (L2 bound)'', which the 2024 slide 13 spells out as $\sup_{0\le t\le T}\E X_t^2<\infty$, as in \cref{ch18:thm-eu}.
\end{coursenote}

% ==================================================================
\section{Solving by coefficient matching: geometric Brownian motion}\label{ch18:sec-cm}

\subsection{The method}
Both lectures show the one systematic recipe they have: \emph{guess that the solution is a function $X_t=f(t,B_t)$ of time and the driving Brownian motion}, differentiate by It\^o's formula \eqref{ch17:eq-ito2}, and match coefficients with the SDE \ts{24}{25}{1:07:37}\ts{13}{21}{6:34}. For $f\in C^{1,2}$,
\[
  \dd X_t=\Bigl(f_t+\tfrac12f_{xx}\Bigr)(t,B_t)\,\dd t+f_x(t,B_t)\,\dd B_t .
\]
Comparing with \eqref{ch18:eq-sde} (and using that $X_t=f(t,B_t)$) we need
\begin{equation}\label{ch18:eq-match}
  \mu\bigl(t,f\bigr)=f_t+\tfrac12f_{xx},\qquad \sigma\bigl(t,f\bigr)=f_x .
\end{equation}
These are two equations for \emph{one} unknown function, so the method can only work for special SDEs; the 2013 lecturer warns that it solves only a few SDEs and that most of the time it will not apply \ts{13}{21}{6:34}\ts{13}{21}{11:01}\ts{13}{21}{19:56}. It is a systematic way to \emph{guess}, and in the end the guess is verified by It\^o's formula.

\begin{example}[Geometric Brownian motion]\label{ch18:ex-gbm}
Let $\mu,\sigma>0$ be constants and $\dd X=\mu X\,\dd t+\sigma X\,\dd B$, $X_0=x_0>0$. Then \eqref{ch18:eq-match} reads $\mu f=f_t+\tfrac12f_{xx}$ and $\sigma f=f_x$. The second equation is an ODE in $x$: $f_x/f=\sigma$, so $f(t,x)=e^{\sigma x+g(t)}$ \ts{24}{25}{1:08:48}\ts{13}{21}{7:45}\ts{13}{21}{8:50}. Inserting in the first,
\[
  \mu f=g'(t)f+\tfrac12\sigma^2f\ \Longrightarrow\ g'(t)=\mu-\tfrac12\sigma^2,\qquad g(t)=\bigl(\mu-\tfrac12\sigma^2\bigr)t+c ,
\]
the factor $f$ dividing out \ts{24}{25}{1:09:56}\ts{24}{25}{1:11:00}\ts{13}{21}{9:55}. The initial condition $f(0,0)=x_0$ fixes $e^c=x_0$, hence
\begin{equation}\label{ch18:eq-gbmsol}
  X_t=x_0\exp\Bigl[\bigl(\mu-\tfrac12\sigma^2\bigr)t+\sigma B_t\Bigr].
\end{equation}
\end{example}
This is the formula \eqref{ch17:eq-gbm} that \cref{ch:stochcalc} obtained from the logarithm; here the same answer comes from a different logic, and by \cref{ch18:thm-eu} it is \emph{the} solution: the coefficients $\mu x,\sigma x$ are Lipschitz and of linear growth. The 2013 lecturer is candid about the status of the method \ts{13}{21}{11:01}\ts{13}{21}{12:01}: we already knew the answer and merely fitted it; the point of the exercise is the algorithm, which works when one cannot guess $f$.

\subsection{The lognormal law and the ``paradox''}
By \eqref{ch18:eq-gbmsol}, $\ln X_t\sim N\bigl(\ln x_0+\mu_*t,\ \sigma^2t\bigr)$ with $\mu_*=\mu-\tfrac12\sigma^2$: $X_t$ is lognormal with $\E X_t=x_0e^{\mu t}$ (the $e^{-\sigma^2t/2}$ from the drift cancels $\E e^{\sigma B_t}=e^{\sigma^2t/2}$) \ts{24}{25}{1:11:00}. The details, the numerical example of the 2024 lecture ($\mu=0.3$, $\sigma=0.4$, half a year: $\E P_T/P_t=1.1618$, standard deviation $0.335\,P_t$) and the annualised return are in \cref{ch17:sec-gbmdist}.

The lecture then points to a seeming paradox \ts{24}{25}{1:12:02}\ts{24}{25}{1:13:03}. Write $\ln X_t=\ln x_0+t\bigl(\mu_*+\sigma B_t/t\bigr)$. Since $B_t/t\to0$ almost surely (strong law of large numbers for Brownian motion),
\begin{equation}\label{ch18:eq-paradox}
  0<\mu<\tfrac12\sigma^2\ \Longrightarrow\ \E X_t=x_0e^{\mu t}\to\infty\quad\text{while}\quad X_t\to0\ \text{a.s.},\ \text{so }\Prob(X_t<c)\to1\ \ \forall c>0 .
\end{equation}
The expected price grows exponentially and yet every path goes to zero. The resolution is that the expectation of a lognormal is dominated by rare, extraordinarily large values. The typical (median) path grows like $x_0e^{\mu_*t}$, which is decreasing here. A numerical illustration ($\mu=5\%$, $\sigma=40\%$, so $\mu_*=-3\%$ per year; own computation):
\begin{center}
\small
\begin{tabular}{@{}rrrr@{}}
\toprule
$t$ (years) & $\E X_t/x_0=e^{\mu t}$ & median $e^{\mu_*t}$ & $\Prob(X_t<x_0)$\\
\midrule
$10$ & $1.65$ & $0.74$ & $0.59$\\
$100$ & $148$ & $0.050$ & $0.77$\\
$1000$ & $5.2\times10^{21}$ & $9\times10^{-14}$ & $0.99$\\
\bottomrule
\end{tabular}
\end{center}
The lecturer connects this with markets: some stocks do extraordinarily well and are noticed, others drift towards zero, and on average there is growth \ts{24}{25}{1:13:03}. (The critical case $\mu=\tfrac12\sigma^2$ is different again: $\ln X_t$ then has no drift and oscillates between $\pm\infty$.)
\begin{intuition}[Ensemble average versus time average]
A physicist would say that the multiplicative process is \emph{non-ergodic}: the ensemble average $\E X_t\propto e^{\mu t}$ and the growth rate of a single realisation, the time average $\lim\frac1t\ln X_t=\mu-\tfrac12\sigma^2$, differ. The gap $\tfrac12\sigma^2$ is the volatility drag, the It\^o correction of \cref{ch:stochcalc}: an investor who holds the asset for a long time experiences $\mu_*$, not $\mu$. (The same quantity, $\mu-\tfrac12\sigma^2$, is the growth rate the Kelly criterion maximises; not discussed in class.)
\end{intuition}
\begin{coursenote}[Erratum: the standard deviation of $\ln X_t$]
Slide 15 writes the lognormal parameters as $\mu_*=(\mu-\sigma^2/2)t$ and $\sigma_*=\sigma t$, and the lecturer says aloud ``our standard deviation is $\sigma t$'' \ts{24}{25}{1:12:02}. The variance of $\sigma B_t$ is $\sigma^2t$, so the standard deviation is $\sigma_*=\sigma\sqrt t$ (equivalently the variance is $\sigma^2t$, as in the transcript at 7:23). The slide's claim about $\E X_t=x_0e^{\mu t}$ and the paradox are correct. The paradox needs $0<\mu<\sigma^2/2$: for $\mu\le0$ the mean does not grow, and there is nothing surprising in $X_t\to0$.
\end{coursenote}

% ==================================================================
\section{Linear SDEs: the Ornstein--Uhlenbeck and Vasicek processes}\label{ch18:sec-ou}

\subsection{The model}
Geometric Brownian motion has no memory of a level: a price that has doubled has no reason to come back. Many quantities do have a preferred level (an interest rate, a credit spread, the velocity of a molecule). Replace the constant drift by a restoring force proportional to the distance from that level:

\begin{definition}[Ornstein--Uhlenbeck process]\label{ch18:def-ou}
For $\kappa>0$ (speed of mean reversion), $m\in\R$ (long-run level) and $\sigma>0$ the \emph{Ornstein--Uhlenbeck} (OU) process solves
\begin{equation}\label{ch18:eq-ou}
  \dd X_t=-\kappa\,(X_t-m)\,\dd t+\sigma\,\dd B_t,\qquad X_0=x_0 .
\end{equation}
\end{definition}
\ts{24}{25}{1:14:09}\ts{24}{25}{1:15:12}\ts{13}{21}{13:16}\ts{13}{21}{14:18}
The drift $-\kappa(X-m)$ is negative when $X>m$ and positive when $X<m$ (slide 16): a force that pulls $X$ back to $m$, proportional to the distance; the local variability $\sigma$ is constant, unlike geometric Brownian motion (\cref{ch18:ex-gbm}), where the noise scales with the level \ts{13}{21}{13:16}. The two lectures use different letters:
\begin{center}
\small
\begin{tabular}{@{}l c c c c c@{}}
\toprule
 & speed & level & noise & source\\
\midrule
these notes, ch.~\ref{ch:timeseries} & $\kappa$ & $m$ ($\mu$) & $\sigma$ ($\sigma_0$) & \\
2024 slides 16--18 & $\alpha$ & $\mu$ & $\sigma$ & \\
2013 note and lecture (level $0$) & $\alpha$ & $0$ & $\sigma$ & \\
2013 PS9 A-3, $\dd R=(\alpha-\beta R)\dd t+\sigma\dd B$ & $\beta$ & $\alpha/\beta$ & $\sigma$ & \\
\bottomrule
\end{tabular}
\end{center}
(The symbol $\mu$ is kept for the \emph{drift} of geometric Brownian motion, so the level of the OU process is called $m$ here; \cref{ch:timeseries} writes $\mu$ for it and $\sigma_0$ for the noise.)

\begin{history}[From Brownian particles to interest rates]
Uhlenbeck and Ornstein (1930) used this equation for the velocity $v$ of a Brownian particle: Langevin's equation $m_pv'=-\gamma v+\sqrt{2\gamma k_BT}\,\xi$ has the form \eqref{ch18:eq-ou} with $m=0$, $\kappa=\gamma/m_p$, $\sigma^2=2\gamma k_BT/m_p^2$; the stationary variance $\sigma^2/(2\kappa)=k_BT/m_p$ is the equipartition theorem. The lecture says the model was first used for the velocity of gas molecules whose average is $m$ \ts{24}{25}{1:14:09}\ts{13}{21}{14:18}. Vasicek (1977) used it for the short interest rate, with $m$ the target or equilibrium rate (for instance the rate targeted by the central bank) \ts{24}{25}{1:15:12}.

\smallskip
{\footnotesize\textbf{References.} P.~Langevin, \emph{C.~R.\ Acad.\ Sci.\ Paris} 146, 530 (1908);
G.~E.~Uhlenbeck and L.~S.~Ornstein, ``On the theory of the Brownian motion'', \emph{Phys.\ Rev.}
36, 823 (1930); O.~Vasicek, ``An equilibrium characterization of the term structure'',
\emph{J.\ Financial Economics} 5, 177--188 (1977).}
\end{history}

\subsection{Solution by the product-process guess}
Take first $m=0$, $\dd X=-\kappa X\,\dd t+\sigma\,\dd B$. Coefficient matching with $f(t,B_t)$ \emph{fails} here \ts{13}{21}{13:16}: \eqref{ch18:eq-match} would require $f_x=\sigma$, so $f=\sigma x+g(t)$, and then $\mu(t,f)=-\kappa(\sigma x+g)$ has to equal $g'$, which depends on $x$. The lectures change the ansatz to a \emph{product process}
\begin{equation}\label{ch18:eq-product}
  X_t=a(t)\Bigl(x_0+\int_0^tb(s)\,\dd B_s\Bigr),\qquad a(0)=1,
\end{equation}
with $a,b$ differentiable deterministic functions \ts{24}{25}{1:15:12}\ts{24}{25}{1:16:17}\ts{13}{21}{14:18}\ts{13}{21}{15:34}. (The 2013 lecturer admits that he cannot motivate the guess beyond experience \ts{13}{21}{15:34}; the reason it works is explained below.) Differentiate with the product rule (\cref{ch17:cor-product}):
\[
  \dd X_t=a'(t)\Bigl(x_0+\int_0^tb\,\dd B\Bigr)\dd t+a(t)\,b(t)\,\dd B_t=\frac{a'(t)}{a(t)}\,X_t\,\dd t+a(t)b(t)\,\dd B_t .
\]
There is \emph{no extra term}: $a$ is deterministic and of finite variation, so $\dd a\,\dd(\int b\,\dd B)=a'\,\dd t\,\dd B=0$; the 2013 lecturer makes this subtle point explicit, noting that once a process is written as an integral, differentiating it returns the integrand and nothing more \ts{13}{21}{16:37}\ts{13}{21}{17:42}. Matching with $\dd X=-\kappa X\,\dd t+\sigma\,\dd B$ gives
\[
  \frac{a'(t)}{a(t)}=-\kappa,\quad a(t)b(t)=\sigma\ \Longrightarrow\ a(t)=e^{-\kappa t},\quad b(t)=\sigma e^{\kappa t}
\]
\ts{24}{25}{1:17:17}\ts{13}{21}{18:42}. For a general level $m$ apply this to $Y=X-m$ (which has the same differential and $Y_0=x_0-m$).

\begin{proposition}[Solution of the OU equation]\label{ch18:prop-ou}
The solution of \eqref{ch18:eq-ou} is
\begin{equation}\label{ch18:eq-ousol}
  X_t=m+(x_0-m)\,e^{-\kappa t}+\sigma\int_0^te^{-\kappa(t-s)}\,\dd B_s .
\end{equation}
\end{proposition}
The \emph{integrating factor} makes the same computation transparent (own remark). Since $e^{\kappa t}$ is deterministic, the product rule gives $\dd\bigl(e^{\kappa t}(X_t-m)\bigr)=e^{\kappa t}\bigl(\dd X_t+\kappa(X_t-m)\dd t\bigr)=\sigma e^{\kappa t}\dd B_t$, and integrating from $0$ to $t$ yields \eqref{ch18:eq-ousol}. This is the stochastic version of the method for a linear ODE $\dot x=-\kappa x+f(t)$, and it also solves the linear SDE with time-dependent coefficients $\dd X=(-\kappa(t)X+f(t))\dd t+\sigma(t)\dd B$ (the Hull--White extension) by replacing $e^{\kappa t}$ with $e^{\int_0^t\kappa}$. That this is a solution can also be checked directly; verifying the formula for the Vasicek form of the model is PS9 Problem~A-3 (\cref{ch18:ex-a3}).

\subsection{Properties}
The stochastic integral in \eqref{ch18:eq-ousol} has a deterministic integrand, so it is a Gaussian random variable with mean $0$ and, by the It\^o isometry, variance $\sigma^2\int_0^te^{-2\kappa(t-s)}\dd s$ (\cref{ch:stochcalc}). This gives:

\begin{proposition}[Moments, law, stationarity]\label{ch18:prop-ouprop}
Let $X$ solve \eqref{ch18:eq-ou}. Then $X$ is a Gaussian process and
\begin{align}
  \E X_t&=m+(x_0-m)e^{-\kappa t}\ \longrightarrow\ m, \label{ch18:eq-oumean}\\
  \Var X_t&=\sigma^2\int_0^te^{-2\kappa(t-s)}\,\dd s=\frac{\sigma^2}{2\kappa}\bigl(1-e^{-2\kappa t}\bigr)\ \longrightarrow\ \frac{\sigma^2}{2\kappa}, \label{ch18:eq-ouvar}\\
  \Cov(X_s,X_t)&=\frac{\sigma^2}{2\kappa}\Bigl(e^{-\kappa(t-s)}-e^{-\kappa(t+s)}\Bigr),\qquad 0\le s\le t. \label{ch18:eq-oucov}
\end{align}
As $t\to\infty$, $X_t\to N\bigl(m,\sigma^2/(2\kappa)\bigr)$. If $X_0\sim N(m,\sigma^2/(2\kappa))$ is independent of $B$, then $X_t\sim N(m,\sigma^2/(2\kappa))$ for all $t$: the process is strictly stationary, with autocorrelation $\rho(\tau)=e^{-\kappa|\tau|}$.
\end{proposition}
\begin{proof}
\eqref{ch18:eq-oumean} and \eqref{ch18:eq-ouvar} were shown above (\ts{24}{25}{1:17:17}, slide 18). For the covariance, $s\le t$: $\Cov(X_s,X_t)=\sigma^2e^{-\kappa(t+s)}\E\bigl[\int_0^se^{\kappa u}\dd B_u\int_0^te^{\kappa u}\dd B_u\bigr]=\sigma^2e^{-\kappa(t+s)}\int_0^se^{2\kappa u}\dd u$, since the part of the second integral beyond $s$ is independent of the first, and this is \eqref{ch18:eq-oucov}. If $X_0$ is independent with variance $\sigma^2/(2\kappa)$, then $\Var X_t=e^{-2\kappa t}\sigma^2/(2\kappa)+\sigma^2(1-e^{-2\kappa t})/(2\kappa)=\sigma^2/(2\kappa)$ and the mean stays $m$; the covariance becomes $\frac{\sigma^2}{2\kappa}e^{-\kappa(t-s)}$, a function of $t-s$; a Gaussian process with stationary mean and covariance is strictly stationary.
\end{proof}
\ts{24}{25}{1:17:17}\ts{13}{21}{19:56}
This is the point of slide 18 and of the transcript: unlike Brownian motion, whose variance grows like $t$, the OU process has a \emph{limiting distribution} with constant mean and variance; it is ergodic (time averages converge to the averages under the stationary law), so the long-run distribution does not depend on the starting point \ts{24}{25}{1:17:17}. The mean decays exponentially to $m$ with relaxation time $1/\kappa$; the \emph{half-life} of a deviation is $\ln2/\kappa$.

\begin{figure}[ht]
\centering
\begin{tikzpicture}
\begin{axis}[width=0.95\linewidth,height=5.4cm,xlabel={$t$},xmin=0,xmax=3,ymin=-2.2,ymax=2.9,
  legend style={at={(0.5,1.03)},anchor=south,legend columns=3,font=\footnotesize},grid=major,grid style={black!10}]
  \addplot[very thick,color=defcol,domain=0:3,samples=60] {2*exp(-2*x)};
  \addplot[dashed,thick,color=thmcol,domain=0:3,samples=60] {2*exp(-2*x)+2*sqrt((1-exp(-4*x))/4)};
  \addplot[dashed,thick,color=thmcol,domain=0:3,samples=60,forget plot] {2*exp(-2*x)-2*sqrt((1-exp(-4*x))/4)};
  \addplot[thin,color=black!55,forget plot] coordinates {(0,2.0) (0.05,1.817) (0.1,1.934) (0.15,2.01) (0.2,1.71) (0.25,1.484) (0.3,1.231) (0.35,1.235) (0.4,1.105) (0.45,1.159) (0.5,0.656) (0.55,0.927) (0.6,0.818) (0.65,0.885) (0.7,0.772) (0.75,0.618) (0.8,0.657) (0.85,0.77) (0.9,0.654) (0.95,0.559) (1.0,0.652) (1.05,0.405) (1.1,0.044) (1.15,0.124) (1.2,-0.031) (1.25,-0.437) (1.3,-0.568) (1.35,-0.614) (1.4,-0.809) (1.45,-1.05) (1.5,-0.942) (1.55,-0.662) (1.6,-0.648) (1.65,-0.745) (1.7,-0.592) (1.75,-0.383) (1.8,-0.411) (1.85,-0.256) (1.9,-0.009) (1.95,-0.052) (2.0,-0.221) (2.05,-0.126) (2.1,-0.061) (2.15,0.179) (2.2,-0.112) (2.25,-0.242) (2.3,-0.397) (2.35,-0.729) (2.4,-0.632) (2.45,-0.46) (2.5,-0.573) (2.55,-0.224) (2.6,-0.028) (2.65,0.109) (2.7,0.184) (2.75,0.37) (2.8,0.051) (2.85,0.177) (2.9,0.288) (2.95,-0.115) (3.0,-0.031)};
  \addplot[thin,color=black!55,forget plot] coordinates {(0,2.0) (0.05,1.756) (0.1,1.756) (0.15,1.495) (0.2,1.349) (0.25,1.294) (0.3,0.984) (0.35,1.018) (0.4,0.899) (0.45,0.918) (0.5,0.719) (0.55,0.882) (0.6,0.927) (0.65,0.801) (0.7,0.859) (0.75,1.046) (0.8,1.327) (0.85,0.866) (0.9,0.972) (0.95,0.978) (1.0,0.865) (1.05,0.569) (1.1,0.782) (1.15,0.439) (1.2,0.518) (1.25,0.746) (1.3,0.334) (1.35,0.238) (1.4,-0.063) (1.45,-0.005) (1.5,0.318) (1.55,0.718) (1.6,0.271) (1.65,0.123) (1.7,0.261) (1.75,0.572) (1.8,0.608) (1.85,0.391) (1.9,0.417) (1.95,0.374) (2.0,0.295) (2.05,0.11) (2.1,0.182) (2.15,0.231) (2.2,0.189) (2.25,0.124) (2.3,-0.162) (2.35,-0.25) (2.4,0.031) (2.45,-0.013) (2.5,-0.318) (2.55,-0.004) (2.6,0.11) (2.65,0.548) (2.7,0.509) (2.75,0.362) (2.8,0.02) (2.85,0.3) (2.9,0.818) (2.95,0.566) (3.0,0.374)};
  \addplot[thin,color=black!55,forget plot] coordinates {(0,2.0) (0.05,1.937) (0.1,1.576) (0.15,1.368) (0.2,1.163) (0.25,1.094) (0.3,1.223) (0.35,1.111) (0.4,1.201) (0.45,0.997) (0.5,0.972) (0.55,0.424) (0.6,0.075) (0.65,0.238) (0.7,0.089) (0.75,0.204) (0.8,0.3) (0.85,0.553) (0.9,0.673) (0.95,0.826) (1.0,0.723) (1.05,0.506) (1.1,0.302) (1.15,0.169) (1.2,-0.087) (1.25,-0.195) (1.3,-0.197) (1.35,-0.124) (1.4,-0.185) (1.45,-0.577) (1.5,-0.537) (1.55,-0.438) (1.6,-0.165) (1.65,-0.027) (1.7,-0.161) (1.75,-0.3) (1.8,0.157) (1.85,0.303) (1.9,0.664) (1.95,1.054) (2.0,0.78) (2.05,0.787) (2.1,0.81) (2.15,0.852) (2.2,0.886) (2.25,0.845) (2.3,0.802) (2.35,0.406) (2.4,0.332) (2.45,0.141) (2.5,0.154) (2.55,0.04) (2.6,0.168) (2.65,0.326) (2.7,0.361) (2.75,0.394) (2.8,0.376) (2.85,0.245) (2.9,0.187) (2.95,0.064) (3.0,0.14)};
  \addplot[thin,color=black!55,forget plot] coordinates {(0,2.0) (0.05,1.813) (0.1,1.764) (0.15,1.313) (0.2,1.377) (0.25,1.084) (0.3,0.824) (0.35,0.704) (0.4,0.517) (0.45,0.53) (0.5,0.357) (0.55,0.095) (0.6,-0.094) (0.65,0.195) (0.7,0.185) (0.75,-0.081) (0.8,-0.071) (0.85,-0.396) (0.9,0.023) (0.95,-0.304) (1.0,-0.173) (1.05,-0.041) (1.1,-0.346) (1.15,-0.249) (1.2,-0.013) (1.25,0.088) (1.3,0.135) (1.35,0.324) (1.4,0.328) (1.45,0.368) (1.5,0.306) (1.55,0.411) (1.6,0.172) (1.65,0.324) (1.7,0.186) (1.75,-0.064) (1.8,0.02) (1.85,0.378) (1.9,0.547) (1.95,0.385) (2.0,0.497) (2.05,0.353) (2.1,0.293) (2.15,0.284) (2.2,-0.236) (2.25,-0.173) (2.3,-0.376) (2.35,-0.488) (2.4,-0.756) (2.45,-1.007) (2.5,-1.112) (2.55,-0.83) (2.6,-0.396) (2.65,-0.364) (2.7,-0.097) (2.75,-0.144) (2.8,-0.537) (2.85,-0.454) (2.9,-0.316) (2.95,-0.377) (3.0,-0.277)};
  \legend{{mean $x_0e^{-\kappa t}$},{mean $\pm2$ s.d.}}
\end{axis}
\end{tikzpicture}
\caption{Four exact simulated paths of the OU process ($\kappa=2$, $m=0$, $\sigma=1$, $x_0=2$; seed~11) with the mean and the band mean $\pm2$ standard deviations from \eqref{ch18:eq-oumean}--\eqref{ch18:eq-ouvar}. The paths are pulled to $m=0$ within a time of order $1/\kappa=0.5$ and then fluctuate with standard deviation $\sigma/\sqrt{2\kappa}=0.5$ around it; compare the growing fan of Brownian paths.}\label{ch18:fig-ou}
\end{figure}

\begin{finance}[Vasicek's model and its limits]\label{ch18:fin-vasicek}
In Vasicek's model the short rate $R$ satisfies $\dd R=(\alpha-\beta R)\dd t+\sigma\dd B$: this is \eqref{ch18:eq-ou} with $\kappa=\beta$, $m=\alpha/\beta$, and \eqref{ch18:eq-ousol} is exactly the formula of PS9 Problem~A-3. Its virtues are tractability (Gaussian rates, closed-form bond prices, \cref{ch18:sec-fk}) and an economic story (rates return to an equilibrium level). Its defects: rates can become negative (the stationary law is normal; with $\sigma/\sqrt{2\kappa}=1\%$ around $m=4\%$ the probability is $\Phi(-4)\approx3\times10^{-5}$, but not zero), the volatility does not depend on the level, and a single factor cannot fit an arbitrary yield curve; the Hull--White extension makes $m$ a function of time to fit the curve, and \cref{ch18:sec-cir} gives a positive alternative (CIR). Historical mean reversion does not survive over decades: the level to which rates revert itself moves (\cref{ch:timeseries}).
\end{finance}

\subsection{Sampling an OU process: the exact AR(1)}\label{ch18:sec-ar1}
\Cref{ch:timeseries} claims that the AR(1) is the discrete-time version of the OU process, with autoregressive coefficient $\phi=e^{-\kappa h}$ and noise variance $\sigma^2(1-e^{-2\kappa h})/(2\kappa)$. Here is the derivation (an addition of these notes).

\begin{proposition}[Exact discretisation]\label{ch18:prop-ar1}
Let $X$ solve \eqref{ch18:eq-ou} and $h>0$. For every $t\ge0$,
\begin{equation}\label{ch18:eq-ar1}
  X_{t+h}=m+\phi\,(X_t-m)+\varepsilon_{t+h},\qquad\phi=e^{-\kappa h},\qquad\varepsilon_{t+h}=\sigma\int_t^{t+h}e^{-\kappa(t+h-s)}\dd B_s,
\end{equation}
where $\varepsilon_{t+h}\sim N\bigl(0,\sigma_\varepsilon^2\bigr)$ with
\begin{equation}\label{ch18:eq-ar1var}
  \sigma_\varepsilon^2=\sigma^2\frac{1-e^{-2\kappa h}}{2\kappa}=\frac{\sigma^2}{2\kappa}\bigl(1-\phi^2\bigr),
\end{equation}
independent of $\mathcal F_t$. Hence $X_{ih}$, $i=0,1,2,\dots$, is \emph{exactly} the Gaussian AR(1) process $X_{i+1}-m=\phi(X_i-m)+\varepsilon_{i+1}$ with i.i.d.\ noise; its stationary variance is $\sigma_\varepsilon^2/(1-\phi^2)=\sigma^2/(2\kappa)$ and its autocorrelation $\phi^j=e^{-\kappa jh}$ is the sampled continuous autocorrelation.
\end{proposition}
\begin{proof}
Apply the integrating factor on $[t,t+h]$ instead of $[0,t]$: $\dd(e^{\kappa u}(X_u-m))=\sigma e^{\kappa u}\dd B_u$, so $e^{\kappa(t+h)}(X_{t+h}-m)-e^{\kappa t}(X_t-m)=\sigma\int_t^{t+h}e^{\kappa u}\dd B_u$; multiply by $e^{-\kappa(t+h)}$ to get \eqref{ch18:eq-ar1}. The noise is a Gaussian Wiener integral of a deterministic function over $[t,t+h]$, so it is independent of $\mathcal F_t$ and normal with variance $\sigma^2\int_t^{t+h}e^{-2\kappa(t+h-s)}\dd s=\sigma^2(1-e^{-2\kappa h})/(2\kappa)$ by isometry. Since $1-\phi^2=1-e^{-2\kappa h}$, the stationary variance of the AR(1) is $\sigma_\varepsilon^2/(1-\phi^2)=\sigma^2/(2\kappa)$, the value of \eqref{ch18:eq-ouvar} at $t=\infty$, \emph{independent of $h$}, as it must be.
\end{proof}
The intercept of the AR(1) in the usual form $X_{i+1}=c+\phi X_i+\varepsilon$ is $c=m(1-\phi)$, and the half-life of a shock, $h\ln2/(-\ln\phi)=\ln2/\kappa$, is the one quoted in \cref{ch09:fin-vasicek}. Three practical consequences:
\begin{itemize}
  \item \emph{Estimation.} Fit an AR(1) to data sampled every $h$ (\cref{ch:timeseries}) and read off the continuous-time parameters: $\hat\kappa=-\ln\hat\phi/h$, $\hat m=\hat c/(1-\hat\phi)$, $\hat\sigma^2=2\hat\kappa\,\hat\sigma_\varepsilon^2/(1-\hat\phi^2)$. This needs $0<\hat\phi<1$: an AR(1) with $\phi\le0$ is not the sampling of any OU process (it would oscillate around the mean). Example: monthly data ($h=1/12$) with $\hat\phi=0.98$ gives $\hat\kappa=12\cdot0.0202=0.24$ per year, half-life $2.9$ years.
  \item \emph{The exact scheme is preferable to Euler.} The Euler--Maruyama step $X_{i+1}=X_i-\kappa(X_i-m)h+\sigma\sqrt hZ$ is also an AR(1), but with $\phi_{\mathrm{E}}=1-\kappa h$ and $\sigma_\varepsilon^2=\sigma^2h$: both are the first-order Taylor expansions of \eqref{ch18:eq-ar1}--\eqref{ch18:eq-ar1var}, and its stationary variance is $\sigma^2/\bigl(\kappa(2-\kappa h)\bigr)$, not $\sigma^2/(2\kappa)$; the scheme is not even stable when $\kappa h>2$. With $\kappa=2$, $h=0.25$ (own computation): exact $\phi=0.607$, noise variance $0.158\,\sigma^2$, stationary variance $0.25\,\sigma^2$; Euler $\phi=0.5$, noise variance $0.25\,\sigma^2$, stationary variance $0.333\,\sigma^2$, a third too large.
  \item \emph{Simulation.} Paths of the OU process on any grid are generated exactly (no discretisation error) by iterating \eqref{ch18:eq-ar1}; this was done for \cref{ch18:fig-ou}. The same holds for geometric Brownian motion, $X_{t+h}=X_t\exp((\mu-\sigma^2/2)h+\sigma\sqrt hZ)$; for most SDEs there is no such scheme (\cref{ch18:sec-num}).
\end{itemize}

% ==================================================================
\section{Systems of SDEs; position and velocity}\label{ch18:sec-sys}

Several quantities driven by several noises are described by a \emph{vector SDE} \ts{24}{25}{1:18:22}\ts{24}{25}{1:19:23}
\begin{equation}\label{ch18:eq-vec}
  \dd\bm X_t=\bm\mu(t,\bm X_t)\,\dd t+\sigma(t,\bm X_t)\,\dd\bm B_t,\qquad\bm X_0=\bm x_0,
\end{equation}
where $\bm X\in\R^m$, $\bm\mu$ is a vector of drifts $\mu_1,\dots,\mu_m$, $\bm B=(B^{(1)},\dots,B^{(m)})$ are \emph{independent} Brownian motions and $\sigma=[\sigma_{ij}]$ is an $m\times m$ matrix (slide 19). The noises in the components are correlated in general: the increments over $\dd t$ have covariance matrix $\sigma\sigma\T\,\dd t$, which is what the lecturer means by ``a covariance matrix reflecting how increments of the components are related'' \ts{24}{25}{1:18:22}. Existence and uniqueness hold under the same Lipschitz and growth conditions with $|\cdot|$ read as Euclidean and matrix norms (quoted).

\subsection{Linear systems}
When $\bm\mu=-A(\bm x-\bm m)$ and $\sigma=\Sigma$ are constant, the integrating factor $e^{At}$ solves the system exactly as in the scalar case (own extension):
\begin{equation}\label{ch18:eq-vecou}
  \bm X_t=\bm m+e^{-At}(\bm x_0-\bm m)+\int_0^te^{-A(t-s)}\Sigma\,\dd\bm B_s,\qquad
  \Cov(\bm X_t)=\int_0^te^{-Au}\Sigma\Sigma\T e^{-A\T u}\dd u .
\end{equation}
The process is stationary in the limit iff all eigenvalues of $A$ have positive real part; the limiting covariance $C_\infty$ solves the Lyapunov equation $AC_\infty+C_\infty A\T=\Sigma\Sigma\T$ (scalar case: $C_\infty=\sigma^2/(2\kappa)$). Sampled every $h$ the process is \emph{exactly} the VAR(1) of \cref{ch:timeseries} (\cref{ch09:sec-var}), $\bm X_{i+1}-\bm m=\Phi(\bm X_i-\bm m)+\bm\varepsilon_{i+1}$ with $\Phi=e^{-Ah}$ and $\Cov(\bm\varepsilon)=\int_0^he^{-Au}\Sigma\Sigma\T e^{-A\T u}\dd u$; the stationarity condition of the VAR (eigenvalues of $\Phi$ inside the unit disc) is precisely $\operatorname{Re}\lambda(A)>0$.

\subsection{The example of the slides: position and velocity}
Slide 20 puts a mean-reverting velocity $V$ and the position $X$ it moves in one system \ts{24}{25}{1:19:23}:
\begin{equation}\label{ch18:eq-posvel}
  \dd X_t=V_t\,\dd t+\sigma_1\,\dd B^{(1)}_t,\qquad \dd V_t=-\kappa(V_t-m)\,\dd t+\sigma_2\,\dd B^{(2)}_t ,
\end{equation}
that is \eqref{ch18:eq-vec} with $\bm\mu=(v,-\kappa(v-m))$ and $\sigma=\operatorname{diag}(\sigma_1,\sigma_2)$. The velocity is an OU process, and the position is its running integral plus independent noise; the position itself is \emph{not} mean-reverting ($A=\begin{psmallmatrix}0&-1\\0&\kappa\end{psmallmatrix}$ has a zero eigenvalue). The slide only sets up the system; we compute its most instructive statistic (own computation). Take $m=0$, $\sigma_1=0$, $X_0=0$ and $V_0\sim N(0,\sigma^2/(2\kappa))$ stationary, $\sigma=\sigma_2$. Then $X_t=\int_0^tV_s\dd s$ and
\begin{equation}\label{ch18:eq-posvar}
  \Var X_t=2\int_0^t(t-s)\,C(s)\,\dd s=\frac{\sigma^2}{\kappa^2}\Bigl(t-\frac{1-e^{-\kappa t}}{\kappa}\Bigr),\qquad C(s)=\frac{\sigma^2}{2\kappa}e^{-\kappa s}.
\end{equation}
Two regimes: for $\kappa t\ll1$, $\Var X_t\approx\frac{\sigma^2}{2\kappa}t^2=\langle v^2\rangle t^2$ (\emph{ballistic} motion: the velocity has no time to change); for $\kappa t\gg1$, $\Var X_t\approx\frac{\sigma^2}{\kappa^2}t=2Dt$ with the diffusion coefficient $D=\int_0^\infty C(s)\dd s=\sigma^2/(2\kappa^2)$ (\emph{diffusive} motion, Einstein's law; this is the Green--Kubo formula). A Monte Carlo check with $\kappa=\sigma=1$, $t=3$ over $2\times10^5$ paths gives $2.046$ against the exact $2.050$ of \eqref{ch18:eq-posvar}. Independent position noise $\sigma_1$ simply adds $\sigma_1^2t$.
\begin{figure}[ht]
\centering
\begin{tikzpicture}
\begin{loglogaxis}[width=0.85\linewidth,height=5.2cm,xlabel={$t$},ylabel={$\Var X_t$},xmin=0.01,xmax=100,ymin=0.00003,ymax=200,
  legend style={at={(0.5,1.03)},anchor=south,legend columns=3,font=\footnotesize},grid=major,grid style={black!10}]
  \addplot[very thick,color=defcol,domain=0.01:100,samples=60] {x-(1-exp(-x))};
  \addplot[dashed,thick,color=thmcol,domain=0.01:10,samples=20] {x*x/2};
  \addplot[dashed,thick,color=intcol,domain=0.3:100,samples=20] {x};
  \legend{{exact \eqref{ch18:eq-posvar}},{ballistic $t^2/2$},{diffusive $t$}}
\end{loglogaxis}
\end{tikzpicture}
\caption{Variance of the position, \eqref{ch18:eq-posvar} with $\kappa=\sigma=1$: slope $2$ at short times, slope $1$ at long times, crossover at $t\sim1/\kappa$.}\label{ch18:fig-posvar}
\end{figure}
\begin{intuition}[How Brownian motion arises]
Let $\kappa\to\infty$ with $\sigma^2/\kappa^2=2D$ fixed: the velocity forgets its past infinitely fast and $X_t\to\sqrt{2D}\,B_t$. Brownian motion is the overdamped limit of the position of a particle with a mean-reverting velocity, the description of \cref{ch11:int-physics}. In finance the same structure appears when a factor (say a trend, or a stochastic mean) drives an asset whose own noise is small.
\end{intuition}
\begin{finance}[Mean-reverting spreads; not discussed in class]
A relative-value (``pairs'') trade models a spread $S_t=\ln P^1_t-\beta\ln P^2_t$ between two related assets as an OU process. Estimating the AR(1) coefficient of the sampled spread gives $\kappa$ and the half-life of \cref{ch18:sec-ar1}; the stationary standard deviation $\sigma/\sqrt{2\kappa}$ sets the scale of the signal (trade when $|S-m|$ exceeds a multiple of it). If $\hat\phi$ is close to $1$ the spread is indistinguishable from a random walk (unit root, \cref{ch:timeseries}) and the strategy has no edge: the OU model is only as good as the mean reversion is real.
\end{finance}

% ==================================================================
\section{Generators, Feynman--Kac and the Kolmogorov equations}\label{ch18:sec-fk}

This section is an addition of these notes: the lectures use the two facts below without proving them (2024 slides 2--3 derive the PDE by hedging \ts{24}{25}{24:16}; both years say that the Black--Scholes PDE is the heat equation, whose solution is the payoff integrated against a Gaussian \ts{24}{25}{32:40}\ts{13}{21}{54:18}, but neither proves the equivalence with an expectation), and \cref{ch:stochcalc,ch:bs} defer them to here.

\subsection{The generator}
For the diffusion \eqref{ch18:eq-sde} define the second-order differential operator
\begin{equation}\label{ch18:eq-gen}
  \mathcal L_tf(x)=\mu(t,x)\,f'(x)+\tfrac12\sigma^2(t,x)\,f''(x).
\end{equation}
It is the \emph{generator} of $X$ (the expansion \eqref{ch11:eq-generator} in \cref{ch:sp2} is the same statement at the level of one small step). It\^o's formula \eqref{ch17:eq-ito3} for $f(t,X_t)$ reads, in this notation,
\begin{equation}\label{ch18:eq-dynkin}
  f(t,X_t)=f(0,x_0)+\int_0^t\bigl(\partial_s+\mathcal L_s\bigr)f(s,X_s)\,\dd s+\int_0^t\sigma\,\partial_xf(s,X_s)\,\dd B_s .
\end{equation}
The last term is a martingale (under an integrability condition, \cref{ch17:thm-ito-int}); so $f(t,X_t)-\int_0^t(\partial_s+\mathcal L_s)f\,\dd s$ is a martingale. This is Dynkin's formula, and \cref{ch17:thm-pde} is its case $\mathcal L=\tfrac12\partial_x^2$: $f(t,B_t)$ is a martingale when $f_t+\tfrac12f_{xx}=0$. The generators of our examples:
\[
  \text{BM: }\tfrac12\partial_x^2;\qquad\text{GBM: }\mu x\,\partial_x+\tfrac12\sigma^2x^2\partial_x^2;\qquad\text{OU: }-\kappa(x-m)\,\partial_x+\tfrac12\sigma^2\partial_x^2 .
\]

\subsection{The Feynman--Kac formula}
\begin{theorem}[Feynman--Kac]\label{ch18:thm-fk}
Let $u\in C^{1,2}\bigl([0,T)\times\R\bigr)\cap C\bigl([0,T]\times\R\bigr)$ solve the terminal value problem
\begin{equation}\label{ch18:eq-fkpde}
  \frac{\partial u}{\partial t}+\mu(t,x)\frac{\partial u}{\partial x}+\tfrac12\sigma^2(t,x)\frac{\partial^2u}{\partial x^2}-r(t,x)\,u=0,\qquad u(T,x)=g(x),
\end{equation}
with $r$ continuous and bounded below, and assume $\E\int_t^T\bigl(e^{-\int_t^sr}\,\sigma\,u_x(s,X_s)\bigr)^2\dd s<\infty$. Then
\begin{equation}\label{ch18:eq-fk}
  u(t,x)=\E\Bigl[\exp\Bigl(-\int_t^Tr(s,X_s)\,\dd s\Bigr)\,g(X_T)\ \Bigm|\ X_t=x\Bigr],
\end{equation}
where $X$ solves \eqref{ch18:eq-sde} started at $x$ at time $t$.
\end{theorem}
\begin{proof}
Fix $(t,x)$ and let $D_s=\exp\bigl(-\int_t^sr(v,X_v)\dd v\bigr)$, a process of finite variation with $\dd D_s=-r(s,X_s)D_s\,\dd s$. By the product rule (\cref{ch17:cor-product}; $\dd D\,\dd X=0$) and It\^o's formula \eqref{ch18:eq-dynkin},
\[
\begin{gathered}
  \dd\bigl(D_su(s,X_s)\bigr)=D_s\Bigl[u_s+\mu u_x+\tfrac12\sigma^2u_{xx}-r\,u\Bigr](s,X_s)\,\dd s+D_s\,\sigma u_x(s,X_s)\,\dd B_s\\
  =D_s\,\sigma u_x(s,X_s)\,\dd B_s ,
\end{gathered}
\]
the bracket vanishing by \eqref{ch18:eq-fkpde}. Under the integrability assumption $D_su(s,X_s)$ is a martingale on $[t,T]$, hence $u(t,x)=\E[D_tu(t,X_t)]=\E[D_Tu(T,X_T)\mid X_t=x]=\E[D_Tg(X_T)\mid X_t=x]$.
\end{proof}
Read the formula in both directions: a PDE solved \emph{backwards} from the terminal condition $g$ is a conditional expectation of the payoff, discounted along the path; conversely every such expectation solves the PDE (under regularity conditions, quoted). Special cases:
\begin{itemize}
  \item $r=0$, $\mu=0$, $\sigma=1$: $u_t+\tfrac12u_{xx}=0$, $u(t,x)=\E\,g(x+B_{T-t})$: the backward heat equation of \cref{ch17:thm-pde}. With $\tau=T-t$ it is the heat equation of \cref{ch18:sec-heat}.
  \item \emph{Black--Scholes.} Take for $X$ the price under the risk-neutral measure, $\dd S=rS\,\dd t+\sigma S\,\dd W$ (\cref{ch17:sec-girsanov}), a constant rate $r$ and $g$ the payoff. Then \eqref{ch18:eq-fkpde} is $u_t+rSu_S+\tfrac12\sigma^2S^2u_{SS}-ru=0$, the Black--Scholes equation \eqref{ch17:eq-bs} that slides 2--3 of the lecture derive by hedging \ts{24}{25}{24:16}, and \eqref{ch18:eq-fk} is the price $e^{-r(T-t)}\E_\Q[g(S_T)\mid S_t]$, the risk-neutral valuation of \cref{ch:bs} (\cref{ch15:eq-rnval}). Two proofs, one equation: the hedging argument shows that the price \emph{must} satisfy the PDE; Feynman--Kac shows that the risk-neutral expectation \emph{does}, and the PDE with its terminal condition has one solution. The real-world drift $\mu$ of $S$ disappears because the expectation is taken under $\Q$, where the drift is $r$ (Girsanov's theorem, \cref{ch17:thm-girsanov2}).
\end{itemize}

\begin{example}[Second moment of the OU process]\label{ch18:ex-fkou}
Let $X$ be the OU process \eqref{ch18:eq-ou} and $g(x)=x^2$, $r=0$, $\tau=T-t$. From \eqref{ch18:eq-ousol}--\eqref{ch18:eq-ouvar}, $u(t,x)=\E[X_T^2\mid X_t=x]=\bigl(m+(x-m)e^{-\kappa\tau}\bigr)^2+\sigma^2\dfrac{1-e^{-2\kappa\tau}}{2\kappa}$. Feynman--Kac says this solves $u_t-\kappa(x-m)u_x+\tfrac12\sigma^2u_{xx}=0$ with $u(T,x)=x^2$; substituting confirms it (checked symbolically for these notes). Similarly for geometric Brownian motion $u=x^p\exp\bigl[(p\mu+\tfrac12p(p-1)\sigma^2)\tau\bigr]$ solves the PDE with $g=x^p$.
\end{example}

\begin{finance}[The Vasicek bond price]\label{ch18:fin-bond}
A zero-coupon bond paying $1$ at $T$ has price $P(t,T)=\E_\Q\bigl[e^{-\int_t^TR_s\dd s}\mid R_t=r\bigr]$ (\cref{ch:rates}). By Feynman--Kac (with $\mu(r)=\kappa(\theta-r)$, the OU drift under $\Q$, and $r(t,x)=x$, $g=1$) it solves $P_t+\kappa(\theta-r)P_r+\tfrac12\sigma^2P_{rr}-rP=0$, $P(T,r)=1$. The affine ansatz $P=\exp\bigl[a(\tau)-b(\tau)r\bigr]$, $\tau=T-t$, gives $-a'+b'r-\kappa(\theta-r)b+\tfrac12\sigma^2b^2-r=0$ for all $r$, hence
\[
  b'=1-\kappa b,\quad a'=-\kappa\theta b+\tfrac12\sigma^2b^2,\quad a(0)=b(0)=0
\]
and, with $b=(1-e^{-\kappa\tau})/\kappa$,
\begin{equation}\label{ch18:eq-vasbond}
  P(t,T)=\exp\Bigl[-b(\tau)\,r+\bigl(b(\tau)-\tau\bigr)\Bigl(\theta-\frac{\sigma^2}{2\kappa^2}\Bigr)-\frac{\sigma^2b(\tau)^2}{4\kappa}\Bigr].
\end{equation}
The yield $-\ln P/\tau$ tends to $\theta-\sigma^2/(2\kappa^2)$ for long maturities, below the level $\theta$ by a convexity term. Example ($\Q$-parameters $\kappa=0.5$, $\theta=4\%$, $\sigma=1\%$, $r_0=3\%$, $\tau=5$): $P=0.8343$, yield $3.62\%$, long yield $3.98\%$; a Monte Carlo average of $e^{-\int R}$ over $2\times10^5$ exactly sampled paths gives $0.83429\pm0.00006$ (own computation; the formula is verified symbolically to satisfy the PDE). Under the real-world measure the drift is $\kappa(\theta_{\mathbb P}-r)$; with a constant market price of risk $\lambda$, Girsanov's theorem (\cref{ch17:thm-girsanov2}) replaces the level by $\theta=\theta_{\mathbb P}-\sigma\lambda/\kappa$ and leaves $\kappa,\sigma$ unchanged. This is a first example of the short-rate models of \cref{ch:hjm} (not discussed in class).
\end{finance}

\subsection{Transition densities: backward and forward equations}
Let $p(s,x;t,y)$ be the density of $X_t$ given $X_s=x$. Feynman--Kac with $r=0$ and $g$ arbitrary says $u(s,x)=\E[g(X_t)\mid X_s=x]=\int g(y)p(s,x;t,y)\dd y$ solves $u_s+\mathcal L_su=0$; hence $p$ satisfies the \emph{Kolmogorov backward equation} in $(s,x)$. There is also an equation in the \emph{forward} variables $(t,y)$:

\begin{proposition}[Fokker--Planck equation]\label{ch18:prop-fp}
If the density $p(t,y)$ of $X_t$ is smooth and decays fast enough,
\begin{equation}\label{ch18:eq-fp}
  \frac{\partial p}{\partial t}=-\frac{\partial}{\partial y}\bigl[\mu(t,y)\,p\bigr]+\frac12\frac{\partial^2}{\partial y^2}\bigl[\sigma^2(t,y)\,p\bigr]=:\mathcal L_t^*p .
\end{equation}
\end{proposition}
\begin{proof}
For a smooth test function $g$ with compact support, taking expectations in \eqref{ch18:eq-dynkin} (the martingale term has mean zero) gives $\frac{\dd}{\dd t}\E\,g(X_t)=\E\,\mathcal L_tg(X_t)=\int\bigl(\mu g'+\tfrac12\sigma^2g''\bigr)p\,\dd y$. Integrating by parts once in the first term and twice in the second (no boundary terms) turns this into $\int g\,\mathcal L_t^*p\,\dd y$, while the left side is $\int g\,\partial_tp\,\dd y$. Since $g$ is arbitrary, \eqref{ch18:eq-fp} follows.
\end{proof}
The Fokker--Planck equation is a conservation law, $\partial_tp+\partial_yJ=0$ with probability current $J=\mu p-\tfrac12\partial_y(\sigma^2p)$: the drift transports probability and the noise diffuses it. It is the equation a physicist derives from the Langevin equation (Kramers--Moyal expansion), and it is the equation $u_t=\lambda u_{xx}$ of the two lectures, with $\lambda=\sigma^2/2$, when $\mu=0$ and $\sigma$ is constant. The backward equation is the adjoint of the forward one: it describes how the \emph{payoff} is smoothed, the forward one how the \emph{density} spreads. Examples:
\begin{itemize}
  \item The OU transition density is the Gaussian $N\bigl(m+(x-m)e^{-\kappa\tau},\sigma^2(1-e^{-2\kappa\tau})/(2\kappa)\bigr)$ (from \eqref{ch18:eq-ar1}); it satisfies \eqref{ch18:eq-fp} in $(\tau,y)$, $\partial_\tau p=\kappa\partial_y[(y-m)p]+\tfrac12\sigma^2\partial_y^2p$, and the backward equation in $(\tau,x)$ (checked symbolically for these notes).
  \item Stationary solution: $\partial_tp=0$ and zero current, $\mu p=\tfrac12(\sigma^2p)'$; for OU, $p'/p=-2\kappa(y-m)/\sigma^2$, so $p_{\mathrm{st}}\propto e^{-\kappa(y-m)^2/\sigma^2}$, the normal law $N(m,\sigma^2/(2\kappa))$ of \cref{ch18:prop-ouprop}. In general, for a drift $-U'(y)$ and constant $\sigma$, $p_{\mathrm{st}}\propto e^{-2U(y)/\sigma^2}$: a Boltzmann distribution at effective temperature $\sigma^2/2$.
  \item Geometric Brownian motion: the lognormal density of \cref{ch:bs} (\cref{ch15:eq-density}).
\end{itemize}

% ==================================================================
\section{The heat equation}\label{ch18:sec-heat}

The 2013 lecture ends and the 2024 lecture begins its second half with the simplest partial differential equation of all, the diffusion (heat) equation \ts{13}{21}{42:14}\ts{24}{25}{32:40}\ts{24}{25}{33:46}
\begin{equation}\label{ch18:eq-heat}
  \frac{\partial u}{\partial t}=\lambda\,\frac{\partial^2u}{\partial x^2},\qquad u(x,0)=u_0(x),\quad x\in\R,\ t>0 ,
\end{equation}
with $\lambda>0$ (2024 slide 5; the 2013 note takes $\lambda=1$). The reason to study it in a finance course is that the Black--Scholes equation \eqref{ch17:eq-bs} reduces to it by a change of variables \ts{13}{21}{43:17}\ts{13}{21}{54:18}\ts{24}{25}{32:40}, made explicit in \cref{ch15:sec-heat}: the payoff is the initial condition, and the price today is the payoff diffused for a time $T-t$. Its physical model is a long, thin, perfectly insulated bar, $u(x,t)$ the temperature at position $x$ \ts{24}{25}{36:58}\ts{13}{21}{43:17}. The heat equation of \cref{ch:sp2} (\cref{ch11:prop-heat}) is the same equation for the transition density of Brownian motion, and by the Fokker--Planck equation \eqref{ch18:eq-fp} $\lambda=\sigma^2/2$.

\subsection{Superposition and the Gaussian kernel}
\begin{itemize}
  \item \emph{Linearity.} If $u_1,u_2$ solve \eqref{ch18:eq-heat}, so does $u_1+u_2$; more generally a superposition $\int_{-\infty}^\infty u_s(x,t)\,c(s)\,\dd s$ of solutions indexed by $s$ is a solution, if one may differentiate under the integral sign \ts{24}{25}{39:01}\ts{13}{21}{44:20}\ts{13}{21}{45:24}. The strategy: solve the ``easy'' initial value problems and superpose \ts{24}{25}{39:01}.
  \item \emph{The easy problem: a Dirac delta.} Put all the heat at one point, $u(x,0)=\delta(x)$ ($\delta=0$ for $x\ne0$, $\int\delta=1$; Dirac's ``function'' that a physicist knows well, and that the lecturer introduces with a note on Dirac himself \ts{24}{25}{40:05}\ts{24}{25}{42:12}\ts{13}{21}{46:28}). Its solution is the Gaussian density with mean $0$ and variance $2\lambda t$,
  \begin{equation}\label{ch18:eq-kernel}
    u_\delta(x,t)=\frac1{\sqrt{4\pi\lambda t}}\,e^{-x^2/(4\lambda t)} ,
  \end{equation}
  which tends to $\delta$ as $t\to0$ \ts{24}{25}{34:54}\ts{24}{25}{40:05}\ts{13}{21}{47:38}\ts{13}{21}{49:59}. (For $\lambda=\tfrac12$ it is the standard $N(0,t)$ density of Brownian motion, slide 7; for $\lambda=1$ it is $\frac1{2\sqrt{\pi t}}e^{-x^2/4t}$, the 2013 note's Observation~2.) Direct substitution shows that it solves \eqref{ch18:eq-heat}: $\partial_tu_\delta=\lambda\partial_x^2u_\delta$ (checked symbolically).
  \item \emph{General initial data.} Since $u_0(x)=\int\delta(x-s)\,u_0(s)\,\dd s$, superposing translated kernels gives
  \begin{equation}\label{ch18:eq-conv}
    u(x,t)=\int_{-\infty}^\infty u_\delta(x-s,t)\,u_0(s)\,\dd s=\E\Bigl[u_0\bigl(x+\sqrt{2\lambda}\,B_t\bigr)\Bigr]
  \end{equation}
  \ts{24}{25}{43:19}\ts{24}{25}{44:27}\ts{13}{21}{47:38}\ts{13}{21}{52:06}\ts{13}{21}{53:09}. For a ``reasonable'' $u_0$ (say bounded and continuous, or of moderate growth such as a call payoff, since the kernel has Gaussian tails) the derivatives $\partial_tu_\delta,\partial_x^2u_\delta$ are integrable and can be taken inside the integral, so $u$ solves the equation, and $u(x,t)\to u_0(x)$ as $t\to0$ \ts{24}{25}{45:27}\ts{24}{25}{46:29}\ts{24}{25}{47:31}\ts{13}{21}{52:06}. The 2013 remark that the integral ``need not be well defined'' for arbitrary $u_0$ is the statement that the class of admissible initial data must be restricted. The form on the right is the probabilistic reading: \emph{the solution at $(x,t)$ is the expectation of the initial datum at the random point reached by a Brownian particle started at $x$}. In the lecture's picture, heat is a cloud of particles released at one point, each moving as a Brownian motion, and their density is normal \ts{13}{21}{47:38}\ts{13}{21}{51:04}. This is Feynman--Kac (\cref{ch18:thm-fk}) for Brownian motion.
\end{itemize}
\begin{example}[Indicator of an interval]\label{ch18:ex-indicator}
Let $u_0=\ind_{(a,b)}$ \ts{24}{25}{48:39}\ts{24}{25}{49:42}. Then by \eqref{ch18:eq-conv}, with $X_t=x+\sqrt{2\lambda}B_t\sim N(x,2\lambda t)$,
\begin{equation}\label{ch18:eq-indic}
  u(x,t)=\Prob\bigl(a<X_t<b\bigr)=\Phi\Bigl(\frac{b-x}{\sqrt{2\lambda t}}\Bigr)-\Phi\Bigl(\frac{a-x}{\sqrt{2\lambda t}}\Bigr)
\end{equation}
\ts{24}{25}{50:44}, slide 9 for $\lambda=\tfrac12$. \Cref{ch18:fig-heat} shows the initial box melting into a bell curve.
\end{example}
\begin{figure}[ht]
\centering
\begin{tikzpicture}
\begin{axis}[width=0.85\linewidth,height=5.2cm,xlabel={$x$},ylabel={$u(x,t)$},xmin=-3,xmax=3,ymin=-0.03,ymax=1.08,
  legend style={at={(0.5,1.03)},anchor=south,legend columns=4,font=\footnotesize},grid=major,grid style={black!10}]
  \addplot[very thick,color=black!70,const plot,forget plot] coordinates {(-3,0) (-1,0) (-1,1) (1,1) (1,0) (3,0)};
  \addplot[thick,color=defcol,smooth] coordinates {(-3.00,0.0000) (-2.75,0.0000) (-2.50,0.0000) (-2.25,0.0000) (-2.00,0.0000) (-1.75,0.0000) (-1.50,0.0002) (-1.25,0.0385) (-1.00,0.5000) (-0.75,0.9615) (-0.50,0.9998) (-0.25,1.0000) (0.00,1.0000) (0.25,1.0000) (0.50,0.9998) (0.75,0.9615) (1.00,0.5000) (1.25,0.0385) (1.50,0.0002) (1.75,0.0000) (2.00,0.0000) (2.25,0.0000) (2.50,0.0000) (2.75,0.0000) (3.00,0.0000)};
  \addplot[thick,color=intcol,smooth] coordinates {(-3.00,0.0000) (-2.75,0.0002) (-2.50,0.0013) (-2.25,0.0062) (-2.00,0.0228) (-1.75,0.0668) (-1.50,0.1587) (-1.25,0.3085) (-1.00,0.5000) (-0.75,0.6912) (-0.50,0.8400) (-0.25,0.9270) (0.00,0.9545) (0.25,0.9270) (0.50,0.8400) (0.75,0.6912) (1.00,0.5000) (1.25,0.3085) (1.50,0.1587) (1.75,0.0668) (2.00,0.0228) (2.25,0.0062) (2.50,0.0013) (2.75,0.0002) (3.00,0.0000)};
  \addplot[thick,color=thmcol,smooth] coordinates {(-3.00,0.0227) (-2.75,0.0400) (-2.50,0.0666) (-2.25,0.1051) (-2.00,0.1573) (-1.75,0.2236) (-1.50,0.3023) (-1.25,0.3891) (-1.00,0.4772) (-0.75,0.5586) (-0.50,0.6247) (-0.25,0.6677) (0.00,0.6827) (0.25,0.6677) (0.50,0.6247) (0.75,0.5586) (1.00,0.4772) (1.25,0.3891) (1.50,0.3023) (1.75,0.2236) (2.00,0.1573) (2.25,0.1051) (2.50,0.0666) (2.75,0.0400) (3.00,0.0227)};
  \legend{{$t=0.02$},{$t=0.25$},{$t=1$}}
\end{axis}
\end{tikzpicture}
\caption{Solution \eqref{ch18:eq-indic} of the heat equation with $\lambda=\tfrac12$ for $u_0=\ind_{(-1,1)}$ (thick black): at time $t$ the profile has width of order $\sqrt t$ and its maximum falls. Each curve is the probability that a Brownian particle started at $x$ is in $(-1,1)$ at time $t$.}\label{ch18:fig-heat}
\end{figure}

\subsection{Solution by similarity and invariance}
The 2024 lecture derives the same solution without the kernel \ts{24}{25}{51:49}. Rescale time and space, $\tau=\alpha t$ and $y=\beta x$, and put $v(y,\tau)=u(x,t)$. The chain rule gives $u_t=\alpha v_\tau$, $u_x=\beta v_y$, $u_{xx}=\beta^2v_{yy}$, so \eqref{ch18:eq-heat} becomes $\alpha v_\tau=\lambda\beta^2v_{yy}$: for $\alpha=\beta^2$, $v_\tau=\lambda v_{yy}$, \emph{the same equation} \ts{24}{25}{52:53}\ts{24}{25}{53:54}\ts{24}{25}{54:58}. The heat equation is invariant under the scaling $t\to\alpha t$, $x\to\sqrt\alpha\,x$ (diffusive scaling). One then looks for solutions with the same invariance, $u(x,t)=u(\sqrt\alpha x,\alpha t)$ for all $\alpha$; taking $\alpha=1/t$,
\[
  u(x,t)=u\bigl(x/\sqrt t,1\bigr)=h\bigl(x/\sqrt t\bigr) ,
\]
a function of \emph{one} variable \ts{24}{25}{56:00}. Substituting $y=x/\sqrt t$ in \eqref{ch18:eq-heat}: $u_t=-\tfrac{y}{2t}h'(y)$ and $u_{xx}=h''(y)/t$, so $-\tfrac12yh'=\lambda h''$, i.e.\ $\dfrac{\dd}{\dd y}\ln h'(y)=-\dfrac{y}{2\lambda}$ \ts{24}{25}{57:04}. Integrating twice,
\[
  h'(y)=c\,e^{-y^2/(4\lambda)},\qquad h(y)=\Phi\Bigl(\frac y{\sqrt{2\lambda}}\Bigr)
\]
(a Gaussian density with variance $2\lambda$, whose antiderivative is the normal distribution function; the constants are fixed by $h(-\infty)=0$, $h(+\infty)=1$) \ts{24}{25}{58:08}. Hence
\begin{equation}\label{ch18:eq-step}
  u(x,t)=\Phi\Bigl(\frac{x}{\sqrt{2\lambda t}}\Bigr)\quad\longrightarrow\quad H(x)=\begin{cases}1&x>0\\ \tfrac12&x=0\\ 0&x<0\end{cases}\quad(t\to0),
\end{equation}
so it solves the heat equation with a \emph{step} initial condition \ts{24}{25}{59:14}\ts{24}{25}{1:00:22}. Translating, $\Phi\bigl((x-a)/\sqrt{2\lambda t}\bigr)$ has a step at $a$; the difference of two of them has the initial condition $\ind_{(a,b)}$ (with value $\tfrac12$ at the endpoints), which reproduces \eqref{ch18:eq-indic} \ts{24}{25}{1:01:25}\ts{24}{25}{1:02:28}\ts{24}{25}{1:03:32}. Every step function, and by limits every reasonable initial condition, is a combination of steps \ts{24}{25}{1:03:32}.
\begin{remark}[The kernel is the derivative of the step solution]
Since $\partial_x$ commutes with the equation, $\partial_xu$ solves it whenever $u$ does, and $\partial_xH=\delta$. Indeed $\partial_x\Phi\bigl(x/\sqrt{2\lambda t}\bigr)=u_\delta(x,t)$ of \eqref{ch18:eq-kernel}: the two derivations of the lecture (delta function, then similarity for the step) are one, in this way (own remark). The 2013 note leaves the derivation of \eqref{ch18:eq-kernel} as its Exercise~3.2, via $\xi=x/\sqrt t$ and $U(\xi)=t^{1/2}u(x,t)$ (\cref{ch18:ex-32}).
\end{remark}
\begin{coursenote}[2013 and 2024 on the heat equation; notation and misprints]
2013 states the kernel and the superposition and asserts the link with Black--Scholes; the reduction is proposed as a final project \ts{13}{21}{54:18}\ts{13}{21}{55:22} (carried out in \cref{ch15:sec-heat}). 2024 derives the solution by similarity and adds the step-function construction. Normalisations differ: 2013 has $u_t=u_{xx}$ ($\lambda=1$), 2024 slide 5 has general $\lambda$ and slides 8--9 specialise to $\lambda=\tfrac12$, i.e.\ Brownian motion with unit variance rate. The 2013 note writes the initial condition sometimes as $u(0,x)$ and sometimes as $u_0(x,0)$; the arguments of $u$ appear in both orders ($u(x,t)$ and $u(0,x)$), which is only a notational slip. In the transcript of 2024 the argument of $\Phi$ is spoken as ``$y$ over $2\sqrt{2\lambda}$'' \ts{24}{25}{58:08}; the slide (and the calculation above) has $y/\sqrt{2\lambda}$.
\end{coursenote}

% ==================================================================
\section{Numerical solution of SDEs}\label{ch18:sec-num}

Most SDEs, like most PDEs, have no closed-form solution, and the closed-form methods above fail for most of them \ts{13}{21}{19:56}\ts{13}{21}{20:56}. Both lectures list the numerical options, finite-difference methods, Monte Carlo methods and tree methods (2024 slide 21, without detail); the 2013 lecture explains them \ts{13}{21}{20:56}.

\subsection{Finite differences for an ODE}
The idea of the finite-difference method is Taylor's formula on a grid. The lecture's example is $u'(x)=5u(x)+2$, $u(0)=0$, to be evaluated at $x=1$ \ts{13}{21}{22:04}. Chop $[0,1]$ into steps of size $h$ and use $u((i+1)h)\approx u(ih)+h\,u'(ih)=u(ih)+h\bigl(5u(ih)+2\bigr)$, which computes each value from the previous one \ts{13}{21}{23:07}\ts{13}{21}{24:12}\ts{13}{21}{25:16}. With $h=\tfrac12$: $u(\tfrac12)\approx1$, $u(1)\approx\tfrac92$ (the note's Example~2.1). The exact solution is $u=\tfrac25(e^{5x}-1)$, $u(1)=58.97$: with two steps the answer is off by a factor $13$, and it improves only as $h\to0$ (own computation of the scheme $u_n=\tfrac25[(1+5h)^n-1]$):
\begin{center}
\small
\begin{tabular}{@{}lrrrrr@{}}
\toprule
$h$ & $1/2$ & $1/10$ & $1/100$ & $1/1000$ & exact\\
\midrule
$u(1)$ & $4.50$ & $22.67$ & $52.20$ & $58.23$ & $58.97$\\
\bottomrule
\end{tabular}
\end{center}
The error is proportional to $h$ (first order): this is the \emph{Euler method}, and the theorem quoted in the lecture is that for a reasonable equation the approximation converges to the true value as $h\to0$ \ts{13}{21}{25:16}. The same idea applies to functions of two variables on a grid, filled layer by layer from the boundary values \ts{13}{21}{26:18}. The lecture also warns about the cost/accuracy trade-off in choosing $h$ \ts{13}{21}{39:02}: a smaller $h$ gives a better answer and costs more time; and a grid fine enough for $X(1)$ may be too coarse for intermediate times close to $0$ where only a few steps have been taken.

\subsection{Why it does not apply directly, and the fix: Monte Carlo along a fixed path}
For the SDE \eqref{ch18:eq-sde} the analogue of $u'=f(u)$ has a random right-hand side: the point at time $2h$ could have been reached from many points at time $h$, and it is unknown which \ts{13}{21}{27:20}\ts{13}{21}{28:21}. The fix is to \emph{fix a Brownian path}: once a sample path $B$ is drawn, every increment $\Delta B_i=B_{(i+1)h}-B_{ih}$ is known and the same Taylor step can be used \ts{13}{21}{28:21}\ts{13}{21}{29:21}\ts{13}{21}{30:30}\ts{13}{21}{31:37}:
\begin{equation}\label{ch18:eq-em}
  X_{i+1}=X_i+\mu(t_i,X_i)\,h+\sigma(t_i,X_i)\,\Delta B_i,\qquad t_i=ih,\quad\Delta B_i=\sqrt h\,Z_i,\ Z_i\sim N(0,1)\ \text{i.i.d.}
\end{equation}
This is the \emph{Euler--Maruyama scheme}. Then repeat with many independent paths (the \emph{Monte Carlo simulation}): each gives one value $X_T^{(k)}$, and the empirical distribution of these values approximates the law of $X_T$ \ts{13}{21}{31:37}\ts{13}{21}{32:42}\ts{13}{21}{33:43}\ts{13}{21}{34:48}. Any intermediate time can be read off the same paths, subject to the accuracy remark above. The steps of the note, \emph{sample a path, solve along it, repeat}, are exactly this. (The obstacle is the underlying randomness of the Brownian motion; once a path is fixed it disappears \ts{13}{21}{33:43}.)

\begin{definition}[Strong and weak convergence; not discussed in class]\label{ch18:def-conv}
A scheme with step $h=T/n$ and output $X^h_T$ has \emph{strong order} $q$ if $\E|X_T-X^h_T|\le Ch^q$ (paths are close: the two are driven by the same $B$), and \emph{weak order} $w$ if $|\E f(X_T)-\E f(X^h_T)|\le Ch^w$ for smooth $f$ (laws are close).
\end{definition}
Euler--Maruyama has strong order $\tfrac12$ and weak order $1$ under the conditions of \cref{ch18:thm-eu} and smooth coefficients (quoted). The order $\tfrac12$ is the price of the noise: over one step the increment is of size $\sqrt h$. The \emph{Milstein} scheme adds the next term of the It\^o--Taylor expansion, $X_{i+1}=X_i+\mu h+\sigma\Delta B_i+\tfrac12\sigma\sigma'(\Delta B_i^2-h)$, and has strong order $1$ (the correction is $0$ when $\sigma$ is constant, so for OU Euler--Maruyama already has order $1$ in the strong sense). Finally the Monte Carlo error itself, of order $N^{-1/2}$ in the number $N$ of paths, adds to the discretisation error; the two should be balanced.

\begin{casestudy}[Convergence of Euler--Maruyama for geometric Brownian motion]\label{ch18:cs-em}
\emph{Goal.} Check the orders with a case where the exact solution \eqref{ch18:eq-gbmsol} is known. \emph{Set-up.} $\dd X=\mu X\dd t+\sigma X\dd B$, $\mu=0.3$, $\sigma=0.4$, $x_0=1$, $T=1$ (the parameters of the 2024 example), $N=2\times10^5$ Brownian paths generated on the finest grid ($n=256$) and coarsened by summing increments, so that all schemes and the exact solution see the \emph{same} path. Code (own; NumPy):
\begin{lstlisting}[language=Python]
import numpy as np
rng = np.random.default_rng(1)
mu, sg, T, N, nmax = 0.3, 0.4, 1.0, 200000, 256
dW = rng.standard_normal((N, nmax)) * np.sqrt(T / nmax)
for n in (4, 8, 16, 32, 64, 128, 256):
    h = T / n
    d = dW.reshape(N, n, nmax // n).sum(axis=2)     # coarse increments
    X = np.ones(N); M = np.ones(N)
    for i in range(n):
        X = X * (1 + mu*h + sg*d[:, i])              # Euler-Maruyama
        M = M * (1 + mu*h + sg*d[:, i] + 0.5*sg**2*(d[:, i]**2 - h))  # Milstein
    exact = np.exp((mu - sg**2/2)*T + sg*d.sum(axis=1))
    print(n, abs(X-exact).mean(), abs(M-exact).mean(), (1+mu*h)**n - np.exp(mu*T))
\end{lstlisting}
\emph{Result} (columns: steps $n$, strong error of Euler, strong error of Milstein, weak error of $\E X_T$, which is deterministic here: $\E X^h_T=(1+\mu h)^n$):
\begin{center}
\small
\begin{tabular}{@{}rccc@{}}
\toprule
$n$ & Euler, $\E|X_T-X_T^h|$ & Milstein & $|\E X_T^h-\E X_T|$\\
\midrule
$4$ & $5.9\times10^{-2}$ & $3.3\times10^{-2}$ & $1.4\times10^{-2}$\\
$16$ & $3.0\times10^{-2}$ & $8.7\times10^{-3}$ & $3.7\times10^{-3}$\\
$64$ & $1.5\times10^{-2}$ & $2.2\times10^{-3}$ & $9.5\times10^{-4}$\\
$256$ & $7.6\times10^{-3}$ & $5.5\times10^{-4}$ & $2.4\times10^{-4}$\\
\bottomrule
\end{tabular}
\end{center}
Halving $h$ divides the Euler strong error by $\sqrt2$ (fitted slopes $0.49$--$0.50$), the Milstein strong error and the weak error by $2$ (slopes $0.94$ rising to $1.00$): strong order $\tfrac12$ (Euler), $1$ (Milstein) and weak order $1$, as stated. The same run shows why the distinction matters for finance: for \emph{prices}, i.e.\ expectations, the weak order is what counts, and Euler is enough; for path-dependent hedging error one needs paths, i.e.\ strong convergence. \Cref{ch18:fig-em} plots the errors. \emph{Lesson.} A convergence experiment against a known solution is the way to validate a simulator before using it where no closed form exists.
\end{casestudy}
\begin{figure}[ht]
\centering
\begin{tikzpicture}
\begin{loglogaxis}[width=0.85\linewidth,height=5.4cm,xlabel={step $h=T/n$},ylabel={error},xmin=0.003,xmax=0.3,ymin=0.0001,ymax=0.1,
  legend style={at={(0.5,1.03)},anchor=south,legend columns=3,font=\footnotesize},grid=major,grid style={black!10}]
  \addplot+[mark=*,thick,color=defcol] coordinates {(0.25,0.05936) (0.125,0.04236) (0.0625,0.03013) (0.03125,0.02141) (0.015625,0.01521) (0.0078125,0.01076) (0.00390625,0.0076)};
  \addplot+[mark=square*,thick,color=thmcol] coordinates {(0.25,0.0326) (0.125,0.01699) (0.0625,0.00867) (0.03125,0.00438) (0.015625,0.0022) (0.0078125,0.0011) (0.00390625,0.00055)};
  \addplot+[mark=triangle*,thick,color=intcol] coordinates {(0.25,0.01439) (0.125,0.007388) (0.0625,0.003744) (0.03125,0.001885) (0.015625,0.000946) (0.0078125,0.000474) (0.00390625,0.000237)};
  \addplot[dashed,black!60,forget plot] coordinates {(0.004,0.0077) (0.25,0.0077*7.905)};
  \legend{{Euler, strong},{Milstein, strong},{Euler, weak ($\E X_T$)}}
\end{loglogaxis}
\end{tikzpicture}
\caption{Errors of \cref{ch18:cs-em} against the step size: Euler strong error (slope $\tfrac12$, the dashed line is a slope-$\tfrac12$ reference), Milstein strong error and the weak error of $\E X_T$ (both slope $1$).}\label{ch18:fig-em}
\end{figure}

\subsection{The tree method}
\emph{Trees} use the other picture of Brownian motion, as the limit of a random walk (\cref{ch:sp2}, Donsker) \ts{13}{21}{34:48}\ts{13}{21}{35:50}: replace $\Delta B_i$ in \eqref{ch18:eq-em} by $\pm\sqrt h$ with probability $\tfrac12$ each. Then $X_T$ takes only finitely many values, and their probabilities are known exactly: after $n$ steps of a recombining tree the node with $j$ up-moves has probability $\binom nj2^{-n}$ (in the lecture's picture with $n=5$ steps, $1,5,10,10,5,1$ over $32$) \ts{13}{21}{36:54}. So one \emph{computes} the expectation instead of sampling it: no statistical error. For geometric Brownian motion the multiplicative form $X_{i+1}=X_i(1+\mu h\pm\sigma\sqrt h)$ does recombine, and \eqref{ch18:eq-em} is Euler--Maruyama with two-point increments. Example (own computation; payoff $(X_T-1)^+$, same parameters, exact value $x_0e^{\mu T}\Phi(d_1)-K\Phi(d_2)=0.41012$ with $d_1=[\ln(x_0/K)+(\mu+\sigma^2/2)T]/(\sigma\sqrt T)$, $d_2=d_1-\sigma\sqrt T$):
\begin{center}
\small
\begin{tabular}{@{}lrrrr@{}}
\toprule
steps $n$ & $16$ & $64$ & $256$ & $1024$\\
\midrule
Euler-type tree $X(1+\mu h\pm\sigma\sqrt h)$, error & $-5.9\times10^{-3}$ & $-1.2\times10^{-3}$ & $-2.2\times10^{-4}$ & $-7.5\times10^{-5}$\\
log-tree $x_0e^{(\mu-\sigma^2/2)T+\sigma\sqrt h(2j-n)}$, error & $-1.1\times10^{-3}$ & $4\times10^{-5}$ & $1.0\times10^{-4}$ & $5\times10^{-6}$\\
\bottomrule
\end{tabular}
\end{center}
The errors are not monotone for the second tree (the strike sits at a different place relative to the nodes as $n$ changes, a known feature of the binomial method for a kinked payoff) but tend to $0$ with $n$. The 2013 speaker adds that accuracy can be increased in two ways, with more possibilities at each step (a trinomial tree with steps $+1,0,-1$ also converges to Brownian motion) or with smaller time steps \ts{13}{21}{41:06}, and remarks that he does not know which of the methods is used in practice, but that Monte Carlo and trees seem the more important ones \ts{13}{21}{36:54}. A comment from the audience in the 2013 lecture gives the practical criterion \ts{13}{21}{40:02}\ts{13}{21}{41:06}: for \emph{path-independent} claims (the payoff depends on $X_T$ only, so its law suffices), Monte Carlo is natural; for \emph{path-dependent} problems, in particular early exercise (American options), a tree is preferable because at every node one can condition on the situation reached and compare exercise value with continuation value.
\begin{finance}[Pricing by simulation and by PDE]
Monte Carlo pricing is Feynman--Kac \eqref{ch18:eq-fk} evaluated by simulation: simulate $S$ under the risk-neutral dynamics $\dd S=rS\dd t+\sigma S\dd W$ (\cref{ch17:sec-girsanov}), average the discounted payoff $e^{-rT}g(S_T)$ over $N$ paths, with standard error $O(N^{-1/2})$ independent of dimension. This is the method of choice for high-dimensional or path-dependent payoffs; PDE methods (finite differences, \cref{ch18:sec-heat,ch:bs}) and trees are faster in one or two dimensions and handle early exercise naturally; the cost of a grid grows exponentially with dimension. Which drift to simulate is the point of \cref{ch18:sec-fk}: the risk-free rate $r$, not the real-world $\mu$ (which matters for risk management and backtests, \cref{ch:var}).
\end{finance}

% ==================================================================
\section{Beyond the lectures: square-root and other diffusions}\label{ch18:sec-cir}

The two examples solved in class have linear coefficients, hence Lipschitz. Two extensions show what changes when they are not (an addition of these notes; none of this is in the lectures).

\paragraph{The CIR process.} The Cox--Ingersoll--Ross model $\dd X=\kappa(\theta-X)\,\dd t+\sigma\sqrt X\,\dd B$ is a mean-reverting process whose noise vanishes at $0$, so it stays nonnegative: the natural model of a positive interest rate or a variance (the Heston model of \cref{ch:volatility} uses it for the variance). The coefficient $\sigma\sqrt x$ is not Lipschitz at $0$, so \cref{ch18:thm-eu} does not apply, but $x\mapsto\sqrt x$ is H\"older-$\tfrac12$ and the Yamada--Watanabe theorem (quoted) still gives a unique solution. Boundary behaviour is governed by the \emph{Feller condition}: if $2\kappa\theta\ge\sigma^2$ the process never reaches $0$; otherwise it hits $0$ and is reflected. The linear drift makes the moments tractable: by Dynkin's formula \eqref{ch18:eq-dynkin}, $\dd\E X_t=\kappa(\theta-\E X_t)\dd t$, so
\[
  \E X_t=\theta+(x_0-\theta)e^{-\kappa t},\qquad
  \Var X_t=x_0\frac{\sigma^2}{\kappa}\bigl(e^{-\kappa t}-e^{-2\kappa t}\bigr)+\theta\frac{\sigma^2}{2\kappa}\bigl(1-e^{-\kappa t}\bigr)^2
\]
(the variance from the second-moment equation $\tfrac{\dd}{\dd t}\E X^2=(2\kappa\theta+\sigma^2)\E X-2\kappa\E X^2$). Unlike OU the law is not normal: $X_t$ is $c$ times a noncentral $\chi^2$ variable with $d=4\kappa\theta/\sigma^2$ degrees of freedom, $c=\sigma^2(1-e^{-\kappa t})/(4\kappa)$ and noncentrality $x_0e^{-\kappa t}/c$, which allows exact sampling. Check ($\kappa=1.5$, $\theta=0.04$, $\sigma=0.2$, $x_0=0.09$, $t=1$; Feller ratio $2\kappa\theta/\sigma^2=3$): the formulas give mean $0.05116$ and variance $7.379\times10^{-4}$; $2\times10^6$ exact samples give $0.05114$ and $7.375\times10^{-4}$ (own computation). The Euler scheme \eqref{ch18:eq-em} may produce negative values, for which $\sqrt X$ is undefined; the standard cure is to truncate ($\sqrt{X^+}$), and the exact sampler above avoids the issue. In the same family, the squared radius $|\bm B|^2$ of an $n$-dimensional Brownian motion (\cref{ch17:sec-multi}) satisfies $\dd Z=n\,\dd t+2\sqrt Z\,\dd W$, a CIR process without mean reversion (the squared Bessel process).

\paragraph{Other state-dependent volatilities.} The choice $\sigma(x)=\sigma x^{\beta}$ (CEV) interpolates between a normal model ($\beta=0$, the OU/Bachelier world) and geometric Brownian motion ($\beta=1$); local- and stochastic-volatility models, in which $\sigma$ is a function of the price or itself a diffusion, lead to SDEs that must be treated numerically (\cref{ch:volatility}). The theory of this chapter (existence, generator, Feynman--Kac, Fokker--Planck, Euler--Maruyama) applies to all of them.

% ==================================================================
\section{The two lectures compared, errata}\label{ch18:sec-compare}

\begin{coursenote}[2013 and 2024 compared]
\emph{Common:} existence and uniqueness under the Lipschitz and growth conditions (stated without proof); coefficient matching for geometric Brownian motion; the Ornstein--Uhlenbeck process solved by the product-process guess $a(t)(x_0+\int b\,\dd B)$; the heat equation solved by superposition of Gaussians, with the link to Black--Scholes. \emph{Only in 2013:} the goal of an SDE as the reverse of differentiation; finite differences for ODEs with the example $u'=5u+2$, the fixed-path Monte Carlo idea, the tree method with its binomial weights, the comments on trinomial trees and path dependence; the Dirac delta and Gaussian kernel as \emph{given} (derivation left as Exercise~3.2). \emph{Only in 2024:} the Black--Scholes equation by hedging (slides 2--4; the 2013 counterpart is a different lecture), the $L^2$ bound in the theorem, the lognormal ``paradox'' of geometric Brownian motion, the properties of the OU process (limit law, stationary start, the case $m\ne0$), Vasicek's model, systems of SDEs and the position--velocity example, the derivation of the heat solution by similarity/invariance and step functions; the numerical methods are only listed (slide 21). Neither year mentions the Feynman--Kac formula, Girsanov's theorem for SDEs, the CIR process or the Euler--Maruyama order of convergence.
\end{coursenote}

\begin{coursenote}[Errata and misprints]
\begin{itemize}
  \item 2024 slide 15: $\sigma_*=\sigma t$ should be $\sigma\sqrt t$ (\cref{ch18:sec-cm}).
  \item \file{lecnote21}, Theorem 1.1 and the paragraph before it: ``for given functions $a$ and $b$'' should be $\mu,\sigma$; the integral equation lacks $x_0$; ``(L2 bound)'' is a placeholder for $\sup_t\E X_t^2<\infty$; ``Lipshictz'' is spelled so in both the note and slide 13.
  \item \file{lecnote21}, \S2.2 says that for a known path $B_t$ ``the equation now becomes a PDE''; for a fixed path it is a recursion in one variable (an ODE-type finite-difference step), as \eqref{ch18:eq-em} shows; the tree and Monte Carlo descriptions write the equation as $\dd f(t,B_t)=g\,\dd B+h\,\dd t$, which is the form of a function of Brownian motion rather than of a general SDE.
  \item The 2013 transcript's remark that OU has ``a drift term \dots\ minus exponential'' at 13:16 is a slip: the drift is \emph{linear}, $-\kappa X$; the exponentials appear only in the solution.
  \item In the 2024 transcript $\Phi\bigl(y/2\sqrt{2\lambda}\bigr)$ is spoken; the slide has $\Phi(y/\sqrt{2\lambda})$, which is right (\cref{ch18:sec-heat}).
  \item PS9 Problem A-1(c): see the note after \cref{ch18:ex-a1}. Everything else in the slides (Black--Scholes derivation, heat kernel, OU moments) was re-derived above and found correct; the numerics of this chapter were checked in Python and, where indicated, symbolically.
\end{itemize}
\end{coursenote}

% ==================================================================
\begin{keyformulas}
\begin{align*}
&\text{SDE:} && \dd X=\mu(t,X)\dd t+\sigma(t,X)\dd B,\quad X_t=x_0+\textstyle\int_0^t\mu\,\dd s+\int_0^t\sigma\,\dd B\\
&\text{Existence, uniqueness:} && |\mu(x)-\mu(y)|^2+|\sigma(x)-\sigma(y)|^2\le K|x-y|^2,\\
& && |\mu|^2+|\sigma|^2\le K(1+|x|^2)\ \Rightarrow\ \exists!\ \text{solution},\ \sup_t\E X_t^2<\infty\\
&\text{Coefficient matching:} && X=f(t,B):\ \mu(t,f)=f_t+\tfrac12f_{xx},\quad\sigma(t,f)=f_x\\
&\text{GBM:} && \dd X=\mu X\dd t+\sigma X\dd B\ \Rightarrow\ X_t=x_0e^{(\mu-\sigma^2/2)t+\sigma B_t},\quad\E X_t=x_0e^{\mu t}\\
& && \mu<\tfrac12\sigma^2:\ X_t\to0\ \text{a.s.};\quad\text{std.\ dev.\ of }\ln X_t=\sigma\sqrt t\\
&\text{OU:} && \dd X=-\kappa(X-m)\dd t+\sigma\dd B\ \Rightarrow\ X_t=m+(x_0-m)e^{-\kappa t}+\sigma\textstyle\int_0^te^{-\kappa(t-s)}\dd B_s\\
& && \E X_t=m+(x_0-m)e^{-\kappa t},\quad\Var X_t=\tfrac{\sigma^2}{2\kappa}(1-e^{-2\kappa t}),\quad X_\infty\sim N\bigl(m,\tfrac{\sigma^2}{2\kappa}\bigr)\\
& && \Cov(X_s,X_t)\to\tfrac{\sigma^2}{2\kappa}e^{-\kappa|t-s|}\ (\text{stationary}),\quad\text{half-life }\ln2/\kappa\\
&\text{Exact AR(1):} && X_{t+h}-m=\phi(X_t-m)+\varepsilon,\quad\phi=e^{-\kappa h},\quad\Var\varepsilon=\tfrac{\sigma^2}{2\kappa}(1-\phi^2)\\
& && \hat\kappa=-\ln\hat\phi/h,\quad\hat\sigma^2=2\hat\kappa\hat\sigma_\varepsilon^2/(1-\hat\phi^2)\\
&\text{Vector OU:} && \bm X_t=\bm m+e^{-At}(\bm x_0-\bm m)+\textstyle\int_0^te^{-A(t-s)}\Sigma\dd\bm B_s,\\
& && \quad\Phi=e^{-Ah}\\
&\text{Generator:} && \mathcal L=\mu\partial_x+\tfrac12\sigma^2\partial_x^2,\quad\dd f=(\partial_t+\mathcal L)f\dd t+\sigma f_x\dd B\\
&\text{Feynman--Kac:} && u_t+\mathcal Lu-ru=0,\ u(T,\cdot)=g\ \Longleftrightarrow\ u=\E\bigl[e^{-\int_t^Tr}\,g(X_T)\mid X_t=x\bigr]\\
&\text{Fokker--Planck:} && \partial_tp=-\partial_y(\mu p)+\tfrac12\partial_y^2(\sigma^2p),\quad p_{\mathrm{st}}\propto e^{-2U/\sigma^2}\ (\mu=-U')\\
&\text{Vasicek bond:} && P=e^{-br+(b-\tau)(\theta-\sigma^2/2\kappa^2)},\\
& && \quad{}-\sigma^2b^2/4\kappa,\quad b=\tfrac{1-e^{-\kappa\tau}}\kappa\\
&\text{Heat equation:} && u_t=\lambda u_{xx}\ \Rightarrow\ u=\textstyle\int\frac{e^{-(x-s)^2/4\lambda t}}{\sqrt{4\pi\lambda t}}u_0(s)\dd s,\\
& && u=\E\,u_0(x+\sqrt{2\lambda}B_t)\\
& && u_0=\ind_{(a,b)}:\ u=\Phi\bigl(\tfrac{b-x}{\sqrt{2\lambda t}}\bigr)\\
& && \quad{}-\Phi\bigl(\tfrac{a-x}{\sqrt{2\lambda t}}\bigr)\\
&\text{Euler--Maruyama:} && X_{i+1}=X_i+\mu h+\sigma\sqrt hZ_i;\ \text{strong order }\tfrac12,\ \text{weak order }1\\
&\text{Milstein:} && +\tfrac12\sigma\sigma'(\Delta B_i^2-h);\ \text{strong order }1
\end{align*}
\end{keyformulas}

% ==================================================================
\section*{Exercises}
\addcontentsline{toc}{section}{Exercises}

\subsection*{From 2013 Problem Set 9 (stochastic differential equations)}
The problem set (\file{pset9.pdf}, due 12/5/13, ``no need to turn in'') has only Part~A, which ``straightforwardly follows from the definition''. There is no 2024 problem set on this topic.

\begin{exercise}[PS9 (2013), Problem A-1: verify solutions]\label{ch18:ex-a1}
Verify that the given process solves the given SDE (use It\^o's formula, \cref{ch:stochcalc}).
\begin{enumerate}[(a)]
  \item $X_t=\exp(B_t)$ solves $\dd X_t=\tfrac12X_t\,\dd t+X_t\,\dd B_t$.
  \item $X_t=B_t/(1+t)$ solves $\dd X_t=-\frac1{1+t}X_t\,\dd t+\frac1{1+t}\,\dd B_t$.
  \item $X_t=\sin B_t$ solves $\dd X_t=-\tfrac12X_t\,\dd t+\sqrt{1-X_t^2}\,\dd B_t$.
\end{enumerate}
\end{exercise}
\begin{coursenote}[Erratum in Problem A-1(c)]
It\^o's formula gives $\dd\sin B_t=-\tfrac12\sin B_t\,\dd t+\cos B_t\,\dd B_t$, and $\sqrt{1-\sin^2B_t}=|\cos B_t|$, which equals $\cos B_t$ only while $\cos B_t\ge0$, i.e.\ up to the first time $|B|$ reaches $\pi/2$. The stated equation therefore holds until that time; afterwards it holds only in law, with $\dd\widetilde B_t=\operatorname{sgn}(\cos B_t)\dd B_t$ (again a Brownian motion, by L\'evy's characterisation). The verification asked in the problem is exactly the computation above.
\end{coursenote}

\begin{exercise}[PS9 (2013), Problem A-2: a cube-root SDE]\label{ch18:ex-a2}
Let $a>0$ and $\dd X_t=\tfrac13X_t^{1/3}\dd t+X_t^{2/3}\dd B_t$, $X_0=a$. Show that $X_t=\bigl(a^{1/3}+\tfrac13B_t\bigr)^3$, $t\ge0$, solves it. (Where do the conditions of \cref{ch18:thm-eu} fail? See \cref{ch18:ex-fail}.)
\end{exercise}

\begin{exercise}[PS9 (2013), Problem A-3: the Vasicek model]\label{ch18:ex-a3}
(Vasicek interest-rate model.) Prove that $R(t)=e^{-\beta t}R(0)+\frac\alpha\beta\bigl(1-e^{-\beta t}\bigr)+\sigma e^{-\beta t}\int_0^te^{\beta s}\dd B_s$ solves the SDE $\dd R(t)=(\alpha-\beta R(t))\,\dd t+\sigma\,\dd B_t$.
\end{exercise}

\subsection*{From the 2013 lecture note}
\begin{exercise}[\file{lecnote21}, Exercise 3.2: derivation of the heat kernel]\label{ch18:ex-32}
Derive $u_\delta(x,t)=\frac1{2\sqrt{\pi t}}e^{-x^2/(4t)}$, the solution of $u_t=u_{xx}$ with $u(x,0)=\delta(x)$, by using the variable $\xi=x/\sqrt t$ and the function $U(\xi)=t^{1/2}u(x,t)$, and restating the heat equation as an ODE for $U$.
\end{exercise}

\subsection*{Additional exercises (not from the course)}
\begin{exercise}[not from the course: OU covariance and estimation]\label{ch18:ex-ou}
\leavevmode
\begin{enumerate}[(a)]
  \item Derive \eqref{ch18:eq-oucov} and check it against \eqref{ch18:eq-ouvar} at $s=t$.
  \item Show that the OU process is the only stationary Gaussian Markov process with continuous paths, up to parameters (hint: $\rho(\tau)$ must satisfy $\rho(\tau+\tau')=\rho(\tau)\rho(\tau')$).
  \item Daily data ($h=1/252$) of a spread give an AR(1) fit with $\hat\phi=0.995$ and residual standard deviation $0.02$. Compute $\hat\kappa$, the half-life in days and the stationary standard deviation of the spread.
\end{enumerate}
\end{exercise}

\begin{exercise}[not from the course: exact versus Euler for OU]\label{ch18:ex-euler}
Find the stationary variance of the Euler--Maruyama chain $X_{i+1}=X_i-\kappa(X_i-m)h+\sigma\sqrt hZ_i$ and compare it with $\sigma^2/(2\kappa)$. For which $h$ is the chain stationary? By what factor is the variance wrong for $\kappa h=0.5$? Write a short program that simulates both schemes and compares the empirical variance.
\end{exercise}

\begin{exercise}[not from the course: Feynman--Kac]\label{ch18:ex-fk}
\leavevmode
\begin{enumerate}[(a)]
  \item For geometric Brownian motion $\dd S=rS\dd t+\sigma S\dd W$ and payoff $g(S)=S^p$ compute $u(t,S)=e^{-r\tau}\E[S_T^p\mid S_t=S]$ and verify \eqref{ch18:eq-fkpde}.
  \item Check that $u=S$ (the stock) and $u=e^{-r\tau}$ (a bond) solve the Black--Scholes equation, and explain in words what this says.
  \item For the OU process solve $u_t-\kappa(x-m)u_x+\tfrac12\sigma^2u_{xx}=0$, $u(T,x)=x$, with the ansatz $u=A(\tau)x+B(\tau)$ and compare with the mean \eqref{ch18:eq-oumean}.
\end{enumerate}
\end{exercise}

\begin{exercise}[not from the course: Fokker--Planck]\label{ch18:ex-fp}
\leavevmode
\begin{enumerate}[(a)]
  \item For $\dd X=-U'(X)\dd t+\sigma\dd B$ with $U(x)=\tfrac14x^4-\tfrac12x^2$ (double well) find the stationary density and show it is normalisable.
  \item Derive the Fokker--Planck equation of geometric Brownian motion and check that the lognormal density solves it.
  \item For $\sigma^2=0.5$ compute the ratio of the stationary density at a well bottom $x=1$ to that at the barrier top $x=0$.
\end{enumerate}
\end{exercise}

\begin{exercise}[not from the course: position and velocity]\label{ch18:ex-posvel}
For the system \eqref{ch18:eq-posvel} with $m=0$, $\sigma_1=0$:
\begin{enumerate}[(a)]
  \item derive \eqref{ch18:eq-posvar} without assuming stationarity of $V_0$, taking $V_0=v_0$ deterministic;
  \item show that the diffusion coefficient in the long-time limit is $D=\sigma^2/(2\kappa^2)$ and relate it to $\gamma$ and $T$ in the physical parametrisation of the intuition box in \cref{ch18:sec-ou};
  \item show that $X_t$ converges in law to $\sqrt{2D}\,B_t$ in the limit $\kappa\to\infty$ with $\sigma^2/\kappa^2=2D$ fixed.
\end{enumerate}
\end{exercise}

\begin{exercise}[not from the course: Vasicek bond]\label{ch18:ex-vasbond}
Derive the ODEs for $a(\tau),b(\tau)$ in \cref{ch18:fin-bond}, solve them and confirm \eqref{ch18:eq-vasbond}. Compute the yield curve for $\tau=1,2,5,10,30$ years for $\kappa=0.5$, $\theta=4\%$, $\sigma=1\%$, $r_0=3\%$. Show that the yield is a decreasing function of $\sigma$ at long maturities and interpret it.
\end{exercise}

\begin{exercise}[not from the course: convergence experiment]\label{ch18:ex-conv}
Repeat \cref{ch18:cs-em} for the OU process (where the exact solution is \eqref{ch18:eq-ar1}) and for the SDE $\dd X=\sin X\,\dd t+\cos X\,\dd B$ (no closed form: use a very fine Euler solution as reference). Fit the strong and weak orders of Euler--Maruyama and Milstein. Why does Milstein give no improvement for OU?
\end{exercise}

\begin{exercise}[not from the course: CIR moments]\label{ch18:ex-cir}
Derive the moment equations of the CIR process from Dynkin's formula and solve for $\E X_t$. Verify the stationary mean $\theta$ and variance $\theta\sigma^2/(2\kappa)$, and check the Feller condition for a short-rate CIR model with $\kappa=0.3$, $\theta=3\%$, $\sigma=0.1$.
\end{exercise}


\part{Additional topics from 2013 (18.S096)}\label{part:2013}
% !TEX root = ../main.tex
\chapter{Value at Risk Models}\label{ch:var}
\begin{paracol}{2}

\begin{sources}
\begin{tabularx}{\linewidth}{@{}l l L@{}}
\toprule
\textbf{Lecture} & \textbf{Speaker} & \textbf{Content}\\
\midrule
2013 L7 & K.~Abbott (Morgan Stanley, Firm Risk Management) & What a risk-management group is for; one-asset VaR as a $1\%$ order statistic; returns instead of levels; the normal multipliers $2.33$ and $1.645$; fat tails; fixed-income positions through the PV01; covariance and correlation, portfolio VaR as $\bm x\T\Sigma\bm x$ (variance--covariance method); data and weighting choices; Monte Carlo VaR through the eigen-decomposition of the covariance matrix (the ``Gaussian copula''). Historical simulation, marginal VaR and back-testing are on the slides only: Abbott announced that he would skip historical simulation.\\
\bottomrule
\end{tabularx}

\smallskip
\textbf{Files.} Slides \file{MIT18_S096F13_lecnote7.pdf} (115 pages, about 100 distinct slides).
\textbf{Problem sets.} None. Abbott mentions a homework on plotting distributions that he usually sets before this lecture; it is not among the course files. (The $t$ distribution that a student asks about \ts{13}{7}{26:29} belongs to \cref{ch:volatility}, where the 2013 problem set on volatility is treated.)
\textbf{Overlap with other chapters.} Coherent risk measures, VaR and expected shortfall, the normal formulas, the square-root-of-time rule, the estimation of VaR by historical simulation and back-testing with the binomial test are in \cref{ch08:sec-riskmeasures} (from the 2024 lecture); VaR with GARCH and $t$ innovations and the estimation of volatility are in \cref{ch13:sec-var,ch:volatility}. This chapter does not repeat them; it treats what the 2013 lecture adds: the \emph{three methodologies} (variance--covariance, historical simulation, Monte Carlo) side by side, the practitioner's choices behind each, decompositions of VaR by position, non-linear positions, and the practice of back-testing and stress testing. Sections or paragraphs marked ``(not discussed in class)'' are our additions.
\end{sources}

\section*{Motivation}
A trading firm holds thousands of positions and must be able to answer, every morning and in a few numbers, ``how much can we lose tomorrow?'' Ken Abbott, who helped set up the risk group of a large bank in the mid-1980s and later became the operating officer of firm risk management at Morgan Stanley, describes the job of risk management as making sure that management knows the risk profile of the firm, protecting it from unacceptable concentrations of risk, and ensuring ``no surprises'' \ts{13}{7}{2:05}. He uses two analogies: a \emph{spotlight}, which points at any risk one wants to look at, and a \emph{colouring book}, in which one then decides whether what is lit is good, bad, too big or too small \ts{13}{7}{5:22}. The second half of the job is communication: the chairman is often not a quant, so a complex idea has to be told in plain prose and simple pictures \ts{13}{7}{3:11}.

Value at risk (VaR) is the number that came out of this need. The methods are elementary statistics dressed in linear algebra:
\begin{itemize}
  \item a VaR is an \emph{order statistic} of the distribution of tomorrow's profit and loss (\cref{ch19:sec-one});
  \item positions in bonds are mapped to yield sensitivities before they enter (\cref{ch19:sec-fi});
  \item under normality the portfolio VaR is $z\sqrt{\bm x\T\Sigma\bm x}$, the \emph{variance--covariance} method (\cref{ch19:sec-vc});
  \item the same order statistic can be read off historical data (\emph{historical simulation}, \cref{ch19:sec-hs}) or off simulated data drawn with the covariance structure of the past (\emph{Monte Carlo}, \cref{ch19:sec-mc});
  \item all three approximate the P\&L of non-linear positions (\cref{ch19:sec-nl}), and all three are checked against realised P\&L by \emph{back-testing} and completed by stress tests (\cref{ch19:sec-bt}).
\end{itemize}
The physicist will recognise most of the mathematics: a quadratic form in a covariance matrix, a diagonalisation, a random-number generator. What is new is the vocabulary and the long list of practical choices, each of which changes the number. Abbott's recurring warning is that whenever a modelling assumption moves the answer materially, ``sirens should go off'' \ts{13}{7}{57:48}.

% ==================================================================
\section{One asset: VaR as an order statistic}\label{ch19:sec-one}
Throughout, VaR is measured on the \emph{loss} $L=-\Delta P$ over a horizon that in this lecture is one day, as in \cref{ch08:def-vares}: for a level $\beta$ (typically $99\%$ or $95\%$), $\mathrm{VaR}_\beta$ is the loss that is exceeded with probability $1-\beta$.
Abbott states it in words: the worst $1\%$ of the possible outcomes \ts{13}{7}{9:34}. The object is an \emph{order statistic}, ``the maximum or the minimum of a set of observations'' in the words of the professor \ts{13}{7}{8:29}, here an intermediate one: the value below which a fraction $1-\beta$ of the ordered outcomes lies. Order statistics have a large literature, so once VaR is recognised as one, its distribution and its estimation error are known \ts{13}{7}{9:34}, as \cref{ch19:sec-hs} will use.

\subsection{Returns, not levels}
The starting point is a time series of prices, say the Mexican IPC stock index over 1995--96 (slide~7). Prices are close to a random walk (\cref{ch:sp1}): the best forecast of tomorrow's level is today's level, the series is not stationary, and a histogram of \emph{levels} is a meaningless multimodal blob that depends on where the walk happened to wander \ts{13}{7}{12:48}. The histogram of one-day \emph{returns} (percentage changes, log changes or, in some markets, absolute changes) is unimodal and roughly symmetric, with the familiar bell shape (slides 8--9) \ts{13}{7}{16:00}. This is the whole reason for working with returns: they are the increments of the walk, approximately stationary with a mean and a variance that do not drift, and it is for them that a distributional assumption can be made. It is the same step from a Brownian motion to its increments that the physicist takes when computing a diffusion coefficient.

\begin{proposition}[One-asset variance method]\label{ch19:prop-one}
Let a position of value $V$ have a one-day return $r\sim N(0,\sigma^2)$. Then for $\beta\in(0,1)$
\[
  \mathrm{VaR}_\beta=z_\beta\,\sigma\,V,\qquad z_\beta=\Phi^{-1}(\beta),
\]
with $z_{0.99}=2.326$ and $z_{0.95}=1.645$. The recipe is: collect prices, compute returns, estimate the variance $\hat\sigma^2=\frac1{n-1}\sum_i(r_i-\bar r)^2$, take the square root, multiply by $z_\beta V$ (slide~14) \ts{13}{7}{30:40}.
\end{proposition}
\begin{proof}
$L=-Vr\sim N(0,\sigma^2V^2)$ and $\Prob(L\le z_\beta\sigma V)=\Phi(z_\beta)=\beta$. This is \cref{ch08:prop-normal} with $\mu=0$; the mean is dropped because for a one-day horizon it is small against $\sigma$.
\end{proof}
The multipliers are \emph{one-sided}: the $99\%$ VaR leaves $1\%$ in a single tail, hence $2.33$ and not the $1.96$ of a two-sided $95\%$ confidence interval, which leaves $2.5\%$ in each tail \ts{13}{7}{10:38}. A student asks about short positions; in this model the answer is that the same number serves for a long or a short position because the density is symmetric, but ``getting your signs right'' is one of Abbott's recurrent themes \ts{13}{7}{24:21}. He also uses the normal quantile as a heuristic on real, skewed data: ``what is a two or three standard-deviation move?'' means something to a manager even when the empirical distribution is lopsided \ts{13}{7}{24:21}.

\begin{example}[Mexican IPC, slides 12--13]\label{ch19:ex-ipc}
The variance of the daily IPC returns from January 1995 to December 1996 is $0.000324$, so $\hat\sigma=0.018012=1.8012\%$. Then $2.33\times1.8012\%=4.1968\%$ (with $z_{0.99}=2.3263$ one gets $4.190\%$): on a position of MXP~$100$ the loss exceeds MXP~$4.2$ only about $1\%$ of the days \ts{13}{7}{22:18}.
In fact the IPC lost more than $4.2\%$ on $8$ days of the sample, ``about $1.5\%$'' of the days, which Abbott reads as a sign of fat tails: large moves are more likely than the normal law implies \ts{13}{7}{26:29}.
\end{example}
\begin{remark}[Eight exceptions are not evidence of fat tails; our computation]\label{ch19:rem-eight}
The sample has about $520$ daily returns (the $8$ exceedances are $1.5\%$ of about $530$). If the normal VaR were exactly right, the number of exceedances would be $\mathrm{Bin}(520,0.01)$ with mean $5.2$ and standard deviation $2.3$, and $\Prob(N\ge8)=15\%$. The likelihood-ratio test of \cref{ch19:sec-bt} gives $p=0.25$. The data are compatible with the normal model at the $99\%$ level: eight exceedances against five are noise. The heavy tails of index returns are real (they are visible in the higher-frequency data of \cref{ch:volatility}), but this count does not show it, which is a small lesson in the very point Abbott stresses next.
\end{remark}

\subsection{Everything is estimated with error}
``Every time you estimate something you estimate it with error'' \ts{13}{7}{17:01}: the \emph{risk is $10$} means it is neither $10\,000$ nor $0$ \ts{13}{7}{18:03}. The sample variance follows a chi-squared law: for i.i.d.\ normal returns $(n-1)\hat\sigma^2/\sigma^2\sim\chi^2_{n-1}$ \ts{13}{7}{19:07} (\cref{ch03:sec-chisq}). Hence it cannot be negative, has a long right tail, and by \cref{ch13:sec-sample} the relative standard error of $\hat\sigma$ is $1/\sqrt{2(n-1)}$: $3.1\%$ for the $520$ observations above, $4.5\%$ for one year of data ($n=252$). The VaR multiplies $\hat\sigma$, so it inherits this error before any model error. Abbott distinguishes the \emph{standard deviation} of a data set from the \emph{standard error} of an estimate \ts{13}{7}{20:10}. The empirical quantile has a larger error still (\cref{ch19:prop-quantile}).

\begin{coursenote}[Spoken claim on the frequency of large moves]\label{ch19:cn-sigma}
Warning against extrapolating the normal law, Abbott says that a three-standard-deviation move occurs ``once every $10\,000$ observations'' \ts{13}{7}{25:25}. For a normal variable $\Prob(|Z|>3)=0.27\%$, i.e.\ one day in $370$ (one in $741$ for one tail), while $\Prob(|Z|>4)=6.3\times10^{-5}$ is one in $15\,800$ (one in $31\,600$ for one tail). The figure quoted lies between the two, and is presumably meant for four standard deviations or for a mixture of both. The remark about eight standard deviations is right in substance: $\Prob(|Z|>8)=1.2\times10^{-15}$, one day in $8\times10^{14}$, about $3\times10^{12}$ years of daily data. Under a normal model such a move ``does not happen'', so when it happens it is the model that has failed, not the market that was unlucky.
\end{coursenote}

\begin{finance}[Clean your data]\label{ch19:fin-data}
``Step 2b: graph your data'' \ts{13}{7}{29:37}. Everybody's data contain errors, missing values, stale quotes, and Abbott's rule is that one may be fired for incompetence but not for using someone else's bad data. Sources named on the slides are the Federal Reserve (the H.15 series for yields and H.10 for exchange rates, free and going back decades), Bloomberg (the cleanest in his experience, precisely because every user reports errors) and commercial vendors (slide~46) \ts{13}{7}{51:31}. The same applies to software: he tells the story of a programmer who claimed that a matrix square root multiplied by itself reproduced the covariance matrix ``to $16$ decimals'', when the code was in fact passing a float to a fixed-point variable \ts{13}{7}{43:13}. The other 2013 case studies (\cref{ch:regression,ch:timeseries}) do the same first step in R.
\end{finance}

\begin{exercise}[not from the course: estimation error of a quantile]
Using \cref{ch19:prop-quantile}, find the number of days of history needed to estimate the $99\%$ VaR of a normal loss with a relative standard error of $5\%$; and the same for the $99.9\%$ VaR. Repeat for a Laplace loss (density $\frac1{2b}e^{-|x|/b}$, quantile $b\ln\frac1{2(1-\beta)}$). (\emph{Answer for the first part:} about $1\,000$ days.)
\end{exercise}

% ==================================================================
\section{Bonds: from yield changes to price changes}\label{ch19:sec-fi}
Equities, currencies and commodities are quoted in prices, and the previous section applies directly. The time series available for bonds are \emph{yields} (slide~16) \ts{13}{7}{31:42}. To turn a yield move into a money move one needs the slope of the price--yield curve at the current yield: a ``poor man's Jacobian'', in Abbott's words, valid for small moves because the curve is non-linear and one uses its tangent \ts{13}{7}{31:42}.

\begin{definition}[PV01 (DV01, PVBP)]
The \emph{PV01} of a position is the change of its value for a rise of the yield by one basis point ($1\,\mathrm{bp}=0.01\%$). Its sign is negative for a long bond. For a bond of price $P$ and modified duration $D$ (\cref{prop:duration}), $\mathrm{PV01}\approx-D\,P\times10^{-4}$.
\end{definition}
If $\sigma_{\mathrm{bp}}$ denotes the standard deviation of the daily yield change in basis points, the P\&L of the position is $\approx\mathrm{PV01}\cdot\Delta y_{\mathrm{bp}}$ and the variance method of \cref{ch19:prop-one} gives
\begin{equation}\label{ch19:eq-fi}
  \mathrm{VaR}_\beta=z_\beta\,\bigl|\mathrm{PV01}\bigr|\,\sigma_{\mathrm{bp}} .
\end{equation}
Abbott's slides express the yield volatility as the volatility $\sigma_y$ of the \emph{relative} changes $\Delta y/y$ (which is the ``percentage change'' of a yield): then $\sigma_{\mathrm{bp}}=100\,y_{\%}\,\sigma_y$ with $y_\%$ the yield in percent. Slide~17 writes this as position $\times$ PV01 $\times$ close $\times\ \sigma_{\mathrm{yield}}\times2.33\times100$ \ts{13}{7}{33:48}. For equities, currencies and commodities the analogue is position $\times\ \sigma_{\mathrm{price}}\times2.33$.

\begin{example}[USD 100 of a 10-year Treasury, slides 18--22]\label{ch19:ex-ust}
The bond has coupon $4.644\%$ (semi-annual), ten years to maturity and yield equal to the coupon, so that its price is $100$ (slide~21). Repricing at $4.654\%$ gives $99.920765$, hence $\mathrm{PV01}=-0.07923$ per $100$ (slide~19; our recomputation gives $-0.079235$). The volatility of the relative changes of the yield is $1.312\%$ per day. The daily yield volatility is therefore $4.644\times100\times0.01312=6.09\,\mathrm{bp}$, and
\[
  \mathrm{VaR}_{99\%}=2.33\times0.07923\times6.093=1.1248\ \text{per }100\text{ of face value}
\]
(slide~22: $1.12479$). The $99\%$ one-day loss is thus $1.12\%$ of the position: a daily yield volatility of $6.1$~bp times a duration of $7.9$ years is $0.48\%$ of price, and $2.33\times0.48\%=1.12\%$.
\end{example}

\begin{coursenote}[Erratum: slides 18--22, the ten-year example]\label{ch19:cn-slide22}
\begin{itemize}
  \item Slide~22 writes $\mathrm{VaR}=-.07923\times0.04644\times.01312\times2.33\times100=1.12479$. With the yield entered as $0.04644$ the product is $0.011248$, not $1.12479$: the factor $100$ that converts percent points of yield to basis points has to be $10^4$ if the yield is a fraction, or the ``close'' has to be entered in percent ($4.644$), as the value $1.12479$ shows and as slide~17 requires.
  \item Slide~21 says that yield and coupon should be divided by $100$ and writes ``$4.644\%=.0644$''; it should read $.04644$.
  \item The volatility on slide~18 is for the period ``since $1/1/03$'' and the DV01 is computed ``at the $12/31/96$ yield'', while the Excel extract of slide~21 uses the date $12/27/2006$: the dates are inconsistent (presumably $1/1/03$ to $12/27/06$, and $12/31/06$).
  \item The cell labelled ``Duration @ Yield=Coupon'' (slide~21) is $7.9273$; that is the \emph{modified} duration, which is what the spreadsheet function returns. The Macaulay duration of this bond is $8.111$ years, and $8.111/(1+0.04644/2)=7.927$. The definition on slide~20 (``weighted average time to cash flows'') is that of the Macaulay duration.
\end{itemize}
\end{coursenote}

\begin{finance}[PV01, duration and units]\label{ch19:fin-pv01}
Duration is the (present-value weighted) average time of the cash flows, measured in years; the PV01 is a money amount per basis point per unit of face value. A duration of $7.9$ years means a PV01 of about $\$790$ per $\$1$~million of face value, ``roughly the same digits but different units'' \ts{13}{7}{33:48}. Bond traders think in PV01, insurance companies and portfolio managers who must match long-dated liabilities think in duration \ts{13}{7}{34:51}. ``Be careful of the units: this is one of the easiest errors to make'' (slide~20) is the most repeated sentence of the lecture. The market quotes yields, so risk is quoted per basis point: nobody computes $(702-701)/702$ in his head \ts{13}{7}{34:51}. See \cref{sec:duration} and, for delta ladders by maturity bucket, \cref{ch06:sec-risk}.
\end{finance}

\paragraph{Credit spread.} For a corporate or asset-backed bond the yield is a risk-free rate plus a credit spread (slide~24: $6\%=4.5\%+1.5\%$ in Abbott's example). One separates the interest-rate PV01 from the spread PV01 (CSPV01) and treats the two as different risk factors with their own volatilities and correlations; slide~23 states that the two coincide when the recovery rate is zero and are otherwise different, and lists calculated spreads, CDS and asset-swap spreads as sources \ts{13}{7}{36:52}. At spreads of $1\,000$--$1\,500$ bp the spread is mainly a default probability, and models of default, loss given default and recovery take over from the smooth sensitivity used here.

\begin{exercise}[not from the course: PV01]
A bond has a semi-annual coupon of $3\%$, five years to maturity, and yield $3\%$. Compute its price at yields $3\%$ and $3.01\%$, the PV01 per $100$ face value, the modified and Macaulay durations, and the one-day $99\%$ VaR of a position of face value $10^7$ if the daily volatility of the relative changes of the yield is $1.5\%$. Compare $\mathrm{PV01}$ with $-D_{\mathrm{mod}}\,P\cdot10^{-4}$.
\end{exercise}

% ==================================================================
\section{The variance--covariance method}\label{ch19:sec-vc}
\subsection{Covariance, correlation and the quadratic form}
For two assets with returns $a,b$ the sample covariance is $\hat\sigma_{ab}=\frac1{n-1}\sum_i(a_i-\bar a)(b_i-\bar b)$ (slide~25) and the correlation is $\hat\rho_{ab}=\hat\sigma_{ab}/(\hat\sigma_a\hat\sigma_b)$. Covariance says by how much and in which direction $b$ moves when $a$ moves away from its mean; it is not unit-free (measure the fertiliser in tons or in ounces and it changes), whereas correlation is a pure number in $[-1,1]$, ``an index of linearity'' \ts{13}{7}{39:01}. If the covariances are known the correlations follow; the converse needs the variances too: $\Sigma=D\,R\,D$ with $D=\diag(\sigma_1,\dots,\sigma_n)$ and $R$ the correlation matrix \ts{13}{7}{41:07}. Both matrices are symmetric; $\Sigma$ has the variances on the diagonal (the covariance of a variable with itself) and $R$ has ones. Abbott convinces himself of $\Var(a+b)=\Var a+\Var b+2\Cov(a,b)$ on a spreadsheet: for the numbers of slide~27, $0.333059+0.469281+2\times0.349660=1.501660$, equal to the variance of $a+b$ computed directly \ts{13}{7}{42:10}. (This is an algebraic identity of the sample covariance, as $\mathrm{Cov}$ is bilinear, so the check tests the spreadsheet rather than the statistics, which is the point of checking.)

\subsection{Portfolio VaR}
With several assets one writes the variance of a combination with $x$ units of $a$ and $y$ units of $b$ as $x^2\Var a+y^2\Var b+2xy\Cov(a,b)$ and so on for three, four, $n$ assets (slides 31--32), a pattern that becomes unmanageable ``very fast'' unless one uses matrices \ts{13}{7}{45:18}. Sum all the entries of the covariance matrix and one gets the variance of the portfolio holding one unit of each asset (slides 35--36); the general statement is \cref{ch03:prop-portvar}.

\begin{proposition}[Variance--covariance (delta-normal) VaR]\label{ch19:prop-vc}
Let $\bm x\in\R^n$ be the \emph{position vector}, $x_i$ being the P\&L of the portfolio per unit return of risk factor $i$ (in money units; for a bond it is the PV01 and the factor is the yield change in basis points), and let the vector of one-day factor changes $\bm r$ be $N(\bm 0,\Sigma)$. Then $\Delta P=\bm x\T\bm r\sim N(0,\sigma_P^2)$ with
\[
  \sigma_P^2=\bm x\T\Sigma\bm x=\sum_{i,j}x_ix_j\Sigma_{ij},\qquad
  \mathrm{VaR}_\beta=z_\beta\sqrt{\bm x\T\Sigma\bm x}\ \ts{13}{7}{47:22}.
\]
Consequently $\mathrm{VaR}_\beta\le z_\beta\sum_i|x_i|\sqrt{\Sigma_{ii}}$, the sum of the stand-alone VaRs, with equality iff the P\&L contributions $x_ir_i$ are pairwise perfectly positively correlated.
\end{proposition}
\begin{proof}
A linear combination of jointly normal variables is normal with variance $\bm x\T\Sigma\bm x$ (\cref{ch03:prop-portvar,ch03:prop-mvn}); the level follows as in \cref{ch19:prop-one}. The bound is the Cauchy--Schwarz inequality, $|\Sigma_{ij}|\le\sqrt{\Sigma_{ii}\Sigma_{jj}}$, applied to each term: $\bm x\T\Sigma\bm x\le\bigl(\sum_i|x_i|\sqrt{\Sigma_{ii}}\bigr)^2$.
\end{proof}
Written with the correlation matrix, $\sigma_P^2=\bm y\T R\bm y$ with $y_i=x_i\sigma_i$ the stand-alone one-standard-deviation P\&L of factor $i$: the standard deviations add in quadrature, corrected by the correlations. The physicist's reading: $\sigma_P$ is the length of $\bm x$ in the metric $\Sigma$; uncorrelated risks add like independent errors, perfectly correlated risks add like lengths, and negative correlations cancel.

\begin{example}[Three currencies, slides 35--42]\label{ch19:ex-fx}
Take positions of $-100$ in CAD, $-50$ in CHF and $-25$ in DEM (USD amounts, written in parentheses on slide~39; on the sign convention see the box below) and the covariance matrix of the daily changes
\[
  \Sigma=\begin{pmatrix}
    0.000037&-0.000018&-0.000017\\-0.000018&0.000321&0.000257\\-0.000017&0.000257&0.000227
  \end{pmatrix}.
\]
The daily volatilities are $0.61\%$, $1.79\%$ and $1.51\%$, and the correlations $\rho_{\mathrm{CAD,CHF}}=-0.17$, $\rho_{\mathrm{CAD,DEM}}=-0.19$, $\rho_{\mathrm{CHF,DEM}}=0.95$ (our computation; unsurprisingly the two European currencies move together). Then
\[
  \bm x\T\Sigma\bm x=1.691875,\qquad\sigma_P=1.3007,\qquad2.33\,\sigma_P=3.0307
\]
(slide~40; $2.3263\,\sigma_P=3.026$ with the exact quantile). The stand-alone VaRs are $1.417$, $2.087$ and $0.878$, adding to $4.382$: the diversification benefit is $1.35$, a third of the sum. Almost none of it comes from the imperfect correlation of CHF and DEM (with $\rho=0.95$ and positions of the same sign their two stand-alone VaRs add to $2.97$ against $2.94$ together, a benefit of $0.03$); the remaining $1.32$ comes from the negative correlation of CAD with both.
\end{example}

\begin{finance}[Signs, units and the position vector]\label{ch19:fin-signs}
Nearly every mistake in this method is a sign or a unit \ts{13}{7}{43:13}: ``like physics, only worse, because if you get the units wrong in physics you blow up, and here it just costs money''. A short gold position is entered as a negative number. The exchange rate dollar/yen is quoted in yen per dollar (say $95$): a dollar investor who owns yen loses when the quote rises to $100$, so the position enters as $-100$, whereas if the covariance matrix measures yen in dollars per yen ($0.0105$) the same position is $+100$ \ts{13}{7}{48:22}. For a bond, similarly, the PV01 of a long position is negative because the yield is the factor \ts{13}{7}{49:25}. Positions must all be in the same currency and in the same scale (units or millions), and the covariance matrix must use the same convention for returns (percent or log changes). Volatilities and correlations must be for the same differencing interval as the VaR horizon.
\end{finance}

\begin{exercise}[not from the course: two assets]
Two positions with daily volatilities $\sigma_1=1\%$, $\sigma_2=2\%$ (P\&L per unit of position value), correlation $\rho$, and values $V_1=100$, $V_2=50$ (long both).
\begin{enumerate}[(a)]
  \item Write the $99\%$ VaR as a function of $\rho$ using \cref{ch19:prop-vc}; for which $\rho$ is it equal to the sum of the stand-alone VaRs?
  \item Find the position $V_2$ (possibly negative) that minimises the VaR for given $V_1$ and $\rho$, and show that it is the minimum-variance hedge $V_2=-\rho\,(\sigma_1/\sigma_2)V_1$.
  \item Compute the component VaRs of the original portfolio for $\rho=0.5$ and check that they add up to the VaR.
\end{enumerate}
\end{exercise}

\subsection{Where does the risk come from? Absolute, incremental and component VaR}
A firm wants to attribute the portfolio VaR to desks or positions. The slides define (slides 59--62) the \emph{absolute VaR} of a bucket as its stand-alone VaR (the position vector with all other entries set to zero; equivalently its sensitivity times its volatility) and the \emph{marginal} (in today's language \emph{incremental}) VaR as the difference between the portfolio VaR and the VaR with that bucket removed, to see ``natural hedges'' and the effect of hedge portfolios \ts{13}{7}{49:25}. For the example: removing the CAD position lowers VaR by $0.096$ only, removing CHF by $1.508$ and removing DEM by $0.709$. These numbers add up to $2.31$, not to $3.03$: incremental VaRs depend on what else is in the book and are not additive. The slides later (77--80, see \cref{ch19:sec-hs}) use another notion, the \emph{component} VaR, which \emph{is} additive.

\begin{proposition}[Component (Euler) VaR; not discussed in class]\label{ch19:prop-euler}
In the variance--covariance method $\mathrm{VaR}(\bm x)=z_\beta\sqrt{\bm x\T\Sigma\bm x}$ is positively homogeneous of degree one, so by Euler's theorem
\[
  \mathrm{VaR}=\sum_ix_i\,\frac{\partial\,\mathrm{VaR}}{\partial x_i}=\sum_iC_i,\qquad
  C_i=z_\beta\,\frac{x_i(\Sigma\bm x)_i}{\sqrt{\bm x\T\Sigma\bm x}}=z_\beta\,\frac{\Cov(x_ir_i,\ \bm x\T\bm r)}{\sigma_P}.
\]
\end{proposition}
\begin{proof}
$\partial\sigma_P/\partial x_i=(\Sigma\bm x)_i/\sigma_P$ and $\sum_ix_i(\Sigma\bm x)_i=\bm x\T\Sigma\bm x=\sigma_P^2$. Also $(\Sigma\bm x)_i=\Cov(r_i,\bm x\T\bm r)$.
\end{proof}
For the currency example $C=(0.425,\ 1.852,\ 0.753)$, adding to $3.031$; the CAD position, nearly uncorrelated with the rest, contributes little, and $\partial\mathrm{VaR}/\partial x_i$ is the change of VaR per unit added of position $i$. A position with negative $C_i$ is a hedge of the rest of the book. Component VaRs are what one compares with limits by desk; the stand-alone VaRs are larger and add to $4.382>3.031$.

\begin{exercise}[not from the course: incremental versus component VaR]
In the variance--covariance method, show that the incremental VaR of position $i$, $\mathrm{VaR}(\bm x)-\mathrm{VaR}(\bm x-x_i\bm e_i)$, is at most the component VaR $C_i$ of \cref{ch19:prop-euler}. (\emph{Hint:} $\bm x\mapsto\sqrt{\bm x\T\Sigma\bm x}$ is a convex function.) Verify it on the numbers of \cref{ch19:ex-fx}: $(0.096,1.508,0.709)$ against $(0.425,1.852,0.753)$. When are the two equal?
\end{exercise}

\subsection{The choices behind the covariance matrix}\label{ch19:sec-assump}
The lecture then goes through the assumptions that hide in $\Sigma$ (slides 45--56). We list them with Abbott's opinions.

\begin{itemize}
  \item \textbf{How much history.} Rule of thumb for time series: at least $100$ observations \ts{13}{7}{53:37}. Regulators require at least a year; firms use one, two or five years. Slide~47 adds that the window should be weighed against data availability, changes of pricing regime (Brazil~1995, the introduction of the euro, the 2008 crisis) and market anomalies that traders might exploit.
  \item \textbf{Weighting.} Equal weights, or weights that decay exponentially into the past: the covariance is
  \[
    \hat\sigma_{xy}=\frac{\sum_{i=0}^{n-1}\omega^i\,(x_i-\bar x)(y_i-\bar y)}{\sum_{i=0}^{n-1}\omega^i},
  \]
  with $i=0$ today, $i=1$ yesterday, and $\omega$ near $0.95$ \ts{13}{7}{55:42}. The weight on the $m$ most recent observations is $(1-\omega^m)/(1-\omega^n)\approx1-\omega^m$; for $\omega=0.94$, $m=76$ observations carry $99.1\%$ of the weight ($99\%$ needs $m=74.4$), as Abbott says \ts{13}{7}{56:44}. This is the RiskMetrics estimator; its effective sample size $(1+\omega)/(1-\omega)\approx32$ observations and its relation to GARCH are in \cref{ch13:sec-ewma,ch13:sec-igarch}. Abbott's opinion is that the choice of $\omega$ has a material effect on the measured volatility (slide~51 compares a rolling six-month window with $\omega=0.92,\dots,0.98$) and that any assumption with such an effect calls for extreme care: one can get whichever volatility one wants (``lies, damn lies, and statistics''). He prefers equal weights out of Occam's razor: not knowing how much more last week matters than three weeks ago, he does not pretend to \ts{13}{7}{53:37}. Slide~49 lists the case for the other side: it smooths the volatility as positions age and reflects a heuristic that recent information matters more.
  \item \textbf{Missing data.} Fill gaps by linear interpolation, the previous day's value, a Brownian bridge (a Monte Carlo between the two neighbours) or regression on other series \ts{13}{7}{54:41}; this is also what protects against non-positive-definite matrices (\cref{ch19:sec-mc}).
  \item \textbf{Update frequency.} Daily updating of $\Sigma$ is overkill; Abbott's firm updates weekly, ``and that is overkill, but tell that to my regulators'' \ts{13}{7}{55:42}. Slide~53: changes of volatility are common in short-term euro-deposit and Japanese rates, changes of correlation are hard to estimate.
  \item \textbf{Percentage or log changes.} $r_{\mathrm{pct}}=e^{r_{\log}}-1=r_{\log}+\tfrac12r_{\log}^2+\cdots$, so the two differ by second-order terms that matter only at the ends of the distribution: a $4\%$ log move is a $4.08\%$ percentage move. Abbott uses percentage changes for convenience, although prices are closer to lognormal; the log form is what one simulates, since it keeps prices positive \ts{13}{7}{30:40}. Even yields, which the lognormal model keeps positive, have gone negative (slide~54).
  \item \textbf{Differencing interval.} Daily data are the norm. One can scale up a daily covariance matrix, but slide~55 warns of asynchronicity (markets close at different times) and serial correlation, whose effect on the $\sqrt T$ rule is quantified in \cref{ch08:prop-autocorr}.
  \item \textbf{Bucketing.} A book cannot carry a risk factor per position; positions are mapped to a set of risk indices (equal-VaR, duration-weighted, hedge-based buckets, slides 52 and 69) at a loss of granularity that should not cost accuracy if done well.
  \item \textbf{Stability of correlations.} A one-day horizon makes correlations less ambiguous, and short horizons mean that linear approximations are less of a problem, but correlations tend to ``swing'' from neutral to directional when markets are under stress (slide~37): the diversification benefit of \cref{ch19:ex-fx} is what disappears first in a crisis.
\end{itemize}
Abbott's summary of the method is the flow chart of slide~44: market data from several sources; returns; tests (graphs, ``$2$ standard deviations''); covariance matrix; positions and PV01s in a vector; the multiplication $\bm x\T\Sigma\bm x$; and then portfolio VaR, absolute and marginal sub-portfolio VaRs and scenario analysis \ts{13}{7}{51:31}. Up to the late 1990s this was how it was done: in 1996 about $80\%$ of the European banks used the variance--covariance method \ts{13}{7}{50:25}. Twenty years earlier the calculation was squeezed into the $256$ columns of a spreadsheet \ts{13}{7}{48:22}.

% ==================================================================
\section{Historical simulation}\label{ch19:sec-hs}
Abbott skipped this part in class (``historical simulation is Monte Carlo without the Monte Carlo part'' \ts{13}{7}{6:26}); it is on slides 63--85, which we summarise.

\subsection{The method}
Historical simulation marks the \emph{current} portfolio to market using the changes of the risk factors observed on each day of a past window, ranks the resulting hypothetical P\&L numbers and reads off the order statistic (slide~63). With $n$ historical scenarios $\bm r_1,\dots,\bm r_n$ (vectors of factor returns, or of changes measured in standard deviations), current factor values $\bm S$ and a valuation function $V$,
\[
  \Delta P_j=V\bigl(\bm S\odot(\bm 1+\bm r_j)\bigr)-V(\bm S)\ \approx\ \bm x\T\bm r_j,\qquad
  \mathrm{VaR}_\beta=-\Delta P_{(k)},\quad k=(1-\beta)n,
\]
$\Delta P_{(k)}$ being the $k$-th smallest P\&L. The steps of slide~67 are: collect position data; map to risk indices; compute sensitivities (deltas, gammas); collect historical data; compute returns; shock the portfolio with the returns; compute P\&L; rank and take the order statistic. The horizontal layout of slide~75 keeps a P\&L column for each desk and for the total, so the desk VaR is the $k$-th worst of the desk column and the total VaR the $k$-th worst of the total column. The rank $k$ depends on the significance level and on the data available:
\begin{center}\footnotesize
\begin{tabular}{@{}lcccc@{}}
\toprule
history & 1 year & 2 years & 3 years & 4 years\\
observations $n$ & $250$ & $500$ & $750$ & $1\,000$\\
$99\%$: worst $k=0.01n$ & $2.5$ (2nd--3rd) & $5$ & $7.5$ (7th--8th) & $10$\\
$95\%$: worst $k=0.05n$ & $12.5$ & $25$ & $37.5$ & $50$\\
\bottomrule
\end{tabular}
\end{center}
(the $99\%$ row is slide~75; the $95\%$ row is ours). It is ``probably sensible'' to compute the $1\%$, $5\%$ and $10\%$ levels together, and to save the detail, because the P\&L at the tail can be driven by very few positions (slide~73). The method is \emph{nonparametric}: it uses no distribution, so skewness, fat tails and the actual dependence between factors, including credit-spread jumps, are in the data; it works equally for long and short positions; aggregation is trivial (add the P\&L columns) (slides 64 and 82).
The weaknesses (slide~65): it needs a lot of data, and history is missing exactly for new products; it assumes that the future is drawn from the same distribution as the window (stationarity) --- slide~66 calls it a limited Monte Carlo, and in statistical language it is a bootstrap in which the past scenarios are resampled with equal probability (slide~66); and the tail estimate rests on very few points. The last point deserves a proof.

\begin{proposition}[Standard error of an empirical quantile; not discussed in class]\label{ch19:prop-quantile}
Let $L_1,\dots,L_n$ be i.i.d.\ losses with distribution function $F$ and density $f$, continuous and positive at $q=F^{-1}(\beta)$, and let $\hat q_n$ be the empirical $\beta$-quantile. Then, as $n\to\infty$,
\[
  \sqrt n\,(\hat q_n-q)\ \Rightarrow\ N\Bigl(0,\ \frac{\beta(1-\beta)}{f(q)^2}\Bigr).
\]
\end{proposition}
\begin{proof}[Sketch]
Let $\hat F_n$ be the empirical distribution function. $n\hat F_n(q)\sim\mathrm{Bin}(n,\beta)$, so $\hat F_n(q)-\beta$ has variance $\beta(1-\beta)/n$ and, by the central limit theorem, is asymptotically normal. By definition $\hat F_n(\hat q_n)\approx\beta$, and $\hat F_n(\hat q_n)-\hat F_n(q)\approx F(\hat q_n)-F(q)\approx f(q)(\hat q_n-q)$ for large $n$ (the empirical process is uniformly close to $F$; this step is the Bahadur representation of the quantile). Hence $\hat q_n-q\approx-(\hat F_n(q)-\beta)/f(q)$, which has the stated variance.
\end{proof}
For the standard normal at $\beta=99\%$, $f(q)=\varphi(2.326)=0.0267$, so $\mathrm{sd}(\hat q_n)=0.236/\sqrt{n/250}$: $10\%$ of the VaR with one year of data, $5\%$ with four years ($n=1\,000$), $3\%$ with ten. A heavy tail is worse, because $f(q)$ is smaller: for the standardised Student $t_4$, $f(q)=0.0123$ and the error is more than twice as large. \cref{ch19:fig-hsnoise} shows the sampling distribution in a simulation.

\begin{figure}[ht]
\centering
\begin{tikzpicture}
\begin{axis}[width=0.9\linewidth,height=5.6cm,xmin=1.5,xmax=3.3,ymin=0,ymax=5.8,
  xlabel={estimated $99\%$ VaR of a $N(0,1)$ loss},ylabel={density},
  tick label style={font=\scriptsize},label style={font=\small},legend pos=north east,legend style={font=\scriptsize}]
\addplot[ideacol,thick,mark=none] coordinates {(1.53,0.0) (1.59,0.003) (1.65,0.02) (1.71,0.03) (1.77,0.087) (1.83,0.213) (1.89,0.437) (1.95,0.63) (2.01,1.013) (2.07,1.477) (2.13,1.677) (2.19,1.94) (2.25,1.9) (2.31,1.657) (2.37,1.467) (2.43,1.303) (2.49,0.997) (2.55,0.713) (2.61,0.373) (2.67,0.243) (2.73,0.193) (2.79,0.143) (2.85,0.067) (2.91,0.023) (2.97,0.03) (3.03,0.017) (3.09,0.007) (3.15,0.003) (3.21,0.003) (3.27,0.0)};
\addplot[fincol,thick,mark=none] coordinates {(1.77,0.003) (1.83,0.013) (1.89,0.07) (1.95,0.243) (2.01,0.577) (2.07,1.09) (2.13,1.697) (2.19,2.26) (2.25,2.58) (2.31,2.263) (2.37,2.187) (2.43,1.447) (2.49,1.0) (2.55,0.62) (2.61,0.33) (2.67,0.14) (2.73,0.08) (2.79,0.053) (2.85,0.01) (2.91,0.0)};
\addplot[thmcol,thick,mark=none] coordinates {(2.07,0.01) (2.13,0.213) (2.19,1.317) (2.25,3.587) (2.31,5.193) (2.37,4.117) (2.43,1.813) (2.49,0.363) (2.55,0.047) (2.61,0.007) (2.67,0.0)};
\legend{$n=250$,$n=500$,$n=2500$}
\draw[dashed] (axis cs:2.326,0) -- (axis cs:2.326,5.8);
\node[font=\scriptsize,anchor=west] at (axis cs:2.36,5.5) {true $2.326$};
\end{axis}
\end{tikzpicture}
\caption{Historical-simulation estimates of the $99\%$ VaR of a standard normal loss: sampling distribution over $5\,000$ independent samples of $n=250$, $500$ and $2\,500$ days (histograms with bin width $0.06$; estimator: linearly interpolated empirical quantile). Standard deviations $0.213$, $0.158$ and $0.073$ against the asymptotic $0.236$, $0.167$ and $0.075$ of \cref{ch19:prop-quantile}; means $2.258$, $2.284$ and $2.318$, so a short window also biases the VaR downwards. Simulation by the authors.}\label{ch19:fig-hsnoise}
\end{figure}

\subsection{Marginal VaR in historical simulation}
Slides 77--81 attribute a historical-simulation VaR along an organisational hierarchy. The VaR engine produces P\&L \emph{vectors}, one entry per date (here $1\,043$ dates, four years of history), and the vector of a parent unit is the sum of the vectors of its children, $y=\sum_ix_i$ (Firm $=$ Equities $+$ Fixed income).

\begin{proposition}[Variance--covariance ratio allocation]\label{ch19:prop-hsmarg}
Let $y=\sum_{i=1}^nx_i$ be the parent P\&L and $x_i$ the child P\&Ls, as random variables over the historical dates. Then
\[
\begin{gathered}
  \sigma_y^2=\sum_{i=1}^n\Cov(y,x_i),\qquad
  \mathrm{MVaR}_i=\frac{\Cov(y,x_i)}{\Var y}\,\mathrm{VaR}_y,\\
  \text{so that}\quad\sum_i\mathrm{MVaR}_i=\mathrm{VaR}_y .
\end{gathered}
\]
\end{proposition}
\begin{proof}
$\Var y=\Cov(y,y)=\Cov\bigl(y,\sum_ix_i\bigr)$ by bilinearity (slides 78--79 write $\sigma_y=\sum_iE[(y-\bar y)(x_i-\bar x_i)]/\sigma_y$). Divide by $\Var y$ and multiply by $\mathrm{VaR}_y$.
\end{proof}
Variance is used as a proxy for risk to split a number that is a quantile. In the variance--covariance model of \cref{ch19:prop-vc} the two are the same thing, $\mathrm{MVaR}_i$ being the component VaR of \cref{ch19:prop-euler}. In a general empirical setting the exact Euler contribution of VaR is a conditional expectation, $\E[x_i\mid y=\mathrm{VaR}]$, which is a noisy average of the very few scenarios near the quantile, and the ratio allocation is a smoothed substitute (a remark of ours). The example on slide~81 has firm VaR $81\,411$ with children marginals $2\,822$, $14\,273$ and $64\,316$ (sum $81\,411$); the second child's stand-alone VaR is $29\,990$ against a marginal of $14\,273$, and the first's $14\,481$ against $2\,822$: the diversifying units carry much less than their stand-alone VaR. The marginals are relative to the immediate parent in the hierarchy, and the five children of the largest unit sum to $69\,020$ against a stand-alone parent VaR of $72\,457$; the slide attributes the gap to a different ``golden run'' of the data, although the text claims that the children's marginals add to the parent's stand-alone VaR.

\subsection{The ratio of the $95\%$ to the $99\%$ VaR}
Slide~85 shows the ratio of the $95\%$ to the $99\%$ VaR of a portfolio over the decade 2001--2011 (the slide does not say whose), with the normal-law value $1.65/2.33$ marked in red. The ratio $z_{0.95}/z_{0.99}=0.7071$ is, to three digits, $1/\sqrt2$. Before 2008 the ratio fluctuates between about $0.63$ and $0.74$, around the normal value; from the end of 2008 it drops to $0.40$--$0.65$. The slide has no comment, but the ratio measures tail thickness: for heavier tails the $99\%$ quantile grows faster than the $95\%$ one. For a standardised Student $t_\nu$ (variance one) the ratio is:
\begin{center}\small
\begin{tabular}{@{}lccccccc@{}}
\toprule
$\nu$ & $3$ & $4$ & $6$ & $10$ & $30$ & $100$ & $\infty$\\
\midrule
$q_{95}$ & 1.359 & 1.507 & 1.587 & 1.621 & 1.640 & 1.644 & 1.645\\
$q_{99}$ & 2.622 & 2.649 & 2.566 & 2.472 & 2.374 & 2.340 & 2.326\\
$q_{95}/q_{99}$ & 0.518 & 0.569 & 0.618 & 0.656 & 0.691 & 0.702 & 0.707\\
\bottomrule
\end{tabular}
\end{center}
so that a ratio of $0.5$--$0.6$ corresponds to a tail like that of a $t_3$--$t_4$ (the table is ours; historical simulation adds sampling noise, \cref{ch19:prop-quantile}). The lesson, again ours: the ratio is a free diagnostic of whether the normal multiplier of \cref{ch19:prop-vc} can be trusted, and after 2008 it could not.

\begin{figure}[ht]
\centering
\begin{tikzpicture}
\begin{axis}[width=0.8\linewidth,height=5.2cm,xmode=log,xmin=2.8,xmax=120,ymin=0.5,ymax=0.73,
  xlabel={degrees of freedom $\nu$ of the standardised $t$},ylabel={$q_{95}/q_{99}$},
  tick label style={font=\scriptsize},label style={font=\small},xtick={3,4,6,10,30,100},xticklabels={3,4,6,10,30,100}]
\addplot[ideacol,thick,mark=*,mark size=1.5pt] coordinates {(3,0.518) (4,0.569) (5,0.599) (6,0.618) (8,0.642) (10,0.656) (15,0.674) (20,0.682) (30,0.691) (50,0.697) (100,0.702)};
\draw[dashed] (axis cs:2.8,0.7071) -- (axis cs:120,0.7071);
\node[font=\scriptsize,anchor=south west] at (axis cs:3,0.7085) {normal: $0.7071$};
\end{axis}
\end{tikzpicture}
\caption{The ratio of the $95\%$ and $99\%$ quantiles for standardised Student-$t$ losses as a function of the degrees of freedom (the observable of slide~85).}\label{ch19:fig-ratio}
\end{figure}

\begin{exercise}[not from the course: the $95/99$ ratio]
Show that for a Laplace distribution the ratio $q_{95}/q_{99}$ does not depend on the scale, and evaluate it: $\ln10/\ln50=0.589$. Compare with the table of \cref{ch19:sec-hs}: which Student-$t$ has the same ratio? What does the ratio become for a Pareto tail $\Prob(L>x)\sim x^{-\alpha}$, and how would you estimate $\alpha$ from the ratio observed on slide~85?
\end{exercise}

% ==================================================================
\section{Monte Carlo simulation}\label{ch19:sec-mc}
Instead of the closed form of \cref{ch19:prop-vc}, one can simulate: draw, say, $10\,000$ possible outcomes for tomorrow, reprice, sort, and take the $100$th worst for the $99\%$ level (the $500$th for $95\%$) \ts{13}{7}{1:00:53}. It is again an order statistic, but now of a sample that we generate, so the assumptions can be relaxed at will: the distribution need not be normal (a $t$, a customised law, mean reversion, jumps) \ts{13}{7}{1:01:58}. The method, named during the Manhattan Project after games of chance, needs only that the system be described by probability density functions (slide~86). Its financial uses (slide~87) are path-dependent products, the risk of convex portfolios in hedge analysis, portfolios with strong jump risk (insurance), credit exposure --- and VaR. The seven steps of slide~89, drawn as a diagram on slide~88, are:
\begin{enumerate}[1.]
  \item generate uniform random numbers;
  \item convert them to standard normals;
  \item build the covariance matrix;
  \item factor it;
  \item create correlated random shocks;
  \item use the shocks to reprice;
  \item sort the P\&L and read the order statistic.
\end{enumerate}

\subsection{Random numbers}
Truly random numbers cannot be produced by a computer; they come from an external source such as radioactive decay \ts{13}{7}{1:03:03}. What software provides are \emph{pseudo-random} numbers, the output of a deterministic algorithm that ``look independent and uniform'' (slide~90; see Knuth), and Abbott recalls, as a young man, seeing two computers print the same ``random'' sequence (the effect of an identical seed) \ts{13}{7}{1:04:09}. \emph{Quasi-random} (low-discrepancy) sequences impose uniformity on purpose, so that a finite draw does not leave a region empty; they are harder to implement (slide~90). Reproducibility is a virtue for an audit trail: fix the seed.

\paragraph{From uniform to normal.} A uniform variable $U$ on $(0,1)$ becomes a standard normal by $Z=\Phi^{-1}(U)$ (the probability integral transform of \cref{ch03:prop-pit}, slide~91) \ts{13}{7}{1:05:13}. Slide~91 also points to the Box--Muller method of \emph{Numerical Recipes}, which avoids inverting $\Phi$:
\[
  Z_1=\sqrt{-2\ln U_1}\,\cos(2\pi U_2),\qquad Z_2=\sqrt{-2\ln U_1}\,\sin(2\pi U_2)
\]
are independent $N(0,1)$ if $U_1,U_2$ are independent uniform (proof, not given in class: if $Z_1,Z_2$ are i.i.d.\ $N(0,1)$ then in polar coordinates the angle is uniform on $(0,2\pi)$, independent of the radius, and $\Prob(Z_1^2+Z_2^2>s)=e^{-s/2}$, so that $U_1=e^{-(Z_1^2+Z_2^2)/2}$ is uniform; invert). Abbott admits that today he trusts the built-in generators of Excel and MATLAB \ts{13}{7}{1:05:13}, and that ``I kind of violate my own rules'' \ts{13}{7}{1:05:13}.

\subsection{Correlated shocks}
Draw independent standard normals and give them the covariance structure of the market. Everything rests on a ``square root'' of $\Sigma$.

\begin{proposition}[Correlated normal shocks]\label{ch19:prop-shocks}
Let $\Sigma$ be a symmetric positive semidefinite matrix and $\bm Z\sim N(\bm0,I_n)$ a column vector. If $A$ is any real matrix with $AA\T=\Sigma$, then $\bm r=A\bm Z\sim N(\bm 0,\Sigma)$. Two standard choices:
\begin{enumerate}[(i)]
  \item \emph{eigenvalue decomposition}: $\Sigma=E\Lambda E\T$ with $E$ orthogonal and $\Lambda=\diag(\lambda_i)$, $\lambda_i\ge0$ (\cref{thm:spectral}); $A=E\Lambda^{1/2}$;
  \item \emph{Cholesky}: $\Sigma=LL\T$ with $L$ lower triangular, which exists for $\Sigma$ positive definite; $A=L$.
\end{enumerate}
The matrix $A$ is not unique: $AQ$ works for every orthogonal $Q$, and all choices give the same distribution.
\end{proposition}
\begin{proof}
$\bm r$ is a linear function of a Gaussian vector and $\Cov\bm r=A\,I\,A\T=\Sigma$, as in \cref{ch03:prop-mvn}. For (i), $AA\T=E\Lambda^{1/2}\Lambda^{1/2}E\T=\Sigma$.
\end{proof}
Written for a \emph{row} vector of shocks (as in a spreadsheet, where each row is a draw) the recipe is $\bm r\T=\bm Z\T\Lambda^{1/2}E\T$ \ts{13}{7}{1:10:40}. The lecture calls $\Lambda^{1/2}E\T$ the ``transformation matrix'', ``a prism'': uncorrelated white noise goes in and correlated noise comes out \ts{13}{7}{1:12:42}. For two assets the Cholesky factor is explicit,
\[
  L=\begin{pmatrix}\sigma_1&0\\ \rho\sigma_2&\sigma_2\sqrt{1-\rho^2}\end{pmatrix},\qquad
  r_1=\sigma_1Z_1,\quad r_2=\sigma_2\bigl(\rho Z_1+\sqrt{1-\rho^2}\,Z_2\bigr),
\]
which shows why $|\rho|\le1$ is necessary. Slides 94--95 list the trade-off: Cholesky is computationally simpler and used by several packages but requires positive \emph{definiteness}; the eigen-decomposition is more complex, more robust (it works for semidefinite matrices) and is what Abbott uses. He notes that this construction is the vaunted \emph{Gaussian copula}, which most people treat as a black box although anyone with a little statistics can open it \ts{13}{7}{1:07:21}. Precisely, the joint law of $\bm r$ is the multivariate normal, whose copula is the Gaussian copula with correlation matrix $R$; the same eigen-decomposition on the correlation matrix is the principal-component analysis of \cref{ch:pca}. Abbott adds in passing that the method ``can really only take two distributions'' \ts{13}{7}{1:18:06}; presumably he means the normal and the Student $t$, the two elliptical laws that share this construction (with one extra step, see \cref{ch19:cs-three}).

\begin{coursenote}[Erratum: slides 95 and 96]\label{ch19:cn-slide95}
\begin{itemize}
  \item Slide~95 states $R\Lambda^{1/2}E'\sim N(0,\Sigma)$ where $R$ is called an $n\times1$ vector of normals; the product as written needs $R$ to be a \emph{row} ($1\times n$) vector, as in the spreadsheet; for a column vector the formula is $E\Lambda^{1/2}R$.
  \item Slide~96 gives, as a test of correlations, $\hat\rho\sim N\bigl(\rho_0,(1-\rho_0^2)/n\bigr)$. The large-sample variance of a sample correlation from $n$ normal pairs is $(1-\rho_0^2)^2/n$, i.e.\ standard deviation $(1-\rho_0^2)/\sqrt n$; equivalently, Fisher's $\operatorname{artanh}\hat\rho$ is $N\bigl(\operatorname{artanh}\rho_0,1/(n-3)\bigr)$. For $\rho_0=0.95$ and $n=1\,000$ the standard deviation is $0.003$, not $0.01$.
  \item The variance test on the same slide, $\hat\sigma_o^2/\sigma^2\sim F(m-1,k-1)$, is a comparison of two independent sample variances from samples of sizes $m$ and $k$ (original and simulated data).
\end{itemize}
\end{coursenote}

\begin{exercise}[not from the course: correlated normals]
\leavevmode
\begin{enumerate}[(a)]
  \item Derive the Cholesky factor of $\Sigma=\begin{pmatrix}4&2.4\\2.4&9\end{pmatrix}$ and give the two-line simulation recipe.
  \item Show that $A=E\Lambda^{1/2}E\T$ (the symmetric square root) also satisfies $AA\T=\Sigma$ and relate it to $E\Lambda^{1/2}$ of \cref{ch19:prop-shocks}.
  \item For three assets with $\rho_{12}=\rho_{23}=\rho$, show that $\rho_{13}$ must lie in $[2\rho^2-1,1]$, and find the eigenvalues of the matrix of \cref{ch19:ex-psd}.
  \item Prove the Box--Muller formulas.
\end{enumerate}
\end{exercise}

\subsection{When the covariance matrix is not positive semidefinite}
A covariance matrix must be positive semidefinite, ``another of those high-school-maths things that suddenly matter'' \ts{13}{7}{1:08:33}. If different entries are estimated from different amounts of data (say $100$ observations for most pairs and $25$ for one of them), the estimated matrix can have a \emph{negative eigenvalue}; then $\Lambda^{1/2}$ is imaginary and the simulation cannot even start. Abbott reports that, in the 1990s, nobody he asked (including a chairman of a statistics department) knew why he was getting negative eigenvalues; the cure is to fill in the missing data so that every entry uses the same sample \ts{13}{7}{1:09:38}. A minimal example shows that it is not a matter of rounding.

\begin{example}[An inconsistent correlation matrix; not from the lecture]\label{ch19:ex-psd}
Let three assets have $\rho_{12}=\rho_{23}=0.9$ and $\rho_{13}=0$, each perfectly plausible if estimated separately:
\[
  R=\begin{pmatrix}1&0.9&0\\0.9&1&0.9\\0&0.9&1\end{pmatrix},\qquad\text{eigenvalues }1-0.9\sqrt2=-0.273,\ 1,\ 1+0.9\sqrt2=2.273 .
\]
For the direction $\bm v=(1,-\sqrt2,1)/2$, $\bm v\T R\bm v=-0.273$: a portfolio with a negative ``variance'', hence an imaginary VaR. In general $\rho_{12}=\rho_{23}=\rho$ forces $\rho_{13}\in[2\rho^2-1,\,1]$: indeed $\det R=-(\rho_{13}-1)(\rho_{13}+1-2\rho^2)$, which must be non-negative (the geometric statement is that three angles $\arccos\rho_{12},\arccos\rho_{23},\arccos\rho_{13}$ between unit vectors satisfy a triangle inequality). For $\rho=0.9$ the lower bound is $0.62$, and $\rho_{13}=0$ is impossible.
\end{example}
Other remedies than filling in the data exist (replacing the matrix by the nearest positive-semidefinite correlation matrix, or setting the negative eigenvalues to zero and rescaling); the slide only lists ``precision, negative eigenvalues, computation time, data storage, error audit trail'' as the potential problems (slide~97).

\subsection{Checking the simulation, and why simulate at all}
Abbott's test of his spreadsheet: $10$ series of yield changes (Treasury, AAA, AA, \dots) with about $800$ observations; a $10\times10$ correlation matrix, for instance $0.83$ between the government $10$-year and the AAA $10$-year; the transformation matrix; $1\,000$ rows of independent normals multiplied by it \ts{13}{7}{1:13:47}. The correlations of the simulated data are close to but not identical with the original ones, ``because of sampling error'': with $1\,000$ draws the standard deviation of an estimated correlation is $(1-\rho^2)/\sqrt{1\,000}$ (\cref{ch19:cn-slide95}), for example $0.02$--$0.03$ for $\rho=0.5$ \ts{13}{7}{1:17:02}. With a million draws the match is much closer. Slide~96 lists the ways of checking: re-create the covariance matrix from the simulated data; test the variances (an $F$ or $\chi^2$ test); test the correlations; Box's $M$ test for the equality of covariance matrices. ``It is difficult to make small mistakes'' (slide~96): a wrong factorisation gives grossly wrong correlations.

Why draw a million scenarios when the data set has a few hundred? Because VaR looks at the tail: with $100$ observations there is one point in the $1\%$ tail, which says little about how fast the distribution drops off, whereas a simulation can put a billion points there \ts{13}{7}{1:17:02}. The price is the assumption of the generating distribution (slide~98 lists ``multivariate normality, serial correlation, skewness, kurtosis'', and the need to adjust for mean reversion and jump diffusion). The simulation error of the estimated quantile is \cref{ch19:prop-quantile} with $n$ replaced by the number $N$ of draws, which is at our disposal: $N=10^4$ gives a standard error of $0.037$ standard deviations of the P\&L, $N=10^6$ one of $0.004$.

\begin{casestudy}[Three methods, one portfolio; our numerical experiment]\label{ch19:cs-three}
Take the three-currency portfolio of \cref{ch19:ex-fx} ($\sigma_P=1.3007$) and compute the one-day $99\%$ VaR in four ways. Simulations use \cref{ch19:prop-shocks} with the eigen-decomposition; the fat-tailed version replaces the normal shock by a multivariate Student $t_4$ with the \emph{same} covariance matrix, $\bm r=\sqrt{(\nu-2)/\nu}\ \bm A\bm Z/\sqrt{W/\nu}$ with $W\sim\chi^2_\nu$ independent of $\bm Z$ (a normal variance mixture, \cref{ch03:prop-mixture}).
\begin{center}\small
\begin{tabular}{@{}lccc@{}}
\toprule
method & $99\%$ VaR & standard error & $99.9\%$ VaR\\
\midrule
variance--covariance (\cref{ch19:prop-vc}), $z_{0.99}\sigma_P$ & $3.026$ & --- & $4.020$\\
Monte Carlo, normal, $N=10^4$ (mean of $300$ runs) & $3.021$ & $0.049$ & --- \\
Monte Carlo, normal, Cholesky instead, $N=10^4$ & $3.022$ & $0.045$ & --- \\
Monte Carlo, normal, $N=10^6$ (one run) & $3.026$ & $\approx0.005$ & --- \\
$t_4$ shocks, exact & $3.446$ & --- & $6.598$\\
$t_4$ shocks, Monte Carlo $N=10^4$ (mean of $300$) & $3.429$ & $0.102$ & --- \\
$t_4$ shocks, Monte Carlo $N=2\cdot10^6$ (one run) & $3.453$ & $\approx0.01$ & --- \\
\bottomrule
\end{tabular}
\end{center}
The normal Monte Carlo converges to the closed form, as it must (the closed form \emph{is} the limit); the standard error agrees with $0.037\,\sigma_P=0.049$ from \cref{ch19:prop-quantile}. The two factorisations are statistically indistinguishable. With the same covariance matrix but heavy tails the $99\%$ VaR is $14\%$ higher and the $99.9\%$ VaR $64\%$ higher: the variance is a summary that hides the tail, and only the simulation (or the history) can see it. Note also that the Monte Carlo error with $10^4$ draws, $0.05$, is already small against this model risk, $0.4$.
\end{casestudy}

% ==================================================================
\section{Non-linear positions}\label{ch19:sec-nl}
The variance--covariance method assumes that the P\&L is linear in the risk factors, $\Delta P=\bm x\T\bm r$. This is exact for stocks and currencies, and for bonds it holds through the PV01 for small moves (\cref{ch19:sec-fi}); for options it is the delta approximation. Slide~70 states the rule of thumb: ``for a daily VaR, delta is probably enough; gamma is needed for a long holding period or for volatile markets'', and that convexity is not a big problem for bonds. Historical simulation and Monte Carlo can do better: they need a way to compute the P\&L for each scenario, either by \emph{exact repricing} of every instrument with the full model (accurate, computationally heavy) or from a \emph{parametric representation} of the price (a Taylor expansion, or a grid of precomputed values interpolated between scenarios, ``pricing grids'', slide~84), which can be very accurate but where cross-partial terms create noise (slides 83 and 99). Slide~83 illustrates it with the P\&L of an option book as a function of a shock: exact curve, coarse grid, refined grid.

The following example, ours, shows that gamma can matter even for one day. Let the risk factor return be $x\sim N(0,1)$ in units of its daily standard deviation, and let the P\&L be
\[
  \Delta P(x)=\delta\,x+\tfrac12\Gamma\,x^2,\qquad\delta>0,\ \Gamma<0
\]
(a short-gamma position, for example a written option; $\delta$ is the P\&L of a one-sigma move to first order, and the units are those of $\delta$). The three approximations below use the same $\delta$ and $\Gamma$.

\begin{proposition}[Exact VaR of a delta--gamma position; ours]\label{ch19:prop-dg}
For $\Gamma<0$ the loss $L=-\Delta P=\tfrac12|\Gamma|x^2-\delta x$ satisfies $L\le\ell$ iff $x_-\le x\le x_+$ with $x_\pm=\bigl(\delta\pm\sqrt{\delta^2+2|\Gamma|\ell}\bigr)/|\Gamma|$. The exact $\mathrm{VaR}_\beta$ is the unique root $\ell$ of $\Phi(x_+)-\Phi(x_-)=\beta$. For $\delta=0$, $\mathrm{VaR}_\beta=\tfrac12|\Gamma|\,\chi^2_{1,\beta}$, with $\chi^2_{1,\beta}$ the $\beta$-quantile of the $\chi^2_1$ law ($6.635$ for $\beta=99\%$).
\end{proposition}
\begin{proof}
Solve the quadratic inequality; the loss is increasing in $\ell$ so the equation has one root. For $\delta=0$, $L=\frac12|\Gamma|x^2$ and $x^2\sim\chi^2_1$.
\end{proof}
\begin{center}\small
\begin{tabular}{@{}lcccc@{}}
\toprule
& delta-normal & delta--gamma, & exact $\mathrm{VaR}_{99\%}$ & exact $\mathrm{ES}_{99\%}$ \\
& $z_{0.99}\,\delta$ & normal moments & (\cref{ch19:prop-dg}) & (simulation) \\
\midrule
$\delta=1,\ \Gamma=-0.5$ & $2.33$ & $2.72$ & $3.68$ & $4.46$ \\
$\delta=0,\ \Gamma=-0.5$ & $0$ & $1.07$ & $1.66$ & $2.11$ \\
\bottomrule
\end{tabular}
\end{center}
The ``delta--gamma normal'' column matches the mean $\frac12\Gamma$ and variance $\delta^2+\frac12\Gamma^2$ of the P\&L to a normal law (see \cref{ch08:prop-simm} for the moments of such a quadratic form, and \cref{ch15:sec-greeks} for the Greeks). The delta-normal VaR is $37\%$ too low for the first position and simply zero for the delta-neutral one, which nevertheless loses $1.66$ one day in a hundred. Even the moment-matched normal misses by a quarter, because the loss is a skewed, bounded-below variable ($\tfrac12|\Gamma|x^2$ is a $\chi^2_1$), and so does not resemble a normal law in the tail; the exact quantile comes from the distribution of the quadratic form, or from simulation.

\begin{figure}[ht]
\centering
\begin{tikzpicture}
\begin{axis}[width=0.85\linewidth,height=6.2cm,xmin=-5,xmax=5,ymin=-12,ymax=4,
  xlabel={risk-factor move $x$ (standard deviations)},ylabel={P\&L},
  tick label style={font=\scriptsize},label style={font=\small},
  legend pos=south east,legend style={font=\scriptsize}]
\addplot[ideacol,thick,domain=-5:5,samples=2] {x};
\addplot[thmcol,thick,domain=-5:5,samples=60] {x-0.25*x*x};
\legend{delta approximation $x$,exact $x-\tfrac14x^2$}
\draw[dashed] (axis cs:-5,-2.326) -- (axis cs:5,-2.326);
\draw[dashed] (axis cs:-5,-3.679) -- (axis cs:5,-3.679);
\node[font=\scriptsize,anchor=south west] at (axis cs:-4.9,-2.326) {delta-normal $-2.33$};
\node[font=\scriptsize,anchor=north west] at (axis cs:-4.9,-3.679) {exact $-3.68$};
\end{axis}
\end{tikzpicture}
\caption{P\&L of the short-gamma position $\delta=1$, $\Gamma=-0.5$ against the risk-factor move. The tangent (delta) sees a loss of $2.33$ at the $99\%$ move $x=-2.33$; the true P\&L there is $-3.68$, which is also the true $99\%$ VaR (the other root $x_+=6.33$ of \cref{ch19:prop-dg} is negligible). For a delta-neutral position the loss is even in $x$ and large moves of \emph{either} sign lose.}\label{ch19:fig-gamma}
\end{figure}

\begin{exercise}[not from the course: gamma]
For the position $\Delta P=\delta x+\frac12\Gamma x^2$ with $x\sim N(0,1)$, $\delta=0$ and $\Gamma<0$, show that $\mathrm{VaR}_\beta=\frac12|\Gamma|\chi^2_{1,\beta}$, and that the expected shortfall is $\mathrm{ES}_\beta=\frac{|\Gamma|}{2(1-\beta)}\bigl(2z_\beta\varphi(z_\beta)+(1-\beta)\bigr)$, where $\varphi$ is the standard normal density and $z_\beta$ the $\frac{1+\beta}2$-quantile. Evaluate for $\Gamma=-0.5$, $\beta=99\%$ and compare with the value $2.11$ of the table. (\emph{Hint:} $\E[x^2\mid|x|>z]=1+z\varphi(z)/(1-\Phi(z))$.)
\end{exercise}

% ==================================================================
\section{Back-testing, stress tests and use of VaR}\label{ch19:sec-bt}
\subsection{Why a bank needs it}
The first paragraph of the motivation has an accounting side that Abbott insists on \ts{13}{7}{13:55}. A balance sheet has assets on one side, financed by borrowing or by the firm's own money (equity); the ratio of the two is the leverage. Every loss reduces the equity, and when it reaches zero the firm is bankrupt. Risk management asks how large a one-day loss can be, with some level of certainty, so that the equity is enough to absorb it \ts{13}{7}{14:57}: \emph{that} is the capital that supports a position. VaR is the number that says how much is needed.

\begin{finance}[The uses of a VaR number]\label{ch19:fin-uses}
Position limits (``how much is too much?'' \ts{13}{7}{4:13}), capital, and the daily report to management. Abbott's caveats, all on slides: ``VaR numbers in isolation are worthless''; one must compute them every day and compare them (slide~74); the actual P\&L, or changes in the risk indices, must be used in the comparison; repricing is preferred to sensitivities; and VaR must be used together with \emph{stress testing} and \emph{scenario analysis}. Slide~76 is a diagram of the data flow: back-office systems provide positions and sensitivities, data providers provide history, and the VaR engine (the ``$1\%$ worst'') feeds limit reports, stress tests, historical scenarios and ad-hoc analyses. Historical scenarios (the 2008 crisis, Brazil in 1995, the introduction of the euro, slide~47) are the natural stress tests: revalue today's book with the moves of a known bad day, an ``automatic scenario analysis'' (slide~82). ``$90\%$ of risk management is knowing what you own and where you own it'' (slide~68).
\end{finance}

\subsection{Back-testing}
A \emph{back-test} compares the VaR computed at time $t$ with the realised P\&L at $t+1$ (slide~101). Reasons: it is required by regulators; it is a reality check on the calculation; it finds errors; it detects changes of the risk profile. The counting test (the binomial law of the number of exceedances, with rejection at $5\%$) is in \cref{ch08:sec-estimation}; we add its likelihood-ratio version, which is the one used in the literature.

\begin{proposition}[Kupiec's unconditional-coverage test; not discussed in class]\label{ch19:prop-kupiec}
If the VaR at level $\beta$ is correct and exceedances are independent, then in $n$ days the number $N$ of exceedances is $\mathrm{Bin}(n,p)$ with $p=1-\beta$. The likelihood-ratio statistic of $H_0\colon p=1-\beta$ against a free exceedance probability, $\hat p=N/n$,
\[
  \mathrm{LR}=-2\ln\frac{(1-p)^{n-N}p^{N}}{(1-\hat p)^{n-N}\hat p^{N}}
\]
is asymptotically $\chi^2_1$ under $H_0$ (Wilks' theorem), and one rejects when it exceeds $3.84$ ($5\%$) or $6.63$ ($1\%$).
\end{proposition}
Applications: the IPC example of \cref{ch19:rem-eight} ($N=8$, $n=520$): $\mathrm{LR}=1.31$, $p$-value $0.25$, not rejected. The $12$ exceedances in $500$ days of the counting test in \cref{ch08:sec-estimation}: $\mathrm{LR}=7.11$, $p=0.008$, rejected, consistent with the $0.5\%$ tail probability found there. Five exceedances in $250$ days: $\mathrm{LR}=1.96$, $p=0.16$: not rejected, although the expected count is $2.5$: a year of data cannot tell a good $99\%$ model from a moderately bad one.

\begin{finance}[The regulatory traffic light; not discussed in class]\label{ch19:fin-basel}
The Basel back-test of the 1990s uses $n=250$ days of $99\%$ VaR. With $p=1\%$, $\Prob(N\le4)=89.2\%$, $\Prob(N\le9)=99.97\%$ and $\Prob(N\ge10)=0.025\%$ (our computation), so the \emph{green zone} is $0$--$4$ exceptions, the \emph{yellow zone} $5$--$9$ and the \emph{red zone} $10$ or more; the capital multiplier applied to the VaR increases from $3$ in the green zone through the yellow zone to $4$ in the red zone (from memory of the 1996 amendment; the reader should check the figures against the Basel text). The 10-day, $99\%$ horizon is obtained from one-day figures by $\sqrt{10}$, with the caveats of \cref{ch08:prop-autocorr}. The later tightening to stressed VaR and expected shortfall is in \cref{ch08:sec-estimation}.
\end{finance}

\begin{exercise}[not from the course: back-test]
A desk reports a $99\%$ one-day VaR and records $7$ exceedances in $250$ days. Compute the likelihood-ratio statistic of \cref{ch19:prop-kupiec} and its $p$-value, the exact binomial tail probability $\Prob(N\ge7)$, and the Basel zone. Do the two tests agree? Why is a test with $250$ days not powerful against a true exceedance probability of $2\%$?
\end{exercise}

\subsection{Practical issues}
Slides 102--103 list the practical issues: back-tests are done down to the desk level, where every outlier must be documented (``need to document outliers''), but one often needs to aggregate, and aggregation hides spikes; \emph{split books} (a position whose P\&L is shared across desks) are a problem; whether one compares with front-office or back-office P\&L, the timing of P\&L recognition, re-organisations, data maintenance, the frequency of updates, and whether upside jumps should count as exceptions (a VaR is one-sided). Slides 104--105 (repeated as 113--114) plot the daily revenue of a business as bars with the $99\%$ and $95\%$ VaR as lines below the axis. In the first, the revenue stays within the lines except for a few isolated days, as it should; in the second the revenue is calm for a long period and then, after a break, the bars fall far outside the VaR lines, which adapt only slowly. The slides on the next pages (106--111, without axis labels or commentary in the lecture) show more series of this kind, of a P\&L band with a horizontal level: they illustrate that the scale of the P\&L can change abruptly (a position is cut, a new product starts from zero, a data error produces a spike) and that the VaR, the limit, or the data process reacts with a delay.

Slide~112 is the cleanest example of a regime change. It shows the spread of a five-year US credit-card asset-backed security over about four years ending in 2008: for three years it stays at about $0.3\%$, and its daily moves at a few basis points; from the summer of 2007 the spread rises to about $4\%$ and the daily moves reach tens of basis points, with the old percentile bands (dashed lines) ending up far inside the new distribution. A $99\%$ VaR estimated from an equally weighted window of one or two years, as Abbott prefers, must fail on such a series for the length of the window, and the exponentially weighted estimate of \cref{ch19:sec-assump} fails for about $1/(1-\omega)$ days: the price of simplicity, stated in Abbott's own terms.

\begin{coursenote}[What the lecture covered and what the slides add]\label{ch19:cn-scope}
Abbott states at the start that he will ``fly through'' the material, which is normally two hours of his university course \ts{13}{7}{6:26}, and announces that he will skip historical simulation to spend the time on Monte Carlo \ts{13}{7}{58:51}. The spoken lecture therefore corresponds to slides 1--62 (one-asset VaR, fixed income, variance--covariance with data issues and marginal VaR) and 86--97 (Monte Carlo); slides 63--85 (historical simulation, marginal VaR by variance ratio, the $95/99$ ratio) and 98--115 (the end of Monte Carlo, back-testing) are not discussed in the recording. He promises to send spreadsheets and papers by email to anyone who asks \ts{13}{7}{7:27}, and jokes, at the end, that the only use he has found for trigonometry is the seasonal decomposition of temperature data in sines and cosines \ts{13}{7}{1:20:10}. He teaches a longer version at NYU and other schools. The 2024 course has no counterpart lecture on these methods: the 2024 Lecture~10 (\cref{ch08:sec-riskmeasures}) takes the estimation of VaR and ES by historical simulation as its starting point and concentrates on the mathematics of risk measures, so the two chapters are complementary.
\end{coursenote}

\subsection{Abbott's principles, in brief}
\begin{intuition}[A physicist's summary]\label{ch19:int-summary}
A VaR is the location of a quantile of a distribution that one never sees and estimates from few data, using a Gaussian approximation whose parameters carry error and whose tails are wrong. The three methods differ in what they take as given: the variance--covariance method a linear map and a normal law (two parameters, a matrix); historical simulation nothing but the stationarity of the past (a sample); Monte Carlo a generating distribution (a model). Each buys a benefit at a cost: speed and transparency against tails; realism against sample size; flexibility against model risk. Abbott's practical advice reads like a laboratory notebook: use the simplest estimator that you understand, look at the data before you trust it, test every step (a matrix square root squared must return the matrix), record units and signs, and treat every number as an estimate with an error bar. Then keep a second, independent estimator, a stress test, to see what the first one cannot.
\end{intuition}

% ==================================================================
\section*{Summary of results}
\begin{keyformulas}
\begin{align*}
&\text{One asset (normal):} && \mathrm{VaR}_\beta=z_\beta\sigma V,\quad z_{0.99}=2.326,\ z_{0.95}=1.645\\
&\text{Bond:} && \mathrm{VaR}_\beta=z_\beta|\mathrm{PV01}|\,\sigma_{\mathrm{bp}},\quad\sigma_{\mathrm{bp}}=100\,y_{\%}\,\sigma_{\Delta y/y}\\
&\text{Portfolio:} && \mathrm{VaR}_\beta=z_\beta\sqrt{\bm x\T\Sigma\bm x},\quad\Sigma=DRD\\
&\text{Component VaR:} && C_i=z_\beta\,\frac{x_i(\Sigma\bm x)_i}{\sqrt{\bm x\T\Sigma\bm x}},\quad\sum_iC_i=\mathrm{VaR}\\
&\text{EWMA covariance:} && \hat\sigma_{xy}=\frac{\sum_i\omega^i(x_i-\bar x)(y_i-\bar y)}{\sum_i\omega^i},\quad\text{weight on last }m\text{ obs.}\approx1-\omega^m\\
&\text{Hist.\ simulation:} && \mathrm{VaR}_\beta=-\Delta P_{(k)},\ k=(1-\beta)n;\quad\mathrm{sd}(\hat q)\approx\frac{\sqrt{\beta(1-\beta)/n}}{f(q)}\\
&\text{Marginal VaR (hist.):} && \mathrm{MVaR}_i=\frac{\Cov(y,x_i)}{\Var y}\mathrm{VaR}_y,\quad\sum_i\mathrm{MVaR}_i=\mathrm{VaR}_y\\
&\text{Monte Carlo:} && \bm r=E\Lambda^{1/2}\bm Z\ \text{or}\ L\bm Z\sim N(\bm0,\Sigma),\quad\Sigma=E\Lambda E\T=LL\T\\
&\text{Box--Muller:} && Z_{1,2}=\sqrt{-2\ln U_1}\,\bigl(\cos,\sin\bigr)(2\pi U_2)\\
&\text{Delta--gamma, }\delta=0\text{:} && \mathrm{VaR}_\beta=\tfrac12|\Gamma|\,\chi^2_{1,\beta}\\
&\text{Kupiec:} && \mathrm{LR}=-2\ln\frac{(1-p)^{n-N}p^N}{(1-N/n)^{n-N}(N/n)^N}\sim\chi^2_1\\
&\text{Correlation error:} && \mathrm{sd}(\hat\rho)\approx(1-\rho^2)/\sqrt n
\end{align*}
\end{keyformulas}

% !TEX root = ../main.tex
\chapter{Regularized Pricing and Risk Models}\label{ch:regularized}
\begin{paracol}{2}

\begin{sources}
\begin{tabularx}{\linewidth}{@{}l l L@{}}
\toprule
\textbf{Lecture} & \textbf{Speaker} & \textbf{Content}\\
\midrule
2013 L10 (1 h 30) & I.~Masyukov (Morgan Stanley) & Guest lecture. Bonds, yield, duration, convexity; interest-rate swaps and the swap rate as an average of forward rates; building a yield curve with splines; bond spread over the swap curve; the cost of hedging a swap book; PCA hedging and its instabilities; a hedging model regularised by smoothness of the forward curve; the pricing/calibration engine and the HJM model; calibration of a two-dimensional volatility surface as an ill-posed inverse problem; SVD analysis, Tikhonov regularisation, truncated SVD; questions on splines, relative-value trading and bond curves.\\
\bottomrule
\end{tabularx}

\smallskip
\textbf{Files.} Slides \file{MIT18_S096F13_lecnote10.pdf} (51 pages, no incremental builds).
\textbf{Problem sets.} None. The exercises (each after the section it practises) are all ``(not from the course)''.
\textbf{Not published.} The talk has no 2024 counterpart. The slides give the numerical tables and figures
of the swap-book example, but no data or code; the toy computations of this chapter (a swap-curve
Jacobian, a one-dimensional calibration) are my own and are marked as such.
\end{sources}

\section*{Motivation}
Most quantities a bank computes every day are the output of an \emph{inverse problem}. A model has
parameters (a yield curve, a volatility surface); a handful of liquid instruments have observed prices;
the parameters must be chosen so that the model reproduces the prices, and only then can it be used to value and
hedge a book of thousands of trades that are \emph{not} among the calibration instruments. The forward map
``parameters $\mapsto$ prices'' is smooth and easy; its inverse is not. Small errors in the prices, or small
differences between instruments, are amplified into wild oscillations of the parameters, and the model that
``fits perfectly'' prices the rest of the book badly. The cure, the subject of this guest lecture, is
\emph{regularisation}: add to the fit a penalty that says what a sensible solution looks like, and give up an exact
fit in exchange for stability.

The chapter follows the lecture. \Cref{ch20:sec-bonds} recalls bonds, swaps and the yield curve. \Cref{ch20:sec-curve}
shows on a toy how calibrating a curve exactly already amplifies a 1~bp change into a ripple of several
basis points. \Cref{ch20:sec-hedge} treats the risk side: hedging with a few liquid swaps by PCA and by a smoothness-regularised
alternative. \Cref{ch20:sec-inverse} develops the mathematics behind both, singular values, Tikhonov regularisation and
truncated SVD, which is the natural continuation of the shrinkage methods of \cref{ch05:sec-shrink}; \cref{ch20:sec-lambda}
discusses how to choose the regularisation parameter (not covered in class). \Cref{ch20:sec-vol} applies everything
to the calibration of a volatility surface for the Heath--Jarrow--Morton model (\cref{ch:hjm}). The linear algebra needed is
the SVD of \cref{sec:svd}.

% ==================================================================
\section{Bonds, swaps and the yield curve}\label{ch20:sec-bonds}

\subsection{Bond price, yield, duration and convexity}\label{ch20:sec-yield}
A fixed-rate bond pays cash flows $C_1,\dots,C_N$ (coupons, plus the face value in $C_N$) at times $t_1<\dots<t_N$
\ts{13}{10}{2:11}. Cash paid later is worth less, so each cash flow is multiplied by a \emph{discount factor}
$\Lambda_i\in(0,1]$ that decreases with $t_i$ \ts{13}{10}{5:16}, and the fair value is
$P=\sum_{i=1}^N\Lambda_iC_i$ \ts{13}{10}{6:20}. For a bond the price $P$ is not something to be computed: it is the
result of a liquid market and is known \ts{13}{10}{7:20}. What is needed is a \emph{model for discounting}. The simplest one uses a single
parameter, the \emph{yield to maturity} $y$ (continuous compounding, one rate for all dates)
\ts{13}{10}{8:27}:
\begin{equation}\label{ch20:eq-yield}
  \Lambda_i=e^{-yt_i},\qquad P(y)=\sum_{i=1}^Ne^{-yt_i}C_i .
\end{equation}
With $C_i>0$ the map $y\mapsto P(y)$ is strictly decreasing, so price and yield determine each other one-to-one
\ts{13}{10}{9:31}. The yield is therefore not a traded quantity but a \emph{convention} to summarise price and
cash flows \ts{13}{10}{10:34}. Differentiating,
\[
  \frac{\dd P}{\dd y}=-\sum_it_ie^{-yt_i}C_i,\qquad
  D=-\frac1P\frac{\dd P}{\dd y}=\sum_{i=1}^Nt_i\,\pi_i,\quad \pi_i=\frac{e^{-yt_i}C_i}{P},\ \ \sum_i\pi_i=1,
\]
\[
  \mathcal C=\frac1P\frac{\dd^2P}{\dd y^2}=\sum_it_i^2\pi_i>0 .
\]
The \emph{duration} $D$ is the present-value-weighted average time of the cash flows (``weighted time'')
\ts{13}{10}{11:40}\ts{13}{10}{12:41}; normalising by $P$ makes it independent of how many bonds are held. A zero-coupon
bond has one weight, so $D=T$; a coupon bond has $D<t_N$, because every $t_i\le t_N$ and at least one is strictly smaller
\ts{13}{10}{13:47}. Convexity is a weighted mean of $t_i^2$, hence positive. These are the definitions of
\cref{sec:duration} in the continuous-compounding convention (\cref{prop:duration}), and \cref{eq:durconv} is the Taylor expansion
$\Delta P/P\approx-D\Delta y+\frac12\mathcal C(\Delta y)^2$.

\begin{coursenote}[Sign and letters in the slides]
Slide~8 defines ``duration'' as $d=\frac1P\frac{\partial P}{\partial y}$, which is \emph{negative} (the speaker says so:
the duration is always negative); here, as in \cref{sec:duration}, $D=-d>0$. Slide~9 names the second derivative $c=\partial^2P/\partial y^2$
and then writes its explicit form as ``$d=\sum t_i^2e^{-yt_i}C_i$'': the letter should be $c$, and this is the
\emph{unnormalised} second derivative, $\mathcal C=c/P$ in the notation used here.
\end{coursenote}

\begin{example}[Duration explains most of the P\&L, convexity the rest (my numbers)]\label{ch20:ex-bond}
Take $N=5$ annual payments $C=(4,4,4,4,104)$ and $y=4\%$. Then $P=99.640$, $D=4.629$ years, $\mathcal C=22.42$. If $y$ rises by $50$~bp the exact
price change is $-2.2868\%$; the duration term alone gives $-2.3146\%$ and adding convexity $-2.2865\%$. For $+200$~bp the figures are
$-8.824\%$ exactly, $-9.258\%$ with duration only and $-8.810\%$ with convexity. The trader who
explains the day's loss with first derivatives only \ts{13}{10}{15:56} is left with an ``unexplained'' term
of order $\frac12\mathcal C(\Delta y)^2$.
\end{example}

All of this assumes one $y$ for all dates, hence that the curve moves in \emph{parallel}. That was an acceptable
approximation before the crisis, but with a steep curve, where short rates sit near zero and long rates are expected to
rise, it is not \ts{13}{10}{14:55}. \Cref{ch20:sec-bondspread} replaces $y$ by a full curve.

\subsection{Interest-rate swaps}\label{ch20:sec-swaps}
In a swap one party pays a fixed rate $C$ and receives a floating rate (three-month LIBOR in the lecture, fixed on the
day and known in advance for the accrual period) \ts{13}{10}{17:02}. With payment dates $T_1<\dots<T_n$, accrual
fractions $\delta_i$ (day-count), discount factors $\Lambda_i$ and forward rates $r_i$ for the period ending at $T_i$,
\[
  PV_{\rm fixed}=\sum_iC\delta_i\Lambda_i=C\sum_iw_i,\qquad PV_{\rm float}=\sum_ir_i\delta_i\Lambda_i=\sum_ir_iw_i,\qquad w_i=\delta_i\Lambda_i,
\]
\ts{13}{10}{18:14}. A swap costs nothing to enter, so the fixed rate is set so that the two legs have the same value
\ts{13}{10}{19:21}:
\begin{equation}\label{ch20:eq-swaprate}
  C=\frac{\sum_ir_iw_i}{\sum_iw_i}\qquad(\text{same payment frequency on both legs}),
\end{equation}
the \emph{swap rate}, a weighted average of forward rates, which is what traders quote \ts{13}{10}{20:27}. If the forward
rates are defined by $\Lambda_{i-1}=(1+r_i\delta_i)\Lambda_i$ (the single-curve convention of \cref{sec:forward}, $\Lambda_0=1$), the floating leg telescopes,
$PV_{\rm float}=\sum_i(\Lambda_{i-1}-\Lambda_i)=1-\Lambda_n$, and $C=(1-\Lambda_n)/\sum_i\delta_i\Lambda_i$ (my remark, not on the slides). The fixed leg is a bond coupon stream, which is why a swap can be
hedged with a bond \ts{13}{10}{20:27}, and why swap risk is measured in \emph{buckets} of par rates: the trader wants to know
the P\&L if the 5-year swap rate moves by 1~bp.

\subsection{Building a curve}\label{ch20:sec-build}
A liquid market gives swap rates for a few maturities, whereas a book of hundreds of thousands of swaps needs a discount factor for every day up to 30--40 years
\ts{13}{10}{22:35}. Constructing the curve means \ts{13}{10}{23:39}:
\begin{enumerate}[1.]
  \item choose the input instruments;
  \item choose
      which quantity is interpolated (daily forward rates, zero rates, discount factors) and the interpolation family
      (piecewise-constant, linear, tension spline, cubic spline) with its \emph{knots};
  \item \emph{calibrate}: solve for the
      spline parameters so that every input instrument is repriced at par (PV of the swap equal to zero)
      \ts{13}{10}{24:42}\ts{13}{10}{28:58}. The lecture's example uses 15 swaps from 1 to 30 years with quotes from $0.33\%$ to $2.67\%$
      (the market of a year earlier; \cref{ch20:tab-risk}); the resulting curve of three-month forwards is far from flat, steep in the first five years
      and with a hump in the 20-year region \ts{13}{10}{30:01}.
\end{enumerate}

\paragraph{Splines and B-splines.}
A \emph{cubic spline} is a function that is a cubic polynomial on each interval between knots and has two continuous derivatives at every knot \ts{13}{10}{25:48}.
The knots in the lecture's picture are $1,10,20,40,80,160,240$ (quarters), denser where instruments are dense
\ts{13}{10}{1:22:09}. All splines with given knots form a vector space, which has a convenient basis of bell-shaped,
locally supported functions \ts{13}{10}{26:49}\ts{13}{10}{27:51}.

\begin{definition}[B-splines, Cox--de Boor recursion]\label{ch20:def-bspline}
For a non-decreasing knot sequence $\tau_0\le\tau_1\le\cdots$ let $B_{i,0}(t)=1$ if $\tau_i\le t<\tau_{i+1}$ and $0$ otherwise, and
\[
  B_{i,k}(t)=\frac{t-\tau_i}{\tau_{i+k}-\tau_i}B_{i,k-1}(t)+\frac{\tau_{i+k+1}-t}{\tau_{i+k+1}-\tau_{i+1}}B_{i+1,k-1}(t)
\]
(a term with zero denominator is dropped) \ts{13}{10}{1:08:13}. Degree $k=3$ gives cubic B-splines.
\end{definition}

\begin{proposition}[Properties of B-splines]\label{ch20:prop-bspline}
$B_{i,k}\ge0$; it vanishes outside $[\tau_i,\tau_{i+k+1}]$; $\sum_iB_{i,k}(t)=1$ on the interior of the knot range; for simple knots
$B_{i,k}\in C^{k-1}$. Hence every $f=\sum_ic_iB_{i,k}$ is a spline of degree $k$ and its values lie between the smallest and the largest $c_i$.
\end{proposition}
\begin{proof}
Non-negativity and support follow by induction from the recursion, whose coefficients are non-negative on the supports involved. For the partition of unity
shift the index in the second term of $\sum_iB_{i,k}$:
\[
  \sum_iB_{i,k}=\sum_iB_{i,k-1}\Bigl[\frac{t-\tau_i}{\tau_{i+k}-\tau_i}+\frac{\tau_{i+k}-t}{\tau_{i+k}-\tau_i}\Bigr]=\sum_iB_{i,k-1},
\]
and the case $k=0$ is $1$. Smoothness: at a simple knot each recursion step raises the differentiability by one, starting from a jump. (Checked numerically: the maximal deviation from $1$ is $4\times10^{-16}$.)
\end{proof}
Local support makes the calibration matrix banded and the response to a change of one parameter local; the convex-combination
property makes the parameters interpretable as ``values of the curve''. This is why the same basis is used for the surface, in two dimensions,
in \cref{ch20:sec-vol}.

\subsection{Bond spread over the curve}\label{ch20:sec-bondspread}
Once a swap curve $\Lambda_i$ exists, a bond can be priced off it with one extra parameter, the \emph{bond spread} $s$,
determined so that the model reproduces the (liquid, known) bond price \ts{13}{10}{31:01}:
\begin{equation}\label{ch20:eq-spread}
  P=\sum_{i=1}^Ne^{-st_i}\Lambda_iC_i .
\end{equation}
Compared with $y$ the correction is small: a curve level of $3\%$ and a spread of a few basis points, about a hundred times smaller
\ts{13}{10}{32:10}. The gain is consistency: sensitivities of the swap book to the curve inputs are inherited by the bond, and one is no longer bound to
the single ``parallel'' mode of duration. The spread contains credit or liquidity components of the bond, and even a
swap-versus-Treasury spread is traded; for a 10-year maturity the most liquid instrument is the bond, then, surprisingly, the spread, and only then the swap
itself, so the curve is built from bond yield plus spread rather than from the raw 10-year swap quote
\ts{13}{10}{33:10}\ts{13}{10}{34:11}\ts{13}{10}{35:12}.

\begin{finance}[Why the curve is built on swaps and not on bonds]
Asked why one does not invert the bond curve instead, the speaker gave two reasons \ts{13}{10}{1:27:24}\ts{13}{10}{1:28:29}. A swap
starting today always exists, and tomorrow's swap starts tomorrow, whereas bonds have a fixed maturity: the newest issue (\emph{on-the-run}) is by far the most liquid
and older ones (\emph{off-the-run}) trade at different levels, so there is no continuous family of maturities and a bond curve can only be a least-squares fit, not an exact one.
The switch of liquidity from one issue to the next creates ``roll effects'' that make the curve unstable. So the swap desk builds the swap curve and \emph{projects} bonds onto it, by the spread
\eqref{ch20:eq-spread}, since bonds are its preferred hedging instruments.
\end{finance}

% ==================================================================
\section{Calibrating a curve exactly is already ill-conditioned}\label{ch20:sec-curve}

The lecture shifts one input of the calibrated curve, the 9-year swap rate, by $1$~bp, keeping all the other 14 quotes fixed, and plots the response of the
three-month forward rates: it is a ripple, zero before about 5 years and after about 20 years, with a swing of about $\pm14$~bp around the 9-year point
\ts{13}{10}{36:16}\ts{13}{10}{37:22}. The curve is exact at every input, so it \emph{must} move in the neighbourhood of the shifted input (the 8- and 10-year par
rates are pinned), and it does so by an oscillation much larger than the shift. The speaker's comment: this is an ill-posed problem, where a
small change of the input causes a large change of the output.

\paragraph{A toy where everything is explicit (not in class).}\label{ch20:par-toy}
Let the instantaneous forward rate be constant on each interval $(T_{i-1},T_i]$, $T_0=0$, between consecutive input maturities, with value $f_i$, and let the par rate of the swap
maturing at $T_i$ be (ignoring discounting) the average forward rate, $S_i=\frac1{T_i}\sum_{k\le i}(T_k-T_{k-1})f_k$, i.e. $S=Mf$ with $M$ lower triangular
and invertible. Inverting,
\[
  f_i=\frac{T_iS_i-T_{i-1}S_{i-1}}{T_i-T_{i-1}},\qquad\text{so}\qquad J=M^{-1}:\quad \Delta f_i=\frac{T_i\,\Delta S_i-T_{i-1}\,\Delta S_{i-1}}{T_i-T_{i-1}} .
\]
A shift $\Delta S_j=1$~bp moves $f_j$ by $+T_j/(T_j-T_{j-1})$ and $f_{j+1}$ by $-T_j/(T_{j+1}-T_j)$: for the 9-year input, $+9$~bp on $(8,9]$ and $-9$~bp on
$(9,10]$; for 10 years $+10$ and $-5$; for 25 years $+5$ and $-5$. Each maturity acts as a lever with arm $T_j$ divided by the spacing
of the instruments, and no amount of care in the method removes it, because the data (an average up to $T_j$) constrain only a
difference of averages. For the 15 maturities of the lecture the condition number of $M$ is $23.4$ (my computation). \Cref{ch20:fig-ripple}
draws the two responses.

\begin{figure}[ht]
\centering
\begin{tikzpicture}
\begin{axis}[width=12cm,height=4.6cm,xlabel={forward time $t$ (years)},ylabel={$\Delta f(t)$ (bp)},xmin=0,xmax=30,ymin=-10.5,ymax=10.5,
  axis lines=left,grid=major,grid style={black!12},legend style={at={(0.5,-0.42)},anchor=north,legend columns=2,draw=none},tick label style={font=\small}]
\addplot[thick,thmcol] coordinates {(0,0)(8,0)(8,9)(9,9)(9,-9)(10,-9)(10,0)(30,0)};
\addlegendentry{$+1$~bp on the 9Y par rate}
\addplot[thick,defcol,dashed] coordinates {(0,0)(20,0)(20,5)(25,5)(25,-5)(30,-5)};
\addlegendentry{$+1$~bp on the 25Y par rate}
\end{axis}
\end{tikzpicture}
\caption{Forward-rate response to a $1$~bp shift of one par rate in the piecewise-constant toy (my computation): the curve calibrated exactly answers with a
ripple whose amplitude is the maturity divided by the spacing of instruments. The lecture's spline curve gives a similar but smoother ripple (\ts{13}{10}{36:16}).}\label{ch20:fig-ripple}
\end{figure}

\begin{remark}[Smoother is not automatically better (my computation)]
Replacing the piecewise-constant forwards by a cubic B-spline with 15 coefficients (interior knots by the averaging rule of de Boor) and again calibrating exactly
gives a condition number of about $6\times10^3$ for the calibration matrix, and the response to the same $1$~bp shift of the 9-year rate has amplitude of
between $10$ and $30$~bp at $9$--$10$~years, about $70$~bp at $12$, $150$~bp at $16$ and $200$~bp at $20$ years. The exact calibration of a few averaged quantities to a flexible smooth
curve is genuinely unstable in the sparse long end; the lecture's curve (ripple confined to $\pm14$~bp) shows that the design of the interpolation matters.
\end{remark}

\begin{exercise}[(not from the course)]
\leavevmode
\begin{enumerate}[(a)]
  \item For a swap with equally spaced payment dates show $C=(1-\Lambda_n)/\sum_i\delta_i\Lambda_i$ and that $C$ lies between the smallest and the largest forward rate.
  \item In the piecewise-constant toy of \cref{ch20:sec-curve} derive $\Delta f_i$ for a shift of the $j$-th par rate and show that its entries are $T_j/(T_j-T_{j-1})$ and $-T_j/(T_{j+1}-T_j)$; what is the largest absolute value over all $j$ for the 15 maturities of the lecture?
  \item The matrix $M$ is lower triangular: compute its determinant.
\end{enumerate}
\end{exercise}

% ==================================================================
\section{Hedging a swap book: PCA and the smoothness-regularised model}\label{ch20:sec-hedge}

\subsection{The cost of hedging bucket by bucket}\label{ch20:sec-cost}
The \emph{risk} of a book is the vector $\bm r\in\R^n$ of P\&L per basis point for a shift of each of the $n$ input rates (here $n=15$ swap maturities);
to first order the P\&L for a shift $\bm u$ of the inputs is $\bm r\T\bm u$. \Cref{ch20:tab-risk} shows the lecture's example \ts{13}{10}{37:22}. The
net risk is tiny ($2{,}266$ in total), but the risk in each bucket is large and of both signs. To be insensitive to \emph{every} bucket the trader must trade
$\pm r_i$ swaps in each of them, paying half the bid--offer of each instrument; the liquid maturities (yellow in the slide: 1, 2, 5, 10, 30 years) cost $0.10$~bp, the
others $0.25$~bp. The charge of a bucket is $|r_i|\times$ bid--offer, and the total is $\sum_i|r_i|b_i=3{,}617{,}629$
\ts{13}{10}{38:25}\ts{13}{10}{39:27}. The gross exposure, not the net one, is what costs money, so traders never hedge every bucket: they hedge a seven-year risk
with the cheaper liquid points around it \ts{13}{10}{40:30}. The question is how to do this systematically.

\begin{table}[ht]
\centering\small
\begin{tabular}{lrrrr}
\toprule
Swap & quote (\%) & raw risk & bid--offer (bp) & charge\\
\midrule
1Y & 0.33 & $-200{,}000$ & 0.10 & 20,000\\
2Y & 0.39 & 1,330,000 & 0.10 & 133,000\\
3Y & 0.49 & $-200{,}000$ & 0.25 & 50,000\\
4Y & 0.64 & 1,200,000 & 0.25 & 300,000\\
5Y & 0.86 & $-722{,}450$ & 0.10 & 72,245\\
6Y & 1.09 & $-35{,}255$ & 0.25 & 8,814\\
7Y & 1.29 & $-537{,}430$ & 0.25 & 134,358\\
8Y & 1.48 & $-3{,}850{,}000$ & 0.25 & 962,500\\
9Y & 1.64 & 1,580,000 & 0.25 & 395,000\\
10Y & 1.79 & 288,751 & 0.10 & 28,875\\
12Y & 2.04 & $-401{,}350$ & 0.25 & 100,338\\
15Y & 2.29 & 50,000 & 0.25 & 12,500\\
20Y & 2.50 & 4,000,000 & 0.25 & 1,000,000\\
25Y & 2.60 & $-1{,}000{,}000$ & 0.25 & 250,000\\
30Y & 2.67 & $-1{,}500{,}000$ & 0.10 & 150,000\\
\midrule
total & & 2,266 & & 3,617,629\\
\bottomrule
\end{tabular}
\caption{The swap book of the lecture (slide~15; the unit of the raw risk is dollars per basis point, presumably). The charge is $|\text{risk}|\times\text{bid--offer}$;
I checked that the column sums to $3{,}617{,}629$.}\label{ch20:tab-risk}
\end{table}

\subsection{Hedging with a few liquid swaps: the general formulation}\label{ch20:sec-hedgeform}
Let $H\in\R^{n\times k}$ have as columns the risk vectors (in the $n$ buckets) of $k$ liquid hedging instruments; if the same instruments calibrate the curve, a hedge instrument
is sensitive only to its own input, and $H$ is made of columns $\bm e_{i(j)}$ of the identity (slide~19). Take $k=5$: the 1, 2, 5, 10 and 30 year swaps. Let $F\in\R^{n\times m}$ collect market
\emph{scenarios} (moves of the inputs) \ts{13}{10}{41:37}. Holding a hedge position that transfers risk $H\bm x$, the residual risk is $\bm r-H\bm x$ and its P\&L in scenario $\bm f$ is $\bm f\T(\bm r-H\bm x)$, so one
chooses
\begin{equation}\label{ch20:eq-hedgeopt}
  \bm x=\arg\min_{\bm x}\ \bigl\|F\T(\bm r-H\bm x)\bigr\|^2 .
\end{equation}
The vector $\bm x$ (the ``translated risk''; the hedge trades are $-\bm x$) is the equivalent risk of the book in the $k$ liquid swaps.

\begin{coursenote}[Sign convention on slides 16 and 19]
Slide~16 writes the objective as $\|F\T(\bm r+H\bm x)\|^2$ and slides 19--23 use $\bm r-H\bm x$ (and $\bm x=R\T\bm r$). I use the latter throughout; the two differ by $\bm x\to-\bm x$.
\end{coursenote}

\subsection{PCA hedging}\label{ch20:sec-pcahedge}
Take for $F$ the first principal components of historical daily changes of the inputs. If $D=USP\T$ is the SVD of the data matrix of changes (rows are days),
$P$ has as columns the principal components, which are eigenvectors of $D\T D$ (\cref{ch:pca}, \cref{thm:svd}) \ts{13}{10}{42:50}; keeping $k$ of them, one for each hedging instrument, is a low-rank
approximation of the market. In the lecture's data the first component is the ``rates do not move now but move in the future'' mode (steep up to 10 years,
flat afterwards, because the Fed is holding), the second is a tilt, the third more complex \ts{13}{10}{43:57}\ts{13}{10}{45:02}.

With $F=P\in\R^{n\times k}$ the objective \eqref{ch20:eq-hedgeopt} can be made \emph{exactly} zero, since $P\T H\in\R^{k\times k}$ is (generically) invertible \ts{13}{10}{45:02}:
\begin{equation}\label{ch20:eq-pcarisk}
  P\T(\bm r-H\bm x)=\bm 0\ \Longrightarrow\ \bm x=(P\T H)^{-1}P\T\bm r=R\T\bm r,\qquad R=P(H\T P)^{-1}\in\R^{n\times k}.
\end{equation}
The matrix $R$ deserves an interpretation (slide~22, \ts{13}{10}{47:12}): it is the unique combination $R=PY$ of the components with $H\T R=I_k$. Its $j$-th column is the
\emph{scenario} (a shift of all 15 inputs, as a combination of principal components) in which the $j$-th hedging swap moves by $1$~bp and the other four do not move at all; and $x_j=\bm r\T R_{\cdot j}$ is the P\&L of the book in that scenario.
Rows of $R$ are the weights with which the risk of each bucket is assigned to the liquid swaps. This is the matrix of slide~20, whose rows have the unit vectors at the hedging tenors (a swap hedges itself)
and, for example, assigns the 4-year risk mainly to the 2-year and 5-year swaps, with weights $0.60$ and $0.70$. The translated risks and their charges are in \cref{ch20:tab-hedge}
(first block): total charge $442{,}036$ against $3{,}617{,}629$, a reduction of a factor $8.2$.

\begin{coursenote}[``Orders of magnitude'']
The speaker says that the charge becomes ``orders of magnitude smaller'', ``something like $400$'' instead of $3.6$ million \ts{13}{10}{47:12}. The slide gives $442{,}036$ for the PCA model
(and $446{,}717$ for the regularised one below), i.e. about $12\%$ of the raw charge: a very large saving, but one order of magnitude, not several.
\end{coursenote}

\paragraph{Weaknesses of PCA hedging} \ts{13}{10}{48:14}\ts{13}{10}{49:19}.
\begin{itemize}
  \item The components are tuned to a historical window, and the coefficients of $R$ are not
      stable, in particular for a short window that reflects the current regime; when they change the hedge must be rebalanced, which costs money.
  \item The instability comes from the noisy components attached to
      small singular values, exactly the mechanism of \cref{ch20:sec-inverse}.
  \item PCA is a least-squares approximation and reacts strongly to outliers: a single day on which rates spiked and reverted has an unwarranted influence.
  \item It overfits history: a model that would have worked for the last three months need not work for the next three.
  \item It makes \emph{no assumption on the shape of the yield curve}: if
      the maturity labels of the swaps were scrambled, the computation would produce the same numbers, unscrambled, although 2 years is close to 1 year and lies between 1 and 5 years. Statistical structure
      without economic structure is the recurrent theme of \cref{ch05:sec-shrink} as well.
\end{itemize}

\subsection{A hedging model regularised by smoothness of the forward curve}\label{ch20:sec-reghedge}
The alternative \ts{13}{10}{50:21}: forget history and impose an \emph{economic assumption}. If nothing is known about what happens in ten years, there is no reason for the forward rate to have a spike at that date, so scenarios should be
smooth in forward-rate space \ts{13}{10}{51:25}. Let $J\in\R^{q\times n}$ be the Jacobian that translates a shift of the $n$ curve inputs into a shift of $q$ forward rates on a fine grid (for the toy above, $q=n$ and $J=M^{-1}$), and $L$ a
smoothness matrix acting on the forward grid, for instance the first-difference matrix, whose rows are $(\dots,1,-1,\dots)$. We look for the smoothest scenarios consistent with the hedging instruments: for each hedge $j$,
\begin{equation}\label{ch20:eq-regR}
  \min_{\bm\rho}\ \|LJ\bm\rho\|^2\quad\text{subject to}\quad H\T\bm\rho=\bm e_j ,
\end{equation}
where $\bm\rho=R_{\cdot j}$ is the input scenario ($1$~bp on hedge $j$, $0$ on the other hedges) and $\|LJ\bm\rho\|^2$ is the roughness of the forward curve it induces. With a soft constraint, penalising
$\|H\T\bm\rho-\bm e_j\|^2+\lambda^2\|LJ\bm\rho\|^2$ and setting the gradient to zero gives, for all columns at once,
\begin{equation}\label{ch20:eq-regRsol}
  R=\bigl(HH\T+\lambda^2\,\mathsf M\bigr)^{-1}H,\qquad \mathsf M=(LJ)\T(LJ),
\end{equation}
with translated risk $\bm x=R\T\bm r$ as before \ts{13}{10}{52:25}.

\begin{proposition}[Soft and hard constraint]\label{ch20:prop-reghedge}
Let $\mathsf M=(LJ)\T LJ$ be invertible. Then $H\T\mathsf M^{-1}H$ is invertible for $H$ of full column rank, the solution of \eqref{ch20:eq-regR} is
$R^\ast=\mathsf M^{-1}H(H\T\mathsf M^{-1}H)^{-1}$, and $(HH\T+\lambda^2\mathsf M)^{-1}H\to R^\ast$ as $\lambda\to0$. If $\mathsf M$ is singular the same limit holds, by the usual penalty-method argument, as long as $\ker\mathsf M\cap\ker H\T=\{0\}$ (as in the toy, where the kernel is the parallel shift).
\end{proposition}
\begin{proof}
Lagrange: $2\mathsf M\bm\rho=H\bm\mu$, so $\bm\rho=\mathsf M^{-1}H\bm\mu/2$ and $H\T\bm\rho=\bm e_j$ fixes $\bm\mu/2=(H\T\mathsf M^{-1}H)^{-1}\bm e_j$. For the limit use the push-through identity
$(\lambda^2\mathsf M+HH\T)^{-1}H=\mathsf M^{-1}H\,(\lambda^2I+H\T\mathsf M^{-1}H)^{-1}$ (multiply both sides by $\lambda^2\mathsf M+HH\T$ and simplify) and let $\lambda\to0$. When $\mathsf M$ is singular the matrix
$HH\T+\lambda^2\mathsf M$ is still invertible under the stated condition, and the minimiser of \eqref{ch20:eq-regR} is unique on the affine set $H\T\bm\rho=\bm e_j$; I do not prove the convergence in that case.
\end{proof}

\begin{coursenote}[Slide 23 leaves out a factor]
Slide~23 shows ``$R\sim(HH\T+\lambda^2(LJ)\T LJ)^{-1}$''. As written the matrix is $n\times n$, not $n\times k$, and does not satisfy $H\T R\approx I$: the correct expression, obtained from the two conditions on the same slide, is \eqref{ch20:eq-regRsol}, with $H$ on the right.
I read ``$\sim$'' as ``is approximately'', the exact constraint being reached for $\lambda\to0$.
\end{coursenote}

With $H\T R\approx I$ up to the size of $\lambda$ (the lecture's matrix has $1.00$ and $0.00$ on the hedging rows, the smoothness being what fills the other rows) the regularised model produces the columns of slide~24 \ts{13}{10}{52:25}. They are
\emph{hat-shaped} (peaked at the hedging tenor and decaying smoothly), whereas the PCA columns follow the historical modes; the results are in the second block of \cref{ch20:tab-hedge}.

\begin{table}[ht]
\centering\footnotesize
\begin{tabular}{lrrrrrr}
\toprule
translated risk on & 1Y & 2Y & 5Y & 10Y & 30Y & total charge\\
\midrule
PCA (slide 20) & $-426{,}757$ & 1,892,770 & $-1{,}538{,}808$ & $-268{,}764$ & 293,261 & 442,036\\
regularised (slide 24) & $-487{,}769$ & 2,082,997 & $-1{,}752{,}958$ & 78,962 & 64,483 & 446,717\\
toy of this chapter & $-331{,}027$ & 1,956,764 & $-1{,}873{,}952$ & $-92{,}944$ & 343,425 & 459,811\\
\bottomrule
\end{tabular}
\caption{Translated risk $\bm x=R\T\bm r$ of the book of \cref{ch20:tab-risk} on the five liquid swaps, and the charge at $0.1$~bp per instrument (the slides print the charge of the PCA
row and I computed the others; the toy, my computation, uses $J=M^{-1}$ of \cref{ch20:par-toy}, first differences for $L$ and $\lambda^2=10^{-4}$). Multiplying the printed two-decimal matrices by the raw risk reproduces the slide values up to the rounding of the entries (differences up to about $25{,}000$ per entry, since raw risks reach $4$ million).}\label{ch20:tab-hedge}
\end{table}

\begin{figure}[ht]
\centering
\begin{tikzpicture}
\begin{axis}[width=12cm,height=5.6cm,xlabel={maturity of the input bucket (years)},ylabel={column of $R$ for the 5Y hedge},xmin=0,xmax=31,ymin=-0.35,ymax=1.1,
  axis lines=left,grid=major,grid style={black!12},legend style={at={(0.5,-0.32)},anchor=north,legend columns=3,draw=none,font=\small},tick label style={font=\small}]
\addplot[thick,thmcol,mark=*,mark size=1.3pt] coordinates {(1,0.00) (2,0.00) (3,0.29) (4,0.70) (5,1.00) (6,0.80) (7,0.54) (8,0.33) (9,0.15) (10,0.00) (12,-0.08) (15,-0.04) (20,-0.08) (25,-0.04) (30,0.00)};
\addlegendentry{PCA (slide 20)}
\addplot[thick,defcol,mark=square*,mark size=1.3pt] coordinates {(1,0.00) (2,0.00) (3,0.44) (4,0.92) (5,1.00) (6,0.81) (7,0.58) (8,0.35) (9,0.15) (10,0.00) (12,-0.17) (15,-0.24) (20,-0.19) (25,-0.09) (30,0.00)};
\addlegendentry{regularised (slide 24)}
\addplot[thick,intcol,dashed,mark=triangle*,mark size=1.5pt] coordinates {(1,-0.00) (2,0.00) (3,0.49) (4,0.86) (5,1.00) (6,0.87) (7,0.63) (8,0.38) (9,0.16) (10,0.00) (12,-0.15) (15,-0.20) (20,-0.16) (25,-0.08) (30,-0.00)};
\addlegendentry{toy (mine)}
\end{axis}
\end{tikzpicture}
\caption{The scenario in which the 5-year hedge moves by $1$~bp and the other four hedges do not: $j=5$Y column of $R$. Both models are equal to $1$ at $5$Y and vanish at the other hedging tenors; the
regularised column is a smooth hat whose negative lobe (rates at 12--25 years move opposite to keep forwards smooth) is more pronounced than PCA's; the toy reproduces the shape. Points are the entries of the
matrices of slides 20 and 24 (two decimals).}\label{ch20:fig-hedge}
\end{figure}

\begin{finance}[Model risk in hedging and relative value]
Two hedges with different assumptions cost the same here (about $445$k) but distribute the risk differently, e.g.\ $10$Y risk of $-268{,}764$ (PCA) versus $+78{,}962$ (regularised). Traders have their own view of the dynamics of the
market, but always \emph{some} model \ts{13}{10}{47:12}. The lecture's argument for the regularised model is that it rests on a \emph{fundamental} assumption (smoothness of forwards) rather than on a window of data
\ts{13}{10}{1:26:24}. The two ideas can also be combined (a PCA regulariser inside the smoothness model).
\end{finance}

\begin{exercise}[(not from the course)]
Let $J=I_4$, $L$ the first-difference matrix ($3\times4$) and $H=(\bm e_1,\bm e_4)$ (hedges in buckets 1 and 4) in \eqref{ch20:eq-regR}. Show that the smoothest scenarios are the linear interpolations $R_{\cdot1}=(1,\tfrac23,\tfrac13,0)\T$ and $R_{\cdot2}=(0,\tfrac13,\tfrac23,1)\T$, compute the translated risk $\bm x=R\T\bm r$ for $\bm r=(0,1,1,0)\T$, and check that $\mathsf M=(LJ)\T LJ$ is singular with kernel spanned by $(1,1,1,1)\T$ (and that \cref{ch20:prop-reghedge} still applies). What would the second-difference matrix give?
\end{exercise}

\begin{exercise}[(not from the course)]
Prove the push-through identity $(\lambda^2\mathsf M+HH\T)^{-1}H=\mathsf M^{-1}H(\lambda^2I+H\T\mathsf M^{-1}H)^{-1}$ used in \cref{ch20:prop-reghedge}, and use it to show that $H\T R_\lambda\to I$ with $R_\lambda=(HH\T+\lambda^2\mathsf M)^{-1}H$ and that $H\T R_\lambda=H\T\mathsf M^{-1}H(\lambda^2I+H\T\mathsf M^{-1}H)^{-1}$ (so the deviation from $I$ is of order $\lambda^2$).
\end{exercise}

% ==================================================================
\section{Linear inverse problems, the SVD and Tikhonov regularisation}\label{ch20:sec-inverse}

The lecture analyses the calibration problem with the tools of linear algebra \ts{13}{10}{1:11:29}. This section develops that analysis in full.

\subsection{The problem and its singular value decomposition}\label{ch20:sec-svdprob}
Let $A\in\R^{M\times N}$ be the linearised pricing map (model parameters $\bm x\in\R^N$ to $M$ market observables), and
\begin{equation}\label{ch20:eq-inv}
  \bm y=A\bm x+\bm\varepsilon
\end{equation}
the market data, with an error $\bm\varepsilon$ (bid--offer, stale quotes, model error) \ts{13}{10}{1:11:29}. Take the thin SVD (\cref{thm:svd}, \cref{eq:svdsum})
$A=\sum_{i=1}^rs_i\bm u_i\bm v_i\T$ with $s_1\ge\dots\ge s_r>0$, $r=\rank A$, orthonormal $\bm u_i\in\R^M$ and $\bm v_i\in\R^N$. In the picture of the slide, $A$ is a rotation, a scaling by $s_i$ and a rotation, and the
inverse is the rotation, the scaling by $1/s_i$, the rotation \ts{13}{10}{1:11:29}.

\begin{proposition}[Minimum-norm least squares]\label{ch20:prop-pinv}
The vectors minimising $\|A\bm x-\bm y\|$ are $\bm x=\sum_{i\le r}\frac{\bm u_i\T\bm y}{s_i}\bm v_i+\bm z$ with $\bm z\in\ker A$; the one of smallest norm is
\begin{equation}\label{ch20:eq-pinv}
  \bm x^\dagger=A^\dagger\bm y=\sum_{i=1}^r\frac{\bm u_i\T\bm y}{s_i}\,\bm v_i=VS^{-1}U\T\bm y .
\end{equation}
It equals $(A\T A)^{-1}A\T\bm y$ if $r=N$ (more data than parameters), and $A\T(AA\T)^{-1}\bm y$, which fits the data exactly, if $r=M<N$.
\end{proposition}
\begin{proof}
Write $\bm x=\sum_j\xi_j\bm v_j+\bm z$ (with $\bm z\perp\bm v_j$). Then $\|A\bm x-\bm y\|^2=\sum_{i\le r}(s_i\xi_i-b_i)^2+\|\bm y\|^2-\sum_{i\le r}b_i^2$, $b_i=\bm u_i\T\bm y$, minimised for $\xi_i=b_i/s_i$; the component $\bm z\in\ker A$ does not enter and the norm is smallest for $\bm z=\bm 0$.
The two closed forms follow by substituting the SVD.
\end{proof}
The slide's formula $\bm x=(A\T A)^{-1}A\T\bm y$ is thus valid only when the columns of $A$ are independent; for the surface calibration, where there are fewer instruments than parameters,
one needs \eqref{ch20:eq-pinv} (the SVD form on the same slide is general).

\subsection{Noise amplification and the condition number}\label{ch20:sec-noise}
Insert \eqref{ch20:eq-inv} into \eqref{ch20:eq-pinv}: $\bm x^\dagger=A^\dagger A\bm x+\sum_i\frac{\bm u_i\T\bm\varepsilon}{s_i}\bm v_i$, the second term being the \emph{propagated error}
\ts{13}{10}{1:12:32}:
\[
  \|\delta\bm x\|^2=\sum_{i=1}^r\frac{(\bm u_i\T\bm\varepsilon)^2}{s_i^2},\qquad
  \E[\|\delta\bm x\|^2]=\sigma^2\sum_{i=1}^r\frac1{s_i^2}\quad\text{if }\Cov\bm\varepsilon=\sigma^2I .
\]
Each mode $i$ of the noise is divided by $s_i$: a small singular value turns an innocent perturbation of the data into a large variation of the parameters. A scale-free measure is the
\emph{condition number} $\kappa=s_1/s_r$ \ts{13}{10}{1:12:32}: for $\bm x$ in the row space of $A$,
\[
  \frac{\|\delta\bm x\|}{\|\bm x\|}\le\kappa\,\frac{\|\bm\varepsilon\|}{\|A\bm x\|},
\]
because $\|\delta\bm x\|\le\|\bm\varepsilon\|/s_r$ and $\|A\bm x\|\le s_1\|\bm x\|$; the bound is attained for $\bm\varepsilon\parallel\bm u_r$, $\bm x\parallel\bm v_1$. A problem is \emph{well posed}
(Hadamard) if a solution exists, is unique and depends continuously on the data; it is \emph{ill-posed} if one of the three fails, and \emph{ill-conditioned} if the discretised problem has a large $\kappa$. The
lecture's calibration problems are underdetermined (uniqueness fails: fewer instruments than parameters, hence $\ker A\ne\{0\}$) \emph{and} unstable (small $s_i$)
\ts{13}{10}{59:53}\ts{13}{10}{1:00:55}.

\begin{intuition}[A physicist's reading: the deconvolution analogue]
The typical forward map averages (the price of a swaption is an integral of the volatility surface over a region), and averaging damps the fine details of the parameters: their contribution to $\bm y$ is $s_i\xi_i$ with $s_i$ small. Inverting means amplifying
back what the data barely contain; when the data have a noise floor $\sigma$, all modes with $s_i|\xi_i|\lesssim\sigma$ are lost, and dividing by their $s_i$ turns them into noise.
This is the same as the inversion of a heat-kernel smoothing, or an image deconvolution, and the standard fix is the same.
\end{intuition}

\paragraph{The discrete Picard picture.} For noiseless data $\bm y=A\bm x$ one has $\bm u_i\T\bm y=s_i\,\bm v_i\T\bm x$: the coefficients $|\bm u_i\T\bm y|$ decay at least as fast as $s_i$ if $\bm x$ is reasonable. With noise, the coefficients level off at the noise floor
$\sigma$ and the ratios $|\bm u_i\T\bm y|/s_i$ that enter \eqref{ch20:eq-pinv} grow. \Cref{ch20:fig-spectral}(b) shows this for the toy calibration of \cref{ch20:sec-lambda}.

\begin{coursenote}[The ``noiseless'' slide]\label{ch20:cn-noiseless}
Slide~46 argues that in the noiseless case ``$A(A\T A)^{-1}A\T\bm y=\bm y$'' so that ``the model appears accurate'', but ``pricing of the actual portfolio with the model $B$ may be unstable''. The point is right; the printed identity is not: slide~46
writes ``$A\T(A\T A)^{-1}A=I$'', which is not even conformable. The correct statement is that $AA^\dagger=\sum_{i\le r}\bm u_i\bm u_i\T$ is the orthogonal projector onto the range of $A$, which is the identity on $\R^M$ exactly when $r=M$, i.e. when
the instruments are independent, for any $N\ge M$. Then \emph{every} data vector, noise included, is repriced exactly: exact repricing certifies nothing, because it is achieved by fitting the noise. What matters is the price $\bm p=B\bm x$ of the
instruments \emph{not} used in the calibration, and $B\bm x^\dagger$ inherits the amplified error $B\,\delta\bm x$. Slide~47 has the same misprint for the regularised repricing matrix, see \eqref{ch20:eq-influence}.
\end{coursenote}

\begin{exercise}[(not from the course)]
Consider $A=\begin{psmallmatrix}1&1\\1&1.01\end{psmallmatrix}$.
\begin{enumerate}[(a)]
  \item Solve $A\bm x=\bm y$ for $\bm y=(2,2.01)\T$ and for $\bm y'=(2,2.02)\T$; by what fraction do data and solution change?
  \item Compute the singular values and $\kappa(A)$ (about $4\times10^2$) and
      check that the change of (a) is compatible with the bound of \cref{ch20:sec-noise}.
  \item Compute the Tikhonov solutions \eqref{ch20:eq-tikh} for $\lambda=0.1$ for both $\bm y$ and $\bm y'$ and compare.
\end{enumerate}
\end{exercise}

\subsection{Truncated SVD and Tikhonov regularisation}\label{ch20:sec-tikh}
Both methods multiply each term of \eqref{ch20:eq-pinv} by a \emph{filter factor} $w_i\in[0,1]$:
\begin{equation}\label{ch20:eq-filter}
  \bm x_w=\sum_{i=1}^rw_i\,\frac{\bm u_i\T\bm y}{s_i}\,\bm v_i .
\end{equation}
\begin{itemize}
  \item \emph{Truncated SVD} (TSVD): $w_i=1$ for $i\le k$ and $0$ for $i>k$ \ts{13}{10}{1:16:49}. One keeps the well-determined modes and discards the rest, which are the part of
        the parameter space that the calibration cannot see (the ``null space'' of the model, in the speaker's words). It is the same idea as the PCA risk model of \cref{ch20:sec-pcahedge} and is principal-components regression in the language of \cref{ch05:sec-shrink}.
        The danger: if the discarded modes do affect the price of the portfolio, the truncation is a bias one pays for.
  \item \emph{Tikhonov regularisation} \ts{13}{10}{1:14:44}:
        \begin{equation}\label{ch20:eq-tikh}
          \bm x_\lambda=\arg\min_{\bm x}\Bigl\{\|A\bm x-\bm y\|^2+\lambda^2\|\bm x\|^2\Bigr\}=(A\T A+\lambda^2I)^{-1}A\T\bm y=\sum_{i=1}^r\underbrace{\frac{s_i^2}{s_i^2+\lambda^2}}_{w_i}\frac{\bm u_i\T\bm y}{s_i}\bm v_i .
        \end{equation}
\end{itemize}
The gradient of the objective is $2A\T(A\bm x-\bm y)+2\lambda^2\bm x$; setting it to zero gives the normal equations $(A\T A+\lambda^2I)\bm x=A\T\bm y$, uniquely solvable for every $\lambda>0$ even if $A$ is rank-deficient (the matrix is
positive definite), and $A\T=V S U\T$ gives the filter form, since $(A\T A+\lambda^2I)^{-1}A\T=V\,\diag\!\bigl(\tfrac{s_i}{s_i^2+\lambda^2}\bigr)U\T$. Hence Tikhonov is a \emph{smooth} filter: $w_i\approx1$ for $s_i\gg\lambda$ and $w_i\approx s_i^2/\lambda^2$ for $s_i\ll\lambda$
(\cref{ch20:fig-spectral}(a)). The essential point is stated by the speaker: one no longer divides by a small number \ts{13}{10}{1:14:44}. Quantitatively, the gain applied to a mode of noise is
$\frac{w_i}{s_i}=\frac{s_i}{s_i^2+\lambda^2}\le\frac1{2\lambda}$ (maximum at $s_i=\lambda$), instead of $1/s_r$ for the exact inverse, and $1/s_k$ for TSVD.

\begin{intuition}[Ridge regression in another language]
This is the ridge regression of \cref{ch05:sec-shrink} with $d_j\to s_i$ and ridge parameter $\lambda_{\rm ridge}=\lambda^2$ (the lecture writes the penalty as $\|\lambda\bm x\|^2$). The smooth filter $s^2/(s^2+\lambda^2)$ is the same shrinkage factor,
and \cref{ch05:prop-ridgebayes} gives its Bayesian reading: a Gaussian prior $\bm x\sim N(\bm 0,\tau^2I)$ with $\lambda^2=\sigma^2/\tau^2$. Early stopping of gradient descent acts as another filter of the same family (\cref{ch16:sec-gd}).
\end{intuition}

The price of stability is that the calibration is no longer exact. The fit and the ``re-pricing'' operator are, by \eqref{ch20:eq-tikh},
\begin{equation}\label{ch20:eq-influence}
  A\bm x_\lambda=H_\lambda\bm y,\qquad H_\lambda=A(A\T A+\lambda^2I)^{-1}A\T=\sum_iw_i\,\bm u_i\bm u_i\T\ne I,
\end{equation}
the corrected form of the ``$A\T(A\T A+\lambda^2I)^{-1}A\neq I$'' of slide~47 \ts{13}{10}{1:15:49}; the residual is $(I-H_\lambda)\bm y$. Its trace $\operatorname{tr}H_\lambda=\sum_iw_i$ is the \emph{effective number of parameters} used by the fit.

\begin{figure}[ht]
\centering
\begin{tikzpicture}
\begin{axis}[name=a,width=6.2cm,height=5cm,xlabel={$\log_{10}(s/\lambda)$},ylabel={filter $w$},xmin=-2,xmax=2,ymin=-0.05,ymax=1.1,
  axis lines=left,grid=major,grid style={black!12},tick label style={font=\small},legend style={at={(0.5,-0.32)},anchor=north,draw=none,legend columns=2,font=\small}]
\addplot[thick,thmcol,domain=-2:2,samples=80] {10^(2*x)/(10^(2*x)+1)};
\addlegendentry{Tikhonov}
\addplot[thick,defcol,dashed] coordinates {(-2,0)(0,0)(0,1)(2,1)};
\addlegendentry{TSVD}
\node[anchor=north west] at (axis cs:-1.95,1.08) {\small(a)};
\end{axis}
\begin{axis}[at={(a.east)},anchor=west,xshift=1.7cm,width=6.2cm,height=5cm,xlabel={mode $i$},ymode=log,ymin=5e-4,ymax=1,xmin=0.5,xmax=10.5,
  axis lines=left,grid=major,grid style={black!12},tick label style={font=\small},legend style={at={(0.5,-0.32)},anchor=north,draw=none,legend columns=3,font=\small}]
\addplot[thick,thmcol,mark=*,mark size=1.3pt] coordinates {(1,0.35564) (2,0.30201) (3,0.27821) (4,0.21286) (5,0.15147) (6,0.10928) (7,0.07985) (8,0.04535) (9,0.02417) (10,0.00281)};
\addlegendentry{$s_i$}
\addplot[thick,defcol,mark=square*,mark size=1.3pt] coordinates {(1,0.45418) (2,0.49752) (3,0.12507) (4,0.00788) (5,0.00641) (6,0.00865) (7,0.00814) (8,0.00214) (9,0.00357) (10,0.00174)};
\addlegendentry{$|\bm u_i\T\bm y|$}
\addplot[thick,intcol,dashed] coordinates {(0.5,0.00189)(10.5,0.00189)};
\addlegendentry{noise}
\node[anchor=south west] at (axis cs:0.8,6e-4) {\small(b)};
\end{axis}
\end{tikzpicture}
\caption{(a) Filter factors of Tikhonov ($w=s^2/(s^2+\lambda^2)$) and of TSVD (cut at $s=\lambda$). (b) Singular values $s_i$ of the toy calibration matrix of \cref{ch20:sec-lambda} (ten instruments, 100 parameters, $\kappa=126$) and the coefficients $|\bm u_i\T\bm y|$ of noisy data
(1\% price noise; my computation). From mode 4 on the coefficients stop tracking $s_i$ and stay within one to four times the noise floor $\|\bm\varepsilon\|/\sqrt M$ (dashed), whereas $s_i$ keeps falling: the terms $\bm u_i\T\bm y/s_i$ of \eqref{ch20:eq-pinv} are then dominated by noise.}\label{ch20:fig-spectral}
\end{figure}

\begin{exercise}[(not from the course)]
Prove the following for Tikhonov regularisation with $L=I$:
\begin{enumerate}[(a)]
  \item $\rho(\lambda)=\|A\bm x_\lambda-\bm y\|$ is non-decreasing and $\|\bm x_\lambda\|$ non-increasing in $\lambda$;
  \item as $\lambda\to0$, $\bm x_\lambda\to A^\dagger\bm y$ (\cref{ch20:prop-pinv}), also when $A$ is rank-deficient;
  \item as $\lambda\to\infty$, $\bm x_\lambda\approx A\T\bm y/\lambda^2$;
  \item $0\le\operatorname{tr}H_\lambda\le r$ and it decreases from $r$ to $0$.
\end{enumerate}
\end{exercise}

\subsection{General form: weights, smoothness matrices and several penalties}\label{ch20:sec-general}
Penalising the amplitude $\|\bm x\|$ is rarely what one wants. The general problem, with a diagonal weight matrix $W$ and a penalty matrix $L$, is
\begin{equation}\label{ch20:eq-gen}
  \bm x_\lambda=\arg\min_{\bm x}\Bigl\{\|W(A\bm x-\bm y)\|^2+\lambda^2\|L\bm x\|^2\Bigr\}=(A\T W^2A+\lambda^2L\T L)^{-1}A\T W^2\bm y .
\end{equation}
The solution is unique iff $\ker(WA)\cap\ker L=\{0\}$. It is an ordinary least-squares problem for the \emph{augmented system}
$\begin{psmallmatrix}WA\\ \lambda L\end{psmallmatrix}\bm x\approx\begin{psmallmatrix}W\bm y\\ \bm 0\end{psmallmatrix}$, which is how it is solved stably in practice (QR, \cref{ch05:sec-qr}). Three remarks.
\begin{itemize}
  \item \emph{Weights.} $W=\diag(1/\sigma_i)$, with $\sigma_i$ the tolerance of instrument $i$ (its bid--offer): the data term is then a $\chi^2$ and a residual of one tolerance costs one unit. Instruments that are
        important or liquid get a small tolerance and are fitted more tightly \ts{13}{10}{1:07:09}.
  \item \emph{Smoothness instead of amplitude.} With $L=I$ the penalty pulls $\bm x$ towards $\bm 0$, which for a volatility surface means a volatility of zero. The speaker's point \ts{13}{10}{1:19:00}: ask first what one does \emph{not} like. A perfectly flat
        surface is acceptable, so the penalty should vanish on it, i.e. $L$ should have the constants in its kernel: the first-difference matrix with rows $(\dots,1,-1,\dots)$ of slide~41 penalises $\sum_i(x_i-x_{i+1})^2$, a discrete $\int|\nabla x|^2$ (in two dimensions the differences of neighbouring nodes along both axes).
  \item \emph{No bias on the kernel of $L$.} If $A\bm x=\bm y$ exactly and $L\bm x=\bm 0$, then $\bm x_\lambda=\bm x$: indeed $\bm x_\lambda=(A\T W^2A+\lambda^2L\T L)^{-1}A\T W^2A\bm x=\bm x-\lambda^2(A\T W^2A+\lambda^2L\T L)^{-1}L\T L\bm x=\bm x$.
        Regularisation biases the solution only in the directions it penalises, so the bias is small when the penalty encodes true prior knowledge \ts{13}{10}{1:17:52}\ts{13}{10}{1:19:00}.
\end{itemize}
Slide~42 combines two penalties, one on the \emph{change} of the surface with respect to a prior $\bm v_0$ (matrix $L_1$ acting on $\bm v-\bm v_0=B\bm x$) and one of Tikhonov type on the coefficients (matrix $L_2$), with parameters $\lambda_1,\lambda_2$ (\cref{ch20:sec-volsolve}); this is \eqref{ch20:eq-gen} with
$\lambda^2L\T L\to\lambda_1^2(L_1B)\T L_1B+\lambda_2^2L_2\T L_2$. The Bayesian reading is the same as for ridge: a Gaussian prior with precision $(\lambda^2/\sigma^2)L\T L$, a random-walk prior for a difference penalty.

\begin{coursenote}[Notation of the slides]
On slide~42 the inverse matrix is called $A$, which is also the model matrix of slides 44--48 (here $K=(J\T W^2J+\dots)^{-1}$ in \cref{ch20:sec-volsolve}); and the penalty matrix $L_2$ is applied to the volatility vector $\bm v$ on slide~40 but to the coefficients $\bm x$ on slide~42.
\end{coursenote}

\subsection{Bias, variance and the meaning of $\lambda$}\label{ch20:sec-biasvar}
For Tikhonov with $L=I$ the error can be computed exactly. Let $\bm y=A\bm x+\bm\varepsilon$, $\E[\bm\varepsilon]=\bm 0$, $\Cov\bm\varepsilon=\sigma^2I$, and $\xi_i=\bm v_i\T\bm x$. From \eqref{ch20:eq-filter}, $\bm x_\lambda=\sum_iw_i(\xi_i+\bm u_i\T\bm\varepsilon/s_i)\bm v_i$, hence
\begin{equation}\label{ch20:eq-mse}
  \E[\|\bm x_\lambda-\bm x\|^2]=\sum_{i=1}^r\Bigl[\underbrace{(1-w_i)^2\xi_i^2}_{\text{bias}^2}+\underbrace{\sigma^2\frac{w_i^2}{s_i^2}}_{\text{variance}}\Bigr]+\sum_{i>r}\xi_i^2 ,\qquad 1-w_i=\frac{\lambda^2}{s_i^2+\lambda^2},\quad\frac{w_i^2}{s_i^2}=\frac{s_i^2}{(s_i^2+\lambda^2)^2}.
\end{equation}
(I checked \eqref{ch20:eq-mse} by Monte Carlo on a random $8\times6$ matrix: $0.7747$ against $0.7750$ from $4\times10^4$ samples.) The last sum is the part of $\bm x$ in the kernel of $A$, which no method can recover.
Bias grows and variance falls with $\lambda$; at $\lambda=0$ the variance $\sigma^2\sum s_i^{-2}$ is the amplification of \cref{ch20:sec-noise}. Minimising each term over $w_i$: $-2(1-w)\xi^2+2w\sigma^2/s^2=0$ gives
\begin{equation}\label{ch20:eq-wiener}
  w_i^\ast=\frac{s_i^2\xi_i^2}{s_i^2\xi_i^2+\sigma^2}=\frac{s_i^2}{s_i^2+\sigma^2/\xi_i^2},
\end{equation}
a Wiener filter, i.e.\ Tikhonov with a mode-dependent $\lambda_i^2=\sigma^2/\xi_i^2$: a mode is kept if the signal it contributes to the data, $s_i|\xi_i|$, exceeds the noise $\sigma$, and suppressed otherwise. TSVD is the $0/1$ version of this filter and a single $\lambda^2=\sigma^2/\tau^2$
(with $\tau^2$ a typical size of $\xi_i^2$) is the ridge--Bayes value of \cref{ch05:prop-ridgebayes}. The optimum cannot be used as it stands, since $\xi_i$ is unknown, but it says what $\lambda$ \emph{is}: a noise-to-signal ratio for the parameters.

\begin{exercise}[(not from the course)]
Take $N=M=1$: $y=ax+\varepsilon$, $\varepsilon\sim N(0,\sigma^2)$. Find the Tikhonov estimator, its bias and variance, and the $\lambda$ that minimises the mean-square error as a function of $x$; compare with \eqref{ch20:eq-wiener}. What happens if one substitutes the ridge--Bayes value $\lambda^2=\sigma^2/\tau^2$ and the true $|x|$ is much larger or much smaller than $\tau$?
\end{exercise}

% ==================================================================
\section{Choosing the regularisation parameter (not covered in class)}\label{ch20:sec-lambda}
The lecture calls $\lambda$ a ``small regularisation parameter'' and says nothing on how to fix it \ts{13}{10}{52:25}. Three standard criteria, then a numerical test. Let $\rho(\lambda)=\|W(A\bm x_\lambda-\bm y)\|$ be the residual.
For $L=I$ and $b_i=\bm u_i\T\bm y$, $\rho^2=\sum_{i\le r}(1-w_i)^2b_i^2+\text{const}$ increases with $\lambda$ and $\|\bm x_\lambda\|^2=\sum_iw_i^2b_i^2/s_i^2$ decreases (each $w_i$ decreases with $\lambda$), so the trade-off is monotone.

\begin{itemize}
  \item \emph{Discrepancy principle} (Morozov): choose $\lambda$ with $\rho(\lambda)=\delta$, the size of the data error, $\delta\approx\sigma\sqrt M$ (with $W=\diag(1/\sigma_i)$: $\rho^2\approx M$). Rationale: fitting closer than the noise fits the noise. This is the lecture's remark that every input has a
        tolerance and calibrating exactly is pointless \ts{13}{10}{1:06:08}; it needs the noise level, here the bid--offer.
  \item \emph{Generalised cross-validation.} Leave-one-out cross-validation predicts each datum from the others. For a linear smoother $\hat{\bm y}=H_\lambda\bm y$ one has $y_i-\hat y_i^{(-i)}=(y_i-\hat y_i)/(1-(H_\lambda)_{ii})$ (\emph{proof}: the fit computed without datum $i$ is also the minimiser when
        $y_i$ is replaced by its own prediction $\hat y_i^{(-i)}$, since that term is then zero; so $\hat y_i^{(-i)}=\sum_{j\ne i}H_{ij}y_j+H_{ii}\hat y_i^{(-i)}$). Replacing $(H_\lambda)_{ii}$ by the average $\operatorname{tr}H_\lambda/M$ gives
        \[
          \mathrm{GCV}(\lambda)=\frac{M\,\rho(\lambda)^2}{\bigl(M-\operatorname{tr}H_\lambda\bigr)^2},
        \]
        which needs no noise level; it is the counterpart of the cross-validation used for ridge in \cref{ch05:sec-shrink}.
  \item \emph{L-curve.} Plot $\log\rho(\lambda)$ against $\log\|L\bm x_\lambda\|$; the corner of the ``L'' separates the regime dominated by noise (small $\lambda$, large seminorm) from the regime dominated by bias.
\end{itemize}

\begin{casestudy}[A toy calibration with regularisation (my computation; not in class)]\label{ch20:cs-toy}
\textbf{Goal.} Mimic the calibration of a volatility slice. \textbf{Set-up.} A smooth ``true'' volatility $v(s)$ on a grid of $N=100$ points, $s\in[0,10]$ (level $0.19$--$0.25$); $M=10$ instruments, each an average of $v$ with a Gaussian window (centres $0.5$--$9.7$, widths growing from $0.46$ to $1.56$, overlapping); prices $y_i=(A\bm v)_i(1+0.01z_i)$ with $z_i$ standard normal, so $\delta=\|\bm\varepsilon\|=0.0060$;
$\kappa(A)=126$ on the ten modes.
\textbf{Methods.}
\begin{enumerate}[(a)]
  \item The lecture's ``formal'' solution: ten piecewise-constant basis functions, square Jacobian, exact inversion.
  \item Tikhonov \eqref{ch20:eq-gen} with the first-difference matrix $L$ of slide~41, $W=I$, on the fine grid.
\end{enumerate}
\textbf{Results} (\cref{ch20:fig-toy}).
\begin{enumerate}[(a)]
  \item Reprices the ten inputs to $5\times10^{-17}$ but the surface swings between $0.002$ and $0.40$ (root-mean-square error $0.106$, roughness $\|L\bm v\|=0.64$), and a $1\%$ change in one input moves it by $16\%$ of its level in the worst place: the skyline of the
      lecture \ts{13}{10}{1:04:01}.
  \item With $\lambda$ from GCV ($0.44$, effective number of parameters $\operatorname{tr}H_\lambda=6.2$): error $0.0031$, roughness $0.0144$ (truth $0.0147$), the same $1\%$ perturbation moves it by $1.3\%$ of its level. The discrepancy principle
      gives $\lambda=0.96$ and error $0.0032$; the best possible value (knowing the truth) is $0.0031$ at $\lambda=0.65$. GCV picks a residual of $0.46\,\delta$, a mild over-fit, typical of GCV.
\end{enumerate}
\textbf{Which penalty?} With the same tuning, the best error is $0.0167$ for $L=I$ (amplitude), $0.0039$ for the second difference and $0.0031$ for the first difference. TSVD of $A$ (minimum-norm solution) reaches at best $0.018$ (with $k=8$ modes) and
$0.063$ with all ten: an unsuitable penalty (here: the minimum norm, which pulls the surface towards zero) gives errors five to six times larger than the suitable one, whereas $\lambda$ can be changed by a factor two with hardly any effect on the error.
\textbf{Lesson.} Regularisation is a matter of choosing \emph{what} to penalise (prior knowledge of the solution) first, and how much second.
\end{casestudy}

\begin{figure}[ht]
\centering
\begin{tikzpicture}
\begin{axis}[name=a,width=6.4cm,height=5.4cm,xlabel={forward time $s$},ylabel={volatility},xmin=0,xmax=10,ymin=-0.02,ymax=0.42,
  axis lines=left,grid=major,grid style={black!12},tick label style={font=\small},legend style={at={(0.5,-0.3)},anchor=north,draw=none,legend columns=1,font=\small}]
\addplot[thick,black,dashed] coordinates {(0.000,0.2309) (0.303,0.2305) (0.606,0.2305) (0.909,0.2312) (1.212,0.2333) (1.515,0.2368) (1.818,0.2416) (2.121,0.2467) (2.424,0.2509) (2.727,0.2528) (3.030,0.2511) (3.333,0.2457) (3.636,0.2371) (3.939,0.2267) (4.242,0.2159) (4.545,0.2063) (4.848,0.1986) (5.152,0.1931) (5.455,0.1895) (5.758,0.1876) (6.061,0.1868) (6.364,0.1868) (6.667,0.1873) (6.970,0.1882) (7.273,0.1894) (7.576,0.1907) (7.879,0.1922) (8.182,0.1937) (8.485,0.1953) (8.788,0.1969) (9.091,0.1984) (9.394,0.1999) (9.697,0.2012) (10.000,0.2024)};
\addlegendentry{truth}
\addplot[thick,thmcol,const plot] coordinates {(0.000,0.2275) (0.101,0.2275) (0.202,0.2275) (0.303,0.2275) (0.404,0.2275) (0.505,0.2275) (0.606,0.2275) (0.707,0.2275) (0.808,0.2275) (0.909,0.2275) (1.010,0.2431) (1.111,0.2431) (1.212,0.2431) (1.313,0.2431) (1.414,0.2431) (1.515,0.2431) (1.616,0.2431) (1.717,0.2431) (1.818,0.2431) (1.919,0.2431) (2.020,0.2528) (2.121,0.2528) (2.222,0.2528) (2.323,0.2528) (2.424,0.2528) (2.525,0.2528) (2.626,0.2528) (2.727,0.2528) (2.828,0.2528) (2.929,0.2528) (3.030,0.2457) (3.131,0.2457) (3.232,0.2457) (3.333,0.2457) (3.434,0.2457) (3.535,0.2457) (3.636,0.2457) (3.737,0.2457) (3.838,0.2457) (3.939,0.2457) (4.040,0.1940) (4.141,0.1940) (4.242,0.1940) (4.343,0.1940) (4.444,0.1940) (4.545,0.1940) (4.646,0.1940) (4.747,0.1940) (4.848,0.1940) (4.949,0.1940) (5.051,0.2413) (5.152,0.2413) (5.253,0.2413) (5.354,0.2413) (5.455,0.2413) (5.556,0.2413) (5.657,0.2413) (5.758,0.2413) (5.859,0.2413) (5.960,0.2413) (6.061,0.0444) (6.162,0.0444) (6.263,0.0444) (6.364,0.0444) (6.465,0.0444) (6.566,0.0444) (6.667,0.0444) (6.768,0.0444) (6.869,0.0444) (6.970,0.0444) (7.071,0.3991) (7.172,0.3991) (7.273,0.3991) (7.374,0.3991) (7.475,0.3991) (7.576,0.3991) (7.677,0.3991) (7.778,0.3991) (7.879,0.3991) (7.980,0.3991) (8.081,0.0020) (8.182,0.0020) (8.283,0.0020) (8.384,0.0020) (8.485,0.0020) (8.586,0.0020) (8.687,0.0020) (8.788,0.0020) (8.889,0.0020) (8.990,0.0020) (9.091,0.2893) (9.192,0.2893) (9.293,0.2893) (9.394,0.2893) (9.495,0.2893) (9.596,0.2893) (9.697,0.2893) (9.798,0.2893) (9.899,0.2893) (10.000,0.2893)};
\addlegendentry{formal (exact repricing)}
\addplot[thick,defcol] coordinates {(0.000,0.2265) (0.303,0.2273) (0.606,0.2297) (0.909,0.2338) (1.212,0.2387) (1.515,0.2436) (1.818,0.2478) (2.121,0.2507) (2.424,0.2520) (2.727,0.2512) (3.030,0.2483) (3.333,0.2431) (3.636,0.2359) (3.939,0.2274) (4.242,0.2183) (4.545,0.2093) (4.848,0.2013) (5.152,0.1946) (5.455,0.1896) (5.758,0.1862) (6.061,0.1842) (6.364,0.1834) (6.667,0.1835) (6.970,0.1842) (7.273,0.1854) (7.576,0.1870) (7.879,0.1891) (8.182,0.1916) (8.485,0.1944) (8.788,0.1973) (9.091,0.2001) (9.394,0.2025) (9.697,0.2043) (10.000,0.2050)};
\addlegendentry{Tikhonov, $\lambda=0.44$}
\node[anchor=north] at (axis cs:8,0.4) {\small(a)};
\end{axis}
\begin{axis}[at={(a.east)},anchor=west,xshift=1.7cm,width=6.4cm,height=5.4cm,xlabel={$\log_{10}\lambda$},xmin=-2,xmax=1,ymode=log,ymin=0.1,ymax=100,
  axis lines=left,grid=major,grid style={black!12},tick label style={font=\small},legend style={at={(0.5,-0.3)},anchor=north,draw=none,legend columns=1,font=\small}]
\addplot[thick,thmcol] coordinates {(-4.000,20.0356) (-3.831,20.0356) (-3.661,20.0324) (-3.492,20.0227) (-3.322,20.0000) (-3.153,19.9547) (-2.983,19.8576) (-2.814,19.6505) (-2.644,19.2071) (-2.475,18.3139) (-2.305,16.6343) (-2.136,13.8900) (-1.966,10.3042) (-1.797,6.8123) (-1.627,4.3398) (-1.458,3.0388) (-1.288,2.4207) (-1.119,2.0097) (-0.949,1.6214) (-0.780,1.2945) (-0.610,1.1003) (-0.441,1.0259) (-0.271,1.0000) (-0.102,1.0032) (0.068,1.1586) (0.237,1.6570) (0.407,2.4207) (0.576,3.1812) (0.746,3.8576) (0.915,4.5955)};
\addlegendentry{error / best error}
\addplot[thick,defcol,dashed] coordinates {(-4.000,5.1023) (-3.831,5.1023) (-3.661,5.1023) (-3.492,5.1023) (-3.322,5.1023) (-3.153,5.1023) (-2.983,5.1023) (-2.814,5.1004) (-2.644,5.0967) (-2.475,5.0911) (-2.305,5.0763) (-2.136,5.0428) (-1.966,4.9740) (-1.797,4.8326) (-1.627,4.5536) (-1.458,4.0644) (-1.288,3.3538) (-1.119,2.5800) (-0.949,1.9438) (-0.780,1.4953) (-0.610,1.1830) (-0.441,1.0000) (-0.271,1.0327) (-0.102,1.5727) (0.068,3.3036) (0.237,7.1168) (0.407,12.9092) (0.576,19.3080) (0.746,26.6555) (0.915,38.3557)};
\addlegendentry{GCV / minimum}
\addplot[thick,intcol] coordinates {(-4.000,0.0000) (-3.831,0.0167) (-3.661,0.0167) (-3.492,0.0500) (-3.322,0.0833) (-3.153,0.2000) (-2.983,0.4333) (-2.814,0.9333) (-2.644,1.9833) (-2.475,4.1333) (-2.305,8.1500) (-2.136,14.7500) (-1.966,23.4500) (-1.797,32.1333) (-1.627,38.7167) (-1.458,42.8667) (-1.288,45.5333) (-1.119,48.2167) (-0.949,52.2000) (-0.780,57.6333) (-0.610,63.7833) (-0.441,71.3500) (-0.271,85.9667) (-0.102,122.8000) (0.068,201.8500) (0.237,330.2000) (0.407,487.9833) (0.576,645.5000) (0.746,809.7833) (0.915,1024.7500)};
\addlegendentry{residual $\rho/\delta$}
\addplot[gray,dotted] coordinates {(-2,1)(1,1)};
\node[anchor=north west] at (axis cs:-1.95,90) {\small(b)};
\end{axis}
\end{tikzpicture}
\caption{The toy calibration of \cref{ch20:cs-toy} (my computation). (a) The formal exact solution is the ``skyline''; the smoothness-regularised one follows the truth to $0.003$. (b) Error, GCV and residual against $\lambda$ (each normalised as indicated; dotted line: $\rho=\delta$, the discrepancy principle). The minima of the error and of GCV are broad and lie within a factor two of each other.}\label{ch20:fig-toy}
\end{figure}

\begin{exercise}[(not from the course)]
Reproduce the toy of \cref{ch20:cs-toy} (any language): choose your own smooth $v$ on $[0,10]$, ten Gaussian-window instruments, $1\%$ noise; compute the formal solution, the Tikhonov solution for a grid of $\lambda$, the discrepancy and GCV parameters. Vary the penalty ($L=I$, first difference, second difference) and the noise level, and report which choice is most sensitive to $\lambda$.
\end{exercise}

% ==================================================================
\section{Calibrating a volatility surface}\label{ch20:sec-vol}

\subsection{The pricing engine and the HJM model}\label{ch20:sec-hjmvol}
The lecture now steps back \ts{13}{10}{53:26}. A \emph{pricing engine} takes the trader's portfolio and the model parameters (curves, volatility surface, \dots) and returns values and risks. To make
the parameters consistent with the market, a \emph{calibration engine} compares the model prices of a set of benchmark instruments with their market prices and adjusts the parameters until they agree
(slide~25); only then is the portfolio priced. The example is the Heath--Jarrow--Morton model (\cref{ch:hjm}), used to price volatility products by Monte Carlo \ts{13}{10}{54:35}:
\begin{equation}\label{ch20:eq-hjm}
  \dd f_{t,s}=\mu_{t,s}\dd t+f_{t,s}^{\beta}\,V(t,s)\,\rho(t,s)\cdot\dd B^{\Q}_t ,
\end{equation}
for the forward rate $f_{t,s}$ observed at calendar time $t$ for forward time $s$: $\mu$ is the drift (fixed by absence of arbitrage, \cref{ch:hjm}), $\beta$ the ``skew factor'' ($\beta=1$ lognormal, $\beta=0$ normal model), $\rho$ the correlation and factor structure, $B^\Q$ a
Brownian motion, and $V(t,s)$ the \emph{volatility surface} that is the subject of the lecture \ts{13}{10}{55:41}. The Monte Carlo advances the whole forward curve from one quarter to the next, and every step needs the value of $V$ on
the corresponding cell of the triangular grid of (calendar time, forward time) \ts{13}{10}{56:47}; the instruments that pin it down, caps and swaptions, are sensitive to \emph{regions} of this triangle (\ts{13}{10}{57:48}; a swaption
depends on the forward rates that will be observed at the option expiry and on their volatility until then), and the regions of different instruments overlap heavily \ts{13}{10}{58:50}.

\subsection{Why the problem is hard}\label{ch20:sec-volhard}
The surface has a value for each of the cells of a triangle with 240 quarters (60 years) on a side, that is $240\cdot241/2=28{,}920$ unknowns \ts{13}{10}{59:53} (the slide says ``$\sim$28k''), whereas only $20$--$50$ swaptions and caps are quoted. The problem is underdetermined
(the kernel of the forward map has dimension at least $28{,}870$), ill-posed because of the overlapping sensitivities, and its unregularised solution is unstable and noisy \ts{13}{10}{1:00:55}. A full $28{,}920\times28{,}920$ matrix of doubles takes $6.7$~GB (my count),
which the speaker rules out as ``no memory''. And the result must be \emph{smooth}: two similar swaptions in the portfolio should get similar prices, and a spike of volatility in some region of the future is either explained by an economic event or is an artefact
\ts{13}{10}{1:00:55}.

\begin{coursenote}[Counting the elements of the triangle]
In the transcript \ts{13}{10}{59:53} the number of triangle elements is first said to be ``200K'' and then a matrix of $28$K$\times28$K is discussed; slide~29 gives ``$\sim$28k'', which is the right count: half of $240^2=57{,}600$ plus the diagonal is $28{,}920$.
\end{coursenote}

\subsection{The ``formal'' solution and why it fails}\label{ch20:sec-volformal}
Represent the surface, flattened to a vector $\bm v$ with one entry per cell, as a prior value plus a combination of basis functions \ts{13}{10}{1:01:56}:
\begin{equation}\label{ch20:eq-vbasis}
  \bm v=\bm v_0+B\bm x,
\end{equation}
where the columns of $B$ are the basis functions (values on the grid) and $\bm x$ the weights. The pricing model gives the sensitivities of the prices of the calibration instruments to a perturbation of the surface along each basis function \ts{13}{10}{1:02:59}; working with
log-prices, $q_i=\ln(q_i^{\rm mdl}/q_i^{0})$, $q_i^{\rm in}=\ln(q_i^{\rm mkt}/q_i^{0})$,
\begin{equation}\label{ch20:eq-jac}
  J_{ij}=\frac{\partial q_i}{\partial x_j},\qquad \bm q\approx J\bm x .
\end{equation}
The \emph{formal} approach takes as many basis functions as instruments ($N=M$), piecewise constant on the region of sensitivity of each instrument, so that $J$ is square and (``if the basis functions are reasonable'') invertible, and solves $\bm x=J^{-1}\bm q^{\rm in}$ iteratively \ts{13}{10}{1:02:59}.
The iteration converges quickly because for at-the-money instruments the price is nearly proportional to the volatility: for the Black--Scholes forward price (\cref{ch:bs}), $C=F\,[2N(\sigma\sqrt T/2)-1]=F\,\sigma\sqrt{T}/\sqrt{2\pi}\,(1-\sigma^2T/24+\dots)$, so $\ln C\approx\ln\sigma+\text{const}$ and $q$ is nearly linear in the volatility
increments (my derivation of the speaker's remark; the calibrated instruments are swaptions in an HJM Monte Carlo, for which the proportionality is heuristic).

The outcome is an exact fit and a meaningless surface: ``Manhattan skyline, with the Hudson river'' \ts{13}{10}{1:04:01} (slide~33, ``Exact, but \dots\ meaningless''). Smoothing the piecewise-constant basis functions gives a better-looking surface but not a good one, and the instability
remains: a $1\%$ change in the price of the $5\text{y}\times10\text{y}$ swaption changes the adjustment of the surface by $4\%$ (slide~35: one tall isolated ``building with an antenna'') \ts{13}{10}{1:05:04}\ts{13}{10}{1:06:08}. This is what \cref{ch20:sec-noise} predicts: $J$ is square, so
$J^\dagger=J^{-1}$ and errors are amplified by up to $\kappa(J)$; the toy of \cref{ch20:cs-toy} shows the same phenomenon (1\% of noise, a factor $16$ on the surface).

\subsection{The regularised design}\label{ch20:sec-volsolve}
The lecture turns ill-posedness to advantage with four changes \ts{13}{10}{1:06:08}\ts{13}{10}{1:07:09} (slide~36).
\begin{itemize}
  \item \emph{Tolerance.} Each input has a bid--offer, so there is no point in calibrating exactly. The fit becomes a weighted least-squares term $\|W(\bm q-\bm q^{\rm in})\|^2$, $W$ diagonal with small weights on instruments of large bid--offer and large ones on important instruments. A
        small loss of accuracy buys a large gain of smoothness.
  \item \emph{Smooth basis functions.} Cubic B-splines (\cref{ch20:def-bspline}) in each direction, and their tensor products $B_a(t)B_b(s)$ as two-dimensional basis functions \ts{13}{10}{1:08:13}. Since B-splines are non-negative and sum to one (\cref{ch20:prop-bspline}), every combination is a
        reasonable surface, and its regularity is inherited from the basis (``every basis function makes sense, so any linear combination will be good enough'').
  \item \emph{More basis functions than instruments} ($N>M$): it is not known in advance which shapes are needed, and the number of unknowns is still tiny compared with the $28{,}920$ cells (with $24$ basis functions per direction, $576$ coefficients, so the matrix of the normal equations is $576\times576$, about $2.7$~MB) \ts{13}{10}{1:07:09}.
        The knots are denser at short forward times, where instruments are dense \ts{13}{10}{1:22:09}.
  \item \emph{Penalties.} The problem becomes \ts{13}{10}{1:09:19}
        \begin{equation}\label{ch20:eq-volpen}
          \|W(\bm q-\bm q^{\rm in})\|^2\to\min,\qquad \|L_1(\bm v-\bm v_0)\|^2\to\min,\qquad \|L_2\bm v\|^2\to\min ,
        \end{equation}
        with a penalty $L_1$ on the \emph{change} of the surface with respect to a prior and a penalty $L_2$ on the surface itself (slide~40). The gradient penalty is a matrix with rows $(\dots,1,-1,\dots)$ for every pair of neighbouring cells (slide~41, \ts{13}{10}{1:10:20}); on a rectangular
        grid of $n_t\times n_s$ nodes it is $L_1=\begin{psmallmatrix}D_t\otimes I\\ I\otimes D_s\end{psmallmatrix}$ with $D$ the one-dimensional first-difference matrices (on the triangle only the pairs of cells that both lie inside are kept).
\end{itemize}

Combining the three terms with weights $\lambda_1,\lambda_2$ gives (slide~42, \ts{13}{10}{1:10:20})
\begin{equation}\label{ch20:eq-volsol}
\begin{gathered}
  \bm x=\arg\min\Bigl\{\|W(J\bm x-\bm q^{\rm in})\|^2+\|\lambda_1L_1B\bm x\|^2+\|\lambda_2L_2\bm x\|^2\Bigr\}=K\,J\T W^2\bm q^{\rm in},\\
  K=\bigl(J\T W^2J+\lambda_1^2(L_1B)\T L_1B+\lambda_2^2L_2\T L_2\bigr)^{-1},
\end{gathered}
\end{equation}
which is \eqref{ch20:eq-gen} with the two penalty terms of \cref{ch20:sec-general} (setting the gradient of the objective to zero gives the normal equations; here $L_1B\bm x=\bm v-\bm v_0$ so the first penalty acts on the change of surface, and $L_2$ acts on the coefficients).
$K$ is invertible when the penalties together have no kernel in common with $WJ$; $\lambda_2>0$ with $L_2=I$ guarantees it. Two computational remarks from the Q\&A \ts{13}{10}{1:24:19}: only $J\T W^2\bm q^{\rm in}$ changes when market quotes tick, whereas the penalty matrices (and, if $J$ is
approximately constant, $J\T W^2J$) can be precomputed, so the calibration runs in real time on live data; and the linear algebra is done in the space of $N\ll28{,}920$ coefficients, never on the full surface. When the pricing map $\bm q(\bm x)$ is not linear, each Newton step solves a regularised problem of this form; the resulting damped iteration is a
Levenberg--Marquardt method (not discussed in class).

\begin{intuition}[Where in the chapter the pieces fit]
Formula \eqref{ch20:eq-volsol} is one object seen three ways. As \emph{statistics}: ridge regression (\cref{ch05:sec-shrink}) with a design matrix $WJ$, observation weights, and a prior that surfaces are smooth and close to $\bm v_0$. As \emph{linear algebra}: the SVD filter \eqref{ch20:eq-filter} on
the singular values of $WJ$ in the metric of the penalty. As \emph{engineering}: what a trader looks at is the residual: the surface is a smooth fit through the market, and what is left over is information.
\end{intuition}

Slide~43 shows the calibrated surface: smooth, with a few gentle peaks and a sharper one at short forward times. Asked about a peak, the speaker explains \ts{13}{10}{1:22:09}\ts{13}{10}{1:23:16} that the grid is uniform in quarters whereas the instruments are at
3~months, 6~months, 1, 2, 5, 10 years, so on a logarithmic axis the feature would not look sharp; the surface has more detail where there are more instruments, and no sharper features can be supported elsewhere.

\begin{finance}[Relative value: the calibrated surface as a benchmark]
Traders watch the surface as it moves with the live quotes \ts{13}{10}{1:24:19}. If an instrument trades away from the smooth surface implied by the others (something that ``is not typical'', a judgement that ``comes with years of experience'') it is a candidate mispricing: a spike is likely
to disappear soon, and one trades (a swaption, or a more exotic product with the appropriate exposure) against it. This is \emph{relative value analysis} and, in the speaker's words, how the desk makes money \ts{13}{10}{1:25:20}. It works only if the model is robust: a regulariser based on a fundamental
principle (smoothness) is more likely to be stable in the future than one estimated from a past window such as PCA \ts{13}{10}{1:26:24}. The same remarks apply to the swap curve.
\end{finance}

\begin{finance}[Interpolation, splines and their limits]
To a question on other fitting techniques the speaker answers that ``spline'' and ``interpolation'' are the same thing for him: one has limited inputs and wants the curve in between \ts{13}{10}{1:20:07}\ts{13}{10}{1:21:07}. The regularised fit goes beyond interpolation: it
does not pass through the data, which is the essential difference between \cref{ch20:sec-build} and \cref{ch20:sec-volsolve}.
\end{finance}

\paragraph{What regularisation buys, and costs} \ts{13}{10}{1:17:52}. It gives stability (indispensable for ill-conditioned problems) and a more realistic solution, at the expense of not fitting the data exactly and possibly of a \emph{biased} solution (for example a smoothness constraint makes the result smoother than reality). The bias
is controlled by \emph{what} is penalised, in the sense of the ``no bias on the kernel of $L$'' remark in \cref{ch20:sec-general}: open the textbook and one finds ``Tikhonov'' and penalises the amplitude, which would push the volatility towards zero; a flat surface is acceptable, so one penalises derivatives instead \ts{13}{10}{1:19:00}. The
slides' references are the Wikipedia pages on the HJM model, the yield curve, inverse problems, Tikhonov regularisation, the SVD, PCA and B-splines.

% ==================================================================
\begin{keyformulas}
\begin{itemize}
  \item Bond: $P=\sum_i\Lambda_iC_i$; yield $\Lambda_i=e^{-yt_i}$; duration $D=\sum_it_i\pi_i$, convexity $\mathcal C=\sum_it_i^2\pi_i$ \eqref{ch20:eq-yield}; bond spread $P=\sum_ie^{-st_i}\Lambda_iC_i$ \eqref{ch20:eq-spread}.
  \item Swap rate $C=\sum_ir_iw_i/\sum_iw_i$, $w_i=\delta_i\Lambda_i$ \eqref{ch20:eq-swaprate}.
  \item Cox--de Boor: $B_{i,k}=\frac{t-\tau_i}{\tau_{i+k}-\tau_i}B_{i,k-1}+\frac{\tau_{i+k+1}-t}{\tau_{i+k+1}-\tau_{i+1}}B_{i+1,k-1}$, $\sum_iB_{i,k}=1$.
  \item PCA hedge: $\bm x=R\T\bm r$, $R=P(H\T P)^{-1}$, $H\T R=I$ \eqref{ch20:eq-pcarisk}; smooth hedge: $R=(HH\T+\lambda^2(LJ)\T LJ)^{-1}H$ \eqref{ch20:eq-regRsol}.
  \item Inverse problem $\bm y=A\bm x+\bm\varepsilon$: $\bm x^\dagger=\sum_i\frac{\bm u_i\T\bm y}{s_i}\bm v_i$, error $\sum_i\frac{\bm u_i\T\bm\varepsilon}{s_i}\bm v_i$, $\kappa=s_1/s_r$ \eqref{ch20:eq-pinv}.
  \item Filters $\bm x_w=\sum_iw_i\frac{\bm u_i\T\bm y}{s_i}\bm v_i$: Tikhonov $w_i=\frac{s_i^2}{s_i^2+\lambda^2}$, $\bm x_\lambda=(A\T A+\lambda^2I)^{-1}A\T\bm y$; TSVD $w_i=\ind_{i\le k}$ \eqref{ch20:eq-tikh}; gain $\le1/(2\lambda)$; repricing $H_\lambda=\sum_iw_i\bm u_i\bm u_i\T\ne I$.
  \item General form $\min\|W(A\bm x-\bm y)\|^2+\lambda^2\|L\bm x\|^2$: $\bm x=(A\T W^2A+\lambda^2L\T L)^{-1}A\T W^2\bm y$; no bias on $\ker L$ \eqref{ch20:eq-gen}.
  \item MSE $=\sum_i[(1-w_i)^2\xi_i^2+\sigma^2w_i^2/s_i^2]+\sum_{i>r}\xi_i^2$; optimal $\lambda_i^2=\sigma^2/\xi_i^2$ \eqref{ch20:eq-mse}--\eqref{ch20:eq-wiener}.
  \item Choosing $\lambda$: discrepancy $\rho(\lambda)=\delta$; $\mathrm{GCV}=M\rho^2/(M-\operatorname{tr}H_\lambda)^2$.
  \item Surface calibration: $\bm v=\bm v_0+B\bm x$, $J=\partial\bm q/\partial\bm x$, $\bm x=KJ\T W^2\bm q^{\rm in}$ \eqref{ch20:eq-volsol}.
\end{itemize}
\end{keyformulas}

\begin{exercise}[(not from the course)]
\leavevmode
\begin{enumerate}[(a)]
  \item Show that the total number of cells of a triangular grid with $n$ quarters on a side is $n(n+1)/2$ and compute the memory (in GB, doubles) of the square matrix of the full problem for $n=240$.
  \item With a tensor-product cubic B-spline basis of $m$ functions per direction, how many coefficients, and what is the size of the matrix $K$ in \eqref{ch20:eq-volsol}?
  \item Explain why $J\T W^2J$ has rank at most the number of instruments, and why $K$ nevertheless exists.
\end{enumerate}
\end{exercise}

% !TEX root = ../main.tex
\chapter{Commodity Models}\label{ch:commodities}
\begin{paracol}{2}

\begin{sources}
\begin{tabularx}{\linewidth}{@{}l l L@{}}
\toprule
\textbf{Lecture} & \textbf{Speaker} & \textbf{Content}\\
\midrule
2013 L13 & A.~Eydeland (Morgan Stanley) & What quants do in commodity markets: contango and the Trafigura trade; storage as a portfolio of calendar-spread options and the storage-optimisation problem; the power plant as a strip of spark-spread options; fat tails and spikes of power prices; why spot models fail; the ``stack'' (supply and demand) model of power prices.\\
\bottomrule
\end{tabularx}

\smallskip
\textbf{Files.} Slides \file{MIT18_S096F13_lecnote13.pdf} (47 pages: 44 of content, then references and disclaimers). \textbf{Problem sets.} None; all exercises at the end
are \emph{not from the course}. \textbf{Year.} Part~II of these notes collects the 2013-only topics; there is no 2024 counterpart. The guest is Alexander Eydeland (the
slides cite his book with K.~Wolyniec, \emph{Energy and Power Risk Management}, Wiley 2002, and a chapter with H.~Geman).
\textbf{Not in the slides.} The whole discussion of the trade, of the bid for storage and of the correlation structure comes from the video; the slides give the formulas
and the pictures. The sections marked \emph{(not discussed in class)} are additions of these notes, needed to make the lecture's claims precise
(Sections~\ref{ch21:sec-carry}, \ref{ch21:sec-schwartz} and part of \ref{ch21:sec-stats}, \ref{ch21:sec-stack}).
\end{sources}

\section*{Motivation}
Every other lecture of the course prices financial claims whose underlying can be bought, held and sold at negligible cost. A commodity is different: oil sits in tanks, gas in
salt caverns, electricity nowhere (it cannot be stored at scale). The quantities that matter are therefore physical: how much can be stored and how fast it can be injected, how
efficiently a plant turns fuel into power, what the supply curve of generators looks like. Eydeland's message is that the mathematics of derivatives
(\cref{ch:bs}) applies, but the \emph{objects} to be valued are physical assets, and each asset turns out to be a portfolio of options on a \emph{spread}:
\begin{center}
\small\begin{tabular}{@{}l l l@{}}
\toprule
asset & underlying spread & option\\
\midrule
oil or gas storage & calendar spread $F(T_2)-F(T_1)$ & \Cref{ch21:sec-spread}\\
gas-fired power plant & power $-$ heat rate $\times$ gas (spark spread) & \Cref{ch21:sec-plant}\\
refinery & products $-$ crude (crack spread) & \Cref{ch21:sec-assets}, \cref{ch09:sec-crack}\\
\bottomrule
\end{tabular}
\end{center}
The plan of the chapter is that of the lecture: a trade in a contango market (\cref{ch21:sec-carry}); the value of storage, the Margrabe spread-option formula and how
to hedge it (\cref{ch21:sec-spread,ch21:sec-storage}); the merchant power plant (\cref{ch21:sec-plant}); why power prices resist the standard spot models
(\cref{ch21:sec-stats,ch21:sec-schwartz}) and the alternative that models supply and demand (\cref{ch21:sec-stack}). The chapter uses the forward and futures
language of \cref{ch:rates} (futures are margined, forwards are not) and of \cref{ch15:prop-fwd}, the Black--Scholes machinery of \cref{ch:bs}, and the
Ornstein--Uhlenbeck process of \cref{ch:sde}; the crack-spread example of \cref{ch09:sec-crack} is the same economic object as the refinery above.

% ==================================================================
\section{Contango, backwardation and cash-and-carry}\label{ch21:sec-carry}

The lecture opens with a news item of March 2009: the trading house Trafigura expects its best results ever \emph{because} oil prices are low and the market is in contango
\ts{13}{13}{1:06}. A year earlier oil had reached \$159 per barrel and traders had been blamed for it; in 2009 it fell to about \$30 and the same traders made record profits
\ts{13}{13}{2:07}. The apparent paradox is resolved by looking at the futures curve.

\subsection{The futures curve}
A \emph{futures contract} lets one buy today, at a price fixed today, a barrel to be delivered at a future date; the prices $F(t,T)$ for all maturities $T$ form the
\emph{futures} (or \emph{forward}) \emph{curve} \ts{13}{13}{3:09}. In commodity markets there is no ``spot price'' as such: what television calls the spot price is the price of the
nearest futures contract \ts{13}{13}{4:10}.

\begin{definition}[Contango and backwardation]\label{ch21:def-contango}
The curve $T\mapsto F(t,T)$ is in \emph{contango} if it is increasing and in \emph{backwardation} if it is decreasing \ts{13}{13}{4:10}. (Loosely, one says that the market is
in contango or backwardation when the curve is so at the front, i.e.\ when $F(t,T)\gtrless S_t$.)
\end{definition}

On 15 January 2009 the WTI curve was in steep contango: about \$35 for the nearest contract (February 2009), \$55 for August, \$58 for December and \$60 for February 2010
(\cref{ch21:fig-wti}; the intermediate points are read off the slide and are approximate).

\begin{figure}[ht]
\centering
\begin{tikzpicture}
\begin{axis}[width=0.82\linewidth,height=6.3cm,xlabel={delivery month},ylabel={\$/bbl},xmin=-0.5,xmax=13.5,ymin=32,ymax=64,
  xtick={0,3,6,9,12},xticklabels={Feb'09,May,Aug,Nov,Feb'10},ymajorgrids,grid style={black!12},tick label style={font=\small},label style={font=\small}]
\addplot[only marks,mark=o,red!70!black,mark size=2pt] coordinates {(0,35.4) (1,43.5) (2,48.1) (3,50.6) (4,52.4) (5,53.8) (6,54.8) (7,55.8) (8,56.7) (9,57.5) (10,58.4) (11,59.1) (12,59.9) (13,60.6)};
\addplot[thick,dashed,blue!60!black,domain=0:13,samples=27] {35.4*exp(0.10*x/12)};
\node[font=\footnotesize,blue!60!black,anchor=west] at (axis cs:7.2,40.7) {cash-and-carry bound, $r=10\%$};
\node[font=\footnotesize,red!70!black,anchor=east] at (axis cs:5.6,60.0) {WTI futures, 15 Jan 2009};
\end{axis}
\end{tikzpicture}
\caption{The WTI futures curve of 15 January 2009 (points read off slide~4 of the lecture) against the upper bound $S\,e^{rT}$ that a holder of free storage would impose with $r=10\%$
(\cref{ch21:prop-carry}). The gap between the two is what the market charges for storage.}\label{ch21:fig-wti}
\end{figure}

\subsection{The Trafigura trade}
Suppose one owns storage. Then on 15 January 2009 one can, at no risk \ts{13}{13}{5:11}\ts{13}{13}{6:16}:
\begin{enumerate}[1.]
  \item borrow \$35 and buy one barrel (the nearby price);
  \item store it for a year;
  \item sell one February-2010 futures contract, fixing the sale price at \$60.
\end{enumerate}
The lock-in is $60-35=\$25$ minus the interest on the loan: \$3.5 at 10\% simple interest, so a profit of \$21.5 per barrel (the lecturer says ``about 21''); with continuous compounding the
interest is $35(e^{0.10}-1)=3.68$ and the profit \$21.3. With 50--60 million barrels of storage the total is enormous and risk-free \ts{13}{13}{6:16}. The trade needs
low prices \emph{because} the interest is proportional to the price paid: at \$125--160 a barrel the same contango would leave about half the profit \ts{13}{13}{7:21}. What it needs
is \emph{storage}: the strategy is ``long the February 2009 contract, short the February 2010 contract'', i.e.\ a long \emph{calendar spread}, plus a tank to receive the
barrel \ts{13}{13}{7:21}.

\begin{proposition}[Cash-and-carry bounds]\label{ch21:prop-carry}
Let the interest rate $r$ and the storage cost rate $u$ (proportional to the price, per unit time) be constant, and let $S$ be the price of the physical commodity, which is
storable. In the absence of arbitrage for an agent that can store,
\[
  F(0,T)\ \le\ S_0\,e^{(r+u)T}.
\]
\end{proposition}
\begin{proof}
If $F(0,T)>S_0e^{(r+u)T}$, borrow $S_0$, buy the commodity, pay the storage cost, and sell the future. At $T$ deliver against $F(0,T)$ and repay the loan with its interest
and the accumulated storage cost $S_0e^{(r+u)T}$: the net cash flow $F(0,T)-S_0e^{(r+u)T}>0$ is certain and costs nothing. This is the argument of \cref{ch15:prop-fwd} with a storage cost.
\end{proof}
The opposite inequality $F\ge S_0e^{(r+u)T}$ does \emph{not} follow, because it would require selling the commodity today and buying it back later: only somebody who holds inventory can do that,
and holding inventory may itself be valuable (a refinery that runs out of crude stops). The shortfall is captured by the \emph{convenience yield} $y$ defined by
\begin{equation}\label{ch21:eq-convenience}
  F(0,T)=S_0\,e^{(r+u-y)T},\qquad y=r+u-\frac1T\ln\frac{F(0,T)}{S_0}\ \ge 0\ \text{ for a storable commodity.}
\end{equation}
(Not discussed in class.) Backwardation corresponds to a large convenience yield, $y>r+u$: inventory is scarce and valuable, spot is expensive against later delivery.

\begin{example}[The 2009 contango]\label{ch21:ex-implied}
With $S_0=35$, $F=60$ and $T=1$ year the implied cost of carry is $\ln(60/35)=54\%$ per year: at $r=10\%$ the bound of \cref{ch21:prop-carry} would require $u\ge44\%$ per year to remove the
arbitrage, i.e.\ $y=r+u-0.539<0$ for any reasonable $u$. The resolution is that the arbitrage is available only to those who own \emph{empty storage}; storage was scarce
(the market was oversupplied) and its shadow price was the difference $60-35e^{0.10}=21.3$ per barrel per year. Contango of this size is the market's bid for tanks.
\end{example}

\begin{finance}[Physical assets and \emph{who} can arbitrage]
The arbitrage of \cref{ch21:prop-carry} is not open to a financial investor: without storage one cannot take the physical leg. Contango therefore pays the owners of scarce
infrastructure (tanks, supertankers used as floating storage, salt caverns for gas) and the shape of the curve is best read as the price of the option to store. The remainder
of the lecture asks the inverse question: given the curve, \emph{what is a tank worth?}
\end{finance}

\begin{intuition}[Physicist's reading]
Think of the curve as a potential energy landscape along the time axis. A steep upward slope in $T$ is a gradient that an owner of storage can ``roll down'' at fixed cost: buy at the
bottom, sell at the top. Backwardation is the reversed slope: the tank owner is paid to empty the tank now and refill it later. The convenience yield is the price of the dissipation that keeps
the slope from exceeding $r+u$ when storage is plentiful.
\end{intuition}

% ==================================================================
\section{The value of storage: a calendar-spread option}\label{ch21:sec-spread}

\subsection{The bid for a tank}
The lecturer inverts the question \ts{13}{13}{7:21}\ts{13}{13}{8:21}. Your boss hands you a credit card and asks you to obtain storage from August to December in an auction, with the same
curve: $F_{\mathrm{Aug}}=55$, $F_{\mathrm{Dec}}=58$. Bid too little and you lose; bid too much and you suffer the winner's curse. The maximal bid is the amount that one can be \emph{sure to recover},
so a bidder needs, before bidding, a strategy that recovers the price paid \ts{13}{13}{9:23}.

\paragraph{The trader's answer.} Buy the August contract and sell the December contract on 1 January; on 1 August take delivery of the barrel and store it; in December deliver against the
December contract. This locks $58-55=3$ dollars per barrel: one can pay up to \$3 (say \$2.99, to keep a penny of profit) \ts{13}{13}{10:25}\ts{13}{13}{11:31}\ts{13}{13}{12:35}. Interest is ignored
``for simplicity'' (slide~9 adds: it should not be) \ts{13}{13}{11:31}. With interest at rate $r$ the locked amount, in December dollars, is $F_{\mathrm{Dec}}-F_{\mathrm{Aug}}e^{r(T_{\mathrm{Dec}}-T_{\mathrm{Aug}})}$, since the \$55 paid in August must be financed until the sale: at $r=10\%$ and four months the financing costs about
$55(e^{0.10/3}-1)\approx1.9$, more than half of the locked \$3, which is why the lecturer says that one must not really forget it.

\paragraph{The quant's answer.} On 1 January sell to somebody an \emph{August--December calendar spread option}, exercised on 31 July, with payoff
\begin{equation}\label{ch21:eq-payoff}
  \bigl[F_{\mathrm{Dec}}(T)-F_{\mathrm{Aug}}(T)\bigr]^+,\qquad T=\text{31 July},
\end{equation}
the futures prices being those quoted on the exercise date \ts{13}{13}{12:35}\ts{13}{13}{13:38}. There is no strike: the ``strike'' is the other futures price, so the option is a function of two random
variables and is more difficult than an option on a single stock \ts{13}{13}{14:39}. Using the formula of \cref{ch21:thm-margrabe} below, the lecture quotes $V=4.4677$ dollars per barrel
\ts{13}{13}{15:39}. So one can bid, for instance, \$4.20 \ts{13}{13}{15:39}: the bid beats the trader's ceiling of \$3, and the margin ($4.4677-4.20=0.27$) is much larger than a penny.

Why is the value larger than the \$3 of the trader? Because $3=F_{\mathrm{Dec}}-F_{\mathrm{Aug}}$ is the \emph{intrinsic value} of the option, its value at zero volatility, and an option is worth more than its
intrinsic value when the volatility is positive; energy volatilities are far larger than in rates, currencies or equity indices \ts{13}{13}{16:44}.

\paragraph{Where is the risk?} Selling the option exposes one to an unbounded payment. Suppose on 31 July $F_{\mathrm{Aug}}=55$ and $F_{\mathrm{Dec}}=80$ (slide~11 uses $65$ and $80$): the seller owes \$25 (slide: \$15)
\ts{13}{13}{17:47}. But he owns the tank: he buys August and sells December \emph{at the same moment} using futures, and later moves the physical barrel, collecting exactly $F_{\mathrm{Dec}}-F_{\mathrm{Aug}}$
\ts{13}{13}{18:48}\ts{13}{13}{19:54}. If instead December falls below August the option pays nothing and the tank is left unused \ts{13}{13}{20:55}. So the tank \emph{is} the option: a full hedge, whatever happens.

\subsection{The spread-option formula}
The value quoted on slide~10 is the formula for an option to exchange one lognormal asset for another (Margrabe, 1978), applied to two futures prices. The natural setting is that of \cref{ch15:sec-rn}:
futures are costless to enter and margined daily, so under the risk-neutral measure $\Q$ a futures price is a martingale (the futures analogue of the statement that the discounted stock is a martingale).

\begin{theorem}[Zero-strike spread option; Margrabe]\label{ch21:thm-margrabe}
Let $F_1,F_2$ be $\Q$-martingales with $\dd F_i=\sigma_iF_i\dd W_i$, $\dd W_1\dd W_2=\rho\dd t$, constant $\sigma_i>0$ and $|\rho|<1$, and let the discount factor to the expiry $T$ be $e^{-rT}$. Then
\begin{equation}\label{ch21:eq-margrabe}
  V=e^{-rT}\,\E_\Q\bigl[(F_2(T)-F_1(T))^+\bigr]=e^{-rT}\bigl[F_2(0)\,\Phi(d_1)-F_1(0)\,\Phi(d_2)\bigr],
\end{equation}
\[
  d_1=\frac{\ln\bigl(F_2(0)/F_1(0)\bigr)+\tfrac12\sigma^2T}{\sigma\sqrt T},\quad d_2=d_1-\sigma\sqrt T,\quad \sigma^2=\sigma_1^2+\sigma_2^2-2\rho\sigma_1\sigma_2,
\]
and $\Phi$ is the standard normal distribution function. (Slide 10, with ``DiscFactor'' $=e^{-rT}$ and $N=\Phi$.)
\end{theorem}
\begin{proof}
Write $(F_2-F_1)^+=F_1(T)\,(X_T-1)^+$ with $X=F_2/F_1$. Define a new measure $\Q^1$ on $\mathcal F_T$ by $\dd\Q^1/\dd\Q=F_1(T)/F_1(0)$, a positive random variable of mean one because $F_1$ is a
$\Q$-martingale. Then
\[
  \E_\Q\bigl[F_1(T)(X_T-1)^+\bigr]=F_1(0)\,\E_{\Q^1}\bigl[(X_T-1)^+\bigr].
\]
Under $\Q^1$ the ratio $X$ is a martingale: for $A\in\mathcal F_t$,
\[
\begin{gathered}
\E_{\Q^1}[X_T\ind_A]=\E_\Q\bigl[F_1(T)F_1(0)^{-1}F_2(T)F_1(T)^{-1}\ind_A\bigr]=\E_\Q[F_2(T)\ind_A]/F_1(0)\\
=\E_\Q[F_2(t)\ind_A]/F_1(0)=\E_{\Q^1}[X_t\ind_A].
\end{gathered}
\]
By Ito's formula $\dd X/X=\sigma_2\dd W_2-\sigma_1\dd W_1+(\text{drift})\dd t$, with quadratic variation $(\sigma_1^2+\sigma_2^2-2\rho\sigma_1\sigma_2)\dd t=\sigma^2\dd t$; the volatility of a process does not change when we change to an equivalent measure
(as recalled after \cref{ch15:eq-rnval}), and the drift is then forced by the martingale property. So $X$ is a driftless lognormal, $X_T=X_0e^{-\sigma^2T/2+\sigma\sqrt T Z}$ under $\Q^1$, and the Black formula with unit strike and zero rate (\cref{ch15:thm-bs})
gives $\E_{\Q^1}[(X_T-1)^+]=X_0\Phi(d_1)-\Phi(d_2)$. Multiplying by $F_1(0)$ gives~\eqref{ch21:eq-margrabe}.
\end{proof}

The proof used only that the \emph{ratio} is lognormal; the numeraire is the asset $F_1$ (this change-of-numeraire trick is the standard proof of Margrabe's formula). Three properties are immediate.

\begin{proposition}[Properties]\label{ch21:prop-props}
\leavevmode
\begin{enumerate}[(i)]
  \item $V\ge e^{-rT}(F_2(0)-F_1(0))^+$, with equality in the limit $\sigma\to0$: the option is worth at least its intrinsic value (Jensen, since $F_2-F_1$ is a martingale and $x\mapsto x^+$ is convex).
  \item $V$ is homogeneous of degree one in $(F_1,F_2)$ and increasing in $\sigma$, and depends on $(\sigma_1,\sigma_2,\rho)$ only through the \emph{spread volatility} $\sigma$; in particular a \emph{higher} correlation $\rho$ \emph{lowers} the value.
  \item The derivatives are $\partial V/\partial F_2=e^{-rT}\Phi(d_1)$ and $\partial V/\partial F_1=-e^{-rT}\Phi(d_2)$.
\end{enumerate}
\end{proposition}
\begin{proof}
(i) is Jensen.

(ii) Degree one is visible in~\eqref{ch21:eq-margrabe}; the vega is $e^{-rT}F_2\varphi(d_1)\sqrt T>0$ ($\varphi$ the normal density; the same computation as the vega of \cref{ch:bs}).

(iii) Differentiating,
$\partial_{F_2}V=e^{-rT}[\Phi(d_1)+(F_2\varphi(d_1)-F_1\varphi(d_2))\,\partial_{F_2}d_1]$ (using $d_2=d_1-\sigma\sqrt T$), and the bracket vanishes because
$\varphi(d_2)/\varphi(d_1)=e^{d_1\sigma\sqrt T-\sigma^2T/2}=F_2/F_1$. The other derivative follows by Euler's identity $V=F_1\partial_1V+F_2\partial_2V$ for degree-one functions.
\end{proof}

\paragraph{The numbers of slide 10.} The slide gives neither the volatilities, nor the correlation, nor the time to expiry. Taking the exercise date 31 July, $T=211/365=0.578$ years from 1 January, and
zero interest, one finds that $\sigma_1=\sigma_2=50\%$ and $\rho=0.95$ reproduce the printed value: $\sigma=0.5\sqrt{2(1-0.95)}=0.1581$ and
\[
  V=58\,\Phi(0.5019)-55\,\Phi(0.3817)=58\cdot0.6921-55\cdot0.6486=4.4677
\]
to four decimals (verified numerically; solving directly for $\sigma$ gives $0.15811$). This is our reconstruction, not something said in the lecture. The value splits as $3.00$ of intrinsic and $1.47$ of time value. The dependence on the correlation
is dramatic (\cref{ch21:tab-rho}): $\rho=1$ gives exactly the intrinsic value, and a fall of $\rho$ from $0.99$ to $0.95$ raises the value by $1.2$ dollars.

\begin{table}[ht]
\centering
\begin{tabular}{@{}l r r r r r r@{}}
\toprule
$\rho$ & $1$ & $0.99$ & $0.98$ & $0.95$ & $0.90$ & $0.80$\\
\midrule
$\sigma$ (spread vol.) & $0$ & $0.0707$ & $0.1000$ & $0.1581$ & $0.2236$ & $0.3162$\\
$V$ (\$/bbl) & $3.000$ & $3.259$ & $3.615$ & $4.468$ & $5.512$ & $7.037$\\
\bottomrule
\end{tabular}
\caption{Aug--Dec spread option, $F_{\mathrm{Aug}}=55$, $F_{\mathrm{Dec}}=58$, $T=211/365$, $r=0$, $\sigma_1=\sigma_2=50\%$. Computed for these notes.}\label{ch21:tab-rho}
\end{table}

\begin{figure}[ht]
\centering
\begin{tikzpicture}
\begin{axis}[width=0.8\linewidth,height=6cm,xlabel={spread volatility $\sigma$},ylabel={$V$ (\$/bbl)},xmin=0,xmax=0.6,ymin=2,ymax=12.5,ymajorgrids,grid style={black!12},tick label style={font=\small},label style={font=\small}]
\addplot[thick,blue!60!black] coordinates {(0.01,3.0000) (0.02,3.0001) (0.04,3.0280) (0.05,3.0792) (0.1,3.6148) (0.15,4.3425) (0.2,5.1305) (0.25,5.9434) (0.3,6.7680) (0.35,7.5985) (0.4,8.4318) (0.45,9.2659) (0.5,10.0995) (0.55,10.9318) (0.6,11.7620)};
\addplot[dashed,red!70!black,domain=0:0.6] {3};
\node[font=\footnotesize,red!70!black,anchor=north west] at (axis cs:0.31,3.0) {intrinsic value $=3$};
\node[font=\footnotesize,blue!60!black,anchor=east] at (axis cs:0.34,9.6) {Margrabe value};
\addplot[only marks,mark=*,mark size=2pt,black] coordinates {(0.15811,4.4677)};
\node[font=\footnotesize,anchor=north west] at (axis cs:0.16,4.3) {slide 10: $4.4677$};
\end{axis}
\end{tikzpicture}
\caption{Value of the August--December spread option as a function of the spread volatility, for $F=(55,58)$, $T=211/365$, no interest. The dot is the number printed on slide~10.}\label{ch21:fig-vol}
\end{figure}

\begin{remark}[Why correlation between contract months is the model's hard part]
The value of a calendar-spread option depends on the joint law of two points of the forward curve, in particular on their correlation (\cref{ch21:tab-rho}). A model that generates all forward prices from one random factor gives
$\rho=1$ and hence an option worth its intrinsic value: it cannot price the storage. This is the content of the slide-31 objection ``cannot model the correlation structure between forward contracts'' and is made
precise in \cref{ch21:sec-schwartz}.
\end{remark}

\begin{coursenote}[Errata and inconsistencies in this part]
\begin{itemize}
  \item Slide~11 repeats the example with $F_{\mathrm{Aug}}=65$, $F_{\mathrm{Dec}}=80$ (owe \$15) while the lecturer speaks of 55 and 80 (owe \$25); both are consistent examples.
  \item Slide~11 says the bid can be raised to \$4.46, the truncation of $4.4677$ (the lecturer says 447 cents and bids
      420).
  \item Slide~5 gives the profit of the Trafigura trade as \$21.5 with 10\% \emph{simple} interest, the speaker says ``21''.
  \item Slide~9 says ``interest rates are ignored for simplicity (should not be)'': formula~\eqref{ch21:eq-margrabe} holds with the discount factor, but the trader's static ceiling of \$3
      needs the financing correction above.
\end{itemize}
\end{coursenote}

% ==================================================================
\section{Hedging the tank, and the real storage problem}\label{ch21:sec-storage}

\subsection{There is no market: replicate}\label{ch21:sec-hedge}
The lecture makes a confession \ts{13}{13}{37:39}: there is no market for spread options, nobody will buy the option that the quant ``sold'' for 4.4677. What survives is the
\emph{method} behind Black--Scholes: the value is also the cost of a dynamic strategy that reproduces the payoff \ts{13}{13}{38:43}\ts{13}{13}{39:49}. The quant therefore does not sell anything; he owns the tank,
whose payoff is that of the option, and delta-hedges it, and he tells the trader what to do every day, and the tank operator when to inject and withdraw \ts{13}{13}{40:49}.

To make this precise (the lecture does not write it down) let $V(F_1,F_2)$ be the value~\eqref{ch21:eq-margrabe}. \emph{Owning} the option (the tank) and hedging means holding $-\Phi(d_1)$ Dec futures and $+\Phi(d_2)$
Aug futures, by \cref{ch21:prop-props}(iii) rebalanced continuously. Futures cost nothing to hold, so the hedge earns $-\Phi(d_1)\dd F_2+\Phi(d_2)\dd F_1$ and the terminal value of the whole position is
\[
\begin{gathered}
  \bigl(F_2(T)-F_1(T)\bigr)^+-\int_0^T\Phi(d_1)\dd F_2+\int_0^T\Phi(d_2)\dd F_1=V(F_1(0),F_2(0))\\
  \text{(continuous rebalancing, } r=0),
\end{gathered}
\]
the standard Black--Scholes replication statement (\cref{ch:bs}): Ito's formula applied to $V(F_1,F_2)$, which solves the two-dimensional Black equation with terminal
condition~\eqref{ch21:eq-payoff}, gives $\dd V=\partial_{F_1}V\dd F_1+\partial_{F_2}V\dd F_2$ (time-derivative and second-order terms cancel).
The hedger who wins the tank at price $P<V$ therefore locks in $V-P$. When the option ends in the money both deltas tend to one: the hedge is ``long August, short December'', exactly the
trader's static position of \cref{ch21:sec-spread}, which the tank then converts into physical barrels. In this sense the trader's strategy is the limit $\sigma\to0$ of the quant's.

\begin{table}[ht]
\centering
\begin{tabular}{@{}l r r r r r@{}}
\toprule
hedging dates $n$ & $1$ & $4$ & $16$ & $64$ & $211$\\
\midrule
mean of final P\&L & $4.462$ & $4.466$ & $4.472$ & $4.469$ & $4.468$\\
standard deviation & $2.033$ & $1.078$ & $0.567$ & $0.290$ & $0.159$\\
\bottomrule
\end{tabular}
\caption{Discrete delta-hedging of the option of \cref{ch21:tab-rho} ($\sigma_1=\sigma_2=50\%$, $\rho=0.95$; $V=4.4677$): terminal payoff plus hedge P\&L. Computed for these notes: $2\cdot10^4$ paths, $n$ equally spaced dates between 1 January and 31 July, no interest.}\label{ch21:tab-hedge}
\end{table}
The simulation (\cref{ch21:tab-hedge}) shows the two things one expects: the average of the final position is the premium $V$ whatever $n$, and the residual risk falls like $n^{-1/2}$, from $2.0$ dollars per barrel with a single rebalancing to $0.16$ with daily hedging.

\begin{remark}
The storage owner is \emph{long} the option, while the lecture speaks of ``selling'' it and covering the short position with the tank. The two views are the same portfolio seen from the two sides: sold option plus tank is riskless with value $V-P$; the hedge of the previous
paragraph reaches the same result without a counterparty. A subtlety: the option is settled on 31 July on the \emph{futures} prices, while the barrels move later, so between the two dates the operator is committed by the futures he has entered, which is why the recipe ``buy August, sell December at the same time'' works.
\end{remark}

\subsection{A portfolio of options under physical constraints}\label{ch21:sec-portfolio}
A real tank is available for two years, not for four months, and one may sell many spread options against it: injection in any month $T_i$ and withdrawal in any later month $T_j$, and also the opposite (withdraw first, refill later), which pays when the curve is
in backwardation \ts{13}{13}{21:58}\ts{13}{13}{27:11}. The first problem is to choose \emph{which portfolio of options} is the most valuable one that the tank can support, and this is an optimisation \ts{13}{13}{23:00}. The constraints are
technical, contractual and legal: injection and withdrawal rates, tank capacity, and above all the requirement that \emph{under no circumstances} the tank is asked to do something impossible \ts{13}{13}{24:02}\ts{13}{13}{25:05}. Slides 13--17 write the problem
as follows. With dates $T_1<\dots<T_N$ and today's futures prices $F_i$ (slide~13):
\begin{itemize}
  \item $S_{i,j}$ ($i<j$) is the value of the option to inject at $T_i$ and withdraw at $T_j$, with payoff $\max\{F_j-F_i-\mathrm{Cost},0\}$ at exercise; $U_{i,j}$ the value of the option to withdraw at $T_i$ and inject at the later
        $T_j$, with payoff $\max\{F_i-F_j-\mathrm{Cost},0\}$ (slide~14) \ts{13}{13}{27:11}\ts{13}{13}{28:16}. The volumes sold against the tank are $x_{i,j}\ge0$ and $\nu_{i,j}\ge0$;
  \item $y_i\ge0$ and $z_j\ge0$ are the volumes committed today by \emph{futures} for injection at $T_i$ and withdrawal at $T_j$ (the trader's static strategy);
  \item the objective is the cash collected now:
\begin{equation}\label{ch21:eq-storage-obj}
\begin{gathered}
  V=\max_{x,\nu,y,z}\Bigl\{\sum_{i<j}x_{i,j}S_{i,j}+\sum_{i<j}\nu_{i,j}U_{i,j}-\sum_iy_iF_i+\sum_jz_jF_j\Bigr\}\\
  \ts{13}{13}{25:05}\ts{13}{13}{26:11}
\end{gathered}
\end{equation}
        (the first two sums are premia received, the last two are the cost of the futures bought to inject and the revenue of the futures sold to withdraw);
  \item Boolean variables record the exercise: $\Omega^S_{i,j}=1$ if $S_{i,j}$ expires in the money (then the tank \emph{must} act), $0$ otherwise, and the same for $\Omega^U_{i,j}$ \ts{13}{13}{28:16}\ts{13}{13}{29:18}.
\end{itemize}
Define the net injection at date $T_i$ (positive means into the tank)
\begin{equation}\label{ch21:eq-qi}
  q_i=\sum_{j>i}x_{i,j}\Omega^S_{i,j}-\sum_{\ell<i}x_{\ell,i}\Omega^S_{\ell,i}+\sum_{\ell<i}\nu_{\ell,i}\Omega^U_{\ell,i}-\sum_{j>i}\nu_{i,j}\Omega^U_{i,j}+y_i-z_i .
\end{equation}
An exercised option of type~$S$ injects $x_{i,j}$ at $T_i$ and withdraws it at $T_j$, so the two terms cancel if both legs fall on the same date (a withdrawal already required cancels an injection required now, \ts{13}{13}{30:19}); for type~$U$ the signs
are reversed. The constraints of slides 16--17 are then
\begin{align}
  &-W_i\ \le\ q_i\ \le\ I_i, && i=1,\dots,N &&\text{(withdrawal and injection rates)}\label{ch21:eq-rates}\\
  &C^{\min}_i\ \le\ C_0+\sum_{k\le i}q_k\ \le\ C^{\max}_i, && i=1,\dots,N &&\text{(capacity)}\label{ch21:eq-cap}
\end{align}
where $C_0$ is the initial inventory. On slide~17 the capacity constraint is written $C_0+\sum_{k\le i}\{\sum_{j>i}x_{k,j}\Omega^S_{k,j}-\sum_{j>i}\nu_{k,j}\Omega^U_{k,j}+y_k-z_k\}$: this is the same thing, because in
$\sum_{k\le i}q_k$ every pair $(k,j)$ with $k<j\le i$ contributes $+x_{k,j}$ at $T_k$ and $-x_{k,j}$ at $T_j$ and cancels, so only the pairs still open at $T_i$ (those with $j>i$) survive.

\begin{remark}[Why the problem is ugly]
The variables $\Omega$ depend on prices at the exercise dates: they are random. The constraints must hold in every price scenario, so \eqref{ch21:eq-rates}--\eqref{ch21:eq-cap} are in effect a \emph{robust} problem, one system for each of the $2^{n_{\mathrm{opt}}}$ exercise patterns.
The unknowns are $x,y,z,\nu$ and, unfortunately, the Boolean $\Omega$, and the constraints are non-linear in them \ts{13}{13}{31:20}. For two years of monthly dates there are $24\cdot23/2=276$ pairs of each type, plus an equally large
set of constraints \ts{13}{13}{32:24}. (The lecturer stumbles over the count; this is the count of pairs $i<j$. The reading of the constraints as a robust problem is ours.)
\end{remark}

\subsection{The linear core: intrinsic value}\label{ch21:sec-intrinsic}
If the option premia are removed and only the futures part is kept, the problem~\eqref{ch21:eq-storage-obj}--\eqref{ch21:eq-cap} with $x=\nu=0$ is a \emph{linear program}: it gives the \emph{intrinsic value} of the tank, what can be locked in with futures alone
(the trader's strategy, made rigorous). (Not discussed in class.) Take twelve monthly dates with the approximate curve of \cref{ch21:fig-wti} (February 2009 to January 2010), a tank with $C_0=0$, $C^{\max}=100$, injection rate $I=30$ and withdrawal rate $W=40$ per month,
no costs. Maximising $\sum_iF_i(z_i-y_i)$ subject to $0\le y_i\le I$, $0\le z_i\le W$ and $0\le\sum_{k\le i}(y_k-z_k)\le C^{\max}$ gives
\begin{center}
\footnotesize
\setlength{\tabcolsep}{3.2pt}
\begin{tabular}{@{}l r r r r r r r r r r r r@{}}
\toprule
month & Feb & Mar & Apr & May & Jun & Jul & Aug & Sep & Oct & Nov & Dec & Jan\\
$F_i$ & 35.4 & 43.5 & 48.1 & 50.6 & 52.4 & 53.8 & 54.8 & 55.8 & 56.7 & 57.5 & 58.4 & 59.1\\
inject $y_i$ & 30 & 30 & 30 & 10 & 0 & 0 & 0 & 0 & 0 & 0 & 0 & 0\\
withdraw $z_i$ & 0 & 0 & 0 & 0 & 0 & 0 & 0 & 0 & 0 & 20 & 40 & 40\\
\bottomrule
\end{tabular}
\end{center}
and an intrinsic value of $5850-4316=1534$ dollars (average purchase price $43.16$, average sale price $58.50$). Without the rate limits one would fill the tank in February for $3540$ and empty it in January for $5910$ (value $2370$): the injection limit alone
costs $836$, and the LP finds the cheapest way around it, filling in the cheap months first. (Computed for these notes with \texttt{scipy.optimize.linprog}.) The quant's value is this plus the value of the flexibility, the time value of the spread options.

\subsection{Methods, robustness and calibration}\label{ch21:sec-methods}
The lecture lists three ways to attack the full problem \ts{13}{13}{32:24}\ts{13}{13}{33:31}: approximate it by a linear or quadratic programme; Monte Carlo simulation of the prices; or \emph{stochastic control}, for which the reference is Carmona and Ludkovski (2010),
``Valuation of energy storage: an optimal switching approach'', \emph{Quantitative Finance} 10(4) (slide~18): the decision inject/withdraw/hold is a control chosen as a function of the current prices and inventory, and the value is the solution of a dynamic-programming equation. None of the methods is ``better'': the
inputs (volatilities, correlations) are known only to a rough accuracy, and an exact method fed with parameters that are 50\% off is still wrong \ts{13}{13}{33:31}\ts{13}{13}{34:35}.

\begin{finance}[Robust, calibration-free models]
Eydeland's desk, in his words, does not use the stochastic-control solution (``we do something different'') and deliberately chooses a method that is \emph{robust}: little changes in the parameters must not change the value much, and each ingredient must be
something that can be traded in the market and verified by outside controllers and regulators, ``in this day and age extremely important'' \ts{13}{13}{34:35}\ts{13}{13}{35:35}. He warns against over-parametrised models (stochastic volatility, regime switching, several jump processes),
which capture the richness of price behaviour at the cost of stability \ts{13}{13}{35:35}. \emph{Calibration}, the fitting of unobservable parameters by least squares to market data, is fragile: the objective is non-convex, the solution set can be very flat, one may stop early
\ts{13}{13}{36:35}\ts{13}{13}{37:39}. Compare the regularisation of ill-posed calibration problems in \cref{ch:regularized} and the least-squares theory of \cref{ch:regression}.
\end{finance}

% ==================================================================
\section{The merchant power plant and the spark spread}\label{ch21:sec-plant}

\subsection{A plant is a strip of daily options}
The same reasoning applies to any physical asset: tankers, refineries, pipelines, power lines \ts{13}{13}{41:53}; only the constraints change (a tanker needs the routes and port limits) \ts{13}{13}{42:59}. The second half of the lecture works out the power plant.

A \emph{merchant} plant decides itself, every morning, whether to run \ts{13}{13}{42:59}\ts{13}{13}{44:00}. The price of power for the day is known; to produce one megawatt-hour it burns fuel. The efficiency is measured by the \emph{heat rate} $h$:
the number of MMBtu of fuel per MWh (typically 7 for a modern gas plant, up to 20 for an inefficient one that runs rarely) \ts{13}{13}{45:05}\ts{13}{13}{46:09}. With the fuel price $G$ (dollars per MMBtu) and the other variable costs $K$ (dollars per MWh: labour, cooling, \dots)
the profit of the day is (slide~20, where the heat rate is written in Btu/kWh, hence the division by 1000)
\begin{equation}\label{ch21:eq-spark}
  \Pi=\max\bigl\{P-hG-K,\ 0\bigr\},\qquad P=\text{power price (\$/MWh)},
\end{equation}
and the plant is run when the argument is positive \ts{13}{13}{47:13}\ts{13}{13}{48:18}. With $h=8$, $G=4$, $K=3$ the plant costs $8\cdot4+3=35$ per MWh: at $P=50$ the profit is 15 and the plant runs; at $P=30$ the owner ``goes back to sleep'' and earns zero
\ts{13}{13}{47:13}\ts{13}{13}{48:18}. So operating a merchant plant is financially equivalent to owning a portfolio of \emph{daily options on the spread} between electricity and fuel, the \emph{spark spread options} (slide~20).
The price the buyer should be willing to pay is the expectation of the sum of these payoffs, taken under the pricing measure, over the remaining life: a two-dimensional distribution of $(P,G)$, with its correlation \ts{13}{13}{50:21}\ts{13}{13}{51:25}.

\subsection{Pricing one daily option}
For $K=0$ the option \eqref{ch21:eq-spark} is a Margrabe option (exchange power for $h$ units of gas): \cref{ch21:thm-margrabe} with $F_2=P$ and $F_1=hG$ prices it. For $K>0$ there is no closed form; a classic approximation (Kirk, 1995; not in the lecture) treats
$hG+K$ as lognormal with volatility $\omega\sigma_G$, $\omega=hG/(hG+K)$, and prices a Margrabe option between $P$ and $hG+K$:
\begin{equation}\label{ch21:eq-kirk}
  V\approx e^{-rT}\bigl[P\,\Phi(d_1)-(hG+K)\Phi(d_2)\bigr],\quad
  \sigma_K^2=\sigma_P^2+\omega^2\sigma_G^2-2\rho\,\omega\sigma_P\sigma_G,
\end{equation}
with $d_{1,2}$ as in \cref{ch21:thm-margrabe} with $F_2/F_1\to P/(hG+K)$ and $\sigma\to\sigma_K$ ($P,G$ here are the forward prices for the delivery date). Numerically (forward prices $P=50$, $G=4$, $h=8$, $K=3$, $\sigma_P=80\%$, $\sigma_G=50\%$, $\rho=0.6$,
three months, $r=0$): intrinsic value $15.00$, Kirk's approximation $15.887$, Monte Carlo (lognormal $P,G$, $2\cdot10^6$ paths) $15.889\pm0.012$. The correlation matters (Kirk, same data): $\rho=0,0.3,0.6,0.9$ give $17.39,\ 16.65,\ 15.89,\ 15.20$. This is why the joint law of power and
gas is a modelling target and not a detail (slide~28: ``if the model does not capture this structure it may misprice spread options (tolling contracts, power plants, etc.)'').

\begin{remark}[What is idealised]
Equation~\eqref{ch21:eq-spark} ignores start-up costs, minimum run times, ramp rates and outages; these turn the strip of independent daily options into a path-dependent (real-option) problem, and it is for these that the Monte Carlo methods of
\cref{ch21:sec-methods} are needed. The lecture keeps to the daily-option picture.
\end{remark}

\begin{intuition}[Options without an underlying you can hold]
Power cannot be stored, so there is no cash-and-carry argument (\cref{ch21:prop-carry}) linking the forward price of power for a given day to the spot price today. The only no-arbitrage requirement is that forward prices be martingales under \emph{some} pricing measure,
$F(0,T)=\E_\Q[P_T]$, and the spot process does not have to be tied to the forward curve by a carry relation. This is the sense of the slide-31 objection ``no no-arbitrage argument: how to price forward contracts and options?''.
\end{intuition}

% ==================================================================
\section{Three assets, three spreads}\label{ch21:sec-assets}
The examples of the lecture are three instances of one construction: a physical asset gives the right, not the obligation, to transform one commodity (or the same commodity at one date) into another with a fixed conversion, so its value is that of an option on the difference of two price indices.

\begin{center}
\begin{tabular}{@{}l l l l@{}}
\toprule
asset & input $\to$ output & spread & exercise decision\\
\midrule
storage & barrel at $T_i$ $\to$ barrel at $T_j$ & $F(T_j)-F(T_i)-\text{cost}$ & inject/withdraw\\
gas power plant & $h$ MMBtu of gas $\to$ 1 MWh & $P-hG-K$ & run or stay off\\
refinery & 1 bbl of crude $\to$ products & $\text{products}-\text{crude}$ & run at what mix\\
\bottomrule
\end{tabular}
\end{center}
The refinery's spread is the crack spread of \cref{ch09:sec-crack}, where the ``cracks'' were the differences between gasoline or heating oil (converted to dollars per barrel with $42$ gallons per barrel) and crude, and where the three futures were found to be cointegrated. The link is useful in
two ways. First, cointegration says that the crack spread is mean reverting, so its variance does \emph{not} grow linearly with time (as in the lognormal spread of \cref{ch21:thm-margrabe}); for long-dated options the Margrabe formula overstates the volatility of a spread that is pulled back to an equilibrium
(not discussed in class). Second, the seasonal pull of heating oil in winter is a seasonal component of this spread, treated in \cref{ch21:sec-season} below.

% ==================================================================
\section{What power prices look like}\label{ch21:sec-stats}

To value the plant one needs the distribution of $P$ and $G$, and the lecture now turns to the data \ts{13}{13}{49:21}\ts{13}{13}{50:21}. The first idea is Brownian motion (or its exponential): the terminal distribution is then normal
(lognormal for the price) \ts{13}{13}{50:21}. The S\&P~500 is often said to have fat tails; the lecturer's point is that by commodity standards it looks normal \ts{13}{13}{51:25}.

\begin{table}[ht]
\centering
\begin{tabular}{@{}l r r r@{}}
\toprule
 & annual volatility & skewness & kurtosis\\
\midrule
Nord Pool (Scandinavian power), daily & $182\%$ & $1.468$ & $26.34$\\
Nord Pool, 6~p.m.\ price & $238\%$ & $2.079$ & $76.82$\\
DAX (equity index) & $23\%$ & $0.004$ & $3.33$\\
\bottomrule
\end{tabular}
\caption{Distribution parameters of returns (slide 23, from A.~Werner, \emph{Risk Management in the Electricity Market}, 2003). Kurtosis is $\E(X-\mu)^4/\sigma^4$, equal to $3$ for a normal law.}\label{ch21:tab-kurt}
\end{table}

The kurtosis of a normal variable is exactly~3; for the DAX it is $3.33$, for power $26$ to $77$ \ts{13}{13}{51:25}\ts{13}{13}{52:27}. And annual volatilities are in the hundreds of percent, against about 10\% for the S\&P at the time of the lecture
\ts{13}{13}{53:35}: intuition built on volatilities of 10--20\% fails at 100--300\% \ts{13}{13}{53:35}. The time series of Texas (ERCOT), New England (NEPOOL) and PJM prices on slides 24--26 show the feature that jumps out: \emph{spikes} \ts{13}{13}{52:27}. The slide-27 list of
what a model must capture is: mean reversion, spikes, high kurtosis, regime switching, lack of data and non-stationarity \ts{13}{13}{54:42}.

\subsection{Jumps produce kurtosis; aggregation removes it}\label{ch21:sec-jumpkurt}
(Not discussed in class.) A minimal fat-tailed model is a diffusion plus compound Poisson jumps. Let a log-return over a horizon $\Delta$ be $X_\Delta=\sigma\sqrt\Delta Z+\sum_{i=1}^{N_\Delta}J_i$ with $Z\sim N(0,1)$, $N_\Delta\sim\mathrm{Poisson}(\lambda\Delta)$ and $J_i\sim N(0,s^2)$, all independent.

\begin{proposition}[Excess kurtosis of jump-diffusion returns]\label{ch21:prop-jumpkurt}
The cumulants of $X_\Delta$ are $\kappa_2=(\sigma^2+\lambda s^2)\Delta$ and $\kappa_4=3\lambda s^4\Delta$, so
\[
  \frac{\E(X_\Delta-\E X_\Delta)^4}{\Var(X_\Delta)^2}-3=\frac{\kappa_4}{\kappa_2^2}=\frac{3\lambda s^4}{(\sigma^2+\lambda s^2)^2}\,\frac1\Delta .
\]
\end{proposition}
\begin{proof}
Cumulants of independent sums add. The Gaussian part contributes only to $\kappa_2$. For the compound Poisson sum, $\ln\E e^{\theta\sum J}=\lambda\Delta\bigl(\E e^{\theta J}-1\bigr)=\lambda\Delta\sum_{n\ge1}\theta^n\E J^n/n!$, so $\kappa_n=\lambda\Delta\,\E J^n$; here $\E J^2=s^2$, $\E J^4=3s^4$. The fourth cumulant of a variable with zero mean is $\E X^4-3(\E X^2)^2$.
\end{proof}
For $\lambda=20$ per year, $s=15\%$, $\sigma=40\%$ the excess kurtosis is $29.8$ for daily returns ($\Delta=1/365$: power trades every calendar day), $4.26$ weekly, $0.99$ monthly and $0.08$ yearly (Monte Carlo: $29.6$, $4.22$, $1.00$, $0.07$): a kurtosis of the size of the Nord Pool figures at the daily scale, disappearing on aggregation as the central limit
theorem requires (\cref{ch:probability}). Fat tails of \emph{daily} power prices therefore do not contradict a normal law for annual returns; the spikes must be modelled at the daily (or hourly) scale, and it is there that the option payoffs~\eqref{ch21:eq-spark} live.

\subsection{Mean reversion and the objections to spot models}
The lecturer's step-by-step argument against the standard toolbox \ts{13}{13}{56:47}\ts{13}{13}{57:52}\ts{13}{13}{58:56}:
geometric Brownian motion has no spikes, no correlation structure, nothing; a spike is a rise followed by a return, so add mean reversion; a rise from \$30 to \$1000 needs a \emph{jump}; the return from \$1000 to \$30 in three or four days needs a very strong reversion speed, but a reversion so strong
would freeze the price at the \$30 level in normal times; so one lets the parameters depend on the level, i.e.\ introduces regime switching, and one ends with 10--12 parameters ``impossible to manage'' \ts{13}{13}{58:56}\ts{13}{13}{1:00:02}. Stochastic volatility, regime switching and multiple jump processes are all used in the industry but
come from other markets and are not what the lecturer recommends \ts{13}{13}{1:00:02}.

A comment (ours): the scale of the reversion is easy to compute. For $\dd X=-\kappa X\dd t+\sigma\dd W$ (\cref{ch:sde}) one has $\mathrm{Corr}(X_t,X_{t+\Delta})=e^{-\kappa\Delta}$, half-life $\ln2/\kappa$ and stationary standard deviation $\sigma/\sqrt{2\kappa}$; a half-life of 1 day is $\kappa=365\ln2=253$ per year (calendar days) and of 4 days
$\kappa=63$ per year, with daily autocorrelation $0.50$ and $0.84$. The tension is really one of \emph{two time scales}: the spike relaxes in days, while the base level of prices, driven by fuel, moves over weeks and months. A one-factor mean-reverting model has a single scale.

\section{Spot models and forward curves}\label{ch21:sec-schwartz}
\subsection{The models of the slides}
Slide 29 lists the spot processes. All are written for the relative change $\dd S_t/S_t$ (here $W$ is a Brownian motion and $q$ a Poisson process of intensity $\lambda$ with jump size $Y_t$, $k=\E(Y-1)$):
\begin{align}
  \text{GBM:}\quad & \dd S_t/S_t=\mu\dd t+\sigma\dd W_t,\\
  \text{with mean reversion:}\quad & \dd S_t/S_t=\kappa(\theta-\ln S_t)\dd t+\sigma\dd W_t,\label{ch21:eq-mr}\\
  \text{with jumps:}\quad & \dd S_t/S_t=(\mu-\lambda k)\dd t+\sigma\dd W_t+(Y_t-1)\dd q_t,\label{ch21:eq-jump}\\
  \text{jumps and mean reversion:}\quad & \dd S_t/S_t=\bigl(\kappa(\theta-\ln S_t)-\lambda k\bigr)\dd t+\sigma\dd W_t+(Y_t-1)\dd q_t.\label{ch21:eq-jumpmr}
\end{align}
The subtraction of $\lambda k$ compensates the jumps so that the expected relative change stays $\mu\,\dd t$ (or zero from mean reversion). Equation~\eqref{ch21:eq-jumpmr} is written on slide~29 as $(\mu-\lambda k-\log S_t)\dd t$, which drops the reversion speed; we give the corrected form.
The more complicated models on slide~30 (stochastic convenience yield, stochastic volatility, regime switching, multiple jump processes, term-structure models) are only listed.

\begin{coursenote}[Erratum: slide 29, last model]
The drift of the last line should be $\kappa(\theta-\ln S_t)-\lambda k$ (or $\mu-\lambda k-\kappa\ln S_t$ with $\mu=\kappa\theta$); as printed, $-\log S_t$ has no coefficient, so it cannot be a speed of reversion.
Also, by Ito, $\dd\ln S=[\kappa(\theta-\ln S)-\sigma^2/2]\dd t+\sigma\dd W$: \eqref{ch21:eq-mr} is an Ornstein--Uhlenbeck process for $\ln S$ with long-run mean $\theta-\sigma^2/(2\kappa)$, not $\theta$.
\end{coursenote}

\subsection{The Schwartz one-factor model and its futures curve}
(Not discussed in class; the lecture only states~\eqref{ch21:eq-mr} and criticises it.) Write $X=\ln S$ and take the mean-reverting log price, in the form of Schwartz (1997),
\begin{equation}\label{ch21:eq-ou}
  \dd X_t=\kappa(\alpha-X_t)\dd t+\sigma\dd W_t,\qquad\kappa,\sigma>0,
\end{equation}
which is~\eqref{ch21:eq-mr} with $\theta=\alpha+\sigma^2/(2\kappa)$. It is solved by $X_T=\alpha+e^{-\kappa T}(X_0-\alpha)+\sigma\int_0^Te^{-\kappa(T-s)}\dd W_s$, which is Gaussian with mean $\alpha+e^{-\kappa T}(X_0-\alpha)$ and variance
$\sigma^2(1-e^{-2\kappa T})/(2\kappa)$ (multiply the equation by $e^{\kappa t}$ and integrate; the variance is the It\^o isometry, \cref{ch:stochcalc}).

A futures price is the expectation of the spot price at delivery under the pricing measure. Because commodity spot prices are not traded assets, the drift under $\Q$ is not fixed by no-arbitrage; assume a constant market price of risk $\lambda_{\mathrm{r}}$, so that $\dd W=\dd W^\Q-\lambda_{\mathrm{r}}\dd t$ and, under $\Q$,
$\dd X=\kappa(\alpha^*-X)\dd t+\sigma\dd W^\Q$ with $\alpha^*=\alpha-\sigma\lambda_{\mathrm{r}}/\kappa$ (the Girsanov change of drift; \cref{ch:sde}).

\begin{theorem}[Futures prices in the one-factor mean-reverting model]\label{ch21:thm-schwartz}
Under~\eqref{ch21:eq-ou} with risk-neutral long-run mean $\alpha^*$, for $\tau=T-t$,
\begin{equation}\label{ch21:eq-schwartz}
  F(t,T)=\E_\Q[S_T\mid\mathcal F_t]=\exp\Bigl\{e^{-\kappa\tau}X_t+(1-e^{-\kappa\tau})\,\alpha^*+\frac{\sigma^2}{4\kappa}\bigl(1-e^{-2\kappa\tau}\bigr)\Bigr\}.
\end{equation}
Moreover $\dd F(t,T)/F(t,T)=\sigma e^{-\kappa(T-t)}\dd W^\Q_t$.
\end{theorem}
\begin{proof}
Under $\Q$, $X_T$ given $\mathcal F_t$ is $N(m,v)$ with $m=\alpha^*+e^{-\kappa\tau}(X_t-\alpha^*)$ and $v=\sigma^2(1-e^{-2\kappa\tau})/(2\kappa)$. For a Gaussian $X_T$, $\E e^{X_T}=e^{m+v/2}$, which is~\eqref{ch21:eq-schwartz}. Since $F(t,T)$ is a $\Q$-martingale in $t$ with $F=\exp(\text{deterministic}+e^{-\kappa(T-t)}X_t)$, its
differential has no drift and the diffusion coefficient $e^{-\kappa(T-t)}\sigma F$ (only $X_t$ is random, and $\dd X_t$ has diffusion $\sigma$).
\end{proof}
Three consequences:
\begin{itemize}
  \item the curve tends, as $T\to\infty$, to $\exp(\alpha^*+\sigma^2/(4\kappa))$;
  \item contango or backwardation depends on where the spot sits;
  \item the volatility of a futures contract is $\sigma e^{-\kappa\tau}$, \emph{decreasing} with maturity (the ``Samuelson effect''): the front of the curve moves violently, the back hardly at all.
\end{itemize}
For contango versus backwardation differentiate~\eqref{ch21:eq-schwartz} in $T$ at $t=0$:
\[
\frac{\partial\ln F(0,T)}{\partial T}=\kappa e^{-\kappa T}(\alpha^*-X_0)+\frac{\sigma^2}{2}e^{-2\kappa T},
\]
so the curve slopes upward at maturity $T$ iff $X_0<\alpha^*+\sigma^2e^{-\kappa T}/(2\kappa)$. Compare with the convenience yield of~\eqref{ch21:eq-convenience}: as $T\to0$, $y-(r+u)=-\partial_T\ln F$ equals $\kappa(X_0-\alpha^*)-\sigma^2/2$, which is large when the spot is high relative to
its long-run level, the ``inventory is scarce'' behaviour. Mean reversion plays the role of a convenience yield that depends on the price level.

\begin{table}[ht]
\centering
\begin{tabular}{@{}l r r r r r@{}}
\toprule
$T$ (years) & $0.25$ & $0.5$ & $1$ & $2$ & $\infty$\\
\midrule
$F(0,T)$, $S_0=35$ & $39.68$ & $43.13$ & $47.36$ & $50.44$ & $51.35$\\
$F(0,T)$, $S_0=70$ & $63.90$ & $59.84$ & $55.28$ & $52.22$ & $51.35$\\
$\sigma_F(T)$ (volatility of the future) & $0.27$ & $0.19$ & $0.09$ & $0.02$ & $0$\\
\bottomrule
\end{tabular}
\caption{One-factor model with $\kappa=1.5$, $\alpha^*=\ln50$, $\sigma=0.4$ (so $\sigma_{\mathrm{spot}}=40\%$): futures curves from a low spot (contango) and a high spot (backwardation). Formula~\eqref{ch21:eq-schwartz} checked against a Monte Carlo of the exact Gaussian transition at $T=1$ ($47.358$ against $47.360$). Computed for these notes.}\label{ch21:tab-schwartz}
\end{table}

\begin{figure}[ht]
\centering
\begin{tikzpicture}
\begin{axis}[width=0.82\linewidth,height=6cm,xlabel={maturity $T$ (years)},ylabel={$F(0,T)$},xmin=0,xmax=3,ymin=30,ymax=75,ymajorgrids,grid style={black!12},tick label style={font=\small},label style={font=\small}]
\addplot[thick,intcol,mark=none] coordinates {(0.0,35.00) (0.2,38.85) (0.4,41.88) (0.6,44.22) (0.8,46.01) (1.0,47.36) (1.2,48.38) (1.4,49.14) (1.6,49.71) (1.8,50.13) (2.0,50.44) (2.2,50.68) (2.4,50.85) (2.6,50.98) (2.8,51.08) (3.0,51.15)};
\addplot[thick,thmcol,mark=none] coordinates {(0.0,70.00) (0.2,64.93) (0.4,61.27) (0.6,58.62) (0.8,56.69) (1.0,55.28) (1.2,54.25) (1.4,53.49) (1.6,52.93) (1.8,52.52) (2.0,52.22) (2.2,51.99) (2.4,51.82) (2.6,51.70) (2.8,51.61) (3.0,51.54)};
\addplot[dashed,black!60,domain=0:3] {51.35};
\node[font=\footnotesize,intcol,anchor=west] at (axis cs:0.9,40.5) {contango ($S_0=35$)};
\node[font=\footnotesize,thmcol,anchor=west] at (axis cs:0.9,64) {backwardation ($S_0=70$)};
\node[font=\footnotesize,black!70,anchor=south east] at (axis cs:3,43) {$e^{\alpha^*+\sigma^2/(4\kappa)}=51.35$};
\end{axis}
\end{tikzpicture}
\caption{Futures curves of the one-factor mean-reverting model (\cref{ch21:thm-schwartz}, parameters of \cref{ch21:tab-schwartz}): the same model gives contango or backwardation depending on the spot.}\label{ch21:fig-schwartz}
\end{figure}

\paragraph{Why one factor cannot price spread options.} Since $\dd F(t,T_i)/F(t,T_i)=\sigma e^{-\kappa(T_i-t)}\dd W^\Q$ involves the \emph{same} Brownian motion for every maturity, all futures returns are perfectly correlated, $\rho=1$, and the spread volatility of \cref{ch21:thm-margrabe} is
$\sigma_{\mathrm{spread}}=|\sigma_1-\sigma_2|=\sigma\lvert e^{-\kappa\tau_1}-e^{-\kappa\tau_2}\rvert$: with the parameters above, two contracts with $\tau=0.6$ and $0.9$ give $0.059$ and hence an option practically at its intrinsic value, against the $0.158$ that reproduces the number of slide~10. The slide-31 list of ``cons'' is now transparent: a spot model
with one factor \emph{cannot model the correlation structure between forward contracts} and cannot produce a complex volatility structure ($\sigma e^{-\kappa\tau}$ is monotone).

\subsection{Stochastic convenience yield: the two-factor model}\label{ch21:sec-gs}
Slide~30 lists ``models with stochastic convenience yield''. The prototype is Gibson and Schwartz (1990): the spot follows $\dd S/S=(\mu-\delta)\dd t+\sigma_1\dd W_1$ with a \emph{stochastic} convenience yield $\delta$, itself mean reverting, $\dd\delta=\kappa(\alpha-\delta)\dd t+\sigma_2\dd W_2$, $\dd W_1\dd W_2=\rho\dd t$.
Under $\Q$ the spot grows at $r-\delta$ and $\delta$ has risk-adjusted level $\hat\alpha$. Define $B(\tau)=(1-e^{-\kappa\tau})/\kappa$.

\begin{proposition}[Futures in the two-factor model]\label{ch21:prop-gs}
With constant $r$ and, under $\Q$, $\dd\ln S=(r-\delta-\sigma_1^2/2)\dd t+\sigma_1\dd W_1$, $\dd\delta=\kappa(\hat\alpha-\delta)\dd t+\sigma_2\dd W_2$,
\begin{equation}\label{ch21:eq-gs}
\begin{gathered}
  \ln F(0,T)=\ln S_0-\delta_0B+\Bigl(r-\hat\alpha+\frac{\sigma_2^2}{2\kappa^2}-\frac{\rho\sigma_1\sigma_2}{\kappa}\Bigr)T\\
  +\Bigl(\hat\alpha-\frac{\sigma_2^2}{\kappa^2}+\frac{\rho\sigma_1\sigma_2}{\kappa}\Bigr)B+\frac{\sigma_2^2}{4\kappa^3}\bigl(1-e^{-2\kappa T}\bigr),\qquad B=B(T),
\end{gathered}
\end{equation}
and $\dd F(t,T)/F(t,T)=\sigma_1\dd W_1-\sigma_2B(T-t)\dd W_2$.
\end{proposition}
\begin{proof}
Integrate: $\delta_s=\delta_0e^{-\kappa s}+\hat\alpha(1-e^{-\kappa s})+\sigma_2\int_0^se^{-\kappa(s-u)}\dd W_2(u)$, so $\int_0^T\delta_s\dd s=\delta_0B+\hat\alpha(T-B)+\sigma_2\int_0^TB(T-u)\dd W_2(u)$ (Fubini). Hence $\ln S_T$ is Gaussian with mean $m=\ln S_0+(r-\sigma_1^2/2)T-\delta_0B-\hat\alpha(T-B)$ and variance
$v=\sigma_1^2T+\sigma_2^2\int_0^TB^2-2\rho\sigma_1\sigma_2\int_0^TB$, where $\int_0^TB(\tau)\dd\tau=(T-B)/\kappa$ and $\int_0^TB^2=\kappa^{-2}\bigl[T-2B+(1-e^{-2\kappa T})/(2\kappa)\bigr]$. Then $F=\E_\Q S_T=e^{m+v/2}$; collecting the coefficients of $T$, of $B$ and the exponential term gives~\eqref{ch21:eq-gs}. The
volatility statement follows from $F(t,T)=S_t\exp(-\delta_tB(T-t)+\text{deterministic})$.
\end{proof}
\emph{Checks} (ours): for $\sigma_2=0$ and constant $\delta$ one recovers $F=S_0e^{(r-\delta)T}$, the cost-of-carry formula~\eqref{ch21:eq-convenience} with $u=0$; and a Monte Carlo of the two SDEs for $S_0=50$, $\delta_0=0.08$, $r=0.03$, $\kappa=1.2$, $\hat\alpha=0.05$, $\sigma_1=0.35$, $\sigma_2=0.5$, $\rho=0.6$, $T=1$ gives $47.302$ against $47.304$ from~\eqref{ch21:eq-gs}
($4\cdot10^5$ paths, 400 Euler steps); the curve is $49.31,\ 48.58,\ 47.30,\ 45.54,\ 42.44$ for $T=0.25,0.5,1,2,5$.

The two Brownian motions decorrelate the front and the back of the curve: from the volatility formula, $\sigma_F(\tau)^2=\sigma_1^2+\sigma_2^2B^2-2\rho\sigma_1\sigma_2B$ and the correlation between contracts at $\tau_1,\tau_2$ is
\[
  \mathrm{Corr}=\frac{\sigma_1^2+\sigma_2^2B(\tau_1)B(\tau_2)-\rho\sigma_1\sigma_2\bigl(B(\tau_1)+B(\tau_2)\bigr)}{\sigma_F(\tau_1)\,\sigma_F(\tau_2)};
\]
with the parameters above this is $0.977$ for $(\tau_1,\tau_2)=(0.6,0.9)$, $0.966$ for $(1,2)$ and $0.77$ for $(0.5,3)$. So calendar spreads acquire a volatility of their own, of the order of the values needed in \cref{ch21:tab-rho}.

\subsection{Seasonality}\label{ch21:sec-season}
(Not discussed in class beyond the observation of heating oil in winter, \cref{ch09:sec-crack}.) Energy prices have a deterministic seasonal pattern: gas and heating oil dearer in winter, power in summer and winter peaks. The simplest device is $\ln S_t=f(t)+X_t$, $f$ a deterministic periodic function (for example a truncated Fourier series in $t$ with period one year) and $X$ a mean-reverting process as
in~\eqref{ch21:eq-ou}. Then $F(t,T)=e^{f(T)}\cdot\E_\Q[e^{X_T}\mid\mathcal F_t]$, so the seasonal shape is multiplied onto the non-seasonal curve of~\eqref{ch21:eq-schwartz}, and $f$ can be chosen so that the model reproduces \emph{exactly} the observed forward curve today. This is the mechanism by which the lecture's constants
$\alpha_1,\alpha_2,\alpha_3$ (next section) ``are chosen to match market data'': shape parameters absorb what the model does not explain and leave the \emph{dynamics} to the model.

% ==================================================================
\section{Why spot models fail for power, and the stack alternative}\label{ch21:sec-stack}

The one- and two-factor models of \cref{ch21:sec-schwartz} are the standard toolkit for oil and gas, but the lecture is explicit that for \emph{power} they are the wrong starting point
\ts{13}{13}{54:12}. Electricity cannot be stored (at grid scale): there is no cash-and-carry argument linking today's price to tomorrow's, no arbitrage forces the spot process to
look like anything in particular, and the empirical price series shows huge, short-lived spikes (\cref{ch21:tab-kurt}) that no diffusion, and not even the simple jump-diffusion
of \eqref{ch21:eq-jump}, reproduces convincingly: a jump-diffusion adds symmetric or one-sided jumps to a smooth base, but real spikes are triggered by a specific, identifiable
mechanism -- the market running out of cheap generation -- and revert almost immediately once it eases \ts{13}{13}{56:30}.

\begin{intuition}[a market with no inventory]
Every other underlying in this course can be held: a bond, a stock, a barrel of oil in a tank. Holding is what forces $F(t,T)\le S_te^{(r+u)(T-t)}$ (\cref{ch21:eq-convenience}):
if the forward is too rich you buy spot, store, and sell forward. Power has no ``store’’: you cannot buy megawatt-hours at 3\,a.m.\ and sell them at 6\,p.m.\ instead. The only
link between today's price and the value of a future delivery is \emph{physical} (what plants will be running then), not financial. This is why the lecture turns from
diffusions to an engineering model of the market itself.
\end{intuition}

\subsection{The merit-order stack}\label{ch21:sec-merit}
The alternative is to model the market micro-structure directly \ts{13}{13}{57:40}. At each instant a set of generating units is available, each with a capacity $q_i>0$ (MW) and
a marginal cost $c_i$ (\$/MWh, the fuel cost divided by the plant's efficiency, i.e.\ its \emph{heat rate} $h_i$ times the fuel price, plus a small operating cost); a system
operator dispatches units in increasing order of $c_i$ (cheapest first: nuclear and hydro, then coal, then gas turbines of increasing heat rate, then oil peakers) until
cumulative capacity meets demand $D$. The market-clearing price is the cost of the \emph{last}, marginal unit dispatched:
\begin{equation}\label{ch21:eq-stack}
P(D)\ =\ c_{(k)},\qquad k=\min\Bigl\{j:\ \textstyle\sum_{i=1}^{j}q_{(i)}\ge D\Bigr\},
\end{equation}
where $c_{(1)}\le c_{(2)}\le\dots$ is the cost-sorted (``stacked'') list of units. Plotting cumulative capacity on the horizontal axis and $c_{(i)}$ on the vertical gives the
\emph{supply stack}, a non-decreasing step function; \eqref{ch21:eq-stack} says the price is where the vertical line at $D$ crosses the stack (not discussed in class beyond the
qualitative picture; the step-function formalisation is ours).

\begin{proposition}[Two sources of randomness, one nonlinear map]\label{ch21:prop-stack}
Let $D_t$ (demand, driven mainly by temperature and the business cycle) and $G_t$ (the marginal fuel price, typically natural gas) be the two state variables. If a
gas unit of heat rate $h$ is marginal, $c=hG_t+m$ ($m$ a small fixed cost); \eqref{ch21:eq-stack} makes $P_t=\Phi(D_t,G_t)$ for a function $\Phi$ that is piecewise linear
in $G_t$ (for fixed $D_t$, as long as the \emph{same} unit stays marginal) but is a non-decreasing step function of $D_t$, jumping up wherever dispatch moves to the next unit on the stack, and an
essentially vertical segment once the stack is exhausted (\emph{price spike}: any further demand can only be met by very expensive last-resort generation, demand-side curtailment,
or a shortage price cap).
\end{proposition}
\begin{proof}[Justification (not discussed in class)]
Immediate from \eqref{ch21:eq-stack}: between two consecutive kink points $D\in[Q_{k-1},Q_k)$ the marginal unit is fixed at $k$, so $P=c_{(k)}(G)$, linear (indeed usually
constant, or linear if unit $k$ itself burns gas); at $D=Q_k$ the marginal unit switches to $k+1$ and $P$ jumps (it is flat in $D$ while a single unit is
marginal, in a discrete stack) by $c_{(k+1)}-c_{(k)}$; if $D>\sum_iq_{(i)}$ (total capacity), \eqref{ch21:eq-stack} has no solution and the price is set
by the market's scarcity mechanism (an administrative cap, e.g.\ \$1{,}000/MWh or higher in some US markets, or unserved load).
\end{proof}

\begin{example}[A numerical stack]\label{ch21:ex-stack}
Take a caricature grid: $400$\,MW nuclear at \$8, $300$\,MW coal at \$18, $400$\,MW of gas units with heat rates $7,7.5,8,8.5$ ($100$\,MW each, cost $hG+3$), $200$\,MW of
peaking gas with heat rates $11,13,15,18$ ($50$\,MW each), and two $40$\,MW oil peakers at \$150 and \$300 (a stand-in for a price cap). With demand $D_t=950+130X_t$
($X_t$ an AR(1) with $\phi=0.8$, roughly a daily load shape) and $\ln G_t$ an OU process ($\phi=0.995$ per period, $\sigma=0.03$, so a slowly wandering gas price around
$G_0=4$), simulating $2\times10^4$ periods gives a price series with mean \$46.5, median \$39.7, excess kurtosis $47$ (against $3$ for a Gaussian) and skewness $4.3$: the
distribution is exactly of the shape of \cref{ch21:tab-kurt}, with a heavy right tail generated \emph{purely} by the nonlinearity of $\Phi$, not by any exogenous jump process
imposed on the price. Conditional on the marginal unit, price and gas price are strongly correlated ($\mathrm{Corr}\approx0.94$ when a gas unit is marginal in the mid-stack); once
demand exceeds the whole stack the price is capped at \$500 regardless of $G_t$: the correlation with gas \emph{breaks down} exactly in the regime that matters most for a hedger,
another symptom of \Cref{ch21:prop-stack}. (Simulation for these notes; the numbers are illustrative, not calibrated to a real grid.)
\end{example}

\begin{figure}[htb]
\centering
\begin{tikzpicture}
\begin{axis}[width=0.95\linewidth,height=5.4cm,xlabel={$t$ (days)},ylabel={price (\$/MWh)},ymode=log,ymin=15,ymax=600,
  grid=major,grid style={black!10}]
  \addplot[thick,color=defcol,mark=none] coordinates {(0,46.52) (1,48.73) (2,48.78) (3,45.54) (4,45.31) (5,40.87) (6,41.07) (7,44.47) (8,47.25) (9,45.21) (10,42.4) (11,44.45) (12,48.02) (13,42.48) (14,18.0) (15,42.8) (16,18.0) (17,45.2) (18,45.38) (19,46.92) (20,47.81) (21,48.85) (22,45.0) (23,49.12) (24,48.8) (25,51.49) (26,66.37) (27,52.64) (28,53.52) (29,51.91) (30,65.06) (31,50.24) (32,81.87) (33,60.49) (34,61.49) (35,48.31) (36,49.8) (37,49.43) (38,47.73) (39,46.8) (40,48.65) (41,47.07) (42,47.57) (43,43.28) (44,46.69) (45,47.39) (46,46.31) (47,41.79) (48,40.53) (49,42.8) (50,44.79) (51,44.45) (52,45.02) (53,42.71) (54,41.49) (55,42.49) (56,43.28) (57,44.81) (58,45.16) (59,45.56) (60,65.9) (61,44.43) (62,42.45) (63,54.04) (64,51.84) (65,57.36) (66,38.46) (67,38.1) (68,38.29) (69,38.99) (70,36.95) (71,35.62) (72,37.57) (73,49.0) (74,38.61) (75,49.15) (76,57.08) (77,37.84) (78,32.31) (79,34.72) (80,43.84) (81,34.16) (82,34.78) (83,32.16) (84,31.78) (85,34.95) (86,33.72) (87,37.13) (88,35.46) (89,38.75) (90,41.71) (91,60.92) (92,53.55) (93,60.73) (94,40.85) (95,38.58) (96,40.18) (97,49.71) (98,39.12) (99,39.6) (100,38.54) (101,34.62) (102,36.01) (103,36.88) (104,37.51) (105,35.17) (106,34.33) (107,33.93) (108,35.32) (109,53.34) (110,71.34) (111,150.0) (112,59.13) (113,44.8) (114,34.87) (115,33.52) (116,30.19) (117,32.15) (118,33.42) (119,32.9) (120,33.12) (121,30.76) (122,30.38) (123,29.29) (124,28.15) (125,29.46) (126,30.48) (127,31.68) (128,31.83) (129,29.65) (130,29.17) (131,30.55) (132,33.22) (133,32.04) (134,35.77) (135,34.63) (136,37.55) (137,40.15) (138,41.65) (139,39.16) (140,37.84) (141,37.81) (142,41.58) (143,56.53) (144,38.89) (145,39.12) (146,39.05) (147,36.59) (148,38.25) (149,38.61) (150,38.52) (151,38.76) (152,37.22) (153,36.69) (154,18.0) (155,40.34) (156,38.73) (157,38.71) (158,39.17) (159,39.75) (160,44.51) (161,44.36) (162,48.64) (163,53.26) (164,80.2) (165,75.71) (166,100.94) (167,71.55) (168,68.9) (169,78.15) (170,80.81) (171,99.38) (172,97.26) (173,300.0) (174,500.0) (175,150.0) (176,500.0) (177,150.0) (178,99.91) (179,79.67)};
\end{axis}
\end{tikzpicture}
\caption{$180$-day window of the simulated stack price of \Cref{ch21:ex-stack} (log scale). Sharp up-spikes occur when demand exceeds the stack (capped at \$500) or dips to a low-cost
segment (the \$18 troughs, an artifact of the toy grid); between spikes the series reverts to the \$30--\$50 range in a few periods, exactly the pattern of \Cref{ch21:tab-kurt}.
Simulated for these notes.}\label{ch21:fig-stack}
\end{figure}

\begin{finance}[trading the stack, not the spot]
Because $\Phi$ is known (approximately) to a utility that owns the plants and knows their costs, a merchant generator does not need a spot \emph{model} of power at all to
value physical assets: \cref{ch21:sec-plant}'s spark-spread option is priced from $F_{\mathrm{power}}$ and $F_{\mathrm{gas}}$, i.e.\ market-observed \emph{forward} prices, and
the stack only enters as an internal simulation tool (to size hedges, to stress-test the effect of an outage, or to price options that depend on the \emph{path} of spikes,
which forward prices alone do not pin down) \ts{13}{13}{58:50}. This is the practical resolution of the lecture's warning that spot models fail: one does not need the spot model
to be ``right’’ everywhere, only to be a useful simulator of the risk that the forward market does not already hedge away.
\end{finance}

\begin{coursenote}[what the slides do and do not formalise]
The slides show the stack picture and assert its consequences (fat tails, spikes, the breakdown of price/fuel correlation at the top of the stack) without deriving them; the
formal statement of \Cref{ch21:prop-stack}, \Cref{ch21:ex-stack} and \cref{ch21:fig-stack} are additions of these notes built to make the qualitative claims precise and checkable.
No specific numerical grid was given in the lecture.
\end{coursenote}

\begin{keyformulas}
\small
\begin{align*}
&\text{Cash-and-carry:} && F(t,T)=S_te^{(r+w-y)(T-t)}\ \ (\text{contango if }r+w>y)\\
&\text{Margrabe (zero-strike spread):} && V=e^{-rT}\bigl[F_2\Phi(d_1)-F_1\Phi(d_2)\bigr],\ \ \sigma=\sqrt{\sigma_1^2+\sigma_2^2-2\rho\sigma_1\sigma_2}\\
&\text{Storage value:} && \max\nolimits_{\{q_i\}}\ \textstyle\sum_iq_iF(t,T_i)-\text{costs}\\
&\text{(rate/capacity limits)} && \text{s.t.\ }\eqref{ch21:eq-rates}\text{--}\eqref{ch21:eq-cap}\\
&\text{Spark spread:} && \Pi=\bigl(P-H\!\cdot\!G-c\bigr)^+,\quad V\ \text{Kirk-approximated,~}\eqref{ch21:eq-kirk}\\
&\text{One-factor Schwartz:} && F(t,T)=S_t\exp\bigl[\text{deterministic function of }X_t,\tau\bigr],\ \ \eqref{ch21:eq-schwartz}\\
&\text{Two-factor (stoch.\ conv.\ yield):} && F(t,T)=S_t\exp\bigl[-\delta_tB(\tau)+A(\tau)\bigr],\ \ \eqref{ch21:eq-gs}\\
&\text{Stack price:} && P(D)=c_{(k)},\ \ \eqref{ch21:eq-stack}\ \ (\text{non-decreasing step function of }D)
\end{align*}
\end{keyformulas}

\section*{Exercises}
\addcontentsline{toc}{section}{Exercises}

\begin{exercise}[not from the course: implied convenience yield]\label{ch21:ex-impliedyield}
On a given day $S_t=\$70$, $r=3\%$, and the one-year futures trades at $\$66$.
\begin{enumerate}[(a)]
  \item Using \eqref{ch21:eq-convenience} with negligible storage cost, find the implied convenience
      yield $y$.
  \item If storage actually costs $w=2\%$ per year, what is $y$ now?
  \item Is the market in contango or backwardation, and is this consistent with a positive convenience
      yield?
\end{enumerate}
\end{exercise}

\begin{exercise}[not from the course: Margrabe with correlated calendar legs]\label{ch21:ex-margrabe-num}
Two futures $F_1=60$ (near), $F_2=63$ (far), $\sigma_1=\sigma_2=45\%$, $T=0.5$, $r=0$.
\begin{enumerate}[(a)]
  \item Compute the spread option value for $\rho=0.3,0.7,0.95$ using
      \Cref{ch21:thm-margrabe} and comment on the monotonicity in $\rho$ found in \cref{ch21:tab-rho}.
  \item Recompute for $\rho=1$: show the value reduces to the intrinsic
      value $(F_2-F_1)^+$ and explain why in terms of \Cref{ch21:prop-props}.
\end{enumerate}
\end{exercise}

\begin{exercise}[not from the course: intrinsic vs extrinsic value of storage]\label{ch21:ex-intrinsic}
A storage facility can inject/withdraw at most $10$ units per month and holds at most $50$ units, currently empty. Given monthly futures prices
$F=(20,24,19,26,21,18,25,30,22,19,23,27)$ (a stylised seasonal curve)
\begin{enumerate}[(a)]
  \item find the intrinsic value \eqref{ch21:eq-storage-obj} by inspection (buy at the cheapest months
      subject to the capacity/rate constraints, sell at the priciest)
  \item explain qualitatively, without solving it, why the \emph{extrinsic} value (\cref{ch21:sec-methods}) would be
      positive even if this particular forward curve were flat.
\end{enumerate}
\end{exercise}

\begin{exercise}[not from the course: heat rate and the spark spread]\label{ch21:ex-heatrate}
A gas plant has heat rate $H=8$ (MMBtu/MWh) and variable cost $c=\$2$/MWh.
\begin{enumerate}[(a)]
  \item At $F_{\mathrm{power}}=\$45$/MWh and $F_{\mathrm{gas}}=\$4$/MMBtu, is the plant in the
      money on a forward basis?
  \item Using Kirk's approximation \eqref{ch21:eq-kirk} with $\sigma_P=60\%$, $\sigma_G=35\%$, $\rho=0.5$, $T=1$\,month, price one spark-spread option
      (one contract = one hour of capacity).
  \item Explain, using \Cref{ch21:prop-stack}, why $\rho$ between power and gas should be expected to fall when demand is very high (the
      top of the stack switches to oil or curtailment, decoupling power from the gas price).
\end{enumerate}
\end{exercise}

\begin{exercise}[not from the course: jumps and aggregation]\label{ch21:ex-jumpagg}
Using \Cref{ch21:prop-jumpkurt}'s formula for the excess kurtosis of a jump-diffusion's returns over an interval $\dd t$,
\begin{enumerate}[(a)]
  \item show that it vanishes as $\dd t\to\infty$ for
      fixed jump intensity $\lambda$ and jump size distribution, i.e.\ that kurtosis is a \emph{high-frequency} phenomenon, consistent with the large \emph{daily} kurtosis of \cref{ch21:tab-kurt} and the aggregation numbers of \cref{ch21:sec-jumpkurt}.
  \item Explain why the stack model's spikes (\cref{ch21:fig-stack}) can average away much more slowly than independent jumps when high demand is persistent (a heat wave lasting weeks keeps the system at the top of the stack), and what this implies
      for someone trying to fit a constant-parameter jump-diffusion to both daily and monthly power returns.
\end{enumerate}
\end{exercise}

\begin{exercise}[not from the course: three-asset crack spread and the stack]\label{ch21:ex-crackstack}
Consider a refinery (crack spread, \cref{ch21:sec-assets}) sitting downstream of a gas-fired power plant that competes for the same natural-gas supply.
\begin{enumerate}[(a)]
  \item Using the
      cointegration result of \cref{ch09:sec-crack}, argue that a shock to the crude price is transmitted to the crack spread only through its \emph{stationary} component.
  \item Using
      \Cref{ch21:prop-stack}, argue instead that a shock to the gas price is transmitted to the power price directly whenever a gas unit is marginal, but not at all once demand pushes
      past the gas units on the stack. Contrast the two propagation mechanisms in one sentence.
\end{enumerate}
\end{exercise}

% !TEX root = ../main.tex
\chapter{Option Prices and Probability Duality}\label{ch:duality}

\begin{sources}
\small
\begin{tabularx}{\linewidth}{@{}l l L@{}}
\toprule
\textbf{Lecture} & \textbf{Speaker} & \textbf{Content}\\
\midrule
2013 L20 & S.~Blythe (Harvard Management Company) & Guest lecture, transcript only (no slides, no notes on OCW). Three payoffs (call, zero-coupon bond, digital); the call spread as a discrete first strike-derivative and its limit, the digital; FTAP and the risk-neutral density as the second strike-derivative of the call price; the call butterfly as a traded proxy for the density; the ``other direction'' (density $\to$ price via FTAP/LOTUS); calls as a spanning set for European payoffs, first by piecewise-linear replication and then by an exact second-order Taylor identity; a remark that the terminal distribution for every maturity does not pin down the stock's stochastic process (Dupire--Derman), with the practical consequence for exotic-option pricing.\\
\bottomrule
\end{tabularx}
\smallskip
\textbf{Files.} None: OCW lists no slides or lecture notes for L20, only the video transcript.
\begin{coursenote}[2013 L22 has no material at all]
The following lecture, 2013 L22, ``Calculus of Variations and its Application in FX Execution''
(E.~Pan, Morgan Stanley), has neither a transcript nor slides on OCW and is not covered anywhere in
these notes.
\end{coursenote}
\textbf{Problem sets.} None. The exercises at the end of the chapter are all ``(not from the course)''.
\textbf{Reconstruction.} Blythe states the Breeden--Litzenberger relation and the spanning formula but
proves neither in full generality (the spanning proof is done only for piecewise-linear $g$, by
inspection). \cref{ch22:sec-bl-recap} is a two-line recap of the proof already given in
\cref{ch03:cor-bl}. \cref{ch22:sec-butterfly} adds the triangle/rectangle area argument only sketched
in class. \cref{ch22:sec-spanning} gives the exact Taylor-with-remainder computation in full (the
lecture writes it down and calls it ``exact'', without proof, and skips the discretisation
$g\to$ piecewise-linear step). All numerics (\cref{ch22:tab-butterfly,ch22:fig-density,ch22:ex-x2})
are computed for these notes, with the Black--Scholes parameters of \cref{ch15:ex-atm}
($S_0=K=80$, $T=2$, $r=5\%$, $\sigma=20\%$) so that the digital and call-spread numbers reproduce
\cref{ch15:eq-digital-limit}--\eqref{ch15:eq-digital}.
\end{sources}

\section*{Motivation}
Chapter~\ref{ch:bs} already used one instance of a strike derivative: the digital option is
$-\partial C/\partial K$ (\cref{ch15:sec-digital}), and \cref{ch03:cor-bl} proved that a call price
viewed as a function of strike encodes the whole distribution of the underlying, $C''(K)=f(K)$. This
chapter is about what that fact \emph{means}. Blythe's title for it, ``option-probability duality'',
is the claim that a call option is not a separate, exotic object requiring a pricing model: it
\emph{is} a piece of probability information, tradeable and observable, and the map between ``the set
of all call prices'' and ``the risk-neutral distribution of the underlying'' goes both ways. Read one
way, market call prices reveal the market's implied probabilities (this is genuinely done by trading
desks, via the \emph{call butterfly} of \cref{ch22:sec-butterfly}). Read the other way, the Fundamental
Theorem of Asset Pricing (\cref{thm:ftap}, generalised here to a continuum of states in
\cref{ch22:sec-ftap}) turns any candidate distribution into a full price system for every European
payoff. The two directions meet in the \emph{spanning} result of \cref{ch22:sec-spanning}: knowing the
call price at every strike is exactly as much information as knowing the whole payoff-to-price map,
because every payoff is a specific, explicit portfolio of zero-coupon bond, stock and calls. The
chapter closes in \cref{ch22:sec-process} with what this does \emph{not} give: the terminal
distribution at every maturity is not enough to pin down the process the stock follows in between,
which is why option markets alone cannot price every exotic derivative and a model (\cref{ch:sde}) is
still needed for path-dependent payoffs.

% ==================================================================
\section{Three payoffs and one assumption}\label{ch22:sec-ftap}

Fix a maturity $T$ and let $S_T$ be the price of the underlying at $T$. Following the lecture
\ts{13}{20}{13:00}, consider three contracts, all defined purely by their payout at $T$:
\begin{itemize}
  \item a \emph{call} with strike $K$: payout $(S_T-K)^+$;
  \item a \emph{zero-coupon bond}: payout $1$;
  \item a \emph{digital} with strike $K$: payout $\ind\{S_T>K\}$.
\end{itemize}
Write $C_K(t,T)$, $Z(t,T)$, $D_K(t,T)$ for their prices at $t\le T$; $Z(T,T)=1$. As in
\cref{def:onepmodel}, a payoff that is a deterministic function of $S_T$ alone is called a
\emph{European} claim, written $g(S_T)$ for a payout function $g$: $g(x)=(x-K)^+$ for the call,
$g\equiv1$ for the bond, $g(x)=\ind\{x>K\}$ for the digital \ts{13}{20}{14:46}.

\begin{finance}[Why call it $C_K(t,T)$]
Blythe is explicit that the notation separates the two roles of time: $K$ and $T$ are contract terms
fixed once and for all, while $t$ is the calendar date at which the price is observed
\ts{13}{20}{15:44}. Determining $C_K(t,T)$ for $t<T$ from a model is exactly the Black--Scholes problem
of \cref{ch:bs}; this chapter instead asks what can be said \emph{without} a model, from the family
$\{C_K(t,T)\}_{K}$ taken as given market data.
\end{finance}

The one assumption of the chapter is the Fundamental Theorem of Asset Pricing, already proved for a
finite state space in \cref{thm:ftap} (positive state prices $\Leftrightarrow$ no arbitrage, price
$=\bm\psi\T\bm C_T=\beta\,\E_\Q[\bm C_T]$). Its continuum version is:

\begin{theorem}[FTAP, continuum of states; not proved here]\label{ch22:thm-ftap}
Let $D(t,T)$ be the price at $t$ of a European claim with payout $g(S_T)$ at $T$. Under no arbitrage
(and suitable regularity, e.g.\ a complete diffusion market) there is a probability measure $\Q$, the
\emph{risk-neutral measure}, such that $D(t,T)/Z(t,T)$ is a $\Q$-martingale with respect to the
filtration generated by $S$; in particular, conditionally on $S_t$,
\begin{equation}\label{ch22:eq-ftap}
  D(t,T) = Z(t,T)\,\E_\Q\bigl[g(S_T)\mid S_t\bigr].
\end{equation}
\end{theorem}
\begin{proof}[Proof]
Not given here. \cref{thm:ftap} is the exact finite-dimensional statement and its complete proof
(Stiemke's lemma plus a separating hyperplane); the continuum, continuous-time version is due to
Harrison and Kreps (1979) and, as Blythe remarks, ``however many times you look at it, you're only
probably going to get through two or three pages before thinking, OK, that's hard''
\ts{13}{20}{51:57}. \cref{ch15:sec-onestep,ch15:sec-tree} prove it in the one-step and binomial-tree
cases, which already contain the essential idea: replicate the claim, then no-arbitrage forces its
price to equal the replicating portfolio's price, \cref{thm:ftap}(iii).
\end{proof}

\eqref{ch22:eq-ftap} is exactly \cref{thm:ftap}(iii), $C_0=\bm\psi\T\bm C_T=\beta\E_\Q[\bm C_T]$, with
the finite sum $\sum_i\psi_i(\cdot)$ replaced by the integral $\int(\cdot)\dd\Q$ and the state-price
vector $\bm\psi$ replaced by its continuum analogue: the density $\psi(x)=Z(t,T)f_\Q(x)$, where
$f_\Q$ is the $\Q$-density of $S_T$ given $S_t$, plays the role of the \emph{Arrow--Debreu price} of
state $x$ (\cref{sec:oneperiod} discusses the discrete Arrow--Debreu security $\bm e_i$ paying
\$1 in state $i$; here the ``state'' is a point $x\in\R_+$ and the security paying \$1 in a
neighbourhood $\dd x$ of $x$ costs $\psi(x)\dd x$). Applying \eqref{ch22:eq-ftap} to the digital
recovers the intuitive statement Blythe leads with: the price of \$1-if-the-coin-is-heads is the
probability of heads, discounted \ts{13}{20}{23:11}--\ts{13}{20}{24:48},
\begin{equation}\label{ch22:eq-digital-ftap}
  D_K(t,T) = Z(t,T)\,\Q(S_T>K\mid S_t).
\end{equation}

% ==================================================================
\section{From call prices to the density}\label{ch22:sec-bl-recap}

\subsection{The digital, twice}\label{ch22:sec-digital-twice}
The digital is priced two different ways: by \eqref{ch22:eq-digital-ftap} above, and as the limit of a
\emph{call spread}, exactly the computation of \cref{ch15:sec-digital}. Long $\lambda$ calls struck at
$K$ and short $\lambda$ calls struck at $K+1/\lambda$ has payout $0$ below $K$, slope $\lambda$ on
$[K,K+1/\lambda]$, and $1$ above $K+1/\lambda$ \ts{13}{20}{17:07}--\ts{13}{20}{19:19}; letting
$\lambda\to\infty$ this payout tends pointwise to $\ind\{S_T>K\}$, and its price tends to
\begin{equation}\label{ch22:eq-cs-limit}
  \lim_{\lambda\to\infty}\lambda\bigl[C_K(t,T)-C_{K+1/\lambda}(t,T)\bigr] = -\frac{\partial C_K}{\partial K}(t,T).
\end{equation}
This is \eqref{ch15:eq-digital-limit} with $\varepsilon=1/(2\lambda)$ and a one-sided spread instead of
a centred one; the limit is the same. Equating \eqref{ch22:eq-cs-limit} with
\eqref{ch22:eq-digital-ftap} gives
\begin{equation}\label{ch22:eq-cdf}
  \Q(S_T>K\mid S_t) = -\frac{1}{Z(t,T)}\frac{\partial C_K}{\partial K}(t,T),
  \qquad
  F_\Q(K):=\Q(S_T\le K\mid S_t) = 1+\frac{1}{Z(t,T)}\frac{\partial C_K}{\partial K}(t,T),
\end{equation}
and one more strike derivative gives the density, restating \cref{ch03:cor-bl} with the discount factor
made explicit:
\begin{equation}\label{ch22:eq-bl}
  f_\Q(K) = \frac{1}{Z(t,T)}\frac{\partial^2 C_K}{\partial K^2}(t,T).
\end{equation}
\eqref{ch22:eq-bl} is the \emph{Breeden--Litzenberger relation}, already proved in
\cref{ch03:cor-bl} (there for $Z\equiv1$, i.e.\ undiscounted expectations) and used in
\cref{ch15:sec-digital} for the digital alone; nothing here is new mathematics, only the
observation, which is the point of the lecture, that \emph{iterating the same one-strike-derivative
trick that gives the digital, one more time, gives the entire density, not just one tail probability}.

\begin{intuition}[One relation, no model]
Nothing about $\Q$ or $S$ entered the derivation beyond FTAP itself: \eqref{ch22:eq-bl} holds for
\emph{whatever process} generates the risk-neutral distribution, log-normal (Black--Scholes) or not
\ts{13}{20}{39:03}--\ts{13}{20}{39:31}. Plugging the Black--Scholes formula \eqref{ch15:eq-call} into
$\partial^2_KC$ recovers the log-normal density of \cref{ch15:eq-gbm}, but \eqref{ch22:eq-bl} is
agnostic to that choice: it converts \emph{observed} call prices into a distribution without ever
positing a model for $S$.
\end{intuition}

\subsection{The call butterfly: trading the density}\label{ch22:sec-butterfly}
A single strike derivative (the call spread) has an obvious trading analogue; Blythe's second
construction shows that the \emph{second} strike derivative does too \ts{13}{20}{30:56}. Consider the
portfolio, long $\lambda$ calls at $K-1/\lambda$, short $2\lambda$ calls at $K$, long $\lambda$ calls at
$K+1/\lambda$ (a \emph{call butterfly}, traded in practice, unlike the idealised limit of
\cref{ch22:sec-digital-twice} \ts{13}{20}{32:22}--\ts{13}{20}{32:37}). Its price is
\begin{equation}\label{ch22:eq-butterfly-price}
  B_K(\lambda;t,T) := \lambda\Bigl[C_{K-1/\lambda}(t,T) - 2C_K(t,T) + C_{K+1/\lambda}(t,T)\Bigr],
\end{equation}
a centred second difference of the call price, scaled by $\lambda$ (one more power of $\lambda$ than
the call spread, matching the second derivative being approximated by
$C(K-\varepsilon)-2C(K)+C(K+\varepsilon)\approx\varepsilon^2C''(K)$ with $\varepsilon=1/\lambda$):
\begin{equation}\label{ch22:eq-butterfly-limit}
  \lim_{\lambda\to\infty}\lambda\,B_K(\lambda;t,T) = \frac{\partial^2 C_K}{\partial K^2}(t,T) = Z(t,T)f_\Q(K).
\end{equation}

\begin{proposition}[The butterfly payoff is a triangle]\label{ch22:prop-triangle}
The payout of the portfolio in \eqref{ch22:eq-butterfly-price} at maturity is the triangular tent
function of base $[K-1/\lambda,K+1/\lambda]$ and height $1$ at $K$,
\[
  h_\lambda(x) = \lambda\Bigl[(x-K+\tfrac1\lambda)^+ - 2(x-K)^+ + (x-K-\tfrac1\lambda)^+\Bigr]
  = \bigl(1-\lambda|x-K|\bigr)^+ .
\]
\end{proposition}
\begin{proof}
Direct case check on the three pieces of the piecewise-linear function; each summand is a hockey stick
with a kink at $K\mp1/\lambda$ or $K$, and the three slopes ($\lambda$, $-2\lambda$, $\lambda$) sum to
$0$ outside $[K-1/\lambda,K+1/\lambda]$ (so the payout is $0$ there) and change by $\mp\lambda$ at each
kink, giving slope $+\lambda$ on $[K-1/\lambda,K]$ and $-\lambda$ on $[K,K+1/\lambda]$, height
$1=\lambda\cdot\frac1\lambda$ at $K$. This is exactly the sketch in class
\ts{13}{20}{36:17}--\ts{13}{20}{36:41}, made explicit.
\end{proof}

This makes \eqref{ch22:eq-butterfly-limit} the exact analogue of the elementary fact
$\Q(K-\delta<S_T<K+\delta)\approx2\delta f_\Q(K)$ for small $\delta$: the triangle of
\cref{ch22:prop-triangle} has base $2/\lambda$ and height $1$, hence area $1/\lambda$, and by FTAP its
discounted expected payout is
\begin{equation}\label{ch22:eq-triangle-area}
  B_K(\lambda;t,T) = Z(t,T)\,\E_\Q[h_\lambda(S_T)\mid S_t] \approx Z(t,T)\cdot\frac1\lambda\, f_\Q(K)
  \quad\text{for large }\lambda,
\end{equation}
which is \eqref{ch22:eq-butterfly-limit} divided by $\lambda$ on both sides
\ts{13}{20}{35:39}--\ts{13}{20}{38:16}. Unlike the call spread, the butterfly is genuinely traded
(interest-rate desks quote butterflies $10$--$20$ basis points wide, as noted in class
\ts{13}{20}{33:24}--\ts{13}{20}{33:37}), so \eqref{ch22:eq-triangle-area} is not a theoretical limit
alone: a market maker who believes $\Q(S_T\approx K)$ is larger than the price of the corresponding
butterfly implies will buy it \ts{13}{20}{38:23}--\ts{13}{20}{38:48}.

\begin{table}[ht]
\centering
\caption{Butterfly price $\lambda\,B_K(\lambda;0,T)$ converging to $C''(K)=e^{-rT}f_\Q(K)$, at $K=80$,
Black--Scholes with $S_0=80$, $r=5\%$, $\sigma=20\%$, $T=2$ (parameters of \cref{ch15:ex-atm}).
Computed for these notes.}\label{ch22:tab-butterfly}
\begin{tabular}{@{}rrr@{}}
\toprule
$\lambda$ & $\varepsilon=1/\lambda$ & $\lambda\,B_{80}(\lambda)$\\
\midrule
$0.05$ & $20$   & $0.014656$\\
$0.1$  & $10$   & $0.015354$\\
$0.2$  & $5$    & $0.015536$\\
$0.5$  & $2$    & $0.015588$\\
$1$    & $1$    & $0.015596$\\
$2$    & $0.5$  & $0.015598$\\
$5$    & $0.2$  & $0.015598$\\
\midrule
\multicolumn{2}{r}{exact $e^{-rT}f_\Q(80)$} & $0.015598$\\
\bottomrule
\end{tabular}
\end{table}

\begin{figure}[ht]
\centering
\begin{tikzpicture}
\begin{axis}[width=0.92\linewidth,height=6cm,axis lines=left,xmin=35,xmax=145,ymin=0,ymax=0.017,
  xlabel={$K$},ylabel={discounted density},tick label style={font=\footnotesize},
  label style={font=\small},
  legend style={at={(0.97,0.97)},anchor=north east,font=\scriptsize,draw=none,fill=none}]
  \addplot[thick,black!45,mark=none] coordinates {
   (40.0,0.00092) (44.0,0.00194) (48.0,0.00347) (52.0,0.00545) (56.0,0.00770) (60.0,0.00999)
   (64.0,0.01208) (68.0,0.01377) (72.0,0.01494) (76.0,0.01554) (80.0,0.01560) (84.0,0.01518)
   (88.0,0.01439) (92.0,0.01333) (96.0,0.01211) (100.0,0.01081) (104.0,0.00950) (108.0,0.00824)
   (112.0,0.00707) (116.0,0.00600) (120.0,0.00504) (124.0,0.00421) (128.0,0.00349) (132.0,0.00287)
   (136.0,0.00235) (140.0,0.00192)};
  \addplot[only marks,mark=o,mark size=1.4pt,defcol] coordinates {
   (40.0,0.00092) (44.0,0.00194) (48.0,0.00347) (52.0,0.00545) (56.0,0.00770) (60.0,0.00999)
   (64.0,0.01208) (68.0,0.01377) (72.0,0.01494) (76.0,0.01554) (80.0,0.01560) (84.0,0.01518)
   (88.0,0.01439) (92.0,0.01333) (96.0,0.01211) (100.0,0.01081) (104.0,0.00950) (108.0,0.00824)
   (112.0,0.00707) (116.0,0.00600) (120.0,0.00504) (124.0,0.00421) (128.0,0.00349) (132.0,0.00287)
   (136.0,0.00235) (140.0,0.00192)};
  \legend{$e^{-rT}f_\Q$ (exact),{$\lambda B_K(\lambda)$, $\lambda=8$}}
\end{axis}
\end{tikzpicture}
\caption{Discounted risk-neutral density recovered from the call-butterfly price
\eqref{ch22:eq-butterfly-price} at $\lambda=8$ (width $2/\lambda=0.25$), against the exact
Black--Scholes log-normal density (same parameters as \cref{ch22:tab-butterfly}); markers overlap the
line to plotting accuracy. Computed for these notes.}\label{ch22:fig-density}
\end{figure}

% ==================================================================
\section{From the density to any price: FTAP the other way}\label{ch22:sec-density-to-price}

\cref{ch22:sec-bl-recap} went from call prices to a distribution. The Fundamental Theorem
\eqref{ch22:eq-ftap} already goes the other way, and the resulting statement is worth writing out
explicitly \ts{13}{20}{49:41}--\ts{13}{20}{54:09}: if $f_\Q(\cdot\mid S_t)$ is known, the price of
\emph{any} European claim with payout $g$ is
\begin{equation}\label{ch22:eq-lotus}
  D(t,T) = Z(t,T)\int_0^\infty g(x)\,f_\Q(x\mid S_t)\,\dd x,
\end{equation}
an application of the law of the unconscious statistician (LOTUS): $\E_\Q[g(S_T)]$ is computed by
integrating $g$ against the density of $S_T$, without first working out the (generally intractable)
density of the transformed variable $g(S_T)$ itself. Taking $g(x)=(x-K)^+$ in \eqref{ch22:eq-lotus}
recovers \cref{ch03:thm-call} (there stated for a general random variable $X$ under the physical
measure; here $X=S_T$ under $\Q$, with the discount factor $Z(t,T)$ attached). So the two directions
of the duality are: call prices $\to$ density via two strike derivatives
(\cref{ch22:sec-bl-recap}), density $\to$ any European price via one integral
(\eqref{ch22:eq-lotus}), and both use nothing beyond FTAP.

% ==================================================================
\section{Calls span every European payoff}\label{ch22:sec-spanning}

The two directions combine into a genuinely new statement, not reducible to
Breeden--Litzenberger alone: knowing $C_K(t,T)$ for every $K$ determines the price of \emph{every}
European claim, and moreover gives an \emph{explicit static replicating portfolio} built only from
the zero-coupon bond, the stock, and calls \ts{13}{20}{55:44}--\ts{13}{20}{56:09}.

\subsection{Piecewise-linear payouts}\label{ch22:sec-pwl}
If $g$ is piecewise linear with kinks at $K_1<K_2<\dots<K_n$, its graph is matched exactly, for
$x\ge K_1$ say, by a portfolio holding $\theta_i$ calls of strike $K_i$, where $\theta_i$ is the
change in slope of $g$ at $K_i$ (a call of strike $K_i$ contributes slope $1$ for $x>K_i$ and $0$
below, so a linear combination reproduces any sequence of slope changes)
\ts{13}{20}{56:47}--\ts{13}{20}{58:13}. Since $g$ and this call portfolio have \emph{identical} payoff
at $T$ for every value of $S_T$, no-arbitrage (\cref{thm:ftap}(iii): a replicable claim has a unique
price, the price of any replicating portfolio) forces their prices to agree at every $t\le T$,
\begin{equation}\label{ch22:eq-pwl-repl}
  D(t,T) = \sum_i \theta_i\, C_{K_i}(t,T), \qquad \theta_i=\Delta g'(K_i).
\end{equation}
A continuous $g$ is a uniform limit of piecewise-linear functions, so \eqref{ch22:eq-pwl-repl} suggests
(but does not by itself prove, because of the limit) that \emph{every} continuous payout is spanned by
calls. \cref{ch22:sec-taylor} makes this precise and quantitative for $g\in C^2$, which is the
practically relevant case and the one the lecture treats.

\begin{coursenote}[A digression worth keeping]
A student asks whether replication at $T$ really forces equality of prices at $t<T$ for maturities
decades away, given that no one can actually hold the position without funding or margin risk in the
meantime \ts{13}{20}{59:44}--\ts{13}{20}{1:02:13}. Blythe's answer: mathematically yes, under
no-arbitrage; empirically, for very long maturities the two prices \emph{have} been observed to drift
apart, because ``no arbitrage'' assumes unlimited risk capital willing to hold the position for the
full horizon, which breaks down since 2008. This is a caveat on the economic content of
\cref{thm:ftap}(iii), not on its mathematics, and it is not pursued further here.
\end{coursenote}

\subsection{The exact Carr--Madan identity}\label{ch22:sec-taylor}
\begin{proposition}[Static replication of a smooth payoff]\label{ch22:prop-carrmadan}
Let $g:\R_+\to\R$ be twice continuously differentiable. For every $x\ge0$,
\begin{equation}\label{ch22:eq-taylor-exact}
  g(x) = g(0) + g'(0)\,x + \int_0^\infty g''(K)\,(x-K)^+\,\dd K.
\end{equation}
\end{proposition}
\begin{proof}
This is the lecture's ``exact Taylor series to second order'' \ts{13}{20}{1:03:33}--\ts{13}{20}{1:04:20},
here proved rather than only asserted. Start from the fundamental theorem of calculus twice,
\[
  g(x) = g(0) + \int_0^x g'(K)\,\dd K, \qquad g'(K) = g'(0) + \int_0^K g''(c)\,\dd c,
\]
substitute and swap the order of integration in the double integral
$\int_0^x\!\int_0^K g''(c)\,\dd c\,\dd K$ over the triangle $0\le c\le K\le x$:
\[
  \begin{gathered}
  \int_0^x\int_0^K g''(c)\,\dd c\,\dd K = \int_0^x g''(c)\int_c^x \dd K\,\dd c = \int_0^x g''(c)(x-c)\,\dd c \\
  = \int_0^\infty g''(c)\,(x-c)^+\,\dd c,
  \end{gathered}
\]
where the last equality extends the integral to $(x,\infty)$ for free because $(x-c)^+=0$ there. Renaming
the dummy variable $c\to K$ gives \eqref{ch22:eq-taylor-exact}. (Equivalently, integrate the right-hand
side of \eqref{ch22:eq-taylor-exact} by parts twice, as the lecture suggests
\ts{13}{20}{1:04:20}--\ts{13}{20}{1:04:31}, which reverses this computation.)
\end{proof}

\begin{remark}
The identity holds pointwise, for \emph{every} $x\ge0$, with no approximation and no limit: this is the
sense in which it is ``exact''. It is the content of \cref{ch22:sec-pwl} made rigorous:
$g''(K)\dd K$ replaces the discrete slope changes $\theta_i=\Delta g'(K_i)$ of a piecewise-linear
function by a continuum of infinitesimal slope changes, each priced by $(x-K)^+$, the call payoff.
Where the lecture stops after the piecewise-linear case and asserts the smooth case by continuity
\ts{13}{20}{1:02:52}--\ts{13}{20}{1:03:03}, \cref{ch22:prop-carrmadan} needs no limiting argument at
all: it is an identity between functions, true for every fixed $x$, and the piecewise-linear intuition
of \cref{ch22:sec-pwl} is recovered by taking $g''$ to be a sum of Dirac masses (formally, outside the
$C^2$ hypothesis, but the calls-with-strikes-at-the-kinks picture is exactly this limit).
\end{remark}

Set $x=S_T$ in \eqref{ch22:eq-taylor-exact} (a pointwise identity between random variables, valid on
every sample path) and apply the discounted-expectation operator $Z(t,T)\,\E_\Q[\cdot\mid S_t]$ to both
sides. The left side becomes $D(t,T)$, the price of the claim with payout $g$. On the right,
linearity of expectation gives three terms: $g(0)$ times the price of $g(0)$ zero-coupon bonds is
$Z(t,T)g(0)$; the term $g'(0)\,\E_\Q[S_T\mid S_t]$, discounted, is $g'(0)$ times the price of \emph{one
unit of stock} (the stock, viewed as its own European claim with payout $S_T$ and replicating portfolio
``hold the stock'', so $Z(t,T)\E_\Q[S_T\mid S_t]=S_t$ by \cref{ch22:thm-ftap} applied to $g(x)=x$
\ts{13}{20}{1:07:25}--\ts{13}{20}{1:07:45}); and the integral term becomes, exchanging expectation and
the $K$-integral (justified when $g''\ge0$ by Tonelli, or under an integrability bound on $g''$ in
general),
\[
  \begin{gathered}
  Z(t,T)\,\E_\Q\Bigl[\int_0^\infty g''(K)(S_T-K)^+\dd K\ \Big|\ S_t\Bigr]
  = \int_0^\infty g''(K)\,Z(t,T)\,\E_\Q\bigl[(S_T-K)^+\mid S_t\bigr]\dd K \\
  = \int_0^\infty g''(K)\,C_K(t,T)\,\dd K.
  \end{gathered}
\]
Collecting the three terms:

\begin{theorem}[Calls span all European payoffs; Carr--Madan]\label{ch22:thm-spanning}
For $g\in C^2(\R_+)$ with the payoff at maturity $g(S_T)$,
\begin{equation}\label{ch22:eq-spanning}
  D(t,T) = g(0)\,Z(t,T) + g'(0)\,S_t + \int_0^\infty g''(K)\,C_K(t,T)\,\dd K.
\end{equation}
Equivalently, the claim is statically replicated at $t$ by $g(0)$ zero-coupon bonds, $g'(0)$ units of
stock, and $g''(K)\dd K$ calls of strike $K$ for every $K>0$.
\end{theorem}

\begin{keyformulas}
\textbf{Call prices $\to$ distribution:} $\displaystyle \Q(S_T>K\mid S_t)=-\frac1{Z}\partial_KC_K$,
\quad $\displaystyle f_\Q(K)=\frac1{Z}\partial^2_KC_K$ \eqref{ch22:eq-cdf}--\eqref{ch22:eq-bl}
(Breeden--Litzenberger, \cref{ch03:cor-bl}). \\
\textbf{Distribution $\to$ any price:} $\displaystyle D(t,T)=Z(t,T)\int g(x)f_\Q(x\mid S_t)\dd x$
\eqref{ch22:eq-lotus} (FTAP/LOTUS). \\
\textbf{Static replication (Carr--Madan):}
$\displaystyle D(t,T)=g(0)Z(t,T)+g'(0)S_t+\int_0^\infty g''(K)C_K(t,T)\dd K$ \eqref{ch22:eq-spanning}.\\
\textbf{Traded proxy for the density:} call butterfly $\lambda B_K(\lambda)\to Z(t,T)f_\Q(K)$
\eqref{ch22:eq-butterfly-limit}.
\end{keyformulas}

\begin{example}[Checking the identity on $g(x)=x^2$]\label{ch22:ex-x2}
Take $g(x)=x^2$, so $g(0)=0$, $g'(0)=0$, $g''(K)=2$ for all $K$. \eqref{ch22:eq-taylor-exact} claims
$x^2=\int_0^\infty 2(x-K)^+\dd K$, and indeed $\int_0^x 2(x-K)\dd K=2\bigl[xK-K^2/2\bigr]_0^x=x^2$: a
direct check of \cref{ch22:prop-carrmadan} with an elementary integral. Numerically (trapezoidal rule,
grid $[0,800]$, $4\times10^5$ points, computed for these notes) the same integral gives $6400.0000$
against the exact $80^2=6400$. A second check with a bounded, non-polynomial $g(x)=e^{-0.01x}$
($g''(K)=10^{-4}e^{-0.01K}$) gives, at $x=80$, exact $g(80)=0.449329$ against the numerical
reconstruction $0.449329$: \eqref{ch22:eq-taylor-exact} is confirmed to $6$ digits both for a payoff
whose second derivative is constant and for one that decays.
\end{example}

% ==================================================================
\section{What the terminal distribution does not determine}\label{ch22:sec-process}

Everything above concerns a single fixed maturity $T$: knowing all call prices $C_K(t,T)$ for that one
$T$ determines the distribution of $S_T$ given $S_t$, and hence the price of every payoff that depends
on $S_T$ alone. The natural next question is whether knowing this for \emph{every} maturity $T$ (a full
implied-volatility surface, in market language) determines the entire stochastic process $(S_u)_{u\ge
t}$ \ts{13}{20}{1:12:11}--\ts{13}{20}{1:13:07}.

\begin{proposition}[Marginals do not determine the process; not proved here]\label{ch22:prop-dupire}
Knowledge of the marginal (risk-neutral) distribution of $S_T$ given $S_t$, for every $T\ge t$, does
not determine the law of the process $(S_u)_{u\in[t,T]}$. If in addition $S$ is assumed to be a
one-dimensional diffusion driven by Brownian motion, $\dd S_u=\mu(u,S_u)\dd u+\sigma(u,S_u)\dd W_u$,
then the marginals \emph{do} determine the process, via the local volatility function
$\sigma_{\mathrm{loc}}(u,x)$ (Dupire, Derman--Kani, early 1990s).
\end{proposition}
\begin{proof}[Discussion, not a proof]
Blythe gives only the intuition, not the construction, and none is attempted here (it belongs to
stochastic calculus proper, beyond \cref{ch:sde}): take a denser and denser grid of maturities with
prescribed, say symmetric, marginals that spread out smoothly in $T$; the stock can still be made to
jump discontinuously between two adjacent grid times while landing in a distribution consistent with
the (denser and denser) marginals at both ends, because \emph{nothing in the one-maturity-at-a-time
data constrains the joint law across maturities} \ts{13}{20}{1:13:34}--\ts{13}{20}{1:14:07}. Ruling
out such jumps, by assuming a diffusion, is what lets Dupire's formula reconstruct a unique
local-volatility function from the full surface of \eqref{ch22:eq-bl}-derived marginals.
\end{proof}

\begin{finance}[Why this matters for a trading desk]
A desk that has priced, and can perfectly hedge with the observed vanilla surface, every European
payoff via \eqref{ch22:eq-spanning}, still cannot claim the same for a path-dependent exotic (a
barrier, an Asian option, or, as Blythe's example, a derivative paying on the number of times $S$
crosses a level in a small interval \ts{13}{20}{1:19:05}--\ts{13}{20}{1:19:10}): its price is not
pinned down by the vanilla surface alone, precisely because \cref{ch22:prop-dupire} says the surface
does not pin down the process. A student's question makes the link to market completeness explicit
\ts{13}{20}{1:16:44}--\ts{13}{20}{1:18:52}: a jump/flip process for which the marginals do not
determine the law corresponds to an \emph{incomplete} market (\cref{def:onepmodel}) with respect to
the vanilla calls, exactly as more states than independent assets leave state prices non-unique in the
one-period model of \cref{sec:oneperiod}. The risk-neutral density $f_\Q(\cdot\mid S_t)$ at each
fixed $T$ is still uniquely pinned down by the calls (Blythe is emphatic on this point when pressed
\ts{13}{20}{1:19:17}--\ts{13}{20}{1:19:48}); what is not determined is the joint law across different
$T$'s, i.e.\ the path.
\end{finance}

\begin{remark}[An alternative spanning set: Arrow--Debreu securities]
A question from the audience asks whether calls are the \emph{only} convenient spanning set
\ts{13}{20}{1:16:16}--\ts{13}{20}{1:17:41}. They are not: knowing all digitals gives the CDF directly
(one derivative instead of two), and knowing the price of every Arrow--Debreu security $\ind\{S_T\in
\dd x\}$ gives the density directly, with no derivative at all. Blythe's preference for calls is that
the identity \eqref{ch22:eq-spanning} is unusually clean, expressing an arbitrary smooth payoff through
an explicit, liquidly-traded static hedge, whereas Arrow--Debreu securities of infinitesimal width are
not traded and digitals are less liquid than calls (\cref{ch15:sec-digital}). This is also exactly the
representation \cref{ch:ross} will use (state prices as a discrete matrix, recovered via
Perron--Frobenius, \cref{sec:pf}) in a setting where the underlying state space is finite rather than
a continuum of strikes.
\end{remark}

% ==================================================================
\section*{Exercises}
\addcontentsline{toc}{section}{Exercises}
All exercises below are (not from the course).

\begin{exercise}[Butterfly width and bias]\label{ch22:ex-width}
For the Black--Scholes parameters of \cref{ch22:tab-butterfly} ($S_0=80$, $r=5\%$, $\sigma=20\%$,
$T=2$), let $\lambda=4$ (butterfly of half-width $1/\lambda=0.25$ around each wing, full base
$0.5$). Compute $B_{80}(4)$ from the closed-form Black--Scholes call price \eqref{ch15:eq-call} and
compare $\lambda B_{80}(4)$ to the exact value in \cref{ch22:tab-butterfly}. Express the leading-order
error as a function of $\lambda$ using a fourth-order Taylor expansion of $C$ around $K=80$, and check
it is consistent with the pattern in the table.
\end{exercise}

\begin{exercise}[Straddle via Carr--Madan]\label{ch22:ex-straddle}
A straddle pays $g(S_T)=|S_T-K_0|$ for a fixed $K_0>0$. This $g$ is not $C^2$ at $K_0$ (it has a
kink), so \eqref{ch22:eq-taylor-exact} must be applied with $g''$ interpreted as a Dirac mass
$2\delta_{K_0}$ (the jump in $g'$ at $K_0$ is $2$). Using \eqref{ch22:eq-pwl-repl} directly instead
(the straddle is piecewise linear with one kink), find the static replicating portfolio of zero-coupon
bonds, stock and calls of a single strike, and check it agrees with formally substituting $g''=2\delta_{K_0}$
into \eqref{ch22:eq-spanning}.
\end{exercise}

\begin{exercise}[Digital from a butterfly, not a call spread]\label{ch22:ex-digital-butterfly}
Show that a digital option cannot be replicated as a \emph{limit of call butterflies} the way it was
replicated as a limit of call spreads in \cref{ch22:sec-digital-twice}: compute
$\lim_{\lambda\to\infty}B_K(\lambda;t,T)$ (not $\lambda B_K(\lambda;t,T)$) and interpret the result in
terms of the payoff triangle of \cref{ch22:prop-triangle} shrinking to a point.
\end{exercise}

\begin{exercise}[Two maturities, same terminal law, different paths]\label{ch22:ex-two-paths}
Let $W$ be a Brownian motion and $Z\sim N(0,1)$. Show that $S_t=1+W_t$ and $\tilde S_t=1+\sqrt t\,Z$
have the same marginal law $N(1,t)$ at every $t\in[0,2]$, but different laws as processes: compute
$\Cov(S_1,S_2)$ and $\Cov(\tilde S_1,\tilde S_2)$. Which of the two is a martingale? This illustrates
\cref{ch22:prop-dupire} concretely: a claim paying $(S_2-S_1)^2$ has different prices under the two
processes although all the terminal laws, hence all the call prices, agree.
\end{exercise}

% !TEX root = ../main.tex
\chapter{Quanto Credit Hedging}\label{ch:quanto}

\begin{sources}
\small
\begin{tabularx}{\linewidth}{@{}l l L@{}}
\toprule
\textbf{Lecture} & \textbf{Speaker} & \textbf{Content}\\
\midrule
2013 L23 & S.~Andreev (Morgan Stanley) & FX spot and forwards, interest-rate parity; the FX betting-game
paradox; sovereign bonds issued in two currencies (Italy) and the credit-spread puzzle; a minimal jump-on-default
FX model; pricing and dynamic hedging of zero-recovery USD/EUR bonds in that model, extension to $R>0$; the
jump-diffusion model left as a final-paper topic; quanto credit default swaps and other applications.\\
\bottomrule
\end{tabularx}

\smallskip
\textbf{Files.} 2013 slides \file{MIT18_S096F13_lecnote23.pdf} (32 slides); accompanying written notes with the
derivations (referenced repeatedly on the slides as ``see class notes'') were not among the materials available
for these notes, so every derivation below is reconstructed from the slides and from the blackboard work described
in the lecture transcript, and re-checked numerically in a Python script written for these notes; reconstructions
are flagged explicitly, and all such numerical checks are marked ``computed for these notes.''
\textbf{Problem sets.} None; the exercises at the end of the chapter are all marked ``not from the course''.
\textbf{Not covered here.} The stochastic-hazard-rate, stochastic-interest-rate, correlated multi-factor version of
the model is only sketched on the last slides (``Real Life''); \cref{ch:hjm} returns to term-structure modelling of
credit.
\end{sources}

\section*{Motivation}
Chapter~\ref{ch:counterparty} showed that a change of numeraire turns a ratio of two traded assets into a
martingale, and listed the FX forward $F_{FX}(t,T)=S_{FX}(t)P^F(t,T)/P^D(t,T)$ as one instance of the pattern
(\cref{ch08:prop-numeraire}). This chapter asks what changes when one of the two ``assets'' can \emph{jump}: a
sovereign issuer defaults, and its currency devalues at the same instant. Two zero-coupon bonds of the same issuer,
one paying 1 USD and one paying 1 EUR, are no longer related by a simple deterministic forward rate, because the
FX rate that converts one payoff into the other is itself correlated with the credit event. Pricing and hedging
one bond with the other -- a \emph{quanto} problem, in the sense that a payoff is being valued in a currency other
than its natural one -- requires modelling the FX jump explicitly.

The chapter builds the smallest possible model that has this feature (a constant-hazard-rate default time with a
deterministic percentage FX jump on default), and carries it through the three-step recipe of
\cref{ch15:sec-onestep}:
\begin{enumerate}[1.]
  \item write the payoffs,
  \item price by risk-neutral expectation and find the replicating
      portfolio,
  \item extract the economic intuition (the sign and size of the ``quanto adjustment'').
\end{enumerate}
The compensator identity used along the way is the jump-process analogue of the risk-neutral drift condition of
\cref{ch15:sec-rn}, and the pricing measure used throughout is exactly the risk-neutral measure $\Q$ of
\cref{ch:bs}; \cref{ch23:sec-numeraire} makes the connection to \cref{ch08:prop-numeraire} precise.

\begin{intuition}[Why symmetric bets are not symmetric]
If I bet you \$100 that a coin lands heads and you bet me 100 EUR that it lands tails, the two bets look
symmetric. They are not, if the coin is Italy defaulting and the EUR devalues whenever it does: winning USD then
happens exactly when EUR is weak, and winning EUR happens exactly when EUR is strong. The currency of the payoff
carries information about the payoff, and a fair price must account for it. This is the whole content of the
chapter, before any formula is written down.
\end{intuition}

\section{FX, interest-rate parity, and the currency of a payoff}\label{ch23:sec-fx}
Write $S_t$ for the spot FX rate USD/EUR, i.e.\ the price in USD of 1 EUR. In the (no-default) Garman--Kohlhagen
model, with domestic (USD) and foreign (EUR) risk-free rates $r_d,r_f$,
\begin{equation}\label{ch23:eq-gk}
  \frac{\dd S_t}{S_t}=(r_d-r_f)\,\dd t+\sigma\,\dd W_t ,\qquad
  S_T=S_0\exp\!\Bigl[(r_d-r_f-\tfrac12\sigma^2)T+\sigma W_T\Bigr],\qquad
  \E[S_T]=S_0e^{(r_d-r_f)T}.
\end{equation}
This is the same lognormal SDE as \cref{ch15:eq-gbm} with $\mu=r_d-r_f$: under $\Q$, a foreign asset earns the
domestic rate once its value is expressed in domestic currency, exactly as a stock earns $r$ under the
risk-neutral measure of \cref{ch15:sec-rn}. The forward FX rate is fixed by no-arbitrage (covered interest-rate
parity), $F(t,T)=S_te^{(r_d-r_f)(T-t)}=\E_t[S_T]$ -- itself an instance of \cref{ch08:prop-numeraire} with
$N_1=1$ (cash) and $N_2=P^D(\cdot,T)$, as recorded in the numeraire table of \cref{ch:counterparty}.

For the rest of the chapter, following the slides, interest rates are set to zero, $r_d=r_f=0$, so that the only
driver of $S_t$ is the credit event: this isolates the ``quanto credit'' effect from the ordinary FX quanto
adjustment of \cref{ch23:eq-gk}, which is unaffected by anything in this chapter and is recovered by simply adding
back $r_d-r_f$ to every drift below.

\begin{finance}[The FX betting game]
Slide 7 poses two \$100-notional bets on whether USD/EUR (currently 1) is above or below 1 in one month. Bet A
pays $\mp\$100$ in USD; Bet B pays $\mp100$ EUR. If the exchange rate is uncorrelated with the outcome the two
bets are worth the same. But if, say, USD/EUR $>1$ (euro strong, dollar weak) exactly in the scenarios where the
bet is lost, then Bet A ($-\$100$ regardless of the rate) is strictly better than Bet B ($-100$ EUR, converted at
an unfavourable rate). Table~\ref{ch23:tab-game} reproduces the slide's numbers: Bet A minus Bet B is always
positive because the loss in Bet B is amplified, and its gain reduced, by exactly the FX move correlated with it. The general
principle: \emph{whenever a payoff and the FX rate used to convert it are correlated, pricing the payoff in the
``wrong'' currency and simply converting at the current spot or forward rate is not valid} -- one must price the
converted payoff directly under a risk-neutral expectation, as done for CDS below.
\end{finance}

\begin{table}[htbp]
\centering
\caption{FX betting game, slide 8: two symmetric-looking \$100/EUR100 bets are worth different amounts once
the payoff currency is correlated with the FX outcome.}\label{ch23:tab-game}
\begin{tabular}{@{}lrrr@{}}
\toprule
Scenario (1 month) & Bet A (USD) & Bet B (EUR, in USD) & A $-$ B\\
\midrule
USD/EUR $=1.25$ & $-100$ & $-125$ & $+25$\\
USD/EUR $=0.75$ & $+100$ & $+75$ & $+25$\\
\bottomrule
\end{tabular}
\end{table}

\section{A reality check: sovereign bonds in two currencies}\label{ch23:sec-italy}
Slides 11--14 use Italy as a running example: a sovereign that issues bonds in EUR and in USD (and other
currencies), with a \emph{cross-default} clause (default on one currency triggers default on all). Two natural
questions: which currency does Italy prefer to borrow in, and which currency do bond investors prefer? The credit
spread (the extra yield over the local risk-free curve) need not be the same in the two currencies, and the slides
show that Italy's EUR and USD five-year credit spreads did in fact move apart during the 2011--2012 sovereign
crisis. Comparing the two spreads directly conflates two different things: the market's view of Italy's
creditworthiness (currency-independent, in a first approximation) and the market's view of what happens to EUR
\emph{if Italy defaults} (which only affects the EUR-denominated bond, since a EUR bond's payoff is naturally
denominated in the currency that would devalue). Section~\ref{ch23:sec-model} builds the smallest model that
captures this second effect, so that the two spreads can be compared on a like-for-like basis.

\subsection{A naive replication attempt, and why it fails}\label{ch23:sec-naive}
Consider two zero-coupon, zero-recovery bonds of the same issuer and maturity $T$: bond $U$ pays 1 USD, bond $E$
pays 1 EUR. With $P_u,P_e$ their prices (in USD and EUR respectively), spot rate $S_t$ and forward rate $F_t$
(USD/EUR, maturity $T$), a naive covered-interest-rate-parity argument suggests buying $100{,}000$ E-bonds, selling
$100{,}000F_t$ U-bonds, and entering a forward to sell $100{,}000$ EUR at $F_t$ (slide 15). If both
bonds survive to $T$, the position pays exactly $0$: the $100{,}000$ EUR received is exchanged at the pre-agreed forward
rate for $100{,}000F_t$ USD, which exactly cancels the short U-bonds. But if the issuer defaults with recovery rate $25\%$
(slide 17), the E-bonds pay $25{,}000$ EUR and the short U-bonds owe $25{,}000F_t$ USD, while the forward contract,
which does not know that the issuer defaulted, still delivers the full $100{,}000$ EUR notional against $100{,}000F_t$ USD.
The strategy is left with a large \emph{unhedged} FX forward position (short $75{,}000$ EUR against $75{,}000F_t$ USD) exactly in the scenario -- default -- in
which EUR is most likely to move sharply (as it did for the Argentine peso in December 2001, slide 18). The naive
strategy is not an arbitrage and is not a valid hedge: the forward is the wrong instrument because it does not
distinguish the default and no-default states. What is needed is a model in which the number of hedging
instruments equals the number of sources of randomness (default timing and the FX jump it triggers) -- a complete
market in the sense of \cref{ch15:sec-onestep} -- so that pricing by replication is again possible.
\section{The minimal jump-on-default FX model}\label{ch23:sec-model}

\subsection{Hazard rate and survival probability}\label{ch23:sec-hazard}
Model the default time $\tau$ of the issuer as the arrival time of the first jump of a Poisson process of constant
intensity $h>0$ (the \emph{hazard rate}): from the perspective of time $t<\tau$,
\begin{equation}\label{ch23:eq-survival}
  \Prob_t(\tau>T)=e^{-h(T-t)},\qquad T\ge t.
\end{equation}
\begin{definition}[Hazard rate]\label{ch23:def-hazard}
$h$ is the instantaneous default intensity: $\Prob_t(\tau\in[t,t+\dd t)\mid\tau>t)=h\,\dd t$.
\end{definition}
\begin{proposition}[Survival probability and default density]\label{ch23:prop-survival}
If $\tau$ has the memoryless survival function \cref{ch23:eq-survival}, then, conditional on $t<\tau$, the density
of $\tau$ at $T>t$ is $\varphi(t,T)=he^{-h(T-t)}$, and the current (\(T\to t\)) default density is $\varphi(t,t)=h$.
\end{proposition}
\begin{proof}
$\Prob_t(\tau\le T)=1-e^{-h(T-t)}$; differentiate in $T$ to get $\varphi(t,T)=he^{-h(T-t)}$; set $T=t$.
\end{proof}
This is exactly the exponential distribution: constant hazard rate $\Leftrightarrow$ $\tau-t\sim\mathrm{Exp}(h)$ given
survival to $t$, the continuous-time analogue of a Poisson arrival process (\cref{ch08:sec-martingale} used the
same intensity idea for the CDS survival measure).

\subsection{FX jump on default}\label{ch23:sec-fxjump}
The model's only new ingredient is that $S_t$ (USD per EUR) jumps by a fixed multiplicative factor at $\tau$:
\begin{equation}\label{ch23:eq-jumpdef}
  S_{\tau}=S_{\tau^-}\,e^{J},\qquad J\in\R,
\end{equation}
so $J<0$ means EUR devalues on default (the Italy/Argentina case) and $J>0$ an appreciation. Let $N_t=\ind_{\{t\ge
\tau\}}$ be the default indicator, a jump process with a single unit jump at $\tau$, $\dd N_t$ the corresponding
point-process differential ($\dd N_t=1$ at $t=\tau$, $0$ elsewhere). Away from the jump, $\log S_t$ is postulated
to drift at some rate $\mu_t$ to be determined:
\begin{equation}\label{ch23:eq-logS}
  \dd\log S_t=\mu_t\,\dd t+J\,\dd N_t .
\end{equation}
No-arbitrage (interest rates being $0$) requires $S_t$ to be a $\Q$-martingale, in particular $\E[S_T]=S_0$ for
every $T$ -- the jump analogue of the risk-neutral drift condition that fixes $\mu=r$ in \cref{ch15:eq-gbm}.

\begin{proposition}[Compensator condition]\label{ch23:prop-compensator}
If $\mu_t=h\bigl(1-e^{J}\bigr)\ind_{\{t<\tau\}}$ in \cref{ch23:eq-logS} (and $\mu_t=0$ once $\tau$ has occurred,
since after default there is nothing left to compensate), then $\E[S_T]=S_0$ for every $T\ge0$.
\end{proposition}
The quantity $h(1-e^J)$ is called the \emph{compensator}: exactly as a Poisson counting process $N_t$ has mean
$ht$ so that $N_t-ht$ is a martingale, here the drift $\mu_t\,\dd t$ must exactly offset, in expectation, the
mean rate of change contributed by the jump term $J\,\dd N_t$, whose expected instantaneous contribution (while
$\tau$ has not yet occurred) is $h\cdot J\,\dd t$ plus a jump-size correction coming from working with $\log S$
rather than $S$ itself; this is made precise in the proof below via It\^o's formula for jump processes.

\begin{lemma}[It\^o's formula for a jump process, reconstruction]\label{ch23:lem-itojump}
Let $X_t$ solve $\dd X_t=\mu_t\,\dd t+J\,\dd N_t$ for a point process $N_t$ with a single jump at $\tau$
(no diffusive part), and let $f\in C^1$. Then
\begin{equation}\label{ch23:eq-itojump}
  \dd f(X_t)=f'(X_t)\,\mu_t\,\dd t+\bigl[f(X_{t^-}+J)-f(X_{t^-})\bigr]\dd N_t .
\end{equation}
\end{lemma}
\begin{proof}
Between jumps $X_t$ is absolutely continuous with $\dot X_t=\mu_t$, so $\dd f(X_t)=f'(X_t)\dot X_t\,\dd t$ by the
ordinary chain rule; at a jump time, $f(X_t)-f(X_{t^-})=f(X_{t^-}+J)-f(X_{t^-})$ by \cref{ch23:eq-jumpdef}-type
jumps of size $J$ in $X$. Summing the continuous drift contribution and the (single) jump contribution gives
\cref{ch23:eq-itojump}. (This is the pure-jump special case of the general jump-diffusion It\^o formula; the
diffusive analogue for Brownian motion is \cref{ch17:thm-ito1}, with the second-derivative term there playing the
role that the finite difference $f(X_{t^-}+J)-f(X_{t^-})$ plays here -- both are ``convexity'' corrections beyond
the naive chain rule, one from quadratic variation, the other from a discrete increment.)
\end{proof}

\begin{proof}[Proof of \cref{ch23:prop-compensator}]
Write $X_t=\log S_t$, so $S_t=e^{X_t}=f(X_t)$ with $f(x)=e^x$. By \cref{ch23:lem-itojump} with $\mu_t$ as in the
statement,
\[
  \dd S_t = S_{t^-}\,h(1-e^J)\ind_{\{t<\tau\}}\,\dd t + S_{t^-}\bigl(e^J-1\bigr)\dd N_t ,
\]
using $f(X_{t^-}+J)-f(X_{t^-})=e^{X_{t^-}}(e^J-1)=S_{t^-}(e^J-1)$. Divide by $S_{t^-}$:
\begin{equation}\label{ch23:eq-dSoverS}
  \frac{\dd S_t}{S_t}=h(1-e^J)\ind_{\{t<\tau\}}\,\dd t+(e^J-1)\,\dd N_t .
\end{equation}
Both terms vanish in expectation instantaneously conditional on survival: on $\{t<\tau\}$, $\E[\dd N_t\mid
\mathcal F_t]=h\,\dd t$, so $\E[\dd S_t/S_t\mid\mathcal F_t,t<\tau]=h(1-e^J)\,\dd t+(e^J-1)h\,\dd t=0$; after
default nothing changes ($\mu_t=0$ and $N_t$ no longer jumps). Hence $S_t$ is a $\Q$-martingale and $\E[S_T]=S_0$.
\end{proof}

\begin{intuition}[Reading the compensator]
If EUR devalues on default ($J<0$, so $e^J<1$), the compensator $h(1-e^J)>0$: between defaults the model must make
EUR \emph{appreciate} against USD, on average, at rate $h(1-e^J)$, precisely so that the (rare, large, negative)
jump at default leaves the unconditional mean unchanged. This upward drift before default is a purely technical
artefact of no-arbitrage in a jump model, not a forecast: it is the exact analogue of the fact that a stock with a
possible downward jump (bankruptcy) must drift upward the rest of the time to keep $\E[S_T]$ fixed at the forward
level, a fact well known from equity jump-diffusion models.
\end{intuition}

\subsection{Explicit solution for \texorpdfstring{$S_t$}{S\_t}}\label{ch23:sec-explicit}
\begin{proposition}[Path of the FX rate]\label{ch23:prop-Spath}
Under \cref{ch23:eq-logS} with the compensator drift of \cref{ch23:prop-compensator},
\begin{equation}\label{ch23:eq-Spath}
  S_T=S_0\exp\Bigl[h\bigl(1-e^J\bigr)(T\wedge\tau)+J\,\ind_{\{\tau\le T\}}\Bigr],
  \qquad T\wedge\tau=\min(T,\tau).
\end{equation}
Equivalently: if $\tau>T$, $S_T=S_0\exp[h(1-e^J)T]$ (pure drift); if $\tau\le T$, $S_T=S_0\exp[h(1-e^J)\tau+J]$
(drift up to default, then the jump).
\end{proposition}
\begin{proof}
Integrate \cref{ch23:eq-logS}: $\log S_T-\log S_0=\int_0^Th(1-e^J)\ind_{\{t<\tau\}}\dd t+J\int_0^T\dd N_t$. The first
integral runs only while $t<\tau$, i.e.\ up to $T\wedge\tau$, contributing $h(1-e^J)(T\wedge\tau)$; the second is
$J$ if a jump occurred by $T$ ($\tau\le T$) and $0$ otherwise. Exponentiate.
\end{proof}

\begin{example}[Numerical check of the martingale property]\label{ch23:ex-mc}
Take $h=0.05$, $J=-0.30$ (a $30\%$ devaluation on default, roughly the order of magnitude of the Argentine peso
move of slide 18), $T=3$, $S_0=1$. Monte Carlo with $6\times10^6$ paths of $\tau\sim\mathrm{Exp}(h)$ and
\cref{ch23:eq-Spath} gives $\widehat\E[S_T]=0.99997$, matching the required value $S_0=1$ to four significant
figures (computed for these notes; ). The exact identity
$\E[S_T]=1$ can also be checked analytically by splitting the expectation on $\{\tau>T\}$ and $\{\tau\le T\}$ and
integrating $\varphi(0,\tau)=he^{-h\tau}$ against \cref{ch23:eq-Spath}; both routes agree.
\end{example}
\section{Pricing and hedging zero-recovery bonds}\label{ch23:sec-pricing}
Consider zero-coupon, zero-recovery bonds of maturity $T$ from the same issuer: $U$ pays 1 USD at $T$ if
$\tau>T$ (and $0$ otherwise); $E$ pays 1 EUR at $T$ under the same condition. Both are priced, in USD, by
risk-neutral expectation (interest rates $0$, so the money-market numeraire is $\beta(t)\equiv1$, as in
\cref{ch08:sec-martingale}):
\begin{equation}\label{ch23:eq-Pu}
  P_u(0,T)=\E\bigl[\ind_{\{\tau>T\}}\bigr]=\Prob(\tau>T)=e^{-hT},
\end{equation}
by \cref{ch23:eq-survival}, exactly as in \cref{ch15:sec-rn}: the price is the expected discounted payoff. For the
EUR bond, the USD-payoff is $S_T\ind_{\{\tau>T\}}$ (converting the 1 EUR payoff at the terminal FX rate), so
\begin{equation}\label{ch23:eq-Pe-setup}
  P_e(0,T)=\E\bigl[S_T\,\ind_{\{\tau>T\}}\bigr].
\end{equation}

\begin{proposition}[Zero-recovery bond prices]\label{ch23:prop-bondprices}
\begin{equation}\label{ch23:eq-bondprices}
  P_u(0,T)=e^{-hT},\qquad\qquad P_e(0,T)=S_0\,e^{-h e^{J}T}.
\end{equation}
\end{proposition}
\begin{proof}
$P_u$ is \cref{ch23:eq-Pu}. For $P_e$, on $\{\tau>T\}$, \cref{ch23:eq-Spath} gives $S_T=S_0e^{h(1-e^J)T}$ (a constant,
not random, given survival), so
\[
  P_e(0,T)=S_0e^{h(1-e^J)T}\,\Prob(\tau>T)=S_0e^{h(1-e^J)T}e^{-hT}=S_0\,e^{hT(1-e^J)-hT}=S_0\,e^{-h e^JT}.
  \qedhere
\]
\end{proof}
Formula \eqref{ch23:eq-bondprices} was checked numerically for these notes: with $h=0.05$, $J=-0.30$, $T=3$,
$S_0=1$, \cref{ch23:eq-bondprices} gives $P_u=0.86071$, $P_e=0.89483$, against Monte Carlo estimates $0.86059$ and
$0.89470$ from $6\times10^6$ paths -- agreement to three decimals.

\begin{intuition}[Reading $P_e$]
$P_e/S_0=e^{-h e^JT}$ rather than $e^{-hT}$: the EUR bond is priced with an \emph{effective hazard rate}
$h_{\mathrm{eff}}=he^J$ once expressed in USD. If EUR devalues on default ($J<0$, $e^J<1$), $h_{\mathrm{eff}}<h$
and $P_e/S_0>P_u$: in USD terms the EUR bond looks \emph{less} risky than the USD bond of the same issuer, because
its payoff is worth relatively little exactly in the states where default has happened -- the market has less to
lose, in USD, from that scenario. This is the quanto adjustment for credit: the raw survival probability $e^{-hT}$
is corrected by $e^J$ to account for the correlation between the default event and the currency of the payoff. It
matches the qualitative content of the Italy example of \cref{ch23:sec-italy}: everything else equal, a devaluing
local-currency bond should trade with a \emph{lower} USD-equivalent implied hazard rate than the hard-currency
bond of the same issuer -- exactly the pattern the slides point to (EUR spread below USD spread once basis-adjusted).
\end{intuition}

\subsection{The replicating portfolio}\label{ch23:sec-replicate}
Sell one $U$ bond and buy $N_E(t)=P_u(t,T)/P_e(t,T)$ $E$-bonds so that the portfolio costs $0$ at every rebalancing
date; by \cref{ch23:eq-bondprices} (evaluated at $t$ instead of $0$, since \cref{ch23:prop-bondprices} only uses
the survival probability from $t$ to $T$, hence holds for any $t<\tau$),
\begin{equation}\label{ch23:eq-hedgeratio}
  N_E(t)=\frac{P_u(t,T)}{P_e(t,T)}=\frac{e^{-h(T-t)}}{S_te^{-he^J(T-t)}}=\frac{1}{S_t}\exp\bigl[h(T-t)(e^J-1)\bigr].
\end{equation}
\begin{proposition}[The portfolio is a hedge]\label{ch23:prop-replicate}
The portfolio $\Pi_t=-P_u(t,T)+N_E(t)P_e(t,T)\equiv0$ for every $t<\tau$; if default occurs, both bonds are worth
$0$ (zero recovery) so $\Pi_\tau=0$ as well. Hence the portfolio replicates the (trivial) zero payoff in every
state, and the hedge ratio \cref{ch23:eq-hedgeratio} is the unique arbitrage-free hedge of one bond against the
other.
\end{proposition}
\begin{proof}
$\Pi_t=-P_u(t,T)+N_E(t)P_e(t,T)=-P_u(t,T)+P_u(t,T)=0$ by construction of $N_E(t)$, for every $t<\tau$ (not only
$t=0$), because \cref{ch23:eq-bondprices} depends on $t$ only through the remaining maturity $T-t$ and through
$S_t$, and $N_E(t)$ was defined by dividing them; the $S_t$ dependence cancels exactly between $P_e(t,T)$'s
explicit $S_t$ factor and $N_E(t)$'s $1/S_t$ factor. At default both bonds default to $0$ recovery, so $\Pi_\tau=0$
trivially. What makes this a genuine hedge, and not a tautology, is that the rebalancing is \emph{self-financing}:
before default $\dd\log P_u(t,T)=h\,\dd t$ and $\dd\log P_e(t,T)=\dd\log S_t+he^J\dd t=h\,\dd t$ (by
\cref{ch23:eq-logS} with the compensator drift), and both bonds jump to $0$ at $\tau$, so
$\dd P_u/P_{u}=\dd P_e/P_{e}=h\,\dd t-\dd N_t$. Hence $-\dd P_u+N_E\,\dd P_e=0$: the gains of the position are
identically zero, and the $N_E(t)$ $E$-bonds, rebalanced at no cost, replicate one $U$-bond.
\end{proof}
This was checked numerically: with the parameters of \cref{ch23:ex-mc} and $T=3$, $N_E(0)=0.96187$
($=e^{h T(e^J-1)}$ since $S_0=1$); stepping $\Delta t=0.4$ forward with no default, both $P_u$ and $N_E(t)P_e(t,T)$
were recomputed from \cref{ch23:eq-bondprices}--\eqref{ch23:eq-hedgeratio} and agreed to machine precision, confirming \cref{ch23:prop-replicate} directly rather than only through the algebraic
cancellation in the proof.

\begin{remark}[Why this is harder than static replication]\label{ch23:rem-dynamic}
$N_E(t)$ depends on $t$ (through the remaining maturity) and on $S_t$ (through the $1/S_t$ factor): the hedge must
be \emph{continuously rebalanced}, unlike a naive static forward hedge. It also depends on $h$ and $J$: how many
EUR bonds are needed to hedge one USD bond depends on both the issuer's credit riskiness and the size of the FX
move it would trigger, not only on the current spot or forward FX rate. This is the chapter's central point: an
FX forward, which prices only the correlation-free part of the FX/credit relationship, cannot replicate a
credit-quanto payoff; only a model with an explicit joint dynamics for $(\tau,S)$ can.
\end{remark}

\section{Non-zero recovery}\label{ch23:sec-recovery}
If the recovery rate is $R\in(0,1)$, the bond payoffs at $T$ become $1$ (USD, resp.\ EUR) if $\tau>T$ and $R$ if
$\tau\le T$.
\begin{proposition}[Bond prices with recovery, reconstruction]\label{ch23:prop-recovery}
\begin{equation}\label{ch23:eq-recovery}
  P_u^R(0,T)=R+(1-R)e^{-hT},\qquad\qquad P_e^R(0,T)=S_0\Bigl[R+(1-R)e^{-he^JT}\Bigr].
\end{equation}
\end{proposition}
\begin{proof}
$P_u^R(0,T)=\E[\ind_{\{\tau>T\}}+R\ind_{\{\tau\le T\}}]=\Prob(\tau>T)+R\Prob(\tau\le T)=e^{-hT}+R(1-e^{-hT})$, which
rearranges to the stated form. For $P_e^R(0,T)/S_0=\E[e^{X_T}\ind_{\{\tau>T\}}+Re^{X_T}\ind_{\{\tau\le T\}}]$ with
$X_T=\log(S_T/S_0)$ from \cref{ch23:eq-Spath}: the first term is $P_e(0,T)/S_0=e^{-he^JT}$ by
\cref{ch23:prop-bondprices}. The second term integrates over the default density $\varphi(0,\tau)=he^{-h\tau}$ of
\cref{ch23:prop-survival}, using $X_T=h(1-e^J)\tau+J$ on $\{\tau\le T\}$:
\[
  R\int_0^Te^{h(1-e^J)\tau+J}\,he^{-h\tau}\,\dd\tau
  =Rhe^J\!\int_0^Te^{-he^J\tau}\,\dd\tau
  =Rhe^J\cdot\frac{1-e^{-he^JT}}{he^J}
  =R\bigl(1-e^{-he^JT}\bigr).
\]
Adding the two terms, $P_e^R(0,T)/S_0=e^{-he^JT}+R(1-e^{-he^JT})=R+(1-R)e^{-he^JT}$.
\end{proof}
The two formulas in \cref{ch23:eq-recovery} are strikingly parallel: each is $R+(1-R)$ times the corresponding
zero-recovery survival probability, with $h$ replaced by $he^J$ in the EUR case. This was checked both by
numerical quadrature of the integral above and by Monte Carlo (parameters as in \cref{ch23:ex-mc}, $R=0.25$):
$P_u^R=0.89553$ (formula) vs.\ $0.89544$ (MC), $P_e^R=0.92112$ (formula, matching the quadrature to 6 digits) vs.\
$0.92102$ (MC).

\begin{remark}[Exact replication is lost]\label{ch23:rem-Rgt0}
With $R>0$ the ratio $P_u^R(t,T)/P_e^R(t,T)$ no longer has the clean multiplicative form of
\cref{ch23:eq-hedgeratio} (the $R$ and $(1-R)e^{-h(\cdot)T}$ terms scale differently with $S_t$ and with $T-t$), so
the zero-cost ratio $P_u^R/P_e^R$ no longer matches the two bonds' jumps at default, and a portfolio of the two bonds
alone is no longer a hedge. The model is still complete (one jump, one risky hedge instrument): holding
$N_E(t)=\Delta P_u^R/\Delta P_e^R$ $E$-bonds, the ratio of the jumps at default
($\Delta P_u^R=R-P_u^R(t,T)$, $\Delta P_e^R=RS_te^J-P_e^R(t,T)$), plus the USD cash needed to match the value,
replicates the $U$-bond, because both prices are $\Q$-martingales and their drifts are fixed by their jumps
(reconstruction, not in the slides). The slides stop here (``even in such a simple model,
hedging is not trivial''); \cref{ch23:sec-jumpdiffusion} sketches the further step (adding a diffusive component)
proposed as a final-project topic.
\end{remark}

\begin{figure}[htbp]
\centering
\begin{tikzpicture}
\begin{axis}[
  width=11cm,height=6.5cm,
  xlabel={maturity $T$ (years)}, ylabel={bond price (USD terms, $S_0=1$)},
  xmin=0,xmax=5, ymin=0.7,ymax=1.01,
  legend pos=south west, legend style={font=\footnotesize},
  grid=both, grid style={gray!15},
]
\addplot[thick,black] coordinates {
(0.00,1.0000) (0.25,0.9876) (0.50,0.9753) (0.75,0.9632) (1.00,0.9512) (1.25,0.9394) (1.50,0.9277)
(1.75,0.9162) (2.00,0.9048) (2.25,0.8936) (2.50,0.8825) (2.75,0.8715) (3.00,0.8607) (3.25,0.8500)
(3.50,0.8395) (3.75,0.8290) (4.00,0.8187) (4.25,0.8086) (4.50,0.7985) (4.75,0.7886) (5.00,0.7788)};
\addlegendentry{$P_u(0,T)=e^{-hT}$}
\addplot[thick,blue,dashed] coordinates {
(0.00,1.0000) (0.25,0.9908) (0.50,0.9816) (0.75,0.9726) (1.00,0.9636) (1.25,0.9548) (1.50,0.9460)
(1.75,0.9372) (2.00,0.9286) (2.25,0.9200) (2.50,0.9116) (2.75,0.9032) (3.00,0.8948) (3.25,0.8866)
(3.50,0.8784) (3.75,0.8703) (4.00,0.8623) (4.25,0.8543) (4.50,0.8465) (4.75,0.8387) (5.00,0.8309)};
\addlegendentry{$P_e(0,T)/S_0$, $J=-0.30$ (devalues)}
\addplot[thick,red,dotted] coordinates {
(0.00,1.0000) (0.25,0.9833) (0.50,0.9668) (0.75,0.9506) (1.00,0.9347) (1.25,0.9191) (1.50,0.9037)
(1.75,0.8886) (2.00,0.8737) (2.25,0.8591) (2.50,0.8447) (2.75,0.8306) (3.00,0.8167) (3.25,0.8030)
(3.50,0.7896) (3.75,0.7764) (4.00,0.7634) (4.25,0.7506) (4.50,0.7381) (4.75,0.7257) (5.00,0.7136)};
\addlegendentry{$P_e(0,T)/S_0$, $J=+0.30$ (appreciates)}
\end{axis}
\end{tikzpicture}
\caption{Zero-recovery bond prices \eqref{ch23:eq-bondprices}, $h=0.05$. A currency that devalues on default
($J<0$) makes the local-currency bond look \emph{less} risky in hard-currency terms ($P_e>P_u$); a currency that
appreciates on default ($J>0$, unusual but not impossible for a safe-haven funding currency) makes it look
\emph{more} risky ($P_e<P_u$). Coordinates computed for these notes in Python.}\label{ch23:fig-bonds}
\end{figure}
\section{Change of numeraire view}\label{ch23:sec-numeraire}
The pricing formula \eqref{ch23:eq-Pe-setup}, $P_e(0,T)/S_0=\E[e^{X_T}\ind_{\{\tau>T\}}]$, can be read as a change
of numeraire in the sense of \cref{ch08:prop-numeraire}. Take $N_1(t)=S_t$ (EUR converted to USD, a traded USD
asset since interest rates are $0$) and $N_2(t)=1$ (USD cash). By \cref{ch08:prop-numeraire},
\[
  \E_0^{\Q}\bigl[Y_T\bigr]=\frac{N_1(0)}{N_2(0)}\,\E_0^{\Q_{S}}\!\left[Y_T\,\frac{N_2(T)}{N_1(T)}\right]
  \quad\Longrightarrow\quad
  \E_0^{\Q}\bigl[S_TY_T\bigr]=S_0\,\E_0^{\Q_S}[Y_T],
\]
where $\Q_S$ is the measure associated with numeraire $S$ (the ``EUR risk-neutral measure expressed in USD units,''
by the Radon--Nikod\'ym derivative $\dd\Q_S/\dd\Q|_{\mathcal F_T}=S_T/S_0$, since $S_t/1$ is itself the
$\Q$-martingale that must be used in \cref{ch08:prop-numeraire}). With $Y_T=\ind_{\{\tau>T\}}$,
\begin{equation}\label{ch23:eq-numeraire-view}
  P_e(0,T)=S_0\,\Prob_0^{\Q}(\tau>T)\Big|_{\text{naive}}\ \ne\ P_e(0,T)=S_0\,\Q_S(\tau>T),
\end{equation}
and \cref{ch23:prop-bondprices} says precisely that $\Q_S(\tau>T)=e^{-he^JT}\ne e^{-hT}=\Q(\tau>T)$: default is
\emph{less} likely under the EUR-numeraire measure $\Q_S$ than under $\Q$ when $J<0$ (EUR weakens on default) --
because $\Q_S$ tilts probability towards paths on which $S_T$ is large, i.e.\ towards the surviving, no-jump
scenario. This is exactly the same mechanism as the FX-forward row of the numeraire table in
\cref{ch08:sec-martingale}, $F_{FX}(t,T)=S_{FX}(t)P^F(t,T)/P^D(t,T)$, except that here the ``foreign bond'' $P^F$ is
itself default-risky and its numeraire property breaks down exactly at $\tau$ (it is only a valid numeraire, i.e.\
strictly positive, on $\{t<\tau\}$) -- the same obstruction \cref{ch08:sec-martingale} met when trying to use the
risky annuity $\tilde A_j$ as a numeraire for the CDS survival measure, and resolved there the same way: work with
$\ind_{\{t<\tau\}}$ times the relevant traded quantity, which \emph{is} a genuine $\Q$-martingale even though the
numeraire itself is not always positive.

\begin{intuition}[Quanto adjustment = change of measure]
Every ``quanto adjustment'' in this course -- FX quanto options, quanto CDS here -- is a change of numeraire from
the natural currency's risk-neutral measure to the payoff currency's risk-neutral measure. When the underlying
(here, the default indicator) is independent of the FX rate, the two measures agree and quanto pricing reduces to
discounting by the forward FX rate. When they are correlated (through the jump $J$), the Radon--Nikod\'ym
derivative $S_T/S_0$ tilts the default probability, and no amount of forward-FX-based reasoning like
\cref{ch23:sec-naive} can substitute for computing the expectation under the correct measure.
\end{intuition}

\section{Quanto credit default swaps and other applications}\label{ch23:sec-quanto-cds}
A credit default swap (CDS) is a bilateral insurance contract on the issuer's default: the protection buyer pays a
periodic spread (the premium leg) until default or maturity, and receives $1-R$ per unit notional from the
protection seller if default occurs (the protection leg). A \emph{quanto CDS} is a CDS whose notional and premium
are denominated in a currency different from the reference obligation's natural currency -- e.g.\ a CDS on Brazil
denominated in Brazilian reais rather than in USD. Exactly the mechanism of this chapter applies: if the reference
entity's local currency devalues on default (as is typical for a sovereign or an issuer whose revenues are tied to
its home currency), then a protection payment fixed in local currency is worth less, in hard-currency terms,
precisely when it is needed -- the quanto CDS spread must be adjusted relative to the hard-currency CDS spread by
a factor analogous to $e^J$ in \cref{ch23:eq-bondprices}, and the correct hedge ratio between the two currencies'
CDS again requires the jump-on-default FX model rather than a simple FX forward. Historically, quanto CDS were
often priced by simply converting the hard-currency CDS spread at the current or forward FX rate, ignoring the
correlation; this reproduces neither the observed market prices nor a valid hedge ratio, which is exactly why a
model with the jump feature became necessary in practice.

\begin{casestudy}[Sovereign quanto CDS baskets]
During the European sovereign crisis (2010--2012), USD- and EUR-denominated CDS on peripheral euro-area sovereigns
(Italy, Spain, Portugal) diverged, giving rise to an actively quoted ``quanto'' basis. Market participants used
models of the type developed here (often with stochastic hazard rates and a jump-diffusion FX process, as sketched
in \cref{ch23:sec-jumpdiffusion}) to convert between the USD and EUR curves and to size cross-currency hedges,
essentially betting on how far the euro would fall if a member state defaulted or exited the currency union.
\end{casestudy}

\section{Towards a more realistic model}\label{ch23:sec-jumpdiffusion}
The pure-jump FX process of \cref{ch23:eq-dSoverS} has no day-to-day volatility except at $\tau$: a natural next
step (posed on the slides as a final-project topic) is to add a diffusive component,
\begin{equation}\label{ch23:eq-jumpdiffusion}
  \frac{\dd S_t}{S_t}=h\bigl(1-e^J\bigr)\ind_{\{t<\tau\}}\,\dd t+\bigl(e^J-1\bigr)\dd N_t+\sigma\,\dd W_t,
\end{equation}
a genuine jump-diffusion in the sense of \cref{ch17:sec-lemma} plus \cref{ch23:sec-fxjump} combined: away from
$\tau$, $S_t$ behaves like the ordinary geometric Brownian motion of \cref{ch23:eq-gk}; at $\tau$, it also jumps by
$e^J$. Pricing bonds in this model requires It\^o's formula with \emph{both} a diffusive and a jump term (the
sum of the Brownian second-derivative correction of \cref{ch17:thm-ito2} and the finite-difference jump correction
of \cref{ch23:lem-itojump}), and hedging now needs (at least) two instruments -- an FX forward for the diffusive
risk and a position sensitive to the jump (e.g.\ a short-dated deep out-of-the-money FX option, or the bonds
themselves) -- since the model has two independent sources of randomness ($W$ and $N$), matching the ``number of
hedging instruments = number of stochastic sources'' completeness criterion of \cref{ch15:sec-onestep}. Once a
closed-form or PDE solution is obtained, it can (and should) be checked against a Monte Carlo simulation that
simulates $\tau$, $W$, and hence $S_T$, jointly, exactly as done at smaller scale in \cref{ch23:ex-mc}. In the more
realistic desk models this course only sketches, both the hazard rate and interest rates are themselves stochastic
and correlated with FX (``Real Life'', slide 29); \cref{ch:hjm} develops the term-structure machinery needed to
make hazard rates stochastic in a no-arbitrage way, and correlating them with FX and with rates is a natural
(unaddressed here) extension of the multi-factor models built there. Because defaults are rare, naive Monte Carlo
is inefficient (most paths never see a jump); practitioners typically integrate the payoff analytically over the
default-time density $\varphi(0,\tau)$ (as done by hand in \cref{ch23:prop-recovery}) and simulate only the
diffusive factors (FX, rates, hazard-rate level) by Monte Carlo, mixing the two techniques.

\section{Girsanov and the diffusive quanto adjustment}\label{ch23:sec-girsanov}
It is worth relating \cref{ch23:eq-jumpdiffusion} back to the purely diffusive quanto adjustment familiar from FX
options (not covered in the 2013 lecture, but a standard consequence of \cref{ch17:sec-girsanov}). If a foreign
asset $V_t$ (say, a EUR equity index) follows a geometric Brownian motion under
the foreign risk-neutral measure $\Q^F$ and the FX rate $S_t$ one under the domestic measure $\Q^D$,
\[
  \frac{\dd V_t}{V_t}=r_f\,\dd t+\sigma_V\,\dd W_t^F,\qquad
  \frac{\dd S_t}{S_t}=(r_d-r_f)\,\dd t+\sigma_S\,\dd \widetilde W_t^D,\qquad
  \dd W^F\,\dd\widetilde W^D=\rho\,\dd t,
\]
then pricing a USD-denominated (\emph{quantoed}) payoff on $V_T$ requires switching from $\Q^F$ to the domestic
measure $\Q^D$. Girsanov's theorem (\cref{ch17:sec-girsanov}) says this changes the drift of $V$ by
$-\rho\sigma_V\sigma_S$: under $\Q^D$, $V$ grows at $r_f-\rho\sigma_V\sigma_S$ rather than at $r_f$. This is the
textbook ``quanto drift adjustment,'' and it is the diffusive sibling of the jump compensator
\cref{ch23:prop-compensator}: in both cases, expressing a foreign payoff in domestic currency under the correlated
FX process shifts the effective drift of the payoff by an amount proportional to the FX/payoff correlation --
$-\rho\sigma_V\sigma_S$ in the Brownian case, $h(e^J-1)$ (via the jump-driven change of measure of
\cref{ch23:sec-numeraire}) in the jump case. \cref{ch23:eq-jumpdiffusion} simply has both effects at once.

\begin{keyformulas}
\begin{align*}
&\text{Survival:} && \Prob_t(\tau>T)=e^{-h(T-t)}\quad\text{(\cref{ch23:eq-survival})}\\
&\text{FX dynamics:} && \dd S_t/S_t=h(1-e^J)\ind_{\{t<\tau\}}\dd t+(e^J-1)\dd N_t\quad\text{(\cref{ch23:eq-dSoverS})}\\
&\text{Zero-recovery bonds:} && P_u(0,T)=e^{-hT},\quad P_e(0,T)=S_0\,e^{-he^JT}\quad\text{(\cref{ch23:eq-bondprices})}\\
&\text{Hedge ratio:} && N_E(t)=\tfrac1{S_t}\exp[h(T-t)(e^J-1)]\quad\text{(\cref{ch23:eq-hedgeratio})}\\
&\text{Bonds, recovery $R$:} && P_u^R(0,T)=R+(1-R)e^{-hT}\\
& && P_e^R(0,T)=S_0[R+(1-R)e^{-he^JT}]\quad\text{(\cref{ch23:eq-recovery})}
\end{align*}
\end{keyformulas}

\section*{Exercises}
\addcontentsline{toc}{section}{Exercises}
\begin{exercise}[No devaluation, not from the course]\label{ch23:ex-J0}
Show that if $J=0$ (no FX jump on default), \cref{ch23:eq-bondprices} collapses to $P_u(0,T)=P_e(0,T)/S_0=e^{-hT}$
and the hedge ratio \eqref{ch23:eq-hedgeratio} collapses to $N_E(t)=1/S_t$ -- the ordinary FX conversion, with no
quanto correction, as expected when default carries no currency information.
\end{exercise}

\begin{exercise}[Sign of the quanto adjustment, not from the course]\label{ch23:ex-sign}
Using \cref{ch23:eq-bondprices}, show that $P_e(0,T)/S_0>P_u(0,T)$ if and only if $J<0$ (devaluation on default),
for every $h>0,T>0$. Interpret: does a devaluing local currency make the local-currency bond, expressed in hard
currency, look safer or riskier than the equivalent hard-currency bond of the same issuer?
\end{exercise}

\begin{exercise}[Small-jump limit, not from the course]\label{ch23:ex-small}
Expand $P_e(0,T)$ of \cref{ch23:eq-bondprices} to first order in $J$ around $J=0$ and show
$P_e(0,T)\approx S_0e^{-hT}(1-hTJ)$. Compare with $P_u(0,T)=e^{-hT}$ and give a one-line rule of thumb for the
size of the quanto adjustment to the credit spread (defined as $-\tfrac1T\log P$) when $J$ is small.
\end{exercise}

\begin{exercise}[Recovery limits, not from the course]\label{ch23:ex-recovery-limits}
Check that \cref{ch23:eq-recovery} reduces to \cref{ch23:eq-bondprices} when $R=0$, and to $P_u^R=P_e^R/S_0=1$ for
every $T$ when $R=1$ (a bond that always pays its notional regardless of default is riskless). Explain why the
second limit no longer depends on $h$ or $J$.
\end{exercise}

\begin{exercise}[Compensator with two jump sizes, not from the course]\label{ch23:ex-twojumps}
Suppose $S_t$ can jump on default to either $S_{\tau^-}e^{J_1}$ (probability $p$) or $S_{\tau^-}e^{J_2}$
(probability $1-p$), independently drawn at $\tau$. Redo the proof of \cref{ch23:prop-compensator} to find the
compensator drift $\mu_t$ that keeps $S_t$ a $\Q$-martingale, and check that it reduces to
\cref{ch23:prop-compensator} when $p=1$.
\end{exercise}

\begin{exercise}[Monte Carlo verification, not from the course]\label{ch23:ex-mc-exercise}
Reproduce \cref{ch23:ex-mc} in Python: simulate $N=10^6$ exponential default times with hazard rate $h=0.08$,
$J=-0.15$, $T=2$, compute $S_T$ from \cref{ch23:eq-Spath}, and check numerically that $\widehat\E[S_T]\approx1$,
$\widehat\E[\ind_{\{\tau>T\}}]\approx P_u(0,T)$ and $\widehat\E[S_T\ind_{\{\tau>T\}}]\approx P_e(0,T)/S_0$ from
\cref{ch23:eq-bondprices}.
\end{exercise}

% !TEX root = ../main.tex
\chapter{The HJM Model for Interest Rates and Credit}\label{ch:hjm}
\begin{paracol}{2}

\begin{sources}
\begin{tabularx}{\linewidth}{@{}l l L@{}}
\toprule
\textbf{Lecture} & \textbf{Speaker} & \textbf{Content}\\
\midrule
2013 L24 & D.~Gorokhov (Morgan Stanley) & Monte Carlo pricing motivated by dynamic hedging of stock options (Black--Scholes recap, Ito's lemma ``under the microscope''); why yield-curve dynamics is harder than a point process; the classic short-rate models (Ho--Lee, Hull--White, CIR) listed without derivation; the HJM forward-rate SDE, the no-arbitrage drift condition derived on the slides via the zero-coupon bond price and Ito's lemma; a practical recipe for pricing with HJM; credit derivatives (survival probabilities, hazard rates, CDS pricing), the hazard-rate HJM analogue and its drift condition left as an exercise on the slides; structured notes as a worked motivation for simulating rates, equity and credit jointly.\\
\bottomrule
\end{tabularx}

\smallskip
\textbf{Files.} \file{MIT18\_S096F13\_lecnote24.pdf} (34 slides).
\textbf{Problem sets.} None assigned on this topic in either year. The exercises (each after the section it practises) are marked \emph{(not from the course)}.
\textbf{Additions of these notes} (marked as such where they occur): the full derivation of the HJM drift condition from Ito's lemma with every step spelled out (the slides state the result and sketch the computation in compressed notation); the explicit recovery of Ho--Lee, Vasicek/Hull--White and the impossibility of an exact CIR recovery as special cases of the HJM volatility structure; the real-world/risk-neutral distinction via Girsanov's theorem (\cref{ch:stochcalc}); a symmetric full derivation of the credit (hazard-rate) drift condition, of which the slides only pose the no-correlation special case as an exercise; and all numerical checks.
\end{sources}

\section*{Motivation}
\Cref{ch:rates} built today's discount curve $Z(0,T)$ from swap quotes and priced linear products from it. \Cref{ch:sde} showed how to write a stochastic differential equation for a single number, the short rate $r_t$, and \cref{ch18:fin-vasicek} used the Ornstein--Uhlenbeck process for this purpose. But a bank does not trade one number: it trades against the whole forward curve, and pricing an option on a swap fifteen years out requires knowing how the \emph{entire curve} can move between now and then, not just its short end. Heath, Jarrow and Morton's 1992 answer is disarmingly simple: write down a stochastic differential equation directly for the forward-rate curve $T\mapsto f(t,T)$, one Brownian increment $\dd B_t$ shared by all maturities but a maturity-dependent volatility $\sigma(t,T)$, and ask what no-arbitrage requires. The answer -- the \emph{HJM drift condition} -- says that the drift of $f(t,T)$ is not free: it is pinned down by $\sigma(t,\cdot)$ alone, exactly as the drift of a stock disappeared from the Black--Scholes equation in \cref{ch:bs} once the stock was used to hedge. Four questions organise the chapter:
\begin{itemize}
  \item \emph{How do you write an SDE for a curve, and what does absence of arbitrage force on its drift?} \Cref{ch24:sec-setup,ch24:sec-drift}.
  \item \emph{Is this new, or have we seen it before?} Ho--Lee, Vasicek/Hull--White and (almost) CIR are the special cases where the whole curve dynamics collapses onto the short rate alone (\cref{ch24:sec-shortrate}).
  \item \emph{Does the condition actually work?} A Monte Carlo simulation of the full curve is checked against the known bond price (\cref{ch24:sec-numerics}).
  \item \emph{Does the same idea price credit risk?} Replace the discount curve by the survival-probability curve and the forward rate by the forward hazard rate; the same drift argument returns, term by term (\cref{ch24:sec-credit}).
\end{itemize}
For a physicist, the HJM equation is a stochastic field theory in one spatial dimension (maturity $T$) driven by a single noise; the no-arbitrage condition is a compatibility constraint between the drift and covariance kernels, of the same flavour as a fluctuation--dissipation relation.

% ==================================================================
\section{The forward-rate curve as a stochastic field}\label{ch24:sec-setup}

\subsection{Notation: two clocks}
\Cref{ch06:sec-curve} built the discount curve $Z(0,T)$ observed \emph{today} ($t=0$) for every future maturity $T$. The instantaneous forward rate seen today is $f(0,T)=-\partial_T\ln Z(0,T)$ (\cref{ch06:prop-cdf}, continuous-time limit). The HJM model asks: as calendar time $t$ moves forward, how does the \emph{whole function} $T\mapsto f(t,T)$, $T\ge t$, evolve? This introduces two clocks, exactly as \cref{ch18:sec-exist} of \cref{ch:sde} introduces $\mu(t,x)$: $t$ is the running (stochastic) time and $T$ is a fixed label, the maturity, carried along unchanged as $t$ increases. This is the one genuinely new feature relative to a stock price $S_t$, which has no second label \ts{13}{24}{45:22}.

\begin{definition}[Forward-rate curve, zero-coupon bond, short rate]\label{ch24:def-fwd}
For each $t\ge0$, $f(t,\cdot):[t,\infty)\to\R$ is the instantaneous forward curve observed at time $t$: $f(t,T)$ is the rate, contracted at $t$, for infinitesimal borrowing starting at $T$. The time-$t$ price of the zero-coupon bond maturing at $T\ge t$ is
\begin{equation}\label{ch24:eq-bond}
  Z(t,T)=\exp\Bigl(-\int_t^Tf(t,s)\,\dd s\Bigr)=:e^{-X_t^T},\qquad X_t^T:=\int_t^Tf(t,s)\,\dd s .
\end{equation}
The short rate is the forward curve evaluated at its own running end, $r_t:=f(t,t)$.
\end{definition}
\ts{13}{24}{46:04}\ts{13}{24}{1:02:41} This is exactly \eqref{ch06:eq-fwd} of \cref{ch:rates} taken to the continuous-time, instantaneous limit, now allowed to be random and re-observed at every $t$; $Z(0,T)$ recovers today's discount curve, and $Z(t,t)=1$ for every $t$ (a bond is worth its notional at its own maturity).

\subsection{The HJM stochastic differential equation}
The starting postulate of the model, parallel to $\dd S_t=\mu S_t\dd t+\sigma S_t\dd B_t$ for a stock, is that for every fixed $T$ the process $t\mapsto f(t,T)$ ($t\le T$) is an Ito process
\begin{equation}\label{ch24:eq-hjmsde}
  \dd f(t,T)=\alpha(t,T)\,\dd t+\sigma(t,T)\,\dd B_t ,\qquad f(0,T)=f_0(T)\ \text{given},
\end{equation}
\ts{13}{24}{45:22} where $B$ is a single standard Brownian motion (one-factor case; \cref{ch24:rem-multi} extends this to several factors) and $\alpha(t,\cdot),\sigma(t,\cdot)$ are, for each $t$, deterministic or path-dependent functions of $T\ge t$, adapted in $t$ and smooth enough in $T$ to differentiate and integrate freely (technical conditions as in \cref{ch18:thm-eu}, not discussed on the slides). One equation \eqref{ch24:eq-hjmsde} for every $T$ is really a continuum of coupled SDEs, all driven by the \emph{same} $B_t$: a single shock moves the whole curve, with an amount and a sign that can differ by maturity through $\sigma(t,T)$.
\begin{intuition}[A field theory with one noise]
Think of $f(t,\cdot)$ as a field on the half-line $T\ge t$ and \eqref{ch24:eq-hjmsde} as its (stochastic) equation of motion. Because a single $B_t$ drives every $T$, the instantaneous correlation between any two points of the curve is $\pm1$ (they are driven by the same shock); what makes the curve twist and not just shift in parallel is that $\sigma(t,T)$ can change sign or magnitude with $T$. A short-lived, front-loaded $\sigma(t,T)$ (large near $T=t$, decaying in $T-t$) gives a curve that mostly moves at the short end, exactly the "level, slope, curvature" factor decomposition of \cref{ch:pca}.
\end{intuition}

\subsection{From the curve SDE to the bond-price SDE}
Everything the model needs to price a derivative -- the growth rate of tradable zero-coupon bonds -- comes from differentiating $Z(t,T)=e^{-X_t^T}$. This requires the stochastic differential of $X_t^T=\int_t^Tf(t,s)\dd s$, an integral over $s$ of a process that is itself stochastic in $t$; this is the one calculation the slides perform in compressed notation \ts{13}{24}{1:02:41}, spelled out in full here.

\begin{lemma}[Differential of the integrated forward curve]\label{ch24:lem-dX}
Under \eqref{ch24:eq-hjmsde} and Leibniz's rule for differentiating under the integral sign (applied formally to the Ito differential, justified by Fubini for the drift and by stochastic Fubini for the martingale part -- both quoted, not proved in class),
\begin{equation}\label{ch24:eq-dX}
  \dd X_t^T=-f(t,t)\,\dd t+\Bigl(\int_t^T\alpha(t,s)\,\dd s\Bigr)\dd t+\Bigl(\int_t^T\sigma(t,s)\,\dd s\Bigr)\dd B_t .
\end{equation}
\end{lemma}
\begin{proof}
Write $X_t^T=\int_0^Tf(t,s)\dd s-\int_0^tf(t,s)\dd s$ is not quite right because the lower limit is also $t$; instead fix $t$ and differentiate $X_{t+h}^T-X_t^T$ directly. Split
\[
  X_{t+h}^T-X_t^T=\underbrace{\int_{t+h}^T\bigl[f(t+h,s)-f(t,s)\bigr]\dd s}_{(\mathrm{I})}\;-\;\underbrace{\int_t^{t+h}f(t,s)\,\dd s}_{(\mathrm{II})} .
\]
For (II), $f(t,\cdot)$ is continuous in $s$ near $s=t$, so $(\mathrm{II})=f(t,t)\,h+o(h)=r_t\,h+o(h)$. For (I), insert \eqref{ch24:eq-hjmsde} pointwise in $s$ and interchange $\int_{t+h}^T(\cdot)\dd s$ with the (finite, over $s\in[t+h,T]$) sum of infinitesimal increments in $t$:
\[
  (\mathrm{I})=\int_t^T\alpha(t,s)\,\dd s\;h+\int_t^T\sigma(t,s)\,\dd s\;\Delta B_t+o(h)+O(h)\cdot h,
\]
where the upper limit was extended from $t+h$ to $t$ at a cost $O(h^2)$, and $\Delta B_t=B_{t+h}-B_t$ is the same shock for every $s$ (the interchange of the $\dd s$-integral and the stochastic increment is the stochastic Fubini theorem: valid here because $\sigma(t,s)$ is bounded and jointly measurable on the compact range $s\in[t,T]$). Collecting the $O(h)$ and $O(\sqrt h)$ terms and passing to the Ito differential gives \eqref{ch24:eq-dX}.
\end{proof}
\begin{coursenote}[Reading the slide's compressed derivation]
The 2013 slide (which writes the drift as $\mu_{tT}$ instead of $\alpha(t,T)$) states $f(t,t)=r_t$ and writes $\dd X_t=-f_{tt}\dd t-\int_t^T\dd f_{ts}\,\dd s$, then substitutes \eqref{ch24:eq-hjmsde} under the integral sign in one line \ts{13}{24}{1:02:41}. The boundary term $-f(t,t)\dd t$ is right (the integration window $[t,T]$ loses its left end as $t$ advances), but the sign in front of the integral is a typo: it should be $+\int_t^T\dd f(t,s)\,\dd s$, as in \eqref{ch24:eq-dX}. The slide's next line, the bond-price SDE, has the correct signs and agrees with \eqref{ch24:eq-dZ}, so the typo does not propagate.
\end{coursenote}

Applying Ito's formula (\cref{ch17:sec-lemma}) to $Z_t^T=e^{-X_t^T}$, a smooth function of the single Ito process $X^T$,
\begin{align}
  \dd Z_t^T&=-e^{-X_t^T}\,\dd X_t^T+\tfrac12e^{-X_t^T}\,(\dd X_t^T)^2\notag\\
  &=Z_t^T\Bigl[-\dd X_t^T+\tfrac12\bigl(\dd X_t^T\bigr)^2\Bigr]\notag\\
  &=Z_t^T\Bigl[r_t\,\dd t-\Bigl(\int_t^T\!\alpha(t,s)\dd s\Bigr)\dd t-\Bigl(\int_t^T\!\sigma(t,s)\dd s\Bigr)\dd B_t+\tfrac12\Bigl(\int_t^T\!\sigma(t,s)\dd s\Bigr)^{\!2}\dd t\Bigr],\label{ch24:eq-dZ}
\end{align}
using $(\dd X_t^T)^2=\bigl(\int_t^T\sigma(t,s)\dd s\bigr)^2\dd t$ (the $\dd t$ and cross terms of \eqref{ch24:eq-dX} do not contribute to the quadratic variation, as in \cref{ch17:sec-lemma}). This is exactly the slide's computation \ts{13}{24}{1:03:52}, reproduced here with every substitution shown.

% ==================================================================
\section{The HJM no-arbitrage drift condition}\label{ch24:sec-drift}

\subsection{Statement and proof}
A bond is a tradable asset with no dividend. By the fundamental theorem of asset pricing (\cref{ch15:sec-rn}, \cref{ch08:prop-numeraire}), under the risk-neutral measure $\Q$ every such asset, discounted by the money-market account $\exp(\int_0^tr_s\dd s)$, is a martingale; equivalently, its instantaneous drift under $\Q$ must equal $r_t$ times its price -- exactly the statement that made the stock's drift $\mu$ disappear from the Black--Scholes equation in \cref{ch:bs}. Reading off the $\dd t$-coefficient of \eqref{ch24:eq-dZ} and setting it equal to $r_tZ_t^T$ gives the condition.

\begin{theorem}[HJM drift condition]\label{ch24:thm-hjm}
Under the risk-neutral measure $\Q$, for the bond price \eqref{ch24:eq-bond} to grow on average at the risk-free rate -- $\dd Z_t^T=r_tZ_t^T\,\dd t+(\cdots)\dd B_t$, no $\dd t$-drift beyond $r_t$ -- the drift and volatility of the forward curve \eqref{ch24:eq-hjmsde} must satisfy, for every $t\le T$,
\begin{equation}\label{ch24:eq-hjmdrift}
  \boxed{\ \alpha^{\Q}(t,T)=\sigma(t,T)\int_t^T\sigma(t,s)\,\dd s\ } .
\end{equation}
\end{theorem}
\begin{proof}
Setting the $\dd t$-coefficient in \eqref{ch24:eq-dZ} equal to $r_t$ (after dividing by $Z_t^T\ne0$) gives
\[
  r_t-\int_t^T\alpha^{\Q}(t,s)\,\dd s+\tfrac12\Bigl(\int_t^T\sigma(t,s)\,\dd s\Bigr)^{\!2}=r_t
  \quad\Longleftrightarrow\quad
  \int_t^T\alpha^{\Q}(t,s)\,\dd s=\tfrac12\Bigl(\int_t^T\sigma(t,s)\,\dd s\Bigr)^{\!2}.
\]
Both sides vanish at $T=t$ (empty integral, and the square of an empty integral), so differentiating both sides in $T$ using the fundamental theorem of calculus on the left and the chain rule on the right,
\[
  \alpha^{\Q}(t,T)=2\cdot\tfrac12\Bigl(\int_t^T\sigma(t,s)\,\dd s\Bigr)\cdot\partial_T\Bigl(\int_t^T\sigma(t,s)\,\dd s\Bigr)=\sigma(t,T)\int_t^T\sigma(t,s)\,\dd s ,
\]
which is \eqref{ch24:eq-hjmdrift}.
\end{proof}
This matches the slide's result exactly \ts{13}{24}{1:04:40}\ts{13}{24}{1:05:15}: ``$\alpha^{\Q}_{tT}=\sigma_{tT}\int_t^T\sigma_{ts}\dd s$''. The logic is identical to Black--Scholes: the model started from an \emph{arbitrary} drift $\alpha$ for the forward curve (there was no reason to guess it), and the requirement that hedged, tradable instruments earn the risk-free rate on average removed that freedom entirely. Only $\sigma(t,\cdot)$, the part that can in principle be read off derivative prices (\cref{ch24:rem-calib}), survives as a free input.

\begin{intuition}[Why the drift is an integral of the volatility]
\eqref{ch24:eq-hjmdrift} says the risk-neutral drift at $T$ equals the covariance between the shock at $T$ and the shock to \emph{every rate between $t$ and $T$}, summed up: $\alpha^{\Q}(t,T)=\Cov\bigl(\dd f(t,T),\,\dd X_t^T\bigr)/\dd t$ up to the sign convention. A bond price is a (negative) exponential of the accumulated curve, so its convexity (the $\tfrac12(\cdot)^2$ term from Ito's lemma, exactly the second-order term of \cref{ch18:ex-fail}'s microscope) must be compensated by a corresponding upward push in the curve's own drift for the bond's growth rate to stay pinned at $r_t$. This is the same mechanism as the $-\tfrac12\sigma^2$ correction between the drift of $\ln S_t$ and the drift of $S_t$ in \cref{ch18:ex-gbm}.
\end{intuition}

\begin{exercise}[not from the course]\label{ch24:ex-driftcheck}
Verify \eqref{ch24:eq-hjmdrift} by direct substitution for the Ho--Lee case $\sigma(t,T)\equiv\sigma_0$: show $\int_t^T\alpha^{\Q}(t,s)\dd s=\tfrac12\sigma_0^2(T-t)^2$ two ways, first from $\alpha^{\Q}(t,s)=\sigma_0^2(s-t)$ integrated in $s$, then from $\tfrac12\bigl(\int_t^T\sigma_0\dd s\bigr)^2$, and confirm they agree.
\end{exercise}

\subsection{The real-world drift and Girsanov's theorem}
\Cref{ch24:thm-hjm} pins down the drift only under $\Q$. Under the physical measure $\Prob$, Girsanov's theorem (\cref{ch17:thm-girsanov2}) allows an arbitrary (adapted, technically admissible) market price of risk $\lambda_t$, with $\dd B_t=\dd B_t^{\Q}+\lambda_t\,\dd t$, so that
\begin{equation}\label{ch24:eq-realdrift}
  \alpha^{\Prob}(t,T)=\alpha^{\Q}(t,T)-\sigma(t,T)\,\lambda_t=\sigma(t,T)\int_t^T\sigma(t,s)\,\dd s-\sigma(t,T)\,\lambda_t ,
\end{equation}
exactly as \cref{ch:commodities} shifts the mean-reversion level of a commodity spot process by a constant market price of risk, and as \cref{ch18:fin-vasicek} shifts $\theta$ in Vasicek's model. Since $\lambda_t$ is unobservable from bond prices alone (it cancels out of every $\Q$-priced payoff), pricing derivatives never needs $\alpha^\Prob$: the whole point of risk-neutral valuation, restated once more, is that only $\sigma(t,\cdot)$ and today's curve $f(0,\cdot)$ enter a price.

\begin{finance}[What this buys a trading desk]\label{ch24:fin-recipe}
The slide's practical summary \ts{13}{24}{1:06:30} is:
\begin{enumerate}[1.]
  \item bootstrap $Z(0,\cdot)$, hence $f(0,\cdot)$, from swap quotes (\cref{ch06:sec-curve});
  \item specify (or calibrate, \cref{ch24:rem-calib}) the volatility structure $\sigma(t,T)$;
  \item compute the unique arbitrage-free drift from \eqref{ch24:eq-hjmdrift};
  \item simulate \eqref{ch24:eq-hjmsde} forward under $\Q$ on a grid of $(t,T)$ pairs, exactly as \cref{ch:bs} simulated $S_t$;
  \item for a derivative paying $g$ at $\tau$, price it as $\E_\Q\bigl[e^{-\int_0^\tau r_s\dd s}\,g\bigr]$, averaging over paths and discounting.
\end{enumerate}
This is the same five-step Monte Carlo recipe of \cref{ch:bs}, with ``simulate the stock'' replaced by ``simulate the curve''.
\end{finance}

\begin{remark}[Several factors]\label{ch24:rem-multi}
With $n$ independent Brownian motions $B^{(1)},\dots,B^{(n)}$ and volatilities $\sigma_1(t,T),\dots,\sigma_n(t,T)$, $\dd f(t,T)=\alpha(t,T)\dd t+\sum_{i=1}^n\sigma_i(t,T)\dd B^{(i)}_t$, the same computation (each $\sigma_i$ contributing its own square, by independence, as in \cref{ch18:eq-vec} of \cref{ch:sde}) gives $\alpha^{\Q}(t,T)=\sum_{i=1}^n\sigma_i(t,T)\int_t^T\sigma_i(t,s)\,\dd s$. Multiple factors (often chosen as the level/slope/curvature loadings of \cref{ch:pca}) let the curve twist in ways a single factor cannot; this is not discussed further in the lecture and not needed for the one-factor numerics of \cref{ch24:sec-numerics}.
\end{remark}

% ==================================================================
\section{Short-rate models as special cases}\label{ch24:sec-shortrate}

The 2013 lecture lists the classic short-rate SDEs \emph{before} deriving the HJM condition, without connecting the two \ts{13}{24}{1:02:13}:
\begin{equation}\label{ch24:eq-shortrates}
\begin{gathered}
  \dd r_t=\theta(t)\,\dd t+\sigma\,\dd B_t\ \text{(Ho--Lee)},\qquad
  \dd r_t=\kappa(\theta(t)-r_t)\,\dd t+\sigma\,\dd B_t\ \text{(Hull--White/Vasicek)},\\
  \dd r_t=\kappa(\theta-r_t)\,\dd t+\sigma\sqrt{r_t}\,\dd B_t\ \text{(CIR)}.
\end{gathered}
\end{equation}
The connection is the content of this section (an addition of these notes): a short-rate model is \emph{equivalent} to the HJM framework precisely when the chosen $\sigma(t,T)$ makes the whole forward curve a deterministic function of $r_t$ alone, so that simulating $r_t$ (a single number) is enough to reconstruct $f(t,T)$ for every $T$ -- the entire motivation, in \cref{ch24:sec-setup}, for needing a curve model.

\begin{proposition}[Ho--Lee is HJM with constant volatility]\label{ch24:prop-holee}
Take $\sigma(t,T)\equiv\sigma_0$ (constant, independent of $t$ and $T$) in \eqref{ch24:eq-hjmsde}. Then \eqref{ch24:eq-hjmdrift} gives $\alpha^{\Q}(t,T)=\sigma_0^2(T-t)$, and
\begin{equation}\label{ch24:eq-holeefwd}
  f(t,T)=f(0,T)+\sigma_0^2t\Bigl(T-\tfrac t2\Bigr)+\sigma_0B_t ,\qquad r_t=f(t,t)=f(0,t)+\tfrac12\sigma_0^2t^2+\sigma_0B_t ,
\end{equation}
which satisfies the Ho--Lee SDE $\dd r_t=\theta(t)\dd t+\sigma_0\dd B_t$ with $\theta(t)=\partial_tf(0,t)+\sigma_0^2t$ (the curve's own slope plus a convexity correction).
\end{proposition}
\begin{proof}
Integrate \eqref{ch24:eq-hjmsde} in $t$ at fixed $T$: $f(t,T)=f(0,T)+\int_0^t\sigma_0^2(T-u)\,\dd u+\sigma_0B_t=f(0,T)+\sigma_0^2t(T-t/2)+\sigma_0B_t$, which is \eqref{ch24:eq-holeefwd}; set $T=t$ for $r_t$. Differentiating $r_t$ in $t$ (Ito, with $B$ contributing $\dd B_t$ and the deterministic part contributing its derivative) gives exactly $\dd r_t=\bigl[\partial_tf(0,t)+\sigma_0^2t\bigr]\dd t+\sigma_0\,\dd B_t$.
\end{proof}
Ho--Lee is the simplest possible HJM model and the one the slide's zero-coupon bond computation implicitly matches when $\sigma$ is taken constant \ts{13}{24}{1:03:00} (the algebra there is written for general $\sigma(t,T)$, but the arithmetic simplifies to Ho--Lee's classic $r_t$-normal, curve-fitting form when $\sigma$ is constant).

\begin{proposition}[Vasicek/Hull--White is HJM with exponential volatility]\label{ch24:prop-vasicek}
Take $\sigma(t,T)=\sigma_0e^{-a(T-t)}$ for constants $\sigma_0,a>0$ (mean-reversion speed $a$). Then $\int_t^T\sigma(t,s)\dd s=\tfrac{\sigma_0}{a}\bigl(1-e^{-a(T-t)}\bigr)$ and \eqref{ch24:eq-hjmdrift} gives $\alpha^{\Q}(t,T)=\sigma(t,T)\cdot\tfrac{\sigma_0}{a}(1-e^{-a(T-t)})$. The short rate $r_t=f(t,t)$ then solves the Hull--White SDE
\begin{equation}\label{ch24:eq-hw}
  \dd r_t=\bigl[\partial_Tf(0,t)+a\bigl(f(0,t)-r_t\bigr)+\tfrac{\sigma_0^2}{2a}\bigl(1-e^{-2at}\bigr)\bigr]\dd t+\sigma_0\,\dd B_t ,
\end{equation}
which is \eqref{ch18:eq-ou} (\cref{ch18:sec-ou}) with a time-dependent, curve-fitting mean level; for a flat initial curve $f(0,T)\equiv r_0$ it approaches, for $t\gg1/a$, the stationary Vasicek SDE $\dd r_t=a(\theta_\infty-r_t)\dd t+\sigma_0\dd B_t$ with $\theta_\infty=r_0+\sigma_0^2/(2a^2)$.
\end{proposition}
\begin{proof}
The volatility factorises as $\sigma(t,T)=e^{-a(T-t)}\sigma_0$, of the separable form $\sigma(t,T)=g(T)/g(t)\cdot\sigma_0$ with $g(u)=e^{au}$: this is exactly the structure that made the Ornstein--Uhlenbeck integrating-factor trick of \cref{ch18:sec-ou} work, and it is why exponential HJM volatility is the one choice that keeps the short rate Markovian (reducible to $r_t$ alone) among mean-reverting models. Writing $f(t,T)=f(0,T)+\int_0^t\alpha^\Q(u,T)\dd u+\sigma_0\int_0^te^{-a(T-u)}\dd B_u$ and specialising to $T=t$ gives, after differentiating in $t$ (details: differentiate $r_t=f(t,t)$ exactly as in the proof of \cref{ch24:lem-dX}, tracking both the moving argument and the stochastic integral's upper limit), the stated SDE; the flat-curve case matches \cref{ch18:fin-vasicek} with $\kappa=a$ after evaluating the $t\to\infty$ limit of the bracket, $a\theta_\infty=\partial_Tf(0,\cdot)|_{\text{flat}}+ar_0+\sigma_0^2/(2a)=ar_0+\sigma_0^2/(2a)$.
\end{proof}

\begin{coursenote}[CIR is not an exact HJM special case]\label{ch24:cn-cir}
The slides list CIR alongside Ho--Lee and Hull--White with no further comment \ts{13}{24}{1:02:13}, but CIR does \emph{not} embed into one-factor HJM the way the other two do: HJM's volatility $\sigma(t,T)$ in \eqref{ch24:eq-hjmsde} is a function of $(t,T)$ only, never of the (random) curve level, whereas CIR's $\sigma\sqrt{r_t}$ (\cref{ch18:sec-cir}) is level-dependent. A level-dependent-volatility HJM model exists (it is the starting point of ``Markov-functional'' and ``quadratic Gaussian'' short-rate models, not discussed here) but the clean closed-form drift \eqref{ch24:eq-hjmdrift}, which relied only on $T$-integrals of a curve known at time $t$, no longer separates from the state in the same way. CIR is best read as a short-rate model in its own right (\cref{ch18:sec-cir}), historically motivated independently of HJM by the wish for a nonnegative rate, rather than as an HJM special case.
\end{coursenote}

\begin{remark}[Calibrating $\sigma(t,T)$]\label{ch24:rem-calib}
As in the equity case (\cref{ch:bs}), the volatility $\sigma$ cannot be chosen freely and then forgotten: it is implied from the market prices of liquid options, here swaptions and caps/floors \ts{13}{24}{1:06:56}. An exponential $\sigma(t,T)=\sigma_0e^{-a(T-t)}$ has only two free parameters and typically cannot fit the whole swaption grid (all expiries and tenors) simultaneously; richer parametrisations, or several factors (\cref{ch24:rem-multi}), trade this off against tractability, exactly the local-versus-global tension of \cref{ch06:sec-curve}'s interpolation choice.
\end{remark}

\begin{exercise}[not from the course]\label{ch24:ex-holeeproof}
Complete the differentiation step skipped in the proof of \cref{ch24:prop-holee}: starting from \eqref{ch24:eq-holeefwd}, apply Ito's formula to $r_t=f(0,t)+\tfrac12\sigma_0^2t^2+\sigma_0B_t$ directly (there is no hidden randomness in the deterministic part) and confirm the stated Ho--Lee SDE, including the exact form of $\theta(t)$ in terms of the initial curve.
\end{exercise}

\begin{exercise}[not from the course]\label{ch24:ex-hw}
Fill in the differentiation step of \cref{ch24:prop-vasicek}: write $r_t=f(0,t)+\int_0^t\alpha^{\Q}(u,t)\dd u+\sigma_0\int_0^te^{-a(t-u)}\dd B_u$ (the $T=t$ specialisation, tracking the moving upper limit of both integrals as in \cref{ch24:lem-dX}), differentiate in $t$, and verify \eqref{ch24:eq-hw}.
\end{exercise}

% ==================================================================
\section{Numerical check of the drift condition}\label{ch24:sec-numerics}

The cleanest test of \cref{ch24:thm-hjm} is direct: simulate the full curve under the computed drift, and check that the resulting bond price matches the one implied by today's curve -- \emph{without ever solving a bond-pricing PDE}, purely by discounting simulated paths, exactly the logic of the Monte Carlo check of the Black--Scholes formula suggested on the slides \ts{13}{24}{38:48} (own computation; not in the lecture).

\begin{example}[Simulating the HJM curve and pricing a bond by Monte Carlo]\label{ch24:ex-mc}
Take the exponential-volatility model of \cref{ch24:prop-vasicek} with $\sigma_0=1.2\%$, $a=0.5$, a flat initial curve $f(0,T)\equiv3\%$, and simulate $f(t,T)$ on a daily grid ($\dd t=1/252$) up to $t=T_h=2$ years, using the closed-form drift $\alpha^{\Q}(t,T)=\sigma(t,T)\cdot\tfrac{\sigma_0}{a}(1-e^{-a(T-t)})$ from \cref{ch24:prop-vasicek} and an Euler--Maruyama step (\cref{ch18:sec-num}) shared by every remaining maturity on the grid at each time step. Over $20{,}000$ paths, the Monte Carlo bond price $\hat Z(0,T_h)=\widehat{\E}_\Q\bigl[e^{-\int_0^{T_h}r_s\dd s}\bigr]$ and its analytic value $Z(0,T_h)=e^{-f_0T_h}$ (flat curve) agree to within Monte Carlo error:
\[
  Z(0,2)_{\text{theory}}=e^{-0.03\times2}=0.941765,\qquad \hat Z(0,2)_{\text{MC}}=0.941796\pm0.000093 .
\]
As a second, independent check, \cref{ch24:prop-vasicek} predicts that the short rate at $t=1$ is Gaussian with variance $\tfrac{\sigma_0^2}{2a}(1-e^{-2a})=9.10\times10^{-5}$; the simulated cross-sectional variance of $r_1$ across the same $20{,}000$ paths is $9.07\times10^{-5}$.
\end{example}

\begin{figure}[ht]
\centering
\begin{tikzpicture}
\begin{axis}[width=0.95\linewidth,height=6cm,xlabel={maturity $T$ (years)},ylabel={forward rate $f(1,T)$ (\%)},
  xmin=1,xmax=5,ymin=0.8,ymax=3.3,grid=major,grid style={black!10},
  legend style={at={(0.5,1.05)},anchor=south,legend columns=2,font=\footnotesize}]
  \addplot[thick,dashed,color=black!55] coordinates {(1.0,3.000) (1.5,3.000) (2.0,3.000) (2.5,3.000) (3.0,3.000) (3.5,3.000) (4.0,3.000) (4.5,3.000) (5.0,3.000)};
  \addplot[thick,color=defcol] coordinates {(1.0,3.097) (1.5,3.079) (2.0,3.063) (2.5,3.051) (3.0,3.040) (3.5,3.032) (4.0,3.025) (4.5,3.020) (5.0,3.015)};
  \addplot[thick,color=thmcol] coordinates {(1.0,1.350) (1.5,1.718) (2.0,2.004) (2.5,2.225) (3.0,2.397) (3.5,2.531) (4.0,2.635) (4.5,2.716) (5.0,2.779)};
  \addplot[thick,color=intcol] coordinates {(1.0,3.084) (1.5,3.068) (2.0,3.055) (2.5,3.044) (3.0,3.035) (3.5,3.028) (4.0,3.022) (4.5,3.017) (5.0,3.013)};
  \addplot[thick,color=fincol] coordinates {(1.0,1.152) (1.5,1.564) (2.0,1.883) (2.5,2.131) (3.0,2.324) (3.5,2.474) (4.0,2.591) (4.5,2.681) (5.0,2.752)};
  \legend{{$f(0,T)$ (initial)},{path A},{path B},{path C},{path D}}
\end{axis}
\end{tikzpicture}
\caption{Four simulated forward curves $T\mapsto f(1,T)$ one year into the simulation of \cref{ch24:ex-mc} ($\sigma_0=1.2\%$, $a=0.5$, flat initial curve at 3\%), against the initial curve $f(0,T)$. All curves are driven by the same Brownian path within a curve but differ across paths; the exponential vol makes the short end move most (as $\sigma(t,T)\to\sigma_0$ when $T\to t$) while the long end is comparatively anchored (as $\sigma(t,T)\to0$ when $T-t\to\infty$), the familiar level/slope pattern of \cref{ch:pca}.}\label{ch24:fig-curves}
\end{figure}

\begin{exercise}[not from the course]\label{ch24:ex-mccode}
Reproduce the Monte Carlo check of \cref{ch24:ex-mc} in your own code: simulate the exponential-volatility HJM curve on a daily grid, compute $\hat Z(0,T_h)$ for $T_h=1,2,5$ years, and compare with $e^{-f_0T_h}$. Then break the drift condition on purpose (set $\alpha\equiv0$, ``drift-free'' forward rates) and show that the simulated bond price no longer matches $Z(0,T_h)$: absence of arbitrage is not automatic, it is exactly what \eqref{ch24:eq-hjmdrift} enforces.
\end{exercise}

% ==================================================================
\section{Extension to credit: hazard rates}\label{ch24:sec-credit}

\subsection{Survival probabilities as a second discount curve}
\Cref{ch08:prop-numeraire} priced default-free cash flows with $Z(t,T)$. A defaultable bond needs a second curve: the market-implied \emph{survival probability} $p(t,T)$, the risk-neutral probability, seen at $t$, that a reference entity has not defaulted by $T$ \ts{13}{24}{1:12:45}\ts{13}{24}{1:13:49}. Exactly as $Z(t,T)$ was parametrised by instantaneous forward rates, $p(t,T)$ is parametrised by instantaneous \emph{forward hazard rates} $h(t,T)$:
\begin{equation}\label{ch24:eq-survival}
  p(t,T)=\exp\Bigl(-\int_t^Th(t,s)\,\dd s\Bigr),\qquad p(t,t)=1,
\end{equation}
the exact analogue of \eqref{ch24:eq-bond}. A zero-recovery defaultable zero-coupon bond -- paying \$1 at $T$ if and only if the reference entity has not defaulted -- is priced (assuming, as the slides do, independence between interest rates and default for this formula) by
\begin{equation}\label{ch24:eq-defbond}
  \tilde Z(t,T)=Z(t,T)\,p(t,T)=\exp\Bigl(-\int_t^T\bigl[f(t,s)+h(t,s)\bigr]\dd s\Bigr) ,
\end{equation}
\ts{13}{24}{1:18:01} shown on the slide as the running example: $\tilde Z_t^T=\ind_{\{\tau>t\}}\exp\bigl(-\int_t^T[f(t,s)+h(t,s)]\dd s\bigr)$, with $\tau$ the (random) default time and the indicator recording that the bond is worthless once default has occurred \ts{13}{24}{1:18:33}.

\subsection{The HJM SDE and drift condition for hazard rates}
Postulate, exactly as in \eqref{ch24:eq-hjmsde}, an Ito process for the hazard curve,
\begin{equation}\label{ch24:eq-hjmcredit}
  \dd h(t,T)=\alpha^h(t,T)\,\dd t+\sigma^h(t,T)\,\dd B^h_t ,
\end{equation}
\ts{13}{24}{1:19:01} where $B^h$ may be correlated with the rates driver $B$ of \eqref{ch24:eq-hjmsde}, $\dd B_t\,\dd B^h_t=\rho\,\dd t$ for a constant correlation $\rho$ (the general case; the slide's exercise asks only for $\rho=0$ \ts{13}{24}{1:19:26}\ts{13}{24}{1:19:46}). The slide poses, but does not answer, the derivation of the drift condition in this case \ts{13}{24}{1:19:26}: ``\emph{Exercise 3: Derive a drift expression for the credit HJM model.}'' The derivation is done in full here, for general $\rho$, and specialises to the slide's hinted answer when $\rho=0$.

\begin{theorem}[HJM drift condition with credit]\label{ch24:thm-hjmcredit}
Under $\Q$, requiring that the pre-default, zero-recovery defaultable bond price $\tilde Z_t^T=\ind_{\{\tau>t\}}\,e^{-Y_t^T}$, $Y_t^T=\int_t^T[f(t,s)+h(t,s)]\dd s$, grows on average at the risk-free rate $r_t$ \emph{plus} the hazard rate ($\ind_{\{\tau>t\}}$ jumps to $0$ at rate $h_t=h(t,t)$, so its compensated growth rate on the event $\{\tau>t\}$ must be $r_t+h_t$ for the discounted, default-compensated price to be a $\Q$-martingale) forces
\begin{equation}\label{ch24:eq-hjmcreditdrift}
  \begin{split}\alpha^h(t,T)={}&\sigma^h(t,T)\int_t^T\sigma^h(t,s)\,\dd s\\&+\rho\Bigl[\sigma^h(t,T)\int_t^T\sigma(t,s)\,\dd s+\sigma(t,T)\int_t^T\sigma^h(t,s)\,\dd s\Bigr].\end{split}
\end{equation}
When $\rho=0$ this is $\alpha^h(t,T)=\sigma^h(t,T)\int_t^T\sigma^h(t,s)\dd s$, matching the slide's hint \ts{13}{24}{1:19:46}: ``\emph{If there are no correlations between IR and hazard rates, the drift is given by the equation similar to the IR case.}''
\end{theorem}
\begin{proof}
Repeat \cref{ch24:lem-dX} verbatim for $Y_t^T$ with $f+h$ in place of $f$ and $(\sigma,\alpha)+(\sigma^h,\alpha^h)$ jointly driving it (the two curves add, so their $\dd t$- and $\dd B$-coefficients add):
\[
\begin{gathered}
  \dd Y_t^T=-\bigl(r_t+h_t\bigr)\dd t+\Bigl(\int_t^T[\alpha(t,s)+\alpha^h(t,s)]\dd s\Bigr)\dd t\\
  {}+\Bigl(\int_t^T\sigma(t,s)\dd s\Bigr)\dd B_t+\Bigl(\int_t^T\sigma^h(t,s)\dd s\Bigr)\dd B^h_t .
\end{gathered}
\]
On $\{\tau>t\}$, $\tilde Z_t^T=e^{-Y_t^T}$ follows Ito's formula for a function of two correlated Brownian drivers (\cref{ch17:sec-multi}):
\[
\begin{gathered}
  \dd\bigl(e^{-Y_t^T}\bigr)=e^{-Y_t^T}\Bigl[-\dd Y_t^T+\tfrac12\,\dd\langle Y^T\rangle_t\Bigr],\\
  \dd\langle Y^T\rangle_t=\Bigl(\int_t^T\!\sigma\dd s\Bigr)^{\!2}\!\dd t+\Bigl(\int_t^T\!\sigma^h\dd s\Bigr)^{\!2}\!\dd t+2\rho\Bigl(\int_t^T\!\sigma\dd s\Bigr)\Bigl(\int_t^T\!\sigma^h\dd s\Bigr)\dd t ,
\end{gathered}
\]
using $\dd\langle B,B^h\rangle_t=\rho\,\dd t$. The jump-to-default piece contributes an extra $-\tilde Z_{t^-}^T\,\dd N_t$ at the default event (with $\dd N_t$ the default indicator's jump, compensated intensity $h_t\dd t$ under $\Q$); requiring the compensated process to be a $\Q$-martingale removes the compensator $h_t\dd t$ exactly against the $+h_t\dd t$ term already present in $-\dd Y_t^T$ (a standard credit-modelling fact, quoted, consistent with pricing formula \eqref{ch24:eq-defbond}). What remains, on $\{\tau>t\}$ and after this cancellation, is the same $\dd t$-matching argument as in \cref{ch24:thm-hjm} applied to $\int_t^T[\alpha+\alpha^h]$ jointly: since the interest-rate piece already satisfies $\int_t^T\alpha=\tfrac12(\int_t^T\sigma)^2$ on its own (rates alone must also be arbitrage-free), what is left is
$\int_t^T\alpha^h(t,s)\dd s=\tfrac12\bigl(\int_t^T\sigma^h\bigr)^2+\rho\bigl(\int_t^T\sigma\bigr)\bigl(\int_t^T\sigma^h\bigr)$, and differentiating in $T$ (product rule on the cross term, which gives two pieces) yields \eqref{ch24:eq-hjmcreditdrift}. This is the result of Sch\"onbucher (1998).
\end{proof}
\begin{intuition}[Why a cross-term appears]
When rates and hazard rates are correlated, a shock that raises $\sigma(t,s)$-driven forward rates on $[t,T]$ also moves (through $\rho$) the hazard curve on the same interval; the defaultable bond price is convex in the \emph{sum} $f+h$, so its quadratic-variation compensation, \eqref{ch24:eq-hjmcreditdrift}'s cross-term, mixes the two volatility kernels exactly as a two-asset Ito's lemma mixes $\sigma_1\sigma_2\rho$ into the drift of a product (\cref{ch17:sec-multi}). Setting $\rho=0$, as the slide does, isolates the credit-only piece and reproduces the interest-rate result \eqref{ch24:eq-hjmdrift} verbatim with $h$ in place of $f$.
\end{intuition}

\begin{exercise}[not from the course]\label{ch24:ex-credit0}
Specialise \eqref{ch24:eq-hjmcreditdrift} to $\rho=0$ and $\sigma^h(t,T)\equiv\sigma_0^h$ constant (a ``Ho--Lee'' hazard-rate model). Derive the resulting SDE for the short hazard rate $h_t=h(t,t)$ and compare its form with \cref{ch24:prop-holee}.
\end{exercise}

\begin{exercise}[not from the course]\label{ch24:ex-corr}
Redo \cref{ch24:ex-mc} with a second, correlated curve $h(t,T)=\sigma_0^he^{-b(T-t)}$-volatility hazard process ($\rho=0.5$ correlation with the rates driver) and the drift \eqref{ch24:eq-hjmcreditdrift}. Check numerically that $\E_\Q[\tilde Z(0,T_h)]$ from simulation matches $Z(0,T_h)\,p(0,T_h)$ from the input curves, as in \eqref{ch24:eq-defbond}.
\end{exercise}

\subsection{Why this matters: callable and structured credit}
The slide motivates the credit HJM model with corporate \emph{callable bonds} \ts{13}{24}{1:14:51}\ts{13}{24}{1:16:56}: an issuer who can refinance at a lower all-in rate (risk-free plus credit spread) saves money by calling old, expensive debt, and pricing that option requires simulating \emph{both} $f(t,\cdot)$ and $h(t,\cdot)$ jointly, exactly the setting \eqref{ch24:eq-hjmcredit}--\eqref{ch24:eq-hjmcreditdrift} was built for.
\begin{example}[Refinancing savings; slide numbers \ts{13}{24}{1:14:51}]
A corporation issues a 10-year \$100M bond at an effective (risk-free $+$ credit) rate of $2.00\%+3.00\%=5.00\%$. Three years later, market rates have fallen to $1.50\%+1.50\%=3.00\%$ with $7$ years remaining. Calling and reissuing saves $(5.00\%-3.00\%)\times\$100\text{M}\times7\text{ years}=\$14\text{M}$ (undiscounted; a real valuation discounts each year's saving and nets the reissuance cost, not shown on the slide).
\end{example}
\begin{casestudy}[Structured notes: coupon enhancement via joint simulation]
The slide's other running example \ts{13}{24}{1:20:05}\ts{13}{24}{1:23:15} is a 10-year, \$100M structured note paying a $10\%$ coupon only while \emph{both} the 30-year swap rate exceeds the 2-year swap rate \emph{and} the S\&P 500 stays above $880$: a joint condition on the rates curve and an equity index (credit enters through the issuer's own default risk on the note, priced with $h(t,\cdot)$ as above). Historically the 30y--2y spread was negative on only a few days in 2005--2006, and the S\&P 500 was above $880$ on $94\%$ of trading days over the prior decade, yet the market-implied (risk-neutral) probabilities were only about $80\%$ and $75\%$ respectively \ts{13}{24}{1:28:28}\ts{13}{24}{1:30:30}: investors are compensated (a higher coupon) for bearing tail risk that is more richly priced than its historical frequency, and the bank prices and hedges the note with exactly the Monte Carlo machinery of \cref{ch24:fin-recipe}, extended to simulate equity, rates and credit together.
\end{casestudy}

% ==================================================================
\begin{keyformulas}
\textbf{Curve and bond.} $Z(t,T)=\exp\bigl(-\int_t^Tf(t,s)\dd s\bigr)$, short rate $r_t=f(t,t)$.\\[2pt]
\textbf{HJM forward-rate SDE.} $\dd f(t,T)=\alpha(t,T)\dd t+\sigma(t,T)\dd B_t$.\\[2pt]
\textbf{HJM drift condition (no-arbitrage, one factor, under $\Q$).}
$\displaystyle\alpha^{\Q}(t,T)=\sigma(t,T)\int_t^T\sigma(t,s)\,\dd s$; with $n$ factors, sum this over $i=1,\dots,n$.\\[2pt]
\textbf{Real-world drift.} $\alpha^{\Prob}(t,T)=\alpha^{\Q}(t,T)-\sigma(t,T)\lambda_t$ (Girsanov, market price of risk $\lambda_t$).\\[2pt]
\textbf{Short-rate special cases.} $\sigma(t,T)\equiv\sigma_0\Rightarrow$ Ho--Lee; $\sigma(t,T)=\sigma_0e^{-a(T-t)}\Rightarrow$ Hull--White/Vasicek; CIR's $\sigma\sqrt{r_t}$ is level-dependent and does not embed in one-factor HJM.\\[2pt]
\textbf{Credit.} $p(t,T)=\exp\bigl(-\int_t^Th(t,s)\dd s\bigr)$, $\dd h(t,T)=\alpha^h(t,T)\dd t+\sigma^h(t,T)\dd B^h_t$,
$\displaystyle\alpha^h(t,T)=\sigma^h(t,T)\int_t^T\sigma^h(t,s)\,\dd s+\rho\Bigl[\sigma^h(t,T)\int_t^T\sigma(t,s)\,\dd s+\sigma(t,T)\int_t^T\sigma^h(t,s)\,\dd s\Bigr]$ ($\rho=\Corr(B,B^h)$; $\rho=0$ recovers the rates-only form with $h$ in place of $f$). Defaultable zero-recovery bond: $\tilde Z(t,T)=Z(t,T)p(t,T)$.
\end{keyformulas}

% !TEX root = ../main.tex
\chapter{The Ross Recovery Theorem}\label{ch:ross}

\begin{sources}
\small
\begin{tabularx}{\linewidth}{@{}l l L@{}}
\toprule
\textbf{Lecture} & \textbf{Speaker} & \textbf{Content}\\
\midrule
2013 L25 & P.~Carr (Morgan Stanley), with J.~Yu & Guest lecture ``Can We Recover?'': Arrow--Debreu
security prices and market beliefs; Ross's Recovery Theorem for finite-state Markov chains via the
Perron--Frobenius theorem; a preference-free re-derivation using a change of numeraire (Long's numeraire
portfolio) for a time-homogeneous diffusion on a bounded state space; successes and failures for
unbounded state spaces (Black--Scholes fails, Cox--Ingersoll--Ross succeeds).\\
\bottomrule
\end{tabularx}

\smallskip
\textbf{Files.} \file{lecnote25} (50 slides, S.~Ross's collaborators P.~Carr and J.~Yu, Morgan Stanley).
There is no 2024 counterpart: the topic is not in the 2024 syllabus.
\textbf{Problem sets.} None (guest lecture, no associated problem-set questions in either year).
\textbf{Not covered / not published.} The slides' Part~I (conditioning and path-dependent
Arrow--Debreu prices, multi-period Markovian consistency) is background that the lecture itself
skipped over quickly (``I'm going to skip these slides''); it is summarised briefly below for
completeness but not derived in full. Ross's original paper (\emph{The Recovery Theorem}, published in
the \emph{Journal of Finance} in 2015, then only forthcoming) is not itself a course file; statements
attributed to it are as quoted or paraphrased on the slides and in the lecture.
\end{sources}

\section*{Motivation}
Every derivative price already encodes a probability measure: the \emph{risk-neutral} measure $\Q$ of
\cref{sec:oneperiod}, under which discounted prices are martingales. But $\Q$ is famously not the market's
actual view of the future --- it is contaminated by risk aversion, so that, say, the risk-neutral chance of
a market crash is typically far larger than anyone's honest assessment of its physical probability. For
decades the going wisdom was that untangling the two required an extra ingredient that no one could agree
on: a representative investor's risk-aversion parameter, supplied from outside the data. Ross's 2011
insight was that a single structural assumption on how risk aversion enters --- \emph{transition
independence} of the pricing kernel --- is enough to pin down the real-world probabilities uniquely, with
no free parameter to calibrate, provided the underlying is a finite-state, time-homogeneous Markov chain
and the matrix of state prices is positive. The mechanism is one this course has already built:
\cref{sec:pf}'s Perron--Frobenius theorem, applied to the matrix of Arrow--Debreu prices instead of a
Markov transition matrix. This chapter states the assumptions precisely, proves the recovery theorem in
full using \cref{thm:pf}, verifies the construction on an explicit numerical example, and closes Part~II of
these notes with a look at what happens (and what is still open) beyond the finite-state setting.

% ==================================================================
\section{Three probability measures: \texorpdfstring{$P$, $\Q$, $R$}{P, Q, R}}\label{ch25:sec-pqr}

Fix a single source of uncertainty $X$ (an index level, a short rate, an exchange rate) and consider
Arrow--Debreu (AD) securities written on it: a claim paying \$1 if $X$ lands in a given state and nothing
otherwise, i.e.\ exactly the state prices $\psi_i$ of \cref{def:onepmodel}, now indexed by both the
\emph{current} state and the terminal state \ts{13}{25}{12:49}. Three probability measures matter here
\ts{13}{25}{5:22}:
\begin{itemize}
  \item $P$, the \emph{physical} (objective, ``God's-eye'') probability measure: the true frequency with
        which $X$ visits each state. No one observes $P$ directly.
  \item $\Q$, the \emph{risk-neutral} measure of \cref{thm:ftap}: it exists under no arbitrage and
        satisfies $A_0=\beta\,\E_\Q[\bm C_T]$ for every replicable claim, but it is not a measure of frequency
        --- it is a repackaging of AD prices, normalised to integrate to one, and under it every asset grows
        at the risk-free rate \ts{13}{25}{9:37}.
  \item $R$, the \emph{recovered} (or \emph{natural}) measure: what one gets by trying to strip $\Q$ of its
        risk-aversion and time-value-of-money content and recover a measure of actual frequencies. Ross
        calls this the market's ``natural probabilities''. $R$ need not equal $P$ --- the market can be
        wrong (2005 mortgage-backed pricing being Carr's stock example) --- so $R$ is honestly only the
        market's \emph{belief} about $P$; a believer in market efficiency is free to set $R=P$
        \ts{13}{25}{10:40}.
\end{itemize}

\begin{intuition}[Why $\Q$ is not a frequency]
A risk-neutral probability of $40\%$ for ``the index is up'' means only that a security paying \$1 in that
event costs \$0.40 today, i.e.\ it is a \emph{price}, not a count of outcomes. Two effects separate this
price from the true chance of the event: discounting (a sure \$1 tomorrow is worth less than \$1 today,
even with no uncertainty at all) and risk aversion (a payoff arriving exactly when the rest of one's wealth
is doing badly is worth more per dollar than one arriving in good times, so its price is inflated relative
to its frequency). Recovering $R$ from $\Q$ means undoing both effects \ts{13}{25}{42:38}.
\end{intuition}

% ==================================================================
\section{The pricing matrix and the pricing kernel}\label{ch25:sec-kernel}

Let $X$ be a Markov chain on a finite state space $\{1,\dots,n\}$, and write $A(x,y)>0$ for the price, at
state $x$ today, of an AD security paying \$1 if $X$ is at state $y$ one period later \ts{13}{25}{40:28}.
Collecting these into a matrix $A\in\R^{n\times n}$, no-arbitrage consistency across horizons forces $X$'s
one-step transition structure under $A$ to be genuinely Markovian and, if desired, time-homogeneous, i.e.\
$A$ does not depend on calendar time (\emph{Appendix, Part~I of the slides}, sketched only: with three
dates $0,1,2$, path-dependent AD prices $A_{jk|i}$ factor as $A_{jk|i}=A_{ji}\cdot A_{k|ji}$, and Markov
means the one-step-ahead conditional price $A_{k|ji}=A_{k|j}$ does not depend on the earlier state $i$
\ts{13}{25}{39:23}).

\begin{coursenote}[Notation clash with \cref{ch:linalg}]
The letter $A$ here denotes the $n\times n$ matrix of \emph{one-step Arrow--Debreu prices between states of
a single Markov chain}, following Ross's and Carr's notation; it is unrelated to the $m\times n$ payoff
matrix $A$ of \cref{def:onepmodel} (assets in columns, states in rows). Both are conventional in their
respective sources and are kept as in the lecture material.
\end{coursenote}

Write $p(x,y)$ for the (unknown) \emph{natural} one-step transition probabilities under $R$: $p(x,y)\ge0$,
$\sum_yp(x,y)=1$ for every $x$. The \emph{pricing kernel} is the entrywise ratio
\begin{equation}\label{ch25:eq-kernel}
  \varphi(x,y) \;=\; \frac{A(x,y)}{p(x,y)},
\end{equation}
capturing exactly the two contaminating effects: time value of money and risk aversion
\ts{13}{25}{47:03}. If interest rates were zero and investors risk-neutral, $\varphi\equiv1$ and $A=p$;
in general $\varphi$ is an unknown positive function of two variables, and the $n^2$ equations \eqref{ch25:eq-kernel} involve
$2n^2-n$ unknowns (the $n^2$ values of $\varphi$ and the $n^2-n$ free entries of $p$, whose rows sum to one),
far too many to pin down $p$ --- this is why decontamination was thought hopeless before Ross
\ts{13}{25}{47:03}\ts{13}{25}{58:40}.

\begin{definition}[Transition independence, Ross's assumption]\label{ch25:def-ti}
The pricing kernel is \emph{transition-independent} if it factors as
\begin{equation}\label{ch25:eq-ti}
  \varphi(x,y) \;=\; \delta\,\frac{z(y)}{z(x)}, \qquad x,y=1,\dots,n,
\end{equation}
for some scalar $\delta>0$ and some function $z:\{1,\dots,n\}\to(0,\infty)$.
\end{definition}

In a world with a representative agent and additively time-separable utility, $z(x)=U'(c(x))$ (marginal
utility of consumption at state $x$) and $\delta$ is the subjective discount factor, so that
$\varphi(x,y)=\delta\,U'(c(y))/U'(c(x))$: risk aversion enters only through how marginal utility differs
\emph{between} the two states, not through the path or timing otherwise \ts{13}{25}{49:10}. Carr's
restatement, avoiding the representative-agent language, is purely dimension-counting
(\cref{ch25:rem-dof}); either reading is the same assumption \eqref{ch25:eq-ti}.

\begin{remark}[Degrees of freedom]\label{ch25:rem-dof}
Without \eqref{ch25:eq-ti}, $\varphi$ has $n^2$ free values. Assumption~\eqref{ch25:eq-ti} reduces this to
$n+1$: the $n$ values $z(1),\dots,z(n)$ (only their ratios matter, so really $n-1$ free numbers) plus
$\delta$. For $n=10$ this is a drop from $100$ to $11$ unknowns, which is what makes identification from
$A$ alone possible \ts{13}{25}{59:42}.
\end{remark}

% ==================================================================
\section{The Recovery Theorem}\label{ch25:sec-thm}

\begin{theorem}[Ross's Recovery Theorem]\label{ch25:thm-ross}
Let $A\in\R^{n\times n}$ be a matrix of one-step AD prices for a time-homogeneous Markov chain, with all
entries $A(x,y)>0$. Under transition independence \eqref{ch25:eq-ti}, there is a unique triple $(\delta,z,p)$
with $\delta>0$, $z\in\R^n$, $z>0$ (up to an overall positive scale factor, which cancels out of $p$ and
$\varphi$), and $p$ a row-stochastic matrix with $p(x,y)>0$, consistent with \eqref{ch25:eq-kernel}--%
\eqref{ch25:eq-ti} and the given $A$. Explicitly, $\delta$ is the Perron--Frobenius eigenvalue of $A$,
$w(x)=1/z(x)$ is (up to scale) its positive eigenvector, and
\begin{equation}\label{ch25:eq-precover}
  p(x,y) \;=\; \frac{1}{\delta}\,A(x,y)\,\frac{w(y)}{w(x)}, \qquad x,y=1,\dots,n.
\end{equation}
\end{theorem}

\begin{proof}
Write $D_z=\diag(z(1),\dots,z(n))$. Assumption \eqref{ch25:eq-ti} in \eqref{ch25:eq-kernel} reads
$A(x,y)=\delta\,p(x,y)\,z(y)/z(x)$, i.e.\ $z(x)A(x,y)=\delta\,p(x,y)z(y)$, or in matrix form
\begin{equation}\label{ch25:eq-matrixform}
  D_zA \;=\; \delta\,p\,D_z .
\end{equation}
Let $w=D_z^{-1}\bm1=(1/z(1),\dots,1/z(n))\T$, where $\bm1=(1,\dots,1)\T$. Multiply \eqref{ch25:eq-matrixform}
on the right by $w=D_z^{-1}\bm1$:
\[
  D_zAw \;=\; \delta\,p\,D_zD_z^{-1}\bm1 \;=\; \delta\,p\,\bm1 \;=\; \delta\bm1,
\]
using that $p$ is row-stochastic ($p\bm1=\bm1$). Multiplying through by $D_z^{-1}$,
\begin{equation}\label{ch25:eq-eigen}
  Aw \;=\; \delta\,D_z^{-1}\bm1 \;=\; \delta\,w .
\end{equation}
So $w$ is an eigenvector of $A$ with eigenvalue $\delta$, and $w>0$ because $z>0$.

\emph{Existence and uniqueness of $(\delta,w)$.} $A$ has all entries strictly positive, so
\cref{thm:pf} (Perron--Frobenius, proved in \cref{ch:linalg}) applies directly: $A$ has a real eigenvalue
$\lambda_0>0$ strictly larger in modulus than every other eigenvalue, with a positive eigenvector unique up
to positive scaling, and no other eigenvalue has a positive eigenvector (every other eigenvector has entries
of mixed sign or is complex). Equation~\eqref{ch25:eq-eigen} exhibits $w>0$ as a positive eigenvector of
$A$, so by \cref{thm:pf}(iii) it must be \emph{the} Perron eigenvector, unique up to scale, and $\delta$
must be \emph{the} Perron eigenvalue $\lambda_0$: no other pair $(\delta',w')$ with $w'>0$ solving
\eqref{ch25:eq-eigen} can exist. This proves $\delta$ and $z=1/w$ (up to the same overall scale) are unique.

\emph{Recovering $p$.} Solve \eqref{ch25:eq-matrixform} for $p$: $p=\delta^{-1}D_zAD_z^{-1}$, i.e.\ entrywise
$p(x,y)=\delta^{-1}z(x)A(x,y)/z(y)=\delta^{-1}A(x,y)w(y)/w(x)$, which is \eqref{ch25:eq-precover}. Positivity
$p(x,y)>0$ is immediate from $A,w,\delta>0$. Row-stochasticity is automatic and does not need to be imposed
separately: by \eqref{ch25:eq-eigen},
\[
\begin{gathered}
  \sum_{y=1}^np(x,y) \;=\; \frac{1}{\delta\,w(x)}\sum_{y=1}^nA(x,y)w(y) \;=\; \frac{(Aw)(x)}{\delta\,w(x)} \\
  \;=\; \frac{\delta\,w(x)}{\delta\,w(x)} \;=\; 1,
\end{gathered}
\]
for every $x$. Since $(\delta,w)$ are the unique Perron data of $A$, $p$ in \eqref{ch25:eq-precover} is the
unique matrix consistent with \eqref{ch25:eq-ti} and the given $A$.
\end{proof}

\begin{coursenote}[What the lecture actually showed]
Carr states Ross's Theorem~1 essentially verbatim (``if the pricing matrix $A$ is positive or irreducible,
then there exists a unique solution... via the Perron--Frobenius theorem'') but says explicitly that he has
no slide on \emph{how} the $n+1$ unknowns are actually computed, referring the audience to Ross's paper
\ts{13}{25}{1:01:49}. The derivation above (via \eqref{ch25:eq-matrixform}--\eqref{ch25:eq-eigen}) is a
reconstruction, not shown on the slides or in the transcript, built from the definitions Carr does give
(the pricing kernel \eqref{ch25:eq-kernel}, transition independence \eqref{ch25:eq-ti}, and the explicit
statement that Perron--Frobenius is the tool used). The theorem also holds, with $|\lambda|\le\delta$ rather
than strict inequality, for \emph{irreducible non-negative} $A$ (some $A(x,y)$ may vanish): this is the
version of \cref{thm:pf} noted in \cref{ch:linalg} as the one needed for Markov chains, and it is the
version Ross actually states.
\end{coursenote}

\begin{finance}[Reading the recovered kernel]
Once $p$ is known, \eqref{ch25:eq-ti} gives back the pricing kernel itself,
$\varphi(x,y)=\delta\,z(y)/z(x)=\delta\,w(x)/w(y)$, fully decontaminated: $\delta$ is the (state-independent)
discount factor and $w(x)/w(y)$ isolates the risk-aversion premium of moving from $x$ to $y$. If the recovered marginal utility
$z=1/w$ is decreasing in the state index (i.e.\ $w$ is increasing) when the index orders states from
``bad'' to ``good'', a transition into a worse state costs relatively more than a fair bet on its recovered
probability alone would suggest --- exactly the insurance-value story of \cref{ch25:sec-pqr}. Unlike the
one-period model of \cref{ch:linalg}, where \emph{any} positive $\bm\psi$ with $A_0=\bm\psi\T A$ is an
admissible state-price vector when the market is incomplete, here positivity of the whole matrix $A$ plus
\eqref{ch25:eq-ti} is what removes all the remaining freedom in one step, without ever specifying a utility
function numerically.
\end{finance}

% ==================================================================
\section{Worked example and numerical verification}\label{ch25:sec-example}

\begin{example}[Recovering a 3-state chain]\label{ch25:ex-3state}
Take three states (say, three ranges of a 1-year index return: down, flat, up) with the ``true'' natural
transition matrix
\[
  p = \begin{pmatrix} 0.6 & 0.3 & 0.1 \\ 0.2 & 0.5 & 0.3 \\ 0.1 & 0.4 & 0.5 \end{pmatrix},
\]
a discount factor $\delta=0.97$, and marginal utilities $z=(1.0,\,1.6,\,2.5)$ (higher in the ``up'' state,
i.e.\ increasing marginal utility of consumption in the better state --- an unusual but not excluded case,
matching Ross's remark that the recovered kernel could be increasing in $y$ \ts{13}{25}{55:28}). By
\eqref{ch25:eq-ti}--\eqref{ch25:eq-kernel}, the (only observable, in this thought experiment) AD price
matrix is
\[
  A(x,y) = \delta\, p(x,y)\, \frac{z(y)}{z(x)}
  = \begin{pmatrix} 0.5820 & 0.4656 & 0.2425 \\ 0.1213 & 0.4850 & 0.4547 \\ 0.0388 & 0.2483 & 0.4850 \end{pmatrix}.
\]
Given only $A$, \cref{ch25:thm-ross} recovers $\delta$ as its Perron eigenvalue and $w=1/z$ as its Perron
eigenvector (fixing the scale by, say, $w(1)=1$); \eqref{ch25:eq-precover} then reconstructs $p$ exactly,
with no knowledge of $z$, $\delta$, or $p$ used in the recovery step itself.
\end{example}

\begin{casestudy}[Numerical check, Python]
The following reconstructs \cref{ch25:ex-3state}: build $A$ from $(p,\delta,z)$, forget $(p,\delta,z)$, and
recover them from $A$ alone via \texttt{numpy.linalg.eig}.
\begin{lstlisting}[language=Python]
import numpy as np

P = np.array([[0.6, 0.3, 0.1],
              [0.2, 0.5, 0.3],
              [0.1, 0.4, 0.5]])
delta, z = 0.97, np.array([1.0, 1.6, 2.5])
A = delta * P * (z[None, :] / z[:, None])     # build the AD price matrix

eigvals, eigvecs = np.linalg.eig(A)           # forget P, delta, z; recover them
i = np.argmax(eigvals.real)                   # Perron eigenvalue = largest one
lam0, w = eigvals[i].real, np.abs(eigvecs[:, i].real)
w /= w[0]                                     # fix the free positive scale

P_recovered = (1 / lam0) * A * (w[None, :] / w[:, None])
print(lam0, w)                 # -> 0.97, [1.0, 0.625, 0.4] = 1/z (up to scale)
print(np.max(np.abs(P_recovered - P)))         # -> ~4e-16
\end{lstlisting}
The recovered eigenvalue is $\lambda_0=0.97=\delta$ to machine precision, the eigenvector is
$w=(1,\,0.625,\,0.4)=1/z$ up to the chosen scale, and $P_{\mathrm{recovered}}$ matches the true $p$ to
within $4\times10^{-16}$ (floating-point roundoff): \cref{ch25:eq-precover} reproduces the natural
transition matrix exactly from the price matrix alone.
\end{casestudy}

\begin{intuition}[Why positivity of $A$ is doing all the work]
Nothing in the recovery step uses probabilistic reasoning about $p$ directly; it is pure linear algebra on
$A$. The only place probability enters is in guaranteeing, via no-arbitrage, that $A$ is a matrix of
strictly positive prices in the first place, and in the modelling assumption \eqref{ch25:eq-ti} that
restricts what $A$ can look like. Once both hold, the Perron eigenvector is forced to be positive and simple
by \cref{thm:pf}, and positivity of the reconstructed $p$ and its automatic row-sum-one property follow for
free --- there is no need to separately check that the answer is a sensible probability matrix.
\end{intuition}

% ==================================================================
\section{Beyond finite-state chains: change of numeraire}\label{ch25:sec-diffusion}

Ross's own proof works in the finite-state, discrete-time setting above. Carr and Yu's contribution is a
preference-free re-derivation, avoiding the representative-agent language, that extends the same conclusion
to a continuous-state diffusion. The key substitute for the discrete pricing kernel is a change of
numeraire.

\begin{definition}[Numeraire, numeraire portfolio]\label{ch25:def-numeraire}
A \emph{numeraire} is a self-financing portfolio whose value is strictly positive at all times (a swap can
be zero or negative, hence disqualified; a stock or a money market account qualifies, as long as it never
hits zero) \ts{13}{25}{25:20}. Long (1990) showed that, given no arbitrage among a money market account and
$n$ risky assets, there is always a self-financing portfolio of them, of value $L_t>0$, the \emph{numeraire
portfolio}, such that every relative price $S^i_t/L_t$ is an $R$-martingale --- i.e.\ expressing P\&L in
units of $L$ turns the (unobservable) natural measure $R$ into the measure under which discounted prices
behave like ordinary risk-neutral martingales \ts{13}{25}{1:16:14}. The same object is also called the
\emph{growth-optimal portfolio} (Kelly, 1956: over an infinite horizon it has the largest mean growth rate
among all portfolios, at the price of being quite risky along the way) or the \emph{natural numeraire}
\ts{13}{25}{19:04}.
\end{definition}

Assume $X$ is a univariate, time-homogeneous diffusion under $\Q$,
\[
  dX_t = b^\Q(X_t)\,dt + a(X_t)\,dW_t,
\]
on a \emph{bounded} state space $(\ell,u)$, driving the short rate $r_t=r(X_t)$ and the value of $n$ risky
assets, all with known $\Q$-coefficients; and that the numeraire portfolio's value is a function
$L_t=L(X_t,t)$ solving $dL_t/L_t=r(X_t)\,dt+\sigma_L(X_t)\,dW_t$ for an unknown volatility function
$\sigma_L$ \ts{13}{25}{1:17:17}. Self-financing forces $L$ to solve the pricing PDE
$\partial_tL+\tfrac12a^2\partial_x^2L+b^\Q\partial_xL=rL$; writing $\sigma_L(x)=a(x)\,\partial_x\ln L(x,t)$
(from It\^o's formula) and separating variables $L(x,t)=\pi(x)e^{\lambda t}$ turns this PDE into a
\emph{regular Sturm--Liouville problem}
\begin{equation}\label{ch25:eq-sl}
  \frac{a^2(x)}{2}\pi''(x) + b^\Q(x)\pi'(x) - r(x)\pi(x) \;=\; -\lambda\,\pi(x), \qquad x\in(\ell,u),
\end{equation}
with separated boundary conditions at $\ell,u$ \ts{13}{25}{1:25:39}.

\begin{finance}[Sturm--Liouville as the continuum Perron--Frobenius]
Equation~\eqref{ch25:eq-sl} is the exact continuous-state analogue of the eigenvalue problem
\eqref{ch25:eq-eigen}: Sturm--Liouville theory guarantees a smallest eigenvalue $\lambda_0$, strictly below
every other eigenvalue, with an associated eigenfunction $\pi_0(x)>0$ everywhere and unique up to positive
scaling, while every other eigenfunction changes sign at least once \ts{13}{25}{1:26:40} --- precisely the
conclusions (i)--(iii) of \cref{thm:pf} carried over from finitely many states to a continuum of them. Once
$\pi_0$ and $\lambda_0$ are known (numerically, in general), $L(x,t)=\pi_0(x)e^{\lambda_0t}$ is known up to
scale, hence so is $\sigma_L(x)=a(x)\,\partial_x\ln\pi_0(x)$, and Girsanov's theorem then gives the $R$-drift
of $X$ as $b^\Q(x)+a(x)\sigma_L(x)=b^\Q(x)+a^2(x)\,\partial_x\ln\pi_0(x)$ and the $R$-dynamics of every asset in the set \ts{13}{25}{1:19:23}.
This is the ``$11$ unknowns'' of \cref{ch25:rem-dof} replaced by one unknown function $\pi_0$ and one
unknown scalar $\lambda_0$.
\end{finance}

\begin{coursenote}[Boundedness is essential, and still an open problem beyond it]
The construction needs $X$'s state space to be an interval $(\ell,u)$ with $\ell,u$ finite (Carr is candid
that some economists object to bounding, say, the S\&P~500 above, and that he personally has no problem
with a bound of \$20 trillion \ts{13}{25}{27:29}). For an \emph{unbounded} state space the argument needs a
Hilbert-space (square-integrability) version of Sturm--Liouville theory, and it does not always go through:
it \emph{fails} for the standard geometric Brownian motion of Black--Scholes (so the natural measure of a
stock price cannot be recovered this way from stock option prices under these assumptions), but it
\emph{succeeds} for the Cox--Ingersoll--Ross short-rate model, recovering an R-dynamics that is again CIR
but with a higher mean-reversion speed. More generally it succeeds for affine diffusions (Linetsky and a
student, cited but not derived in the lecture). At the time of the 2013 lecture, necessary and sufficient
conditions for recovery on an unbounded state space were an open research question
\ts{13}{25}{1:27:41}.
\end{coursenote}

% ==================================================================
\begin{keyformulas}
\begin{align*}
&\text{Pricing kernel:} && \varphi(x,y)=A(x,y)/p(x,y)\\
&\text{Transition independence:} && \varphi(x,y)=\delta\,z(y)/z(x)\iff D_zA=\delta\,p\,D_z\\
&\text{Recovery ($w=1/z$):} && Aw=\delta w \ \text{(Perron eigenpair of }A\text{, \cref{thm:pf})}\\
&\text{Recovered transitions:} && p(x,y)=\delta^{-1}A(x,y)\,w(y)/w(x),\quad\textstyle\sum_yp(x,y)=1\\
&\text{Diffusion analogue:} && \tfrac12a^2\pi''+b^\Q\pi'-r\pi=-\lambda\pi\ \text{(Sturm--Liouville)}\\
  &&& \sigma_L(x)=a(x)\partial_x\ln\pi_0(x),\quad b^R(x)=b^\Q(x)+a(x)\sigma_L(x)
\end{align*}
\end{keyformulas}

% ==================================================================
\section*{Exercises}
\addcontentsline{toc}{section}{Exercises}

\begin{exercise}[Recovering a 2-state chain by hand (not from the course)]\label{ch25:ex-x-2state}
Let $A=\bigl(\begin{smallmatrix}0.45&0.30\\0.20&0.40\end{smallmatrix}\bigr)$.
\begin{enumerate}[(a)]
  \item Find the Perron eigenvalue
      $\delta$ and a positive eigenvector $w$ of $A$ by solving the $2\times2$ characteristic equation directly
      (no software).
  \item Use \eqref{ch25:eq-precover} to recover $p$, and check it is row-stochastic.
  \item Recover
      $z=1/w$ up to scale and the pricing kernel $\varphi(x,y)=\delta z(y)/z(x)$; is $\varphi$ larger for a
      transition into state 1 or state 2? Which state carries the ``insurance'' premium?
\end{enumerate}
\end{exercise}

\begin{exercise}[Recovery is not just normalising rows (not from the course)]\label{ch25:ex-x-notnorm}
For the $A$ of \cref{ch25:ex-x-2state}, compute the naive row-normalisation
$\hat p(x,y)=A(x,y)/\sum_zA(x,z)$.
\begin{enumerate}[(a)]
  \item Verify $\hat p\ne p$ from \cref{ch25:ex-x-2state}(b).
  \item Explain, in
      terms of \eqref{ch25:eq-matrixform}, why row-normalising $A$ recovers the natural probabilities only in the
      special case $z\equiv\mathrm{const}$ (no risk aversion), and identify what $\hat p$ actually equals in
      general (a $\Q$-like object, not $R$).
\end{enumerate}
\end{exercise}

\begin{exercise}[Sensitivity of the recovered kernel (not from the course)]\label{ch25:ex-x-sensitivity}
Using the Python code of \cref{ch25:ex-3state}, perturb one entry of $A$ (e.g.\ multiply $A(1,3)$ by $1.05$)
while leaving the others fixed, keeping $A$ positive.
\begin{enumerate}[(a)]
  \item Recompute $\delta$, $w$ and $p$; by how much does
      each change relative to the unperturbed values?
  \item Which entries of $p$ are most sensitive to a change in
      $A(1,3)$, and does this match the intuition that Perron--Frobenius data depend on the \emph{whole} matrix,
      not just one entry?
  \item Repeat with a $5\times5$ random positive matrix and comment on whether the recovered
      $p$ remains close to row-stochastic under floating-point rounding alone.
\end{enumerate}
\end{exercise}

\begin{exercise}[Irreducible but not positive (not from the course)]\label{ch25:ex-x-irreducible}
Let $A=\bigl(\begin{smallmatrix}0&0.5&0.3\\0.4&0&0.6\\0.2&0.4&0\end{smallmatrix}\bigr)$ (zero diagonal, but
irreducible: every state reachable from every other in finitely many steps).
\begin{enumerate}[(a)]
  \item Check numerically that $A$
      still has a positive Perron eigenvector, consistent with the irreducible case of \cref{thm:pf} mentioned in
      \cref{ch:linalg}.
  \item Recover $p$ via \eqref{ch25:eq-precover} and check it is a valid (row-stochastic,
      non-negative) transition matrix with some zero entries.
  \item Explain financially what $A(x,x)=0$ means (no
      one-period AD security paying off in the same state you start in) and why this does not break recovery.
\end{enumerate}
\end{exercise}

\begin{exercise}[From the Sturm--Liouville problem to the discrete one (not from the course)]\label{ch25:ex-x-discretize}
Discretise \eqref{ch25:eq-sl} on a grid $x_1<\dots<x_n$ in $(\ell,u)$ using centred differences for
$\pi''(x_i)\approx(\pi_{i+1}-2\pi_i+\pi_{i-1})/h^2$ and $\pi'(x_i)\approx(\pi_{i+1}-\pi_{i-1})/(2h)$, with
absorbing (Dirichlet) boundary conditions $\pi_0=\pi_{n+1}=0$.
\begin{enumerate}[(a)]
  \item Write the resulting eigenvalue problem as
      $M\bm\pi=-\lambda\bm\pi$ for an explicit tridiagonal $M$ in terms of $a(x_i)$, $b^\Q(x_i)$, $r(x_i)$, $h$.
  \item For constant $a,b^\Q,r$, compute $M$'s eigenvalues in closed form and identify the smallest one and its
      eigenvector's sign pattern; check it matches \cref{thm:pf}'s conclusion that only the extremal eigenvalue
      has an everywhere-positive eigenvector.
  \item Implement this in Python for $a=0.2$, $b^\Q(x)=0.1(0.05-x)$,
      $r(x)=0.03$, on $(\ell,u)=(-0.1,0.2)$ with $n=50$ grid points, and plot $\pi_0$; is it unimodal?
\end{enumerate}
\end{exercise}

\begin{remark}[End of Part~II]
This closes Part~II of these notes. The recovery theorem was a fitting last stop: it leans on nearly every
tool built up along the way --- state prices and no-arbitrage (\cref{ch:linalg}), the risk-neutral measure
(\cref{ch:bs}), Markov chains and Perron--Frobenius (\cref{ch:linalg} again), and, in its diffusion form,
stochastic calculus and PDEs (\cref{ch:stochcalc,ch:sde}). The appendices that follow are practical rather
than mathematical: how to run the course's \textsf{R} material, and the logistics of both years of the
course.
\end{remark}


\appendix
% !TEX root = ../main.tex
\chapter{Map of the Course Materials}\label{app:materials}

Every file downloaded from the two OCW sites, with the chapter where it is used. Paths are relative to
the course folder; \file{static_resources/} is abbreviated \file{sr/}. Status: \textbf{done} = already
integrated in a written chapter, \emph{partly} = some of the problems or chapters that use it are written, \emph{planned} = assigned to a chapter still in preparation.

\newcommand{\done}{\textbf{done}}
\newcommand{\planned}{\emph{planned}}
\newcommand{\partly}{\emph{partly}}

\section*{18.642, Fall 2024 (\file{18.642-fall-2024/})}
\begin{center}\footnotesize\renewcommand{\file}[1]{\texttt{\detokenize{#1}}}
\begin{longtable}{@{}p{0.6\linewidth} p{0.27\linewidth} l@{}}
\toprule
\textbf{File} & \textbf{Used in} & \textbf{Status}\\
\midrule\endhead
\file{sr/mit18_642_f24_lec01_1.pdf} & \cref{ch:markets} & \done\\
\file{sr/mit18_642_f24_lec01_2.pdf} & \cref{ch:markets}, \cref{app:rsetup} & \done\\
\file{sr/mit18_642_f24_lec02_1.pdf} & \cref{ch:linalg} & \done\\
\file{sr/mit18_642_f24_lec02_2.pdf} & \cref{ch:linalg} & \done\\
\file{sr/mit18_642_f24_lec02_3.pdf} & \cref{ch:linalg}, \cref{app:rsetup} & \done\\
\file{sr/mit18_642_f24_lec04_1.pdf}, \file{lec04_2.pdf} & \cref{ch:probability} & \done\\
\file{sr/mit18_642_f24_lec05.pdf} & \cref{ch:sp1} & \done\\
\file{sr/mit18_642_f24_lec06_1.pdf} -- \file{lec06_4.pdf} & \cref{ch:regression} & \done\\
\file{sr/mit18_642_f24_lec07.pdf} & \cref{ch:rates} & \done\\
\file{sr/mit18_642_f24_lec08.pdf} & \cref{ch:regression} & \done\\
\file{sr/mit18_642_f24_lec09.pdf} & \cref{ch:pca} & \done\\
\file{sr/mit18_642_f24_lec10.pdf} & \cref{ch:counterparty} & \done\\
\file{sr/mit18_642_f24_lec11-12.pdf} & \cref{ch:timeseries} & \done\\
\file{sr/mit18_642_f24_lec13.pdf} & \cref{ch:portfolio} & \done\\
\file{sr/mit18_642_f24_lec14_1.pdf} -- \file{lec14_9.pdf} & \cref{ch:sp2} & \done\\
\file{sr/mit18_642_f24_lec17_1.pdf}, \file{lec17_2.pdf} & \cref{ch:volatility} (2024 L19) & \done\\
\file{sr/mit18_642_f24_lec19_1.pdf}, \file{lec19_2.pdf} & \cref{ch:stochcalc} (2024 L24) & \done\\
\file{sr/mit18_642_f24_lec21.pdf} & \cref{ch:bs} & \done\\
\file{sr/mit18_642_f24_lec23.pdf} & \cref{ch:ml} & \done\\
\file{sr/mit18_642_f24_lec24.pdf} & \cref{ch:sde} (2024 L25) & \done\\
\midrule
\file{sr/mit18_642_f24_pset1.pdf} & \cref{ch:linalg} (all problems) & \done\\
\file{sr/mit18_642_f24_pset2.pdf} & \cref{ch:probability,ch:sp1} (P1--2), \cref{ch:pca} (P3), \cref{app:logistics} (P4) & \done\\
\file{sr/mit18_642_f24_pset3.pdf} & \cref{ch:probability} (Stein's lemma), \cref{ch:portfolio}, \cref{ch:pca} & \done\\
\file{sr/mit18_642_f24_pset4.pdf} & \cref{ch:sp2} (Brownian motion), \cref{ch:timeseries} (AR, ARMA) & \done\\
\file{sr/mit18_642_f24_pset5.pdf} & \cref{ch:sp1,ch:stochcalc}, \cref{ch:volatility} & \done\\
\file{sr/mit18_642_f24_pset_analysis.pdf}, \file{pset_rcode.pdf} & \cref{ch:pca} (PCA of Treasury yields) & \done\\
\midrule
\file{sr/mit18_642_f24_rstudio.zip} (\file{FM_Intro1.Rmd}) & \cref{ch:markets} & \done\\
\file{sr/mit18_642_f24_equal_weighted_portfolios.zip} & \cref{ch:linalg} & \done\\
\file{sr/mit18_642_f24_etf_casestudy.zip} & \cref{ch:regression} & \done\\
\file{sr/mit18_642_f24_rstudio_pset5.zip} & \cref{ch:volatility} & \done\\
\file{sr/fedbondyielddata_20230911.csv}, \file{sr/feddataanalysis.ipynb} & \cref{ch:rates}, \cref{ch:pca} & \done\\
\midrule
\file{sr/*_transcript.pdf} (23 files) & PDF copies of the \file{.txt} transcripts; the \file{.txt} files are used instead & --\\
\file{sr/21m383-s23-lecture-28-apr-28-v3_transcript.pdf} & to be identified & \planned\\
\file{pages/syllabus}, \file{calendar}, \file{assignments}, \file{problem-sets} & \cref{app:logistics} & \done\\
\file{pages/investment-game} & \cref{ch:markets} (\cref{ex:game}), \cref{app:logistics} & \done\\
\file{pages/group-project}, \file{pages/final-paper} & \cref{app:logistics} & \done\\
\file{pages/week-1} -- \file{week-15} & mapping lectures $\leftrightarrow$ files (this appendix) & \done\\
\file{static_shared/}, \file{*.html}, \file{*.js}, \file{*.css}, fonts, images & website infrastructure, no content & --\\
\bottomrule
\end{longtable}
\end{center}

\section*{18.S096, Fall 2013 (\file{18.s096-fall-2013/})}
File names in \file{sr/} start with a hash, omitted here.
\begin{center}\footnotesize\renewcommand{\file}[1]{\texttt{\detokenize{#1}}}
\begin{longtable}{@{}p{0.6\linewidth} p{0.27\linewidth} l@{}}
\toprule
\textbf{File} & \textbf{Used in} & \textbf{Status}\\
\midrule\endhead
\file{MIT18_S096F13_lecnote1.pdf} & \cref{ch:markets} & \done\\
\file{lecnote2.pdf} & \cref{ch:linalg} & \done\\
\file{lecnote3.pdf} & \cref{ch:probability} & \done\\
\file{lecnote5.pdf} & \cref{ch:sp1} & \done\\
\file{lecnote6.pdf} & \cref{ch:regression} & \done\\
\file{lecnote7.pdf} & \cref{ch:var} & \done\\
\file{lecnote8.pdf}, \file{lecnote11.pdf}, \file{lecnote12.pdf} & \cref{ch:timeseries} & \done\\
\file{lecnote9.pdf} & \cref{ch:volatility} & \done\\
\file{lecnote10.pdf} & \cref{ch:regularized} & \done\\
\file{lecnote13.pdf} & \cref{ch:commodities} & \done\\
\file{lecnote14.pdf}, \file{lecnote16.pdf} & \cref{ch:portfolio} & \done\\
\file{lecnote15.pdf} & \cref{ch:pca} & \done\\
\file{lecnote17.pdf} & \cref{ch:sp2} & \done\\
\file{lecnote18.pdf} & \cref{ch:stochcalc} & \done\\
\file{lecnote19.pdf} & \cref{ch:bs} & \done\\
\file{lecnote21.pdf} & \cref{ch:sde} & \done\\
\file{lecnote23.pdf} & \cref{ch:quanto} & \done\\
\file{lecnote24.pdf} & \cref{ch:hjm} & \done\\
\file{lecnote25.pdf} & \cref{ch:ross} & \done\\
\file{lecnote26.pdf} & \cref{ch:counterparty} & \done\\
\midrule
\file{CaseStudy1.pdf}, \file{CaseStudy2.pdf}, \file{fm_casestudy_1_rcode.r}, \file{fm_casestudy_1_0.r}, \file{fm_casestudy_2_rcode.r}, \file{fm_casestudy_fx_1.r} & \cref{ch:regression} & \done\\
\file{CaseStudy3.pdf}, \file{CaseStudy5.pdf} & \cref{ch:timeseries} & \done\\
\file{CaseStudy4.pdf} & \cref{ch:volatility} & \done\\
\file{CaseStudy6.pdf}, \file{Smltn_TwoAst.pdf}, \file{ETF_pridA_15/30.pdf}, \file{ETF_pridB_15/30.pdf} & \cref{ch:portfolio} & \done\\
\file{CaseStudy7.pdf} & \cref{ch:pca} & \done\\
\file{fm_casestudy_0_InstallOrLoadLibraries*.r}, \file{R_Rstudio_Instructions_CaseStudy1/2.txt} & \cref{app:rsetup} & \done\\
\midrule
\file{pset1.pdf} & \cref{ch:linalg} & \done\\
\file{pset2.pdf} & \cref{ch:probability,ch:sp1} & \done\\
\file{pset3.pdf} & \cref{ch:regression} & \done\\
\file{pset4.pdf} & \cref{ch:timeseries} & \done\\
\file{pset5.pdf} & \cref{ch:volatility} & \done\\
\file{pset6.pdf} & \cref{ch:timeseries} (VAR), \cref{ch:portfolio} & \done\\
\file{pset7.pdf} & \cref{ch:pca} & \done\\
\file{pset8.pdf} & \cref{ch:stochcalc} & \done\\
\file{pset9.pdf} & \cref{ch:sde} & \done\\
\midrule
\file{<youtube-id>.pdf} (24 files) & PDF copies of the \file{.txt} transcripts & --\\
\file{pages/syllabus}, \file{calendar}, \file{lecture-notes}, \file{case-studies}, \file{instructor-insights}, \file{projects}, \file{related-resources} & \cref{app:logistics} & \done\\
\bottomrule
\end{longtable}
\end{center}

\section*{Transcripts}
\file{Lectures 2024 transcript/} (22 files) and \file{Lectures 2013 transcript/} (24 files): used in every
chapter for explanations given orally and for the video timestamps \ts{24}{1}{m:ss}. Lectures with no
recording: 2024 L3, L15--L17 (student presentations), L22, L26 (presentations); 2013 L4, L22.

% !TEX root = ../main.tex
\chapter{R and RStudio Setup}\label{app:rsetup}
\begin{paracol}{2}

Several case studies in both courses are distributed as runnable R code rather than only as PDF
write-ups: the 2013 case studies (\cref{ch:regression}) ship as plain \file{.r} scripts, and the 2024
case studies (\cref{ch:markets,ch:linalg,ch:volatility}) ship as RStudio projects with R Markdown
(\file{.Rmd}) files. This appendix collects, from the actual install scripts and instruction files
found in the course materials, what is needed to run them and which chapter each file belongs to. It
is a practical reference, not a tutorial on R itself.

\section{2013 case studies: plain R scripts}

The 18.S096 (2013) case studies used in \cref{ch:regression} come with a small family of files whose
names are hash-prefixed in \file{static\_resources/} (the hash is omitted below, as in
\cref{app:materials}): a library-install script, one or two data/plotting scripts per case study, and
a plain-text instructions file.

\begin{coursenote}[Install script pattern]
\file{fm\_casestudy\_0\_InstallOrLoadLibraries.r} (used for Case Study 1) and
\file{fm\_casestudy\_0\_InstallOrLoadLibraries Case\_Study\_2.r} (Case Study 2) both follow the same
pattern: a boolean flag \texttt{ind.install0} guards a block of \texttt{install.packages(...)} calls,
followed by \texttt{library(...)} calls that load the same packages. Set \texttt{ind.install0 <- TRUE}
the first time the script runs on a machine (or to update the packages), then switch it back to
\texttt{FALSE} on later runs so the script only loads what is already installed. The packages required
are:
\begin{center}
\texttt{quantmod}, \texttt{tseries}, \texttt{vars}, \texttt{fxregime}, \texttt{moments},
\texttt{rugarch}, \texttt{xts}
\end{center}
(\texttt{xts} appears only in the Case Study 1 variant of the install script, and only in its
\texttt{install.packages} block, not among the \texttt{library} calls; the Case Study 2 variant installs and
loads the other six.)
\end{coursenote}

The two \file{R\_Rstudio\_Instructions\_CaseStudy1/2.txt} files give the following steps (paraphrased;
both files agree except for the file lists in step 2):
\begin{enumerate}[1.]
\item Install R and RStudio (\url{http://www.r-project.org/} and
  \url{http://www.rstudio.com/ide/download/desktop}).
\item Create a project directory (e.g.\ \file{fm\_casestudy1} or \file{fm\_casestudy2}) and copy into
  it the case study's \file{.r} script files (for Case Study 2, also the \file{fred\_fxrates.txt} data
  file).
\item In RStudio: Project $\to$ Create Project $\to$ Create Project from Existing Directory, and select
  that directory.
\item Open the script files, listed in the lower-right Files pane, in the order given by the
  instructions file.
\item Source (or run with Ctrl+Enter) the install/load script first, with \texttt{ind.install0 <-
  TRUE} on the very first run; then source the data-download/plotting script; then the case study's
  \texttt{\_rcode.r} script, which reproduces the R commands shown in the case study's PDF.
\item Restart RStudio between steps if prompted (e.g.\ after installing packages), answering \texttt{n}
  to "save workspace" -- the scripts reload the data explicitly and do not rely on a saved
  \file{.RData} workspace.
\end{enumerate}

\begin{center}\footnotesize
\begin{tabular}{@{}p{0.3\linewidth}p{0.62\linewidth}@{}}
\toprule
\textbf{Case study} & \textbf{Files} \\
\midrule
Case Study 1 (regression, FX \& rates) & \file{fm\_casestudy\_0\_InstallOrLoadLibraries.r},
  \file{fm\_casestudy\_1\_0.r}, \file{fm\_casestudy\_1\_rcode.r} \\
Case Study 2 (FX regime) & \file{fm\_casestudy\_0\_InstallOrLoadLibraries Case\_Study\_2.r},
  \file{fm\_casestudy\_fx\_1.r}, \file{fm\_casestudy\_2\_rcode.r} \\
\bottomrule
\end{tabular}
\end{center}

Both case studies are used in \cref{ch:regression}; see \cref{app:materials} for the exact
correspondence with the PDF write-ups (\file{CaseStudy1.pdf}, \file{CaseStudy2.pdf}).

\section{2024 case studies: RStudio projects}

The 18.642 (2024) case studies distribute a full RStudio project per case study (a \file{.Rproj} file,
one or more \file{.Rmd} files, and any data the case study needs as an \file{.RData} or \file{.csv}
file), rather than loose scripts. Opening the \file{.Rproj} file in RStudio sets the working directory
and loads the project automatically; there is no separate install-libraries script in these archives,
so the required packages must be installed with \texttt{install.packages(...)} once, before knitting
the \file{.Rmd} file, if not already present.

\begin{center}\footnotesize
\begin{tabular}{@{}p{0.36\linewidth}p{0.13\linewidth}p{0.42\linewidth}@{}}
\toprule
\textbf{Archive / project} & \textbf{Used in} & \textbf{R packages (\texttt{library(...)})} \\
\midrule
\texttt{\small mit18\_642\_f24\_\allowbreak rstudio.zip} (\file{FM\_Intro1.Rmd}) & \cref{ch:markets} &
  \texttt{tidyquant}, \texttt{tidyverse}, \texttt{ggplot2} \\
\texttt{\small mit18\_642\_f24\_\allowbreak equal\_weighted\_\allowbreak portfolios.zip} & \cref{ch:linalg} &
  \texttt{tidyquant}, \texttt{tidyverse}, \texttt{ggplot2} \\
\texttt{\small mit18\_642\_f24\_\allowbreak etf\_casestudy.zip} & \cref{ch:regression} &
  \texttt{tidyquant}, \texttt{tidyverse}, \texttt{car}, \texttt{broom} \\
\texttt{\small mit18\_642\_f24\_\allowbreak rstudio\_pset5.zip} (ARKK/S\&P 500 volatility) & \cref{ch:volatility} &
  \texttt{tidyquant}, \texttt{tidyverse}, \texttt{fpp2}, \texttt{ggfortify}, \texttt{glmnet} \\
\bottomrule
\end{tabular}
\end{center}

(This is the union of \texttt{library(...)} calls actually found across the \file{.Rmd} files in these
four archives; a given \file{.Rmd} may only use a subset of the packages listed for its row. No
combined install script for the 2024 material was found in the static resources -- unlike the 2013
scripts, package installation here is left implicit and must be done once per package with
\texttt{install.packages("<name>")} before the first knit.)

\begin{coursenote}
Two of these are financial-data packages worth calling out: \texttt{tidyquant} and
\texttt{quantmod} both pull historical price/rate series from the internet (typically Yahoo Finance or
FRED), so an internet connection is needed the first time a script or \file{.Rmd} runs, even though the
2013 scripts also cache the result in an \file{.RData} workspace for later offline reruns.
\end{coursenote}

\section{Try it yourself}

\begin{itemize}
\item Run \file{FM\_Intro1.Rmd} (\cref{ch:markets}) end to end in RStudio with "Knit" once the
  packages above are installed, and compare the output to \file{FM\_Intro1.pdf} in the same archive.
\item Change \texttt{ind.install0} back and forth in
  \texttt{\small fm\_casestudy\_0\_\allowbreak InstallOrLoadLibraries.r} and confirm that with it set
  to \texttt{FALSE} the script still runs (using the already-installed packages) but skips the
  \texttt{install.packages} calls.
\end{itemize}

% !TEX root = ../main.tex
\chapter{Course Logistics}\label{app:logistics}
\begin{paracol}{2}

This appendix collects the practical, non-mathematical information about how the two courses were
run: syllabus and grading, calendar structure, problem-set and assignment logistics, the 2024
investment game, the 2024 group project and final paper, and the 2013 instructor interview. It is
assembled from the \file{pages/} subdirectories of the two OCW course sites (\file{syllabus},
\file{calendar}, \file{assignments}, \file{problem-sets}, \file{investment-game},
\file{group-project}, \file{final-paper} for 2024; \file{syllabus}, \file{instructor-insights},
\file{projects}, \file{related-resources} for 2013), reporting only what those pages state.

\section{18.642, Fall 2024}

\subsection{Syllabus and grading}

Two 1.5-hour lectures per week. Prerequisites: 18.03 (Differential Equations), 18.05 or 18.600
(Probability), and 18.06 (Linear Algebra); prior finance knowledge is not required. Mathematics
lectures were given by MIT mathematicians (mainly Peter Kempthorne), applications lectures by
industry professionals (a mix of investment banks, FinTech and hedge funds).

The class had five problem sets, an optional investment game, a mid-semester group project (a
written lecture note plus an in-class presentation), and, in place of a final exam, a final paper on
a math finance topic of the student's choice. The grading page states the weights explicitly:
\begin{center}
\begin{tabular}{@{}ll@{}}
\toprule
Component & Weight\\
\midrule
Homework (problem sets) & 40\%\\
Group project presentation & 20\%\\
Participation & 10\%\\
Final paper & 30\%\\
\bottomrule
\end{tabular}
\end{center}

\subsection{Calendar structure}

The calendar page lists 26 lecture slots, alternating math lectures (mostly Kempthorne, covering the
topics of \cref{part:2024}) with applications lectures by outside speakers (e.g.\ Jeff Shen of
BlackRock on quantitative equity investing, L3; Andrew Gunstensen of Mizuho on rates, L7; Stefan
Andreev of Two Sigma on PCA, L9; Andrew W.~Lo on AI in biomedical portfolios, L18; John Hull on
machine learning, L23). Problem sets were due at lectures 3, 7, 9, 18 and 22. Lectures 15--17 were
student presentations of the group project, and lecture 26 was the final paper presentations.

\subsection{Problem sets}

Five problem sets, each covering the math lecture topics up to that point (see
\cref{app:materials} for the file-to-chapter mapping). The problem-sets page lists two accompanying
readings: Gron, ``Optimal Portfolio Choice and Stochastic Volatility'' (\emph{Applied Stochastic
Models in Business and Industry}, 2012), referenced by Problem 2 of Problem Set~2, and, for the PCA
problem, a note on principal component analysis of interest rates plus sample R code (both used in
\cref{ch:pca}); Problem Set~5 references Parkinson, ``The Extreme Value Method for Estimating the
Variance of the Rate of Return'' (\emph{The Journal of Business}, 1980) and ships a companion R
Markdown volatility case study.

\subsection{Investment game}\label{appC:game}

The investment game (ungraded) is stated in full detail in \cref{ex:game} of \cref{ch:markets}; it
is not repeated here. Two points from the course page are purely administrative and not part of the
exercise: students had to submit their choice of instrument to the instructors before 14 September
2024, and their optional switch before 12 October 2024, and at the end of the term the instructors
announced the top three performers by sharing the tracking spreadsheets.

\subsection{Group project}

Groups worked on an introductory, written lecture note aimed at fellow students with no prior
knowledge of the topic; each group member submitted their own 5--7 page draft on their section. The
first submission (``Draft 1'') was graded only on evidence of individual, focused effort (5+ pages
of draft text), and closed with a reflection paragraph asking for specific feedback. It was followed
by a peer review of another group's draft, revision, and a 15--20 minute oral presentation in class
(lectures 15--17 on the calendar).

The topics offered in 2024 were: How Human Fallacies Lead to Market Anomalies; Regression Modeling
of UK Property Prices; The Use of Machine Learning in Predicting Future Interest Rates; Alternative
Data; Jump-Diffusion Model; Maximizing Utility While Minimizing Risk; Fourier Analysis and Its
Applications in the Stock Market; Black Scholes Merton Model; Implied Volatility; Partial
Differential Equation Models for Option Pricing; Manipulating the Sports Betting Market; Time
Series Modeling; Analyzing Volatility Curves in Option Pricing; The Black-Scholes Model; and
Delta-Hedging.

\subsection{Final paper}

Every student wrote their own final paper, individually graded even when it continued a group
project topic; sections of a joint paper had to be clearly attributed to their author. Papers were
due on the last day of lectures, typically 10--15 pages, graded in part on mathematical content,
writing quality, and organization. As with the group project, a Draft~1 was required first (3--4
polished pages plus 1--2 pages of motivation, or 5--7 pages of developing notes, with a closing
reflection section), assessed only on evidence of focused effort, followed by peer review the next
week.

The final-paper page lists 33 topics proposed for 2024, spanning market microstructure (credit
expansion and land prices, election effects on volatility), derivatives pricing (SABR, Heston,
jump-diffusion Monte Carlo, comparative Black-Scholes/Monte Carlo studies), portfolio and risk
methods (Black-Litterman, PCA, tail-risk hedging), machine learning applications (LSTM volatility
forecasting, reinforcement learning for Greeks hedging, AI-driven trading), and a few lighter
entries (sports-betting binary options, ``The Wolf of Wall Street'' and stock-market cinematography,
DEI-focused ETFs). The full list is not reproduced here; consult the source page for all 33 titles.

\section{18.S096, Fall 2013}

\subsection{Syllabus, goals and grading}

Same meeting pattern (two 1.5-hour lectures/week) and a longer prerequisite list: 18.01, 18.02,
18.03, 18.05 or 18.440, and 18.06. The course had no exams: problem sets were due two weeks after
each math lecture, and the final grade was 75\% homework, 25\% final paper. The syllabus page lists
ten explicit course goals, later repeated verbatim on the instructor-insights page, e.g.\ deriving
the price-yield relationship and convexity, bootstrapping a yield curve, computing standard VaR,
estimating option volatility, deriving Black--Scholes via risk-neutral arguments, PCA as matrix
decomposition, limiting theorems, It\^o's lemma, and Girsanov's theorem and change of measure --
i.e.\ essentially the syllabus of \cref{part:2024,part:2013} combined. The course also offered an
optional field trip to Morgan Stanley's New York offices.

\subsection{Instructor insights (2013)}

The course grew out of Jake Xia's and Vasily Strela's wish, as MIT alumni working in industry, to
give back to the department; it was first taught in Fall 2012 as invited industry lectures only, and
the pure-math component was added the following year to better integrate theory and practice.
Enrollment was about 50 students, mostly undergraduates with a few graduate students, mainly
mathematics majors with some from EECS, economics, and aero/astro; about half had no prior finance
background. Students were expected to spend about 12 hours/week on the course, split between the two
lecture types and, outside class, problem sets (due two weeks after each lecture) and the final
paper.

\begin{coursenote}
The instructors' own reflection: designing problem sets for a mixed-background cohort was an
acknowledged difficulty (``Future offerings will distinguish core versus advanced problems''), which
is one motivation, together with the addition of the applications lectures, for how the course
evolved from 2013 to the 2024 edition mapped in \cref{app:materials}.
\end{coursenote}

\subsection{Final paper topics (2013)}\label{appC:2013paper}

Twenty-five percent of the grade was the final paper (see above); the projects page gives sample
topics -- portfolio risk management, regime-shift modeling, low-volatility investing, modeling
financial bubbles (citing Didier Sornette's work), the Black--Scholes/heat-equation correspondence,
hybrid FX/rate products, HJM versus short-rate models, and Ross recovery -- plus a list of topics
actually chosen by students the previous year, mostly variations on the same themes (Black-Scholes
to heat-equation derivations, HJM/Ho-Lee models, FX jump-diffusion bond pricing, Asian options,
minimal-entropy martingale measures, PCA of energy/currency swap curves around the 2008 crisis, and
Monte Carlo/Heston option pricing).

\subsection{Related resources (2013)}

The related-resources page lists six textbooks (Hull's \emph{Options, Futures, and Other
Derivatives}; Baxter and Rennie's \emph{Financial Calculus}; Wilmott, Howison and Dewynne's
\emph{The Mathematics of Financial Derivatives}; Grinold and Kahn's \emph{Active Portfolio
Management}; Fabozzi, Focardi and Kolm's \emph{Financial Modeling of the Equity Market}; and Tsay's
\emph{Analysis of Financial Time Series}) together with a financial glossary of about 40 terms
(alpha, beta, the Greeks, VaR, swaptions, zero-coupon bonds, and so on). The glossary substantially
overlaps with the terminology already introduced across \cref{part:2024,part:2013}; \cref{ex:glossary}
in \cref{ch:markets} is the course's own suggested exercise for building such a glossary, so it is
not reproduced again here.

\section{Sources not covered}

The 2024 \file{pages/week-1} through \file{week-15} pages are a lecture-by-lecture index used to
build \cref{app:materials} and are not repeated here. No \file{case-studies} or
\file{lecture-notes} index page beyond what is already listed in \cref{app:materials} added further
logistics content for the 2013 course.


\end{document}
